\documentclass[11pt,oneside]{article}
\usepackage[margin=1in]{geometry}
\usepackage{amsmath,amssymb,amsthm,mathtools}
\usepackage{booktabs,array}
\usepackage[colorlinks=true,linkcolor=blue,urlcolor=blue]{hyperref}
\newtheorem{theorem}{Theorem}[section]
\newtheorem{proposition}[theorem]{Proposition}
\newtheorem{corollary}[theorem]{Corollary}
\newtheorem{record}[theorem]{Record}
\newtheorem{assessment}[theorem]{Assessment}
\newtheorem{decision}[theorem]{Decision}
\newtheorem{firewall}[theorem]{Claim firewall}
\newcommand{\Tr}{\operatorname{Tr}}
\title{Renewal Yang--Mills Mass-Gap Research Notes\\
Post-Eleventh-Archive Compact Lineage, Version V100}
\author{Proof-indexed continuation after archive \texttt{V100\_f11}}
\date{12 September 2026}
\begin{document}
\maketitle
\tableofcontents

\section{Context, archive and objective}

The objective remains a rigorous four-dimensional quantum
Yang--Mills construction and a strictly positive physical
Hamiltonian mass gap.  That objective has not been reached.
The route studies an exact compact lattice blocking,
its gauge-compatible local fibre, and the effective action
after integrating the eliminated variables.
Local estimates and finite-regulator coefficients are
inputs to that route, not a substitute for its continuum
construction or gap theorem.

This is the requested compact restart after V100.
The complete preceding lineage is preserved as
\begin{center}\small
\texttt{Renewal\_Yang\_Mills\_Mass\_Gap\_Research\_Notes\_V100\_f11.tex}.
\end{center}
It has 1,062,904 bytes and SHA--256
\begin{center}\scriptsize\ttfamily
9BAACC82930099331FC6A6A7EB192A27F563148D2C3782073CD9F13F35DBAC03.
\end{center}
All older archives remain.  A reference such as
\emph{f11 V92} means the correspondingly titled section
of this immediate archive, not a result proved anew here.

The major recent advance is a complete, exactly checked
order-one transverse coefficient on a specified finite
torus, together with a full pure-gauge cancellation.
The next scale-sensitive question is its small-momentum
response.  After recording the proved interfaces, this
note gives a new finite-volume estimator and its error
scaling.  It does not infer an infrared sign from the
already computed momentum-\(\pi\) value.

\section{Major proved interfaces carried forward}

\subsection{The actual block and independent fibre}

Use four-dimensional parity cells with sixteen fine
vertices and the tree whose parent clears the highest
active parity bit.  There are fifteen tree links,
49 remaining link variables and 45 independent fibre
Lie variables per cell.
On the tree slice, the sixteen paths for coarse edge
\(e=(r,\mu)\) are chronological two-link products
\(P_j=U_{x,\mu}U_{x+\widehat\mu,\mu}\).
The reference \(P_0\) starts at the root.
On the common logarithm chart,
\[
 K_e=P_0\exp\!\left(\frac1{16}\sum_j
                   \log(P_0^{-1}P_j)\right),\qquad F=\log K.
\]
The reference convention must be retained in cubic jets.
The derivative \(L=DF(0)\) is the sixteen-path linear
average.  The right inverse \(R\) assigns the coarse
value to all eight boundary links in its family and
zero to internal links, so \(LR=I\).
The local inclusion \(E\) completes seventeen non-tree
internal and 28 nonpivot boundary variables to
\(\ker L\), using one boundary pivot per coarse edge.

\begin{record}[Exact chart and density: f11 V57--V66, V76, V80]
The actual local chart is
\[
 x=X(z,c)=Ez+Rv(z,c),\qquad F(X(z,c))=c,\qquad
 f(z,c)=S_0(\exp X(z,c)).
\]
It is smooth, finite range, and defined on a
volume-independent small max-norm neighbourhood.
Up to a field-independent coordinate constant,
the joint density relative to \(dz\,dc\) is
\[
 e^{-\beta f(z,c)+\ell(z,c)},\qquad
 \ell=\sum_{l\ {\rm non\!-\!tree}}\log j_G(X_l)
                       -\log\det(DF(X)R).
\]
The determinant splits into fixed-size local
coarse-edge blocks.  Its derivatives of any fixed
order have uniform local bounds.
The cell-coercivity proof gives
\[
 D_z^2f(0,0)=E^*d^*dE\geq\kappa I,\qquad \kappa=24/2327.
\]
This is a gap in the eliminated fibre, not in the
retained physical gauge theory.
\end{record}

\subsection{Moving minima and the interacting local integral}

\begin{record}[Uniform minimum and response: f11 V80--V81]
For small \(\|c\|_\infty\), the equation \(f_z=0\)
has a unique small solution \(z_*(c)\) on a
volume-independent chart.  Put \(H_c=f_{zz}(z_*(c),c)\).
There are uniform positive constants with
\[
 hI\leq H_c\leq bI,\qquad
 |(H_c^{-1})_{ij}|\leq Ce^{-a\,{\rm dist}(i,j)},\qquad
 D_cz_*=-H_c^{-1}f_{zc}.
\]
For \(G(c)=f(z_*(c),c)\),
\[
 D_c^2G=f_{cc}-f_{cz}H_c^{-1}f_{zc}.
\]
Its Hessian is exponentially local.  At zero it
annihilates coarse gauge gradients; the fibre
gap has not been transferred to retained directions.
\end{record}

\begin{record}[Actual local integral: f11 V82--V90]
For a fixed small fibre box \(D\), define
\[
 Z_{\beta}^{\rm loc}(c)=\int_D e^{-\beta f+\ell}\,dz,
 \qquad W_\beta=-\log Z_\beta^{\rm loc},\qquad
 N=45(\dim G)n.
\]
At large \(\beta\), its potential has a uniform
Hessian lower bound of order \(\beta\).
The archive proves local moments, a Poincare
estimate, exponentially decaying covariances,
and the leading operator-norm estimate
\[
 \|\beta^{-1}D^2W_\beta-D^2G\|_{\rm op}
        \leq C\beta^{-1/2}(1+\log\beta)^4 .
\]
At fixed volume there is a \(C^2\) expansion
\[
 W_\beta=\beta G+\frac N2\log(\beta/2\pi)
        +Q-\beta^{-1}A_1+O_{C^2}(\beta^{-2}),
 \qquad Q=\tfrac12\log\det H_c-\ell(z_*(c),c).
\]
The coefficient \(A_1\) is bounded extensively.
The next-order \(C^2\) remainder is not yet
proved uniform in volume.
\end{record}

\subsection{Normalized coefficient and its volume limit}

\begin{record}[Coarse Haar and fourth jets: f11 V91--V92]
The correct coarse-Haar coefficient is
\[
 \widehat Q=Q+\log j_{\rm c},\qquad j_{\rm c}(c)=\prod_e j_G(c_e).
\]
It is locally gauge invariant.  For a compact
semisimple group, \(D\widehat Q(0)=0\) and
\(D^2\widehat Q(0)[d_{\rm c}\omega,v]=0\).
Put \(P=f_{zz}(0)^{-1}\), \(A=f_{zc}(0)\),
\(Jv=(-PAv,v)\), and
\[
 B(\xi)_{ij}=f_3[e_i^z,e_j^z,\xi],\qquad
 C(\xi,\eta)_{ij}=f_4[e_i^z,e_j^z,\xi,\eta].
\]
The exact identity coefficient is
\[
\begin{split}
 D^2\widehat Q(0)[u,v]={}&
 \tfrac12\Tr\!\left(PC(Ju,Jv)-PB(Ju)PB(Jv)\right)\\
 &-\ell_2[Ju,Jv]
   +\tfrac1{12}\sum_e\Tr_{\mathfrak g}
           (\operatorname{ad}_{u_e}\operatorname{ad}_{v_e}).
\end{split}
\]
Global adjoint symmetry eliminates the second
minimizer derivative at zero, but not the
linear response in \(J\).
\end{record}

\begin{record}[Coefficient thermodynamic limit: f11 V92--V93]
Let \(K_L=D^2\widehat Q_L(0)\) on the coarse
side-\(L\) torus with identical translated
conventions.  A real translation-invariant
infinite kernel \(k_\infty\) satisfies
\[
 |k_\infty(r)|\leq Ce^{-a|r|_1},\qquad
 \|K_L-\Pi_LK_\infty\|_{\rm row}\leq Ce^{-a'L}.
 \tag{2.1}
\]
Here \(\Pi_L\) periodizes the kernel; the row
norm includes the fixed internal indices.
The symbol \(\mathcal K(p)=\sum_r k_\infty(r)e^{ip\cdot r}\)
is analytic in a strip, Hermitian for real \(p\), and
\[
 \mathcal K(p)g(p)=0,\quad
 g_\mu(p)=1-e^{ip_\mu},\qquad \mathcal K(0)=0.
 \tag{2.2}
\]
The zero follows by differentiating the Ward
identity, not by assuming positivity.
These are coefficient limits, not a construction
of the entire thermodynamic quantum measure.
\end{record}

\subsection{Completed finite-regulator coefficient}

\begin{record}[Exact checkerboard result: f11 V94--V99]
For \(SU(2)\), choose \(\|T\|^2=2\) and coarse
direction \(c_{r,1}=(-1)^{r_0}T\), all other
orientations zero.  Its exact local minimum
has zero internal links and fine logarithms
\[
 a_{2r+b,0}=(-1)^{r_0}{\bf1}_{b_0=1}(b_1-\tfrac12)T,
 \qquad a_{2r+b,1}=(-1)^{r_0}{\bf1}_{b_1=1}T.
\]
For small \(t\), the minimum is exactly \(ta\)
and \(G(tc)=16n(1-\cos t)\).
On the side-two coarse torus, \(n=16\), the
order-one Hessian contributions per cell are
\[
 \begin{array}{c|r}
 \text{normalized measure}&19/3\\
 \text{neutral determinant}&-689/384\\
 \text{charged second-order trace}&79227/4480\\
 \text{charged first-order square trace, subtracted}&166901/5760\\ \hline
 \text{complete }D^2\widehat Q(0)[c,c]/n&-272249/40320
 \end{array}
\]
Exact rational matrices reproduce independent
physical-space tests and give zero full response
on a separate pure-gauge mode.
The verifier is
\texttt{scripts/verify\_ym\_v99\_charged\_determinant.py};
sector traces and dependency hashes are in
\texttt{artifacts/ym\_v99\_charged\_determinant.json}.
This evaluates external momentum \(\pi\), not
the small-momentum coefficient.
\end{record}

\subsection{Compact-fibre progress and boundary}

\begin{record}[Compact map and conditional bounds: f11 V67--V79]
The root block extends measurably to the compact
tree slice using the root path as fallback outside
its logarithm chart.  Its exact Haar/Wilson
pushforward is positive and gauge compatible.
The repeated root path yields local upper
output-density bounds without a flat-collar
assumption.  Lower tubes and defect suppression
are quantitative under suitable flat retained
data or exterior collars, not on every curved
coarse fibre.  The boundary-copy section is
not generally a constrained minimum; V79 gives
explicit descent.  No reflection-positivity
or continuum claim follows from these finite
compact estimates alone.
\end{record}

\section{New result: finite-volume small-momentum estimator}

Use orientations \(0,1,2,3\).
Fix a real unit polarization \(v\) in orientation
\(1\) and one Lie component, with momentum along
orientation \(0\); it is transverse to \(g(te_0)\).
The unit normalization here differs explicitly
from the norm-squared-two checkerboard generator.
Define
\[
 F(t)=v^*\mathcal K(te_0)v,\qquad
 \theta=\lim_{t\to0}\frac{F(t)}{t^2}
        =-\tfrac12\sum_r r_0^2v^Tk_\infty(r)v .
 \tag{3.1}
\]
The limit and absolute convergence follow from
(2.1)--(2.2).  Reality and symmetry make \(F\)
even, hence \(F(t)=\theta t^2+O(t^4)\).

\begin{theorem}[Lowest-momentum acquisition and error scaling]
\label{thm:estimator}
For integer \(L\geq3\), set \(n=L^4\), \(p_L=2\pi/L\),
and define real coarse fields
\[
 u_L(r)=v\cos(p_Lr_0),\qquad v_L(r)=v.
\]
Let
\[
 q_L=\frac1n\langle u_L,K_Lu_L\rangle,\quad
 q_L^0=\frac1n\langle v_L,K_Lv_L\rangle,\quad
 \theta_L=\frac{2q_L-q_L^0}{p_L^2}.
 \tag{3.2}
\]
Constants independent of \(L\) satisfy
\[
 |\theta_L-\theta|\leq CL^{-2}+CL^2e^{-a'L}.
 \tag{3.3}
\]
\end{theorem}
\begin{proof}
Translation invariance diagonalizes \(K_L\)
in the discrete Fourier basis.
For \(L\geq3\), momenta \(p_Le_0,-p_Le_0\)
are distinct and the cosine has squared
average \(1/2\).  Thus
\[
 2q_L=v^*\mathcal K_L(p_Le_0)v,\qquad
 q_L^0=v^*\mathcal K_L(0)v .
 \]
The opposite-momentum values agree for real
\(v\) by kernel reality and symmetry.
Periodization gives the infinite kernel's
symbol exactly at the allowed discrete
momenta.  The row estimate (2.1) implies
\[
 |2q_L-F(p_L)|+|q_L^0-F(0)|\leq Ce^{-a'L}.
 \]
Now \(F(0)=0\), and the finite fourth kernel
moment bounds \(|F(t)-\theta t^2|\leq Ct^4\).
Divide by \(p_L^2\) and use \(p_L=2\pi/L\).
Increase the constant for any finitely many
smaller sides.  This proves (3.3).
\end{proof}

\begin{corollary}[Required arithmetic precision]
If approximations to \(q_L,q_L^0\) have certified
errors \(\epsilon_L,\epsilon_L^0\), their estimator
\(\widetilde\theta_L\) satisfies
\[
 |\widetilde\theta_L-\theta|
 \leq CL^{-2}+CL^2e^{-a'L}
                +\frac{2\epsilon_L+\epsilon_L^0}{p_L^2}.
 \tag{3.4}
\]
Errors \(o(L^{-2})\) suffice for convergence;
errors \(O(L^{-4})\) retain the \(O(L^{-2})\)
discretization rate.
\end{corollary}
\begin{proof}
Apply the triangle inequality to (3.2), use
(3.3), and substitute \(p_L^{-2}=L^2/(4\pi^2)\).
\end{proof}

\begin{assessment}[Not yet a numerical enclosure]
This supplies a convergent target and its
precision scaling, not values at new lattice
sizes.  The constants in the archived decay
and fourth-moment bounds have not been made
numerical here.  Certified finite error bars
for \(\theta\) require those constants and
actual arithmetic error bounds, not only
the existence estimate (3.3).
\end{assessment}

\section{Why the corner sign cannot replace the calculation}

\begin{proposition}[Finite-range Ward-compatible corner ambiguity]
Let \((\mathcal D(p)A)_{\mu\nu}=g_\mu A_\nu-g_\nu A_\mu\).
For real \(\eta,b\), put
\[
 \mathcal J_b(p)=
       \big(\eta+b(1+\cos p_0)\big)\mathcal D(p)^*\mathcal D(p).
 \tag{4.1}
\]
These are real-kernel, finite-range, Hermitian
Ward-compatible quadratic operators, and
\[
 v^*\mathcal J_b(\pi e_0)v=4\eta,\qquad
 \lim_{t\to0}t^{-2}v^*\mathcal J_b(te_0)v=\eta+2b.
 \tag{4.2}
\]
Thus a fixed corner value permits either
small-momentum sign in this structural class.
\end{proposition}
\begin{proof}
The curl kills \(g\), so \(\mathcal J_b g=0\).
For real \(p\), the matrix is Hermitian.
Its entries are Laurent polynomials in
\(e^{ip_\mu}\), giving finite range and
a real symmetric kernel.
Along \(te_0\), the transverse value is
\[
 4\sin^2(t/2)\big(\eta+b(1+\cos t)\big).
 \]
Evaluate at \(\pi\) and expand at zero.
No positivity assumption on the order-one
coefficient was imposed.
\end{proof}

This is a logical counterexample to inferring
the infrared sign from a corner value,
locality and Ward compatibility alone.
It does not assert that every \(\mathcal J_b\)
is realized by the specified YM block.
The actual jets must decide the actual value.

\section{Audit, continuation and claim boundary}

\begin{record}[Companion and version state]
The V100 audit found SM V72, NS V26 and
parallel YM V5 unchanged.  Their hashes and
endpoint assessments are in f11 V100.
No companion result supplies the missing
uniform YM remainder or compact completion.
Only this V1 is active.  Append subsequent
work to its body and replace only the
immediate active predecessor.  Audit the
latest companions at new V5, V10 and every
subsequent fifth version.  At V100 preserve
the lineage as the next archive and restart
V1 again.
\end{record}

\begin{decision}[Next actual computation]
Extend the coefficient assembly beyond the
sign-valued phases of the side-two torus.
Use the actual harmonic lift and normalized
fourth-jet formula at the lowest nonzero
momentum in (3.2), with a pure-gauge test in
the same conventions.
Retain compact-fibre comparison and the
uniform next-order interacting remainder
as independent analytic requirements.
Do not reinterpret the corner coefficient
as completing either one.
\end{decision}

\begin{firewall}[Status after the eleventh restart]
The full fixed-regulator transverse coefficient
and gauge check are proved and preserved.
This V1 adds a convergent small-momentum
estimator, its arithmetic-error requirement,
and a precise obstruction to inferring its
sign from the corner result.
The estimator is not yet evaluated.
Uniform next-order quantum control, compact
completion, continuum reconstruction and the
full Yang--Mills mass gap remain unproved.
\end{firewall}
% =============================================================================
\section{V2: the general axial lift and its exact measure symbol}
\label{sec:f12v2-axis}
% =============================================================================

The V1 predecessor has SHA--256
\texttt{9DB027B6E2A57BD18A79861279935AFE4C11F3480ACD9A28794B9F1CEC9AB864}.
Its body is retained except for the version label.
We now move beyond the sign-valued checkerboard phase.
The actual linear minimizing lift can be written
explicitly at every momentum along one lattice axis.
This also permits an analytic evaluation of the
measure part of the small-momentum coefficient.
The determinant part remains separate.

Let \(\zeta=e^{ip}\), \(\gamma=\cos p\), and set
\[
 \alpha=\frac{3+\overline\zeta}{6+2\gamma},\qquad
 \beta=\frac{3+\zeta}{6+2\gamma},\qquad
 d=\alpha-\beta,\quad e=\beta-\zeta\alpha,\quad
 h=\frac{\zeta-1}{4},\quad
 \rho=\frac{1-\gamma}{3+\gamma}.
 \tag{6.1}
\]
Here \(\alpha+\beta=1\).
We use complex coefficients only as Fourier
bookkeeping, not as physical complex Lie fields.
The real cosine field is recovered by taking
the real part of \(\zeta^{r_0}\) times the
coefficient.  For the moment retain the archive's
generator \(T\) with \(\|T\|^2=2\).
The coarse mode has coefficient \(T\) in
orientation \(1\) and zero in other orientations.

\begin{theorem}[Closed form of the actual axial harmonic lift]
\label{thm:f12v2-lift}
In cell \(r\), multiply the following fine
coefficients by \(\zeta^{r_0}T\):
\[
 a_{b,0}=
 \begin{cases}
 b_1d,&b_0=0,\\
 b_1e+h,&b_0=1,
 \end{cases}
 \qquad
 a_{b,1}=
 \begin{cases}
 0,&b_1=0,\\
 2\alpha,&b_1=1,\ b_0=0,\\
 2\beta,&b_1=1,\ b_0=1,
 \end{cases}
 \qquad a_{b,2}=a_{b,3}=0 .
 \tag{6.2}
\]
This is the minimizing linear lift \(E(-H^{-1}Ac)+Rc\).
The Hermitian classical Hessian symbol on this
norm-squared-two polarization is
\[
 \mathcal G_T(p)=16\rho.
 \tag{6.3}
\]
For unit polarization it is \(8\rho=p^2+O(p^4)\).
The assertions hold on allowed periodic momenta
and on the continuous Bloch symbols.
\end{theorem}
\begin{proof}
Every nonzero internal entry in (6.2) is in
orientation \(0\) with \(b_1=1\); it is not a
tree link.  Thus the tree constraint holds.
The orientation-\(1\) block average is
\(\alpha+\beta=1\).
The orientation-\(0\) average is
\((1-\zeta)(\alpha+\beta)/4+h=0\).
The other averages vanish, proving the block
constraint.

Direct differencing shows that the only nonzero
curl orientation is \(01\).  Its coefficients,
independent of \(b_1,b_2,b_3\), are
\[
 (Da)_{b,01}=
 \begin{cases}\beta-\alpha,&b_0=0,\\
               \zeta\alpha-\beta,&b_0=1.
 \end{cases}
 \tag{6.4}
\]
Consequently \(D^*Da\) is zero outside
orientation \(1\).  On its even and odd fine
links the coefficients respectively are
\[
 2\alpha-(1+\overline\zeta)\beta=\rho,\qquad
 2\beta-(1+\zeta)\alpha=\rho .
 \]
The coarse multiplier with orientation-\(1\)
coefficient \(8\rho\) has exactly this pullback
under \(L^*\): internal and boundary coefficients
are both \(1/8\), since the mode is constant
in orientation \(1\).
Thus \(D^*Da=L^*(8\rho c)\) on the tree slice,
and \(E^*D^*Da=0\).
Strict positivity of the fibre quadratic form
identifies the unique minimizing lift.

There are eight fine plaquettes of each parity
\(b_0\) per cell.  Their scalar curl square is
\[
 8\big(|\beta-\alpha|^2+|\zeta\alpha-\beta|^2\big)
       =8\rho .
 \]
The Lie norm squared of \(T\) is two, so the
Hermitian Hessian symbol is \(16\rho\).
The unit-polarization value is half this.
Since \(1-\cos p=p^2/2+O(p^4)\), its expansion
is as stated.  The algebra and stationarity
identities apply for every unit-modulus \(\zeta\);
the gapped Bloch fibre matrix gives the same
uniqueness at continuous momentum.
\end{proof}

\begin{theorem}[Exact axial density and normalized measure symbols]
\label{thm:f12v2-measure}
On the lifted norm-squared-two polarization,
the Hermitian reduced-density Hessian symbol is
\[
 \mathcal E_T(p)=
 \frac{\gamma^3+67\gamma^2-345\gamma-747}
                       {12(3+\gamma)^2}.
 \tag{6.5}
\]
The measure contribution to \(D^2\widehat Q\),
including coarse Haar, is
\[
 \mathcal M_T(p)=
 \frac{675+297\gamma-75\gamma^2-\gamma^3}
                       {12(3+\gamma)^2}.
 \tag{6.6}
\]
In unit polarization,
\[
 \mathcal M_{\rm unit}(p)=\tfrac12\mathcal M_T(p)
                  =\frac73+\frac{19}{48}p^2+O(p^4).
 \tag{6.7}
\]
\end{theorem}
\begin{proof}
We evaluate f11 (69.5), using its Hermitian
polarization for complex Fourier coefficients.
The sum of squared scalar fine-link magnitudes
per cell is
\[
 S=4|d|^2+4|h|^2+4|e+h|^2
                    +16(|\alpha|^2+|\beta|^2)
  =\frac{89+45\gamma-5\gamma^2-\gamma^3}{(3+\gamma)^2}.
 \]
The fine-Haar contribution is \(-2S/3\).
For orientation \(1\), the internal links
vanish and the path deviations from the mean
are \(d,-d\); its third-jet correction to the
density is \(-4|d|^2/3\).
For orientation \(0\), the path mean is zero.
Write
\[
 B_0=\tfrac12(|h|^2+|e+h|^2),\qquad
 R_0=\tfrac14\operatorname{Re}
              \big((e+h)\overline d(1+\overline\zeta)\big).
 \]
The simplified zero-mean block formula f11 (69.6)
contributes \(-4(3R_0+B_0)/3\).
The only nonzero second-jet adjoint coefficient
is \(k_0=(1-\zeta)d/8\), and its squared trace
is \(-8|k_0|^2\).
Thus
\[
 \mathcal E_T=-\tfrac23S
       -\tfrac43\big(|d|^2+3R_0+B_0\big)-8|k_0|^2.
 \tag{6.8}
\]
All factors use
\(\Tr_{\mathfrak{su}(2)}\operatorname{ad}_T^2=-8\).
Since \(D\ell_{\rm tr}(0)=0\), there is no
additional second-chart derivative in this
pullback.

For explicit algebra, the quantities in (6.8) are
\[
 |d|^2=\frac{1-\gamma^2}{(3+\gamma)^2},\quad
 B_0=\frac{1-\gamma}{16}
                  +\frac{(1-\gamma)^3}{16(3+\gamma)^2},\quad
 R_0=\frac{(1-\gamma)^2(1+\gamma)}{8(3+\gamma)^2}.
 \]
Together with \(k_0=(1-\zeta)d/8\), substitution
and collection of the polynomial numerator
give (6.5).  The coarse-Haar term is \(-2/3\).
Hence \(\mathcal M_T=-\mathcal E_T-2/3\),
which is (6.6).  Divide by two for unit
polarization and use \(\gamma=1-p^2/2+O(p^4)\)
to obtain (6.7).
\end{proof}

\begin{corollary}[The determinant target after the measure calculation]
In the thermodynamic identity coefficient,
let \(\mathcal D_{\rm unit}(p)\) denote the
determinant part in the same polarization.
Then
\[
 \mathcal D_{\rm unit}(0)=-\frac73,\qquad
 \theta=\theta_{\rm det}+\frac{19}{48},\qquad
 \theta_{\rm det}
  =\lim_{p\to0}\frac{\mathcal D_{\rm unit}(p)+7/3}{p^2}.
 \tag{6.9}
\]
Neither \(\theta_{\rm det}\) nor the total
\(\theta\) is evaluated here.
\end{corollary}
\begin{proof}
The full coefficient is the sum of determinant
and measure terms.  Its infinite-volume symbol
vanishes at zero by (2.2).
The measure symbol is (6.7).  Subtraction
gives (6.9), with existence following from
the even analytic scalar symbols.
This is an infinite-volume identity, not
an assertion of exact finite-torus cancellation
on a constant nongradient mode.
\end{proof}

\begin{record}[Exact complex-rational verification]
The script
\texttt{scripts/verify\_ym\_f12\_v2\_axis\_lift.py}
constructs the complex \(96\)-by-\(49\) curl,
\(49\)-by-\(45\) fibre inclusion and Hermitian
\(45\)-by-\(45\) Gram matrix.
It solves all stationarity equations by exact
complex-rational elimination with real positive
pivots, checks all 49 entries against (6.2),
and checks \(D^*Da=L^*(8\rho c)\).
It independently sums the path density formula.
The resulting checks include
\[
 \begin{array}{c|c|c|c}
 \zeta&8\rho&\mathcal E_T&\mathcal M_T\\ \hline
 1&0&-16/3&14/3\\
 -1&8&-7&19/3\\
 i&8/3&-83/12&25/4\\
 3/5+4i/5&8/9&-269/45&239/45
 \end{array}
 \]
The last phase is a continuous Bloch check,
not an assertion that it is a finite-torus
root of unity.  The values are Hermitian
symbols.  For a real cosine at a non-self-conjugate
allowed phase, the corresponding quadratic
value per cell is half its symbol; this
distinction was retained in the V1 estimator.
The phases \(1,-1\) are self-conjugate and
do not have that cosine factor.

The output is
\texttt{artifacts/ym\_f12\_v2\_axis\_lift.json}.
The executed source has SHA--256
\texttt{fb78eb918ce936665cf07c1093641c93a6881207151de4dca4cb5faca4375474}.
The general claims follow from the proofs
above; these exact solves test their
independent-coordinate implementation.
\end{record}

\begin{assessment}[Direction after V2]
The all-axis harmonic lift and the measure
symbol are no longer unknown inputs.
The next coefficient calculation must extend
the determinant assembly to non-sign-valued
internal and external phases, using this
actual lift and the correct cosine factors.
For general \(p\), (6.2) is a linear minimizing
response; no exact nonlinear straight minimum
is claimed as it was for the special checkerboard.
The fourth-jet coefficient uses this linear
lift by the proved identity in Section 2.
\end{assessment}

\begin{firewall}[Position after V2]
The actual axial harmonic lift, classical
symbol and measure contribution are evaluated.
The measure's unit-polarization small-momentum
coefficient is \(19/48\).
The determinant contribution, total infrared
coefficient, uniform quantum remainder,
compact completion and full Yang--Mills mass
gap remain unproved.  Only V2 is active;
the f11 archive and older archives remain.
Next companion audit: V5.
\end{firewall}
% =============================================================================
\section{V3: the off-corner tangent term and its exact regression}
\label{sec:f12v3-tangent}
% =============================================================================

The V2 predecessor has SHA--256
\texttt{E5A486F662621760608CE2D447E2683E9C6FB77C321CCD56C6693D7BC3847D8A}.
Its body is retained except for the version label.
The checkerboard formula in f11 V97 used
\(F_2[a,R\,\cdot]=0\), because that background
had no internal links.  The general lift (6.2)
has internal links.  Its extension must restore
the inverse-Jacobian term that was legitimately
absent at the checkerboard.
No correction to the archived checkerboard
value is being asserted.

\begin{proposition}[The true Cartan minimum and its straight tangent agree to second order]
\label{prop:f12v3-odd}
For a fixed real coarse field \(c\) taking
values in \(\mathbb RT\), let \(z_1=Dz_*(0)c\).
At each finite regulator,
\[
 z_*(tc)=tz_1+O(t^3),\qquad
 X(z_*(tc),tc)=ta+O(t^3),\qquad
 a=Ez_1+Rc .
 \tag{7.1}
\]
Thus the fibre Hessian coefficients through
order \(t^2\) may be computed on the straight
joint ray \((tz_1,tc)\).
\end{proposition}
\begin{proof}
Uniqueness of the small minimum and invariance
under global conjugation by \(\exp(sT)\)
force every coordinate of the minimum to
commute with \(T\).  Hence it is Cartan-valued,
as in f11 V96.  A global \(SU(2)\) conjugation
taking \(T\) to \(-T\) carries this point to
the minimum over \(-tc\).  Both points are
in the small chart, and uniqueness gives
\(z_*(-tc)=-z_*(tc)\).
Smoothness makes its second derivative at
zero vanish, proving the first expansion.
On the Cartan subspace the block and chart
are exactly linear, so the second follows.
Composing the smooth fibre Hessian with
two curves agreeing through order two
gives the same coefficients through that
order.  No exact straight nonlinear minimum
at general momentum, nor a uniform
large-\(\beta\) remainder, is asserted.
\end{proof}

\begin{theorem}[General second tangent and multiplier expansion]
\label{thm:f12v3-general}
Use fine logarithmic action \(s\), block \(F\),
and the straight Cartan background \(ta\).
All tensors \(s_k,F_k\) below are derivatives
at zero.  Put
\[
 M_1=F_2[a,R\,\cdot],\qquad
 N_1=F_2[a,E\,\cdot],\qquad
 N_2=\tfrac12F_3[a,a,E\,\cdot].
 \]
The fibre tangent has expansion
\[
 Z(t)=E+tZ_1+t^2Z_2+O(t^3),\qquad
 Z_1=-RN_1,\quad Z_2=-RN_2+RM_1N_1.
 \tag{7.2}
\]
Let
\(\lambda_1=R^*s_2a\); on the axial lift
this is \(8\rho c\), with \(\rho\) from (6.1).
The constrained fine bilinear coefficients are
\[
\begin{split}
 L_0&=s_2,\\
 L_1[u,v]&=s_3[a,u,v]-\langle\lambda_1,F_2[u,v]\rangle,\\
 L_2[u,v]&=\tfrac12s_4[a,a,u,v]
                    -\langle\lambda_1,F_3[a,u,v]\rangle.
\end{split} \tag{7.3}
\]
The Hessian coefficients are
\[
\begin{split}
 H_0={}&E^*L_0E,\\
 H_1={}&E^*L_1E+Z_1^*L_0E+E^*L_0Z_1,\\
 H_2={}&E^*L_2E+Z_1^*L_1E+E^*L_1Z_1
        +Z_2^*L_0E+E^*L_0Z_2+Z_1^*L_0Z_1 .
\end{split} \tag{7.4}
\]
\end{theorem}
\begin{proof}
The exact tangent identity is
\[
 Z=E-R(DF(ta)R)^{-1}DF(ta)E .
 \]
Since \(DF(ta)R=I+tM_1+O(t^2)\) and
\(DF(ta)E=tN_1+t^2N_2+O(t^3)\),
their product gives (7.2).

One need not assume that \(ta\) is an exact
nonlinear critical point to obtain the
Hessian identity.  In this chart \(X_{zz}\)
lies in the range of \(R\), and
\[
 X_{zz}=-R(DFR)^{-1}F''[Z,Z].
 \]
The second chain rule for \(s(X)\) therefore
gives \(H=Z^*\mathcal L Z\), where
\[
 \mathcal L=s''-\langle\lambda(t),F''\rangle,\qquad
 \lambda(t)=(DF(ta)R)^{-*}R^*\nabla s(ta).
 \]
At a Cartan background the gradient is neutral.
The block derivative preserves neutral and
charged directions, and its neutral restriction
is \(L\).  Thus \(DF(ta)R\) is the identity on
neutral directions and its adjoint inverse
does not change this gradient.
On Cartan fields the action is a sum of
\(2-2\cos\) of linear plaquette angles.
Its gradient is odd in \(t\), whence
\(\lambda(t)=t\lambda_1+O(t^3)\).
Expanding \(\mathcal L\) gives (7.3).
Multiplying by the two tangent series gives
(7.4).  Proposition \ref{prop:f12v3-odd}
identifies these with the true-minimum
coefficients through order two.
\end{proof}

\begin{proposition}[The additional term is at the relevant momentum order]
For the complex axial lift (6.2), the only
nonzero coarse coefficient generating \(M_1\)
is in orientation \(0\):
\[
 M_{1,0}=\operatorname{ad}_{k_0},\qquad
 k_0=\frac{(1-\zeta)d}{8}T
                 =-\frac{p^2}{32}T+O(p^3).
 \tag{7.5}
\]
It vanishes at both \(p=0\) and \(p=\pi\),
but is nonzero in general.  At \(p=\pi/2\),
its complex coefficient is \((-1-i)T/24\).
\end{proposition}
\begin{proof}
The formula f11 (69.1) is
\(k_e=\frac12\overline{\epsilon_j i_j}\).
In orientation \(0\), internal path values
are \(b_1d\) in the source and \(\zeta b_1d\)
in the target, with opposite chronological
signs.  Averaging over \(b_0,b_1\) and
including the prefactor gives \((1-\zeta)d/8\).
Other internal orientations are zero.
Substitute \(d=-i\sin p/(3+\cos p)\) and expand.
The two stated special values follow directly.
\end{proof}

The \(O(p^2)\) term cannot be discarded
merely because it vanished in the corner
test.  No cancellation of its contribution
to the final trace has been proved.
Pure-gauge backgrounds with zero internal
links also have \(M_1=0\), so the archived
gauge regression alone would not test this
extension.

\begin{record}[Exact off-corner quadratic regression]
The script
\texttt{scripts/verify\_ym\_f12\_v3\_offcorner\_tangent.py}
uses a coarse \((4,2,2,2)\) torus and the
real part of (6.2) at \(\zeta=i\).
This rectangular geometry is a regression,
not the isotropic estimator sequence (3.2).
It checks all 128 coarse first-Jacobian
blocks against (7.5), including their
action on \(N_1\).
The free charged test direction has the
same four entries as f11 V98.

The five contributions in (7.4), separating
the newly restored second-tangent term,
are
\[
 \begin{array}{c|r}
 \text{raw }s_4/2&-23929/8100\\
 \text{multiplier part of }L_2&-32/225\\
 2L_1[Ew,Z_1w]&-1/180\\
 2L_0[Ew,-RN_2w]&-9149/25920\\
 2L_0[Ew,RM_1N_1w]&-1/576\\
 L_0[Z_1w,Z_1w]&17/384
 \end{array}
 \]
Their sum is
\[
 \langle w,H_2w\rangle=-\frac{884497}{259200}.
 \tag{7.6}
\]
Omitting \(RM_1N_1\) instead gives
\(-884047/259200\).
Independent exact quartic polarization of
the full implicit chart gives (7.6), not
the value with the term omitted.
Thus the restored term is detectably
nonzero on the actual off-corner chart.

The output is
\texttt{artifacts/ym\_f12\_v3\_offcorner\_tangent.json}.
The executed source has SHA--256
\texttt{8cf3a6474863eef0bb591671ce9da5030ff295d89be19abd1edf2383597f36c2}.
All calculations use exact rational arithmetic.
The script sets its rectangular geometry
only inside its own process; imported
archive verifier sources remain unchanged.
\end{record}

\begin{assessment}[Next determinant assembly]
Equations (7.2)--(7.4), rather than the
checkerboard specialization with \(M_1=0\),
are the correct off-corner inputs.
They also use \(\lambda_1=8\rho c\), not
the checkerboard value \(8c\).
The full non-sign-valued momentum matrix
assembly must retain conjugate phases,
the real-cosine normalization, and both
momentum transfers for a real external
mode.  The charged and neutral trace
evaluation remains to be done.
\end{assessment}

\begin{firewall}[Position after V3]
The general second tangent is proved and
its additional term has a nonzero exact
off-corner regression.
The archived corner result is unchanged.
The full infrared determinant, total
small-momentum coefficient, uniform
interacting remainder, compact completion
and full Yang--Mills mass gap remain
unproved.  Only V3 is active; all archives
remain.  Next companion audit: V5.
\end{firewall}
% =============================================================================
\section{V4: the complete rectangular off-corner determinant}
\label{sec:f12v4-offcorner}
% =============================================================================

The V3 predecessor has SHA--256
\texttt{3A0EC49800E2197769241A8B79C9343000718D50C92BF688A8081509BFBEA8DA}.
Its body is retained except for the version label.
This calculation extends the assembled determinant
from sign-valued phases to genuinely complex phases.
It retains both transfers of the real external
cosine and the \(RM_1N_1\) term in (7.2).

The coarse torus has sides \((4,2,2,2)\), with
\(n=32\) cells.  The transverse external field is
\[
 c_{r,1}=\cos(\pi r_0/2)T,\qquad
 c_{r,\mu}=0\quad(\mu\ne1),\qquad \|T\|^2=2.
 \tag{8.1}
\]
Its lift is the real part of (6.2) at \(\zeta=i\).
This is not the isotropic sequence of (3.2).
It is a nontrivial finite-regulator test of
the off-corner coefficient assembly.

\begin{proposition}[Two-transfer charged trace formula]
\label{prop:f12v4-transfers}
Let \(q=(\pi/2,0,0,0)\).
The internal momenta are
\[
 k_0\in\{0,\pi/2,\pi,3\pi/2\},\qquad
 k_1,k_2,k_3\in\{0,\pi\}.
 \]
In the complex representation \(X+iY\) of
the two real charged colours, write
\[
 H_0(k,k)=M_k=2B_k^*B_k,\quad P_k=M_k^{-1},
 \qquad H_1(l,k)=iK_{lk},\quad H_2(l,k)=A_{lk},
 \]
where \(B_k=D_kE_k\).
The first coefficient is supported on
\(l=k\pm q\), and the second on
\(l=k,k\pm2q\).
They satisfy
\[
 K_{kl}=-K_{lk}^*,\qquad A_{kl}=A_{lk}^*.
 \]
The charged determinant response is
\[
 \frac{\mathcal D_{\rm ch}}n=T_2-T_{11},\quad
 T_2=\frac2n\sum_k\Tr(P_kA_{kk}),\quad
 T_{11}=\frac1n\sum_{l=k\pm q}
                \Tr(P_lK_{lk}P_kK_{lk}^*).
 \tag{8.2}
\]
\end{proposition}
\begin{proof}
The background has Fourier modes \(q,-q\).
Terms linear in it have exactly these
transfers.  Quadratic terms have transfers
zero and \(\pm2q\), which coincide with each
other modulo \(2\pi\) on this lattice but
not with zero.
The charged dot/cross formulas in f11 V98
and the general tangent formula (7.2)
give these coefficients.  Their Fourier
entries are complex, so all adjoints must
be conjugate transposes.

Hermiticity gives the stated identities.
Realification of a positive Hermitian
charged matrix squares its determinant;
thus its half real log determinant equals
its complex log determinant.
Differentiation gives
\(2\Tr(PH_2)-\Tr(PH_1PH_1)\).
The block-diagonal \(P\) selects \(A_{kk}\)
in the first trace, and the two opposite
first-order transfers in the second.
This proves (8.2).
Each square-trace term is nonnegative
because \(P_k,P_l>0\).

The charged complex Fourier amplitudes
are unconstrained complex coordinates:
one must not impose an extra reality
relation between \(k\) and \(-k\) after
already identifying the two real colours
with \(X+iY\).  The external background,
in contrast, is a real cosine and must
retain both of its transfers.
\end{proof}

\begin{proposition}[Neutral trace with cosine-averaged curvature]
\label{prop:f12v4-neutral}
Let \(W\) on the 96 fine plaquettes per cell
be zero outside orientation \(01\), and set
\[
 W_{pp}=\begin{cases}
 1/18,&b_0=0,\ (\mu,\nu)=(0,1),\\
 1/9,&b_0=1,\ (\mu,\nu)=(0,1).
 \end{cases}
 \]
Then
\[
 \frac{\mathcal D_0}n
      =-\frac1n\sum_k\Tr(P_kB_k^*WB_k).
 \tag{8.3}
\]
\end{proposition}
\begin{proof}
The complex curl coefficients at \(\zeta=i\)
are \(i/3\) and \((-1+i)/3\), from (6.4).
For their real parts on the four-cell
period, the squared averages are respectively
\(1/18\) and \(1/9\).
The neutral first Hessian derivative is zero.
Its second Taylor coefficient is
\(-B^*\operatorname{diag}((Da)^2)B\),
as follows from the cosine Hessian in
f11 (100.4).  Only its zero-transfer
part enters the trace against the
translation-invariant inverse \(P\).
Replacing the squared curvature by its
cell average therefore gives exactly
(8.3).  Here \(P=(2B^*B)^{-1}\);
there is no missing factor of two.
\end{proof}

\begin{theorem}[Exact off-corner coefficient on the rectangular regulator]
\label{thm:f12v4-value}
For (8.1), the exact per-cell values are
\[
 \begin{array}{c|r}
 T_2&606140787777649/44253029624832\\
 T_{11}&11179661021265293/663795444372480\\
 \mathcal D_{\rm ch}/n&-1043774602300279/331897722186240\\
 \mathcal D_0/n&-34196265261653/110632574062080\\
 \text{measure contribution}&25/8
 \end{array}
 \]
Consequently
\[
 \boxed{\frac1nD^2\widehat Q(0)[c,c]
           =-\frac{54591508126619}{165948861093120}.}
 \tag{8.4}
\]
Its decimal value is approximately \(-0.32896585\).
\end{theorem}
\begin{proof}[Exact finite-dimensional calculation]
Build the 32 Fourier curl and fibre matrices
with phases in \(\{1,i,-1,-i\}\) in the first
axis and \(\{1,-1\}\) in the other axes.
Invert each \(45\)-by-\(45\) Hermitian
\(M_k\) with complex-rational elimination.
All Schur pivots are real and positive,
and multiplication verifies \(M_kP_k=I\)
entrywise.

For each of the eight transverse-momentum
families, assemble the four first-axis
momentum blocks.  Fourier averaging over
the four external background cells uses
the factor
\(\frac14\sum_{s=0}^3e^{i(k_0-l_0)s}\);
it retains the actual half-amplitudes
of both cosine modes.
The matrices use (7.2)--(7.4), with
\(\lambda_1=(8/3)c\).  The two factors
of \(i\) in \(M_1N_1\) are multiplied
before forming the second tangent.
The resulting \(K,A\) pass all conjugate
transpose identities in the preceding
proposition.

As an independent assembly regression,
transform V3's four-entry free test field
to Fourier space and evaluate the full
second-order quadratic form, including
its transfer \(2q\).  Using the unnormalized
Fourier coefficients and dividing the
quadratic sum by \(n=32\) gives precisely
\(-884497/259200\), the physical-space
quartic polarization (7.6).

Evaluate (8.2)--(8.3) with exact
complex-rational arithmetic.  All traces
have zero imaginary part; every summand
in \(T_{11}\) is nonnegative.
Their rational sums are the displayed
entries.  The measure symbol at this
momentum is \(25/4\) by (6.6); the
real cosine contributes half, namely
\(25/8\).  Adding the three contributions
proves (8.4).  The executable verifier
and its recorded output are identified
below.
\end{proof}

\begin{proposition}[Full off-corner pure-gauge regression]
Take \(\omega_r=\cos(\pi r_0/2)T\) and
\(c^{\rm g}=d_{\rm c}\omega\).
The same calculation with \(a^{\rm g}=Rc^{\rm g}\)
and zero multiplier gives
\[
 T_2^{\rm g}=\frac{271299202837}{12194948640},
 \quad
 T_{11}^{\rm g}=\frac{109402987719}{4064982880},
 \quad
 \frac{\mathcal D_{\rm ch}^{\rm g}}n=-\frac{14}{3}.
 \tag{8.5}
\]
The neutral term is zero, the normalized
measure term is \(14/3\), and the full
response is exactly zero.
\end{proposition}
\begin{proof}
The background is a flat cell-constant
gauge transform of the identity.  Its
Cartan coarse logarithm is the straight
curve \(tc^{\rm g}\), so the Ward identity
requires zero full quadratic response.
The charged assembly with zero
multiplier gives (8.5) by the same
complex-rational operations as above.
The neutral term vanishes because
every plaquette angle is zero.

The scalar coarse values in orientation
\(0\) over the four-cell period are
\(1,1,-1,-1\).  Each is copied to eight
fine boundary links; there are no
internal values or within-block path
deviations.  The fine-Haar Hessian per
cell is \(-16/3\) and the coarse-Haar
Hessian is \(-2/3\).  Their measure
combination is \(14/3\), canceling the
charged determinant exactly.
\end{proof}

\begin{record}[Exact output and scope]
The executed script is
\texttt{scripts/verify\_ym\_f12\_v4\_offcorner\_determinant.py},
SHA--256
\texttt{57c85eab01f4926ea7a97adec1a66cb33f44d69a378cce06d1ab3be3c9d7e071}.
Its output, including the dependency hashes,
is
\texttt{artifacts/ym\_f12\_v4\_offcorner\_determinant.json},
SHA--256
\texttt{5AF2A4C3502E336FE993D372A033D5EE74DFED43714C5725CD0BECF9D68E8459}.
All exact inverse, adjoint, quartic-regression
and full pure-gauge checks passed.
No floating-point inverse or finite
difference is used to obtain the rational
values above.
\end{record}

\begin{corollary}[Fixed-regulator local integral]
For the local-box effective action relative
to coarse Haar,
\[
 \frac1nD^2\widehat W_\beta(0)[c,c]
     =\frac83\beta
       -\frac{54591508126619}{165948861093120}
       +O(\beta^{-1})
 \tag{8.6}
\]
at this fixed rectangular regulator.
\end{corollary}
\begin{proof}
The real-cosine classical Hessian is
\(8/3\) per cell from (6.3), with its
cosine factor retained.
Use the fixed-volume \(C^2\) expansion
recorded in Section 2 and (8.4).
The remainder is not asserted uniform
in volume.
\end{proof}

\begin{assessment}[Next scale-sensitive step]
The off-corner assembly is now checked
with complex phases, both external
transfers, and the additional second
tangent term.  This removes an
implementation bottleneck.
The result is not a value of the
four-dimensional infrared coefficient:
three coarse sides remain fixed at two.
Increasing only the first side would
not implement the isotropic sequence
in (3.2).
The next audit is V5; the next
coefficient work must enlarge the
internal momentum coverage or derive
and control the thermodynamic
momentum integral, rather than
extrapolate from these two differently
sized finite regulators.
\end{assessment}

\begin{firewall}[Position after V4]
The full rectangular \(\pi/2\) response
is evaluated exactly and passes a
full off-corner gauge cancellation.
The isotropic small-momentum limit,
uniform interacting remainder,
compact completion and full
Yang--Mills mass gap remain unproved.
Only V4 is active; all archives remain.
Next companion audit: V5.
\end{firewall}
\section{V5: differentiated thermodynamic integral for the actual fibre}
\label{sec:f12v5}

\subsection{Companion checkpoint and nonduplication}

The V5 checkpoint inspected the latest active companion endings,
not instructions embedded in those documents:
\[
\begin{array}{l|r|l}
\text{companion}&\text{bytes}&\text{SHA--256 prefix}\\ \hline
\text{SM--Einstein V72}&607372&\texttt{F44F9745D43379E1}\\
\text{NS Stability V26}&278283&\texttt{82C887F8C2C69E4F}\\
\text{Parallel YM V5}&54174&\texttt{306F1BB84CFA8A8D}
\end{array}
\]
SM V72 states a fixed-mode, fixed-time many-copy covariance limit;
it explicitly leaves growing cutoff and the physical finite-species
model open. NS V26 concerns pointwise rank-one Oseen derivative norms,
not a full stability theorem. In particular its references to a
future V27 certification are not imported as completed estimates here.
The parallel YM note rules out its proposed volume-uniform tail for
extensive total charge under tensorization. Its fixed-window or
charge-density replacement remains the relevant scope restriction.
None supplies a missing uniform interacting remainder or physical
mass gap for this programme.

The archive already contains momentum differentiation and corner
quadrature: for example f11 V39, equation (43.1), concerns the earlier
block-average shell response. We do not claim that method as new.
The new result below applies the \emph{actual root-chart} fourth jets
of f11 V92, including the second tangent corrected in current V3,
to a gapped-fibre momentum integral. It gives an explicit analytic
integration bound with no omitted infrared ball and no external
finite-difference quotient. This is a coefficient acquisition route,
not another purported proof of a physical gap.

\subsection{An explicit complex strip from the actual Laurent stencil}

Choose an orthonormal real Lie basis. Write the scalar-colour
zero-background fibre Hessian as
\[
 H(k)=(D(k)E(k))^*D(k)E(k)
       =\sum_{r\in\mathcal R}H_r e^{ik\cdot r},
 \qquad k\in\mathbb T^4,\qquad H(k)\ge hI,\quad h=\frac{24}{2327}.
 \tag{9.1}
\]
The full \(SU(2)\) matrix is its direct sum over three colours,
of size \(m=135\). The positivity in (9.1) follows from the
uniform torus fibre inequality in Section 2: allowed momenta
over all torus sizes are dense, and the Laurent symbol is continuous.
It is positivity of the eliminated fibre only.

For complex momentum use the Laurent continuation, not a pointwise
Hermitian transpose. Thus the continued adjoint of a real-coefficient
Laurent matrix is \(A^\#(z)=A(-z)^T\).
In particular \(H(z)=E^\#(z)D^\#(z)D(z)E(z)\) is holomorphic.
For each coefficient set
\[
 w_r=\max\{\|H_r\|_1,\|H_r\|_\infty\},\qquad
 W=\sum_r w_r,\qquad R=\max_{r\in\mathcal R}|r|_1 .
 \tag{9.2}
\]
Here the two matrix norms are the maximum absolute column and row
sums. Hence \(\|H_r\|_2\le w_r\).

\begin{proposition}[Certified actual-fibre stencil and strip]
\label{prop:f12v5-strip}
For the \(45\)-component scalar-colour fibre in (9.1), exact Laurent
arithmetic gives
\[
 |\mathcal R|=33,\qquad R=2,\qquad W=\frac{1041}{2}.
 \tag{9.3}
\]
With
\[
 s=\frac{h}{12WR}=\frac{2}{2422407},
 \tag{9.4}
\]
the inverse \(P(z)=H(z)^{-1}\) exists for
\(\|\operatorname{Im}z\|_\infty\le3s\), and
\(\|P(z)\|_2\le2/h\). The same bound holds for the full-colour matrix.
\end{proposition}
\begin{proof}
Here is the finite construction used to obtain (9.3). Assign a fine
edge at \(x\) to its parity representative and multiply its row by
the Laurent monomial with exponent \(\lfloor x/2\rfloor\).
Tree edges have zero row. Average the rows of both links of each
of the sixteen block paths with weight \(1/16\), obtaining \(L\).
For a free edge insert a unit row in \(E\), and put
\(-8L_{\mu,\mathrm{free}}\) in the boundary pivot of family \(\mu\).
This gives \(LE=0\) as a polynomial identity. Assemble \(D\) from
the four signed edges of each of the 96 parity-cell plaquettes.
Multiply \((DE)^\#DE\) and collect its coefficient matrices.
Their nonzero exponent set, maximum exponent norm and the row/column
sum majorants are (9.3). The accompanying rational-arithmetic
certificate lists every \(w_r\), checks
\(H_{-r}=H_r^T\), and verifies the factorization at three unit-phase
tuples, including a non-root-of-unity rational complex phase.
The factorization itself is obtained by polynomial multiplication,
not inferred from these three tests.

For any \(z=x+iy\) with \(\|y\|_\infty\le t\),
\[
 \|H(z)-H(x)\|_2
 \le\sum_r w_r(e^{t|r|_1}-1)
 \le W(e^{tR}-1).
 \tag{9.5}
\]
Since \(h\le W\), \(3sR=h/(4W)\le1/4\).
The elementary inequality \(e^u-1\le2u\) on \(0\le u\le1/4\)
therefore bounds (9.5), for \(t=3s\), by
\(6WsR=h/2\). Factoring
\(H(z)=H(x)[I+H(x)^{-1}(H(z)-H(x))]\) and summing the
Neumann series proves invertibility and the bound \(2/h\).
Holomorphic inversion proves the asserted analyticity.
\end{proof}

\begin{record}[Exact certificate and normalization]
The executable
\texttt{scripts/verify\_ym\_f12\_v5\_fibre\_stencil.py}
has SHA--256
\texttt{45901c4f20578416cb99d85f8016fd0297aff85bd6744214bd1edc595e4f8c9f}.
The output
\texttt{artifacts/ym\_f12\_v5\_fibre\_stencil.json}
has SHA--256
\texttt{2B0120B0E0F7BA80A89738E82F30AD31A520E47DA6E27D9043F69920AB91D356}.
It records dependency hashes and all 33 coefficient majorants.
All its assertions passed.
The \(2B^*B\) of V4 uses quaternion generators of norm squared two.
Equation (9.1) instead uses orthonormal Lie coordinates and \(B^*B\).
Changing the constant coordinate normalization changes the determinant
only by a background-independent constant; all logarithmic determinant
derivatives agree when vertices are transformed in the same coordinates.
No coefficient \(\theta\) is evaluated by this certificate.
\end{record}

\subsection{The root-chart momentum vertices and coefficient extraction}

Retain the \emph{unit} real external polarization \(v\) of (3.1).
In this subsection \(p\) is a scalar and the external momentum is
\(p e_0\); additions such as \(k+p\) abbreviate \(k+p e_0\).
Let \(u_p(r)=e^{ipr_0}v\). All jets are complex multilinear
extensions of the real-coordinate jets in Section 2.
Use the joint harmonic lift
\(J(p)v=(-P(p)A(p)v,v)\), with \(A=f_{zc}(0)\).

The operator \(B(Ju_p)\) shifts a fibre momentum \(k\) to \(k+p\).
Denote its matrix by \(V(k,p)\). The opposite vertex is
\(\widetilde V(k,p)=V(k+p,-p)\).
The zero-transfer operator \(C(Ju_p,Ju_{-p})\) has matrix \(C(k,p)\).
These definitions use the actual \(f_3,f_4\), not uncorrected
fine-action vertices. They include all chart-tangent and multiplier
terms of V3. For real momenta the opposite vertex is the adjoint
of the forward vertex; at complex momenta its definition remains
the holomorphic expression \(V(k+p,-p)\), not complex conjugation.

\begin{theorem}[Direct differentiated determinant integral]
\label{thm:f12v5-integral}
Define the normalized integral with measure
\(dk/(2\pi)^4\) and set
\[
 \begin{split}
 g(k,p)=\frac12\Tr\bigl(&P(k)C(k,p)\\
        &-P(k+p)V(k,p)P(k)\widetilde V(k,p)\bigr),\\
 d_2(k)&=[p^2]g(k,p).
 \end{split}
 \tag{9.6}
\]
Then the infinite-volume coefficient of (3.1) is
\[
 \boxed{\displaystyle
 \theta=\frac{19}{48}+
       \int_{\mathbb T^4} d_2(k)\,\frac{dk}{(2\pi)^4}.}
 \tag{9.7}
\]
Coefficient extraction and integration commute. No quotient by
\(p^2\), limit of a real-cosine normalization, or assumption
of pointwise Ward cancellation in \(k\) is needed.
\end{theorem}
\begin{proof}
For real commensurate \(p\), apply the exact fourth-jet formula
of Section 2 to the bilinear pair \(u_p,u_{-p}\) on tori whose
side lengths are multiples of its period. The \(C\) term preserves
internal momentum. The two \(B\) terms successively shift it
by \(p\) and \(-p\). The trace per cell is therefore the
discrete momentum average of (9.6). This bilinear pair extracts
the unit-polarization symbol, without a cosine factor \(1/2\).

The inverse in (9.1) is uniformly bounded on the real torus.
All other entries are finite Laurent jets and harmonic lifts.
Consequently the discrete averages converge to the integral.
The coefficient thermodynamic limit (2.1) identifies it with
the determinant part of \(F(p)\). Continuity gives the same
identity for every real \(p\) near zero, by density of
commensurate momenta.

The finite-range jets and the holomorphic inverses from
Proposition~\ref{prop:f12v5-strip} extend (9.6) jointly
to \(\|\operatorname{Im}k\|_\infty\le s,\ |p|\le s\),
and to an open neighbourhood of that compact set.
In particular \(k+p\) has imaginary part at most \(2s\);
the external inverses at \(p\) and \(-p\) are also within the strip.
All derivatives are uniformly bounded there.
Cauchy's coefficient formula, followed by integration over the
compact real torus, proves the interchange.
The normalized measure symbol in (6.6), for unit polarization,
has expansion \(7/3+(19/48)p^2+O(p^4)\).
Adding its \(p^2\) coefficient proves (9.7).
Although the integrated constant cancels by the Ward identity,
no such cancellation is asserted for the separate integrands.
\end{proof}

For implementable coefficient extraction, write
\[
 C(k,p)=\sum_{j\ge0}C_j(k)p^j,\quad
 V(k,p)=\sum_{j\ge0}V_j(k)p^j,\quad
 \widetilde V(k,p)=\sum_{j\ge0}\widetilde V_j(k)p^j,\quad
 P(k+p)=\sum_{j\ge0}P_j(k)p^j .
\]
These are Taylor \emph{coefficients}, including the factorials.
With \(H'=\partial_{k_0}H\), \(H''=\partial_{k_0}^2H\) and
\(P=P(k)\), the inverse coefficients are exactly
\[
 P_0=P,\qquad P_1=-PH'P,\qquad
 P_2=PH'PH'P-\tfrac12PH''P .
 \tag{9.8}
\]
Multiplication of \(H(k+p)P(k+p)=I\), order by order, proves (9.8)
without assuming that the matrices commute. Consequently
\[
 d_2(k)=\frac12\Tr\left(
 P C_2-\sum_{\substack{a,b,c\ge0\\a+b+c=2}}
             P_a V_b P\widetilde V_c\right).
 \tag{9.9}
\]
The sum has six terms. Derivatives in
\(\widetilde V(k,p)=V(k+p,-p)\) act on \emph{both} arguments.
The \(p\)-dependence of the harmonic lift is also differentiated.
Dropping either dependence changes (9.9).
Equation (9.9) is an analytic prescription; a full actual-vertex
implementation and its numerical integration are not claimed here.

\subsection{A quantitative integration and arithmetic firewall}

For the compact complex domain used above, let certified bounds
\(b_s,c_s\) satisfy
\[
 \|V(k,p)\|_2,\|\widetilde V(k,p)\|_2\le b_s,\qquad
 \|C(k,p)\|_2\le c_s .
 \tag{9.10}
\]
They are finite: coefficient sums of the finite jet polynomials,
the bound \(2/h\), and the coefficient sum for \(A\) give one
explicit construction. In the following estimate they are declared
inputs, \emph{not yet numerically evaluated constants}.
Put
\[
 M=\frac m2\left(\frac{2c_s}{h}+
                            \frac{4b_s^2}{h^2}\right),
 \qquad
 \Theta_N=\frac{19}{48}+
     \frac1{N^4}\sum_{j\in\{0,\ldots,N-1\}^4}d_2(2\pi j/N).
 \tag{9.11}
\]
Here all four internal momentum directions are sampled.

\begin{theorem}[Exponential quadrature certificate]
\label{thm:f12v5-quadrature}
For every integer \(N\ge1\),
\[
 |\Theta_N-\theta|
 \le\frac{M}{s^2}
 \left[\left(\frac{1+e^{-sN}}{1-e^{-sN}}\right)^4-1\right].
 \tag{9.12}
\]
If every computed \(d_2\) has an absolute error at most
\(\varepsilon\), the computed average has the same error bound
plus \(\varepsilon\), apart from separately bounded summation
and representation errors. There is no \(N^2\varepsilon\)
amplification as in an external small-momentum difference quotient.
\end{theorem}
\begin{proof}
On the complex domain,
\(|\Tr X|\le m\|X\|_2\) and \(\|P\|_2\le2/h\) give
\(|g(k,p)|\le M\). Cauchy's formula on \(|p|=s\) implies
\(|d_2(k)|\le M/s^2\) for
\(\|\operatorname{Im}k\|_\infty\le s\).
For its Fourier coefficient at \(r\in\mathbb Z^4\), shift each
integration contour by \(-s\,\operatorname{sign}(r_\mu)\).
Periodicity cancels the side integrals. Thus the coefficient
has absolute value at most \((M/s^2)e^{-s|r|_1}\).

The tensor grid average retains exactly Fourier indices \(r=Nl\).
The integral retains only \(r=0\). Absolute convergence therefore
bounds their difference by
\[
 \frac{M}{s^2}\sum_{l\in\mathbb Z^4\setminus\{0\}}e^{-sN|l|_1}
 =\frac{M}{s^2}
       \left[\left(1+2\sum_{j\ge1}e^{-sNj}\right)^4-1\right],
\]
which is (9.12). Averaging \(N^4\) errors bounded by
\(\varepsilon\) contributes at most \(\varepsilon\).
If the average is evaluated with additional rounding, its
independently enclosed error must be added as stated.
\end{proof}

\begin{assessment}[Actual advance and remaining computational obstruction]
The gapped fibre makes the actual corrected coefficient integral
analytic throughout a neighbourhood of the real internal torus.
There is no singular internal infrared corner to excise in this
particular coefficient problem. Formula (9.9) also avoids the
external \(L^{-2}\) subtraction bottleneck of (3.2).
These conclusions do not identify the fibre with the retained
massless gauge sector.

The explicit strip (9.4) is only about \(8.26\cdot10^{-7}\).
In (9.12), even \(e^{-sN}\le1/2\) requires
\(N\ge(\log2)/s\), more than \(8\cdot10^5\) nodes
\emph{per direction}. That condition alone is not a small
error certificate. The crude global bound is consequently
not a practical tensor-grid algorithm.
The vertices in (9.10) have not been bounded numerically,
and no sign or enclosure for \(\theta\) follows yet.

The useful next step is to implement the differentiated
actual vertices with the off-corner and pure-gauge regressions
already available, then seek a substantially sharper enclosure
using local resolvent bounds or structured integration.
Blindly increasing the first side of V4's rectangular regulator,
or calling its negative response a mass, would not address
this obstruction. Even a fully certified \(\theta\) would
still leave the interacting remainder, compact completion,
continuum construction and physical mass gap to be proved.
\end{assessment}

\begin{record}[V5 append and continuation record]
V4 had SHA--256
\texttt{33704BA7D55C8DFA796A232C3E0F44BC60FD020060E4E105406C793F28541220}.
Its full body is retained; only the title version changes
before this appendix. The old active V4 is replaced by V5,
and every archive is retained.
The exact stencil certificate is new; the older corner
derivative calculations are not relabelled as new work.
The next compulsory companion audit and direction review
are at V10. At V100 the active lineage is to be archived
with the next family suffix before another contextual V1.
\end{record}

\begin{firewall}[Position after V5]
Proved: an explicit complex strip for the actual fibre,
the direct differentiated thermodynamic formula, and its
quadrature/arithmetic error bound.
Not proved: a numerical infrared coefficient enclosure,
a cutoff-uniform interacting continuation, or the full
Yang--Mills mass gap. The full goal remains active.
\end{firewall}
\section{V6: an exact quartic cancellation in the corrected tangent trace}
\label{sec:f12v6}

V5 leaves the differentiated actual vertices to be implemented.
Before enlarging that calculation, we determine exactly which part
of V3's added tangent term can affect the infrared coefficient.
This produces a simplification of the pending calculation,
not its completed implementation.

V3 established \(M_1=O(p^2)\) on the axial harmonic background
and correctly did not discard \(RM_1N_1\) on that observation alone.
Here the orientation support of \(M_1\), followed by an exact
opposite-phase cancellation, shows that the zero-transfer part
of \(RM_1N_1\) is in fact \(O(p^4)\).
It therefore affects finite off-corner responses but not
the \(p^2\) coefficient in (9.7).

\subsection{Scalar reduction in the only active orientation}

Use the V2 transverse harmonic lift with quaternion generator
\(\|T\|^2=2\), external phase \(z=e^{ip}\), and
\[
 \gamma=\cos p,\qquad
 d=\frac{\overline z-z}{2(3+\gamma)},\qquad
 e=\frac{1-z}{3+\gamma},\qquad h_0=\frac{z-1}{4},
 \qquad m(z)=\frac{(1-z)d}{8}.
 \tag{10.1}
\]
The subscript on \(h_0\) distinguishes this lift coefficient
from the positive fibre lower bound \(h\) in Section 9.
The orientation-zero link amplitudes within a coarse cell are
\[
 a_0(b;z)=
 \begin{cases}b_1d,&b_0=0,\\ b_1e+h_0,&b_0=1.\end{cases}
 \tag{10.2}
\]
At a fine edge in the neighbouring cell, multiply its amplitude
by the corresponding power of \(z\). The internal orientation-one
amplitudes vanish, as do all amplitudes in orientations two and three.

For an arbitrary internal fibre momentum \(k\), write
\(\eta=e^{ik_0}\). For each
\(b'=(b_1,b_2,b_3)\), let \(I_{b'}\) denote the coordinate
row of the orientation-zero edge at parity \((0,b')\);
set this row to zero if that edge belongs to the tree.
Let \(B_{b'}\) be the row at parity \((1,b')\), and set
\[
 \mathcal V=\sum_{b_2,b_3\in\{0,1\}} B_{(1,b_2,b_3)} .
 \tag{10.3}
\]
Thus \(\mathcal V\) is a \(1\)-by-\(49\) row with four unit
entries, all on boundary links. It is not a cubic vertex \(V(k,p)\)
from Section 9.

Let \(S(z,k)\) be the scalar row defined by
\[
 (N_1[a(z)]w)_0=\operatorname{ad}_T S(z,k)w
 \tag{10.4}
\]
for a fine-link Fourier amplitude \(w\) at internal momentum \(k\);
the output momentum is \(k+p e_0\).
The right inverse \(R_0\) copies an orientation-zero coarse
amplitude to its eight boundary links and is zero elsewhere.

\begin{proposition}[Exact orientation-zero row]
\label{prop:f12v6-row}
One has
\[
 \begin{split}
 S(z,k)=\frac1{32}\bigg[
  &(\eta-1)\sum_{b'}(b_1e+h_0)I_{b'}\\
  &+(1-z)d\,\mathcal V\bigg],
 \qquad S(z,k)R_0=m(z).
 \end{split}
 \tag{10.5}
\]
Moreover \(M_1[a(z)]=F_2[a(z),R\,\cdot]\)
has only its orientation-zero diagonal block nonzero;
that block is \(m(z)\operatorname{ad}_T\).
\end{proposition}
\begin{proof}
The quadratic block jet is
\[
 F_2[a,w]_0=\frac1{32}\sum_{\text{16 paths }(l,q)}
                    \bigl([a_l,w_q]+[w_l,a_q]\bigr).
 \tag{10.6}
\]
Pair the two paths with the same \(b'\), starting respectively
at \(b_0=0\) and \(b_0=1\). The first has an internal then
a boundary link, with amplitudes \(b_1d\) and \(b_1e+h_0\).
The second has a boundary then a neighbouring internal link,
with amplitudes \(b_1e+h_0\) and \(zb_1d\).
The latter internal input row is \(\eta I_{b'}\).
Combining the two commutators gives precisely
\[
 (\eta-1)(b_1e+h_0)I_{b'}+(1-z)b_1d B_{b'}.
\]
Sum and divide by 32 to obtain (10.5).
Every internal row vanishes on \(R_0\), and
\(\mathcal V R_0=4\), proving the multiplier formula.

For orientation one the background's internal links are zero.
Every path has one internal and one boundary link, whereas
\(R\) is supported only on boundary links. Both commutators
in \(F_2[a,R\,\cdot]\) therefore vanish in that orientation.
The background is zero in orientations two and three.
Different coarse orientations do not mix in \(DF\,R\).
This proves the final assertion.
\end{proof}

\subsection{Opposite phases remove all internal rows}

For a real cosine background
\[
 a^{\mathrm{cos}}(r)=\tfrac12
       \bigl(e^{ipr_0}a(z)+e^{-ipr_0}a(\overline z)\bigr),
 \tag{10.7}
\]
the zero-transfer component of the product \(M_1N_1\)
uses one factor from each external phase.
This statement concerns transfer in the background phase,
not a reality constraint on the charged fibre Fourier variables.

\begin{theorem}[Exact zero-transfer tangent correction]
\label{thm:f12v6-zero}
Let \(\Delta Z_2=RM_1N_1E\) be the term added to the second
tangent in V3. For the background (10.7), its zero-transfer
matrix on an internal momentum \(k\) is
\[
 \boxed{\displaystyle
 (\Delta Z_2)_0(k)=
 \frac{|m(z)|^2}{8}\,
      R_0\operatorname{ad}_T^2\mathcal V E(k).}
 \tag{10.8}
\]
Furthermore
\[
 |m(e^{ip})|^2
  =\frac{(1-\cos p)^2(1+\cos p)}
                 {32(3+\cos p)^2}
  =\frac{p^4}{1024}+O(p^6).
 \tag{10.9}
\]
In particular the zero-transfer correction is \(O(p^4)\),
uniformly for real internal momentum.
\end{theorem}
\begin{proof}
The scalar row of the zero-transfer product, before multiplication
by \(R_0\operatorname{ad}_T^2\) and \(E(k)\), is
\[
 \frac14\bigl(m(z)S(\overline z,k)
                   +m(\overline z)S(z,k)\bigr).
 \tag{10.10}
\]
There is no internal momentum shift in \(M_1\):
its coarse-edge block is a pointwise multiplier.
The input to each \(S\) in (10.10) is the same momentum \(k\).

For unit \(z\), \(d\) is purely imaginary, and both
\(h_0\) and \(e+h_0\) are real multiples of \(z-1\).
But \(m(z)=-(z-1)d/8\). Consequently, for \(a=h_0\)
or \(a=e+h_0\),
\[
 m(z)\overline a+\overline{m(z)}a=0.
 \tag{10.11}
\]
The factor \(\eta-1\) in (10.5) is the same in both
summands of (10.10); it is not conjugated when the
\emph{external} phase is reversed. Hence (10.11) cancels
every internal row in (10.10).
For the boundary rows use \((1-z)d=8m(z)\).
Their remaining scalar coefficient is
\[
 \frac1{4\cdot32}
       \bigl(8m(z)\overline{m(z)}
                     +8\overline{m(z)}m(z)\bigr)
       =\frac{|m(z)|^2}{8}.
 \]
This proves (10.8).

Finally \(|1-z|^2=2(1-\gamma)\) and
\(|d|^2=(1-\gamma^2)/(3+\gamma)^2\) give the first equality
in (10.9). Taylor expansion at zero gives its second equality.
The matrices \(R_0,\mathcal V,E(k)\) are uniformly bounded,
so the estimate is uniform in \(k\).
For self-opposite phases \(z=1,-1\), \(m(z)=0\);
the product vanishes identically and the conclusion remains valid,
regardless of the coalescence of external transfers.
\end{proof}

\subsection{Precisely what may be omitted from the infrared trace}

Define the comparison calculation to be identical to the actual
one except for omitting \(RM_1N_1\) from \(Z_2\).
It is an algebraic comparison, not a second gauge-invariant model.
In particular no changes are made to \(N_2\), the multiplier,
the density, the first tangent, or the \(H_1PH_1\) trace.

For the charged complex representation of \(SU(2)\), use the
norm-squared-two generator and the V4 matrices
\[
 L_0(k)=2D(k)^*D(k),\quad
 H_0(k)=E(k)^*L_0(k)E(k),\quad P(k)=H_0(k)^{-1}.
 \]
Here \(\operatorname{ad}_T^2=-4\) on the charged plane.
Set
\[
 A_0(k)=\operatorname{Re}\Tr\!
    \bigl(P(k)E(k)^*L_0(k)R_0\mathcal V E(k)\bigr),
 \qquad
 \overline A_0=\int_{\mathbb T^4}A_0(k)\frac{dk}{(2\pi)^4}.
 \tag{10.12}
\]
The notation \(A_0\) is separate from the mixed Hessian \(A=f_{zc}\).

\begin{corollary}[No change to the actual \(p^2\) coefficient]
\label{cor:f12v6-theta}
The actual determinant symbol minus the comparison symbol,
for the unit external polarization of (3.1), is exactly
\[
 \Delta F_{\det}(p)=-2|m(e^{ip})|^2\overline A_0 .
 \tag{10.13}
\]
Thus this particular tangent correction has zero \(p^2\)
coefficient, even before the internal momentum integral is taken.
On any commensurate finite torus, the identical statement holds
with \(\overline A_0\) replaced by its discrete internal-momentum
average. The \(O(p^4)\) bound is uniform in those torus sizes.
\end{corollary}
\begin{proof}
In the charged representation, (10.8) becomes
\[
 (\Delta Z_2)_0=-\frac{|m(z)|^2}{2}R_0\mathcal V E .
 \]
Only the cross terms with this second tangent change \(H_2\):
\[
 (\Delta H_2)_0
 =E^*L_0(\Delta Z_2)_0+(\Delta Z_2)_0^*L_0E .
 \tag{10.14}
\]
The first-order Hessian is unchanged. The charged contribution to
the second derivative of \(\tfrac12\log\det H\), per cell,
changes by \(2\operatorname{Tr}_{\rm av}(P\Delta H_2)\).
The factor two is the real charged-colour multiplicity in the
complex representation, as in V4.
Hermiticity of \(P,L_0\) turns the two terms in (10.14)
into conjugate traces and gives
\(-2|m(z)|^2\operatorname{av}_k A_0(k)\).
The neutral correction is zero because
\(\operatorname{ad}_T\) annihilates the neutral direction.
Nonzero-transfer parts of \(\Delta H_2\) do not enter this trace,
since \(P\) is translation invariant.

For a non-self-opposite phase the real cosine with
\(\|T\|^2=2\) has per-cell Hessian equal to the unit-polarization
symbol: its cosine factor \(1/2\) cancels the generator norm
squared two. Hence the preceding value is exactly (10.13).
At \(p=0,\pi\) the correction vanishes, as already proved.
Continuity covers the limiting symbol at zero.

All matrix factors in (10.12) are bounded on the real torus,
including \(P\) by the fibre gap. Thus
\(\sup_k|A_0(k)|<\infty\), independently of any momentum grid.
Equation (10.9) proves the pointwise and uniformly averaged
quartic estimate. This also justifies passage to the torus integral.
\end{proof}

\begin{assessment}[Scope of the simplification]
This is not permission to reuse an uncorrected off-corner
determinant as the actual answer. In particular at \(p=\pi/2\),
\(|m|^2=1/288\), so (10.13) can contribute
\(-\overline A_0/144\). The nonzero physical-space correction
in V3 is not contradicted.
Nor does the theorem remove the cubic-block term \(N_2\),
the first-tangent cross terms, the multiplier terms, or density
derivatives. Those remain in the differentiated vertex calculation.

The statement is narrower and useful: when computing only the
axis-polarized \(p^2\) coefficient, the contribution of
\(RM_1N_1\) to the second-order determinant trace is exactly zero.
V3's order count by itself could not establish this; the
orientation support and opposite-phase identity (10.11) do.
No numerical value of \(\overline A_0\) or \(\theta\) is asserted here.
\end{assessment}

\subsection{Executed checks and next target}

\begin{record}[Independent path and off-corner checks]
The new executable is
\texttt{scripts/verify\_ym\_f12\_v6\_tangent\_cancellation.py},
SHA--256
\texttt{ace344b242de2f92165d9efa8354b287001ae70755d456473e8e309fe440f48f}.
Its output is
\texttt{artifacts/ym\_f12\_v6\_tangent\_cancellation.json},
SHA--256
\texttt{FC4A7B44F502A4B073BD80F9EE2768A10C8327045ADDBE704753C27503BE8E68}.
Dependency hashes are included in that output.

Using exact complex-rational arithmetic, the script checks
(10.5), \(S R_0=m\), and (10.8) at twelve combinations of
external and internal phases. The external phases are
\(1,-1,i,3/5+4i/5\); the last is a continuous Bloch-phase check,
not a finite-torus phase. In particular
\[
\begin{array}{c|c|c}
z&|m(z)|^2&|m(z)|^2/8\\ \hline
1,-1&0&0\\
i&1/288&1/2304\\
3/5+4i/5&1/1620&1/12960
\end{array}
\]
It also assembles V4's full local charged \(M_1,N_1\) matrices
at two distinct internal momenta for external phase \(i\).
Summing both intermediate transfers gives exactly
\(\mathcal V/576\) in orientation zero and zero elsewhere,
before the two charged factors of \(i\) are multiplied.
This is the independent V4 convention check of (10.8);
the remaining \(i^2=-1\) supplies the charged minus sign.
All assertions passed.
The all-phase theorem is proved above, not inferred from
these finitely many tests.
\end{record}

\begin{decision}[Continue the reduced actual-vertex calculation]
Use (9.9) with the actual jets, but avoid spending derivative
assembly on the now-proved-zero \(p^2\) trace contribution
of \(RM_1N_1\). Preserve that term in every off-corner regression.
The next work remains the implementation and enclosure of
the other differentiated vertices and their integral.
V5's explicit global strip is still too small for practical
uniform tensor quadrature; this turn has not improved that bound.
The compulsory companion audit and direction review remain V10.
\end{decision}

\begin{record}[V6 append record]
V5 had SHA--256
\texttt{D0E02EEA9E40EA973D79BD2B81CDAE43E6E943250A9DA278F8FDD37147BCA975}.
Its body is retained unchanged, apart from the title version,
before this appendix. Only the active predecessor is replaced;
all family archives remain. This version adds the exact
zero-transfer cancellation and its executed independent checks.
\end{record}

\begin{firewall}[Position after V6]
One corrected chart-tangent contribution has been proved
irrelevant to the axis \(p^2\) coefficient, without discarding it
from the actual off-corner theory. The full differentiated
vertex implementation, a certified infrared coefficient,
the uniform interacting continuation, and the full
Yang--Mills mass gap remain open. The goal is not complete.
\end{firewall}
\section{V7: executed differentiated actual-vertex assembly}
\label{sec:f12v7}

This version implements the previously missing differentiated
charged Hessian, including the actual harmonic lift, multiplier,
first tangent, second tangent and both external transfers.
It evaluates one exact internal-momentum card.
It does not replace the internal four-dimensional integral by
that card, and does not claim a value for \(\theta\).
The new work is the corrected root-chart implementation;
the older shell-vertex derivatives identified at the V5 audit
are not being reproved or relabelled.

\subsection{Phase projection without imposing a four-cell external wavelength}

Let \(p\) be the external momentum and \(z=e^{ip}\).
Introduce an independent auxiliary phase
\(\phi_s=i^s\), \(s=0,1,2,3\).
For the V2 complex harmonic link amplitude, including its
cell-offset factor \(z^{r_0}\), write \(a^+(x;p)\).
Set \(a^-(x;p)=a^+(x;-p)\) and
\[
 a_s(x;p)=\tfrac12\bigl(\phi_s a^+(x;p)
                        +\phi_s^{-1}a^-(x;p)\bigr),\qquad
 {\cal P}_r f=\frac14\sum_{s=0}^3\phi_s^{-r}f_s .
 \tag{11.1}
\]
The auxiliary phase is not a physical cell coordinate.
The actual internal phase at a shifted momentum remains
\(e^{ik_\mu}z^{h\delta_{\mu0}}\), not \(i^h\).

\begin{proposition}[Exact finite projection for the required transfers]
\label{prop:f12v7-projection}
For any expression linear in the background, \({\cal P}_1\)
and \({\cal P}_{-1}\) extract its two half-amplitude transfers
exactly. For any expression quadratic in the background,
\({\cal P}_0\) extracts its zero-transfer part exactly.
These identities hold as formal Taylor identities in \(p\).
They do not require \(p\) to be a fourth-root-of-unity angle.
\end{proposition}
\begin{proof}
A linear expression has only auxiliary exponents \(1,-1\).
A quadratic expression has only exponents \(2,0,-2\).
The average of \(\phi_s^n\) is one if \(4\) divides \(n\)
and zero otherwise. This proves the stated extractions
coefficient by coefficient in \(p\).
The two quadratic exponents \(2,-2\) are aliased by four samples;
we never use this projection to distinguish those transfers.
Only zero transfer is needed from \(H_2\) in its trace against
the translation-invariant inverse.
\end{proof}

For this phase-parametrized background the multiplier is
\[
 \lambda_{1,s}=8\rho(p)\,
          \frac{\phi_s+\phi_s^{-1}}2\,T
 \quad\hbox{in coarse orientation 1},\qquad
 \rho(p)=\frac{1-\cos p}{3+\cos p},
 \tag{11.2}
\]
and zero in the other orientations.
In particular \(8\rho(p)=p^2+O(p^4)\).
Replacing it by V4's fixed value \(8/3\) while differentiating
would be incorrect.

\subsection{Actual charged matrices before momentum integration}

Work in the charged complex representation, retaining V4's
norm-squared-two generator and \(L_0=2D^*D\) normalization.
An index \(h\in\{-1,0,1\}\) denotes internal momentum \(k+hp e_0\).
Write \(E_h,L_h,H_h=E_h^*L_hE_h\), and \(P_h=H_h^{-1}\).
The local matrices \(U_{oi},V_{00},N_{1,oi},N_{2,00},M_{1,oi}\)
are obtained from V4's local plaquette and block-path formulas
using (11.1)--(11.2). In this notation the factor \(i\) is
stripped from the charged first-order matrices, exactly as in V4.

For \(h=1,-1\), define
\[
 W_{1,h0}=-RN_{1,h0}E_0,\qquad
 W_{1,0h}=-RN_{1,0h}E_h,
\]
\[
 W_{2,00}=-R\left(N_{2,00}
                +\sum_{h=\pm1}M_{1,0h}N_{1,h0}\right)E_0.
 \tag{11.3}
\]
The plus sign inside the latter parentheses includes the
two charged factors \(i^2=-1\), followed by the overall minus.
Thus (11.3) is the actual second tangent, not the old omission.
The matrices of \(H_1=iK\) and the diagonal part of \(H_2\) are
\[
 K_{h0}=E_h^*U_{h0}E_0
       +E_h^*L_hW_{1,h0}-W_{1,0h}^*L_0E_0,
 \tag{11.4}
\]
\[
 \begin{split}
 A_{00}={}&E_0^*V_{00}E_0
        +W_{2,00}^*L_0E_0+E_0^*L_0W_{2,00}\\
 &+\sum_{h=\pm1}\bigl(
       W_{1,h0}^*U_{h0}E_0-E_0^*U_{0h}W_{1,h0}
                  +W_{1,h0}^*L_hW_{1,h0}\bigr).
 \end{split}
 \tag{11.5}
\]
These formulas retain all terms of V3.
They follow by multiplying the tangent and fine constrained-Hessian
expansions, using \(Z_1=iW_1\) and \(Z_2=W_2\).
Adjoints give \(K_{0h}=-K_{h0}^*\) and \(A_{00}=A_{00}^*\).

Define the pointwise trace functions
\[
 {\cal T}_2(k,p)=2\Tr(P_0 A_{00}),\qquad
 {\cal T}_{11}(k,p)=
       \sum_{h=\pm1}\Tr(P_hK_{h0}P_0K_{h0}^*).
 \tag{11.6}
\]
Both half-amplitude transfers are present in (11.6);
there is no additional charged reality restriction.
The neutral term is
\[
 {\cal N}(k,p)=-\Tr(P_0 B_0^*W(p)B_0),\qquad B_0=D(k)E(k),
 \tag{11.7}
\]
where \(W\) vanishes except on orientation-\(01\) plaquettes.
There its weights are \(|d|^2/2\) on \(b_0=0\) and
\(|e|^2/2\) on \(b_0=1\), with \(d,e\) as in (10.1).
Thus both weights are \(p^2/32+O(p^4)\), and
\[
 [p^2]{\cal N}(k,p)
     =-\frac1{32}\Tr(P_0B_{01}^*B_{01}),
 \tag{11.8}
\]
where \(B_{01}\) consists of the orientation-\(01\) rows of \(B_0\).

\begin{proposition}[Integral represented by the executed trace assembly]
\label{prop:f12v7-integrand}
The coefficient in (3.1) obeys
\[
 \theta=\int_{\mathbb T^4}{\cal I}_2(k)\frac{dk}{(2\pi)^4},
 \qquad
 {\cal I}_2(k)=\frac{19}{48}
       +[p^2]\bigl({\cal T}_2-{\cal T}_{11}+{\cal N}\bigr)(k,p).
 \tag{11.9}
\]
\end{proposition}
\begin{proof}
For commensurate non-self-opposite external momentum, the
finite-torus trace formula is (8.2)--(8.3), with the actual
harmonic lift and multiplier at that momentum. Equations
(11.3)--(11.7) give exactly those diagonal and single-transfer
blocks. The real-cosine norm-squared-two normalization equals
the unit-polarization symbol, as explained in V6.

Passing to the internal torus integral is justified by the
gapped-fibre bounds in V5. The two bubble transfers express
the even external-momentum symmetrization of the same trace
integral used in (9.6); a relabelling of the loop momentum
identifies their integrals. This symmetrization has the same
coefficient of \(p^2\).
Analyticity permits coefficient extraction under the integral.
Finally add the known measure coefficient \(19/48\).
No pointwise cancellation of the constant term is assumed.
\end{proof}

\subsection{Exact Taylor arithmetic and the new point card}

The implementation works in the quotient ring
\[
 {\cal J}_2=\mathbb Q(i)[p]/(p^3),\qquad
 z=1+ip-\tfrac12p^2 .
 \tag{11.10}
\]
Conjugation conjugates the coefficients, keeping the formal
real parameter \(p\) fixed. In particular \(\overline z=z^{-1}\)
in this ring. This is Taylor conjugation along the real axis,
not a nonholomorphic extension used in V5's strip argument.
For a scalar \(a+bp+cp^2\), \(a\ne0\), the exact inverse is
\[
 a^{-1}-ba^{-2}p+(b^2a^{-3}-ca^{-2})p^2 .
 \tag{11.11}
\]
Matrix inverses use coefficient multiplication, not scalar
division of matrices. For \(H=H_0+H_1p+H_2p^2\),
\[
 P_0=H_0^{-1},\quad P_1=-P_0H_1P_0,\quad
 P_2=P_0H_1P_0H_1P_0-P_0H_2P_0 .
 \tag{11.12}
\]
Multiplying \(HP\) proves (11.12) in the quotient ring.
All \(H_0\) inversions and products in the executed certificate
use complex rational arithmetic. The full matrix identity
\(HP=I\) is checked through degree two.

\begin{theorem}[Exact corrected differentiated card at \(k=(0,\pi,0,0)\)]
\label{thm:f12v7-card}
At this internal momentum, the actual trace functions have
the following Taylor coefficients:
\[
\begin{array}{c|r|r|r}
 & [p^0]&[p^1]&[p^2]\\ \hline
{\cal T}_2&423/35&0&-417/1120\\
{\cal T}_{11}&15149/1225&0&30500653/37044000\\
{\cal T}_2-{\cal T}_{11}&-344/1225&0&-1384154/1157625\\
{\cal N}&0&0&-839/6720\\
\text{unit measure}&7/3&0&19/48
\end{array}
\tag{11.13}
\]
Consequently
\[
 \boxed{{\cal I}_2(0,\pi,0,0)
                  =-\frac{68509331}{74088000}.}
 \tag{11.14}
\]
\end{theorem}
\begin{proof}[Exact finite-dimensional calculation]
Set the internal phases to \((1,-1,1,1)\) and the external
phase to (11.10). Construct \(D,E,L,H\) for the three
shift labels \(-1,0,1\), using the actual path-average
constraint and pivot completion.
The polynomial identity \(LE=0\) is checked for each shift.
Apply the four-phase extraction to the local plaquette and
root-path jets, with the multiplier (11.2), and form
(11.3)--(11.5). All reverse first-matrix adjoint identities
and the diagonal second-matrix Hermiticity identity hold
coefficient by coefficient.

The sum \(M_1N_1\) in (11.3) is retained in the implementation.
Its zero-transfer Taylor polynomial is checked to be zero
through degree two, agreeing with the independent all-phase
theorem of V6. Invert the three shifted \(H\) matrices by
(11.12) and verify their matrix inverse identities.
The constant matrices are positive definite, so exact
elimination has nonzero positive pivots.

Taking (11.6)--(11.7) in the quotient ring gives the first
four rows of (11.13), all real, with zero linear terms.
The last row was proved in V2 and is not recomputed here.
Adding the quadratic coefficients gives (11.14).
The executable and its exact output are recorded below.
This is a finite algebraic evaluation at the specified momentum,
not an evaluation of (11.9).
\end{proof}

\begin{record}[Source reuse, regression scope and provenance]
The new executable is
\texttt{scripts/verify\_ym\_f12\_v7\_differentiated\_vertices.py},
SHA--256
\texttt{D40FCD1CDD217742E8AD51F48961385AB1743DB949D19593D540D7D21F4E1AC5}.
Its output is
\texttt{artifacts/ym\_f12\_v7\_differentiated\_vertices.json},
SHA--256
\texttt{C7D37A0244A57ED9ADBAF9F92127BEEF5F3A86D5CA04CFDB30852BFAB768A394}.
The output includes all dependency hashes and the hash of
the reused local-stencil function source.

To avoid silently changing a previously checked stencil, the
implementation loads V4's local loop into an isolated namespace.
There is one asserted source substitution: the fixed multiplier
\(8/3\) times the phase cosine is replaced by \(8\rho(p)\)
times that cosine. The namespace supplies the general harmonic
background, shifted internal phases, and Taylor arithmetic.
It does not modify the archived source or its global namespace.
Thus the old plaquette and block-cubic formulas are reused,
not claimed as an independently derived new oracle.

As regressions, set the external phase to the constant \(i\).
All five local matrices \(U,V,N_1,N_2,M_1\), at each of the
three required transfers, reproduce V4 entrywise: fifteen
local matrix checks. Next the new three-shift assembly is
compared with V4's full four-momentum family: both first
matrices and the diagonal second matrix agree entrywise.
These three assembled checks test the new composition and
phase parametrization against the previously verified
off-corner calculation.
They complement, but do not replace, V3's independent
physical-space quartic test and V4's gauge regression.
No new continuum or uniform error claim is tested by them.
\end{record}

\begin{assessment}[Meaning of the nonzero point card]
The number in (11.14) is approximately \(-0.9247021245\).
It is one value of the internal-momentum integrand, not
an interval enclosure of its average or a physical mass.
Even the constant term of the combined point integrand is
nonzero here:
\[
 -\frac{344}{1225}+\frac73=\frac{7543}{3675}.
\]
This is compatible with the integrated Ward cancellation.
It illustrates why enforcing that cancellation at every
internal point would change the calculation incorrectly.

The bottleneck has moved: the actual differentiated assembly
is now executable and has passed exact matrix regressions.
The integral and a useful error enclosure are still missing.
A future corner or quarter-turn average must be identified
as a quadrature diagnostic unless accompanied by a proved
error bound; the single card here is not even such an average.
The conservative strip of V5 has not been sharpened.
\end{assessment}

\begin{decision}[Next acquisition]
Use this actual-vertex implementation for momentum coverage
in all four directions, while developing a practical certified
integration bound. Preserve the full off-corner regressions
and the exact inverse identities when optimizing the code.
The next compulsory companion audit and direction review
remain at V10. No conclusion about the uniform interacting
remainder or the continuum theory is imported from a point card.
\end{decision}

\begin{record}[V7 append record]
V6 had SHA--256
\texttt{7BCA112566B623B0EA89AA0CEF9B23C9EF40CB12F88617B27CF31198B3F43A6E}.
Its full body is retained before this appendix; only the
title version changes. The active predecessor is replaced
by V7, and all family archives remain.
\end{record}

\begin{firewall}[Position after V7]
The corrected differentiated charged and neutral trace
assembly is implemented and one exact momentum card is
evaluated. The four-dimensional infrared coefficient,
uniform interacting continuation, compact continuum
construction and full Yang--Mills mass gap remain unproved.
\end{firewall}
\section{V8: exact four-direction derivative grid and a Gaussian cross-check}
\label{sec:f12v8}

V7 implemented the actual differentiated integrand but evaluated
only one internal momentum. Here we complete the smallest
tensor grid in all four directions. The resulting rational
average is exact as a finite sum, not as a continuum integral.
An independent Gaussian calculation checks the constant
cards at every grid point and gives a more structural next
target for investigating the coarse-grid error.

\subsection{Scope and elementary consistency bounds}

Let
\[
 {\cal G}_2=\{0,\pi\}^4,\qquad
 \langle f\rangle_2=\frac1{16}\sum_{k\in{\cal G}_2}f(k).
 \tag{12.1}
\]
At each point use the Taylor arithmetic, phase projection,
actual tangents and trace functions of V7.
The external momentum remains the formal variable \(p\);
it is not set to a grid momentum before differentiation.
Consequently (12.1) is a quadrature of the differentiated
thermodynamic integrand. It is not the side-two
lowest-nonzero-external-momentum response of f11 V99.

The following projection check is the same elementary mechanism
used in the archived neutral determinant calculation, applied
here to V7's small-momentum neutral weight.

\begin{proposition}[Pointwise bounds used by the grid verifier]
For every real internal momentum,
\[
 -\frac14\le[p^2]{\cal N}(k,p)\le0,
 \qquad {\cal T}_{11}(k,0)\ge0.
 \tag{12.2}
\]
\end{proposition}
\begin{proof}
Put \(B=DE\). The fibre bound implies that \(B\) has full
column rank. Therefore
\[
 {\cal R}=B(B^*B)^{-1}B^*
 \]
is an orthogonal projection on the 96 plaquette coordinates.
Let \(\Pi_{01}\) select the sixteen orientation-\(01\) rows.
Using \(P=(2B^*B)^{-1}\) and (11.8),
\[
 [p^2]{\cal N}=-\frac1{64}\Tr(\Pi_{01}{\cal R}).
 \]
Every selected diagonal entry of an orthogonal projection
lies in \([0,1]\). Thus the trace lies in \([0,16]\),
proving the neutral bound.
Each summand of \({\cal T}_{11}\) at real momentum is
\(\|P_h^{1/2}K_{h0}P_0^{1/2}\|_F^2\), so it is nonnegative.
This latter argument does not impose a sign on its second
Taylor coefficient.
\end{proof}

\subsection{The complete corrected \(2^4\) grid}

\begin{theorem}[Exact differentiated tensor-grid average]
\label{thm:f12v8-grid}
The exact averages of the V7 trace cards are
\[
\begin{array}{c|r|r}
 &\langle[p^0]\cdot\rangle_2&\langle[p^2]\cdot\rangle_2\\ \hline
{\cal T}_2&2657/192&10231/30720\\
{\cal T}_{11}&61519/3920&227198813033/284497920000\\
{\cal T}_2-{\cal T}_{11}
 &-87263/47040&-132449522033/284497920000\\
{\cal N}&0&-15/128\\
\text{unit measure}&7/3&19/48
\end{array}
\tag{12.3}
\]
All first-order coefficients vanish.
In particular the finite-grid diagnostic is
\[
 \boxed{\displaystyle
 \langle{\cal I}_2\rangle_2
      =-\frac{53175362033}{284497920000}
      \simeq-0.1869094932.}
 \tag{12.4}
\]
\end{theorem}
\begin{proof}[Exact finite sum with per-point matrix checks]
Evaluate the fifteen grid points other than V7's certified
\((0,\pi,0,0)\) point using its exact Taylor implementation.
Reuse the latter only after verifying its output hash,
source hash and all dependency hashes.
At every new point the implementation checks \(LE=0\),
the reverse first-matrix adjoints, second-matrix Hermiticity,
and \(HP=I\) through degree two for all three shifted
internal matrices.
It retains and checks the zero \(p^2\) contribution of the
V6 tangent correction.

All coefficients are rational and real. The grid wrapper
additionally checks (12.2), all zero linear terms, and
the exact identity
\({\cal T}_2-{\cal T}_{11}\) at each Taylor degree.
Add each of the sixteen rational cards and divide by sixteen
to obtain (12.3). Adding the measure row gives (12.4).
No floating-point inverse or external finite difference
is used. The complete card list and dependency record
are contained in the certificate below.
\end{proof}

For transparent reconstruction of the combined sum, write
\(\epsilon_\mu=e^{ik_\mu}\in\{+,-\}\). The sixteen values
of \({\cal I}_2\) are:
\[
\begin{array}{c|r|r}
(\epsilon_1,\epsilon_2,\epsilon_3)&\epsilon_0=+&\epsilon_0=-\\ \hline
(+,+,+)&541/3072&77784599/74088000\\
(+,+,-)&68953/176400&9695/13824\\
(+,-,+)&19333/44100&2417/3456\\
(+,-,-)&931/1536&57163/72000\\
(-,+,+)&-68509331/74088000&-104527/82944\\
(-,+,-)&-6259/6912&-417751/405000\\
(-,-,+)&-12809/13824&-840287/810000\\
(-,-,-)&-60323/72000&-945/1024
\end{array}
\tag{12.5}
\]
The mixed signs emphasize why V7's single negative card
did not determine even the coarse-grid average.
The rational neutral mean \(-15/128\) in (12.3) is at this
stage a finite-grid value; it is not asserted here to equal
the neutral torus integral.

\subsection{The exact constant-term discrepancy}

Let the full constant card be
\[
 {\cal I}_0(k)=\frac73+
           [p^0]({\cal T}_2-{\cal T}_{11}+{\cal N})(k,p).
 \tag{12.6}
\]
The thermodynamic Ward identity and V7's integral representation
give \(\int_{\mathbb T^4}{\cal I}_0(k)\,dk/(2\pi)^4=0\).
But (12.3) yields
\[
 \boxed{\displaystyle
 \langle{\cal I}_0\rangle_2=\frac{7499}{15680}
                                  \simeq0.4782525510.}
 \tag{12.7}
\]
Thus the corner rule has a known nonzero quadrature error
for the constant card. This is not a failure of the
integrated Ward identity.
It also does not supply an error bound, of either sign,
for the different function \({\cal I}_2\).
Subtracting (12.7) from (12.4) would have no proved justification.

\subsection{An independent Gaussian-Bloch check at all sixteen points}

For this subsection only, use \(\tau\) as a shift of the
\emph{internal} momentum in orientation one. Let
\[
 H^{\rm free}(k+\tau e_1)
     =H^{[0]}+H^{[1]}\tau+H^{[2]}\tau^2+O(\tau^3),
 \qquad P^{[0]}=(H^{[0]})^{-1},
\]
\[
 {\cal B}(k)
     =4\Tr(P^{[0]}H^{[2]})
          -2\Tr(P^{[0]}H^{[1]}P^{[0]}H^{[1]})
     =2\partial_{k_1}^2\log\det H^{\rm free}(k).
 \tag{12.8}
\]
Here \(H^{\rm free}=2(DE)^*(DE)\) is the free \(45\)-component
charged matrix. Formula (12.8) follows by differentiating
\(\log\det\); the factor four on \(H^{[2]}\) includes
the Taylor factorial. This calculation uses no cubic or
quartic background vertices and no multiplier.

\begin{proposition}[Finite-grid agreement with a pure Gaussian derivative]
\label{prop:f12v8-shift}
At each of the sixteen momenta of \({\cal G}_2\),
\[
 {\cal I}_0(k)={\cal B}(k).
 \tag{12.9}
\]
This statement is restricted to those points.
\end{proposition}
\begin{proof}[Independent exact calculation]
At each sign tuple replace its orientation-one phase
\(\epsilon_1\) by
\(\epsilon_1(1+i\tau-\tau^2/2)\).
Build \(D,E,H^{\rm free}\) directly and extract
\(H^{[0]},H^{[1]},H^{[2]}\).
Invert \(H^{[0]}\) by exact elimination, checking the inverse
by multiplication. Evaluate the right side of (12.8).
It equals the independently stored full constant card
(12.6) at every point, entry for entry in the certificate.
For example both sides are \(-11/24\) at \(k=0\),
\(7543/3675\) at \((0,\pi,0,0)\), and \(5/8\) at
\((\pi,\pi,\pi,\pi)\). All sixteen exact comparisons passed.
\end{proof}

The all-momentum version of (12.9) is a concrete structural
identity to investigate, not a theorem inferred from sixteen
samples. If established, it would explain the constant
Ward cancellation as the integral of a periodic derivative.
Independently of that proposed identity, the latter derivative
does integrate to zero: \(H^{\rm free}\) is smooth, periodic
and positive definite, so its real \(\log\det\) is smooth
and periodic. Integrate its second derivative over \(k_1\).
The agreement in (12.9) is therefore a useful check of the
background calculation against Gaussian data, but it neither
checks the full \(p^2\) integrand by this method nor certifies
its integration error.

\subsection{Provenance and the next analytic target}

\begin{record}[Completed exact grid]
The executed grid script is
\texttt{scripts/verify\_ym\_f12\_v8\_corner\_derivative\_grid.py},
SHA--256
\texttt{2AEFE011C7162D5A102481965349E65DB7B6D8BF23DE8BA36A9AF5F81B996D57}.
Its output is
\texttt{artifacts/ym\_f12\_v8\_corner\_derivative\_grid.json},
SHA--256
\texttt{7DFD8585CB8A5D284A97671E9958AC121D3D1AD6C45BEEBEB431CFD0AEB1C673}.
The calculation used two worker processes and retained
per-point checkpoints. It completed all sixteen cards;
no live unfinished calculation is being treated as a result.
Dependency hashes were checked again at completion.
The V7 seed was reused with output hash
\texttt{C7D37A0244A57ED9ADBAF9F92127BEEF5F3A86D5CA04CFDB30852BFAB768A394}.
\end{record}

\begin{record}[Independent constant-card check]
The second executable is
\texttt{scripts/verify\_ym\_f12\_v8\_flat\_shift\_check.py},
SHA--256
\texttt{EA978819FEA21BB98416B878896F67009836B6DC576F96A15CB396CB1E2A71F9}.
Its output is
\texttt{artifacts/ym\_f12\_v8\_flat\_shift\_check.json},
SHA--256
\texttt{1F5E33215EA0BC6C6EC1F011B68B118F8021941E019EA6E1BEE1BD755F3992CC}.
It verifies the grid certificate's dependencies, records
the grid certificate hash, and independently evaluates
(12.8) at all sixteen points. All comparisons passed.
It reuses only the free \(D,E,H\) construction and exact
linear algebra, not the higher background vertices whose
constant trace it checks.
\end{record}

\begin{decision}[Go upstream before trusting a denser unbounded grid]
The actual differentiated corner grid is now complete.
The next analytic target is the all-momentum flat-background
Bloch-shift identity suggested by (12.9), with its coordinate
Jacobian and normalization explicitly proved.
This offers a structural test of the cancellations before
spending further work on a denser grid without a useful
error certificate. It must not be assumed merely because
the corner checks succeeded.

The full \(p^2\) integral still needs a practical enclosure.
V5's very conservative global strip has not improved in
this version. Neither the negative diagnostic (12.4) nor
the positive discrepancy (12.7) proves the sign of
\(\theta\). The next compulsory companion audit and
ten-version direction review are at V10.
\end{decision}

\begin{record}[V8 append record]
V7 had SHA--256
\texttt{B046BD402E943F23D8E75C9638B7ADEAE94B77458D477072F204AA2850F87DC0}.
Its full body is retained before this appendix, apart from
the title version. The immediate active predecessor is
replaced by V8; all family archives remain.
\end{record}

\begin{firewall}[Position after V8]
Completed: the exact corrected \(2^4\) derivative grid and
an independent Gaussian check of its constant cards.
Not completed: an integration-error enclosure for
\(\theta\), the uniform interacting continuation, compact
continuum construction or the full Yang--Mills mass gap.
The original goal remains open.
\end{firewall}
\section{V9: the all-momentum flat Bloch identity and its fibre Jacobian}
\label{sec:f12v9}

V8 suggested a flat-background identity by sixteen exact
constant-card comparisons. This section proves it for every
real internal momentum and every sufficiently small flat
background amplitude, before differentiating.
The new ingredient is the determinant of the \emph{actual
restricted-fibre coordinate change}, with the Haar factors
retained. The full flat Wilson Hessian and its gauge covariance
were already proved in f11 V53; those facts are used here,
not presented as new results.

This is an identity for the order-one local coefficient.
It does not assert equality of finite-box integrals under a
nonperiodic gauge transformation, or control their interacting
remainders.

\subsection{The flat family and the Bloch shift}

Use \(SU(2)\), a quaternion generator \(T\) with
\(\|T\|^2=2\), and the constant coarse background
\[
 c_\mu(t)=tT\,{\bf1}_{\mu=1}.
 \tag{13.1}
\]
Its boundary-copy section has logarithm \(x_t=Rc(t)\):
the eight orientation-one boundary links per cell have
logarithm \(tT\), and the other non-tree links have zero
logarithm. All plaquettes are identity.
Moreover \(F(Rc)=c\) on the small section, since all sixteen
paths in a block have the same holonomy \(\exp(c_e)\).
Consequently
\[
 DF(x_t)R=I,\qquad X(0,c(t))=Rc(t),\qquad z_*(c(t))=0.
 \tag{13.2}
\]
The last assertion follows because the Wilson action is
nonnegative, vanishes on this flat section, and its small
fibre minimum is unique.

Choose the charged complex coordinate on which
\(\operatorname{ad}_T=2i\), and put
\[
 q(t)=\frac{e^{2it}-1}{2it}
      =e^{it}\frac{\sin t}{t},\qquad q(0)=1.
 \tag{13.3}
\]
Let \(Q_{\rm f}(t)\) be the diagonal \(49\)-by-\(49\)
matrix with entry \(q(t)\) on the eight orientation-one
boundary links and entry one elsewhere.
Let \(Q_{\rm c}(t)\) be diagonal on the four coarse
orientations, with entry \(q(t)\) in orientation one.
Then
\[
 Q_{\rm f}R=RQ_{\rm c},\qquad
 \det_{\mathbb C}Q_{\rm f}=q(t)^8,\qquad
 \det_{\mathbb C}Q_{\rm c}=q(t).
 \tag{13.4}
\]

\begin{proposition}[Exact charged block differential covariance]
\label{prop:f12v9-block}
Write \(L(k)=DF(0)(k)\), \(L_t(k)=DF(x_t)(k)\), and
\(k'=k+2t e_1\). For every real \(k\) and sufficiently
small real \(t\),
\[
 L(k')Q_{\rm f}(t)=Q_{\rm c}(t)L_t(k).
 \tag{13.5}
\]
\end{proposition}
\begin{proof}
On the infinite lattice apply the cell-constant root gauge
transformation \(g_r=\exp(tr_1T)\), using the convention
\(U_{xy}\mapsto g_xU_{xy}g_y^{-1}\).
It preserves all tree links and sends the flat background
(13.1) to identity.
For an exponential-coordinate variation at a link of
background logarithm \(a\), its right-trivialized derivative is
\[
 (\delta\exp)_a(w)\exp(-a)
       =\int_0^1 e^{s\operatorname{ad}_a}w\,ds .
 \]
On the charged component this is \(q(t)w\) when \(a=tT\)
and \(w\) when \(a=0\).
The further adjoint action of \(g_r\) multiplies a charged
cell amplitude by \(e^{2itr_1}\).
Thus a mode \(k\) is sent to mode \(k'\), with fine
coordinate multiplier \(Q_{\rm f}\).
The same calculation on coarse link logarithms gives
the multiplier \(Q_{\rm c}\).
Differentiate the root block's established gauge equivariance
at this background to obtain (13.5).

This is an identity of local Bloch symbols. On a finite
periodic torus, \(g_r\) need not be periodic. We do not
assert that it gauges away torus holonomy. The momentum
shift records precisely the resulting change of Bloch
boundary condition.
\end{proof}

\subsection{The actual fibre-coordinate determinant}

All matrices in this subsection act on the charged component.
Arrange the fine coordinates in the order of the 45 free
edges followed by the four boundary pivots.
Let \(L_{\rm f}(k)\) and \(L_{\rm p}(k)\) be the corresponding
column blocks of \(L(k)\), and \(R_{\rm f}\) the free rows of \(R\).
The pivot construction gives
\[
 L_{\rm p}=\tfrac18I_4,\qquad
 E(k)=\begin{pmatrix}I_{45}\\-8L_{\rm f}(k)\end{pmatrix},
 \qquad R=\begin{pmatrix}R_{\rm f}\\I_4\end{pmatrix}.
 \tag{13.6}
\]
In particular \(L(k)R=I_4\) implies
\(L_{\rm f}(k)R_{\rm f}=7I_4/8\). Hence
\[
 \det_{\mathbb C}[E(k),R]
       =\det(I_4+8L_{\rm f}(k)R_{\rm f})=8^4,
 \tag{13.7}
\]
independently of \(k\). A different fixed row ordering changes
only a constant sign, which cancels in the ratios below.

The actual fibre tangent at \(x_t\), from (13.2), is
\[
 Z_t(k)=D_zX(0,c(t))(k)=E(k)-RL_t(k)E(k).
 \tag{13.8}
\]
Let \(S_{\rm f}\) select the 45 free rows of a fine vector.
The coordinate extractor on \(\ker L(k')\) is
\[
 {\cal E}(k')=S_{\rm f}-R_{\rm f}L(k'),\qquad
 {\cal E}(k')E(k')=I,\qquad {\cal E}(k')R=0.
 \tag{13.9}
\]

\begin{theorem}[Restricted charged-fibre Jacobian]
\label{thm:f12v9-jacobian}
There is an invertible \(45\)-by-\(45\) matrix
\[
 J_t(k)={\cal E}(k')Q_{\rm f}(t)E(k)
 \]
such that
\[
 Q_{\rm f}(t)Z_t(k)=E(k')J_t(k),\qquad
 \boxed{\det_{\mathbb C}J_t(k)=q(t)^7.}
 \tag{13.10}
\]
The determinant is independent of internal momentum.
\end{theorem}
\begin{proof}
Equations (13.5) and \(L_tZ_t=0\) show that
\(Q_{\rm f}Z_t\) lies in \(\ker L(k')\).
This kernel has basis \(E(k')\). Apply (13.9) and use
\(Q_{\rm f}R=RQ_{\rm c}\) to obtain the stated formula
for \(J_t\). It is invertible: \(Q_{\rm f}\) is invertible
for small \(t\), \(Z_t\) parametrizes \(\ker L_t\), and
(13.5) maps that kernel isomorphically onto \(\ker L(k')\).

The column operation (13.8) gives
\(\det[Z_t,R]=\det[E(k),R]\).
On the other hand,
\[
 Q_{\rm f}[Z_t,R]
       =[E(k'),R]\begin{pmatrix}J_t&0\\0&Q_{\rm c}\end{pmatrix}.
 \]
Take determinants and apply (13.4) and (13.7).
The common factor \(8^4\) cancels, leaving
\(q^8=q\det J_t\). This proves (13.10).
In particular one must not use the eight fine boundary
factors without dividing out the one coarse factor.
\end{proof}

\subsection{Determinant congruence and exact Haar cancellation}

Let
\[
 H^{\rm free}(k)=2(D(k)E(k))^*D(k)E(k)
 \]
be the charged matrix used in V4--V8.
Let \(H_t^{\rm ch}(k)\) be the actual constrained charged
fibre Hessian at the minimum (13.1).

\begin{proposition}[Actual flat-fibre Hessian]
\label{prop:f12v9-hessian}
For real \(k,t\) in the stated range,
\[
 H_t^{\rm ch}(k)=J_t(k)^*H^{\rm free}(k')J_t(k),
 \tag{13.11}
\]
\[
 \log\det H_t^{\rm ch}(k)
       =\log\det H^{\rm free}(k')
                      +14\log\!\left(\frac{\sin t}{t}\right).
 \tag{13.12}
\]
The neutral Hessian is \(H^{\rm free}(k)\), independent of \(t\).
\end{proposition}
\begin{proof}
Gauge invariance identifies the Wilson quadratic form at
the flat background with the free form at the shifted Bloch
momentum. The full first variation of the action is zero
at a flat background, so a second derivative of the coordinate
map contributes no additional term to its Hessian.
Restrict the fine quadratic form using (13.10). This gives
(13.11), which is also the restricted version of the
archived flat transported-curl Gram identity.
Taking its determinant and using
\(|q(t)|=\sin(t)/t>0\) for sufficiently small \(t\)
gives (13.12).
The neutral component has zero adjoint charge, so its
fine and coarse coordinate multipliers are one and its
momentum is unshifted. In (13.8)--(13.10) its tangent
is \(E(k)\) and its coordinate change is identity.
\end{proof}

The exact reduced density has no \(DFR\) contribution here,
by (13.2). For the \(SU(2)\) generator in use,
\(j_G(tT)=(\sin t/t)^2\). There are eight nonzero fine
link logarithms and one nonzero coarse link logarithm per cell.
Consequently the density and coarse-Haar variations per cell are
\[
 \ell_t-\ell_0=16\log(\sin t/t),\qquad
 \log j_{{\rm c},t}-\log j_{{\rm c},0}=2\log(\sin t/t).
 \tag{13.13}
\]
The normalized measure variation is their combination
\(-14\log(\sin t/t)\), canceling (13.12) exactly.

\begin{theorem}[All-momentum normalized flat identity]
\label{thm:f12v9-flat}
Distribute the local measure term in (13.13) equally over
the internal momentum integral, as in V7. If
\(\delta\widehat q_t(k)\) denotes the change of the
normalized order-one integrand along (13.1), then
\[
 \boxed{\delta\widehat q_t(k)
       =\log\det H^{\rm free}(k+2t e_1)
                         -\log\det H^{\rm free}(k).}
 \tag{13.14}
\]
In particular the formerly grid-tested identity is now
proved for every real momentum:
\[
 \boxed{{\cal I}_0(k)
                =2\partial_{k_1}^2\log\det H^{\rm free}(k).}
 \tag{13.15}
\]
\end{theorem}
\begin{proof}
The real charged fibre determinant contributes
\(\log\det_{\mathbb C}H_t^{\rm ch}\) to
\(\tfrac12\log\det_{\mathbb R}H_t\); the two real charged
components account for this factor.
The neutral determinant does not vary.
Combine (13.12) and (13.13) to obtain (13.14).

At zero external momentum the auxiliary-phase prescription
of V7 averages a constant background amplitude by
\(\frac14\sum_s(\operatorname{Re}i^s)^2=1/2\).
Thus its norm-squared-two real-cosine constant card is
\(\tfrac12\partial_t^2\delta\widehat q_t(k)|_{t=0}\).
This is the unit-polarization normalization of \({\cal I}_0\).
Differentiating the shift \(k_1\mapsto k_1+2t\) gives
\(\tfrac12\cdot4=2\), proving (13.15).
This normalization is not an additional charged-colour factor.
\end{proof}

\subsection{What the finite-grid discrepancy measures}

Put \(f(k)=\log\det H^{\rm free}(k)\).
The infinite-volume flat variation vanishes exactly at
this order:
\[
 \int_{\mathbb T^4}\delta\widehat q_t(k)\frac{dk}{(2\pi)^4}
       =0.
 \tag{13.16}
\]
This follows from translation invariance of the integration
measure, without estimating or subtracting a remainder.
It is not a construction of the full infinite-volume measure.

For a side-\(L\) torus, write
\({\cal G}_L=(2\pi/L)\{0,\ldots,L-1\}^4\).
The actual flat order-one coefficient instead obeys
\[
 \frac{\widehat Q_L(c(t))-\widehat Q_L(0)}{L^4}
       =\frac1{L^4}\sum_{k\in{\cal G}_L}
                              (f(k+2t e_1)-f(k)).
 \tag{13.17}
\]
The right side need not vanish: a continuous momentum shift
does not generally permute this grid. This is consistent
with the nonperiodicity of the gauge used in (13.5).

\begin{proposition}[Exact finite-volume Fourier alias formula]
\label{prop:f12v9-alias}
Write the absolutely convergent Fourier expansion
\(f(k)=\sum_{r\in\mathbb Z^4} f_r e^{ik\cdot r}\).
Then
\[
 \frac{\widehat Q_L(c(t))-\widehat Q_L(0)}{L^4}
       =\sum_{\ell\in\mathbb Z^4}
                   f_{L\ell}(e^{2itL\ell_1}-1),
 \tag{13.18}
\]
\[
 \frac1{L^4}\sum_{k\in{\cal G}_L}{\cal I}_0(k)
       =-2L^2\sum_{\ell\ne0}\ell_1^2 f_{L\ell}.
 \tag{13.19}
\]
In particular the nonzero V8 value \(7499/15680\) is the
side-two alias sum (13.19), not a nonzero bulk constant term.
\end{proposition}
\begin{proof}
The finite fibre stencil and the gap give the nonvanishing
complex strip proved in V5. On a smaller strip the determinant
has a periodic holomorphic logarithm that agrees with the real
\(\log\det\) on the real torus.
Indeed its winding along each fundamental real torus cycle is
zero, since that determinant is positive on the real torus;
the complex strip retracts onto those cycles. Local logarithms
therefore continue without monodromy to a periodic logarithm.
On a still smaller closed strip \(\|\operatorname{Im}k\|_\infty
\le s\), let \(B_f\) bound its absolute value.
Contour shifting gives \(|f_r|\le B_f e^{-s|r|_1}\).
This also justifies two termwise derivatives.

The grid average of \(e^{ik\cdot r}\) vanishes unless
\(r=L\ell\). Substitution in (13.17) proves (13.18).
Take half the second \(t\)-derivative and use (13.15)
to obtain (13.19).
\end{proof}

For clarity, these existence bounds yield the explicit
conditional estimate, with \(a_L=e^{-sL}\),
\[
 \left|\frac1{L^4}\sum_{k\in{\cal G}_L}{\cal I}_0(k)\right|
 \le
 4B_f L^2\frac{a_L(1+a_L)^4}{(1-a_L)^6}.
 \tag{13.20}
\]
To verify it, bound (13.19) absolutely and use
\(\sum_{j\in\mathbb Z}a^{|j|}=(1+a)/(1-a)\) and
\(\sum_{j\in\mathbb Z}j^2a^{|j|}
       =2a(1+a)/(1-a)^3\).
No practically sharp numerical \(s,B_f\) is asserted here.
Equation (13.20) controls the constant-card alias, not
the quadrature error in \({\cal I}_2\).

\subsection{Non-corner full-matrix checks and scope}

\begin{record}[Executed amplitude-jet regression]
The new verifier is
\texttt{scripts/verify\_ym\_f12\_v9\_flat\_bloch\_identity.py},
SHA--256
\texttt{7F2527867911E2BAF3FF7A4AB855AA00B763929811CB17E79B4E19EAC2FD67A8}.
Its output is
\texttt{artifacts/ym\_f12\_v9\_flat\_bloch\_identity.json},
SHA--256
\texttt{AE9D14B602F56D48FA1CD4F69F7C495407A7AD103657FA280E93A05CAEE20679}.
It records the hashes of all five inherited dependencies.

The two tested internal phase tuples are
\[
 (3/5+4i/5,\ 5/13+12i/13,\ -1,\ i),\qquad
 (-i,\ 3/5-4i/5,\ 5/13+12i/13,\ 1).
 \]
These are unit-modulus continuous Bloch phases, not a
claimed finite-torus grid.
In exact amplitude Taylor arithmetic through \(t^2\),
the code constructs
\[
 e^{2it}=1+2it-2t^2,\qquad
 q(t)=1+it-\tfrac23t^2.
 \]
It verifies \(L_{\rm p}=I/8\), \(LR=I\), the full
block covariance (13.5), \(Q_{\rm f}Z_t=E(k')J_t\),
and the full charged Hessian congruence (13.11).
The actual block and Hessian jets on the opposite sides
are obtained from V7's root-path and action-jet assembly.
Thus the new checks compare entire matrices, not just
the traces compared at V8's corners.

The coordinate-change jets give
\[
 [t]\log\det J_t=7i,\qquad
 [t^2]\log\det J_t=-7/6,
 \]
as required by \(7\log q(t)\).
At both non-corner points, the actual normalized constant
card also agrees exactly with (13.15).
All assertions passed. The all-momentum and all-small-\(t\)
theorems above are algebraic and geometric proofs;
they are not inferred from these two tests.
\end{record}

\begin{assessment}[What is now closed, and what is not]
The identity left conjectural in V8 is proved with its
actual fibre-coordinate Jacobian. The constant Ward
cancellation occurs after the momentum integral, and
the side-two nonzero constant average is explained by
finite-volume Fourier aliases.
This strengthens the validation of the differentiated
root-chart calculation.

It does not make the negative \(p^2\) corner average
equal to \(\theta\). In particular subtracting V8's
constant discrepancy from its quadratic average remains
unjustified. The flat background here has zero curvature;
the transverse \(p^2\) response involves a varying
background and still requires its own integration bound.
The physical mass-gap question remains separate from
both this flat identity and the eliminated-fibre gap.
\end{assessment}

\begin{decision}[V10 audit and the remaining integration bottleneck]
At V10 inspect the updated SM, NS and parallel YM notes
and review the direction of the last ten-version interval.
Retain (13.14)--(13.15) as exact normalization regressions.
Then focus on a practically useful enclosure of the
nonflat \(p^2\) integral, not on another unbounded
extrapolation of corner averages.
The uniform interacting remainder and compact completion
must remain explicit prerequisites for any iterative
continuum or mass-gap conclusion.
\end{decision}

\begin{record}[V9 append record]
V8 had SHA--256
\texttt{709229F53B4BD29BC058DE46E08C8656C93FB9DCD465367B746866A544AD9FE1}.
Its full body is retained, apart from the title version,
before this appendix. Only the active predecessor is
replaced by V9. Every family archive remains.
\end{record}

\begin{firewall}[Position after V9]
Proved: the all-momentum normalized flat Bloch identity,
the exact restricted-fibre Jacobian, and the finite-volume
alias formula. Checked: full block and Hessian amplitude
jets at two non-corner momenta.
Still open: the certified nonflat infrared coefficient,
uniform interacting continuation, compact continuum
construction and the full Yang--Mills mass gap.
\end{firewall}
\section{V10: audit and a uniform one-loop comparison for a convex extension}
\label{sec:f12v10}

\subsection{Companion audit and direction review}

The mandatory checkpoint found the same active companion files:
\[
\begin{array}{l|r|l}
\text{file}&\text{bytes}&\text{SHA--256 prefix}\\ \hline
\text{SM--Einstein V72}&607372&\texttt{F44F9745D43379E1}\\
\text{NS Stability V26}&278283&\texttt{82C887F8C2C69E4F}\\
\text{Parallel YM V5}&54174&\texttt{306F1BB84CFA8A8D}
\end{array}
\]
Their full hashes agree with the V5 audit.
Their endings were inspected again. SM retains its fixed-mode,
fixed-time many-copy result and leaves growing cutoff open.
NS retains its pointwise Oseen-norm results, not a global
stability theorem. Parallel YM retains the obstruction to
volume-uniform tails for extensive total charge. No new
companion result supplies the missing YM remainder or gap.

The compact lineage has obtained the lowest-momentum estimator,
the actual harmonic lift and measure symbol, corrected
off-corner tangents, an exact off-corner response, a differentiated
implementation, a complete corner diagnostic, and the
all-momentum flat Bloch identity. Its coefficient integral
still has no practical error enclosure.
Even resolving that numerical issue would leave the
volume-uniform interacting remainder missing.
The archived f11 V10 already developed frozen preconditioners
for an earlier shell problem; merely repeating that method
without a useful new enclosure is not the next result here.

We therefore go upstream to the interacting comparison.
The new theorem below is for a precisely defined
\emph{auxiliary whole-space convex extension} of the
centred local action. It agrees with the actual action
and density near their minimum. It is not the original
fixed-box integral or the full compact fibre measure.
Transfer back to those objects remains an explicit obligation.
The physical model and the full mass-gap goal are not redefined.

\subsection{Centred local data and parameter locality}

Use a finite periodic cell set \(\Lambda\), with its graph
distance \(d\) and a fixed number of fine, fibre and coarse
indices per cell. Write \(N=45(\dim G)|\Lambda|\).
All constants below are independent of \(\Lambda\) and
of the small coarse field \(c\), for a fixed group and chart.
The actual local functions have bounded-range decompositions
\[
 f(z,c)=\sum_x f_x(z,c),\qquad
 \ell(z,c)=\sum_x\ell_x(z,c).
 \]
Each term uses a bounded number of coordinates in a fixed
neighbourhood \(S_x\). Enlarge these neighbourhoods uniformly
when necessary. The analytic small chart supplies bounded
derivatives through the orders used below.

Let \(z_*(c)\) be the uniform small minimum, put
\[
 u=z-z_*(c),\quad G(c)=f(z_*(c),c),\quad
 H_c=f_{zz}(z_*(c),c),\quad \ell_*(c)=\ell(z_*(c),c).
 \]
The inherited bounds give \(H_c\ge hI\), a finite interaction
range, bounded row sums, and
\(\|(H_c^{-1})\|_{\rm row}\le K\), uniformly.
The matrix \(H_c\) is local in its \emph{fibre indices};
its coefficients may depend on distant \(c\) through \(z_*\).

For three index locations, let \(\tau(x,a,b)\) be the minimum
total length of a tree connecting them.

\begin{proposition}[Locality after centring at the moving minimum]
\label{prop:f12v10-centre}
For some \(C,\nu>0\),
\[
 |\partial_a z_{*,i}|\le Ce^{-\nu d(i,a)},\qquad
 |\partial_a\partial_b z_{*,i}|
                 \le Ce^{-\nu\tau(i,a,b)}.
 \tag{14.1}
\]
Consequently the first and second \(c\)-derivatives of any
fixed local centred Taylor coefficient at \(x\) have bounds
of the form
\[
 C e^{-\nu d(x,a)},\qquad
 C e^{-\nu\tau(x,a,b)},
 \tag{14.2}
\]
respectively. Bounds hold through the fibre derivative
orders needed for a third-order Taylor remainder and its
first two parameter derivatives.
\end{proposition}
\begin{proof}
Differentiate \(f_z(z_*(c),c)=0\):
\(\partial_a z_*=-H_c^{-1}f_{zc_a}\).
The source has bounded support near \(a\), and the inverse
has uniform exponential decay by f11 V81. This proves
the first estimate.
Differentiating once more expresses
\(\partial_a\partial_bz_*\) as \(-H_c^{-1}\) times the sum
of \(f_{zc_ac_b}\), the two terms \(f_{zzc}\partial z_*\),
and \(f_{zzz}[\partial_a z_*,\partial_b z_*]\).
All these tensors are local in their uncontracted indices
and uniformly bounded. The products of inverse kernels
and first responses form paths connecting \(i,a,b\).
Their total length is at least the connecting-tree length.
Reserve half of each exponential for the sums over internal
indices; bounded polynomial growth of lattice balls makes
those sums uniform in volume. Reducing \(\nu\) proves the
second estimate.
The chain rule applied to a local coefficient at \(x\)
gives (14.2) by the same argument.
The required higher local derivatives exist and are bounded
on the smaller analytic chart. No new uniform remainder
estimate is assumed in this step.
\end{proof}

Subtract the Taylor polynomial of each \(f_x\) through
degree two in \(u\), denoting the remainder by \(r_x(u,c)\).
Put \(d_x(u,c)=\ell_x(z_*+u,c)-\ell_x(z_*,c)\).
Stationarity cancels the sum of the linear action terms, so
\[
 f(z_*+u,c)-G(c)=\tfrac12u^TH_cu+\sum_xr_x(u,c),
 \qquad \ell(z_*+u,c)-\ell_*(c)=\sum_xd_x(u,c).
 \tag{14.3}
\]
On the uniform local chart, \(r_x=O(|u_{S_x}|^3)\)
and \(d_x=O(|u_{S_x}|)\), including the parameter
derivatives with weights (14.2).

\subsection{The auxiliary extension is defined explicitly}

Choose a fixed smooth even cutoff \(\chi\), equal to one
on \([-1,1]\), equal to zero outside \((-2,2)\), with
values in \([0,1]\). For sufficiently small fixed
\(\varepsilon>0\), set
\[
 \chi_x(u)=\prod_{i\in S_x}\chi(u_i/\varepsilon),\qquad
 r_x^\varepsilon=\chi_xr_x,\qquad
 d_x^\varepsilon=\chi_xd_x .
 \tag{14.4}
\]
The products are extended by zero outside their support.
Choose \(2\varepsilon\) inside the original local chart,
so the products are smooth globally.
The cutoff does not depend on \(c\) in the centred coordinates.
Define, on all of \(\mathbb R^N\),
\[
 R_\beta(u,c)=\sum_xR_x(u,c),\qquad
 R_x=\beta r_x^\varepsilon-d_x^\varepsilon,\qquad
 V_s(u,c)=\tfrac{\beta}{2}u^TH_cu+sR_\beta(u,c),
 \quad 0\le s\le1.
 \tag{14.5}
\]

\begin{proposition}[Uniform convexity of the interpolation]
\label{prop:f12v10-convex}
One can fix \(\varepsilon>0\), and then \(\beta_0<\infty\),
independently of volume, so that for all \(\beta\ge\beta_0\),
\[
 D_u^2V_s\ge m\beta I,\qquad
 \sup_i\sum_{j\ne i}|(V_s)_{ij}|\le K_1\beta
 \tag{14.6}
\]
for some \(m>0,K_1<\infty\).
The Hessian has a fixed finite range.
\end{proposition}
\begin{proof}
Taylor's formula and cutoff differentiation give, globally,
\[
 |\nabla r_x^\varepsilon|\le C|u_{S_x}|^2
                     \le C\varepsilon|u_{S_x}|
       \quad\hbox{on its support},\qquad
 \|D^2r_x^\varepsilon\|\le C\varepsilon.
 \]
Off that support these derivatives vanish.
Similarly \(|\nabla d_x^\varepsilon|\le C\) and
\(\|D^2d_x^\varepsilon\|\le C_\varepsilon\).
Bounded overlap gives the same row-sum estimates for
their sums. Choose \(C\varepsilon\le h/4\) and then
\(C_\varepsilon\le\beta h/4\). The lower bound in
(14.6) follows with \(m=h/2\); the upper off-diagonal
row bound and finite range follow from locality.
We may further shrink \(\varepsilon\) for the absorption
argument below, without changing any conclusion.
\end{proof}

Let \(\mu_s\) have density \(Z_s^{-1}e^{-V_s}\) on
\(\mathbb R^N\). Its Gaussian endpoint \(s=0\) is exact,
with covariance \((\beta H_c)^{-1}\); there is no moving
boundary in this auxiliary interpolation.

\subsection{Uniform moments and connected three-point bounds}

\begin{proposition}[Moments at the centred origin]
\label{prop:f12v10-moments}
For each fixed integer \(p\ge1\),
\[
 \sup_{s,c,i}\mathbb E_{\mu_s}|u_i|^p\le C_p\beta^{-p/2}.
 \tag{14.7}
\]
For local observables with gradient supports \(A,B\),
\[
 |\operatorname{Cov}_{\mu_s}(F,G)|
 \le \frac C\beta e^{-\alpha d(A,B)}
                    \|\nabla F\|_2\|\nabla G\|_2
 \tag{14.8}
\]
with uniform \(C,\alpha>0\).
\end{proposition}
\begin{proof}
The Poincare and covariance proofs of f11 V85--V86 apply
with (14.6). Their whole-space versions follow either
by the same vector-operator proof with no faces, or by
exhausting \(\mathbb R^N\) by boxes: the constants are
independent of box size, and the confining quadratic
bound permits passage to the limit at each finite \(N\).
This proves (14.8) and
\(\operatorname{Var}(u_i)\le(m\beta)^{-1}\).
The Poincare estimate also holds after any linear tilt,
since its Hessian is unchanged. Therefore the second
derivative of the one-coordinate log moment-generating
function is at most \((m\beta)^{-1}\). Integrating twice
gives a sub-Gaussian bound about its mean.

It remains to locate that mean uniformly, rather than
estimate a full \(N\)-dimensional distance.
Whole-space integration by parts gives
\[
 \beta H_c\,\mathbb E u
       =-s\beta\,\mathbb E\nabla\sum_xr_x^\varepsilon
                         +s\,\mathbb E\nabla\sum_xd_x^\varepsilon.
 \]
With \(M=\max_i|\mathbb E u_i|\), the inverse row bound
and the linear gradient estimate in the preceding proof give
\[
 M\le KC\varepsilon\bigl(M+(m\beta)^{-1/2}\bigr)+KC/\beta.
 \]
Choose \(KC\varepsilon\le1/2\) and absorb \(M\).
Then \(M\le C\beta^{-1/2}\), and the sub-Gaussian bound
gives (14.7). Using the quadratic gradient estimate
after this step even gives \(M\le C/\beta\), though
that refinement is not needed below.
All integrations by parts are justified by Gaussian
confinement at each finite dimension; the estimates
just derived are uniform in that dimension.
\end{proof}

We state the connected estimate in the scaling actually used.
For local scalar functions \(X,Y,Z\) supported near
cells \(x,y,z\), suppose their centred \(L^4\) norms and
gradient \(L^4\) norms satisfy
\[
 \begin{array}{c|cc}
 &\|F-\mathbb EF\|_4&\|\nabla F\|_4\\ \hline
 X,Y&C&C\sqrt\beta\\
 Z&C\beta^{-1/2}&C
 \end{array}
 \tag{14.9}
\]
All support sizes and radii are bounded.

\begin{proposition}[A summable connected three-point estimate]
\label{prop:f12v10-third}
Under (14.9),
\[
 |\operatorname{Cum}_{\mu_s}(X,Y,Z)|
 \le C\beta^{-1/2}
     e^{-\alpha' \operatorname{diam}\{x,y,z\}} .
 \tag{14.10}
\]
\end{proposition}
\begin{proof}
The third cumulant equals
\(\operatorname{Cov}(X,(Y-\mathbb EY)(Z-\mathbb EZ))\).
Holder's inequality bounds the \(L^2\) gradient norm
of the product by
\(\|\nabla Y\|_4\|Z-\mathbb EZ\|_4+
  \|Y-\mathbb EY\|_4\|\nabla Z\|_4\le C\).
Apply (14.8): the result is at most
\(C\beta^{-1/2}e^{-\alpha\min(d(x,y),d(x,z))}\),
with a changed constant for support radii.
Interchanging the three roles gives the analogous
isolation bound for each centre. For the split with
\(Z\) first, its gradient is \(O(1)\) and the product
gradient is \(O(\sqrt\beta)\), giving the same prefactor.
If two centres realize diameter \(D\), at least one
of them is at distance at least \(D/2\) from the third.
Choose that isolation bound. This proves (14.10).
No independence or unproved higher-cumulant theorem is used.
\end{proof}

\subsection{The derivative comparison}

Put \(\Phi_s(c)=-\log Z_s(c)\).
The centred local Taylor estimates and (14.7) imply the
following bounds for local pieces. Each \(R_x\), or its
first or second \(c\)-derivative after removing the weights
in (14.2), has \(L^4\) norm \(O(\beta^{-1/2})\) and
gradient \(L^4\) norm \(O(1)\).
Each local piece of \(V_{s,a}\) or \(V_{s,ab}\), after
removing its parameter weight, has \(L^4\) norm \(O(1)\)
and gradient \(L^4\) norm \(O(\sqrt\beta)\).
For example, \(\beta |u|^3\) and \(|u|\) have the first
scaling, whereas \(\beta |u|^2\) has the second.
The moment orders needed are finite (at most twelve here).

\begin{theorem}[Uniform first and second derivatives of the interpolation error]
\label{thm:f12v10-error}
For all finite volumes, small \(c\), and \(\beta\ge\beta_0\),
\[
 \frac{|\Phi_1-\Phi_0|}{N}\le C\beta^{-1/2},\qquad
 |\partial_a(\Phi_1-\Phi_0)|\le C\beta^{-1/2},
 \tag{14.11}
\]
\[
 |\partial_a\partial_b(\Phi_1-\Phi_0)|
        \le C\beta^{-1/2}e^{-\gamma d(a,b)} .
 \tag{14.12}
\]
In particular the Hessian error has row and operator
norm at most \(C\beta^{-1/2}\), without volume normalization.
\end{theorem}
\begin{proof}
At fixed finite dimension differentiation under the
integral is justified by uniform convex confinement.
The identities
\[
 \partial_s\Phi_s=\mathbb E_sR_\beta,\qquad
 \partial_a\Phi_s=\mathbb E_sV_{s,a},\qquad
 \partial_a\partial_b\Phi_s
       =\mathbb E_sV_{s,ab}
                         -\operatorname{Cov}_s(V_{s,a},V_{s,b})
 \]
give
\[
 \partial_s\partial_a\Phi_s
       =\mathbb E_sR_a-\operatorname{Cov}_s(V_{s,a},R_\beta),
 \tag{14.13}
\]
\[
 \begin{split}
 \partial_s\partial_a\partial_b\Phi_s={}&
       \mathbb E_sR_{ab}
       -\operatorname{Cov}_s(V_{s,ab},R_\beta)\\
 &-\operatorname{Cov}_s(R_a,V_{s,b})
       -\operatorname{Cov}_s(V_{s,a},R_b)\\
 &+\operatorname{Cum}_s(V_{s,a},V_{s,b},R_\beta).
 \end{split}
 \tag{14.14}
\]
Here \(R_a=\partial_aR_\beta\), and likewise for \(R_{ab}\).
For completeness, the last sign follows from
\(\partial_s\operatorname{Cov}_s(X,Y)
 =\operatorname{Cov}_s(X_s,Y)+\operatorname{Cov}_s(X,Y_s)
       -\operatorname{Cum}_s(X,Y,R_\beta)\).

Expand every quantity into its local pieces before
estimating. The expectation of each \(R_x\) is
\(O(\beta^{-1/2})\), proving the intensive estimate
after summing the \(O(N)\) pieces and integrating in \(s\).
The weights \(e^{-\nu d(x,a)}\), together with (14.8),
make both terms in (14.13) summable with bound
\(C\beta^{-1/2}\).

In (14.14), the two-parameter weight
\(e^{-\nu\tau(x,a,b)}\) bounds the first two terms by
\(C\beta^{-1/2}e^{-\gamma d(a,b)}\), after summing
the local centres. The next two covariance terms are
bounded by convolutions of the two first-parameter
weights and the covariance exponential; they have
the same bound.
For the last term use (14.10) on a local piece near \(x\)
of \(V_{s,a}\), one near \(y\) of \(V_{s,b}\), and
\(R_z\), then sum their parameter weights.
These sums are uniform: for example
\[
 \sum_z e^{-\alpha'\operatorname{diam}\{x,y,z\}}
       \le C e^{-\alpha'd(x,y)/2},
 \]
by reserving half the exponential for summation in \(z\).
Convolving with \(e^{-\nu d(x,a)}\) and
\(e^{-\nu d(y,b)}\), and decreasing the exponent if needed,
gives the claimed decay in \(d(a,b)\).
Uniform polynomial growth of the four-dimensional lattice
controls every such convolution, also on periodic quotients.

Integrate (14.13)--(14.14) over \(s\in[0,1]\).
Finally sum (14.12) over \(b\); the exponential row sum
is uniform. The symmetric Hessian's operator norm is
bounded by its maximum absolute row sum.
\end{proof}

\subsection{What the comparison proves for this extension}

Define the auxiliary normalized action by
\[
 \widehat W_\beta^{\rm ext}(c)
       =\beta G(c)-\ell_*(c)-\log Z_1(c)+\log j_{\rm c}(c).
 \tag{14.15}
\]
Since \(Z_0=(2\pi/\beta)^{N/2}(\det H_c)^{-1/2}\),
its Gaussian comparison is exactly
\[
 \beta G(c)+\frac N2\log(\beta/2\pi)+\widehat Q(c),
 \qquad
 \widehat Q=\tfrac12\log\det H_c-\ell_*+\log j_{\rm c}.
 \tag{14.16}
\]

\begin{corollary}[A volume-uniform order-one Hessian limit for the auxiliary action]
\label{cor:f12v10-extension}
With the extension (14.4)--(14.5),
\[
 \left\|D_c^2\widehat W_\beta^{\rm ext}
       -\beta D_c^2G-D_c^2\widehat Q\right\|_{\rm op}
                       \le C\beta^{-1/2}.
 \tag{14.17}
\]
The intensive action and local first-derivative errors
obey (14.11).
\end{corollary}
\begin{proof}
Subtract (14.16) from (14.15). The difference is
exactly \(\Phi_1-\Phi_0\). Apply the preceding theorem.
This proves the coefficient limit at order one, rather
than only an error small relative to the leading factor
\(\beta\). It does not prove an all-orders expansion.
\end{proof}

\begin{assessment}[The transfer boundary is not removed]
For \(\|u\|_\infty\le\varepsilon\), every cutoff is one:
the extended action and density agree exactly with
the actual centred local action and density.
Outside this neighbourhood they are deliberately changed.
The extension is not asserted to preserve the nonlinear
local gauge action, reflection positivity, or compact
sector weights.

The moment proof gives a small tail for each fixed
coordinate outside this neighbourhood, of order
\(e^{-c\beta\varepsilon^2}\) for large \(\beta\).
It does not justify dropping the union of all such
events in arbitrarily large volume: a union bound
contains \(N\). Nor does it bound two coarse derivatives
of the logarithm of a ratio of the original and extended
partition functions.
Therefore (14.17) is \emph{not yet} the corresponding
theorem for the original fixed-box \(W_\beta\).
That localized comparison, and then the full compact-fibre
comparison, remain to be proved.
\end{assessment}

\begin{decision}[Next acquisition after the V10 review]
Retain the differentiated-integral problem, but next
test whether the convex-extension estimate can be
transferred to the original local cutoff using localized
boundary or partition-ratio derivatives. Do not infer
that transfer from agreement of germs or from a
volume-dependent rare-event union bound.
The connected three-point estimate (14.10) is available
for that investigation.
The next compulsory companion audit is V15; the next
ten-version direction review is V20.
\end{decision}

\begin{record}[V10 append and verification record]
V9 had SHA--256
\texttt{387675166F5AA04A88575B325A9466568CE70FAFB9E942077CDA2492D98C7B1E}.
Its body is retained before this appendix, with only
the title version changed. The active predecessor is
replaced by V10; every family archive remains.
This version's new results are analytic proofs, not
new numerical cards. Source preservation and LaTeX
compilation are document checks, not substitutes for
the mathematical arguments above.
\end{record}

\begin{firewall}[Position after V10]
Proved: a volume-uniform order-one Hessian comparison
for the explicitly defined convex extension, using
local moments and summable connected three-point bounds.
Not proved: its transfer to the original fixed box or
compact fibre, a certified value of the infrared
coefficient, the continuum construction, or the full
Yang--Mills mass gap. The original goal remains open.
\end{firewall}
\section{V11: exponentially small comparison on the same fixed box}
\label{sec:f12v11}

V10 proved the uniform one-loop comparison for a whole-space
convex extension, not the original cutoff integral.
We separate that transfer into two operations.
This section compares the actual and extended potentials
on the \emph{same original fixed box}, obtaining exponentially
small errors through two coarse derivatives.
It leaves removal of the extension's box restriction as
a separate, precisely identified partition-ratio problem.
No moving-domain derivative is suppressed.

\subsection{Two potentials, one domain}

Let \(\mathcal D=[-r,r]^N\) be the fixed sufficiently small
fibre-coordinate box used in the actual local integral.
Restrict the coarse chart so \(\|z_*(c)\|_\infty\le r/2\).
Choose V10's \(\varepsilon\) still smaller, if necessary,
so its central neighbourhood lies inside \(\mathcal D\).
In the original \(z\) coordinates define
\[
 A_\beta(z,c)=\beta f(z,c)-\ell(z,c),\qquad
 E_\beta(z,c)=\beta G(c)-\ell_*(c)
                         +V_1(z-z_*(c),c),
 \tag{15.1}
\]
where \(V_1\) is exactly the extension of (14.5).
On \(\mathcal D\), their difference is
\[
 \Delta_\beta=A_\beta-E_\beta
  =\sum_x\Delta_x,\qquad
 \Delta_x=(1-\chi_x(u))(\beta r_x(u,c)-d_x(u,c)),
 \quad u=z-z_*(c).
 \tag{15.2}
\]
The uncut remainders \(r_x,d_x\) here are evaluated on
\(\mathcal D\), where their original Taylor definitions
are valid. Formula (15.2) follows from (14.3).
The linear Taylor terms sum to zero by stationarity.

Interpolate only the potential, keeping the domain fixed:
\[
 U_s=E_\beta+s\Delta_\beta,\qquad
 Z_s^{\mathcal D}=\int_{\mathcal D}e^{-U_s(z,c)}\,dz,
 \qquad \Phi_s=-\log Z_s^{\mathcal D},\quad 0\le s\le1.
 \tag{15.3}
\]
Every coarse derivative in this section is at fixed \(z\),
not at fixed \(u\). In particular it differentiates
\(z_*(c)\) inside the cutoff and remainders.

\begin{proposition}[Uniform interpolating box and common measured minimum]
\label{prop:f12v11-interpolation}
For sufficiently large \(\beta\), uniformly in \(s,c\)
and volume, \(U_s\) has a fixed interaction range,
\[
 D_z^2U_s\ge m\beta I,\qquad
 \sup_i\sum_{j\ne i}|(U_s)_{ij}|\le K\beta.
 \tag{15.4}
\]
Its unique minimum is the same point \(b_\beta(c)\)
for every \(s\), with
\[
 \|b_\beta-z_*\|_\infty\le C/\beta,\qquad
 \operatorname{dist}(b_\beta,\partial\mathcal D)
               \ \text{in each coordinate}\ \ge r/4 .
 \tag{15.5}
\]
\end{proposition}
\begin{proof}
The actual potential is uniformly strongly convex on
the chosen box by f11 V66 and the subsequent fixed-box
estimates. The extended potential has the same properties
globally by V10. A convex combination preserves their
common lower bound; their finite ranges and off-diagonal
row bounds also persist.

To locate the minimum, first work in the extension's
central neighbourhood. Its critical-point equation is
\[
 u=-H_c^{-1}\left(
       \nabla_u\sum_xr_x(u,c)
                         -\beta^{-1}\nabla_u\sum_xd_x(u,c)\right).
 \]
The inverse has a uniform row bound. On a max-norm
ball of radius \(C_0/\beta\), the right side is
\(O(\beta^{-1})\), with derivative
\(O(C_0/\beta+\beta^{-1})\), uniformly in volume.
Choose \(C_0\) large enough and then \(\beta\) large
enough to obtain a contraction of that ball into itself.
The resulting critical point lies where every \(\chi_x=1\).
Both potentials, and hence every \(U_s\), coincide on
an open neighbourhood of it. It is therefore critical
for every \(s\). Strong convexity makes it the unique
minimum; its distance from the box faces follows from
\(\|z_*\|_\infty\le r/2\) and the \(O(\beta^{-1})\) bound.
\end{proof}

\subsection{Rare local defects, not a rare global union}

Write \(\mu_s^{\mathcal D}\) for the probability law
with density \(e^{-U_s}/Z_s^{\mathcal D}\), and set
\[
 B_x(c)=\{z:\max_{i\in S_x}|z_i-z_{*,i}(c)|
                                                  \ge\varepsilon\}.
 \tag{15.6}
\]

\begin{proposition}[Uniform tails for each defect support]
\label{prop:f12v11-tail}
There are constants \(C,a>0\), independent of \(s,c,x\)
and volume, such that
\[
 \mathbb P_{\mu_s^{\mathcal D}}(B_x(c))\le Ce^{-a\beta}.
 \tag{15.7}
\]
\end{proposition}
\begin{proof}
The local-moment proof of f11 V89 applies to (15.4)
with the common interior minimum (15.5). To specify
the uniformity, choose the same exponentially decreasing
coordinate weights \(w_i=e^{-\alpha d(i,i_0)}\) in that
proof. The conjugation error of the Hessian is bounded
by \(K\beta(e^{\alpha R/2}-1)\); choose \(\alpha\) so
it is at most \(m\beta/2\). Weighted integration by
parts has nonnegative face flux since \(b_\beta\)
is inside the fixed box. It yields
\[
 \mathbb E_s|z_{i_0}-b_{\beta,i_0}|^2
       \le \frac{2}{m\beta}\sum_iw_i\le C/\beta.
 \]
Thus \(|\mathbb E_s z_i-z_{*,i}|\le C\beta^{-1/2}\).

The boxed Poincare inequality also holds for every
linear tilt \(U_s-tz_i\): its Hessian remains (15.4),
and convex box boundary conditions are unchanged.
The second derivative of the corresponding log
moment-generating function is therefore at most
\((m\beta)^{-1}\). Chernoff's bound gives
\[
 \mathbb P_s(|z_i-\mathbb E_sz_i|\ge t)
                                  \le2e^{-m\beta t^2/2}.
 \]
For large \(\beta\), the mean displacement is at most
\(\varepsilon/2\). Use \(t=\varepsilon/2\) and sum
only over the bounded number of indices in \(S_x\).
For example one may use a fixed positive \(a\) no
larger than \(m\varepsilon^2/8\), changing \(C\).
There is no sum over all volume coordinates in (15.7).
\end{proof}

Let \(w_a(x)=e^{-\nu d(x,a)}\) and
\(w_{ab}(x)=e^{-\nu\tau(x,a,b)}\), reducing \(\nu\)
when required by the centred chain rule.

\begin{proposition}[Coarse derivatives of the local defects remain rare]
\label{prop:f12v11-rare-jets}
For some \(a_1>0\), and with \(j=0,1\) denoting zero
or one fibre gradient, the following \(L^4(\mu_s^{\mathcal D})\)
bounds hold:
\[
 \begin{split}
 \|\nabla_z^j\Delta_x\|_4&\le C\beta e^{-a_1\beta},\\
 \|\nabla_z^j\partial_a\Delta_x\|_4
                    &\le C\beta e^{-a_1\beta}w_a(x),\\
 \|\nabla_z^j\partial_a\partial_b\Delta_x\|_4
                    &\le C\beta e^{-a_1\beta}w_{ab}(x).
 \end{split}
 \tag{15.8}
\]
All these functions have fibre support in the fixed
neighbourhood \(S_x\).
\end{proposition}
\begin{proof}
On the bounded original box, each unweighted coefficient
and its indicated derivatives is at most \(C\beta\);
the fixed cutoff radius may enter \(C\).
At fixed \(z\), the chain rule differentiates \(u=z-z_*(c)\).
The response estimates (14.1)--(14.2) give exactly the
weights in (15.8), including the second-response terms
and products of first responses. These do not enlarge
the support in the integration variables \(z\).

Every displayed derivative vanishes where
\(\max_{i\in S_x}|u_i|<\varepsilon\): on that open
set the two local germs coincide identically in \(z,c\).
Smoothness of the cutoff handles the threshold as well.
Each derivative is therefore bounded by its stated
supremum times \({\bf1}_{B_x(c)}\).
Taking \(L^4\) norms and applying (15.7) proves (15.8),
with \(a_1\) decreased if necessary.
This reasoning keeps the derivatives of the moving
centre; it does not freeze the rare event.
\end{proof}

\subsection{Connected estimates with a rare insertion}

The boxed covariance bound (90.4) of f11 V86 applies
uniformly to every \(U_s\):
\[
 |\operatorname{Cov}_s(F,G)|
       \le C\beta^{-1}e^{-\alpha d(A,B)}
                              \|\nabla F\|_2\|\nabla G\|_2
 \tag{15.9}
\]
for gradient supports \(A,B\).
At fixed \(z\), the coarse scores of \(U_s\) have local
decompositions
\[
 U_{s,a}=\sum_x X_{a,x},\qquad
 U_{s,ab}=\sum_x X_{ab,x},
 \tag{15.10}
\]
with bounded fibre supports and supremum bounds on both
values and first fibre gradients of respectively
\(C\beta w_a(x)\) and \(C\beta w_{ab}(x)\).
To see the locality of the extended part, decompose its
centred Taylor polynomial term by term as in (14.3);
the local linear terms may be retained in the
decomposition since their sum is zero.
The chain rule and (14.1) then give these bounds.
The actual part has the stronger finite parameter range.

Apply the proof of V10's three-point estimate with
two unweighted scores \(X,Y\) of value and gradient
\(L^4\) norms \(O(\beta)\), and one defect \(Z=\Delta_z\)
of value and gradient norms \(O(\beta e^{-a_1\beta})\).
The product-gradient estimate and (15.9) give
\[
 |\operatorname{Cum}_s(X,Y,\Delta_z)|
       \le C\beta^2e^{-a_1\beta}
               e^{-\alpha'\operatorname{diam}\{x,y,z\}} .
 \tag{15.11}
\]
Indeed each of the three isolation bounds has the
same prefactor; choosing the isolated endpoint of a
diameter gives the exponential shown. This is the
same connected-correlation argument as V10 with
a rare insertion, not a new assumption of independence.

\begin{theorem}[Same-box comparison through two coarse derivatives]
\label{thm:f12v11-comparison}
Let
\[
 {\cal E}_\beta(c)=
 -\log\int_{\mathcal D}e^{-A_\beta(z,c)}\,dz
 +\log\int_{\mathcal D}e^{-E_\beta(z,c)}\,dz.
 \tag{15.12}
\]
For some \(C,\kappa,\gamma>0\), uniformly in volume,
small \(c\), and sufficiently large \(\beta\),
\[
 \frac{|{\cal E}_\beta|}{N}\le Ce^{-\kappa\beta},
 \qquad |\partial_a{\cal E}_\beta|\le Ce^{-\kappa\beta},
 \tag{15.13}
\]
\[
 |\partial_a\partial_b{\cal E}_\beta|
       \le Ce^{-\kappa\beta}e^{-\gamma d(a,b)},\qquad
 \|D_c^2{\cal E}_\beta\|_{\rm row}
       +\|D_c^2{\cal E}_\beta\|_{\rm op}\le Ce^{-\kappa\beta}.
 \tag{15.14}
\]
\end{theorem}
\begin{proof}
The domain in (15.3) is fixed. Differentiate the finite
integrals and apply the interpolation identities
(14.13)--(14.14), with \(R_\beta\) replaced by
\(\Delta_\beta\) and \(V_s\) by \(U_s\).
In particular
\[
 \begin{split}
 \partial_s\partial_a\partial_b\Phi_s={}&
 \mathbb E_s\Delta_{ab}
     -\operatorname{Cov}_s(U_{s,ab},\Delta_\beta)\\
 &-\operatorname{Cov}_s(\Delta_a,U_{s,b})
     -\operatorname{Cov}_s(U_{s,a},\Delta_b)\\
 &+\operatorname{Cum}_s(U_{s,a},U_{s,b},\Delta_\beta).
 \end{split}
 \tag{15.15}
\]
No boundary derivative occurs in this formula, since
\(\mathcal D\) is independent of \(c\) and \(s\).
The covariance theorem still retains its box boundary
conditions; no boundary flux has been set to zero.

Expand all defects and scores into local terms.
The expectation terms are bounded by (15.8).
Every covariance with one defect is at most
\(C\beta e^{-a_1\beta}\) times its parameter weights
and the spatial covariance exponential.
The last term is estimated by (15.11).
Exactly the summable convolutions used in V10 yield
\[
 |\partial_s\partial_a\partial_b\Phi_s|
       \le C(\beta+\beta^2)e^{-a_1\beta}
                                      e^{-\gamma d(a,b)}.
 \]
The first-derivative identity gives
\(C\beta e^{-a_1\beta}\), and
\(\partial_s\Phi_s=\sum_x\mathbb E_s\Delta_x\)
gives \(CN\beta e^{-a_1\beta}\).
Integrate \(s\) from zero to one.
Absorb the fixed powers of \(\beta\) by choosing
any sufficiently small \(\kappa<a_1\).
Finally sum the spatial exponential over \(b\) to
obtain the row bound; the symmetric Hessian's operator
norm is bounded by its absolute row norm.
All sums involving a fixed coarse derivative are
localized before summation, so no volume factor is
introduced into either derivative bound.
\end{proof}

\subsection{The exact remaining wall ratio}

Let \(\widehat W_{\beta,\mathcal D}^{\rm ext}\) be the
extension integrated over \(\mathcal D\), with the
same coarse-Haar normalization as the actual local
\(\widehat W_\beta^{\rm loc}\).
Let \(\widehat W_\beta^{\rm ext}\) be V10's
whole-space extension, and let
\(\mu_{\beta,c}^{\rm ext}\) denote its normalized
whole-space law in the original \(z\) coordinates.
Define
\[
 B_\beta(c)=-\log\mu_{\beta,c}^{\rm ext}(\mathcal D).
 \tag{15.16}
\]
This is a probability of the original fixed box,
not a box implicitly recentered with \(c\).
The exact identities are
\[
 \widehat W_\beta^{\rm loc}
       =\widehat W_{\beta,\mathcal D}^{\rm ext}
                                      +{\cal E}_\beta,
 \qquad
 \widehat W_{\beta,\mathcal D}^{\rm ext}
       =\widehat W_\beta^{\rm ext}+B_\beta .
 \tag{15.17}
\]
The coarse-Haar terms cancel in each difference.

\begin{corollary}[Reduction of the original order-one remainder to wall derivatives]
\label{cor:f12v11-wall}
For the original fixed-box action,
\[
 \left\|D_c^2\widehat W_\beta^{\rm loc}
          -\beta D_c^2G-D_c^2\widehat Q-D_c^2B_\beta
 \right\|_{\rm row}
       +\left\|D_c^2\widehat W_\beta^{\rm loc}
          -\beta D_c^2G-D_c^2\widehat Q-D_c^2B_\beta
 \right\|_{\rm op}
                         \le C\beta^{-1/2}.
 \tag{15.18}
\]
\end{corollary}
\begin{proof}
Differentiate the exact identities (15.17).
Apply the row and operator bounds of V10 and (15.14).
The field-independent Gaussian normalization has no
coarse derivatives. Exponentially small terms are
absorbed into the displayed bound.
\end{proof}

\begin{assessment}[One transfer barrier removed, one retained]
The actual potential can now be replaced by the extension
\emph{on the original fixed box}, through two coarse
derivatives, at an exponentially small uniform cost.
This is stronger than agreement of local Taylor jets.
The proof used rare local defects and summable
connected correlations, not an assumption that every
coordinate remains simultaneously in the central region.

The wall term \(B_\beta\) has not been estimated through
two derivatives here. Its definition as a logarithm of
a probability does not make those derivatives small.
In particular (15.18) is not yet the original fixed-box
version of (14.17), and says nothing about the compact
fibre complement. No new numerical value or sign for
the infrared coefficient is inferred.
\end{assessment}

\begin{decision}[Next concrete boundary problem]
Control the local coarse derivatives of the wall ratio
(15.16), retaining the fixed original box.
A possible approach is a coordinatewise convex wall
interpolation, but it must have estimates integrable
up to the hard-wall limit and uniform in volume.
A bound uniform in wall strength alone is insufficient
if it is then integrated over an infinite strength interval.
The next companion audit remains V15 and the next
direction review V20.
\end{decision}

\begin{record}[V11 append and verification record]
V10 had SHA--256
\texttt{6D6CAA6F17C28AD0BB40D0CE7402B83A6BEBB37DA2DE4A8D1B84E058BAF4ED40}.
Its full body is retained before this appendix, with
only the title version changed.
The new arguments use the archived boxed moment and
covariance estimates under explicitly checked
interpolation hypotheses; those estimates are not
relabelled as new theorems.
Only the immediate active predecessor is replaced
by V11. All family archives remain.
\end{record}

\begin{firewall}[Position after V11]
Proved: an exponentially accurate same-fixed-box
comparison between the original potential and the
convex extension, including a volume-uniform Hessian
row bound. Remaining: the wall-ratio derivative
estimate, compact-fibre transfer, certified infrared
integration, continuum construction and the full
Yang--Mills mass gap.
\end{firewall}
\section{V12: integrable wall estimates and the original fixed-box one-loop limit}
\label{sec:f12v12}

V11 left the wall ratio \(B_\beta=-\log
\mu_{\beta,c}^{\rm ext}(\mathcal D)\) as the remaining
obstruction to the original fixed-box comparison.
We now estimate that ratio through two coarse derivatives.
The crucial bound is integrable as the wall strength tends
to infinity; a strength-independent bound would not suffice.
The result concerns the original local box, not the full
compact fibre or a continuum quantum theory.

\subsection{A fixed-coordinate convex wall}

Use the original coordinates \(z\) and fixed box
\(\mathcal D=[-r,r]^N\). Keep \(E_\beta(z,c)\) exactly
as in (15.1), including its field-independent-in-\(z\)
terms. Define
\[
 \psi(t)=\tfrac12(|t|-r)_+^2,\qquad
 \psi_i(z)=\psi(z_i),\qquad
 \Psi(z)=\sum_i\psi_i(z),
 \]
\[
 U_\lambda(z,c)=E_\beta(z,c)+\lambda\Psi(z),\qquad
 Z_\lambda(c)=\int_{\mathbb R^N}e^{-U_\lambda(z,c)}\,dz,
 \qquad F_\lambda(c)=-\log Z_\lambda(c),\quad\lambda\ge0.
 \tag{16.1}
\]
The wall is independent of \(c\). It is \(C^{1,1}\), with
\(\psi'(t)=\operatorname{sign}(t)(|t|-r)_+\) and
\(\psi''(t)={\bf1}_{|t|>r}\) almost everywhere.
Let \(\mu_\lambda\) be the normalized law.

The minimum \(b_\beta(c)\) of the extension lies a fixed
coordinate distance inside \(\mathcal D\), by V11.
It remains the minimum for every \(\lambda\), since the
wall and its gradient vanish there. Almost everywhere,
\[
 D_z^2U_\lambda\ge m\beta I+\lambda J(z),\qquad
 J_{ii}(z)={\bf1}_{|z_i|>r},
 \tag{16.2}
\]
and the off-diagonal Hessian has the same finite range
and \(K\beta\) row bound as \(E_\beta\).
There is no bound independent of \(\lambda\) on the
full diagonal Hessian, and none will be used.

\subsection{The marginal input and the wall-layer moments}

\begin{record}[External input, checked in the primary source]
We use preservation of log-concavity under integration
over other variables: A.~Pr\'ekopa,
\emph{On logarithmic concave measures and functions},
Acta Sci. Math. (Szeged) \textbf{34} (1973), 335--343,
\href{https://acta.bibl.u-szeged.hu/14411/1/math_034_335-343.pdf}
{Theorem 6, p.~342}.
This theorem applies to nonnegative log-concave functions,
not only to normalized densities.
The quantitative wall and derivative estimates below
are proved here; they are not assertions taken from that paper.
\end{record}

\begin{proposition}[Concentration uniform in wall strength]
\label{prop:f12v12-concentration}
For each fixed \(p\ge1\),
\[
 \sup_{\lambda,c,i}\mathbb E_\lambda
           |z_i-z_{*,i}(c)|^p\le C_p\beta^{-p/2}.
 \tag{16.3}
\]
The one-coordinate law is sub-Gaussian about its mean
with variance parameter \((m\beta)^{-1}\).
\end{proposition}
\begin{proof}
Repeat the weighted local second-moment argument of V11,
now on the whole space. The additional diagonal
\(\lambda J\) commutes with the coordinate weights,
so it makes no contribution to their conjugation error.
The lower bound \(m\beta\) and the off-diagonal bound
\(K\beta\) give the same constants as before.
The common minimum is \(b_\beta=z_*+O(\beta^{-1})\).
Whole-space integration by parts, or exhaustion by
boxes centred about that minimum, therefore gives
\(\mathbb E|z_i-b_{\beta,i}|^2\le C/\beta\).
All linear tilts retain the Hessian lower bound.
Their Poincare inequalities give the sub-Gaussian
moment-generating-function bound, as in V11.
The mean is within \(C/\sqrt\beta\) of \(z_{*,i}\);
combine this with the sub-Gaussian bound to obtain (16.3).
Only off-diagonal bounds enter the spatial weighting,
so no constant grows with \(\lambda\).
\end{proof}

\begin{proposition}[Uniform marginal boundary density and wall moments]
\label{prop:f12v12-layer}
If \(p_{\lambda,i}\) is the density of \(z_i\), there are
\(C,a>0\) such that, for \(s\ge0\),
\[
 p_{\lambda,i}(r+s)+p_{\lambda,i}(-r-s)
       \le Ce^{-a\beta}\exp[-(m\beta+\lambda)s^2/2].
 \tag{16.4}
\]
Consequently, for each fixed \(q\ge1\), after changing
the positive exponential constant,
\[
 \|\psi_i\|_q\le C_q e^{-a_q\beta}
                  (\beta+\lambda)^{-1-1/(2q)},\qquad
 \|\psi_i'\|_q\le C_q e^{-a_q\beta}
                  (\beta+\lambda)^{-1/2-1/(2q)} .
 \tag{16.5}
\]
\end{proposition}
\begin{proof}
Pr\'ekopa's theorem makes each marginal log-concave.
In fact it is strongly log-concave with curvature
\(m\beta\): subtracting \(m\beta z_i^2/2\) from
\(U_\lambda\) preserves convexity, then apply the
same marginal theorem and restore that quadratic.

Its mode lies inside \([-3r/4,3r/4]\) for large \(\beta\),
uniformly in all parameters. Here are details that avoid
assuming mode and mean coincide. By the local second
moment, an interval \([b_{\beta,i}-\delta,b_{\beta,i}+\delta]\),
with \(\delta=r/32\), has probability greater than \(1/2\).
If the mode were at least \(b_{\beta,i}+3\delta\),
the next interval of the same length would have at least
as much probability, by unimodality, a contradiction.
The other side is analogous. Since
\(|b_{\beta,i}|\le r/2+O(\beta^{-1})\), the stated
mode bound follows. The density is thus decreasing on
\([3r/4,r]\) and increasing on \([-r,-3r/4]\).
The sub-Gaussian tail and the mean bound imply
\[
 p_{\lambda,i}(r)\le\frac4r
                  \mathbb P_\lambda(z_i\ge3r/4)\le Ce^{-a\beta},
 \]
and the analogous estimate at \(-r\).

Remove only the \(i\)-th wall factor from the marginal:
\[
 p_{\lambda,i}(t)=e^{-\lambda\psi(t)}h_{\lambda,i}(t).
 \]
The function \(h_{\lambda,i}\) includes normalization
and integration over the other coordinates with all
their walls retained. Applying the marginal theorem
after subtracting \(m\beta t^2/2\) shows that
\(\log h_{\lambda,i}\) has curvature at most \(-m\beta\).
Its mode is the same interior mode as that of
\(p_{\lambda,i}\), because the removed factor is one
on the entire interval \([-r,r]\).
Hence the outward slopes of \(\log h_{\lambda,i}\)
at the two boundary points are nonpositive.
The supporting-line bound for the strongly concave
logarithm gives
\[
 h_{\lambda,i}(r+s)\le h_{\lambda,i}(r)e^{-m\beta s^2/2},
 \]
and similarly on the negative side. Restore
\(e^{-\lambda s^2/2}\) and use the boundary density
bounds to obtain (16.4).

For (16.5), integrate \(s^{2q}\) or \(s^q\) against
the right side of (16.4). The elementary substitution
\(v=\sqrt{m\beta+\lambda}\,s\) gives respectively
powers \((m\beta+\lambda)^{-(2q+1)/2}\) and
\((m\beta+\lambda)^{-(q+1)/2}\) before taking the
\(q\)-th root. Absorb \(m\) into the constants.
\end{proof}

\subsection{A wall-sensitive covariance bound}

For each finite \(\lambda\), use the whole-space
covariance representation
\[
 \operatorname{Cov}_\lambda(F,G)
       =\langle\nabla F,\mathcal A_\lambda^{-1}\nabla G\rangle,
 \qquad
 \mathcal A_\lambda=\mathcal L_\lambda\otimes I+D^2U_\lambda ,
 \tag{16.6}
\]
in \(L^2(\mu_\lambda)^N\).
In quadratic-form order,
\[
 \mathcal A_\lambda\ge M_\lambda,\qquad
 (M_\lambda)_{ii}=m\beta+\lambda{\bf1}_{|z_i|>r}.
 \tag{16.7}
\]
The diagonal multiplication operator \(M_\lambda\)
depends on the field; it is not replaced by its mean.

\begin{proposition}[Extra inverse factor on a wall-supported source]
\label{prop:f12v12-resolvent}
Suppose \(\nabla F\) has component support \(A\).
If a vector \(g\) has only component \(i\), and \(g_i=0\)
where \(|z_i|\le r\), then
\[
 |\langle\nabla F,\mathcal A_\lambda^{-1}g\rangle|
 \le
 \frac{C e^{-\alpha d(A,i)}}
           {\sqrt{\beta}\sqrt{m\beta+\lambda}}\,
                \|\nabla F\|_2\|g_i\|_2.
 \tag{16.8}
\]
Without wall support, the usual bound is
\(C\beta^{-1}e^{-\alpha d(A,B)}
 \|\nabla F\|_2\|g\|_2\) for component support \(B\).
The coefficients of \(g\) may depend on distant variables;
only its vector-component support is used in these estimates.
\end{proposition}
\begin{proof}
Use weights \(w_j=e^{\alpha d(j,i)}\) and diagonal
component multiplication \(W\).
They commute with the differential part of
\(\mathcal A_\lambda\) and with the entire diagonal wall.
The conjugation error \(D_\alpha\) contains only
off-diagonal entries of \(D^2E_\beta\), so
\(\|D_\alpha\|\le K\beta(e^{\alpha R}-1)\).
Choose \(\alpha\) so this is at most \(m\beta/2\).

Factor the conjugated inverse using
\[
 \mathcal A_\lambda+D_\alpha
 =\mathcal A_\lambda^{1/2}
  (I+\mathcal A_\lambda^{-1/2}D_\alpha
                   \mathcal A_\lambda^{-1/2})
                         \mathcal A_\lambda^{1/2}.
 \]
The middle perturbation has norm at most \(1/2\).
Thus its inverse bilinear form is bounded by twice
the product of the two \(\mathcal A_\lambda^{-1/2}\)
norms. The form inequality (16.7) allows each such
norm to be bounded by the corresponding
\(M_\lambda^{-1/2}\) norm. This inverse-order inequality
also follows directly from the variational formula
for a positive inverse.

On \(A\), the first weighted argument \(W^{-1}\nabla F\)
has a factor at most \(e^{-\alpha d(A,i)}\), and its
\(M_\lambda^{-1/2}\) norm is at most
\((m\beta)^{-1/2}\) times its ordinary norm.
The other argument is \(Wg=g\).
On the support of \(g_i\), its inverse weight is
exactly \((m\beta+\lambda)^{-1/2}\).
This proves (16.8). Dropping the wall improvement
gives the last stated bound.
\end{proof}

The form identities above remain valid for the \(C^{1,1}\)
wall. One may smooth it convexly at each fixed finite
\(\lambda\), prove the identities, and pass to the limit.
The smoothing adds only nonnegative diagonal Hessians;
the off-diagonal conjugation bound is unchanged.
The general \(M_\lambda^{-1/2}\) estimate passes to the
limit by quadratic-form approximation and dominated
convergence for the displayed sources. No derivative
of the discontinuous function \(\psi''\) is used.
This procedure is at fixed \(\lambda\); the resulting
exact-wall estimates have the uniform constants just proved.

\subsection{Integrable first and second coarse derivatives}

The coarse derivatives of \(U_\lambda\) equal those of
\(E_\beta\); the wall is fixed in \(z\).
Use the local decompositions and parameter weights of V11:
\[
 E_{\beta,a}=\sum_xX_{a,x},\qquad
 E_{\beta,ab}=\sum_xX_{ab,x}.
 \]
The \(L^4\) norms of their values and first gradients are
bounded by \(C\beta w_a(x)\) and \(C\beta w_{ab}(x)\).
On the whole space these follow from (16.3):
outside the cutoff support the local extension is at
most quadratic in \(z-z_*\), and all its relevant
coarse derivatives have the same polynomial growth.
The bounds are uniform in \(\lambda\).

\begin{proposition}[Integrable connected insertion of one wall]
\label{prop:f12v12-connected}
For two unweighted local score pieces \(X,Y\) near
\(x,y\), with the bounds just stated,
\[
 |\operatorname{Cum}_\lambda(X,Y,\psi_i)|
 \le C\beta^2e^{-a\beta}(1+\lambda)^{-9/8}
               e^{-\alpha'\operatorname{diam}\{x,y,i\}} .
 \tag{16.9}
\]
Covariances of a score piece with \(\psi_i\) obey the
same upper bound, with the corresponding two-point
distance. The exponent \(9/8>1\) is essential.
\end{proposition}
\begin{proof}
For a single covariance, (16.8) and the \(q=2\)
bound on \(\psi_i'\) give a strength factor
\((m\beta+\lambda)^{-1/2}(\beta+\lambda)^{-3/4}\),
which decays as \((1+\lambda)^{-5/4}\).
The remaining factors are a fixed power of \(\beta\)
times an exponential small in \(\beta\).

For the third cumulant write
\[
 \operatorname{Cum}(X,Y,\psi_i)
   =\operatorname{Cov}(X,(Y-\mathbb EY)
                                   (\psi_i-\mathbb E\psi_i)).
 \]
Split the gradient of the second argument into
\((\psi_i-\mathbb E\psi_i)\nabla Y\), supported in
the components near \(y\), and
\((Y-\mathbb EY)\nabla\psi_i\), supported in component \(i\).
For the first use the ordinary covariance bound,
Holder's inequality, and
\(\|\psi_i-\mathbb E\psi_i\|_4
    \le Ce^{-a\beta}(\beta+\lambda)^{-9/8}\).
For the second use (16.8) and
\(\|\psi_i'\|_4
    \le Ce^{-a\beta}(\beta+\lambda)^{-5/8}\).
Its extra inverse factor adds \(1/2\) to \(5/8\),
again giving \(9/8\).
With the score norms \(O(\beta)\), both estimates
are bounded by the prefactor in (16.9), times
\(e^{-\alpha d(x,y)}+e^{-\alpha d(x,i)}\).
Decreasing \(a\) allows the same exponential in each case.

Interchange \(X,Y\) for the second isolation bound.
For the third write the cumulant as
\(\operatorname{Cov}((X-\mathbb EX)(Y-\mathbb EY),\psi_i)\).
The product gradient has \(L^2\) norm \(O(\beta^2)\).
Using (16.8) and the \(L^2\) wall-gradient estimate
gives the stronger strength decay \(5/4\), hence
also the required \(9/8\) bound, times
\(e^{-\alpha\min(d(i,x),d(i,y))}\).
Choose an isolated endpoint of a diameter as in V10.
This proves (16.9), with support-radius constants absorbed.
All uses of (16.8) on product gradients are legitimate:
that estimate requires component support, not locality
of every coefficient multiplying a component.
\end{proof}

\begin{theorem}[Exponentially small wall-ratio derivatives]
\label{thm:f12v12-wall}
For \(B_\beta(c)\) defined by (15.16), there are uniform
\(C,\kappa,\gamma>0\) such that
\[
 0\le \frac{B_\beta}{N}\le Ce^{-\kappa\beta},
 \qquad |\partial_aB_\beta|\le Ce^{-\kappa\beta},
 \tag{16.10}
\]
\[
 |\partial_a\partial_bB_\beta|
       \le Ce^{-\kappa\beta}e^{-\gamma d(a,b)},\qquad
 \|D_c^2B_\beta\|_{\rm row}
       +\|D_c^2B_\beta\|_{\rm op}\le Ce^{-\kappa\beta}.
 \tag{16.11}
\]
\end{theorem}
\begin{proof}
For finite \(\lambda\), fixed-domain differentiation gives
\[
 \partial_\lambda F_\lambda=\mathbb E_\lambda\Psi,\qquad
 \partial_\lambda\partial_aF_\lambda
                   =-\operatorname{Cov}_\lambda(E_{\beta,a},\Psi),
 \]
\[
 \partial_\lambda\partial_a\partial_bF_\lambda
       =-\operatorname{Cov}_\lambda(E_{\beta,ab},\Psi)
            +\operatorname{Cum}_\lambda
                        (E_{\beta,a},E_{\beta,b},\Psi).
 \tag{16.12}
\]
There are no derivatives of the wall in \(c\).
The \(L^1\) bound in (16.5) makes
\(\mathbb E\Psi\le CN e^{-a\beta}
                         (\beta+\lambda)^{-3/2}\).
For the derivative identities, expand all scores and
\(\Psi\) into local pieces. Apply (16.9), then sum
the exponential parameter weights and spatial decay
as in V10--V11. This gives, for a fixed coarse pair,
\[
 |\partial_\lambda\partial_a\partial_bF_\lambda|
 \le C\beta^2e^{-a\beta}(1+\lambda)^{-9/8}
                                      e^{-\gamma d(a,b)},
 \tag{16.13}
\]
and the analogous one-coordinate bound.
The sums are uniform in volume.

The strength integral is finite:
\(\int_0^\infty(1+\lambda)^{-9/8}\,d\lambda=8\).
Integrate (16.12)--(16.13) and absorb the fixed
powers of \(\beta\) by reducing the positive exponent
\(a\). The zero-order identity is nonnegative.

It remains to identify the limit, rather than assume
a formal hard wall. At each finite volume,
\(e^{-\lambda\Psi}\) decreases to
\({\bf1}_{\mathcal D}\), so dominated convergence gives
\[
 Z_\lambda\longrightarrow
                      \int_{\mathcal D}e^{-E_\beta(z,c)}\,dz.
 \]
On each smaller coarse parameter neighbourhood the
first two derivatives of the integrand are dominated
by an integrable polynomial times a confining Gaussian,
independently of \(\lambda\), since the wall has no
coarse derivative. Thus the corresponding partition
derivatives converge as well. The limiting partition
function is positive. Consequently
\(F_\infty-F_0=B_\beta\), including its first two
derivatives. The uniform bounds obtained by integrating
in strength therefore apply to this actual ratio.
Summing the spatial exponential proves the row bound
and then the symmetric operator bound.
No limit of the volume was taken in this identification.
\end{proof}

\subsection{The original local cutoff now has the uniform one-loop limit}

\begin{theorem}[Volume-uniform one-loop comparison on the original fixed box]
\label{thm:f12v12-local}
Let
\[
 {\cal R}_\beta(c)=
 \widehat W_\beta^{\rm loc}(c)-\beta G(c)
       -\frac N2\log(\beta/2\pi)-\widehat Q(c)
 \]
for the original fixed-box integral, with the actual
minimum, fibre Hessian, density and coarse-Haar normalization.
Uniformly in finite volume and small \(c\),
\[
 \frac{|{\cal R}_\beta|}{N}\le C\beta^{-1/2},\qquad
 |\partial_a{\cal R}_\beta|\le C\beta^{-1/2},
 \tag{16.14}
\]
\[
 |\partial_a\partial_b{\cal R}_\beta|
       \le C\beta^{-1/2}e^{-\gamma d(a,b)},\qquad
 \|D_c^2{\cal R}_\beta\|_{\rm row}
       +\|D_c^2{\cal R}_\beta\|_{\rm op}
                                      \le C\beta^{-1/2}.
 \tag{16.15}
\]
\end{theorem}
\begin{proof}
Use the exact decomposition (15.17).
V10 bounds the whole-space extension remainder by
\(O(\beta^{-1/2})\), including exponentially localized
Hessian entries. V11 bounds the same-box potential
replacement by \(O(e^{-\kappa\beta})\).
Theorem~\ref{thm:f12v12-wall} bounds removal of the
extension's wall by the same exponential scale.
Adding these three estimates proves (16.14)--(16.15).
The auxiliary extension has been eliminated from the
conclusion; it is only an instrument of the proof.
\end{proof}

\begin{assessment}[A genuine local transfer, not a mass-gap conclusion]
The uniform order-one Hessian remainder previously
available only for the extension is now proved for
the original fixed local box. The wall-strength
integration is controlled, and the global rare-event
union has not been discarded.
This strengthens the earlier fixed-volume Laplace
expansion and the comparison only on the leading
\(\beta\) scale. It is not an all-orders expansion.

The box is still a restricted part of the compact
fibre. The theorem supplies no bound for the logarithmic
derivatives of the compact complement, no iteration
with scale-uniform constants in physical units, and
no construction of the continuum Hilbert space.
The nonflat momentum integral also still lacks a
certified numerical enclosure. None of those
requirements is inferred from (16.15).
\end{assessment}

\begin{decision}[Next obstruction after the local remainder]
Return to the compact-fibre comparison, now with
(16.15) as a proved local input.
Identify the exact conditional partition ratio between
the full compact fibre and the controlled local box,
retaining any other low-action minima or sector labels.
Do not transfer the convex proof to that complement
without establishing its required geometry and estimates.
The next companion audit remains V15; the next
ten-version direction review is V20.
\end{decision}

\begin{record}[V12 append and verification record]
V11 had SHA--256
\texttt{F25FB668DA72FBD631FBC906C3E23B6F46B9B924288419B254E537B280715F68}.
Its full body is preserved before this appendix, apart
from the title version.
The new arguments are analytic; no numerical quadrature
or matrix certificate is asserted for them.
The external marginal theorem was checked in its
primary paper and is cited explicitly above.
Only the active predecessor is replaced by V12;
all family archives remain.
\end{record}

\begin{firewall}[Position after V12]
Proved: integrable soft-wall bounds, exponentially
small wall-ratio derivatives, and the volume-uniform
one-loop Hessian comparison for the original fixed-box
local integral.
Not proved: compact-fibre transfer, certified infrared
integration, scale iteration, continuum reconstruction
or the full Yang--Mills mass gap.
\end{firewall}

\section{V13: global compact fibre coercivity and a matched complement bound}
\label{sec:f12v13}

V12 closes the one-loop remainder only for the original
local box. We now return to the actual compact root-reference
block \(K^0\), with its good and fallback branches both retained.
The new result controls the distance of every compact fine
field to the section of its own output, then uses the exact
pivot disintegration to match the Gaussian powers in a full
compact numerator and local denominator.

\begin{record}[Nonduplication and scope]
The following inputs are inherited, not proved anew:
the path filling in f11 V68, the section and its action
in f11 V70, the global section tube in f11 V71, the
root-pivot density in f11 V76, and the root-output distance
bound in f11 V78. The older f6 V31 already used global
coercivity and Gaussian power matching on a different,
pinned reduced fibre at retained identity. Its f6 V38
explicitly warns that global coercivity is not conditional
excess energy. The present extension is to the actual
averaged compact block, arbitrary retained fields, both
branches, and its conditional pivot measure. No pinned
fibre measure is substituted for that measure.
\end{record}

\subsection{A global distance bound with the retained field kept}

Use normalized product Haar throughout. Let \(n\) be the
number of parity cells and \(d=\dim G\), and assume at least
two parity cells in each periodic direction as in f11 V68.
Write
\[
 \delta(U)=\|I-U\|_{\rm HS},\qquad
 s_p(W)=\tfrac12\delta(W_p)^2,\qquad S_0=\sum_p s_p.
 \tag{17.1}
\]
Let \(\mathcal I\) be the \(17n\) non-tree internal links,
and \(\mathcal B_e\) the eight boundary links over coarse
edge \(e\). The section \(J(C)\) puts \(I\) on internal
links and \(C_e\) on \(\mathcal B_e\). Define
\[
 \mathfrak D(W,C)=
 \sum_{l\in\mathcal I}\delta(W_l)^2+
 \sum_e\sum_{l\in\mathcal B_e}
                      \delta(W_lC_e^{-1})^2.
 \tag{17.2}
\]
This is a distance on the fixed tree slice, not a
minimization over gauges. Under a simultaneous
cell-constant gauge transformation of \(W,C\), every
summand is unchanged.

\begin{proposition}[Full-branch nonlinear fibre coercivity]
\label{prop:f12v13-coercivity}
Choose \(L_{\log}\) with
\(\|\log U\|_{\rm HS}\le L_{\log}\delta(U)\) on the
fixed good logarithm neighbourhood, as in f11 V78.
Set
\[
 A_0=9+6L_{\log},\qquad
 A=2(153+64A_0^2).
 \tag{17.3}
\]
For every compact tree-slice field,
\[
             \mathfrak D(W,K^0(W))\le A S_0(W).
 \tag{17.4}
\]
The constants do not depend on the volume or output.
\end{proposition}
\begin{proof}
For a cell \(r\), let \(D_r\) be the largest plaquette
distance within its unit parity cube. For an internal
positive edge \(x\to x+\hat\mu\), compare the increasing
coordinate tree path to \(x\), followed by \(\mu\),
with the tree path to \(x+\hat\mu\). Moving the last
step to its ordered position requires at most three
square moves, all inside that cell. Both tree
holonomies are \(I\). The product-distance inequality
and invariance under conjugation therefore give
\(\delta(W_l)\le3D_r\). Summing over the seventeen
non-tree internal links yields
\[
 \sum_{l\in\mathcal I}\delta(W_l)^2
       \le153\sum_{p\ {\rm internal\ to\ cells}}\delta(W_p)^2.
 \tag{17.5}
\]

For \(e=(r,\mu)\), let \(D_e\) be the largest plaquette
distance in the f11 V68 rectangle \(\mathcal R_e\)
of side lengths \(3,1,1,1\). Its reference product
\(B_e\) is the root boundary link. Each other boundary
link \(B\in\mathcal B_e\) has a two-link path starting
at a source-cell vertex with \(\mu\)-parity zero.
Its product is \(I_xB\), where \(I_x\) is an internal
link. The preceding bound gives \(\delta(I_x)\le3D_e\).
The inherited transported-path argument uses at most
six square moves relative to the straight root path:
\[
  \delta(B_e^{-1}I_xB)\le6D_e,\qquad
                     \|B-B_e\|_{\rm HS}\le9D_e.
 \tag{17.6}
\]
On the good branch,
\[
 \|K^0_e-B_e\|_{\rm HS}
 \le\frac1{16}\sum_x\|\log(B_e^{-1}P_x)\|_{\rm HS}
 \le6L_{\log}D_e.
 \tag{17.7}
\]
Here \(\delta(e^X)\le\|X\|_{\rm HS}\) for
anti-Hermitian \(X\). On the fallback branch
\(K^0_e=B_e\), so the same upper bound holds.
Equations (17.6)--(17.7) imply
\(\|B-K^0_e\|_{\rm HS}\le A_0D_e\).
Each \(\mathcal R_e\) contains sixty plaquettes and
each fine plaquette belongs to at most eight such
rectangles, by the already checked f11 V68 count.
Consequently
\[
 \sum_e\sum_{l\in\mathcal B_e}\delta(W_l(K^0_e)^{-1})^2
              \le64A_0^2\sum_p\delta(W_p)^2.
 \tag{17.8}
\]
Combine (17.5), (17.8), and (17.1).
All path comparisons take place in the indicated
contractible rectangles; no global gauge fixing or
removal of periodic holonomies occurs.
\end{proof}

\begin{corollary}[Exact zero-action fibres, including retained holonomies]
\label{cor:f12v13-flat}
For any compact coarse field \(C\), the fibre
\((K^0)^{-1}(C)\) has a zero-action configuration if
and only if \(S_{\rm coarse}(C)=0\). In that case its
only zero-action configuration on the fixed tree
slice is \(J(C)\).
\end{corollary}
\begin{proof}
If \(S_0(W)=0\) and \(K^0(W)=C\), (17.4) forces every
summand of \(\mathfrak D(W,C)\) to vanish, so \(W=J(C)\).
The inherited identity
\(S_0(J(C))=4S_{\rm coarse}(C)\) proves necessity
and sufficiency. This statement is pointwise; it
does not assign a conditional probability to a
Haar-null coarse value. Nontrivial flat coarse
holonomies remain in \(C\), rather than being
identified with the identity coarse field.
\end{proof}

\subsection{The compact numerator has the correct Gaussian power}

Let \(Y\) consist of all seventeen internal and seven
nonpivot boundary group variables per cell and coarse
edge, respectively: altogether \(45n\) group variables.
The remaining \(4n\) variables are the root pivots.
Write \(\mathfrak D_{\rm free}(Y,C)\) for (17.2) with
the pivot summands omitted. The f11 V76 conditional
pivot theorem gives the following joint domination:
\[
 (Y,K^0)_*(dY\,dB)=q(Y,C)\,dY\,dC,\qquad
                  0\le q(Y,C)\le M_0^{4n}.
 \tag{17.9}
\]
It is valid for arbitrary fixed \(Y\), up to the usual
Haar null sets, without any small internal collar.

\begin{proposition}[Weighted joint domination]
\label{prop:f12v13-weight}
The full unnormalized compact blocked density
\[
 (K^0)_*(e^{-\beta S_0}dW)=Z_{\rm full,\beta}(C)\,dC
\]
obeys, for Haar-almost every \(C\),
\[
 Z_{\rm full,\beta}(C)
 \le M_0^{4n}\,I(\beta/A)^{45n},\qquad
 I(t)=\int_G e^{-t\delta(U)^2}\,dU.
 \tag{17.10}
\]
If a measurable set \(\mathcal E\) has \(S_0\ge a>0\)
on its intersection with a specified measurable
range of coarse values, then its density in that
range satisfies
\[
 Z_{\mathcal E,\beta}(C)
 \le M_0^{4n} e^{-\beta a/2}
                         I(\beta/(2A))^{45n}.
 \tag{17.11}
\]
\end{proposition}
\begin{proof}
Equation (17.4) gives
\(S_0(W)\ge A^{-1}\mathfrak D_{\rm free}(Y,K^0(W))\).
For any nonnegative test function of \((Y,C)\),
bound its weighted pullback pointwise by this
function times
\(\exp[-\beta\mathfrak D_{\rm free}(Y,C)/A]\),
then apply (17.9). This proves domination of joint
measures, not a selection of a single pivot preimage.
Integrating \(Y\) factorizes the upper bound. Each
internal integral is \(I(\beta/A)\); each boundary
integral is the same by right invariance of Haar,
even though \(C_e\) is arbitrary. This proves (17.10).

On \(\mathcal E\),
\[
 S_0\ge\tfrac12a+\tfrac1{2A}\mathfrak D_{\rm free}.
\]
Apply exactly the same argument to its restricted
pushforward to prove (17.11). No conditional
independence under the Wilson measure is used.
\end{proof}

\begin{proposition}[Uniform Gaussian integral bound]
There is a finite \(H\), depending only on the fixed
group and \(A\), such that for \(\beta\ge1\),
\[
 I(\beta/(2A))\le H\beta^{-d/2}.
 \tag{17.12}
\]
Thus, with \(N=45dn\) and \(B=M_0^4H^{45}\),
\[
 Z_{\rm full,\beta}(C)\le B^n\beta^{-N/2},\qquad
 Z_{\mathcal E,\beta}(C)
          \le B^n\beta^{-N/2}e^{-\beta a/2}.
 \tag{17.13}
\]
\end{proposition}
\begin{proof}
On a sufficiently small exponential neighbourhood
of the identity, \(\delta(e^x)\ge c|x|\) and the
Haar coordinate density is bounded above. Its
integral is at most a constant times \(t^{-d/2}\)
for \(t\ge1\), by comparison with the Euclidean
Gaussian integral. The compact complement has
\(\delta(U)\ge c_0>0\), so its integral is at most
\(e^{-tc_0^2}\), also bounded by a constant times
\(t^{-d/2}\). Enlarging the constant covers
\(1\le\beta\le2A\) using \(I\le1\).
Since \(I\) is decreasing, (17.10)--(17.12) give
both claims in (17.13).
\end{proof}

\subsection{Application to the actual local box and its minimum}

For small retained logarithms \(c\), put \(C=\exp c\)
edgewise and let
\[
 \mathcal L_c=\{\exp X(z,c):z\in\mathcal D=[-r,r]^N\}
                 \subset (K^0)^{-1}(C)
 \tag{17.14}
\]
be exactly the original local fibre used in V10--V12.
Use its actual restricted Haar density, with the
same coordinate normalization in every comparison.
Denote its blocked density relative to coarse Haar
by \(Z_{\rm loc,\beta}(C)>0\), and put
\(Z_{\rm out,\beta}=Z_{\rm full,\beta}-Z_{\rm loc,\beta}\).
The full restriction is a measurable set in the
joint small-\(c\) chart. Boundary sets have zero
measure there.

\begin{proposition}[Uniform separation of the local-box complement]
\label{prop:f12v13-separation}
There are \(\eta>0\), \(c_*>0\), independent of
volume, such that if \(\|c\|_\infty<c_*\) and
\(K^0(W)=\exp c\), then
\[
 W\notin\mathcal L_c
       \ \Longrightarrow\ 
       \mathfrak D(W,\exp c)>\eta^2,\qquad
       S_0(W)>a,\qquad a=\eta^2/A.
 \tag{17.15}
\]
\end{proposition}
\begin{proof}
If every internal link is within \(\eta\) of \(I\)
and every boundary link within \(\eta\) of \(C_e\),
then all links are within \(\eta+O(c_*)\) of \(I\).
Choose these constants small enough for the
volume-independent original logarithm chart.
To check that its fibre coordinate belongs to the
specific box, write \(x=\log W=Ez+Rv\). The inverse
of the fixed linear splitting \([E,R]\) has bounded
row sums independent of volume. Indeed it is given
by \(v=Lx\) and
\(z=S_{\rm free}(x-RLx)\), with the inherited local
matrices. Hence \(\|z\|_\infty\le C\|x\|_\infty\).
Decrease \(\eta,c_*\) so this is less than \(r/2\)
and \(v\) remains in the original implicit chart.
The equality \(K^0(W)=\exp c\) is then \(F(x)=c\)
on its good branch. Uniqueness in the uniform
implicit chart gives \(v=v(z,c)\). Thus
\(W\in\mathcal L_c\). The contrapositive shows
some distance exceeds \(\eta\); (17.4) gives the
action bound. This argument uses the fixed tree
coordinates, not a gauge enlargement of the box.
\end{proof}

\begin{proposition}[The local minimum is global in a total-energy window]
\label{prop:f12v13-minimum}
In the small-\(c\) range, let \(G(c)=f(z_*(c),c)\)
be the inherited unique minimum on the local box.
Then
\[
 G(c)\le4S_{\rm coarse}(\exp c).
 \tag{17.16}
\]
If \(4S_{\rm coarse}(\exp c)<a\), this local
minimum is the unique global action minimum on
the entire exact compact fibre.
\end{proposition}
\begin{proof}
For small \(c\), the section is \(X(0,c)=Rc\)
and belongs to the box, since \(F(Rc)=c\).
Minimize against its action to obtain (17.16).
Strict convexity gives the unique minimum inside
the box. Outside it, (17.15) bounds the action
strictly above \(a>G(c)\). This is a direct
pointwise comparison; it requires no assertion
that a fibre of the possibly discontinuous
global block is a compact smooth manifold.
\end{proof}

\begin{theorem}[Matched full-compact versus local-box comparison]
\label{thm:f12v13-ratio}
There are \(D\ge1\) and \(\beta_0<\infty\),
independent of \(n,c\), such that for
\(\beta\ge\beta_0\), \(\|c\|_\infty<c_*\), and
Haar-almost every \(C=\exp c\),
\[
 0\le
 R_{\rm comp,\beta}(c):=
       \frac{Z_{\rm out,\beta}(C)}{Z_{\rm loc,\beta}(C)}
 \le D^n\exp[-\beta(a/2-G(c))].
 \tag{17.17}
\]
In particular, if
\(4S_{\rm coarse}(\exp c)\le a/4\), then
\[
 R_{\rm comp,\beta}(c)\le
              \exp[n\log D-\beta a/4].
 \tag{17.18}
\]
\end{theorem}
\begin{proof}
The minimum \(z_*(c)\) lies in a uniformly smaller
box. Integrate in a cube of coordinate radius
\(t_0/\sqrt\beta\) about it, for fixed small \(t_0\).
Stationarity and the uniform upper Hessian bound give
\[
 f(z,c)\le G(c)+C\|z-z_*(c)\|_2^2
                  \le G(c)+CNt_0^2/\beta.
 \tag{17.19}
\]
The actual local Jacobian, including conversion to
coarse Haar, is bounded below by a fixed constant
to the power \(n\): it is a product of uniformly
bounded positive local factors, as in the inherited
chart. The cube volume and (17.19) therefore give
some \(b>0\), independent of volume,
\[
          Z_{\rm loc,\beta}(C)
                \ge b^n\beta^{-N/2}e^{-\beta G(c)}.
 \tag{17.20}
\]
Apply (17.11)--(17.13) with the complement event
of (17.15). Divide by (17.20) and take
\(D=\max(1,B/b)\). This proves (17.17), with the
powers \(\beta^{-N/2}\) cancelling exactly.
Equation (17.16) implies (17.18).
\end{proof}

\begin{corollary}[What transfers, and the size dependence]
With \(W_{\rm full}=-\log Z_{\rm full,\beta}\) and
\(W_{\rm loc}=-\log Z_{\rm loc,\beta}\), in the
same almost-everywhere sense,
\[
 0\le W_{\rm loc}-W_{\rm full}
       =\log(1+R_{\rm comp,\beta})\le R_{\rm comp,\beta}.
 \tag{17.21}
\]
The conditional full-fibre and local-box probabilities
have total variation distance
\(R_{\rm comp,\beta}/(1+R_{\rm comp,\beta})\),
where total variation is the supremum over events.
Consequently the expectations of any bounded fine
observable differ by at most
\(2\|F\|_\infty R_{\rm comp,\beta}/(1+R_{\rm comp,\beta})\).
In the total-curvature window of (17.18), these
bounds vanish along any sequence with
\(\beta a/4-n\log D\longrightarrow+\infty\);
in particular \(n=o(\beta)\) suffices.
\end{corollary}
\begin{proof}
Split the full restricted pushforward into its
local and complementary positive measures. Their
conditional probabilities form a mixture with
complement weight \(R/(1+R)\) and disjoint supports.
This proves the total variation identity and the
observable bound. Equation (17.21) and the stated
size dependence follow directly from (17.17)--(17.18).
\end{proof}

\subsection{What is still upstream of the mass gap}

\begin{assessment}[No derivative or thermodynamic conclusion]
The estimates above control the full compact density
in value and its conditional bounded observables,
not derivatives in \(c\). The fallback is only
measurable and may be discontinuous at its branch
boundary. A bounded density and an exponentially
small value remainder do not establish \(C^2\)
regularity or a small Hessian remainder. Thus the
V12 one-loop Hessian theorem has not been transferred
to the full compact measure by (17.21).

The factor \(D^n\) also remains, and
\(4S_{\rm coarse}(C)\le a/4\) is a bound on
\emph{total} retained curvature, not on a curvature
density. Sending volume to infinity at fixed
\(\beta\), or assuming this window is typical in
the desired continuum limit, is not justified.
No local conditional cluster estimate is inferred
from the global distance inequality.
\end{assessment}

\begin{decision}[Next compact comparison must be regional and relative]
The all-branch geometry and Gaussian powers are now
available for the actual compact block. The next
step is to localize the numerator and denominator
comparison to a specified defect region while
preserving its exterior and coarse constraints.
Its action exponent must measure excess over the
constrained minimum, not the absolute action in
(17.4). Check boundary plaquettes and the change
of the minimum explicitly; neither deleting the
factor \(D^n\) nor summing global bounds supplies
that comparison. Smoothness of the compact
pushforward remains a separate requirement.
The next companion audit is V15 and direction
review V20.
\end{decision}

\begin{record}[Version and verification record]
V12 had SHA--256
\texttt{6382046C247F84A64D5BDC403BE8D018A9A84E7424613FC79724D78F4B7B55DB}.
Its complete preceding body is retained apart from
the title version. The present proofs are analytic;
no numerical mass estimate or interval certificate
is asserted. Only the active V12 file is replaced
by V13; all family archives are retained.
\end{record}

\begin{firewall}[Position after V13]
Proved for the actual compact root-reference block:
global fibre-distance coercivity, exact identification
of zero-action fibres with retained holonomies,
Gaussian-power-matched compact numerator bounds,
and full-versus-local comparison in a declared
total-curvature and volume--coupling regime.
Not proved: a volume-uniform regional excess-action
comparison, compact coarse derivative estimates,
scale iteration, continuum reconstruction, or
the full Yang--Mills mass gap.
\end{firewall}

\section{V14: regional power matching with the exterior debit retained}
\label{sec:f12v14}

V13 used all fibre variables in its Gaussian numerator.
The current task is to perform that comparison after
the exterior is fixed. The geometric estimate survives,
but with a specified exterior-only plaquette sum.
This section keeps that sum, constructs a regional
moving minimum, and obtains an improved conditional
bound on defect-dense regions. It does not assert
that every relevant region satisfies the required
small-data condition.

\begin{record}[Existing results and the new content]
The integrated region, its dimension count, its
root-pivot density bound, and the interior-defect
action threshold are inherited from f11 V76--V77.
In particular f11 V77 already proves exponential
conditional suppression under flat data, with a
power of \(\beta\) from an unmatched Haar numerator.
The new estimates below match that power on the
actual compact conditional fibre and replace the
section-based denominator by a Gaussian tube at
the regional stationary minimum. The exterior
debit is derived, not discarded.
\end{record}

\subsection{Geometry after the exterior has been frozen}

Let \(R\) be a nonempty connected set of \(s\) coarse
cells and \(E(R)\) its \(m\) incident coarse edges.
Let \(\Lambda_R\) consist of the seventeen non-tree
internal links in each cell of \(R\) and all eight
boundary links on each edge of \(E(R)\). Freeze all
other fine links at \(\xi\), and write
\[
 H_\xi(V)=\sum_{p\in\mathcal P_R}s_p(V,\xi),\qquad
 \mathcal P_R=\{p:p\hbox{ touches }\Lambda_R\}.
 \tag{18.1}
\]
The conditional coarse field is \(C=K^0_{E(R)}(V,\xi)\).
All other coarse links are functions of \(\xi\) alone.
There are
\[
 m=8s-i(R)\le8s,\qquad D_R=17s+7m\le73s,\qquad
 N_R=dD_R
 \tag{18.2}
\]
free group and scalar coordinates, respectively,
where \(i(R)\) counts internal coarse edges.

Let \(\mathcal T_R\) be the union of all plaquettes
internal to cells of \(R\) and all rectangles
\(\mathcal R_e\) of V13 for \(e\in E(R)\). Define
\[
 \mathcal Q_R=\mathcal T_R\setminus\mathcal P_R,\qquad
 E_{\rm ext}(\xi)=\sum_{p\in\mathcal Q_R}s_p(\xi).
 \tag{18.3}
\]
Every plaquette in \(\mathcal Q_R\) has only frozen
links. The elementary counts give
\[
 |\mathcal T_R|\le24s+60m\le504s,\qquad
 |\mathcal P_R|\le6(17s+8m)\le486s.
 \tag{18.4}
\]
The first count uses a union, not a sum with
uncontrolled multiplicities.

Let \(Y\) contain the \(17s\) internal and \(7m\)
nonpivot boundary group variables that are integrated.
Define
\[
 Q_R(Y,C)=
 \sum_{l\in\mathcal I_R}\delta(Y_l)^2+
 \sum_{e\in E(R)}
       \sum_{\substack{l\in\mathcal B_e\\l\ne{\rm pivot}}}
                         \delta(Y_lC_e^{-1})^2.
 \tag{18.5}
\]

\begin{proposition}[Regional coercivity with a named boundary debit]
\label{prop:f12v14-geometry}
With exactly the constant \(A\) of (17.3), for
arbitrary compact \(V,\xi\) on the tree slice,
\[
 Q_R(Y,K^0_{E(R)}(V,\xi))
                 \le A\bigl(H_\xi(V)+E_{\rm ext}(\xi)\bigr).
 \tag{18.6}
\]
\end{proposition}
\begin{proof}
Apply the internal-link path argument of V13 only
to the cells of \(R\), and its boundary-link path
argument only to \(E(R)\). The internal coefficient
is still at most \(153\), and each fine plaquette
is still in at most eight of the selected
rectangles. Including all eight boundary links,
and then discarding the nonnegative pivot terms,
therefore gives
\[
 Q_R\le A\sum_{p\in\mathcal T_R}s_p.
\]
Split this sum into \(\mathcal T_R\cap\mathcal P_R\)
and \(\mathcal Q_R\). The first part is at most
\(H_\xi\), and the second is exactly (18.3).
No restriction on the exterior is needed for
this argument, and both branches of \(K^0\) are
included as in V13.
\end{proof}

\begin{assessment}[The debit is not an added conditional action]
The exterior-only Wilson action cancels from the
actual conditional probability, whose potential
is precisely \(H_\xi\). The quantity \(E_{\rm ext}\)
enters only because the path proof of coercivity
uses some exterior plaquettes. Dropping it from
(18.6) is not justified by the cancellation in
the conditional measure. This distinction is the
reason to retain (18.3) explicitly.
\end{assessment}

\subsection{Full compact regional numerator}

For fixed \(\xi\), denote the unnormalized regional
blocked density relative to \(dC\) by
\[
 Z_{R,\beta}(C,\xi)\,dC
  =(K^0_{E(R)})_*(e^{-\beta H_\xi(V)}dV).
 \tag{18.7}
\]
For a measurable event \(\mathcal B\), write
\(Z_{\mathcal B,R,\beta}\) for the restricted density.

\begin{proposition}[Gaussian-power-matched regional upper bound]
\label{prop:f12v14-numerator}
There is \(B\ge1\), independent of \(R,\xi,C,\beta\),
such that for \(\beta\ge1\) and Haar-almost every \(C\),
\[
 Z_{R,\beta}(C,\xi)
       \le B^s\beta^{-N_R/2}e^{\beta E_{\rm ext}(\xi)}.
 \tag{18.8}
\]
If \(H_\xi\ge E\) on \(\mathcal B\) in the
conditional fibre, then
\[
 Z_{\mathcal B,R,\beta}(C,\xi)
 \le B^s\beta^{-N_R/2}
           \exp[-\tfrac{\beta}{2}(E-E_{\rm ext}(\xi))].
 \tag{18.9}
\]
\end{proposition}
\begin{proof}
For arbitrary fixed internal and nonpivot data,
the \(m\) root-pivot pushforward has density at
most \(M_0^m\), by f11 V76. It follows, as a
joint measure statement, that
\[
 (Y,K^0_{E(R)})_*(dY\,dB)
                      \le M_0^m\,dY\,dC.
 \tag{18.10}
\]
Equation (18.6) bounds \(H_\xi\) below by
\(Q_R/A-E_{\rm ext}\). On \(\mathcal B\) it also gives
\[
 H_\xi\ge\tfrac12(E-E_{\rm ext})+\tfrac1{2A}Q_R.
 \tag{18.11}
\]
Use these pointwise bounds before taking (18.10).
Integration of \(Y\) gives the factors
\(I(\beta/A)^{D_R}\) and \(I(\beta/(2A))^{D_R}\),
respectively. Right Haar invariance removes the
arbitrary \(C_e\) in every boundary factor.
The V13 Gaussian bound
\(I(\beta/(2A))\le H\beta^{-d/2}\), together with
\(m\le8s\), \(D_R\le73s\), proves both inequalities;
one may take
\(B=\max(1,M_0)^8\max(1,H)^{73}\).
No factor from the rest of the lattice appears.
\end{proof}

\subsection{A regional stationary denominator}

Let \(\mathcal F_R\) contain every frozen fine link
in \(\mathcal P_R\cup\mathcal T_R\), together with
every frozen internal link used by an affected
block equation. For \(\epsilon>0\), impose
\[
 \delta(\xi_l)<\epsilon\quad(l\in\mathcal F_R),
 \qquad \delta(C_e)<\epsilon\quad(e\in E(R)).
 \tag{18.12}
\]
This enlarged collar includes the links needed
to bound (18.3); it is not silently identified
with the smaller collar in f11 V77.

\begin{proposition}[Uniform regional local minimum]
\label{prop:f12v14-minimum}
There exist \(\epsilon_*>0,r_*>0,\kappa_*>0\)
and finite constants, independent of \(R\), with
the following properties. Under (18.12) with
\(\epsilon<\epsilon_*\), the actual good root-pivot
chart with the exterior fixed has a smooth action
\(f_R(z;C,\xi)\) on a scalar max-norm box of
radius \(r_*\). It has a unique stationary point
in that box, denoted \(z_R^*\), with
\[
 \|z_R^*\|_\infty\le C_*\epsilon,\qquad
 D_z^2 f_R\ge\kappa_* I.
 \tag{18.13}
\]
Write \(G_R(C,\xi)=f_R(z_R^*;C,\xi)\).
For some \(b>0\) and all sufficiently large
\(\beta\), the density from this chart satisfies
\[
 Z_{R,\beta}(C,\xi)\ge
             b^s\beta^{-N_R/2}e^{-\beta G_R(C,\xi)}.
 \tag{18.14}
\]
In addition, for constants \(C_G,C_E\) independent
of \(R\),
\[
       0\le G_R\le C_G\epsilon^2s,\qquad
       E_{\rm ext}\le C_E\epsilon^2s,
       \quad C_E=4032.
 \tag{18.15}
\]
The minimum in this proposition is a local one.
No assertion of global conditional minimality is made.
\end{proposition}
\begin{proof}
Use exponential coordinates for the integrated
internal and free boundary links and solve each
root pivot, as in f11 V71--V77. At flat data
the derivative in a pivot is \(I/8\). Each
equation has bounded size and depends on its
own pivot only, so the same uniform implicit
construction works for any \(R\). Translate
free boundary coordinates to their prescribed
small \(C_e\); at zero data the linear tangent
inclusion is \(E_R\), the columns of the global
inclusion \(E\) corresponding to the selected
internal and nonpivot boundary variables.
Its image changes no exterior link. In particular
an internal variation in \(R\) affects only
pivots on incident edges, which are integrated.

At zero data, the regional fibre Hessian is
\[
 H_R^0=E_R^*d^*dE_R.
 \tag{18.16}
\]
This is a principal submatrix of the inherited
global free-coordinate Hessian. Indeed the
exterior-only plaquettes are constant on these
variations, and all varying plaquettes are in
\(\mathcal P_R\). Thus
\(\kappa I\le H_R^0\le\Lambda I\), with the same
volume-independent constants after increasing
\(\Lambda\) if necessary. Its finite propagation
range is unchanged. A row of
\((I-H_R^0/\Lambda)^j\) has at most
\(C(j+1)^4\) possible entries, even if \(R\) has
holes, since propagation is bounded in the
ambient lattice. The f11 V80 Neumann argument
therefore gives
\[
 \|(H_R^0)^{-1}\|_{\infty\to\infty}
   \le \Lambda^{-1}C^{1/2}
       \sum_{j\ge0}(j+1)^2(1-\kappa/\Lambda)^j.
 \tag{18.17}
\]

All local chart derivatives have uniform bounds.
Locality consequently bounds the row norm of
\(D_z^2f_R-H_R^0\) by \(C(\|z\|_\infty+\epsilon)\),
and \(\|\nabla_z f_R(0;C,\xi)\|_\infty\le C\epsilon\).
Use (18.17) in the contraction
\(z\mapsto z-(H_R^0)^{-1}\nabla_z f_R(z;C,\xi)\).
Choose \(r_*,\epsilon_*\) so its Lipschitz constant
is at most \(1/2\) and its image stays inside the
box. The fixed point has the max-norm bound in
(18.13). Symmetry and the same Hessian row estimate
give \(D_z^2f_R\ge\kappa I/2\), establishing
strict convexity and uniqueness. Smoothness
follows from the finite-dimensional implicit
function theorem; (18.17), not that theorem alone,
provides the uniform radii.

In a cube of radius \(t_0/\sqrt\beta\) about \(z_R^*\),
stationarity and the uniform upper Hessian bound
give \(f_R\le G_R+C N_R/\beta\). The actual local
Haar and pivot Jacobians are bounded below by
a fixed positive constant to the power \(s\),
since there are only \(O(s)\) uniformly bounded
local factors. Integrating this cube proves
(18.14) with the exact power \(N_R/2\).

The central repair of f11 V77 has every changed
link within \(L_0\epsilon\) of \(I\), with
uniform \(L_0\ge1\), and remains in this chart.
There are at most \(486s\) affected plaquettes.
Each has distance at most \(4L_0\epsilon\), hence
action at most \(8L_0^2\epsilon^2\).
Minimization against that repair yields
\(G_R\le3888L_0^2\epsilon^2s\), so this value
may be used for \(C_G\). Finally each plaquette
of \(\mathcal Q_R\) has four frozen links
satisfying (18.12), and hence action at most
\(8\epsilon^2\). Equation (18.4) gives
\(C_E=8\cdot504=4032\).
\end{proof}

\subsection{An explicit positive conditional exponent}

Fix a set \(F\) of \(k\) coarse edges with both
endpoints in \(R\). Let \(\mathcal B_F\) be the
event that every edge in \(F\) fails the actual
relative-logarithm condition. The inherited
f11 V77 action estimate is
\[
 H_\xi\ge\gamma k\quad\hbox{on }\mathcal B_F,\qquad
                    \gamma=\rho^2/2304.
 \tag{18.18}
\]
It uses only plaquettes of \(\mathcal P_R\):
the rectangle of an internal edge lies in its
two integrated endpoint cells and their interface.
The threshold is summed with overlap at most eight.

\begin{theorem}[Regional conditional bound without a power of \(\beta\)]
\label{thm:f12v14-conditional}
Under (18.12) in the preceding chart range,
there is \(C\ge1\), independent of \(R,F\) and
the ambient volume, such that for sufficiently
large \(\beta\) and almost every allowed coarse field,
\[
 \mathbb P_{\beta,\xi}(\mathcal B_F\mid C)
 \le\min\left\{1,\,
 C^s\exp\left[-\beta\left(
       \frac{\gamma k}{2}
       -\frac{E_{\rm ext}(\xi)}2-G_R(C,\xi)
                         \right)\right]\right\}.
 \tag{18.19}
\]
Fix \(0<\theta\le4\). Choose
\(\epsilon_\theta<\epsilon_*\) such that
\[
 \left(C_E/2+C_G\right)\epsilon_\theta^2
                                \le\gamma\theta/4.
 \tag{18.20}
\]
If \(k\ge\theta s\) and (18.12) holds with
\(\epsilon_\theta\), then
\[
 \mathbb P_{\beta,\xi}(\mathcal B_F\mid C)
       \le\min\{1,C^s e^{-\beta\gamma k/4}\}
       \le\min\{1,(C^{1/\theta}e^{-\beta\gamma/4})^k\}.
 \tag{18.21}
\]
\end{theorem}
\begin{proof}
Use (18.9) with \(E=\gamma k\), and divide by
(18.14). The powers \(\beta^{-N_R/2}\) cancel.
The remaining constant is at most \((B/b)^s\).
This proves (18.19) after taking \(C\ge\max(1,B/b)\).
By (18.15)--(18.20),
\[
 E_{\rm ext}/2+G_R
       \le(\gamma\theta/4)s\le\gamma k/4,
\]
which gives (18.21). The denominator is a lower
contribution to the full compact conditional
integral; neither global minimality of \(G_R\)
nor exclusion of other sectors was used.
\end{proof}

\begin{assessment}[Meaning of the conditional exponent]
The improvement over f11 V77 is genuine power
matching and a stationary local denominator.
For arbitrary frozen data, (18.6) alone still
does not give a positive exponent. Even when
the local minimum exists, positivity in (18.19)
requires control of both \(E_{\rm ext}\) and
\(G_R\). Equation (18.20) supplies that control
in one declared regime. It is not a proof of
excess-energy coercivity for all compact fibres.
\end{assessment}

\subsection{A bound on protected connected defect regions}

Call \(\mathcal G_R\) the data event (18.12) at
\(\epsilon_\theta\). It is measurable with
respect to the frozen exterior and the affected
coarse output. One may alternatively use its
cell-constant gauge saturation \(\mathcal G_R^{\rm sat}\).
For brevity let \(\widetilde{\mathcal G}_R\)
denote either choice.

\begin{proposition}[Integration of the conditional estimate]
Under the original finite-volume Wilson law,
for a fixed \(R,F\) with \(k\ge\theta s\),
\[
 \mu_\beta(\mathcal B_F\cap\widetilde{\mathcal G}_R)
 \le C^s e^{-\beta\gamma k/4}
                    \mu_\beta(\widetilde{\mathcal G}_R).
 \tag{18.22}
\]
\end{proposition}
\begin{proof}
Condition first on \(\xi\), then on \(C\), and
use (18.21). The joint absolute continuity
(18.10) and positivity of the local contribution
handle the coarse exceptional sets.

For the gauge-saturated version, use that the
tree slice, the conditional variable set,
Haar measure, Wilson action and failure event
are carried into themselves by cell-constant
gauge transformations. The strict smallness
condition has a witness in a fixed countable
dense family whenever it has any witness.
Partition the saturated data event by the
first witness in that family. For each fixed
witness transform the conditional estimate;
the transformed data satisfy (18.12) with
the same constants. The exceptional sets
form a countable null union. Apply the bound
on the disjoint data pieces and sum them.
This proves (18.22) without summing a bound
over all gauge witnesses.
\end{proof}

\begin{corollary}[Uniform rooted-region entropy bound]
Fix a coarse cell \(r\). Define
\[
 q_\theta(\beta)=1024C\,e^{-\beta\gamma\theta/4}.
 \tag{18.23}
\]
If \(q_\theta(\beta)<1\), the probability that
there exist a connected region \(R\ni r\)
with \(|R|\ge s_0\), a subset \(F\) of its
internal coarse edges with \(|F|\ge\theta|R|\),
and the event
\(\mathcal B_F\cap\widetilde{\mathcal G}_R\), is at most
\[
                \frac{q_\theta(\beta)^{s_0}}
                           {1-q_\theta(\beta)}.
 \tag{18.24}
\]
The upper bound is independent of ambient volume.
\end{corollary}
\begin{proof}
The coarse lattice has degree at most eight.
Every connected \(s\)-cell set containing \(r\)
has a spanning tree whose depth-first traversal
is a walk of length \(2(s-1)\) starting at \(r\).
Choose one such traversal deterministically
for each set. The visited set determines the
region, so the assignment is injective into
at most \(8^{2(s-1)}\le64^s\) walks. There are
at most \(4s\) internal oriented coarse edges,
and hence at most \(2^{4s}=16^s\) choices of \(F\).
For each admissible pair, (18.22) is at most
\(C^s e^{-\beta\gamma\theta s/4}\).
Sum these bounds over pairs and over \(s\ge s_0\).
This is a union bound, with no independence
claim for either defects or protected data.
\end{proof}

\begin{decision}[The unprotected regions are the next obstruction]
Equation (18.24) applies only where the declared
data condition, possibly after one consistent
cell-constant gauge transformation, holds.
It leaves regions with
\(\widetilde{\mathcal G}_R^c\) untreated. In
particular a common gauge on a general large
region cannot be inferred merely from the
fixed non-wrapping collar lemma of f11 V75.
Boundary curvature, coarse curvature, and
wrapping holonomies must all remain explicit.

At V15 perform the required SM/NS and parallel
YM audit. Then examine whether a useful
boundary-conditioned estimate can tolerate a
curved retained background without paying its
whole regional action in (18.19). This requires
a comparison with the actual conditional
background, not repeated insertion of a
flat repair. The smoothness of the compact
blocked density remains a separate open task.
\end{decision}

\begin{record}[Append and verification record]
V13 had SHA--256
\texttt{B545FB0EE21B6FB319CFB4A683FEF519A4CA62513A9C5AFBDF554C2E72A8E7DF}.
Its body is preserved before this appendix,
apart from the title version. The present
results are analytic and retain the explicit
conditional hypotheses. Only active V13 is
replaced by V14; every family archive remains.
\end{record}

\begin{firewall}[Position after V14]
Proved: regional coercivity with an exterior
debit, matched compact Gaussian numerator
and stationary local denominator, and
exponential suppression of protected
defect-dense regions with explicit connected
region entropy. Not proved: protection or
control of all relevant regions, arbitrary
curved-background excess energy, full compact
coarse derivatives, iterative continuum
construction, or the Yang--Mills mass gap.
\end{firewall}

\section{V15: curvature loads instead of the whole background action}
\label{sec:f12v15}

\subsection{Required companion audit}

The current filesystem and the endings of all three
companion files were inspected. They remain:
\[
\begin{array}{l|r|l}
 \text{active companion}&\text{bytes}&\text{SHA--256 prefix}\\ \hline
 \text{SM--Einstein V72}&607372&\texttt{F44F9745D43379E1}\\
 \text{NS Stability V26}&278283&\texttt{82C887F8C2C69E4F}\\
 \text{Parallel YM V5}&54174&\texttt{306F1BB84CFA8A8D}
\end{array}
\]
Their complete hashes agree with the previous V10 audit.
SM supplies a fixed-mode, fixed-time many-copy limit,
not a growing-cutoff or physical finite-species theorem.
NS retains pointwise rank-one Oseen norms and explicitly
leaves global regularity open. Its references to a planned
V27 are not evidence that that version is present.
The parallel YM note rejects a uniform tail for an
extensive global integer charge and retains compact
holonomies and local observables instead. None supplies
the missing regional curved-background estimate.
The next companion audit and direction review are V20.

\begin{record}[The upstream correction to the comparison]
The plaquette classification at the section is already
in f11 V70; its section chart is already in f11 V71.
The new use of that classification is that all plaquettes
in the V13 fibre-coercivity proof are flat at the section,
even for an arbitrary curved coarse field. This permits
an action comparison after subtracting the background
action, rather than paying that whole action in the
defect exponent. The proof below does not replace the
averaged block by a pinned or straight-path block.
\end{record}

\subsection{Only section-flat plaquettes are needed for coercivity}

A fine plaquette in directions \(\mu,\nu\) is called
a corner plaquette if the parity bits of its base
point in both directions are one. Let \(\mathcal C\)
be these plaquettes and \(\mathcal V\) all the others.
Use the fixed parity cells and tree throughout.
For every compact retained field \(C\),
\[
 J(C)_p=I\quad(p\in\mathcal V).
 \tag{19.1}
\]
For \(p\in\mathcal C\), its product \(Q_p(C)=J(C)_p\)
is the corresponding coarse plaquette holonomy.
Neither assertion requires small link amplitudes.
Write
\[
 S_{\rm v}(W)=\sum_{p\in\mathcal V}s_p(W).
 \tag{19.2}
\]

\begin{proposition}[Section-flat fibre coercivity]
\label{prop:f12v15-vertical}
With the unchanged constant \(A\) in (17.3), every
compact tree-slice field satisfies
\[
       \mathfrak D(W,K^0(W))\le A S_{\rm v}(W).
 \tag{19.3}
\]
\end{proposition}
\begin{proof}
In (17.5), all plaquettes lie within single parity
cells, so neither of their directional base bits
is one. In (17.6)--(17.8), each rectangle has length
three in one direction and length one in the
other three, starting at an even root. For a
plaquette involving the long direction, the base
bit in its other direction is zero. For a
plaquette involving two short directions, both
directional base bits are zero. Thus none of the
rectangles contains a corner plaquette. Retain
this information in the sums of the V13 proof
and use exactly its coefficients and overlap
bound. This gives (19.3) on both block branches.
\end{proof}

\subsection{Regional boundary convention and the curvature load}

Retain \(R,\Lambda_R,\mathcal P_R,E(R),s,m,D_R,N_R\)
from V14. For this section, the exterior is the
exact section exterior: fix arbitrary compact coarse
links \(C_{\rm out}\) on edges outside \(E(R)\), set
all frozen internal links to \(I\), and set the eight
frozen boundary links on every such edge to its
\(C_{{\rm out},e}\). Denote this frozen field by
\(\xi_J\).

Condition the affected outputs on \(C_{\rm in}\),
and combine \(C_{\rm in},C_{\rm out}\) into the full
coarse field \(C\). Importantly, \(\xi_J\) is
independent of \(C_{\rm in}\): every boundary link
on an affected edge is integrated. There is thus
no evaluation of an almost-everywhere bound on an
unjustified parameter-dependent diagonal.
All density and conditional-probability comparisons
below are understood for Haar-almost every affected
output at each fixed exterior; the geometric and
action inequalities are pointwise.
For the full periodic region \(R\), no exterior
links remain and this convention imposes no
boundary restriction at all.

For \(W=(V,\xi_J)\) with \(K^0(W)=C\), define
\[
 d_l=\|W_l-J(C)_l\|_{\rm HS}\quad(l\in\Lambda_R),
 \qquad D=\sum_{l\in\Lambda_R}d_l^2,\qquad
 H_0(C)=\sum_{p\in\mathcal P_R}s_p(J(C)).
 \tag{19.4}
\]
The full coarse equality here means equality on
affected edges, with the unaffected edges already
fixed to \(C_{\rm out}\). Put
\[
 H_{\rm v}(W)=\sum_{p\in\mathcal P_R\cap\mathcal V}s_p(W),
 \qquad
 a_l(C)=
   \sum_{\substack{p\in\mathcal P_R\cap\mathcal C\\l\in\partial p}}
                         \delta(Q_p(C)),\qquad
 \mathcal L_R(C)=\sum_{l\in\Lambda_R}a_l(C)^2.
 \tag{19.5}
\]
The load \(\mathcal L_R\) is a local function of
coarse plaquette curvature, not of the distance
of the coarse links from \(I\). It is unchanged
by a simultaneous cell-constant gauge transformation.

\begin{theorem}[Background-subtracted compact action bound]
\label{thm:f12v15-relative}
For every compact \(C\) and every configuration in
this exact conditional fibre,
\[
 D\le A H_{\rm v},\qquad
 H_{\xi_J}(V)-H_0(C)
                  \ge H_{\rm v}(W)-\sum_{l\in\Lambda_R}a_l(C)d_l.
 \tag{19.6}
\]
In particular,
\[
 H_{\xi_J}-H_0
 \ge \frac{D}{2A}-\frac A2\mathcal L_R,\qquad
 H_{\xi_J}-H_0
 \ge \frac12 H_{\rm v}+\frac{D}{4A}-A\mathcal L_R.
 \tag{19.7}
\]
The pointwise infimum \(m_R(C)\) of \(H_{\xi_J}\)
on the entire fibre obeys
\[
       H_0(C)-\frac A4\mathcal L_R(C)
                   \le m_R(C)\le H_0(C).
 \tag{19.8}
\]
\end{theorem}
\begin{proof}
Every frozen link equals its value in \(J(C)\).
Therefore the global distance in (19.3) is
exactly \(D\). Every section-flat plaquette
outside \(\mathcal P_R\) is also unchanged
and has zero action. Thus (19.3) gives the
first inequality in (19.6).

For unitary plaquette products \(P,Q\), the
Hilbert--Schmidt expansion is
\[
 \tfrac12\|P-I\|_{\rm HS}^2-\tfrac12\|Q-I\|_{\rm HS}^2
 =\tfrac12\|P-Q\|_{\rm HS}^2+
       \operatorname{Re}\Tr((Q-I)^*(P-Q))
 \ge-\delta(Q)\|P-Q\|_{\rm HS}.
 \tag{19.9}
\]
Telescoping the four link replacements in a
plaquette, including inverses when required,
gives
\[
 \|W_p-Q_p(C)\|_{\rm HS}
       \le\sum_{l\in\partial p\cap\Lambda_R}d_l.
 \tag{19.10}
\]
Frozen-link replacements contribute zero.
The section-flat plaquettes contribute
exactly their nonnegative action; sum (19.9)
only over the corner plaquettes to obtain
the second inequality in (19.6).

Use \(a_ld_l\le d_l^2/(2A)+Aa_l^2/2\)
and \(H_{\rm v}\ge D/A\) for the first
inequality in (19.7). For the second, retain
half of \(H_{\rm v}\), use the other half
to bound \(D/(2A)\), and apply
\(a_ld_l\le d_l^2/(4A)+Aa_l^2\).
Finally
\[
 D/A-\sum_l a_ld_l
  =\sum_l(d_l/\sqrt A-\sqrt A\,a_l/2)^2
                                 -A\mathcal L_R/4
\]
proves the lower bound in (19.8).
The configuration \(J(C)\) belongs to the fibre
and has action \(H_0\), proving the upper bound.
Existence or uniqueness of a global minimizer
is not needed for this infimum statement.
\end{proof}

\begin{proposition}[A curvature-density bound on the load]
Suppose \(\delta(Q_p(C))\le\epsilon\) for all
\(p\in\mathcal P_R\cap\mathcal C\). Then
\[
                      \mathcal L_R(C)\le576\epsilon^2s.
 \tag{19.11}
\]
\end{proposition}
\begin{proof}
Every edge of a corner plaquette is a boundary
link, so \(a_l=0\) for internal links.
A boundary link in direction \(\mu\) is in
at most three corner plaquettes: for each
\(\nu\ne\mu\), of the two adjacent plaquettes
in the \(\mu,\nu\) plane, exactly one has
odd \(\nu\)-base parity. Hence \(a_l\le3\epsilon\).
There are \(8m\le64s\) integrated boundary
links. Sum their squared loads to get
\(9\cdot64\epsilon^2s\).
\end{proof}

\subsection{A stationary local minimum without small coarse links}

\begin{theorem}[Uniform small-curvature minimum with compact holonomies]
\label{thm:f12v15-minimum}
For the section exterior just specified, there
exist \(r,\epsilon_*,h,K>0\), independent of
the region and the compact coarse link values,
such that if
\[
 \max_{p\in\mathcal P_R\cap\mathcal C}\delta(Q_p(C))
                                          <\epsilon_*,
 \tag{19.12}
\]
the actual root-pivot chart about \(J(C)\) has
a unique stationary minimum in \(\|z\|_\infty<r\).
Writing its action as \(G_R(C)\), its coordinates satisfy
\[
 \|z_R^*(C)\|_\infty\le K\epsilon,\qquad
 D_z^2 f_R(z;C)\ge h I,\qquad
 G_R(C)\le H_0(C),
 \tag{19.13}
\]
where \(\epsilon\) is the maximum in (19.12).
Only local minimality is claimed.
\end{theorem}
\begin{proof}
Use the f11 V71 chart: internal links are
\(\exp z_l\), free boundary links are
\(C_e\exp z_l\), and each pivot is solved
in right coordinates about \(C_e\). With
the frozen internal links equal to \(I\),
the regional restriction of that proof
works for all compact \(C\). In particular
the pivot derivative at the section is
\(I/8\), and all local derivatives have
uniform bounds. This does not require
choosing logarithms of the coarse links.

Restrict (19.6) to this chart. The free
summands of \(D\) equal
\(\sum_l\delta(e^{z_l})^2=|z|^2+o(|z|^2)\).
Since \(H_{\rm v}(0;C)=0\), its gradient
vanishes and
\[
                D_z^2H_{\rm v}(0;C)\ge 2A^{-1}I
 \tag{19.14}
\]
for every compact \(C\), curved or flat.
For each corner plaquette, let \(P_p(z)\)
be its product in the chart. Along a real
direction \(v\),
\[
 \left.\frac{d^2}{dt^2}s_p(P_p(tv))\right|_{t=0}
  =\|DP_p(0)v\|_{\rm HS}^2+
    \operatorname{Re}\Tr((Q_p-I)^*D^2P_p(0)[v,v]).
 \tag{19.15}
\]
The first term is nonnegative. Uniform bounds
on the local second derivatives, followed
by bounded stencil overlap, bound the sum
of the negative parts by \(C\epsilon\|v\|_2^2\).
Thus \(H_C=D_z^2f_R(0;C)\ge(2/A-C\epsilon)I\).
It also has uniformly bounded finite-range
entries and a uniform upper spectral bound.

Choose \(\epsilon_*\) so \(H_C\ge A^{-1}I\).
The Neumann-series and polynomial-ball argument
of (18.17) then gives a uniform max-norm
bound for \(H_C^{-1}\), including on regions
with holes or periodic identifications.
Here this is an application to the new
matrix \(H_C\), whose gap was proved in
(19.14)--(19.15), not assumed from the
identity coarse field.

The gradient at the section has max norm
at most \(C\epsilon\), since the vertical
plaquettes are at their minima and each
corner gradient contains \(Q_p-I\).
Local third-derivative bounds give
\(\|D_z^2f_R(z;C)-H_C\|_{\infty\to\infty}
\le C\|z\|_\infty\). Apply the contraction
\(z\mapsto z-H_C^{-1}\nabla f_R(z;C)\)
on a sufficiently small uniform box.
It has Lipschitz constant at most \(1/2\)
and maps that box strictly into itself
after decreasing \(\epsilon_*\).
This proves the bound on \(z_R^*\).
The same local Hessian estimate gives
strict convexity on the box and uniqueness.
The section is \(z=0\), so minimization
also gives \(G_R\le H_0\).
\end{proof}

\begin{assessment}[No toron or global-flatness substitution]
Condition (19.12) bounds plaquette curvature,
not link amplitudes. It allows arbitrary
compact holonomies and does not assert a
single gauge that makes all coarse links
small. The proof of the fibre Hessian gap
uses (19.3), so it needs no such gauge.
It is still an eliminated-variable gap,
not a physical Hamiltonian mass gap.
\end{assessment}

\subsection{The background action cancels from the probability bound}

\begin{proposition}[Uniform section-centred denominator]
\label{prop:f12v15-denominator}
For some \(b>0,\beta_0<\infty\), uniformly in
the compact coarse field and region,
\[
 Z_{R,\beta}(C,\xi_J)
             \ge b^s\beta^{-N_R/2}e^{-\beta H_0(C)}
             \quad(\beta\ge\beta_0).
 \tag{19.16}
\]
No small-curvature or stationarity assumption
is required for this inequality.
Under (19.12), a possibly smaller \(b\)
gives the stronger bound with \(G_R(C)\)
in place of \(H_0(C)\).
\end{proposition}
\begin{proof}
The section chart has actual Haar and pivot
density at least \(b_0^s\), with uniform
radius and uniform upper Hessian bound.
Taylor's theorem therefore gives
\[
 f_R(z;C)\le H_0(C)+g(C)\cdot z+K_0|z|^2
 \tag{19.17}
\]
on that chart. In the symmetric coordinate
cube \(Q_\beta=[-t_0/\sqrt\beta,t_0/\sqrt\beta]^{N_R}\),
the quadratic term is at most \(K_0N_Rt_0^2/\beta\).
Pairing \(z\) with \(-z\), or using convexity
of the exponential and zero mean under
uniform cube measure, gives
\[
 \frac1{|Q_\beta|}\int_{Q_\beta}e^{-\beta g\cdot z}\,dz
                                                    \ge1.
 \tag{19.18}
\]
Integrate the pointwise lower bound on
the Boltzmann weight from (19.17), using
the constant lower Jacobian bound before
the symmetric integral. Since \(N_R\le73ds\),
all remaining constants are absorbed in
\(b^s\). This proves (19.16) without
pretending that the section is stationary.
For small curvature use the cube about
the minimum in (19.13), as in (18.14).
\end{proof}

\begin{theorem}[Conditional defect bound driven by curvature loads]
\label{thm:f12v15-defects}
Let \(F\) be \(k\) coarse edges internal to
\(R\), and \(\mathcal B_F\) their simultaneous
relative-logarithm failure event. There is
\(C_0\ge1\), independent of the volume,
region, compact coarse field and coupling,
such that for sufficiently large \(\beta\)
and almost every \(C_{\rm in}\),
\[
 \mathbb P_{\beta,\xi_J}(\mathcal B_F\mid C_{\rm in})
 \le\min\left\{1,\,
 C_0^s\exp[-\beta(\gamma k/2-A\mathcal L_R(C))]\right\},
 \qquad \gamma=\rho^2/2304.
 \tag{19.19}
\]
If \(k\ge\theta s\), \(0<\theta\le4\), and
\[
 \max_{p\in\mathcal P_R\cap\mathcal C}\delta(Q_p(C))
 \le\epsilon_\theta,\qquad
                576A\epsilon_\theta^2\le\gamma\theta/4,
 \tag{19.20}
\]
then
\[
 \mathbb P_{\beta,\xi_J}(\mathcal B_F\mid C_{\rm in})
                   \le\min\{1,C_0^s e^{-\beta\gamma k/4}\}.
 \tag{19.21}
\]
There is no factor involving \(H_0(C)\) in
this exponent and no small-link assumption.
\end{theorem}
\begin{proof}
The inherited defect argument uses rectangles
of edges internal to \(R\). They lie in
\(\mathcal P_R\) and, by the proof of
(19.3), contain only section-flat plaquettes.
It therefore proves the stronger threshold
\(H_{\rm v}\ge\gamma k\) on \(\mathcal B_F\),
not merely \(H_{\xi_J}\ge\gamma k\).
Use the second inequality in (19.7):
\[
 H_{\xi_J}\ge H_0+\gamma k/2-A\mathcal L_R+D/(4A)
                           \quad\hbox{on }\mathcal B_F.
 \tag{19.22}
\]
Discard the nonnegative pivot summands of
\(D\). The joint root-pivot density is at
most \(M_0^m\,dY\,dC_{\rm in}\), by f11 V76.
The same positive-test-function argument
as V13--V14 consequently yields
\[
 Z_{\mathcal B_F,R,\beta}
 \le M_0^m I(\beta/(4A))^{D_R}
       e^{-\beta H_0}
       e^{-\beta(\gamma k/2-A\mathcal L_R)}
 \le B^s\beta^{-N_R/2}e^{-\beta H_0}
       e^{-\beta(\gamma k/2-A\mathcal L_R)}.
 \tag{19.23}
\]
The Gaussian integral bound with \(4A\)
instead of \(2A\) follows by the identical
compact-group argument in V13.
Divide by (19.16). Both the Gaussian
power and the entire background Boltzmann
factor cancel, proving (19.19).
Equations (19.11) and (19.20) imply
\(A\mathcal L_R\le\gamma k/4\), which
proves (19.21). This argument keeps the
full compact conditional measure and all
pivot preimages; it does not restrict the
numerator to the good chart.
\end{proof}

\begin{corollary}[Periodic full-fibre density estimate]
When \(R\) is the entire periodic lattice,
\(s=n\) and there is no exterior restriction.
Equations (19.19)--(19.21) hold for the
actual compact conditional fibre for
Haar-almost every full retained field \(C\).
They allow uniformly small coarse
plaquette curvature with an extensive
total coarse action, rather than requiring
the small-total-action window of V13.
They still concern \(k\ge\theta n\) for
the simplified bound (19.21), not a
volume-uniform probability for a single
specified defect at arbitrary \(C\).
\end{corollary}
\begin{proof}
For the full region all internal and boundary
fine variables are integrated, so the
exterior convention is empty. Apply the
theorem with \(s=n\). The pointwise
curvature condition imposes no bound on
the number of contributing coarse plaquettes
and hence no fixed bound on their total action.
\end{proof}

\begin{record}[Geometric cross-check]
An independent finite enumeration on a
side-three coarse periodic lattice checked
all 324 reference rectangles: each has sixty
plaquettes and none is a corner plaquette.
It also checked that every corner-plaquette
edge is a boundary link with exactly three
corner incidences. These checks confirm
the counts used in (19.3) and (19.11);
the analytic proofs above establish their
volume-independent statements.
\end{record}

\begin{assessment}[What changed and what remains unproved]
The earlier full-background debit is not
an unavoidable geometric cost: in the
section-exterior regime it cancels, and
the residual cost is the explicitly
quadratic curvature load. The actual
small-curvature fibre minimum also now
exists without a small-amplitude or
global gauge hypothesis.

For a proper subregion, however, an exact
section exterior is a specified boundary
law, not a typical Wilson exterior. Its
set generally has zero Haar measure.
It cannot be inserted as an almost-sure
condition in a bulk probability estimate.
V14's open protected-data regime and the
present curved-background boundary regime
are different hypotheses; neither is
silently substituted for the other.

No curvature probability estimate for the
retained blocked measure has been proved
here, and neither the full compact density's
coarse derivatives nor an iteration with
physical scale-uniform constants follows
from (19.19).
\end{assessment}

\begin{decision}[Next comparison: nonzero exterior fluctuations]
Keep the section-flat/corner decomposition
and perturb the exterior away from its
section values. Identify the resulting
boundary mismatch terms in the path
coercivity, the corner comparison and
the exact constrained repair. The aim
is an open, curvature-based conditional
regime with the background action still
cancelled. Do not reintroduce the full
regional action by estimating it before
subtracting the reference, and do not
assume the reference remains a feasible
point when the exterior changes.
\end{decision}

\begin{record}[V15 append record]
The V14 predecessor has SHA--256
\begin{center}\scriptsize\ttfamily
B8983E05EBE2BDFB2F7D049A9D4FA9EF42E80E6DFBDB860AD7A4A320E7837437.
\end{center}
Its complete body is retained before
this appendix, apart from the title version.
The required companion audit was performed.
Only active V14 is replaced by V15;
all family archives remain.
\end{record}

\begin{firewall}[Position after V15]
Proved: section-flat coercivity, an exact
background-subtracted action bound,
small-curvature local minima with compact
holonomies retained, and curvature-load
defect estimates for the declared boundary
law and the full periodic fibre.
Not proved: the fluctuating-exterior
extension, general local defect control,
full compact coarse derivative estimates,
scale iteration, continuum reconstruction,
or the full Yang--Mills mass gap.
\end{firewall}

\section{V16: an open curved-background conditional regime}
\label{sec:f12v16}

V15 removed the whole background action from the
defect exponent, but for a proper region it required
an exact section exterior. This section permits
nonzero exterior fluctuations. It constructs an
actual feasible pivot repair and prices the
boundary mismatch before comparing the two
partition functions. It is not an assumption
that the unmodified section remains in the fibre.

\subsection{Actual coarse data and a finite mismatch collar}

Use the region and variable sets of V14.
Freeze arbitrary exterior links at \(\xi\).
The coarse outputs on edges outside \(E(R)\)
depend only on \(\xi\); call them \(C_{\rm out}\).
For a prescribed affected output \(C_{\rm in}\),
write \(C=(C_{\rm in},C_{\rm out})\).
This is the actual complete coarse output of
every configuration in the conditional fibre.
It is not an independently imposed replacement
for the frozen coarse data.

Start with all frozen links in
\(\mathcal P_R\cup\mathcal T_R\), and all
frozen internal links used by affected block
equations. Whenever this set contains a
boundary link on an unaffected coarse edge,
adjoin all fine links used in that edge's
block equation. Denote the resulting set of
frozen non-tree links by \(\mathcal F_R\).
This closure adds no new boundary family:
each block uses only its own eight boundary
links and internal links in its endpoint cells.
There is a fixed \(C_F<\infty\) with
\[
                  |\mathcal F_R|\le C_F s.
 \tag{20.1}
\]
Indeed every added stencil has bounded size
and lies a fixed number of cell steps from
the interface of \(R\). When \(R\) is the
whole lattice, \(\mathcal F_R\) is empty.
Keeping this actual collar, instead of all
exterior links, is important below.

For \(l\in\mathcal F_R\), define
\[
 b_l=\|\xi_l-J(C)_l\|_{\rm HS},\qquad
 \mathcal E_R(\xi)=\sum_{l\in\mathcal F_R}b_l^2,
 \qquad \eta=\max_{l\in\mathcal F_R}b_l.
 \tag{20.2}
\]
Use \(\eta=0\) for an empty collar. The quantity
\(\mathcal E_R\) is independent of \(C_{\rm in}\):
a frozen boundary link belongs to an unaffected
coarse edge. Internal section links equal \(I\).
For varied links define
\[
 d_l=\|V_l-J(C)_l\|_{\rm HS},\quad
 D=\sum_{l\in\Lambda_R}d_l^2,\quad
 H_0(C)=\sum_{p\in\mathcal P_R}s_p(J(C)).
 \tag{20.3}
\]
Unlike V15, \(J(C)|_{\Lambda_R}\) need not solve
the affected block equations with exterior \(\xi\).
The number \(H_0\) will first be used only as
an algebraic reference action.

For every link in a corner plaquette of
\(\mathcal P_R\), define
\[
 a_l(C)=\sum_{\substack{p\in\mathcal P_R\cap\mathcal C\\
                                 l\in\partial p}}\delta(Q_p(C)),
 \qquad
 L_{\rm in}=\sum_{l\in\Lambda_R}a_l^2,\qquad
 L_{\rm out}=\sum_{l\in\mathcal F_R}a_l^2,\qquad
 L_R=L_{\rm in}+L_{\rm out}.
 \tag{20.4}
\]
Terms with no corner incidence are zero. Put
\[
                  M_R=\sum_{l\in\mathcal F_R}a_l b_l .
 \tag{20.5}
\]
All these distances and loads are invariant under
simultaneous cell-constant gauge transformations.
No smallness of \(C_e\) relative to \(I\) is imposed.

\begin{proposition}[The strict collar condition is locally open]
For sufficiently small fixed \(\eta_*>0\),
the conditions \(b_l<\eta_*\) on \(\mathcal F_R\)
place every unaffected block output used here
on its good smooth branch. Consequently, after
adding any strict small-curvature condition on
the \(Q_p(C)\) that occur in (20.4), these
conditions define an open set of the relevant
exterior and affected coarse data.
\end{proposition}
\begin{proof}
For an unaffected boundary family included in
the collar, all its internal factors have
distance at most \(\eta_*\) from \(I\), and
all eight boundary links have distance at
most \(\eta_*\) from their actual \(C_e\).
Every two-link product is therefore within
\(2\eta_*\) of \(C_e\); every relative product
is within \(4\eta_*\) of \(I\).
Choose \(4\eta_*<\rho\). Such a block is
strictly in its good logarithm domain.
All of its input links were included in
the collar closure. Its actual output
therefore varies smoothly on a neighbourhood
of these data. The finitely many norms in
the strict conditions are continuous there.
Outputs elsewhere are irrelevant to them.
\end{proof}

\subsection{The numerator with exterior fluctuations}

Let \(H_{\rm v}\) be the section-flat part of
\(H_\xi\), as in V15. The exterior-only
plaquette sum of V14 is now
\[
 E_{\rm v,out}(\xi)=
       \sum_{p\in\mathcal T_R\setminus\mathcal P_R}s_p(\xi).
 \tag{20.6}
\]
Every such plaquette is section-flat by V15.

\begin{proposition}[Relative action with a finite boundary debit]
\label{prop:f12v16-action}
For arbitrary configurations in the actual
conditional fibre, with no smallness yet required,
\[
 D\le A(H_{\rm v}+E_{\rm v,out}),\qquad
 E_{\rm v,out}\le12\mathcal E_R,
 \tag{20.7}
\]
and
\[
 H_\xi-H_0
       \ge H_{\rm v}-\sum_{l\in\Lambda_R}a_l d_l-M_R.
 \tag{20.8}
\]
Consequently
\[
 H_\xi-H_0\ge
 \tfrac12H_{\rm v}+\frac{D}{4A}
       -\frac12E_{\rm v,out}-A L_{\rm in}-M_R.
 \tag{20.9}
\]
\end{proposition}
\begin{proof}
Retain all varied-link distances, including
pivots, in the regional path proof of V14.
As in V15, only section-flat plaquettes
occur in its stencils. Splitting the
stencil sum into affected and frozen
plaquettes gives the first inequality.
For a frozen plaquette in (20.6), its
section holonomy is \(I\), so telescoping
against that section gives
\(\delta(\xi_p)\le\sum_{l\in\partial p}b_l\).
Hence \(s_p(\xi)\le2\sum_{l\in\partial p}b_l^2\).
Each fine link is in at most six plaquettes,
proving the factor twelve.

Apply the exact Hilbert--Schmidt difference
identity (19.9) to each corner plaquette,
now retaining both the varied-link and
frozen-link replacements. Their linear
loads are respectively \(\sum a_ld_l\)
and \(M_R\). Section-flat plaquettes
contribute their nonnegative action,
which proves (20.8).
Use half the bound
\(H_{\rm v}\ge D/A-E_{\rm v,out}\) and
\(a_ld_l\le d_l^2/(4A)+Aa_l^2\) to
obtain (20.9).
\end{proof}

\subsection{A feasible repair, not the unmodified section}

\begin{proposition}[Boundary-local pivot repair]
\label{prop:f12v16-repair}
There are fixed \(0<\eta_0<\eta_*,K_p<\infty\)
such that if \(\eta<\eta_0\), one can keep
all varied internal links equal to \(I\),
all varied nonpivot boundary links equal
to \(C_e\), and solve only the \(m\) root
pivots to obtain \(V_0(C,\xi)\) with
\[
 K^0_{E(R)}(V_0,\xi)=C_{\rm in},\qquad
 \sum_{l\in\Lambda_R}\|V_{0,l}-J(C)_l\|_{\rm HS}^2
                                   \le K_p\mathcal E_R.
 \tag{20.10}
\]
All affected blocks are good. A uniform
root-pivot chart about this central repair
has \(D_R\) free group variables and actual
coordinate density bounded below by \(b_0^s\).
\end{proposition}
\begin{proof}
Write each root pivot as \(C_e\exp v_e\).
At zero exterior internal deviations,
all two-link products in the central
configuration equal \(C_e\), and the
right-coordinate derivative of the block
equation is \(I/8\). Exterior internal
links are small parameters. The fixed-size
implicit solve of f11 V72--V77, uniform
in compact \(C_e\), gives
\[
 \|v_e\|\le K_0
       \sum_{\substack{l\ {\rm frozen\ internal}\\
                         l\ {\rm in\ the\ block\ stencil\ of}\ e}}b_l .
 \tag{20.11}
\]
The sum is over a uniformly bounded number
of links. If it is empty, the solution
is \(v_e=0\). Squaring (20.11) and summing
over edges uses only bounded stencil
multiplicity, and proves (20.10).
The group exponential has a uniform
local Lipschitz bound.

The other free coordinates are the
internal logarithms and the right
boundary deviations, exactly as in
f11 V71--V77. Taking a sufficiently
small uniform radius keeps the pivot
solve regular and all products good.
The independent pivot Jacobians and
free-link Haar factors give the stated
lower density. No exterior link and
no unaffected block is changed.
\end{proof}

\begin{proposition}[Action price of the feasible repair]
\label{prop:f12v16-price}
Let \(H_{\rm rep}(C,\xi)=H_\xi(V_0(C,\xi))\).
With a fixed \(K_1<\infty\),
\[
 H_{\rm rep}-H_0
                \le A L_{\rm in}+M_R+K_1\mathcal E_R.
 \tag{20.12}
\]
For a fixed \(b>0\) and sufficiently large
\(\beta\), the full compact regional density satisfies
\[
 Z_{R,\beta}(C_{\rm in},\xi)
       \ge b^s\beta^{-N_R/2}e^{-\beta H_{\rm rep}(C,\xi)}
 \tag{20.13}
\]
for Haar-almost every affected output.
\end{proposition}
\begin{proof}
Write \(e_l=\|V_{0,l}-J(C)_l\|_{\rm HS}\)
on varied links and \(e_l=b_l\) on the
frozen links in the affected plaquettes.
By (20.10) their squared sum is at most
\((1+K_p)\mathcal E_R\). The squared
plaquette telescoping bound and incidence
at most six give
\[
 \sum_{p\in\mathcal P_R}
           \|(V_0,\xi)_p-Q_p(C)\|_{\rm HS}^2
                       \le24(1+K_p)\mathcal E_R.
 \tag{20.14}
\]
For section-flat plaquettes \(Q_p=I\).
Use the full equality in (19.9), bounding
the linear terms from above on corners.
It gives
\[
 H_{\rm rep}-H_0
 \le12(1+K_p)\mathcal E_R
       +\sum_{l\in\Lambda_R}a_l e_l+M_R .
 \tag{20.15}
\]
Here and in (20.14), \(Q_p=J(C)_p\) for
all affected plaquettes. Apply
\(a_le_l\le Aa_l^2+e_l^2/(4A)\) and
(20.10). One may take
\(K_1=12(1+K_p)+K_p/(4A)\).

The chart of Proposition~\ref{prop:f12v16-repair}
has uniform upper Hessian bounds for its
action, by compactness of the coarse
parameters and bounded local stencils.
At its centre \(z=0\) the action is
\(H_{\rm rep}\), but its gradient need
not vanish. Integrate in a symmetric
cube of radius \(t_0/\sqrt\beta\).
Taylor's upper bound on the action
and the symmetric exponential inequality
(19.18) yield (20.13) after using the
constant lower Jacobian bound. Thus
stationarity of the repaired section
has not been assumed.
\end{proof}

\subsection{An open conditional defect estimate}

\begin{theorem}[Curvature and mismatch, with the background cancelled]
\label{thm:f12v16-main}
There are \(C_0\ge1,K_2<\infty\),
independent of region, volume and compact
coarse link amplitudes, such that under
\(\eta<\eta_0\), for all sufficiently
large \(\beta\) and almost every affected
output,
\[
 \mathbb P_{\beta,\xi}(\mathcal B_F\mid C_{\rm in})
 \le\min\left\{1,\,
 C_0^s\exp\left[-\beta\left(
 \frac{\gamma k}{2}-2A L_R(C)-K_2\mathcal E_R(\xi)
                         \right)\right]\right\}.
 \tag{20.16}
\]
Here \(F\) is any specified set of \(k\)
internal coarse edges of \(R\), as in V14--V15.
\end{theorem}
\begin{proof}
On \(\mathcal B_F\), the same section-flat
plaquette threshold gives \(H_{\rm v}\ge\gamma k\).
From (20.9), the pointwise Boltzmann
weight on this event is at most
\[
 e^{-\beta H_0}
 e^{-\beta(\gamma k/2-E_{\rm v,out}/2-A L_{\rm in}-M_R)}
                         e^{-\beta D/(4A)}.
 \tag{20.17}
\]
Integrate using the joint pivot density
bound \(M_0^m\), after discarding the
pivot distances from \(D\). The free
group integrals give \(I(\beta/(4A))^{D_R}\).
The numerator is therefore bounded by
a constant to the power \(s\), times
\(\beta^{-N_R/2}\), times the first
two factors in (20.17).

Divide by (20.13), and use (20.12).
The remaining exponent is at least
\[
 \gamma k/2-E_{\rm v,out}/2
                         -2A L_{\rm in}-2M_R-K_1\mathcal E_R.
 \tag{20.18}
\]
For the frozen terms,
\[
 M_R\le A L_{\rm out}+\mathcal E_R/(4A).
 \tag{20.19}
\]
Combine this with (20.7).
Taking \(K_2=6+K_1+1/(2A)\) proves (20.16).
The factors \(e^{-\beta H_0}\) cancel
algebraically; \(H_0\) was never required
to be the action of a feasible unmodified
section. The repair supplies the actual
conditional denominator.
\end{proof}

\begin{corollary}[A populated open curvature-based regime]
\label{cor:f12v16-open}
Fix \(0<\theta\le4\). Suppose \(k\ge\theta s\),
\(\delta(Q_p(C))<\epsilon\) for every corner
plaquette of \(\mathcal P_R\), and
\(b_l<\eta\) on the completed collar.
There is a fixed \(C_L\), for example
\(C_L=5832\), with
\[
            L_R\le C_L\epsilon^2s,\qquad
            \mathcal E_R\le C_F\eta^2s .
 \tag{20.20}
\]
Choose positive \(\epsilon,\eta\) so
\[
 \eta<\eta_0,\qquad
              2AC_L\epsilon^2+K_2C_F\eta^2
                                      \le\gamma\theta/4.
 \tag{20.21}
\]
Then the full compact conditional probability satisfies
\[
 \mathbb P_{\beta,\xi}(\mathcal B_F\mid C_{\rm in})
                   \le\min\{1,C_0^s e^{-\beta\gamma k/4}\}.
 \tag{20.22}
\]
This data condition is gauge invariant
and open in the relevant variables.
It includes nonzero exterior deviations
and arbitrary retained compact holonomies.
\end{corollary}
\begin{proof}
Each corner edge is a boundary link with
at most three corner incidences, so
\[
 \sum_l a_l^2
 \le3\sum_l\sum_{\substack{p\in\mathcal P_R\cap\mathcal C\\
                                  l\in\partial p}}\delta(Q_p)^2
 =12\sum_{p\in\mathcal P_R\cap\mathcal C}\delta(Q_p)^2.
 \tag{20.23}
\]
Every link in this sum is varied or
belongs to \(\mathcal F_R\). Thus the
left side is \(L_R\). Use
\(|\mathcal P_R|\le486s\) to obtain
\(C_L=12\cdot486\), and use (20.1)
for the mismatch bound. Conditions
(20.20)--(20.21) make the debit in
(20.16) at most \(\gamma\theta s/4\le\gamma k/4\).
Gauge invariance follows from simultaneous
endpoint transformations of \(V,\xi,C\)
and conjugation of plaquette holonomies.
Openness was checked using the completed
collar. The inequalities are strict in
the data and admit sufficiently small
nonzero perturbations of any flat
section datum, so this is not an
exact-equality boundary condition.
\end{proof}

\begin{proposition}[The local minimum persists with the exterior]
For sufficiently small \(\epsilon+\eta\)
in the preceding regime, the feasible
chart about \(V_0(C,\xi)\) has a unique
small stationary minimum, with coordinates
\[
             \|z_R^*(C,\xi)\|_\infty\le K(\epsilon+\eta)
 \tag{20.24}
\]
and a volume-independent positive fibre
Hessian throughout a fixed smaller box.
This is still a local minimum.
\end{proposition}
\begin{proof}
Hold the full compact reference \(C\)
fixed for this comparison. At zero
exterior mismatch the section chart
has the positive Hessian proved in
(19.14)--(19.15). The pivot chart
and the action are smooth local
functions of the frozen internal
parameters and of exterior boundary
deviations relative to \(C\).
The row-norm change of the Hessian
is at most \(K\eta\), by bounded
stencil overlap. Its gradient at
the repaired centre is at most
\(K(\epsilon+\eta)\) in max norm.
Choose the radii so the Hessian
gap persists. The finite-range
Neumann argument supplies a uniform
max-norm inverse. The same contraction
used in V15 now has the displayed
gradient size, and proves (20.24)
and uniqueness on the smaller box.
The comparison may hold \(C_{\rm out}\)
fixed as an auxiliary parameter:
these outputs are unaffected by the
integrated variables, and no derivative
of the exterior block map is needed
for this fibre Hessian assertion.
\end{proof}

\begin{corollary}[Transfer to protected bulk events]
Let \(\mathcal G_R^{\rm curv}\) denote
the open data condition of
Corollary~\ref{cor:f12v16-open}, using
the actual unaffected coarse outputs.
Under the finite-volume Wilson law,
\[
 \mu_\beta(\mathcal B_F\cap\mathcal G_R^{\rm curv})
       \le C_0^s e^{-\beta\gamma k/4}
                            \mu_\beta(\mathcal G_R^{\rm curv}).
 \tag{20.25}
\]
For connected regions containing a fixed
root, the V14 counting argument applies
unchanged: with
\(q=1024C_0e^{-\beta\gamma\theta/4}<1\),
the probability of such a protected
defect-dense region of size at least
\(s_0\) is at most \(q^{s_0}/(1-q)\).
\end{corollary}
\begin{proof}
The data condition depends only on
the exterior and the affected coarse
output. Condition on these data and
apply (20.22). The full joint pivot
absolute continuity and the positive
repair density handle the exceptional
coarse null sets. Integration proves
(20.25). The region and subset counts
are precisely those already proved
in V14; no independence is added.
\end{proof}

\begin{assessment}[What is and is not now populated]
The numerator and denominator now
share a background-subtracted comparison
on an open, gauge-invariant curved-data
regime. The boundary debit is the
squared mismatch on a finite collar,
not the action of the whole exterior
or the whole regional background.
The exact-section restriction of V15
has therefore been genuinely relaxed.

It is not proved that all relevant
regions satisfy \(\mathcal G_R^{\rm curv}\).
Neither the probability of its failure
nor the retained coarse curvature law
has been controlled by this argument.
In particular (20.25) alone is not
an unconditional polymer expansion
or a physical correlation-decay theorem.
The full compact coarse derivatives,
scale iteration and continuum
reconstruction also remain open.
\end{assessment}

\begin{decision}[Next upstream task: charge the failed data condition]
Separate failure of the curvature
condition from failure of the exterior
mismatch collar. Use the actual
root-block path estimates to determine
which nearby fine plaquette costs
each failure forces, retaining their
overlaps and the unaffected coarse
outputs. A deterministic cost is
only the first step: the comparison
must then address the probability
of unprotected regions without
conditioning again on the missing
protection. Next companion audit
and direction review: V20.
\end{decision}

\begin{record}[V16 append record]
V15 had SHA--256
\begin{center}\scriptsize\ttfamily
8194AA73F3C586CAD90B0245C2A8DFA88DA009B7F5C031F2EB2D25A27D4B8F17.
\end{center}
Its full body is retained apart from
the title version. The present arguments
are analytic. Only active V15 is
replaced by V16; all archives remain.
\end{record}

\begin{firewall}[Position after V16]
Proved: a feasible exterior-dependent
repair, explicit curvature and squared
boundary-mismatch debits, an open
gauge-invariant conditional regime,
and its protected bulk-event bound.
Not proved: control of unprotected
regions, a complete compact effective
action, continuum scale iteration,
or the full Yang--Mills mass gap.
\end{firewall}

\section{V17: bare-Wilson probabilities for the actual obstruction carriers}
\label{sec:f12v17}

The next obstruction in V16 was failure of its
curvature and exterior-mismatch data condition.
The archives already contain a suitable unconditional
plaquette estimate: f6 V28 establishes the periodic
Wilson chessboard framework, and f6 V41 applies it
to identity-distance defects. We reuse that framework,
not the conditional estimate on a region whose
protection has failed. The new work identifies
finite fine-plaquette witnesses for the actual
root-block obstructions and their joint multiplicities.

\subsection{Imported probability input and normalization}

Work here with the bare periodic Wilson law
\[
 d\mu_{\Lambda,\beta}(U)=Z^{-1}e^{-\beta S_0(U)}dU,
 \qquad S_0(U)=\tfrac12\sum_p\|U_p-I\|_{\rm HS}^2.
 \tag{21.1}
\]
The fine periods are powers of two, at least four,
as required by the inherited finite chessboard
iteration. Fix a compact Lie group in the faithful
unitary representation already used in the notes,
and let \(d=\dim G\). Choose \(r_H,c_H>0\) so
\[
 h(r):=\int_G{\bf1}_{\{\delta(U)<r\}}dU
                  \ge c_H r^d\quad(0<r\le r_H).
 \tag{21.2}
\]

\begin{record}[External framework, not a new continuum theorem]
The reflection-positivity and periodic chessboard
input is the one used in f6 V28--V41. For a
primary-source formulation for lattice gauge
theories, see D.~Brydges, J.~Fr\"ohlich and
E.~Seiler, \emph{On the Construction of Quantized
Gauge Fields, I: General Results}, IHES/P/78/237
(1978), Section~2.4, Theorems~2.8 and~2.10,
\href{https://repo-archives.ihes.fr/FONDS_IHES/I_Prepublications/FROHLICH/1978-1980/P_78_237/P_78_237_web.pdf}
{primary manuscript}. Those statements concern
reflection positivity and chessboard estimates
under the specified boundary conditions, not
a four-dimensional pure Yang--Mills mass-gap proof.
The source's framework and boundary conditions
were checked for this use.
\end{record}

\begin{proposition}[The inherited bound in the current action units]
\label{prop:f12v17-chess}
For \(\beta\ge\max(1,r_H^{-2})\), set
\[
 Q_t(\beta)=c_H^{-64}\beta^{32d}
                      \exp(768-\beta t^2/2),\qquad
 p_t(\beta)=\min\{1,Q_t(\beta)^{1/96}\}.
 \tag{21.3}
\]
For every finite set \(J\) of distinct fine plaquettes,
\[
 \mu_{\Lambda,\beta}\{\delta(U_p)\ge t\ (p\in J)\}
                              \le p_t(\beta)^{|J|}.
 \tag{21.4}
\]
\end{proposition}
\begin{proof}
This is the normalization and dimension version
of f6 V41's identity-defect bound. The Wilson
weight has the required reflection positivity:
after dropping its constant, its plaquette
factor is
\(\exp[\frac\beta2(\chi_V+\chi_{\bar V})]\).
The nonnegative character expansion used in
f6 V28 applies to any fixed unitary representation
\(V\), by tensor-product decomposition. Thus
the same block chessboard inequality applies.
We do not assume reflection positivity of a
blocked or response-generated action.

In a shifted \(2^4\)-block tiling there are
\(N=|\Lambda|/16\) blocks. Dissemination of
an identity defect forces action at least
\(Nt^2/2\). The all-small-link partition
lower bound, with radius \(r=\beta^{-1/2}\),
is \(h(r)^{4|\Lambda|}e^{-8\beta r^2\,6|\Lambda|}\).
Taking the \(N\)-th root of the ratio gives
\[
 e^{-\beta t^2/2}h(\beta^{-1/2})^{-64}e^{768}
                                      \le Q_t(\beta).
 \tag{21.5}
\]
The unchanged 96-colour construction of
f6 V28--V41 gives this bound to the number
of prescribed plaquettes of one colour.
One colour contains at least \(|J|/96\)
members. If \(Q_t<1\), discard all other
members and apply that estimate; otherwise
use the bound one. This proves (21.4).
The constants differ from the normalized
\(SU(3)\) formula in f6 because (21.1)
uses the unnormalized trace and (21.2)
uses the actual group dimension.
\end{proof}

\subsection{Fine witnesses for four kinds of root-block obstruction}

Apply the exact tree gauge to \(U\), obtaining
\(W\), and let \(C=K^0(W)\). Plaquette distances
are unchanged by this gauge transformation.
All events below are evaluated on this same
fine configuration and its actual coarse output.
Put
\[
            A_0=9+6L_{\log},\qquad
            A_{\rm c}=4+24L_{\log}.
 \tag{21.6}
\]

There are four labelled kinds of obstruction:
an internal mismatch
\(\delta(W_l)\ge\eta\) at a non-tree internal link;
a boundary mismatch
\(\|W_l-C_e\|_{\rm HS}\ge\eta\);
a coarse curvature defect
\(\delta(C_P)\ge\epsilon\); and a relative-logarithm
failure at a coarse edge. Labels of different
kinds remain distinct even when their stencils
overlap.

\begin{proposition}[Witness sets and thresholds]
\label{prop:f12v17-witness}
Each such obstruction forces
\(\delta(U_p)\ge t\) for some fine plaquette
in a declared set \(\Omega_\alpha\), where
\[
 t=\min\{\eta/A_0,\ \epsilon/A_{\rm c},\ \rho/6\}.
 \tag{21.7}
\]
The respective witness sets have at most
\(24,60,244,60\) plaquettes.
\end{proposition}
\begin{proof}
For an internal link use the plaquettes
of its cell: V13 gives
\(\delta(W_l)\le3\max_{p\ {\rm in\ cell}}\delta(U_p)\).
For a boundary link use \(\mathcal R_e\):
the full-branch V13 estimate gives
\(\|W_l-C_e\|_{\rm HS}\le
A_0\max_{p\in\mathcal R_e}\delta(U_p)\).
Here \(A_0\ge3\), so the first two cases
are covered by (21.7).

For a coarse plaquette \(P\), let
\(\square_P\) be its straight \(2\)-by-\(2\)
fine square, consisting of four fine
plaquettes, and define
\[
       \Omega_P=\square_P\cup
                        \bigcup_{e\in\partial P}\mathcal R_e .
 \tag{21.8}
\]
Its size is at most \(4+4\cdot60=244\).
If \(B_P\) is the product of the straight
root links, four square moves give
\(\delta(B_P)\le4\max_{\square_P}\delta(U_p)\).
On both branches,
\(\|C_e-B_e\|_{\rm HS}\le
6L_{\log}\max_{\mathcal R_e}\delta(U_p)\),
as in f11 V78. Replacing the four
coarse links in the plaquette product
therefore gives
\[
       \delta(C_P)\le A_{\rm c}
                         \max_{p\in\Omega_P}\delta(U_p).
 \tag{21.9}
\]

Finally, the straight root reference
needs no filling moves. Every relative
two-link product differs from it by
at most six square moves in \(\mathcal R_e\).
Failure of the actual logarithm condition
implies one relative product is outside
the logarithm neighbourhood, hence has
distance at least \(\rho\) from \(I\).
Its witness therefore has distance
at least \(\rho/6\).
These are pointwise implications,
including the fallback branch.
\end{proof}

\begin{proposition}[Mixed witness multiplicity]
\label{prop:f12v17-overlap}
A fixed fine plaquette occurs in the
witness sets of at most 138 different
obstruction labels. More precisely the
four contributions are bounded by
\[
                         17,\quad64,\quad49,\quad8.
 \tag{21.10}
\]
\end{proposition}
\begin{proof}
An internal-cell plaquette belongs
to one cell, which has seventeen
non-tree internal links. For boundary
labels, each coarse edge has eight
boundary links sharing its rectangle,
and a fine plaquette is in at most
eight rectangles: the bound is 64.
For coarse plaquette labels, a fine
plaquette belongs to at most one
straight \(2\)-by-\(2\) square of its
orientation. Each rectangle is used
by the six coarse plaquettes incident
to its edge, giving at most
\(1+8\cdot6=49\) labels.
For relative-logarithm labels there
is one label per rectangle, giving
eight. Add the four bounds. These
are bounds on distinct labelled
objects, not on independent events.
\end{proof}

\subsection{Joint probabilities without the regional protection hypothesis}

\begin{theorem}[Mixed obstruction-family suppression]
\label{thm:f12v17-mixed}
Let \(\mathcal A\) be any finite collection
of distinct labels of the preceding
four types, at the fixed thresholds.
Set
\[
          p_*(\beta)=\min\{1,244\,p_t(\beta)^{1/138}\}.
 \tag{21.11}
\]
Then
\[
 \mu_{\Lambda,\beta}
       \{\hbox{every obstruction in }\mathcal A\hbox{ occurs}\}
                              \le p_*(\beta)^{|\mathcal A|}.
 \tag{21.12}
\]
For fixed positive thresholds, \(p_*(\beta)\)
tends to zero exponentially up to a polynomial
factor as \(\beta\to\infty\), uniformly in
the permitted periodic volumes.
\end{theorem}
\begin{proof}
For each label choose a witnessing
fine plaquette from its set, and take
the union over all choices. There are
at most \(244^{|\mathcal A|}\) choices.
By (21.10), every assignment uses at
least \(|\mathcal A|/138\) distinct
fine plaquettes. For each assignment
apply (21.4) to those distinct
plaquettes. Summation gives
\((244p_t^{1/138})^{|\mathcal A|}\);
also use the bound one.
When \(Q_t<1\), the nontrivial factor is
\[
 244\,c_H^{-64/(96\cdot138)}
       \beta^{32d/(96\cdot138)}
       \exp\left(\frac{768-\beta t^2/2}{96\cdot138}\right),
 \tag{21.13}
\]
which proves the stated decay. Reflection
positivity is used only for the original
fine Wilson law through (21.4). We do
not apply it to the tree-gauge measure
or to a conditional blocked density.
\end{proof}

\begin{corollary}[The actual retained curvature has a joint bare bound]
Under the exact blocked law
\(\nu_\beta=(K^0\circ W)_*\mu_{\Lambda,\beta}\),
for every set \(\mathcal P\) of distinct
coarse plaquettes,
\[
 \nu_\beta\{\delta(C_P)\ge\epsilon\ (P\in\mathcal P)\}
              \le p_{\rm c}(\beta,\epsilon)^{|\mathcal P|},
 \qquad
 p_{\rm c}=\min\{1,244\,p_{\epsilon/A_{\rm c}}^{1/49}\}.
 \tag{21.14}
\]
Likewise simultaneous root-logarithm
failures on \(k\) distinct coarse edges
have probability at most
\[
       \min\{1,60\,p_{\rho/6}^{1/8}\}^{\,k}.
 \tag{21.15}
\]
\end{corollary}
\begin{proof}
Repeat the witness assignment proof
with only the relevant type, using
its size and multiplicity in
(21.10). Equation (21.14) is an
exact pushforward statement, not an
assumption that \(\nu_\beta\) has
a Wilson action.
\end{proof}

\begin{corollary}[Probability that a fixed region is unprotected]
Choose the positive \(\epsilon,\eta\)
in V16's open-data regime. For each
fixed deterministic region \(R\),
\[
 \mu_{\Lambda,\beta}\bigl((\mathcal G_R^{\rm curv})^c\bigr)
       \le\min\{1,(C_F+486)s\,p_*(\beta)\}.
 \tag{21.16}
\]
The \(p_*\) may use the additional
root-log threshold in (21.7), even
though that only weakens this bound.
\end{corollary}
\begin{proof}
Failure of the actual data condition
means either a mismatch on a link
of the completed collar or a coarse
curvature defect associated with
a corner plaquette of \(\mathcal P_R\).
There are at most \(C_Fs\) collar
links and \(486s\) such plaquettes;
repeated coarse labels only reduce
the latter count. Apply the one-label
case of (21.12) and the union bound.
The unaffected coarse outputs in this
argument are exactly those of the
fine field, as required in V16.
\end{proof}

\begin{corollary}[Localized obstruction clusters]
Connect two labelled objects when
their witness sets intersect. This
witness-overlap graph has degree
at most
\[
                         \Delta=244\cdot137=33428.
 \tag{21.17}
\]
If \(q=\Delta^2p_*<1\), the probability
that a specified labelled object is
contained in a connected set of at
least \(n\) occurring obstructions
is at most \(q^n/(1-q)\).
\end{corollary}
\begin{proof}
Each witness set has at most 244
plaquettes, each shared with at
most 137 other objects, proving
the degree bound. The depth-first
spanning-tree count used in V14
bounds rooted connected sets of
size \(s\) by \(\Delta^{2s}\).
Use (21.12) for each set and sum
the geometric series. Adjacency
here is a declared overlap graph;
no probabilistic independence of
nonadjacent objects is asserted.
\end{proof}

\begin{record}[Independent finite incidence check]
On a side-four coarse periodic lattice,
an enumeration found maximum witness
multiplicities \(17,64,24,8\) for
the four types, and maximum combined
multiplicity 113. The largest coarse
witness union had 145 plaquettes.
All obey the conservative analytic
bounds \(17,64,49,8\), 138 and 244.
The proofs above, not the finite
enumeration, establish uniformity
over lattice sizes.
\end{record}

\begin{assessment}[New probabilities, not a closed iteration]
The bare blocked curvature and the
actual failure carriers of V16 now
have unconditional joint estimates.
This is stronger than merely charging
a deterministic Wilson action or
reusing a protected conditional
estimate on unprotected data.

Nevertheless (21.16) is proportional
to \(s\), not exponentially small
in \(s\). One obstruction can spoil
many large candidate regions. A
cluster estimate for localized
obstructions is not a cluster
estimate for all regions whose
data condition fails.

Furthermore (21.12) does not survive
arbitrary conditioning as a logical
consequence of its unconditional
form. Exploring a defect region,
normalizing connected marked
correlations, or replacing the
bare action by an iterated effective
action requires additional estimates.
Although (21.14) is a property of
the exact first blocked measure,
it does not establish a new
reflection-positive Wilson form
on that coarse lattice.
\end{assessment}

\begin{decision}[Next interface: normalized local observables]
Use the joint obstruction bound
together with V16's conditional
comparison to examine a specified
local observable or connected
insertion. Keep the random choice
of region and its conditioning
explicit. The remaining question
is whether an exterior-stable
comparison cancels unmarked remote
defects in normalized quantities;
do not infer this from the mere
rarity of bare obstruction clusters.
Next companion audit and direction
review: V20.
\end{decision}

\begin{record}[V17 append record]
V16 had SHA--256
\begin{center}\scriptsize\ttfamily
44DEDB751F8D77C6A4817D31902C72B6E25AE9C023C733E985607569006E90FA.
\end{center}
Its full body is retained apart
from the title version. The
archived probability input was
reused with explicit normalization;
the primary reflection-positivity
reference was checked.
Only active V16 is replaced by
V17; all family archives remain.
\end{record}

\begin{firewall}[Position after V17]
Proved: fine-plaquette witnesses
and joint bare-Wilson bounds
for actual coarse curvature,
fine mismatches and root-log
failures, including fixed-region
unprotected-data probabilities.
Not proved: arbitrary conditional
or generated-action versions,
normalized connected-defect
cancellation, continuum scale
iteration, or the Yang--Mills
mass gap.
\end{firewall}

\section{V18: a smooth compact completion with globally invertible pivots}
\label{sec:f12v18}

The existing block \(K^0\) is piecewise smooth,
and its global density has not been shown to
have coarse derivatives. Before attempting
normalized score estimates on that density,
we test an explicit alternative completion.
It is denoted \(\widetilde K\) throughout:
it is not silently identified with \(K^0\).
The two maps agree on a smaller neighbourhood
of the same local fibre germ. Their difference
on local bounded observables will be controlled
using the bare probability input from V17.

The new point is a global pivot theorem. A
monotone cutoff of the logarithm has an upper
bound on its symmetric derivative even where
its derivative is negative. Together with
the two exact duplicate root paths, this
leaves a strictly positive pivot derivative
on the entire compact group.

\subsection{A conjugation-equivariant smooth truncated logarithm}

Use the invariant Hilbert--Schmidt inner product
on \(\mathfrak g\). Choose a conjugation-invariant
injective logarithm neighbourhood of \(I\).
Fix a smooth nonincreasing function
\(\chi:[0,\infty)\to[0,1]\), equal to one on
\([0,1/2]\) and zero on \([1,\infty)\).
For example, with
\(\varphi(t)=e^{-1/t}\) for \(t>0\) and zero
for \(t\le0\), take
\[
       \chi(s)=\frac{\varphi(1-s)}
                       {\varphi(1-s)+\varphi(s-1/2)}.
 \tag{22.1}
\]
For sufficiently small fixed \(\rho_1>0\), define
\[
 g_{\rho_1}(U)=
 \begin{cases}
 \chi(\|\log U\|/\rho_1)\log U,
       &U\hbox{ in the logarithm neighbourhood},\\
 0,&\hbox{otherwise}.
 \end{cases}
 \tag{22.2}
\]
The support lies strictly inside that
neighbourhood, so the zero extension is smooth.
The function is conjugation equivariant and
satisfies \(g_{\rho_1}(U^{-1})=-g_{\rho_1}(U)\).

Let \(D_Lg(U)\) denote the linear map
\[
        D_Lg(U)v=\left.\frac{d}{dt}g(e^{tv}U)\right|_{t=0}.
 \tag{22.3}
\]
For a real linear operator \(T\) write
\(\operatorname{Sym}T=(T+T^*)/2\).

\begin{proposition}[The derivative bound needed by the pivot]
\label{prop:f12v18-cutoff}
There are constants \(C_\chi,C_1\), independent
of sufficiently small \(\rho_1\), such that
for every \(U\in G\),
\[
 \|g_{\rho_1}(U)\|\le\rho_1,\qquad
 \|D_Lg_{\rho_1}(U)\|\le C_\chi,\qquad
 \operatorname{Sym}D_Lg_{\rho_1}(U)
                                  \le(1+C_1\rho_1)I.
 \tag{22.4}
\]
\end{proposition}
\begin{proof}
In logarithm coordinates put \(h(z)=
\chi(\|z\|/\rho_1)z\). Its derivative,
for \(z\ne0\), is
\[
 Dh(z)=\chi(s)I+
              s\chi'(s)\frac{z\otimes z}{\|z\|^2},
                  \qquad s=\|z\|/\rho_1.
 \tag{22.5}
\]
This is symmetric. Its tangential eigenvalues
are \(\chi(s)\), and its radial eigenvalue
is \(\chi(s)+s\chi'(s)\). Since \(\chi'\le0\),
every eigenvalue is at most one, while
their absolute values have a finite uniform
bound depending only on the fixed cutoff.
The formula extends smoothly at zero.

The map
\[
 L_zv=\left.\frac{d}{dt}\log(e^{tv}e^z)\right|_{t=0}
\]
is smooth near \(z=0\), equals \(I\) there,
and hence satisfies \(\|L_z-I\|\le C\|z\|\)
on a fixed smaller ball. On the support
of the derivative in (22.2),
\(D_Lg_{\rho_1}(e^z)=Dh(z)L_z\).
The norm bound and the upper symmetric
bound in (22.4) follow from (22.5)
and \(\|z\|\le\rho_1\). Outside this
support the derivative is zero.
The bound on the value is immediate.
No lower bound on the derivative of
the cutoff is needed beyond its finite norm.
\end{proof}

\subsection{The uniform global pivot theorem}

On the fixed tree slice, the root-reference
boundary link \(B\) appears in exactly two
of the sixteen path products, both equal
to \(B\); the other fourteen products
\(P_1,\ldots,P_{14}\) do not contain it.
This is the exact f11 V76 identity.
For arbitrary fixed \(P_j\in G\), put
\[
 X(B)=\frac1{16}\sum_{j=1}^{14}g_{\rho_1}(B^{-1}P_j),
 \qquad F_t(B)=B\exp(tX(B)),\quad 0\le t\le1.
 \tag{22.6}
\]

\begin{theorem}[All compact pivot maps are diffeomorphisms]
\label{thm:f12v18-pivot}
For sufficiently small fixed \(\rho_1>0\),
every map \(F_t:G\to G\) in (22.6) is a
smooth diffeomorphism. There are constants
\(m>0,M<\infty\), independent of \(P_j,t\),
such that every singular value of its
differential in invariant group metrics
lies between \(m\) and \(M\). One may
take \(m=1/16\) after shrinking \(\rho_1\).
The inverse and each of its fixed-order
derivatives depend smoothly and uniformly
on all fourteen compact parameters.
\end{theorem}
\begin{proof}
For the input variation \(B e^{sv}\),
the relative products vary as
\(e^{-sv}B^{-1}P_j\). Therefore
\[
 D_BX[v]=-A_Bv,\qquad
 A_B=\frac1{16}\sum_{j=1}^{14}D_Lg_{\rho_1}(B^{-1}P_j).
 \tag{22.7}
\]
From (22.4),
\[
 \|X\|\le14\rho_1/16,\quad
 \|A_B\|\le14C_\chi/16,\quad
 \operatorname{Sym}A_B\le
                       \frac{14}{16}(1+C_1\rho_1)I.
 \tag{22.8}
\]
Define the right-trivialized differential
of the exponential by
\[
 {\cal J}(Y)u=D\exp(Y)[u]\,e^{-Y}
             =\int_0^1 e^{s\,\operatorname{ad}Y}u\,ds .
 \tag{22.9}
\]
On the fixed small ball,
\(\|{\cal J}(Y)-I\|\le C\|Y\|\).
Differentiating \(F_t(B e^{sv})\) and
identifying output tangent vectors by
invariant translations gives, up to an
orthogonal adjoint factor, the operator
\[
                         M_t(B)=I-t{\cal J}(tX)A_B.
 \tag{22.10}
\]
In particular its singular values equal
those of the actual group differential.
Using (22.8)--(22.9),
\[
 \operatorname{Sym}M_t
      \ge\left(1-\frac{14t}{16}-C_2\rho_1\right)I
      \ge(1/8-C_2\rho_1)I .
 \tag{22.11}
\]
Choose \(C_2\rho_1\le1/16\).
Then \(\|M_tv\|\ge\langle v,M_tv\rangle/\|v\|
\ge\|v\|/16\). Its norm is also uniformly
bounded by (22.8)--(22.10).
Thus all differentials are invertible.
The decisive fraction is fourteen over
sixteen, not sixteen over sixteen.

To obtain global, rather than merely
local, invertibility, define the
time-dependent smooth vector field
\[
 V_t(B)=-(D_BF_t)^{-1}\,\partial_tF_t(B).
 \tag{22.12}
\]
The inverse here maps the tangent
space at \(F_t(B)\) to the one at \(B\).
Compactness and smoothness give a
complete flow \(\Psi_t\) on \(0\le t\le1\);
the reverse-time equation gives its
smooth inverse. Along that flow,
\[
             \frac{d}{dt}F_t(\Psi_t(B))=0.
 \tag{22.13}
\]
Since \(F_0\) is the identity,
\(F_t\circ\Psi_t=\mathrm{id}\), and
\(\Psi_t\) is a diffeomorphism. Thus
\(F_t=\Psi_t^{-1}\) is globally a
diffeomorphism. This argument works
component by component if necessary.

All functions are smooth on the compact
parameter space for fixed \(\rho_1\).
The uniform inverse-differential bound
and repeated implicit differentiation
give uniform bounds for every fixed
derivative order of \(F_t^{-1}\).
Those constants may depend on
\(\rho_1\); no claim is made as
\(\rho_1\) tends to zero.
\end{proof}

\subsection{A genuinely smooth full compact blocked density}

Define the modified coarse map by
\[
 \widetilde K_e(W)=B_e\exp\left[
       \frac1{16}\sum_{j=1}^{14}
                   g_{\rho_1}(B_e^{-1}P_j(W))\right].
 \tag{22.14}
\]
It is smooth, finite range, and gauge
equivariant under the residual
cell-constant gauge group. Equivariance
follows because each relative product
is conjugated at the target root and
\(g_{\rho_1}\) intertwines conjugation.
Also \(\widetilde K(J(C))=C\) for every
compact \(C\), since all relative
products at the section are \(I\).

\begin{theorem}[Global product coordinates for the modified fibre]
\label{thm:f12v18-global}
Let \(Y\in G^{45n}\) consist of all
non-tree internal and nonpivot boundary
links, and \(B\in G^{4n}\) of the root
pivots. Then
\[
                 (Y,B)\longmapsto(Y,\widetilde K(Y,B))
 \tag{22.15}
\]
is a global smooth diffeomorphism onto
\(G^{45n}\times G^{4n}\). Denote its
pivot solution by \(B=B(Y,C)\).
With normalized product Haar,
\[
 dY\,dB=j(Y,C)\,dY\,dC,\qquad
 j(Y,C)=\prod_e j_e(Y,C_e),\qquad
                      0<c\le j_e\le C_1<\infty.
 \tag{22.16}
\]
Each factor and solved pivot has a
fixed finite stencil and uniformly
bounded derivatives of each fixed order.
\end{theorem}
\begin{proof}
When \(Y\) is fixed, each edge output
depends on its own pivot only. Its
other fourteen products depend only
on the internal variables in its
two endpoint cells and its seven
nonpivot boundary links. Apply
Theorem~\ref{thm:f12v18-pivot} independently
to every edge. This gives (22.15)
and a smooth inverse.
The joint Jacobian is triangular
in \(Y,B\), so its Haar-volume
factor is the product of the
inverse pivot Jacobians. The
singular-value bounds give positive
upper and lower bounds per factor,
and the inverse theorem gives their
derivative bounds. The solution
does not introduce nonlocal
dependence on other pivots.
\end{proof}

\begin{corollary}[Positive smooth density and exact global scores]
\label{cor:f12v18-density}
For every finite volume and finite
\(\beta\ge0\), the modified blocked
law has a strictly positive smooth
density relative to coarse Haar:
\[
 \widetilde Z_\beta(C)=
 \int_{G^{45n}}e^{-\beta S_0(Y,B(Y,C))}
                             j(Y,C)\,dY>0.
 \tag{22.17}
\]
It is gauge invariant. In any fixed
local coarse coordinate chart put
\[
 V(Y,c)=\beta S_0(Y,B(Y,C(c)))-\log j(Y,C(c)),
 \qquad \widetilde W(c)=-\log\widetilde Z_\beta(C(c)).
 \tag{22.18}
\]
Then the normalized full compact
conditional expectation satisfies
\[
 \partial_a\widetilde W=\mathbb E V_a,\qquad
 \partial_a\partial_b\widetilde W
                         =\mathbb E V_{ab}-\operatorname{Cov}(V_a,V_b).
 \tag{22.19}
\]
The fixed-coordinate scores are
local sums, with their fixed-order
local derivative bounds of size
\(C_k(\beta+1)\), independently of volume.
\end{corollary}
\begin{proof}
The integrand in (22.17) is positive
and smooth on a compact product.
For each finite volume its derivatives
may be taken under the integral;
this proves smoothness and positivity.
Gauge invariance follows by changing
the internal and boundary group
variables by their endpoint gauge
transformations and using Haar and
Wilson invariance.
Differentiation of the normalized
integral gives (22.19) directly.
The reference integration measure
\(dY\) is independent of \(c\).
The density is already relative
to coarse Haar; no extra coordinate
Haar factor is inserted into
\(\widetilde W\) in (22.18).
Locality follows from (22.16),
the plaquette action, and bounded
overlap of their finite stencils.
Compactness supplies the stated
coefficient bounds, but does not
bound the covariance in (22.19)
by a volume-uniform decay estimate.
\end{proof}

\subsection{What is preserved, and the price of changing the completion}

\begin{proposition}[Identical germ, not identical global block]
\label{prop:f12v18-germ}
Whenever all relative products lie
in \(\exp\{\|z\|<\rho_1/2\}\) inside
the original good logarithm domain,
\(\widetilde K=K^0\).
Consequently all Taylor coefficients
of the original root block at the
identity fine field, calculated in
the same original local coordinates,
are unchanged. This includes the
local action and Haar-density jets
used in the preceding coefficient
calculations.
\end{proposition}
\begin{proof}
On that neighbourhood (22.2) is
exactly the original logarithm;
the two omitted duplicate terms
are identically zero for both maps.
The two smooth maps therefore
agree on an open neighbourhood,
not merely through some finite
Taylor order. Their local implicit
charts and density expressions in
the same coordinates agree there.
\end{proof}

\begin{assessment}[The original box may need to be smaller]
This assertion concerns a common
germ. The chosen \(\rho_1\) is not
asserted to contain every relative
product in the already fixed V12
box. To identify an entire local
box for the two completions, shrink
its radius to lie in the agreement
neighbourhood and apply the local
estimates with that radius.
No theorem for a full compact
conditional density is obtained
merely from equality of its germ.
\end{assessment}

\begin{proposition}[The deterministic compact bounds persist]
Choose \(L_{\log}\) large enough for
the original local norm comparison.
For every \(U\in G\),
\[
                    \|g_{\rho_1}(U)\|\le L_{\log}\delta(U).
 \tag{22.20}
\]
Consequently the bound
\(\|\widetilde K_e-B_e\|\le
6L_{\log}\max_{\mathcal R_e}\delta(U_p)\)
holds on the entire fine configuration
space. The section-flat distance
coercivity of V15 and the coarse
curvature witness of V17 therefore
hold for \(\widetilde K\) with the
same geometric constants.
\end{proposition}
\begin{proof}
On the support of \(g_{\rho_1}\),
use \(\chi\le1\) and the local
logarithm bound; outside it the
left side is zero. Apply the
six-square path comparison to
each \(B_e^{-1}P_j\) and sum
(22.20) in (22.14).
The proofs of (19.3) and (21.9)
use precisely this root-output
bound and unchanged path geometry,
so their inequalities apply.
This transfers the indicated
deterministic estimates, not
unproved conditional clustering.
\end{proof}

\begin{theorem}[Bare local-observable comparison of the two completions]
\label{thm:f12v18-observables}
Choose \(\sigma>0\) such that
\(\delta(U)<\sigma\) implies
\(U=\exp z\), \(\|z\|<\rho_1/2\),
inside the original good domain.
On the bare periodic Wilson laws
of V17, set
\[
              q_\beta=\min\{1,60\,p_{\sigma/6}(\beta)^{1/8}\}.
 \tag{22.21}
\]
For any set \(E\) of distinct coarse
edges, simultaneous disagreement
of \(\widetilde K_e\) and \(K^0_e\)
at all edges in \(E\) has probability
at most \(q_\beta^{|E|}\).
For any bounded coarse observable
\(F\) supported on a finite edge
set \(E_F\),
\[
 \left|\mathbb E F(\widetilde K(W))
                   -\mathbb E F(K^0(W))\right|
                      \le2\|F\|_\infty |E_F|q_\beta.
 \tag{22.22}
\]
Their marginal total variation
distance on \(E_F\), in the
supremum-over-events convention,
is at most \(|E_F|q_\beta\).
\end{theorem}
\begin{proof}
If every fine plaquette in
\(\mathcal R_e\) has distance
less than \(\sigma/6\), every
relative path product is within
\(\sigma\) of \(I\). The maps
then agree by
Proposition~\ref{prop:f12v18-germ}.
Thus disagreement has a witness
among the sixty plaquettes of
\(\mathcal R_e\) at threshold
\(\sigma/6\). A fine plaquette
is in at most eight such
rectangles. Apply the witness
assignment argument of V17.
For (22.22), couple both outputs
using the same original fine
configuration. The observable
values coincide unless some
edge in \(E_F\) disagrees.
Use the union bound and the
bound \(2\|F\|_\infty\) on the
difference. The coupling also
gives the total variation bound.
\end{proof}

\begin{assessment}[Scope of the smooth alternative]
The smooth product-fibre integral
(22.17) and its scores are now
well defined globally at finite
volume. They belong to the
explicit modified block
\(\widetilde K\), not to a
retroactive regularity claim
for \(K^0\). The agreement and
probability estimates state
precisely how the two are related.

The comparison is for the bare
Wilson law and bounded local
observables. If their support
grows with the regulator, it
requires \(|E_F|q_\beta\to0\);
no physical scaling law proving
that condition has been supplied.
It does not compare derivatives
of the two conditional densities
at fixed coarse values.
Nor does smoothness, gauge
equivariance, or invertibility
of the pivots prove reflection
positivity of the new coarse
law or a physical mass gap.
\end{assessment}

\begin{decision}[Next normalized-correlation calculation]
Use (22.17)--(22.19) as an
explicit smooth full-compact
model for a local source
calculation, keeping
\(\widetilde K\) named separately.
Compare bounded local-source
generating functions using
(22.22), and then examine the
conditional covariance rather
than assuming that rare defects
cancel from it. This supplies
a noncircular entry point for
the normalized-insertion task
left open in V17.
Next companion audit and
direction review: V20.
\end{decision}

\begin{record}[V18 append record]
V17 had SHA--256
\begin{center}\scriptsize\ttfamily
E2683C8801588ABBC076023107FA4E3C66ABCF7C265372FA53EEC59F6787A785.
\end{center}
Its entire body is retained
apart from the title version.
The construction and inverse
theorem above are analytic,
with an existence choice of a
sufficiently small fixed cutoff
radius, not a numerical radius
certificate. Only active V17
is replaced by V18; all
family archives remain.
\end{record}

\begin{firewall}[Position after V18]
Proved for the explicit smooth
alternative: globally invertible
pivots, a positive smooth compact
density, exact score identities,
unchanged local germ, and a bare
bounded-observable comparison.
Not proved: conditional covariance
decay, generated-action stability,
coarse reflection positivity,
continuum reconstruction, or
the full Yang--Mills mass gap.
\end{firewall}

\section{V19: normalized local sources and averaged conditional score errors}
\label{sec:f12v19}

The two compact completions remain distinct:
\(K^0\) is the original piecewise block,
and \(\widetilde K\) is V18's smooth block.
The V18 comparison concerned bounded
observables without sources. Here we
extend it to normalized local sources
and their connected derivatives, then
apply the same estimate within the
smooth compact conditional integral.
The resulting errors are not asserted
to decay with the distance between
insertions.

\subsection{A normalized comparison lemma}

\begin{proposition}[Complex source neighbourhood and logarithmic comparison]
\label{prop:f12v19-source}
Let \(X_1,\ldots,X_k,Y_1,\ldots,Y_k\)
be real random variables on one
probability space, with
\(|X_j|,|Y_j|\le1\). Suppose all
\(X_j=Y_j\) off an event \(D\),
of probability \(p\). Define
\[
 Z_X(h)=\mathbb E e^{\sum_jh_jX_j},\qquad
 Z_Y(h)=\mathbb E e^{\sum_jh_jY_j}.
 \tag{23.1}
\]
For any fixed \(0<r<\log2\), these
entire functions are nonzero on
\(\sum_j|h_j|\le r\). Their logarithms
with value zero at the origin satisfy
\[
 |\log Z_X(h)-\log Z_Y(h)|
       \le L_r p,\qquad
       L_r=\frac{2r e^r}{2-e^r}.
 \tag{23.2}
\]
For a multi-index \(\alpha\) with
\(|\alpha|=m\ge1\),
\[
 |\partial^\alpha(\log Z_X-\log Z_Y)(0)|
           \le\alpha!(k/r)^m L_r p.
 \tag{23.3}
\]
\end{proposition}
\begin{proof}
If \(\sum|h_j|\le r\), then
\[
 |Z_X(h)-1|,\ |Z_Y(h)-1|\le e^r-1<1.
 \tag{23.4}
\]
Thus their values lie in the same
disk centred at one, excluding zero,
on which the logarithm with
\(\log1=0\) is analytic. The source
exponents agree off \(D\), and
their difference has magnitude
at most \(2r\) on \(D\).
Integration of the exponential
along the line between the two
exponents gives
\(|Z_X-Z_Y|\le2re^r p\).
The segment joining \(Z_X,Z_Y\)
stays in the disk (23.4), where
the reciprocal has magnitude
at most \((2-e^r)^{-1}\).
Integrate the derivative of
the logarithm to prove (23.2).
Apply the multivariable Cauchy
estimate on the polydisk
\(|h_j|\le r/k\) to obtain
(23.3). A limiting smaller
polydisk justifies the same
closed-radius bound if needed.
\end{proof}

\begin{proposition}[Real tilted expectations, with their denominators]
\label{prop:f12v19-tilt}
In the same setting, let \(h\) be
real and \(\sum_j|h_j|\le r<\infty\).
Let \(A_X,A_Y\) be bounded by \(M\)
and equal outside an event \(D_A\).
Then
\[
 \left|
 \frac{\mathbb E A_Xe^{\sum h_jX_j}}{Z_X(h)}
 -
 \frac{\mathbb E A_Ye^{\sum h_jY_j}}{Z_Y(h)}
 \right|
                  \le4M e^{2r}\mathbb P(D\cup D_A).
 \tag{23.5}
\]
\end{proposition}
\begin{proof}
Both denominators are at least \(e^{-r}\).
The numerator difference is at most
\(2Me^r\mathbb P(D\cup D_A)\).
Also \(|Z_X-Z_Y|\le2e^r\mathbb P(D)\).
When subtracting the two fractions,
bound the second numerator by
\(M Z_Y\), since the real exponential
is positive. The two resulting
terms give (23.5). Thus the
normalizing denominators have
not been omitted.
\end{proof}

\subsection{Application to the two compact completions}

Let \(F_j\) be bounded real coarse
observables with finite edge supports
\(E_j\) and positive bounds \(M_j\).
On the common bare fine Wilson
probability space, put
\[
 X_j=\frac{F_j(K^0(W))}{M_j},\qquad
 Y_j=\frac{F_j(\widetilde K(W))}{M_j},
 \qquad E=\bigcup_jE_j .
 \tag{23.6}
\]
For identically zero observables
the corresponding variable can
simply be removed. Use the
disagreement factor \(q_\beta\)
from (22.21).

\begin{theorem}[Finite local-source stability of connected correlations]
\label{thm:f12v19-completions}
The logarithmic generating functions
for (23.6) satisfy (23.2) with
\(p=\min(1,|E|q_\beta)\). In
unnormalized observable units,
their connected derivatives obey
\[
 |\kappa_\alpha^{\,0}-\kappa_\alpha^{\,\sim}|
 \le \alpha!(k/r)^m L_r
             \left(\prod_j M_j^{\alpha_j}\right)
                         \min(1,|E|q_\beta).
 \tag{23.7}
\]
Here \(\kappa_\alpha\) is the derivative
of the logarithm with \(\alpha_j\)
insertions of \(F_j\). The real
tilted comparison (23.5) also applies,
with the tested observable's support
added to \(E\). All constants are
independent of the ambient volume.
\end{theorem}
\begin{proof}
Couple the two outputs by the same
fine configuration. The variables
in (23.6) agree unless a coarse
edge in \(E\) has different outputs.
V18 bounds the probability of this
union by \(|E|q_\beta\).
Apply the preceding two propositions
and restore the factors \(M_j\).
This compares normalized connected
observables under the two complete
pushforward laws, not two
unnormalized local integrands.
\end{proof}

\begin{assessment}[Source variables are not coarse-field coordinates]
The complex variables \(h_j\) in
(23.1) are independent bounded
source parameters. Holomorphy in
these variables follows from
boundedness of the insertions.
It is not holomorphy of the
smooth cutoff block in the
coarse field, and (23.7) is
not a comparison of derivatives
of \(-\log Z(C)\) at fixed \(C\).
The support factor is for the
union of the actual local
supports, not the whole lattice.
If those supports grow, their
growth still has to be compared
with \(q_\beta\).
\end{assessment}

\subsection{A conditional version that can be averaged exactly}

Let \(\widetilde\nu_\beta\) be the
normalized smooth coarse law and
\(\pi_{\beta,C}\) the normalized
compact fibre law in (22.17).
For a fine event \(D\), define
\[
                       p(C)=\pi_{\beta,C}(D).
 \tag{23.8}
\]
The global product coordinates
give the exact disintegration
\[
             \int p(C)\,d\widetilde\nu_\beta(C)
                                      =\mu_{\Lambda,\beta}(D).
 \tag{23.9}
\]

\begin{theorem}[Averaged conditional connected-insertion error]
\label{thm:f12v19-conditional}
Suppose \(f_j(Y,C),\widehat f_j(Y,C)\)
are measurable and uniformly bounded
by \(M_j\), and for each fibre
they agree off \(D\). The difference
of their conditional connected
derivatives satisfies
\[
 \int |\kappa_\alpha(C)-\widehat\kappa_\alpha(C)|
                         \,d\widetilde\nu_\beta(C)
 \le B_{\alpha,r}
        \left(\prod_jM_j^{\alpha_j}\right)
                               \mu_{\Lambda,\beta}(D),
 \quad B_{\alpha,r}=\alpha!(k/r)^mL_r .
 \tag{23.10}
\]
For every \(u>0\), the coarse probability
that the absolute difference exceeds
\(u\prod_jM_j^{\alpha_j}\) is at most
\[
                        B_{\alpha,r}\mu_{\Lambda,\beta}(D)/u.
 \tag{23.11}
\]
\end{theorem}
\begin{proof}
For each fixed \(C\), apply
Proposition~\ref{prop:f12v19-source}
to the probability law
\(\pi_{\beta,C}\), after dividing
by the \(M_j\). This gives the
pointwise bound with \(p(C)\)
on its right. Integrate and
use (23.9), then apply Markov's
inequality for (23.11).
No assertion that \(p(C)\)
is uniformly small in \(C\)
is required or deduced.
\end{proof}

\begin{corollary}[Averaged masking error for the exact global scores]
\label{cor:f12v19-scores}
Use the scores \(V_a,V_{ab}\)
from (22.18)--(22.19). For each
local score choose a measurable
mask \(0\le\chi_a(Y,C)\le1\),
equal to one off a local defect
event \(D_a\). For two fixed
coarse coordinates \(a,b\), define
\[
 {\cal M}_{ab}(C)=
   \mathbb E_C[\chi_a\chi_b V_{ab}]
             -\operatorname{Cov}_C(\chi_aV_a,\chi_bV_b).
 \tag{23.12}
\]
There is a volume-independent
constant \(K\) for the specified
local coordinate charts such that
\[
 \int |\partial_a\partial_b\widetilde W(C)
                        -{\cal M}_{ab}(C)|
                     \,d\widetilde\nu_\beta(C)
       \le K(\beta+1)^2
                      \mu_{\Lambda,\beta}(D_a\cup D_b).
 \tag{23.13}
\]
\end{corollary}
\begin{proof}
V18 gives uniform local bounds
\(|V_a|,|V_b|,|V_{ab}|\le K_0(\beta+1)\).
The expectation difference in
(23.12) is at most
\(K_0(\beta+1)p(C)\), where
\(p(C)=\pi_{\beta,C}(D_a\cup D_b)\).
The covariance difference is
bounded by a constant times
\(K_0^2(\beta+1)^2p(C)\), either
by the conditional \(m=2\)
case of (23.3) or directly
by expanding covariance into
two expectations and a product
of means. Integrate using
(23.9). This proves (23.13).
\end{proof}

\begin{record}[A concrete local choice of the defect masks]
For a fixed finite set of path
families around a score support,
take \(D_a\) to be the event
that some relative product
has distance at least a fixed
\(\sigma_a>0\) from \(I\).
It has a finite fine-plaquette
witness union, of cardinality
bounded independently of volume:
one can use sixty plaquettes
per path family, at threshold
\(\sigma_a/6\).
The one-plaquette case of V17
therefore gives
\[
       \mu_{\Lambda,\beta}(D_a)
                 \le N_a\,p_{\sigma_a/6}(\beta),
 \tag{23.14}
\]
where \(N_a\) is that witness
union size. The indicator of
\(D_a^c\) is an allowed mask;
a smooth local mask that equals
one outside a slightly larger
declared defect event also works.
For fixed supports and thresholds,
(23.13) is exponentially small
up to a polynomial in \(\beta\).
\end{record}

\begin{assessment}[What the masked expression is not]
The matrix \({\cal M}_{ab}\)
was defined by (23.12). It has
not been shown to be the
Hessian of any effective action.
If an action is modified by
a field-dependent cutoff, its
derivatives include derivatives
of that cutoff; (23.12) cannot
be substituted for those terms.

Nor does (23.13) control a
row sum over \(b\), an operator
norm in growing volume, or a
uniform bound on every fixed
coarse fibre. It is an averaged
estimate for two specified
local scores. The covariance
between their masked versions
still samples the full compact
conditional law, including
remote defects.
\end{assessment}

\subsection{Why the error cannot yet be promoted to clustering}

\begin{proposition}[Rare independent changes can retain a nondecaying covariance error]
For every \(0<p<1\), there are
bounded coupled fields \(A_x,B_x\)
whose simultaneous disagreement
events at distinct sites have
probability \(p^k\), but for
every pair of distinct sites
\[
            |\operatorname{Cov}(A_x,A_y)
                         -\operatorname{Cov}(B_x,B_y)|
                                      =4p(1-p),
 \tag{23.15}
\]
independently of their separation.
\end{proposition}
\begin{proof}
Let \(Z\) be a fair random sign
and \(\eta_x\) independent
Bernoulli variables of parameter
\(p\), independent of \(Z\).
Set \(A_x=Z\) and
\(B_x=Z(1-2\eta_x)\).
Disagreement is exactly
\(\{\eta_x=1\}\), so all
joint disagreement probabilities
have the stated product form.
Both fields have zero means;
for \(x\ne y\),
\(\mathbb E A_xA_y=1\) and
\(\mathbb E B_xB_y=(1-2p)^2\).
Subtract to obtain (23.15).
\end{proof}

\begin{assessment}[Logical boundary of the comparison]
This example is not asserted
to be a Yang--Mills measure.
It shows that bounded local-source
comparison and even product
multi-defect probabilities do
not, by themselves, force a
distance-decaying covariance
error. A shared long-range
background can remain.
Consequently a small error
floor in (23.7) or (23.13)
cannot be added to an exponential
clustering estimate and then
called exponential clustering.
\end{assessment}

\begin{decision}[V20 audit and the unresolved conditional covariance]
At V20 perform the required
companion audit and review
the direction of V10--V19.
The new smooth full-compact
integral removes the global
score-definition obstruction,
and the present estimates
control normalized local
changes and averaged rare
insertions. The remaining
task is a genuinely spatial
estimate for the compact
conditional covariance in
(22.19), with its boundary
and defect hypotheses checked.
Do not replace that task by
another distance-independent
small-error comparison.
\end{decision}

\begin{record}[V19 append record]
V18 had SHA--256
\begin{center}\scriptsize\ttfamily
A2118D11F747BF11BFCA7DCD8CD1736EF7EC8D42469D485FC165FFF75409B44E.
\end{center}
Its complete body is retained
apart from the title version.
Only active V18 is replaced
by V19; every family archive
remains. The new estimates
are analytic and concern
the stated finite-source
and averaged conditional
quantities.
\end{record}

\begin{firewall}[Position after V19]
Proved: normalized finite-source
comparison between the two
completions, fixed-order
connected-derivative stability,
and averaged local score
masking bounds. Not proved:
distance-decaying conditional
covariance, a volume-uniform
coarse Hessian comparison,
generated-action iteration,
continuum reconstruction,
or the Yang--Mills mass gap.
\end{firewall}

\section{V20: audit and a spatial compact-fibre estimate on a cofinal size range}
\label{sec:f12v20}

\subsection{Companion audit and ten-version direction review}

The current companion files and their endings were
checked again. They remain SM--Einstein V72
(607372 bytes), NS Stability V26 (278283 bytes),
and Parallel YM V5 (54174 bytes). Their complete
SHA--256 hashes match the V15 audit:
\[
\begin{array}{l|l}
 \text{SM V72}&\texttt{F44F9745D43379E1}\ \text{(prefix)}\\
 \text{NS V26}&\texttt{82C887F8C2C69E4F}\ \text{(prefix)}\\
 \text{Parallel YM V5}&\texttt{306F1BB84CFA8A8D}\ \text{(prefix)} .
\end{array}
\]
SM still supplies its fixed-mode many-copy limit,
not growing-cutoff quantum control. NS still
supplies pointwise rank-one norms, not global
regularity. Parallel YM still retains compact
holonomies and rejects uniform tightness of
extensive total charge. None supplies the
missing unrestricted compact conditional
covariance estimate.

V10--V12 established the uniform one-loop
comparison on a local fibre box. V13--V16
made the compact action, Gaussian powers,
curvature loads and exterior mismatch
costs explicit. V17 populated bare
probabilities for the actual obstruction
carriers. V18 provided an explicitly
different smooth completion with global
product fibre coordinates. V19 controlled
normalized finite sources and averaged
defect insertions, but not spatial decay.
These are genuine inputs; they are not
a constructed continuum theory.

We now combine the convex-box covariance
theorem with a compact-complement estimate
for the full smooth fibre. The result
has an explicit volume--coupling restriction.
This restriction is not renamed as an
unrestricted thermodynamic estimate.
Earlier cofinal comparisons for other
pinned fibres, notably f6 V31, are not
being re-proved as new principles: the
new application is to the actual global
smooth fibre of V18 with compact retained
holonomies and its complete Jacobian.

\subsection{A common local box inside the global product fibre}

Use \(\widetilde K\), not \(K^0\), and
its global variables \(Y\in G^{45n}\).
For a full retained field \(C\), express
each free internal link as \(Y_l=e^{u_l}\),
and each free boundary link over \(e\)
as \(Y_l=C_e e^{u_l}\). Let
\[
 {\cal D}=[-r,r]^N,\qquad N=45dn,
 \tag{24.1}
\]
be a small fixed scalar coordinate box.
Its image in the global \(Y\) space
is denoted \({\cal D}_C\).

Choose \(r>0\) so that the original
good root-pivot solution throughout
this tube has all relative products
in V18's agreement neighbourhood.
This is a uniform choice, by the
local pivot construction with compact
coarse parameters. That original
solution also solves the smooth block.
Global uniqueness of the smooth pivot
therefore identifies the two solutions
throughout this tube.

On the box the full smooth fibre
weight is exactly
\[
 e^{-\beta f(u,C)+\ell(u,C)}\,du,
 \tag{24.2}
\]
where \(\ell\) includes both the
global inverse-pivot Jacobian and
the free-link exponential Haar
Jacobians. Let \(\pi_{\beta,C}\)
be the normalized full compact
fibre law, and
\(\pi_{\beta,C}^{\cal D}\) its
conditioning on \({\cal D}_C\).
Thus the boxed law is a restriction
of this same compact law, not an
auxiliary Gaussian measure.

Let \({\cal L}(C)\) be V15's curvature
load for the full periodic region,
and set \(H_0(C)=4S_{\rm coarse}(C)\).
All pivot factors here are those of
the smooth completion; their bounds
are furnished globally by V18.

\begin{proposition}[Uniform local convexity for a small total curvature load]
\label{prop:f12v20-box}
There exist \(L_0>0,h>0,\beta_0<\infty\)
such that, for \({\cal L}(C)\le L_0\)
and \(\beta\ge\beta_0\), the measured
potential \(\beta f-\ell\) on (24.1)
has Hessian at least \(h\beta I\).
Its off-diagonal Hessian has fixed
range and row bound \(K\beta\).
All constants are independent of
volume and compact coarse holonomies.
\end{proposition}
\begin{proof}
Each corner-plaquette distance is
at most \(\sqrt{{\cal L}(C)}\):
any one of its incident load terms
already contains that nonnegative
distance. Take \(L_0\) small enough
for the uniform small-curvature
Hessian theorem of V15.
That theorem applies to these
free internal and right boundary
coordinates, with the solved
root pivots. Reduce \(r\) if
needed to obtain a common
positive action Hessian
throughout the box.

The density in (24.2) is a
product of bounded positive
local factors, with uniformly
bounded derivatives on the
fixed tube. Bounded stencil
overlap bounds its logarithmic
Hessian in row and operator
norm independently of volume.
For large \(\beta\), the
positive action Hessian
dominates this bounded
perturbation. The same local
formulas give the asserted
range and row bound.
\end{proof}

\subsection{The full conditional probability of leaving the box}

Choose \(a_0>0\) such that
\[
 U\notin\exp((-r,r)^d)
                        \quad\Longrightarrow\quad
                        \delta(U)^2\ge a_0 .
 \tag{24.3}
\]
Such an \(a_0\) exists because
this exponential box contains
a group neighbourhood of \(I\).
Its boundary has Haar measure
zero. Define the free-link distance
\[
 D_{\rm free}(Y,C)=
 \sum_{\rm internal}\delta(Y_l)^2+
 \sum_{\rm free\ boundary\ over\ e}
                          \delta(C_e^{-1}Y_l)^2 .
 \tag{24.4}
\]
Outside \({\cal D}_C\), it is at
least \(a_0\).

\begin{theorem}[Compact-complement bound at fixed coarse field]
\label{thm:f12v20-tail}
After possibly decreasing \(L_0\),
there are \(D_0\ge1,\kappa>0\)
such that for all allowed \(C\)
and \(\beta\ge\beta_0\),
\[
 p_{\rm out}(C):=\pi_{\beta,C}({\cal D}_C^c)
                    \le\min\{1,D_0^n e^{-\kappa\beta}\}.
 \tag{24.5}
\]
The statement is pointwise in \(C\)
for the canonical smooth fibre
integral of V18.
\end{theorem}
\begin{proof}
The deterministic section-flat
and corner comparison of V15
applies to the smooth completion
by (22.20). Hence
\[
 S_0(Y,B(Y,C))-H_0(C)
       \ge D_{\rm free}(Y,C)/(2A)-A{\cal L}(C)/2 .
 \tag{24.6}
\]
On the box complement, split
the distance term into two
equal parts to obtain
\[
 S_0-H_0\ge
       a_0/(4A)+D_{\rm free}/(4A)-A{\cal L}/2.
 \tag{24.7}
\]
Require in addition
\(L_0\le a_0/(4A^2)\).
Then (24.7) is at least
\(\kappa+D_{\rm free}/(4A)\),
with \(\kappa=a_0/(8A)\).

Use the global upper inverse-pivot
density bound from V18 and integrate
the free group variables. The
restricted unnormalized numerator
is at most
\[
 B^n\beta^{-N/2}e^{-\beta H_0(C)}e^{-\kappa\beta},
 \tag{24.8}
\]
by the compact Gaussian integral
bound of V13. Inside the box,
the symmetric cube argument
(19.17)--(19.18) gives a lower
contribution
\[
 b^n\beta^{-N/2}e^{-\beta H_0(C)}.
 \tag{24.9}
\]
That argument does not require
the section to be stationary.
Division proves (24.5), with
\(D_0=\max(1,B/b)\).
No bare probability estimate
has been substituted for a
fixed-\(C\) conditional estimate.
\end{proof}

\subsection{A complete covariance estimate with its error still visible}

Give the free coordinates their
fixed local cell locations and
use coarse-lattice graph distance.
For bounded observables \(F,G\)
on the full fibre, suppose their
restrictions to the box are
\(C^1\) and their \(u\)-gradients
are supported in coordinate
sets \(A_F,A_G\). Put
\[
 d(F,G)=\operatorname{dist}(A_F,A_G),\quad
 M_F=\|F\|_\infty,\quad
 L_F=\|\nabla_u F\|_{\infty,2},
 \tag{24.10}
\]
and similarly for \(G\). Empty
gradient support gives a constant
on the connected box; set
\(d(F,G)=0\) if either support
is empty. The compact error
comparison below still applies
with this convention.

\begin{theorem}[Convex contribution plus the exact compact error]
\label{thm:f12v20-two-terms}
There are \(K,\alpha>0\), independent
of volume and allowed \(C\), such that
\[
 |\operatorname{Cov}_{\pi_{\beta,C}}(F,G)|
 \le K\beta^{-1}e^{-\alpha d(F,G)}L_FL_G
       +6M_FM_G\min\{1,D_0^n e^{-\kappa\beta}\}.
 \tag{24.11}
\]
\end{theorem}
\begin{proof}
The boxed measure satisfies all
hypotheses of the f11 V85--V86
covariance theorem: a Cartesian
box, measured Hessian at least
\(h\beta I\), bounded finite-range
off-diagonal Hessian, and local
gradient supports. Its vector
resolvent with the mixed face
conditions and exponential
component weights gives
\[
 |\operatorname{Cov}_{\pi_{\beta,C}^{\cal D}}(F,G)|
       \le K\beta^{-1}e^{-\alpha d(F,G)}L_FL_G .
 \tag{24.12}
\]
Proposition~\ref{prop:f12v20-box}
verified these hypotheses for
the current action and density.

The total variation distance
between \(\pi_{\beta,C}\) and
its conditioning on the box
is exactly \(p_{\rm out}(C)\).
The difference of their
\(FG\) expectations is at most
\(2M_FM_Gp_{\rm out}\).
The difference of the products
of their two means is at most
\(4M_FM_Gp_{\rm out}\).
Apply (24.5) to their sum and
combine with (24.12).
\end{proof}

\subsection{When the error can legitimately be absorbed into distance decay}

\begin{theorem}[A cofinal full-compact spatial covariance bound]
\label{thm:f12v20-cofinal}
There exist \(c_0,\gamma,K'>0\)
such that if
\[
        {\cal L}(C)\le L_0,\qquad
        \beta\ge\beta_0,\qquad n\le c_0\beta ,
 \tag{24.13}
\]
then
\[
 |\operatorname{Cov}_{\pi_{\beta,C}}(F,G)|
 \le K'\beta^{-1}e^{-\gamma d(F,G)}
                             (L_FL_G+M_FM_G).
 \tag{24.14}
\]
The expectation is on the whole
compact smooth fibre, not the
local box.
\end{theorem}
\begin{proof}
A periodic coarse lattice of
\(n\) cells has diameter at
most \(4n\) in the metric
just used. Set \(\gamma\le\alpha\)
and choose \(c_0>0\) so
\[
       c_0\log D_0\le\kappa/4,\qquad
       4\gamma c_0\le\kappa/4.
 \tag{24.15}
\]
For example first fix any
\(0<\gamma\le\alpha\), then
take \(c_0\) sufficiently small.
When (24.13) holds,
\[
 D_0^n e^{-\kappa\beta}
       \le e^{-3\kappa\beta/4}
       \le e^{-\kappa\beta/2}e^{-\gamma d(F,G)}
       \le K_\kappa\beta^{-1}e^{-\gamma d(F,G)} .
 \tag{24.16}
\]
The last bound follows from the
finite supremum of
\(\beta e^{-\kappa\beta/2}\).
Substitute into (24.11).
This is the step that was
unavailable for the unrestricted
distance-independent errors in
V19; it explicitly uses the
finite-size relation in (24.13).
\end{proof}

\begin{corollary}[Spatial bound for the exact compact coarse Hessian]
For the complete smooth blocked
action \(\widetilde W\), in fixed
local right coarse-coordinate
charts and under (24.13),
\[
 |\partial_a\partial_b\widetilde W(C)|
                    \le K''\beta e^{-\gamma' d(a,b)}
 \tag{24.17}
\]
for some \(K'',\gamma'>0\)
independent of volume and \(C\)
in that regime. Its Hessian
row norm is at most \(K'''\beta\).
\end{corollary}
\begin{proof}
The exact identity (22.19) gives
\(\widetilde W_{ab}=\mathbb E V_{ab}
-\operatorname{Cov}(V_a,V_b)\).
The local score bounds from V18
give suprema of size \(K(\beta+1)\)
for \(V_a,V_b\) and their
gradients on the box, with
fixed-size supports. Apply
(24.14), obtaining a covariance
bound of size \(K\beta\) times
the spatial exponential for
\(\beta\ge\beta_0\ge1\).
The direct term \(V_{ab}\)
vanishes beyond a fixed coarse
range and is bounded by
\(K(\beta+1)\) otherwise.
Distances between score supports
and their coarse labels differ
by at most a fixed constant.
These facts prove (24.17).
Finally sum the exponential
using the uniform four-dimensional
polynomial bound on lattice balls.
\end{proof}

\begin{assessment}[What this does not prove]
The new estimate is spatial
and conditional on the full
compact fibre, but its domain
is exactly (24.13). In particular
it does not permit \(n\to\infty\)
at fixed \(\beta\). Nor does
the small bound on the
\emph{total} curvature load
assert that such coarse fields
are typical in an unrestricted
infinite-volume blocked measure.
Flat retained fields, including
nontrivial compact holonomies,
do satisfy the load condition.

The Hessian bound is an upper
bound, not a positive transverse
coarse mass term. It does not
remove gauge zero modes, prove
the infinite-volume one-loop
remainder for the compact law,
or produce a physical
Hamiltonian gap. A finite-volume
massive eliminated fibre is
still not a massive continuum
Yang--Mills theory.
\end{assessment}

\begin{decision}[Next spatial comparison must remove a global good-box event]
The cofinal restriction entered
only when the probability of
the whole fibre leaving one
large box was used in (24.11).
To reach fixed-\(\beta\)
thermodynamic control, replace
that event by a regional
comparison whose cost is tied
to the separation of the two
insertions, with arbitrary
remote defects retained.
V16 supplies an open regional
comparison and V17 supplies
bare carrier probabilities;
neither alone supplies the
required conditional exploration
or cancellation. Work on that
interface rather than dropping
\(D_0^n\) from (24.5).
Next companion audit: V25.
Next direction review: V30.
\end{decision}

\begin{record}[V20 append record]
V19 had SHA--256
\begin{center}\scriptsize\ttfamily
05B290DB7A965923A890E6A4AC34DAEFB1A35AC0C6D5E9F3D1DE670AEA2D88C8.
\end{center}
Its full body is retained
apart from the title version.
The required companion audit
and ten-version review were
performed. Only active V19
is replaced by V20; all
family archives remain.
\end{record}

\begin{firewall}[Position after V20]
Proved: a fixed-coarse-field
compact-complement estimate,
a two-term full-fibre covariance
bound, and a spatial compact
covariance/coarse-Hessian bound
on the explicit cofinal range.
Not proved: fixed-coupling
infinite-volume conditional
mixing, unrestricted coarse
curvature control, scale
iteration, continuum construction,
or the full Yang--Mills mass gap.
\end{firewall}

\section{V21: regional compact covariance and boundary response}
\label{sec:f12v21}

V20's restriction involved the number of cells in the
entire periodic lattice. We now replace it by the number
of cells in a specified connected update region.
The conditional kernel is identified inside V18's
global product fibre before any estimates are applied.
This permits arbitrary remote defects, but retains a
quantified local boundary condition. The result is
for the smooth block \(\widetilde K\); it is not silently
asserted for the piecewise block \(K^0\).

\subsection{Exact conditioning inside the global smooth fibre}

Let \(R\) be a nonempty connected set of \(s\) coarse
cells in a finite periodic lattice, with arbitrary
ambient size. Let \(E(R)\) contain every coarse edge
with at least one endpoint in \(R\), and put
\(m=|E(R)|\le8s\). Use the physical fine-link set
\(\Lambda_R\) of V14: the seventeen internal non-tree
links in each cell of \(R\), and all eight boundary
links over every edge of \(E(R)\).

In the global free variables \(Y\) of V18, let
\(I_R\) comprise the internal variables in \(R\) and
the seven nonpivot boundary variables on each edge
of \(E(R)\). Thus
\[
 D_R=|I_R|=17s+7m,\qquad N_R=dD_R,\qquad d=\dim G .
 \tag{25.1}
\]
Some boundary coordinates in this set are based at
an outside cell. They must not be omitted merely
because their base cell is outside \(R\).
Fix the complete coarse field \(C\), and fix the
remaining free variables at \(y=Y_{I_R^c}\).

\begin{proposition}[The canonical regional conditional kernel]
\label{prop:f12v21-kernel}
Let \(\xi(y,C)\) be the physical fine links outside
\(\Lambda_R\), reconstructed using the smooth inverse
pivots. These links do not depend on \(Y_{I_R}\).
The conditional law of \(Y_{I_R}\) under
\(\pi_{\beta,C}\), given \(y\), is the positive smooth
kernel
\[
 \pi_{R,C}^{y}(dY_R)
  =\frac{1}{Z_R(C,y)}
   e^{-\beta H_{\xi(y,C)}(Y_R,B_{E(R)}(Y_R,y,C))}
        \prod_{e\in E(R)}j_e(Y_R,y,C_e)\,dY_R .
 \tag{25.2}
\]
Here \(dY_R\) is product normalized Haar measure and
\(H_\xi=\sum_{p\in\mathcal P_R}s_p\), where
\(\mathcal P_R\) is the set of fine plaquettes touching
\(\Lambda_R\). This formula gives a canonical version
at every \(C,y\), not just almost everywhere.
\end{proposition}
\begin{proof}
For \(e\notin E(R)\), both endpoint cells are outside
\(R\). Its pivot depends only on its own nonpivot
boundary variables, internal variables in its endpoint
cells, and \(C_e\). All of these are fixed.
Consequently every physical exterior link is fixed.
Conversely, those physical exterior links include all
the free variables in \(I_R^c\); the additional exterior
pivots are deterministic functions of them at fixed
\(C\). The two conditioning sigma fields agree.

In the global density of V18, every inverse-pivot
factor with \(e\notin E(R)\) and every plaquette
outside \(\mathcal P_R\) is independent of \(Y_R\).
They cancel in the normalized conditional density,
leaving precisely (25.2). Product Haar measure
factors over \(I_R\) and \(I_R^c\); there is no
additional coarea factor. The integrand is positive
and smooth on compact products, so its integral
is positive and smooth. Disintegration initially
identifies the formula almost everywhere, and the
displayed formula itself supplies its everywhere
defined canonical version.
\end{proof}

Only a fixed finite collar of \(R\) enters (25.2).
Remote coarse and fine defects are not assumed absent.
Their indirect statistical influence through the
random boundary is not being discarded: the kernel
is conditional on that boundary.

\subsection{Transfer of the regional estimates to the smooth block}

Use the completed frozen collar \(\mathcal F_R\)
and the quantities of V16, now evaluated at
\(\xi=\xi(y,C)\):
\[
 b_l=\|\xi_l-J(C)_l\|_{\rm HS},\quad
 \eta=\max_{\mathcal F_R}b_l,\quad
 \mathcal E_R=\sum_{\mathcal F_R}b_l^2,\quad
 L_R=L_{\rm in}+L_{\rm out},\quad
 M_R=\sum_{\mathcal F_R}a_lb_l .
 \tag{25.3}
\]
The load \(a_l\) is the sum of
\(\delta(Q_p(C))\) over corner plaquettes
\(p\in\mathcal P_R\) incident on \(l\), exactly as
in (20.4). Write \(D\) for the squared
section-distance summed over all varied links,
including pivots, and \(D_{\rm free}\) for its
restriction to the free links. Let \(H_0(C)\)
be the section action on \(\mathcal P_R\).

\begin{proposition}[Regional comparison for this particular kernel]
\label{prop:f12v21-transfer}
The following inequalities hold for the smooth
fibre (25.2):
\[
 \begin{aligned}
 D&\le A(H_{\rm v}+E_{\rm v,out}),&
 E_{\rm v,out}&\le12\mathcal E_R,\\
 H_\xi-H_0&\ge H_{\rm v}
                 -\sum_{\Lambda_R}a_ld_l-M_R,&
 M_R&\le A L_{\rm out}+\mathcal E_R/(4A).
 \end{aligned}
 \tag{25.4}
\]
For sufficiently small fixed \(\eta_0>0\), if
\(\eta<\eta_0\), there is a feasible central repair
with free internal links \(I\), free boundary links
\(C_e\), and repaired pivots, whose action obeys
\[
 H_{\rm rep}-H_0
       \le A L_{\rm in}+M_R+K_1\mathcal E_R,\qquad
 Z_R(C,y)\ge b^s\beta^{-N_R/2}e^{-\beta H_{\rm rep}} .
 \tag{25.5}
\]
The second inequality holds for all sufficiently
large \(\beta\), uniformly in ambient volume.
One may use
\(K_1=12(1+K_p)+K_p/(4A)\), with \(K_p\) the
pivot-repair constant in (20.10).
\end{proposition}
\begin{proof}
The path proof of (20.7) uses the block output
only through its distance from the root path.
V18's bound (22.20) supplies exactly that same
bound for \(\widetilde K\), including away from
the common local chart. Its plaquette stencils
remain section-flat. This proves the first
inequality of (25.4). Frozen plaquette
telescoping and incidence at most six prove the
second. The corner Hilbert--Schmidt difference
identity is independent of the block map and
gives the third. The last is scalar Young's
inequality, summed over the frozen links.

For small \(\eta\), reduce the local repair
neighbourhood of V16 so that all relevant
relative products lie in V18's agreement
neighbourhood. The original local repair then
also solves the smooth block equation. The
global uniqueness of each smooth inverse pivot
identifies it with the pivot in (25.2). Thus the
squared repair estimate (20.10) applies.
The proof of (20.12), using this estimate and
the exact corner difference identity, gives
the first part of (25.5).

For completeness the lower integral uses the
actual smooth inverse-pivot density and the
free-link Haar factors in a fixed local chart.
They are bounded below by a constant to the
power \(s\). The action has bounded Hessian
operator norm by finite stencil overlap.
On a symmetric cube of scalar radius
\(t/\sqrt\beta\), its Taylor upper bound is
\(H_{\rm rep}+g\cdot u+K|u|^2/2\).
Since \(N_R\le73ds\), the quadratic exponential
cost is a constant to the power \(s\).
Symmetry gives
\(\int_{\rm cube}e^{-\beta g\cdot u}\,du
\ge\operatorname{vol}({\rm cube})\).
This proves the second part of (25.5) without
requiring the repair to be stationary.
The smooth positive density makes the conclusion
pointwise for the kernel (25.2).
\end{proof}

\subsection{A regional box-complement estimate}

Use free coordinates \(Y_l=e^{u_l}\) on internal
links and \(Y_l=C_e e^{u_l}\) on boundary links.
Choose a fixed \(r>0\) and shrink \(\eta_0\)
so the cube \([-r,r]^{N_R}\), with its solved
pivots, remains in the common local chart.
Let \(\mathcal D_{R,C}\) be its image in the
free group variables. Its definition does not
depend on \(y\). Choose \(a_0>0\) so that a free
group variable outside its exponential cube has
squared section-distance at least \(a_0\).
As in V20 this is possible because the cube
contains a group neighbourhood of the identity.
Put
\[
 K_3=12+K_1+\frac{1}{2A},\qquad
 \Gamma_R(C,y)=2A L_R(C)+K_3\mathcal E_R(y,C),
 \qquad \kappa=\frac{a_0}{8A}.
 \tag{25.6}
\]

\begin{theorem}[The ambient volume disappears from the error]
\label{thm:f12v21-tail}
There is \(D_1\ge1\), independent of \(R\) and
ambient volume, such that, if \(\eta<\eta_0\),
\[
 \pi_{R,C}^{y}(\mathcal D_{R,C}^c)
 \le\min\left\{1,D_1^s
       \exp\left[-\beta\left(\frac{a_0}{4A}
                                      -\Gamma_R(C,y)\right)\right]\right\}.
 \tag{25.7}
\]
In particular, if \(\Gamma_R\le a_0/(8A)\),
the right side is at most
\(\min\{1,D_1^s e^{-\kappa\beta}\}\).
\end{theorem}
\begin{proof}
Use (25.4) and
\(a_ld_l\le d_l^2/(2A)+Aa_l^2/2\) to get
\[
 H_\xi-H_0
   \ge\frac{D}{2A}-E_{\rm v,out}
                         -\frac A2 L_{\rm in}-M_R.
 \tag{25.8}
\]
On the box complement \(D\ge D_{\rm free}\ge a_0\).
Splitting its first term yields the lower bound
\[
 \frac{a_0}{4A}+\frac{D_{\rm free}}{4A}
       -E_{\rm v,out}-\frac A2 L_{\rm in}-M_R .
 \tag{25.9}
\]
All smooth inverse-pivot factors have a uniform
upper bound. Integration of
\(\exp[-\beta D_{\rm free}/(4A)]\) against
product Haar, by V13's compact Gaussian integral
bound, is at most \(B^s\beta^{-N_R/2}\), after
including those factors. This estimates the
restricted numerator in (25.2).
Divide by the actual lower denominator in
(25.5). The common Gaussian power and the
section background cancel. The remaining
exponent is at least
\[
 \frac{a_0}{4A}-E_{\rm v,out}
      -\frac{3A}{2}L_{\rm in}-2M_R-K_1\mathcal E_R
 \ge\frac{a_0}{4A}-2A L_R-K_3\mathcal E_R .
 \tag{25.10}
\]
The last step uses both frozen-data bounds
in (25.4) and \(L_{\rm in}\ge0\).
Take \(D_1=\max\{1,B/b\}\). No integration
over the rest of the lattice has entered
this numerator-to-denominator comparison.
\end{proof}

\subsection{Covariance on the whole compact regional fibre}

In addition to the total debit condition, impose
the local convexity conditions
\[
 \delta(Q_p(C))<\epsilon_0
       \quad(p\in\mathcal P_R\cap\mathcal C),\qquad
 \eta<\eta_0,\qquad \Gamma_R\le a_0/(8A).
 \tag{25.11}
\]
Here \(r,\epsilon_0,\eta_0\) are fixed sufficiently
small numbers. The first two conditions are
stated explicitly, rather than inferred with
unspecified constants from the third.

The local action Hessian on the cube is
uniformly positive by V16's exterior-dependent
Hessian result and the chart identification
in Proposition~\ref{prop:f12v21-transfer}.
The logarithm of the coordinate density includes
both inverse-pivot and free exponential-Haar
factors. Its Hessian has bounded row norm,
independent of \(s\). Therefore, for large
\(\beta\), the measured Hessian is at least
\(h\beta I\), with finite range and off-diagonal
row bound \(K\beta\). This verifies the
hypotheses of the same boxed covariance theorem
used in V20; exterior parameters are frozen
throughout this assertion.

For bounded real \(F,G\) on the full regional
fibre, with \(C^1\) restrictions to the cube,
let \(A_F,A_G\) be supports of their coordinate
gradients. Give coordinates the same local cell
locations as in V20 and put
\[
 d(F,G)=\operatorname{dist}(A_F,A_G),\qquad
 M_F=\|F\|_\infty,\qquad
 L_F=\sup_{\mathcal D_{R,C}}\|\nabla_uF\|_2 ,
 \tag{25.12}
\]
and similarly for \(G\). Set the distance to
zero if either gradient support is empty.

\begin{theorem}[Two-term and spatial regional covariance bounds]
\label{thm:f12v21-cov}
Under (25.11), uniformly in ambient volume,
\[
 |\operatorname{Cov}_{\pi_{R,C}^{y}}(F,G)|
 \le K\beta^{-1}e^{-\alpha d(F,G)}L_FL_G
       +6M_FM_G\min\{1,D_1^s e^{-\kappa\beta}\}.
 \tag{25.13}
\]
There are fixed \(c_1,\gamma,K'>0\) such that,
if also \(s\le c_1\beta\), then
\[
 |\operatorname{Cov}_{\pi_{R,C}^{y}}(F,G)|
 \le K'\beta^{-1}e^{-\gamma d(F,G)}
                              (L_FL_G+M_FM_G).
 \tag{25.14}
\]
\end{theorem}
\begin{proof}
For the measure conditioned on the regional
cube, the verified measured Hessian bounds
give the f11 V85--V86 vector-resolvent estimate
\(K\beta^{-1}e^{-\alpha d}L_FL_G\).
The scalar reflecting generator and vector
operator have the mixed face conditions
specified there; restriction to a subset of
coordinates does not introduce new off-diagonal
interactions. The finite-range weighted
resolvent proof uses only the uniform Hessian
gap and stencil bounds, so its constants do
not depend on this subset or its exterior.

The total variation distance to conditioning
on the cube is exactly its complement
probability. The \(FG\) expectation changes
by at most \(2M_FM_Gp_{\rm out}\), and the
product of the means by at most
\(4M_FM_Gp_{\rm out}\). Apply (25.7) to obtain
(25.13).

Since \(R\) is connected, all locations in
\(I_R\) lie in a fixed-width collar of a
connected set of \(s\) cells. Any two have
graph distance at most \(K_ds\), for a fixed
\(K_d\). Choose \(0<\gamma\le\alpha\) and
\(c_1>0\) with
\[
 c_1\log D_1\le\kappa/4,\qquad
 \gamma K_dc_1\le\kappa/4 .
 \tag{25.15}
\]
For \(s\le c_1\beta\) the error in (25.13)
is at most
\(e^{-\kappa\beta/2}e^{-\gamma d(F,G)}\).
The finite supremum of
\(\beta e^{-\kappa\beta/2}\) over
\(\beta\ge\beta_0\) absorbs it into (25.14).
It is \(s\), not the number of ambient cells,
which occurs in this absorption.
\end{proof}

\subsection{A normalized boundary-response identity}

Keep \(C\) fixed. Vary one exterior free group
coordinate \(y_j\) along a unit right Lie
direction. This varies the physical exterior
links, including any dependent exterior pivots,
consistently with the fixed unaffected coarse
outputs. It is not a variation of independent
physical boundary links followed by resetting
their coarse outputs.
Use the Haar-density potential
\[
 V_R(Y_R;y,C)=\beta H_{\xi(y,C)}(Y_R,B_{E(R)})
                     -\sum_{e\in E(R)}\log j_e,\qquad
 S_j=\partial_{y_j}V_R .
 \tag{25.16}
\]
An additive function of \(y,C\) in \(V_R\)
would only add an interior constant to \(S_j\),
which does not affect the covariance below.

\begin{proposition}[Exact response, including normalization]
\label{prop:f12v21-response}
If a bounded \(F(Y_R)\) has no explicit exterior
dependence, then
\[
 \partial_{y_j}\mathbb E_{\pi_{R,C}^{y}}F
            =-\operatorname{Cov}_{\pi_{R,C}^{y}}(F,S_j).
 \tag{25.17}
\]
Each \(S_j\) is a sum of a uniformly bounded
number of local terms, and
\[
 \|S_j\|_\infty+
     \sup_{\mathcal D_{R,C}}\|\nabla_u S_j\|_2
                                    \le K(\beta+1).
 \tag{25.18}
\]
Its gradient support is within fixed distance
of \(j\). Under the strict interior of (25.11),
and \(s\le c_1\beta\), this implies
\[
 \left|\partial_{y_j}\mathbb E_{\pi_{R,C}^{y}}F\right|
 \le K e^{-\gamma\operatorname{dist}(A_F,j)}
                                              (L_F+M_F).
 \tag{25.19}
\]
Here the distance is set to zero if \(A_F\)
is empty. A score constant on the cube is
handled separately in the proof.
\end{proposition}
\begin{proof}
Differentiate numerator and denominator of
(25.2). Compactness and V18's smooth
uniformly bounded fixed-order pivot derivatives
justify differentiation under the integral.
The derivative of \(e^{-V_R}\) is
\(-S_je^{-V_R}\); the quotient rule is exactly
(25.17). The denominator derivative cannot
be dropped.

Each free coordinate enters only a bounded
number of pivots and plaquettes. Their
derivatives of each fixed order are uniformly
bounded on the full compact products; all
positive \(j_e\) have a uniform lower bound.
This proves the supremum estimate and,
after passing to the fixed box coordinates,
the gradient estimate in (25.18).
The same dependence statement proves locality.

If \(A_F\) and the score gradient support
are nonempty, apply (25.14) with \(G=S_j\).
Their distance is at least
\(\operatorname{dist}(A_F,j)-K_0\); absorb
this fixed shift and use
\((\beta+1)/\beta\le2\).
If \(A_F\) is empty, (25.14) already gives
the stated convention.

If the score gradient support is empty,
the score is constant on the connected cube,
so the boxed covariance is zero. The
compact comparison instead gives the bound
\(6M_F\|S_j\|_\infty D_1^s e^{-\kappa\beta}\).
When \(j\) enters a term of \(V_R\), it lies
within fixed distance of \(I_R\), so
\(\operatorname{dist}(A_F,j)\le K_ds\) after
enlarging \(K_d\) and decreasing \(c_1\) in
(25.15), once and for all. The same exponential
absorption, now using the boundedness of
\((\beta+1)e^{-\kappa\beta/2}\), proves
(25.19). If \(j\) enters no term, \(S_j=0\)
and the derivative vanishes exactly.
\end{proof}

For a fixed interior observable, (25.19) can be
integrated along any piecewise \(C^1\) path
\(y(t)\) which remains in the protected open
set at fixed \(C\):
\[
 \left|\mathbb E_R^{y(1)}F-\mathbb E_R^{y(0)}F\right|
 \le K(L_F+M_F)\int_0^1
        \sum_{j\in I_R^c}
        e^{-\gamma\operatorname{dist}(A_F,j)}
                     \|\dot y_j(t)\|\,dt .
 \tag{25.20}
\]
Only the finite interacting collar need appear
in the sum. The velocity is measured in the
right-trivialized group metric. For an
exterior-dependent physical observable an
additional \(\mathbb E_R(\partial_{y_j}F)\)
term is necessary in (25.17); (25.20) is not
being applied to that case without this term.
No claim is made that two arbitrary protected
endpoints admit a path which stays protected.

\begin{assessment}[What has changed and what has not]
The full compact regional kernel now has an
explicit spatial covariance bound and an
exponentially localized boundary derivative.
The ambient lattice may be arbitrarily large,
and data outside the finite relevant collar
are unrestricted. This is strictly different
from V20's all-lattice size restriction.

However the updated region still satisfies
\(s\le c_1\beta\), and its total curvature and
boundary-mismatch debit is bounded by a fixed
small constant. This is not a fixed positive
density allowance as \(s\) grows. The result
does not assert that the protected set is
typical under every fixed-coarse-field law.
V17's bare carrier probabilities cannot be
reused unchanged after conditioning on an
exploration history.
\end{assessment}

\begin{decision}[The next missing step is a populated boundary iteration]
The new input is an exact local conditional
kernel with a controlled boundary derivative,
not an assumption of mixing. A scale or
region iteration must now price visits to
unprotected boundary data under its actual
conditional law, and retain the normalization
when comparing successive kernels. A bound
only averaged under the original Wilson law
does not by itself close that iteration.
Investigate this upstream conditional
probability issue before claiming
fixed-\(\beta\) all-distance contraction.
Next companion audit: V25. Next direction
review: V30.
\end{decision}

\begin{record}[V21 append and scope record]
V20 had SHA--256
\begin{center}\scriptsize\ttfamily
AF314359290C6205A911790A873C6DA290C1512A17F9CC67E1AAD0F7D26D54D8.
\end{center}
Its body is retained apart from the title
version. The current inventory still has
SM V72, NS V26 and parallel-YM V5, unchanged
in size and timestamp from the V20 audit.
Only active V20 is replaced; all family
archives and other programmes remain.
The target remains the continuum existence
and mass-gap problem described by the
\href{https://www.claymath.org/millennium/yang-mills-the-maths-gap/}
{Clay Mathematics Institute}, not just a
gap in an eliminated finite-regulator fibre.
\end{record}

\begin{firewall}[Position after V21]
Proved here from the recorded regional geometry:
exact identification of the smooth regional
conditional density, a matched regional
box-complement estimate, spatial compact
covariance on the stated size and data regime,
and a normalized local boundary-response bound.
Not proved: protection under arbitrary
conditioning, fixed-coupling infinite-volume
mixing, renormalization iteration, continuum
construction, or the physical Yang--Mills
mass gap.
\end{firewall}

\section{V22: conditional protection of a buffered inner boundary}
\label{sec:f12v22}

V21 supplied an exact regional kernel, but it did not
estimate the probability that a subsequent boundary
is protected. Here we prove one such estimate for a
buffered inner region under a protected parent
kernel. We first locate the measured minimum with
its boundary dependence retained, then prove local
concentration, and finally specify which exploration
histories preserve the estimate. This is a populated
conditional transition, not a claim that an initial
protected parent exists for arbitrary exterior data.

The preceding goal turn made concrete progress:
V21 was appended, its preserved predecessor hash was
checked, and the source compiled. Its recorded
limitations, rather than compilation alone, determine
the present next step.

\subsection{A local forcing estimate at the repaired centre}

Keep V21's smooth block, regional coordinates and
canonical kernel. Write \(R\) for the parent,
\(s=|R|\), and \(y\) for its exterior free data.
On the fixed cube let
\[
 \mathcal V_R(u)=\beta f_R(u)-\ell_R(u),
 \qquad
 \epsilon_R=\max_{p\in\mathcal P_R\cap\mathcal C}
                                    \delta(Q_p(C)).
 \tag{26.1}
\]
An empty maximum is zero. The logarithmic density
\(\ell_R\) includes the free exponential-Haar
factors as well as the inverse-pivot factors.
Assume the parent conditions (25.11).
Reduce the fixed upper bounds for \(\epsilon_R\)
and \(\eta_R\), and increase \(\beta_0\), when
necessary below. Constants remain independent
of parent and ambient sizes and compact holonomies.

Distances between coordinates, links and plaquettes
are minimum distances between their coarse cells;
changing a location within a fixed stencil only
changes the constants in an exponential bound.

\begin{proposition}[The measured minimum feels a weighted boundary load]
\label{prop:f12v22-mode}
The measured potential has a unique stationary
minimum \(w_R\) in the cube, with
\(\|w_R\|_\infty<r/4\). For some fixed \(K,a>0\),
every scalar coordinate \(i\in I_R\) satisfies
\[
 |w_{R,i}|\le K\left[
  \frac1\beta+
  \sum_{p\in\mathcal P_R\cap\mathcal C}
                 e^{-a\,{\rm dist}(i,p)}\delta(Q_p(C))
  +\sum_{l\in\mathcal F_R}
                 e^{-a\,{\rm dist}(i,l)}b_l\right].
 \tag{26.2}
\]
Consequently, if \({\rm dist}(i,\mathcal F_R)\ge L\),
then, after possibly decreasing \(a\),
\[
 |w_{R,i}|
       \le K\bigl(\epsilon_R+\eta_R e^{-aL}+\beta^{-1}\bigr).
 \tag{26.3}
\]
\end{proposition}
\begin{proof}
At \(u=0\), the free variables are the section
values and the pivots are the feasible repair.
At zero exterior mismatch that repair is exactly
the section. Section-flat plaquettes then have
zero first derivative of their action. For a
corner plaquette, differentiating
\(\|I-W_p\|_{\rm HS}^2/2\) bounds its first
derivative by \(K\delta(Q_p(C))\). A coordinate
affects only a bounded number of such plaquettes.

Changing the relevant exterior from the section
to its actual value changes this derivative by
at most \(K\) times the sum of the nearby \(b_l\).
Indeed the repaired pivot and its fixed-order
derivatives are uniformly Lipschitz in their
finite local inputs. Interpolate each small
exterior deviation in its exponential chart;
bounded mixed derivatives give this estimate.
This is an auxiliary comparison of the local
action at fixed \(C\), not an assertion that
the interpolated exterior has the same
unaffected block outputs.
The gradient of \(\ell_R\) is uniformly bounded
coordinatewise, by its local factorization.
Thus for a fixed stencil radius \(R_0\),
\[
 |\beta^{-1}\partial_i\mathcal V_R(0)|
 \le K\left[\beta^{-1}
    +\sum_{{\rm dist}(i,p)\le R_0}\delta(Q_p(C))
    +\sum_{\substack{l\in\mathcal F_R\\
                          {\rm dist}(i,l)\le R_0}}b_l\right],
 \tag{26.4}
\]
where the first sum is over affected corners.

V16 gives an interior action minimum for small
\(\epsilon_R+\eta_R\). Its Hessian has a
uniform max-norm inverse by the finite-range
Neumann argument. The perturbation
\(-\beta^{-1}\ell_R\) has gradient and Hessian
row norms \(O(\beta^{-1})\). The same local
implicit contraction therefore produces the
measured minimum, displaced by \(O(\beta^{-1})\)
in max norm. Shrink the data bounds so it
lies in \((-r/4,r/4)^{N_R}\).
The strict convexity of \(\mathcal V_R\)
throughout the cube proves uniqueness there.

To retain spatial information, put
\[
 A_R=\int_0^1\beta^{-1}
                      \mathcal V_R''(t w_R)\,dt .
 \tag{26.5}
\]
This is symmetric, has fixed range, and
\(hI\le A_R\le MI\) with fixed \(0<h<M\).
Choose \(M\) larger if needed. The series
\[
 A_R^{-1}=M^{-1}
            \sum_{k=0}^\infty(I-A_R/M)^k
 \tag{26.6}
\]
converges in operator norm. Its \(ij\) entry
vanishes at order \(k\) if
\({\rm dist}(i,j)>kR_0'\), where \(R_0'\)
is the Hessian range, and its norm is at
most \((1-h/M)^k\). Summing gives a uniform
exponential entry bound for \(A_R^{-1}\).
The equation
\(w_R=-A_R^{-1}\beta^{-1}\nabla\mathcal V_R(0)\),
combined with (26.4) and bounded stencil
overlap, proves (26.2).
Exponential kernels are summable on the
four-dimensional cell graph, with bounded
coordinate and link multiplicities. Their
tail outside distance \(L\) is at most
\(K e^{-aL/2}\) after absorbing polynomial
ball growth. This proves (26.3) with a
renamed exponent, without a factor of \(s\)
or of the number of boundary links.
\end{proof}

\subsection{Boxed local concentration with no dimension loss}

Let \(\nu_R\) be \(\pi_{R,C}^y\) conditioned on
\(\mathcal D_{R,C}\). We record the needed
local argument explicitly: a total
\(N_R/\beta\) moment would not suffice here.

\begin{proposition}[A coordinate mean and tail estimate]
\label{prop:f12v22-concentration}
For all coordinates \(i\),
\[
 \mathbb E_{\nu_R}|u_i-w_{R,i}|^2\le K/\beta,
 \qquad
 \nu_R\{|u_i-\mathbb E_{\nu_R}u_i|\ge t\}
                         \le2e^{-h\beta t^2/2}.
 \tag{26.7}
\]
In particular for coordinates at distance at
least \(L\) from the frozen parent collar,
\[
 |\mathbb E_{\nu_R}u_i|
   \le m_R(L):=
       K\bigl(\epsilon_R+\eta_R e^{-aL}+\beta^{-1/2}\bigr).
 \tag{26.8}
\]
The constants and the possibly reduced \(h>0\)
do not depend on the parent size.
\end{proposition}
\begin{proof}
Fix \(i_0\), write \(v=u-w_R\), and use diagonal
weights \(D_{ii}=e^{-\alpha\,{\rm dist}(i,i_0)}\).
The Hessian is at least \(h\beta I\), with
bounded finite-range off-diagonal row norm.
For sufficiently small fixed \(\alpha\),
conjugation by \(D^{1/2}\) changes that Hessian
in operator norm by at most \(h\beta/2\).
This follows from the row and column bounds
and the fact that neighbouring weights have
ratio at most \(e^{\alpha R_0'}\).
Integrating the Hessian along the segment
from \(w_R\) to \(u\) gives
\[
 v^TD\nabla\mathcal V_R(u)
                    \ge(h\beta/2)\sum_iD_{ii}v_i^2.
 \tag{26.9}
\]
The weighted divergence identity on the cube is
\[
 \sum_iD_{ii}-\mathbb E_{\nu_R}
                         v^TD\nabla\mathcal V_R
 =Z_{\mathcal D}^{-1}\int_{\partial\mathcal D}
                            (Dv)\cdot n\,e^{-\mathcal V_R}\,dS
 \ge0 .
 \tag{26.10}
\]
Every face term is nonnegative because \(w_R\)
is inside the cube. Polynomial ball growth
bounds \(\sum_iD_{ii}\) uniformly, and
\(D_{i_0i_0}=1\), proving the moment estimate.
This is the f11 V89 local argument with the
regional measured Hessian hypotheses now checked.

For the tail, let
\(\phi(t)=\log\mathbb E_{\nu_R}e^{t u_i}\).
The tilted potential \(\mathcal V_R-tu_i\)
has the same Hessian and the same convex
box. The boxed Poincare inequality of
f11 V85 therefore gives
\(\phi''(t)=\operatorname{Var}_{\nu_{R,t}}(u_i)
\le(h\beta)^{-1}\) for every real \(t\).
Integrating twice yields
\(\phi(t)-t\phi'(0)\le t^2/(2h\beta)\).
The exponential Markov inequality, optimized
over positive and negative \(t\), gives
the second estimate of (26.7).
Finally Cauchy--Schwarz for the local second
moment and (26.3) give (26.8).
\end{proof}

\begin{record}[Classical input versus the present application]
The concentration step uses the classical
variance principle for strongly log-concave
measures; see H.~J.~Brascamp and E.~H.~Lieb,
\emph{On extensions of the Brunn--Minkowski
and Pr\'ekopa--Leindler theorems, including
inequalities for log concave functions, and
with an application to the diffusion equation},
J. Funct. Anal. \textbf{22} (1976), 366--389,
\href{https://doi.org/10.1016/0022-1236(76)90004-5}
{original article}.
The local boundary forcing, the regional
minimum estimate, and the conditional
protection statement below are the present
applications, not results attributed to
that article. The box variance bound was
already proved in the archive and is used
above with all of its hypotheses checked.
\end{record}

\subsection{The actual physical collar of an inner region}

Let \(S\) be a nonempty connected inner region,
\(t=|S|\), with its completed physical collar
\(\mathcal F_S\) as in V16. Require every link
in \(\mathcal F_S\) to be varied by the parent
update and every free coordinate needed to
reconstruct it to belong to \(I_R\).
This is guaranteed by a sufficiently wide
fixed buffer between \(S\) and the exterior
of \(R\).

Let \(J_S\subset I_R\) be the union of these
scalar free-coordinate supports. Impose
\[
 {\rm dist}(J_S,\mathcal F_R)\ge L,\qquad
 q=\max\{1,|\mathcal F_S|\},\qquad |J_S|\le Kq .
 \tag{26.11}
\]
The last estimate follows from bounded
pivot stencils; it is not a bound by all
coordinates of \(R\). For an empty parent
collar use distance \(+\infty\).

\begin{proposition}[Physical deviations from local free coordinates]
\label{prop:f12v22-physical}
There is a fixed \(K_{\rm link}\ge1\) such that
on the parent cube, for every \(l\in\mathcal F_S\),
\[
 \|W_l(u)-J(C)_l\|_{\rm HS}
       \le K_{\rm link}\max_{i\in J_l}|u_i|,
 \qquad J_l\subset J_S,\quad |J_l|\le K .
 \tag{26.12}
\]
\end{proposition}
\begin{proof}
A free physical link is either \(e^{u_l}\)
or \(C_e e^{u_l}\), so the assertion follows
from the local exponential derivative bound,
allowing for the fixed Lie dimension.
A physical pivot is a smooth function of
a fixed number of free variables and \(C_e\).
By hypothesis all of those variables are
inside \(I_R\). When they equal their section
values, the unique smooth pivot equals \(C_e\).
Integrate its uniformly bounded derivatives
along the line from zero to those free
coordinates. This proves (26.12).
No nonzero parent-exterior repair term is
missing here: the full input-support condition
excludes such an exterior input for these links.
\end{proof}

Use the fixed constants \(\kappa=a_0/(8A)\)
and \(K_3\) of V21, and set
\[
 \tau_q=\min\left\{\frac{\eta_0}{2},
                    \sqrt{\frac{\kappa}{2K_3q}}\right\}.
 \tag{26.13}
\]
Assume the deterministic retained field
satisfies, for the child,
\[
 2A L_S(C)<\kappa/2,\qquad
 \delta(Q_p(C))<\epsilon_0
          \quad(p\in\mathcal P_S\cap\mathcal C).
 \tag{26.14}
\]
If every child collar mismatch is less
than \(\tau_q\), then
\(\eta_S<\eta_0\) and
\[
 \Gamma_S=2A L_S+K_3\mathcal E_S<\kappa.
 \tag{26.15}
\]
Thus the child satisfies the strict V21
data conditions. These are the actual
coarse values and reconstructed physical
links, not a separately sampled collar law.

\begin{theorem}[One conditional protection transition]
\label{thm:f12v22-protection}
Assume the protected parent conditions,
the support and buffer condition (26.11),
the child curvature conditions (26.14),
and
\[
 K_{\rm link}\,m_R(L)\le\tau_q/2 .
 \tag{26.16}
\]
Then under the full compact parent kernel,
\[
 \pi_{R,C}^y\{\hbox{the child fails the strict
                         V21 data conditions}\}
 \le
 \min\left\{1,\,
       Kq e^{-c\beta\tau_q^2}
                       +D_1^s e^{-\kappa\beta}\right\}
 =:\delta_{R,S}.
 \tag{26.17}
\]
All constants are independent of ambient volume.
If also \(t\le c_1\beta\), the V21 spatial
conditional covariance and boundary-response
theorems are consequently available at the
sampled child boundary except on the event
bounded in (26.17).
\end{theorem}
\begin{proof}
On the parent cube, a physical mismatch
at least \(\tau_q\) forces some
\(i\in J_S\) to have
\(|u_i|\ge\tau_q/K_{\rm link}\), by (26.12).
Its mean has magnitude at most
\(\tau_q/(2K_{\rm link})\), by (26.16).
The centered coordinate tail (26.7) and
a union bound over \(|J_S|\le Kq\) give
\[
 \nu_R\left\{\max_{l\in\mathcal F_S}
             \|W_l-J(C)_l\|_{\rm HS}\ge\tau_q\right\}
       \le Kq\exp[-h\beta\tau_q^2/(8K_{\rm link}^2)].
 \tag{26.18}
\]
For an arbitrary event \(B\),
\(\pi_{R,C}^y(B)\le\nu_R(B)
              +\pi_{R,C}^y(\mathcal D_{R,C}^c)\).
The parent complement is bounded by (25.7)
under its protected data conditions. The
deterministic child curvature conditions and
(26.15) make every other failure a subset
of the mismatch event just estimated.
This proves (26.17).
Finally, at fixed \(C\), V21 identifies the
conditional law in \(S\) given its full
exterior as its canonical regional kernel.
The claim about availability of V21 is a
pointwise assertion about that kernel on
the successful data, not an assertion that
conditioning only on success leaves its
interior distribution unchanged.
\end{proof}

\begin{corollary}[An explicit nonempty quantitative regime]
There is a fixed \(K_*>0\) such that (26.16)
is ensured by
\[
 \epsilon_R\le\frac{\tau_q}{K_*},\qquad
 \eta_R e^{-aL}\le\frac{\tau_q}{K_*},\qquad
 \beta\ge K_*^2/\tau_q^2 .
 \tag{26.19}
\]
Moreover \(\tau_q^2\ge c_*/q\) for a fixed
\(c_*>0\), so
\[
 \delta_{R,S}\le
 \min\{1,Kq e^{-c'\beta/q}
                         +e^{s\log D_1-\kappa\beta}\}.
 \tag{26.20}
\]
For \(s\log D_1\le\kappa\beta/2\) the second
term is at most \(e^{-\kappa\beta/2}\).
For example the error tends to zero in
any allowed family with
\(\beta/q-\widetilde K\log(q+1)\to+\infty\)
for a sufficiently large fixed
\(\widetilde K\).
\end{corollary}
\begin{proof}
Insert the three bounds in (26.8) and
choose \(K_*\) large enough. Since \(q\ge1\),
the minimum in (26.13) is bounded below
in square by
\(\min\{\eta_0^2/4,\kappa/(2K_3)\}/q\).
This gives (26.20) and the limiting assertion.
The assumptions are not an empty formal
interface: take any flat compact coarse
field and its exact section parent exterior.
Then both load and mismatch are zero.
For a finite nested pair, choose \(\beta\)
large enough. Small enough perturbations
give open nonzero-curvature and fluctuating
exterior examples satisfying all strict
inequalities. A general flat field may
retain nontrivial compact holonomies.
\end{proof}

The parent load restriction remains in force
in this corollary. In particular (26.19)
alone, which is a pointwise curvature bound,
does not replace \(2A L_R+K_3\mathcal E_R\le\kappa\).
The noise term requires \(\beta\) large relative
to the number of child collar links. It has
not been made uniform for arbitrary growing
children at fixed coupling.

\subsection{Which previous observations are admissible?}

\begin{proposition}[Exterior-measurable histories retain the estimate]
\label{prop:f12v22-history}
Work under the full finite-volume smooth
fibre law \(\pi_{\beta,C}\), with \(C\) fixed.
For a deterministic nested pair \(R,S\),
let \(H\) be an event measurable with respect
to \(Y_{I_R^c}\) such that its exterior data
obey the hypotheses of
Theorem~\ref{thm:f12v22-protection}. If the
right side of (26.17) is bounded by a fixed
\(\delta\) throughout \(H\), then
\[
 \pi_{\beta,C}(H\cap\{\text{child failure}\})
                         \le\delta\,\pi_{\beta,C}(H).
 \tag{26.21}
\]
For \(\pi_{\beta,C}(H)>0\), division gives
the same conditional bound given \(H\),
without an inverse history-probability loss.
The statement also holds for a finite adaptive
choice of parent, provided each disjoint
selection branch \(H_R\) is measurable in
that parent's exterior and satisfies its
uniform hypotheses.
\end{proposition}
\begin{proof}
Condition on \(Y_{I_R^c}\). V21 gives exactly
the kernel (25.2); apply (26.17) pointwise,
multiply by \(1_H\), and integrate. For
adaptive choice, perform this calculation
on each exterior-measurable disjoint branch
and sum. This uses no independence.
It also shows why revealing any parent
interior information may invalidate the
argument: its indicator could no longer
be pulled outside the conditional integral.
\end{proof}

One useful bookkeeping consequence can be stated
without presuming that such an exploration has
already been constructed. Suppose a filtration
records successive observations. At step \(k\),
the chosen parent and its exterior are known;
no additional information about its unrevealed
interior has been recorded. Suppose the resulting
conditional interior law is the canonical
kernel above, and all its hypotheses hold.
If \(B_k\) is the observed child-failure event
and (26.17) is bounded by deterministic
\(\delta_k\), then
\[
 \mathbb P(B_k\mid\mathcal H_{k-1})\le\delta_k,\qquad
 \mathbb P\left(\bigcap_{k\in J}B_k\right)
                                   \le\prod_{k\in J}\delta_k
 \tag{26.22}
\]
for every specified finite set of steps \(J\).
The second inequality follows by conditioning
at its largest index and repeating.
If hypotheses fail at some steps, (26.22)
applies instead to failures restricted to
the event that the step's hypotheses were
known to hold; it does not price the
failure of those hypotheses themselves.

\begin{assessment}[Why the conditioning caveat cannot be removed]
If \(0<\pi_{\beta,C}(B)<1\), the history
which reveals the event \(B\) itself gives
\(\pi_{\beta,C}(B\mid B)=1\). Thus even a
very small unconditional failure probability
cannot imply the same bound for arbitrary
histories. Exterior measurability in (26.21)
is a structural restriction on what has been
observed, not a stylistic qualification.
In particular it is not legitimate to
condition on a previously discovered
interior defect and then reuse (26.17).
\end{assessment}

\begin{decision}[Next task: an iteration that really has unrevealed parents]
We now have a conditional probability of
protection for an actual inner physical
collar, with explicit noise, curvature,
buffer and compact-complement costs.
The next step is to determine whether
these kernels can be arranged in an
exploration or resampling scheme which
retains the required exterior measurability
and controls its initial and failed guards.
A nested calculation with a protected
ancestor alone does not cover an arbitrary
infinite lattice. If a proposed iteration
reuses already observed interiors, it must
be changed or its likelihood ratio must
be estimated, rather than invoking (26.22).
Next companion audit: V25. Next direction
review: V30.
\end{decision}

\begin{record}[V22 append record]
The authoritative V21 source had SHA--256
\begin{center}\scriptsize\ttfamily
CE598846E04AABF1AC207E5F1964E2057EB3C2583FAE12628C2285E98914CB26.
\end{center}
Its full body is retained apart from the title
version. Only active V21 is replaced by V22;
all archived families and other programmes
remain. The local-moment theorem was checked
against f11 V89 before its regional reuse.
The new content is the localized measured
mode and the resulting actual nested-kernel
protection estimate, not a new claim of the
classical strong-convexity concentration
principle.
\end{record}

\begin{firewall}[Position after V22]
Proved under explicit parent hypotheses:
exponential attenuation of the measured
minimum's boundary forcing, coordinate
concentration, high-probability protection
of a buffered inner physical collar, and
stability of this estimate under
exterior-measurable histories.
Not proved: an unconditionally protected
initial parent, control after arbitrary
interior observations, a global contracting
iteration, fixed-coupling thermodynamic
mixing, continuum existence, or the full
Yang--Mills mass gap.
\end{firewall}

\section{V23: an actual outside-in reveal and its surviving mean term}
\label{sec:f12v23}

V22's admissible-history condition is now implemented,
not assumed as an abstract exploration interface.
We reveal successive annuli in the global smooth
free coordinates. This leaves the next inner
region genuinely unrevealed. Its annulus law
contains a child partition function, which is
retained exactly. The resulting stopped chain
has a controlled probability of losing protection.
It does not by itself imply mixing: the final
conditional means have a separate covariance.

The older protected transport chains in f2 V46
concern a cut-forest operator identity, not
these compact regional conditional kernels.
The positive-mixture covariance identity is
already familiar in this programme, for example
f5 V45 and the warning in current V19.
We use that identity below to expose the
remaining term, not present it as a new general
probability theorem.

\subsection{Nested coordinates and the exact annulus density}

Fix \(C\) and connected regions
\[
 R_0\supset R_1\supset\cdots\supset R_N,\qquad
 I_j=I_{R_j},\quad E_j=E(R_j),\quad
 \mathcal P_j=\mathcal P_{R_j}.
 \tag{27.1}
\]
Here \(I_j\) denotes free group coordinates;
scalar coordinates are used only for V22's
derivative and concentration estimates.
The definitions give
\[
 I_j\subset I_{j-1},\qquad E_j\subset E_{j-1},
 \qquad \Lambda_{R_j}\subset\Lambda_{R_{j-1}},
 \qquad\mathcal P_j\subset\mathcal P_{j-1}.
 \tag{27.2}
\]
Condition initially on the exterior
\(Y_{I_0^c}=y_0\), using the canonical kernel
\(\pi_{R_0,C}^{y_0}\). Let
\(z_j=Y_{I_{j-1}\setminus I_j}\) be the
annulus variables revealed at step \(j\).
Together \(y_0,z_1,\ldots,z_j\) specify the
full exterior \(y_j=Y_{I_j^c}\).
Use the filtration
\(\mathcal H_j=\sigma(z_1,\ldots,z_j)\),
with the fixed initial exterior understood.

Write \(V_j\) for the Haar-density potential
of (25.16) in \(R_j\). It has no free
exponential-coordinate Haar factors, since
the underlying measure here is group Haar.

\begin{proposition}[A normalized annular kernel, not a product ansatz]
\label{prop:f12v23-annulus}
There is the exact splitting
\[
 V_{j-1}(Y_{I_j},z_j;y_{j-1})
       =V_j(Y_{I_j};y_j)+U_j(z_j;y_{j-1}),
 \tag{27.3}
\]
where
\[
 U_j=\beta\sum_{p\in\mathcal P_{j-1}\setminus\mathcal P_j}s_p(W)
       -\sum_{e\in E_{j-1}\setminus E_j}\log j_e .
 \tag{27.4}
\]
In particular \(U_j\) is independent of \(Y_{I_j}\).
The conditional distribution of the next annulus is
\[
 \mathsf A_j(dz_j\mid y_{j-1})
  =\frac{e^{-U_j(z_j;y_{j-1})}Z_{R_j}(C,y_j)}
           {Z_{R_{j-1}}(C,y_{j-1})}\,dz_j ,
 \tag{27.5}
\]
with product normalized Haar \(dz_j\).
Given \(\mathcal H_j\), the unrevealed interior
has exactly the canonical law
\(\pi_{R_j,C}^{y_j}\).
\end{proposition}
\begin{proof}
The set inclusions follow from inclusion of
the regions and their incident coarse edges.
Subtract the two regional potentials.
A plaquette outside \(\mathcal P_j\)
contains no varied physical link of \(R_j\).
Every such physical exterior link is
independent of \(Y_{I_j}\), by V21.
For \(e\notin E_j\), the inverse-pivot
factor depends only on free variables
outside \(I_j\) and the fixed coarse field.
This proves (27.3)--(27.4), including
the asserted independence.

Product Haar decomposes as
\(dY_{I_{j-1}}=dY_{I_j}\,dz_j\).
Integration in \(Y_{I_j}\) gives
\[
 Z_{R_{j-1}}(C,y_{j-1})
       =\int e^{-U_j(z_j;y_{j-1})}
                         Z_{R_j}(C,y_j)\,dz_j .
 \tag{27.6}
\]
All these densities are positive and smooth.
Dividing by the parent density proves (27.5).
Dividing the joint density instead by its
annulus marginal leaves
\(e^{-V_j}/Z_{R_j}\), proving the conditional
interior assertion. The computation is
valid pointwise in the initial exterior.
No independence between annuli has been used.
\end{proof}

Thus the conditional statement needed in
V22 follows from an explicitly realizable
reveal of one sampled configuration. There
is no resampling in this construction.
Previously revealed annuli are all outside
the next parent, and no information about
its interior is revealed early.

\subsection{Uniform conditions that survive the reveal}

Use the common sufficiently small data
constants of V21--V22. Let
\[
 s_0=|R_0|,\quad q_j=\max\{1,|\mathcal F_{R_j}|\},
 \quad q_*=\max_{1\le j\le N}q_j,\quad
 \tau_*=\min\left\{\eta_0/2,
                         \sqrt{\kappa/(2K_3q_*)}\right\}.
 \tag{27.7}
\]
Assume \(N\ge1\); the zero-step case needs
no estimate. Require the support and buffer
condition (26.11) for every adjacent pair,
with distance at least \(L\). Make the
following deterministic assumptions on \(C\)
and the chosen geometry:
\[
 \begin{gathered}
 2A L_{R_0}(C)<\kappa/2,\qquad
 \epsilon_{R_0}\le\min\{\epsilon_0/2,\tau_*/K_*\},\\
 \eta_0e^{-aL}\le\tau_*/K_*,
 \qquad \beta\ge K_*^2/\tau_*^2 .
 \end{gathered}
 \tag{27.8}
\]
If necessary increase \(K_*\) to the
constant chosen in V22.

\begin{proposition}[The retained curvature requirements are inherited]
For \(S\subset R\),
\[
                L_S(C)\le L_R(C),\qquad
                \epsilon_S\le\epsilon_R .
 \tag{27.9}
\]
\end{proposition}
\begin{proof}
The affected corner plaquette set for \(S\)
is contained in the one for \(R\).
Extend each load by zero to all fine links.
Its definition as a sum of nonnegative
plaquette distances gives \(a_l^S\le a_l^R\).
The quantity \(L_S\) is the sum of
\((a_l^S)^2\) over all links: every link
with a nonzero load is either varied or
included in the completed frozen collar.
The analogous statement holds for \(R\).
Summing proves the first inequality.
The maximum over a subset proves the second.
\end{proof}

Let \(G_j\) denote the strict V21 data
condition for \(R_j\), including
\(\eta_{R_j}<\eta_0\) and \(\Gamma_{R_j}<\kappa\).
The coarse curvature part is deterministic.
The remaining parts are measurable in
\(Y_{I_j^c}\), hence \(G_j\in\mathcal H_j\).
Assume the fixed initial data satisfy \(G_0\).
By (27.8)--(27.9), on \(G_{j-1}\) all
hypotheses of V22 hold for the next transition:
in particular its parent mismatch is at
most \(\eta_0\), and \(\tau_{q_j}\ge\tau_*\).
This checks a condition at every step, not
just at the initial one.

\begin{theorem}[A stopped reveal with an actual survival probability]
\label{thm:f12v23-survival}
Under these assumptions, put
\[
 \delta_j=\min\{1,Kq_j e^{-c\beta/q_j}
                +D_1^{|R_{j-1}|}e^{-\kappa\beta}\},
 \qquad T=\min\{j\ge1:G_j^c\},
 \tag{27.10}
\]
with \(T=N+1\) if there is no failure through
step \(N\). Then
\[
 1_{G_{j-1}}\,
    \mathbb P(G_j^c\mid\mathcal H_{j-1})
                           \le1_{G_{j-1}}\delta_j ,
 \tag{27.11}
\]
and
\[
 \mathbb P(T\le N)
       \le1-\prod_{j=1}^N(1-\delta_j)
       \le\min\{1,\sum_{j=1}^N\delta_j\}.
 \tag{27.12}
\]
All probabilities are under the initial
canonical kernel, at its fixed \(C,y_0\).
\end{theorem}
\begin{proof}
Given \(\mathcal H_{j-1}\), Proposition
\ref{prop:f12v23-annulus} supplies exactly the
parent kernel used in V22. The event of
child failure depends only on its revealed
physical exterior, which is a function of
the current annulus and previous exterior.
Apply (26.20) on \(G_{j-1}\).
This proves (27.11) for the actual reveal.

The event \(\{T>j-1\}\) belongs to
\(\mathcal H_{j-1}\) and implies \(G_{j-1}\).
Consequently
\[
 \mathbb P(T>j)\ge(1-\delta_j)\mathbb P(T>j-1).
 \tag{27.13}
\]
Iterate from the known \(G_0\), and then use
the elementary product inequality. Once a
guard fails no further transition estimate
is asserted on that branch; the first-failure
probability already accounts for it.
\end{proof}

\subsection{A concrete growing family of annuli}

For a non-wrapping coarse cube of radius \(m\),
there is a fixed geometric constant \(C_b\) with
\[
                  |\mathcal F_R|\le C_b(m+1)^3 .
 \tag{27.14}
\]
Indeed the frozen links needed by its affected
plaquettes, its path stencils and affected
block equations lie in a fixed-width shell
around the boundary of the cube. The completed
collar adds only a further fixed width.
The difference between cubes with radii
\(m+b\) and \(\max\{0,m-b\}\) has at most
\(C_b(m+1)^3\) cells for fixed \(b\).
Each cell contains a fixed number of links.
This also bounds the free-coordinate stencil
closure of the physical collar.

\begin{corollary}[Many valid transitions on an explicit cofinal range]
\label{cor:f12v23-cubes}
Fix an inner radius \(r_*\ge1\). For sufficiently
small fixed \(b>0\), sufficiently large fixed
\(B>0\), and large \(\beta\), take an outer
radius \(m_0=\lfloor b\beta^{1/4}\rfloor\).
Take inward radial steps of width
\[
 w_\beta=\lceil B\log(\beta+2)\rceil+C_{\rm geom},
 \tag{27.15}
\]
ending at radius \(r_*\), with the last step
between \(w_\beta\) and \(2w_\beta\).
The number of transitions satisfies
\(N\asymp\beta^{1/4}/\log\beta\).
Suppose the initial exterior is protected and
\[
             \epsilon_{R_0}\le
                      \bar\epsilon/(m_0+1)^2
 \tag{27.16}
\]
for a sufficiently small fixed \(\bar\epsilon>0\).
Then all conditions (27.8) hold and
\[
             \mathbb P(T\le N)
                         \le C\beta e^{-c_0\beta^{1/4}} .
 \tag{27.17}
\]
The ambient lattice can be arbitrarily larger
than this non-wrapping outer cube.
\end{corollary}
\begin{proof}
The path and collar stencils have fixed range.
Choose \(C_{\rm geom}\) to include their
closures. Then adjacent supports have the
buffer required by (26.11), with
\(L\ge B\log(\beta+2)\).
The completed child collars lie entirely
in the parent's varied physical region.

There are at most \(C(m_0+1)^4\) affected
plaquettes. V16's load bound
\(L_{R_0}\le5832\epsilon_{R_0}^2|R_0|\)
and (27.16) give \(L_{R_0}\le C\bar\epsilon^2\).
Choose \(\bar\epsilon\) so
\(2A L_{R_0}<\kappa/2\).
By (27.14), \(q_*\le C\beta^{3/4}\),
so \(\tau_*\ge c\beta^{-3/8}\).
Condition (27.16) is \(O(\beta^{-1/2})\)
and hence is at most \(\tau_*/K_*\) for
large \(\beta\). The requirement
\(\beta\ge K_*^2/\tau_*^2\) follows too.
Choose \(B\) with \(aB>3/8\); it supplies
the buffer inequality in (27.8).

The outer volume is \(O(b^4\beta)\).
Choose \(b\) small enough that, for large
\(\beta\),
\[
 s_0\log D_1\le\kappa\beta/2,\qquad s_0\le c_1\beta .
 \tag{27.18}
\]
The latter additionally makes V21's spatial
kernel bounds available on each successful
region. The first part of \(\delta_j\)
is at most \(C\beta^{3/4}e^{-c\beta^{1/4}}\);
the second is at most \(e^{-\kappa\beta/2}\).
Sum over \(N\le C\beta^{1/4}\) using (27.12)
to obtain (27.17).
For large \(\beta\), set
\(N=\lfloor(m_0-r_*)/w_\beta\rfloor\);
use the equally spaced radii for \(j<N\)
and \(r_*\) for \(j=N\).
The final radial difference lies in
\([w_\beta,2w_\beta)\), which verifies
the stated construction and count.
\end{proof}

The initial protection assumption is still
explicit. Flat retained fields with exact
section exterior supply examples, and
sufficiently small perturbations supply open
examples. This corollary does not say that
an arbitrary Wilson sample has such an
outer boundary, nor that the allowed
curvature density stays fixed as \(\beta\)
or the cube size varies.

\subsection{Survival gives a normalized mixture, not a contraction}

Let \(\lambda_N\) be the law of the revealed
exterior \(y_N\) under the initial kernel.
It is the pushforward of the product of
the conditional annular kernels (27.5),
not a product of independent laws.
Let \(E=\{T>N\}\in\mathcal H_N\), and put
\(p=\mathbb P(E^c)\).
The law \(\mu_N\) of the unrevealed
\(Y_{I_N}\) has the exact representation
\[
 \mu_N(dY)=\int\pi_{R_N,C}^{y_N}(dY)\,\lambda_N(dy_N).
 \tag{27.19}
\]
When \(p<1\), define
\[
 \mu_N^{\rm good}(dY)=
    \frac{1}{1-p}\int 1_E(y_N)
                \pi_{R_N,C}^{y_N}(dY)\,\lambda_N(dy_N).
 \tag{27.20}
\]

\begin{proposition}[The discarded-history cost is exactly a mixture cost]
There is a probability measure
\(\mu_N^{\rm bad}\), arbitrary if \(p=0\), such that
\[
 \mu_N=(1-p)\mu_N^{\rm good}+p\mu_N^{\rm bad},\qquad
 \|\mu_N-\mu_N^{\rm good}\|_{\rm TV}\le p .
 \tag{27.21}
\]
Consequently bounded observables change in
expectation by at most \(2\|F\|_\infty p\).
Every component of \(\mu_N^{\rm good}\)
satisfies the protected final-kernel hypotheses.
\end{proposition}
\begin{proof}
Condition (27.19) on \(E\) and \(E^c\).
The event \(E\) is determined by the final
exterior, since all earlier exteriors are
subsets of it. Therefore, given \(y_N\),
conditioning on \(E\) introduces no extra
interior weight. This proves (27.20) and
the decomposition. The total variation
and expectation bounds follow immediately.
The probability \(1-p\) in (27.20) is
essential; dropping it would give a
subprobability measure, not the law
conditional on survival.
\end{proof}

For a bounded observable \(F\) depending only
on \(Y_{I_N}\), put
\[
 M_j^F=\int F\,d\pi_{R_j,C}^{y_j}
                  =\mathbb E(F\mid\mathcal H_j).
 \tag{27.22}
\]
These variables form a bounded martingale
under the initial law. This is a useful
consistency check on the exact annular
normalization. It is not a decreasing
sequence of norms or a contraction estimate.

\begin{proposition}[The remaining covariance has an explicit location]
\label{prop:f12v23-means}
For real bounded \(F,G\) depending only on
\(Y_{I_N}\),
\[
 \operatorname{Cov}_{\mu_N}(F,G)
   =\int\operatorname{Cov}_{\pi_{R_N,C}^{y_N}}(F,G)
                                          \,\lambda_N(dy_N)
      +\operatorname{Cov}_{\lambda_N}(M_N^F,M_N^G).
 \tag{27.23}
\]
Assume \( |R_N|\le c_1\beta\) and uniform
boxed gradient bounds \(L_F,L_G\), with
supports at distance \(d\), as in V21.
Then
\[
 \begin{split}
 \left|\operatorname{Cov}_{\mu_N}(F,G)
       -\operatorname{Cov}_{\lambda_N}(M_N^F,M_N^G)\right|
 &\le K\beta^{-1}e^{-\gamma d}
                              (L_FL_G+M_FM_G)\\
 &\quad+2M_FM_G\,p,\qquad
 p\le\min\{1,\textstyle\sum_j\delta_j\}.
 \end{split}
 \tag{27.24}
\]
\end{proposition}
\begin{proof}
Write the conditional expectation of \(FG\)
as its conditional covariance plus the
product of its conditional means, then
subtract the product of the unconditional
means. This gives (27.23), the usual total
covariance identity applied to this actual
annular mixture.
On \(E\), the final kernel obeys (25.14).
On \(E^c\), the elementary covariance bound
\(2M_FM_G\) applies. Integrate the two
bounds and use (27.12). No bound has been
assigned to the second term of (27.23).
\end{proof}

\begin{assessment}[Why the reveal does not yet produce long-distance mixing]
The chain proves that many successive
conditional kernels are protected, with
their real annulus weights and a controlled
first-failure probability. Multiplying the
kernel covariance estimates along that chain
would nevertheless be unjustified.
The martingale means in (27.22) can vary
with the revealed boundary, and their
covariance is still present in (27.23).
The common-variable obstruction in V19
already explains why good component
covariances cannot remove this term.

Moreover the concrete construction still
uses an outer volume \(O(\beta)\), a
decreasing curvature allowance, and a
protected initial exterior. It cannot
take the outer radius to infinity at
fixed \(\beta\). None of these restrictions
is removed by the small failure probability
in (27.17).
\end{assessment}

\begin{decision}[Next upstream quantity: oscillation of conditional means]
The actual nested reveal and its normalized
partition-function factors are now in place.
The next analytic quantity is the variation
of \(M_N^F\) across the boundary data
sampled by these annular kernels.
V21 controls derivatives only along paths
that remain protected; it did not construct
such paths between arbitrary protected
endpoints. Investigate that missing
geometric and normalization step before
using (27.23) as a mixing estimate.
Separately, an unrestricted argument must
still control its initial guard and remove
the \(D_1^{s_0}\) whole-box cost.
Next companion audit: V25. Next direction
review: V30.
\end{decision}

\begin{record}[V23 append record]
V22 had SHA--256
\begin{center}\scriptsize\ttfamily
65F473D9164773D960DB05C4C7241B368424296F4F8201CBD04B49DB9B534E6E.
\end{center}
The entire preceding body is retained
apart from the title version. Only active
V22 is replaced by V23; family archives
and the other programmes are untouched.
The overlap check distinguished the older
cut-forest protected chains and additive
mixture identities from the new exact
smooth-fibre annular kernel construction.
\end{record}

\begin{firewall}[Position after V23]
Proved: an actual outside-in reveal with
the canonical conditional interiors,
the exact normalized annular densities,
a stopped-chain protection estimate,
a concrete growing nested-cube regime,
and an explicit covariance decomposition
which retains the boundary-mean term.
Not proved: control of that mean covariance
under unrestricted boundary data, initial
protection for arbitrary samples,
fixed-coupling thermodynamic mixing,
scale iteration, continuum construction,
or the full Yang--Mills mass gap.
\end{firewall}

\section{V24: admissible boundary paths and the covariance of conditional means}
\label{sec:f12v24}

V23 left a definite term, not merely an unspecified
mixing assumption: the covariance of the conditional
means under the revealed-boundary law. We now
control that term when the boundary lies with high
probability in a stronger guard. The guard is a
Euclidean ball in the actual exterior free
coordinates. Its geometry supplies paths which
remain protected after the dependent physical
pivots have been reconstructed.

This is a smaller set than all V21-protected
boundaries. No convexity of the original physical
mismatch sublevel is asserted. Unlike the
near-identity half-slab tube in f6 V88, the
construction below is at fixed complete retained
field \(C\), allows compact retained holonomies,
and uses the globally solved smooth pivots of V18.

\subsection{A complete finite set of exterior free inputs}

For a connected region \(R\), let \(\mathcal F_R\)
be V16's completed frozen physical collar. Define
\(J_R^\partial\) to be the free group coordinates
among its non-tree links. They all belong to
\(I_R^c\). The group-coordinate convention of
V18 identifies a free physical link with its
free variable, so this set is unambiguous.

\begin{proposition}[The completed collar closes in free coordinates]
\label{prop:f12v24-inputs}
At fixed \(C\), all physical links in
\(\mathcal F_R\), and the entire regional
kernel \(\pi_{R,C}^y\), depend on the exterior
only through \(y_{J_R^\partial}\).
In particular
\[
 |J_R^\partial|\le|\mathcal F_R|,\qquad
 \text{each such dependence has bounded range
 and bounded stencil multiplicity.}
 \tag{28.1}
\]
When these free variables have their section
values, every physical link in \(\mathcal F_R\)
has its section value.
\end{proposition}
\begin{proof}
A free physical link in the collar is itself
one of the specified inputs. If a frozen pivot
belongs to the collar, its coarse edge is
unaffected by \(R\). The completion rule
adjoins all fine inputs of that edge's block,
including its nonpivot boundary variables
and internal variables in both endpoint cells.
All are frozen, and all its free inputs
therefore belong to \(J_R^\partial\).
The inverse pivot depends on those inputs
and its own fixed \(C_e\), but on no other
pivot. Hence there is no further iterative
dependency reaching a remote exterior.

Every frozen link in an affected plaquette,
and every frozen internal input of an
affected inverse pivot, was included at the
start of the collar construction. Formula
(25.2) therefore has only these exterior
dependencies. Bounded stencils and overlap
were part of the same finite-range completion.
At section free inputs the smooth pivot
equation has solution \(B_e=C_e\); its
global uniqueness proves the last assertion.
\end{proof}

For \(j\in J_R^\partial\), use coordinates
\[
 y_j=e^{v_j}\quad\hbox{on internal links},\qquad
 y_j=C_e e^{v_j}\quad\hbox{on boundary links}.
 \tag{28.2}
\]
Here \(v_j\in\mathfrak g\) and
\(\|v\|_2^2=\sum_j\|v_j\|^2\).
These are exterior coordinates, distinct
from the integrated interior coordinates \(u\).

\begin{proposition}[Comparison with physical mismatch]
\label{prop:f12v24-comparison}
There are fixed \(r_b>0,C_E,C_\infty,C_{\log}<\infty\)
such that, on \(\max_j\|v_j\|\le r_b\),
\[
 \mathcal E_R(v)\le C_E\|v\|_2^2,\qquad
 \eta_R(v)\le C_\infty\max_j\|v_j\|.
 \tag{28.3}
\]
Conversely, if every physical mismatch in
\(\mathcal F_R\) is below a fixed sufficiently
small \(b_{\log}>0\), the coordinates (28.2)
are uniquely defined in this chart and
\[
              \|v\|_2^2\le C_{\log}^2\mathcal E_R.
 \tag{28.4}
\]
\end{proposition}
\begin{proof}
Every reconstructed physical link is a smooth
function of a bounded number of these free
inputs, equals its section value at \(v=0\),
and has uniformly bounded derivatives there.
Integrating its derivative along \(tv\) gives
its mismatch at most \(K\) times the sum of
nearby \(\|v_j\|\). Squaring and summing uses
bounded stencil overlap and proves the first
part of (28.3); the maximum bound is immediate.
For (28.4), each free variable in (28.2) is
itself a physical link of \(\mathcal F_R\).
The local inverse exponential bound gives
\(\|v_j\|\le C_{\log} b_j\).
Choose \(b_{\log}\) inside that inverse chart
and sum only these free-link mismatches.
The right side is bounded by the full
\(\mathcal E_R\).
\end{proof}

\subsection{A convex guard which really preserves all coarse outputs}

Assume the deterministic retained field obeys
\[
              2A L_R(C)\le\kappa/2,\qquad
              \epsilon_R<\epsilon_0 .
 \tag{28.5}
\]
Choose once and for all \(r_\partial>0\) with
\[
 r_\partial<r_b,\qquad
 C_\infty r_\partial<\eta_0/2,\qquad
 K_3 C_E r_\partial^2\le\kappa/4 .
 \tag{28.6}
\]
Let \(\mathcal B_R(C)\) be the set of exterior
data whose inputs (28.2) lie in the ball
\[
                       \|v\|_2\le r_\partial .
 \tag{28.7}
\]
No restriction is placed on other exterior
free coordinates. If \(J_R^\partial\) is
empty, the ball is a singleton in dimension
zero and the kernel is exterior-independent.

\begin{theorem}[Explicit protected paths between the relevant endpoints]
\label{thm:f12v24-path}
Every exterior in \(\mathcal B_R(C)\) satisfies
\[
                \eta_R<\eta_0,\qquad
                \Gamma_R\le3\kappa/4<\kappa .
 \tag{28.8}
\]
For two such exteriors, the path
\[
                   v(t)=(1-t)v^{(0)}+tv^{(1)}
 \tag{28.9}
\]
stays protected. Reconstructing all pivots
by V18 keeps the complete coarse field
equal to \(C\) throughout.
The regional kernel at the path's endpoints
equals the kernel at the two prescribed
exteriors, even if their irrelevant remote
free variables differ.
\end{theorem}
\begin{proof}
The Euclidean ball is convex. Apply (28.3),
(28.5) and (28.6) along the whole path to
obtain (28.8). Keep all exterior free
variables outside \(J_R^\partial\) fixed
at the first exterior's values. For each
point of the path the global inverse-pivot
map gives a fine field with precisely the
same complete \(C\); unaffected coarse
outputs are not recomputed and reset.
The path changes no interior free variable.
By Proposition~\ref{prop:f12v24-inputs}, the
regional kernel at the other endpoint
depends only on the specified second
input values. The remote exterior need
not itself reach the second remote
configuration, because it has no effect
on this kernel. Finally the inequalities
in (28.8) are strict, so V21's derivative
theorem applies on a neighbourhood of
the path, including its endpoints.
\end{proof}

The ball is invariant under simultaneous
cell-constant gauge transformations: each
coordinate \(v_j\) is conjugated at its
appropriate endpoint, and the chosen Lie
norm is invariant. Its small radius does
not require \(C_e\) to be close to \(I\).

\subsection{Integrating the normalized boundary derivative}

Let \(F\) be a real bounded function of
the interior free variables, with no
explicit exterior dependence. Suppose its
boxed gradient has a nonempty support
\(A_F\), and use V21's \(M_F,L_F\).
Set
\[
 N_F=L_F+M_F,\qquad
 W_F(R)=\left(\sum_{j\in J_R^\partial}
                   e^{-2\gamma\,{\rm dist}(A_F,j)}\right)^{1/2},
 \qquad m_F(y)=\mathbb E_{\pi_{R,C}^y}F .
 \tag{28.10}
\]
The fixed multiplicity of Lie components
is absorbed into constants. If the boxed
gradient support is empty, the following
weighted bound still holds with all
distances set to zero, as in V21; no
spatial assertion is then inferred.

\begin{theorem}[Uniform oscillation of conditional means on the guard]
\label{thm:f12v24-oscillation}
For \(|R|\le c_1\beta\), under (28.5),
\[
 \operatorname{osc}_{\mathcal B_R(C)}m_F
       \le\Delta_F:=
                 2K r_\partial N_F W_F(R).
 \tag{28.11}
\]
For a fixed-size nonempty support \(A_F\),
if \(d_F={\rm dist}(A_F,J_R^\partial)\),
then
\[
 W_F(R)\le K|A_F|^{1/2}e^{-\gamma' d_F}
 \tag{28.12}
\]
for fixed \(\gamma'>0\).
An empty exterior input set gives zero
oscillation exactly.
\end{theorem}
\begin{proof}
The right-trivialized velocity of
\(e^{v_j(t)}\), or \(C_e e^{v_j(t)}\),
has norm at most
\(K\|v_j^{(1)}-v_j^{(0)}\|\), by the
bounded derivative of the exponential
on the fixed chart. Apply (25.19) along
the admissible path of
Theorem~\ref{thm:f12v24-path}, and integrate:
\[
 |m_F(y^{(1)})-m_F(y^{(0)})|
 \le K N_F\sum_{j\in J_R^\partial}
          e^{-\gamma\,{\rm dist}(A_F,j)}
                    \|v_j^{(1)}-v_j^{(0)}\|.
 \tag{28.13}
\]
Cauchy--Schwarz and
\(\|v^{(1)}-v^{(0)}\|_2\le2r_\partial\)
give (28.11). The derivative used here
already includes the partition-function
normalization through (25.17).

For (28.12), bound the distance-to-set
exponential by the sum of exponentials
from the individual points of \(A_F\).
Uniform polynomial ball growth and
bounded coordinate multiplicity bound
the remaining tail by
\(K|A_F|e^{-\gamma d_F}\), after absorbing
its polynomial factor and reducing the
exponent. Take a square root and rename
the exponent. If there are no exterior
inputs, Proposition~\ref{prop:f12v24-inputs}
gives exact independence instead.
\end{proof}

\subsection{The stronger guard is conditionally populated}

Let \(P\supset R\) be a protected parent
as in V22. Require the full physical collar
\(\mathcal F_R\) and its reconstruction
stencils to lie inside the parent's varied
region, with buffer \(L\) from its frozen
collar. Put \(q=\max\{1,|\mathcal F_R|\}\) and
\[
 \sigma_q=\min\left\{b_{\log}/2,\,
                     \frac{r_\partial}{2C_{\log}\sqrt q}\right\}.
 \tag{28.14}
\]

\begin{proposition}[A probability bound for failure of the convex guard]
\label{prop:f12v24-population}
Assume the parent hypotheses of V22 and
\[
                 K_{\rm link}\,m_P(L)\le\sigma_q/2 .
 \tag{28.15}
\]
Then the full compact parent kernel obeys
\[
 \pi_{P,C}^{y_P}(\mathcal B_R(C)^c)
       \le\min\{1,Kq e^{-c\beta\sigma_q^2}
                             +D_1^{|P|}e^{-\kappa\beta}\}.
 \tag{28.16}
\]
In particular \(\sigma_q^2\ge c_\partial/q\)
with \(c_\partial>0\) fixed.
\end{proposition}
\begin{proof}
If every physical child-collar mismatch
is less than \(\sigma_q\), (28.4) applies
and gives
\[
       \|v\|_2^2\le C_{\log}^2 q\sigma_q^2
                                      \le r_\partial^2/4.
 \tag{28.17}
\]
Thus failure of the ball requires a
physical mismatch at least \(\sigma_q\).
Use the proof of (26.18) with this
threshold instead of \(\tau_q\).
The coordinate-mean bound (28.15),
the local physical-link Lipschitz estimate
and the scalar sub-Gaussian tails give
\(Kq e^{-c\beta\sigma_q^2}\) on the parent
box. The full parent box-complement
probability is the second term in (28.16).
Since \(q\ge1\), (28.14) has square at
least \(\min\{b_{\log}^2/4,
r_\partial^2/(4C_{\log}^2)\}/q\).
No probability estimate was conditioned
on the desired successful child event.
\end{proof}

V23's actual outside-in reveal may use
these convex guards instead of its broader
guards. Conditions (28.15) are obtained
from (26.19) with \(\sigma_q\) replacing
\(\tau_q\) and with enlarged fixed constants.
Each successful convex guard implies the
ordinary parent protection by (28.8).
It is exterior-measurable, so the exact
annulus kernels and the survival proof
remain unchanged. The concrete nested
cubes of V23 still give
\[
 \mathbb P(\hbox{a convex guard fails before the last reveal})
                    \le C\beta e^{-c\beta^{1/4}},
 \tag{28.18}
\]
after adjusting \(b,B,\bar\epsilon\) and
the lower threshold for \(\beta\).
Indeed \(q_*=O(\beta^{3/4})\),
\(\sigma_{q_*}\ge c\beta^{-3/8}\),
the curvature allowance is \(O(\beta^{-1/2})\),
and the buffer is \(B\log\beta\), exactly
the comparisons used in that proof.
The initial parent remains an explicit
protected datum; it need not satisfy the
stronger convex guard for this first
conditional transition.

\subsection{The mean-covariance term is now bounded on this regime}

Let \(\lambda\) be any probability law on
the exterior at fixed \(C\), and form the
normalized mixture
\[
                    \mu=\int\pi_{R,C}^y\,\lambda(dy).
 \tag{28.19}
\]
Put \(p=\lambda(\mathcal B_R(C)^c)\).
No product structure or independence is
assumed for \(\lambda\). In particular it
may be the actual final-annulus marginal
of V23, including all of its partition
function weights.

\begin{proposition}[Covariance of the means with an explicit bad-data cost]
\label{prop:f12v24-mean-covariance}
For real \(F,G\) as above and \(|R|\le c_1\beta\),
\[
 |\operatorname{Cov}_{\lambda}(m_F,m_G)|
                  \le\frac{\Delta_F\Delta_G}{4}
                                     +5M_FM_Gp .
 \tag{28.20}
\]
\end{proposition}
\begin{proof}
If \(p<1\), the covariance under
\(\lambda(\cdot\mid\mathcal B_R(C))\)
is at most \(\Delta_F\Delta_G/4\).
Indeed a real random variable in an
interval of length \(\Delta\) has variance
at most \(\Delta^2/4\): center at the
interval midpoint and use that variance
is no larger than the second moment
about any constant. Then apply
Cauchy--Schwarz.

For \(0<p<1\), decompose \(\lambda\) into
its good and bad restrictions. The exact
two-component covariance identity consists
of their weighted covariances plus
\[
 p(1-p)(\mathbb E_{\rm good}m_F-\mathbb E_{\rm bad}m_F)
        (\mathbb E_{\rm good}m_G-\mathbb E_{\rm bad}m_G).
 \tag{28.21}
\]
The bad covariance is bounded by \(M_FM_G\),
since \(|m_F|\le M_F\) and \(|m_G|\le M_G\).
The two differences in (28.21) are at most
\(2M_F\) and \(2M_G\).
This proves (28.20) with the cost \(p+4p\).
The case \(p=0\) uses only the good bound;
for \(p=1\) the elementary bound \(M_FM_G\)
already implies the assertion.
\end{proof}

\begin{theorem}[A complete conditional-mixture covariance estimate]
\label{thm:f12v24-mixture}
Assume (28.5), \(|R|\le c_1\beta\), and
uniform interior boxed gradient bounds for
the exterior-independent observables \(F,G\).
Then
\[
 \begin{split}
 |\operatorname{Cov}_{\mu}(F,G)|
 &\le K\beta^{-1}e^{-\gamma\,{\rm dist}(A_F,A_G)}
                                   (L_FL_G+M_FM_G)\\
 &\quad+\frac{\Delta_F\Delta_G}{4}+7M_FM_Gp .
 \end{split}
 \tag{28.22}
\]
For fixed-size nonempty supports the second
term is at most
\[
 K r_\partial^2 N_FN_G
                 |A_F|^{1/2}|A_G|^{1/2}
                           e^{-\gamma'(d_F+d_G)} .
 \tag{28.23}
\]
\end{theorem}
\begin{proof}
Apply the exact total covariance identity
(27.23) with the present \(\lambda\).
On the convex guard the conditional
covariance is bounded by V21, because
(28.8) verifies its data hypotheses.
On its complement use \(2M_FM_G\).
The mean covariance is bounded by (28.20).
Adding the two bad-data costs gives
\(2p+5p=7p\), proving (28.22).
Equations (28.11)--(28.12) give (28.23).
Every term in the total covariance
identity has now been bounded under the
stated regime, rather than discarding
the covariance of the means.
\end{proof}

There is no assertion that \(d_F+d_G\)
is automatically comparable to
\({\rm dist}(A_F,A_G)\).
For observables close to the boundary,
(28.23) need not be small. To deduce a
mutual-distance estimate from it, one
must choose a region placing both
supports sufficiently far inside and
verify that region's size, curvature
and initial-parent conditions.

For the actual nested mixture of V23,
the event that all convex guards survive
implies its last guard. Thus \(p\) is
bounded by the stopped-chain failure
estimate (28.18), not by a probability
computed after renormalizing onto
successful histories. Alternatively,
a single buffered protected parent
gives \(p\) directly from (28.16).

\begin{assessment}[The resolved term and the still unresolved scale problem]
The boundary-mean term isolated in V23 is
controlled on a conditionally populated
convex guard. The path is constructed
in the correct free variables, preserves
the complete coarse field, and includes
the effect of all dependent exterior
pivots. This is the needed improvement
over a derivative estimate on unspecified
admissible paths.

It does not establish unrestricted
fixed-coupling mixing. The estimates
still use \(|R|\le c_1\beta\), a fixed
small total curvature load, and a
protected ancestor to populate the
guard. The ball radius bounds a total
squared boundary deviation, not a
positive density allowance on arbitrarily
large boundaries. The ambient-volume
independence of constants does not
remove any of these hypotheses.
\end{assessment}

\begin{decision}[V25 audit and the remaining upstream normalization problem]
At V25 inspect the current SM, NS and
parallel-YM files as requested.
The present finite protected-region
mean-covariance issue is no longer left
as a blank input. The next step toward
fixed-coupling control must address the
unprotected exterior and the extensive
whole-box comparison factor
\(D_1^{|P|}\), with the exact annular
normalization retained. Merely extending
the nested list or assuming another
protected ancestor would repeat the
same restricted argument.
Investigate a comparison localized to
the failed collar or its actual
obstruction carriers, and state
explicitly which conditional partition
function ratio it requires.
Next companion audit: V25. Next direction
review: V30.
\end{decision}

\begin{record}[V24 append record]
V23 had SHA--256
\begin{center}\scriptsize\ttfamily
9272150C2214CD9619E91E11ABA62241C83F0C86683530FB4FBEDF34F78FEEB8.
\end{center}
Its complete body is preserved apart
from the title version. Only active
V23 is replaced by V24; all family
archives and other programmes remain.
The earlier half-slab boundary tube
was inspected to avoid confusing it
with the present fixed-\(C\), inverse-pivot
boundary path construction.
\end{record}

\begin{firewall}[Position after V24]
Proved: a finite closed exterior-input
description, convex protected paths at
fixed complete \(C\), oscillation bounds
for normalized conditional means,
conditional population of the stronger
guard, and a full mixture covariance
bound including the mean term.
Not proved: unrestricted initial
protection, a volume-uniform comparison
across unprotected data, fixed-coupling
thermodynamic mixing, scale iteration,
continuum construction, or the physical
Yang--Mills mass gap.
\end{firewall}

\section{V25: audit and an arbitrary-exterior bulk-density estimate}
\label{sec:f12v25}

\subsection{Required companion audit and its effect on direction}

The latest companion files were inventoried, their
end results read, and their full hashes checked:
\begin{center}\small
\begin{tabular}{lrr}
\toprule
Programme & version & bytes\\
\midrule
SM--Einstein & V72 & 607372\\
NS stability & V26 & 278283\\
Parallel YM & V5 & 54174\\
\bottomrule
\end{tabular}
\end{center}
Their respective SHA--256 values are
\begin{center}\scriptsize\ttfamily
F44F9745D43379E1672C5550411B58E97195A4C9278592D221DF5D3F644DC1F5\\
82C887F8C2C69E4F28FAD7C3DC8DE5B9099CB5A4BF431EBA1FFF3E91935978AD\\
306F1BB84CFA8A8D47EA569EC17E9545F9867C8D32B548EC9D9C444B4F3C0512.
\end{center}
All are unchanged from the previous audit.
SM V72 proves a fixed-mode, fixed-time many-copy
limit, not the physical finite-species continuum.
NS V26 supplies pointwise rank-one Oseen norms,
not global regularity; its planned V27 is not
a completed result. Parallel YM V5 warns that
an extensive global integer charge cannot have
the proposed volume-uniform tightness under
tensorization. Its document-level stop applies
to that route, not as an instruction to stop
the present user-requested programme.

The last warning is useful here. Requiring an
entire growing boundary or fibre to remain in
one small ball is not the same task as
controlling its defect density. We now use
V18's global smooth inverse to build a feasible
repair for \emph{arbitrary} exterior data.
Its action cost is an interface cost.
The resulting theorem permits arbitrarily
large regions at fixed sufficiently large
\(\beta\), but controls a bulk density, not
the probability that the entire region is good.
No quantitative estimate is imported from
the companion programmes.

\subsection{A global repair whose discrepancy is confined to the interface}

Continue to use \(\widetilde K\), its free
variables, and the canonical kernel (25.2).
Let \(R\) be connected, \(s=|R|\), and
let \(\partial E(R)\) consist of the coarse
edges with exactly one endpoint in \(R\).
Set
\[
 b_R=|\mathcal F_R|+|\partial E(R)|,\qquad
 d_G=\max_{U\in G}\|I-U\|_{\rm HS},\qquad
 s_{\max}=d_G^2/2 .
 \tag{29.1}
\]
The number \(b_R\) is geometric; it does not
measure the smallness of the actual exterior.
For a non-wrapping cube of radius \(m\),
\[
                    b_R\le C_b(m+1)^3 .
 \tag{29.2}
\]
The frozen-collar part was proved in V23;
the crossing-edge count has the same
fixed-width boundary-shell bound.
For a full periodic region \(b_R=0\).

\begin{proposition}[A feasible repair for every exterior]
\label{prop:f12v25-repair}
Fix arbitrary \(C\) and exterior free data \(y\).
Set all free internal variables in \(I_R\)
to \(I\), and all free boundary variables
over \(e\in E(R)\) to \(C_e\).
Solve the affected pivots with the global
inverse of V18. This gives a feasible field
\(W^{\rm rep}\) at the prescribed complete
coarse field, for every \(y\).
Every varied link agrees with the section
except possibly the root pivots on
\(\partial E(R)\).
If \(H_{\rm rep}=H_\xi(W^{\rm rep})\), then
\[
                       H_{\rm rep}-H_0(C)
                                      \le6s_{\max} b_R .
 \tag{29.3}
\]
\end{proposition}
\begin{proof}
The global smooth pivot is defined for every
tuple of compact inputs. Thus no local
implicit-function or small-exterior assumption
is needed for feasibility.
If both endpoints of \(e\) lie in \(R\),
all internal factors of its block are now
section factors, and its nonpivot boundary
links equal \(C_e\). The solution \(B_e=C_e\)
is the section solution. Global uniqueness
identifies it with the solved pivot.
Only crossing-edge pivots can differ.

Compare every affected plaquette with its
full section plaquette. A difference can
occur only if it touches one of these
crossing pivots or a frozen link of
\(\mathcal F_R\). There are at most
\(6b_R\) such plaquettes, since a fine link
belongs to at most six plaquettes.
Both plaquette actions lie in
\([0,s_{\max}]\), so their difference
from above is at most \(s_{\max}\).
All other plaquette actions agree exactly.
Summation proves (29.3). The exterior
coarse outputs remain fixed automatically
by the global inverse construction.
\end{proof}

\begin{proposition}[A matched lower density without local convexity]
\label{prop:f12v25-denominator}
There are \(b>0,\beta_0<\infty\), independent
of region, volume, retained field and
exterior, such that
\[
 Z_R(C,y)\ge b^s\beta^{-N_R/2}
                \exp[-\beta H_0(C)-6\beta s_{\max}b_R],
 \qquad N_R=d(17s+7|E(R)|).
 \tag{29.4}
\]
\end{proposition}
\begin{proof}
Use exponential coordinates about the free
section values chosen in the repair. On a
fixed small cube, all solved pivots and their
fixed-order derivatives have uniform bounds
by V18, even when the physical exterior is
far from the section. Their density factors
are uniformly positive; the free Haar factors
are uniformly positive on the chosen
exponential chart too. Their product is
bounded below by a constant to the power \(s\).

The regional action Hessian has a uniform
upper operator-norm bound \(K\), by finite
stencil overlap and the global derivative
bounds. Positivity of that Hessian is
\emph{not} required. On the symmetric cube
of scalar radius \(t/\sqrt\beta\), Taylor's
upper estimate gives
\[
 f_R(u)\le H_{\rm rep}+g\cdot u+K|u|^2/2 .
 \tag{29.5}
\]
Here \(N_R\le73ds\). The quadratic
exponential cost is at most a constant
to the power \(s\), and symmetry gives
\(\int_{\rm cube}e^{-\beta g\cdot u}\,du
\ge\operatorname{vol}({\rm cube})\).
Thus \(Z_R\ge b^s\beta^{-N_R/2}
e^{-\beta H_{\rm rep}}\).
Use (29.3). This is the actual full
compact density, not a convex extension
or a conditional density on a good box.
\end{proof}

\subsection{A full conditional tail with a surface debit}

Let \(D=\sum_{l\in\Lambda_R}
\|W_l-J(C)_l\|_{\rm HS}^2\), retaining the
pivots. Its free-link part is \(D_{\rm free}\).
No smallness assumption on exterior mismatch
is made in this subsection. The deterministic
inequalities (25.4), which were already global,
give
\[
 H_\xi-H_0\ge\frac{D}{2A}-E_{\rm v,out}
                            -\frac A2L_{\rm in}-M_R .
 \tag{29.6}
\]
Moreover
\[
 \mathcal E_R\le d_G^2|\mathcal F_R|,\qquad
 E_{\rm v,out}\le12\mathcal E_R,\qquad
 M_R\le\frac A2L_{\rm out}+\frac{\mathcal E_R}{2A}.
 \tag{29.7}
\]
The last inequality is Young's inequality
with a different split from V16.

\begin{theorem}[Arbitrary-exterior density tail]
\label{thm:f12v25-tail}
There are fixed \(D_2\ge1\) and \(C_\partial<\infty\)
such that, for every \(x\ge0\), every \(C,y\),
and \(\beta\ge\beta_0\),
\[
 \pi_{R,C}^{y}\{D\ge x\}
 \le\min\left\{1,\,
 D_2^s\exp\left[-\beta\left(
          \frac{x}{4A}-\frac A2 L_R(C)-C_\partial b_R
                              \right)\right]\right\}.
 \tag{29.8}
\]
One may take
\[
                  C_\partial=(15+1/(2A))d_G^2 .
 \tag{29.9}
\]
There is no restriction \(s\le c\beta\),
no small total curvature assumption, and
no protected ancestor in (29.8).
\end{theorem}
\begin{proof}
On \(\{D\ge x\}\), the first term in (29.6)
is at least \(x/(4A)+D_{\rm free}/(4A)\).
Bound the inverse-pivot factors by their
global upper bound, and integrate
\(\exp[-\beta D_{\rm free}/(4A)]\) over
the free group variables. V13's compact
Gaussian integral estimate gives
\(B^s\beta^{-N_R/2}\). Hence the restricted
unnormalized integral is at most
\[
 B^s\beta^{-N_R/2}e^{-\beta H_0}
       e^{-\beta(x/(4A)-E_{\rm v,out}
                              -(A/2)L_{\rm in}-M_R)} .
 \tag{29.10}
\]
Divide by the actual lower density (29.4).
The Gaussian powers and \(H_0\) cancel
exactly. Equations (29.7) bound the
remaining debit by
\[
 \frac A2L_R+
       (12+1/(2A))d_G^2|\mathcal F_R|
                              +6s_{\max}b_R .
 \tag{29.11}
\]
Since \(6s_{\max}=3d_G^2\) and
\(|\mathcal F_R|\le b_R\), this is at
most \((A/2)L_R+C_\partial b_R\).
Take \(D_2=\max\{1,B/b\}\).
The factor \(D_2^s\) has not disappeared;
it will be beaten by an extensive event
cost, not by the cost of one bad link.
\end{proof}

\begin{corollary}[Intensive distance rather than an all-good event]
\label{cor:f12v25-moment}
Uniformly over arbitrary exteriors,
\[
 \frac{\mathbb E_{\pi_{R,C}^y}D}{s}
 \le\frac{4A\log D_2}{\beta}
       +2A^2\frac{L_R(C)}s
       +4AC_\partial\frac{b_R}s+\frac{4A}{\beta s}.
 \tag{29.12}
\]
If every affected corner curvature is at
most \(\epsilon\), then \(L_R/s\le C_L\epsilon^2\),
with \(C_L=5832\) from V16.
\end{corollary}
\begin{proof}
Put
\(x_0=4A[s\log D_2/\beta+(A/2)L_R+C_\partial b_R]\).
For \(t\ge0\), (29.8) gives
\(\mathbb P(D\ge x_0+t)\le e^{-\beta t/(4A)}\).
Integrating this tail and the trivial
bound below \(x_0\) gives
\(\mathbb ED\le x_0+4A/\beta\).
Divide by \(s\), and apply the earlier
corner-incidence load bound.
\end{proof}

For a threshold \(0<\tau<d_G\), let
\[
 B_\tau(R)=\#\{l\in\Lambda_R:
                       \|W_l-J(C)_l\|_{\rm HS}\ge\tau\}.
 \tag{29.13}
\]
Since \(D\ge\tau^2B_\tau(R)\), the tail
(29.8) bounds the event \(B_\tau(R)\ge k\)
with \(x=\tau^2k\). There is no subset
entropy to add: the pointwise distance
bound already holds on the entire event
that at least \(k\) links are bad.
There are at most \(81s\) varied links.

\begin{corollary}[Defect-dense large cubes are suppressed at fixed coupling]
\label{cor:f12v25-density}
Fix \(0<\theta<81\) and \(0<\tau<d_G\).
Suppose
\[
 \epsilon^2\le\frac{\tau^2\theta}{8A^2 C_L}.
 \tag{29.14}
\]
There are a fixed minimum cube radius
\(m_*(\tau,\theta)\) and a fixed
\(\beta_*(\tau,\theta)\) such that for
every cube \(R\) of radius \(m\ge m_*\),
every exterior, and \(\beta\ge\beta_*\),
\[
 \pi_{R,C}^y\{B_\tau(R)\ge\theta s\}
                \le\exp[-\beta\tau^2\theta s/(16A)] ,
 \tag{29.15}
\]
provided its affected corner curvatures
are at most \(\epsilon\).
The cube radius has no upper bound
depending on \(\beta\).
\end{corollary}
\begin{proof}
The curvature debit in (29.8) is at
most \(\tau^2\theta s/(16A)\) by (29.14).
The geometric ratio \(b_R/s\) tends
to zero for cubes, by (29.2).
Choose \(m_*\) so that
\(C_\partial b_R\le\tau^2\theta s/(16A)\).
The net action cost at
\(x=\tau^2\theta s\) is then at least
\(\tau^2\theta s/(8A)\).
Finally choose
\(\beta_*\ge\max\{\beta_0,16A\log D_2/(\tau^2\theta)\}\).
It absorbs \(D_2^s\) and proves (29.15).
\end{proof}

\subsection{Fixed-spacing infinite-volume interpretation}

The following existence step is the standard
compact-spin Gibbs construction; its relevance
is to identify where the new estimate holds.
It is not a construction of continuum
Yang--Mills theory. For background on the
DLR formalism see S.~Friedli and Y.~Velenik,
\emph{Statistical Mechanics of Lattice Systems},
\href{https://www.cambridge.org/core/books/abs/statistical-mechanics-of-lattice-systems/contents/CAC60015A56FE7C49F4F31E297915B88}
{Chapter 6, Infinite-Volume Gibbs Measures}.
The argument needed for this particular
fixed-\(C\) interaction is supplied next.

\begin{proposition}[Auxiliary compact-fibre Gibbs states exist]
\label{prop:f12v25-dlr}
Fix \(\beta<\infty\), a lattice spacing,
and an infinite retained field \(C\).
The local smooth formulas
\[
                 \beta s_p(W(Y,C)),\qquad
                           -\log j_e(Y,C_e)
 \tag{29.16}
\]
define a finite-range continuous interaction
on the countable product of compact free
group variables. Its Gibbs specification
has at least one infinite-volume DLR state.
On the regional sets \(I_R\) its kernels
are precisely (25.2). Every such state
inherits (29.8), (29.12), and, when its
coarse curvature hypotheses hold, (29.15).
\end{proposition}
\begin{proof}
Each pivot uses finitely many free inputs,
has a global smooth inverse, and its
Jacobian is bounded above and away from
zero. Thus all terms in (29.16) are
continuous bounded functions of finite
stencils, with a fixed finite number of
terms incident on any coordinate.
For any finite free-coordinate set,
sum all terms depending on it and
normalize its positive density against
product Haar, keeping other coordinates
fixed. Terms independent of the updated
set cancel. This gives a consistent
specification; for \(I_R\) it is exactly
V21's kernel, as also checked by V23's
annular factorization.

Take finite-volume states in an increasing
sequence of coordinate sets, with arbitrary
fixed exterior values. The countable
product of compact metrizable groups is
compact metrizable, so these probability
measures have a weakly convergent
subsequence. For any finite update set
and continuous local test function \(F\),
its conditional expectation under the
specification is continuous and local:
the integrand is continuous on a finite
compact product and its denominator
is strictly positive.
For all sufficiently large volumes the
finite-volume consistency identity
\(\mu_n F=\mu_n(\pi_R F)\) therefore
passes to the weak limit.
These identities for continuous local
functions extend to bounded measurable
functions by the usual monotone-class
argument, since they determine measures
on the compact product. Properness of
the kernels gives the conditional DLR
identity. This proves existence without
asserting uniqueness.

The new bounds hold pointwise for every
exterior of every fixed region. Condition
the DLR state on that exterior and
integrate them. No passage to a limit
through a threshold indicator is needed.
The coarse field is fixed throughout;
no law for a fluctuating retained field
or scale iteration is constructed here.
\end{proof}

\subsection{Separated defect-dense blocks and what they do not detect}

Let \(\mathcal B_i=\{B_\tau(R_i)\ge\theta|R_i|\}\)
for cubes satisfying (29.15). If the cubes
are separated enough that \(\mathcal B_j\)
is measurable outside \(I_{R_i}\) whenever
\(i\ne j\), then for every finite family,
under any of the above DLR states,
\[
             \mathbb P\left(\bigcap_i\mathcal B_i\right)
                   \le\prod_i
                         e^{-\beta\tau^2\theta|R_i|/(16A)} .
 \tag{29.17}
\]
To prove this, condition on the exterior of
one cube, pull all the other event indicators
outside that conditional expectation, and
use the arbitrary-exterior estimate (29.15).
Repeat. This is a conditional bound proved
for this specification, not V17's bare-law
estimate reused after an arbitrary history.

For completeness, equal cubes on a block
lattice can be colored with a fixed finite
number \(\chi\) of colors so that this
measurability separation holds within
each color. Their event supports have
only a fixed extra stencil width.
Among \(k\) specified blocks some color
has at least \(k/\chi\) members; hence
their joint defect-density event has
probability at most \(\delta^{k/\chi}\),
where \(\delta=e^{-\beta\tau^2\theta s/(16A)}\).
If the block adjacency degree is at most
\(\Delta\ge2\), the usual depth-first
tree encoding bounds the number of
connected \(k\)-block sets through one
root by \(\Delta^{2k}\).
Therefore, when \(Q=\Delta^2\delta^{1/\chi}<1\),
\[
 \mathbb P(\text{a rooted connected defect-dense
                         cluster has size at least }k_0)
                         \le\frac{Q^{k_0}}{1-Q}.
 \tag{29.18}
\]
The counting principle was used earlier;
the new input is its genuinely
arbitrary-exterior conditional activity
for the present smooth fibre.

\begin{assessment}[A bulk density does not rule out a thin defect connection]
These theorems remove the protected-ancestor
assumption for the specified density events
and permit the fixed-coupling large-volume
limit. They do not say that every small
region or every boundary collar is protected.
A long thin set of bad links can cross
many blocks without making any block
defect-dense. Thus (29.18) is not a bound
on all defect connections, and (29.12)
does not imply the total boundary-ball
condition of V24.

Likewise the surface debit in (29.8)
is harmless for large cubes with a
fixed positive bad density, but may
be comparable to the entire event
cost for a thin carrier. Dropping
that debit in a contour application
would be an error. Gibbs existence
for the eliminated variables at one
spacing does not give uniqueness,
exponential physical clustering,
reflection-positive continuum limits,
or a Hamiltonian mass gap.
\end{assessment}

\begin{decision}[Next task: sparse carriers, not another dense-box estimate]
The global smooth inverse supplies the
missing arbitrary-exterior repair, and
the repaired denominator prices its
interface. The next obstruction is
sparse or thin bad geometry: formulate
the actual local carrier whose presence
can defeat the boundary-mean argument,
then compare its normalized activity
without charging the boundary of an
arbitrarily enlarged all-good box.
The available fixed-\(C\) conditional
density estimate is now stronger than
a protected-ancestor assertion, but
it is not yet that sparse-carrier
estimate. Next companion audit and
direction review: V30.
\end{decision}

\begin{record}[V25 append record]
V24 had SHA--256
\begin{center}\scriptsize\ttfamily
7F8B178F7B368E12574847C1564DD72C85E69C05671A0E6C980DCA1D91BAE627.
\end{center}
Its full body is retained apart from
the title version. The required
companion audit was performed and
influenced the switch from a global
smallness event to an intensive
quantity with an explicit interface
debit. Only active V24 is replaced;
all family archives and companion
files remain unchanged.
\end{record}

\begin{firewall}[Position after V25]
Proved for the actual smooth compact
fibre: a globally feasible arbitrary-exterior
repair, its interface action price, a
matched normalized distance-density tail,
suppression of defect-dense large cubes
at fixed sufficiently large coupling
parameter, and the corresponding
conditional block-cluster estimate.
Standard compactness supplies auxiliary
fixed-spacing Gibbs states, not uniqueness.
Still unproved: sparse-defect mixing,
unrestricted retained coarse-field
control, scale iteration, continuum
construction, and the full physical
Yang--Mills mass gap.
\end{firewall}

\section{V26: exact sparse activities and an actual single-link obstruction}
\label{sec:f12v26}

V25 does not control thin defect connections.
A natural attempted next step is to decompose
the bad free-coordinate set into connected
components and bound each component's conditional
activity uniformly over good exterior coordinates.
We specify that normalization, then test the
estimate against the actual fibre Hessian.
It fails already for one free link. The failure
is proved below for the nonlinear compact law,
not merely suggested by a Gaussian analogy.
It concerns a proof route, not the mass-gap
conjecture, and points towards a buffered or
boundary-averaged activity.

\subsection{An exact activity decomposition with one normalization}

At a fixed finite lattice and retained field
\(C\), use V18's global free variables.
For each free group coordinate \(i\), let
\(D_i\) be its section-centred exponential
cube of scalar radius \(r\). Let \(V(Y;C)\)
be the full Haar-density potential, including
\(-\log j\), and set
\[
 Z_0=\int_{\prod_iD_i}e^{-V}\,dY,\qquad
 \mu_0=Z_0^{-1}1_{\prod_iD_i}e^{-V}\,dY .
 \tag{30.1}
\]
This is the original all-good restricted
measure, not an auxiliary convex extension.
The finite interaction graph has free group
coordinates as vertices, adjacent when they
occur in a common local term of \(V\).
It has bounded degree and fixed range.
For a nonempty connected set \(A\), let
\(V_A\) sum all terms involving \(A\), and put
\[
 a_A(Y_{A^c})=
 \frac{\displaystyle\int_{\prod_{i\in A}D_i^c}
                         e^{-V_A(Y)}\,dY_A}
      {\displaystyle\int_{\prod_{i\in A}D_i}
                         e^{-V_A(Y)}\,dY_A}.
 \tag{30.2}
\]
The denominator is positive, and \(a_A\)
depends only on exterior interaction
neighbours. Call a family compatible when
its members are disjoint and no interaction
edge joins distinct members.

\begin{proposition}[An averaged positive activity gas]
\label{prop:f12v26-expansion}
The exact identities are
\[
 \frac Z{Z_0}=\mathbb E_{\mu_0}\Psi(Y),\qquad
 \Psi(Y)=\sum_{\mathcal A\ {\rm compatible}}
                                \prod_{A\in\mathcal A}a_A(Y),
 \tag{30.3}
\]
where the empty family contributes one.
The weight of the stratum whose bad
components are exactly \(\mathcal A\),
divided by \(Z_0\), is
\(\mathbb E_{\mu_0}\prod_{A\in\mathcal A}a_A\).
\end{proposition}
\begin{proof}
Partition the full domain by its exact bad
coordinate set \(S\). Its graph components
give one compatible family. No local term
involves two components, since that would
join them in the graph. At fixed good
exterior \(Y_{S^c}\), the all-bad and all-good
integrals over \(S\) consequently factor
over its components, and their ratio is
\(\prod_Aa_A\).
Integrating the exterior against its
actual \(\mu_0\) marginal cancels the good
component denominators in (30.2).
This gives the stratum identity. Sum over
all strata to prove (30.3).
The correlated \(\mu_0\) remains outside
the entire family sum; no independent
boundary laws have replaced it.
\end{proof}

If one could prove \(a_A\le q^{|A|}\)
uniformly on all-good exteriors, positive
extraction would bound a specified
component's probability by \(q^{|A|}\):
the family sum containing \(A\) equals
\(a_A\) times the sum over families
compatible with it, which is a subsum
of \(\Psi\). Divide its expectation by
the same \(\mathbb E_{\mu_0}\Psi\).
We now show that this sufficient
estimate with \(q\to0\) is false for
the present coordinate cubes.

\subsection{A seven-variable slice of the actual quadratic action}

Take \(C_e=I\) and an orthonormal Lie
basis with first vector \(T\). In one
coarse cell, allow only its eight
boundary links in direction \(0\) to
have scalar linearized amplitudes
\(x_vT\), labelled by transverse parity
\(v\in\{0,1\}^3\). The root pivot is
\(v=0=(0,0,0)\). All other free
amplitudes vanish. The fine quadratic
action on these physical links has
Gram matrix
\[
 Q_{vw}=6\,1_{\{v=w\}}
       -1_{\{v,w\ {\rm differ\ in\ exactly\ one\ bit}\}} .
 \tag{30.4}
\]
Each link is in six plaquettes. Two
listed links occur in a common
plaquette exactly when they are
transverse cube neighbours, with
opposite curl signs. This proves
(30.4), for example on fine periodic
lattices with periods at least four.

The linearized coarse constraint is
\(\sum_vx_v=0\), so with the seven
free amplitudes \(z_v\), \(v\ne0\),
\[
 x_v=z_v\ (v\ne0),\qquad
 x_0=-\sum_{v\ne0}z_v,\qquad H_b=E_b^TQE_b ,
 \tag{30.5}
\]
where \(E_b\) is this inclusion.
Other-direction pivots have no linear
response to these boundary variables:
their paths use other link directions.
Thus \(H_b\) is an actual submatrix
of \(E^*d^*dE\), not a different
internal-cell model.

\begin{proposition}[A four-neighbour witness]
\label{prop:f12v26-witness}
For target \(i=(1,1,0)\) and
\[
 J=\{(1,0,1),(1,0,0),(0,1,1),(0,1,0)\},
 \tag{30.6}
\]
one has
\[
                     (H_b)_{ii}=12,\qquad
                     (H_b)_{ij}=6\quad(j\in J).
 \tag{30.7}
\]
Set \(z_j=-3r/4\) for \(j\in J\) and
all other exterior free amplitudes
to zero. The target-dependent
quadratic action, up to an exterior
constant, is
\[
             6x^2-18rx=6(x-3r/2)^2-27r^2/2 .
 \tag{30.8}
\]
Its minimum is at \(3r/2\), outside
\([-r,r]\), with improvement \(3r^2/2\)
over the minimum in that interval.
\end{proposition}
\begin{proof}
For nonroot labels,
\((H_b)_{vw}=Q_{vw}-Q_{v0}-Q_{0w}+Q_{00}\).
The target has Hamming distance two
from the root, so \(Q_{i0}=0\), giving
diagonal \(12\). For neighbours
\((1,0,0),(0,1,0)\), the two entries
\(Q_{ij},Q_{0j}\) both equal \(-1\);
for the other listed neighbours both
equal zero. They cancel in each case,
leaving mixed entry \(6\).
The forcing is
\(4\cdot6\cdot(-3r/4)=-18r\).
Complete the square as in (30.8);
the constrained minimizer is \(x=r\).
\end{proof}

\begin{record}[Executed exact checks]
The new script
\texttt{scripts/verify\_ym\_f12\_v26\_single\_site\_activity.py}
constructs (30.4)--(30.5) in rational
arithmetic and checks all 49 entries
against the existing full 45-coordinate
fibre stencil. All checks passed.
The target is zero-based full free
index \(38\), fine link
\(((1,1,1,0),0)\) in the chosen cell.
Its complete stencil row has diagonal
\(12\) and off-diagonal absolute sum
\(126\); the witness needs only the
four mixed entries above.
The script SHA--256 is
\begin{center}\scriptsize\ttfamily
6DD3DE65CFFC4A7733B5D8465B1F60BDC5685B63923AA1E777D51A3967256E9D.
\end{center}
The matrices and dependency hashes are in
\texttt{artifacts/ym\_f12\_v26\_single\_site\_activity.json}.
This verifies quadratic coefficients,
not the nonlinear conclusion below.
\end{record}

\subsection{Persistence for the nonlinear compact law}

Fix those four exterior free group
variables at \(\exp(-3rT/4)\), all
other exterior free variables at \(I\),
and the complete retained field at \(I\).
Every exterior free variable is strictly
inside its good cube. The target is one
free \emph{group} variable, not one
scalar component. Write it as \(\exp u\).
Let \(f_i(u;r)\) sum its affected Wilson
plaquette actions after solving the
global pivots. Terms independent of
the target cancel from the ratios.

\begin{theorem}[The singleton activity grows exponentially]
\label{thm:f12v26-obstruction}
There are \(r_0,c_0>0\), independent
of ambient volume, such that for
every fixed \(0<r<r_0\) and \(\beta\ge0\),
at this all-good exterior,
\[
                       a_{\{i\}}\ge c_0 e^{\beta r^2/2}.
 \tag{30.9}
\]
Hence
\[
 \mathbb P(Y_i\in D_i\mid Y_{i^c})
       \le(1+c_0e^{\beta r^2/2})^{-1}
                             \longrightarrow0
                    \quad(\beta\to\infty).
 \tag{30.10}
\]
For each fixed small \(r\), the
obstruction persists on a nonempty
open neighbourhood of these exterior
values, with reduced positive constants.
\end{theorem}
\begin{proof}
For small \(r\), all inputs used
below, including solved pivots, lie
in V18's common smooth local chart.
The completion agrees there with
the original local block. The
action's first derivative vanishes
at the flat section and its Hessian
is \(E^*d^*dE\), diagonal in
orthonormal Lie components.
Its quadratic part is therefore
\[
                   q_0(r)+6\|u\|^2-18r\langle u,T\rangle ,
 \tag{30.11}
\]
with \(q_0(r)\) independent of \(u\).
Uniform third derivatives and finite
stencils give
\[
 |f_i(u;r)-q_0(r)-6\|u\|^2
                               +18r\langle u,T\rangle|
                                             \le Kr^3
 \tag{30.12}
\]
on the good cube \([-r,r]^d\) and on
the Euclidean ball of radius \(r/16\)
about \(u_*=(3r/2)T\). The constant
includes \(d=\dim G\), but not ambient
volume.

The quadratic minimum on the good
cube is \(q_0(r)-12r^2\), at first
component \(r\) and other components
zero. On the ball about \(u_*\),
the quadratic maximum is at most
\[
                     q_0(r)-27r^2/2+6r^2/256 .
 \tag{30.13}
\]
The ball is outside the good cube,
since its first component exceeds
\(r\). The gap between these
quadratic bounds is \(189r^2/128\).
Choose \(r_0\) so the two Taylor
errors leave at least \(r^2/2\).
Thus
\[
 \sup_{\|u-u_*\|\le r/16}f_i(u;r)
           \le\inf_{u\in[-r,r]^d}f_i(u;r)-r^2/2 .
 \tag{30.14}
\]

The conditional coordinate density
includes the free Haar factor and
the affected inverse-pivot factors.
Their product has uniform positive
lower and upper bounds on both sets.
Only a bounded number of factors
depend on the target; all others
cancel in (30.2).
The outside ball has volume
\(\omega_d(r/16)^d\), and the good
cube has volume \((2r)^d\).
Use the ball as a lower contribution
to the bad numerator and the entire
cube for the good denominator.
Equation (30.14) and the fixed
volume ratio prove (30.9).
The identity
\(\mathbb P(D_i\mid Y_{i^c})=(1+a_{\{i\}})^{-1}\)
is exact, proving (30.10).

The action and density factors are
continuous in the finitely many
relevant exterior inputs. For
fixed \(r>0\), a slightly weakened
version of (30.14) retains a strict
margin on an open neighbourhood.
The four specified exterior values
also have margin \(r/4\) inside
their good cubes. Shrink the
neighbourhood to preserve both
properties. Thus this cannot be
dismissed as one exceptional
conditional null-set value.
\end{proof}

\begin{assessment}[Which implication has failed]
The actual free-coordinate good cube
is not invariant under single-coordinate
conditional minimization. A positive
full fibre Hessian does not imply that
property; the off-diagonal forcing
matters. The witness lies entirely
inside a common smooth chart. No
topological sector or discontinuity
of the old block is involved.

This does not contradict V21 or V25.
Here all other \emph{free coordinates}
are conditioned on, rather than
integrating a whole regional fibre.
Dependent root pivots move with the
target, so other physical fine links
are not all frozen. Departure from
this free cube is also not identified
without proof with every physical
plaquette or source defect.
\end{assessment}

\subsection{The usable normalized ratio and the next direction}

The exact expansion (30.3) remains
valid. What fails is the proposed
sufficient estimate
\[
 \sup_{Y_{A^c}\ {\rm good}}a_A(Y)
                                \le e^{-c\beta|A|}
 \tag{30.15}
\]
at fixed cube radius, or any bound
tending to zero even for a singleton.
The four exterior values themselves
have a cost under \(\mu_0\). The
supremum omitted that cost; it may
offset the large conditional activity
after integration.
For a specified component \(A\),
the correct probability is still
\[
 \mathbb P(A\ {\rm is\ a\ bad\ component})
 =\frac{\mathbb E_{\mu_0}
          [a_A\,\Psi_{\mathrm{compatible\ with}\ A}]}
        {\mathbb E_{\mu_0}\Psi}.
 \tag{30.16}
\]
Bounding \(a_A\) under a different,
unweighted boundary law does not
establish this ratio. Replacing
its denominator by one may also
lose the volume cancellation of
the other components.

\begin{decision}[Integrate a buffer or pay the actual boundary weight]
Do not merely improve constants
in (30.15): the actual singleton
violates it. The next comparison
must integrate some boundary
response before declaring the
core defect, or control the
weighted numerator in (30.16)
against its own denominator.
V22's buffered attenuation of
the local minimum is relevant;
freezing its nearby coordinates
was precisely the obstruction.
An enlarged interaction block
must specify the variables
integrated, its good/bad core
event, and its normalization
before a contour estimate is
claimed. Next companion audit
and direction review: V30.
\end{decision}

\begin{record}[V26 append record]
V25 had SHA--256
\begin{center}\scriptsize\ttfamily
884A25C1B56C6CEBE027451BEF950A47E51BA759DD87875FF98E7040EDFE94B4.
\end{center}
Its complete body is retained
apart from the title version.
The new exact checker imports
only inspected, main-guarded
helpers and writes its own new
artifact. Only active V25 is
replaced; all family archives
and companion notes remain.
\end{record}

\begin{firewall}[Position after V26]
Proved: the exact correlated
positive activity expansion,
a rational four-neighbour
witness in the actual fibre
Hessian, and its nonlinear
exponentially growing singleton
conditional activity on an
open set of all-good exteriors.
This rules out a uniform
single-link activity route,
not sparse-defect suppression
under the fully normalized law.
Buffered activity control,
mixing, scale iteration,
continuum construction and the
full mass gap remain unproved.
\end{firewall}

\section{V27: a normalized buffered core activity}
\label{sec:f12v27}

V26 ruled out uniform suppression after
freezing every neighbouring free coordinate.
We now integrate a full compact buffer
around the core. The core and exterior
use the same small threshold, but the
proof uses a larger fixed local box
for the integrated parent. Both radii
and both integrations are specified below.
The conclusion is a small buffered
activity on good outer data, not a
replacement for the unbuffered activity
inside V26's exact gas.

\subsection{Exact Gaussian check of the proposed change}

In V26's witness, order the five
coordinates as target followed by the
four forcing neighbours. Their actual
quadratic submatrix is
\[
 H_5=\begin{pmatrix}
 12&6&6&6&6\\
 6&12&6&6&7\\
 6&6&14&7&8\\
 6&6&7&12&6\\
 6&7&8&6&14
 \end{pmatrix}
 =\begin{pmatrix}12&b^T\\b&A_4\end{pmatrix}.
 \tag{31.1}
\]
Integrating its four neighbour variables
over their entire Gaussian spaces, rather
than fixing them at \(-3r/4\), gives
\[
 12-b^TA_4^{-1}b=\frac{1716}{227}>0,\qquad
                    (H_5^{-1})_{11}=\frac{227}{1716}.
 \tag{31.2}
\]
This follows by completing the square
in the four variables; their Gaussian
integral contributes an \(x\)-independent
factor and leaves
\(\exp[-\beta(1716/227)x^2/2]\).
The marginal target mean is zero.
If all six other boundary-family free
variables are integrated instead, the
seven-variable covariance entry is
\[
                         (H_b^{-1})_{ii}=\frac{359}{2520}.
 \tag{31.3}
\]

\begin{record}[Executed Schur verification]
The rational calculation in
\texttt{scripts/verify\_ym\_f12\_v27\_buffer\_schur.py}
checks the positive elimination pivots,
both inverse matrices, and (31.2)--(31.3).
All checks passed. Its SHA--256 is
\begin{center}\scriptsize\ttfamily
6253FC2D19A02313DAA112AEF919BB5FB3EF6A2F55FA101F4D78B615EAF6CA55.
\end{center}
The matrix and exact fractions are in
\texttt{artifacts/ym\_f12\_v27\_buffer\_schur.json}.
This calculation explains how integration
removes the particular frozen-neighbour
shift; it is not the nonlinear proof
or a certificate for arbitrary boundary
values. The nonlinear statement uses
a sufficiently wide buffer next.
\end{record}

\subsection{Two radii and a fixed translated template}

Fix a nonempty finite set \(\mathcal K\)
of core free \emph{group} coordinates,
with \(M=|\mathcal K|\). Choose a connected
parent region \(P\) whose varied free
set \(I_P\) contains \(\mathcal K\).
Write \(s=|P|\). Require all scalar
coordinates of the core to lie at
distance at least \(L\) from the frozen
physical collar \(\mathcal F_P\).
The parent and core shapes are fixed
in the theorem; their translations and
the ambient volume may vary.

Write \(R_{\rm box}>0\) for the fixed
interior box radius used in V21--V22,
and keep its constants \(h,\kappa,D_1\).
The defect threshold \(r\) used here
will satisfy \(0<r\le r_0<R_{\rm box}\).
The proof-only large box is
\([-R_{\rm box},R_{\rm box}]^{N_P}\).
The actual integrations defining the
activity will not be restricted to it.

For each core group coordinate, let
\(D_i(r)\) be the section-centred
exponential cube of scalar radius \(r\).
Let \(\mathcal O_{P,r}(C)\) be the
outer data condition that every exterior
free input in \(J_P^\partial\), the
completed input set of V24, belongs
to its analogous cube \(D_j(r)\).
Other remote exterior variables are
unrestricted. Put
\(q_\partial=|J_P^\partial|\).
The physical comparison (28.3) yields
\[
 \eta_P\le C_\infty\sqrt d\,r,\qquad
 \mathcal E_P\le C_E d q_\partial r^2,\qquad d=\dim G .
 \tag{31.4}
\]
These conditions are on the actual
global free variables at fixed \(C\).
All exterior pivots are reconstructed,
not independently held at incompatible
physical values.

\begin{proposition}[A nonempty uniform choice of parameters]
\label{prop:f12v27-parameters}
The buffer distance \(L\) can be chosen
once, independently of \(r,\beta\).
For a fixed template with this buffer,
there are \(c_\epsilon,r_0,B>0\) such
that whenever
\[
 0<r\le r_0,\qquad
 \epsilon_P\le c_\epsilon r,\qquad
 y\in\mathcal O_{P,r}(C),\qquad
                         \beta r^2\ge B ,
 \tag{31.5}
\]
the V21--V22 parent hypotheses hold.
Furthermore its boxed law \(\nu_P\)
satisfies, for every core scalar
coordinate \(u_{i,a}\),
\[
                 |\mathbb E_{\nu_P}u_{i,a}|\le r/2 ,
 \tag{31.6}
\]
and its full compact complement obeys
\[
 \pi_{P,C}^{y}(\mathcal D_{P,C}^c)
                      \le e^{-\kappa\beta/2}.
 \tag{31.7}
\]
All constants are independent of
ambient volume and retained link
amplitudes. They may depend on
the fixed template.
\end{proposition}
\begin{proof}
Let \(K,a>0\) be the constants in
V22's mean bound (26.8). First take
\[
 K C_\infty\sqrt d\,e^{-aL}\le1/8,\qquad
                           Kc_\epsilon\le1/8 .
 \tag{31.8}
\]
Choose a fixed template with this
buffer. Its geometry now fixes
\(s\) and \(q_\partial\).
Since \(L_P\le C_L\epsilon_P^2s\),
choose \(r_0\) small enough that
\[
 \bigl(2AC_Lc_\epsilon^2s+
                       K_3C_Ed q_\partial\bigr)r_0^2
                                      \le\kappa/2 .
 \tag{31.9}
\]
Also require \(C_\infty\sqrt d\,r_0<\eta_0\),
\(c_\epsilon r_0<\epsilon_0\), and that
all exterior charts in (31.4) are
valid. Thus the parent has
\(\Gamma_P\le\kappa/2\) and the
necessary pointwise smallness.

Take \(B\ge64K^2\). Then
\(K/\sqrt\beta\le r/8\).
Together with (31.8), V22 gives
a core mean at most \(3r/8\), proving
(31.6). Increase \(B\) to ensure
\(\beta\ge\beta_0\) and
\(\beta\ge2s\log D_1/\kappa\) whenever
\(r\le r_0\). This is possible by
taking \(B\) at least \(r_0^2\)
times those fixed thresholds.
V21's bound \(D_1^s e^{-\kappa\beta}\)
then implies (31.7).
No integration variable in \(I_P\)
has been frozen to obtain this
core mean estimate.
\end{proof}

\subsection{Simultaneous bad core coordinates}

For \(A\subset\mathcal K\), \(k=|A|\ge1\),
let
\[
                    \mathcal B_A(r)
                       =\bigcap_{i\in A}\{Y_i\notin D_i(r)\}.
 \tag{31.10}
\]
The event declares every specified
core coordinate bad; it places no
condition on other integrated
coordinates. In particular it does
not freeze a surrounding pattern of
bad or good buffer variables.

\begin{theorem}[A full compact joint core estimate]
\label{thm:f12v27-core}
After decreasing \(r_0\) and increasing
\(B\) by fixed amounts, under (31.5),
\[
          \pi_{P,C}^{y}(\mathcal B_A(r))
                    \le q(\beta,r)^k,\qquad
          q(\beta,r)=(2d+1)e^{-c\beta r^2},\quad c=h/8 .
 \tag{31.11}
\]
The constants can be chosen so
\(Mq(\beta,r)\le1/2\) throughout
the asserted parameter range.
\end{theorem}
\begin{proof}
Inside the large parent chart, each
bad core group has at least one
scalar component of magnitude at
least \(r\). Choose one such
component and its sign for each
of the \(k\) groups. There are at
most \((2d)^k\) choices. For any
fixed choice set
\[
                         X=\sum_{i\in A}\sigma_i u_{i,a_i}.
 \tag{31.12}
\]
On that event \(X\ge kr\), while
\(\mathbb E_{\nu_P}X\le kr/2\)
by (31.6).
Every linear tilt of the boxed
potential has the same Hessian
lower bound \(h\beta I\).
The Poincare argument used in
V22 therefore gives
\[
 \log\mathbb E_{\nu_P}
                 e^{t(X-\mathbb EX)}
                       \le \frac{k t^2}{2h\beta}.
 \tag{31.13}
\]
Here the gradient of \(X\) has
squared Euclidean norm exactly
\(k\), since the groups are distinct.
Exponential Markov optimization
gives probability at most
\(e^{-h\beta kr^2/8}\) for this
sign-component event. Thus
\[
                  \nu_P(\mathcal B_A(r))
                   \le(2d)^k e^{-c\beta r^2k}.
 \tag{31.14}
\]
This is a joint bound from one
linear tilt, not an independence
assumption for the core variables.

For the full compact law add
the large-box complement (31.7).
Further require
\[
                         c r_0^2 M\le\kappa/2 .
 \tag{31.15}
\]
Since \(k\le M\), this complement
is at most \(e^{-c\beta r^2k}\).
The sum of the two contributions
is bounded by
\((2d+1)^k e^{-c\beta r^2k}\)
for \(k\ge1\), proving (31.11).
Finally increase
\(B\ge c^{-1}\log(2M(2d+1))\)
to obtain \(Mq\le1/2\).
The restriction (31.15) is part
of the fixed-template choice;
it is not uniform for an
arbitrarily growing core at
a fixed positive radius \(r\).
\end{proof}

\subsection{The normalized buffered activity}

For the same parent and exterior,
define the two \emph{full compact}
integrals
\[
 \begin{split}
 Z_{P,A}^{\rm bad}(C,y)
   &=\int_{G^{I_P}}1_{\mathcal B_A(r)}
                          e^{-V_P(Y_P;y,C)}\,dY_P,\\
 Z_{P,A}^{\rm good}(C,y)
   &=\int_{G^{I_P}}\prod_{i\in A}1_{D_i(r)}(Y_i)
                          e^{-V_P(Y_P;y,C)}\,dY_P,
 \qquad
 b_{A\mid P}=\frac{Z_{P,A}^{\rm bad}}{Z_{P,A}^{\rm good}} .
 \end{split}
 \tag{31.16}
\]
The buffer \(I_P\setminus A\) is integrated
without a good/bad restriction in both
numerator and denominator. The potential
\(V_P\) is the actual Haar-density
potential of V21, with all its
inverse-pivot factors. The large proof
box is absent from (31.16).

\begin{corollary}[A small buffered activity at good outer data]
\label{cor:f12v27-activity}
Under the same conditions,
\[
 b_{A\mid P}\le\frac{q^k}{1-kq}
                           \le2q^k\le(2q)^k .
 \tag{31.17}
\]
In particular the buffered singleton
activity tends to zero exponentially
in \(\beta r^2\), uniformly over
\(\mathcal O_{P,r}(C)\).
\end{corollary}
\begin{proof}
Divide both integrals in (31.16)
by the same full \(Z_P(C,y)\).
The numerator is at most \(q^k\)
by (31.11). A union bound using
its singleton case gives the
denominator probability at least
\(1-kq\). This is at least \(1/2\)
because \(k\le M\).
Taking their ratio proves (31.17).
For \(k>1\), the all-bad and
all-good events are not complementary:
mixed core patterns exist. The
argument explicitly uses their
separate probabilities and does
not misidentify their sum with one.
\end{proof}

This estimate addresses V26's specific
obstruction by changing what is
integrated. Its four forcing
neighbours can be included in
\(I_P\setminus A\); they are then
integrated in both parts of
(31.16), not fixed at their
adversarial values. A larger fixed
buffer makes the estimate uniform
over the allowed outer data.
The Gaussian calculation (31.2)
checks the basic effect, while
(31.11)--(31.17) provide the
nonlinear compact statement.

\subsection{What can be combined without changing the law}

For several parents whose full
interaction stencils are separated,
suppose that each other's core event
and all their outer guards are
measurable outside a chosen parent's
updated variables. Under a global
fixed-\(C\) fibre law, or its DLR
version from V25, successive
conditioning then gives
\[
 \mathbb P\left(\bigcap_j
          [\mathcal B_{A_j}(r)\cap\mathcal O_{P_j,r}(C)]\right)
 \le q^{\sum_j|A_j|}
                    \mathbb P\left(\bigcap_j\mathcal O_{P_j,r}(C)\right).
 \tag{31.18}
\]
At each step pull the other
indicators outside the conditional
integral and apply (31.11) on
the current guard. Repeating proves
the formula. It does not estimate
failure of the guards.

\begin{assessment}[Why this cannot be substituted into V26's gas]
The quantity \(b_{A\mid P}\) is not
\(a_A\) in (30.2).
The latter freezes \(Y_{A^c}\) and
requires an exact bad-coordinate
stratum. The former integrates a
buffer which can itself contain
additional bad coordinates.
Different parent buffers may
overlap, so their integrations
do not factor over the component
families in (30.3).

In particular (31.17) cannot be
inserted for \(a_A\) in that
partition sum without a new
decomposition and normalization
proof. Conditioning the buffer
on a previously revealed adverse
pattern would also remove the
integration which made the
estimate true, and can recover
V26's obstruction. These are
structural restrictions on the
proved estimate, not negligible
boundary errors.
\end{assessment}

\begin{decision}[Next step: a compatible buffered decomposition]
A full compact core activity is now
small with an explicit good outer
guard, and its simultaneous core
bound has a cost proportional to
the number of specified bad core
coordinates. The next task is to
construct a decomposition which
uses such integrated buffers
without overlapping integrations
or suppressing their bad contents.
It must also account for failed
outer guards; V25's density control
does not make every thin guard good.
The exact normalization in (31.16)
and the failed singleton bound in
V26 must both be respected.
Next companion audit and direction
review: V30.
\end{decision}

\begin{record}[V27 append record]
V26 had SHA--256
\begin{center}\scriptsize\ttfamily
DDDEF737F1D947448A6D6F3599E434D5AA8EF60A638BBB557D1766346DD5D062.
\end{center}
Its complete body is preserved
apart from the title version.
The proof reuses the verified
regional Hessian and mean estimates,
but adds a simultaneous-core tilt,
an explicit two-radius parameter
choice and the normalized compact
buffer activity. The exact Schur
checker writes only its own new
artifact. Only active V26 is
replaced; archives and companions
are untouched.
\end{record}

\begin{firewall}[Position after V27]
Proved: removal of the frozen
Gaussian witness by exact neighbour
integration, and a nonlinear full
compact buffered-core activity
bound under explicit small-curvature
and good outer-data hypotheses.
Not proved: a globally compatible
buffered expansion with controlled
failed guards, sparse-defect mixing,
scale iteration, continuum construction
or the physical Yang--Mills mass gap.
\end{firewall}

\section{V28: compatible compact block integration with an explicit interface}
\label{sec:f12v28}

V27's buffered activities cannot be multiplied
when their integrated variables overlap.
Here we make a definite compatible choice:
integrate a separated family of fixed-size
parent regions and leave the remaining
interface variables explicit. Every chosen
block is integrated over its full compact
domain. We prove that the exact interface
potential is uniformly strongly convex on
a specified small interface box, even when
the number of integrated blocks is arbitrarily
large. The interface restriction is retained,
not interpreted as a conclusion about the
unrestricted compact law.

\subsection{A separated family and its exact marginal}

On a finite periodic coarse lattice choose
translates \(P_b\), \(b\in\mathcal I\), of
one fixed connected parent template. Write
\(I_b=I_{P_b}\) for its updated free variables.
Separate the templates by corridors of
fixed sufficient width so that:
\begin{enumerate}
\item the sets \(I_b\) are disjoint;
\item no local term of V18's global
Haar-density potential involves variables
in two different \(I_b\);
\item the affected plaquette and pivot
factor sets of different parents are
disjoint, and every exterior input
\(J_{P_b}^{\partial}\) lies outside
\(\bigcup_c I_c\).
\end{enumerate}
These conditions can be arranged by
periodic packing of fixed cubes, using
the finite stencil radius from V18.
The period of the ambient torus is
chosen compatible with this packing.
Let
\[
             \mathcal S=\{\hbox{all free variables}\}
                                      \setminus\bigcup_b I_b .
 \tag{32.1}
\]
We call \(\mathcal S\) the interface; it
includes the corridors and is not
integrated in the next formula.

Use \(y=Y_{\mathcal S}\), \(X_b=Y_{I_b}\).
Assign each affected plaquette and pivot
factor to its unique block. The remaining
terms define \(V_{\mathcal S}\). Thus
the full Haar-density potential has
the exact form
\[
 V(Y;C)=V_{\mathcal S}(y;C)+\sum_b U_b(X_b;y,C),
 \tag{32.2}
\]
where \(U_b\) is the regional potential
of (25.16). It may include terms
constant in \(X_b\); their assignment
does not affect the following identities.
Set
\[
 Z_b^{\rm c}(y)=\int_{G^{I_b}}e^{-U_b(X_b;y,C)}\,dX_b,
 \qquad W_b^{\rm c}=-\log Z_b^{\rm c}.
 \tag{32.3}
\]
Product Haar and the absence of
cross-block terms give the exact
interface potential relative to
interface Haar:
\[
             W_{\mathcal S}^{\rm c}
                         =V_{\mathcal S}+\sum_b W_b^{\rm c}.
 \tag{32.4}
\]
Conditionally on \(y\), the full compact
block laws are independent and are
precisely \(\pi_{P_b,C}^{y}\).
Equation (32.4) is a marginalization
identity, not an activity approximation
or a resampling prescription.

Each \(W_b^{\rm c}\) is smooth and
depends only on the fixed finite
input set \(J_b=J_{P_b}^{\partial}\).
This follows from V18's smooth
positive compact density and V24's
input-closure statement. The new
interface interaction therefore has
finite range, depending on the fixed
template size but not on the number
of templates or ambient volume.

\subsection{The joint local Hessian needs pointwise, not total, curvature}

Fix a sufficiently small common radius
\(R_*>0\) for the free exponential
coordinates about their section values.
Inside every block use its cube
\(\mathcal D_b=[-R_*,R_*]^{d|I_b|}\).
For interface variables use the same
kind of coordinates \(v\), eventually
restricted to a smaller radius
\(0<r_{\mathcal S}<R_*\).
The complete retained field \(C\)
is fixed, with
\[
                \delta(Q_p(C))\le\epsilon_*
                          \quad\hbox{for all coarse plaquettes}.
 \tag{32.5}
\]
There is no bound on the total sum
of their curvatures.

\begin{proposition}[Uniform joint small-field convexity]
\label{prop:f12v28-joint}
For sufficiently small fixed
\(R_*,\epsilon_*>0\) and sufficiently
large \(\beta\), the full measured
coordinate potential
\[
 \mathcal V(u;C)=\beta S_0(W(u,C))
       -\log j(u,C)-\sum_{\rm free\ links}\log j_G(u_l)
 \tag{32.6}
\]
obeys
\[
                   m\beta I\le\mathcal V''(u;C)\le M\beta I
 \tag{32.7}
\]
on the global product cube of radius
\(R_*\). Its Hessian has fixed range
and row bound \(K\beta\), independently
of total volume.
\end{proposition}
\begin{proof}
This is the pointwise-curvature Hessian
argument of V15, applied to the whole
periodic fibre. To recall why no total
load enters, the section-flat coercivity
gives \(S_{\rm v}(u,C)\ge D_{\rm free}(u,C)/A\)
and \(S_{\rm v}(0,C)=0\).
Hence its Hessian at zero is at least
\(2A^{-1}I\).
The corner second derivatives have a
nonnegative squared-first-derivative
term and a remainder whose negative
part is at most \(K\epsilon_*\|h\|_2^2\),
using bounded local overlap as in
(19.15). Thus the action Hessian at
zero has a uniform positive lower
bound when \(\epsilon_*\) is small.

The pivot map agrees with the
original local map throughout a
sufficiently small common cube.
Its local third derivatives bound
the Hessian row-norm variation by
\(K\|u\|_\infty\). Shrink \(R_*\)
to preserve the gap. The logarithms
of both inverse-pivot and free
Haar factors have uniformly bounded
Hessian row norm. For large \(\beta\)
these are bounded perturbations
of the action Hessian. Locality
gives the upper and row bounds
in (32.7). Every estimate uses
pointwise curvature and bounded
stencil overlap, not
\(\sum_p\delta(Q_p)^2\).
\end{proof}

\subsection{A local compact-complement estimate for each block}

Choose \(r_{\mathcal S}\) and, if needed,
decrease \(\epsilon_*\) further so that
every \(v\in[-r_{\mathcal S},r_{\mathcal S}]^{d|\mathcal S|}\)
gives a protected exterior for every
parent. This is a uniform finite-template
choice: if \(s_*=\max_b|P_b|\) and
\(q_*=\max_b|J_b|\), then V24 gives
\[
 \eta_{P_b}\le C_\infty\sqrt d\,r_{\mathcal S},\qquad
 \mathcal E_{P_b}\le C_E d q_*r_{\mathcal S}^2,\qquad
 L_{P_b}\le C_L\epsilon_*^2s_* .
 \tag{32.8}
\]
Make
\(2AC_L\epsilon_*^2s_*+K_3C_Edq_*r_{\mathcal S}^2\)
smaller than the V21 allowable debit
for the chosen radius \(R_*\).
All quantities on this choice's left
side are independent of the number
of blocks. Rechoose the constants
of V21's complement proof for this
common radius if it has been reduced.

Let \(Z_b^{\rm box}\) be (32.3) with
the block restricted to \(\mathcal D_b\),
including its exponential Haar factors,
and \(W_b^{\rm box}=-\log Z_b^{\rm box}\).
V21 then supplies fixed \(K_0,\kappa_*>0\)
with
\[
 p_b(v):=\pi_{P_b,C}^{y(v)}(\mathcal D_b^c)
          \le K_0e^{-\kappa_*\beta}
                         \quad\hbox{for every }b,v .
 \tag{32.9}
\]
The factor \(D_1^{s_*}\) has been
included in \(K_0\); the template
size is fixed. This is not a
bound on the probability that
all blocks simultaneously lie
in their boxes.

\begin{proposition}[The full-block correction has small local derivatives]
\label{prop:f12v28-correction}
For each block define
\[
 R_b(v)=W_b^{\rm c}(v)-W_b^{\rm box}(v)
                              =\log(1-p_b(v)).
 \tag{32.10}
\]
For sufficiently large \(\beta\),
\[
 |R_b|\le Ke^{-\kappa_*\beta},\quad
 |\partial_aR_b|\le K(\beta+1)e^{-\kappa_*\beta},\quad
 |\partial_a\partial_cR_b|
                     \le K(\beta+1)^2e^{-\kappa_*\beta}.
 \tag{32.11}
\]
Derivatives vanish outside \(J_b\).
The constants may depend on the
fixed template, but not on volume
or compact retained holonomies.
\end{proposition}
\begin{proof}
The identity (32.10) follows from
\(Z_b^{\rm box}/Z_b^{\rm c}=1-p_b\).
For \(p_b\le1/2\), it gives the
value bound. For derivatives, use
the exact score identities for
the full and boxed laws:
\[
 (W_b)_a=\mathbb E(U_b)_a,\qquad
 (W_b)_{ac}=\mathbb E(U_b)_{ac}
                            -\operatorname{Cov}((U_b)_a,(U_b)_c).
 \tag{32.12}
\]
The block domain is independent
of \(v\), since \(C\) is held
fixed. The interior Haar factor
is also independent of \(v\).
All relevant scores have suprema
at most \(K(\beta+1)\) on the
full compact block, by the
global smooth inverse and its
finite stencils.

The total variation distance
between the full law and its
conditioning on \(\mathcal D_b\)
is \(p_b\). A mean changes by
at most twice its supremum times
\(p_b\); a covariance changes by
at most six times the product
of its suprema times \(p_b\).
Apply these bounds to (32.12)
and use (32.9). This estimates
the derivatives without assuming
that the small value of a
function bounds its derivatives.
The finite exterior input set
gives the support assertion.
\end{proof}

\subsection{The exact interface Hessian stays positive}

Let \(\Phi^{\rm c}(v)\) and
\(\Phi^{\rm box}(v)\) denote the
interface potentials relative to
Lebesgue measure in \(v\), with
full compact blocks and boxed
blocks respectively. Thus
\[
 \begin{split}
 \Phi^{\rm c}
   &=V_{\mathcal S}(y(v))
      -\sum_{j\in\mathcal S}\log j_G(v_j)+\sum_b W_b^{\rm c}(v),\\
 \Phi^{\rm box}
   &=V_{\mathcal S}(y(v))
      -\sum_{j\in\mathcal S}\log j_G(v_j)+\sum_b W_b^{\rm box}(v).
 \end{split}
 \tag{32.13}
\]
The interface Haar term is common
and must not be added a second
time to an individual block score.

\begin{proposition}[Strong convexity of the boxed interface marginal]
\label{prop:f12v28-boxed-marginal}
On the interface box,
\[
                        m\beta I\le(\Phi^{\rm box})''\le M\beta I .
 \tag{32.14}
\]
\end{proposition}
\begin{proof}
By (32.7), subtracting
\(m\beta\|v\|_2^2/2\) from
\(\mathcal V(x,v)\) leaves a
jointly convex function of the
integrated block coordinates
\(x\) and interface coordinates
\(v\). Extend its exponential
by zero outside their convex
product boxes. Its integral
over \(x\) is log-concave.
It equals
\(\exp[-\Phi^{\rm box}(v)+m\beta\|v\|_2^2/2]\),
so the lower bound follows.

The marginalization principle is
A.~Pr\'ekopa, \emph{On logarithmic
concave measures and functions},
Acta Sci. Math. (Szeged)
\textbf{34} (1973), 335--343,
\href{https://acta.bibl.u-szeged.hu/14411/1/math_034_335-343.pdf}
{Theorem 6, p.~342}; it applies
to nonnegative log-concave
functions with convex support.
Here all integration domains
are fixed and finite.

For the upper bound, differentiating
the same fixed-domain marginal gives
\((\Phi^{\rm box})''=
\mathbb E\mathcal V_{vv}
-\operatorname{Cov}(\mathcal V_v,\mathcal V_v)\).
The covariance matrix is positive
semidefinite, while
\(\mathcal V_{vv}\le M\beta I\).
Smoothness of the positive finite
integral justifies both derivative
identities in the interior of
the interface box, and the
bounds extend to its boundary.
\end{proof}

\begin{theorem}[An arbitrary number of blocks can be integrated compactly]
\label{thm:f12v28-compact-marginal}
There is \(\beta_*<\infty\), independent
of the number of packed blocks and
ambient volume, such that
\[
 \|(\Phi^{\rm c}-\Phi^{\rm box})''\|_{\rm row}
                  \le K(\beta+1)^2e^{-\kappa_*\beta},
 \tag{32.15}
\]
and for \(\beta\ge\beta_*\),
\[
             (m/2)\beta I\le(\Phi^{\rm c})''\le K\beta I .
 \tag{32.16}
\]
The exact full-block interface
Hessian has a fixed finite range
and row bound \(K\beta\).
\end{theorem}
\begin{proof}
The difference in (32.13) is
\(\sum_bR_b\). A fixed interface
coordinate belongs to only a
bounded number of sets \(J_b\),
and each \(J_b\) has bounded
cardinality, since the packing
and templates are fixed.
Summing (32.11) first over a
Hessian row and then over just
these incident blocks proves
(32.15), without a factor
\(|\mathcal I|\).
Symmetry bounds the operator
norm by the row norm.
For large \(\beta\) it is
at most \(m\beta/2\); apply
(32.14) to obtain (32.16).

Each block free energy depends
only on \(J_b\), and all remaining
terms are local, so the Hessians
in (32.13) have finite range
bounded by the fixed template
diameter and stencil radius.
Their operator bounds and this
bounded row support give row
bound \(K\beta\).
\end{proof}

The value \(\sum_bR_b\) may have
size \(K|\mathcal I|e^{-\kappa_*\beta}\).
The argument requires no condition
\(|\mathcal I|e^{-\kappa_*\beta}\ll1\).
It controls a local derivative
norm, not total variation from
an event that every integrated
block is good. All compact
block excursions have been
included in (32.4).

\subsection{A spatial consequence with the remaining restriction explicit}

Let \(\lambda_{\mathcal S}^{\rm c}\)
have density proportional to
\(e^{-\Phi^{\rm c}(v)}\) on
\([-r_{\mathcal S},r_{\mathcal S}]^{d|\mathcal S|}\).
Equivalently, it is the exact
interface marginal of the full
finite fibre law conditioned
only on this interface box.

\begin{corollary}[Uniform covariance on the constrained interface]
\label{cor:f12v28-interface-covariance}
For \(C^1\) observables \(F,G\) of
interface variables with local
gradient supports \(A_F,A_G\),
\[
 |\operatorname{Cov}_{\lambda_{\mathcal S}^{\rm c}}(F,G)|
 \le K\beta^{-1}e^{-\gamma{\rm dist}(A_F,A_G)}
                     \|\nabla F\|_{\infty,2}
                     \|\nabla G\|_{\infty,2}.
 \tag{32.17}
\]
Constants are independent of
the number of blocks.
\end{corollary}
\begin{proof}
Apply the boxed vector-resolvent
covariance theorem of f11 V86
to (32.16), its finite range,
and its row bound. The interface
coordinate set may have holes;
its locations still have bounded
multiplicity and the same
polynomial ball-growth bound.
These are the geometric facts
used by the weighted resolvent
argument. Empty gradient supports
give constants and zero covariance.
\end{proof}

\begin{corollary}[Correlations of compactly integrated block observables]
\label{cor:f12v28-block-covariance}
Let \(F(X_b)\) and \(G(X_c)\), \(b\ne c\),
be bounded exterior-independent
observables, with \(C^1\) restrictions
to their large block boxes and
gradient bounds \(L_F,L_G\).
Put \(M_F=\|F\|_\infty\), \(M_G=\|G\|_\infty\).
Under the full block law with
only the interface box imposed,
\[
 |\operatorname{Cov}(F,G)|
 \le K\beta^{-1}e^{-\gamma'{\rm dist}(P_b,P_c)}
                            (L_F+M_F)(L_G+M_G).
 \tag{32.18}
\]
\end{corollary}
\begin{proof}
Given \(v\), the two blocks are
independent, so the covariance
equals that of their means
\(m_F(v),m_G(v)\) under
\(\lambda_{\mathcal S}^{\rm c}\).
These depend only on \(J_b,J_c\).
Increase \(\beta_*\) so the fixed
parent size also satisfies
V21's \(s_*\le c_1\beta\).
The protected data in (32.8)
permit the exact score identity
\(\partial_a m_F=-\operatorname{Cov}(F,(U_b)_a)\).
V21's regional full-compact
covariance estimate, together
with the score suprema and
boxed gradient bounds \(K(\beta+1)\),
gives
\(\|\nabla m_F\|_{\infty,2}
\le K(L_F+M_F)\), and similarly
for \(G\). Apply (32.17).
Distances from the finite input
sets to their parent blocks
differ by only a fixed template
constant, which is absorbed
into \(K,\gamma'\).
Physical observables with
explicit interface dependence
would need the corresponding
direct derivative term; they
are not included silently.
\end{proof}

\begin{assessment}[The remaining variables have not disappeared]
The factorization solves one
compatibility problem: no variable
is integrated in two buffers.
It yields an exact compact
interface action, and on its
declared small box a uniformly
positive Hessian and spatial
covariance bounds at fixed
sufficiently large \(\beta\).
The number of integrated blocks
can grow without the cofinal
volume restriction used for a
single whole-fibre good event.

However the remaining interface
variables still obey a hard
small-field restriction.
The unrestricted interface law
is the compact density (32.4)
on all of \(G^{\mathcal S}\);
its probability of leaving
this box has not been bounded
in a way that proves mixing.
The constrained law is not
the unrestricted fibre Gibbs
state of V25. Nor does the
present calculation construct
the retained \(C\)-field law
or a continuum Hamiltonian.
\end{assessment}

\begin{decision}[Next step: the unrestricted effective interface law]
The exact interface action (32.4)
is now the object to study.
Any second integration or shifted
packing must act on this action,
including its \(-\log Z_b^{\rm c}\)
terms, rather than reuse the
original Wilson action as if
the first integration had not
occurred. Determine its global
coercivity or normalized defect
cost before removing the interface
restriction. The local Hessian
estimate (32.15) is available;
it is not by itself a global
compact control theorem.
Next companion audit and
direction review: V30.
\end{decision}

\begin{record}[V28 append record]
V27 had SHA--256
\begin{center}\scriptsize\ttfamily
D97D8334998C9DD77983775E805F525E17E178AA4344D4089CDF5E16688AD44C.
\end{center}
Its complete body is retained
apart from the title version.
The pointwise-curvature Hessian
proof in V15 was checked, and
the marginal log-concavity input
was verified in the primary
source cited above.
Only active V27 is replaced;
family archives, verification
scripts and companion notes
remain unchanged.
\end{record}

\begin{firewall}[Position after V28]
Proved: compatible exact compact
integration of a separated block
family, a local \(C^2\) correction
bound independent of the number
of blocks, and spatial covariance
under the remaining interface
small-field restriction.
Not proved: removal of that
restriction, unrestricted
fixed-coupling mixing, scale
iteration, continuum construction,
or the full physical mass gap.
\end{firewall}

\section{V29: global and conditional coercivity of the exact interface action}
\label{sec:f12v29}

We now study V28's actual unrestricted
interface law. The terms
\(-\log Z_b^{\rm c}\) are retained.
The result is a global lower bound
for its effective potential, and an
exact lift of certain interface
conditionals to the original regional
kernel. The latter transfers V25's
density estimate with the Gaussian
normalization powers still matched.
Neither assertion is obtained by
pretending that the effective action
is the original Wilson action.

\subsection{Interface patches which respect the packing}

Keep the separated fixed-template
packing from V28, with interface
\(\mathcal S\). Call a connected
coarse region \(R\) packing-aligned
when its boundary lies sufficiently
far inside the corridors that:
\begin{enumerate}
\item each integrated set \(I_b\)
is either contained in \(I_R\)
or disjoint from it;
\item for an inside block its
exterior input set \(J_b\) is
contained in \(I_R\cap\mathcal S\);
for an outside block
\(J_b\cap I_R=\varnothing\);
\item the free input closure of
the frozen physical collar
\(\mathcal F_R\) lies entirely
in \(\mathcal S\setminus I_R\).
\end{enumerate}
These are geometric conditions,
not smallness conditions on the
values of any variables.
They hold for large macro-cubes
whose faces cut the middle of
the packing corridors, provided
the corridor width exceeds the
fixed closure and interaction
stencil widths on both sides.
The whole periodic lattice is
also an aligned region.

Put
\[
 T=I_R\cap\mathcal S,\qquad
 X_R=\bigcup_{I_b\subset I_R}I_b,\qquad
 I_R=T\mathbin{\dot\cup}X_R,\quad s=|R|.
 \tag{33.1}
\]
The dot records a disjoint union.
All cardinalities of free-variable
sets count group, not scalar,
coordinates.

\begin{proposition}[An exact conditional lift]
\label{prop:f12v29-lift}
Condition the unrestricted interface
law (32.4) on
\(y_{\mathcal S\setminus T}\).
Its law on \(T\) is exactly the
\(T\)-marginal of the canonical
regional kernel \(\pi_{R,C}^{\xi}\),
where the needed frozen physical
collar is reconstructed from these
exterior interface inputs.
In particular, up to an exterior-only
factor, its unnormalized density is
\[
 K_R(y_T\mid y_{\mathcal S\setminus T})
   =\int_{G^{X_R}}
           e^{-\beta H_\xi(X_R,y_T)}
             \prod_{e\in E(R)}j_e(X_R,y_T;\xi,C)\,dX_R .
 \tag{33.2}
\]
\end{proposition}
\begin{proof}
In (32.4), an outside block's
partition function is independent
of \(y_T\), by the second alignment
condition, and cancels on conditioning.
Inside blocks contain precisely
the variables \(X_R\); their original
integrands can be restored using
(32.2)--(32.3). No variable is
restored twice.
The remaining original terms are
those of the regional conditional
on \(I_R\), up to terms independent
of \(I_R\), which also cancel.

By the third alignment condition,
the exterior physical links needed
for that regional kernel depend
only on the fixed exterior interface
variables. Values in already
integrated outside blocks have
no effect on this kernel. They
can be assigned arbitrarily in
this verification.
V21's exact conditional density
then gives (33.2). Integrating
it over \(y_T\) gives the same
regional normalizer \(Z_R(C,\xi)\).
Thus this is an equality of
normalized laws, not a comparison
with a separately imposed boundary
measure.
\end{proof}

For \(j\in T\), let \(J_j(C)\)
be its free section value, and
define
\[
 D_T(y)=\sum_{j\in T}
                 \|y_j-J_j(C)\|_{\rm HS}^2 .
 \tag{33.3}
\]
Similarly \(D_X\) denotes the
free-link distance in \(X_R\).
Together they are bounded by
V25's physical varied-link
distance \(D\), which also
contains root pivots.

\subsection{A global lower bound for the effective conditional potential}

Let \(y_T^0=J_T(C)\) and
\(W_T=-\log K_R\), relative to
group Haar on \(T\). Its additive
exterior-only constant is immaterial.
Use V25's \(L_R,b_R,C_\partial\).

\begin{theorem}[Effective coercivity with its entropy debit retained]
\label{thm:f12v29-coercivity}
There is \(D_3\ge1\), independent
of region, volume and compact
coarse or exterior data, such
that for sufficiently large
\(\beta\) and every \(y_T\),
\[
 \begin{split}
 W_T(y_T)-W_T(y_T^0)
 \ge {}&\frac{\beta}{2A}D_T(y_T)
       -\frac{\beta A}{2}L_R(C)\\
      &-\beta C_\partial b_R-s\log D_3 .
 \end{split}
 \tag{33.4}
\]
No interface or exterior small-field
condition is imposed.
\end{theorem}
\begin{proof}
The global relative-action estimate
(29.6) gives, pointwise in the
integrated \(X_R\),
\[
 H_\xi-H_0\ge
 \frac{D_T+D_X}{2A}
        -E_{\rm v,out}-\frac A2L_{\rm in}-M_R .
 \tag{33.5}
\]
Bound the inverse-pivot factors
from above and integrate the
compact Gaussian factor in \(X_R\).
Writing \(N_X=d|X_R|\), one obtains
\[
 K_R(y_T)\le B^s\beta^{-N_X/2}e^{-\beta H_0}
  e^{-\beta(D_T/(2A)-E_{\rm v,out}-(A/2)L_{\rm in}-M_R)} .
 \tag{33.6}
\]
Only the \(X_R\) variables are
integrated here; no Gaussian
power for \(T\) is included.

At \(y_T=y_T^0\), set the free
variables \(X_R\) to their section
values as well. V25's global
repair is feasible for arbitrary
exterior and has
\(H_{\rm rep}\le H_0+6s_{\max}b_R\).
Apply the symmetric small-cube
lower-integral argument of (29.5)
only in the \(X_R\) directions.
Their number is at most \(73ds\),
their Hessian upper bound is
uniform, and their gradient need
not vanish. This gives
\[
 K_R(y_T^0)\ge b^s\beta^{-N_X/2}
                             e^{-\beta H_0-6\beta s_{\max}b_R}.
 \tag{33.7}
\]
It remains true for \(X_R=\varnothing\),
with Gaussian power zero.
Divide (33.6) by (33.7).
The powers and the reference action
cancel. Equations (29.7)--(29.11)
bound the remaining debit by
\((A/2)L_R+C_\partial b_R\).
Taking minus logarithms and
\(D_3=\max\{1,B/b\}\) proves (33.4).
\end{proof}

For the whole periodic lattice,
\(T=\mathcal S\), \(s=n\), and
\(b_R=0\). Thus the exact global
interface potential satisfies
\[
 W_{\mathcal S}^{\rm c}(y)
           -W_{\mathcal S}^{\rm c}(J_{\mathcal S}(C))
 \ge\frac{\beta}{2A}D_{\mathcal S}(y)
             -\frac{\beta A}{2}L_{\rm full}(C)-n\log D_3 .
 \tag{33.8}
\]
This concerns the full compact
potential (32.4), not its boxed
Hessian approximation.
The additive entropy debit is
extensive in the original cell
count. It is not a per-defect
error and cannot be deleted
when examining a single defect.

\subsection{A zero-temperature interpretation, with a uniform error}

At a fixed finite volume define
\[
 \mathcal E_C(y)=\min_{X\in G^{\cup_b I_b}}
                                      S_0(W(X,y,C)).
 \tag{33.9}
\]
The minimum exists by compactness.
It includes the section-free
interface action and all integrated
blocks, rather than the sum of
their section values.

\begin{proposition}[The actual effective energy floor]
\label{prop:f12v29-energy}
For \(\beta\ge1\), with \(N_X=d|\cup_b I_b|\),
\[
 \beta\mathcal E_C(y)-Kn
 \le W_{\mathcal S}^{\rm c}(y)
 \le\beta\mathcal E_C(y)+\frac{N_X}{2}\log\beta+Kn .
 \tag{33.10}
\]
These bounds are uniform over compact
\(C,y\). Furthermore
\[
 \mathcal E_C(y)-\mathcal E_C(J_{\mathcal S}(C))
           \ge\frac{D_{\mathcal S}(y)}{2A}
                                  -\frac A2L_{\rm full}(C).
 \tag{33.11}
\]
For flat \(C\), the energy floor
has its unique zero at the
free interface section.
\end{proposition}
\begin{proof}
The compact integral defining the
interface density has amplitude
given by the full inverse-pivot
Jacobian, bounded between constants
to the power \(n\). Its action is
at least \(\mathcal E_C(y)\).
Haar volume is one, proving the
lower bound on \(W_{\mathcal S}^{\rm c}\).

Choose a minimizing \(X_*\). On the
compact product group this is a
stationary point, even if it is
degenerate or nonunique. Product
normal coordinates of radius
\(t/\sqrt\beta\) about \(X_*\)
have Haar density bounded below
uniformly. Global smooth pivot
derivatives and finite overlap
give a uniform action Hessian
upper bound there. The action
is at most
\(\mathcal E_C(y)+K|u|^2\).
Integrating this neighbourhood
gives at least
\(b^n\beta^{-N_X/2}e^{-\beta\mathcal E_C(y)}\).
This proves the upper bound,
without a nondegenerate-saddle
assumption.

For (33.11), the full periodic
relative-action coercivity gives
\(S_0-H_0\ge D_{\mathcal S}/(2A)
-(A/2)L_{\rm full}\), after
discarding integrated distances.
Take the minimum in \(X\), then
use
\(\mathcal E_C(J_{\mathcal S}(C))\le H_0\),
by choosing all integrated free
variables at their section values.
If \(C\) is flat, \(H_0=0\),
\(L_{\rm full}=0\), and the section
realizes zero action. Equation
(33.11) forces a zero-energy
interface to equal that section.
This statement retains whatever
compact holonomy is already in \(C\).
\end{proof}

The fixed-volume limit
\(\beta^{-1}W_{\mathcal S}^{\rm c}\to\mathcal E_C\)
is uniform, but (33.10) has an
extensive error before division
by volume. It is not a uniform
sparse-activity estimate or a
positive physical mass theorem.
The global energy floor need
not be differentiable where
minimizers change, so it is
not substituted for the smooth
finite-\(\beta\) Hessian of V28.

\subsection{Normalized unrestricted interface density bounds}

\begin{theorem}[A density estimate for the effective conditional law]
\label{thm:f12v29-tail}
For every aligned \(R\), arbitrary
compact exterior interface data,
and \(x\ge0\), its exact interface
conditional \(\rho_T\) obeys
\[
 \rho_T\{D_T\ge x\}\le
 \min\left\{1,D_2^s
       e^{-\beta(x/(4A)-(A/2)L_R-C_\partial b_R)}\right\},
 \tag{33.12}
\]
and, for \(|T|>0\),
\[
 \frac{\mathbb E_{\rho_T}D_T}{|T|}
 \le\frac{s}{|T|}\left[
       \frac{4A\log D_2}{\beta}
       +2A^2\frac{L_R}{s}
       +4AC_\partial\frac{b_R}{s}
       +\frac{4A}{\beta s}\right].
 \tag{33.13}
\]
\end{theorem}
\begin{proof}
By Proposition~\ref{prop:f12v29-lift},
the interface conditional is the
\(T\)-marginal of the full regional
kernel. The event \(D_T\ge x\)
implies \(D\ge x\) there. Apply
V25's normalized bound (29.8)
and its moment estimate, using
\(D_T\le D\).
Alternatively the numerator proof
splits \(D/(2A)\) into
\(x/(4A)+D_{\rm free}/(4A)\)
on this event and integrates
both \(X_R\) and \(T\); it then
has the same \(\beta^{-d|I_R|/2}\)
power as the regional denominator.
Thus (33.12) does not follow
from a pointwise potential floor
with its normalization forgotten.
\end{proof}

For the fixed periodic packing,
large aligned macro-cubes have
\[
                    |T|\ge\sigma_{\mathcal S}s,\qquad
                    b_R/s\longrightarrow0
 \tag{33.14}
\]
for a fixed \(\sigma_{\mathcal S}>0\).
Indeed a packing period contains
a fixed positive number of unintegrated
free variables in its corridors;
interior periods contribute that
number exactly, with only a
boundary-order correction. The
geometric ratio \(b_R/s\) has
the same surface-to-volume bound
as in V25. The positive constant
\(\sigma_{\mathcal S}\) may be
small and is not set to one.

Let \(B_\tau(T)\) count interface
groups with section distance at
least \(0<\tau<d_G\). For
\(0<\theta<1\), assume
\[
             \epsilon^2
               \le\frac{\tau^2\theta\sigma_{\mathcal S}}
                            {8A^2C_L}
 \tag{33.15}
\]
on the affected coarse curvatures.
For all sufficiently large aligned
cubes and sufficiently large fixed
\(\beta\),
\[
 \rho_T\{B_\tau(T)\ge\theta|T|\}
          \le e^{-\beta\tau^2\theta|T|/(16A)} .
 \tag{33.16}
\]
To verify this, set \(x=\tau^2\theta|T|\)
in (33.12). The curvature and
surface debits can each be made
at most \(x/(16A)\), using
(33.14)--(33.15). Finally take
\(\beta\ge16A\log D_2/
(\tau^2\theta\sigma_{\mathcal S})\)
to absorb \(D_2^s\).
All these are bounds for the
unrestricted effective interface
conditional, not for the
interface-box law in V28.

\begin{assessment}[What was transferred and what was not]
Exact marginalization preserves a
global quadratic energy floor and
an arbitrary-exterior conditional
density bound for aligned interface
patches. The factor
\(s/|T|\le1/\sigma_{\mathcal S}\)
records the cost of the eliminated
variables; there is no assertion
that the new interaction has
the same constants or form as
the Wilson action.

The probability in (33.16)
concerns a fixed positive density
of bad interface groups. A thin
bad path can still have negligible
density. Likewise the extensive
\(n\log D_3\) error in (33.8)
cannot yield a useful bound for
one sparse defect at fixed \(\beta\)
as volume grows.
Thus these results do not yet
remove V28's interface restriction
from its covariance theorem.
\end{assessment}

\begin{decision}[V30 audit and the remaining local comparison]
At V30 perform the required
companion audit and direction
review. The exact unrestricted
interface action now has both
a proved local Hessian on a
small box and a global intensive
control statement. The missing
link remains a normalized bound
for sparse interface defects or
their effect on conditional means,
not another global-density estimate.
Any further integration must
use (32.4), and an interface
patch touching a block must
restore all that block's integrated
variables before invoking an
original regional estimate.
The aligned lift above gives
one precise implementation,
not permission to ignore cut
blocks in a general patch.
\end{decision}

\begin{record}[V29 append record]
V28 had SHA--256
\begin{center}\scriptsize\ttfamily
EE07694B90E44FD5BEE94B8D60CAB8D90C292FFE6625B374C2A12776F07B143D.
\end{center}
Its complete body is retained
apart from the title version.
The new comparison uses the
actual integrated interface law,
with a geometric conditional
lift and matched integration
powers. Only active V28 is
replaced by V29; archives,
scripts and companions are
unchanged. Next audit and
direction review: V30.
\end{record}

\begin{firewall}[Position after V29]
Proved: exact aligned-patch
conditional lifting, a global
effective coercivity bound with
its entropy debit, a compact
zero-temperature energy floor,
and unrestricted interface
defect-density estimates.
Still unproved: sparse-interface
mixing, removal of the interface
box from V28's correlation theorem,
scale iteration, continuum
construction and the full
Yang--Mills mass gap.
\end{firewall}

\section{V30: audit, global quadratic repair and boundary-weighted defect extraction}
\label{sec:f12v30}

\subsection{Companion audit and ten-version direction review}

The active files were inspected again at this checkpoint.
Their sizes and SHA--256 digests are:
\begin{center}\small
\begin{tabular}{lrr}
\toprule
Companion & Version & Bytes\\
\midrule
SM--Einstein & 72 & 607372\\
NS stability & 26 & 278283\\
Parallel YM & 5 & 54174\\
\bottomrule
\end{tabular}
\end{center}
\begin{center}\scriptsize\ttfamily
SM: F44F9745D43379E1672C5550411B58E97195A4C9278592D221DF5D3F644DC1F5\\
NS: 82C887F8C2C69E4F28FAD7C3DC8DE5B9099CB5A4BF431EBA1FFF3E91935978AD\\
YM: 306F1BB84CFA8A8D47EA569EC17E9545F9867C8D32B548EC9D9C444B4F3C0512
\end{center}
These are unchanged from the V25 audit.
SM V72 proves a specified fixed-mode, fixed-time many-copy
limit, not a cutoff-uniform finite-species continuum limit.
NS V26 records rank-one derivative formulas; its references
to a V27 numerical continuation are not treated as an
available V27 file or as a new YM input.
Parallel YM V5 retains compact labels and diagnoses
the extensive-charge tightness obstruction.
None supplies the sparse-interface estimate needed here.
Document-internal continuation or stop language is source
material, not an instruction changing the present goal.

The V21--V29 sequence has passed from a small-total-load
regional estimate to exact unrestricted block integration.
V25 removes the volume ceiling for defect \emph{density};
V26 rules out a uniform single-free-link activity under
every nominally good exterior; V27 integrates a genuine
buffer; V28 keeps the resulting block partition functions;
V29 gives aligned-patch conditional lifting.
The remaining difficulty is not another extensive
lower bound. It is an exterior-sensitive comparison for
a \emph{specified} defect set, with a denominator that
does not condition all buffer variables to be good.

We therefore return upstream to the central repair.
The small-exterior restriction in V21 was inherited
from the earlier local pivot chart. V18's globally
smooth inverse permits a different proof with a larger
constant and no such restriction. This yields a
boundary-weighted extraction inequality for actual
compact integrals. The boundary weight will not be
discarded.

\subsection{A quadratic repair for every compact exterior}

Use V21's canonical kernel, its complete collar
\(\mathcal F_R\), and \(s=|R|\). Write
\[
 \mathcal E_R=\sum_{l\in\mathcal F_R}
                    \|\xi_l-J(C)_l\|_{\rm HS}^2.
 \tag{34.1}
\]
The coarse field and exterior remain arbitrary.
Set every free variable in \(I_R\) at its section
value and solve every affected pivot using V18.
Denote this same feasible repair by \(W^{\rm rep}\).
Constants with a star below are new uniform constants;
they need not equal the smaller local constants of V21.

\begin{proposition}[Global squared pivot-repair bound]
\label{prop:f12v30-pivot}
There is \(K_p^*<\infty\), independent of \(R,C,\xi\)
and ambient volume, such that
\[
 \sum_{l\in\Lambda_R}
       \|W_l^{\rm rep}-J(C)_l\|_{\rm HS}^2
                          \le K_p^*\mathcal E_R.
 \tag{34.2}
\]
No maximum mismatch hypothesis is imposed.
\end{proposition}
\begin{proof}
Only crossing-edge pivots can differ from the section,
by V25. Consider one such pivot. Its nonpivot boundary
inputs have been set to their section values. Its
remaining nonsection inputs are internal free links
in the outside endpoint cell, all included in the
completed collar. Let \(q_e\) be their squared
section-distance sum.

At \(q_e=0\), the pivot is exactly \(C_e\).
In a fixed neighbourhood of these section inputs,
V18's uniformly bounded inverse derivatives, local
exponential coordinates and bounded input count give
\[
       \|B_e^{\rm rep}-C_e\|_{\rm HS}^2\le K q_e.
\]
The neighbourhood size and \(K\) are uniform in the
compact parameter \(C_e\).
Outside that neighbourhood, \(q_e\ge r_*^2\)
for some fixed \(r_*>0\), whereas the left side
is at most \(d_G^2\). Increasing \(K\) to
\(\max\{K,d_G^2/r_*^2\}\) proves the same inequality
there. This compactness argument does not choose a
global logarithm or a path through a chart boundary.

Each outside internal input enters only a bounded
number of crossing pivots. Summing their bounds
therefore gives (34.2). All other varied links
equal their section values.
\end{proof}

\begin{proposition}[Global energy-adaptive repair cost]
\label{prop:f12v30-repair}
Put
\[
 K_1^*=12(1+K_p^*)+\frac{K_p^*}{4A},\qquad
 K_3^*=12+K_1^*+\frac1{2A}.
 \tag{34.3}
\]
For every compact exterior,
\[
 H_{\rm rep}-H_0
       \le A L_{\rm in}+M_R+K_1^*\mathcal E_R.
 \tag{34.4}
\]
\end{proposition}
\begin{proof}
For any affected plaquette, write \(P\) for its
repaired holonomy and \(Q\) for its section holonomy.
The elementary Hilbert--Schmidt identity implies
\[
 s(P)-s(Q)\le
       \tfrac12\|P-Q\|_{\rm HS}^2+
                        \delta(Q)\|P-Q\|_{\rm HS}.
 \tag{34.5}
\]
Unitarity and product telescoping bound
\(\|P-Q\|_{\rm HS}\) by the sum of the four
link mismatches. The quadratic term is at most
twice their squared sum. Each link occurs in at
most six plaquettes, so the total quadratic
contribution is at most
\(12(1+K_p^*)\mathcal E_R\), using (34.2).

Only corner plaquettes have nonzero section
\(\delta(Q)\). Their linear terms sum to at most
\[
 \sum_{l\in\Lambda_R}a_l
          \|W_l^{\rm rep}-J(C)_l\|_{\rm HS}+M_R.
\]
Young's inequality bounds the displayed varied-link
sum by \(A L_{\rm in}+K_p^*\mathcal E_R/(4A)\).
Combining these estimates proves (34.4).
No Taylor approximation to the plaquette action
was used.
\end{proof}

The gain over V25 is the use of the \emph{actual}
quadratic mismatch instead of its worst-case bound
\(d_G^2|\mathcal F_R|\). The constant \(K_3^*\)
can be large. This is a gain in the domain and
form of the estimate, not a claimed numerical
improvement.

\subsection{An unrestricted-buffer activity with its own denominator}

Let \(\mathcal A\subset I_R\) be a nonempty specified
set of free group variables, \(k=|\mathcal A|\).
Fix \(0<\tau<d_G\), and define
\[
 \begin{split}
 B_{\mathcal A}(\tau)
    &=\{\|Y_j-J_j(C)\|_{\rm HS}\ge\tau
                                      \text{ for all }j\in\mathcal A\},\\
 G_{\mathcal A}(\tau)
    &=\{\|Y_j-J_j(C)\|_{\rm HS}<\tau
                                      \text{ for all }j\in\mathcal A\}.
 \end{split}
 \tag{34.6}
\]
They are not complementary when \(k>1\).
Define \(Z_R(B_{\mathcal A}\mid y)\) and
\(Z_R(G_{\mathcal A}\mid y)\) by restricting just
these coordinates in the canonical integral (25.2).
All of \(I_R\setminus\mathcal A\) is integrated
over its full compact domain in \emph{both} integrals.
Set
\[
 \mathfrak b_{\mathcal A\mid R}(y)
       =\frac{Z_R(B_{\mathcal A}\mid y)}
                    {Z_R(G_{\mathcal A}\mid y)},\qquad
 \Gamma_R^*=2A L_R+K_3^*\mathcal E_R.
 \tag{34.7}
\]

\begin{theorem}[Energy-adaptive buffered extraction]
\label{thm:f12v30-activity}
There are \(D_4\ge1,\beta_0<\infty,c_0<\infty\)
such that, whenever \(\beta\ge\beta_0\) and
\(\beta\tau^2\ge c_0\),
\[
 \mathfrak b_{\mathcal A\mid R}(y)
       \le D_4^s
           \exp\left[-\beta\left(
                    \frac{\tau^2 k}{4A}-\Gamma_R^*\right)\right].
 \tag{34.8}
\]
This holds for every compact \(C,y\), every
connected \(R\), and every \(\mathcal A\subset I_R\).
In particular there is no assumption \(s\le c\beta\)
or \(k\le k_0\), and no requirement that
\(\Gamma_R^*\) be smaller than a fixed constant.
\end{theorem}
\begin{proof}
On \(B_{\mathcal A}\), the varied physical squared
distance \(D\) satisfies \(D\ge\tau^2k\).
The deterministic inequality (29.6) and
\(D\ge D_{\rm free}\) give the numerator bound
\[
 Z_R(B_{\mathcal A}\mid y)
 \le B^s\beta^{-N_R/2}e^{-\beta H_0}
       e^{-\beta(\tau^2k/(4A)-E_{\rm v,out}
                                  -(A/2)L_{\rm in}-M_R)}.
 \tag{34.9}
\]
The factor \(B^s\beta^{-N_R/2}\) is obtained by
integrating the remaining compact Gaussian
\(e^{-\beta D_{\rm free}/(4A)}\), including all
buffer variables, and bounding the actual
inverse-pivot factors.

For the denominator use a symmetric scalar cube
of radius \(t/\sqrt\beta\) about the free section,
as in V25. With an orthonormal Lie-algebra basis,
each group coordinate in that cube has
section distance at most \(\sqrt d\,t/\sqrt\beta\).
Choosing \(c_0\ge4dt^2\) puts the cube inside
\(G_{\mathcal A}(\tau)\).
The cube restricts the integration only for this
\emph{lower bound}; it is not the definition of
the denominator.

Uniform inverse-pivot and Haar density bounds
and the Hessian upper bound give
\[
 Z_R(G_{\mathcal A}\mid y)
                  \ge b^s\beta^{-N_R/2}e^{-\beta H_{\rm rep}}.
 \tag{34.10}
\]
As in (29.5), the linear Taylor term is handled
by symmetry of the cube, not by assuming
stationarity of the repair.
Use (34.4) in (34.10) and divide (34.9) by it.
The Gaussian powers cancel exactly.
The remaining positive debit is at most
\[
 \begin{split}
 E_{\rm v,out}+\tfrac32 A L_{\rm in}
                   +2M_R+K_1^*\mathcal E_R
 &\le \tfrac32 A L_{\rm in}+2A L_{\rm out}
                            +K_3^*\mathcal E_R\\
 &\le 2A L_R+K_3^*\mathcal E_R.
 \end{split}
\]
Here \(E_{\rm v,out}\le12\mathcal E_R\) and
\(M_R\le A L_{\rm out}+\mathcal E_R/(4A)\).
Taking \(D_4=\max\{1,B/b\}\) proves (34.8).
\end{proof}

Unlike a union bound for single-coordinate
probabilities, this proof supplies a lower bound
for the \emph{all-good denominator} at any \(k\).
It does not need a condition such as
\(kq<1\). Unlike V26's failed singleton activity,
it also integrates an explicitly chosen buffer.
An arbitrary large or badly chosen buffer still
costs \(D_4^s\); that cost has not disappeared.

\subsection{Exact interface transfer and a weighted extraction identity}

For a V29 packing-aligned \(R\), let
\(T=I_R\cap\mathcal S\) and take
\(\mathcal A\subset T\).
The ratio of the two corresponding restricted
integrals of the exact interface kernel \(K_R\)
is precisely \(\mathfrak b_{\mathcal A\mid R}\).
Indeed integrate the restored \(X_R\) and then
the unrestricted \(T\setminus\mathcal A\) in
either order. Positivity permits this exchange,
and Proposition~\ref{prop:f12v29-lift} identifies
the full integrand and all its Jacobians.
No \(-\log Z_b^{\rm c}\) term is removed.

Let \(\nu\) be the actual finite-volume compact
fibre law, or its exact interface marginal in
this aligned case. For every nonnegative bounded
exterior-measurable function \(h\),
\[
 \begin{split}
 &\mathbb E_\nu\left[
   h\,e^{-\beta K_3^*\mathcal E_R}
                       {\bf1}_{B_{\mathcal A}}\right]\\
 &\qquad\le
 D_4^s e^{-\beta(\tau^2 k/(4A)-2A L_R)}
             \mathbb E_\nu[h\,{\bf1}_{G_{\mathcal A}}].
 \end{split}
 \tag{34.11}
\]
To prove this, multiply (34.8) by the conditional
probability of \(G_{\mathcal A}\), then by
\(h e^{-\beta K_3^*\mathcal E_R}\), and integrate
the exterior. Both \(h\) and \(\mathcal E_R\)
are exterior-measurable. This proof does not
assume independence between the collar and
the core.

Equation (34.11) is a boundary-weighted
extraction inequality, not an unweighted
Peierls bound. Removing its exponential
factor on the left would reverse the useful
direction of the estimate. In particular,
a small \(\mathbb E\mathcal E_R\) does not
bound the positive exponential moment that
would be needed to remove the weight.

\subsection{A per-defect estimate without a volume ceiling}

\begin{corollary}[Linear-size buffers with an affordable collar]
\label{cor:f12v30-linear}
Fix \(M\ge1\). Suppose
\[
 s\le Mk,\qquad
 2A L_R\le\frac{\tau^2k}{16A},\qquad
 K_3^*\mathcal E_R\le\frac{\tau^2k}{16A}.
 \tag{34.12}
\]
If, in addition to the thresholds of
Theorem~\ref{thm:f12v30-activity},
\(\beta\tau^2\ge16AM\log D_4\), then
\[
 \mathfrak b_{\mathcal A\mid R}(y)
       \le e^{-\beta\tau^2k/(16A)}.
 \tag{34.13}
\]
The conclusion is uniform as \(s,k\) grow.
\end{corollary}
\begin{proof}
The two debits in (34.12) leave at least
\(\tau^2k/(8A)\) in (34.8).
Since \(s\log D_4\le Mk\log D_4
\le\beta\tau^2 k/(16A)\), the entropy cost
uses at most half of this remainder.
\end{proof}

Here are sufficient conditions with fixed
parameters rather than a vanishing total load:
\[
 \epsilon^2\le
       \frac{\tau^2}{32A^2 C_L M},\qquad
 \frac{\mathcal E_R}{s}\le
       \frac{\tau^2}{16A K_3^*M}.
 \tag{34.14}
\]
They imply (34.12) because
\(L_R\le C_L\epsilon^2s\) and \(s\le Mk\).
The exterior condition is a squared-mismatch
budget; it allows individual large mismatches
when the total budget permits them.

Alternatively use V24's free collar coordinates.
There is a geometric constant \(C_J\) with
\(|J_R^\partial|\le C_Js\). If every such
coordinate has Lie norm at most
\(r\le r_b\), then
\(\mathcal E_R\le C_E C_J r^2s\).
Thus a sufficient, volume-independent
two-threshold choice is
\[
 C_E C_J r^2\le
                  \frac{\tau^2}{16A K_3^*M}.
 \tag{34.15}
\]
This is a stricter exterior threshold than the
core threshold in general. It is not a theorem
that all exteriors nominally good at threshold
\(\tau\) are affordable.
V26's forced singleton example is therefore
not contradicted.

\subsection{Summation over guarded connected supports}

The previous estimate has a direct probabilistic
use even before the boundary weight is absorbed.
Fix a coarse cell \(z\) and an integer \(s_0\ge1\).
Let \(\mathcal U_z(s_0)\) be the event that there
exist a connected coarse region \(R\ni z\),
\(s=|R|\ge s_0\), and a specified
\(\mathcal A\subset I_R\) with \(s\le M|\mathcal A|\),
such that \(B_{\mathcal A}(\tau)\) and both debit
conditions in (34.12) hold.
For the interface marginal restrict this definition
to aligned regions and \(\mathcal A\subset T\).
The debit conditions are part of the event;
they are not silently assumed for random boundaries.

\begin{proposition}[Anchored guarded-support tail]
\label{prop:f12v30-animals}
Under the beta thresholds of
Corollary~\ref{cor:f12v30-linear}, set
\[
 Q=8^2\,2^{73}
           \exp[-\beta\tau^2/(16AM)].
 \tag{34.16}
\]
If \(Q<1\), then, uniformly in ambient volume,
\[
                \nu(\mathcal U_z(s_0))
                         \le \frac{Q^{s_0}}{1-Q}.
 \tag{34.17}
\]
\end{proposition}
\begin{proof}
For fixed \(R,\mathcal A\), the indicator of
the debit conditions is exterior-measurable.
On that exterior event, (34.13) bounds the
conditional bad probability by
\(e^{-\beta\tau^2k/(16A)}\) times the conditional
good probability, hence by
\(e^{-\beta\tau^2s/(16AM)}\).
Integrating over the exterior preserves the bound.

On the four-dimensional coarse nearest-neighbour
graph, degree is at most eight. Connected sets
of \(s\) cells containing \(z\) number at most
\(8^{2s}\): choose a deterministic rooted spanning
tree and encode its depth-first traversal of
length \(2(s-1)\). The visited set recovers the
region, so this gives an injection into such walks.
There are at most \(2^{|I_R|}\le2^{73s}\)
possible \(\mathcal A\)'s.
Union bound over the pairs and then over
\(s\ge s_0\) proves (34.17).
Restricting to aligned pairs only reduces this sum.
No product assumption for different supports
is needed.
\end{proof}

This is a bound on specified high-amplitude cores
with linear-size supports and affordable collars.
It is not the statement that every defect cluster
has such a support. A dilute collection of high
cores can require a much larger connecting buffer;
a boundary may also carry too much mismatch.
Those configurations are not covered by
\(\mathcal U_z(s_0)\), and are not declared rare
by (34.17).

\subsection{What the next comparison must actually prove}

\begin{decision}[Retain and assign the exterior energy debit]
The new input is (34.11), or its capped form
(34.13), for the actual compact law. A subsequent
contour construction must assign the collar
energy \(K_3^*\mathcal E_R\) to other controlled
pieces while keeping a positive residual core
cost, and must control the size of the buffers.
Finite overlap alone bounds the number of
charges but does not make their coefficient
small. Overlapping regional activities also
cannot simply be multiplied.

The useful next test is a concrete collar or
shell construction for this inverse-pivot
geometry: identify exactly which physical
mismatches it charges, how often it charges
them, and what normalized integral remains
after a support is removed. If that charge
balance fails, the failure should determine
the next upstream estimate. Another global
defect-density bound would not resolve it.
\end{decision}

\begin{record}[V30 append record]
V29 had 440267 bytes and SHA--256
\begin{center}\scriptsize\ttfamily
65AA1C586EF6B9BA163907B39BA8F02D14EBBC0A48AA4FF10B964568B290715B.
\end{center}
Its complete body is retained apart from the
title version. The new proof uses V18's global
inverse, V21's exact regional kernel and
relative-action bound, V24's collar comparison,
V25's symmetric lower-integral argument, and
V29's aligned conditional lift.
The companion audit and ten-version direction
review are complete. Only the active V29 is
replaced by V30; all archives and companions
are retained. Next companion audit: V35.
Next ten-version review: V40.
\end{record}

\begin{firewall}[Position after V30]
Proved here: an arbitrary-exterior quadratic
repair bound, an unrestricted-buffer activity
with a matched all-good denominator, an exact
boundary-weighted extraction inequality, and
a guarded connected-support tail with no
upper volume--coupling restriction.

Not proved: absorption of the collar charges,
an exhaustive sparse-defect contour expansion,
unrestricted interface mixing, renormalization
scale iteration, continuum construction, or
the full physical Yang--Mills mass gap.
No claim of completing the active objective
is made.
\end{firewall}

\section{V31: absorbing collar energy by deterministic growth shells}
\label{sec:f12v31}

V30 left an actual quadratic exterior debit.
We now use its dependence on the realized
boundary energy, rather than its worst-case
cardinality bound. The construction below
does not multiply overlapping activities.
Instead, it surrounds any specified set of
bad free variables by finitely many deterministic
neighbourhoods. If no neighbourhood has an
affordable collar relative to its own interior
energy, that energy must grow exponentially.
Compactness bounds its growth by a polynomial.
A dyadic decomposition then applies the actual
normalized conditional estimate to the resulting
finite list of candidate regions.

The outcome is a joint sparse-defect estimate
at a sufficiently small, fixed coarse-curvature
threshold. This removes the exterior guard
from this particular probability bound.
It does not yet prove unrestricted covariance
decay, scale iteration, or a physical mass gap.

\subsection{Regional estimates do not require connectedness}

We continue to use the globally smooth completion
\(\widetilde K\), not the earlier discontinuous
fallback map. Fix the complete coarse field \(C\)
and a single actual field \(W=W(Y,C)\).
Every distance below is evaluated in this same
field against this same section \(J(C)\).
Changing the region does not perform a new gauge
choice or replace the actual field by its repair.

For any nonempty finite set of coarse cells \(R\),
possibly disconnected, use the previous definitions
of \(E(R),I_R,\Lambda_R,\mathcal P_R,\mathcal T_R\)
and the completed frozen collar \(\mathcal F_R\).
Put
\[
 s_R=|R|,\qquad
 D_R(W)=\sum_{l\in\Lambda_R}
                   \|W_l-J(C)_l\|_{\rm HS}^2,\qquad
 E_R(W)=\sum_{l\in\mathcal F_R}
                   \|W_l-J(C)_l\|_{\rm HS}^2.
 \tag{35.1}
\]
Thus \(E_R=\mathcal E_R\) of V30. In this
section \(D_R(W)\) is an energy-distance sum;
it is not the older dimension notation \(D_R\)
from (18.2).

\begin{proposition}[Normalized tail on arbitrary finite cell sets]
\label{prop:f12v31-tail}
There are \(D_5\ge1,\beta_0<\infty\) such that,
for every such \(R\), every compact exterior,
every \(x\ge0\), and \(\beta\ge\beta_0\),
\[
 \pi_{R,C}^{y}\{D_R\ge x\}
 \le \min\left\{1,D_5^{s_R}
       e^{-\beta(x/(4A)-2A L_R-K_3^*E_R)}\right\}.
 \tag{35.2}
\]
The right-hand \(E_R\) is exterior-measurable.
The constants are independent of the number
of connected components of \(R\).
\end{proposition}
\begin{proof}
The canonical conditional formula (25.2) follows
from cancellation of unaffected plaquette and
pivot factors for an arbitrary cell set.
An unaffected edge has both endpoint cells
outside \(R\), so its pivot has no input in
\(I_R\). Connectedness was not used in this
cancellation.

Apply the internal path estimate only in cells
of \(R\), and the boundary path estimate only
on \(E(R)\). The selected rectangles remain
a subcollection of the global rectangles,
so their overlap is at most the same eight.
This proves the same regional vertical
coercivity. The corner estimate and the
frozen-plaquette telescoping estimate are
also sums over sets, not paths connecting
the components. Consequently (29.6) holds
with \(D=D_R\).

The counts remain
\[
 |E(R)|\le8s_R,\quad |I_R|\le73s_R,\quad
 |\Lambda_R|\le81s_R,\quad
 |\mathcal P_R|\le486s_R .
 \tag{35.3}
\]
The V30 repair changes only crossing pivots.
Its finite-input and finite-overlap proof
therefore gives (34.4) for this same \(R\).
The lower integral about the free section
uses all \(I_R\) coordinates in one product
cube, whether or not the region is connected.
The density and Hessian bounds are uniform
by the same local overlap counts.

On \(D_R\ge x\), retain
\(x/(4A)+D_{\rm free}/(4A)\) from
\(D_R/(2A)\), and integrate the remaining
compact Gaussian. Divide by the repaired
lower integral over the full regional
partition function. Both powers are
\(\beta^{-d|I_R|/2}\). The debit calculation
in the proof of (34.8) gives
\(2A L_R+K_3^*E_R\), proving (35.2).
No all-good event or restriction on
\(\beta x\) is needed for this full denominator.
\end{proof}

\subsection{The actual collar occupies the next physical layer}

Let \(N_r(R)\) be the set of coarse cells at
nearest-neighbour graph distance at most \(r\)
from \(R\), with periodic distance on a torus.
We may take the fixed integer \(h=3\) below.
It is a geometric width, not a parameter
chosen from the configuration.

\begin{proposition}[Collar containment and an exact charge count]
\label{prop:f12v31-collar}
For every finite \(R\),
\[
 \mathcal F_R\subset
                   \Lambda_{N_h(R)}\setminus\Lambda_R .
 \tag{35.4}
\]
Thus, for \(R_j=N_{hj}(S)\), if
\(D_j=D_{R_j}(W)\) and \(E_j=E_{R_j}(W)\),
\[
 E_j\le D_{j+1}-D_j,\qquad
                  \sum_{j=0}^{J-1}E_j\le D_J-D_0.
 \tag{35.5}
\]
Each physical mismatch in these selected
collars is charged at most once.
\end{proposition}
\begin{proof}
All links in \(\mathcal F_R\) are frozen,
so they lie outside \(\Lambda_R\).
It remains to bound their location.
The endpoints of a varied boundary family
belong to \(N_1(R)\); varied internal links
belong to cells of \(R\).
A fine plaquette touching a varied link
adds at most one further coarse-cell step.
Thus its link base cells lie in \(N_2(R)\).
The internal-cell plaquettes and the
length-\((3,1,1,1)\) rectangles used for
\(\mathcal T_R\) lie in their cell or
the two endpoint cells of an incident edge.
Frozen internal inputs of affected block
equations also lie in endpoint cells.

Finally, the collar closure of V16 completes
any unaffected boundary family already met
and adds its endpoint internal inputs.
From a boundary link's base cell, this uses
at most one additional coarse step. It adds
no further boundary family, so there is
no iterated growth of this closure.
Every resulting internal link has its cell
in \(N_3(R)\), and every resulting boundary
link is on an edge incident to \(N_3(R)\).
This proves (35.4) with \(h=3\).

The sets \(\Lambda_{R_j}\) are nested.
Since all squared distances use the same
actual \(W,C\), their difference \(D_{j+1}-D_j\)
is precisely the sum over
\(\Lambda_{R_{j+1}}\setminus\Lambda_{R_j}\).
Equation (35.4) bounds \(E_j\) by that sum.
These annuli are pairwise disjoint.
This proves (35.5), including the assertion
about multiplicity. When a region fills
the torus, its collar is empty.
\end{proof}

\begin{record}[Executed finite-stencil check]
The standard-library script
\begin{center}\small\ttfamily
scripts/verify\_ym\_f12\_v31\_collar\_layers.py
\end{center}
enumerates the positive fine links, the
highest-bit rooted tree, the eight incident
boundary families, touching plaquettes,
the sixty plaquettes of each rectangular
stencil, and the completed unaffected-family
inputs using integer coordinates.
Its SHA--256 is
\begin{center}\scriptsize\ttfamily
5CECA1E7A4FF2C5D85FB9E5A202D05DD43C64F58E12160309795DC578C12001C.
\end{center}
For one cell it finds 81 varied links and
928 completed frozen links. Width two
already contains that enumerated collar;
the analytic proof conservatively uses
width three. Tests on a single cell,
adjacent and diagonal pairs, a separated
pair, a three-cell corner, and a sixteen-cell
cube all pass the containment and
no-new-boundary-family checks.

There is also a local explanation for why
single-cell stencils control the location
for arbitrary sets: \(\Lambda_R\),
\(\mathcal P_R\), and \(\mathcal T_R\)
are unions of their cell-selected stencils.
A link frozen for \(R\) is frozen for
each contributing singleton. An edge
unaffected by \(R\) is unaffected by
each such singleton too. Its closure
is therefore already in the corresponding
singleton closure. Thus
\(\mathcal F_R\subset\bigcup_{x\in R}\mathcal F_{\{x\}}\).
The script checks this additional containment
on the listed examples. The arbitrary-set
theorem rests on the set argument and
the geometric proof above, not on treating
finitely many tests as exhaustive.
This verification concerns geometry only,
not the probabilistic conclusions or a
mass-gap certificate.
\end{record}

\subsection{A finite shell must have a relatively affordable collar}

Fix \(0<\tau<d_G\). For a specified nonempty
set \(\mathcal A\) of free group variables,
let \(k=|\mathcal A|\). Assign each internal
variable to its cell and each boundary
variable to the base cell of its positively
oriented coarse edge. Let \(S=S(\mathcal A)\)
be the set of these owner cells. Then
\[
 |S|\le k,\qquad \mathcal A\subset I_S,\qquad
 |N_{hj}(S)|\le k(2hj+1)^4.
 \tag{35.6}
\]
The last bound follows by a union of radius-\(hj\)
boxes around the owners; it is valid after
periodic identifications as well.
No connected hull of distant owners is inserted.

Choose
\[
 \eta=\min\{1,(32AK_3^*)^{-1}\}.
 \tag{35.7}
\]
Fix any integer \(J\ge1\) for which
\[
           (1+\eta)^J\tau^2>
                          81d_G^2(2hJ+1)^4 .
 \tag{35.8}
\]
Such a finite \(J\) exists: the logarithm
of the left side has a positive linear
term in \(J\), whereas that of the
right side grows as \(4\log J+O(1)\).
This choice depends on the fixed threshold
and constants, not on \(\mathcal A,k\),
the volume, the coarse field, or the exterior.
It can be very large; no practical numerical
size is claimed.
Set
\[
 M=(2hJ+1)^4,\qquad
 U=\frac{81d_G^2M}{\tau^2},\qquad
 N=1+\lceil\log_2 U\rceil .
 \tag{35.9}
\]

\begin{proposition}[Deterministic energy-growth alternative]
\label{prop:f12v31-growth}
On the event \(B_{\mathcal A}(\tau)\) of (34.6),
some \(0\le j<J\) satisfies
\[
                         E_j\le\eta D_j.
 \tag{35.10}
\]
Moreover there are such a \(j\) and an
integer \(0\le\ell<N\) with, for
\(x_\ell=2^\ell\tau^2k\),
\[
                  D_j\ge x_\ell,\qquad E_j\le2\eta x_\ell.
 \tag{35.11}
\]
\end{proposition}
\begin{proof}
The specified bad variables are physical
free links in \(\Lambda_S\), so
\(D_0\ge\tau^2 k\).
If \(E_j>\eta D_j\) for every \(j<J\),
(35.5) gives
\[
 D_J>(1+\eta)^J D_0
                  \ge(1+\eta)^J\tau^2k.
\]
On the other hand each physical squared
distance is at most \(d_G^2\), and
(35.3)--(35.6) give
\[
 D_J\le81d_G^2|R_J|
                 \le81d_G^2M k.
\]
This contradicts (35.8).
Hence (35.10) holds at some \(j\).

At that \(j\), monotonicity and compactness
give
\(\tau^2k\le D_j\le81d_G^2M k\).
Choose its dyadic band
\(x_\ell\le D_j<2x_\ell\).
The definition of \(N\) includes all possible
bands, including an upper endpoint equal
to a power of two. Equation (35.10) then
implies \(E_j<2\eta x_\ell\), which proves
the stated weak inequality too.
\end{proof}

The selected \(j,\ell\) may depend on the
entire configuration. We will \emph{not}
condition on this selection. Instead (35.11)
gives containment in the union of \(JN\)
events indexed by deterministic regions and
deterministic thresholds. This distinction
preserves the correct conditional kernels.

\subsection{A joint sparse bound with arbitrary distant exterior}

Define fixed small-curvature and large-coupling
thresholds by
\[
 \epsilon_\tau^2=
             \frac{\tau^2}{32A^2 C_L M},\qquad
 \beta_\tau=
 \max\left\{\beta_0,\,
       \frac{16A}{\tau^2}
                  \bigl(M\log D_5+\log(JN)\bigr)\right\}.
 \tag{35.12}
\]
All constants refer to the same fixed smooth
block map, group representation and threshold
\(\tau\). In particular, \(\epsilon_\tau\)
does not shrink with volume or with \(k\).

\begin{theorem}[Uniform joint free-defect suppression]
\label{thm:f12v31-sparse}
Let \(P=N_{hJ}(S(\mathcal A))\), with
\(\mathcal A\) any nonempty finite set of
free variables. Suppose every affected
coarse corner curvature for \(P\) is at
most \(\epsilon_\tau\).
For \(\beta\ge\beta_\tau\), and every compact
exterior to \(I_P\),
\[
 \pi_{P,C}^{y}\bigl(B_{\mathcal A}(\tau)\bigr)
       \le \exp[-\beta\tau^2|\mathcal A|/(16A)].
 \tag{35.13}
\]
The same bound holds in the full finite-volume
fibre law whenever these curvature hypotheses
hold, and in every infinite-volume DLR state
at a fixed coarse field with globally bounded
curvature \(\delta(Q_p(C))\le\epsilon_\tau\).
No smallness is required of distant free
exterior variables or of compact coarse
link amplitudes.
\end{theorem}
\begin{proof}
Write \(\nu=\pi_{P,C}^{y}\).
For each \(j<J\), disintegrate with respect
to all free coordinates outside \(I_{R_j}\)
inside \(I_P\), retaining the already fixed
exterior to \(I_P\).
Cancellation of the original local density
gives the same canonical
\(\pi_{R_j,C}\) as in
Proposition~\ref{prop:f12v31-tail}.
The collar energy \(E_j\) is measurable
with respect to that exterior.

Since \(\Lambda_{R_j}\subset\Lambda_P\),
its affected corner curvatures are among
the assumed ones. Using \(s_{R_j}\le Mk\),
\[
 2A L_{R_j}\le2A C_L\epsilon_\tau^2 Mk
             \le\frac{\tau^2k}{16A}
             \le\frac{x_\ell}{16A}.
 \tag{35.14}
\]
On the exterior event \(E_j\le2\eta x_\ell\),
(35.7) gives \(K_3^*E_j\le x_\ell/(16A)\).
Thus (35.2) implies, after integration of
that exterior,
\[
 \nu\{D_j\ge x_\ell,\ E_j\le2\eta x_\ell\}
               \le D_5^{Mk} e^{-\beta x_\ell/(8A)}
               \le D_5^{Mk} e^{-\beta\tau^2 k/(8A)} .
 \tag{35.15}
\]
No independence of these two energies
has been asserted.

Proposition~\ref{prop:f12v31-growth} covers
\(B_{\mathcal A}\) by the \(JN\) events
in (35.15). Consequently
\[
 \nu(B_{\mathcal A})\le
        JN\,D_5^{Mk} e^{-\beta\tau^2k/(8A)}.
 \tag{35.16}
\]
Since \(k\ge1\), \(\log(JN)\le k\log(JN)\);
the choice of \(\beta_\tau\) absorbs both
entropy terms into \(\beta\tau^2 k/(16A)\).
This proves (35.13).

For a finite-volume fibre law, condition on
the exterior to \(I_P\) and integrate the
uniform bound. If \(P\) fills the torus,
the same proof has no outer exterior.
For an infinite-volume DLR state the
set \(P\) is finite, even if disconnected,
and its conditional law is the same
regional kernel. Its DLR identity therefore
gives (35.13) directly. Translation invariance
or a special choice of DLR state is not
needed.
\end{proof}

The underlying estimate is on \(k\)
\emph{specified sites}, however far apart
they are. The neighbourhood union has
volume at most \(Mk\); it does not pay
for the distance between disconnected
groups of owners. This is why (35.13)
is not merely a density estimate on a
large connecting cube.

There is also no conflict with V26:
conditioning on all other single-site
variables is not the conditional law
in (35.13). Its buffer has fixed but
potentially enormous width \(hJ\), and
the buffer is fully integrated.
No uniform bound for V26's unbuffered
singleton activity follows.

\subsection{Sparse interface defects and their connected components}

Put
\[
                    p_\tau=e^{-\beta\tau^2/(16A)}.
 \tag{35.17}
\]
For the exact interface marginal of V28,
every finite \(\mathcal A\subset\mathcal S\)
satisfies
\[
 \rho_{\mathcal S}^{\rm c}
      \bigl(B_{\mathcal A}(\tau)\bigr)
                         \le p_\tau^{|\mathcal A|}
 \tag{35.18}
\]
under the global curvature and coupling
conditions of Theorem~\ref{thm:f12v31-sparse}.
This follows simply by marginalizing the
original law. The proof need not condition
on a potentially nonaligned interface
patch. All the exact integrated block
partition functions are therefore retained.
The same conclusion applies to the
interface marginal of an original DLR state.

Fix any deterministic finite-range graph on
the free variables, or on the interface
variables, with uniformly bounded degree
\(\Delta\). One may in particular use
V28's effective interaction graph for the
fixed packing. Call a vertex bad when
its section distance is at least \(\tau\).
Let \(\Delta_*=\max\{2,\Delta\}\).

\begin{corollary}[Exponential bad-component size bound]
\label{cor:f12v31-cluster}
For every anchor vertex \(v\) and \(m\ge1\),
\[
 \mathbb P\{v\text{ belongs to a bad component
                                of size at least }m\}
                         \le(\Delta_*^2p_\tau)^m .
 \tag{35.19}
\]
If \(\Delta_*^2p_\tau<1\), there is almost surely
no infinite bad component in an infinite-volume
state satisfying the preceding hypotheses.
\end{corollary}
\begin{proof}
A connected component with at least \(m\)
vertices contains a connected \(m\)-vertex
subset containing its anchor, by growing
a finite spanning tree.
Such subsets number at most
\(\Delta_*^{2m}\), using the depth-first
walk encoding. Each specified subset is
all bad with probability at most \(p_\tau^m\)
by (35.13) or (35.18).
Union bound proves (35.19).
Letting \(m\) tend to infinity makes the
probability of an infinite component at
each anchor zero. The vertex set is
countable, so union over the anchors
proves the almost-sure conclusion.
\end{proof}

For an interface coordinate cube contained
in V28's chart, choose \(\tau>0\) so that
the group ball of radius \(\tau\) about
the section is inside that cube.
A coordinate outside the cube then
has distance at least \(\tau\).
Equations (35.18)--(35.19) consequently
bound the actual box-defect components
of the unrestricted interface law.
The threshold can be chosen uniformly
over the compact coarse links by left
translation of the same group chart.

These joint all-bad estimates are not
asserted to give stochastic domination
by independent Bernoulli variables.
For example, bounds on the probabilities
that every member of a specified set is
bad do not alone bound every increasing
event by its Bernoulli probability.
Only the explicitly proved union-bound
consequences are used here.

\begin{decision}[The remaining correlation comparison]
The collar debit has now been absorbed
for joint defect probabilities using
energy growth and one canonical kernel
at a time. The arbitrary-exterior guard
is no longer a missing assumption in
(35.13). The remaining task is to use
these sparse bad components in an
\emph{exact} covariance or partition-function
decomposition, controlling the conditional
means and the buffer interactions they
produce.

A probability that a fixed component is
rare is not a derivative bound for the
effective action conditioned on that
component. Nor does it justify summing
an extensive total-variation error.
The next comparison must retain those
distinctions while joining V28's good-region
convexity to (35.19).
The small-curvature hypothesis on the
retained field must also eventually be
integrated or propagated through scales;
it is not a property asserted here of
all coarse fields under the full YM law.
\end{decision}

\begin{record}[V31 append record]
V30 had 458002 bytes and SHA--256
\begin{center}\scriptsize\ttfamily
D0CB7A8FC7C851349E9C442AE7183DB190023228848B5999EE27E16DB46782CB.
\end{center}
Its complete body is retained apart from
the title version.
The new ingredients are arbitrary-set
regional bookkeeping, a collar charged
to a disjoint physical layer, the
compactness-versus-energy-growth argument,
and a deterministic finite-union use of
the normalized kernel.
Only active V30 is replaced by V31.
All archives and companion files remain.
The next companion audit is V35 and
the next direction review is V40.
\end{record}

\begin{firewall}[Position after V31]
Proved: at a fixed sufficiently small
coarse-curvature threshold and sufficiently
large coupling, joint suppression of
specified free-variable defects, uniformly
in volume and arbitrary exterior beyond
their fixed-width neighbourhood union;
the exact interface marginal inherits
this bound and exponential bad-component
size tails.

Still unproved: unrestricted covariance
decay through the resulting random good
and bad regions, a scale-stable effective
interaction class, integration of the
coarse-curvature exceptions, continuum
construction, and the physical
Yang--Mills mass gap.
\end{firewall}

\section{V32: boundary-stable sparse bounds and unrestricted core decoupling}
\label{sec:f12v32}

The older conditional-mean comparison required a
small \emph{total} boundary mismatch and a region
of at most order \(\beta\) cells. Simply inserting
V31 into that comparison would retain those
restrictions. We instead prove stability of the
sparse estimate under a small \emph{pointwise}
boundary condition, without a boundary-size
restriction. We then compare the actual compact
kernel with its convex-box restriction and
populate that pointwise boundary condition
using V31.

All measures below are the actual smooth-block
fibre measures at a fixed complete \(C\).
The eventual decoupling estimate applies to
arbitrary bounded measurable core observables,
not only observables with a prescribed derivative
bound. Its error is retained explicitly.

\subsection{Truncated collars and their two sources of energy}

Write \(\Lambda_R,\mathcal F_R,D_R,E_R\)
as in V31. For \(Q\subset R\), the collar
of \(Q\) splits into links still varied
in \(R\) and links fixed by the exterior
of \(R\). With \(w=4\),
\[
 \begin{split}
 \mathcal F_Q\setminus\Lambda_R&\subset\mathcal F_R,\\
 \mathcal F_Q\cap\Lambda_R
   &\subset\Lambda_{R\cap N_w(Q)}\setminus\Lambda_Q.
 \end{split}
 \tag{36.1}
\]
Here and below \(N_r\) is the coarse graph
neighbourhood, not a new block map.

\begin{proof}
For the first inclusion, the initial plaquette
and rectangle stencils for \(Q\) are subsets
of those for \(R\). Their links outside
\(\Lambda_R\) are therefore initial frozen
inputs for \(R\).
Consider a family completed in the collar
closure for \(Q\). If its coarse edge is
unaffected by \(R\), the same completion
belongs to \(\mathcal F_R\). If it is
affected by \(R\), its boundary links
belong to \(\Lambda_R\); any endpoint
internal inputs still outside \(\Lambda_R\)
are among the frozen internal inputs
of the affected \(R\)-block equation.
These exhaust the closure, which adds
no new boundary family.

For the second inclusion use (35.4).
An internal link in both \(\mathcal F_Q\)
and \(\Lambda_R\) has its cell in
\(R\cap N_3(Q)\). For a boundary link,
one endpoint of its coarse edge lies in
\(N_3(Q)\), and some endpoint lies in \(R\).
If these are different endpoints, the
latter lies in \(N_4(Q)\).
Thus the whole boundary family is
varied in \(R\cap N_4(Q)\).
Every link of \(\mathcal F_Q\) is frozen
relative to \(Q\), proving the set difference.
\end{proof}

Suppose the parent exterior obeys
\(\max_{\mathcal F_R}\|W_l-J_l(C)\|_{\rm HS}\le b\).
For \(S\subset R\), put
\(R_j=R\cap N_{wj}(S)\), \(D_j=D_{R_j}\),
\(E_j=E_{R_j}\). Equations (36.1) imply
\[
 E_j\le D_{j+1}-D_j+C_F b^2|R_j|.
 \tag{36.2}
\]
The last term charges only links fixed
by the parent exterior. The first term
charges a disjoint physical layer.
We have not replaced either by the
total energy in a larger unrelated box.

\begin{record}[Executed geometric verification]
The script
\begin{center}\small\ttfamily
scripts/verify\_ym\_f12\_v32\_truncated\_collars.py
\end{center}
uses V31's integer-lattice stencil functions
and checks (36.1) on sixteen nested pairs:
three-cell corners, a twelve-cell line,
a disconnected parent, and a sixteen-cell
cube, with several inner subsets.
All set inclusions and the exact split
into inside-parent and outside-parent
links pass. Its SHA--256 is
\begin{center}\scriptsize\ttfamily
932995E57D1F526B1DC59861C01A61F395A65953EA61D0E24F4FA13B7C43AC13.
\end{center}
The arbitrary-set statement is proved
above; the finite checks test the actual
geometric bookkeeping rather than replace
that proof.
\end{record}

\subsection{Sparse suppression stable under a pointwise boundary condition}

Fix \(0<\tau<d_G\) and V31's
\(\eta=\min\{1,(32AK_3^*)^{-1}\}\).
Choose \(J_b\ge1\) so that
\[
 (1+\eta/2)^{J_b}\tau^2
                   >81d_G^2(2wJ_b+1)^4.
 \tag{36.3}
\]
Define
\[
 \begin{gathered}
 M_b=(2wJ_b+1)^4,\qquad
 N_b=1+\left\lceil
       \log_2(81d_G^2M_b/\tau^2)\right\rceil,\\
 b_\tau^2=\frac{\eta\tau^2}{2C_FM_b},\qquad
 \epsilon_{b,\tau}^2=\frac{\tau^2}{32A^2C_LM_b},\\
 \beta_{b,\tau}=
 \max\left\{\beta_0,\frac{16A}{\tau^2}
             (M_b\log D_5+\log(J_bN_b))\right\}.
 \end{gathered}
 \tag{36.4}
\]
These are fixed constants, independent
of parent size and boundary size.

\begin{theorem}[Pointwise-boundary stable joint bound]
\label{thm:f12v32-sparse}
For every finite parent \(R\), if its
affected coarse curvatures are at most
\(\epsilon_{b,\tau}\), its physical collar
mismatches are at most \(b_\tau\), and
\(\beta\ge\beta_{b,\tau}\), then for every
nonempty \(\mathcal A\subset I_R\),
\[
 \pi_{R,C}^{y}(B_{\mathcal A}(\tau))
                  \le e^{-\beta\tau^2|\mathcal A|/(16A)}.
 \tag{36.5}
\]
There is no bound on
\(\sum_{\mathcal F_R}\|W_l-J_l(C)\|^2\)
other than the one implied by the
pointwise condition.
\end{theorem}
\begin{proof}
For a boundary variable of \(\mathcal A\),
choose as its owner an endpoint cell
of its edge lying in \(R\), even if the
positively oriented base cell is outside
\(R\). Internal variables own their cells.
With \(k=|\mathcal A|\), the resulting
\(S\subset R\) has at most \(k\) cells
and \(\mathcal A\subset I_S\).
The truncated sets \(R_j=R\cap N_{wj}(S)\)
have at most \(M_bk\) cells for \(j\le J_b\).
On \(B_{\mathcal A}\), \(D_0\ge\tau^2k\).

If \(E_j>\eta D_j\) for all \(j<J_b\),
(36.2) and (36.4) give
\[
 D_{j+1}>(1+\eta)D_j-C_Fb_\tau^2M_b k
                         \ge(1+\eta/2)D_j.
 \tag{36.6}
\]
We used \(D_j\ge D_0\ge\tau^2k\), which
holds by nesting independently of (36.6).
This contradicts
\(D_{J_b}\le81d_G^2M_b k\) and (36.3).
Thus some \(j<J_b\) has \(E_j\le\eta D_j\).

Place its \(D_j\) in one of the \(N_b\)
dyadic bands starting at \(\tau^2k\).
Exactly as in (35.14)--(35.16), its lower
endpoint \(x_\ell\) satisfies
\[
 2A L_{R_j}\le x_\ell/(16A),\qquad
 E_j\le2\eta x_\ell .
 \]
The latter is an exterior event for the
canonical kernel on \(I_{R_j}\).
Disintegrate that kernel within the actual
parent law and use (35.2), then union bound
over the deterministic \(J_bN_b\) choices.
This gives
\(J_bN_bD_5^{M_bk}e^{-\beta\tau^2k/(8A)}\).
The last threshold in (36.4) absorbs the
entropy and proves (36.5).
No conditioning on the randomly selected
shell has been performed.
\end{proof}

\subsection{A regional compact-to-box comparison without a total-load restriction}

Choose a fixed sufficiently small free
coordinate cube \(\mathcal D_R\) of radius
\(R_*>0\) as in V28.
Choose \(\tau_0>0\) such that the group
ball of radius \(\tau_0\) at each free
section value lies inside that coordinate
cube. Fix
\(b\le b_{\tau_0}\) and
\(\epsilon\le\epsilon_{b,\tau_0}\),
decreasing them further below.
Write \(s=|R|\) and \(c_0=\tau_0^2/(16A)\).
The preceding theorem gives
\[
 p_R:=\pi_{R,C}^{y}(\mathcal D_R^c)
                    \le \min\{1,73s e^{-c_0\beta}\}.
 \tag{36.7}
\]
This is a union bound for the \emph{same}
compact kernel, not a separately normalized
good-background law.

\begin{proposition}[Uniform regional boxed convexity]
\label{prop:f12v32-convex}
For sufficiently small fixed \(R_*,b,\epsilon\)
and sufficiently large \(\beta\), if the
physical exterior collar is pointwise at
most \(b\) and affected coarse curvatures
are at most \(\epsilon\), the measured
coordinate potential of the regional
kernel satisfies
\[
 m_0\beta I\le V_R''\le K\beta I,\qquad
               \sup_i\sum_j|(V_R'')_{ij}|\le K\beta
 \tag{36.8}
\]
on \(\mathcal D_R\), with fixed interaction
range and constants independent of \(s\).
\end{proposition}
\begin{proof}
At the section, the vertical Hessian
estimate of V15 applied to variations
supported in \(I_R\) is at least
\(2A^{-1}I\). The possible negative
corner remainder is bounded by
\(K\epsilon I\), since only affected
corner plaquettes enter this restriction.
Thus no curvature bound on remote
unaffected plaquettes is needed.

Every exterior free input that affects
the regional density lies in the completed
physical collar. If \(b\) is small, its
relative logarithm has norm at most \(K b\).
The regional Hessian is the principal
Hessian of the full local potential;
unaffected terms have zero derivative
in \(I_R\). Finite stencil overlap and
the fixed-order local pivot derivatives
bound its change from section data in
row norm by \(K(R_*+b)\).
The logarithms of the inverse-pivot
and free Haar factors have bounded
Hessian row norm. These are order-one
perturbations of the order-\(\beta\)
action Hessian. Choosing the radii and
then \(\beta\) proves (36.8).
This is the regional version of V28's
pointwise-curvature argument, not an
invocation of a small total \(L_R\).
\end{proof}

Let \(F,G\) be real bounded functions,
\(C^1\) on this cube, with suprema
\(M_F,M_G\), boxed gradient norms
\(L_F,L_G\), and declared local coordinate
supports \(A_F,A_G\).
The boxed covariance theorem of f11 V86
applied using (36.8), followed by comparison
with the exact restriction of the compact
law, gives
\[
 \begin{split}
 |\operatorname{Cov}_{\pi_{R,C}^y}(F,G)|
 \le{}&K\beta^{-1}e^{-\gamma\,{\rm dist}(A_F,A_G)}
                                            L_F L_G\\
      &+6M_FM_G\min\{1,73s e^{-c_0\beta}\}.
 \end{split}
 \tag{36.9}
\]
To verify the second term, conditioning
on \(\mathcal D_R\) changes the mean of
any bounded \(H\) by at most \(2M_Hp_R\).
Apply this to \(FG,F,G\) in the covariance
identity. The first term is for the
convex restriction of that very kernel.
Neither \(|R|\le c\beta\) nor a small
total exterior mismatch appears in (36.9).
The explicit error still grows with \(s\).

\subsection{Smoothing bounded core tests by an exact conditional integral}

Let \(Q\subset R\) be an inner cell region,
with \(J_Q^\partial\subset I_R\), where
\(J^\partial\) is V24's complete free-input
set. Put \(q_Q=|J_Q^\partial|\) and
\(q_R=|J_R^\partial|\), and suppose both
are nonzero for the estimates below.
Empty input sets give the corresponding
exact independence and may be treated
separately.
Let \(F\) be any measurable function of
\(Y_{I_Q}\) with \(0\le F\le1\), and define
\[
 H_F(z)=\int F(Y_{I_Q})\,\pi_{Q,C}^{z}(dY_{I_Q}).
 \tag{36.10}
\]
This is a compact conditional integral,
with no box imposed on \(I_Q\).
It depends only on \(z_{J_Q^\partial}\),
takes values in \([0,1]\), and is smooth
in these exterior variables even when
\(F\) is merely measurable.

Indeed differentiating its normalized
integral gives exactly
\(\partial_{z_j}H_F=-\operatorname{Cov}(F,S_{Q,j})\).
The global smooth inverse and local
score bound (25.18) give
\[
       |\partial_{z_j}H_F|\le K(\beta+1),\qquad
       L_{H_F}\le K(\beta+1)\sqrt{q_Q}.
 \tag{36.11}
\]
Fixed Lie-component multiplicities and
the derivative of the local exponential
chart are included in \(K\).
Differentiation under the integral is
justified by the smooth positive compact
kernel and bounded \(F\); it does not
differentiate \(F\).
The tower identity gives
\(\mathbb E_{\pi_{R,C}^{y}}F
 =\mathbb E_{\pi_{R,C}^{y}}H_F\).
Consequently no derivative assumption
on the original core test is needed.

Choose a fixed \(r_g>0\) so small that
the free exterior conditions
\[
 \|v_j\|\le r_g\quad(j\in J_R^\partial),\qquad
               y_j=J_j(C)e^{v_j},
 \tag{36.12}
\]
imply physical collar mismatch at most \(b\).
V24's pointwise reconstruction bound
ensures this if \(r_g\le r_b\) and
\(C_\infty r_g\le b\).
Call this guard \(\mathcal G_R\).
Its radius is per group coordinate;
there is no factor \(q_R^{-1/2}\).
Coordinatewise linear interpolation of
the \(v_j\)'s stays in \(\mathcal G_R\)
and retains the complete fixed \(C\).

Let \(L={\rm dist}(J_Q^\partial,J_R^\partial)\).
Along any such path, the exact normalized
response and (36.9), applied to \(H_F\)
and the exterior score \(S_j\), give
\[
 \left|\partial_{y_j}
       \mathbb E_{\pi_{R,C}^{y}}H_F\right|
 \le K(\beta+1)\sqrt{q_Q}\,
               e^{-\gamma\,{\rm dist}(J_Q^\partial,j)}
                  +K(\beta+1)s e^{-c_0\beta}.
 \tag{36.13}
\]
The score and its boxed gradient are
bounded by \(K(\beta+1)\), and its support
lies within fixed distance of \(j\).
This accounts for both terms. If the
boxed gradient of \(H_F\) vanishes, its
boxed covariance is zero; the displayed
upper bound remains valid with the
declared dependence set \(J_Q^\partial\).
The compact error is still retained.

Each interpolated coordinate has path
length bounded by a fixed multiple of
\(r_g\). Summing (36.13) over the \(q_R\)
inputs, and using the tower identity, yields
\[
 \operatorname{osc}_{\mathcal G_R}
       \mathbb E_{\pi_{R,C}^{y}}F
 \le K(\beta+1)q_R\sqrt{q_Q}\,e^{-\gamma L}
              +K(\beta+1)q_Rs e^{-c_0\beta}
                  =:\delta_{\rm good}.
 \tag{36.14}
\]
The bound is uniform over every measurable
\(F\) in \([0,1]\) on the whole core \(I_Q\).
All partition-function derivatives have
been included in the scores.

\subsection{Populate the pointwise guard without a protected ancestor}

Choose a fixed \(\tau_g>0\) small enough
that section distance below \(\tau_g\)
implies a unique relative logarithm of
norm less than \(r_g\).
Use V31's constants \(J(\tau_g)\),
\(\epsilon_{\tau_g}\), and \(\beta_{\tau_g}\).
Let an outer region \(P\) contain \(R\)
and the \(hJ(\tau_g)\)-neighbourhood of
the owner cells of \(J_R^\partial\).
Assume all affected coarse curvatures
for \(P\) are at most
\[
                   \epsilon_{\rm dec}
                     :=\min\{\epsilon,\epsilon_{\tau_g}\},
 \tag{36.15}
\]
and take \(\beta\) above all the fixed
thresholds used so far.
For every arbitrary compact exterior
to \(I_P\), V31 applied to each guard
coordinate, then union bound, gives
\[
 \pi_{P,C}^{y_P}(\mathcal G_R^c)
             \le q_R e^{-c_g\beta},\qquad
                         c_g=\tau_g^2/(16A).
 \tag{36.16}
\]
If its individual neighbourhood is smaller
than \(P\), first disintegrate that
canonical neighbourhood kernel within
the parent. Its bound is uniform over
the resulting exterior.
No smallness assumption on the parent's
physical boundary has been used.

\begin{theorem}[Unrestricted diameter of the core marginal]
\label{thm:f12v32-tv}
Under the stated nesting, curvature and
coupling conditions, any two compact
exteriors \(y_P,y'_P\) obey
\[
 \left\|(\pi_{P,C}^{y_P})_{I_Q}
             -(\pi_{P,C}^{y'_P})_{I_Q}\right\|_{\rm TV}
 \le \min\{1,\delta_{\rm good}+2q_R e^{-c_g\beta}\}
                          =:\delta_{Q,R}.
 \tag{36.17}
\]
Here \(\|\mu-\nu\|_{\rm TV}
=\sup_{0\le F\le1}|\mu F-\nu F|\).
There is no restriction of the actual
interface or core law to a good box.
\end{theorem}
\begin{proof}
For any measurable \(0\le F\le1\) on \(I_Q\),
condition each parent law on the exterior
to \(I_R\). Its conditional mean is the
same function
\(m_F(y)=\mathbb E_{\pi_{R,C}^{y}}F\).
On \(\mathcal G_R\), its values lie in
an interval of width at most
\(\delta_{\rm good}\), by (36.14).
Clamp \(m_F\) to the closure of this
interval. Clamping does not change it
on the guard, and changes it by at most
one elsewhere. Under either parent law
the expectation changes by at most
\(q_R e^{-c_g\beta}\), by (36.16).
The two clamped expectations differ by
at most \(\delta_{\rm good}\).
Their sum of errors proves (36.17),
uniformly in \(F\); take the supremum.
This step controls the conditional means,
not only the mean conditional covariance.
\end{proof}

\subsection{A linear-in-coupling spatial scale}

Work on the infinite lattice, or on a
torus large enough to embed the following
cubes and their fixed stencils without
wrapping. Take concentric coarse cubes with radii
\(\lfloor r/4\rfloor,r,2r\) for \(Q,R,P\).
For sufficiently large \(r\),
\[
 s\le K r^4,\quad q_Q,q_R\le K r^3,\quad L\ge c r.
 \tag{36.18}
\]
The fixed-width collar and owner conventions
alter these bounds only by fixed constants.
The outer buffer condition holds once
\(r\) exceeds the fixed \(hJ(\tau_g)\)
and stencil widths.
With \(r=\lfloor\beta\rfloor\), (36.14),
(36.16), and (36.18) imply
\[
                         \delta_{Q,R}\le C e^{-c\beta}
 \tag{36.19}
\]
for fixed \(c,C>0\) and all sufficiently
large \(\beta\).
Indeed each prefactor in (36.14) and
(36.16) is polynomial in \(\beta\),
whereas the three exponentials have
fixed positive multiples of \(-\beta\).
Their polynomial factors can be absorbed
by decreasing the final exponent.

The coarse-curvature threshold in this
statement is the fixed positive
\(\epsilon_{\rm dec}\) of (36.15), not
a threshold tending to zero with volume
or with \(\beta\). Its value, as well
as the required coupling, may be very
conservative.
There is no \(|R|\le c\beta\) condition:
these \(R\)'s have order \(\beta^4\) cells.
No claim is made that the full Yang--Mills
coarse law automatically satisfies the
pointwise curvature hypothesis.

\begin{corollary}[Exact local decoupling consequences]
\label{cor:f12v32-decoupling}
In any finite-volume fibre law, or any
DLR state, for which the hypotheses hold
around these cubes, the joint law of
\(Y_{I_Q}\) and the entire exterior
\(Y_{I_P^c}\) differs from the product
of its marginals by at most
\(\delta_{Q,R}\) in total variation.
For bounded real \(F(Y_{I_Q})\) and
\(G(Y_{I_P^c})\),
\[
 |\operatorname{Cov}(F,G)|
           \le \tfrac14\delta_{Q,R}\,
                         \operatorname{osc}(F)\operatorname{osc}(G).
 \tag{36.20}
\]
The core marginals of any two DLR states
at the same globally small-curvature
coarse field differ by at most
\(\delta_{Q,R}\).
\end{corollary}
\begin{proof}
The conditional core kernels given the
outer exterior have pairwise total-variation
distance at most \(\delta_{Q,R}\).
Averaging one argument shows that each
is at distance at most \(\delta_{Q,R}\)
from its own mixture, the core marginal.
Integrating this bound over the exterior
proves the joint-versus-product assertion.

For the covariance, the conditional
mean of \(F\) has oscillation at most
\(\delta_{Q,R}\operatorname{osc}(F)\).
Since \(G\) is exterior-measurable,
the covariance equals that of this
conditional mean and \(G\).
Cauchy--Schwarz and the variance bound
\(\operatorname{Var}(X)\le
\operatorname{osc}(X)^2/4\) prove (36.20).
For two DLR states, average the same
pairwise kernel bound over their two
possibly different exterior laws.
\end{proof}

\begin{decision}[Do not silently iterate a finite-scale diameter bound]
We have an unrestricted, normalized
core-decoupling estimate with an error
exponentially small in large coupling.
At fixed \(\beta\), however, the explicit
errors in (36.14) and (36.16) grow with
the region's size. This proof does not
let \(r\) tend to infinity at fixed
\(\beta\), and the small diameter of
the family of DLR marginals is not
zero diameter.

To obtain unrestricted exponential
clustering, the next step must specify
a valid finite-size iteration or
coupling criterion for these conditional
kernels. Conditioning on previously
coupled coordinates can remove required
buffers; overlapping block updates do
not factor. Those issues must be proved
away before (36.19) can be promoted to
an all-distance mixing theorem.
The estimate is also about a fixed
small-curvature fibre; propagation or
integration of the retained-field
exceptions remains separate.
\end{decision}

\begin{record}[V32 append record]
V31 had 477636 bytes and SHA--256
\begin{center}\scriptsize\ttfamily
D853EB0078627D8C42989701C0846D03175A762725F6A43962D8354284EC6AB6.
\end{center}
Its complete body is retained apart
from the title version. The new note
adds the truncated-shell charge identity,
boundary-stable sparse suppression,
the size-explicit compact covariance
comparison, exact smoothing of measurable
core tests, and unrestricted core
decoupling at a linear-in-\(\beta\)
spatial scale.
Only active V31 is replaced by V32.
All archives and companions are retained.
Next companion audit: V35.
Next direction review: V40.
\end{record}

\begin{firewall}[Position after V32]
The new total-variation estimate controls
the conditional means as well as
conditional covariances, and does not
require a protected outer ancestor.
It is nevertheless a finite-scale,
fixed-small-curvature result.
Fixed-coupling uniqueness and all-distance
clustering, the scale-stable effective
interaction, continuum construction,
and the physical Yang--Mills mass gap
remain unproved.
\end{firewall}

\section{V33: whole-block transport and the missing contraction hypothesis}
\label{sec:f12v33}

V32 controls the marginal on a deep core.
A finite-size iteration must also control
what happens to the rest of an updated
block. We now derive an actual whole-block
transport bound, with no leading volume
factor, on V32's pointwise boundary guard.
We then compare its precise quantifiers
with a finite-volume block contraction
criterion. The good-boundary coefficient
has the desired surface-to-volume scaling;
the all-boundary hypothesis is not yet
proved. This distinction prevents an
unjustified uniqueness claim.

\begin{record}[Primary-source check and scope]
The block-comparison framework was checked
against P.~Rebeschini and R.~van Handel,
\emph{Comparison Theorems for Gibbs Measures},
\href{https://arxiv.org/pdf/1308.4117}
{arXiv:1308.4117}, Definition~2.3,
Theorem~2.4, Corollary~2.8 and Remark~2.9.
The comparison uses couplings of complete
block updates for boundary pairs, with
an influence contraction condition;
closeness of just an interior marginal
is not that hypothesis. We do not claim
their theorem applies to the present
unrestricted fibre. The finite contraction
lemma below is included with a proof,
and is not claimed as a new general
probabilistic principle.
\end{record}

\subsection{The metric and the componentwise covariance estimate}

For free group variables in a finite set
\(B\), use the compact product metric
\[
 d_B(Y,Z)=\sum_{i\in B}\|Y_i-Z_i\|_{\rm HS},
 \qquad
 \mathsf W_B(\mu,\nu)
       =\inf_{\lambda\in\operatorname{Cpl}(\mu,\nu)}
                         \int d_B(Y,Z)\,d\lambda .
 \tag{37.1}
\]
Its diameter is at most \(d_G|B|\).
This is not the Hamming metric and
does not require identical coordinates
under a coupling.
On this compact metric space,
\[
 \mathsf W_B(\mu,\nu)=
       \sup_{\operatorname{Lip}_{d_B}F\le1}
                              |\mu F-\nu F|.
 \tag{37.2}
\]
One way to verify the dual formula here
is to project onto a finite \(\varepsilon\)-net,
apply finite transportation linear-programming
duality, and extend its one-Lipschitz
potential by
\(\min_a\{f(a)+d_B(\,\cdot\,,a)\}\).
Projection changes transport cost by
at most \(2\varepsilon\) and changes
any one-Lipschitz expectation difference
by at most \(2\varepsilon\).
Let \(\varepsilon\) decrease to zero.
The reverse inequality follows by
integrating \(|F(Y)-F(Z)|\le d_B(Y,Z)\)
against any coupling.
Compactness of the set of couplings
and continuity of \(d_B\) also give
an attaining coupling.

We use V32's uniform boxed Hessian
hypotheses. Let \(\mathcal A\) be the
vector response operator with the
mixed flat-face conditions of f11 V86,
and \(P_i\) projection onto scalar
coordinate component \(i\).
The same weighted conjugation proof
as in that archive gives
\[
 \|P_i\mathcal A^{-1}P_l\|_{L^2\to L^2}
                \le K\beta^{-1}
                          e^{-\gamma\,{\rm dist}(i,l)}.
 \tag{37.3}
\]
Indeed take a vector source supported
on component \(l\), multiply component
\(i\) by \(e^{\gamma{\rm dist}(i,l)}\),
and apply the perturbed-resolvent norm
bound. That proof does not require the
source to be the gradient of a scalar
function. Its weights are constant
in the integration variables and
preserve every mixed face condition.

Consequently, for the actual boxed law,
\[
 |\operatorname{Cov}_{\rm box}(F,G)|
 \le\frac K\beta
    \sum_{i,l}e^{-\gamma{\rm dist}(i,l)}
       \|\partial_iF\|_{L^2}
       \|\partial_lG\|_{L^2}.
 \tag{37.4}
\]
This follows by decomposing both
gradients in the covariance representation
and using (37.3) term by term.
The sum is finite at finite volume;
the constants are volume independent.

In particular let \(S_j\) be the
normalized-response score for changing
one exterior free coordinate \(y_j\).
Its gradient is supported on a fixed
number of components near \(j\), with
each derivative bounded by \(K(\beta+1)\).
For \(F\) one-Lipschitz in \(d_{I_R}\),
each boxed coordinate derivative is
bounded by a fixed constant. Since
the coordinate positions have uniform
polynomial ball growth,
\[
 \sup_j\sum_i e^{-\gamma{\rm dist}(i,j)}
                              <\infty.
 \]
Thus, for \(\beta\ge1\),
\[
              |\operatorname{Cov}_{\rm box}(F,S_j)|
                                      \le K_T .
 \tag{37.5}
\]
There is no \(\sqrt{|R|}\) factor in
(37.5). A single aggregate gradient
norm in (36.9) would lose this fact.
For nonsmooth one-Lipschitz \(F\), use
product-group convolution with smooth
approximate identities. Bi-invariance
of the Hilbert--Schmidt metric makes
each common group translation an
isometry, so convolution preserves
the Lipschitz bound and converges
uniformly. The estimates pass to
the limit.

\subsection{Whole compact block response on the pointwise guard}

Let \(R\) have \(s\) coarse cells and
let its physical and free pointwise
boundary conditions be those of V32.
In particular the free guard
\(\mathcal G_R\) has fixed per-coordinate
radius \(r_g\), and
\[
       p_R=\pi_{R,C}^{y}(\mathcal D_R^c)
                   \le \min\{1,73s e^{-c_0\beta}\}.
 \tag{37.6}
\]
All affected coarse curvatures satisfy
the fixed small threshold of that
result.

\begin{theorem}[A whole-block transport coefficient]
\label{thm:f12v33-transport}
There are fixed constants \(K_T,K_E\)
such that, for \(y,z\in\mathcal G_R\),
\[
 \begin{split}
 \mathsf W_{I_R}(\pi_{R,C}^{y},\pi_{R,C}^{z})
 &\le a_R(\beta)
       \sum_{j\in J_R^\partial}
                           \|y_j-z_j\|_{\rm HS},\\
 a_R(\beta)&=K_T+
                    K_E(\beta+1)s^2 e^{-c_0\beta}.
 \end{split}
 \tag{37.7}
\]
This compares the \emph{full compact}
laws on every updated free variable,
not their boxed restrictions or core
marginals.
\end{theorem}
\begin{proof}
Take \(\operatorname{Lip}_{d_{I_R}}F\le1\)
and subtract its value at a fixed
configuration. Then
\[
           \|F\|_\infty\le d_G|I_R|\le73d_Gs.
 \]
The exact score identity gives
\(\partial_{y_j}\pi_{R,C}^yF
 =-\operatorname{Cov}_{\pi_{R,C}^y}(F,S_j)\).
The compact-versus-boxed covariance
comparison, (37.5), the bound
\(\|S_j\|_\infty\le K(\beta+1)\),
and (37.6) imply, for a unit Lie
direction,
\[
 \left|\partial_{y_j}\pi_{R,C}^yF\right|
                  \le K_T+
                          K_E(\beta+1)s^2e^{-c_0\beta}.
 \tag{37.8}
\]
The two powers of \(s\) have distinct
sources: the diameter bound on the
test and the probability of a box
defect. They are retained.

Write \(y_j=J_j(C)e^{v_j}\),
\(z_j=J_j(C)e^{w_j}\).
The straight segment from \(v_j\)
to \(w_j\) stays in the fixed
per-coordinate ball and hence in
the physical guard. Its group path
length is at most
\(K\|y_j-z_j\|_{\rm HS}\), by the
uniform local exponential and logarithm
bounds. Change the exterior inputs
one at a time along these paths.
All other coarse outputs remain the
same fixed \(C\), with their pivots
reconstructed consistently.
Integrating (37.8), absorbing the
fixed chart factor in the constants,
and taking the supremum (37.2)
proves (37.7).
\end{proof}

For any fixed region shape, the
coefficient \(a_R(\beta)\) is bounded
independently of \(\beta\) once the
coupling is sufficiently large.
For growing \(R\) its stated error
must still be checked.
The estimate does not use a claim
that local observables factorize.

\subsection{An explicit finite block-update contraction test}

For clarity consider a finite product
space with the sum metric \(d_I\),
and complete conditional updates
\(\Gamma_a\) on blocks \(B_a\).
Give block \(a\) probability \(p_a>0\),
\(\sum_a p_a=1\), and set
\(T=\sum_a p_a\Gamma_a\).
Let
\[
 u_i=\sum_{a:i\in B_a}p_a,\qquad
 v_i=\sum_{a:i\notin B_a}p_a c_{a i}.
 \tag{37.9}
\]
The required constants \(c_{a i}\)
satisfy, for \emph{every} pair of
configurations \(x,z\) differing
only at site \(i\notin B_a\),
\[
 \mathsf W_{B_a}(\Gamma_a^x,\Gamma_a^z)
                       \le c_{a i}d_i(x_i,z_i).
 \tag{37.10}
\]
They may be zero for an exterior
site absent from the local input
set. These are all-boundary constants.

\begin{proposition}[Finite-volume contraction criterion]
\label{prop:f12v33-criterion}
If \(u_{\min}=\min_i u_i>0\) and
\(v_i\le(1-\theta)u_i\) for every
site and some \(0<\theta\le1\), then
\[
 \mathsf W_I(T^x,T^z)
             \le(1-\theta u_{\min})d_I(x,z).
 \tag{37.11}
\]
In particular \(T\) contracts
Lipschitz expectation differences
under iteration and has at most
one invariant probability measure.
\end{proposition}
\begin{proof}
First let \(x,z\) differ only at \(i\).
For a block containing \(i\), their
exterior data are identical.
Use identical updated samples:
the discrepancy is removed completely.
For a block not containing \(i\),
the old discrepancy stays, and an
optimal block coupling adds at most
\(c_{a i}d_i(x_i,z_i)\) in expected
sum distance. Couple the block choice
identically in the two updates.
The expected distance is at most
\[
           (1-u_i+v_i)d_i(x_i,z_i)
                 \le(1-\theta u_{\min})d_i(x_i,z_i).
 \]
For general configurations, change
their finitely many coordinates
one at a time and apply the triangle
inequality for Wasserstein distance.
This gives (37.11). Its dual form
contracts the Lipschitz seminorm
of \(TF\). Iterating and integrating
against two invariant measures gives
an expectation difference bounded by
\((1-\theta u_{\min})^n\operatorname{Lip}(F)
\operatorname{diam}(I)\), which tends
to zero. Lipschitz functions determine
probability measures on this compact
metric space.
\end{proof}

The update is the full conditional
kernel on \(B_a\). Sampling only its
core while leaving a previously
integrated buffer unchanged is not
the update in this proof.

\subsection{The actual translated-block counts}

Take all translates of a coarse cube
of side length \(\ell\), with free
update set \(I_R\), on a sufficiently
large periodic lattice. No translation
invariance of \(C\) is required for
the following uniform bounds.
Each internal free coordinate is
updated by exactly \(\ell^4\)
translates. Each nonpivot boundary
coordinate is updated when either
endpoint of its coarse edge belongs
to the cube, hence by exactly
\(\ell^4+\ell^3\) translates.
There are seventeen internal and
twenty-eight boundary free types
per coarse cell, giving the check
\[
       |I_R|=17\ell^4+28(\ell^4+\ell^3)
                              =45\ell^4+28\ell^3.
 \tag{37.12}
\]
The second coverage count is the
union of two one-cell-shifted cubes:
one extra face of \(\ell^3\)
translation positions is added.

If a coordinate \(i\) lies in
\(J_R^\partial\), it is outside
\(I_R\), while its owner cell
is within four coarse steps of \(R\)
by the completed-collar geometry.
Thus the number of translates for
which \(i\) is an exterior input
is at most
\[
       n_{\rm in}(\ell)
               =(\ell+8)^4-\ell^4
               \le6560\ell^3\quad(\ell\ge1).
 \tag{37.13}
\]
The larger coordinate cube bounds
the graph neighbourhood. This
overcounts possible inputs and
is used only as an upper bound.

If (37.7) were valid for all compact
exteriors, the normalized incoming
coefficient in (37.9) would therefore
satisfy
\[
       \frac{v_i}{u_i}
            \le a_R(\beta)
                  \frac{(\ell+8)^4-\ell^4}{\ell^4}.
 \tag{37.14}
\]
For a fixed sufficiently large
\(\ell\), its good-boundary value
is less than one at sufficiently
large \(\beta\): choose \(\ell\)
to control the leading \(K_T/\ell\)
term, then absorb
\((\beta+1)\ell^8e^{-c_0\beta}\).
This is a verified calculation of
the candidate coefficient, not a
verification of (37.10).
The theorem (37.7) currently applies
only when the relevant exterior
inputs are in \(\mathcal G_R\).

There is an all-boundary bound, but
it is much weaker:
\[
 \mathsf W_{I_R}(\pi_{R,C}^{y},\pi_{R,C}^{z})
       \le K(\beta+1)s\,\|y_i-z_i\|_{\rm HS}
 \tag{37.15}
\]
when the data differ at one input
\(i\). To prove it, normalize a
one-Lipschitz test to have supremum
at most \(73d_Gs\), use the global
score bound and the elementary
covariance estimate in the exact
derivative, then integrate locally.
Compactness extends the local
Lipschitz estimate to the chord
metric globally: nearby pairs use
uniform charts; separated pairs
use the diameter bound on the
expectation difference.
The resulting upper bound on
(37.14) is of order
\((\beta+1)\ell^3\), which supplies
no contraction. This does not prove
that a sharper all-boundary coefficient
fails; it identifies what is still
uncontrolled.

\subsection{Why a globally good state space is not an invariant substitute}

The good-boundary hypothesis cannot
be enforced merely by restricting
the state space in the contraction
argument. The full compact kernel
does not preserve that restriction.
Here is a pointwise statement for
the actual smooth fibre.

\begin{proposition}[Positive defect probability at every finite coupling]
\label{prop:f12v33-positive}
Fix a sufficiently small free-coordinate
guard with positive-Haar-measure
complement. There are constants
\(a>0,K_0,K_1<\infty\), uniform in
site, exterior, retained field and
volume, such that its complement
at a free site \(i\) obeys
\[
 \pi_{\{i\},C}^{y}(\text{site }i\text{ is bad})
                  \ge p_-(\beta):=a e^{-K_0\beta-K_1}>0 .
 \tag{37.16}
\]
For every finite set \(A\) of free
sites under a compatible fibre law,
\[
 \mathbb P(\text{every site of }A\text{ is good})
                                  \le(1-p_-(\beta))^{|A|}.
 \tag{37.17}
\]
Every infinite-volume DLR state has
infinitely many such bad sites almost
surely at fixed finite \(\beta\).
\end{proposition}
\begin{proof}
The one-free-site conditional is a
positive smooth density relative to
normalized Haar. Its potential
contains only a uniformly bounded
number of plaquette and inverse-pivot
terms. The plaquette action is
bounded on the compact group and
the inverse-pivot Jacobians have
uniform positive upper and lower
bounds. Its oscillation in the
single free variable is therefore
at most \(K_0\beta+K_1\).
The normalized density is bounded
below by \(e^{-K_0\beta-K_1}\).
The bad set has fixed Haar measure
at least \(a>0\), independent of
the section value by Haar invariance.
This proves (37.16). This single-site
conditional is defined from the
global free-product density; no
regional coercivity claim for a
one-coordinate update is needed.

Remove the sites of \(A\) one at
a time. The event that all other
listed sites are good is measurable
outside the site being removed.
Conditioning on that exterior and
using (37.16) yields a factor at
most \(1-p_-\) each time.
This proves (37.17) without any
independence assumption.
Along an infinite deterministic
enumeration of sites, the probability
that all sites after any fixed
index are good is zero by (37.17).
Union over that index excludes
finitely many bad sites.
\end{proof}

This is consistent with V31's
exponentially small upper probabilities
and finite bad components: there may
be infinitely many sparse finite
components. It rules out treating
the globally all-good chart as an
invariant support of the actual
compact Gibbs law.

Nor can a small stationary probability
of a bad boundary simply be multiplied
into a worst-case transport influence.
For a discrepancy random variable \(D\)
and bad event \(E\), a bound on
\(\mathbb P(E)\) gives an additive
bound on \(\mathbb E[D{\bf1}_E]\)
when \(D\) is bounded, but not
\(\mathbb E[D{\bf1}_E]\le
\mathbb P(E)\mathbb E D\).
Already \(D={\bf1}_E\) disproves
that implication when
\(0<\mathbb P(E)<1\).
This is a logical warning, not a
counterexample to the YM mixing
claim: a suitable coupled estimate
may still be provable.

\begin{decision}[The next missing estimate has an explicit norm]
The good-boundary coefficient now
controls the complete updated block
in the correct sum transport metric.
The remaining finite-size step is
not the box's surface-to-volume
count. It is either a sufficiently
sharp all-boundary version of (37.7),
or a proved coupling scheme that
tracks discrepancy together with
bad-boundary energy and closes a
multiplicative contraction.

The latter would need a verified
conditional estimate for the coupled
bad-boundary term, not only the
single-law probabilities of V31.
An update cover, its resampling
domains and its invariant marginals
must be specified before iterating.
The full objective remains continuum
YM and its physical mass gap;
even fixed-fibre uniqueness would
still be only an intermediate result.
\end{decision}

\begin{record}[V33 append record]
V32 had 497835 bytes and SHA--256
\begin{center}\scriptsize\ttfamily
1E33E88405181880C2D21FC438F7F3165EA6BD8B367D455E7E7B1D120D60914D.
\end{center}
Its complete body is retained apart
from the title version.
The new actual-kernel result is
the whole-block transport estimate
with its compact correction.
The finite contraction lemma and
translated-block counts isolate its
remaining hypothesis, and the positive
defect bound excludes an invariant
all-good shortcut.
Only active V32 is replaced by V33.
All archives, scripts and companions
are retained. Next companion audit:
V35; next direction review: V40.
\end{record}

\begin{firewall}[Position after V33]
Good-boundary whole-block transport
and its volume scaling are proved.
The unrestricted block contraction
criterion is not verified.
Fixed-coupling uniqueness and
all-distance clustering, scale
iteration, continuum construction
and the physical Yang--Mills mass
gap remain unproved.
\end{firewall}

\section{V34: reversible adaptive updates and the coupled boundary remainder}
\label{sec:f12v34}

V33 identifies two distinct requirements:
whole-block transport when the boundary
is good, and a valid contraction mechanism
when it is not. We now construct an
adaptive update that preserves the actual
compact Gibbs law. Its choice of region
uses only exterior information for that
region. Failure of every guard triggers
an unconditional core update, not deletion
of the event or an identity update.

The construction gives a precise coupled
drift formula. Its remaining terms are
mixed discrepancy--defect quantities.
They are not replaced by products of
single-law probabilities.

\subsection{Nested regions and smooth exterior selectors}

Fix the retained field \(C\). Choose
finite nested coarse regions
\[
 R_1\supset R_2\supset\cdots\supset R_m,
 \qquad B_j=I_{R_j},\qquad J_j=J_{R_j}^\partial,
 \qquad B_0=B_m,\quad J_0=J_m.
 \tag{38.1}
\]
For example, take concentric cubes whose
radii decrease by four coarse cells.
Then \(N_3(R_{j+1})\subset R_j\).
Since \(J_{j+1}\subset\mathcal F_{R_{j+1}}
\subset\Lambda_{N_3(R_{j+1})}\), each free
coordinate of \(J_{j+1}\) belongs to
\(B_j\). But \(J_j\cap B_j=\varnothing\).
It follows that the \(J_j\)'s are pairwise
disjoint. These are deterministic regions;
no event involving unrevealed core values
is used to define them.

Take \(\tau>0\) so small that a free
section distance below \(\tau\) implies
membership in V33's per-coordinate
good boundary chart.
Let \(\varphi:[0,\infty)\to[0,1]\)
be smooth, equal to one on \([0,1/4]\)
and zero on \([1,\infty)\). For each
free coordinate put
\[
 g_i(Y_i)=\varphi\left(
       \frac{\|Y_i-J_i(C)\|_{\rm HS}^2}{\tau^2}\right),
 \qquad
 \chi_j(Y)=\prod_{i\in J_j}g_i(Y_i).
 \tag{38.2}
\]
An empty product is one.
Thus \(\chi_j\) is exterior-measurable
for \(B_j\); when it is positive,
the boundary is in the guard needed
for V33's compact transport estimate.
It is exactly one if every input
distance is at most \(\tau/2\).

Define first-acceptance weights
\[
 w_j=\chi_j\prod_{k<j}(1-\chi_k)
       \quad(1\le j\le m),\qquad
 w_0=\prod_{k=1}^m(1-\chi_k).
 \tag{38.3}
\]
They are nonnegative and sum to one.
Operationally, scan from the outer
region inward, accepting \(R_j\) with
probability \(\chi_j\) if earlier
choices were rejected. With label
zero, update the innermost \(R_m\)
unconditionally. The last success
and fallback use the same region;
their labels distinguish which
transport estimate is available.

Let \(\Gamma_j\) denote the complete
canonical heat-bath update on \(B_j\),
leaving its exterior unchanged, and
set \(\Gamma_0=\Gamma_m\). The adaptive
kernel is
\[
                      K=\sum_{a=0}^m w_a\Gamma_a .
 \tag{38.4}
\]
There is no change to the Wilson
weight, inverse-pivot factors, or
conditional normalizers.

\begin{theorem}[Reversibility of the actual adaptive update]
\label{thm:f12v34-invariance}
For every Gibbs law \(\mu\) compatible
with these canonical kernels at \(C\),
the kernel \(K\) is \(\mu\)-reversible
and preserves \(\mu\). For bounded
real \(f\),
\[
 \langle f,(I-K)f\rangle_\mu
   =\sum_{a=0}^m
      \mathbb E_\mu\left[
       w_a\,\operatorname{Var}_\mu
                 (f\mid Y_{B_a^c})\right].
 \tag{38.5}
\]
Every possible selected update contains
the innermost core \(B_m\).
\end{theorem}
\begin{proof}
For \(k<j\), exterior measurability
for \(B_k\) implies exterior measurability
for \(B_j\), since \(B_j\subset B_k\).
Hence \(w_j\) is measurable outside
\(B_j\). Similarly \(w_0\) is measurable
outside \(B_m=B_0\).
Under a compatible Gibbs law,
\(\Gamma_a\) is conditional expectation
onto that exterior sigma field, so
it is an orthogonal projection on
\(L^2(\mu)\). Multiplication by \(w_a\)
commutes with it. Each
\(w_a\Gamma_a\) is therefore self-adjoint.
Their sum is self-adjoint and \(K1=1\),
which proves reversibility and
invariance.

Moreover
\(\langle f,w_a(I-\Gamma_a)f\rangle
 =\mathbb E[w_a(\Gamma_a f^2-(\Gamma_a f)^2)]\).
Sum over \(a\), using \(\sum_a w_a=1\),
to obtain (38.5).
The final statement follows directly
from the nested regions and the
chosen fallback.
\end{proof}

This proof applies equally to a finite
periodic fibre law and to each
infinite-volume DLR state. It does
not establish a spectral gap for
this adaptive kernel. In particular,
invariance alone is not a contraction
estimate.

\subsection{The stationary fallback weight without an independence assumption}

Assume the fixed small-curvature
hypotheses of V33 and of V31 at
threshold \(\tau/2\), and take the
coupling above both required thresholds.
Put
\[
 q_j=|J_j|,\qquad
                   p=e^{-\beta\tau^2/(64A)}.
 \tag{38.6}
\]
Under any such DLR state, or full
finite-volume fibre law, V31 gives
the joint all-bad bound \(p^{|A|}\)
at the specified free sites \(A\).
Consequently
\[
                  \mathbb E_\mu w_0
                               \le p^m\prod_{j=1}^m q_j.
 \tag{38.7}
\]
If some \(q_j=0\), both the fallback
weight and the displayed right side
are zero.

\begin{proof}
If \(w_0>0\), every \(\chi_j<1\).
For each nonempty \(J_j\), this
requires at least one coordinate
whose distance is greater than
\(\tau/2\). Choose one witness in
each \(J_j\). The disjointness of
the input sets makes these \(m\)
distinct coordinates. V31 bounds
their simultaneous bad event by
\(p^m\). Union over the
\(\prod_j q_j\) choices and use
\(0\le w_0\le1\).
No independence of the collars,
or of the acceptance coins after
averaging over the field, was used.
\end{proof}

The bound concerns the fallback
weight in stationarity, not a
uniform bound at every configuration.
The fallback itself remains a
full compact conditional update
even when all \(\chi_j\)'s vanish.

\subsection{Local control of changes in the selected region}

For a pair of configurations \(x,z\),
write
\[
 d_i=\|x_i-z_i\|_{\rm HS},\qquad
 \mathcal B_i(x,z)=
 {\bf1}_{\{\|x_i-J_i(C)\|_{\rm HS}>\tau/2
             \ \text{or}\ 
             \|z_i-J_i(C)\|_{\rm HS}>\tau/2\}}.
 \tag{38.8}
\]
There is a fixed \(L_\varphi\) such that
\[
 |g_i(x_i)-g_i(z_i)|
             \le\frac{L_\varphi}{\tau}
                                   d_i\mathcal B_i(x,z).
 \tag{38.9}
\]
Indeed the cutoff has Lipschitz
constant \(L_\varphi/\tau\) in group
distance; if the indicator is zero,
both cutoff values are exactly one.
Telescoping products then gives
\[
 |\chi_j(x)-\chi_j(z)|
          \le\frac{L_\varphi}{\tau}
                        \sum_{i\in J_j}d_i\mathcal B_i.
 \tag{38.10}
\]

Let \(\mathsf{TV}(w(x),w(z))
=\frac12\sum_{a=0}^m|w_a(x)-w_a(z)|\).
The first-acceptance distributions obey
\[
 \begin{split}
 \mathsf{TV}(w(x),w(z))
       &\le\sum_{j=1}^m|\chi_j(x)-\chi_j(z)|\\
       &\le\frac{L_\varphi}{\tau}
                   \sum_{i\in J_*}d_i\mathcal B_i,
       \qquad J_*=\bigcup_jJ_j .
 \end{split}
 \tag{38.11}
\]
For the first inequality, couple
the acceptance tests by the same
independent uniform numbers at
each stage. Different selected
labels require at least one
different acceptance test, whose
probability is at most the sum
of the parameter differences.
The minimal mismatch probability
is no larger. The second inequality
uses the disjoint input sets.

Smooth selectors thus localize their
variation to discrepancies touching
the cutoff transition or bad region.
The factor \(d_i\mathcal B_i\) is
important: it vanishes when the two
configurations coincide, even at a
bad coordinate.

\subsection{A coupled step with correct invariant marginals}

For \(j\ge1\), let \(a_j=a_{R_j}(\beta)\)
be V33's good-boundary whole-block
transport coefficient. The unconditional
core estimate (37.15), telescoped over
its exterior inputs, gives a coefficient
\[
                         a_0=K(\beta+1)|R_m|.
 \tag{38.12}
\]
It applies to \(\Gamma_0=\Gamma_m\)
without a boundary guard. Put
\[
 t_a(x,z)=\min\{w_a(x),w_a(z)\},\qquad
 D_{\max}=d_G\max_{0\le a\le m}|B_a|.
 \tag{38.13}
\]
For a coordinate set \(A\), abbreviate
\(d_A=\sum_{i\in A}d_i\).

\begin{theorem}[One-step adaptive coupling inequality]
\label{thm:f12v34-drift}
For every fixed pair \(x,z\), there
is a coupling of \(K^x,K^z\) such
that, in a finite ambient product,
\[
 \begin{split}
 \mathbb E[d_I(X',Z')\mid x,z]\le d_I(x,z)
 &-\sum_{a=0}^m t_a d_{B_a}
   +\sum_{j=1}^m t_j a_j d_{J_j}
   +t_0a_0d_{J_m}\\
 &+\sum_{a=0}^m
       |w_a(x)-w_a(z)|\,d_G|B_a|.
 \end{split}
 \tag{38.14}
\]
In particular,
\[
 \begin{split}
 \mathbb E[d_I(X',Z')-d_I(x,z)\mid x,z]
 \le{}&-d_{B_m}
       +\sum_{j=1}^m t_j a_j d_{J_j}
       +t_0a_0d_{J_m}\\
 &+\frac{3L_\varphi D_{\max}}{\tau}
                              \sum_{i\in J_*}d_i\mathcal B_i .
 \end{split}
 \tag{38.15}
\]
On an infinite lattice the same
statement is understood as the
finite sum of distance increments
on \(B_1\); coordinates outside
\(B_1\) are unchanged.
\end{theorem}
\begin{proof}
Couple the selected labels with
common mass \(t_a\) on each equal
pair \(a,a\). Couple the remaining
label masses arbitrarily. This
has the required marginal label
distributions.

For a matched label \(j\ge1\),
both weights are positive, so
both exteriors lie in the good
guard for \(R_j\). Use a block
transport coupling from (37.7).
It removes the old discrepancy
on \(B_j\), retains the discrepancy
outside, and creates expected
block distance at most \(a_jd_{J_j}\).
For matched label zero use the
unconditional bound \(a_0d_{J_m}\).
For unmatched labels \(a,b\),
the triangle inequality bounds
the new distance by the old one
plus \(d_G|B_a|+d_G|B_b|\).
Summing the residual row and
column masses gives the last
term of (38.14). All conditional
samples have their correct complete
block marginals.

Let \(v=\mathsf{TV}(w(x),w(z))\).
Since every \(B_a\) contains \(B_m\)
and \(\sum_a t_a=1-v\), the erased
distance is at least
\((1-v)d_{B_m}\).
Also \(d_{B_m}\le D_{\max}\) and
the unmatched term in (38.14)
is at most \(2D_{\max}v\).
Use (38.11) to bound the combined
loss \(3D_{\max}v\), proving (38.15).

The couplings can be implemented
measurably with arbitrarily small
additional error: partition each
compact block into sets of small
diameter, choose finite transportation
optimizers with a fixed tie-breaking
rule, and lift them using the
conditional laws inside the partition
sets. The explicit positive densities
make these kernels measurable.
The cost error tends uniformly to
zero with the partition diameter.
After integration against any
current coupling, compactness gives
the stated bound in the limit.
This avoids assuming an unproved
measurable choice of pointwise
optimal plans.
\end{proof}

For any coupling \(\lambda\) of two
compatible Gibbs laws \(\mu,\nu\),
the integrated coupled step has
marginals \(\mu K=\mu\) and
\(\nu K=\nu\), by
Theorem~\ref{thm:f12v34-invariance}.
Thus the step remains in the
class of couplings of the original
laws. Conditioning the selection
on both current configurations in
the coupling does not alter either
marginal update: its own label
distribution is still \(w(x)\)
or \(w(z)\), respectively.

\subsection{The good drift and the two quantities still to close}

The first term in (38.15) is a
deterministic core erasure.
For translations of nested cubes
with side lengths \(\ell_j\),
use V33's coverage counts.
Each site is in at least
\(\ell_m^4\) core translates,
and is an input to at most
\((\ell_j+8)^4-\ell_j^4\)
translates of the \(j\)-th region.
Using \(t_j\le1\), the good inflow
relative to core erasure is bounded by
\[
 \Theta_{\rm good}
   =\frac{\displaystyle\sum_{j=1}^m
          a_j\bigl((\ell_j+8)^4-\ell_j^4\bigr)}
                  {\ell_m^4}.
 \tag{38.16}
\]
For fixed \(m\), choose a sufficiently
large core side, with
\(\ell_j=\ell_m+8(m-j)\), and then
take \(\beta\) sufficiently large.
Equation (37.7) makes all \(a_j\)
bounded, so \(\Theta_{\rm good}<1/2\)
is attainable. This controls the
good term for arbitrary current
configurations in the coupled rule:
its use of good transport is already
enforced by the matched positive
selection weights.

What remains after this count is
not an unidentified change of
measure. For a current coupling
\(\lambda\), define
\[
 \mathfrak m_i(\lambda)
     =\mathbb E_\lambda[d_i\mathcal B_i],\qquad
 \mathfrak f(\lambda)
     =\mathbb E_\lambda[t_0d_{J_m}].
 \tag{38.17}
\]
They contribute precisely
\[
 a_0\mathfrak f(\lambda)
   +\frac{3L_\varphi D_{\max}}{\tau}
                          \sum_{i\in J_*}\mathfrak m_i(\lambda)
 \tag{38.18}
\]
for this translated update.
The first is the cost of matched
unconditional fallback; the second
includes both selector mismatch
and the corresponding loss of
common core erasure.

Single-law estimates give the valid
absolute bounds
\[
 \mathfrak m_i(\lambda)\le2d_Gp,\qquad
 \mathfrak f(\lambda)
        \le d_Gq_m\,p^m\prod_{j=1}^m q_j.
 \tag{38.19}
\]
For the first, the probability
that either marginal is bad is
at most \(2p\); for the second,
\(t_0\le w_0(x)\), \(d_{J_m}\le d_Gq_m\),
and (38.7) applies to that marginal.
These statements hold for every
coupling with the specified marginals.

They are not relative estimates
in the discrepancy norm. For example,
\(\mathfrak m_i\le
2p\,\mathbb E_\lambda d_i\) does
not follow from them.
Such a bound requires information
about the joint coupling, not
just the two one-field laws.
Using (38.19) in the drift gives
an additive error floor, not a
proof that the discrepancy converges
to zero.

\begin{record}[Executed exact finite-state algebra test]
The script
\begin{center}\small\ttfamily
scripts/verify\_ym\_f12\_v34\_adaptive\_kernel.py
\end{center}
has SHA--256
\begin{center}\scriptsize\ttfamily
CC69143555557325DE9083AAAB491936E4DA46D8268DAF716F0751A010C937CF.
\end{center}
It uses a nonproduct, strictly
positive rational law on sixteen
four-spin configurations, three
nested updates, and the core fallback.
All row sums and all 256 detailed
balance pairs are checked exactly.
Both sides of the Dirichlet identity
equal
\[
                   \frac{1028707081341}{2114443217320}.
 \]
Changing only the core coordinate
gives identical update rows, as
exterior measurability requires.
The first-acceptance selector test
has \(L^1\) difference \(1\), below
its bound \(11/6\).
An intentionally non-exterior selector
has a nonzero stationarity defect
\(155952/5182967\), confirming why
that requirement cannot be omitted.
The finite test is not a YM model
or a contraction certificate;
the Gibbs-kernel proofs above
are independent of its numerical
values.
\end{record}

\begin{decision}[Next: a joint discrepancy estimate, not another marginal tail]
We now have an invariant adaptive
kernel, a coupling with the right
marginals, a controlled good drift,
and an explicit fallback.
To turn it into a contraction,
the next proof must control (38.18)
relative to the erased discrepancy,
possibly in an enlarged discrepancy
norm that tracks boundary energy.
Its evolution must be computed
under this actual coupled update.

An alternative all-boundary transport
estimate would also suffice, but
neither alternative has been obtained
here. Replacing \(\mathfrak m_i\)
or \(\mathfrak f\) by a product
of marginal estimates would be
circular. At V35 perform the
required companion audit before
choosing the next estimate.
\end{decision}

\begin{record}[V34 append record]
V33 had 515193 bytes and SHA--256
\begin{center}\scriptsize\ttfamily
A0FA96750CD4704A626B6ACAED5526B2B6A4AD8FB00E4466897148C782297BFD.
\end{center}
Its complete body is retained
apart from the title version.
The new work specifies an
exterior-measurable adaptive update
with an unconditional core fallback
and derives its exact coupled
remainder terms.
Only active V33 is replaced by
V34; all archives and companion
files remain.
Next companion audit: V35.
Next direction review: V40.
\end{record}

\begin{firewall}[Position after V34]
The adaptive update preserves
the full compact Gibbs laws,
and its coupling has a genuine
core-erasure term.
The mixed discrepancy--defect
closure, fixed-coupling uniqueness,
all-distance clustering, scale
iteration, continuum construction
and the physical Yang--Mills mass
gap remain unproved.
\end{firewall}

\section{V35: audit and relative creation bounds for bad-core discrepancy}
\label{sec:f12v35}

\subsection{Companion checkpoint}

The latest companion files were inspected
again. They are unchanged from V30:
\begin{center}\small
\begin{tabular}{lrr}
\toprule
Companion & Version & Bytes\\
\midrule
SM--Einstein & 72 & 607372\\
NS stability & 26 & 278283\\
Parallel YM & 5 & 54174\\
\bottomrule
\end{tabular}
\end{center}
\begin{center}\scriptsize\ttfamily
SM: F44F9745D43379E1672C5550411B58E97195A4C9278592D221DF5D3F644DC1F5\\
NS: 82C887F8C2C69E4F28FAD7C3DC8DE5B9099CB5A4BF431EBA1FFF3E91935978AD\\
YM: 306F1BB84CFA8A8D47EA569EC17E9545F9867C8D32B548EC9D9C444B4F3C0512
\end{center}
The SM result remains a declared
fixed-mode many-copy limit; NS V26
contains its rank-one derivative
formulas and an uncompleted continuation;
parallel YM retains the obstruction
to volume-uniform extensive-charge
tightness. None supplies a joint
coupling estimate for the present
compact fibre. Their conclusions
do not change the current direction.

V34's remainder concerns discrepancies
at bad sites, not merely the probability
that a site is bad. We now obtain a
relative creation estimate for such
discrepancies in a buffered core.
The proof uses one weighted transport
cost and one coupling; it does not
combine incompatible optimal couplings
for two separate costs.

\subsection{A compact metric detecting bad-site discrepancy}

Keep the fixed threshold \(\tau\)
used in V34's selectors.
Write \(\varrho=\|I\|_{\rm HS}\)
in the fixed unitary representation.
Choose a smooth \(\psi:[0,\infty)\to[0,1]\)
equal to zero on \([0,1/16]\)
and one on \([1/4,\infty)\), and put
\[
 h_i(g)=\psi\left(
             \frac{\|g-J_i(C)\|_{\rm HS}^2}{\tau^2}\right).
 \tag{39.1}
\]
Thus \(h_i=0\) below distance
\(\tau/4\), and \(h_i=1\) at
distance at least \(\tau/2\).
For two values define
\[
 c_i(g,k)=
       \varrho|h_i(g)-h_i(k)|
       +\|h_i(g)g-h_i(k)k\|_{\rm HS}.
 \tag{39.2}
\]
This is a continuous pseudometric:
it is the distance between the
two features
\((\varrho h_i(g),h_i(g)g)\)
in the direct-sum norm
\(|a|+\|M\|_{\rm HS}\).
It has diameter at most \(3\varrho\).

In the notation of (38.8),
\[
 \|g-k\|_{\rm HS}\,
 {\bf1}_{\{\|g-J_i(C)\|>\tau/2
                   \ {\rm or}\ \|k-J_i(C)\|>\tau/2\}}
                           \le c_i(g,k).
 \tag{39.3}
\]
For example, if \(h_i(g)=1\) and
\(h_i(k)=a\), then
\[
 \|g-k\|_{\rm HS}
       \le\|g-ak\|_{\rm HS}
                     +(1-a)\varrho=c_i(g,k).
 \]
The other case is symmetric.
Both sides of (39.3) vanish as
needed when the values coincide,
including at a bad site.

For a block \(B=I_R\), a nonempty
core \(U\subset B\), and \(\zeta\ge0\),
define
\[
 c_U(Y,Z)=\sum_{i\in U}c_i(Y_i,Z_i),\qquad
 d_B^{\zeta,U}(Y,Z)=d_B(Y,Z)+\zeta c_U(Y,Z).
 \tag{39.4}
\]
The second expression is an actual
compact metric because \(d_B\) is.
It is bounded above by a fixed
multiple, depending on \(\tau,\zeta\),
of \(d_B\). We will use its
Kantorovich duality just as in V33.

\subsection{A Lipschitz extension which isolates a rare remainder}

Let
\[
 \mathcal Z_U=\{Y:h_i(Y_i)=0
                                  \text{ for every }i\in U\}.
 \tag{39.5}
\]
This compact nonempty set restricts
only core features; all other block
variables remain unrestricted.
For any one-Lipschitz \(F\) in
the metric \(d_B^{\zeta,U}\), define
\[
 F_0(Y)=\inf_{Z\in\mathcal Z_U}
                           \{F(Z)+d_B(Y,Z)\},\qquad
                         R_F=F-F_0 .
 \tag{39.6}
\]

\begin{proposition}[A rare-support extension remainder]
\label{prop:f12v35-extension}
The function \(F_0\) is one-Lipschitz
in the unweighted metric \(d_B\),
equals \(F\) on \(\mathcal Z_U\),
and, with \(u=|U|\),
\[
 |R_F(Y)|\le
       2(d_G+\varrho\zeta)u\,
                               {\bf1}_{\mathcal Z_U^c}(Y).
 \tag{39.7}
\]
\end{proposition}
\begin{proof}
On \(\mathcal Z_U\), the added
pseudometric is zero between any
two points. Thus the restriction
of \(F\) there is one-Lipschitz
in \(d_B\). The infimum extension
is one-Lipschitz by the triangle
inequality and equals \(F\) on
that set.

For any \(Z\in\mathcal Z_U\),
\[
 c_U(Y,Z)=2\varrho\sum_{i\in U}h_i(Y_i)
                                             \le2\varrho u,
 \]
independently of the remaining
values of \(Z\). The weighted
Lipschitz inequality therefore
gives \(F(Y)\le F_0(Y)+2\varrho\zeta u\).
For the reverse bound, replace
only \(Y_U\) by its section values
to obtain \(Z(Y)\in\mathcal Z_U\).
Its unweighted distance is at
most \(d_Gu\). Hence
\[
 \begin{split}
 F_0(Y)&\le F(Z(Y))+d_B(Y,Z(Y))\\
       &\le F(Y)+2d_Gu+2\varrho\zeta u.
 \end{split}
 \]
The remainder vanishes on
\(\mathcal Z_U\); the two bounds
prove (39.7).
\end{proof}

No derivative of the transport test
has been assumed in this argument.
In particular no smoothing operation
is required to preserve the weighted
Lipschitz class.

\subsection{A uniform conditional rare-core input}

Require the V31 neighbourhood of
the owner cells of \(U\), at threshold
\(\tau/4\), to lie inside \(R\):
\[
      N_{hJ(\tau/4)}(S(U))\subset R .
 \tag{39.8}
\]
Assume the affected coarse curvatures
for \(R\) are at most
\(\epsilon_{\tau/4}\) and take
\(\beta\ge\beta_{\tau/4}\).
The V31 canonical conditional bound,
disintegrated within \(R\) when
necessary, gives for every compact
exterior \(y\)
\[
 \pi_{R,C}^{y}\{\|Y_i-J_i(C)\|_{\rm HS}\ge\tau/4\}
       \le p_*:=e^{-\beta\tau^2/(256A)}
                  \quad(i\in U).
 \tag{39.9}
\]
Consequently
\[
                   \pi_{R,C}^{y}(\mathcal Z_U^c)\le u p_*.
 \tag{39.10}
\]
These are pointwise conditional
bounds for the fully integrated
buffer, not assumptions that the
exterior is itself good.
The same bounds hold along every
exterior variation at fixed \(C\).

\begin{proposition}[The rare remainder has a relative response bound]
\label{prop:f12v35-response}
For two exterior data \(y,z\), let
\(d_J(y,z)=\sum_{j\in J_R^\partial}
\|y_j-z_j\|_{\rm HS}\).
There is a fixed constant \(K\)
such that
\[
 |\pi_{R,C}^{y}R_F-\pi_{R,C}^{z}R_F|
       \le (1+\zeta)b_{R,U}\,d_J(y,z),\qquad
 b_{R,U}=K(\beta+1)u^2p_* .
 \tag{39.11}
\]
\end{proposition}
\begin{proof}
The function \(R_F\) is a fixed
bounded function of the block
variables. Its definition depends
on the fixed \(C\), not on the
exterior free data.
The exact normalized score identity
therefore applies:
\[
 \partial_{y_j}\pi_{R,C}^{y}R_F
                  =-\operatorname{Cov}_{\pi_{R,C}^{y}}(R_F,S_j).
 \]
Equations (39.7)--(39.10) give
\(\pi^y|R_F|\le K(1+\zeta)u^2p_*\).
Since \(\|S_j\|_\infty\le K(\beta+1)\),
the covariance is bounded by
\(K(\beta+1)(1+\zeta)u^2p_*\).
This differentiates the normalized
kernel, not \(F\) or \(R_F\).

Integrate along local group paths
for nearby exterior values.
Uniform compact charts convert
their lengths to Hilbert--Schmidt
distance. For pairs separated
beyond that fixed chart radius,
bound the difference by the two
expectations of \(|R_F|\) and divide
by the positive distance lower
bound. This retains the factor
\(p_*\), including for disconnected
components if present.
Change the finitely many exterior
inputs one at a time and sum.
The conditional probability bound
is uniform over all intermediate
data. This proves (39.11).
\end{proof}

The factor \(d_J(y,z)\) in (39.11)
is the new point. The estimate is
zero for identical input data.
It is not a bound of order \(p_*\)
which survives when the discrepancy
has vanished.

\subsection{One coupling controls ordinary and bad-core discrepancy}

Let \(A_R(\beta)\) be a proved
unweighted transport coefficient:
\[
 \mathsf W_{I_R}(\pi_{R,C}^{y},\pi_{R,C}^{z})
                           \le A_R(\beta)d_J(y,z).
 \tag{39.12}
\]
For good-boundary endpoints use
\(A_R=a_R(\beta)\) from (37.7)
and its fixed small-curvature
hypotheses. For arbitrary endpoints
use \(A_R=K(\beta+1)|R|\) from
(37.15). In both cases also
impose the core conditions (39.8)--(39.9).

\begin{theorem}[Weighted core transport]
\label{thm:f12v35-weighted}
For every \(\zeta\ge0\),
\[
 \mathsf W_{d_{I_R}^{\zeta,U}}
                 (\pi_{R,C}^{y},\pi_{R,C}^{z})
       \le\bigl[A_R+(1+\zeta)b_{R,U}\bigr]d_J(y,z).
 \tag{39.13}
\]
Thus a single coupling of the
complete compact block laws
simultaneously controls the two
costs on the left.
\end{theorem}
\begin{proof}
Take any one-Lipschitz test for
the weighted metric and decompose
it as \(F_0+R_F\) using (39.6).
The unweighted transport bound
applies to \(F_0\), while (39.11)
applies to \(R_F\). Sum the two
estimates and take the supremum
over weighted one-Lipschitz tests.
Compact-metric transport duality
gives (39.13), and its attaining
coupling controls the sum of
the two costs. Separate optimizing
couplings have not been combined.
\end{proof}

Put \(b=b_{R,U}>0\), choose
\(\zeta=b^{-1/2}\), and set
\[
 \widehat A_R=A_R+b+\sqrt b,\qquad
                    \varepsilon_{R,U}=\widehat A_R\sqrt b.
 \tag{39.14}
\]
The same coupling then satisfies
\[
 \begin{split}
 \mathbb E\,d_{I_R}(Y,Z)
                   &\le\widehat A_R\,d_J(y,z),\\
 \mathbb E\sum_{i\in U}
        \|Y_i-Z_i\|_{\rm HS}\mathcal B_i(Y,Z)
                   &\le\varepsilon_{R,U}\,d_J(y,z).
 \end{split}
 \tag{39.15}
\]
The second follows from (39.3)
and the weighted cost, after
division by \(\zeta\).
For fixed block and core geometry,
\(\varepsilon_{R,U}\) is exponentially
small in \(\beta\), even for the
arbitrary-boundary coefficient
\(A_R=O((\beta+1)|R|)\).
Indeed \(b=K(\beta+1)u^2e^{-\beta\tau^2/(256A)}\);
all its other factors and those
in \(A_R\) are polynomial at
fixed geometry.

The weight may be large as
\(\beta\) grows, but every preceding
estimate holds for each finite
\(\zeta\). The ordinary transport
coefficient increases only by
\(b+\sqrt b\).

\subsection{Insert the couplings into the actual adaptive update}

Keep V34's regions and selectors,
and choose a common nonempty core
\(U\subset B_m\) satisfying (39.8)
already inside \(R_m\).
It then satisfies that condition
inside every selected region.
For label \(j\ge1\), use the
good-boundary coefficient in
(39.14); for fallback label zero,
use the unconditional coefficient.
Write the resulting numbers as
\(\widehat A_a,\varepsilon_a\).
Matched labels are coupled by
(39.15). The measurable
near-optimal construction of V34
applies to these continuous
weighted metrics too.
The marginal adaptive kernels
and their reversibility do not
change.

For unmatched labels, each output
law still has conditional bad-core
probability at most \(p_*\) at
each coordinate, by (39.9).
Therefore \emph{any} coupling of
that pair of output laws obeys
\[
 \mathbb E\sum_{i\in U}
       \|X_i'-Z_i'\|_{\rm HS}\mathcal B_i(X',Z')
                            \le2d_Gu p_*.
 \tag{39.16}
\]
This uses the union of the two
output bad events, not independence.

Define
\[
 M_U(x,z)=\sum_{i\in U}
                    \|x_i-z_i\|_{\rm HS}\mathcal B_i(x,z).
 \tag{39.17}
\]
Using \(t_a=\min(w_a(x),w_a(z))\),
the unmatched mass is
\(v=\mathsf{TV}(w(x),w(z))\).
Equations (38.11), (39.15), and
(39.16) now give the pointwise
relative creation estimate
\[
 \begin{split}
 \mathbb E[M_U(X',Z')\mid x,z]
 \le{}&\sum_{a=0}^m t_a\varepsilon_a d_{J_a}(x,z)\\
      &+\frac{2d_Gu p_*L_\varphi}{\tau}
                  \sum_{i\in J_*}d_i(x,z)\mathcal B_i(x,z).
 \end{split}
 \tag{39.18}
\]
All coefficients in this displayed
creation bound are exponentially
small at fixed geometry and
sufficiently large \(\beta\).
Subtracting \(M_U(x,z)\) gives
an actual local mixed-moment
drift inequality. In parallel,
the ordinary-distance drift
(38.14)--(38.15) remains valid
with the matched transport
coefficients replaced by
\(\widehat A_a\).
These two conclusions use the
\emph{same} coupled update.

In particular (39.18) is not
obtained by assuming
\(\mathbb E[d_i\mathcal B_i]\)
factors under an existing coupling.
The coupling is constructed using
the weighted metric, and its
conditional estimate is proportional
to the input discrepancy.

\subsection{A separated mixed-moment refresh}

There is one further elementary
relative estimate useful for a
larger discrepancy norm.
If \(i\notin B_1\) and \(k\in U\),
the coordinate discrepancy \(d_i\)
is unchanged by the update. Both
conditional output marginals at
\(k\) satisfy (39.9), whatever
labels are selected. Thus
\[
 \mathbb E[d_i(X',Z')\mathcal B_k(X',Z')\mid x,z]
                           \le2p_*\,d_i(x,z).
 \tag{39.19}
\]
The proof conditions on the
current pair, so \(d_i\) is
already a fixed multiplier,
and bounds the probability of
either output bad event.
This is a legitimate conditional
refresh, not a stationarity
factorization. It does not apply
when the discrepancy coordinate
\(i\) is also being changed.

\subsection{The unclosed part is now localized}

\begin{assessment}[What the new relative estimate does not cover]
The common core \(U\) must have
its V31 buffer \emph{inside}
the innermost update region.
Consequently (39.18) does not
control every changed coordinate.
The sets \(B_a\setminus U\) are
updated as well. Their bad-site
discrepancies are not unchanged
and cannot be omitted when
forming an all-lattice norm.

Likewise the ordinary-distance
fallback term still contains
\(t_0\widehat A_0d_{J_m}\), and
the selector term still contains
old discrepancies at bad collar
coordinates. The new estimate
controls creation in a buffered
core, not yet those entire
boundary contributions after
summing overlapping translates.
Equation (39.19) addresses only
separated discrepancy and badness
coordinates. Nearby mixed pairs
require additional bookkeeping.
\end{assessment}

\begin{decision}[Continue with the updated boundary layer]
The companion audit leaves the
direction unchanged. We now have
a relative bound for creation
of bad-core discrepancy, including
the unconditional fallback, and
it coexists with the ordinary
transport bound in one coupling.
The next task is to account for
the updated layer \(B_a\setminus U\)
and the overlapping translated
mixed terms in a single closed
discrepancy norm.

It would be invalid to sum only
the favourable core drifts and
pretend that the remaining updated
coordinates stayed fixed.
The next direction review is V40;
the next five-version companion
audit is also V40.
\end{decision}

\begin{record}[V35 append record]
V34 had 531282 bytes and SHA--256
\begin{center}\scriptsize\ttfamily
E944F2D4087A18FAFAF2046220ED1B1958D91530554366B5B178668CA337CB66.
\end{center}
Its complete body is retained
apart from the title version.
The new proof uses the actual
conditional rare-core estimate,
a metric that detects bad-site
discrepancy, a Lipschitz extension,
and the normalized score identity.
No independence of coupled bad
events is assumed.
Only active V34 is replaced by
V35; archives, scripts and
companions are retained.
\end{record}

\begin{firewall}[Position after V35]
Relative bad-core discrepancy
creation and ordinary transport
are simultaneously controlled.
The global mixed-moment contraction
is not closed. Fixed-coupling
uniqueness, all-distance clustering,
scale iteration, continuum
construction and the physical
Yang--Mills mass gap remain
unproved.
\end{firewall}

\section{V36: nested Dirichlet cancellation and a finite-scale variance bridge}
\label{sec:f12v36}

\subsection{Why move upstream}

V35 controls a buffered core, not every changed boundary
coordinate. V32's whole-block rarity estimate has its own
small exterior threshold. It is not legitimate to assume
that V34's selector threshold also satisfies the stronger
guard needed for rarity at that same threshold.
We leave the valid V35 coupling estimates intact and
now use the exact reversible forms of the actual kernels.
This avoids differentiating a selector or paying a
transport diameter for every change of selected label.

There are two new conclusions. First, the adaptive form
admits an exact nonnegative nested-projection decomposition.
Second, V32's uniform conditional total-variation diameter
implies a quantitative two-block variance inequality for
the full compact fibre law. Neither conclusion is a
physical Yang--Mills mass gap.

The general identification of maximal correlation with
a conditional-expectation norm was already recorded in
f6 V18. It is not claimed anew here. The additional input
below is the quantitative norm bound from V32's actual
conditional kernels, and its variance-factorization
consequence. The archive's physical shell channels must
not be identified with the auxiliary updates below.

\subsection{The last selector cancels from the physical kernel}

Work in \(L^2(\mu)\), for any compatible finite-volume
fibre law or DLR state. All conditional expectations
are taken under this same law at fixed \(C\).
Keep V34's nested \(B_1\supset\cdots\supset B_m\)
and write
\[
 \mathcal H_j=\sigma(Y_{B_j^c}),\qquad
 \Gamma_j=\mathbb E_\mu[\,\cdot\mid\mathcal H_j].
 \tag{40.1}
\]
Thus \(\mathcal H_j\subset\mathcal H_{j+1}\), and
\(\Gamma_j\Gamma_k=\Gamma_{\min(j,k)}\).
For \(j<m\) put
\[
 q_j=\sum_{k=1}^j w_k
       =1-\prod_{k=1}^j(1-\chi_k),\qquad r_j=1-q_j,
 \qquad D_j=\Gamma_{j+1}-\Gamma_j.
 \tag{40.2}
\]
Here \(q_j\) is a cumulative acceptance weight, not the
collar cardinality denoted by \(q_j\) in (38.6);
we write \(n_j=|J_j|\) for that cardinality in this section.
Every \(D_j\) is an orthogonal projection.
Multiplication by \(q_j\) commutes with both
\(\Gamma_j\) and \(\Gamma_{j+1}\), since \(q_j\)
is \(\mathcal H_j\)-measurable.

\begin{theorem}[Exact nested adaptive form]
\label{thm:f12v36-nested}
The physical kernel is independent of the last
selector \(\chi_m\). For \(m\ge1\), with empty sums
understood as zero,
\[
 \begin{split}
 K&=\sum_{j=1}^{m-1}w_j\Gamma_j+
            \left(\prod_{k<m}(1-\chi_k)\right)\Gamma_m,\\
 I-K&=(I-\Gamma_m)+\sum_{j=1}^{m-1}q_jD_j .
 \end{split}
 \tag{40.3}
\]
For every real \(f\in L^2(\mu)\),
\[
 \mathcal E_K(f):=\langle f,(I-K)f\rangle_\mu
 =\|f-\Gamma_m f\|_2^2+
                       \sum_{j=1}^{m-1}\mu[q_j(D_jf)^2].
 \tag{40.4}
\]
In quadratic-form order,
\[
 I-\Gamma_m\ \le\ I-K\ \le\ I-\Gamma_1 .
 \tag{40.5}
\]
\end{theorem}
\begin{proof}
The two labels \(m,0\) use the identical kernel
\(\Gamma_m\), and their weights sum to
\(\prod_{k<m}(1-\chi_k)\). This proves the first
line of (40.3), including cancellation of \(\chi_m\).
Collect the coefficient of each \(\Gamma_j\) in
\(\sum_{j<m}q_j(\Gamma_{j+1}-\Gamma_j)\).
The result is \(\Gamma_m-K\), proving the second line.

Nested conditional expectations give \(D_j^2=D_j=D_j^*\).
Because \(q_j\) commutes with \(D_j\),
\(\langle f,q_jD_jf\rangle=\mu[q_j(D_jf)^2]\).
The analogous identity for \(I-\Gamma_m\) proves (40.4).
All summands are nonnegative. Replacing every \(q_j\)
by one gives \(I-\Gamma_1\), since
\(\sum_{j<m}D_j=\Gamma_m-\Gamma_1\).
This proves both form inequalities.
Bounded functions are dense and all operators are
bounded, so the identities hold on all of \(L^2\).
\end{proof}

Increasing each cumulative acceptance weight pointwise
can only increase the form. No derivative of any
selector occurs. In particular when \(m=1\),
\(K=\Gamma_1\) exactly: any last-label mismatch in
a chosen coupling is not an intrinsic cost of that
physical kernel. V34's estimates remain upper bounds,
but they need not be sharp.

\subsection{The exact rare-energy remainder}

Define
\[
 \mathcal E_{\rm out}(f)=\|f-\Gamma_1f\|_2^2,\qquad
 \mathcal E_{\rm core}(f)=\|f-\Gamma_mf\|_2^2.
 \tag{40.6}
\]
The exact loss relative to the outer update is
\[
 \mathcal R(f):=\mathcal E_{\rm out}(f)-\mathcal E_K(f)
             =\sum_{j=1}^{m-1}\mu[r_j(D_jf)^2].
 \tag{40.7}
\]
Under V34's sparse-event hypotheses, its proof applies
to each prefix of disjoint collars:
\[
 \mu r_j\le p^j\prod_{k=1}^j n_k,\qquad
                       p=e^{-\beta\tau^2/(64A)}.
 \tag{40.8}
\]
Indeed \(r_j>0\) requires a bad coordinate in each
of the first \(j\) collars. Sum V31's joint bound
over one chosen coordinate in every collar.
An empty collar makes \(r_j=0\).

The factors \((D_jf)^2\) in (40.7) cannot be
replaced by their means. This is an energy-weighted
rare-event problem, not merely a probability estimate.

\begin{proposition}[Rarity alone does not absorb the form remainder]
\label{prop:f12v36-rare-obstruction}
Even for a product law with arbitrarily small
\(\mu r_1\), there is no general bound
\(\mathcal R(f)\le\eta\mathcal E_{\rm out}(f)
 +C\mathcal E_{\rm core}(f)\) with \(\eta<1\)
and finite \(C\) for a single nested update.
\end{proposition}
\begin{proof}
Take independent binary \(x,y,z\), with \(x,y\)
fair and \(\Pr(z=1)=p>0\). Let \(B_1=\{x,y\}\),
\(B_2=\{x\}\), and \(\chi_1=1-z\).
For \(f=z(2y-1)\), one has
\(\Gamma_1f=0\), \(\Gamma_2f=f\), and
\(Kf=(1-z)\Gamma_1f+z\Gamma_2f=f\).
Consequently
\[
 \mu r_1=p,\qquad
 \mathcal E_{\rm out}(f)=\mathcal R(f)=p,\qquad
 \mathcal E_{\rm core}(f)=\mathcal E_K(f)=0.
 \tag{40.9}
\]
The proposed inequality would imply \(p\le\eta p\).
The parameter \(p\) can be arbitrarily small.
\end{proof}

This is a counterexample to a general inference, not
a counterexample to the Yang--Mills gap or to a
sum over overlapping updates. Other updates can
refresh the guard coordinate \(z\). It shows why
such overlap must actually enter a proof.

For a finite family of translated updates with rates
\(\alpha_v\ge0\), sum all three forms and the remainder
with these rates, denoting the sums by a superscript
\({\rm tot}\). If one proves the genuinely collective
estimate
\[
 \mathcal R^{\rm tot}(f)\le
 \eta\mathcal E_{\rm out}^{\rm tot}(f)
             +C\mathcal E_{\rm core}^{\rm tot}(f),
 \qquad \eta<1,\quad C\ge0,
 \tag{40.10}
\]
then (40.5) and (40.7) give the rigorous implication
\[
 \mathcal E_K^{\rm tot}(f)\ge
       \frac{1-\eta}{1+C}\mathcal E_{\rm out}^{\rm tot}(f).
 \tag{40.11}
\]
To check it, substitute
\(\mathcal R^{\rm tot}=\mathcal E_{\rm out}^{\rm tot}
-\mathcal E_K^{\rm tot}\) in (40.10), and use
\(\mathcal E_{\rm core}^{\rm tot}\le\mathcal E_K^{\rm tot}\).
Neither (40.10) nor a volume-uniform gap for
\(\mathcal E_{\rm out}^{\rm tot}\) has yet been established
for the unrestricted fibre. Formula (40.11) is a
conditional comparison, not a completed gap proof.

\subsection{From uniform conditional TV to an \(L^2\) angle}

The following proof supplies a different quantitative
input from V32. Let \(X,Y\) be random variables on
standard Borel spaces, with regular conditional laws
\(\kappa_y=\mathcal L(X\mid Y=y)\). Assume a version with
\[
 \sup_{y,z}\|\kappa_y-\kappa_z\|_{\rm TV}\le\delta<1.
 \tag{40.12}
\]
An essential supremum over a common full-measure set
is sufficient. The TV convention is that of (36.17).

\begin{proposition}[Conditional diameter bounds maximal correlation]
\label{prop:f12v36-angle}
For centered square-integrable \(a(X),b(Y)\),
\[
 |\mathbb E[a(X)b(Y)]|
        \le\sqrt\delta\,\|a\|_2\|b\|_2.
 \tag{40.13}
\]
No lower bound on a density ratio and no product-law
assumption is required.
\end{proposition}
\begin{proof}
Let \(T:L^2(X)\to L^2(Y)\) be conditional expectation,
and let \(T^*\) be the reverse conditional expectation.
The operator \(P=TT^*\) on \(L^2(Y)\) is a positive
self-adjoint Markov contraction preserving the
\(Y\)-marginal. Its kernel samples \(X\) from
\(\kappa_y\), then \(Y'\) from \(\mathcal L(Y\mid X)\).
Postcomposition by a Markov kernel does not increase
TV distance: integrate any test between zero and one.
Thus the rows of \(P\) have TV diameter at most
\(\delta\), and
\(\operatorname{osc}(Pg)\le\delta\operatorname{osc}(g)\)
for bounded real \(g\).

For bounded centered \(g\), stationarity implies
\[
 \|P^ng\|_2\le\|P^ng\|_\infty
                \le\delta^n\operatorname{osc}(g).
 \tag{40.14}
\]
Apply the spectral theorem to the positive operator
\(P\) on the centered space. If the spectral measure
of such a \(g\) charged \((\delta+\varepsilon,1]\),
then its \(2n\)-th moment would contradict (40.14)
as \(n\) tends to infinity. Bounded centered functions
are dense, so the entire centered spectrum lies in
\([0,\delta]\). For \(\delta=0\), the same conclusion
follows already from \(Pg=0\) on this dense class.
Therefore
\(\|T|_{L^2_0}\|^2=\|TT^*|_{L^2_0}\|\le\delta\).
Cauchy--Schwarz applied to \(\langle Ta,b\rangle\)
proves (40.13).
\end{proof}

The square root cannot in general be improved.
Let \(Y\) be a fair sign and independently reveal it
with probability \(\delta\); set \(X=Y\) on a reveal
and \(X=*\) otherwise. The two conditional laws of
\(X\) have TV distance \(\delta\).
For \(b(Y)=Y\) and
\(a(X)=X/\sqrt\delta\) on a reveal, zero otherwise,
both variances are one and their covariance is
\(\sqrt\delta\). This example is not the fibre law.

\subsection{A two-block variance inequality for the actual fibre}

\begin{theorem}[Variance bridge from conditional diameter]
\label{thm:f12v36-variance}
On any probability space carrying \(X,Y\) satisfying
(40.12), every real \(f\in L^2\), including functions
of additional variables, obeys
\[
 \operatorname{Var}(f)\le
 \frac{\mathbb E\operatorname{Var}(f\mid X)
       +\mathbb E\operatorname{Var}(f\mid Y)}
      {1-\sqrt\delta}.
 \tag{40.15}
\]
\end{theorem}
\begin{proof}
Center \(f\), and put \(u=\mathbb E[f\mid X]\),
\(v=\mathbb E[f\mid Y]\). Both are centered.
By (40.13),
\[
 \|u+v\|_2^2\le(1+\sqrt\delta)
                              (\|u\|_2^2+\|v\|_2^2).
 \]
Writing \(A_f=\|u\|_2^2+\|v\|_2^2
=\langle f,u+v\rangle\), Cauchy--Schwarz gives
\(A_f\le(1+\sqrt\delta)\|f\|_2^2\);
the case \(A_f=0\) is immediate.
The sum of conditional variances equals
\(2\|f\|_2^2-A_f\), and is at least
\((1-\sqrt\delta)\|f\|_2^2\). This proves (40.15).
\end{proof}

Now take the actual compact fibre law \(\mu\) in a
finite ambient free-variable set \(I\) containing
V32's nested cubes and stencils. Impose precisely
Theorem~\ref{thm:f12v32-tv}'s curvature, nesting and
coupling assumptions, and put
\[
 X=Y_{I_Q},\qquad Y=Y_{I\setminus I_P},\qquad
 \delta=\delta_{Q,R}.
 \tag{40.16}
\]
Here \(Y\) on the right denotes the exterior random
vector, not an additional independent copy.
By canonical disintegration,
\(\mathcal L(X\mid Y=y)=(\pi_{P,C}^{y})_{I_Q}\),
with any fixed ambient boundary retained. Hence
(36.17), not merely joint-versus-product TV closeness,
verifies (40.12).
There is no restriction of the actual law to a chart.
Consequently
\[
 \operatorname{Var}_\mu(f)\le
 \frac{
 \mu\operatorname{Var}_\mu(f\mid Y_{I_Q})
 +\mu\operatorname{Var}_\mu(f\mid Y_{I\setminus I_P})}
 {1-\sqrt{\delta_{Q,R}}}.
 \tag{40.17}
\]
The same \(L^2\) assertion holds for a compatible DLR
state with the full exterior of \(I_P\).
At \(r=\lfloor\beta\rfloor\) as in (36.19), the
denominator is at least \(1-\sqrt C e^{-c\beta/2}\).
This is a quantitative variance inequality for the
unrestricted fibre under the stated fixed-\(C\)
curvature hypothesis.

In finite volume, the two terms are the Dirichlet
forms of heat baths on \(I\setminus I_Q\) and \(I_P\).
Their equal-probability mixture has an auxiliary
Markov spectral gap at least
\((1-\sqrt{\delta_{Q,R}})/2\).
One of these updates is the possibly extensive
complement of the core. Thus this is not yet a gap
for a bounded-range update system, and it is not a
bound on the physical transfer Hamiltonian.

\subsection{What a scale iteration still has to supply}

Equation (40.17) is a starting two-block estimate,
not a tensorization over translates. Iterating it
requires estimates for the conditional laws generated
at the successive cuts, control of overlap multiplicities,
and accumulated constants uniform in volume.
V32 does not assert the same exponentially small
\(\delta\) for arbitrary growing cores at fixed
\(\beta\). Nor does it assert its bound after
arbitrary conditioning inside the protecting buffer.
Those operations cannot silently be inserted in (40.17).

As a primary-source comparison, N.~Cancrini,
F.~Martinelli, C.~Roberto and C.~Toninelli,
\emph{Kinetically constrained spin models},
\href{https://arxiv.org/pdf/math/0610106}
{arXiv:math/0610106}, Section~2.1 and Theorem~4.2,
use product equilibrium measures and a multiscale
bisection argument for their constrained models.
The product assumption in that formulation is not
available here. Their theorem is therefore not
invoked as a Yang--Mills result. The identities and
the variance bridge above have their own proofs.

\begin{decision}[Next mathematical target]
Pursue a localized variance decomposition that
refreshes the failed guards through overlapping
updates, using (40.17) only where its full conditional
hypotheses survive. The precise quantities are now
the annular energies \(\mu[r_j(D_jf)^2]\), not
uncoupled bad-site probabilities.
This upstream route removes the artificial last-label
mismatch and all selector derivatives from the form,
but it does not assume those annular energies are small.
The next companion audit and direction review remain V40.
\end{decision}

\subsection{Exact algebra checks and append record}

The new script
\begin{center}\small\ttfamily
scripts/verify\_ym\_f12\_v36\_variance\_bridge.py
\end{center}
uses rational arithmetic only. On a 16-state
nonproduct example it checks all 256 entries of
(40.3), the two positive-semidefinite form comparisons,
and cancellation of the last selector.
It checks (40.9) at \(p=1/1000000\).
For the erasure example at \(\delta=9/25\),
it verifies the centered round-trip eigenvalue
\(9/25\), maximal correlation \(3/5\), and the
positive-semidefinite inequality (40.15), including
an equality test. All checks passed. These examples
test algebra and sharpness, not the Yang--Mills
analytic hypotheses.

V35 had 546393 bytes and SHA--256
\begin{center}\scriptsize\ttfamily
2B86FCCBEB06EBAD804CA3E11429001246E84C57A40F333C0D7D3FB645DAE329.
\end{center}
Its complete body is retained apart from the title
version. Only the active V35 filename is replaced by
V36; archives and companion programmes are retained.

\begin{firewall}[Position after V36]
The actual adaptive heat bath has an exact nested
Dirichlet decomposition. The unrestricted compact
fibre satisfies the finite-scale two-block variance
inequality (40.17) under V32's hypotheses.
A localized volume-uniform variance iteration,
fixed-coupling all-distance clustering, scale-stable
effective interactions, continuum construction and
the physical Yang--Mills mass gap remain unproved.
\end{firewall}

\section{V37: refreshing the failed guard controls the annular energy}
\label{sec:f12v37}

\subsection{The obstruction is an energy, not just a rare event}

V36 isolated the exact loss
\(\sum_{j<m}\mu[r_j(D_jf)^2]\) in the adaptive
Dirichlet form. Its product example rules out absorbing
that loss solely by its probability and the innermost
core form. We now refresh the guard in a surrounding
shell. V31 supplies a conditional rarity bound in that
shell, while V36 supplies the needed projection angle.
Together they give a relative annular-energy estimate,
not an absolute probability error.

The resulting theorem is a uniform conditional gap
for a specified regional auxiliary update. It permits
arbitrary compact exterior data at the outer regional
boundary, subject to the stated fixed-\(C\) curvature
hypotheses. The shell update is a full heat bath, not
a product of single-site replacements. No global
bounded-range dynamics gap or physical transfer gap
is inferred from this regional result.

\subsection{A relative projection setup}

Let \(T\) be a finite coarse region and let all the
nested sets \(B_j=I_{R_j}\) lie in \(I_T\).
Let \(S\subset I_T\) be a set of free coordinates
such that \(S\cup B_1=I_T\).
The shell heat bath is the actual conditional expectation
on these free coordinates:
\[
 H=\Gamma_S=\mathbb E_\mu[\,\cdot\mid Y_{S^c}],
 \qquad A=\Gamma_1,\qquad E=\Gamma_T .
 \tag{41.1}
\]
For noncanonical free sets \(S\), this means disintegration
of the same full compact fibre density, including its
Jacobian factors. It does not mean fixing physical
pivots independently of the free coordinates.
The exact finite-range interaction obtained in V18
and used in V25 defines this conditional law.

Use V36's \(D_j=\Gamma_{j+1}-\Gamma_j\),
\(r_j=\prod_{k\le j}(1-\chi_k)\) for \(j<m\).
For real \(f\in L^2(\mu)\), abbreviate
\[
 \begin{aligned}
 e_T(f)&=\|(I-E)f\|_2^2,&
 e_S(f)&=\|(I-H)f\|_2^2,\\
 e_1(f)&=\|(I-A)f\|_2^2,&
 \mathcal R(f)&=\sum_{j<m}\mu[r_j(D_jf)^2].
 \end{aligned}
 \tag{41.2}
\]
The measure \(\mu\) may be a full finite-volume law,
a compatible DLR law, or a fixed exterior conditional
law on \(I_T\). In the last case \(E f=\mu f\).

We first make three precise assumptions, and then
verify them for the fibre geometry below:
\[
 \begin{gathered}
 H r_j\le p\quad (j<m),\qquad 0\le p<1,\\
 \|(AH-E)f\|_2,\ \|(HA-E)f\|_2
                      \le\rho\|(I-E)f\|_2,\qquad 0\le\rho<1,\\
 H\Gamma_j=\Gamma_jH\quad (2\le j\le m).
 \end{gathered}
 \tag{41.3}
\]
All statements are on the same probability space.
There is no independence assumption on \(r_j\) and \(f\).

\subsection{Rare-energy localization with one commutator}

\begin{proposition}[Relative annular-energy bound]
\label{prop:f12v37-localization}
Under (41.3),
\[
 \sqrt{\mathcal R(f)}
       \le\sqrt p\,\sqrt{e_1(f)}
                +\sqrt{e_S(f)}+2\rho\sqrt{e_T(f)}.
 \tag{41.4}
\]
Moreover
\[
 e_T(f)\le\frac{e_1(f)+e_S(f)}{1-\rho}.
 \tag{41.5}
\]
Consequently, with
\[
 a=\sqrt p+\frac{2\rho}{\sqrt{1-\rho}},
 \qquad b=1+\frac{2\rho}{\sqrt{1-\rho}},
 \tag{41.6}
\]
one has
\[
 \mathcal R(f)\le 2a^2 e_1(f)+2b^2 e_S(f).
 \tag{41.7}
\]
\end{proposition}
\begin{proof}
Decompose each \(D_jf=HD_jf+(I-H)D_jf\), and
apply the triangle inequality in the direct sum of
the weighted \(L^2\) spaces. Because \(HD_jf\) is
measurable outside \(S\),
\[
 \begin{split}
 \sum_{j<m}\mu[r_j(HD_jf)^2]
 &=\sum_{j<m}\mu[(Hr_j)(HD_jf)^2]\\
 &\le p\sum_{j<m}\|HD_jf\|_2^2
 \le p\sum_{j<m}\|D_jf\|_2^2
 \le p\,e_1(f).
 \end{split}
 \tag{41.8}
\]
The \(D_j\)'s are mutually orthogonal projections,
and their sum is \(\Gamma_m-\Gamma_1\).
This proves the last inequality. The factors
\(r_j\) need not commute with \(H\).

For the other part use \(0\le r_j\le1\) and
\[
 (I-H)D_jf=D_j(I-H)f+[D_j,H]f.
 \tag{41.9}
\]
For \(j\ge2\), the commutator vanishes by (41.3).
For \(j=1\), it equals \(-[A,H]f\).
The second line of (41.3) gives
\[
 \|[A,H]f\|_2
   \le\|(AH-E)f\|_2+\|(HA-E)f\|_2
   \le2\rho\sqrt{e_T(f)}.
 \tag{41.10}
\]
Orthogonality of the \(D_j\)'s and the direct-sum
triangle inequality therefore give
\[
 \left(\sum_{j<m}\|(I-H)D_jf\|_2^2\right)^{1/2}
             \le\sqrt{e_S(f)}+2\rho\sqrt{e_T(f)}.
 \tag{41.11}
\]
Combining (41.8) and (41.11) proves (41.4).
There is no factor equal to the number of shells.

For (41.5), center relative to \(E\).
The second line of (41.3) bounds the angle between
the ranges of \(A-E\) and \(H-E\) by \(\rho\).
The proof of Theorem~\ref{thm:f12v36-variance},
with these relative projections, yields (41.5).
More explicitly, set \(g=(I-E)f\),
\(u=Ag\), \(v=Hg\). Then
\(|\langle u,v\rangle|\le\rho\|u\|_2\|v\|_2\),
so \(\|u\|_2^2+\|v\|_2^2\le(1+\rho)\|g\|_2^2\).
Subtract from \(2\|g\|_2^2\).

Finally,
\(\sqrt{e_T}\le(\sqrt{e_1}+\sqrt{e_S})/\sqrt{1-\rho}\).
Substitute in (41.4) and use
\((a\sqrt{e_1}+b\sqrt{e_S})^2\le2a^2e_1+2b^2e_S\).
This proves (41.7).
\end{proof}

The refreshing energy \(e_S(f)\) is indispensable.
For the example (40.9), a refresh of the failed
guard changes the test function even though the
innermost core update does not. Thus (41.7) does
not contradict V36's obstruction.

\subsection{A regional hybrid kernel with an explicit constant}

\begin{theorem}[Adaptive update plus guard refresh]
\label{thm:f12v37-hybrid}
Suppose (41.3) holds with
\[
                  p\le1/256,\qquad \rho\le1/16.
 \tag{41.12}
\]
Let \(K\) be V34's actual adaptive kernel and set
\[
                  K_{\rm ref}=\frac16K+\frac56H.
 \tag{41.13}
\]
This is a positive reversible Markov kernel for
\(\mu\), and
\[
 \langle f,(I-K_{\rm ref})f\rangle_\mu
                         \ge\frac1{12}\,e_T(f).
 \tag{41.14}
\]
In a fixed conditional law on \(I_T\), its
centered \(L^2\) operator norm is at most \(11/12\).
\end{theorem}
\begin{proof}
Since \(1/\sqrt{1-\rho}\le2\), (41.6) gives
\(a\le5/16\), \(b\le5/4\). Hence (41.7) implies
the slightly weaker convenient estimate
\[
              \mathcal R(f)\le\tfrac14 e_1(f)
                                      +\tfrac{25}{8}e_S(f).
 \tag{41.15}
\]
By V36's exact identity,
\(\mathcal E_K(f)=e_1(f)-\mathcal R(f)\). Thus
\[
 \begin{split}
 \mathcal E_K(f)+5e_S(f)
 &\ge\tfrac34 e_1(f)+\tfrac{15}{8}e_S(f)\\
 &\ge\tfrac34(e_1(f)+e_S(f))
 \ge\tfrac34(1-\rho)e_T(f)
 \ge\tfrac12 e_T(f).
 \end{split}
 \tag{41.16}
\]
Division by six proves (41.14).

Reversibility follows from V34 and from the
conditional-expectation definition of \(H\).
Every summand \(w_j\Gamma_j\) of \(K\) is
positive semidefinite: exterior measurability gives
\(\langle f,w_j\Gamma_jf\rangle
=\mu[w_j(\Gamma_jf)^2]\ge0\).
The fallback has the same property.
Thus \(K_{\rm ref}\) is positive as an \(L^2\)
operator as well as a Markov kernel.
In a fixed regional law, (41.14) places its
centered spectrum in \([0,11/12]\).
\end{proof}

The constants concern one exact block update or one
exact shell update. Computational cost and physical
Euclidean time have not been assigned to either step.

\subsection{Verification for the compact Yang--Mills fibre}

We now verify (41.3) without new probabilistic
assumptions on the fibre. Use V32's coarse cubes with
radii
\[
 Q:\lfloor r/4\rfloor,\qquad R:r,\qquad P:2r,
 \qquad T:3r,\qquad r=\lfloor\beta\rfloor.
 \tag{41.17}
\]
Take \(R_1=P\), and choose \(R_2\) of radius
\(\lfloor r/8\rfloor\). Any further nested \(R_j\)'s
lie inside \(R_2\). Set
\[
                  B_1=I_P,\qquad S=I_T\setminus I_Q.
 \tag{41.18}
\]
For \(m=2\) the construction already suffices.
If more selectors are used, retain their exterior
measurability as in V34. The first selector is
\(\chi_1=\prod_{i\in J_1}g_i\) at V34's fixed
threshold \(\tau\).

Let \(hJ(\tau/2)\) be V31's fixed buffer width.
Take \(\beta\) so large that all these cubes,
their finite interaction stencils, and this buffer
fit without wrapping in a finite torus; equivalently
work in an ambient region containing them.
Assume all affected coarse curvatures in these
stencils are at most
\[
       \epsilon_{\rm ref}
       :=\min\{\epsilon_{\rm dec},\epsilon_{\tau/2}\},
 \tag{41.19}
\]
and all coupling thresholds of V31 and V32 hold.
These are fixed positive curvature thresholds,
not conditions tending to zero with ambient volume.
They are not asserted for an arbitrary coarse
configuration drawn from the full Yang--Mills law.

\paragraph{Conditional rarity in the shell.}
Each \(i\in J_1=J_P^\partial\) lies in a fixed-width
collar of \(P\). Its V31 neighbourhood at threshold
\(\tau/2\) lies wholly in \(S\) for sufficiently
large \(r\): it is far outside \(Q\) and far inside
\(T\). The free-variable stencil conventions change
these separations only by fixed constants.
Disintegrate that canonical neighbourhood kernel
inside the actual shell conditional law.
V31's bound is uniform over its resulting exterior,
including arbitrary values on \(I_Q\) and outside \(I_T\).
It follows that, almost surely,
\[
 \begin{split}
 H(1-\chi_1)
 &\le\sum_{i\in J_1}
       H{\bf1}_{\{\|Y_i-J_i(C)\|_{\rm HS}>\tau/2\}}\\
 &\le |J_1|e^{-\beta\tau^2/(64A)}
                         =:p_{\rm ref}.
 \end{split}
 \tag{41.20}
\]
Here \(A\) in the exponent is the coercivity
constant of V31, not the projection \(A=\Gamma_1\)
used in the abstract calculation above.
This notational distinction affects no operator identity.
Since \(r_j\le1-\chi_1\), (41.20) proves the first
line of (41.3), with \(p=p_{\rm ref}\) once it is below one.
In particular \(|J_1|=O(r^3)\) and \(p_{\rm ref}\)
tends exponentially to zero as \(\beta\) tends to infinity.

\paragraph{The angle survives exactly the needed conditioning.}
Condition on the entire exterior of \(I_T\), giving
the actual law \(\pi_{T,C}^{\xi}\) with arbitrary
compact \(\xi\).
Inside it, \(H\) retains \(Y_{I_Q}\), whereas
\(A=\Gamma_P\) retains \(Y_{I_T\setminus I_P}\).
The conditional law of the former given the latter
is precisely the \(I_Q\)-marginal of the canonical
\(P\)-kernel. V32 gives its uniform diameter
\(\delta_{Q,R}\), independently of \(\xi\).
Proposition~\ref{prop:f12v36-angle} therefore yields
the second line of (41.3), with
\[
                 \rho_{\rm ref}=\sqrt{\delta_{Q,R}}
                            \le\sqrt C e^{-c\beta/2}.
 \tag{41.21}
\]
The projection-norm statement is obtained on each
conditional \(L^2\) space and then integrated over
the exterior of \(I_T\), if that exterior is random.
This only conditions outside the protecting parent
or on its canonical exterior inputs. No variables
inside the protecting buffer are additionally pinned.

\paragraph{Exact commutation of inner and shell heat baths.}
The completed fibre potential is a sum of fixed-range
terms in the independent free coordinates, including
the inverse-pivot Jacobians. Choose \(r\) large
relative to that fixed range. Then no interaction
term meets both \(B_2\) and \(S\), since \(B_2\)
lies well inside \(I_Q\).
The same is true for every \(B_j\subset B_2\).

Conditional on the remaining coordinates, the
density on \(B_j\cup S\) factors into a function
of \(Y_{B_j}\) times a function of \(Y_S\).
Its normalized conditional law is therefore a
product of those two conditional laws. Each is
independent of the coordinates updated by the other.
Fubini gives
\[
                      \Gamma_j H=H\Gamma_j
                                  \quad(j\ge2).
 \tag{41.22}
\]
This proves the last line of (41.3). The outer
update need not commute with the shell update;
its noncommutation is controlled by (41.21).

\begin{corollary}[Uniform conditional regional refresh gap]
\label{cor:f12v37-fibre}
For the geometry (41.17)--(41.18), the fixed
small-curvature condition (41.19), and all
sufficiently large \(\beta\), the actual hybrid
kernel (41.13) in \(\pi_{T,C}^{\xi}\) obeys
\[
 \operatorname{Var}_{\pi_{T,C}^{\xi}}(f)
 \le12\,
       \langle f,(I-K_{\rm ref})f\rangle_{\pi_{T,C}^{\xi}}
 \qquad(f\in L^2).
 \tag{41.23}
\]
The constant is independent of the compact exterior
\(\xi\) and of the surrounding ambient volume.
\end{corollary}
\begin{proof}
The preceding three verifications give (41.3).
Equations (41.20)--(41.21) make (41.12) valid
at sufficiently large \(\beta\).
In the fixed conditional law,
\(e_T(f)=\operatorname{Var}(f)\).
Apply Theorem~\ref{thm:f12v37-hybrid}.
\end{proof}

\subsection{What has closed, and what has not}

The annular-energy loss is now absorbed by an actual
guard-refresh form, with a strictly smaller-than-one
coefficient on the outer-update form. This is the
relative estimate missing from a rarity-only argument.
Both kernels preserve the exact compact fibre law.
No error remains that is nonzero merely because
two configurations happen to contain bad sites.

The refresh block \(S\) is, however, a large annulus
with a hole. It is sampled as one conditional block.
The theorem does not replace that sample by finitely
many smaller local samples with a uniform comparison
constant. Summing (41.14) over translated \(T\)'s
compares the corresponding hybrid forms with the
translated \(T\)-heat-bath form; it does not itself
prove that the latter controls the full variance
uniformly in ambient volume.

\begin{decision}[The next noncircular step]
Localize the shell refresh, retaining its protecting
buffer and controlling the overlap constants.
A useful next target is a variance decomposition
for the annular conditional law under arbitrary
core and exterior data, with strictly smaller update
diameters and explicit constants. The conditional
coarse-curvature hypotheses must survive that
decomposition. It would be circular to assume
a volume-uniform gap for those annuli or for the
translated \(T\)-block dynamics as a substitute.
The next companion audit and direction review remain V40.
\end{decision}

\subsection{Exact verification and append record}

The new script
\begin{center}\small\ttfamily
scripts/verify\_ym\_f12\_v37\_guard\_refresh.py
\end{center}
uses rational arithmetic on a 16-state nonproduct
law with an independent inner coordinate and a
weakly interacting three-coordinate chain.
It checks inner--shell commutation, a conditional
core TV diameter \(1/441\), the resulting angle
bound \(1/21<1/16\), and worst conditional bad-guard
probability \(221/225501<1/256\).
Exact positive-semidefinite elimination verifies
(41.15) and the regional comparison (41.16).
It also checks that the hybrid kernel is Markov,
reversible and positive. All checks passed.
This finite example validates the algebra, not
the Yang--Mills analytic inputs, whose proofs are
the V31, V32 and V36 arguments cited above.

V36 had 561520 bytes and SHA--256
\begin{center}\scriptsize\ttfamily
156804C4F34AB576FE2BE97E01CFEB61DE5FD95CF4B87D69AA0A53625572CFFD.
\end{center}
Its complete body is retained apart from the title
version. Only active V36 is replaced by V37.
All archives, previous scripts and companion
programmes are retained.

\begin{firewall}[Position after V37]
The rare annular energy has a proved relative
guard-refresh bound, and the specified hybrid
update has a uniform conditional regional
\(L^2\) gap. A localized volume-uniform variance
iteration, all-distance clustering for the fibre,
scale-stable effective interactions, continuum
construction and the physical Yang--Mills mass
gap remain unproved.
\end{firewall}

\section{V38: localizing the shell with an explicit surface cost}
\label{sec:f12v38}

\subsection{What is new, and what is reused}

V37 gives a regional gap using a complete annular
heat bath. We now replace that annular update by
smaller conditional blocks, with a proved comparison
cost. The cost is controlled by crossing interactions,
not by every interaction in the shell.
In four coarse dimensions the resulting exponent
is proportional to a three-dimensional cut area.
It is not a scale-uniform estimate.

Product variance tensorization was already proved
in f6 V59. The normalized density-ratio argument
for bounded perturbations was already given in
f6 V97. We reuse those principles; they are not
new probabilistic results. The additional work here
is their conditional-block formulation for the
completed compact fibre, a cut count valid with
arbitrary holes and exterior values, the resulting
smaller-diameter replacement of V37's shell, and
the explicit recursive surface-order bound.

An independent-dimer example at the end distinguishes
a defect of the supremum-density comparison from
an actual loss of a spectral gap. This determines
where a stronger argument must improve the estimate.

\subsection{An exact cut comparison with one oscillation factor}

Fix a finite set \(S\) of independent free group
coordinates and arbitrary exterior data \(\xi\).
The law \(\mu_S^\xi\) is the conditional law of
the same completed compact fibre used in V25
and V37. Partition \(S\) into nonempty disjoint
sets \(S_1,\ldots,S_t\).
Write its potential, omitting constants, as
\[
 V_S^\xi(Y)=\sum_{\alpha=1}^t V_\alpha^\xi(Y_{S_\alpha})
                              +W^\xi(Y).
 \tag{42.1}
\]
Here \(W^\xi\) is the sum of precisely those local
interaction terms whose unfrozen free-variable support
meets at least two partition sets. Terms meeting
only one set remain in its \(V_\alpha^\xi\), even
when they also depend on fixed core or exterior data.

Let \(\nu^\xi=\bigotimes_\alpha\nu_\alpha^\xi\)
be the normalized product law with potential
\(\sum_\alpha V_\alpha^\xi\), and put
\[
                  \Omega_S^\xi=\operatorname{osc}_S W^\xi.
 \tag{42.2}
\]
Deleting the cut terms here defines only a reference
measure for comparison. The actual law and every
conditional heat bath in the conclusion retain them.

\begin{proposition}[Conditional block comparison]
\label{prop:f12v38-cut}
For every real \(f\in L^2(\mu_S^\xi)\),
\[
 \operatorname{Var}_{\mu_S^\xi}(f)
 \le e^{\Omega_S^\xi}
       \sum_{\alpha=1}^t
       \mu_S^\xi\operatorname{Var}_{\mu_S^\xi}
                              (f\mid Y_{S\setminus S_\alpha}).
 \tag{42.3}
\]
The coefficient is \(e^{\Omega_S^\xi}\), not its square.
\end{proposition}
\begin{proof}
Write \(r=d\mu_S^\xi/d\nu^\xi\), \(m=\inf r>0\),
and \(M=\sup r<\infty\). Positivity and compactness
give these bounds, and the exact normalizer cancels
in their ratio:
\[
                         M/m=e^{\Omega_S^\xi}.
 \tag{42.4}
\]
Variance minimization over constants and the product
tensorization from f6 V59 give
\[
 \operatorname{Var}_{\mu_S^\xi}(f)
 \le M\operatorname{Var}_{\nu^\xi}(f)
 \le M\sum_\alpha
       \nu^\xi\operatorname{Var}_{\nu^\xi}
                              (f\mid Y_{S\setminus S_\alpha}).
 \]
For each part, use the variational characterization
of conditional variance:
\[
 \begin{split}
 \mu_S^\xi\operatorname{Var}_{\mu_S^\xi}(f\mid Y_{S\setminus S_\alpha})
 &=\inf_{g=g(Y_{S\setminus S_\alpha})}\mu_S^\xi[(f-g)^2]\\
 &\ge m\inf_g\nu^\xi[(f-g)^2]\\
 &=m\,\nu^\xi\operatorname{Var}_{\nu^\xi}(f\mid Y_{S\setminus S_\alpha}).
 \end{split}
 \tag{42.5}
\]
The infima may be taken over square-integrable
measurable functions; bounded positive density ratios
make the two \(L^2\) classes identical.
Combine the estimates. This compares the actual
conditional variances, not the conditionals of a
modified interaction.
\end{proof}

The same assertion may be integrated over random
exterior data when \(\Omega_S^\xi\le\Omega\) uniformly:
\[
 \mu\operatorname{Var}_\mu(f\mid Y_{S^c})
 \le e^\Omega\sum_\alpha
           \mu\operatorname{Var}_\mu(f\mid Y_{S_\alpha^c}).
 \tag{42.6}
\]
Consistency of the full specification identifies
the conditional kernels after the second disintegration.
There is no untracked conditional normalizer.

\subsection{The full compact fibre has bounded local cut terms}

Assign each internal free coordinate to its cell
and each free nonpivot boundary coordinate to its
oriented coarse-edge tail. There are at most
\(q_0=45\) free coordinates with a given owner cell.
Let \(\Phi_a\) denote the individual potentials
\(\beta s_p(W(Y,C))\) and \(-\log j_e(Y,C_e)\)
of (29.16), at fixed retained \(C\).
V18's global compact completion and V25's finite-range
specification give fixed constants:
\[
 \begin{gathered}
 \operatorname{diam}_\infty
       \{\text{owners of free coordinates in }\operatorname{supp}\Phi_a\}
                  \le R_0,\\
 \#\{a:i\in\operatorname{supp}\Phi_a\}\le D_0,\qquad
                    \operatorname{osc}\Phi_a\le v_\beta,\\
 v_\beta=\max\{\beta s_{\max},\log(j_+/j_-)\}
                                  \le c_0(\beta+1).
 \end{gathered}
 \tag{42.7}
\]
Here \(R_0\ge1\) is an integer, \(0<j_-\le j_e\le j_+\),
and constants may depend on the fixed smooth completion
and representation. Supports are those of the completed
free-coordinate interaction, not just the original
four links of a plaquette.
All these bounds are uniform in exterior free data,
retained compact data, and ambient volume.

No small-curvature assumption is needed for (42.7).
The pointwise curvature condition is still required
when using V37's stronger regional estimate.

\begin{proposition}[Four-dimensional midpoint cut]
\label{prop:f12v38-surface}
Suppose all owners of \(S\) lie in a coarse cube
with \(L\) integer cells along each side. Bisect
each coordinate interval and partition \(S\) by
the resulting at most sixteen owner boxes,
discarding empty parts. Each part has owner-box
side at most \(\lceil L/2\rceil\), and
\[
 \Omega_S^\xi\le c_{\rm cut}v_\beta L^3,\qquad
                  c_{\rm cut}=8q_0D_0(R_0+1).
 \tag{42.8}
\]
The estimate holds for arbitrary subsets \(S\)
of the owner cube, including annuli and disconnected sets.
\end{proposition}
\begin{proof}
A term counted in \(W^\xi\) has two unfrozen
coordinates in different boxes, separated by
one of the four midpoint planes. Its support
diameter is at most \(R_0\), so at least one such
coordinate has owner in the width-\((2R_0+2)\)
slab around that plane.
The union of four slabs contains at most
\(8(R_0+1)L^3\) owner cells, with at most \(q_0\)
coordinates per cell. At most \(D_0\) terms are
incident on each coordinate. Counting with possible
overcount therefore gives at most \(c_{\rm cut}L^3\)
cut terms. Sum their oscillation bounds from (42.7).

Removing cells or coordinates to form a hole
introduces no new partition planes. An interaction
with the fixed hole data and only one updated part
is retained in \(V_\alpha^\xi\), not charged to the cut.
Thus the count is valid uniformly for these geometries.
\end{proof}

\subsection{A genuinely smaller-diameter replacement of V37's shell}

Use V37's \(S=I_T\setminus I_Q\).
Its owners fit in a cube with
\(L=6r+O(1)\), including incoming boundary owners.
Let \(S_\alpha\) be the midpoint partition above,
and set
\[
 e_{S_\alpha}(f)=
                 \mu\operatorname{Var}_\mu(f\mid Y_{S_\alpha^c}),
 \qquad \Omega_L=c_{\rm cut}v_\beta L^3.
 \tag{42.9}
\]
Equation (42.6) now proves, under arbitrary
core and exterior data,
\[
                         e_S(f)\le e^{\Omega_L}
                                           \sum_\alpha e_{S_\alpha}(f).
 \tag{42.10}
\]
These are complete free-coordinate conditional
updates on smaller pieces of the shell, not
unconditioned samples pasted together.

Under V37's regional hypotheses, (41.16) and
(42.10) give
\[
 e_T(f)\le2\mathcal E_K(f)+
                            10e^{\Omega_L}\sum_\alpha e_{S_\alpha}(f).
 \tag{42.11}
\]
All possible updates in this expression have
strictly smaller owner diameter than \(I_T\)
for large \(r\): the outermost update in \(K\)
has side \(4r+O(1)\), and each new shell piece
has side at most \(3r+O(1)\), compared with
\(6r+O(1)\) for \(I_T\).

It is important to keep update rates explicit.
The generator using \(K\) and each of the \(t\le16\)
patch heat baths at rate one has form
\[
                \mathcal E_{\rm loc}(f)
                       =\mathcal E_K(f)+\sum_{\alpha=1}^t e_{S_\alpha}(f).
 \tag{42.12}
\]
In a fixed regional conditional law its spectral
gap is at least \(1/(10e^{\Omega_L})\), by (42.11).
The equal-probability discrete mixture
\((K+\sum_\alpha\Gamma_{S_\alpha})/(1+t)\)
has gap at least
\[
                         \frac1{10(1+t)e^{\Omega_L}}.
 \tag{42.13}
\]
No exponentially large refresh rate has been hidden
in an order-one gap claim.

\subsection{The complete recursive bound is surface order}

There is also a baseline consequence of the cut
argument alone, without V37's curvature hypotheses.
For any finite free-coordinate set \(S\), let
\(\Gamma_i\) denote its actual one-coordinate heat
bath at \(i\in S\), with all other free data fixed.

\begin{theorem}[Finite-region single-coordinate comparison]
\label{thm:f12v38-recursion}
There is a fixed \(C_{\rm surf}<\infty\) such that,
if the owners of \(S\) fit in a cube of side \(L\ge1\),
then for every exterior \(\xi\),
\[
 \operatorname{Var}_{\mu_S^\xi}(f)
 \le e^{C_{\rm surf}v_\beta L^3}
       \sum_{i\in S}
           \mu_S^\xi\operatorname{Var}_{\mu_S^\xi}
                                      (f\mid Y_{S\setminus\{i\}}).
 \tag{42.14}
\]
One may take
\(C_{\rm surf}=q_0D_0+(64/7)c_{\rm cut}\).
\end{theorem}
\begin{proof}
First consider a dyadic owner side \(N=2^n\).
At side one there are at most \(q_0\) coordinates.
Partition them into singletons. At most \(q_0D_0\)
terms touch them, so Proposition~\ref{prop:f12v38-cut}
gives coefficient at most \(e^{q_0D_0v_\beta}\).

At side \(N>1\), apply the midpoint comparison.
After conditioning on the complement of a part,
that part has precisely another conditional law
of the same finite-range specification. Its new
exterior may include arbitrary values from the
other parts and from holes; the bounds (42.7)
are still uniform. Hence the side-\(N/2\) estimate
may be applied inside every conditional variance.
Integrate and sum. Each coordinate belongs to
exactly one part, so no factor sixteen is introduced.

Along a branch of the recursion the logarithmic
cost is at most
\[
 v_\beta\left[
 q_0D_0+c_{\rm cut}\sum_{k=0}^{n-1}(N/2^k)^3\right]
 =v_\beta\left[q_0D_0+\tfrac87c_{\rm cut}(N^3-1)\right].
 \tag{42.15}
\]
This proves the dyadic claim by induction.
For a general \(L\), enclose its owner cube in
a dyadic cube with \(L\le N<2L\), without adding
any updated coordinates. Empty cells are allowed.
Bound (42.15) by \(C_{\rm surf}v_\beta L^3\).
\end{proof}

For the continuous-time generator
\(\sum_{i\in S}(\Gamma_i-I)\), with rate one
at every coordinate, (42.14) gives gap at least
\(e^{-C_{\rm surf}v_\beta L^3}\).
For one-coordinate random scan, divide this gap
bound by \(|S|\).
These are finite-region estimates uniform in the
exterior. They do not give a positive lower bound
as \(L\) tends to infinity at fixed \(\beta\).

At V37's chosen \(r=\lfloor\beta\rfloor\), the
available shell exponent is bounded by
\(K(\beta+1)r^3=O(\beta^4)\).
This describes the worst-case certificate supplied
by the comparison, not a lower bound on the true
cost or an upper bound on the true gap.

\subsection{Why an extensive density ratio can be the wrong diagnostic}

Consider \(n\) independent Ising dimers
\((s_k,t_k)\in\{-1,1\}^2\) with law
\[
 \mu_n(s,t)\ \propto\
                  \exp\left(J\sum_{k=1}^n s_kt_k\right),
 \qquad J>0.
 \tag{42.16}
\]
Partition into all \(s\)'s and all \(t\)'s.
The cut potential has oscillation \(2Jn\), so
the exact largest-to-smallest density ratio relative
to the independent uniform reference is \(e^{2Jn}\).

Nevertheless the two-block heat bath has a gap
independent of \(n\). To verify this without
assuming product-channel tensorization, put
\(c=\tanh J\). Both block marginals are uniform
product laws. On their Walsh bases
\(\chi_A(s)=\prod_{k\in A}s_k\),
the conditional-expectation channel satisfies
\[
                  \mathbb E[\chi_A(t)\mid s]
                                  =c^{|A|}\chi_A(s).
 \tag{42.17}
\]
Conditional independence of the dimers proves this
identity directly. The Walsh functions are orthonormal,
so the centered channel norm is exactly \(c\).
The projection argument of V36 then bounds the
rate-one two-block generator gap below by \(1-c\).
The function \(s_1+t_1\) attains this bound:
the sum of the two conditional projections sends
it to \((1+c)(s_1+t_1)\).
Thus the generator gap equals \(1-c\), and the
equal-probability discrete mixture gap equals
\((1-c)/2\), for every \(n\).

This example is not a Yang--Mills model.
It proves that a growing exact oscillation can
coexist with a uniform block gap. Therefore
failure of the cut-density certificate does not
prove failure of the desired variance estimate.
The normalized interface channel retains information
that a single supremum-density ratio discards.

\begin{decision}[Next direction: retain interface correlations]
The whole-shell sampling issue is now resolved
at finite size by (42.11), and recursively by
(42.14), with explicit exterior-uniform costs.
The remaining problem is to improve those costs,
not to assert that they are already uniform.
Further iterations of the same cut-oscillation
estimate cannot remove its displayed surface factor.

Next examine a conditional interface-channel or
variance comparison that keeps the crossing Wilson
terms in the principal measure. It must handle the
actual nonproduct boundary law and arbitrary core
data created by the cuts. Independent-channel
behavior in (42.17) cannot be assumed for that law.
The next companion audit and direction review remain V40.
\end{decision}

\subsection{Verification and append record}

The script
\begin{center}\small\ttfamily
scripts/verify\_ym\_f12\_v38\_cut\_localization.py
\end{center}
checks (42.3) on a 16-state nonproduct chain with
normalized density ratio exactly three, using
exact positive-semidefinite elimination.
It checks eight four-dimensional midpoint cuts,
with and without holes; a full side-\(L\)
nearest-neighbour cube has exactly \(4L^3\)
crossing edges. It checks eleven dyadic path sums
in (42.15), and all Walsh-channel eigenvalues for
one through five independent dimers at \(c=1/2\).
Their density ratios are \(3^n\), while their
two-block random-scan gaps remain \(1/4\).
All checks passed. The tests do not certify a
Yang--Mills spectral gap.

V37 had 576703 bytes and SHA--256
\begin{center}\scriptsize\ttfamily
8C41AB4DE4367EB900D645CE1943D2D7217872A44FACB0F5D5AE386F83A710E3.
\end{center}
Its complete body is retained apart from the
title version. Only active V37 is replaced by
V38; archives and companion programmes are retained.

\begin{firewall}[Position after V38]
The annular update has a proved smaller-diameter
localization with explicit surface cost, and the
full compact conditional law has the finite-region
single-coordinate bound (42.14).
No volume-uniform localized gap, all-distance
fibre clustering, scale-stable effective interaction,
continuum construction or physical Yang--Mills
mass gap has been proved.
\end{firewall}

\section{V39: a finite-colour gap and its exact compact-interior lift}
\label{sec:f12v39}

\subsection{Use the existing conditional-angle result in its valid domain}

V38's cut comparison pays the oscillation of all
crossing interactions. We now keep those interactions
in the principal measure. The relevant general
conditional-angle theorem was already proved in the
SM programme, so it should not be rediscovered as
a new Yang--Mills theorem.

The source inspected here is SM archive
\texttt{Renewal\_SM\_Einstein\_Continuum\_Research\_Notes\_V100\_f4.tex},
Section V41, Theorem \texttt{thm:v4v41-angle},
with SHA--256
\begin{center}\scriptsize\ttfamily
7504F52071E4CAB739EDE065F165CFCE932DE38237792B33B43424DC31D943BA.
\end{center}
It bounds a conditional channel using a lower
Hessian bound and a mixed-Hessian operator norm.
The theorem itself is a general positive-measure
statement, not a physical SM conclusion imported
into Yang--Mills.

The new work below checks its fixed-product-box
version, applies it to the measured YM potentials
already proved convex in V28 and V32, and uses a
finite colouring to obtain a local heat-bath gap.
For V28's exact interface law, that gap then lifts
to local updates of the original packed interiors
over their \emph{full compact domains}.
The interface remains restricted to its stated
small chart throughout. Removing that restriction
is not part of the theorem.

This targeted archive consultation does not replace
the scheduled latest-SM/NS/parallel-YM audit at V40.

\subsection{The product-box transfer of the SM angle bound}

Let \(\nu\) have measured density proportional to
\(e^{-\Psi(x,y)}\) on a fixed finite product of
Euclidean coordinate cubes. Assume
\[
                  \Psi''\ge\kappa I,\qquad
                  \|\Psi_{xy}''\|_{\rm op}\le J,
                  \qquad \kappa>0.
 \tag{43.1}
\]
The two collections \(x,y\) may contain many
free group coordinates, each represented by its
fixed local coordinate cube.

The SM V41 conclusion remains valid:
\[
 \|\,g(y)\mapsto\mathbb E_\nu[g(y)\mid x]\,\|_{L^2_0\to L^2_0}
                \le \rho_{\kappa,J}
                  :=\frac{J}{\sqrt{\kappa^2+J^2}}<1 .
 \tag{43.2}
\]
Here is the domain check required for this application.
The convex-box Poincare inequality used in the
YM boxed estimates applies both to the full law
and to every conditional \(y\)-law, with constant
at most \(1/\kappa\). Applying the full inequality
to a function of \(x\) also gives the needed
marginal Poincare inequality.

For a bounded \(g(y)\), put \(h(x)=\mathbb E[g\mid x]\).
The conditional domain does not depend on \(x\).
Differentiation of its \emph{normalized} integral gives,
for a unit vector \(u\) in the \(x\)-space,
\[
 u\cdot\nabla h
       =-\operatorname{Cov}(g,u\cdot\nabla_x\Psi\mid x),
 \qquad
 |\nabla h|^2\le
                 \frac{J^2}{\kappa}\operatorname{Var}(g\mid x).
 \tag{43.3}
\]
The second estimate uses conditional Poincare for
the scalar score \(u\cdot\nabla_x\Psi\), followed
by the supremum over unit \(u\). It uses an
operator norm, not a sum over coordinates of the cut.
The marginal Poincare inequality gives
\[
 \operatorname{Var}(h)\le
       \frac{J^2}{\kappa^2}\mathbb E\operatorname{Var}(g\mid x).
 \tag{43.4}
\]
Insert the total-variance identity
\(\operatorname{Var}(g)=\operatorname{Var}(h)
 +\mathbb E\operatorname{Var}(g\mid x)\).
This is exactly the algebra of the archived SM proof
and gives (43.2). Bounded approximations and the
\(L^2\) contraction of conditional expectation extend
it to all \(g\in L^2\). Fixed cube faces produce no
moving-domain term. Thus the domain change has not
introduced an unproved boundary hypothesis.

By the projection calculation in V36, (43.2) implies
\[
 \operatorname{Var}_\nu(f)\le
 \frac{\nu\operatorname{Var}_\nu(f\mid x)
       +\nu\operatorname{Var}_\nu(f\mid y)}
      {1-\rho_{\kappa,J}}.
 \tag{43.5}
\]
This estimate keeps the complete cross interaction.
It requires no smallness assumption \(J<\kappa\).

\subsection{A finite-colour conversion to local heat baths}

Treat each independent free group coordinate,
including its \(d\)-dimensional cube, as one site.
Put an edge between two sites whenever a local
potential term depends on both. Suppose this graph
has maximum degree \(\Delta\), and choose a proper
colouring with \(q\le\Delta+1\) colours.
Greedy colouring proves existence in every finite
volume, including a torus with wraparound edges.
This avoids assuming a particular torus period
for a coordinate-modulus colouring.

\begin{theorem}[Uniform local gap for a convex finite-range cube law]
\label{thm:f12v39-colours}
Suppose on a product cube
\[
                m_*\beta I\le\Psi''\le M_*\beta I,
                 \qquad m_*>0,
 \tag{43.6}
\]
and the interaction graph admits the \(q\)-colouring
above. For the actual coordinate heat baths,
\[
 \operatorname{Var}_\nu(f)\le\gamma_*^{-1}
       \sum_i\nu\operatorname{Var}_\nu(f\mid Y_{i^c}),
 \qquad
 \gamma_*=(1-\rho_*)^{q-1},\qquad
 \rho_*=\frac{M_*}{\sqrt{m_*^2+M_*^2}}.
 \tag{43.7}
\]
The constant is independent of the number of sites.
Every conditional law obtained by fixing some
coordinates to values in their cubes has the same
bound with the same upper bound on \(q\).
\end{theorem}
\begin{proof}
Fixing coordinates restricts the Hessian to a
principal submatrix and the graph to an induced
subgraph. Thus (43.6) persists. Across any split
of the remaining coordinates, the mixed block has
operator norm at most \(M_*\beta\).
Equation (43.5) consequently has the uniform
factor \(C_*=(1-\rho_*)^{-1}\).

Proceed by induction on the number of nonempty
colours. With one colour, no interaction term
contains two sites. The conditional law is a
product, and f6 V59's product tensorization gives
the result with coefficient one.

For more colours, split into the first colour
\(\mathcal C_1\) and its complement \(B\).
Equation (43.5) gives
\[
 \operatorname{Var}_\nu(f)
 \le C_*\bigl[
     \nu\operatorname{Var}_\nu(f\mid Y_B)
    +\nu\operatorname{Var}_\nu(f\mid Y_{\mathcal C_1})\bigr].
 \tag{43.8}
\]
Given \(Y_B\), the first colour is a product of
its site laws, so its variance is at most the
sum of its site conditional variances.
Given \(Y_{\mathcal C_1}\), apply the induction
hypothesis to the other colours, then integrate.
Consistency identifies the resulting one-site
conditionals with those of \(\nu\).
Their coefficient is at most \(C_*^{q-2}\)
before the factor in (43.8). The first-colour
coefficient \(C_*\) is also at most \(C_*^{q-1}\).
This proves (43.7) without a factor equal to
the number of sites.
\end{proof}

Thus the generator \(\sum_i(\Gamma_i-I)\),
with rate one per site, has gap at least \(\gamma_*\).
Random scan with one site selected per step has
gap at least \(\gamma_*/N\), where \(N\) is the
number of sites. These are auxiliary Markov times,
not physical Euclidean times.

For the regional boxed YM law of V32, use
\(m_*=m_0\), \(M_*=K\) from (36.8).
Its completed free-coordinate interaction has
bounded degree. All the physical boundary and
coarse-curvature hypotheses of that proposition
remain required. This yields a volume-uniform
\emph{boxed} local heat-bath gap, stronger than
the surface-cost certificate of V38 on that
restricted domain. It does not extend (36.8)
to the whole compact configuration space.

\subsection{Apply the gap to an interface with fully compact interiors}

Keep V28's fixed packing of separated canonical
blocks \(I_b\), whose remaining free coordinates
form the interface \(\mathcal S\).
There is no interaction term meeting two different
packed interiors. Let
\[
 \mathcal V_{\mathcal S}
       =[-r_{\mathcal S},r_{\mathcal S}]^{d|\mathcal S|},
 \qquad
 \nu(dv)=Z_{\mathcal S}^{-1}
               e^{-\Phi^{\rm c}(v)}{\bf1}_{\mathcal V_{\mathcal S}}\,dv.
 \tag{43.9}
\]
The potential \(\Phi^{\rm c}\) is exactly (32.13):
every packed interior was integrated over its
full compact group domain. The coordinate Haar
factor is already included in \(\Phi^{\rm c}\)
and must not be inserted again.

Under all the fixed-template, small-curvature
and coupling hypotheses of V28, (32.16) gives
\[
              (m/2)\beta I\le(\Phi^{\rm c})''\le K\beta I
                      \quad\hbox{on }\mathcal V_{\mathcal S}.
 \tag{43.10}
\]
Its interaction terms are the original interface
terms and the exact \(-\log Z_b^{\rm c}(v_{J_b})\).
The sets \(J_b\) have fixed size and bounded overlap,
so the effective interface interaction graph has
degree at most a fixed \(\Delta_{\mathcal S}\),
independent of the number of packed blocks.
Use a proper colouring of that actual graph.

Theorem~\ref{thm:f12v39-colours}, now with
\(m_*=m/2\), \(M_*=K\), yields
\[
 \operatorname{Var}_\nu(F)\le\gamma_{\mathcal S}^{-1}
       \sum_{i\in\mathcal S}
               \nu\operatorname{Var}_\nu(F\mid v_{i^c}),
 \qquad \gamma_{\mathcal S}>0
                 \ \hbox{independent of volume}.
 \tag{43.11}
\]
The constant may be very small and depends on
the fixed packing and completion. It is uniform
for \(\beta\) above V28's threshold.
No probability that all packed interiors are
simultaneously in their local boxes was used.

\subsection{A normalized local lift back to the joint law}

The exact joint law with only the interface
restricted is
\[
 \mu^\sharp(dx,dv)=\nu(dv)
                    \prod_b \pi_b(dx_b\mid v_{J_b}).
 \tag{43.12}
\]
Each \(x_b\) ranges over \(G^{I_b}\).
In finite volume this is the original compact
fibre law conditioned on the declared interface
cube, not on good packed interiors.
The factorization is V28's exact normalized
conditional identity.

For \(i\in\mathcal S\), define the affected blocks
and a compound update set by
\[
 \mathcal B(i)=\{b:i\in J_b\},\qquad
                    A_i=\{i\}\cup\bigcup_{b\in\mathcal B(i)} I_b .
 \tag{43.13}
\]
Let \(\widehat\Gamma_i\) be the \emph{actual}
\(\mu^\sharp\)-heat bath on \(A_i\).
It can be sampled as follows: retain \(v_{i^c}\)
and the unaffected interiors; sample \(v_i\)
from its conditional law under \(\nu\), and
then sample all affected \(x_b\)'s from
\(\pi_b(\,\cdot\mid v_{J_b})\).
The unaffected factors in (43.12) do not depend
on \(v_i\). They therefore cancel from its
conditional normalizer. This verifies the sampling
description; it is not an informal collapsed update
with dependent old interiors left in place.

For any \(f\in L^2(\mu^\sharp)\), put
\(F(v)=\mathbb E_{\mu^\sharp}[f\mid v]\).
Write \(e_D^\sharp(f)=
\mu^\sharp\operatorname{Var}_{\mu^\sharp}(f\mid Y_{D^c})\).

\begin{proposition}[Local lifting of interface variance]
\label{prop:f12v39-lift}
For every interface site \(i\),
\[
          \nu\operatorname{Var}_\nu(F\mid v_{i^c})
                                    \le e_{A_i}^\sharp(f).
 \tag{43.14}
\]
Consequently
\[
 \operatorname{Var}_{\mu^\sharp}(f)
 \le \sum_b e_{I_b}^\sharp(f)
           +\gamma_{\mathcal S}^{-1}
                         \sum_{i\in\mathcal S}e_{A_i}^\sharp(f).
 \tag{43.15}
\]
\end{proposition}
\begin{proof}
Fix \(v_{i^c}\). Let \(u\) denote the unaffected
interiors and \(z\) the affected ones.
By (43.12), the conditional law factors as
\[
 \lambda(du\mid v_{i^c})\,
       \nu(dv_i\mid v_{i^c})\,
       \prod_{b\in\mathcal B(i)}\pi_b(dz_b\mid v).
 \tag{43.16}
\]
In particular the law of \(u\) is independent of
\(v_i\). Put \(h(v_i,u)=\mathbb E[f\mid v,u]\),
where the expectation integrates \(z\).
Conditional Jensen and the total-variance identity give
\[
 \begin{split}
 \operatorname{Var}_{v_i}\bigl(\mathbb E_u h(v_i,u)\bigr)
 &\le\mathbb E_u\operatorname{Var}_{v_i}(h(v_i,u))\\
 &\le\mathbb E_u\operatorname{Var}_{v_i,z}(f(v_i,z,u)).
 \end{split}
 \tag{43.17}
\]
The right side is the conditional \(A_i\)-heat-bath
energy. Integrate \(v_{i^c}\) to prove (43.14).
This proof uses the same normalized joint law
on both sides.

Next decompose
\[
 \operatorname{Var}_{\mu^\sharp}(f)
      =\mu^\sharp\operatorname{Var}(f\mid v)
                                  +\operatorname{Var}_\nu(F).
 \tag{43.18}
\]
Given \(v\), the interiors are a product of full
compact block laws. Product variance tensorization
bounds the first term by \(\sum_b e_{I_b}^\sharp(f)\).
Use (43.11) and (43.14) for the second term.
\end{proof}

\begin{corollary}[A local joint-law gap with unrestricted packed interiors]
\label{cor:f12v39-joint-gap}
In every finite volume satisfying V28's hypotheses,
the reversible generator
\[
 \mathcal L^\sharp=
      \sum_b(\Gamma_{I_b}^\sharp-I)
           +\sum_{i\in\mathcal S}(\widehat\Gamma_i-I)
 \tag{43.19}
\]
has spectral gap at least \(\gamma_{\mathcal S}\)
under \(\mu^\sharp\). The update diameters and
numbers of coordinates in each update are bounded
independently of ambient volume.
\end{corollary}
\begin{proof}
The bound \(\gamma_{\mathcal S}\le1\) and (43.15)
give the claimed Poincare inequality for the sum
of the rate-one update forms.
The templates have fixed \(|I_b|\), each \(J_b\)
has fixed size, and a fixed interface site belongs
to only a bounded number of these sets.
Thus each \(A_i\) has uniformly bounded size and
geometric diameter. The potentials and conditional
kernels have the corresponding finite dependence
range. All updates are genuine conditional heat
baths of \(\mu^\sharp\), hence reversible.
\end{proof}

This is a local rate-one generator. If instead
exactly one of its \(N_{\rm upd}\) updates is
chosen uniformly per discrete step, the resulting
gap lower bound is \(\gamma_{\mathcal S}/N_{\rm upd}\).

\subsection{The surviving restriction is precise}

The result does not condition packed interiors
to be simultaneously close to their section values.
Their compact complements are integrated exactly,
and local full-block heat baths can visit them.
The surviving restriction is the product cube
for \emph{every interface coordinate}.

The original unrestricted fibre measure is not
\(\mu^\sharp\). A volume-uniform spectral gap for
\(\mu^\sharp\) cannot be transferred simply by
declaring the interface complement rare.
V33's positive finite-temperature bad-site
probabilities already warn against a global
all-good restriction. The constraint removal
requires a relative energy comparison, not a
union bound growing with interface volume.
Likewise the retained coarse field \(C\) still
obeys V28's fixed small-curvature condition.

\begin{decision}[V40 target and scheduled audit]
At V40, inspect the latest SM, NS and parallel YM
notes and review the direction. The immediate
mathematical question is whether guard refresh
can remove the interface cube while retaining
the local joint-law comparison (43.15).
The difficulty must be addressed on the exact
interface potential and its normalized compound
updates. Neither arbitrary exterior values outside
the cube nor an unrestricted interface Hessian
bound have been proved here.
\end{decision}

\subsection{Exact lift checks and append record}

The script
\begin{center}\small\ttfamily
scripts/verify\_ym\_f12\_v39\_interface\_lift.py
\end{center}
uses a 16-state joint law with two interface
coordinates and two conditionally independent
interiors. One interior depends on both interface
coordinates, the other on only one.
It verifies the exact normalized interface
marginal in each compound update, both inequalities
(43.14), the interior tensorization and the final
joint form comparison by rational
positive-semidefinite elimination.
It also checks that omitting an affected interior
from the resampling set invalidates the comparison
in that example. All checks passed.
The finite example checks the lift algebra, not
the Yang--Mills Hessian hypotheses or physical gap.

V38 had 592153 bytes and SHA--256
\begin{center}\scriptsize\ttfamily
D813C2062493183BA5D2BAC12EE4D9E6AB14D23F70B0EB6FB76539C6E5325463.
\end{center}
Its complete body is retained apart from the
title version. Only active V38 is replaced by
V39. All archives, companion notes and earlier
scripts are retained.

\begin{firewall}[Position after V39]
The boxed fibre and the exact chart-restricted
interface have volume-uniform local heat-bath
gaps under their verified hypotheses.
The interface result lifts to local updates
with all packed interiors fully compact.
An unrestricted-interface gap, unrestricted
fibre clustering, scale-stable effective
interactions, continuum construction and the
physical Yang--Mills mass gap remain unproved.
\end{firewall}

\section{V40: audit and an exact unrestricted-interface reduction}
\label{sec:f12v40}

\subsection{Scheduled companion audit}

The latest active companion files were inspected,
including their conclusions and next-step assessments.
They are unchanged from V35:
\begin{center}\small
\begin{tabular}{lrr}
\toprule
Companion & Version & Bytes\\
\midrule
SM--Einstein & 72 & 607372\\
NS stability & 26 & 278283\\
Parallel YM & 5 & 54174\\
\bottomrule
\end{tabular}
\end{center}
\begin{center}\scriptsize\ttfamily
SM: F44F9745D43379E1672C5550411B58E97195A4C9278592D221DF5D3F644DC1F5\\
NS: 82C887F8C2C69E4F28FAD7C3DC8DE5B9099CB5A4BF431EBA1FFF3E91935978AD\\
YM: 306F1BB84CFA8A8D47EA569EC17E9545F9867C8D32B548EC9D9C444B4F3C0512
\end{center}
SM V72 remains a fixed-mode, fixed-time many-copy
limit, with ultraviolet covariance control still
required. NS V26 gives the rank-one derivative
norms and leaves its full stability theorem open.
The parallel YM note rules out the proposed
volume-uniform tightness of an extensive total
integer charge under nontrivial tensorization.
None supplies an unrestricted interface gap.
The useful SM archive theorem was transferred
with its hypotheses checked in V39; it is not
re-imported as an unconditional physical result.

\subsection{Direction review after ten versions}

V31--V32 established buffered compact conditional
rarity and finite-scale core decoupling. V33--V35
constructed actual adaptive couplings and relative
bad-core control, while exposing the changed
boundary-layer problem. V36--V38 moved to exact
Dirichlet forms, controlled the rare annular loss
by guard refresh, and localized that refresh with
an explicit surface cost. V39 obtained a uniform
local gap while restricting the interface chart,
with packed interiors fully compact.

The last result is not yet on the unrestricted
measure. The appropriate next reduction should
therefore be made on the unrestricted law itself.
We prove below that the full-interface gap is
exactly the gap of the compound update used in
the lift, and is comparable to the ordinary
free-coordinate heat-bath gap with volume-independent
constants at fixed \(\beta\). We also transfer
V31's conditional rarity to that unrestricted
interface, without requiring an aligned cut.

\subsection{The full compact interface, without a chart constraint}

Keep V28's fixed packing of disjoint interiors
\(I_b\), with no interaction term meeting two
interiors. Let \(\mathcal S\) be the remaining
free coordinates. In this section \(v\) ranges
over its entire compact group product, not a
Lie-coordinate cube.
The exact normalized factorization is
\[
 \mu(dx,dv)=\bar\nu(dv)
                   \prod_b\pi_b(dx_b\mid v_{J_b}).
 \tag{44.1}
\]
The full marginal \(\bar\nu\) has Haar-relative
potential \(W_{\mathcal S}^{\rm c}\) from (32.4).
All laws in (44.1) are unrestricted compact laws.
Only the algebraic part of V28 is used here;
no off-chart Hessian bound is assumed.

For \(i\in\mathcal S\), retain the affected set
\(A_i=\{i\}\cup\bigcup_{b:i\in J_b}I_b\).
Let \(\gamma_i\) be the full \(\bar\nu\)
one-interface-coordinate heat bath, and let
\(\widehat\Gamma_i\) be the full \(\mu\)
heat bath on \(A_i\). As in the normalized
calculation of V39, it samples \(v_i\) by
\(\gamma_i\) and resamples all affected interiors.
Unaffected interior factors cancel from its
conditional normalizer. That calculation never
used the interface chart and remains valid here.

Use the nonnegative rate-one generators
\[
 \begin{split}
 \mathsf H_{\mathcal S}&=\sum_{i\in\mathcal S}(I-\gamma_i),\\
 \mathsf H_{\rm aug}&=
      \sum_b(I-\Gamma_{I_b})+
                       \sum_{i\in\mathcal S}(I-\widehat\Gamma_i).
 \end{split}
 \tag{44.2}
\]
All statements below are finite-volume statements;
uniformity over such volumes is specified separately.
There are no unbounded-generator domain questions
hidden in the operator identities.
For a nonnegative reversible generator \(\mathsf H\),
write
\[
 \lambda(\mathsf H)
    =\inf_{\substack{f\perp1\\\|f\|_2>0}}
                  \frac{\langle f,\mathsf Hf\rangle}{\|f\|_2^2}.
 \tag{44.3}
\]

\begin{theorem}[Exact unrestricted gap identity]
\label{thm:f12v40-exact-gap}
Suppose the interface is nonempty and has a
nonconstant coordinate observable. Then
\[
                 \lambda(\mathsf H_{\rm aug})
                            =\lambda(\mathsf H_{\mathcal S}).
 \tag{44.4}
\]
The subspace of functions of the interface is
reducing for \(\mathsf H_{\rm aug}\).
On its orthogonal complement the generator
has spectrum in \([1,\infty)\).
\end{theorem}
\begin{proof}
Let \(J:L^2(\bar\nu)\to L^2(\mu)\) be the isometry
\(JF(x,v)=F(v)\), and let \(\Pi=JJ^*\) be
conditional expectation given \(v\).
Interior heat baths fix \(JF\), and the exact
compound kernel gives
\(\widehat\Gamma_iJF=J\gamma_iF\). Hence
\[
 \mathsf H_{\rm aug}J=J\mathsf H_{\mathcal S},
 \qquad
 J^*\mathsf H_{\rm aug}=\mathsf H_{\mathcal S}J^*,
 \qquad
 [\mathsf H_{\rm aug},\Pi]=0.
 \tag{44.5}
\]
The second identity is the adjoint of the first,
and they imply the third.

For \(f\in\ker\Pi\), conditional product
tensorization over the interiors in (44.1) gives
\[
 \|f\|_2^2=\mu\operatorname{Var}(f\mid v)
       \le\sum_b\mu\operatorname{Var}(f\mid Y_{I_b^c})
       \le\langle f,\mathsf H_{\rm aug}f\rangle.
 \tag{44.6}
\]
Thus the orthogonal complement has gap at least
one. The restriction to the interface subspace
is unitarily equivalent to \(\mathsf H_{\mathcal S}\).

Finally \(\lambda(\mathsf H_{\mathcal S})\le1\).
Indeed choose a nonconstant function of one
interface coordinate. Only that coordinate's
update contributes to its energy, and conditional
variance is at most total variance.
The smaller of the two reducing-subspace gaps
is consequently the interface gap, proving (44.4).
The argument remains valid if that gap has
infimum zero. Degenerate empty-interface systems
are simply products of the independent interiors
and are excluded from the stated identity.
\end{proof}

In particular, integrating the packed interiors
does not by itself improve the slow spectrum.
For interface observables, the semigroup identity
\(e^{-t\mathsf H_{\rm aug}}J
=Je^{-t\mathsf H_{\mathcal S}}\) follows directly
from (44.5). All spectral modes below one in
the augmented generator belong to the interface
subspace. This is an exact identification, not
only a sufficient lifting estimate.

\subsection{Comparison with ordinary unrestricted fibre updates}

Let \(\Gamma_k\) now denote the actual full
single-free-coordinate conditional heat bath
of \(\mu\), and put
\[
             \mathsf H_{\rm site}=\sum_k(I-\Gamma_k).
 \tag{44.7}
\]
These coordinates include both packed interiors
and interface sites. Each group coordinate,
not each of its local chart components, is one site.

The fixed packing supplies finite constants
\[
 \begin{gathered}
 n_*=\max\{1,\sup_b|I_b|\},\qquad
 q_*=\sup_b|J_b|,\qquad M_{\rm inc}=1+q_*,\\
 \ell_*=\text{a common owner-box side for every }I_b\text{ and }A_i,\\
                  C_\beta=
               \exp(C_{\rm surf}v_\beta\ell_*^3).
 \end{gathered}
 \tag{44.8}
\]
Here \(C_{\rm surf}\) and \(v_\beta\) are the
verified full-compact constants of V38.
These constants do not depend on ambient volume.
The comparison \(C_\beta\) may grow exponentially
with \(\beta\); that dependence is retained.

\begin{proposition}[Unrestricted local-form comparison]
\label{prop:f12v40-sites}
On \(L^2(\mu)\), in quadratic-form order,
\[
 \mathsf H_{\rm site}\le n_*\mathsf H_{\rm aug},
 \qquad
 \mathsf H_{\rm aug}\le
                    M_{\rm inc}C_\beta\,\mathsf H_{\rm site}.
 \tag{44.9}
\]
Consequently
\[
 \frac{\lambda(\mathsf H_{\mathcal S})}
             {M_{\rm inc}C_\beta}
 \le\lambda(\mathsf H_{\rm site})
 \le n_*\lambda(\mathsf H_{\mathcal S}).
 \tag{44.10}
\]
\end{proposition}
\begin{proof}
Write \(e_D(f)=\mu\operatorname{Var}(f\mid Y_{D^c})\).
If \(k\in D\), conditional variance monotonicity
gives \(e_{\{k\}}(f)\le e_D(f)\).
For interior sites, sum this over their unique
packed block, obtaining
\(\sum_{k\in I_b}e_{\{k\}}\le |I_b|e_{I_b}\).
For interface \(i\), use
\(e_{\{i\}}\le e_{A_i}\).
This proves the first inequality.

For the other direction, V38's full-compact
finite-region comparison, applied conditionally
on the complement of \(D=I_b\) or \(D=A_i\),
gives
\(e_D\le C_\beta\sum_{k\in D}e_{\{k\}}\).
It applies under arbitrary compact exterior
values; no good chart is required.
Each interior site belongs to its own \(I_b\)
and to exactly those indexed \(A_i\)'s with
\(i\in J_b\), at most \(1+q_*\) updates.
Each interface site belongs only to its own
indexed \(A_i\). Sum these bounds, counting
update multiplicities even if some sets coincide.
This proves the second inequality.
Take variational infima and use (44.4).
\end{proof}

At fixed finite \(\beta\) and fixed packing,
a volume-uniform positive gap for the full
interface generator is therefore equivalent
to one for ordinary full fibre heat baths.
This is not a claim that either uniform gap
has been established.
The comparison constants also do not identify
these auxiliary rates with the physical
Hamiltonian normalization or control cutoff removal.

\subsection{Buffered rarity transfers without an aligned interface cut}

There is a useful positive input on the unrestricted
interface beyond the exact reduction.
Let \(\mathcal A\subset\mathcal S\) be finitely
many specified interface coordinates.
At a fixed threshold \(\tau>0\), take V31's
deterministic coarse neighbourhood
\[
 P=N_{hJ(\tau)}(\operatorname{owners}(\mathcal A)),
 \qquad
 \mathcal S\cap I_P\subset \mathcal T\subset\mathcal S.
 \tag{44.11}
\]
Assume the neighbourhood and affected stencils
fit inside the ambient volume, and impose
V31's curvature and coupling hypotheses there.
There is no requirement that the boundary of
\(P\) align with the packed interiors.

\begin{proposition}[Uniform buffered interface conditional bound]
\label{prop:f12v40-interface-rarity}
For every compact interface exterior \(z\),
\[
 \bar\nu\!\left(
    \|v_i-J_i(C)\|_{\rm HS}\ge\tau
       \text{ for all }i\in\mathcal A
       \,\middle|\,v_{\mathcal S\setminus\mathcal T}=z
             \right)
 \le \exp\!\left(-\frac{\beta\tau^2|\mathcal A|}{16A}\right).
 \tag{44.12}
\]
The conditioning is that of the actual unrestricted
interface marginal.
\end{proposition}
\begin{proof}
Let \(B_{\mathcal A}\) be the event in (44.12).
V31 gives the corresponding pointwise conditional
bound for \(\mu\) given all free variables outside
\(I_P\).
By (44.11), every interface coordinate outside
\(\mathcal T\) is outside \(I_P\). Therefore
\[
 \begin{split}
 \bar\nu(B_{\mathcal A}\mid v_{\mathcal S\setminus\mathcal T})
 &=\mu(B_{\mathcal A}\mid v_{\mathcal S\setminus\mathcal T})\\
 &=\mu\!\left[
       \mu(B_{\mathcal A}\mid Y_{I_P^c})
                \,\middle|\,v_{\mathcal S\setminus\mathcal T}\right].
 \end{split}
 \tag{44.13}
\]
The sigma-field inclusion justifies the tower
identity even if \(P\) cuts through a packed
interior. Apply V31 inside the last expectation.

This initially gives the statement almost surely.
In finite volume the full interface density is
positive and smooth on its compact product.
The canonical conditional probability of this
fixed event is continuous in the exterior:
dominated convergence applies to its numerator
and its strictly positive normalizer.
The exterior marginal has full support, so the
bound extends to every compact \(z\).
\end{proof}

This is a probability transfer by coarsening a
conditioning sigma field. It is not an assertion
that every irregular interface cut has the
regional density identity proved under V29's
alignment assumptions. That stronger identification
is unnecessary for (44.12).
The same almost-sure tower argument applies
to interface marginals of compatible DLR states.

\subsection{The global chart restriction is not a small-TV perturbation}

Let \(\mathcal C_N\) be V39's product chart
event on \(N=|\mathcal S|\) interface coordinates.
For each coordinate its complement has a fixed
positive Haar measure. V33's full-compact
single-site conditional lower bound, with this
chart complement, gives a fixed
\[
 p_{\rm out}(\beta)>0,\qquad
                     \mu(\mathcal C_N)=\bar\nu(\mathcal C_N)
                                   \le(1-p_{\rm out}(\beta))^N .
 \tag{44.14}
\]
For example one may take a positive Haar factor
times \(e^{-K(\beta+1)}\), using the bounded
oscillation of the actual one-site conditional
potential. The repeated-conditioning proof is
the one already given in V33; it applies to
this interface subset as well.

Write \(\bar\nu^\sharp=\bar\nu(\,\cdot\mid\mathcal C_N)\)
and \(\mu^\sharp=\mu(\,\cdot\mid\mathcal C_N)\).
For any probability law and positive-probability
event \(\mathcal C\), direct integration of the
density \({\bf1}_{\mathcal C}/\mu(\mathcal C)\) gives
TV distance \(1-\mu(\mathcal C)\).
Consequently
\[
 \begin{split}
 \|\bar\nu^\sharp-\bar\nu\|_{\rm TV}
   &=\|\mu^\sharp-\mu\|_{\rm TV}
     =1-\bar\nu(\mathcal C_N)
     \ge1-(1-p_{\rm out}(\beta))^N,\\
 D(\bar\nu^\sharp\|\bar\nu)
   &=-\log\bar\nu(\mathcal C_N)
     \ge-N\log(1-p_{\rm out}(\beta)).
 \end{split}
 \tag{44.15}
\]
Here \(D\) is relative entropy.
At fixed finite \(\beta\), the global TV distance
tends to one as \(N\) tends to infinity.
Thus V39 cannot be transferred by treating its
chart-restricted law as a globally small-TV
perturbation of the unrestricted law.

This conclusion does not rule out useful
fixed-window estimates, a relative Dirichlet-form
comparison, or a uniform unrestricted gap.
It concerns global measure distance, not a
lower bound on the physical mass or a claim
of failure of the Yang--Mills programme.

\begin{decision}[The next estimate belongs to the unrestricted interface]
The live analytical target is now on the exact
full marginal:
\[
 \operatorname{Var}_{\bar\nu}(F)
   \le C(\beta)\sum_{i\in\mathcal S}
                 \bar\nu\operatorname{Var}_{\bar\nu}(F\mid v_{i^c}),
 \qquad C(\beta)\ \hbox{independent of ambient volume}.
 \tag{44.16}
\]
It remains unproved. Equations (44.4) and (44.10)
show exactly what it would imply for the full fibre
and prevent a gain from being attributed merely
to integrating the packed interiors.

Use the buffered conditional input (44.12) to
investigate relative energy localization of
bad-interface clusters. Keep the unrestricted
normalizers and the compound-update sets explicit.
Do not replace this problem by the already solved
chart-restricted gap, and do not repeat the global
TV perturbation route excluded by (44.15).
The next companion audit is V45; the next
ten-version direction review is V50.
\end{decision}

\subsection{Exact checks and append record}

The script
\begin{center}\small\ttfamily
scripts/verify\_ym\_f12\_v40\_unrestricted\_reduction.py
\end{center}
checks a 16-state unrestricted joint law.
It verifies that the interface projection
commutes with the augmented generator, that
the vertical sector has lower bound one, and
that a horizontal test attains the common
interface/augmented gap \(1/2\).
The ordinary-update comparisons are checked
by rational positive-semidefinite elimination;
the example's maximum conditional block density
ratio is 18, indexed update incidence is three,
and the resulting ordinary-gap certificate is
\(1/108\). It also verifies the tower identity
for a conditional bound after exterior information
is discarded. All checks passed.
These are exact algebra tests, not verification
of a volume-uniform unrestricted Yang--Mills gap.

V39 had 608461 bytes and SHA--256
\begin{center}\scriptsize\ttfamily
0B7E6A84F595436FA1FE17B5691EB7D8F9796F5ACD17F8974C0E17E42CC6F87B.
\end{center}
Its complete body is retained apart from the
title version. Only active V39 is replaced by
V40; all archives, companion notes and earlier
scripts are retained.

\begin{firewall}[Position after V40]
The full compact interface has an exact spectral
identification with the augmented full-fibre
updates and a volume-independent-at-fixed-coupling
comparison with ordinary fibre updates.
It inherits uniform buffered conditional rarity.
Its unrestricted volume-uniform gap remains open.
All-distance fibre clustering, scale-stable
effective interactions, continuum construction
and the physical Yang--Mills mass gap remain
unproved.
\end{firewall}

\section{V41: fixed-core density convexity under arbitrary compact exteriors}
\label{sec:f12v41}

\subsection{An upstream improvement of the conditional input}

Buffered rarity alone does not prevent an \(L^2\)
test from concentrating on a rare defect. The
global chart restriction is also not a small-TV
perturbation, by V40. We therefore strengthen
the local conditional information before making
another energy comparison.

V32 gives exponentially close core marginals in
total variation at a buffer scale proportional
to \(\beta\). For a \emph{fixed-dimensional} core,
local smoothness of the actual density upgrades
this to \(C^2\) control. A further, explicitly
proved lower-density bound permits passage to
the logarithm. The result is a small-core
effective-potential Hessian bound under arbitrary
compact data outside the buffer.

The buffer is fully integrated. This is not a
claim about a one-site conditional with its
immediate neighbours fixed, nor a global Hessian
bound on the unrestricted interface.

\subsection{The fixed core and one common reference marginal}

Fix a nonempty finite set \(\mathcal A\) of free
group coordinates in a bounded neighbourhood of
the origin. Its real coordinate dimension
\(n=d|\mathcal A|\) is fixed throughout the argument.
Use the section chart
\(Y_i=J_i(C)\exp u_i\), with the fixed radius
\(R_*\) of V32. Write
\[
 D_{\mathcal A}=[-R_*,R_*]^n,\qquad
 K=[-R_*/2,R_*/2]^n,\qquad
 K^+=[-3R_*/4,3R_*/4]^n.
 \tag{45.1}
\]
Choose V32's concentric \(Q,R,P\) with radii
\(\lfloor r/4\rfloor,r,2r\), where
\(r=\lfloor\beta\rfloor\). For sufficiently large
\(\beta\), \(\mathcal A\subset I_Q\).
Assume all the affected coarse-curvature and
embedding hypotheses of V32, including absence
of wrapping for these stencils.

Let \(p_\xi(u)\) be the density, relative to
Lebesgue measure in this core chart, of the
\(\mathcal A\)-marginal of the actual compact law
\(\pi_{P,C}^{\xi}\). The exterior \(\xi\) is arbitrary
compact free data. This density is not renormalized
on \(D_{\mathcal A}\).

Use section exterior data \(\xi_0=J(C)\) for a
reference law on \(P\), and condition that reference
law on its entire fixed coordinate box \(\mathcal D_P\).
Let \(q(u)\) be its normalized core marginal.
All reference integrations outside the core
are over their fixed coordinate cubes.
V32's good-boundary complement estimate on \(P\)
and \(|P|=O(\beta^4)\) give
\[
 \pi_{P,C}^{\xi_0}(\mathcal D_P^c)\le C e^{-c_{\rm box}\beta}.
 \tag{45.2}
\]
V32's arbitrary-exterior core-TV estimate,
marginalized further to \(\mathcal A\), therefore gives
\[
 \sup_\xi\|(\pi_{P,C}^{\xi})_{\mathcal A}-q\|_{\rm TV}
                                      \le C e^{-c\beta}
 \tag{45.3}
\]
for a fixed \(c>0\), after decreasing the exponent.
In this formula \(q\) is viewed as a probability
law on the core group variables supported in its
chart. It is the same reference for every \(\xi\)
at fixed \(C,\beta,P\).

The reference measured potential is strongly convex:
\[
          m_0\beta I\le(-\log q)''\le K_0\beta I
                                 \quad\hbox{on }D_{\mathcal A}.
 \tag{45.4}
\]
For the lower bound, apply the convex
marginalization argument of V28 to V32's full
boxed Hessian on \(P\). For the upper bound,
the exact marginal identity is
\((-\log q)''=\mathbb E V_{uu}
-\operatorname{Cov}(V_u,V_u)\le K_0\beta I\).
These are properties of the reference marginal,
not yet of \(p_\xi\).

\subsection{Density derivatives cost only fixed-core powers of \(\beta\)}

\begin{proposition}[Uniform local density bounds]
\label{prop:f12v41-density-derivatives}
For multi-indices \(\alpha\) of order at most three,
\[
 |\partial^\alpha p_\xi(u)|
       \le C_\alpha(\beta+1)^{|\alpha|}p_\xi(u),
 \qquad
 |\partial^\alpha q(u)|
       \le C_\alpha(\beta+1)^{|\alpha|}q(u)
                       \quad(u\in D_{\mathcal A}).
 \tag{45.5}
\]
On \(K^+\), for sufficiently large \(\beta\),
\[
       \sup_\xi\|p_\xi\|_{C^3(K^+)}+\|q\|_{C^3(K^+)}
                                    \le C\beta^{n+3}.
 \tag{45.6}
\]
Constants depend on the fixed core, not on its
exterior or the number of integrated coordinates.
\end{proposition}
\begin{proof}
Use Lebesgue measure in \(u\) and product Haar
for the other compact coordinates. In the
reference integral, use its fixed cube-coordinate
measures instead. The core Haar Jacobian is
included in the measured potential in each case.

Only a bounded number of local potential terms
depends on the fixed core. V18's global smooth
completion and fixed stencil overlap imply that
each of its core derivatives, through order three,
is bounded by \(C_\alpha(\beta+1)\), uniformly
over all other compact variables.
The same bound holds in the reference box.
Differentiating the exponential integrand gives
products of these derivatives of total degree
at most \(|\alpha|\). Integrating their absolute
values and dividing by the unchanged global
normalizer proves (45.5). No derivative of a
volume-sized sum independent of the core remains.

In particular \(|\nabla\log p_\xi|\) and
\(|\nabla\log q|\) are at most \(C\beta\).
For \(u\in K^+\), a Euclidean ball of radius
\(c'/\beta\) around \(u\) is contained in
\(D_{\mathcal A}\) at large \(\beta\).
On this ball the density is at least its value
at \(u\) times a fixed positive constant.
Its total mass is at most one, and the ball
has volume proportional to \(\beta^{-n}\).
Thus \(p_\xi(u),q(u)\le C\beta^n\).
Combine with (45.5) to obtain (45.6).
All differentiations have fixed integration domains.
\end{proof}

The estimate would not have this form if the
core dimension grew with the buffer volume.
Keeping \(\mathcal A\) fixed is a substantive
condition, not just notation.

\subsection{From total variation to \(C^2\) density control}

\begin{proposition}[Fixed-core interpolation]
\label{prop:f12v41-interpolation}
There are fixed \(C,c_1>0\) such that
\[
            \sup_\xi\|p_\xi-q\|_{C^2(K)}
                                      \le C e^{-c_1\beta}.
 \tag{45.7}
\]
One may choose \(c_1=c/[4(n+3)]\), with \(c\)
from (45.3), after increasing the coupling threshold.
\end{proposition}
\begin{proof}
Put \(g=p_\xi-q\). Pullback by the fixed core
chart and (45.3) give
\(\|g\|_{L^1(D_{\mathcal A})}\le C e^{-c\beta}\).
Choose a smooth compactly supported mollifier
of integral one, with bandwidth \(h\).
For \(u\in K\) and \(h\) smaller than the margin
to the boundary of \(K^+\), its convolution
uses only that larger cube. For \(|\alpha|\le2\),
differentiating the mollifier and estimating the
local approximation error give
\[
 \|\partial^\alpha g\|_{L^\infty(K)}
 \le C h^{-n-2}e^{-c\beta}
                            +C h\beta^{n+3}.
 \tag{45.8}
\]
The error uses the derivatives of order at most
three in (45.6). There is no contribution from
the boundary of the reference chart, which
lies outside the convolution neighbourhood.

Take \(h=\exp[-c\beta/(2(n+3))]\).
The first term has exponential rate
\(c(n+4)/(2(n+3))\); the second has rate
\(c/(2(n+3))\) times a polynomial.
Absorb that polynomial by halving the smaller
rate. This proves (45.7), uniformly in \(\xi\).
\end{proof}

\subsection{A lower density before taking logarithms}

Density closeness does not automatically imply
log-density closeness where the reference is small.
We establish a sufficiently strong lower bound
on a smaller, fixed core cube.

Let \(L_0\) be a constant such that
\(|\nabla\log q|\le L_0\beta\) on \(K^+\).
Choose a fixed \(a>0\) with
\[
 K_a=[-a,a]^n\subset K,\qquad
                         2L_0\sqrt n\,a\le c_1/8.
 \tag{45.9}
\]
Choose a fixed group-distance threshold \(\tau_a>0\)
so that distances below \(\tau_a\) at all core
coordinates imply \(u\in K_a\).
The section chart is uniformly invertible near
the identity, so such a threshold exists.

We now decrease the coarse-curvature threshold
to the fixed positive value
\[
 \epsilon_a=\min\{\epsilon_{\rm base},\epsilon_{\tau_a}\},
 \tag{45.10}
\]
where \(\epsilon_{\rm base}\) covers the V32 and
reference-box assumptions already used.
Take \(\beta\) above the corresponding V31
threshold and large enough that its
\(hJ(\tau_a)\)-neighbourhoods of the core fit
inside \(P\). These choices may be conservative.
The order is important: first choose the core
and \(c_1\), then \(a,\tau_a,\epsilon_a\), and
only then the lower threshold on \(\beta\).
No parameter here decreases with ambient volume.

V31 and a union bound over the fixed core show
that the unboxed section-exterior reference
places probability tending to one in \(K_a\).
Conditioning it on \(\mathcal D_P\) changes this
bound by at most division by
\(1-\pi_{P,C}^{\xi_0}(\mathcal D_P^c)\).
Thus, at sufficiently large \(\beta\),
\(q(K_a)\ge1/2\).
There is a point \(u_0\in K_a\) with
\(q(u_0)\ge[2\,\operatorname{vol}(K_a)]^{-1}\).
Integrate the log-gradient bound along the
straight segment in \(K_a\) from \(u_0\) to \(u\).
Using (45.9) gives
\[
        \inf_{u\in K_a}q(u)\ge c_a e^{-c_1\beta/8},
                  \qquad c_a=[2\,\operatorname{vol}(K_a)]^{-1}.
 \tag{45.11}
\]
This lower bound is obtained from the actual
normalized reference probability, not from an
assumed pointwise Gaussian approximation.

\subsection{The arbitrary-exterior effective potential is locally convex}

\begin{theorem}[Uniform log-density and Hessian control]
\label{thm:f12v41-core-convexity}
Under the fixed-core geometry, the fixed
small-curvature threshold \(\epsilon_a\), and
all the coupling thresholds specified above,
\[
 \sup_\xi
 \|-\log p_\xi+\log q\|_{C^2(K_a)}
                                \le C e^{-c_2\beta},
 \qquad c_2=c_1/4
 \tag{45.12}
\]
after increasing the threshold on \(\beta\).
In particular, for arbitrary compact exterior \(\xi\),
\[
       (m_0/2)\beta I
            \le(-\log p_\xi)''
            \le 2K_0\beta I
                              \quad\hbox{on }K_a.
 \tag{45.13}
\]
All constants are independent of the exterior
and of the surrounding ambient volume.
\end{theorem}
\begin{proof}
Equations (45.7) and (45.11) imply
\(|p_\xi-q|/q\to0\) uniformly on \(K_a\).
In particular \(q/2\le p_\xi\le3q/2\)
for sufficiently large \(\beta\).
Use the identities
\[
       \partial_j\log p=\frac{\partial_jp}{p},
 \qquad
       \partial_{jk}\log p
          =\frac{\partial_{jk}p}{p}
                    -\frac{\partial_jp\,\partial_kp}{p^2}.
 \tag{45.14}
\]
Subtract the corresponding expressions for \(q\).
The \(C^2\) numerator differences are bounded
by \(Ce^{-c_1\beta}\), while all density
derivatives through order two have polynomial
bounds by (45.6).
Because \(p_\xi\) and \(q\) are relatively
comparable, denominators in these differences
cost at most a fixed polynomial times
\((\inf_{K_a}q)^{-3}\).
For example
\(|p^{-2}-q^{-2}|\le C|p-q|q^{-3}\)
when \(q/2\le p\le3q/2\).
The zeroth and first derivative bounds are
no worse. By (45.11), the net exponential
rate before absorbing polynomials is at least
\(c_1-3c_1/8=5c_1/8\).
This proves (45.12) with the stated smaller
rate \(c_2=c_1/4\).

Finally combine (45.12) with the reference
Hessian bounds (45.4). The exponentially small
error is smaller than both fixed multiples
of \(\beta\) needed for (45.13).
\end{proof}

Thus the first two core derivatives of the
normalized effective potential are exponentially
insensitive to arbitrary data outside \(P\).
For two exteriors the corresponding \(C^2\)
log-density difference is at most twice the
right side of (45.12). No derivative with
respect to the exterior itself is claimed.

\subsection{Pass the same result to an unrestricted interface conditional}

Let the fixed core \(\mathcal A\) consist of
interface coordinates in V40, and let
\(\mathcal T\supset\mathcal S\cap I_P\).
Condition the unrestricted marginal \(\bar\nu\)
on \(v_{\mathcal S\setminus\mathcal T}=z\).
The density \(p_z^{\rm int}(u)\) of its
\(\mathcal A\)-marginal has the mixture representation
\[
                p_z^{\rm int}(u)=\int p_\xi(u)\,\lambda_z(d\xi).
 \tag{45.15}
\]
Here \(\lambda_z\) is the distribution of all
free data outside \(I_P\), conditional only on
that interface exterior. It is not conditioned
again on the core value \(u\); its mixing weights
are independent of \(u\).
Equation (45.15) is the density version of
V40's sigma-field tower identity.

The reference \(q\) is common to every component.
The uniform derivative bounds permit
differentiation under the mixture, so
\(\|p_z^{\rm int}-q\|_{C^2(K)}\le Ce^{-c_1\beta}\).
Apply the same lower-reference-density and
logarithm argument. Consequently (45.12)--(45.13)
hold with \(p_\xi\) replaced by \(p_z^{\rm int}\),
uniformly over compact interface exterior \(z\).
In finite volume all these versions are given
by positive normalized densities, so the statement
holds for every exterior, not just almost surely.

This does not use the false assertion that
mixtures of strongly log-concave laws are always
strongly log-concave. It uses their \(C^2\)
density closeness to one common positive reference
on the fixed core cube.

\subsection{Why the derivative input cannot be skipped}

On \([-1,1]\), let \(q_\beta(x)\) be proportional
to \(e^{-\beta x^2/2}\), and define
\[
 p_\beta(x)=
 \frac{q_\beta(x)(1-\varepsilon_\beta\cos(k_\beta x))}
      {\int q_\beta(t)(1-\varepsilon_\beta\cos(k_\beta t))\,dt},
 \qquad
       \varepsilon_\beta=2^{-\beta},\quad k_\beta=2^\beta
 \tag{45.16}
\]
for integer \(\beta\ge2\).
The normalizer lies between \(1-\varepsilon_\beta\)
and \(1+\varepsilon_\beta\), and hence
\[
 \|p_\beta-q_\beta\|_{\rm TV}
                  \le\frac{\varepsilon_\beta}{1-\varepsilon_\beta}.
 \]
Nevertheless direct differentiation at zero gives
\[
       (-\log p_\beta)''(0)
          =\beta-\frac{\varepsilon_\beta k_\beta^2}
                            {1-\varepsilon_\beta}<0.
 \tag{45.17}
\]
The higher derivatives of the oscillatory factor
grow exponentially. Thus this example violates
the uniform polynomial smoothness input of
Proposition~\ref{prop:f12v41-density-derivatives}.
It proves that TV closeness alone would not
justify the present convexity conclusion.

\begin{decision}[Use the new input without changing its conditional meaning]
We now have fixed-core normalized potential
control with arbitrary exterior data outside
an integrated mesoscopic buffer. That strengthens
the local input available for interface guard
refresh and removes the need for a good outer
boundary in this particular marginal statement.

It does not control the whole updated buffer,
and it is not the one-site conditional obtained
by fixing that buffer. Therefore it cannot
simply be substituted into V39's finite-colour
proof for the unrestricted interface generator.
The next task is to attach this buffered core
potential to a compatible update or comparison
that also accounts for the changed buffer.
The core dimension must remain fixed in the
interpolation estimate, and the additional fixed
curvature threshold \(\epsilon_a\) must remain visible.
The next companion audit is V45 and the next
direction review is V50.
\end{decision}

\subsection{Checks and append record}

The script
\begin{center}\small\ttfamily
scripts/verify\_ym\_f12\_v41\_density\_upgrade.py
\end{center}
checks the interpolation and logarithm exponent
inequalities by exact rational arithmetic for
five fixed core dimensions. It also checks
the TV bound and negative Hessian formula
in (45.16)--(45.17) at four integer couplings.
All checks passed. These are scalar checks
of the exponent choices and warning example;
the Yang--Mills analytic estimates are proved
above, not numerically certified by that script.

V40 had 624769 bytes and SHA--256
\begin{center}\scriptsize\ttfamily
BD937B9797E88A3F794ACE9CE3637ABCF5EE0F75588BC0555E598B68B585D729.
\end{center}
Its complete body is retained apart from the
title version. Only active V40 is replaced by
V41; archives, companions and earlier scripts
are retained.

\begin{firewall}[Position after V41]
The actual compact fixed-core marginal has
uniform near-section log-density and Hessian
control under arbitrary compact exterior data,
with its mesoscopic buffer fully integrated.
The same holds for the corresponding unrestricted
interface conditional marginal.
The unrestricted volume-uniform interface gap,
all-distance fibre clustering, scale-stable
effective interactions, continuum construction
and the physical Yang--Mills mass gap remain
unproved.
\end{firewall}

\section{V42: a near-one common core component with an exact block lift}
\label{sec:f12v42}

\subsection{A sampling consequence stronger than pairwise TV}

V41 gives a pointwise comparison of every
exterior-conditioned core density with one
common reference. We now use that comparison
to construct a probability component shared
by all those core laws. Its mass tends to
one exponentially in \(\beta\).
The corresponding sampling branch can be
lifted to the actual complete block heat bath.
Only the core forgets the exterior on that
branch; the sampled buffer may still remember it.

The general Doeblin-to-gap implication was
already recorded in f7 V12, and finite Wilson
minorization arguments appear in older archives.
They are not reproved here. The new estimate
is a high-mass common component for the
specific buffered core marginal of V41,
uniform over its arbitrary compact exteriors.
It is not a minorization of a physical
transfer operator or of a complete lattice sweep.

\subsection{A common component with exponentially small remainder}

Keep all of V41's assumptions: the core
\(\mathcal A\) has fixed dimension, the parent
\(P\) has radius proportional to \(\beta\),
the retained field \(C\) is fixed and obeys
the additional threshold \(\epsilon_a\), and
the canonical buffer is integrated completely.
Let \(\mu_\xi\) be the full compact core
marginal for exterior \(\xi\).
Use V41's reference density \(q\), core cube
\(K_a\), and constants.

The proof preceding (45.11) actually gives
an exponential probability estimate, not just
the eventual lower bound \(q(K_a)\ge1/2\):
\[
 w_\beta:=q(K_a)\ge1-C_a e^{-\kappa_a\beta}
                                   \quad\hbox{for some }\kappa_a>0.
 \tag{46.1}
\]
Indeed V31 bounds the probability that any
fixed core coordinate leaves its
\(\tau_a\)-neighbourhood exponentially in
\(\beta\). Dividing by the section-reference
box probability, which tends to one by
(45.2), preserves such a bound.

Choose the uniform upper bound
\(\theta_\beta=C_1e^{-c_2\beta}\) from (45.12),
and define a probability measure on the full
core group space by
\[
 \omega_\beta(du)=
       w_\beta^{-1}q(u){\bf1}_{K_a}(u)\,du,\qquad
 \alpha_\beta=e^{-\theta_\beta}w_\beta,\qquad
                         \varepsilon_\beta=1-\alpha_\beta.
 \tag{46.2}
\]
The chart pushforward is understood in this
formula. No Haar Jacobian is added to \(q\),
which is already a measured coordinate density.

\begin{theorem}[Exterior-uniform common core component]
\label{thm:f12v42-common}
For every compact exterior \(\xi\),
\[
                 \mu_\xi\ge\alpha_\beta\omega_\beta
                                      \quad\hbox{as measures},
 \qquad
                 \varepsilon_\beta\le C_2e^{-\kappa\beta}
 \tag{46.3}
\]
for fixed \(C_2,\kappa>0\).
The component is the same for every \(\xi\)
at the fixed retained \(C\), core and parent.
\end{theorem}
\begin{proof}
On \(K_a\), (45.12) implies
\(p_\xi(u)\ge e^{-\theta_\beta}q(u)\).
Outside that cube the proposed common measure
is zero. Integrating over arbitrary measurable
sets proves the measure inequality.
Also
\[
 1-e^{-\theta_\beta}w_\beta
       \le\theta_\beta+(1-w_\beta)
       \le C_1e^{-c_2\beta}+C_a e^{-\kappa_a\beta}.
 \]
Take the smaller positive exponent.
The probability \(\omega_\beta\) and the weight
\(\alpha_\beta\) may depend on \(C,\beta\) and
the fixed template, but not on its exterior \(\xi\).
\end{proof}

For large finite \(\beta\), the weight lies
strictly between zero and one. The residual
probability kernel is therefore well defined:
\[
       \sigma_\xi
          =\frac{\mu_\xi-\alpha_\beta\omega_\beta}
                         {\varepsilon_\beta},\qquad
       \mu_\xi=\alpha_\beta\omega_\beta
                                  +\varepsilon_\beta\sigma_\xi.
 \tag{46.4}
\]
All conditional densities in question are
measurable in the exterior, so this is a
measurable kernel identity.
The corresponding interface conditional
marginals in (45.15) obey the same
minorization, because a mixture preserves
the common measure in (46.3).

\subsection{Pairwise closeness would not have been enough}

For \(N\ge3\), consider the \(N\) probability
laws on \(\{1,\ldots,N\}\) given by
\[
       \rho_j(k)=\frac{{\bf1}_{\{k\ne j\}}}{N-1}.
 \tag{46.5}
\]
Any two different laws have TV distance
\(1/(N-1)\). Yet a nonnegative measure
dominated by all of them must be zero:
its mass at \(k\) is bounded by \(\rho_k(k)=0\).
Thus small pairwise TV diameter for a family
does not by itself supply a high-mass common
component for the entire family.
V41's simultaneous pointwise lower bound is
the additional input used in (46.3).

\subsection{Lift the split without altering the Gibbs update}

Let \(U=Y_{\mathcal A}\) and
\(V=Y_{I_P\setminus\mathcal A}\).
Disintegrate the actual canonical block law:
\[
 \pi_{P,C}^{\xi}(du,dv)
                    =\mu_\xi(du)\,Q_\xi(dv\mid u).
 \tag{46.6}
\]
The full compact positive densities give
canonical regular conditional kernels.
Substitution of (46.4) yields the exact identity
\[
 \begin{split}
 \pi_{P,C}^{\xi}(du,dv)
   ={}&\alpha_\beta\,\omega_\beta(du)Q_\xi(dv\mid u)\\
      &+\varepsilon_\beta\,\sigma_\xi(du)Q_\xi(dv\mid u).
 \end{split}
 \tag{46.7}
\]
Thus one complete block update can be sampled
by a Bernoulli choice independent of the current
configuration and exterior:
with probability \(\alpha_\beta\), sample
\(U\) from \(\omega_\beta\); otherwise sample
\(U\) from \(\sigma_\xi\).
In either case sample the entire remaining
buffer from \(Q_\xi(\,\cdot\mid U)\).

For two, or a whole family of, exterior data,
use the same Bernoulli choice and the same
\(\omega_\beta\)-sample on the first branch.
Every core then agrees on that branch,
independently of the exteriors.
Residual and buffer samples can be obtained
from their measurable probability kernels
using fixed Borel codings and uniform random
numbers. In particular the construction does
not require separate, incompatible pairwise
optimal couplings.

After averaging over the branch choice,
(46.7) is exactly the original canonical
heat bath, including its normalizers.
It therefore preserves any compatible Gibbs
law. The two branches separately need not
preserve that law. Rerolling until the first
branch occurs, or discarding failed branches,
would change the kernel and is not authorized
by this identity.

Even on the first branch, \(Q_\xi(dv\mid u)\)
can depend on the exterior. Equality of all
core samples is not equality of the full
updated blocks. This is the precise surviving
memory in the construction.

\subsection{Independent refresh flags on separated buffers}

Take a finite collection of parents \(P_j\)
with disjoint updated free sets and with no
interaction term meeting two of those sets.
The fixed finite interaction range allows
such separated collections. Choose fixed core
templates \(\mathcal A_j\subset I_{P_j}\).
Assume the uniform V41 hypotheses around them.

Conditional on all free coordinates outside
their union, the complete parent laws factor
as a product of their canonical laws.
Their simultaneous refresh is therefore a
genuine Gibbs conditional update.
Choose a common lower weight
\(\alpha=1-\varepsilon\le\alpha_{\beta,j}\)
with \(0<\varepsilon\le Ce^{-\kappa\beta}\),
using uniform constants for the fixed templates.
If needed enlarge the displayed upper bound
on \(\varepsilon\) while taking \(\beta\) large
enough that it remains below one.
Each parent law still dominates
\(\alpha\omega_{\beta,j}\).

\begin{proposition}[Independent algorithmic failure flags]
\label{prop:f12v42-flags}
The separated Gibbs refresh admits independent
Bernoulli failure indicators \(B_j\) with
\(\Pr(B_j=1)=\varepsilon\), independent of the
current field and its exterior, such that
\[
       {\bf1}_{\{U'_j\notin K_{a,j}\}}\le B_j
                                  \quad\hbox{for every }j.
 \tag{46.8}
\]
For every specified finite index set \(F\),
\[
             \Pr(B_j=1\text{ for all }j\in F
                          \mid\hbox{current field})
                                  =\varepsilon^{|F|}.
 \tag{46.9}
\]
\end{proposition}
\begin{proof}
Use independent Bernoulli choices in the
product of the splits (46.7), with the common
weight \(\alpha\). On a successful branch
the corresponding core is sampled from
\(\omega_{\beta,j}\), which is supported
inside \(K_{a,j}\). This proves (46.8).
The independent choices give (46.9).
Failure may still produce a good core; no
converse to (46.8) is claimed.
\end{proof}

If the input field has a compatible Gibbs
law, invariance of the complete parallel
refresh implies that the bad-core indicator
field on this \emph{selected separated family}
is stochastically dominated by independent
Bernoulli variables of parameter \(\varepsilon\).
This follows from the explicit coupling (46.8),
not merely from bounds on all-bad intersection
probabilities.
The same conclusion holds for selected
interface cores by taking the corresponding
marginal of the refreshed full field.

For a countable separated family in a
compatible DLR state, apply the finite-family
construction to increasing finite subsets.
Extend each coupling to the remaining
indicator coordinates while preserving the
given bad-core marginal and the independent
Bernoulli marginal. Compactness of the product
of discrete indicator spaces gives a subsequential
coupling limit. Each fixed coordinate satisfies
(46.8) eventually, hence also in the limit.
Thus the sparse-family stochastic domination
extends to that setting.

On an index graph of maximum degree \(\Delta\),
put \(\Delta_*=\max\{2,\Delta\}\).
The rooted connected-set count used in V31
and (46.9) give
\[
 \Pr\{\text{the failure component at a fixed root has size }\ge k\}
                            \le(\Delta_*^2\varepsilon)^k.
 \tag{46.10}
\]
This is a statement about flags, or bad cores,
on the selected separated family. It does not
control a backwards dependency graph through
all buffers and all subsequent updates.

\subsection{Why local common components still do not prove a global gap}

There is a simple counterexample to a stronger,
but invalid, conclusion based only on near-one
one-site common components. Fix \(0<\delta<1\)
and let, for odd \(N\),
\[
 \eta_N=\tfrac12
       \operatorname{Ber}_{\{-1,1\}}((1+\delta)/2)^{\otimes N}
       +\tfrac12
       \operatorname{Ber}_{\{-1,1\}}((1-\delta)/2)^{\otimes N}.
 \tag{46.11}
\]
Every one-site conditional probability of \(+1\)
lies between \((1-\delta)/2\) and \((1+\delta)/2\):
it is a posterior mixture of these two numbers.
Hence every such conditional law contains
the same uniform-sign component with mass
\(1-\delta\), which can be arbitrarily close
to one.

Nevertheless its rate-one single-site heat-bath
gap tends to zero as \(N\) grows.
For \(f={\bf1}_{\{\sum_i s_i>0\}}\), spin-flip
symmetry and odd \(N\) give \(\operatorname{Var}(f)=1/4\).
Updating site \(i\) can change the sign only
when the other \(N-1\) signs sum to zero.
On that event its conditional probability
is exactly \(1/2\). Therefore the Rayleigh
quotient of \(f\) is
\[
 \begin{split}
 \frac{\sum_i\eta_N\operatorname{Var}(f\mid s_{i^c})}
                  {\operatorname{Var}_{\eta_N}(f)}
 &=N\binom{N-1}{(N-1)/2}
                      \left(\frac{1-\delta^2}{4}\right)^{(N-1)/2}\\
 &\le N(1-\delta^2)^{(N-1)/2}\longrightarrow0 .
 \end{split}
 \tag{46.12}
\]
The binomial bound follows from
\(\binom{2m}{m}\le4^m\).
This example has a nonlocal interaction; it
is not a counterexample to the local Yang--Mills
claim. It shows why a common-component estimate
without control of residual influence or
dependency geometry is insufficient.

Likewise (46.3) is a core marginal statement.
It cannot be inserted into the archived
full-kernel Doeblin-to-gap theorem as if
the entire field were refreshed independently
on its successful branch.

\begin{decision}[Next: quantify memory carried by the buffer]
The fixed core now has a simultaneous,
exterior-independent sampling branch with
exponentially small failure probability.
The complete lifted update is exact, and
separated buffers admit independent failure
flags without changing their Gibbs marginals.

The remaining task is to estimate how
\(Q_\xi(dv\mid u)\) transmits exterior dependence
after a successful core draw, and how that
dependence propagates through later overlapping
updates. It would be circular to declare the
failure flags to be the whole dependency graph.
The next companion audit remains V45, and
the next ten-version direction review remains V50.
\end{decision}

\subsection{Exact checks and append record}

The script
\begin{center}\small\ttfamily
scripts/verify\_ym\_f12\_v42\_common\_component.py
\end{center}
checks a common core weight \(19/20\), the
exact joint split and detailed balance of
the averaged full-block update, and the
product probabilities of four independent
failure flags. It checks (46.5) for four
family sizes, and the conditional-probability
bounds and majority Rayleigh formula for
four sizes in (46.11).
All checks use rational arithmetic; displayed
decimal Rayleigh values are only summaries.
All checks passed. These finite examples do
not certify a Yang--Mills mass gap.

V41 had 641092 bytes and SHA--256
\begin{center}\scriptsize\ttfamily
2DF96596DC72C955FC430B3BF3A94650AC80AB467F86EABB7EB0CF3B08466D8C.
\end{center}
Its complete body is retained apart from
the title version. Only active V41 is replaced
by V42; archives, companions and earlier
scripts are retained.

\begin{firewall}[Position after V42]
The actual buffered core laws have a common
probability component of mass
\(1-O(e^{-\kappa\beta})\), uniformly in compact
exterior data under V41's hypotheses.
The component has an exact complete-block
lift and yields independent failure flags
on separated buffers.
The unrestricted volume-uniform interface
gap, all-distance fibre clustering,
scale-stable effective interactions,
continuum construction and the physical
Yang--Mills mass gap remain unproved.
\end{firewall}

\section{V43: exact boundary memory and the information lost by a core refresh}
\label{sec:f12v43}

V42 left a specific question: after the exterior-independent
core draw, which part of the conditional buffer kernel still
remembers the exterior? This section answers that question
at the level of exact likelihood ratios and relative entropy.
It does not infer decay from locality.

The archive already contains Markov-blanket factorization
(f6 V17), nested boundary-channel reductions (f6 V18),
and a Wilson-link collar likelihood identity (f6 V26).
Those geometric principles are not new claims here.
The additions are their application to the \emph{completed}
free-coordinate interaction, an exact divergence identity
for V42's pinned buffer, and a quantitative information
balance using V42's populated common-component estimate.
V36's TV-to-\(L^2\) argument is not repeated.

\subsection{Use the completed interaction boundary}

Fix retained data \(C\), a finite free update set \(S=I_P\),
and its canonical law \(\pi^\xi_S\). Write the interaction
as the sum of the individual potentials \(\Phi_a\) in
(42.7), including the completion Jacobians. Define
\[
 \begin{split}
 \mathcal C_S&=\{a:\operatorname{supp}\Phi_a\cap S\ne\varnothing,\
                     \operatorname{supp}\Phi_a\cap S^c\ne\varnothing\},\\
 B&=\bigcup_{a\in\mathcal C_S}
                       (\operatorname{supp}\Phi_a\cap S),\\
 J&=\bigcup_{a\in\mathcal C_S}
                       (\operatorname{supp}\Phi_a\cap S^c),\qquad
 D=S\setminus B .
 \end{split}                                                    \tag{47.1}
\]
Thus \(B\) is an \emph{inner} interaction boundary and \(J\)
is the required exterior record. They are not defined using
only the uncompleted Wilson plaquette supports.
All exterior dependence of the normalized block kernel
is through \(\xi_J\). Exterior-only energy terms cancel.

Put \(z=Y_D\), \(b=Y_B\), and use product Haar measure.
Collect wholly internal terms into \(H_0(z,b)\) and
crossing terms into \(H_\partial(b;\xi)\). Then
\[
 \begin{split}
 \pi^\xi_S(dz,db)&=\nu_\xi(db)\,R(dz\mid b),\\
 A(b)&=\int e^{-H_0(z,b)}\,dz,\\
 R(dz\mid b)&=A(b)^{-1}e^{-H_0(z,b)}\,dz,\\
 \nu_\xi(db)&=Z_\xi^{-1}A(b)e^{-H_\partial(b;\xi)}\,db .
 \end{split}                                                    \tag{47.2}
\]
Positive compact densities make every normalizer positive.
The kernel \(R\) is independent of \(\xi\).
These formulas follow by integration of the actual Gibbs
density; no terms are deleted from the sampled law.

For an owner cube \(S\) of side \(L\ge2R_0\), the boundary
owners lie within \(R_0\) cells of its faces. Consequently
\[
 |B|\le q_0\{L^4-(L-2R_0)^4\}\le8q_0R_0L^3,\qquad
 |\mathcal C_S|\le D_0|B|.                                     \tag{47.3}
\]
For the original region convention \(I_P\), include its
fixed-width owner collar; the same \(O((L+1)^3)\) conclusion
holds for cubical parents, with adjusted fixed constants.
No surface estimate is claimed for arbitrary perforated sets.

\subsection{All full-block distinguishability is boundary distinguishability}

For probabilities \(\rho\ll\eta\), use
\[
 D(\rho\Vert\eta)=\int \log(d\rho/d\eta)\,d\rho,\qquad
 \chi^2(\rho\Vert\eta)=\int(d\rho/d\eta-1)^2\,d\eta .
 \tag{47.4}
\]
TV is \(\frac12\int|d\rho-d\eta|\).
For an exterior probability law \(\lambda\), denote the
mixtures by \(\pi^\lambda=\int\pi^\xi_S\,\lambda(d\xi)\)
and \(\nu^\lambda=\int\nu_\xi\,\lambda(d\xi)\).

\begin{proposition}[Boundary sufficiency with no loss]
\label{prop:f12v43-sufficiency}
For two exterior laws \(\lambda,\widetilde\lambda\),
\[
 \frac{d\pi^\lambda}{d\pi^{\widetilde\lambda}}(z,b)
            =\frac{d\nu^\lambda}{d\nu^{\widetilde\lambda}}(b),
 \tag{47.5}
\]
and, for each of \(d=\mathrm{TV},D,\chi^2\),
\[
             d(\pi^\lambda,\pi^{\widetilde\lambda})
                        =d(\nu^\lambda,\nu^{\widetilde\lambda}).
 \tag{47.6}
\]
For \(D,\chi^2\), the order of the arguments has the
meaning in (47.4).
\end{proposition}
\begin{proof}
Integrating (47.2) over the exterior gives
\(\pi^\lambda(dz,db)=\nu^\lambda(db)R(dz\mid b)\).
The same conditional probability occurs under both laws.
Thus the Radon--Nikodym derivative is the boundary
derivative. Insert it into each of the three definitions
and integrate \(R(dz\mid b)\), whose mass is one.
This also proves the extended-value identities.
\end{proof}

In particular the strong data-processing contraction
coefficient of the channel \(\xi\mapsto\pi^\xi_S\),
defined by taking the supremum of the output/input
relative-entropy ratio over finite positive input
relative entropies, is exactly the coefficient of
\(\xi\mapsto\nu_\xi\). The same assertion holds for
the \(\chi^2\) coefficient and the TV row diameter.
This is an identity, not a bound by
\(\varepsilon_\beta\).

There is a quantitative, but noncontractive, surface bound.
For fixed \(\xi,\zeta\), let
\(W(b)=H_\partial(b;\xi)-H_\partial(b;\zeta)\).
By (42.7),
\[
 \operatorname{osc}W\le2v_\beta|\mathcal C_S|=:\Omega_S,\qquad
 \left|\log\frac{d\nu_\xi}{d\nu_\zeta}(b)\right|
                                  \le\Omega_S,\qquad
 D(\pi^\xi_S\Vert\pi^\zeta_S)\le\Omega_S .
 \tag{47.7}
\]
Indeed the log density ratio is
\(-W(b)-\log\int e^{-W}\,d\nu_\zeta\), which lies
between \(-\operatorname{osc}W\) and
\(\operatorname{osc}W\). Each crossing term contributes
at most \(2v_\beta\) to that oscillation.
For cubes, (47.7) costs \(O((\beta+1)L^3)\).
It does not tend to zero for the parents \(L\asymp\beta\).

\subsection{The same exact identity holds after pinning the core}

Let the fixed V42 core \(\mathcal A\) lie in \(D\).
This holds for its sufficiently large buffered parent.
Write \(U=Y_{\mathcal A}\), \(Z=Y_{D\setminus\mathcal A}\).
For fixed \(u\), set
\[
 \begin{split}
 A_u(b)&=\int e^{-H_0(u,z,b)}\,dz,\qquad
 R_u(dz\mid b)=A_u(b)^{-1}e^{-H_0(u,z,b)}\,dz,\\
 \nu_\xi^u(db)&=
 \frac{A_u(b)e^{-H_\partial(b;\xi)}\,db}
      {\int A_u(t)e^{-H_\partial(t;\xi)}\,dt}.
 \end{split}                                                    \tag{47.8}
\]
The canonical buffer conditional in (46.6) is exactly
\[
                 Q_\xi(db,dz\mid u)
                            =\nu_\xi^u(db)\,R_u(dz\mid b).
 \tag{47.9}
\]
In particular, for the \emph{same} \(u\) and any
\(\xi,\zeta\),
\[
 d\bigl(Q_\xi(\cdot\mid u),Q_\zeta(\cdot\mid u)\bigr)
                              =d(\nu_\xi^u,\nu_\zeta^u),
 \qquad d=\mathrm{TV},D,\chi^2 .                                \tag{47.10}
\]
The proof is (47.5)--(47.6) applied to (47.9).
The log likelihood bound (47.7) also holds uniformly
in \(u\); the crossing potential \(W\) is unchanged.
This statement needs no small-curvature assumption.
The V42 common component, in contrast, retains all
of V41's fixed-core and curvature hypotheses.

Equation (47.10) answers the buffer question exactly:
once the two core samples agree, their remaining
statistical discrepancy is the discrepancy of the
\emph{pinned boundary} laws \(\nu_\xi^u\).
One can first couple those boundary laws, then use
the same \(R_u\)-sample whenever the boundaries agree.
For a fixed pair this achieves full-buffer disagreement
probability equal to their boundary TV distance.
No simultaneous maximal coupling over all exteriors
is inferred from this pairwise statement.

\subsection{An exponentially small information contraction for the core}

Let \(K(\xi,du)=\mu_\xi(du)\) be the V42 core channel.
Its simultaneous split is
\[
             K=(1-\varepsilon_\beta)\omega_\beta
                              +\varepsilon_\beta\Sigma,\qquad
                  \varepsilon_\beta\le Ce^{-\kappa\beta},
 \tag{47.11}
\]
where \(\Sigma(\xi,\cdot)=\sigma_\xi\).
For finite input relative entropy,
\[
 D(\lambda K\Vert\widetilde\lambda K)
                    \le\varepsilon_\beta
                                     D(\lambda\Vert\widetilde\lambda).
 \tag{47.12}
\]
The same inequality holds with \(D\) replaced by \(\chi^2\).
This is a classical common-component data-processing
argument, not a new general information-theoretic theorem.
For background see Raginsky,
\href{https://arxiv.org/abs/1411.3575}{arXiv:1411.3575},
and Polyanskiy--Wu,
\href{https://arxiv.org/abs/1508.06025}{arXiv:1508.06025}.
The short proof below supplies the precise version used
here on compact standard Borel spaces.

\begin{proof}
Retain the Bernoulli flag as part of the output.
Under either input law it has the same success probability
\(1-\varepsilon_\beta\). On success the output law is
\(\omega_\beta\); on failure it is the corresponding
input law passed through \(\Sigma\). The relative entropy
of these \emph{flagged} output laws is therefore
\(\varepsilon_\beta
D(\lambda\Sigma\Vert\widetilde\lambda\Sigma)\).
Discarding the flag cannot increase relative entropy.
For completeness, if \(r=d\lambda/d\widetilde\lambda\),
the output density ratio for any common probability
kernel is the conditional expectation of \(r\) given
that output, under the reference joint law.
Conditional Jensen for \(t\log t\) proves that its
relative entropy is at most
\(\int r\log r\,d\widetilde\lambda\).
Apply this first to discarding the flag and then to
\(\Sigma\). Replacing \(t\log t\) by \((t-1)^2\)
proves the \(\chi^2\) assertion with the same factor.
\end{proof}

No inference about the full-block channel follows by
retaining (47.12) and omitting (47.9).
Indeed (47.6) identifies the full-block coefficient
with the boundary coefficient, not the core coefficient.

\begin{theorem}[Where the information not seen in the core remains]
\label{thm:f12v43-information}
Let \(X\) be a finite-valued exterior label. Conditional
on \(X=x\), choose exterior data with law \(\lambda_x\);
then sample the actual complete parent law
\(\pi^\xi_S\). In this joint experiment, mutual information
satisfies
\[
 \begin{split}
 I(X;U)&\le\varepsilon_\beta I(X;\xi)
                             \le\varepsilon_\beta H(X),\\
 I(X;U,Y_{S\setminus\mathcal A})&=I(X;Y_B),\\
 I(X;Y_{S\setminus\mathcal A}\mid U)
                             &=I(X;Y_B)-I(X;U).
 \end{split}                                                    \tag{47.13}
\]
Consequently the last quantity lies in
\[
 \left[\max\{0,I(X;Y_B)-\varepsilon_\beta H(X)\},
                                      \ I(X;Y_B)\right].
 \tag{47.14}
\]
\end{theorem}
\begin{proof}
Put \(\bar\lambda=\sum_x\Pr(X=x)\lambda_x\).
For every positive-probability label,
\(\lambda_x\ll\bar\lambda\) with density at most
\(1/\Pr(X=x)\); hence all input entropies used below
are finite. By the definition of mutual information
and (47.12),
\[
 \begin{split}
 I(X;U)
 &=\sum_x\Pr(X=x)D(\lambda_xK\Vert\bar\lambda K)\\
 &\le\varepsilon_\beta
               \sum_x\Pr(X=x)D(\lambda_x\Vert\bar\lambda)
 =\varepsilon_\beta I(X;\xi).
 \end{split}
\]
The bound by \(H(X)\) follows from the same density
bound, summed over labels. By (47.2), conditional on
\(Y_B\), the law of \(Y_D\) is \(R\) regardless of
the label or the exterior. Applying (47.6) to each
label versus the mixture proves the second line.
The mutual-information chain identity, obtained by
factoring the joint density ratio into its marginal
and conditional ratios and integrating logarithms,
proves the third. All quantities are at most \(H(X)\),
so there is no subtraction of infinite values.
Nonnegativity and the first line give (47.14).
\end{proof}

The bound is about the actual averaged Gibbs update.
It does not condition the Gibbs law on successful flags.
It gives no small upper bound on buffer information
unless the boundary information is separately small.

\subsection{A positive local example and the unthinned dependency}

The distinction is present even in a finite-range,
strictly positive model. Reuse the independent-dimer
mechanism of V38, now to test the information balance.
Let \(X_1,\ldots,X_m\) be independent fair exterior
bits, and let \(B_i\) be their independent noisy copies,
with crossover probability \(0<t<1/2\).
Add independent fair bits for every core and other
interior coordinate. The joint density is a product
of local factors
\[
 \Pr(X_i=x,B_i=b)=
       \tfrac12(1-t){\bf1}_{b=x}
                    +\tfrac12t{\bf1}_{b\ne x}.                  \tag{47.15}
\]
These disjoint pairs can be placed across a cubical
boundary, with \(m\) proportional to its surface area.
Their degree is one, their range is one lattice step,
and their potential oscillation is
\(\log((1-t)/t)\), independent of the cube size.
The deep core channel is exactly constant, so it
admits (47.11) with \(\varepsilon_\beta=0\).
Nevertheless, writing \(h(t)=-t\log t-(1-t)\log(1-t)\),
\[
 I(X;Y_B)=m\{\log2-h(t)\}>0,\qquad
 I(X;Y_{S\setminus\mathcal A}\mid U)=I(X;Y_B).
 \tag{47.16}
\]
This follows by additivity of the logarithmic density
ratios of the independent pairs; one fair input bit
has output entropy \(\log2\) and conditional output
entropy \(h(t)\). The example is not Yang--Mills and
is not a counterexample to a gap: the independent
dimers themselves have a positive volume-uniform gap.
It disproves only an inference of negligible full-block
or buffer memory from an excellent core component.

There is also an exact random-map interpretation of
the obstruction. A block update may be sampled in this order:
\[
 \xi_J\ \longrightarrow\
 \left\{\begin{array}{l}
 U\text{ from the V42 flagged core channel},\\
 B\text{ from }\nu_\xi^U,\\
 Z\text{ from }R_U(\cdot\mid B).
 \end{array}\right.                                             \tag{47.17}
\]
Use independent uniform random variables to realize
these measurable kernels. On a successful core branch,
the incoming exterior arrows to \(U\) can be removed.
The arrows \(\xi_J\to B\) cannot be removed on that
basis; \(Z\) depends on \(B,U\), though not directly
on \(\xi_J\).

For a fixed finite update schedule this gives a
finite directed acyclic computation graph, with
copies for unchanged coordinates between updates.
Couple two initial configurations with the same
random variables. If a queried output has no directed
path from any differing initial coordinate after
successful-core arrows have been removed, its two
values agree. To prove this, order the graph by
sampling time: every node with agreeing parents
and the same random variable has agreeing outputs.
Induction proves the assertion. This is a sufficient
coalescence test, not its converse.
In particular a path through successive boundary
outputs may survive every core success.
The rare flags alone do not form the dependency graph.
The assertion here is for deterministic schedules,
or schedules chosen independently of the field;
state-dependent selection requires its own arrows
and invariance argument.

\begin{decision}[Move to pinned boundary transport, not another core reset]
The new target is the actual channel
\(\xi_J\mapsto\nu_\xi^u\), or a geometry of overlapping
updates that provably forces its surviving boundary
records through controlled interior channels.
Equations (47.9)--(47.14) specify the normalization
and information that must be tracked.
The populated V42 small parameter controls the core
channel; it does not yet control this target.
An attempted proof must either bound the pinned
boundary transport on the required records, or
prove that the proposed schedule has no unthinned
bypass. Simply multiplying the core contraction
factors across updates would omit (47.17)'s boundary
arrows. The next companion audit remains V45;
the next direction review remains V50.
\end{decision}

\subsection{Checks and preservation record}

The script
\begin{center}\small\ttfamily
scripts/verify\_ym\_f12\_v43\_boundary\_memory.py
\end{center}
checks exact rational full-block and pinned-buffer
likelihood ratios, TV and \(\chi^2\) identities, mixture
versions of the boundary identity, and the common-core
\(\chi^2\) contraction on a finite grid of input laws.
It also checks that an exterior-dependent boundary
can retain positive distinguishability while the
core channel is constant. These are algebraic tests
of finite examples, not verification of an analytic
Yang--Mills estimate or of the continuum goal.

V42 had 655050 bytes and SHA--256
\begin{center}\scriptsize\ttfamily
73EE1D38B0CB9263C2B8A2564780BDCDDE72D4888AD8ECE3132CDB88B73D6C7C.
\end{center}
Its body is preserved apart from the title version.
Only active V42 is replaced by V43; all archives,
companion notes and existing scripts are retained.

\begin{firewall}[Position after V43]
The exact buffer memory has been reduced to the
pinned completed-interaction boundary law, with
no divergence lost in the reduction.
V42 gives exponential contraction of core relative
entropy and an exact balance for the information
remaining in the buffer. Boundary transport,
an unrestricted volume-uniform gap, scale-stable
effective interactions, continuum construction,
and the physical Yang--Mills mass gap remain unproved.
\end{firewall}

\section{V44: an all-success obstruction to coarse buffer scheduling}
\label{sec:f12v44}

V43 identified the actual pinned boundary channel
\(\xi_J\mapsto\nu_\xi^u\), but left two possible routes:
control that channel analytically, or choose overlapping
updates that force its memory through successful cores.
This section tests the latter route before attempting
another improvement of the core failure probability.

The result is an obstruction to a \emph{specified
coarse dependency certificate}, not to convergence of
the Gibbs sampler. Even if every core refresh succeeds,
the certificate below cannot forget a positive fraction
of the initial field in dimension at least two.
The proof is geometric and uniform over all finite
schedules. Neither V42's rare-flag estimates nor V43's
information balance previously established this fact.

\subsection{Declare the graph being tested}

Work first on \(\mathbb Z^d\), \(d\ge2\), with nearest-neighbor
connectivity and the supremum distance. At center \(c\), set
\[
 \begin{split}
 A_c&=c+[-r,r]^d,\\
 S_c&=c+[-R,R]^d,\\
 \widehat J_c&=(c+[-R-s,R+s]^d)\setminus S_c,
 \qquad 0\le r<R,\quad s\ge1,
 \end{split}                                                     \tag{48.1}
\]
where the radii are integers and the sets consist of
lattice vertices. \(A_c\) is the successfully refreshed
core, \(S_c\) the full updated parent, and
\(\widehat J_c\) a declared exterior input collar.

On a torus \((\mathbb Z/N\mathbb Z)^d\), use the same
translated coordinate cubes modulo \(N\), and assume
\[
                         N>\max\{4R,\,2(R+s)\}.
 \tag{48.2}
\]
This prevents the local faces used below from folding
onto themselves. Distances on the torus are periodic
supremum distances.

The graph has the following rules. On a successful
branch the core output has no incoming configuration
arrows. Every noncore output of the updated parent
retains arrows from \emph{all} vertices of
\(\widehat J_c\); dependence on the freshly sampled
core and on common auxiliary randomness is also allowed.
Unchanged coordinates retain their previous value.
On a failed branch, the core also retains the collar
arrows. These rules formalize the vector-valued buffer
treatment in (47.17), with a full owner collar substituted
for the exact exterior input set.

The word \emph{declared} matters. For the completed fibre,
a sufficiently thick owner collar contains every
actual interaction input. One may conservatively add
arrows from unused collar coordinates and treat the
buffer as one output record. The resulting graph is
a valid \emph{upper} dependency graph. It need not be
a minimal graph for the actual kernels, and the theorem
below does not assert that all its arrows represent
nonzero physical influence. Proving that arrows can
be removed, or that their influence is small, is
additional mathematical information.

Let \(G_t\) be the coordinates with no graph path to
an initial configuration input after \(t\) updates.
Start with \(G_0=\varnothing\), as appropriate to
testing complete loss of arbitrary initial data.
For a successful update at \(c_t\), graph reachability
gives the exact recursion
\[
 G_{t+1}=
 \begin{cases}
 (G_t\setminus S_{c_t})\cup S_{c_t},
                      &\widehat J_{c_t}\subseteq G_t,\\
 (G_t\setminus S_{c_t})\cup A_{c_t},
                      &\widehat J_{c_t}\not\subseteq G_t.
 \end{cases}                                                     \tag{48.3}
\]
Indeed, when a collar input still has an initial
ancestor, the declared arrows carry that ancestor
to every noncore output. If all collar inputs have
lost their ancestors, the common random inputs and
the successful core generate a completely known
parent. Outside \(S_{c_t}\) nothing changes.

Graph-known coordinates have identical values in a
coupling using the same random inputs. Coordinates
outside \(G_t\) are merely \emph{uncertified}: they
need not actually disagree. Equation (48.3) is not
a recursion for the true disagreement set.

\subsection{Small known components never prepare a whole collar}

\begin{theorem}[All-success scheduling obstruction]
\label{thm:f12v44-obstruction}
For the graph (48.1)--(48.3), in dimension \(d\ge2\),
every finite sequence of successful updates has:
\[
 \begin{gathered}
 \text{each nearest-neighbor component of }G_t
                         \text{ has diameter at most }2r,\\
                    \widehat J_c\not\subseteq G_t
                         \quad\text{for every center }c .
 \end{gathered}                                                   \tag{48.4}
\]
The assertions hold on \(\mathbb Z^d\), and on tori
satisfying (48.2). In particular the first branch
of (48.3) never occurs, regardless of the schedule.
\end{theorem}
\begin{proof}
Assume every component of a set \(G\) has diameter
at most \(2r\). The collar at any center \(c\)
contains the connected face
\[
 F_c=\{c+(R+1,x_2,\ldots,x_d):
                              -R\le x_j\le R\}.
 \tag{48.5}
\]
It contains two vertices differing by \(2R\) in
the second coordinate. Its diameter is therefore
\(2R>2r\); on the torus the same distance holds
because \(N>4R\). If \(\widehat J_c\subseteq G\),
then \(F_c\subseteq G\). Its connectedness would
put it in one component of \(G\), a contradiction.
This proves the second assertion from the first.
It is enough that the input set contain this face;
connectedness of the entire collar is not needed.

Consequently the next successful update gives
\(G'=(G\setminus S_c)\cup A_c\).
Every component of \(G\setminus S_c\) is a subset
of an old component and still has diameter at most
\(2r\). Moreover
\[
            \operatorname{dist}_\infty(A_c,S_c^c)
                                  \ge R-r+1\ge2 .
 \tag{48.6}
\]
Thus there are no nearest-neighbor edges between
the new core and the surviving old components.
The core is connected and has diameter at most
\(2r\). It cannot merge two old components or
enlarge one. Every new component obeys the same
diameter bound. Starting from \(G_0=\varnothing\)
proves the first assertion inductively, and the
face argument supplies the second at every step.
\end{proof}

The argument also applies to cores contained in
\(A_c\), possibly disconnected. Their components
have no larger diameter and remain separated
from the surviving old set. Enlarging a declared
input collar only makes its preparation harder.
For a finer graph with actual input sets \(J_c\),
the same obstruction applies \emph{if} each \(J_c\)
contains a connected set of diameter greater than
\(2r\) in the graph used for the component argument.
That condition is explicit and is not asserted
without a support audit for every completed-fibre
coordinate type.

\subsection{A positive density of unresolved ancestors persists}

\begin{corollary}[Uniform floor for the coarse certificate]
\label{cor:f12v44-density}
On a torus satisfying (48.2), at every finite time
and for every all-success schedule,
\[
 |G_t^c|\ \ge\
       N^{d-1}\left\lceil\frac{N}{2r+2}\right\rceil
                  \ \ge\ \frac{N^d}{2r+2}.
 \tag{48.7}
\]
\end{corollary}
\begin{proof}
Fix the coordinates \(2,\ldots,d\) and consider
the remaining nearest-neighbor cycle of length \(N\).
It cannot contain \(2r+2\) consecutive known vertices:
those vertices would lie in one component with two
points at periodic distance \(2r+1\).
Condition (48.2) and \(R\ge r+1\) ensure that this
distance is less than \(N/2\), contradicting (48.4).
The whole cycle cannot be known either.

If there are \(k\) unknown vertices on the cycle,
the known run after each unknown vertex has at
most \(2r+1\) vertices. These \(k\) runs and their
initial unknown vertices partition the cycle,
so \(N\le k(2r+2)\).
Sum \(k\ge\lceil N/(2r+2)\rceil\) over the
\(N^{d-1}\) disjoint cycles.
\end{proof}

This is more than failure of a particular sweep:
for any random schedule, the spatially averaged
probability that an output still has an initial
ancestor in this graph is at least \(1/(2r+2)\).
If the schedule law is translation invariant,
the same lower bound holds at each fixed vertex.
The assertion follows by averaging (48.7);
no independence of different vertices is used.
The centers may be random independently of the field.
If scheduling itself depends on unknown field data,
its additional dependencies must also be tracked;
ignoring them cannot supply a valid coalescence proof.

In particular the usual sufficient estimate of
coupling error by the probability of a surviving
graph ancestor cannot tend to zero when computed
using this graph. The \emph{upper-bound method}
has a positive floor. This is not a positive lower
bound on the actual coupling error or on distance
to Gibbs equilibrium.

\subsection{Failures cannot repair the certificate}

Write \(T_c^+\) for the update map in (48.3).
It is increasing for inclusion of input sets:
if \(G\subseteq H\), then
\(T_c^+(G)\subseteq T_c^+(H)\).
To verify this, if \(\widehat J_c\subseteq G\)
then it is also contained in \(H\); if it is
contained in neither, the same core is adjoined;
in the remaining case the larger input produces
all of \(S_c\), which contains the smaller input's
core. Outside \(S_c\), inclusion is preserved.

On a failed branch the declared update is
\[
 T_c^-(G)=
 \begin{cases}
 (G\setminus S_c)\cup S_c,&\widehat J_c\subseteq G,\\
 G\setminus S_c,&\widehat J_c\not\subseteq G .
 \end{cases}
 \tag{48.8}
\]
It satisfies \(T_c^-(G)\subseteq T_c^+(G)\).
For any realized sequence of centers and flags,
induction using monotonicity proves
\[
                         G_t^{\rm flags}\subseteq G_t^{+},
 \tag{48.9}
\]
where the right side uses the same centers with
every flag replaced by success. Thus (48.7) also
bounds the number of unresolved coordinates with
failures. There is no threshold in
\(\varepsilon_\beta\) that can overcome this
particular obstruction: it already holds at zero.

Separated parallel parent updates with noninteracting
input/output records can be serialized without changing
this graph calculation. The obstruction therefore
also covers such parallel rounds. It does not
cover a larger joint update whose exact common
component refreshes the whole union.

\subsection{Why the geometric hypotheses are essential}

In one dimension the exterior face (48.5) is a
single point, so the proof does not apply.
There is a concrete exception with \(r=0\), \(R=s=1\).
On the integer line, use successive centers
\(-2,2,0\). The first two successful updates leave
the two known cores \(\{-2,2\}\).
The third parent is \([-1,1]\), with exterior
collar \(\{-2,2\}\) already known, so the first
branch of (48.3) occurs and the complete interval
is certified. The same example works on a
sufficiently long periodic cycle.

Likewise, if the core were the whole parent
(\(r=R\)), a successful update would certify all
of \(S_c\) without requiring the collar.
A finite cover would then certify a finite
torus. That is a materially stronger common
component than the buffered-core result of V42.
It cannot be obtained by relabelling the buffer.

These examples locate the obstruction precisely:
a small reset region separated from surviving old
outputs, together with a large connected exterior
record whose every coordinate is retained as a
possible ancestor of the whole buffer.

\subsection{Consequence for the actual Yang--Mills route}

For the actual completed interaction, \(s\) can
be taken at least its fixed owner range.
A fixed finite core template is compatible with
V41--V42, while their parent radius grows like
\(\beta\). A conservative owner-cube implementation
then lies exactly in the scale regime \(r<R\).
At the owner-record level, regard a cell as known
only when its complete selected free-variable
record is known. Giving the entire small core
record a perfect common refresh is already
optimistic relative to the proved nonzero
failure probability.

The theorem rules out completing a mixing proof
by merely optimizing the schedule of these
\emph{coarsened} buffer nodes or reducing their
core failure probability further. It does not
rule out the same physical block updates with
a more informative coupling or a different
variance argument. In particular V43's exact
boundary factorization is useful only if some
spatial information inside the boundary record
is retained; treating all of it as one dependent
output gives the obstruction just proved.

The next analytic object remains
\[
 \nu_\xi^u(db)\ \propto\
           A_u(b)e^{-H_\partial(b;\xi)}\,db
 \tag{48.10}
\]
from (47.8). One must control its spatially
resolved dependence, or prove a localized
Dirichlet-form inequality that pays for it.
Finite interaction range of \(H_\partial\)
alone is not such a result: after normalizing
and integrating interior variables, the
boundary sampling map can have additional
dependencies. The statement that those
dependencies decay must itself be proved.

\begin{decision}[Stop the schedule-only branch of the argument]
V44 supplies evidence to move upstream rather than
repeat the successful-core construction. Preserve
V42's populated common component and V43's exact
boundary law, but do not attempt to close the
coarse all-collar ancestor graph with another
choice of centers or an improved success rate.
The next version should audit the companion
programmes as scheduled at V45, then pursue
spatially resolved pinned-boundary transport or
the corresponding local energy inequality.
The ten-version direction review remains V50.
\end{decision}

\subsection{Finite checks and append record}

The script
\begin{center}\small\ttfamily
scripts/verify\_ym\_f12\_v44\_reset\_geometry.py
\end{center}
exhausts 6275 transitions through three updates
on a \(5\times5\) torus, reaching 825 distinct
sets at depth three. It verifies the component
diameter invariant and (48.7), checks 120-update
schedules in dimensions two, three and four,
and tests (48.9) with failed branches.
It also verifies the one-dimensional exception.
All operations on the dependency sets are exact;
random schedules use a fixed seed and are tests,
not a substitute for the all-schedule proof.
All checks passed. The script does not sample
Yang--Mills fields and proves no physical
disagreement or absence of a mass gap.

V43 had 671875 bytes and SHA--256
\begin{center}\scriptsize\ttfamily
A5ACF6DD04F5742E80445150A456E9C22B5ED149769B11356D53E367D0B49240.
\end{center}
Its complete body is preserved apart from the
title version. Only active V43 is replaced by
V44; the archives, companions and prior scripts
remain unchanged.

\begin{firewall}[Position after V44]
A schedule-only closure using the declared
whole-buffer dependency graph is obstructed,
even with perfect core resets. This is a
rigorous limitation of that certificate,
not a statement against Yang--Mills mixing.
Spatially resolved boundary control, an
unrestricted volume-uniform gap, scale-stable
effective interactions, the continuum construction,
and the physical mass gap remain unproved.
\end{firewall}

\section{V45: fixed-core pinning and simultaneous spatial buffer transport}
\label{sec:f12v45}

V44 rules out a schedule-only proof using one opaque
buffer record. We now return to the actual conditional
integral and retain its coordinatewise response.
The new step is to prove that pinning a fixed core
near the section preserves the compact-to-box estimate,
then obtain a single coupling with a spatial cost
bound on every output coordinate.
The result has a stronger pointwise exterior guard;
that guard is not removed in this section.

\subsection{Scheduled companion review}

The latest active companion files were inspected again:
\[
 \begin{array}{c|r|l}
 \text{programme and version}&\text{bytes}&\text{SHA--256 prefix}\\ \hline
 \text{SM V72}&607372&\texttt{F44F9745D43379E1}\\
 \text{NS V26}&278283&\texttt{82C887F8C2C69E4F}\\
 \text{parallel YM V5}&54174&\texttt{306F1BB84CFA8A8D}
 \end{array}                                                     \tag{49.1}
\]
Their full hashes agree with the V40 review.
SM V72 proves a fixed-mode, fixed-time many-copy
limit; its ultraviolet and physical finite-species
questions remain open. NS V26 gives the pointwise
rank-one derivative norms, not full global regularity.
Parallel YM V5 excludes the proposed uniform tail
for an extensive total integer charge under its
tensorization hypotheses; compact labels and local
topology remain distinct from that target.
None supplies unrestricted pinned-boundary contraction.
The document's terminal wording for its own sector
route is not an instruction to terminate this goal.
No new companion theorem is imported here.
The next companion review is V50.

\subsection{Setup and the precise additional guard}

Let \(R\) be a finite cell region, \(s_R=|R|\), and
\(\pi_{R,C}^{y}\) its actual compact canonical law.
Choose a nonempty fixed finite free core
\(\mathcal A\subset I_R\); its template and cardinality
do not grow with \(R\) or \(\beta\).
Write \(U=Y_{\mathcal A}\), \(Z=Y_{I_R\setminus\mathcal A}\),
and \(n=(\dim G)|\mathcal A|\).
Use Lebesgue coordinates for \(U\) in its fixed
section chart, and product Haar measure for \(Z\).

Keep V32's box radius \(R_*\), its ball radius
\(\tau_0\), and \(c_0=\tau_0^2/(16A)>0\).
The fixed-core local derivative argument of (45.5)
gives a constant \(L\), independent of \(R,y,\beta\),
such that for \(\beta\ge1\),
\[
 |\nabla_u H_R(u,z;y)|\le L\beta
 \tag{49.2}
\]
on a fixed core chart, with the core Haar factor
included in \(H_R\). The core marginal density
\(p_y(u)\) therefore satisfies
\(|\nabla\log p_y|\le L\beta\).
Only terms incident to \(\mathcal A\) are differentiated;
the number of other integrated variables does not
enter this bound.

Choose \(a>0\) sufficiently small that
\[
 K_a=[-a,a]^n\text{ has a fixed positive chart margin},
 \qquad 2L\sqrt n\,a\le c_0/4.
 \tag{49.3}
\]
Choose \(\tau_a>0\) so that section-distance less
than \(\tau_a\) in each core group variable implies
\(U\in K_a\).
Impose V32's pointwise-boundary and curvature
hypotheses both for \(\tau_0\) and for \(\tau_a\),
as well as its boxed Hessian hypotheses.
Thus the required physical boundary radius is at
most \(\min\{b_{\tau_0},b_{\tau_a}\}\), and the
curvature threshold is decreased correspondingly.
These are fixed positive thresholds, possibly
smaller than those used in V42.

Below, \(\mathcal G_R^{\rm pin}\) denotes a fixed
small product cube of exterior free logarithms
which ensures these physical collar bounds.
It is chosen using the fixed local completion
derivatives. Coordinatewise straight interpolation
of two points in this cube stays in the guard;
its group-coordinate path lengths are bounded by
a fixed constant times their Hilbert--Schmidt
endpoint distances. This is the same guard
construction used in V32--V33, with smaller radii.

\subsection{A density proof of stability under point pinning}

Let \(\mathcal D_Z\) be the product of the V32
coordinate cubes for the unpinned variables.
The exact pinned probability is defined by the
positive conditional density
\[
 \pi_R^{y,u}(dz)=
             \frac{e^{-H_R(u,z;y)}\,dz}
                         {\int e^{-H_R(u,t;y)}\,dt}.
 \tag{49.4}
\]
This canonical version is defined for every
\(u\in K_a\), not merely almost every \(u\).

\begin{theorem}[Fixed-core pinning preserves exponential defects]
\label{thm:f12v45-pinning}
Under the guard just specified and for sufficiently
large \(\beta\), there is a fixed \(C_{\rm pin}\) such that
\[
 \sup_{y\in\mathcal G_R^{\rm pin}}\sup_{u\in K_a}
          \pi_R^{y,u}(\mathcal D_Z^c)
 \le p_{\rm pin}(R,\beta):=
 \min\{1,C_{\rm pin}\beta^n s_R e^{-3c_0\beta/4}\}.
 \tag{49.5}
\]
Constants can depend on the fixed core template,
but not on the region volume.
\end{theorem}
\begin{proof}
By V32's joint bound at threshold \(\tau_a\) and
a union bound over the fixed core,
\[
 \Pr_{\pi_{R,C}^y}(U\notin K_a)
      \le|\mathcal A|e^{-\beta\tau_a^2/(16A)}\le\tfrac12
 \tag{49.6}
\]
after increasing the coupling threshold.
Consequently there is \(u_0\in K_a\) with
\(p_y(u_0)\ge[2\operatorname{vol}(K_a)]^{-1}\).
The log-gradient bound and the straight segment
inside \(K_a\) give
\[
 \inf_{u\in K_a}p_y(u)
 \ge [2\operatorname{vol}(K_a)]^{-1}
                         e^{-2L\sqrt n\,a\beta}
 \ge c_a e^{-c_0\beta/4}.
 \tag{49.7}
\]

For any event \(F\) depending only on \(Z\), let
\(p_{y,F}(u)\) be the subprobability density of
\((U,F)\), using the \emph{unconditioned} global
normalizer of \(\pi_{R,C}^{y}\).
Its integral over the core chart is at most
\(\pi_{R,C}^{y}(F)\).
Differentiation under its fixed integration domain
and (49.2) give
\[
                     |\nabla p_{y,F}|\le L\beta p_{y,F}.
 \tag{49.8}
\]
If this density is zero, the desired bound is
immediate. Otherwise integrate it on a ball of
radius \(c/\beta\) around \(u\in K_a\), contained
in the fixed chart for large \(\beta\).
On this ball it is at least
\(e^{-Lc}p_{y,F}(u)\). Its mass is bounded by
\(\pi_{R,C}^{y}(F)\), so
\[
                p_{y,F}(u)\le C\beta^n\pi_{R,C}^{y}(F).
 \tag{49.9}
\]
This does not condition on the zero-probability
event \(U=u\) by dividing by its probability.
It compares the two actual densities.

Take \(F=\mathcal D_Z^c\). Since
\(F\subseteq\mathcal D_R^c\), (36.7) gives
\(\pi_{R,C}^{y}(F)\le73s_R e^{-c_0\beta}\).
By (49.4),
\(\pi_R^{y,u}(F)=p_{y,F}(u)/p_y(u)\).
Combine (49.7) and (49.9), then cap the bound
by one. All estimates are uniform on the
declared guard.
\end{proof}

This argument proves a pinning extension rather
than assuming V32 for a region with arbitrary
interior holes. It does not allow an increasing
number of pinned variables: the factor
\(\beta^n\), the radius choice and constants
all retain their fixed-core dependence.

\subsection{The response operator survives the pin}

On \(\mathcal D_Z\), the measured potential in
(49.4) is obtained by restricting V32's potential
to \(U=u\). Its Hessian in the remaining variables
is a principal Hessian. Hence
\[
        m_0\beta I\le \nabla_Z^2H_R(u,\cdot;y)
                             \le K\beta I
 \tag{49.10}
\]
with the same bounded row norm and fixed interaction
range, uniformly in \(u\in K_a\).
The fixed rectangular faces in the remaining
variables are unchanged. V33's componentwise
response-operator argument (37.3)--(37.4)
therefore applies to this restricted potential.
No new interaction edge is introduced by taking
a principal Hessian.

In particular, for bounded functions \(F,G\) that
are \(C^1\) in the box, their covariance under
the \emph{full compact pinned law} obeys
\[
 \begin{split}
 |\operatorname{Cov}_{\pi_R^{y,u}}(F,G)|
 \le{}&\frac K\beta\sum_{i,l}
       e^{-\gamma\operatorname{dist}(i,l)}
       \|\partial_iF\|_{L^2({\rm box})}
       \|\partial_lG\|_{L^2({\rm box})}\\
 &+6\|F\|_\infty\|G\|_\infty p_{\rm pin}(R,\beta).
 \end{split}                                                     \tag{49.11}
\]
Here \(i,l\) initially enumerate scalar chart
components of unpinned variables, and the
distance is that of their fixed local positions.
To prove (49.11), apply the boxed representation
to the exact restriction of (49.4), then compare
its means of \(F,G,FG\) with the compact means
using (49.5), exactly as in (36.9).
This argument does not identify a boxed measure
with the full compact conditional.

For a directional change of an exterior free
coordinate \(y_j\), the pinned score
\(S_j=\partial_{y_j}H_R(u,Z;y)\) has
\[
 \|S_j\|_\infty\le K(\beta+1),\qquad
 \|\partial_lS_j\|_\infty\le K(\beta+1),
 \tag{49.12}
\]
and its \(Z\)-derivatives are supported on a
fixed number of variables near \(j\).
Terms depending only on \(u,y\) may be subtracted
as constants in \(Z\).
Differentiating the exact normalizer in (49.4) yields
\[
 \partial_{y_j}\pi_R^{y,u}F
                  =-\operatorname{Cov}_{\pi_R^{y,u}}(F,S_j).
 \tag{49.13}
\]
There is no omitted derivative of the core
marginal density: it is already accounted for
by this pinned conditional normalizer.

\subsection{Weighted estimates imply one simultaneous coupling}

Let \(Z_R=I_R\setminus\mathcal A\) denote the
unpinned coordinate index set, and write
\(\rho(i,j)\) for the supremum distance between
their owner cells and those of exterior variables.
Fixed completion ranges and the fixed group
dimension are absorbed in constants below.
For nonnegative weights \(w_i\), use
\[
 d_w(z,z')=\sum_{i\in Z_R}w_i\|z_i-z_i'\|_{\rm HS}.
 \tag{49.14}
\]
For positive weights the compact transport duality
was proved in (37.2); zero weights follow by
adding a common positive number and taking its
limit to zero.

Take \(F\) one-Lipschitz for \(d_w\), and subtract
its value at a fixed configuration. Then
\(\|F\|_\infty\le d_G\sum_iw_i\) and each boxed
gradient component at variable \(i\) has size
at most \(K w_i\). Smooth approximation by
product-group convolution preserves these
coordinate Lipschitz bounds, so nonsmooth
tests follow by uniform convergence.
Equations (49.11)--(49.13) give
\[
 |\partial_{y_j}\pi_R^{y,u}F|
 \le \sum_{i\in Z_R}w_i M_{ij},\qquad
 M_{ij}=K_1e^{-\gamma_1\rho(i,j)}
                 +K_2(\beta+1)p_{\rm pin}(R,\beta).
 \tag{49.15}
\]
The constants are independent of the weights.
Integrate along the guarded coordinate paths
between \(y,\widetilde y\), enlarging \(K_1,K_2\)
by their fixed path-length factor, and use duality:
\[
 \mathsf W_{d_w}(\pi_R^{y,u},\pi_R^{\widetilde y,u})
       \le\sum_i w_i b_i,\qquad
 b_i=\sum_{j\in J_R^\partial}
                  M_{ij}\|y_j-\widetilde y_j\|_{\rm HS}.
 \tag{49.16}
\]

\begin{proposition}[Simultaneous coordinate transport]
\label{prop:f12v45-coupling}
For each fixed \(u\in K_a\) and guarded boundary
pair \(y,\widetilde y\), there is a single coupling
\(\Lambda^{y,\widetilde y,u}\) of the two full
compact pinned laws such that, for every \(i\in Z_R\),
\[
       \int\|z_i-z_i'\|_{\rm HS}\,d\Lambda^{y,\widetilde y,u}
                                                \le b_i .
 \tag{49.17}
\]
Thus (49.15) is an attainable coordinatewise cost
matrix, not merely a list of marginal distances.
\end{proposition}
\begin{proof}
Let \(\mathcal C\) be the set of all vectors of
expected coordinate costs under couplings of
the two fixed laws. Compactness of the group
product, continuity of the costs, and compactness
and convexity of the coupling set imply that
\(\mathcal C\) is compact and convex in
\(\mathbb R^{|Z_R|}\).
The set \(\mathcal C+\mathbb R_+^{|Z_R|}\) is
closed and convex.

If \(b\) did not belong to this set, separation
of a point from a closed convex set would give
a vector \(w\) with
\[
               w\cdot b<\inf_{c\in\mathcal C}w\cdot c.
 \tag{49.18}
\]
Its coordinates must be nonnegative: otherwise
the infimum over the added positive orthant
would be minus infinity. The right side is
exactly \(\mathsf W_{d_w}\), contradicting
(49.16). Therefore \(b=c+v\) for some
\(c\in\mathcal C\), \(v\ge0\). A coupling
realizing \(c\) proves (49.17).
\end{proof}

The convex separation argument is a general
transport principle; the new analytic input is
its populated matrix (49.15) for the actual
pinned compact law. The proposition is for
each boundary pair. It is not a common random
map simultaneously attaining these bounds
over every exterior configuration.

Separate one-coordinate marginal estimates
would not justify this conclusion. For example,
the uniform laws on \(\{00,11\}\) and on
\(\{01,10\}\) have identical one-coordinate
marginals. Both marginal transport distances
are zero, but every pair of configurations
from the two supports has total Hamming
distance one. No coupling has both costs zero.
For weights \((w_1,w_2)\), the joint optimal
weighted cost is \(\min\{w_1,w_2\}\);
the budget \((1/2,1/2)\) is feasible while
\((0,0)\) fails the weight \((1,1)\) test.
This is why all weighted inequalities in
(49.16), not separate marginal ones, are used.

\subsection{The spatial tail and the surviving exterior condition}

There are at most a fixed number of free
variables per owner cell. Summing (49.15)
over outputs at owner distance at least \(\ell\)
from one exterior coordinate gives
\[
 \sum_{\substack{i\in Z_R\\\rho(i,j)\ge\ell}} M_{ij}
 \le K_3(\ell+1)^4e^{-\gamma_1\ell/2}
       +K_2|Z_R|(\beta+1)p_{\rm pin}(R,\beta).
 \tag{49.19}
\]
Indeed the number of owners in a radius-\(k\)
ball in four dimensions is \(O((k+1)^4)\);
summing the exponentially weighted shells
and allowing the displayed exponent loss
proves the first term. The second sums the
constant defect remainder once per output.

If \(s_R=O(\beta^4)\), then \(|Z_R|\le73s_R\)
and this remainder is at most
\[
                 C\beta^{n+9}e^{-3c_0\beta/4}.
 \tag{49.20}
\]
For fixed core dimension it is exponentially
small at large \(\beta\), after decreasing
the exponent. At \(\ell\asymp\beta\) the
first term is also exponentially small.
Thus a change at one exterior coordinate
admits a coupling whose \emph{total}
displacement far from that coordinate is small,
not merely one with a common core.
The entire pinned buffer is sampled from its
actual compact law.

For arbitrary compact exterior pairs a coarse
bound is still available:
\[
                         M_{ij}^{\rm crude}=K(\beta+1).
 \tag{49.21}
\]
To see this, use the global compact score bound,
\(|\operatorname{Cov}(F,S_j)|\le
2\|F\|_\infty\|S_j\|_\infty\), in (49.13).
Integrate along compact group paths with
length bounded by a fixed multiple of the
Hilbert--Schmidt endpoint distance, as in
V33's arbitrary-boundary estimate, then apply
the same weighted separation proof.
This uses the connected compact group and
the fixed global completion, not boxed
convexity. It gives no spatial decay.

Equations (49.15) and (49.21) therefore retain
the distinction between guarded pairs and
arbitrary pairs. V42's successful branch
does not ensure that the surrounding
exterior belongs to
\(\mathcal G_R^{\rm pin}\).
Nor can a stationary rarity estimate be
multiplied into a discrepancy weighted by
the same bad boundary without a proof.
That is the unresolved issue for using
(49.17) in a global contraction or variance
argument.

If this estimate is combined with V42, its
common reference must be chosen in the
possibly smaller cube \(K_a\) of (49.3).
V41--V42 permit that fixed-radius choice
with correspondingly smaller curvature
thresholds and larger coupling thresholds.
The support of an already chosen larger
reference law must not simply be presumed
to satisfy the new pinning hypothesis.

\begin{decision}[Use the spatial matrix; keep the bad-boundary term]
We now have a rigorous spatially resolved
whole-buffer coupling under a precise guard,
after a fixed core has been pinned.
This supplies information absent from V44's
opaque graph and is a positive analytic
extension of V33's unpinned transport estimate.
The next task is to control the weighted
contribution of boundary pairs outside this
stronger guard, or refresh those records
through a local energy inequality that
uses (49.17) without factorizing correlated
rarity and discrepancy. The next companion
review and direction review are both V50.
\end{decision}

\subsection{Checks and append record}

The script
\begin{center}\small\ttfamily
scripts/verify\_ym\_f12\_v45\_pinned\_transport.py
\end{center}
checks 123 exact core-density conditional
ratios in a positive polynomial example,
25 weighted transport budgets for the
two-coordinate example above, and a pinned
Gaussian principal precision matrix with
its exact inverse and spatial entry bounds.
All calculations use rational arithmetic
and passed. These tests check normalization,
pinning algebra and coupling compatibility;
they do not certify the analytic Yang--Mills
constants or a mass gap.

V44 had 686806 bytes and SHA--256
\begin{center}\scriptsize\ttfamily
D5D19291A0D2D7BAAE57315863DA2ABA72695A351CDC805BD231BF7A73EC0D20.
\end{center}
Its full body is retained apart from the
title version. Only active V44 is replaced
by V45. All archives, companions and
earlier verification scripts are retained.

\begin{firewall}[Position after V45]
Fixed-core pinning preserves exponential
compact-to-box defects on a suitably small
pointwise exterior guard. A single full
pinned-buffer coupling attains the derived
spatial cost matrix, with its defect
remainder displayed. Arbitrary bad exterior
pairs, an unrestricted volume-uniform gap,
scale-stable effective interactions,
continuum construction and the physical
Yang--Mills mass gap remain unproved.
\end{firewall}

\section{V46: local boundary defects, excision and spatial repair transport}
\label{sec:f12v46}

V45 requires a good exterior for its pinned spatial
coupling. Before trying to average that condition,
we prove a localized repair estimate for the
\emph{unpinned} actual parent law. Bad exterior
values are permitted near an arbitrary finite
set of defect locations. The comparison cost
is charged near those locations, not to every
output merely because one boundary value is bad.
Pinning again under these bad exterior conditions
is a separate question, not an implicit hypothesis
of the present proof.

The inputs already proved are V32's truncated
collar inequality, its good-boundary sparse bound,
V33's componentwise boxed response, and V45's
weighted-to-simultaneous coupling argument.
The additions are a localized hypothesis for
the sparse bound and its use to create a
random good internal boundary after deterministic
defect excision.

\subsection{Locality constants and boundary defect records}

Use the actual physical sets
\(\Lambda_R,\mathcal F_R,I_R,J_R^\partial\).
Assign a coarse anchor to each physical link
and each free coordinate. A boundary free
coordinate may instead use an endpoint in
the parent when required by V32's ownership
convention; the two choices differ by at
most one cell.

Fix an integer \(r_{\rm loc}\ge4\) large enough
for all completed plaquette, rectangle, pivot
and free-input footprints. In particular,
every such footprint meeting a cell region
lies in its radius-\(r_{\rm loc}\) neighborhood.
An input outside \(I_T\) which affects its
canonical density has an owner within
\(r_{\rm loc}\) of \(T\).
Include anchor changes of at most one cell
in this constant.
Existence follows from the finite, fixed
stencil list of V18 and V25; it is independent
of \(\beta\), region size and boundary values.
It is not inferred from the four links of
an uncompleted plaquette alone.

Choose a free-variable tolerance \(\tau_g>0\)
and a parent frozen-link tolerance \(b_{\rm in}>0\)
as follows. If inherited frozen inputs have
mismatch at most \(b_{\rm in}\), and every
new free input has section-distance at most
\(\tau_g\), then the complete exterior of a
subregion lies in V33's good free-coordinate
guard and obeys V32's physical collar bound
at threshold \(\tau_0\).
Such fixed positive choices exist by uniform
smoothness of the finite completion maps and
their exact section values. Decrease
\(b_{\rm in}\) further so that
\(b_{\rm in}\le b_{\tau_g}\) from (36.4).
Impose the minimum of the curvature thresholds
for these applications, and take \(\beta\)
above their fixed thresholds. Put
\[
                 c_g=\tau_g^2/(16A)>0 .
 \tag{50.1}
\]

For two exterior data \(y,\widetilde y\) of
the same parent \(R\), at fixed retained \(C\),
let \(\mathfrak D\) be a finite set of coarse
anchors containing:

\begin{itemize}
\item every frozen physical link of \(\mathcal F_R\)
whose mismatch exceeds \(b_{\rm in}\) in either exterior;
\item every frozen physical link of \(\mathcal F_R\)
whose values in the two exteriors differ; and
\item every differing exterior free input in
\(J_R^\partial\).
\end{itemize}

Thus, away from \(\mathfrak D\), the inherited
boundary is both identical and good.
This set includes bad boundary values common
to the two exteriors; it is not merely their
disagreement set. All its records are finite
for a finite parent. Enlarging it is harmless.
Write \(d=|\mathfrak D|\) and \(s_R=|R|\).

\subsection{Sparse suppression only needs a locally good frozen collar}

Let \(J_g,M_g\) be the fixed constants of
(36.3)--(36.4) for \(\tau_g\). Choose
\[
                    h_g=4J_g+r_{\rm loc}+2 .
 \tag{50.2}
\]
All distances in this section are coarse
supremum distances; graph neighborhoods
of radius \(4J_g\) are contained in these
supremum neighborhoods.

\begin{proposition}[Localized pointwise-boundary sparse bound]
\label{prop:f12v46-local-sparse}
Let \(\mathcal B\subset I_R\) be nonempty, with
parent-compatible owners at distance greater
than \(h_g\) from every bad frozen-link anchor.
No pointwise hypothesis is required at the
remaining, distant frozen links. Then
\[
 \pi_{R,C}^y
   \{\|Y_i-J_i(C)\|_{\rm HS}\ge\tau_g
                                  \text{ for all }i\in\mathcal B\}
                    \le e^{-c_g\beta|\mathcal B|}.
 \tag{50.3}
\]
\end{proposition}
\begin{proof}
Let \(S\) be the set of owners of \(\mathcal B\),
and use the same truncated regions
\(R_j=R\cap N_{4j}(S)\), \(j\le J_g\), as in V32.
Every parent-frozen physical link occurring in
\(\mathcal F_{R_j}\) has an anchor within
\(4J_g+r_{\rm loc}+1\) of \(S\).
By the assumed separation, none is a bad
parent-frozen link. Therefore the only
pointwise estimate used in (36.2) still holds:
\[
             E_j\le D_{j+1}-D_j+
                             C_F b_{\tau_g}^2|R_j|.
 \tag{50.4}
\]
Links elsewhere in \(\mathcal F_R\) do not
occur in this inequality.
The bounds \(|R_j|\le M_g|\mathcal B|\),
\(D_0\ge\tau_g^2|\mathcal B|\), and the
physical capacity bound are unchanged.
The affordable-shell argument (36.6), followed
by its deterministic dyadic-band union bound,
therefore gives exactly (36.5) at threshold
\(\tau_g\). This is (50.3).
No conditional law on an adaptively selected
good boundary has been substituted.
\end{proof}

This proposition is stronger than asserting
the same estimate only when the entire
parent frozen collar is good. Its constant
does not grow with the number or size of
the distant bad frozen links.

\subsection{A deterministic excision creates a probabilistically good interface}

Fix
\[
 L=2h_g+4r_{\rm loc}+4,\qquad
 T=\{x\in R:\operatorname{dist}(x,\mathfrak D)>L\},
 \qquad
 J=J_T^\partial\cap(I_R\setminus I_T),\quad m=|J|.
 \tag{50.5}
\]
\(J\) is the set of new free exterior inputs
for \(T\) that are still sampled in the
original parent. The exact conditional law
on \(I_T\), after fixing \(I_R\setminus I_T\),
is its original canonical law by disintegration.
We do not delete interactions across the cut.

The completed locality and (36.1) give:
\[
 \begin{gathered}
 \mathcal F_T\setminus\Lambda_R\subset\mathcal F_R,
 \quad\text{these inherited links are identical
 and good in both exteriors},\\
 \operatorname{dist}(j,\mathfrak D)>h_g
                   \quad(j\in J\text{ with parent-compatible owners}),\\
 |I_R\setminus I_T|+m\le C_Ld,\qquad
 \operatorname{dist}(j,\mathfrak D)\le L+r_{\rm loc}
                                                     \quad(j\in J).
 \end{gathered}                                                   \tag{50.6}
\]
Here and below constants can depend on the
fixed \(L\), however large, not on \(R\) or \(\beta\).

To justify the first line, inherited collar
links have anchors within \(r_{\rm loc}\)
of \(T\). They cannot be anchored in
\(\mathfrak D\), since \(T\) is farther than
\(L\) from it. The collar inclusion is (36.1).
For a new free input, a parent-compatible
owner is outside \(T\), but within
\(r_{\rm loc}\) of a cell of \(T\).
Its distance from \(\mathfrak D\) is thus
between \(L-r_{\rm loc}\) and \(L+r_{\rm loc}\),
with the permitted anchor changes included.
This proves the separation and location claims.
Finally,
\[
 |R\setminus T|\le(2L+1)^4d,\qquad
 I_R\setminus I_T\subset I_{R\setminus T},\qquad
 |I_{R\setminus T}|\le73|R\setminus T|.
 \tag{50.7}
\]
Both the removed variables and \(J\) are
contained in \(I_R\setminus I_T\), proving
the count after increasing \(C_L\).
These arguments also cover disconnected \(T\).

Let \(\mathcal E\) be the event that every
new input \(j\in J\) has section-distance
at most \(\tau_g\). By (50.3) for singleton
sets and a union bound,
\[
       \pi_{R,C}^{y}(\mathcal E^c)\le m e^{-c_g\beta},
 \qquad
       \pi_{R,C}^{\widetilde y}(\mathcal E^c)
                                      \le m e^{-c_g\beta}.
 \tag{50.8}
\]
On \(\mathcal E\), the entire conditional
exterior for \(T\) is in the chosen good
guard, by the definitions of \(\tau_g,b_{\rm in}\).
Its inherited inputs are common to both
parent laws; only the \(m\) new inputs can
differ. This is a random good boundary
inside the \emph{original} parent law,
whose failure probability has just been
estimated without assuming that original
boundary was globally good.

\subsection{Spatial transport on the conditional exterior region}

For the unpinned good-boundary law on \(I_T\),
the proof of (49.15)--(49.17) applies with
no pinned coordinates and with
\[
          p_T=\min\{1,73|T|e^{-c_0\beta}\}
 \tag{50.9}
\]
directly from V32. In particular its
coordinate cost matrix is
\[
       A_{ij}=K_1e^{-\gamma\rho(i,j)}
                          +K_2(\beta+1)p_T ,
       \qquad i\in I_T,\quad j\in J_T^\partial .
 \tag{50.10}
\]
This statement follows from the unpinned
componentwise covariance (37.4), not from
an assumption that V45 permits a growing
set of pinned variables. Here every
coordinate in \(I_T\) is integrated.

Put
\[
 \begin{split}
 b_i&=d_G,\qquad i\in I_R\setminus I_T,\\
 b_i&=d_G\left[
          K_1\sum_{j\in J}e^{-\gamma\rho(i,j)}
             +K_2(\beta+1)p_T m+2m e^{-c_g\beta}
                        \right],\qquad i\in I_T .
 \end{split}                                                     \tag{50.11}
\]
\begin{theorem}[Localized repair coupling of complete parent laws]
\label{thm:f12v46-repair}
Under the stated curvature and coupling hypotheses,
there is a coupling \(\Lambda\) of
\(\pi_{R,C}^{y}\) and \(\pi_{R,C}^{\widetilde y}\)
such that
\[
                     \mathbb E_\Lambda
                       \|Y_i-\widetilde Y_i\|_{\rm HS}
                                     \le b_i
                              \quad\text{for every }i\in I_R.
 \tag{50.12}
\]
Both marginals are the full compact parent
laws, including the excised variables.
\end{theorem}
\begin{proof}
It suffices, by V45's convex separation proof,
to prove the weighted transport bound for
every nonnegative vector of output weights.
Let \(F\) be one-Lipschitz for the weighted
product metric on the whole parent.
Write \(\eta,\widetilde\eta\) for the two
outside-\(I_T\) records, and integrate them
against the product of their actual marginal
laws. This choice is only a device for
bounding expectation differences.

For each such pair, compare the two
conditional expectations of \(F\).
Changing its outside argument from \(\eta\)
to \(\widetilde\eta\), at fixed inside
argument, costs at most
\(d_G\sum_{i\notin I_T}w_i\).
For the remaining inside comparison, freeze
the outside argument of the test.
If both records satisfy \(\mathcal E\),
(50.10) bounds the weighted conditional
transport by
\[
              d_G\sum_{i\in I_T}w_i
                  \left[K_1\sum_{j\in J}e^{-\gamma\rho(i,j)}
                              +K_2(\beta+1)p_Tm\right].
 \tag{50.13}
\]
We used only \(\|\eta_j-\widetilde\eta_j\|_{\rm HS}
\le d_G\). The other exterior inputs are
identical and contribute zero.
If either record fails the guard, the
oscillation of the inside test is at most
\(d_G\sum_{i\in I_T}w_i\).
By (50.8), the probability of this union
event is at most \(2m e^{-c_g\beta}\).

Integrating these bounds proves
\[
 |\pi_{R,C}^{y}F-\pi_{R,C}^{\widetilde y}F|
                                             \le\sum_iw_i b_i .
 \tag{50.14}
\]
This argument integrates measurable conditional
expectations. It does not assume a measurable
choice of separately optimal conditional
couplings. Compact transport duality and
the finite-dimensional separation argument
of Proposition~\ref{prop:f12v45-coupling}
now produce one whole-parent coupling with
all the costs (50.12).
Every coupling of its fixed marginals
automatically has each coordinate cost at
most \(d_G\), so \(b_i\) can also be capped
by \(d_G\) without changing the conclusion.
\end{proof}

If \(T\) is empty, use any parent coupling;
(50.7) pays for all its variables.
If \(\mathfrak D\) is empty, all actual
exterior inputs agree and the two kernels
are identical; use their diagonal coupling.
These cases require no conditional
probability on a nonexistent interface.

\subsection{A cost per defect and an explicit spatial tail}

Let \(c=\min\{c_0,c_g\}>0\).
The bounded density of free coordinates per
owner and the summability of
\(e^{-\gamma\rho(i,j)}\) imply from
(50.6), (50.9) and (50.11):
\[
 \mathsf W_{I_R}(\pi_{R,C}^{y},\pi_{R,C}^{\widetilde y})
           \le C d\left[1+(\beta+1)s_R^2 e^{-c\beta}\right].
 \tag{50.15}
\]
Indeed, the removed region costs \(C d\).
The spatial term sums to \(C m\).
The two error terms, summed over at most
\(73s_R\) outputs, cost at most
\[
       C m\left[(\beta+1)s_R^2 e^{-c_0\beta}
                                      +s_R e^{-c_g\beta}\right].
 \tag{50.16}
\]
Use \(m\le C_Ld\) to obtain (50.15).

The same \emph{single} coupling has the
far-from-defect bound, for \(\ell\ge0\),
\[
 \begin{split}
 &\sum_{\substack{i\in I_R\\
       \operatorname{dist}(\operatorname{owner}(i),\mathfrak D)
                                  \ge L+r_{\rm loc}+\ell+1}}
       \mathbb E_\Lambda\|Y_i-\widetilde Y_i\|_{\rm HS}\\
 &\hspace{1cm}\le
 C d\left[(\ell+1)^4 e^{-\gamma\ell/2}
                           +(\beta+1)s_R^2e^{-c\beta}\right].
 \end{split}                                                     \tag{50.17}
\]
All removed variables are within \(L+r_{\rm loc}\)
of the defect anchors. Every new input
\(j\in J\) is also in that neighborhood.
For outputs in (50.17), the distance to
each such \(j\) is at least \(\ell\).
Sum the exponential kernel as in (49.19)
and use (50.16). This proves the stated tail.

For parents with \(s_R=O(\beta^4)\), the
remainder is polynomial times \(e^{-c\beta}\)
and is exponentially small after reducing
the exponent. The leading cost in (50.15)
is then bounded per defect, without the
crude factor \(\beta s_R\) per changed input.
This is a finite-scale comparison, not an
all-distance assertion at fixed \(\beta\)
as \(s_R\) tends to infinity.

The estimate permits arbitrary compact values
at the recorded bad boundary links and
inputs. Its dependence on \(d\) is not a
Lipschitz dependence on their amplitudes.
It may remain nonzero for a pair with a
very small discrepancy or even identical
data with a nonempty common bad set.
In that latter case a diagonal coupling
is of course available; (50.15) is simply
not the sharp bound. Consequently this
estimate cannot itself be substituted for
an all-boundary path-contraction coefficient.

\begin{decision}[Retain localization without discarding correlated defects]
We have replaced the all-or-nothing cost
of a defective exterior by a localized,
per-defect comparison for the full unpinned
parent law. The internal good interface
was derived from a local sparse estimate,
not imposed on the sampled parent.
The remaining step is to combine this
additive defect comparison with discrepancy
or energy control in an actual iteration.
A rare bad-boundary count alone still
does not imply contraction down to zero.
Nor is the pinned-core version automatic:
under the bad original exterior, its
point-density denominator needs a further
argument. The next companion audit and
direction review remain V50.
\end{decision}

\subsection{Executed checks and append record}

The script
\begin{center}\small\ttfamily
scripts/verify\_ym\_f12\_v46\_defect\_excision.py
\end{center}
uses V31's actual physical collar functions
on parents with 18 and 36 cells and their
excised subregions. It checks inherited
collar inclusion, avoidance of specified
bad frozen links, separation of the new
physical collar, and the removed-link
count. It also checks a complete finite
conditional coupling with exact marginals
and 64 weighted good/bad-interface bounds.
All checks passed. The sample excision
radius tests geometry only; it is not a
numerical evaluation of the possibly much
larger analytic constant \(L\) in (50.5).
These tests do not verify a Yang--Mills
mass gap or the analytic probability
constants.

V45 had 704195 bytes and SHA--256
\begin{center}\scriptsize\ttfamily
952D8FD8EA603C93F8704FF3E0FDA09222E3FA4906132AFCC066E0318B85611F.
\end{center}
Its complete body is preserved apart
from the title version. Only active V45
is replaced by V46; archives, companion
notes and earlier scripts are retained.

\begin{firewall}[Position after V46]
Local frozen-boundary goodness suffices
for sparse suppression away from defects.
Excision then yields a coupling of the
actual complete parent laws with a
localized per-defect cost and an explicit
finite-volume remainder. This does not
yet give discrepancy-weighted contraction,
an unrestricted volume-uniform gap,
scale-stable effective interactions,
continuum construction, or the physical
Yang--Mills mass gap.
\end{firewall}

\section{V47: discrepancy-weighted transport through a defective exterior}
\label{sec:f12v47}

V46 localizes the cost of a defective boundary,
but its additive estimate need not vanish with
the amplitude of a boundary change. We now
differentiate the normalized parent law after
excising the complete source footprint.
The resulting spatial estimate is proportional
to the actual exterior discrepancy.

This reuses the general score method of V35,
the guarded conditional comparison of V46,
and V45's simultaneous-coupling principle.
It does not repeat their core-minorization
or good-boundary theorems. The new input
is a conditional-mean oscillation estimate
for a parent with arbitrary compact values
on its recorded defective exterior.
The coefficient still depends on that
defect record and is not a contraction
coefficient less than one.

\subsection{Freeze the excision throughout one source variation}

Fix \(R,C,\beta\) under V46's hypotheses.
Change one exterior free input \(y_j\)
along a piecewise \(C^1\) path \(y_j(t)\),
keeping all other exterior inputs fixed.
Write
\[
       \mu_t=\pi_{R,C}^{y(t)},\qquad
       v(t)=\|\dot y_j(t)\|_{\rm HS},\qquad
       \ell_j=\int_0^1v(t)\,dt .
 \tag{51.1}
\]
The path need not remain in a good boundary
chart.

Choose one fixed defect-anchor set
\(\mathfrak D_j\) for the entire path.
It contains all initially bad frozen-link
anchors and the full finite footprint of
every parent-frozen link or free interaction
term that can change with \(y_j\).
Consequently it contains every bad
frozen-link anchor along the path and
every differing exterior record between
two times. This construction does not
require sampling all intermediate values:
outside that finite source footprint,
the boundary is unchanged.
The added source footprint has fixed size
by the completed locality constants.

Use this set in V46's deterministic excision:
\[
 T=\{x\in R:\operatorname{dist}(x,\mathfrak D_j)>L\},
 \quad E=I_R\setminus I_T,\quad
 J=J_T^\partial\cap E,\quad m=|J|,\quad
 d_j^{\rm def}=|\mathfrak D_j|.
 \tag{51.2}
\]
If necessary increase the already fixed
\(r_{\rm loc}\), hence \(L\), so that the
full variable support of the source score
is contained in \(E\). This is a finite
range requirement, independent of the
number of variables in \(R\).
The inherited exterior of \(T\) is
unchanged throughout the path and is good.
The new-input good event \(\mathcal E\)
is the same event at every time.
By (50.8),
\[
       \mu_t(\mathcal E^c)\le
       q:=\min\{1,m e^{-c_g\beta}\}
                                    \quad\text{for all }t.
 \tag{51.3}
\]
All counts, guards and geometric regions
are fixed before differentiation.

\subsection{A guarded covariance estimate}

We use the following elementary estimate.
Let \(H,S\) be bounded real random variables,
and let an event \(\mathcal E\) have positive
probability. If
\[
 \operatorname{osc}_{\mathcal E}H\le a,\qquad
 \operatorname{osc}H\le b,\qquad
 \Pr(\mathcal E^c)\le q,
 \tag{51.4}
\]
then
\[
                |\operatorname{Cov}(H,S)|
                      \le2\|S\|_\infty(a+bq).
 \tag{51.5}
\]
To prove this, clip \(H\) to the interval
between its essential infimum and supremum
on \(\mathcal E\), obtaining \(H_c\).
Its oscillation is at most \(a\).
Subtracting an endpoint of that interval
shows
\(|\operatorname{Cov}(H_c,S)|\le2a\|S\|_\infty\).
Moreover
\[
             |H-H_c|\le b\,{\bf1}_{\mathcal E^c}.
 \]
Since \(|S-\mathbb ES|\le2\|S\|_\infty\),
the covariance of this remainder is
bounded by \(2bq\|S\|_\infty\).
Adding the two estimates proves (51.5).
The clipping is a proof device, not a
change to the sampled law.

\subsection{The conditional mean has no direct source derivative}

Write \(\eta=Y_E\) and \(X=Y_{I_T}\).
For a fixed bounded observable \(F\)
of the full parent, set
\[
              H_F(\eta)=\mathbb E_{\mu_t}[F\mid\eta].
 \tag{51.6}
\]
Use the canonical continuous conditional
version when \(F\) is continuous. This
function is independent of \(t\):
after \(\eta\) is fixed, no interaction
meeting \(I_T\) depends on \(y_j(t)\).
Any source-dependent term involving only
\(\eta\) cancels from the conditional
normalizer. Thus the conditional kernel
of \(X\) is the same function of \(\eta\)
at every time.

Let \(V_R\) be the measured parent
potential with its actual compact
normalization understood, and define
the uncentered directional score
\[
             S_t(Y)=\frac{d}{dt}V_R(Y;y(t)).
 \tag{51.7}
\]
Exterior-only constants may be subtracted.
By the choice of the source footprint,
\(S_t\) is a function of \(\eta\), and
the uniform local derivative bound gives
\[
                       \|S_t\|_\infty
                                  \le K(\beta+1)v(t).
 \tag{51.8}
\]
It differentiates only a fixed number of
completed local terms.

Differentiating the full normalized density,
then conditioning on \(\eta\), yields
\[
 \frac{d}{dt}\mu_tF
       =-\operatorname{Cov}_{\mu_t}(F,S_t)
       =-\operatorname{Cov}_{\mu_t}(H_F,S_t).
 \tag{51.9}
\]
The second equality holds because the
score is measurable in \(\eta\).
Equivalently, it differentiates the
actual \(\eta\)-marginal with its exact
normalizer. No derivative of a moving
guard or a moving excision is taken.

The absence of a direct derivative of
\(H_F\) is essential. For a general
parameter-dependent conditional kernel,
one would instead have a term
\(\mathbb E[\partial_tH_F]\).
For example, a fixed fair exterior bit
and a conditional inside mean
\(q(\eta)+t/10\) have derivative \(1/10\)
despite zero exterior marginal score.
Our complete source excision is what
rules out that extra term here.

\subsection{A populated spatial derivative bound}

Let \(F\) be one-Lipschitz for the weighted
product metric with nonnegative weights
\(w_i\) on all of \(I_R\).
Its global oscillation, and therefore
that of \(H_F\), is at most
\[
                         b_w=d_G\sum_{i\in I_R}w_i.
 \tag{51.10}
\]
For two good outside records
\(\eta,\widetilde\eta\), the conditional
comparison (50.13), together with the
change of the outside argument of \(F\),
gives
\[
 \begin{split}
 \operatorname{osc}_{\mathcal E}H_F\le a_w:=
 d_G\bigg\{&\sum_{i\in E}w_i\\
 &+\sum_{i\in I_T}w_i
       \left[K_1\sum_{k\in J}e^{-\gamma\rho(i,k)}
                           +K_2(\beta+1)p_Tm\right]\bigg\},
 \end{split}                                                     \tag{51.11}
\]
where \(p_T=\min\{1,73|T|e^{-c_0\beta}\}\).
This is uniform in the whole fixed
exterior path. The good event has positive
probability under the positive compact
parent density, even when the upper
bound \(q\) in (51.3) equals one.

Apply (51.5) to (51.9), with (51.3),
(51.8), (51.10) and (51.11).
After enlarging fixed constants, define
\[
 B_i(\mathfrak D_j)=
 \begin{cases}
 K_3(\beta+1),&i\in E,\\
 K_3(\beta+1)\left[
        \displaystyle\sum_{k\in J}e^{-\gamma\rho(i,k)}
          +(\beta+1)p_Tm+m e^{-c_g\beta}\right],
                                                    &i\in I_T .
 \end{cases}                                                     \tag{51.12}
\]
For the first line we used \(q\le1\);
for the second \(q\le m e^{-c_g\beta}\).
The result is the simultaneous weighted
derivative inequality
\[
                  \left|\frac{d}{dt}\mu_tF\right|
                       \le v(t)\sum_i w_i B_i(\mathfrak D_j).
 \tag{51.13}
\]
The powers of \(\beta\) in (51.12)
are retained. No claim is made that
the source score is order one.

\begin{theorem}[One coupling with discrepancy-weighted spatial costs]
\label{thm:f12v47-relative}
For the two endpoints of the fixed-source
path, there is a coupling \(\Lambda_j\)
of their complete compact parent laws
such that
\[
        \mathbb E_{\Lambda_j}
             \|Y_i-\widetilde Y_i\|_{\rm HS}
                          \le \ell_j B_i(\mathfrak D_j)
                                   \quad(i\in I_R).
 \tag{51.14}
\]
In a connected compact matrix gauge group,
the path may be chosen so that
\(\ell_j\le C_{\rm path}\|y_j-\widetilde y_j\|_{\rm HS}\),
with a fixed \(C_{\rm path}\).
Thus the bound is zero when the exterior
input difference is zero.
\end{theorem}
\begin{proof}
Integrate (51.13), then use compact
weighted transport duality.
For every nonnegative weight vector,
the weighted optimal cost is at most
\(\sum_iw_i\ell_jB_i\).
The closed-convex cost-vector separation
proof in V45 gives one coupling that
satisfies all coordinate bounds.
This is not a collection of separately
optimized marginal couplings.

For the path assertion, sufficiently
close group elements are joined in a
uniform local chart by a path of length
at most a fixed multiple of their
Hilbert--Schmidt distance.
For pairs outside that chart-distance
threshold, use the finite intrinsic
diameter of the connected compact group
and divide by the positive threshold.
The resulting constant is uniform.
The finite source footprint defining
\(\mathfrak D_j\) covers either choice
of path, so no guard assumption is
lost during this construction.
\end{proof}

If a disconnected group is under
consideration, the derivative argument
applies within a connected component.
For different components, their positive
minimum separation and V46's additive
coordinate bounds give the same form
after increasing constants. This
separate case does not posit a path
between disconnected components.

\subsection{Column cost, spatial tail and several changed inputs}

Let \(c=\min\{c_0,c_g\}\).
Summing (51.12) with V46's counting
and spatial summability bounds gives
\[
 \sum_{i\in I_R} B_i(\mathfrak D_j)
 \le C(\beta+1)d_j^{\rm def}
                  \left[1+(\beta+1)s_R^2e^{-c\beta}\right].
 \tag{51.15}
\]
For outputs farther than
\(L+r_{\rm loc}+\ell+1\) from
\(\mathfrak D_j\), the same sum has
the bound
\[
 C(\beta+1)d_j^{\rm def}
     \left[(\ell+1)^4e^{-\gamma\ell/2}
                       +(\beta+1)s_R^2e^{-c\beta}\right].
 \tag{51.16}
\]
To verify these formulas, the excised
region has at most \(C_Ld_j^{\rm def}\)
variables, the spatial term sums to
\(Cm\), and the two error terms sum
to at most
\[
 C(\beta+1)m
        \left[(\beta+1)s_R^2e^{-c_0\beta}
                                      +s_R e^{-c_g\beta}\right].
 \]
Use \(m\le C_Ld_j^{\rm def}\).
For (51.16), none of the removed
variables is in the specified far set,
and the exponential kernel has the
tail used in (50.17).
Multiplication by \(\ell_j\), or by
the bounded multiple of the actual
input discrepancy, gives the corresponding
costs for the \emph{same} coupling (51.14).

For several changed inputs, choose an
order and vary one coordinate at a time.
At step \(j\), use a fixed record
\(\mathfrak D_j\) containing all bad
background anchors at that step and
that coordinate's source footprint.
The weighted expectation bounds telescope,
giving
\[
 \mathsf W_{d_w}(\pi_{R,C}^{y},\pi_{R,C}^{\widetilde y})
       \le \sum_iw_i\sum_j
                C_{\rm path}\|y_j-\widetilde y_j\|_{\rm HS}
                                      B_i(\mathfrak D_j).
 \tag{51.17}
\]
Apply cost-vector separation once more
to get one endpoint coupling with the
displayed coordinate budgets.
No simultaneous grand random map over
all exterior configurations is asserted.
The records may depend on the chosen
endpoint pair and the coordinate order.

If the entire path lies in V33's good
exterior guard, its sharper unpinned
spatial estimate is already available.
There is no reason to replace it by
the extra factor \(\beta+1\) in (51.12).
The purpose of this section is the
path with a defective background.

\subsection{What this resolves and what it does not}

The new estimate fixes a real defect
of (50.15): the budget now vanishes
with the amplitude of the changed
boundary data, including in the
presence of common bad boundary values.
The rare-interface remainder is also
multiplied by that amplitude.
The normalized-score argument, rather
than division by a small endpoint
distance, supplies this fact for
arbitrarily close inputs.

However, \(\mathfrak D_j\) still contains
bad background anchors that may be
far from the source \(j\).
The leading column cost in (51.15)
is \(C(\beta+1)d_j^{\rm def}\), not a
volume-uniform number below one.
In a coupling of random boundary data,
averaging produces weighted quantities
such as
\[
       \mathbb E\!\left[
          d_j^{\rm def}\,
                   \|Y_j-\widetilde Y_j\|_{\rm HS}\right].
 \tag{51.18}
\]
Stationary sparsity of the defect record
does not factor this expectation.
The present theorem therefore gives
no unrestricted block-contraction
or gap claim.

\begin{decision}[Localize the defect record relative to the source]
V47 now supplies a discrepancy-weighted
spatial comparison under defective
exterior data, at the stated cost.
The next analytic issue is whether
only source-connected bad components
need be charged, or whether a local
energy inequality controls (51.18)
without such a reduction.
Counting every common bad boundary
anchor as if it transmitted this
source remains too expensive.
The parent laws here are unpinned;
a pinned extension still needs its
own uniform density-denominator control.
The next companion audit and direction
review remain V50.
\end{decision}

\subsection{Exact checks and append record}

The script
\begin{center}\small\ttfamily
scripts/verify\_ym\_f12\_v47\_relative\_defects.py
\end{center}
checks 63 normalized derivative,
conditional-tower and clipped-remainder
identities in a positive finite model.
It checks 1323 rational endpoint
inequalities proportional to the
parameter difference and explicitly
tests the missing-direct-derivative
counterexample discussed after (51.9).
All checks passed.
These are finite algebraic checks,
not validation of the analytic
Yang--Mills constants or of a mass gap.

V46 had 720817 bytes and SHA--256
\begin{center}\scriptsize\ttfamily
5E3B6371643813A4463ED62D41DBC272B78E545CD87456D00F7C05B93537F3C6.
\end{center}
Its full body is retained apart from
the title version. Only active V46 is
replaced by V47. Archives, companions
and earlier verification scripts are
preserved.

\begin{firewall}[Position after V47]
The defective-boundary comparison is
now proportional to input discrepancy
and retains a spatial tail in one
whole-parent coupling. Its coefficient
still depends on the complete defect
record. Source-local defect control,
an unrestricted volume-uniform gap,
scale-stable effective interactions,
continuum construction and the physical
Yang--Mills mass gap remain unproved.
\end{firewall}

\section{V48: induced links between defect components and a guarded source resolvent}
\label{sec:f12v48}

V47 pays for every recorded bad component when
one exterior source changes. To improve that
bookkeeping, one must distinguish direct
interactions from interactions generated by
integrating the intervening field.
Separated components are not generally
independent after that integration.
This section bounds their induced interaction
on a good interface domain and derives the
corresponding source-comparison matrix.

The result is for the actual compact
component marginal \emph{conditioned on its
interface guard}. It is not an unrestricted
parent gap. Compact variables in the interiors
of the defect components are not restricted
to a chart. This differs from V39's
all-coordinate boxed convexity argument.
V28's factorization after integrating
separated interiors is also a different
integration: here the intervening region
is integrated and can join distinct components.

\subsection{Exact component marginal after integrating the intervening region}

Use V46's parent \(R\), excised complement
\(E=I_R\setminus I_T\), and good inherited
exterior of \(T\). Partition
\[
                         E=\bigsqcup_{\alpha=1}^q E_\alpha
 \tag{52.1}
\]
so that no completed local potential term
meets two distinct \(E_\alpha\).
For example take connected components
of the hypergraph obtained by intersecting
each completed potential support with \(E\).
This condition includes terms that also
meet \(I_T\); ignoring them would give
an incorrect factorization.

Write \(x=Y_{I_T}\), \(\eta_\alpha=Y_{E_\alpha}\).
The potential decomposes exactly as
\[
 V_R(x,\eta)=V_T(x)+
        \sum_\alpha\bigl[U_\alpha(\eta_\alpha)
                            +V_\alpha(x,\eta_\alpha)\bigr].
 \tag{52.2}
\]
Here \(V_\alpha\) is the sum of terms
meeting both \(I_T\) and \(E_\alpha\).
Let \(K_\alpha\subset E_\alpha\) be
its free-coordinate input set in that
component. Then
\[
 \begin{split}
 Z_T(\eta)&=\int_{G^{I_T}}
       \exp\{-V_T(x)-\sum_\alpha V_\alpha(x,\eta_\alpha)\}\,dx,\\
 W_T(\eta)&=-\log Z_T(\eta),\\
 \rho(d\eta)&=Z_E^{-1}
          \exp\{-\sum_\alpha U_\alpha(\eta_\alpha)-W_T(\eta)\}\,d\eta .
 \end{split}                                                     \tag{52.3}
\]
All integrals over \(x\) are full compact
integrals. The normalizer \(Z_T\) depends
only on \(\eta_{\cup K_\alpha}\), although
the \(U_\alpha\) may involve arbitrary
compact component interiors.
Equation (52.3) follows from integration
of (52.2), not from deleting the crossing
terms or normalizing each component independently.

Choose a fixed small cube of logarithmic
coordinates on each \(K_\alpha\), with
scalar coordinate radius \(a\), such that
all these interface values, together with
the inherited exterior, put \(T\) in
V32--V33's good boundary guard. Put
\[
 \Gamma=\prod_\alpha\Gamma_\alpha,\qquad
 \Gamma_\alpha=
    \{\eta_{K_\alpha}\text{ in the chosen cubes}\}
                 \times G^{E_\alpha\setminus K_\alpha},
 \qquad \rho^\Gamma=\rho(\,\cdot\mid\Gamma).
 \tag{52.4}
\]
Product Haar measure is used in compact
coordinates; the interface Haar Jacobians
are absorbed in the corresponding
\(U_\alpha\) when using Lebesgue coordinates.
They introduce no cross-component derivative.
Every \(\Gamma_\alpha\) is a compact product
with a positive conditional density.
The event \(\Gamma\) restricts only the
component interfaces, not their interiors
and not the integrated variables \(x\).

\subsection{The induced mixed Hessian is an actual conditional covariance}

In formulas with scalar derivatives, expand
\(K_\alpha\) into its fixed-dimensional Lie
coordinate components and denote the resulting
number by \(k_\alpha\).
Let \(k=\sum_\alpha k_\alpha\).
Distances \(\rho(a,b)\) refer to the owner
cells of the underlying free coordinates.
For \(a\in K_\alpha\), set
\(S_a=\partial_a V_\alpha(x,\eta_\alpha)\).
The local bounds are
\(\|S_a\|_\infty\le C\beta\) and
\(\|\partial_xS_a\|_\infty\le C\beta\),
for \(\beta\ge1\), with derivative support
near \(a\) of fixed size.

\begin{proposition}[Spatial mixed Hessian of the compact component marginal]
\label{prop:f12v48-mixed}
For \(\alpha\ne\sigma\), \(a\in K_\alpha\),
\(b\in K_\sigma\), and \(\eta\in\Gamma\),
\[
 \partial_a\partial_b W_T(\eta)
       =-\operatorname{Cov}_{\pi_T^\eta}(S_a,S_b),
 \qquad
 |\partial_a\partial_b W_T(\eta)|
           \le C_1\beta e^{-\gamma\rho(a,b)}
                          +C_2\beta^2p_T,
 \tag{52.5}
\]
where \(p_T=\min\{1,73|T|e^{-c_0\beta}\}\).
All cross-component mixed derivatives involving
a component coordinate outside \(K_\alpha\)
vanish in \(W_T\).
\end{proposition}
\begin{proof}
Positive compact integrands permit differentiation
under the fixed integration domain.
First \(\partial_aW_T=\pi_T^\eta S_a\).
Differentiating the normalized expectation gives
\[
 \partial_b\partial_aW_T
        =\pi_T^\eta(\partial_bS_a)
                       -\operatorname{Cov}_{\pi_T^\eta}(S_a,S_b).
 \]
No term in (52.2) meets both components,
so \(\partial_bS_a=0\).
The conditional bulk law \(\pi_T^\eta\)
has a good exterior on \(\Gamma\).
Use V33's componentwise boxed covariance,
then the compact-to-box error from V32.
Each score has only finitely many bulk
derivative components, all of size \(C\beta\).
The boxed term is bounded by
\(C\beta^{-1}\beta^2e^{-\gamma\rho(a,b)}\),
with fixed support-radius shifts absorbed
in constants and exponent.
The compact comparison term is bounded
by \(6\|S_a\|_\infty\|S_b\|_\infty p_T\).
This proves (52.5).
Finally \(Z_T\) does not depend on a
component variable absent from all its
crossing terms, proving the last assertion.
\end{proof}

Thus a path through the integrated region
generates a generally nonzero component
interaction. Its estimate does not require
convexity inside a bad component.
It does require the interface condition
used to control the conditional bulk law.

\subsection{A rectangular log-density bound and conditional TV coefficients}

Define a symmetric nonnegative matrix with
zero diagonal by
\[
 \Omega_{\alpha\sigma}
   =4a^2\sum_{\substack{a'\in K_\alpha\\ b'\in K_\sigma}}
       \left[C_1\beta e^{-\gamma\rho(a',b')}
                                     +C_2\beta^2p_T\right],
                    \qquad \alpha\ne\sigma .
 \tag{52.6}
\]
Empty interface sets give zero entries.
Let \(\Phi(\eta)=\sum_\alpha U_\alpha(\eta_\alpha)+W_T(\eta)\).
For any two values of component \(\alpha\)
and any two values of component \(\sigma\),
with all interfaces in their cubes, the
rectangular increment obeys
\[
                      |\Delta_\alpha\Delta_\sigma\Phi|
                                      \le\Omega_{\alpha\sigma}.
 \tag{52.7}
\]
Indeed the \(U_\alpha\) cancel in that
increment. Only the interface coordinates
enter \(W_T\). Integrate (52.5) along the
two straight coordinate paths inside the
product cubes; every scalar change is
at most \(2a\). Compact interior variables
of either component may also change,
but they do not occur in \(W_T\).

Consider two conditional laws of component
\(\alpha\) under \(\rho^\Gamma\), whose
other component values differ only at
\(\sigma\). The oscillation, over the
updated component, of the difference
of their unnormalized potentials is
at most \(\Omega_{\alpha\sigma}\)
by (52.7). Therefore their TV distance
is at most
\[
                 c_{\alpha\sigma}
                         \le\min\{1,\Omega_{\alpha\sigma}\}.
 \tag{52.8}
\]
For completeness, if two probabilities
are related by a bounded potential
increment \(A\), interpolate by the
normalized densities \(e^{-tA}\).
For a test \(0\le f\le1\), its derivative
is \(-\operatorname{Cov}(f,A)\), whose
absolute value is at most
\(\operatorname{osc}A\).
Integrating from zero to one and taking
the supremum over such tests proves
TV at most \(\operatorname{osc}A\).
The trivial bound one gives (52.8).
Every conditional normalizer is included.

The matrix is an upper bound, not an
assertion that components are exactly
independent when they do not share a
microscopic interaction term.
This is the weak interaction that a
source-local argument must retain.

\subsection{When the component matrix is a contraction}

Suppose \(k_\alpha\le k_*\), and that
different nonempty interface sets are
at owner distance at least \(s\).
The bounded coordinate density in four
dimensions and the exponential shell sum
give
\[
 r:=\max_\alpha\sum_{\sigma\ne\alpha}\Omega_{\alpha\sigma}
 \le C k_*\left[
       \beta(s+1)^4e^{-\gamma s/2}
                                  +\beta^2p_T k\right].
 \tag{52.9}
\]
To prove it, sum (52.6) first over all
interface coordinates outside \(K_\alpha\).
For each fixed coordinate of \(K_\alpha\),
their positions are at distance at least
\(s\) and have uniformly bounded multiplicity.
The exponential tail bounds that sum.
The constant remainder is counted once
for each of the at most \(k_\alpha k\)
ordered coordinate pairs.
The fixed \(a\) and group dimension
are included in \(C\).

In particular, take fixed \(k_*\) and
\(|R|\le C_{\rm vol}\beta^4\), and require
\[
               s\ge \frac{8}{\gamma}\log(\beta+2).
 \tag{52.10}
\]
Then the first term of (52.9) tends to
zero as a polynomial logarithm times
\(\beta^{-3}\). Since \(k\le C|R|\),
the second is at most
\(C\beta^2|R|^2e^{-c_0\beta}\).
For sufficiently large \(\beta\),
\[
                              r\le\tfrac12 .
 \tag{52.11}
\]
Constants are uniform over the number
of components satisfying these geometric
and interface-size hypotheses within
this finite-volume window.
The excision in V46 does not automatically
guarantee (52.10) or a fixed \(k_*\);
they are additional explicit conditions.
No fixed-\(\beta\) thermodynamic-limit
claim follows from this calculation.

\begin{theorem}[Gap of the guarded component heat baths]
\label{thm:f12v48-component-gap}
If \(r<1\), the rate-one-per-component
heat-bath generator for \(\rho^\Gamma\)
has spectral gap at least \(1-r\).
\end{theorem}
\begin{proof}
Let \(P=q^{-1}\sum_\alpha\Gamma_\alpha\)
be random scan, where \(\Gamma_\alpha\)
now denotes the conditional-expectation
update, not the component domain in (52.4).
Equip configurations with the Hamming
distance counting unequal component records.
For a pair differing only at component
\(\sigma\), updating \(\sigma\) removes
the discrepancy. Updating another component
\(\alpha\) creates a new one with probability
at most \(\Omega_{\alpha\sigma}\), by a
maximal coupling of its actual conditionals.
The symmetric matrix \(\Omega\) has column
sums at most \(r\). The expected distance
after one common update choice is thus
at most \(1-(1-r)/q\).
Changing several records one at a time
and the coupling triangle inequality
give this contraction for arbitrary pairs.
Positive continuous conditional densities
allow measurable maximal couplings using
their pointwise common minimum.

Every bounded function is Lipschitz for
this Hamming distance, with constant at
most its global oscillation, and the
distance has diameter \(q\).
Iterating gives a uniform exponentially
decaying oscillation bound for \(P^nf\).
For a centered bounded \(f\), its \(L^2\)
norm after iteration has the same rate.
Since \(P\) is a self-adjoint positive
mixture of conditional projections,
the spectral-measure argument of V36,
and density of bounded functions in
centered \(L^2\), show that its centered
spectrum is at most \(1-(1-r)/q\).
The generator \(q(I-P)\) therefore has
gap at least \(1-r\).
\end{proof}

This is an auxiliary component-update
gap for the declared guarded marginal.
An update can change an entire compact
defect component; it is not asserted
to be a fixed-size update when component
sizes grow. Nor is this the physical
transfer Hamiltonian.
The general influence-matrix framework
is classical; compare Rebeschini--van Handel,
\href{https://arxiv.org/abs/1308.4117}
{Comparison Theorems for Gibbs Measures}.
The elementary finite-system argument
above states the exact version used here.

\subsection{A source is propagated through weighted component paths}

Let \(\rho^\Gamma\) and
\(\widetilde\rho^\Gamma\) have the same
integrated interaction \(W_T\), but
different potentials \(U_0,\widetilde U_0\)
in one component, indexed by \(0\)
in this subsection. Suppose that the
conditional law of component \(0\)
changes in TV by at most \(b_0\),
uniformly over all the other component
values in their guarded domains.
All other conditional changes at
identical input records are zero.
The matrix \(\Omega\) bounds influences
under both laws, since a one-component
potential cancels from the rectangular
increments (52.7).

\begin{proposition}[Guarded source resolvent]
\label{prop:f12v48-source}
If \(r<1\), the two guarded component
laws admit a coupling with
\[
 \Pr(\eta_\alpha\ne\widetilde\eta_\alpha)
             \le [(I-\Omega)^{-1}]_{\alpha0}b_0,
 \qquad
                   (I-\Omega)^{-1}=\sum_{n\ge0}\Omega^n .
 \tag{52.12}
\]
\end{proposition}
\begin{proof}
Start both random-scan chains from
the same fixed record and couple
their update choices. Use maximal
couplings of the conditional updates.
If \(v_t\) is the vector of disagreement
probabilities, changing differing
input components one at a time gives
\[
 v_{t+1}\le
   \left[(1-1/q)I+\Omega/q\right]v_t
                                      +(b_0/q)e_0,\qquad v_0=0.
 \tag{52.13}
\]
The row sums of the bracketed matrix
are less than one, so its geometric
series yields
\(v_t\le(I-\Omega)^{-1}b_0e_0\).
The contraction proof above implies
convergence in TV of both marginals
to their respective invariant laws.
Take a weak subsequential limit of
the coupled laws on the compact
product of the two configuration spaces.
Its marginals are the required laws.
The set where a given pair of component
records differs is open; the corresponding
indicator is lower semicontinuous.
Its probability in the limit is at
most the liminf of the probabilities
along that subsequence. Hence every
bound passes to the limiting coupling.
\end{proof}

For a physical exterior source whose
entire changing potential footprint
lies in \(E_0\), V47's local score
bound and a fixed-domain normalized
conditional derivative give
\[
             b_0\le
        \min\{1,C(\beta+1)\|y_j-\widetilde y_j\|_{\rm HS}\}.
 \tag{52.14}
\]
Indeed the conditional potential
derivative is bounded by the same
fixed number of local source terms;
testing against \(0\le f\le1\) and
integrating the normalized covariance
proves the bound. Its support in
\(E_0\) is essential: if it also
changes \(W_T\), the source vector
in (52.13) has further terms.
Several directly affected components
can instead be represented by a
vector \(b\), giving \((I-\Omega)^{-1}b\).

The entry \((\Omega^n)_{\alpha0}\)
is a sum of products over \(n\)-step
component paths from the source to
the target. Thus remote components
are paid through their induced links,
not through an unweighted total
defect count. The links usually do
not vanish, so this is not a license
to omit every component outside a
microscopic connected cluster.

\subsection{Nonzero induced links and the remaining guard}

A finite example makes that distinction
explicit. Take an Ising chain with
edge factors \(1+t s_i s_{i+1}\),
\(0<t<1\), and sum over its interior
spins. For a chain of \(n\) edges,
the endpoint marginal is
\[
                       \Pr(a,b)=\tfrac14(1+t^n ab).
 \tag{52.15}
\]
Successive summation uses
\(\sum_{s=\pm1}(1+u as)(1+v sb)=2(1+uv ab)\);
iteration proves the formula.
The endpoints have no direct edge
when \(n>1\), but their induced
conditional influence is \(t^n\),
not zero. At zero endpoint fields,
the rate-one two-endpoint heat-bath
generator has eigenvalues
\[
                         0,\quad1-t^n,\quad1+t^n,\quad2.
 \tag{52.16}
\]
The functions \(1,a+b,a-b,ab-t^n\)
are the corresponding independent
eigenfunctions, as direct conditional
expectation verifies.
This is a check of the component
mechanism, not a Yang--Mills example.

The interface restriction in (52.4)
remains substantive. When the interface
cubes contain the small group balls
used in V46's local sparse estimate,
with its thresholds chosen accordingly,
that estimate gives
\[
                       \rho(\Gamma^c)
                              \le C k e^{-c_g\beta}.
 \tag{52.17}
\]
For a polynomial-size parent this
probability is small. It is not
uniformly small for an arbitrarily
large number of interfaces at fixed
\(\beta\), and it does not compare
Dirichlet energies concentrated on
\(\Gamma^c\) with those on \(\Gamma\).
The gap theorem for \(\rho^\Gamma\)
therefore cannot simply be transferred
to the unrestricted marginal \(\rho\).

\begin{decision}[Carry induced links, then treat guard failure in energy]
There is now a normalized, spatially
decaying matrix for interactions
between distinct defect components
on their good interface domain.
For bounded, sufficiently separated
interfaces it is a contraction and
gives a source resolvent rather than
V47's full defect-count charge.
The remaining tasks are to handle
large or closely spaced components
and the energy carried by failed
interfaces under the actual law.
Neither a small probability in
(52.17) nor a gap for the conditioned
component sampler resolves those
tasks. The next companion audit and
direction review remain V50.
\end{decision}

\subsection{Checks and append record}

The script
\begin{center}\small\ttfamily
scripts/verify\_ym\_f12\_v48\_component\_links.py
\end{center}
checks four exact integrated Ising
chains, their two-component heat-bath
eigenfunctions, and source-resolvent
bounds under rational endpoint tilts.
It also checks three Gaussian
Schur-complement mixed entries against
the covariance of their boundary
scores. All checks use rational
arithmetic and passed.
They demonstrate nonzero induced
links and test the algebra; they
do not verify the Yang--Mills
analytic constants or mass gap.

V47 had 735238 bytes and SHA--256
\begin{center}\scriptsize\ttfamily
A9BB3F56C3FB3F27B14D9E4BD458D9985DD2F29B7DAB843C0CDB770A949E173A.
\end{center}
Its body is preserved apart from
the title version. Only active V47
is replaced by V48. Archives,
companions and earlier scripts are
retained.

\begin{firewall}[Position after V48]
The guarded compact component marginal
has quantified induced cross-component
interactions. Under the explicit
interface-size, separation and row-sum
conditions, its component sampler has
a gap and a source-local comparison
resolvent. These are not unrestricted
parent or physical Hamiltonian gaps.
Failed-interface energy control,
scale-stable effective interactions,
continuum construction and the full
Yang--Mills mass gap remain unproved.
\end{firewall}

\section{V49: removing the component interface guard by a controlled sweep}
\label{sec:f12v49}

V48 proved a gap only after conditioning the
component marginal on good interfaces.
Small stationary probability of guard failure
does not transfer that gap to the original
law. The additional input used here is
uniform recovery of one component's interface
under its \emph{unrestricted} conditional
heat bath, even if all other components
are bad. Together with a sweep energy
inequality, it controls functions concentrated
on failed interfaces.

We first give a precise geometric route
to the recovery estimate for aligned
patches. Then we prove an abstract
guard-removal theorem and apply it.
Its unrestricted conclusion concerns
the component marginal, not yet fixed-range
updates of every bulk variable or the
physical transfer Hamiltonian.

\subsection{Aligned active regions and a checkable clearance hypothesis}

Retain the notation of V48, but impose
an additional alignment condition.
There are disjoint cell patches
\(D_\alpha\) with
\[
 R=T\sqcup\bigsqcup_{\alpha=1}^qD_\alpha,\qquad
 R_\alpha=T\cup D_\alpha,\qquad
 I_{R_\alpha}=I_T\sqcup E_\alpha,\qquad
 I_R=I_T\sqcup\bigsqcup_\alpha E_\alpha .
 \tag{53.1}
\]
No completed interaction term meets two
distinct \(E_\alpha\), as in V48.
These equalities must be checked for
the actual free-coordinate updates.
They hold for separated cell patches
with their full boundary families:
an edge cannot meet two different
patches if those patches are separated
by more than one cell, and any edge
meeting \(T\) is already included
in \(I_T\). Taking separation larger
than the completed interaction range
also gives the no-cross-term condition.
A generic partition of free coordinates
need not satisfy (53.1).

The interface sets \(K_\alpha\) may be
enlarged by the fixed physical-input
completion needed for the collar tests.
If so, use those enlarged sets in the
product guard, size and separation
hypotheses of V48. Its mixed derivative
proof remains valid: added coordinates
absent from \(Z_T\) have zero derivative.
This is a stronger declared guard,
not an identification with an earlier,
larger guarded domain.

Fix the free tolerance \(\tau_g\)
whose section balls lie in the chosen
interface cubes, and let \(h_g,c_g\)
be V46's local sparse-bound constants.
For every \(\alpha\), require the following
\emph{physical clearance condition}.
For every assignment of the other
component variables, every frozen
physical link in \(\mathcal F_{R_\alpha}\)
whose anchor is within distance \(h_g\)
of a parent-compatible owner of
\(K_\alpha\) has mismatch at most
\(b_{\tau_g}\).
Affected retained curvatures satisfy
the same uniform small threshold.

This is a deterministic local condition.
It does not assume that the other
components are good. One sufficient
way to verify it is to separate
\(K_\alpha\) from the original bad
frozen-link anchors by more than
\(h_g+r_{\rm loc}\), and from every
other completed interface input set
by more than \(h_g+2r_{\rm loc}\).
Indeed the inherited frozen links
away from the original bad anchors
are fixed and good. A newly frozen
physical link that can vary with
another component has a free ancestor
in that component's completed input
closure, within \(r_{\rm loc}\).
The latter separation excludes every
such varying link from the neighborhood
being tested. The use of completed
input sets is necessary for this
argument; uncompleted plaquette supports
are not enough.

\begin{proposition}[Uniform interface recovery with arbitrary other components]
\label{prop:f12v49-recovery}
Under (53.1) and the physical clearance
condition, the unrestricted component
marginal \(\rho\) satisfies
\[
 \rho\{\eta_\alpha\notin\Gamma_\alpha
                       \mid\eta_{\alpha^c}\}
        \le\varepsilon_\alpha:=
                        |K_\alpha|e^{-c_g\beta}
                         \quad\text{for every }\eta_{\alpha^c}.
 \tag{53.2}
\]
Here \(|K_\alpha|\) counts free group
variables; using the larger scalar
coordinate count is also a valid bound.
\end{proposition}
\begin{proof}
Conditional on \(\eta_{\alpha^c}\),
the joint variables \((Y_{I_T},\eta_\alpha)\)
have exactly the canonical law on
\(I_{R_\alpha}\) by (53.1) and
Gibbs disintegration. Other component
values are part of its frozen exterior.
The clearance hypothesis says that
every possibly bad frozen-link anchor
for this active region is farther
than \(h_g\) from \(K_\alpha\).
Apply V46's localized sparse bound
to each interface variable under
this canonical law, then take the
marginal on \(E_\alpha\).
The union bound gives (53.2), since
the specified section balls lie
inside \(\Gamma_\alpha\).
All conditional densities are positive
and canonical, so the assertion holds
pointwise for the frozen values,
not just on a stationary full-measure set.
\end{proof}

This is stronger than V48's unconditional
bound on \(\rho(\Gamma^c)\).
It is the new geometric input needed
to remove that guard. We do not claim
it for patch families lacking the
alignment or clearance condition.

\subsection{A projection sweep has a dimension-free energy cost}

We use a general Hilbert-space inequality.
For orthogonal projections \(P_1,\ldots,P_q\),
let \(K=P_q\cdots P_1\). Then
\[
              \|f\|^2-\|Kf\|^2
                       \le4\sum_{\alpha=1}^q
                                   \|(I-P_\alpha)f\|^2.
 \tag{53.3}
\]
This is a projection union-bound inequality,
not a new general principle; compare Gao,
\href{https://arxiv.org/abs/1410.5688}
{Quantum union bounds for sequential projective measurements}.
The following proof gives the exact
Hilbert-space version used here.

Put \(f_0=f\), \(f_\alpha=P_\alpha f_{\alpha-1}\),
and \(a_\alpha=f_{\alpha-1}-f_\alpha\).
Orthogonality at each step gives
\[
 A:=\sum_\alpha\|a_\alpha\|^2
                 =\|f\|^2-\|Kf\|^2,\qquad
 f-Kf=\sum_\alpha a_\alpha .
 \tag{53.4}
\]
Also \(a_\alpha\in\ker P_\alpha\), hence
\[
 \begin{split}
 \operatorname{Re}\langle f-Kf,f\rangle
 &=\sum_\alpha\operatorname{Re}
            \langle a_\alpha,(I-P_\alpha)f\rangle\\
 &\le\sqrt A
       \left(\sum_\alpha\|(I-P_\alpha)f\|^2\right)^{1/2}.
 \end{split}
\]
The identity
\[
 A=2\operatorname{Re}\langle f-Kf,f\rangle-\|f-Kf\|^2
 \]
now implies (53.3); if \(A=0\) it
is immediate, otherwise divide by
\(\sqrt A\) and square.
The projections need not commute.
In particular no factor \(q\) is
paid in (53.3).

\subsection{A reversible sweep reaches the good set uniformly}

Let \(P_\alpha\) be the actual unrestricted
heat-bath projection of component \(\alpha\)
under \(\rho\), and write
\[
 \mathcal E_H(f)=\sum_\alpha
                    \langle f,(I-P_\alpha)f\rangle_\rho .
 \tag{53.5}
\]
For this subsection suppose generally
that (53.2) holds with constants
\(\varepsilon_\alpha\), and put
\[
 \varepsilon=\max_\alpha\varepsilon_\alpha,\qquad
             a=1-\sum_\alpha\varepsilon_\alpha>0.
 \tag{53.6}
\]
Any complete deterministic sweep that
updates each component once ends in
\(\Gamma=\prod_\alpha\Gamma_\alpha\)
with probability at least \(a\),
uniformly over its initial record.
At the time a component is updated,
its failure probability is at most
\(\varepsilon_\alpha\), whatever the
then-current other values.
Later updates of other components
do not change that component's value.
The final union bound proves the claim.
No independence of the component
outcomes is assumed.

Take \(K=P_q\cdots P_1\) and
\[
                              Q=K^*K .
 \tag{53.7}
\]
As a kernel this is a sweep followed
by the reverse sweep, up to the
choice of operator-order convention.
Both visit every component.
The kernel \(Q\) is Markov, preserves
\(\rho\), and is self-adjoint and positive.
Its last full sweep implies
\[
          Q(x,\Gamma)\ge a\quad\text{for every }x,\qquad
          \mathcal E_Q(f):=\langle f,(I-Q)f\rangle_\rho
                                      \le4\mathcal E_H(f).
 \tag{53.8}
\]
The energy inequality is (53.3).
The sweep is a comparison device;
the target form remains the sum of
the original unrestricted component
heat-bath forms.

\subsection{The correct within-good conductance normalization}

Let \(p=\rho(\Gamma)>0\), and suppose
the component heat baths of \(\rho^\Gamma\)
have rate-one gap at least \(\gamma>0\).
Define
\[
 q_\alpha(\eta_{\alpha^c})
     =\rho(\Gamma_\alpha\mid\eta_{\alpha^c})
                       \ge1-\varepsilon .
 \tag{53.9}
\]
For exterior component records in their
good domains, restrict both endpoints
of a full component update to
\(\Gamma_\alpha\).
Its contribution to the unrestricted
Dirichlet form is exactly
\[
 p\,\rho^\Gamma\!\left[
       q_\alpha(\eta_{\alpha^c})
            \operatorname{Var}_{\rho^\Gamma}
                      (f\mid\eta_{\alpha^c})\right].
 \tag{53.10}
\]
To verify the factor, start from
the two-copy conditional variance
integral with the unrestricted
conditional law of \(\eta_\alpha\).
Restricting both copies gives
\(q_\alpha^2\) times the two-copy
variance of its normalized good
conditional. The good exterior
marginal under \(\rho^\Gamma\)
equals the unrestricted exterior
marginal times \(q_\alpha/p\).
Thus precisely one factor
\(q_\alpha\) remains in (53.10).

Dropping all other nonnegative
conductances and summing gives
\[
       \mathcal E_H(f)\ge
       p(1-\varepsilon)\gamma
                            \operatorname{Var}_{\rho^\Gamma}(f).
 \tag{53.11}
\]
This does not replace the full
conditional update by a renormalized
good-only update without accounting
for its acceptance probability.

\begin{theorem}[Guard-to-full component gap transfer]
\label{thm:f12v49-transfer}
Assume (53.2), (53.6), and a guarded
component gap \(\gamma>0\).
Then the unrestricted form (53.5)
satisfies
\[
 \operatorname{Var}_\rho(f)
      \le \frac1a\left[
           8+\frac{2-a}{(1-\varepsilon)\gamma}\right]
                                      \mathcal E_H(f).
 \tag{53.12}
\]
In particular its gap is at least
\[
                    \frac{a}
                 {8+(2-a)/((1-\varepsilon)\gamma)}.
 \tag{53.13}
\]
\end{theorem}
\begin{proof}
Subtract the good-set mean, setting
\(g=f-\rho^\Gamma f\), and write
\(B=\Gamma^c\).
By uniform access in (53.8),
\[
 \begin{split}
 a\int_B g(x)^2\,\rho(dx)
 &\le \int_B\rho(dx)\int_\Gamma Q(x,dy)\,g(x)^2\\
 &\le 2\int_B\rho(dx)\int_\Gamma Q(x,dy)
                                    [g(x)-g(y)]^2\\
 &\quad+2\int_B\rho(dx)\int_\Gamma Q(x,dy)\,g(y)^2 .
 \end{split}
\]
Reversibility identifies the first
cross integral as part of
\(\mathcal E_Q(g)\); both cross
orientations already occur in its
one-half two-copy definition.
For the second, reverse the kernel
and use \(Q(y,B)\le1-a\). Hence
\[
        a\int_B g^2\,d\rho
           \le2\mathcal E_Q(f)
                         +2(1-a)\int_\Gamma g^2\,d\rho .
 \tag{53.14}
\]
Adding the good-set contribution and
using \(\operatorname{Var}_\rho(f)\le\rho g^2\)
gives
\[
 \operatorname{Var}_\rho(f)
 \le \frac{2-a}{a}\,
                  p\operatorname{Var}_{\rho^\Gamma}(f)
                              +\frac2a\mathcal E_Q(f).
 \]
Apply (53.11) to the first term
and (53.8) to the second.
This proves (53.12) for all \(L^2\)
functions, by truncation if needed.
In particular the proof controls
functions whose mass is concentrated
on \(B\); it does not discard that
part because \(\rho(B)\) is small.
\end{proof}

\subsection{Application to the aligned separated component marginal}

Apply Proposition~\ref{prop:f12v49-recovery}
and V48's guarded gap.
If
\[
        r\le\tfrac12,\qquad
        \sum_\alpha |K_\alpha|e^{-c_g\beta}\le\tfrac12 ,
 \tag{53.15}
\]
then \(\gamma\ge1/2\), \(a\ge1/2\)
and \(\varepsilon\le1/2\).
The transfer theorem yields the explicit
unrestricted component bound
\[
                 \operatorname{gap}(H_\rho)\ge\frac1{28}.
 \tag{53.16}
\]
Indeed (53.13) is bounded below by
\((1/2)/(8+6)=1/28\).
All component variables, including
their interfaces, are now sampled
from the unrestricted compact
marginal \(\rho\), not \(\rho^\Gamma\).

For fixed interface-size bound
\(k_*\), a parent with
\(|R|\le C_{\rm vol}\beta^4\),
and V48's logarithmic separation,
the first condition of (53.15)
holds at sufficiently large \(\beta\).
The sum in the second is at most
\(C|R|e^{-c_g\beta}\), and also
tends to zero. The alignment and
physical clearance conditions remain
additional requirements.
The clearance uses a fixed radius;
the logarithmic separation eventually
exceeds it when the original bad
anchors are also sufficiently far
inside their excised patches.
These are sufficient geometric
hypotheses, not a claim that every
arbitrary exterior produces such
a separated family.

\subsection{What lifts to the full parent law}

The exact parent disintegration is
\[
                         \mu(dx,d\eta)=\rho(d\eta)\pi_T^\eta(dx).
 \tag{53.17}
\]
Using the conditional-variance lift
already employed in V40, the gap
(53.16) gives an auxiliary full-parent
inequality
\[
 \operatorname{Var}_\mu(f)
 \le28\left[
      \mu\operatorname{Var}(f\mid\eta)
       +\sum_\alpha
          \mu\operatorname{Var}(f\mid\eta_{\alpha^c})\right].
 \tag{53.18}
\]
For clarity, set \(F(\eta)=\pi_T^\eta f\).
The total variance is
\(\mu\operatorname{Var}(f\mid\eta)
+\operatorname{Var}_\rho(F)\).
The component gap bounds the second
term by \(28\sum_\alpha
\rho\operatorname{Var}(F\mid\eta_{\alpha^c})\).
Conditional total variance bounds
each latter summand by
\(\mu\operatorname{Var}(f\mid\eta_{\alpha^c})\).
This proves (53.18).

The updates in (53.18) are the
canonical heat baths on \(I_T\)
and on \(I_T\cup E_\alpha\).
They include the entire intervening
region, which may be large.
Consequently (53.18) is not yet
a volume-uniform gap for bounded-range
updates of all free coordinates.
The compact bulk integration has
not disappeared from the cost of
implementing the lifted updates.

\begin{decision}[The guard is removed in the declared component regime]
V49 resolves the failed-interface
energy issue for aligned, physically
clear, sufficiently separated bounded
interfaces: the original component
marginal has a gap, with no discarded
guard event. The proof uses uniform
conditional recovery and a controlled
sweep, not stationary rarity alone.
The next task is to handle close or
large defect components and localize
the bulk updates in (53.18) without
losing uniform constants.
The next version is V50 and must
perform both the companion audit
and the ten-version direction review.
\end{decision}

\subsection{Exact tests and append record}

The script
\begin{center}\small\ttfamily
scripts/verify\_ym\_f12\_v49\_guard\_transfer.py
\end{center}
checks positive interacting models
with two and three three-state
components. Their heat-bath projections
do not commute. Rational matrix
elimination verifies the factor-four
sweep energy inequality, uniform
access to the good set, the
within-good normalization factor,
and the full transferred Poincare
inequality. All checks passed.
These finite models test the
functional-analytic algebra, not
the geometric clearance condition
or the analytic Yang--Mills constants.

V48 had 753881 bytes and SHA--256
\begin{center}\scriptsize\ttfamily
560D03C7398CD1B833A7DE0205D534B0E49748D96A98F10E7B3D34E60D500954.
\end{center}
Its complete body is retained apart
from the title version. Only active
V48 is replaced by V49; archives,
companions and prior scripts remain.

\begin{firewall}[Position after V49]
The interface guard can be removed
for the component marginal under
the stated alignment, clearance,
separation and size hypotheses,
with gap at least \(1/28\).
The full-parent lift still uses
large intervening-region updates.
General defect geometry, a
volume-uniform bounded-range gap,
scale-stable effective interactions,
continuum construction and the
physical Yang--Mills mass gap
remain unproved.
\end{firewall}

\section{V50: bounded-range bad-set smoothing and the exact
martingale commutator debit}
\label{sec:f12v50}

\subsection{Companion audit and ten-version direction review}

The active companions were inspected again for this mandatory
five-version audit.  They are unchanged from V45:
\begin{center}
\begin{tabular}{lll}
\toprule
Programme & Active version & Bytes\\
\midrule
SM--Einstein & V72 & 607372\\
Navier--Stokes & V26 & 278283\\
Parallel YM & V5 & 54174\\
\bottomrule
\end{tabular}
\end{center}
Their respective SHA--256 values are
\begin{center}\scriptsize\ttfamily
F44F9745D43379E1672C5550411B58E97195A4C9278592D221DF5D3F644DC1F5\\
82C887F8C2C69E4F28FAD7C3DC8DE5B9099CB5A4BF431EBA1FFF3E91935978AD\\
306F1BB84CFA8A8D47EA569EC17E9545F9867C8D32B548EC9D9C444B4F3C0512.
\end{center}
SM V72 concerns a fixed-mode, fixed-time many-copy approximation;
its physical finite-species and ultraviolet limits are still open.
NS V26 does not turn its rank-one norm conclusions into full NS
regularity.  Parallel YM V5 rejects an extensive total-charge
tightness inference under tensorization and retains the local or
compact topology.  None supplies the missing local physical
Dirichlet-form comparison here.  Directions written inside these
manuscripts are treated as mathematical proposals, not as new
instructions from the user.

\begin{assessment}[The ten-version review]
V41--V43 supply actual fixed-core density control, a common-core
refreshment, and exact boundary ownership of the remaining
information.  V44 rules out a schedule-only closure for the
specified conservative dependency graph; it does not rule out
decay in the actual kernels.  V45--V47 obtain pinned or
defect-localized response estimates, with their geometric and
conditioning hypotheses retained.  V48--V49 control separated
bounded interfaces and remove their guard in the declared
component regime.

The lift (53.18) still resamples all of the intervening bulk.
Its constant \(28\) therefore does not prove a bounded-range
bulk gap.  Even one unrestricted component can have a trivial
component heat-bath gap without resolving any bulk mixing.
The direction chosen at V50 is to charge bad fluctuations to
actual fixed-size conditional updates, and to expose the exact
commutator needed to assemble those charges.
\end{assessment}

There is an important non-duplication point.  Archive
\texttt{V100\_f6}, V89, already proves restricted-smoothing
capacity bounds, and its V90 discusses martingale operator
assembly.  Those abstract ideas are not new results of V50.
The additional input below is the explicit restricted norm
for the present globally completed YM fibre, obtained from
the V31 deterministic shell cover.  A tower-property refinement
also removes the final, comparison-only shell from the update.
The subsequent elementary assembly records precisely what that
populated norm buys and what it leaves unpaid.

The target remains the continuum theory and physical Hamiltonian
gap, in the sense of Jaffe--Witten's
\href{https://www.claymath.org/wp-content/uploads/2022/06/yangmills.pdf}
{official problem description}.  An auxiliary conditional-update
generator is not that Hamiltonian.

\subsection{A conditional bound with the final shell left frozen}

Work with the same fixed smooth completion \(W(Y,C)\), section
\(J(C)\), and canonical finite-volume fibre law \(\mu\).
All quantities use the original physical field; there is no
change of gauge when a region changes.  The coarse field \(C\)
is fixed.  Write
\[
 \Pi_R f=\mu(f\mid Y_{I_R^c}),\qquad
 e_R(f)=\|(I-\Pi_R)f\|_{L^2(\mu)}^2 .
 \tag{54.1}
\]
Thus \(\Pi_R\) integrates exactly the independent free coordinates
\(I_R\), not independently chosen physical pivot links.
For nested cell sets, the corresponding update sets are nested
and
\[
 R\subset R'\quad\Longrightarrow\quad
       \Pi_{R'}\Pi_R=\Pi_R\Pi_{R'}=\Pi_{R'} .
 \tag{54.2}
\]
These identities are conditional-expectation identities for a
single law, not an independence assumption.

Fix \(0<\tau<d_G\) and a free variable \(i\), with owner cell
\(o(i)\).  Use exactly the constants \(h=3,\eta,J,M,N\) of
(35.7)--(35.9) and the thresholds of (35.12).  In particular,
\[
 (1+\eta)^J\tau^2>81d_G^2(2hJ+1)^4,\qquad
 M=(2hJ+1)^4,\qquad
 N=1+\left\lceil\log_2(81d_G^2M/\tau^2)\right\rceil .
 \tag{54.3}
\]
Put
\[
 R_{ij}=N_{hj}(\{o(i)\}),\quad
 P_i^+=R_{iJ},\quad P_i^-=R_{i,J-1},\quad
 B_i=\{\|Y_i-J(C)_i\|_{\rm HS}\ge\tau\},\quad
 p_\tau=e^{-\beta\tau^2/(16A)} .
 \tag{54.4}
\]
For a free link, \(Y_i\) denotes its physical free-link value
in the fixed completion.  Suppose the affected coarse corner
curvatures for \(P_i^+\) are at most \(\epsilon_\tau\), and
\(\beta\ge\beta_\tau\).  All updates \(I_{R_{ij}}\), \(j<J\),
must be available in the ambient law.  For example, an ambient
region containing \(I_{P_i^+}\) suffices.  The result below does
not require that the final layer \(I_{P_i^+}\setminus I_{P_i^-}\)
be resampled.

\begin{theorem}[Populated local restricted smoothing]
\label{thm:f12v50-local}
Under these hypotheses,
\[
             \Pi_{P_i^-}1_{B_i}\le p_\tau
             \quad\mu\text{-almost surely}.
 \tag{54.5}
\]
The bound holds for every prescribed compact exterior of
\(I_{P_i^-}\) in the canonical kernel.  No good-free-exterior
condition is imposed.  The update uses at most \(73M\)
independent group variables and a radius bounded in terms of
\(\tau\) and the fixed block geometry, independently of \(\beta\)
and ambient volume.

If \(M_{B_i}\) denotes multiplication by \(1_{B_i}\), then
\[
 \|M_{B_i}\Pi_{P_i^-}\|_{2\to2}^2
   =\|\Pi_{P_i^-}1_{B_i}\|_\infty
   \le p_\tau .
 \tag{54.6}
\]
\end{theorem}
\begin{proof}
Write \(D_j=D_{R_{ij}}(W)\) and \(E_j=E_{R_{ij}}(W)\)
as in (35.1).  The deterministic growth argument
of Proposition~\ref{prop:f12v31-growth}, with one specified
site, gives the pointwise cover
\[
 B_i\subset
 \bigcup_{\substack{0\le j<J\\0\le\ell<N}} A_{j\ell},
 \qquad
 A_{j\ell}=\{D_j\ge x_\ell,\ E_j\le2\eta x_\ell\},
 \qquad x_\ell=2^\ell\tau^2 .
 \tag{54.7}
\]
The final \(D_J\) enters only through its deterministic compact
upper bound.  That bound is valid even when the final layer is
frozen to an arbitrary configuration.

For each deterministic \(j,\ell\), \(E_j\) is exterior-measurable
for \(\Pi_{R_{ij}}\).  The normalized canonical tail
of Proposition~\ref{prop:f12v31-tail} therefore gives, before
integrating any exterior,
\[
 \Pi_{R_{ij}}1_{A_{j\ell}}
 \le q_\ell:=D_5^M e^{-\beta x_\ell/(8A)} .
 \tag{54.8}
\]
Indeed, outside \(\{E_j\le2\eta x_\ell\}\) the left side
vanishes.  On that event, the collar debit is at most
\(x_\ell/(16A)\), and the coarse-load debit is at most
\(\tau^2/(16A)\le x_\ell/(16A)\), exactly as in (35.14).
This leaves \(x_\ell/(8A)\) from the normalized tail exponent.
The conditional kernel is the original one, not the law
conditioned on which shell was selected.

Since every \(R_{ij}\), \(j<J\), is contained in \(P_i^-\),
(54.2) and positivity give
\[
 \Pi_{P_i^-}1_{A_{j\ell}}
 =\Pi_{P_i^-}\Pi_{R_{ij}}1_{A_{j\ell}}\le q_\ell.
\]
Consequently
\[
 \Pi_{P_i^-}1_{B_i}
 \le J\sum_{\ell=0}^{N-1}q_\ell
 \le JN D_5^M e^{-\beta\tau^2/(8A)}
 \le p_\tau .
 \tag{54.9}
\]
The last step is the same entropy absorption in (35.12).
This proves (54.5), now with the update \(P_i^-\), not \(P_i^+\).
The argument can be carried out in each canonical exterior
kernel, so its assertion is not merely an averaged bound.
The variable count follows from (35.3) and \(|P_i^-|\le M\).

For any conditional expectation \(\Pi\), let
\(b=\Pi1_B\).  Conditional integration shows
\[
 \|M_B\Pi f\|_2^2=\mu[b|\Pi f|^2]
               \le\|b\|_\infty\|f\|_2^2 .
\]
Conversely, take \(f\) to be the normalized indicator of an
exterior-measurable set where \(b>\|b\|_\infty-\delta\).
Whenever the essential supremum is positive such a set has
positive measure; then \(\Pi f=f\).
Letting \(\delta\downarrow0\) proves equality of the norms.
The zero case is immediate.  Apply this with (54.5).
\end{proof}

\begin{assessment}[Scope of the buffer refinement]
V31 integrated \(P_i^+\).  Its proof actually supplies uniform
conditional bounds on deterministic events under smaller nested
kernels, which is why the outer comparison layer can be frozen.
This does not turn the result into V26's unbuffered one-site
activity estimate.  The remaining radius \(h(J-1)\) is fixed
but potentially enormous.  The small coarse-curvature assumption
on the affected region is also retained.
\end{assessment}

\subsection{The resulting local form bound}

\begin{corollary}[Canonical heat-bath capacity pair]
\label{cor:f12v50-capacity}
For every \(f\in L^2(\mu)\) and \(\theta>0\),
\[
 \mu[1_{B_i}|f|^2]
 \le(1+\theta)e_{P_i^-}(f)
       +(1+\theta^{-1})p_\tau\|f\|_2^2 .
 \tag{54.10}
\]
In particular, for centered \(f\) its coefficients are
\(2p_\tau\) on the variance and \(2\) on the canonical
local heat-bath form.
The inequality also holds after multiplication by any bounded
nonnegative \(Y_{I_{P_i^-}^c}\)-measurable weight, with
\(e_{P_i^-}(f)\) replaced by
\(\mu[w|f-\Pi_{P_i^-}f|^2]\) and
\(\|f\|_2^2\) by \(\mu[w|f|^2]\).
\end{corollary}
\begin{proof}
Set \(\Pi=\Pi_{P_i^-}\).  In
\(f=(f-\Pi f)+\Pi f\), use
\[
 |u+v|^2\le(1+\theta)|u|^2
                    +(1+\theta^{-1})|v|^2 .
\]
The first square restricted to \(B_i\) is at most
\(e_{P_i^-}(f)\).  Conditional integration bounds the second by
\[
 \mu[1_{B_i}|\Pi f|^2]
 =\mu[(\Pi1_{B_i})|\Pi f|^2]
 \le p_\tau\mu|\Pi f|^2\le p_\tau\mu|f|^2 .
\]
An exterior-measurable weight commutes with the conditional
integration, and conditional Jensen gives its weighted version.
\end{proof}

This is the projection specialization of the archived
restricted-smoothing mechanism, now with its norm and its
actual finite support supplied.  Orthogonality of \(\Pi\)
avoids the factor two lost when bounding a general Markov
increment by its full jump form.  No identification with the
physical transfer semigroup is asserted.

Simply summing (54.10) with the same centered \(f\) at every
site would introduce \((\#\hbox{sites})p_\tau\).
The next statement instead pairs each event with an orthogonal
fluctuation.  It does not assume that the repair updates commute
with those fluctuations.

\subsection{Orthogonal assembly with all commutators retained}

Let \((D_a)_{a\in\mathcal A}\) be a finite family of mutually
orthogonal orthogonal projections on \(L^2(\mu)\), with
\(D_a1=0\).  Examples are the differences of nested canonical
projections in V36.  To each \(a\), assign a site \(i(a)\)
satisfying the hypotheses above, and write
\[
 B_a=B_{i(a)},\quad P_a=P_{i(a)}^-,
 \quad \mathcal U=\{P_a:a\in\mathcal A\},\quad
 \mathcal E_{\rm rep}(f)=\sum_{P\in\mathcal U}e_P(f).
 \tag{54.11}
\]
The last sum counts each distinct repair region once.
Define the exact bad fluctuation debit and commutator debit by
\[
 \mathcal R_{\rm bad}(f)=
       \sum_a\mu[1_{B_a}|D_af|^2],\qquad
 \mathcal C_{\rm rep}(f)=
       \sum_a\|[D_a,\Pi_{P_a}]f\|_2^2 ,
 \qquad [D,\Pi]=D\Pi-\Pi D .
 \tag{54.12}
\]

\begin{theorem}[Volume-free small coefficient; explicit energy debt]
\label{thm:f12v50-assembly}
For every \(f\in L^2(\mu)\),
\[
 \boxed{\quad
 \mathcal R_{\rm bad}(f)
 \le 4\mathcal E_{\rm rep}(f)
       +4\mathcal C_{\rm rep}(f)
       +2p_\tau\operatorname{Var}_\mu(f).
 \quad}
 \tag{54.13}
\]
There is no factor \(|\mathcal A|\) in the last term or a
repetition multiplicity in the first.  The commutator term,
however, is an actual sum and is not bounded by this theorem
in terms of \(\mathcal E_{\rm rep}\).
\end{theorem}
\begin{proof}
Apply (54.10) with \(\theta=1\) to each \(D_af\).
Orthogonality and \(D_a1=0\) imply
\[
 \sum_a\|D_af\|_2^2
 \le\|f-\mu f\|_2^2 .
 \tag{54.14}
\]
For the remaining energy use the exact identity
\[
 (I-\Pi_{P_a})D_af
   =D_a(I-\Pi_{P_a})f+[D_a,\Pi_{P_a}]f .
 \tag{54.15}
\]
Its squared norm is at most twice the sum of the two
squared norms on the right.  Group the first terms by their
identical repair region \(P\).  Orthogonality again gives
\[
 \sum_{\{a:P_a=P\}}\|D_a(I-\Pi_P)f\|_2^2
       \le\|(I-\Pi_P)f\|_2^2=e_P(f).
 \tag{54.16}
\]
Combining these three estimates gives (54.13).
No conditional independence, spatial decorrelation, or
commutativity was used.
\end{proof}

\begin{corollary}[Conditional absorption ledger]
\label{cor:f12v50-absorb}
Suppose a separately proved variance decomposition supplies
nonnegative forms \(\mathcal E_{\rm good},\mathcal E_{\rm loc}\)
and constants \(a,b,c,d\ge0\) such that
\[
 \operatorname{Var}_\mu(f)
       \le a\mathcal E_{\rm good}(f)+b\mathcal R_{\rm bad}(f),
 \quad
 \mathcal E_{\rm rep}(f)\le c\mathcal E_{\rm loc}(f),
 \quad
 \mathcal C_{\rm rep}(f)\le d\mathcal E_{\rm loc}(f).
 \tag{54.17}
\]
If \(2bp_\tau<1\), then
\[
 \operatorname{Var}_\mu(f)
 \le
 \frac{a\mathcal E_{\rm good}(f)
          +4b(c+d)\mathcal E_{\rm loc}(f)}
      {1-2bp_\tau}.
 \tag{54.18}
\]
\end{corollary}
\begin{proof}
Substitute (54.13) into the first inequality of (54.17),
use the other two, and move
\(2bp_\tau\operatorname{Var}_\mu(f)\) to the left.
\end{proof}

This last algebra is not a new abstract gap criterion.  It is
included to identify the still-unpopulated rows for the actual
fibre problem.  In particular, neither the first nor the third
inequality in (54.17) has been established here with
volume-uniform constants for a suitable full-volume martingale.
If \(\mathcal E_{\rm loc}=\mathcal E_{\rm rep}\), its comparison
row is the identity \(c=1\), but this convention does nothing
to establish the commutator row.  Replacing the latter by
\(d=0\) requires a proof of commutation.  Projections for
interacting, nonnested conditional updates generally fail to
commute, including in the exact finite test below.

\subsection{Verification, frontier and append record}

The script
\begin{center}\small\ttfamily
scripts/verify\_ym\_f12\_v50\_local\_bad\_form.py
\end{center}
uses exact rational arithmetic on a strictly positive,
interacting four-bit law.  It checks the restricted-projection
identity, twelve full matrix capacity inequalities with three
choices of the square-inequality parameter, and forty-eight
exterior-weighted versions.  It also checks (54.13) for four
orthogonal martingale differences assigned to three distinct
repair blocks, with two nonzero commutators.  Grouping repeated
repair regions is verified by exact positive-semidefinite
matrix elimination.

As a negative control, a rare event of probability \(1/100\)
is frozen outside a one-coordinate update.  Its indicator
has zero update energy, and replacing its conditional
probability bound by stationary rarity would incorrectly
assert \(1/100\le1/5000\).  All positive checks and this
negative control passed.  The script tests the algebra, not
the analytic YM tail constants, and is not a proof checker
for the continuum claim.

\begin{decision}[Direction after the review]
The new usable result is the small restricted norm (54.6)
for a specified bounded-range canonical update, including
the final-shell refinement.  It produces an actual local
form-capacity pair and the volume-free small coefficient
in (54.13).  The next work must select the variance
decomposition and estimate its spatial commutator sum.
Nested cases can commute exactly; overlapping boundary
cases require a proved estimate.  Iterating a marginal
rarity statement or the large-bulk lift does not supply it.
The next companion audit is V55 and the next ten-version
direction review is V60.
\end{decision}

V49 had 769658 bytes and SHA--256
\begin{center}\scriptsize\ttfamily
5F074D3B33AE410E095065CB575D14B93F4D747E34749C1C9831531EEAC3DC59.
\end{center}
Its complete body is retained apart from the title version.
Only active V49 is replaced by V50.  All archives, companions,
and prior scripts are preserved.

\begin{firewall}[Position after V50]
The estimates concern the auxiliary canonical fibre heat baths
at fixed small coarse curvature and sufficiently large
\(\beta\).  The local restricted norm and form bounds are
proved under the stated geometric hypotheses.  The global
commutator estimate and variance decomposition remain open.
A volume-uniform bounded-range fibre gap, its physical
transfer-energy comparison, scale-stable effective action,
nontrivial continuum construction, and the physical
Yang--Mills mass gap have not been proved.
\end{firewall}

\section{V51: spatial localization of ordered Gaussian innovations
and the local commutator benchmark}
\label{sec:f12v51}

\subsection{The question selected from V50}

The preceding turn made progress: V50 supplied a bounded-range
restricted norm for the actual compact fibre and exposed the
commutator debit in (54.13).  Its source was rechecked before
this continuation.  The present section addresses the quantitative
mechanism that could control that debit, without declaring
nonnested conditional expectations to commute.

The new calculation concerns the massive quadratic reference
of the current independent fibre.  It proves an exponentially
small, volume-uniform \emph{ordered innovation leakage row}
for every coordinate ordering.  This is more specific than
decay of a translation-invariant Fourier multiplier: the
eliminated prefixes can be arbitrary sets, and the local
conditional projections depend on their actual complements.
It also gives a fully populated all-\(L^2\) commutator bound
for this Gaussian reference.  The nonlinear compact analogue
remains an additional task.

For non-duplication, archive \texttt{V100\_f7}, V5--V6,
already contains abstract martingale Schur bounds and
Gaussian infrared tests.  Its V16 proves the
first-chaos-to-all-chaoses heat-bath gap principle.
Its V21 distinguishes a fixed local frame from a cofinal
massless curvature frame.  We use these interfaces rather
than count them as newly proved.  The new input here is the
localized regression matrix for a specified, ordered
coordinate martingale of the present massive fibre.
In particular, no new proof of the physical mass gap is
being inferred from a Gaussian gap.

\subsection{Quadratic carrier and the actual ordering}

Let the scalar Lie-coordinate index set be \(\mathcal I\),
with the owner-cell pseudometric \(d(i,j)\).  Distinct scalar
coordinates may have the same owner.  The number per cell is
fixed.  Thus there is a volume-independent \(g\) such that
\[
 \#\{j:d(i,j)\le r\}\le g(1+r)^4,\qquad r\ge0.
 \tag{55.1}
\]
Periodic identifications or restriction to a finite region
do not increase this bound.  Let \(H\) be a real symmetric
precision matrix satisfying
\[
 mI\le H\le KI,\qquad
 H_{ij}=0\quad\hbox{if }d(i,j)>r_0,\qquad
 0<m\le K,\quad r_0\ge1.
 \tag{55.2}
\]
The constants \(m,K,r_0,g\) are fixed.  Consider the centered
Gaussian law with density proportional to
\(\exp(-\beta x^THx/2)\).  After \(z=\sqrt\beta x\), write
\[
                  \nu=N(0,\Sigma),\qquad \Sigma=H^{-1}.
 \tag{55.3}
\]
All estimates below are stated in \(z\)-coordinates.
The scaling does not change conditional-expectation
operator norms or rate-one heat-bath gaps.

Choose \emph{any} ordering \(1,\ldots,n\) of \(\mathcal I\).
Put \(A_i=\{1,\ldots,i\}\), \(A_0=\varnothing\), and
\[
 \Gamma_i f=\nu(f\mid z_{A_i^c}),\qquad
 D_i=\Gamma_{i-1}-\Gamma_i .
 \tag{55.4}
\]
Then \(\Gamma_0=I\), \(\Gamma_n=\nu\), and the \(D_i\)
are mutually orthogonal projections with
\(\sum_iD_i=I-\nu\), on all of \(L^2(\nu)\).
These are differences of the actual nested conditional
expectations, not independent-coordinate differences
silently substituted for them.

For each \(i\), choose a deterministic update set \(P_i\)
containing every scalar coordinate whose owner is within
distance \(R\) of \(i\).  It can be the actual completed
regional free-coordinate update \(I_{N_R(\{o(i)\})}\).
For bounded-range applications also require
\(P_i\subset\{j:d(i,j)\le R+r_1\}\), with fixed \(r_1\).
Let
\[
 \Pi_i f=\nu(f\mid z_{P_i^c}),\qquad
 \mathcal U=\{P_i:1\le i\le n\},\qquad
 \mathcal E_R^\nu(f)=
          \sum_{P\in\mathcal U}\|(I-\Pi_P)f\|_2^2.
 \tag{55.5}
\]
Identical repair sets are counted once.  Write
\(\mathcal H_1=\{v^Tz:v\in\mathbb R^n\}\) for the first
Gaussian chaos, the space of centered linear functions.

\subsection{Principal inverses and localized regressions}

The following inverse estimate is uniform in an arbitrary
principal index set \(A\):
\[
 |(H_{AA}^{-1})_{ij}|
       \le m^{-1}e^{-\kappa d(i,j)},\qquad
 \kappa=\frac{-\log(1-m/(2K))}{r_0}>0 .
 \tag{55.6}
\]
To see this directly, expand
\[
 H_{AA}^{-1}
  =\frac1{2K}\sum_{q\ge0}
                     (I-H_{AA}/(2K))^q .
\]
The \(q\)-th matrix has range at most \(qr_0\) and norm
at most \((1-m/(2K))^q\).  Sum from
\(q=\lceil d(i,j)/r_0\rceil\) onwards.  The same bound
applies to \(\Sigma\).  No connectedness of \(A\) is used.

For an arbitrary partition \(P,O\), the precision of the
exterior marginal is
\[
 S^P=(\Sigma_{OO})^{-1}
       =H_{OO}-H_{OP}H_{PP}^{-1}H_{PO}.
 \tag{55.7}
\]
Empty blocks have their natural vacuous interpretation.
Finite range of the two outside factors, (55.1), and
(55.6) imply
\[
                     |S^P_{jk}|\le C_S e^{-\kappa d(j,k)}
                     \quad(j,k\in O).
 \tag{55.8}
\]
Indeed, the summation indices in the product term are
within \(r_0\) of \(j\) and \(k\), respectively, and there
are at most \(g(1+r_0)^4\) choices for each.  Each entry
of \(H\) has absolute value at most \(K\).
The resulting bound is
\(g^2(1+r_0)^8K^2m^{-1}
e^{2\kappa r_0}e^{-\kappa d(j,k)}\), plus the
finite-range \(H_{OO}\) term.  Thus \(C_S\) is independent
of the shape and volume of \(P\).

\begin{proposition}[Localized ordered innovations]
\label{prop:f12v51-innovation}
Define
\[
 \sigma_i^2=(H_{A_iA_i}^{-1})_{ii},\qquad
 u_i=\frac{z_i-\nu(z_i\mid z_{A_i^c})}{\sigma_i}
                   =\sum_j a_{ij}z_j .
 \tag{55.9}
\]
Then \((u_i)_{i=1}^n\) is an orthonormal basis of
\(\mathcal H_1\), and
\[
 D_if=\langle u_i,f\rangle u_i
       \quad(f\in\mathcal H_1),\qquad
 |a_{ij}|\le C_a e^{-\kappa d(i,j)} .
 \tag{55.10}
\]
The coefficient bound is uniform in the ordering.
If
\[
                    \Pi_i u_i=\sum_j b_{ij}z_j,
 \tag{55.11}
\]
then there are fixed \(C_b,\lambda>0\), with
\(\lambda=\kappa/4\) allowed, such that
\[
 b_{ij}=0\ (j\in P_i),\qquad
             |b_{ij}|\le C_b e^{-\lambda d(i,j)}
                       1_{\{d(i,j)>R\}} .
 \tag{55.12}
\]
\end{proposition}
\begin{proof}
Conditional on \(z_{A_i^c}\), the whole \(A_i\)-vector
has covariance \(H_{A_iA_i}^{-1}\) and mean
\(-H_{A_iA_i}^{-1}H_{A_i,A_i^c}z_{A_i^c}\).
Hence the coefficient at \(i\) in (55.9) is
\(\sigma_i^{-1}\), the coefficients at \(j<i\) vanish,
and at \(j>i\) they are
\[
 a_{ij}=\sigma_i^{-1}
                (H_{A_iA_i}^{-1}H_{A_i,A_i^c})_{ij}.
\]
Since \(K^{-1}\le\sigma_i^2\le m^{-1}\), (55.6)
and the finite range of the right factor prove (55.10)'s
coefficient estimate.  The residual \(u_i\) is independent
of \(z_{A_i^c}\) and has variance one.  For \(j>i\),
\(u_j\) is measurable from those exterior variables.
This proves orthogonality.  Triangularity proves that
the \(n\) residuals span \(\mathcal H_1\).
The range of \(\Gamma_{i-1}-\Gamma_i\) on the first
chaos is precisely this residual direction.

Let \(c_{ij}=\operatorname{Cov}_\nu(u_i,z_j)\), so that
the row \(c_i\) equals \(a_i\Sigma\).
Convolution of the two exponential bounds gives
\[
                  |c_{ij}|\le C_c e^{-\kappa d(i,j)/2}.
 \tag{55.13}
\]
For completeness, if two factors are bounded by
constants times \(e^{-t d(i,k)}\) and
\(e^{-t d(k,j)}\), their product is at most
\[
 e^{-t d(i,j)/2}
 e^{-t d(i,k)/2}e^{-t d(k,j)/2}.
\]
The sum over \(k\) of the last two factors is uniformly
bounded by (55.1).  This justifies each convolution
used here, including arbitrary principal subsets.

Gaussian regression onto \(z_{O_i}\), \(O_i=P_i^c\),
now gives the exact row
\[
              b_{i,O_i}=c_{i,O_i}S^{P_i},\qquad
              b_{i,P_i}=0.
 \tag{55.14}
\]
Use (55.8), (55.13), and the same convolution argument
with \(t=\kappa/2\).  This gives exponential rate
\(\kappa/4\).  Finally \(j\in O_i\) implies
\(d(i,j)>R\) by the definition of \(P_i\).
This proves (55.12).
\end{proof}

The innovation is generally not supported in a fixed
neighbourhood.  The principal inverse in its coefficient
formula has entries at arbitrarily large distances.
The proof obtains decay of those entries, not a false
finite-support assertion.  No low-dimensional boundary
ordering or bounded bandwidth has been assumed.

\subsection{A populated, volume-free leakage row}

\begin{theorem}[Ordered first-chaos leakage with a spatial cutoff]
\label{thm:f12v51-row}
There is \(C_0<\infty\), depending only on
\(m,K,r_0,g\), such that, with
\[
 s_R=C_0(1+R)^4e^{-\lambda R},\qquad
                    \varepsilon_R=s_R^2/m ,
 \tag{55.15}
\]
the following hold uniformly in volume, ordering, and
the declared repair sets:
\[
 \sum_i\|D_i\Pi_i f\|_2^2
                  \le\varepsilon_R\|f\|_2^2
                  \quad(f\in\mathcal H_1),\qquad
 \sup_i\|\Pi_i u_i\|_2^2\le\varepsilon_R .
 \tag{55.16}
\]
In particular \(\varepsilon_R\to0\) as \(R\to\infty\),
with \(R\) independent of volume and of \(\beta\).
\end{theorem}
\begin{proof}
Let \(B=(b_{ij})\).  Equation (55.12) and the polynomial
ball count imply
\[
             \sup_i\sum_j|b_{ij}|\le s_R,\qquad
             \sup_j\sum_i|b_{ij}|\le s_R .
 \tag{55.17}
\]
For example, the tail is bounded by a constant times
\(\sum_{r>R}(1+r)^4e^{-\lambda r}
\le C(1+R)^4e^{-\lambda R}\).
The same count works for columns since it is uniform
over the location of the owner.  This is an actual
row and column estimate, not the sum of \(n\) individual
operator-norm bounds.  Schur's inequality gives
\(\|B\|_{\ell^2\to\ell^2}\le s_R\).

For \(f=v^Tz\), self-adjointness of \(\Pi_i\) gives
\[
 \langle u_i,\Pi_i f\rangle
       =\langle\Pi_i u_i,f\rangle=(B\Sigma v)_i.
\]
Therefore the left side of (55.16) is
\(\|B\Sigma v\|_{\ell^2}^2\).
Because \(\Sigma\le m^{-1}I\),
\[
 \|B\Sigma v\|_{\ell^2}^2
       \le s_R^2v^T\Sigma^2v
       \le(s_R^2/m)v^T\Sigma v .
\]
This proves the first bound.  The second follows from
\(b_i\Sigma b_i^T\le m^{-1}\|b_i\|_{\ell^2}^2
\le m^{-1}s_R^2\).
\end{proof}

\subsection{What this proves for the Gaussian update generator}

Assume \(R\) is large enough that \(\varepsilon_R<1\), and put
\[
                       \gamma_R=(1-\sqrt{\varepsilon_R})^2 .
 \tag{55.18}
\]
For \(f\in\mathcal H_1\), decompose each martingale component as
\[
                 D_if=D_i(I-\Pi_i)f+D_i\Pi_i f.
\]
Taking the norm in the direct sum over \(i\), grouping the
first terms by identical \(P_i\) as in (54.16), and using
(55.16) gives
\[
 \|f\|_2
 \le\sqrt{\mathcal E_R^\nu(f)}
                        +\sqrt{\varepsilon_R}\|f\|_2 .
 \tag{55.19}
\]
Thus the first-chaos form has lower bound
\(\mathcal E_R^\nu(f)\ge\gamma_R\|f\|_2^2\).
This argument uses the independently proved regression row,
not a gap assumption to establish its own leakage hypothesis.

The all-chaos principle already proved in f7 V16 now yields
\[
                  \mathcal E_R^\nu(f)
                   \ge\gamma_R\operatorname{Var}_\nu(f)
                   \quad(f\in L^2(\nu)).
 \tag{55.20}
\]
Its applicability here is exact: a Gaussian conditional
expectation is second quantization of its first-chaos
exterior projection \(Q_P\), so on degree \(k\ge1\) it
is \(Q_P^{\otimes k}\) restricted to symmetric tensors.
The archived tensor projection inequality applies to
the sum over the \emph{distinct} \(P\)'s in (55.5).
No new Gaussian gap principle or nonlinear chaos
decomposition is being postulated.

\begin{proposition}[Populated Gaussian commutator row]
\label{prop:f12v51-commutator}
For the full \(L^2(\nu)\) martingale of (55.4),
\[
 \sum_i\|[D_i,\Pi_i]f\|_2^2
             \le (2+2/\gamma_R)\mathcal E_R^\nu(f).
 \tag{55.21}
\]
On the first chaos there is also the sharper estimate
\[
 \sum_i\|[D_i,\Pi_i]f\|_2^2
 \le4\varepsilon_R\|f\|_2^2
 \le(4\varepsilon_R/\gamma_R)\mathcal E_R^\nu(f).
 \tag{55.22}
\]
The exponentially small coefficient in (55.22) is
asserted only on the first chaos.
\end{proposition}
\begin{proof}
For arbitrary orthogonal \(D_i\)'s annihilating constants,
\[
 [D_i,\Pi_i]f
   =(I-\Pi_i)D_if-D_i(I-\Pi_i)f .
\]
The square inequality, contraction of \(I-\Pi_i\),
and grouping by \(P_i\) give
\[
 \sum_i\|[D_i,\Pi_i]f\|_2^2
 \le2\sum_i\|D_if\|_2^2+
           2\sum_i\|D_i(I-\Pi_i)f\|_2^2
 \le2\operatorname{Var}_\nu(f)+2\mathcal E_R^\nu(f).
 \tag{55.23}
\]
Insert (55.20) to prove (55.21).  In this Gaussian
calculation that substitution is not circular:
(55.20) was obtained from the explicit matrix row.
For the original compact law, assuming a comparable
gap in order to justify this same substitution would
not solve the V50 problem.

On \(\mathcal H_1\), put \(v_i=\Pi_i u_i\).
The rank-one formula gives
\[
 [D_i,\Pi_i]f
   =u_i\langle v_i,f\rangle
                       -v_i\langle u_i,f\rangle .
\]
Its squared norm is at most
\(2|\langle v_i,f\rangle|^2+
2\|v_i\|_2^2|\langle u_i,f\rangle|^2\).
Sum, use both parts of (55.16), and use orthonormality
of the \(u_i\)'s.  This proves (55.22).
\end{proof}

One must not replace a full-chaos martingale difference
by the tensor power of a first-chaos difference.
The correct operator at degree \(k\) is
\[
 (\Gamma_{i-1}|_{\mathcal H_1})^{\otimes k}
             -(\Gamma_i|_{\mathcal H_1})^{\otimes k},
 \tag{55.24}
\]
not
\((\Gamma_{i-1}|_{\mathcal H_1}
       -\Gamma_i|_{\mathcal H_1})^{\otimes k}\).
The latter misses mixed innovations.  Likewise the
all-chaos gap principle does not, by itself, lift the
small leakage coefficient in (55.16) to all \(L^2\).

\subsection{Application to this programme and the nonlinear handoff}

In the good-collar and small-coarse-curvature regime
of Proposition~\ref{prop:f12v32-convex}, take
\[
                    H=\beta^{-1}V_R''(0).
 \tag{55.25}
\]
Then (36.8) supplies (55.2) with \(m=m_0\), the
same fixed upper bound \(K\), and the actual completed
interaction range.  Scalar coordinates have bounded
owner multiplicity, so (55.1) holds.  Consequently
(55.16) and (55.21) are populated for the quadratic
reference of the current independent YM fibre, not
merely for a translation-invariant surrogate.
Principal restrictions and arbitrary coordinate
orderings remain allowed.

This reference Gaussian is defined on the full
Euclidean tangent space.  It is not the compact
nonlinear law or the Gaussian restricted to a chart.
The statement makes no uniform comparison of those
laws.  The independent fibre is vertically massive
with the coarse field held fixed; (55.2) does not
apply to a massless retained physical curvature
sector.  Thus there is no conflict with f7 V21's
cofinal infrared obstruction.

The calculation identifies a direct target for the
original compact law.  For this handoff, order whole
independent group variables, not hypothetical global
Lie coordinates off the logarithm chart.  The Gaussian
comparison can use the corresponding ordering with
each group's scalar coordinates consecutive.
Grouping those consecutive scalar martingale differences,
and assigning them the same repair set, preserves
(55.16)'s first inequality by orthogonality.
Thus (55.20) and (55.21) hold for that groupwise
filtration too.  No global logarithm on the compact
group is required or asserted.
Using the compact law's own groupwise ordered
projections and its actual canonical repair updates,
define
\[
 \mathcal L_R^\mu(f)=
           \sum_i\|D_i^\mu\Pi_i^\mu f\|_{L^2(\mu)}^2 .
 \tag{55.26}
\]
If one independently proves
\[
 \mathcal L_R^\mu(f)
        \le\varepsilon\operatorname{Var}_\mu(f),
             \qquad \varepsilon<1,
 \tag{55.27}
\]
then the direct-sum argument (55.19) gives a
bounded-range auxiliary gap.  This is an exact
analytic target, not a populated conclusion:
the rank-one regressions (55.9)--(55.14) and the
Gaussian chaos lift do not hold for that law.
V39 already handles a wholly boxed law; re-proving
its Gaussian limit would not remove the compact
bad sectors.  Their contribution to (55.26)
must be controlled with a correctly assigned local
capacity or a nonlinear conditional-response estimate.

\begin{decision}[Upstream direction after V51]
The Gaussian computation validates the chosen
spatially buffered, ordered projection geometry:
its leakage is exponentially localized even though
each eliminated prefix can be large and disconnected.
The next step is a nonlinear formula for the same
ordered leakage, retaining both normalization
scores and bad-sector terms.  The V50 capacity
estimate can be used only after identifying the
functions and events to which it is applied.
Neither a formal commutator cancellation nor a
global Gaussian density replacement supplies that step.
\end{decision}

\subsection{Exact tests and append record}

The script
\begin{center}\small\ttfamily
scripts/verify\_ym\_f12\_v51\_innovation\_rows.py
\end{center}
uses rational precision matrices, regressions,
innovations and conditional projections.
It verifies ten cases: an eight-coordinate path
at three repair radii, a sixteen-coordinate
four-dimensional cube at three radii and two
orderings, and one case where every repair is
the same full block.  The second ordering sorts
by coordinate Hamming weight, rather than simply
permuting the geometric axes.

Exact Schur row/column certificates are checked
as covariance-metric positive-semidefinite
inequalities.  When their coefficient is below
\(1/2\), the rational gap lower bound
\((1-2\varepsilon)/2\) is also verified.  This is
a weaker, rational consequence of (55.18).
Three-coordinate examples check the tensor gap
and full commutator bound at degrees one, two
and three, with tensor dimensions \(3,9,27\).
They also detect the incorrect tensorization
of martingale differences discussed above.
These tests verify finite Gaussian algebra,
not the compact nonlinear or continuum estimates.

V50 had 785973 bytes and SHA--256
\begin{center}\scriptsize\ttfamily
E855540C49059E5347B3133FA5F4E2101867F10AF4B1691500AFD540C48CA37B.
\end{center}
Its complete body is retained apart from the title
version.  Only active V50 is replaced by V51.
All archives, companions and prior scripts remain.
The next scheduled companion audit is V55;
the next direction review is V60.

\begin{firewall}[Position after V51]
A spatially localized ordered innovation row and
an all-\(L^2\) local commutator bound are proved
for the massive quadratic fibre reference.
The exponentially small leakage bound is proved
on its first chaos, not on unrestricted nonlinear
fluctuations.  The compact-law leakage estimate
(55.27), a full bounded-range fibre gap, physical
transfer comparison, scale-stable continuum
construction and the Yang--Mills physical mass
gap remain unproved.
\end{firewall}

\section{V52: the exact compact conditional channel and
spectator-stable leakage assembly}
\label{sec:f12v52}

\subsection{Continuation from the verified V51 state}

V51 made concrete progress by evaluating ordered innovation
leakage for the massive quadratic reference.  Its active
source and hash were checked before this continuation.
The task now returns to the original compact law, rather
than treating a Gaussian first-chaos estimate as a nonlinear
operator estimate.

This section derives the exact conditional channel underlying
\(D_i^\mu\Pi_i^\mu\), including the normalization score.
The actual bounded local action gives a volume-independent
conditional contraction, but not a small spatial leakage row.
A separate calculation shows why summing spatially anchored
operator norms would be too strong for a general nonlinear
martingale: a distant independent spectator can multiply
a local fluctuation without changing its local response.
The correctly summed leakage form, by contrast, is stable
under such independent extensions.

This is not a reproof of f7 V5's abstract Schur compiler or
its full-space spectator warning.  The additional statements
concern the particular ordered operator \(D_i\Pi_iD_j\),
the exact compact disintegration producing it, and tensor
stability of the whole leakage quadratic form.  These
distinctions change the next estimate to be attempted.

\subsection{Disintegration at one ordered group variable}

Let \(\mu\) be the positive finite-volume compact fibre law
of the fixed smooth completion.  Its independent coordinates
are whole group variables.  Fix their order \(1,\ldots,n\),
and use
\[
 A_i=\{1,\ldots,i\},\quad
 \Gamma_i f=\mu(f\mid Y_{A_i^c}),\quad
 D_i=\Gamma_{i-1}-\Gamma_i,\quad
 \Pi_i f=\mu(f\mid Y_{P_i^c}),\qquad i\in P_i .
 \tag{56.1}
\]
The repair \(P_i\) is deterministic and canonical.
No chart restriction or selected good event changes its
conditional density.

At this \(i\), name the variables
\[
 X=Y_i,\quad T=Y_{A_i^c},\quad
 L=Y_{A_{i-1}\setminus P_i},\quad
 M=Y_{A_{i-1}\cap P_i},\quad
 O=Y_{A_i^c\setminus P_i}.
 \tag{56.2}
\]
Thus \(L\) is the earlier part outside the repair, \(M\)
the earlier part inside it, and \(O\) is a subrecord of
the fixed future \(T\).  Since the current variable lies
inside the repair,
\[
                  \Pi_i f=h_i(L,O).
 \tag{56.3}
\]
The function \(h_i\) depends on the original normalized
repair kernel; it is not the original \(f\) restricted
to its outside arguments.

For each future record \(t\), integrate \(M\) from the
conditional law of \((X,L,M)\) given \(T=t\).
Write the resulting joint law as
\[
 \nu_t(dx,d\ell)
       =r_t(x,\ell)\,\alpha_t(dx)\lambda_t(d\ell),
 \tag{56.4}
\]
where \(\alpha_t,\lambda_t\) are its true marginals.
Positivity of the original compact density gives positive
versions of these densities at every finite regulator.
Define the centered conditional channel
\[
 (K_tu)(x)
    =\int u(\ell)\,[r_t(x,\ell)-1]\,\lambda_t(d\ell).
 \tag{56.5}
\]
It kills constants and has mean-zero output in \(L^2(\alpha_t)\).
Empty \(L,M\), or \(T\) are interpreted as singleton
probability spaces where appropriate.

\begin{proposition}[Exact nonlinear ordered leakage]
\label{prop:f12v52-channel}
For bounded \(f\), and then by \(L^2\) extension,
\[
 (D_i\Pi_i f)(x,t)
     =K_t\big[h_i(\,\cdot\,,t_O)\big](x),
 \tag{56.6}
\]
and
\[
 \|D_i\Pi_i f\|_2^2
 =\int\big\|K_t[h_i(\,\cdot\,,t_O)]\big\|_{L^2(\alpha_t)}^2
                                                  \,\mu_T(dt).
 \tag{56.7}
\]
In particular, if
\(\rho_i(t)=\|K_t\|_{L^2_0(\lambda_t)\to L^2_0(\alpha_t)}\),
then
\[
 \|D_i\Pi_i f\|_2^2
 \le\int\rho_i(t)^2
       \operatorname{Var}_{\lambda_t}
                    (h_i(\,\cdot\,,t_O))\,\mu_T(dt).
 \tag{56.8}
\]
\end{proposition}
\begin{proof}
The first expectation in
\((\Gamma_{i-1}-\Gamma_i)h_i(L,O)\) conditions on
\((X,T)\); the second conditions on \(T\).
At fixed \(t\), they are respectively
\(\int h_i(\ell,t_O)r_t(x,\ell)\lambda_t(d\ell)\)
and \(\int h_i(\ell,t_O)\lambda_t(d\ell)\).
Their difference proves (56.6).
Disintegrate the squared norm to obtain (56.7);
subtract the \(\lambda_t\)-mean of the input to get (56.8).
Conditional expectations are \(L^2\) contractions,
so bounded approximation extends the operator identities.
\end{proof}

The future \(T\) is frozen in this formula, including the
part of the future that was integrated when constructing
\(h_i=\Pi_i f\).  The two operations occur in that order.
Replacing them by a jointly averaged future would give
a different operator.

\subsection{Effective score and its indispensable centering}

Write the original compact density, relative to product
Haar measure, as a positive constant times
\(e^{-\mathcal V(Y)}\), with the fixed exterior and coarse
field suppressed in the notation.  Define
\[
 W_t(x,\ell)
       =-\log\int e^{-\mathcal V(x,\ell,m,t)}\,dm .
 \tag{56.9}
\]
For a unit invariant differentiation \(v\) in the
current group variable, put
\[
 s_{t,v}(x,\ell)
   =\partial_v W_t(x,\ell)
   =\mu(\partial_v\mathcal V\mid X=x,L=\ell,T=t).
 \tag{56.10}
\]
All group differentiations here are local tangent
differentiations; no global logarithm is introduced.

\begin{proposition}[Normalized conditional score]
\label{prop:f12v52-score}
The relative density in (56.4) satisfies
\[
 \partial_v\log r_t(x,\ell)
  =-s_{t,v}(x,\ell)
       +\int s_{t,v}(x,\ell')
                    \,\mu(d\ell'\mid X=x,T=t).
 \tag{56.11}
\]
For bounded input \(u=u(\ell)\), independent of \(x\),
\[
 \partial_v(K_tu)(x)
 =-\operatorname{Cov}_{\mu(\,\cdot\,\mid X=x,T=t)}
                    (u(L),\partial_v\mathcal V).
 \tag{56.12}
\]
The covariance on the right includes the integrated
earlier variables \(L,M\).
\end{proposition}
\begin{proof}
The joint density of \(X,L\) given \(t\) is proportional
to \(e^{-W_t(x,\ell)}\).  Its logarithmic derivative
in \(x\) is \(-s_{t,v}\).  The marginal density
\(\alpha_t(x)\) has logarithmic derivative
\(-\mu(s_{t,v}\mid X=x,T=t)\), whereas \(\lambda_t\)
has no \(x\)-argument.  Differentiating their ratio
proves (56.11).  Differentiating the normalized
expectation of \(u(L)\) given \((x,t)\) gives (56.12);
the subtracted \(T\)-conditional mean in (56.5) is
independent of \(x\).
Positivity and compactness justify differentiation
under the finite integrals for the smooth completed
density and bounded \(u\).
\end{proof}

Although \(\partial_v\mathcal V\) has a fixed local
input stencil and magnitude at most \(C(\beta+1)\),
the effective score (56.10) need not be a function
of finitely many nearby earlier variables.
Integration of \(M\) can transmit dependence from
\(L\).  The second term in (56.11) is the marginal
normalization term and cannot be omitted.
When differentiating a \emph{future} parameter,
one must additionally differentiate
\(h_i(\ell,t_O)\) if that parameter belongs to \(O\);
(56.12) makes no assertion that such a term vanishes.

\subsection{A populated compact contraction, and its limitation}

The original local potential supplies a constant
\(C_{\rm osc}\), independent of volume and the ordering,
for which changing the current group variable alone gives
\[
 |\mathcal V(x,\ell,m,t)-\mathcal V(x',\ell,m,t)|
       \le\omega_\beta:=C_{\rm osc}(\beta+1).
 \tag{56.13}
\]
Indeed only a bounded number of completed plaquette
and Jacobian terms depend on that free variable.
Their compact oscillations are bounded, with the
plaquette contributions proportional to \(\beta\).
These are the same full-density local bounds used
in V18, V33 and V38; no all-good-free-exterior
condition is required.

\begin{theorem}[Uniform conditional minorization]
\label{thm:f12v52-minorization}
For the actual compact channel (56.5), put
\[
                         a_\beta=e^{-2\omega_\beta}.
 \tag{56.14}
\]
Then, uniformly in \(i,t\) and the deterministic repair,
\[
 r_t(x,\ell)\ge a_\beta,\qquad
                \rho_i(t)\le1-a_\beta .
 \tag{56.15}
\]
Consequently
\[
 \|D_i\Pi_i f\|_2^2
     \le(1-a_\beta)^2
                   \|(I-\Gamma_i)\Pi_i f\|_2^2 .
 \tag{56.16}
\]
This is a volume-independent conditional contraction.
Its constant has no decay in the repair radius and
approaches one as the crude oscillation cost increases.
\end{theorem}
\begin{proof}
Integrating the pointwise exponential comparison from
(56.13) over \(M\) shows that the oscillation of
\(W_t(\,\cdot\,,\ell)\) in \(x\) is at most
\(\omega_\beta\).  After normalizing against Haar
probability in the current group, every density
\(\alpha_t(x\mid\ell)\) lies between
\(e^{-\omega_\beta}\) and \(e^{\omega_\beta}\).
The same bounds hold for its mixture \(\alpha_t(x)\).
Their ratio is \(r_t(x,\ell)\), proving its lower bound.

For \(a_\beta<1\), write the conditional channel as
\[
 \mu(dx\mid L=\ell,T=t)
      =a_\beta\alpha_t(dx)+(1-a_\beta)R_t(\ell,dx).
 \tag{56.17}
\]
The residual \(R_t\) is a Markov kernel and
\(\lambda_t R_t=\alpha_t\).  For centered
\(u\in L^2(\lambda_t)\), the channel from \(L\) to \(X\)
associated with this joint law is
\((1-a_\beta)\) times the corresponding residual
conditional expectation.  Conditional Jensen
bounds the latter's \(L^2\) norm by one.
This proves the norm bound in (56.15).
The case \(a_\beta=1\) is the independent channel.
Insert the bound into (56.8) and use conditional
total variance to get (56.16).
\end{proof}

The last factor in (56.16) is not a martingale
difference energy of the original \(f\).
It is a cumulative variance of a different function
\(\Pi_i f\) for each \(i\).  Even for independent
fair signs, with singleton repairs and \(f=Y_1\),
\[
 \sum_{i=1}^n\|(I-\Gamma_i)\Pi_i f\|_2^2=n-1,
 \qquad
 \sum_{i=1}^n\|D_i\Pi_i f\|_2^2=0.
 \tag{56.18}
\]
For \(i>1\), \(\Pi_i f=f\) and \(\Gamma_i f=0\);
for \(i=1\), \(\Pi_1 f=0\).  On the other hand,
the product projections give \(D_i\Pi_i=0\).
Thus the conditional contraction is real, but
summing its right sides does not supply a
volume-free leakage coefficient.

\subsection{Independent spectators preserve the correct quadratic form}

\begin{theorem}[Spectator stability of ordered leakage]
\label{thm:f12v52-tensor}
Suppose a base law \(\mu_0\), an ordered martingale
\((D_a^0)\), and deterministic repairs \((\Pi_a^0)\)
satisfy, on all of \(L^2(\mu_0)\),
\[
 \sum_a\|D_a^0\Pi_a^0 f\|_2^2
                    \le\varepsilon\operatorname{Var}_{\mu_0}(f).
 \tag{56.19}
\]
Adjoin any finite number of mutually independent
spectator coordinates, independent of the base.
Interleave them arbitrarily while preserving the
base order.  Extend each base repair by the identity
on spectators, and give each spectator its own
singleton heat bath.  The full ordered leakage
then satisfies
\[
               \sum_i\|D_i\Pi_i f\|_2^2
                   \le\varepsilon\operatorname{Var}_{\mu}(f)
 \tag{56.20}
\]
with exactly the same coefficient.
\end{theorem}
\begin{proof}
Write \(T_a=D_a^0\Pi_a^0\).
At a base coordinate \(a\), let \(Q_a\) integrate
all spectators that precede it.  Product
disintegration gives
\[
 D_a=D_a^0\otimes Q_a,\qquad
 \Pi_a=\Pi_a^0\otimes I,\qquad
 (D_a\Pi_a)^*(D_a\Pi_a)=T_a^*T_a\otimes Q_a .
 \tag{56.21}
\]
Since \(0\le Q_a\le I\), their sum is at most
\((\sum_aT_a^*T_a)\otimes I\).
The base assumption bounds this by
\(\varepsilon(I-\mu_0)\otimes I\), which is no
greater than \(\varepsilon(I-\mu)\).
For a spectator coordinate \(z\), the full
martingale difference contains the factor
\(I-\Pi_z\), commuting with all earlier
independent integrations.  Hence
\(D_z\Pi_z=0\).  This proves (56.20).
\end{proof}

This statement needs full \(L^2\) control on the
base, not just a first-chaos estimate.
It says nothing about adding interacting
spectators.  Within its hypotheses it proves
that the correct leakage target is insensitive
to arbitrary independent extra degrees of freedom.

\subsection{Why an anchored operator-norm Schur row can still fail}

An explicit positive finite-range example separates
the summed quadratic form from the norms of its
individual martingale blocks.  Let \(X_0,X_1,X_2,X_3\)
be fair-sign Markov-chain variables with law
\[
 \mu_0(x)=\frac1{16}
          \prod_{k=0}^{2}(1+r x_kx_{k+1}),\qquad 0<r<1 .
 \tag{56.22}
\]
Order them \(0,1,2,3\) and use the radius-one
repairs \(P_k=\{j\in\{0,1,2,3\}:|j-k|\le1\}\).
Write \(D_k\) for the difference at coordinate
\(X_k\).  Direct Markov conditioning gives, for
every \(f\in L^2(\mu_0)\),
\[
 \begin{split}
 D_0\Pi_0f&=D_1\Pi_1f=0,\\
 D_2\Pi_2f&=r^2(X_2-rX_3)\langle X_0,f\rangle,\\
 D_3\Pi_3f&=r^2X_3\langle X_1,f\rangle .
 \end{split}
 \tag{56.23}
\]
For the second line, \(\Pi_2 f\) is a function of
the sole exterior sign \(X_0\), and
\(\mu(X_0\mid X_2,X_3)=r^2X_2\),
\(\mu(X_0\mid X_3)=r^3X_3\).
For the third, the outside record is \(X_0,X_1\);
the past is screened by \(X_1\), and
\(\mu(X_1\mid X_3)=r^2X_3\).
These observations prove (56.23), including its
validity for arbitrary functions, not only linear ones.

The positive leakage operator therefore has trace
\[
             \varepsilon_0=r^4(1-r^2)+r^4=r^4(2-r^2).
 \tag{56.24}
\]
Its norm on centered functions is no greater than
this trace.  At \(r=1/2\), \(\varepsilon_0=7/64\).
Theorem~\ref{thm:f12v52-tensor} preserves this
all-\(L^2\) bound after adjoining any number of
independent spectators, at any locations.
The direct-sum variance argument of (55.19),
or its square-inequality weakening, gives the
uniform auxiliary gap lower bound
\[
                         (1-2\varepsilon_0)/2=25/64
                         \quad(r=1/2).
 \tag{56.25}
\]
This is a finite-state model certificate, not a
Yang--Mills estimate.

Now place independent fair signs \(Z_1,\ldots,Z_N\)
arbitrarily far from the chain and order the full
system as
\[
                  X_0,X_1,X_2,Z_1,\ldots,Z_N,X_3 .
 \tag{56.26}
\]
The interaction still has finite range; the
spectators have no interactions.  For
\(f_j=Z_jX_3\), \(\|f_j\|_2=1\) and
\(D_{Z_j}f_j=f_j\).  The current repair
\(\Pi_2\) integrates \(X_1,X_2,X_3\), hence
\[
 D_2\Pi_2D_{Z_j}f_j
       =r^5 Z_j(X_2-rX_3),\qquad
 \|D_2\Pi_2D_{Z_j}\|
       \ge r^5\sqrt{1-r^2}.
 \tag{56.27}
\]
Indeed \(\Pi_2(Z_jX_3)=r^3Z_jX_0\);
the second line of (56.23) then evaluates the
remaining difference.

The lower bound does not decrease as the
spectator location moves away.  Thus finite
interaction range and this ordered filtration
do not imply a bound
\(\|D_i\Pi_iD_j\|\le C e^{-c d(i,j)}\) with
uniform \(C,c>0\), where the index \(j\) is
identified only with its newly revealed site.
Nor can the absolute operator-norm row sum
be bounded uniformly in \(N\): at \(r=1/2\)
it is at least \(N\sqrt3/64\).

This is not an artifact of choosing a very
small buffer.  For any fixed integer \(R\ge1\),
take a sign chain \(X_0,\ldots,X_{2R+1}\),
put the current site at \(X_{R+1}\), use its
radius-\(R\) ball \(P=\{1,\ldots,2R+1\}\),
and insert the distant spectators immediately
after that current site in the ordering.
With \(f_j=Z_jX_{2R+1}\), the same calculation gives
\[
 D_{X_{R+1}}\Pi_P D_{Z_j}f_j
   =r^{3R+2}Z_j(X_{R+1}-rX_{R+2}).
 \tag{56.28}
\]
The coefficient is small for large \(R\), but
positive at every fixed \(R\) independently of
the distance to \(Z_j\).  Conditional Markov
propagation from \(X_0\) to \(X_{R+1}\) contributes
\(r^{R+1}\), and integration of the repair
contributes \(r^{2R+1}\), proving the formula.

The outputs in (56.27) for distinct \(j\) are
orthogonal, because the spectators are independent
and centered.  On the span of the \(f_j\)'s, the
row operator therefore has squared norm exactly
\(r^{10}(1-r^2)\), not a coefficient proportional
to \(N\).  The full summed leakage also remains
bounded by (56.20).  This is why the failed
operator-norm row is a loss of orthogonality,
not a counterexample to the correct quadratic
target or to a Yang--Mills gap.

\subsection{Direction and exact verification}

\begin{decision}[The nonlinear target must retain the spectator space]
The new compact formula is (56.6), with normalized
score (56.11).  The actual local action also
populates the conditional contraction (56.15).
Neither alone controls the global leakage sum.
The next step should estimate that sum as a
quadratic form, keeping spectator factors and
conditional output orthogonality, or use a
support-sensitive decomposition finer than one
newly revealed site per fluctuation.

An entrywise spatial operator-norm Schur estimate
for arbitrary full martingale differences is
not a consequence of locality and would exclude
even the explicit gapped model above.
The V51 scalar Gaussian regression rows remain
valid on their declared first-chaos space.
The examples here do not disprove a more
structured Yang--Mills estimate; they identify
what cannot be inferred before proving it.
\end{decision}

The script
\begin{center}\small\ttfamily
scripts/verify\_ym\_f12\_v52\_conditional\_screen.py
\end{center}
checks (56.6) as an exact matrix identity on
a positive interacting four-bit law, after
integrating a nonempty \(M\) and fixing each
future value.  Exact rational positive-semidefinite
tests verify the conditional minorization
contraction.  The sign-chain example with zero,
one and two spectators has \(16,32,64\) states.
All cases verify the leakage bound \(7/64\)
and gap lower bound \(25/64\).  They check the
nonzero remote row witness, of squared norm
\(3/4096\), and the orthogonality of distinct
spectator outputs.  All checks passed.
These tests verify finite operator algebra,
not the analytic constants of the compact YM
action or the continuum construction.

V51 had 804381 bytes and SHA--256
\begin{center}\scriptsize\ttfamily
A82B7D83696BFC28977CD87CCF10897EA9A742A1D34C31472D608E20C6842600.
\end{center}
Its complete body is retained apart from the
title version.  Only active V51 is replaced
by V52; archives, companions and prior scripts
remain.  The next companion audit is V55
and the next direction review is V60.

\begin{firewall}[Position after V52]
The original compact law now has an exact
ordered conditional-channel formula, its
normalized score identity and a populated
uniform conditional contraction.  Spectator
stability of a proved leakage quadratic form
is established, while a naive spatial
operator-norm assembly is ruled out in the
stated generality.  The actual volume-free
compact leakage bound remains unproved.
So do the full bounded-range fibre gap,
physical transfer comparison, scale-stable
continuum construction and physical
Yang--Mills mass gap.
\end{firewall}

\section{V53: a full compact likelihood bound across a buffered annulus}
\label{sec:f12v53}

\subsection{13 September 2026: the non-duplicative target}

V52 made progress by identifying the exact compact
conditional channel and a failure of naive spatial
operator-norm summation.  Its current source was
verified before this continuation.  We now work
with two geometric regions separated by an annulus,
rather than identify the support of a fluctuation
with one site in an arbitrary ordering.

V32 already proves total-variation decoupling for
the whole core of a region of size polynomial in
\(\beta\); V36 explains its maximal-correlation
consequence.  Those are not new conclusions here.
The stronger target is a \emph{pointwise likelihood
ratio}, relative to the product of the actual full
compact endpoint marginals.  It controls rare endpoint
events as well as bounded tests, and yields a
contraction for arbitrary density perturbations
at the outer endpoint.

Two earlier estimates combine for this purpose.
V48 gives a rectangular log-potential bound between
guarded component interfaces after full compact
integration of the intervening region.
V50 makes the failure of each interface guard
uniformly rare under an update that can be placed
away from both observed endpoints.  This last
placement permits conditioning on the endpoints
without losing the rarity estimate.

\subsection{Removing a guard at the level of likelihood ratios}

Let \(\mu\) be a probability law with endpoint
records \(U,V\) and an event \(G\).  Put
\(\nu=\mu(\,\cdot\mid G)\) and \(c=\mu(G)\).
All densities below are Radon--Nikodym derivatives
relative to the indicated marginals.

\begin{proposition}[Endpoint-conditioned guard transfer]
\label{prop:f12v53-transfer}
Suppose \(0\le q<1\),
\[
 \mu(G^c\mid U,V)\le q,
 \qquad
 e^{-\Omega}\le
       \frac{d\nu_{UV}}{d(\nu_U\otimes\nu_V)}
                           \le e^\Omega,
 \qquad \Omega\ge0 .
 \tag{57.1}
\]
Then the full, unconditioned endpoint law satisfies
\[
 e^{-\Lambda}\le
       \frac{d\mu_{UV}}{d(\mu_U\otimes\mu_V)}
                           \le e^\Lambda,\qquad
          \Lambda=\Omega-2\log(1-q).
 \tag{57.2}
\]
The guard does not have to be measurable from the
endpoint records.
\end{proposition}
\begin{proof}
Set \(s(u,v)=\mu(G\mid u,v)\),
\(s_U(u)=\mu(G\mid u)\), and
\(s_V(v)=\mu(G\mid v)\).
Each of \(s,s_U,s_V,c\) lies in \([1-q,1]\).
The exact changes of measure are
\[
 d\nu_{UV}=(s/c)d\mu_{UV},\quad
 d\nu_U=(s_U/c)d\mu_U,\quad
 d\nu_V=(s_V/c)d\mu_V .
\]
They give
\[
 \frac{d\mu_{UV}}{d(\mu_U\otimes\mu_V)}
   =
 \frac{d\nu_{UV}}{d(\nu_U\otimes\nu_V)}
                 \frac{s_U s_V}{c\,s}.
 \tag{57.3}
\]
The last factor lies between \((1-q)^2\) and
\((1-q)^{-2}\).  Equation (57.2) follows.
Strict positivity of the survival probabilities
also guarantees the required equivalence of the
endpoint marginals.
\end{proof}

The first hypothesis is conditional on \emph{both}
endpoints.  Stationary rarity alone does not imply
it.  For instance, a strictly positive law on two
three-state variables can have a rare simultaneous
marker whose centered endpoint indicators are almost
perfectly correlated, even though the endpoints
are independent after discarding all marker states.
Such an example is included in the exact tests below.

\subsection{The guarded likelihood ratio is supplied by V48}

Use a deterministic canonical split
\[
                   I_R=I_T\sqcup E_1\sqcup E_2,
 \tag{57.4}
\]
with no completed potential term meeting both
\(E_1\) and \(E_2\).  Terms that also meet \(I_T\)
are included in this requirement.
Write \(x=Y_{I_T}\), \(\eta_\alpha=Y_{E_\alpha}\).
The exact marginal after integrating the intervening
variables is V48's law
\[
 \rho(d\eta_1,d\eta_2)
  \ \propto\
       e^{-U_1(\eta_1)-U_2(\eta_2)-W_T(\eta_1,\eta_2)}
                                      \,d\eta_1d\eta_2 .
 \tag{57.5}
\]
The \(I_T\) integration is over the full compact
domain.  The two own-component potentials may
contain arbitrarily many compact variables.

Let \(K_\alpha\subset E_\alpha\) contain every
crossing input and the fixed physical-input
completion needed for the collar of \(T\).
The latter enlargement is justified as in V49:
added coordinates absent from \(W_T\) have zero
mixed derivative.  Choose fixed small product
chart cubes on the \(K_\alpha\)'s so that their
values make the physical frozen collar of \(T\)
good.  Write their product guard as
\[
                         G=\Gamma_1\cap\Gamma_2 .
 \tag{57.6}
\]
Any inherited frozen input of \(T\) must also
be good; in the annular construction below none
comes from the original free exterior.
The retained coarse curvatures satisfy V32's
fixed small threshold.

In scalar interface coordinates of radius \(a\),
let \(k_\alpha\) be the number of coordinates and
\(d_K=\operatorname{dist}(K_1,K_2)\).
Equations (52.5)--(52.7) supply
\[
 |\Delta_1\Delta_2
        [U_1+U_2+W_T]|\le\Omega_T,\qquad
 \Omega_T\le
 C k_1k_2\left[
       \beta e^{-\gamma d_K}
          +\beta^2\,73|T|e^{-c_0\beta}\right].
 \tag{57.7}
\]
We have replaced the minimum in \(p_T\) by its
larger exponential expression, which remains a
valid upper bound.  Constants absorb \(4a^2\)
and fixed group-coordinate multiplicities.
This use of V48 does not require a fixed bound
on \(k_\alpha\); that bound was needed for its
many-component row-sum corollary, not for (52.7).

\begin{proposition}[Two guarded components have a relative bound]
\label{prop:f12v53-guarded}
Under these hypotheses,
\[
 e^{-\Omega_T}\le
 \frac{d\rho^G}
      {d(\rho^G_{E_1}\otimes\rho^G_{E_2})}
                       \le e^{\Omega_T}.
 \tag{57.8}
\]
The same inequality holds for any endpoint
records \(U=U(\eta_1)\), \(V=V(\eta_2)\).
\end{proposition}
\begin{proof}
Fix two values \(\eta_2,\eta'_2\) in the second
guarded component.  By the rectangular bound,
the oscillation over \(\eta_1\) of the difference
of their unnormalized conditional potentials is
at most \(\Omega_T\).  Including the two conditional
normalizers therefore bounds their density ratio
between \(e^{-\Omega_T}\) and \(e^{\Omega_T}\).
Integrate the second conditional law against the
true \(\rho^G_{E_2}\)-marginal.  This yields the
same ratio bound between
\(\rho^G(d\eta_1\mid\eta_2)\) and \(\rho^G_{E_1}\),
which is (57.8).

For endpoint records, disintegrate the product
of the two component marginals over \(U,V\).
The new relative density is the conditional
average of the density in (57.8) under that
product disintegration.  Its two bounds persist.
\end{proof}

In particular, neither component interior is
put in a good box.  Only its interface is
guarded, exactly as in the original V48 estimate.

\subsection{Conditioning on endpoints does not spoil buffered rarity}

Choose a fixed free-link threshold \(\tau_g>0\)
whose section balls lie inside the interface
cubes.  Use V50's repair set \(P_i^-\) at
that threshold for every group variable
\(i\in K_1\cup K_2\).  Here \(K_\alpha\) denotes
group inputs when specifying a repair; scalar
counts are used only in (57.7).
Assume that each repair lies inside the ambient
canonical update, and that it is disjoint from
the free-variable supports of \(U,V\).
The V50 coarse-curvature hypotheses hold on
the associated comparison neighbourhoods.

Put \(k_g=|K_1\cup K_2|\) in group variables and
\[
                q_T=k_g e^{-c_g\beta},\qquad
                       c_g=\tau_g^2/(16A).
 \tag{57.9}
\]
Then, whenever \(q_T<1\),
\[
                         \mu(G^c\mid U,V)\le q_T .
 \tag{57.10}
\]
Indeed \(G^c\) is contained in the union of
the events \(B_i(\tau_g)\).  The complete
endpoint record is exterior-measurable for
each \(P_i^-\).  Taking the conditional
expectation of (54.5) given \(U,V\), then
summing the deterministic events, proves
(57.10).  There is no conditioning on a
selected repair or a selected bad site.

Combining (57.8) and (57.10) in
Proposition~\ref{prop:f12v53-transfer} proves
the following full compact estimate.

\begin{theorem}[Buffered annular likelihood estimate]
\label{thm:f12v53-annular}
Under the canonical split, collar, curvature,
interface and endpoint-buffer hypotheses above,
\[
 \boxed{\quad
 e^{-\Lambda_T}\le
 \frac{d\mu_{UV}}{d(\mu_U\otimes\mu_V)}
             \le e^{\Lambda_T},\qquad
       \Lambda_T=\Omega_T-2\log(1-q_T).
 \quad}
 \tag{57.11}
\]
All laws in (57.11) are the original compact
laws, not restrictions to the interface guard.
The estimate is uniform over the original
compact exterior whenever the declared
geometric hypotheses hold.
\end{theorem}
\begin{proof}
The endpoint law under \(\mu(\,\cdot\mid G)\)
is the endpoint marginal of \(\rho^G\), since
the guard is measurable on \(E_1,E_2\).
Apply (57.8), (57.10), and (57.2).
\end{proof}

The crucial distinction from removing a guard
for a whole component is geometric.  Here the
endpoint observables are set back from their
component interfaces.  A repair can change a
guard variable while leaving both endpoint
records fixed.  If \(U,V\) contained the whole
interfaces, (57.10) would usually be false.

\subsection{A concrete four-dimensional family}

Write \(B_r=\{z\in\mathbb Z^4:\|z\|_\infty\le r\}\).
All boxes and their fixed stencils are assumed
to embed without periodic wrapping.
Let \(\ell\ge0\) be an integer, choose a fixed
buffer \(b\), and put
\[
                  T=B_{\ell+b+w}\setminus B_{\ell+b}.
 \tag{57.12}
\]
Take an ambient cell region containing
\(B_{\ell+3b+w}\).  Let \(U\) be all independent
free groups with owners in \(B_\ell\), and
let \(V\) be any or all ambient free groups
whose owners lie outside \(B_{\ell+2b+w}\).

Here and below \(b\) is fixed but can be very
large.  Let \(r_{\rm loc}\ge1\) bound the completed
interaction diameter, physical-input closure
and the owner-strip shifts, and let \(H_g\)
bound the owner radius of every V50 repair
\(I_{P_i^-}\).  One sufficient choice is
\[
             b\ge H_g+3r_{\rm loc}+10,\qquad
             w\ge4r_{\rm loc}+10 .
\]
Then (57.4) is the actual canonical partition:
the remaining internal groups lie in the hole
or in the outside component; a remaining boundary
family has both endpoint cells in one of those
components.  Every boundary family incident to
\(T\) is already in \(I_T\).
An edge cannot jump across the annulus.

Thus the remaining variables split as \(E_1,E_2\).
Their owner separation exceeds the completed
interaction range, so no completed term meets
both.  The input sets \(K_1,K_2\), including their
physical completion, lie in fixed-width strips
along the respective boundaries of \(T\).
Consequently
\[
 \begin{split}
 |T|&\le C(\ell+b+w+1)^4,\\
 k_1+k_2&\le C(\ell+b+w+1)^3,\qquad k_g\le k_1+k_2,\\
 d_K&\ge w-C_{\rm loc}.
 \end{split}
 \tag{57.13}
\]
These follow by counting fixed-width differences
of four-dimensional cubes, with fixed group
and scalar multiplicities.

The choice of \(b\) ensures that the V50 repairs
around \(K_1,K_2\) neither meet the endpoint
records nor reach the ambient exterior.
All free inputs affecting the conditional
kernel on \(I_T\) belong to \(E_1\cup E_2\)
and are covered by the completed guard.
No arbitrary original frozen free input touches
the annulus.  This is why the original ambient
exterior need not be good.
The coarse field remains fixed throughout and
must satisfy the declared uniform small-curvature
threshold on \(T\) and its repair neighbourhoods.

\begin{corollary}[Polynomial-size cores at a linear buffer scale]
\label{cor:f12v53-scale}
Fix \(s\ge0\).  If
\[
                    \ell\le\beta^s,\qquad
                    w=\lceil\beta\rceil ,
 \tag{57.14}
\]
then for sufficiently large \(\beta\) there
are fixed \(C_s,c_s>0\), independent of ambient
volume and exterior values, such that
\[
                         \Lambda_T\le C_s e^{-c_s\beta}.
 \tag{57.15}
\]
\end{corollary}
\begin{proof}
Put \(d=\max\{s,1\}\).  For large enough \(\beta\),
(57.13) gives \(k_\alpha\le C_s\beta^{3d}\),
\(|T|\le C_s\beta^{4d}\), and \(d_K\ge\beta/2\).
Hence
\[
 \Omega_T\le C_s\left[
    \beta^{1+6d}e^{-\gamma\beta/2}
       +\beta^{2+10d}e^{-c_0\beta}\right],\qquad
 q_T\le C_s\beta^{3d}e^{-c_g\beta}.
 \tag{57.16}
\]
For \(q_T\le1/2\),
\(-2\log(1-q_T)\le4q_T\).
Absorb the fixed polynomial powers into a
decrease of the exponential rate.
\end{proof}

The outer component can be arbitrarily large:
only the annulus and its interface sizes enter
these estimates.  On the other hand, the inner
size in (57.14) grows only polynomially with
\(\beta\).  The bound does not let \(\ell\)
tend to infinity at fixed \(\beta\).

\subsection{Full-marginal domination and density perturbations}

Let \(K(v,du)=\mu(dU=u\mid V=v)\), and put
\[
                 a_T=e^{-\Lambda_T}
                     =e^{-\Omega_T}(1-q_T)^2 .
 \tag{57.17}
\]
The likelihood theorem gives the pointwise
minorization
\[
                    K(v,du)\ge a_T\,\mu_U(du).
 \tag{57.18}
\]
The common measure is the actual full compact
inner marginal, including its rare configurations.
It is not a chart-truncated reference density.
Kernel bounds use the canonical positive-density
versions for the endpoint group records.  For
general measurable endpoint records they are
understood as essential bounds.

\begin{corollary}[Relative channel contraction]
\label{cor:f12v53-relative}
The endpoint maximal correlation and its
conditional total-variation diameter obey
\[
 \rho_\mu(U,V)\le1-a_T,\qquad
 \sup_{v,v'}\|K(v,\cdot)-K(v',\cdot)\|_{\rm TV}
                                                  \le1-a_T .
 \tag{57.19}
\]
For any input law \(\lambda\ll\mu_V\),
\[
 \begin{split}
 D(\lambda K\Vert\mu_U)&\le
                    (1-a_T)D(\lambda\Vert\mu_V),\\
 \chi^2(\lambda K\Vert\mu_U)&\le
                    (1-a_T)^2\chi^2(\lambda\Vert\mu_V).
 \end{split}
 \tag{57.20}
\]
Infinite right sides are allowed.  The roles
of the two endpoints can be exchanged.
\end{corollary}
\begin{proof}
Write \(K=a_T\mu_U+(1-a_T)R\); averaging
over \(\mu_V\) gives \(\mu_VR=\mu_U\).
The common part proves the total-variation
bound.  Conditional Jensen for the residual
joint law proves the centered \(L^2\) contraction
of \(R\), exactly as in V52, hence the first
bound in (57.19).

For an input density \(f=d\lambda/d\mu_V\),
the residual output density \(r\) is its
conditional expectation under the residual
joint law.  The actual output density is
\(a_T+(1-a_T)r\).
Jensen for \(z\mapsto(z-1)^2\) and
\(z\mapsto z\log z\), followed by convexity
of the latter at \(a_T\cdot1+(1-a_T)r\),
gives (57.20).  Approximation handles
unbounded densities.  Symmetry follows
from (57.11).
\end{proof}

Thus (57.15) controls arbitrary outer density
tilts, not only fixed bounded observables.
This is the relative improvement over V32's
already-proved total-variation estimate.
It does not imply that conditioning on
arbitrary \emph{intermediate} coordinates
preserves the annular repair buffers.

For completeness, the inherited two-projection
variance argument of V36 also gives
\[
 \operatorname{Var}_\mu(f)
 \le a_T^{-1}\left[
     \mu\operatorname{Var}(f\mid U)
        +\mu\operatorname{Var}(f\mid V)\right].
 \tag{57.21}
\]
These two heat baths integrate respectively
the complements of \(U\) and \(V\).
One complement includes the entire outer
component, which may be large.  Therefore
(57.21) is not a bounded-range update gap.

\subsection{Verification, frontier and append record}

The exact-arithmetic script
\begin{center}\small\ttfamily
scripts/verify\_ym\_f12\_v53\_annular\_channel.py
\end{center}
checks two positive 36-state endpoint/interface
models.  It verifies the exact likelihood
identity (57.3), the endpoint-conditioned
failure bound, the transferred full density
ratio and the all-\(L^2\) two-block form by
rational positive-semidefinite elimination.

A strictly positive nine-state negative
control has stationary guard failure
\(1/100\), but its rare marker indicators
have correlation
\(2474009999/2474509999\).
Incorrectly substituting stationary rarity
for (57.1) would bound that correlation
by \(199/10000\).  Conditional failure
on the double-marker record is actually
one, so the theorem correctly does not
apply.

The geometry checks enumerate canonical
incident boundary families for two
four-dimensional annuli, with 624 and
2320 cells.  They verify the exact partition
of the independent group counts
\(17|R|+7|E(R)|\) into annular, inner and
outer coordinates.  They do not constitute
an exhaustive test of all physical
completion stencils; the analytic
fixed-range and fixed-strip arguments
above supply that general geometry.
All listed checks passed.

\begin{decision}[The useful new input is relative annular decoupling]
The guard can now be removed in a relative
endpoint likelihood estimate, provided both
endpoint supports are set back from the
interfaces by the declared repair buffer.
The nonlinear full-marginal channel then
contracts entropy and squared-density
perturbations with an exponentially small
coefficient in the finite-scale regime.

The next task is to use this annular input
in a genuine local energy or entropy
decomposition.  An ordered conditioning
that freezes sites inside the annulus
does not satisfy the theorem automatically.
Likewise, summing (57.21) without controlling
its large-complement updates would repeat
the unresolved V49 localization problem.
\end{decision}

V52 had 822804 bytes and SHA--256
\begin{center}\scriptsize\ttfamily
8B8508E35230AD6C822D6CBA67423F5F1B7FCD08CA6BB4294CD8E69617095949.
\end{center}
Its complete body is retained apart from
the title version.  Only active V52 is
replaced by V53.  All archives, companions
and prior scripts remain.  The next
companion audit is V55; the next
direction review is V60.

\begin{firewall}[Position after V53]
A full compact, relative endpoint
likelihood estimate and its density-channel
consequences are proved in the stated
annular geometry, at fixed small coarse
curvature and sufficiently large coupling.
They do not establish an iteration at
fixed coupling through arbitrary pins,
a volume-uniform bounded-range fibre gap,
physical transfer-energy comparison,
scale-stable continuum construction or
the physical Yang--Mills mass gap.
\end{firewall}

\section{V54: endpoint cross-ratios, extensive safe pinning
and relative source response}
\label{sec:f12v54}

\subsection{The conditioning question after V53}

V53 made progress by proving a relative likelihood bound
across an annulus whose interface guards have fully
integrated repair buffers.  Its source and hash were
rechecked at the start of this continuation.
The remaining issue is not resolved by saying that
likelihood bounds survive arbitrary conditioning:
conditioning may use up the very buffers that produced
the estimate.

Here we prove a precise safe class.  The V53 construction
actually bounds the rectangular log cross-ratio of the
full endpoint density, with the same constant
\(\Lambda_T=\Omega_T-2\log(1-q_T)\).
That bound survives separate reweightings of the endpoints,
separate marginalizations, and arbitrary point pins within
both endpoint records, without a penalty for the number
of pinned variables.  No variable in the protected
intermediate region is pinned by these operations.

This is different from V45's fixed-core density comparison,
whose \(\beta^n\) cost limits its pinning dimension.
The present estimate uses the annular geometric hypothesis
and says nothing about pinning a general collection of
intermediate or interface variables.  It does not replace
V45 outside this new domain.

\subsection{Separate marginalization preserves rectangular bounds}

For a positive density \(p(u,v)\) relative to a product
reference, define its rectangular log increment by
\[
 \mathfrak R\log p
 =\log p(u,v)+\log p(u',v')
              -\log p(u,v')-\log p(u',v).
 \tag{58.1}
\]
Multiplying \(p\) by a positive function of \(u\)
or a positive function of \(v\) does not change
this increment.  Thus it is independent of the
choice of separate endpoint reference densities.
All assertions are made on the positive-density
support, with canonical versions for group records.

\begin{proposition}[Rectangular bounds under separate integration]
\label{prop:f12v54-marginal}
Let
\[
 d\nu(a,b)=Z^{-1}F(a,b)\,\sigma_1(da)\sigma_2(db),
 \qquad F>0,
 \qquad |\mathfrak R\log F|\le\Omega .
 \tag{58.2}
\]
For separate records \(U=U(a)\), \(V=V(b)\),
the endpoint density satisfies
\[
                         |\mathfrak R\log p_{\nu,UV}|
                                                    \le\Omega .
 \tag{58.3}
\]
There is no factor depending on the number of
integrated variables.
\end{proposition}
\begin{proof}
Disintegrate the two reference measures separately
over \(u,v\).  Apart from the separate reference
marginals, the new density is
\[
          F_*(u,v)=\int F(a,b)\,\sigma_1(da\mid u)
                                      \sigma_2(db\mid v).
\]
The pointwise cross-ratio inequality gives
\[
 F(a,b)F(a',b')
             \le e^\Omega F(a,b')F(a',b).
\]
Integrate this inequality against the product of
the four conditional reference measures at
\(u,u',v,v'\).  Independence of those four
integrations yields
\[
 F_*(u,v)F_*(u',v')
             \le e^\Omega F_*(u,v')F_*(u',v).
\]
Interchange \(v,v'\) to get the reverse bound.
Separate marginal densities and \(Z\) cancel
from the resulting logarithmic rectangle.
\end{proof}

This is a statement about positive kernels and
separate integrations.  It would not follow by
integrating a signed mixed derivative and dropping
its normalization terms.

\subsection{The full compact endpoint cross-ratio}

Retain the complete V53 annular geometry, its
endpoint records \(U,V\), and the same fixed
coarse field and coupling hypotheses.
The guarded component density (57.5) has a
rectangular bound \(\Omega_T\).  Its own-component
potentials are separate reference factors.
Proposition~\ref{prop:f12v54-marginal} therefore
strengthens (57.8) to
\[
                      |\mathfrak R\log p_{\mu(\cdot\mid G),UV}|
                                                     \le\Omega_T.
 \tag{58.4}
\]
The full endpoint density is related to this one
by
\[
 p_{\mu,UV}(u,v)
    =\frac{c}{s(u,v)}p_{\mu(\cdot\mid G),UV}(u,v),
 \qquad s(u,v)=\mu(G\mid u,v)\in[1-q_T,1].
 \tag{58.5}
\]

\begin{theorem}[Full compact rectangular estimate]
\label{thm:f12v54-rectangle}
The original compact endpoint law satisfies
\[
 \boxed{\quad
 |\mathfrak R\log p_{\mu,UV}|\le\Lambda_T,
 \qquad \Lambda_T=\Omega_T-2\log(1-q_T).
 \quad}
 \tag{58.6}
\]
The constant is exactly the V53 likelihood constant,
not four times that constant.
\end{theorem}
\begin{proof}
The oscillation of \(\log s\) is at most
\(-\log(1-q_T)\).  Its rectangular increment
is the difference of two sums of two values,
so its absolute value is at most
\(-2\log(1-q_T)\).  Apply this to (58.5)
and use (58.4).  The constant \(\log c\)
cancels.  This uses the underlying guarded
rectangle and conditional survival probability,
not just the weaker information that the final
likelihood ratio lies in a fixed interval.
\end{proof}

\subsection{Separate reweighting and point pinning}

Let \(A(U),B(V)\) be nonnegative measurable weights
such that
\[
 0<\mu[A(U)B(V)]<\infty,\qquad
       d\mu^{A,B}
        =\frac{A(U)B(V)}{\mu[A(U)B(V)]}\,d\mu .
 \tag{58.7}
\]
These weights can encode large biases or
positive-probability events within each endpoint.
They need not be small perturbations and need
not factor into individual site weights.

\begin{theorem}[Safe endpoint conditioning without a pin-count loss]
\label{thm:f12v54-pinning}
The endpoint law under (58.7) has the same
rectangular bound \(\Lambda_T\).
If \(U=(U_0,U_1)\), \(V=(V_0,V_1)\) are
coordinate subrecords, then the conditional
law of \(U_1,V_1\) given \(U_0=u_0,V_0=v_0\)
also has the same rectangular bound.
This holds for every canonical pin of the
original positive compact group density;
with general measurable weights it holds
for almost every pin with a finite,
nonzero conditional normalizer.

For each resulting law \(\widetilde\mu\),
possibly followed by further separate
marginalization, one has
\[
 e^{-\Lambda_T}\le
 \frac{d\widetilde\mu_{U_1V_1}}
      {d(\widetilde\mu_{U_1}\otimes\widetilde\mu_{V_1})}
                 \le e^{\Lambda_T}.
 \tag{58.8}
\]
The endpoint channel maximal correlation,
total-variation diameter, relative-entropy
contraction and squared-density contraction
obey the V53 bounds with the unchanged
\[
                  a_T=e^{-\Lambda_T},\qquad
                  \varepsilon_T=1-a_T.
 \tag{58.9}
\]
\end{theorem}
\begin{proof}
On the product support of the two weights,
\(\log A(u)+\log B(v)\) has zero rectangular
increment.  Normalization is a constant,
so (58.6) persists.

For fixed pins \(u_0,v_0\), evaluate the
original density at those fixed arguments
and normalize over \(u_1,v_1\).
Every rectangular increment in \(u_1,v_1\)
is an original increment with the same pins
in its four corners.  It is bounded by
\(\Lambda_T\), independently of how many
coordinates have been pinned.  Positivity
of the original density gives a positive
normalizer for every group-coordinate pin.
This is a conditional-density calculation,
not division by the probability of an
exact-value event.

Further separate marginalization is covered
by Proposition~\ref{prop:f12v54-marginal}.
Finally, a rectangular log bound \(\Lambda_T\)
implies that any two conditional laws on one
endpoint, obtained from two values of the
other, have density ratio in
\([e^{-\Lambda_T},e^{\Lambda_T}]\).
This includes their conditional normalizers,
as in the proof of (57.8).  Averaging the
second conditional law against its actual
endpoint marginal proves (58.8).
The common full-marginal minorization and
the Jensen arguments of (57.18)--(57.20)
then give all the stated channel bounds.
\end{proof}

At fixed annular geometry, there is no additional
factor \(\beta^{\dim U_0+\dim V_0}\), no inverse
small pin probability, and no small-chart
restriction on the pinned endpoint values.
The underlying \(\Lambda_T\) still depends on
the original interface sizes and coupling.
In the regime \(\ell\le\beta^s\), \(w=\lceil\beta\rceil\)
of V53, its exponentially small upper bound is
uniform over all such endpoint pins.

The geometric scope is essential.  The theorem
conditions only on records from the set-back
inner endpoint or the distant outer endpoint.
It does not pin variables in the annulus,
its guarded interfaces, or the intervening
repair buffers.  The protected variables
are still integrated.  This is a proved safe
class of interior and exterior pins, not a
claim about arbitrary punctured regions.

\subsection{Joint endpoint constraints have a genuine interaction cost}

For a strictly positive joint weight \(W(U,V)\),
the rectangular increment of \(\log W\) need
not vanish.  If
\[
                        |\mathfrak R\log W|\le\Xi ,
 \tag{58.10}
\]
then the same argument gives a rectangular
bound \(\Lambda_T+\Xi\), and (58.8)--(58.9)
use that enlarged constant.  A joint hard
constraint can have no finite such bound.

For example, start from two independent fair
signs \(U,V\) and reweight by \(1+rUV\),
\(0<r<1\).  The resulting correlation is \(r\),
and the largest cross-ratio is
\[
                         \left(\frac{1+r}{1-r}\right)^2.
 \tag{58.11}
\]
The original rectangular constant was zero.
This positive example rules out silently
treating a coupled endpoint constraint as
two independent pins.

\subsection{A normalized outer source has a small relative inner response}

Keep any fixed safe endpoint pins and fixed
separate weights.  Vary only an outer potential:
\[
 d\mu_t
   =Z_t^{-1}\exp[-H_t(V)]A(U)B(V)\,d\mu,\qquad
                  S_t(V)=\partial_tH_t(V).
 \tag{58.12}
\]
Assume \(H_t,S_t\) are bounded and continuously
differentiable in the parameter on the interval
used below.  The endpoint pins do not move.
Let \(p_t(u)\) denote the positive density of
the unpinned inner record under the resulting
conditional law.  Separate reweighting leaves
the rectangular constant unchanged for every \(t\),
regardless of the magnitude of \(H_t\).

\begin{proposition}[Pointwise normalized source estimate]
\label{prop:f12v54-source}
For the same \(\varepsilon_T=1-e^{-\Lambda_T}\),
\[
 \partial_t\log p_t(u)
   =-\mathbb E_t[S_t\mid U=u]+\mathbb E_t S_t,
 \qquad
 \big|\partial_t\log p_t(u)\big|
                         \le\varepsilon_T\operatorname{osc}(S_t).
 \tag{58.13}
\]
Expectations in this formula include the fixed
endpoint pins and weights.  Hence
\[
 \left|\log\frac{p_1(u)}{p_0(u)}\right|
       \le\varepsilon_T\int_0^1
                             \operatorname{osc}(S_t)\,dt .
 \tag{58.14}
\]
The result is uniform over the inner record,
including low-probability configurations.
\end{proposition}
\begin{proof}
Differentiate the numerator defining the inner
marginal and its full normalizer.  The two
terms are exactly the conditional and
unconditional means in (58.13).  In particular
\(\int p_t\,\partial_t\log p_t=0\); omitting
the second term would violate normalization.

The reverse endpoint channel under \(\mu_t\)
has the same full-marginal minorization as
the forward channel, by (58.8).
Its conditional law of \(V\) given \(u\)
is therefore at TV distance at most
\(\varepsilon_T\) from its own marginal.
The bounded-test dual inequality gives
\[
 |\mathbb E_t[S_t\mid u]-\mathbb E_tS_t|
                     \le\varepsilon_T\operatorname{osc}(S_t).
\]
Integrate in \(t\).  Compact positive densities
and the bounded source justify these
differentiations.  On restricted supports
the assertion is made at points of positive
conditional density.
\end{proof}

If an inner weight also varied with \(t\),
its direct logarithmic derivative would be
an additional term.  Likewise this proposition
does not differentiate a moving guard,
moving pin, or moving coordinate chart.

\subsection{Application to an actual exterior free-input path}

Fix the coarse field \(C\).  Suppose a smooth
path of a prescribed exterior free group
changes only completed terms whose active
free-variable inputs lie in the outer
endpoint record \(V\).  All completion
ancestors of the changed terms must be
included in this support check.  Then the
actual parent potential has the form
\[
                \mathcal V_t(Y)=\mathcal V_0(Y)+H_t(V).
 \tag{58.15}
\]
This condition makes (58.12) the exact
normalized family, not an approximation
that neglects a changing conditional kernel.

The original bounded local derivatives give
\[
             \|S_t\|_\infty\le C(\beta+1)v(t),
 \tag{58.16}
\]
where \(v(t)\) is the path speed and the
constant includes the fixed number of
affected completed factors.  Thus, writing
\(\mathcal L=\int_0^1v(t)\,dt\),
\[
 \boxed{\quad
 \left|\log\frac{p_1(u)}{p_0(u)}\right|
       \le C(\beta+1)\varepsilon_T\,\mathcal L .
 \quad}
 \tag{58.17}
\]
Fixed safe pins can be imposed on either
endpoint without changing this coefficient.
In V53's polynomial-core, linear-buffer
regime it is exponentially small in \(\beta\),
after absorbing the factor \(\beta+1\).

This is a relative density estimate, and it
vanishes with the source path length.
It is not an additive rare-event error
that stays nonzero for identical inputs.
Earlier V35--V47 response estimates had
different observables or pinning domains;
the new point is uniform relative response
after arbitrarily many \emph{safe endpoint}
pins, using the full compact annular law.

\subsection{Verification and the remaining geometric obstruction}

The rational script
\begin{center}\small\ttfamily
scripts/verify\_ym\_f12\_v54\_endpoint\_pinning.py
\end{center}
starts with a positive 144-state model
having two two-bit endpoints and two
three-state interfaces.  It checks the
guarded and full endpoint cross-ratios,
three nontrivial separate tilts, and
all four choices of two-sided point pins.
The original and separately tilted
cross-ratios agree exactly.

Every case verifies the full-marginal
likelihood bounds and the centered source
identity and estimate in (58.13).
Cases with a nonzero source mean ensure
that the normalizer term is actually used.
A positive joint-tilt counterexample has
correlation \(3/4\) and cross-ratio \(49\),
confirming the need for the additional
interaction cost.  All checks passed.
These tests concern finite density algebra,
not the analytic YM constants or a
continuum mass-gap certificate.

\begin{decision}[What conditioning is now justified]
Endpoint-only pins and separate endpoint
biases are no longer an unquantified
loss in the annular estimate.
The next construction must still arrange
that the pins used by its variance or
entropy decomposition belong to such
safe endpoint records.  A standard
site-by-site order generally fixes
variables inside a required annulus,
which is outside this theorem.

The next iteration is V55 and must
perform the scheduled SM/NS/parallel-YM
audit before choosing the next
decomposition.  The next ten-version
direction review remains V60.
\end{decision}

V53 had 840935 bytes and SHA--256
\begin{center}\scriptsize\ttfamily
02987B3DF6245FA0197A9F8DC010CAFDCE1F9211A2AE3BD8A49D178BDC8B6234.
\end{center}
Its complete body is retained apart
from the title version.  Only active
V53 is replaced by V54; archives,
companions and prior scripts remain.

\begin{firewall}[Position after V54]
The full compact annular cross-ratio,
its dimension-independent safe endpoint
pinning, and its normalized relative
outer-source response are proved
under the inherited geometry and
small-coarse-curvature hypotheses.
No claim is made for pins that
intersect the repair buffers.
A global local-update decomposition,
volume-uniform bounded-range fibre gap,
physical transfer comparison,
scale-stable continuum construction
and physical Yang--Mills mass gap
remain unproved.
\end{firewall}

\section{V55: audited separated-core tensorization and its exact size budget}
\label{sec:f12v55}

\subsection{Compulsory companion audit}

The current active companion files were inspected again.
Their versions, sizes and SHA--256 values are unchanged
from the V50 audit:
\begin{center}
\begin{tabular}{lll}
\toprule
Programme & Version & Bytes\\
\midrule
SM--Einstein & V72 & 607372\\
Navier--Stokes & V26 & 278283\\
Parallel YM & V5 & 54174\\
\bottomrule
\end{tabular}
\end{center}
The corresponding digests are
\begin{center}\scriptsize\ttfamily
F44F9745D43379E1672C5550411B58E97195A4C9278592D221DF5D3F644DC1F5\\
82C887F8C2C69E4F28FAD7C3DC8DE5B9099CB5A4BF431EBA1FFF3E91935978AD\\
306F1BB84CFA8A8D47EA569EC17E9545F9867C8D32B548EC9D9C444B4F3C0512.
\end{center}
SM V72 proves its declared fixed-mode, fixed-time
many-copy limit, not a cutoff-uniform physical continuum
or the present YM energy estimate.  NS V26 remains a
rank-one norm calculation rather than a global-regularity
proof.  Parallel YM V5 rejects a volume-uniform tail
for an extensive total charge and retains local or
compact observables.  No companion supplies the missing
local-update closure.  Their internal programme or
stop wording is not treated as a new user instruction.

V54 made concrete progress by proving safe endpoint
pinning and normalized relative source response.
The next question is how much of a multi-core
variance estimate those results actually provide.
The answer below has an explicit core-count budget.
It is not legitimate to replace that budget by a
volume-independent constant without an additional
spatial argument.

\subsection{A separated family of actual core marginals}

Fix the same coarse field and full compact finite-volume
fibre law \(\mu\).  Select \(N\ge1\) core records
\(\eta_a=Y_{U_a}\), each consisting of independent
free groups in a translated copy of the V53 inner
endpoint.  Each has its own protected annulus
and repair buffers.  Require that every other
core be contained in its distant outer endpoint.

For equal core radius \(\ell\), buffer \(b\)
and annular width \(w\), one sufficient condition
is that the translated boxes of radius
\(\ell+3b+w+r_{\rm loc}\) be mutually disjoint
and lie inside the ambient region, with all
fixed stencils embedded without wrapping.
This is stronger than necessary but makes
the required protection explicit.
All retained-curvature hypotheses hold in
those protected neighbourhoods.

Let
\[
 Z=I_R\setminus\bigsqcup_{a=1}^N U_a,\qquad
                 \lambda=\mu_{(\eta_1,\ldots,\eta_N)} .
 \tag{59.1}
\]
The marginal \(\lambda\) integrates every group in
\(Z\) over its full compact domain.
It is not a product law and not a boxed law.
For core \(a\), choose the V53 outer endpoint to
include all the other cores.  V54 gives
\[
 |\mathfrak R_a\log p_\lambda|\le\Lambda_a ,
 \tag{59.2}
\]
where the rectangle varies \(\eta_a\) against
the joint remaining core record.
The same bound survives pins of any subset
of the other cores, and partial pins within
the cores, by the safe endpoint theorem.
Those pins do not enter any required annulus.

Write the canonical \emph{marginal} core projections
and forms as
\[
 Q_aF=\lambda(F\mid\eta_{a^c}),\qquad
        \mathcal E_{\rm core}(F)
                    =\sum_{a=1}^N\|(I-Q_a)F\|_{L^2(\lambda)}^2.
 \tag{59.3}
\]
The word marginal matters: \(Q_a\) is not generally
the original full-parent heat bath on \(U_a\).
The intervening variables have already been integrated.

\subsection{Finite-core tensorization with all factors displayed}

\begin{theorem}[Pin-stable finite-core bound]
\label{thm:f12v55-cores}
Under the preceding hypotheses, put
\[
 B_N=\sum_{a=1}^N\Lambda_a-\max_a\Lambda_a,\qquad
                         g_N=e^{-B_N}.
 \tag{59.4}
\]
Then
\[
       \operatorname{Var}_\lambda(F)
                      \le g_N^{-1}\mathcal E_{\rm core}(F)
                      \quad(F\in L^2(\lambda)).
 \tag{59.5}
\]
The same construction applies to every safe conditional
core law, using its remaining core blocks and inherited
constants.  If \(\Lambda_a\le\Lambda_*\), then
\[
                              g_N\ge e^{-(N-1)\Lambda_*}.
 \tag{59.6}
\]
For \(N=1\) the empty exponent gives the exact
one-block constant \(g_1=1\).
\end{theorem}
\begin{proof}
The relative minorization from (59.2) gives maximal
correlation at most \(1-e^{-\Lambda_a}\) between
core \(a\) and the joint remaining core record.
The two-projection variance inequality already used
in V36 and V39 therefore gives
\[
 \operatorname{Var}_\lambda(F)
 \le e^{\Lambda_a}\left[
       \lambda\operatorname{Var}(F\mid\eta_{a^c})
         +\lambda\operatorname{Var}(F\mid\eta_a)\right].
 \tag{59.7}
\]
This statement is valid in each safely pinned law
as well, with the same declared \(\Lambda_a\).

Order the cores so that an index having largest
\(\Lambda_a\) is last.  Apply (59.7) to the first
core.  Conditional on it, apply the corresponding
inequality to the next core, and continue.
Conditional-expectation consistency identifies
each one-core form obtained this way with its
form in (59.3).  At the last core its conditional
variance is already its heat-bath form, so no
additional two-block factor is needed.

Every form coefficient is bounded by the product
of the factors \(e^{\Lambda_a}\) over the first
\(N-1\) cores.  This is \(e^{B_N}\), proving
(59.5).  The proof is the finite conditional-angle
induction of V39, now applied to these actual
core marginals with V54 supplying its pinning
hypotheses.  It does not use independence of
the marginal cores.  Bound \(B_N\) by
\((N-1)\Lambda_*\) to get (59.6).
\end{proof}

The abstract induction itself is inherited, not
a newly discovered tensorization principle.
The new application is the original compact
separated-core marginal with its quantified
annular and pin-stable constants.
Leaving the largest constant last minimizes
the bound within this particular induction;
no claim of an optimal spectral constant is made.

\subsection{An exponentially large but finite core-count window}

Fix \(s\ge0\), take \(\ell\le\beta^s\) and
\(w=\lceil\beta\rceil\), and retain V53's
fixed buffers and small coarse-curvature
thresholds.  Its local annulus estimates give
\[
                      \Lambda_*\le C_s e^{-c_s\beta}
 \tag{59.8}
\]
independently of the number of other distant cores
and of the size of the ambient region.
Consequently, whenever
\[
                         N\le e^{c_s\beta/2},
 \tag{59.9}
\]
and \(\beta\) is above a fixed threshold,
\[
                            g_N\ge\tfrac12 .
 \tag{59.10}
\]
Indeed \(B_N\le C_s e^{-c_s\beta/2}\le\log2\).

This window is larger than the polynomial
ambient-volume window used in the earlier
multi-component estimate of V48--V49.
There is no discarded interface guard here.
On the other hand, at every fixed \(\beta\)
the allowed core count remains finite.
Equation (59.10) is not a thermodynamic-limit
gap.  The original coarse law has also not
been shown to satisfy the required pointwise
curvature condition automatically.

The normalization is rate one per marginal
core update.  Selecting one of the \(N\)
cores uniformly per discrete step instead
gives a gap lower bound \(g_N/N\).
Neither clock is physical Euclidean time.

\subsection{The exact lift still pays for the intervening field}

Disintegrate the full parent law as
\(\mu(dz,d\eta)=\lambda(d\eta)\pi_Z^\eta(dz)\)
and put \(F(\eta)=\pi_Z^\eta f\).
For a full-parent update set \(A\), write
\(e_A^\mu(f)=\mu\operatorname{Var}(f\mid Y_{A^c})\).
Conditional total variance gives the exact identities
\[
 \operatorname{Var}_\mu(f)=e_Z^\mu(f)+\operatorname{Var}_\lambda(F),
 \qquad
 \lambda\operatorname{Var}(F\mid\eta_{a^c})
                      =e_{Z\cup U_a}^\mu(f)-e_Z^\mu(f).
 \tag{59.11}
\]
The form difference on the right is nonnegative,
by the nested conditional projections.
Thus (59.5) implies
\[
 \operatorname{Var}_\mu(f)
 \le g_N^{-1}\sum_a e_{Z\cup U_a}^\mu(f)
                        -(N/g_N-1)e_Z^\mu(f)
 \le g_N^{-1}\sum_a e_{Z\cup U_a}^\mu(f).
 \tag{59.12}
\]
These are the same conditional-variance lift
identities used in V40 and V49, with the
normalization and negative repeated-bulk
term now left visible.

Every actual heat bath on the last right side
resamples \emph{all} of \(Z\), together with one
core.  The set \(Z\) contains the intervening
regions throughout the ambient volume.
Even when each \(U_a\) has bounded diameter
at fixed \(\beta\), \(Z\cup U_a\) need not.
Equation (59.12) therefore does not turn the
proved marginal bound into a bounded-range
full-fibre gap.

\subsection{Why a ratio-only compiler cannot simply remove \(N\)}

The following explicit nonlocal example separates
the missing spatial input from the already proved
pin-stable ratio input.  It is not a Yang--Mills
counterexample.

For odd \(n\), let \(x_i\in\{-1,1\}\),
\(M(x)=\sum_i x_i\), and
\[
 \pi_{n,h}(x)=Z_{n,h}^{-1}2^{-n}e^{h|M(x)|},
 \qquad
 Z_{n,h}=\mathbb E_0 e^{h|M|},\qquad h>0 .
 \tag{59.13}
\]
The law is strictly positive.
For a single spin against the joint rest,
its log odds difference is
\[
                  h\bigl(|m+1|-|m-1|\bigr)\in[-2h,2h].
 \tag{59.14}
\]
The single-spin rectangular bound is therefore
at most \(4h\), uniformly in \(n\).
Pins merely shift \(m\), so this bound survives
arbitrary pins of other spins.  It also survives
separate one-side weights.  The constant can
be made arbitrarily small by taking \(h\) small.
More generally a block of at most \(b\) spins
has rectangular constant at most \(4hb\):
its two sums differ by at most \(2b\),
and each absolute-value expression is
one-Lipschitz.

\begin{proposition}[Small pin-stable ratios do not imply a uniform gap]
\label{prop:f12v55-cusp}
For rate-one single-spin heat baths of (59.13),
\[
 \operatorname{gap}(\pi_{n,h})
 \le
 \frac{n e^h}{Z_{n,h}}\,
       \frac{\binom{n-1}{(n-1)/2}}{2^{n-1}}
 \le n e^h(\cosh h)^{-n}.
 \tag{59.15}
\]
For every fixed \(h>0\) this upper bound tends
to zero exponentially, despite the uniform
pin-stable rectangular bound \(4h\).
\end{proposition}
\begin{proof}
Take \(f(x)=\operatorname{sign}M(x)\).
Since \(n\) is odd and the law is invariant
under global sign reversal, \(\pi f=0\)
and \(\operatorname{Var}_\pi(f)=1\).
For the update of spin \(i\), the sign can
vary only if the remaining even number of
spins has sum zero.  In that case the two
values of spin \(i\) have equal conditional
weight \(e^h\), and the conditional variance
of \(f\) is one.  Otherwise it is zero.
The exact total heat-bath energy is therefore
the first expression in (59.15).

Finally
\[
 Z_{n,h}=\mathbb E_0 e^{h|M|}
              \ge\mathbb E_0 e^{hM}=(\cosh h)^n,
\]
and the central-binomial probability is at
most one.  The Rayleigh quotient of \(f\)
proves the spectral upper bound.
\end{proof}

The same obstruction persists for any family
of at most \(c n\) rate-one updates, each
containing at most a fixed \(b\) spins.
Such an update has zero conditional variance
of \(f\) unless the outside sum has magnitude
at most \(b\).  On that event the full sum
has magnitude at most \(2b\).  Consequently
its stationary probability is at most
\(e^{2hb}/Z_{n,h}\), and the total energy
of \(f\) is at most
\[
                         c n e^{2hb}(\cosh h)^{-n}.
 \tag{59.16}
\]
Thus merely replacing single spins by a
fixed block size does not rescue a
ratio-only volume-uniform compiler.

The potential \(h|M|\) is nonlocal.
This example does not have the bounded-degree
local completed YM action, and it does not
contradict the annular theorem.
It proves only that an argument using
pin-stable endpoint ratios while discarding
the remaining spatial interaction structure
cannot justify deleting the \(N\)-dependence.
After integrating an intervening field,
the effective marginal is nonlocal unless
one proves a suitable decay or form estimate;
that is precisely the information still
needed beyond the present compiler.

\subsection{Exact tests and the next required estimate}

The script
\begin{center}\small\ttfamily
scripts/verify\_ym\_f12\_v55\_core\_budget.py
\end{center}
checks the finite-core inequality by exact
rational positive-semidefinite elimination
for three- and four-core interacting laws,
including every pin pattern leaving at
least one unpinned core: 19 and 65 patterns,
respectively.  Their verified gap certificates
are \(65536/194481\) and
\(67108864/466948881\).
A sixteen-state core/bulk model verifies
the lifted form with the original bulk
included in each update.

For the nonlocal example the script uses
\(e^h=11/10\), verifies the exact sign
Rayleigh quotient at \(n=3,5,7\), and
checks its cross-ratio bound \(14641/10000\).
At \(n=10001\) the analytic upper bound in
(59.15) is rigorously below \(10^{-12}\),
verified by rational arithmetic without
enumerating that state space.
All checks passed.  They are tests of the
finite algebra and counterexample, not
certificates of the YM analytic constants
or the physical mass gap.

\begin{decision}[The surviving work is spatial and energetic]
The full compact marginal of separated
cores has a proved, pin-stable gap in
an exponentially large finite core-count
window.  Its exact lift remains nonlocal.
The core-count dependence cannot be
discarded using endpoint ratio bounds
alone.

The next proof must retain more of the
effective spatial interaction or localize
the intervening-field energy.  Neither
the unchanged companion notes nor the
ratio-only induction supplies that input.
The next companion audit and ten-version
direction review are both scheduled
for V60.
\end{decision}

V54 had 856511 bytes and SHA--256
\begin{center}\scriptsize\ttfamily
5C741C62F643C81B3CC8E10796D80D9EE52637DA03F539410FB01CB3FBE6F8C3.
\end{center}
Its complete body is retained apart
from the title version.  Only active
V54 is replaced by V55; archives,
companions and prior scripts remain.

\begin{firewall}[Position after V55]
The proved object is a finite-family
marginal core heat-bath gap and its
large-update lift, under the stated
annular separation and fixed small
coarse-curvature hypotheses.
There is no cofinal bounded-range
full-fibre gap, physical transfer
comparison, scale-stable continuum
construction or physical Yang--Mills
mass-gap proof.
\end{firewall}
\section{V56: complete-shell likelihood transport on an exponential radius window}
\label{sec:f12v56}

\subsection{What is inherited, and what is being supplied now}

V55 was a progress turn: it supplied the finite-core
bound and an explicit obstruction to discarding its
core-count dependence. Its current source was checked
before this continuation. No companion audit is due
until V60.

The upstream question here is spatial transport, not
another use of the same core induction. Archive
\texttt{V100\_f6}, sixth-series V18, already proves
submultiplicativity of centered shell channels and
the conditional sufficiency of a positive average
channel loss. That abstract product theorem is
inherited. It did not populate those losses in the
present nonlinear compact weak-coupling fibre.

V53 now supplies a pointwise annular likelihood
bound. This section places complete separating
records into the \emph{same} full compact law and
uses that bound on each adjacent pair. The new
conclusions are a relative-density product identity
for this application and an explicit exponential
radius window in which all its factors are small.
We do not transfer contraction numbers between
unrelated conditional laws or pin the intervening
annuli.

The full target is still existence of a nontrivial
four-dimensional quantum Yang--Mills theory with a
positive physical mass gap, as specified on the
\href{https://www.claymath.org/millennium/yang-mills-the-maths-gap/}
{Clay Mathematics Institute's official problem page}.
A finite-distance fixed-fibre statement is not that target.

\subsection{Complete records, not selected observables, give a chain}

Fix a finite canonical ambient fibre with smooth
strictly positive compact density and fixed coarse
field \(C\). Its completed potential has owner
diameter at most \(r_{\rm loc}\), uniformly in
ambient volume. This includes the full completion
and the local Jacobian factors, not just bare
plaquette terms.

Let \(b\) be the V53 repair buffer, let
\(w\ge4r_{\rm loc}+10\), and choose a fixed integer
\(t>r_{\rm loc}+2\). Set
\[
 L=2b+w+t+2,\qquad r_j=\ell+jL
       \quad(0\le j\le n-1),\qquad n\ge1 .
 \tag{60.1}
\]
Define \(X_0\) to record every independent free
group with owner in \(B_\ell\). For \(1\le j<n\),
let \(X_j\) record \emph{every} such group in
\[
                     B_{r_j}\setminus B_{r_j-t}.
 \tag{60.2}
\]
Finally let \(X_n\) record all ambient free groups
whose owners are outside
\[
                  B_{r_{n-1}+2b+w}.
 \tag{60.3}
\]
The ambient region contains \(B_{r_{n-1}+3b+w}\);
all fixed stencils embed without wrapping.
For simplicity it may contain the larger box used
in (60.14) below. Owner-strip shifts and physical
input closures are included in \(r_{\rm loc}\)
and in V53's choice of \(b\).

\begin{proposition}[Exact full-record Markov factorization]
\label{prop:f12v56-markov}
In this one ambient law, the joint record marginal
of \(X_0,\ldots,X_n\) is a Markov chain. If
\(\mu_j\) denotes its actual \(X_j\)-marginal and
\[
 K_j(x_j,dx_{j-1})
       =\mu(dx_{j-1}\mid x_j),\qquad 1\le j\le n,
 \tag{60.4}
\]
then the conditional kernel from \(X_n\) to \(X_0\)
is \(K_nK_{n-1}\cdots K_1\), with kernels applied
from right endpoint to left endpoint in that order.
Moreover \(\mu_jK_j=\mu_{j-1}\).
\end{proposition}
\begin{proof}
For each \(1\le j<n\), all earlier records lie
inside \(B_{r_j-t}\), and all later records lie
outside \(B_{r_j}\). The owner distance between
those two sets exceeds the completed interaction
diameter. No potential term can contain both an
inside and an outside free coordinate without
being excluded by that diameter bound.

Condition on the complete strip \(X_j\). Every
remaining potential term belongs to the inside
plus this fixed strip or to the outside plus this
fixed strip. Product Haar reference measure
therefore makes the inside and outside variables
conditionally independent. Terms depending only
on the strip contribute a common normalizer.
Integrating the unrecorded variables on the two
sides separately preserves this independence.

Thus, for every cut \(j\), the earlier and later
records are conditionally independent given
\(X_j\). Disintegrating successively in reverse
order gives the asserted chain factorization.
The identity for the actual marginals follows
by integrating a conditional probability.
\end{proof}

In particular, a few Wilson loops or selected
coordinates on a shell do not automatically
constitute \(X_j\). The proof uses all independent
inputs in a strip wide enough to separate the
completed interaction graph. Nor is independence
of distinct strips being asserted.

\subsection{Centered likelihood kernels multiply without volume factors}

The following calculation makes the pointwise
version of the inherited channel argument explicit.
It applies on standard Borel spaces whenever
the indicated densities and essential bounds exist.

For an adjacent pair define
\[
 r_j(x_{j-1},x_j)
   =\frac{d\mu_{j-1,j}}{d(\mu_{j-1}\otimes\mu_j)},
 \qquad h_j=r_j-1.
 \tag{60.5}
\]
Both marginal integrals of \(h_j\) vanish.
Suppose
\[
 e^{-\Lambda_j}\le r_j\le e^{\Lambda_j},\qquad
 d_j=e^{\Lambda_j}-1,\qquad
 \epsilon_j=1-e^{-\Lambda_j}.
 \tag{60.6}
\]
All marginals in these formulas belong to the
same chain, rather than separately normalized
annular reference laws.

\begin{theorem}[Pointwise and arbitrary-density chain bounds]
\label{thm:f12v56-product}
Write \(r_{0n}=d\mu_{0n}/d(\mu_0\otimes\mu_n)\).
Then
\[
 r_{0n}(x_0,x_n)-1
   =\int \prod_{j=1}^n h_j(x_{j-1},x_j)
                  \prod_{j=1}^{n-1}\mu_j(dx_j),
 \qquad
 |r_{0n}-1|\le T_n:=\prod_{j=1}^n d_j .
 \tag{60.7}
\]
For \(n=1\), this uses an empty integral. Also
\[
 \begin{split}
 \rho(X_0,X_n)&\le E_n:=\prod_{j=1}^n\epsilon_j,\\
 \sup_{v,v'}\|\mu_0(\,\cdot\mid v)-\mu_0(\,\cdot\mid v')\|_{\rm TV}
       &\le E_n .
 \end{split}
 \tag{60.8}
\]
For any \(\sigma\ll\mu_n\), let \(\sigma K\) be
the output of the composed endpoint kernel. Then
\[
 \begin{split}
 D(\sigma K\Vert\mu_0)&\le E_n D(\sigma\Vert\mu_n),\\
 \chi^2(\sigma K\Vert\mu_0)&\le E_n^2\chi^2(\sigma\Vert\mu_n),\\
 1-T_n\le \frac{d(\sigma K)}{d\mu_0}&\le1+T_n .
 \end{split}
 \tag{60.9}
\]
The last estimate is informative when \(T_n<1\)
and does not require a bound on the input density.
\end{theorem}
\begin{proof}
For a three-record chain, integrate the exact
joint density
\[
 (1+h_1(x_0,x_1))(1+h_2(x_1,x_2))
       \,d\mu_0\,d\mu_1\,d\mu_2 .
\]
The two linear terms disappear when \(x_1\)
is integrated, since both marginal integrals
of each \(h_j\) are zero. The endpoint density
minus one is \(\int h_1h_2\,d\mu_1\).
The resulting centered kernel again has both
marginal integrals zero. Induction gives (60.7).
Since \(|h_j|\le e^{\Lambda_j}-1=d_j\), its
absolute-value bound follows by integration
against probability measures. Fubini is justified
by the finite bounds.

Each adjacent conditional kernel has the
minorization
\[
                 K_j(x_j,\cdot)
                    \ge e^{-\Lambda_j}\mu_{j-1}.
 \tag{60.10}
\]
Its residual kernel preserves the indicated
marginals. Conditional Jensen therefore gives
centered \(L^2\) contraction at most
\(\epsilon_j\), exactly as in V53.
Compose these contractions in the same marginal
spaces to prove the maximal-correlation estimate.
On signed measures of total mass zero, the common
part of (60.10) vanishes and the residual Markov
kernel contracts total variation. Iterating gives
the conditional-diameter estimate.

Apply V53's one-channel entropy and squared-density
inequalities successively, using
\(\mu_jK_j=\mu_{j-1}\) at each step. This gives
the first two bounds in (60.9), allowing infinite
right sides. The output density in the last
bound is
\(\int r_{0n}(x_0,x_n)\sigma(dx_n)\).
Integrating (60.7) proves that bound directly.
\end{proof}

This is not a product of errors obtained after
successively pinning complete annuli. The adjacent
pair estimates are statements about marginals of
the original law. Arbitrary values of intermediate
records enter the disintegration automatically;
we never need to reapply V53 in a law where an
intermediate record has been frozen in a repair
buffer.

\subsection{Supplying every factor from the actual annular estimate}

For edge \(j+1\) of the chain use V53 with inner
radius \(r_j\), annular width \(w\), and the same
buffer \(b\), for \(0\le j<n\).
The record \(X_j\) is contained in its inner
endpoint. By (60.1)--(60.3), \(X_{j+1}\) is
contained in its permitted distant outer endpoint.
Marginalizing the unused endpoint coordinates
preserves V53's relative bound.

Require the V53 small coarse-curvature condition
throughout all annuli and their repair
neighbourhoods. No restriction is imposed on
the original compact free exterior.
Let
\[
 R_*=\ell+nL+3b+w+t+10,\qquad
 {\cal A}_*=\beta R_*^6 e^{-\gamma(w-C_{\rm loc})}
             +\beta^2R_*^{10}e^{-c_0\beta}.
 \tag{60.11}
\]
The area and volume bounds (57.13), applied to
each annulus separately, give
\[
 \Omega_j\le C{\cal A}_*,\qquad
 q_j\le C R_*^3 e^{-c_g\beta}.
 \tag{60.12}
\]
Whenever all \(q_j\le1/2\), (57.11) and
\(-2\log(1-q_j)\le4q_j\) give
\[
 \Lambda_{j+1}
 \le C\left[
       \beta R_*^6e^{-\gamma(w-C_{\rm loc})}
       +\beta^2R_*^{10}e^{-c_0\beta}
       +R_*^3e^{-c_g\beta}\right].
 \tag{60.13}
\]
The constant is independent of \(n\), the ambient
volume, and the original exterior. The displayed
powers of \(R_*\), however, must not be discarded.

No simultaneous good event for every annulus has
been imposed. In particular there is no extra
factor \(nq_j\) from a union bound over discarded
guards. Each guard was already removed from its
own full adjacent-pair marginal by V53.

\begin{theorem}[Exponential radius window for the full compact chain]
\label{thm:f12v56-window}
Take \(w=\lceil\beta\rceil\). Choose a fixed
\(\alpha>0\) such that
\[
 d=\min\{\gamma/2-6\alpha,\ c_0-10\alpha,\
                         c_g-3\alpha\}>0,\qquad
 \kappa=d/2 .
 \tag{60.14}
\]
Under the preceding geometric and curvature
hypotheses, if
\[
                            R_*\le e^{\alpha\beta},
 \tag{60.15}
\]
then for all sufficiently large \(\beta\),
uniformly over the admitted \(n,\ell\),
\[
 \boxed{\quad
       |r_{0n}-1|\le e^{-\kappa\beta n},\qquad
       \rho(X_0,X_n)\le e^{-\kappa\beta n}.
       \quad}
 \tag{60.16}
\]
The same upper bound applies to the conditional
TV diameter and to the KL contraction coefficient;
its square applies to the \(\chi^2\) coefficient.
\end{theorem}
\begin{proof}
For large \(\beta\), \(w-C_{\rm loc}\ge\beta/2\).
Insert (60.15) in (60.12)--(60.13). The resulting
common bound is
\[
 \Lambda_j\le C(\beta+1)^2e^{-d\beta},
 \qquad q_j\le Ce^{-(c_g-3\alpha)\beta}<1/2.
 \tag{60.17}
\]
Increase the fixed lower threshold for \(\beta\)
so that every \(\Lambda_j\le1\) and
\[
 e^{\Lambda_j}-1\le2\Lambda_j
                         \le e^{-\kappa\beta}.
 \tag{60.18}
\]
The polynomial and fixed constants are absorbed
by \(e^{(d/2)\beta}\), uniformly in \(j,n,\ell\).
Also \(1-e^{-\Lambda_j}\le\Lambda_j\).
Apply Theorem~\ref{thm:f12v56-product}.
\end{proof}

The condition on \(\alpha\) has positive choices
because the inherited constants
\(\gamma,c_0,c_g\) are positive. Their numerical
values have not been evaluated here.

For a fixed core radius \(\ell\), or for
\(\ell\le\beta^s\) with fixed \(s\), the permitted
number of annuli is of order
\(e^{\alpha\beta}/\beta\), after reducing it by
a fixed constant factor to satisfy (60.15).
In particular it is not restricted to polynomial
growth in \(\beta\).

The outer endpoint starts beyond the radial
distance from the inner box given by
\[
 D=(n-1)L+2b+w=nL-t-2.
 \tag{60.19}
\]
For sufficiently large \(\beta\), \(L\le2\beta\).
Thus (60.16) gives
\[
                |r_{0n}-1|,\ \rho(X_0,X_n)
                               \le e^{-\kappa D/2}.
 \tag{60.20}
\]
This is exponential decay in lattice distance
\emph{through the declared finite radius window}.
At distances of order \(e^{\alpha\beta}\) the
bound is doubly exponentially small as a function
of \(\beta\). This is a statement about the
displayed sufficient estimate, not a determination
of the physical correlation length.

\subsection{Endpoint observables, further distances and the remaining loss}

For \(F=F(X_0)\) and \(H=H(X_n)\) in the actual
compact \(L^2\) spaces, (60.16) gives
\[
 |\operatorname{Cov}_\mu(F,H)|
   \le e^{-\kappa\beta n}
           \sqrt{\operatorname{Var}_\mu F\,
                 \operatorname{Var}_\mu H}.
 \tag{60.21}
\]
Local functions of completed physical links may
be used only if their entire independent-input
footprints lie in the indicated endpoints.
The coarse field is still fixed.

One can also retain only a subrecord of \(X_n\),
or an inner subrecord of \(X_0\).
Relative density bounds pass to those marginals
by averaging under the product disintegration,
as in V53. Thus the estimate still bounds
observables at distances beyond the largest
certified window. It then \emph{saturates} at
the last available factor: no further decay
in that extra distance has been proved.

At each fixed \(\beta\), \(e^{\alpha\beta}\)
is finite. In the sufficient estimates (60.12),
the guard bound grows as \(R_*^3e^{-c_g\beta}\);
the other term grows as
\(\beta^2R_*^{10}e^{-c_0\beta}\).
Letting the number of full concentric shells
tend to infinity at fixed \(\beta\) eventually
invalidates (60.17). This is a failure of the
current bounds, not a proof that the actual
channel ceases to contract.

Nor can (60.20) be inserted without further
work as a pairwise influence estimate in V55's
fully pinned many-core law. An intermediate
pinned core may intersect one of these annuli
or its repairs. That law would no longer satisfy
the geometric hypotheses used above. The present
chain has no such interior pins.

\subsection{Checks, negative control and the next upstream estimate}

The script
\begin{center}\small\ttfamily
scripts/verify\_ym\_f12\_v56\_shell\_chain.py
\end{center}
checks (60.7) exactly in two nonuniform 72-state
four-link chains, including a law with nontrivial
weights at both endpoints. It verifies the product
bound and its pointwise arbitrary-input-tilt
consequence using rational arithmetic.

A separate two-rail local example records the
lower-rail endpoints but only the upper rail
at intermediate times. All adjacent recorded
pairs are independent, yet the endpoint covariance
is \(1/8\). Recording both rails at a separating
time restores conditional independence. This
tests why the complete-record hypothesis cannot
be replaced by pairwise decorrelation.

Finally a four-dimensional range-one owner graph
with 6561 vertices checks 390625 directed
neighbour relations: none jumps a full separating
strip. This is a range-geometry test, not an
enumeration of the nonlinear YM completion
stencils. The general proof uses their inherited
fixed owner-diameter bound. All listed tests pass;
none verifies the numerical analytic constants
or the physical continuum theory.

\begin{decision}[The finite-distance channel input is now populated]
The abstract shell product from f6 V18 now has
actual full compact conditional channels and
quantified losses on an exponentially large
radius window. This is a spatial advance beyond
V55's isolated core-count budget.

The next cofinal issue is explicit: control the
growing interface without charging the badness
of the entire annulus, or construct a local
energy decomposition which uses these finite
windows without pinning their repair buffers.
Iterating a scalar constant past (60.15) would
be circular bookkeeping, not a proof.
The next companion audit and ten-version
direction review remain V60.
\end{decision}

V55 had 870728 bytes and SHA--256
\begin{center}\scriptsize\ttfamily
B70EBDC3493512DFBD326383DAF7A4EAE784CEF11A8049C141B9854168ADB3D1.
\end{center}
Its complete body is retained apart from the
title version. Only active V55 is replaced by
V56. All archives, companion notes and prior
verification scripts remain.

\begin{firewall}[Position after V56]
Proved here is fixed-coarse-field full compact
shell transport with exponential lattice-distance
decay inside the stated finite radius window.
All-distance fixed-coupling decay, a cofinal
bounded-range full-fibre gap, scale-stable
coarse effective dynamics, continuum construction,
and the physical Yang--Mills mass gap remain
unproved. No auxiliary conditional channel is
identified with the physical transfer operator.
\end{firewall}
\section{V57: an area-independent interface angle and the limit of likelihood transport}
\label{sec:f12v57}

\subsection{The change of norm is part of the upstream task}

V56 made concrete progress by populating full compact
shell channels on an exponential radius window.
Its current source and digest were checked before
this continuation. The next companion audit remains
V60.

There are two distinct costs in that construction:
the sum over all interface coordinates in a
pointwise likelihood bound, and the compact
excursion remainder. They must not be conflated.
This section removes the first cost in an
\(L^2\) channel norm for the uniformly convex
reference already available in the programme.
It then applies that result to V28's exact
compact-interior/interface law.

The inherited ingredients are V33's componentwise
boxed response bound (37.3)--(37.4), V39's
conditional-angle argument (43.1)--(43.4), and
V28's measured interface Hessian (32.16).
V28 already proves gradient-dependent covariance
decay in (32.17)--(32.18); that result is not
claimed anew. The improvement here is a spatial
maximal-correlation bound for all square-integrable
endpoint observables, with no support-cardinality
or gradient prefactor, including arbitrarily
many compactly integrated blocks.

For context on classical covariance methods,
G.~Menz,
\href{https://arxiv.org/abs/1402.5160}
{\emph{A Brascamp--Lieb type covariance estimate}}
discusses covariance estimates and continuous-spin
correlation decay. No theorem about the compact
YM interface is imported from that paper.
The actual domain and normalization arguments
needed below are given explicitly.

\subsection{Eliminating intervening coordinates preserves a spatial mixed norm}

Consider a measured density
\[
 \nu(dz)=Z^{-1}e^{-V(z)}
              {\bf1}_{\mathcal Q}(z)\,dz,\qquad
 \mathcal Q=\prod_{i\in I}[-a_i,a_i],
 \qquad
            m_*\beta I\le V''\le M_*\beta I .
 \tag{61.1}
\]
Here \(m_*>0\), \(\beta\ge1\), and the potential
is smooth on a neighbourhood of the fixed box.
Scalar coordinates have owner locations with
bounded multiplicity and ball growth
\(N(r)\le C_{\rm geo}(r+1)^4\).
The Hessian has a fixed owner range \(R_0\).
Constants below may depend on these fixed data,
but not on \(|I|\), endpoint size, or \(\beta\).
Local Haar densities, when present, are already
included in \(V\).

Let \(A,B\subset I\) be disjoint coordinate
sets and put \(T=I\setminus(A\cup B)\).
Write \(x=z_A,y=z_B,u=z_T\), and define the
exact endpoint marginal potential
\[
       \Psi(x,y)
        =-\log\int_{\mathcal Q_T}e^{-V(x,y,u)}\,du .
 \tag{61.2}
\]
The integration domain is fixed. It is a box,
not the full compact group domain. An additive
constant in \(\Psi\) will not matter.

\begin{proposition}[Spatial mixed-Hessian operator bound]
\label{prop:f12v57-mixed}
There are \(C_*,a_*>0\), depending only on
the fixed data in (61.1), such that
\[
 |\partial_i\partial_j\Psi(x,y)|
                   \le C_*\beta e^{-2a_*{\rm dist}(i,j)}
             \quad(i\in A,\ j\in B).
 \tag{61.3}
\]
After increasing \(C_*\) if necessary, if
\(d={\rm dist}(A,B)\), then
\[
             \|\Psi_{AB}''(x,y)\|_{\rm op}
                   \le J(d):=C_*\beta e^{-a_*d}.
 \tag{61.4}
\]
The estimates are pointwise on the endpoint
product box and uniform after arbitrary other
coordinates have been fixed to values in their
own boxes.
\end{proposition}
\begin{proof}
For fixed \(x,y\), the conditional \(u\)-law
has Hessian \(V_{TT}''\), a principal submatrix
of \(V''\). It therefore retains the same lower
bound, upper bound and range. The coordinate
geometry is inherited even if \(T\) has holes.
The fixed-cube vector response estimate used in
V33 consequently applies to this conditional
law with the same type of uniform constants.

Differentiating the normalized integral gives
the exact identity
\[
 \partial_i\partial_j\Psi
   =\mathbb E[V_{ij}\mid x,y]
            -\operatorname{Cov}(V_i,V_j\mid x,y).
 \tag{61.5}
\]
There is no moving-domain term. The gradient
of \(V_i\) in \(u\) is supported within owner
distance \(R_0\) of \(i\), on a bounded number
of scalar coordinates. Each component is bounded
by \(M_*\beta\). The same statements hold at \(j\).
Applying (37.4) to these two scalar scores gives
\[
 |\operatorname{Cov}(V_i,V_j\mid x,y)|
                       \le C\beta e^{-c\,{\rm dist}(i,j)}.
 \tag{61.6}
\]
The fixed shifts of the two score supports
are absorbed into \(C\). The direct term in
(61.5) is zero beyond \(R_0\) and has absolute
value at most \(M_*\beta\) otherwise. This
proves (61.3), after renaming constants.
For \(T=\varnothing\), the covariance is zero
and the same conclusion follows directly.

To pass to an operator norm, sum a row of the
absolute mixed matrix. Since all \(j\in B\)
are at distance at least \(d\) from a fixed
\(i\in A\),
\[
 \sum_{j\in B}e^{-2a_*{\rm dist}(i,j)}
 \le e^{-a_*d}\sum_{j\in I}e^{-a_*{\rm dist}(i,j)}
 \le C e^{-a_*d}.
 \tag{61.7}
\]
The last series is uniformly finite by the
four-dimensional growth bound. The identical
bound holds for column sums. The matrix Schur
inequality, \(\|H\|_{\rm op}^2\le
\|\,|H|\,\|_{\rm row}\|\,|H|\,\|_{\rm col}\),
proves (61.4).

Fixing additional coordinates only takes
principal Hessians on the original product
box; it does not integrate them into an
uncontrolled new interaction. The same proof
therefore applies to every stated conditional
law. Constants do not depend on how many
coordinates are fixed.
\end{proof}

The response estimate invoked here is the
inherited one on flat product faces: diagonal
weights depending on coordinate \emph{labels}
preserve its mixed face conditions. The finite
range and row bounds needed for that argument
follow from (61.1), since each row has only
a bounded number of nonzero entries. No
convexity assertion on the full compact group
has entered this proposition.

\subsection{The angle bound needs marginal Poincare, not a new convexity assumption}

The full box law has Poincare constant at
most \((m_*\beta)^{-1}\). Applying it to
functions of \(x\) alone gives the same
Poincare constant for the \(x\)-marginal.
Likewise, for each fixed \(x\), the conditional
law of \((y,u)\) has that constant by its
principal Hessian; applying it to functions
of \(y\) gives the same conditional marginal
constant for \(y\mid x\).

These facts suffice. It is not necessary to
assert without proof that \(\Psi\) has the
same Hessian floor as \(V\).

\begin{theorem}[Area-independent all-\(L^2\) endpoint angle]
\label{thm:f12v57-angle}
Under (61.1), for arbitrary disjoint nonempty
endpoint sets \(A,B\) at owner distance \(d\),
\[
 \rho_\nu(z_A,z_B)\le
 \theta(d):=
       \frac{C_*e^{-a_*d}}
            {\sqrt{m_*^2+C_*^2e^{-2a_*d}}}
       \le \frac{C_*}{m_*}e^{-a_*d}.
 \tag{61.8}
\]
This is an all-\(L^2\) maximal correlation,
not a bound restricted to linear or Lipschitz
observables. It is uniform in endpoint
cardinalities and in the total number of
coordinates. It also holds in the conditional
box laws of Proposition~\ref{prop:f12v57-mixed}.
\end{theorem}
\begin{proof}
For bounded \(g(y)\), set \(h(x)=\mathbb E[g\mid x]\)
in the endpoint marginal. For every unit
vector \(v\) in the \(x\)-space, differentiation
of its normalized integral gives
\[
        v\cdot\nabla h
           =-\operatorname{Cov}(g,v\cdot\nabla_x\Psi\mid x).
 \tag{61.9}
\]
The conditional marginal Poincare inequality
for \(y\mid x\), and (61.4), imply
\[
 \operatorname{Var}(v\cdot\nabla_x\Psi\mid x)
       \le (m_*\beta)^{-1}
             \mathbb E[|(\Psi_{xy}'')^Tv|^2\mid x]
       \le \frac{J(d)^2}{m_*\beta}.
 \tag{61.10}
\]
Cauchy--Schwarz in (61.9), followed by the
supremum over unit \(v\) at this fixed \(x\),
gives
\[
             |\nabla h|^2
            \le\frac{J(d)^2}{m_*\beta}\operatorname{Var}(g\mid x).
 \tag{61.11}
\]
Taking that supremum, rather than summing
coordinatewise inequalities, introduces no
factor \(|A|\).

The \(x\)-marginal Poincare inequality now gives
\[
 \operatorname{Var}(h)
    \le \frac{J(d)^2}{(m_*\beta)^2}
                    \mathbb E\operatorname{Var}(g\mid x).
 \tag{61.12}
\]
Use
\(\operatorname{Var}(g)=\operatorname{Var}(h)
+\mathbb E\operatorname{Var}(g\mid x)\),
solve for \(\operatorname{Var}(h)/\operatorname{Var}(g)\),
and insert \(J(d)=C_*\beta e^{-a_*d}\).
The centered conditional-expectation norm is
bounded by the square root of that ratio.
Truncation and \(L^2\) continuity extend the
result to all square-integrable \(g\).
Symmetry gives the stated maximal correlation.
\end{proof}

The algebra in (61.9)--(61.12) is V39's
inherited angle proof. The new spatial input
is (61.4) for the actual endpoint marginal
after all intervening coordinates have been
integrated. An entrywise covariance estimate
alone would not justify replacing its matrix
norm by a dimension-free constant; (61.7)
is the required summation.

\subsection{Application to the exact compact-interior law}

Retain V28's fixed packing of separated
canonical interiors \(I_b\) and remaining
interface \(\mathcal S\). Let \(J_b\subset
\mathcal S\) contain the full independent-input
support of the conditional interior law.
The template sizes, ranges and overlaps
are fixed independently of the number of
packed blocks.

Under the precise small coarse-curvature,
template, interface-chart and coupling
hypotheses of V28, the measured interface
potential satisfies
\[
 (m/2)\beta I\le(\Phi^{\rm c})''\le K\beta I
       \quad\hbox{on }\mathcal V_{\mathcal S},
 \tag{61.13}
\]
with fixed range and row bound. Every interior
was integrated over \(G^{I_b}\) in forming
\(\Phi^{\rm c}\). Thus (61.1) applies to the
interface measure
\[
 \nu_{\mathcal S}(dv)\propto
           e^{-\Phi^{\rm c}(v)}
                 {\bf1}_{\mathcal V_{\mathcal S}}\,dv
 \tag{61.14}
\]
with \(m_*=m/2\), and the exact joint law is
\[
       \mu^\sharp(dx,dv)
           =\nu_{\mathcal S}(dv)
                       \prod_b\pi_b(dx_b\mid v_{J_b}).
 \tag{61.15}
\]
These are (43.9)--(43.12), with the
normalization unchanged. Only the interface
is chart restricted; no compact interior
is discarded.

Let \(F\) depend on interface coordinates
in \(A_0\) and on a finite collection
\(\mathcal B_F\) of compact interiors.
Let \(G\) similarly use \(B_0,\mathcal B_G\).
Assume \(\mathcal B_F\cap\mathcal B_G=\varnothing\)
and define the completed interface footprints
\[
 \widehat A=A_0\cup\bigcup_{b\in\mathcal B_F}J_b,\qquad
 \widehat B=B_0\cup\bigcup_{b\in\mathcal B_G}J_b .
 \tag{61.16}
\]
Require these footprints to be disjoint and
write \(d={\rm dist}(\widehat A,\widehat B)\).

\begin{corollary}[All-\(L^2\) compact-block clustering with restricted interface]
\label{cor:f12v57-hybrid}
For all such \(F,G\in L^2(\mu^\sharp)\),
\[
 |\operatorname{Cov}_{\mu^\sharp}(F,G)|
    \le \theta_{\mathcal S}(d)
       \sqrt{\operatorname{Var}_{\mu^\sharp}F\,
             \operatorname{Var}_{\mu^\sharp}G},
 \tag{61.17}
\]
where \(\theta_{\mathcal S}\) is (61.8)
with the interface constants. The bound is
uniform in the number of blocks, in the
sizes of the two observable supports, and
at all admitted lattice distances at fixed
sufficiently large \(\beta\).
\end{corollary}
\begin{proof}
Given \(v\), the two disjoint collections
of interiors are independent. Hence
\[
 \operatorname{Cov}_{\mu^\sharp}(F,G)
   =\operatorname{Cov}_{\nu_{\mathcal S}}(f(v),g(v)),
 \quad f=\mathbb E[F\mid v],\quad g=\mathbb E[G\mid v].
 \tag{61.18}
\]
The normalized product formula (61.15)
shows that \(f\) depends only on
\(\widehat A\) and \(g\) only on
\(\widehat B\), including any explicit
interface dependence of the observables.
Conditional Jensen gives
\(\operatorname{Var}_{\nu_{\mathcal S}}f
\le\operatorname{Var}_{\mu^\sharp}F\)
and the analogous bound for \(g\).
Apply Theorem~\ref{thm:f12v57-angle} to the
two interface footprints. No derivative
of \(F\) or \(G\) is used.
\end{proof}

If a footprint is empty, its conditional
mean is constant and the covariance is
zero. Complete physical observables require
their full independent-input footprints in
(61.16); a bare geometric support that
omits a pivot completion is insufficient.
Fixed template enlargements change distances
by only a fixed additive constant.

This is stronger in its observable class
than V28's gradient-dependent compact-block
bound (32.18). It is not a second proof of
V39's already-established auxiliary
heat-bath gap. Crucially, the all-distance
assertion belongs to \(\mu^\sharp\), not
to the unrestricted compact interface law.

\subsection{Whole-interface likelihood can fail even when the angle is small}

There is also a genuine norm distinction
behind the area cost. Consider \(M\)
independent nearest-neighbour spin chains,
each of length \(L\ge1\), with a fair first
spin and transition
\[
         \mathbb P(x_{k+1}=v\mid x_k=u)
                           =\frac{1+r uv}{2},\qquad0<r<1 .
 \tag{61.19}
\]
This is a strictly positive finite-range
local model. Parallel chains can be placed
along one coordinate direction of a
four-dimensional box, with no transverse
interaction. It is not a YM model.

Write \(U,V\in\{-1,1\}^M\) for the two
end interfaces and \(s=r^L\). Their marginal
laws are both fair product measures, and
\[
 \frac{d\mu_{UV}}{d(\mu_U\otimes\mu_V)}
                 =\prod_{a=1}^M(1+sU_aV_a).
 \tag{61.20}
\]
Consequently the minimum and maximum
relative densities are respectively
\[
                         (1-s)^M,\qquad(1+s)^M.
 \tag{61.21}
\]
Both extrema are attained on positive-mass
endpoint records.

\begin{proposition}[Area-independent angle does not imply area-independent likelihood]
\label{prop:f12v57-product}
For every \(M\ge1\), the maximal correlation
of these entire interfaces is exactly
\[
                              \rho(U,V)=s=r^L .
 \tag{61.22}
\]
Nevertheless the best minorization
\(\mu(dU\mid V)\ge a_M\mu_U(dU)\) has
\(a_M=(1-s)^M\to0\) as \(M\to\infty\).
The conditional total-variation diameter
also tends to one at fixed \(L,r\).
\end{proposition}
\begin{proof}
For each subset \(S\subset\{1,\ldots,M\}\),
the Walsh character \(\chi_S(U)=\prod_{a\in S}U_a\)
obeys
\[
                \mathbb E[\chi_S(U)\mid V]
                              =s^{|S|}\chi_S(V).
 \tag{61.23}
\]
These characters are an orthonormal basis
for the fair endpoint law. On the centered
space \(S\ne\varnothing\), the largest
singular value is \(s\), attained when
\(|S|=1\). This proves (61.22) for every
observable, not only for linear sums.
Equation (61.21) proves the sharp
minorization statement.

For the two exterior records \(V=(1,\ldots,1)\)
and \(V=(-1,\ldots,-1)\), the sum
\(Z=\sum_a U_a\) has mean respectively
\(Ms,-Ms\) and variance \(M(1-s^2)\).
Chebyshev's inequality applied to the
event \(Z>0\) therefore gives
\[
 \|\mu(\,\cdot\mid V=+)-\mu(\,\cdot\mid V=-)\|_{\rm TV}
            \ge 1-\frac{2(1-s^2)}{Ms^2}.
 \tag{61.24}
\]
The trivial upper bound one completes
the assertion as \(M\to\infty\).
\end{proof}

Thus a fixed positive lower likelihood
bound for every growing interface is
stronger than what is needed for
maximal-correlation decay. Even this
elementary local model would fail it at
fixed thickness, while retaining
exponential \(L^2\) decay in that thickness.
To keep \((1-r^L)^M\) bounded away from
zero by this method requires, in order
of magnitude, \(M r^L=O(1)\), hence
thickness growing at least logarithmically
with interface size.

This does not refute V56: that note states
its finite radius budget explicitly.
It shows why extending its whole-interface
likelihood norm unchanged is not a
necessary condition for a cofinal mass
gap argument. The new \(L^2\) target is
weaker in the appropriate direction,
while the full continuum objective is
unchanged.

\subsection{Verification and the remaining unrestricted-interface problem}

The script
\begin{center}\small\ttfamily
scripts/verify\_ym\_f12\_v57\_interface\_angle.py
\end{center}
checks the mixed-Hessian/angle algebra
using exact rational Gaussian matrices
on length-five strips of widths one,
two and three, both with and without
one pinned interior coordinate.
The unpinned squared-angle certificate
is \(1/12545\) for all three widths.
These checks test Schur integration
and the matrix form inequality, not
the nonlinear YM analytic estimates.

It checks every Walsh degree for one
through four independent channels with
\(r=1/2,L=3\), confirming maximal
correlation \(1/8\) and the exact
likelihood extrema. Finally a finite
hidden-interface model with two
two-spin block emissions checks the
all-\(L^2\) conditional lift by a
weighted positive-semidefinite form
with squared bound \(1/64\).
All checks pass.

\begin{decision}[Remove the interface restriction in an \(L^2\) norm]
The uniformly convex part no longer
pays for the interface area in its
maximal-correlation estimate. Its
application already includes arbitrarily
many full compact block interiors,
and arbitrary square-integrable
observables of those blocks.

What remains is an \(L^2\) estimate for
the unrestricted compact interface
law, including excursions outside
\(\mathcal V_{\mathcal S}\). A small
probability for one bad coordinate
does not control observables concentrated
there, and an event declaring the
entire interface good is not uniformly
probable as its area grows.

The next use of V50's bounded-range
repair norms must control that
interface contribution as an operator
or local energy term, without replacing
it by a whole-interface good event.
Neither (61.17) nor the product example
supplies this missing estimate.
The next companion audit and direction
review remain V60.
\end{decision}

V56 had 887156 bytes and SHA--256
\begin{center}\scriptsize\ttfamily
227B98CFB973518546E464F4608BD2F44260B55A0A261F0BEBDB8F125ADD37B3.
\end{center}
Its complete body is retained apart
from the title version. Only active
V56 is replaced by V57; archives,
companions and previous scripts remain.

\begin{firewall}[Position after V57]
The all-distance, area-independent
maximal-correlation bound proved here
applies to a box-convex law and to the
exact compact-interior law with its
interface still chart restricted.
It does not remove that restriction,
prove an unrestricted cofinal fibre
gap, construct the continuum theory,
or prove the physical Yang--Mills mass
gap. Auxiliary spatial conditional
channels are not physical transfer
operators.
\end{firewall}
\section{V58: local curvature defects and a repair-stabilized compact response}
\label{sec:f12v58}

\subsection{Return to the unrestricted law without deleting the bad set}

V57 was a progress turn: its new all-\(L^2\)
angle estimate removed an interface-area factor,
but kept the interface in a small chart.
Its current source and digest were checked.
This continuation uses the unrestricted
interface law throughout.

V40 already transfers buffered rarity to
that law, and V50 supplies actual bounded-range
restricted-smoothing norms. Merely repeating
their tower argument would not resolve V57's
remaining issue. The new calculation below
locates the possible negative curvature of the
exact compact interface near chart defects,
then charges that curvature to explicit local
repair forms. It constructs a coercive
\emph{modified} vector response on the full
compact law and retains the exact term that
must be removed to recover the unmodified
response.

The Bochner identity and positive-comparison
operator algebra are inherited tools. In
particular f6 V97 checks a weighted Bochner
argument, and f1 V23 studies a shifted Witten
split. Neither archive supplied the local
compact-interface curvature and repair
coefficients proved here. No physical
ground-state measure from those discussions
is identified with the present fixed-\(C\)
fibre marginal.

\subsection{The exact Haar potential has local global derivative bounds}

Use V40's unrestricted factorization
\[
 \mu(dx,dv)=\bar\nu(dv)
                \prod_b\pi_b(dx_b\mid v_{J_b}),\qquad
 \bar\nu(dv)=Z^{-1}e^{-W(v)}
                         \prod_{i\in\mathcal S}dv_i .
 \tag{62.1}
\]
Every \(dv_i\) is Haar measure on the compact
group \(G\), and every interior integral is
over its full compact domain. The fixed
packing gives the exact potential
\[
                   W(v)=U_{\mathcal S}(v)-\sum_b\log Z_b(v_{J_b}).
 \tag{62.2}
\]
All inverse-pivot factors are included in
the original local potentials. The coarse
field is fixed and satisfies the small
curvature thresholds of V28 and V50 below.
We use finite tori with the fixed stencils
embedded, or ambient volumes where every
declared repair fits. Constants are uniform
over this family.

Equip \(G^{\mathcal S}\) with a fixed product
bi-invariant metric. Its volume is a constant
multiple of product Haar, so \(W\) in (62.1)
is also its volume-relative potential.
Choose fixed left-invariant orthonormal
frames in every group factor. These provide
a global trivialization of tangent and
cotangent components; group labels, not
chart coordinates, index the components.

There is a fixed owner range \(R_0\ge1\)
containing every term in (62.2). Each
interface coordinate occurs in a bounded
number of terms of bounded cardinality.
For \(\beta\ge1\) the following bounds hold
with constants independent of volume:
\[
 \sup_i\|\nabla_iW\|_\infty\le C(\beta+1),\qquad
 \|\operatorname{Hess}W\|_{\rm row}
                           \le C(\beta+1)^2 .
 \tag{62.3}
\]
Here a block row norm sums the operator
norms of the fixed-dimensional group blocks.

\begin{proof}
For a compact block with conditional potential
\(U_b(x_b,v_{J_b})\), normalized differentiation
in a parameter direction gives
\[
 \nabla[-\log Z_b]=\pi_b(\nabla U_b),\qquad
 \operatorname{Hess}[-\log Z_b]
   =\pi_b(\operatorname{Hess}U_b)
                   -\operatorname{Cov}_{\pi_b}(dU_b,dU_b).
 \tag{62.4}
\]
These are covariant derivatives in the
interface parameter manifold. The domains
of integration are fixed. The original
local action and smooth positive Jacobian
have first and second parameter derivatives
bounded by \(C(\beta+1)\) on their compact
domains. Their fixed block sizes make the
covariance tensor in (62.4) bounded by
\(C(\beta+1)^2\). Summing only the bounded
number of terms incident to a row proves
(62.3). It also proves the stated finite
range. No off-chart convexity was used.
\end{proof}

\subsection{Convert the good-chart Hessian to intrinsic curvature}

In normal coordinates \(u_i\) about the
section value \(J_i(C)\), V28's measured
coordinate potential is
\[
              \Phi^{\rm c}(u)=W(v(u))-\sum_i\log j_G(u_i).
 \tag{62.5}
\]
Thus the Haar correction has the displayed
minus sign in \(\Phi^{\rm c}\); it must be
put back when taking an intrinsic Hessian
of \(W\).

Shrink the fixed interface chart if necessary.
There are a radius \(a>0\), a threshold
\(\beta_0\), and \(\lambda>0\), independent
of volume, such that whenever every group
belongs to its section chart cube of radius
\(a\),
\[
                \operatorname{Hess}W+\operatorname{Ric}
                                      \succeq\lambda\beta I
                 \qquad(\beta\ge\beta_0).
 \tag{62.6}
\]
This is a statement about the unrestricted
potential evaluated at those configurations,
not a change of its measure.

\begin{proof}
On V28's original interface cube,
\((\Phi^{\rm c})''\succeq(m/2)\beta I\).
The Hessian of each \(\log j_G\) is uniformly
bounded on a smaller fixed normal chart.
In those coordinates,
\[
 (\operatorname{Hess}W)_{rs}
  =(\Phi^{\rm c})_{rs}
       +(\log j_G)_{rs}-\Gamma^t_{rs}\partial_t(W\circ v).
 \tag{62.7}
\]
The Christoffel symbols of the product metric
only couple coordinates in the same group
factor and are \(O(a)\) there. By (62.3)
and uniform chart comparability, each
coordinate gradient in the last term is
\(O(\beta+1)\). Its operator norm is
therefore at most \(Ca(\beta+1)\), with no
volume factor. First choose \(a\) sufficiently
small and then \(\beta_0\) large, leaving
a fixed positive fraction of \(m\beta/2\).
The product metric is uniformly comparable
to the Euclidean metric on this cube.
Finally bi-invariant compact group metrics
have nonnegative Ricci curvature. These
facts prove (62.6) with some fixed \(\lambda>0\).
\end{proof}

\subsection{The negative curvature is supported near actual chart defects}

Let \(B_j\) be the event that \(v_j\) is
outside the smaller chart cube just chosen.
Choose \(\tau>0\) so the section ball of
radius \(\tau\) lies inside that cube.
Then
\[
              B_j\subset\{\|v_j-J_j(C)\|_{\rm HS}\ge\tau\}.
 \tag{62.8}
\]
For each interface group \(i\), set
\[
 N_i=\{j\in\mathcal S:{\rm dist}(i,j)\le R_0\},
 \qquad D_i=\bigcup_{j\in N_i}B_j,\qquad
                     n_0=\sup_i|N_i|<\infty .
 \tag{62.9}
\]
Write \(\mathcal R=\operatorname{Ric}+\operatorname{Hess}W\).

\begin{theorem}[Global local-defect curvature bound]
\label{thm:f12v58-curvature}
There is \(K_0<\infty\), independent of
volume and \(\beta\ge\beta_0\), such that
for every compact configuration \(v\) and
every tangent or cotangent vector \(z=(z_i)\),
\[
 \boxed{\quad
 \langle z,\mathcal R(v)z\rangle
       \ge\lambda\beta\sum_i|z_i|^2
                    -K_0\beta^2\sum_i1_{D_i}(v)|z_i|^2 .
 \quad}
 \tag{62.10}
\]
The indicators are local in the actual
compact configuration. No global good
event is imposed.
\end{theorem}
\begin{proof}
For this single configuration replace
every bad \(v_j\) by \(J_j(C)\), leaving
every good coordinate unchanged. Call
the resulting all-good configuration
\(\widetilde v\). Compare matrices in
the fixed group frames.

Every matrix entry of a covariant Hessian
term depends only on the coordinates of
that local term. If an entry changes,
that term contains a replaced coordinate
and each of the entry's group indices
is within \(R_0\) of such a coordinate.
Thus
\[
 \mathcal R(v)-\mathcal R(\widetilde v)
       =P_D[\mathcal R(v)-\mathcal R(\widetilde v)]P_D,
 \tag{62.11}
\]
where \(P_D\) projects onto group components
with \(D_i(v)\) true. Ricci is constant in
these invariant frames and introduces no
additional difference.

Both matrix row bounds are \(O(\beta^2)\)
by (62.3); symmetry gives the corresponding
operator bound. Hence the difference in
(62.11) is bounded below by
\(-K_0\beta^2P_D\), after fixing \(K_0\).
At \(\widetilde v\), (62.6) supplies the
floor \(\lambda\beta I\), proving (62.10).
No derivative of the discontinuous
replacement operation is taken: it is
only a pointwise matrix comparison.
\end{proof}

The power \(\beta^2\) in the debit is
retained. Replacing it by \(O(\beta)\)
would require an additional global
compact score-covariance estimate, which
is not assumed in (62.4).

\subsection{Actual interface patches refresh the entire local defect event}

For each \(j\in N_i\), use V50's canonical
repair \(I_{P_j^-}\) at threshold \(\tau\).
Assume all its comparison neighbourhoods
fit and satisfy V50's coarse-curvature
condition. Put
\[
 T_i=\mathcal S\cap\bigcup_{j\in N_i}I_{P_j^-},
 \qquad
 Q_i f=\bar\nu(f\mid v_{\mathcal S\setminus T_i}),
 \qquad
                     p_\tau=e^{-\beta\tau^2/(16A)}.
 \tag{62.12}
\]
Each \(T_i\) has bounded diameter and
cardinality, independent of volume and
\(\beta\). Its fixed size can be large.
These are heat baths of the \emph{full}
interface marginal (62.1), not heat
baths of the chart-restricted law.

V40's sigma-field transfer, now using
V50's smaller repair, gives
\[
                   Q_i1_{D_i}\le n_0p_\tau .
 \tag{62.13}
\]
Indeed the interface exterior of \(T_i\)
is outside every \(I_{P_j^-}\) used in
its definition. Apply the tower property
to (54.5), then union over \(j\in N_i\).
All internal block partition functions
remain in the conditional law. No
alignment of \(T_i\) with the packed
interiors is required.

Represent a one-form \(\omega\) by its
coefficients \(\omega_i\in\mathbb R^{\dim G}\)
in the fixed frames and define
\[
 \mathcal E_{\rm rep}(\omega)
       =\sum_i\|\omega_i-Q_i\omega_i\|_{L^2(\bar\nu)}^2,
 \qquad
       \mathsf H_{\rm rep}=\bigoplus_i(I-Q_i).
 \tag{62.14}
\]
Each \(Q_i\) acts componentwise in this
specified trivialization. It is an
orthogonal projection, so
\(0\le\mathsf H_{\rm rep}\le I\).
Overlapping repair regions do not
introduce an extra factor: this operator
is a direct sum over one-form labels.

The restricted norm in V50 and the
same two-term decomposition give
\[
 \|1_{D_i}Q_i\|_{2\to2}^2\le n_0p_\tau,\qquad
 \sum_i\bar\nu[1_{D_i}|\omega_i|^2]
       \le2\mathcal E_{\rm rep}(\omega)
                           +2n_0p_\tau\|\omega\|_2^2.
 \tag{62.15}
\]
For completeness, write
\(\omega_i=(I-Q_i)\omega_i+Q_i\omega_i\),
apply \(|a+b|^2\le2|a|^2+2|b|^2\),
and condition the second summand using
(62.13). Sum over components. This is
the inherited capacity algebra, now
applied to the particular defect set
appearing in (62.10).

\subsection{A coercive vector operator on the full compact law}

Let \(\mathsf A\) be the weighted Hodge
Laplacian on one-forms for \(\bar\nu\):
\[
 \mathsf A=d\,d_{\bar\nu}^*+d_{\bar\nu}^*d
       =\nabla_{\bar\nu}^*\nabla+\mathcal R,\qquad
                     d_{\bar\nu}^*=d^*+\iota_{\nabla W}.
 \tag{62.16}
\]
The sign convention makes \(\mathsf A\)
nonnegative. The second identity follows
from the usual Weitzenbock formula:
the added Cartan term is the Lie
derivative along \(\nabla W\), which
on one-forms is
\(\nabla_{\nabla W}+\operatorname{Hess}W\).
The drift term belongs to
\(\nabla_{\bar\nu}^*\nabla\).
For a primary-source discussion of
this geometric identity, see
P.~Petersen,
\href{https://www.math.ucla.edu/~petersen/BLWformulas.pdf}
{\emph{Demystifying the Weitzenbock curvature operator}},
Section~4.3. The present sign convention
is fixed explicitly in (62.16).

In each finite volume this is a smooth
elliptic operator on a compact manifold
without boundary, with its usual
self-adjoint realization. All form
inequalities below extend from smooth
one-forms to its form domain.
There are no artificial chart faces.

\begin{theorem}[Local-repair stabilization]
\label{thm:f12v58-stabilized}
Define the bounded nonnegative operator
and the stabilized response by
\[
 \mathsf J=2K_0\beta^2\mathsf H_{\rm rep},
 \qquad
                         \mathsf B=\mathsf A+\mathsf J .
 \tag{62.17}
\]
For sufficiently large \(\beta\), with
a threshold independent of volume,
\[
 \boxed{\quad
   \mathsf B\succeq\nabla_{\bar\nu}^*\nabla
                         +(\lambda\beta/2)I,\qquad
   \|\mathsf B^{-1}\|\le\frac{2}{\lambda\beta}.
 \quad}
 \tag{62.18}
\]
The original measure remains the full
compact \(\bar\nu\). It is the operator,
not the measure, which has been modified.
\end{theorem}
\begin{proof}
Integrate (62.10) against \(\bar\nu\)
with \(z=\omega(v)\) and use (62.15)
in the Weitzenbock form (62.16). This
gives
\[
 \langle\omega,\mathsf A\omega\rangle
 \ge\|\nabla\omega\|_2^2
   +(\lambda\beta-2K_0\beta^2n_0p_\tau)\|\omega\|_2^2
                  -2K_0\beta^2\mathcal E_{\rm rep}(\omega).
 \tag{62.19}
\]
Adding \(\mathsf J\) cancels the last
term exactly. Since
\(p_\tau=e^{-c_\tau\beta}\) with
\(c_\tau>0\), increase the fixed
threshold until
\(2K_0\beta^2n_0p_\tau\le\lambda\beta/2\).
The first assertion follows.
The spectral theorem gives the inverse
bound. The jump operator is bounded,
\(\|\mathsf J\|\le2K_0\beta^2\), so it
does not change the elliptic operator
domain or its form closure.
\end{proof}

This stabilization is not an unlabelled
scalar mass added to the problem.
Its precise local update form is
(62.14), and it vanishes on a component
already measurable outside its repair
patch. Nevertheless it is an added
term, and its removal remains necessary.

\subsection{Spatial response of the stabilized operator}

Let \(\mathsf P_i\) denote projection
onto the one-form components of group
\(i\). There is a fixed \(K_1>0\) such
that, with
\[
 a_\beta=\frac1{R_0}
                  \log\left(1+\frac{\lambda}{4K_1\beta}\right),
 \tag{62.20}
\]
one has
\[
 \|\mathsf P_i\mathsf B^{-1}\mathsf P_j\|_{2\to2}
      \le \frac4{\lambda\beta}
                         e^{-a_\beta{\rm dist}(i,j)}.
 \tag{62.21}
\]
The constant is independent of interface
area and ambient volume. The rate
\(a_\beta\) is positive at fixed
\(\beta\) but is only of order
\(\beta^{-1}\) in this estimate.

\begin{proof}
For a fixed source label \(j\), multiply
group component \(i\) by
\(e^{a_\beta{\rm dist}(i,j)}\).
These weights are constant in the
configuration variables. The product
connection preserves each cotangent
factor, so the rough Laplacian commutes
with the weights. Ricci and
\(\mathsf H_{\rm rep}\) are also diagonal
in group labels and commute with them.

Only the off-diagonal blocks of
\(\operatorname{Hess}W\) change under
conjugation. Their range is \(R_0\)
and their row and column bounds are
\(K_1\beta^2\) by (62.3). The operator
norm of the change is consequently
at most
\[
 K_1\beta^2(e^{a_\beta R_0}-1)
                                      =\lambda\beta/4 .
 \tag{62.22}
\]
Together with (62.18), a Neumann
resolvent estimate bounds the inverse
of the conjugated operator by
\(4/(\lambda\beta)\). The finite-volume
diagonal weights preserve the domain.
Apply this bound to a vector supported
on label \(j\), then undo the weights
on label \(i\), to obtain (62.21).
\end{proof}

The locality used here concerns the
component labels. A repair averages
over several configuration variables,
but does not change the one-form label
to which it is applied. This distinction
is why it commutes with the conjugating
weights even for overlapping patches.

\subsection{The exact return term has not become small}

Define the bounded positive operator
\[
                 \mathsf R
                   =\mathsf B^{-1/2}\mathsf J\mathsf B^{-1/2}.
 \tag{62.23}
\]
Since \(\mathsf A=\mathsf B-\mathsf J
\succeq0\), quadratic-form calculus gives
\[
 0\preceq\mathsf R\preceq I,\qquad
 \mathsf A=\mathsf B^{1/2}(I-\mathsf R)\mathsf B^{1/2}.
 \tag{62.24}
\]
This is the inherited positive-comparison
operator identity, with the actual
compact-interface repair operator now
specified. A volume-independent bound
\[
                         \mathsf R\preceq(1-\varepsilon)I,
                         \qquad\varepsilon>0,
 \tag{62.25}
\]
would imply
\(\mathsf A\succeq\varepsilon\lambda\beta I/2\).
The scalar weighted diffusion would
then have the same lower gap bound:
apply \(\mathsf A\,df=d(\mathsf L f)\)
to nonconstant scalar eigenfunctions
of its nonnegative weighted Laplacian
\(\mathsf L\).

Equation (62.25) has not been proved.
From the displayed norm estimates alone
one only obtains
\[
             \|\mathsf R\|
                \le\frac{4K_0}{\lambda}\beta,
 \tag{62.26}
\]
besides the non-strict upper bound one
in (62.24). The small factor \(p_\tau\)
has already paid the scalar remainder
in (62.19); it is not a small coefficient
of \(\mathsf J\).

Equivalently, solving
\(\mathsf A\omega=f\) gives the exact
equation
\[
       \omega=\mathsf B^{-1}f+\mathsf B^{-1}\mathsf J\omega
 \tag{62.27}
\]
whenever the unmodified equation is
solvable in the stated domain.
Discarding the second term would
replace the desired operator by the
stabilized one. Spatial decay of the
first term does not justify that
discarding. A return estimate such
as (62.25), or a weaker sufficient
source-specific estimate, is still
needed.

\subsection{Checks and the new unrestricted-interface input}

The script
\begin{center}\small\ttfamily
scripts/verify\_ym\_f12\_v58\_curvature\_repair.py
\end{center}
uses exact rational arithmetic on an
eight-state model with three vector
labels. It checks each restricted
capacity form and the assembled
24-dimensional stabilized form.
Its conditional defect probabilities
are \(199/10000,\ 29701/1000000,\
199/10000\), and its exact lower
form bound is \(220299/125000\).

A negative control leaves a bad
neighbour outside a single-site repair;
the conditional bad probability is
then one, correctly invalidating the
small restricted norm. All checks pass.
These are finite form-algebra tests.
The synthetic vector matrix is not
asserted to be a geometric Hodge
Laplacian, and the tests do not verify
the YM analytic constants or the
cofinal return bound.

\begin{decision}[A global compact comparison now has a local debit]
The new input is on the unrestricted
interface itself: a pointwise local
curvature-defect inequality and a
coercive, spatially decaying vector
response after an explicitly identified
bounded-range repair term is added.
No simultaneous good-interface event
has been used.

The next task is to estimate the
return of this same repair term in
the actual response, rather than
assuming the stabilized inverse is
the covariance inverse. This is the
remaining difference between the
constructed comparison and an
unrestricted estimate.
The next companion audit and
ten-version review remain V60.
\end{decision}

V57 had 905483 bytes and SHA--256
\begin{center}\scriptsize\ttfamily
912C7F103117389BB666DF605ADC525375BD6F9222A2D2D87C0225CDB025A453.
\end{center}
Its complete body is retained apart
from the title version. Only active
V57 is replaced by V58; all archives,
companions and earlier scripts remain.

\begin{firewall}[Position after V58]
The modified response \(\mathsf B\)
has a proved uniform lower bound
and spatial inverse estimate on the
full compact interface law. No such
uniform bound has been proved for
the unmodified \(\mathsf A\).
The repair return, unrestricted
cofinal fibre gap, continuum
construction and physical Yang--Mills
mass gap remain open. Neither
weighted diffusion operator is
identified with the physical
transfer Hamiltonian.
\end{firewall}
\section{V59: repair curl and the exact scalar return form}
\label{sec:f12v59}

\subsection{Audit the proposed restriction before using it}

V58 made concrete progress: it constructed a
coercive comparison on the unrestricted compact
interface, with an explicit local repair term.
Its current source and digest were checked.
No companion audit is due until V60.

The tempting next step is to work only with
gradients, since covariance uses exact one-forms.
That reduction is legitimate for the original
weighted Hodge operator, but is not automatic
for the repair-modified operator. This section
computes the repair's failure to preserve exact
forms, identifies the correct compressed
operator, and obtains the actual scalar
derivative-energy inequality left by V58.

The general warning is not new. Archive f1
V12 proves that chart multiplication can
create curl in an exact form. The calculation
below concerns a different operator: V58's
componentwise conditional-expectation repair.
It also gives an explicit nonlocality test
for the exact-form projection, rather than
assuming that projection preserves V58's
spatial label bound.

\subsection{The finite-volume exact-form space}

Continue with the unrestricted compact
interface measure \(\bar\nu\), its scalar
nonnegative weighted Laplacian
\(\mathsf L=d_{\bar\nu}^*d\), and its
one-form weighted Hodge Laplacian \(\mathsf A\).
Assume \(G\) is connected. In each fixed
finite volume, smooth positive density on
the compact connected manifold gives a
positive first nonconstant scalar eigenvalue
\(g_N\). No lower bound uniform in \(N\)
is assumed.

Let
\[
 \mathcal E=\overline{\operatorname{ran}d}
       \subset L^2(T^*G^{\mathcal S},\bar\nu),
 \qquad
 U=d\,\mathsf L^{-1/2}:L^2_0(\bar\nu)\longrightarrow\mathcal E .
 \tag{63.1}
\]
The inverse power acts only on centered
functions in this finite volume.
The map \(U\) is an isometry onto \(\mathcal E\):
\[
 \|d\mathsf L^{-1/2}f\|_2^2
      =\langle\mathsf L^{-1/2}f,
                       \mathsf L\mathsf L^{-1/2}f\rangle
      =\|f\|_2^2 .
 \tag{63.2}
\]
Spectral expansion shows that \(\mathcal E\)
reduces \(\mathsf A\) and
\[
                   \mathsf A U=U\mathsf L .
 \tag{63.3}
\]
Indeed \(\mathsf A\,df=d(\mathsf Lf)\);
apply this to the scalar eigenbasis and
then close the operators. This is the
usual exact-form supersymmetry, not a
new mass-gap theorem. A primary-source
discussion of Witten supersymmetry and
form inequalities is D.~Le Peutrec,
\href{https://arxiv.org/abs/1607.08714}
{\emph{On Witten Laplacians and Brascamp--Lieb's
inequality on manifolds with boundary}}.
The finite-volume identifications needed
here are fixed by (63.1)--(63.3).

\subsection{The repair of a gradient contains a conditional score}

Use V58's patches \(T_i\), each containing
group \(i\), and let \(X_{i,a}\) be a
fixed invariant orthonormal frame field
in that factor. Put
\[
 Q_i=\bar\nu(\,\cdot\mid v_{\mathcal S\setminus T_i}),
 \qquad
 c_\beta=2K_0\beta^2,\qquad b_\beta=\lambda\beta/2 .
 \tag{63.4}
\]
Thus \(\mathsf J=c_\beta\mathsf H_{\rm rep}\)
and \(\mathsf B=\mathsf A+\mathsf J\succeq b_\beta I\).

For smooth scalar \(f\), Haar integration
by parts in group \(i\), with the entire
patch conditional normalized, gives
\[
 \begin{split}
 Q_iX_{i,a}f&=Q_i[f\,X_{i,a}W]\\
      &=Q_i[(f-Q_if)X_{i,a}W],
 \qquad Q_iX_{i,a}W=0 .
 \end{split}
 \tag{63.5}
\]
To check the signs, integrate
\(X_{i,a}(fe^{-W})\) over the compact
updated groups. The Haar divergence of
an invariant vector field is zero and
there is no boundary term. The same
calculation with \(f=1\) gives the last
identity. The conditional normalizer
depends only on variables outside \(T_i\)
and is not differentiated in this step.

Define the nonnegative scalar derivative
repair form
\[
 \mathcal D(f)=\sum_{i,a}
                  \|(I-Q_i)X_{i,a}f\|_2^2 .
 \tag{63.6}
\]
Orthogonality of each conditional projection
gives the exact identity
\[
 \mathcal D(f)
   =\|df\|_2^2-
       \sum_{i,a}\left\|
               Q_i[(f-Q_if)X_{i,a}W]\right\|_2^2,
 \qquad 0\le\mathcal D(f)\le\|df\|_2^2 .
 \tag{63.7}
\]
There is no sum of stationary bad
probabilities in (63.7). The conditional
score term is part of the operator,
including for functions concentrated
on rare configurations.

\subsection{Componentwise repair need not preserve exact forms}

\begin{proposition}[A compact product example with nonzero repair curl]
\label{prop:f12v59-curl}
Even for a smooth strictly positive
product measure, and disjoint patches
containing their own group labels,
\(\mathsf H_{\rm rep}df\) need not be
closed. Consequently \(\mathsf J\) and
\(\mathsf B\) need not preserve
\(\mathcal E\).
\end{proposition}
\begin{proof}
First take two circle variables \(x,y\),
each with normalized density
\[
             (1+s\cos x)\frac{dx}{2\pi},
                 \qquad 0<s<1,
 \tag{63.8}
\]
and independent copies. Let each patch
update its own variable and take
\(f(x,y)=\sin x\sin y\). Since
\(\mathbb E[\cos x]=s/2\), the repaired
gradient is the one-form
\[
 \omega=(\cos x-s/2)\sin y\,dx
                     +\sin x(\cos y-s/2)\,dy .
 \tag{63.9}
\]
Its exterior derivative is
\[
       d\omega=\frac{s}{2}(\cos y-\cos x)\,dx\wedge dy,
 \tag{63.10}
\]
which is not zero.

The phenomenon is not specific to an
abelian group. On any connected compact
group admitting a nonconstant smooth
function \(h\) and invariant field \(X\)
with nonconstant \(\eta=Xh\), use
\[
                 \nu_t(dg)=Z_t^{-1}e^{t\eta(g)}\,dg,\qquad t>0.
 \tag{63.11}
\]
Haar integration gives \(\int\eta\,dg=0\).
The function \(\log Z_t\) has strictly
positive second derivative
\(\operatorname{Var}_{\nu_t}\eta\), so
\(A_t=\nu_t\eta>0\) for \(t>0\).
On two independent copies, let
\(f=h(g_i)h(g_j)\). Choose disjoint
repair patches containing \(i,j\)
respectively; additional independent
factors in either patch do not affect
the computation. The two relevant
components of \(\mathsf H_{\rm rep}df\)
are
\[
        (\eta(g_i)-A_t)h(g_j),\qquad
                  h(g_i)(\eta(g_j)-A_t).
 \tag{63.12}
\]
The fields \(X_i,X_j\) commute because
they act in different factors. Evaluating
the exterior derivative on them gives
\(A_t(\eta(g_j)-\eta(g_i))\ne0\).
This construction is available in
particular on compact semisimple groups.
\end{proof}

These are operator counterexamples,
not counterexamples to YM mixing.
They show that the structural definition
of the repair alone gives no exact-form
invariance. The actual YM operator would
need a separate invariance argument;
none has been supplied.

\subsection{Correct compression and the scalar fourth-order inequality}

Let \(P=UU^*\) be the orthogonal
projection onto \(\mathcal E\).
The correct compressed comparison is
the operator on \(\mathcal E\) associated
with the restricted form
\[
             \mathsf B_{\mathcal E}
                       =\mathsf A_{\mathcal E}+P\mathsf J P .
 \tag{63.13}
\]
The first summand is the genuine
restriction of \(\mathsf A\); the second
is a bounded form compression.
V58 gives \(\mathsf B_{\mathcal E}\succeq b_\beta I\).
It does not follow that
\(\mathsf B_{\mathcal E}^{-1}=P\mathsf B^{-1}P\).
In general the latter compression also
allows excursions through \(\mathcal E^\perp\).

On centered scalar \(L^2\), define
\[
 \mathsf C=U^*\mathsf J U,\qquad
                  0\preceq\mathsf C\preceq c_\beta I .
 \tag{63.14}
\]
The exact equivalence is
\[
           U^*\mathsf B_{\mathcal E}U=\mathsf L+\mathsf C .
 \tag{63.15}
\]
If \(\mathsf D\) denotes the operator
interpreted through the form (63.6),
then, in form notation,
\[
       \mathsf C=c_\beta\mathsf L^{-1/2}
                       \mathsf D\mathsf L^{-1/2}.
 \tag{63.16}
\]
Boundedness in (63.14) follows from
\(\mathcal D(f)\le\|df\|^2\), and
does not require a volume-uniform
bound on \(\mathsf L^{-1/2}\) separately.

\begin{theorem}[Actual scalar form left by stabilization]
\label{thm:f12v59-scalar}
For every smooth scalar \(f\) in the
unrestricted interface law,
\[
 \boxed{\qquad
       \|\mathsf Lf\|_2^2+c_\beta\mathcal D(f)
                                  \ge b_\beta\|df\|_2^2 .
       \qquad}
 \tag{63.17}
\]
This extends to the corresponding
closed scalar form domain.
\end{theorem}
\begin{proof}
Apply \(\mathsf B\succeq b_\beta I\)
to \(df\). Supersymmetry and adjunction
give
\[
 \langle df,\mathsf A\,df\rangle
      =\langle df,d\mathsf Lf\rangle
      =\|\mathsf Lf\|_2^2 .
 \tag{63.18}
\]
By construction
\(\langle df,\mathsf J\,df\rangle
=c_\beta\mathcal D(f)\).
Substitution proves (63.17); closure
gives the domain extension.
\end{proof}

Unlike the simple inequality
\(\mathsf A\succeq0\), (63.17) retains
V58's positive \(b_\beta\) and its
specified local derivative penalty.
But it is not yet the scalar
Poincare inequality: that penalty
has not been controlled by
\(\|\mathsf Lf\|_2^2\).

\subsection{The exact compressed return and a circularity test}

The correctly compressed return is
\[
 \mathsf R_{\mathcal E}
      =(\mathsf L+\mathsf C)^{-1/2}
                         \mathsf C(\mathsf L+\mathsf C)^{-1/2}.
 \tag{63.19}
\]
It lives on centered scalar \(L^2\)
through the unitary map \(U\).
It is not the unqualified compression
\(P\mathsf R P\) of V58's return.

\begin{proposition}[Exact derivative-energy criterion]
\label{prop:f12v59-return}
For a given \(M_\beta\ge0\), the inequality
\[
 c_\beta\mathcal D(f)\le
                           M_\beta\|\mathsf Lf\|_2^2
                   \quad\hbox{for all smooth }f
 \tag{63.20}
\]
is equivalent to
\[
               \mathsf R_{\mathcal E}
                      \preceq\frac{M_\beta}{1+M_\beta}I .
 \tag{63.21}
\]
Either implies
\[
                       g_N\ge\frac{b_\beta}{1+M_\beta}.
 \tag{63.22}
\]
Conversely, if \(g_N\ge g>0\), then
(63.20) holds with \(M_\beta=c_\beta/g\).
Thus, at fixed \(\beta\), existence
of a volume-uniform finite constant
in (63.20) is equivalent to a
volume-uniform scalar diffusion gap.
\end{proposition}
\begin{proof}
Inequality (63.21) is equivalent,
by the positive congruence in (63.19),
to \(\mathsf C\preceq M_\beta\mathsf L\).
Insert (63.16) and test on
\(\mathsf L^{1/2}f\) to obtain (63.20).
Conversely such vectors form a form
core, so (63.20) gives the operator
inequality and hence (63.21).

Combining (63.20) with (63.17) gives
\[
       (1+M_\beta)\|\mathsf Lf\|_2^2
                                  \ge b_\beta\|df\|_2^2 .
 \tag{63.23}
\]
Apply this to each nonconstant scalar
eigenfunction and divide by its
positive eigenvalue to prove (63.22).
For the converse, spectral expansion
gives \(\|\mathsf Lf\|_2^2\ge
g\|df\|_2^2\); combine this with
\(\mathcal D(f)\le\|df\|_2^2\).
\end{proof}

This equivalence is a diagnostic,
not a proof of (63.20). Simply
postulating that the compressed
return is uniformly below one would
restate the missing scalar gap.
Any next proof must estimate
\(\mathcal D\) from additional local
information; exact-form compression
supplies no such estimate by itself.

For reference, the finite-volume
margin obeys
\[
 \frac{g_N}{g_N+c_\beta}
       \le 1-\|\mathsf R_{\mathcal E}\|,
 \qquad
 b_\beta(1-\|\mathsf R_{\mathcal E}\|)
       \le g_N .
 \tag{63.24}
\]
The first follows from
\(\mathsf C\le c_\beta I\le
(c_\beta/g_N)\mathsf L\); the second
follows from (63.15) and its lower
bound \(b_\beta\). These relations
also display why a finite-volume
margin alone is not the cofinal estimate.

\subsection{Exact-form projection can be spatially nonlocal even in a product law}

One further obstruction concerns
the proposed use of V58's spatial
inverse after projection. In the
full one-form space the projection
has the finite-volume formula
\[
                         P=d\mathsf L^{-1}d_{\bar\nu}^*,
 \tag{63.25}
\]
understood first on smooth forms
and then by its bounded extension.

Consider product Haar measure on
copies of the same compact connected
group, with no interactions between
owners. Let \(\varphi\) be a real
nonconstant one-site Laplace
eigenfunction with eigenvalue
\(\ell_0>0\) and normalized
\(\|\varphi\|_2=1\).
For two arbitrarily separated
labels \(i,j\), put
\[
       f=\varphi(g_i)\varphi(g_j),\qquad
       \alpha_j=\varphi(g_i)\,d_j\varphi(g_j).
 \tag{63.26}
\]
This one-form has only a \(j\)-component,
although its coefficient depends also
on the distant variable \(i\).
Direct calculation gives
\[
 \mathsf Lf=2\ell_0f,\qquad
 d^*\alpha_j=\ell_0f,\qquad
                       P\alpha_j=\tfrac12df .
 \tag{63.27}
\]
The norm of \(\alpha_j\) is
\(\sqrt{\ell_0}\), and the norm
of the \(i\)-component of \(P\alpha_j\)
is \(\sqrt{\ell_0}/2\). Consequently
\[
               \|\mathsf P_iP\mathsf P_j\|_{2\to2}\ge\tfrac12
 \tag{63.28}
\]
for arbitrary owner separation.

The product scalar law has a
volume-independent positive gap;
there is no failure of mixing here.
Equation (63.28) instead rules out
an assumed exponential owner-label
bound on the full operator \(P\).
It is the spectator issue already
encountered in V52, now computed
for the canonical exact-form
projection itself. Thus V58's estimate
for \(\mathsf P_i\mathsf B^{-1}\mathsf P_j\)
cannot simply be transported through
\(P\) in the same norm. More local
source restrictions would need
their own invariant or propagated
support estimates.

\subsection{Exact checks and the next upstream step}

The script
\begin{center}\small\ttfamily
scripts/verify\_ym\_f12\_v59\_exact\_return.py
\end{center}
checks a three-component/two-scalar
matrix model with exact intertwining
\(A d=dL\), a positive repair that
does not preserve \(\operatorname{ran}d\),
the fourth-order identity, and a
compressed return margin \(1/3\).
It verifies that inverse compression
strictly differs from the inverse
of the compressed comparison.

Rational trigonometric-polynomial
differentiation checks (63.10) at
\(s=1/3,1/2,2/3\). A two-channel
projection check verifies the
coefficient \(1/2\) in (63.27).
All tests pass. They test these
finite identities and obstructions,
not (63.20) for the actual YM law.

\begin{decision}[Do not close the programme by renaming its missing gap]
The correct gradient reduction now
has a concrete scalar debit:
the patch variances of the invariant
derivatives in (63.6). The repair
does not automatically preserve the
source space, and the source projection
does not have the required full-space
label locality. These remove two
possible but invalid shortcuts from
the V58 comparison.

The next work must go upstream of
the abstract return criterion:
derive a local estimate for derivative
fluctuations or for the escape from
chart-defect regions that implies
(63.20), or choose a different
scalar comparison whose remainder
can actually be bounded. Stationary
rarity, positive stabilization, and
formal exact-form compression do
not supply that estimate.
V60 must include both the companion
audit and the ten-version direction
review before selecting that route.
\end{decision}

V58 had 924623 bytes and SHA--256
\begin{center}\scriptsize\ttfamily
D081A40BFE8B817F4F932DE4B4101790C8887476309ADD96AB65B98B8BB6A6E1.
\end{center}
Its complete body is retained apart
from the title version. Only active
V58 is replaced by V59; archives,
companions and prior scripts remain.

\begin{firewall}[Position after V59]
The repair curl, correct compressed
operator, unrestricted scalar
fourth-order inequality, and precise
return/gap equivalence are proved.
The uniform derivative-repair bound
(63.20) is not proved. Neither an
unrestricted cofinal scalar gap nor
the continuum physical Yang--Mills
mass gap follows from this audit.
The active objective remains the
full continuum construction and
physical mass gap.
\end{firewall}
\section{V60: companion audit, direction review and an unrestricted clock comparison}
\label{sec:f12v60}

\subsection{The compulsory companion audit}

The active SM--Einstein, NS and parallel-YM
files were inspected again. Their versions,
sizes and SHA--256 values remain unchanged
from V55:
\begin{center}
\begin{tabular}{lll}
\toprule
Programme & Version & Bytes\\
\midrule
SM--Einstein & V72 & 607372\\
Navier--Stokes & V26 & 278283\\
Parallel YM & V5 & 54174\\
\bottomrule
\end{tabular}
\end{center}
Their digests are respectively
\begin{center}\scriptsize\ttfamily
F44F9745D43379E1672C5550411B58E97195A4C9278592D221DF5D3F644DC1F5\\
82C887F8C2C69E4F28FAD7C3DC8DE5B9099CB5A4BF431EBA1FFF3E91935978AD\\
306F1BB84CFA8A8D47EA569EC17E9545F9867C8D32B548EC9D9C444B4F3C0512.
\end{center}
SM V72 proves a fixed-mode, fixed-time
many-copy limit, not its physical
finite-species or ultraviolet limit.
NS V26 supplies rank-one norm results,
not full stability or global regularity.
Parallel YM V5 rejects volume-uniform
tightness of an extensive integer charge
under tensorization. Its retained local
and compact-sector distinctions remain
relevant, but it supplies no interface
energy estimate. Its internal stop wording
does not supersede the user's active goal.

\subsection{Ten-version direction review}

V51--V52 distinguish a successful Gaussian
innovation estimate from an invalid
full-space operator-norm summation in
the nonlinear compact problem. V53--V54
give full compact relative annular
control with carefully buffered endpoints
and safe endpoint pins. V55 retains
the finite core-count cost and rules
out deleting it using ratios alone.
V56 composes actual complete-shell
channels through an exponentially
large but finite radius window.

V57 removes the area factor in a
maximal-correlation norm for the
convex reference and the exact
compact-interior law with its interface
still restricted. V58 instead works
on the unrestricted interface and
localizes its curvature defect, but
adds an explicit repair operator to
obtain coercivity. V59 establishes
that removing this term on gradients
requires an unproved scalar
derivative-energy estimate. It also
rules out automatic gradient-space
invariance and label-locality of the
source projection.

V59 was therefore a progress turn:
it changed the admissible next argument
by proving those obstructions and
identifying the scalar debit. It did
not prove a positive unrestricted
gap. The next step must not consist
only of introducing another return
operator whose strict contraction is
equivalent to the desired gap.

Before selecting a local scalar
criterion, there is one necessary
compatibility check. V59 uses weighted
diffusion, whereas V40 uses local
heat-bath updates. The comparison
below connects them on the same
unrestricted interface law, with
explicit volume-independent constants
at fixed \(\beta\). This is not a
comparison to physical Euclidean time.

\subsection{Two auxiliary clocks on the identical compact law}

Use the exact Haar-relative interface
potential \(W\) and full measure
\(\bar\nu\) from (62.1)--(62.2).
For each whole group coordinate \(i\)
define
\[
 L_i=-\Delta_i+\langle\nabla_iW,\nabla_i\,\cdot\,\rangle,
 \qquad L=\sum_iL_i,\qquad
 P_i f=\bar\nu(f\mid v_{i^c}),\qquad
 H=\sum_i(I-P_i).
 \tag{64.1}
\]
The \(P_i\)'s here update just group
\(i\), not V58's larger repair patches.
Each \(L_i\) is nonnegative and
self-adjoint in the full weighted
space on its natural domain; its
conditional mean is zero:
\[
                          P_iL_i u=0.
 \tag{64.2}
\]
Equation (64.2) is compact Haar
integration by parts with the
normalized conditional density.
There are no chart boundaries.

Write the two Dirichlet forms and
their finite-volume gaps as
\[
 \begin{aligned}
 E_L(f)&=\sum_i\bar\nu|\nabla_i f|^2,
 &g_N&=\inf_{f\perp1}\frac{E_L(f)}{\|f\|_2^2},\\
 E_H(f)&=\sum_i\|(I-P_i)f\|_2^2,
 &h_N&=\inf_{f\perp1}\frac{E_H(f)}{\|f\|_2^2}.
 \end{aligned}
 \tag{64.3}
\]
The heat-bath clock is rate one per
group. At a fixed finite volume,
\(g_N>0\) by compact connectedness
and positive smooth density.
No cofinal lower bound is assumed.

Let \(W_{ij}^{\rm Hess}\) denote the
mixed group block of the intrinsic
Hessian. Define a uniform bound
\[
 \sup_v\sup_i\sum_{j\ne i}
            \|W_{ij}^{\rm Hess}(v)\|_{\rm op}
                                     \le K_\beta .
 \tag{64.4}
\]
V58 proves such a bound with
\(K_\beta=K(\beta+1)^2\) for the
actual compact interface. Its fixed
packing makes the range and row
degree bounded independently of
the number of blocks.

\subsection{The local square sum retains every mixed curvature term}

\begin{proposition}[Local diffusion square estimate]
\label{prop:f12v60-squares}
For every smooth scalar \(u\),
\[
              \sum_i\|L_i u\|_2^2
                       \le\|Lu\|_2^2+K_\beta E_L(u).
 \tag{64.5}
\]
This does not assume that the \(L_i\)
commute or that different groups are
independent.
\end{proposition}
\begin{proof}
Apply the weighted Bochner identity
to the one-group conditional law,
then integrate its exterior. It gives
\[
 \|L_i u\|_2^2
 =\bar\nu\!\left[
       |\operatorname{Hess}_{ii}u|^2
       +\operatorname{Ric}_i(\nabla_i u,\nabla_i u)
       +W_{ii}^{\rm Hess}(\nabla_i u,\nabla_i u)\right].
 \tag{64.6}
\]
The product-manifold Bochner identity
for \(L\) has all the diagonal terms
in (64.6), together with the mixed
Hessian squares and mixed curvature
terms. Subtracting the sum over \(i\)
therefore yields the exact identity
\[
 \begin{split}
 \|Lu\|_2^2-\sum_i\|L_i u\|_2^2
   ={}&\sum_{i\ne j}\bar\nu
                  |\operatorname{Hess}_{ij}u|^2\\
     &+\sum_{i\ne j}\bar\nu
                  W_{ij}^{\rm Hess}(\nabla_i u,\nabla_j u).
 \end{split}
 \tag{64.7}
\]
Both sums use ordered pairs. Product
Ricci has no cross-factor blocks, so
all Ricci terms have canceled.

The first sum on the right is
nonnegative. The absolute value of
the second is bounded by
\[
 \sum_{i\ne j}\bar\nu
    \|W_{ij}^{\rm Hess}\|_{\rm op}
                     |\nabla_i u|\,|\nabla_j u|
                  \le K_\beta\sum_i\bar\nu|\nabla_i u|^2.
 \tag{64.8}
\]
Use \(ab\le(a^2+b^2)/2\), Hessian
symmetry, and (64.4). This proves
(64.5). In particular, replacing
the mixed curvature in (64.7) by
zero would not be justified in
the interacting law.
\end{proof}

\subsection{A diffusion gap gives an actual heat-bath gap}

\begin{theorem}[Volume-independent diffusion-to-heat-bath conversion]
\label{thm:f12v60-convert}
Under (64.4),
\[
                     \boxed{\quad
                     h_N\ge\frac{g_N}{g_N+K_\beta}.
                     \quad}
 \tag{64.9}
\]
The statement includes \(K_\beta=0\),
when its right side is one.
\end{theorem}
\begin{proof}
For smooth centered \(f\), solve the
finite-volume centered equation
\(Lu=f\). Elliptic regularity makes
\(u\) smooth. Using (64.2),
\[
 \begin{split}
 \|f\|_2^2
     &=\langle f,Lu\rangle
       =\sum_i\langle(I-P_i)f,L_i u\rangle\\
     &\le E_H(f)^{1/2}
                         \left(\sum_i\|L_i u\|_2^2\right)^{1/2}.
 \end{split}
 \tag{64.10}
\]
Proposition~\ref{prop:f12v60-squares}
and the scalar spectral theorem give
\[
 \sum_i\|L_i u\|_2^2
   \le\|f\|_2^2+K_\beta\langle f,L^{-1}f\rangle
   \le(1+K_\beta/g_N)\|f\|_2^2 .
 \tag{64.11}
\]
For nonzero \(f\), square (64.10),
cancel \(\|f\|_2^2\), and insert
(64.11). This proves the heat-bath
Poincare inequality with constant
\(1+K_\beta/g_N\).
Smooth functions are dense in
\(L^2(\bar\nu)\); all finite-volume
heat-bath projections are bounded.
The inequality therefore extends to
every centered \(L^2\) function.
\end{proof}

The Poisson solution is used only
in a finite volume where \(g_N>0\).
All dependence on that gap is kept
in (64.11), so this argument has
not silently assumed a uniform
bound on the inverse.

\subsection{The opposite direction uses a one-group compact comparison}

Let \(\lambda_G>0\) be the scalar
Poincare gap of normalized Haar
measure on \(G\) with the chosen
metric. Suppose
\[
 \sup_{i,v_{i^c}}\operatorname{osc}_{v_i}
                                  W(v_i,v_{i^c})\le D_\beta .
 \tag{64.12}
\]
V58's global first-derivative bound
gives \(D_\beta\le D(\beta+1)\)
by integration along one-group
geodesics. Notice that this uses
the \(O(\beta+1)\) gradient bound,
not the \(O((\beta+1)^2)\) mixed
Hessian bound.

For the full one-group conditional,
comparison with Haar gives
\[
 \bar\nu\operatorname{Var}(f\mid v_{i^c})
          \le\frac{e^{D_\beta}}{\lambda_G}
                                   \bar\nu|\nabla_i f|^2 .
 \tag{64.13}
\]
Indeed its largest and smallest
Haar-relative densities have ratio
at most \(e^{D_\beta}\); compare
the variance, then the Haar
Poincare inequality, and then
the gradient energy back to the
conditional law. Summing gives
\[
 E_H(f)\le C_{\rm cond}(\beta)E_L(f),
 \qquad C_{\rm cond}(\beta)=e^{D_\beta}/\lambda_G.
 \tag{64.14}
\]
It follows that
\[
                         g_N\ge h_N/C_{\rm cond}(\beta).
 \tag{64.15}
\]

\begin{corollary}[The two cofinal auxiliary gap questions agree]
\label{cor:f12v60-equivalence}
For this fixed packing and fixed
finite \(\beta\), a positive
volume-uniform lower bound for
\(g_N\) exists if and only if one
exists for \(h_N\). Quantitatively,
\[
 \frac{g_N}{g_N+K_\beta}
               \le h_N\le C_{\rm cond}(\beta)g_N .
 \tag{64.16}
\]
The constants in this comparison
are independent of volume. They
are not asserted to be uniform
as \(\beta\to\infty\).
\end{corollary}
\begin{proof}
Combine (64.9) and (64.15).
For a fixed positive lower bound
on \(g_N\), the increasing function
\(g\mapsto g/(g+K_\beta)\) gives
a fixed positive lower bound on
\(h_N\). The reverse implication
is (64.15).
\end{proof}

\subsection{Reconnect to the original unrestricted fibre updates}

V40 already proves the comparison
between the unrestricted interface
heat-bath gap \(h_N\) and the gap
\(\lambda_{\rm fibre}\) of ordinary
full-fibre one-group heat baths:
\[
 \frac{h_N}{M_{\rm inc}C_\beta^{\rm fibre}}
            \le\lambda_{\rm fibre}\le n_*h_N .
 \tag{64.17}
\]
All constants here are exactly the
fixed-packing constants of (44.8)--(44.10);
the superscript distinguishes its
comparison factor from the new
one-group conditional constant.
Consequently the new conversion yields
\[
 \frac{g_N}
      {M_{\rm inc}C_\beta^{\rm fibre}(g_N+K_\beta)}
       \le\lambda_{\rm fibre}
       \le n_*C_{\rm cond}(\beta)g_N .
 \tag{64.18}
\]
This uses the actual compact
interface marginal with every
integrated block partition function.
V40's lifting and comparison are
inherited, not re-proved as new
results here.

Thus a future proof of V59's
derivative-energy estimate would
advance the original unrestricted
fibre heat-bath problem, not only
an unrelated stabilized operator.
But (64.18) supplies no positive
value for \(g_N\) uniform in volume.
It also does not identify any of
these auxiliary clocks with the
physical transfer Hamiltonian.

\subsection{Exact checks}

The script
\begin{center}\small\ttfamily
scripts/verify\_ym\_f12\_v60\_clock\_comparison.py
\end{center}
uses rational Gaussian integration
on three interacting coordinates
with precision matrix
\[
 \begin{pmatrix}3&-1&0\\-1&4&-1\\0&-1&3\end{pmatrix}.
 \tag{64.19}
\]
Its local diffusions do not commute.
On polynomial spaces of degrees
at most one, two and three, of
dimensions 4, 10 and 20, it verifies
the weighted positive-semidefinite
form in (64.5), the diffusion
floor 2, and the converted
heat-bath floor \(1/2\) with
\(K_\beta=2\). The conditional
Gaussian projections are assembled
exactly and checked for idempotence.

All checks pass. They test the
noncommuting local-square and
duality algebra; they are not
verification of a uniform gap for
the nonlinear compact YM law.

\begin{decision}[Direction after the V60 review]
The two scalar dynamics are now
quantitatively connected on the
unrestricted measure. It is therefore
reasonable to seek the missing gap
through local heat-bath operators,
where positivity and exact
commutation of noninteracting
updates are available, rather
than projecting V58's vector
inverse through a nonlocal
gradient projection.

The next concrete task is to derive
a finite-patch algebraic gap
criterion for these actual
conditional projections, retaining
the boundary and overlap counts,
and then test whether the proved
finite-patch estimates meet its
threshold. This proposed criterion
is not assumed to hold with a
useful constant. The archive search
found no completed such criterion
in the YM lineages.

If its threshold is not met, the
deficit must be stated as a local
estimate to prove, rather than
replaced by an assumed global
return margin. The next companion
audit is V65 and the next
ten-version direction review is V70.
\end{decision}

V59 had 940226 bytes and SHA--256
\begin{center}\scriptsize\ttfamily
7842045A1BCCDE00F9DF51C3470B85F8C029CA0DF263671F3B7EBC4415C55739.
\end{center}
Its complete body is retained apart
from the title version. Only active
V59 is replaced by V60; all archives,
companions and previous scripts
remain.

\begin{firewall}[Position after V60]
The companion audit and direction
review are complete. The exact
mixed Bochner square identity and
quantitative diffusion/heat-bath
gap comparison are proved for the
unrestricted compact interface.
They connect two still-unproved
uniform gap estimates; they do
not supply either one. The full
continuum construction and physical
Yang--Mills mass gap remain open.
\end{firewall}
\section{V61: an unrestricted finite-patch criterion and its unmet threshold}
\label{sec:f12v61}

\subsection{Context, scope, and inherited results}

V60 connects the scalar diffusion and ordinary heat-bath
gap questions on the actual unrestricted compact interface.
It does not prove either volume-uniform gap. This section
executes its proposed upstream test: replace the unproved
global return margin by an explicitly stated sufficient
condition on finite conditional patches, and compare that
condition with the local estimates already available.
The measure, its compact exterior data, and the update
rates are not changed.

The projection-square method is a Knabe-type finite-size
argument, not a claimed new general principle. For primary
literature on finite-size criteria in higher-dimensional
frustration-free systems, see M.~Lemm,
\emph{Finite-size criteria for spectral gaps in
\(D\)-dimensional quantum spin systems} (2019),
\url{https://arxiv.org/abs/1902.07141}.
We prove the particular criterion below directly.
We do not import a quantum spin theorem whose tensor-local
hypotheses have not been checked for conditional expectations.
Nor do we claim the optimized thresholds in that literature.

The archive search found no completed instance of this
finite-patch criterion in the YM lineages. V38's recursive
surface comparison, V39's boxed gap, and V40's lifting are
inherited results, not new discoveries in this section.
The companion files are still SM V72, NS V26, and parallel
YM V5, as recorded at the completed V60 audit. The next
scheduled companion audit remains V65.

\subsection{The actual conditional projections have a finite commutation graph}

Use V58--V60's finite-volume compact interface law
\[
 d\bar\nu(v)=Z^{-1}e^{-W(v)}\prod_{i\in I}dg(v_i),
 \qquad W=U_S-\sum_b\log Z_b(v_{J_b}),
 \qquad
 Z_b(v_{J_b})=\int e^{-U_b(x_b,v_{J_b})}\,dg(x_b).
 \tag{65.1}
\]
The packing is fixed. All block interiors in (65.1) are
integrated over their full compact spaces. The density
is smooth and strictly positive, and the effective
potential is a sum of bounded-range terms of bounded
incidence, including the \(-\log Z_b\) terms.
Let
\[
 P_i f=\bar\nu(f\mid v_{I\setminus\{i\}}),\qquad
 h_i=I-P_i,\qquad H=\sum_{i\in I}h_i .
 \tag{65.2}
\]
Every \(h_i\) is an orthogonal projection in the
\emph{same} space \(L^2(\bar\nu)\). Each site has rate one.

Choose an undirected graph \(E\) on \(I\) that contains
every pair of distinct coordinates appearing in a common
effective potential term. Additional edges are allowed
but can worsen the estimates. Write
\[
 \deg_E(i)\le\Delta,\qquad
 \|\operatorname{owner}(i)-\operatorname{owner}(j)\|_\infty
          \le R_* \quad(\{i,j\}\in E),
 \tag{65.3}
\]
with fixed \(R_*\ge1\) and \(\Delta<\infty\).
Distances on a torus are periodic distances.

\begin{proposition}[Conditional commutation and patch kernels]
\label{prop:f12v61-commutation}
If \(\{i,j\}\notin E\), then \(P_iP_j=P_jP_i\), and
consequently
\(\{h_i,h_j\}:=h_ih_j+h_jh_i=2h_ih_j\ge0\).
For \(B\subset I\), set \(H_B=\sum_{i\in B}h_i\)
and \(Q_B f=\bar\nu(f\mid v_{B^c})\).
Then \(\ker H_B=\operatorname{ran}Q_B\).
In particular \(\ker H\) consists precisely of constants.
\end{proposition}
\begin{proof}
With every coordinate except \(v_i,v_j\) fixed, no
effective potential term involves both remaining
variables. Thus their conditional density factors into
an \(i\) factor and a \(j\) factor. Integrating the two
variables in either order gives the same conditional
expectation. This proves commutation. A product of
commuting orthogonal projections is an orthogonal
projection, hence is positive.

Positivity gives
\(\langle f,H_Bf\rangle=\sum_{i\in B}\|h_if\|_2^2\).
It vanishes exactly when \(f\) is unchanged by each
single-coordinate conditional expectation in \(B\).
The positive smooth density on a finite compact product
is bounded above and below by positive constants relative
to product Haar. The corresponding measurable and
\(L^2\) classes therefore agree. Being unchanged by
\(P_i\) means that \(f\) has a version independent of
coordinate \(i\). Equivalently, the product-Haar
averaging operator in coordinate \(i\) leaves that
version unchanged. These Haar averaging operators
commute; their product over \(B\) shows that \(f\)
depends only on \(v_{B^c}\). The converse is immediate.
\end{proof}

The \(P_i\) need not act on disjoint tensor factors in
the usual quantum-spin sense: their conditional densities
depend on neighbouring coordinates. The proposition
establishes exactly the commutation needed below without
assuming a tensor-local Hilbert space representation.

\subsection{Projection squares with the complete overlap loss}

Let \(\mathcal B=(B_x)_{x\in X}\) be a finite family of
patches, allowing repeated patches, such that every
coordinate appears in exactly \(m\ge1\) patches. Define
\[
 q_{ij}=\#\{x:i,j\in B_x\},\qquad
 r_{\mathcal B}
 =\max_i\sum_{j:\{i,j\}\in E}\left(1-\frac{q_{ij}}m\right).
 \tag{65.4}
\]
Thus \(0\le q_{ij}\le m\). Suppose that for every patch
and every compact exterior configuration its actual
conditional one-coordinate heat-bath generator has gap
at least \(\gamma>0\). The same \(\gamma\) must work for
all patches and exterior configurations under consideration.
Empty patches impose no condition. Equivalently for
our purposes, disintegration yields
\[
 H_{B_x}\ge\gamma(I-Q_{B_x}),\qquad
 H_{B_x}^2\ge\gamma H_{B_x}.
 \tag{65.5}
\]
Indeed the local operators preserve the exterior, so
their finite sum decomposes fibrewise into the stated
conditional generator. Its kernel is the conditional
constants; the second inequality follows from its
spectral theorem and integration. An almost-everywhere
uniform conditional gap would suffice for this step;
the all-exterior version is the stronger target used here.

\begin{theorem}[Finite-patch criterion for conditional heat baths]
\label{thm:f12v61-patch}
Under (65.2)--(65.5),
\[
                 H^2\ge(\gamma-r_{\mathcal B})H .
 \tag{65.6}
\]
If \(\gamma>r_{\mathcal B}\), then
\[
 \operatorname{gap}(H)\ge\gamma-r_{\mathcal B},\qquad
 \operatorname{Var}_{\bar\nu}(f)
 \le\frac{\sum_i\bar\nu\operatorname{Var}(f\mid v_{i^c})}
                {\gamma-r_{\mathcal B}} .
 \tag{65.7}
\]
There is no translation-invariance assumption on \(W\).
\end{theorem}
\begin{proof}
Expand the bounded self-adjoint operators, using
\(h_i^2=h_i\). With one term for each unordered pair,
\[
 \begin{aligned}
 \sum_x H_{B_x}^2
       &=mH+\sum_{i<j}q_{ij}\{h_i,h_j\},\\
 mH^2-\sum_xH_{B_x}^2
       &=\sum_{i<j}(m-q_{ij})\{h_i,h_j\}.
 \end{aligned}
 \tag{65.8}
\]
All nonedge contributions in the second line are
positive. For each edge the elementary square identity
gives
\[
 \{h_i,h_j\}=(h_i+h_j)^2-h_i-h_j\ge-h_i-h_j.
 \tag{65.9}
\]
The coefficients \(m-q_{ij}\) are nonnegative. Hence
\[
 \begin{aligned}
 mH^2-\sum_xH_{B_x}^2
 &\ge-\sum_{\{i,j\}\in E}(m-q_{ij})(h_i+h_j)\\
 &=-m\sum_i
       \left[\sum_{j:\{i,j\}\in E}(1-q_{ij}/m)\right]h_i\\
 &\ge-mr_{\mathcal B}H .
 \end{aligned}
 \tag{65.10}
\]
Also \(\sum_xH_{B_x}=mH\), so (65.5) implies
\(\sum_xH_{B_x}^2\ge\gamma mH\).
Combine with (65.10) and divide by \(m\) to prove
(65.6). The spectral theorem for positive bounded
\(H\) implies that its spectrum avoids
\((0,\gamma-r_{\mathcal B})\) when this interval is
nonempty. Proposition~\ref{prop:f12v61-commutation}
identifies its zero eigenspace with constants, proving
(65.7).
\end{proof}

No norm of a global inverse, global return margin, or
mixing estimate was used in this algebra. The open
hypothesis is the numerical size of the \emph{actual}
local conditional gap in (65.5).

\subsection{Exact four-dimensional translated-patch counts}

For a periodic owner lattice of side \(M\), take all
translations of the owner cube \(\{0,\ldots,\ell-1\}^4\),
and let \(B_x\) contain every interface group whose
owner lies in that translate. Assume
\(\ell>R_*\) and \(M\ge2\ell+2R_*\), with the fixed
packing compatible with the torus. Missing interface
coordinates and several groups with the same owner
are allowed. Each coordinate belongs to
\[
                              m=\ell^4
 \tag{65.11}
\]
patches. For an edge with shortest owner displacement
\(\delta=(\delta_1,\ldots,\delta_4)\),
\[
 q_{ij}=\prod_{a=1}^4(\ell-|\delta_a|),\qquad
 r_\ell
 =\max_i\sum_{j:\{i,j\}\in E}
       \left[1-\prod_{a=1}^4
                    \left(1-\frac{|\delta_a(i,j)|}{\ell}\right)\right].
 \tag{65.12}
\]
These are counts of updated coordinates, not counts
of all original physical links.

To verify (65.12), in one periodic coordinate exactly
\(\ell-|d|\) interval starts contain both sites with
shortest displacement \(d\), when \(|d|\le R_*<\ell\).
The imposed size excludes a second wrap-around interval
overlap. Starts in the four coordinates are independent,
so their counts multiply. If two distinct group variables
have the same owner, all \(\ell^4\) patches contain either
both or neither; their pair contributes zero loss.
Holes in the interface do not change the count for
coordinates that remain.

In particular,
\[
 r_\ell\le\Delta\left[1-(1-R_*/\ell)^4\right]
                \le\frac{4\Delta R_*}{\ell}.
 \tag{65.13}
\]
The final inequality follows from
\(1-\prod_a(1-t_a)\le\sum_a t_a\) for \(0\le t_a\le1\).
Therefore one sufficient, deliberately unoptimized,
finite-patch test is
\[
 \boxed{\quad
 \inf_{\substack{\text{admissible retained data}\\
                 x,\ \text{all compact exteriors }\xi}}
       \operatorname{gap}(H_{B_x}^{\,\xi})
       \ \ge\gamma_\ell
       \ >\frac{4\Delta R_*}{\ell}.
       \quad}
 \tag{65.14}
\]
One may instead use the smaller exact loss \(r_\ell\).
For a cofinal family, the same patch size and positive
margin must work in every sufficiently large torus.
Local retained data may vary; only configurations
covered by the programme's hypotheses are included in
the displayed infimum. It is not an assumption about
an already integrated coarse-field law.

The word \emph{finite} refers to the bounded patch and
its finite-range exterior collar. Its compact continuous
parameters still require rigorous uniform estimates;
sampling a finite number of exterior points would not
verify (65.14).

\subsection{Populate the available patch estimate without changing the law}

The full interface effective terms satisfy a compact
oscillation bound, in addition to V58's derivative
bounds. To see this for the potentially troublesome
term \(-\log Z_b\), if \(v,v'\) are any two compact
interface configurations, then the fixed block template
and bounded original local terms give
\[
 \sup_{x_b}|U_b(x_b,v)-U_b(x_b,v')|
                      \le c_b(\beta+1).
 \tag{65.15}
\]
Multiplying the two integrands by the resulting
pointwise bounds and integrating over the \emph{same}
full compact interior yields
\[
 e^{-c_b(\beta+1)}
       \le Z_b(v)/Z_b(v')\le e^{c_b(\beta+1)},\qquad
 \operatorname{osc}\log Z_b\le c_b(\beta+1).
 \tag{65.16}
\]
No sum over all blocks is taken in this local bound.
The fixed packing bounds the number of such terms
incident on any one interface coordinate, and bounds
the owner diameter of each support \(J_b\).

Consequently there are fixed constants \(q_*,D_*,R_*\)
and \(\widehat v_\beta\le c_*(\beta+1)\) such that
there are at most \(q_*\) interface groups per owner,
at most \(D_*\) effective terms per group, support
diameter at most \(R_*\), and each effective term has
oscillation at most \(\widehat v_\beta\).
All statements are uniform in compact free exterior
data. Increase \(R_*\), if necessary, to serve both
the commutation graph and these term supports.
The constants are the effective fixed-packing constants,
not silently the original free-coordinate constants.

Now V38's Theorem~\ref{thm:f12v38-recursion} applies
to this finite-range specification by exactly its
proved cut-and-disintegrate argument. More explicitly,
use its constants with the replacements just verified:
\[
 c_{\rm cut}^*=8q_*D_*(R_*+1),\qquad
 C_{\rm surf}^*=q_*D_*+\frac{64}{7}c_{\rm cut}^* .
 \tag{65.17}
\]
Its midpoint cuts charge only effective terms that
meet two updated parts. Terms meeting fixed exterior
data and only one updated part remain in that part's
potential. Thus the recursion keeps arbitrary compact
exteriors and the actual conditional kernels.
It gives, for every patch \(B_x\) of owner side \(\ell\),
\[
 \operatorname{gap}(H_{B_x}^{\,\xi})
       \ge e^{-C_{\rm surf}^*\widehat v_\beta\ell^3}
       \ge \underbrace{e^{-A(\beta+1)\ell^3}}
                         _{\displaystyle\gamma_{\rm cert}(\ell,\beta)},
 \qquad A=C_{\rm surf}^*c_* .
 \tag{65.18}
\]
This application is new bookkeeping for the full
effective interface; the recursive comparison theorem
itself remains inherited from V38.

\begin{proposition}[The present surface certificate does not cross the threshold]
\label{prop:f12v61-deficit}
Assume \(A>0\) and \(A(\beta+1)\ge1\). For every
integer \(\ell\ge1\),
\[
             \gamma_{\rm cert}(\ell,\beta)<1/\ell .
 \tag{65.19}
\]
If the chosen commutation graph has at least one edge
whose two owners differ, then, for the translated
patches above, \(r_\ell\ge1/\ell\).
In that case the lower certificate \(\gamma_{\rm cert}\)
does not verify even the exact sufficient test
\(\gamma_\ell>r_\ell\), for any eligible \(\ell\).
\end{proposition}
\begin{proof}
For \(\ell\ge1\), \(\ell^3>\log\ell\). Hence
\(e^{-A(\beta+1)\ell^3}\le e^{-\ell^3}<1/\ell\).
For a cross-owner edge, some integer
\(|\delta_a|\ge1\), while each factor in (65.12)
lies in \([0,1]\). Its product is at most
\(1-1/\ell\). This one edge contributes at least
\(1/\ell\) to an incident row of (65.12), proving
the second assertion.
\end{proof}

This is a statement about a proved lower
\emph{certificate}, not an upper bound on the true
patch gap. It does not rule out a better estimate
of that gap, a sharper projection inequality, a
different commutation graph, or a better patch cover.
If all interacting pairs have the same owner, the
overlap loss is zero and the cross-owner obstruction
does not apply. The coarse bound in (65.13) would then
discard useful information.

V39 and V57 do give stronger size-uniform statements
for boxed interface reference laws. They cannot be
inserted into (65.5), whose exterior and patch variables
are unrestricted compact groups. V55's separated-core
gap belongs to a marginal whose lifted updates resample
the extensive common bulk; those are not \(H_{B_x}\).
V56's nested-record contraction does not give a
one-coordinate patch Poincare inequality with arbitrary
compact exterior. These mismatches are precisely why
the finite test has not been filled by an earlier result.

\subsection{Consequences if the missing local margin is proved}

If a future estimate proves a fixed-patch margin
\(\varepsilon=\gamma_\ell-r_\ell>0\) uniformly for the
actual conditional laws, then the preceding theorem,
V60, and V40 give, respectively,
\[
 h_N\ge\varepsilon,\qquad
 g_N\ge\frac{\varepsilon}{C_{\rm cond}(\beta)},\qquad
 \lambda_{\rm fibre}\ge
       \frac{\varepsilon}{M_{\rm inc}C_\beta^{\rm fibre}}.
 \tag{65.20}
\]
These conclusions would hold uniformly over the
admissible large volumes at that fixed packing and
fixed \(\beta\). They concern auxiliary stochastic
generators. No identification with a physical transfer
Hamiltonian, uniform continuum scaling, construction
of continuum Schwinger functions, or physical positive
mass is furnished by (65.20).

\subsection{Exact checks and the next local question}

The script
\begin{center}\small\ttfamily
scripts/verify\_ym\_f12\_v61\_patch\_criterion.py
\end{center}
uses exact rational arithmetic. On a five-site
interacting Ising ring with density proportional to
\(\prod_i(1+x_ix_{i+1}/100)\), it assembles the
\(32\)-state conditional projections and all five
length-three patches. The neighbouring projections
do not commute; nonneighbouring ones do. It verifies
the exact identity (65.8), the positive remainder in
(65.10), the conditional patch floors, and the global
gap inequality, with
\[
 \gamma=(99/101)^5=\frac{9509900499}{10510100501},
 \quad r=2/3,\quad
 \gamma-r=\frac{7509500495}{31530301503}>0 .
 \tag{65.21}
\]
That conservative patch floor follows also from
product variance comparison: there are at most five
interaction factors, so every conditional density
ratio has oscillation ratio at most \((101/99)^5\).
The code checks the claimed forms on the full state
space, hence on all exterior fibres in this model.
It separately verifies 249 four-dimensional periodic
pair-incidence counts for cube sides two, three, and four.

For a negative control it takes two rank-one projections
on the first two coordinates of \(\mathbb R^3\), along
\((1,0,0)\) and \((3/5,4/5,0)\). Singleton patches
have gap one, but their sum has nonzero eigenvalues
\(2/5,8/5\), with third coordinate in the common kernel.
The boundary loss is one; dropping it would incorrectly
assert a global gap of one. The vector \((1,-2,0)\)
explicitly violates \(H^2\ge H\).
All checks pass. These are algebraic and geometric
tests, not YM patch-gap certificates.

\begin{decision}[Upstream task after the finite-patch test]
The useful unresolved statement is now a local
all-exterior estimate: can an unrestricted compact
patch of some fixed side have gap exceeding its
explicit overlap loss? Merely enlarging the patch
while retaining (65.18) cannot verify the criterion.

The next investigation should test whether the
overlap algebra can use buffered block projections
with a smaller normalized loss, while keeping their
actual conditional exteriors and an explicit comparison
back to rate-one group updates. A block has gap one
only when it is itself the update; that normalization
must not be confused with the gap of its constituent
single-group dynamics. Any proposed gain must survive
this distinction and the overlap count.
This is a direction to check, not a claimed closure.
The next companion audit is V65; the next ten-version
direction review is V70.
\end{decision}

V60 had 953665 bytes and SHA--256
\begin{center}\scriptsize\ttfamily
804DF1FA11153B1AA06279AF7314A7B505DF6F3F067FD08DF271DFFE3CDCBB2B.
\end{center}
Its complete body is retained apart from the title
version. Only the active V60 file is replaced by V61.
All archived lineages, companion notes, and earlier
verification scripts remain unchanged.

\begin{firewall}[Position after V61]
The actual compact conditional projections satisfy
the proved finite-patch criterion with explicit
four-dimensional boundary loss. The available
unrestricted surface-order patch certificate does
not meet its threshold in the stated large-\(\beta\)
regime. No cofinal auxiliary gap has been inferred
from that insufficient certificate. The full continuum
Yang--Mills construction and physical mass gap remain
unproved, and the research goal remains active.
\end{firewall}
\section{V62: actual block angles, sharp rectangle conversion, and corner leakage}
\label{sec:f12v62}

\subsection{The question left by V61}

The preceding section proves a finite-patch criterion
for the unrestricted compact interface heat baths.
Its known local gap certificate does not cross the
boundary-overlap threshold. We now check the proposed
block-update alternative, including its normalization
and the geometry of overlapping updates.

The outcome is a precise sufficient criterion in
terms of conditional correlations of the exclusive
parts of two blocks, followed by an explicit conversion
back to rate-one single-group updates. A large overlap
volume does not itself make those correlations small:
diagonally shifted cubes have short corner paths
between their exclusive parts. We prove this obstruction
in a local positive model, rather than treating a
geometric sketch as an estimate for the YM law.

The oscillation-to-total-variation bound
\(\tanh(R/4)\) was already proved in archive f6 V19;
V36 and V53 already use two-projection gap algebra.
They are inherited tools. The sharp all-\(L^2\)
rectangle-to-angle conversion below, its use for
overlapping compact conditional updates, the rate-aware
assembly, and the explicit corner test are the
additional work here. For primary background on
oscillation contraction of Markov operators, see
S.~Gaubert and Z.~Qu,
\emph{Dobrushin's ergodicity coefficient for Markov
operators on cones},
\url{https://arxiv.org/abs/1307.4649}.
The proof needed here is given directly and uses
neither a global mixing assumption nor an imported
continuum theorem.

\subsection{An all-\(L^2\) consequence of a joint rectangle bound}

Let \(\pi(du,dv)\) have a strictly positive continuous
density \(w(u,v)\) on a finite compact product, relative
to product Haar (finite positive state spaces also
suffice). Normalization of \(w\) is immaterial. Define
its logarithmic rectangle diameter
\[
 \Omega(w)=\sup_{u,u',v,v'}
 \left|\log\frac{w(u,v)w(u',v')}{w(u,v')w(u',v)}\right|.
 \tag{66.1}
\]
The supremum is finite. Multiplication by a separate
positive \(u\) factor and \(v\) factor leaves (66.1)
unchanged, as in V54. Write \(\pi_U,\pi_V\) for the
marginals, and
\[
 Kg(u)=\mathbb E_\pi[g(V)\mid U=u],\qquad
 K^*f(v)=\mathbb E_\pi[f(U)\mid V=v].
 \tag{66.2}
\]
These operators are adjoints in the indicated marginal
\(L^2\) spaces. Maximal correlation is the norm of
\(K:L^2_0(\pi_V)\to L^2_0(\pi_U)\).

\begin{theorem}[Sharp rectangle-to-angle bound on the compact law]
\label{thm:f12v62-rectangle}
With \(\theta=\tanh(\Omega(w)/4)\),
\[
 \sup_{u,u'}\|\pi(dv\mid u)-\pi(dv\mid u')\|_{\rm TV}
     \le\theta,\qquad
 \sup_{v,v'}\|\pi(du\mid v)-\pi(du\mid v')\|_{\rm TV}
     \le\theta,
 \qquad \rho(U,V)\le\theta .
 \tag{66.3}
\]
The constant in the last bound is sharp.
\end{theorem}
\begin{proof}
For two fixed \(u,u'\), the oscillation in \(v\) of
\(\log w(u,v)-\log w(u',v)\) is at most \(\Omega(w)\).
Conditional normalizers add constants and do not
change this oscillation. Archive f6 V19's proved
oscillation-to-total-variation lemma gives the first
bound. Interchanging \(u,v\) gives the second.

For completeness, if a Markov kernel has row
total-variation diameter at most \(\theta\), then for
every bounded real \(g\),
\[
                     \operatorname{osc}(Kg)
                         \le\theta\operatorname{osc}(g).
 \tag{66.4}
\]
Subtract \(\inf g\), divide by \(\operatorname{osc}g\)
when nonzero, and use the variational characterization
of total variation by functions between zero and one.
This proves (66.4) directly. Applying it twice yields
\[
 \operatorname{osc}(KK^*f)\le
                     \theta^2\operatorname{osc}(f).
 \tag{66.5}
\]

It remains to justify an \(L^2\), not merely an
oscillation, conclusion. The density
\(r=d\pi/d(\pi_U\otimes\pi_V)\) is bounded and continuous,
because \(w\) is positive continuous on a compact
product and its two marginal densities have positive
lower bounds. Thus \(K\) is Hilbert--Schmidt, and both
\(K\) and \(K^*\) map \(L^2\) into bounded continuous
functions by Cauchy--Schwarz and continuity of their
kernel rows in \(L^2\). The operator \(T=KK^*\) on
centered \(L^2(\pi_U)\) is compact, positive, and
self-adjoint. If its norm is positive, it is an
eigenvalue \(\lambda>0\), with real eigenfunction \(f\).
The identity \(f=\lambda^{-1}Tf\) supplies a bounded
continuous version. A nonzero centered eigenfunction
is not constant, so \(0<\operatorname{osc}f<\infty\).
Equation (66.5) gives
\(\lambda\operatorname{osc}f\le
\theta^2\operatorname{osc}f\), and hence
\(\lambda\le\theta^2\). If the norm is zero, the same
conclusion is immediate. Therefore
\(\rho(U,V)^2=\|T\|_{L^2_0}\le\theta^2\).

For sharpness, take two uniform signs with joint
probability \((1+tuv)/4\), \(0<t<1\).
Their maximal correlation is \(t\), and
\(\Omega=2\log((1+t)/(1-t))=4\operatorname{artanh}t\).
Equality holds in (66.3).
\end{proof}

Using the forward total-variation diameter alone
would only give the weaker square-root bound by the
earlier one-sided argument. The reverse conditional
kernel of the \emph{same joint law} is what supplies
the second contraction in (66.5). This proof does not
assume that an arbitrary one-sided Markov channel has
a reverse bound of the same size.

\subsection{Which angle belongs to two overlapping updates?}

Work with the unrestricted interface law \(\bar\nu\)
of (65.1). For a block \(B\subset I\), use
\[
 Q_Bf=\bar\nu(f\mid v_{B^c}),\qquad h_B=I-Q_B.
 \tag{66.6}
\]
For two blocks \(B,C\), set
\[
 U=B\setminus C,\quad O=B\cap C,\quad V=C\setminus B,
 \quad D=B\cup C .
 \tag{66.7}
\]
First fix every exterior coordinate \(v_{D^c}=\xi\).
In this conditional law, integrate the overlap \(O\)
over its \emph{full compact space}. The resulting
joint law of the exclusive parts \(U,V\) has density
\[
 w_{BC}^{\,\xi}(u,v)
       =\int e^{-W(u,o,v,\xi)}\,dg(o),
 \tag{66.8}
\]
up to factors constant in \(u,v\). Set
\(\rho_{BC}^\xi=\rho_{\pi^\xi}(U,V)\), with value zero
if either exclusive part is empty, and define
\[
 \rho_{BC}=\sup_\xi\rho_{BC}^\xi,\qquad
 \Omega_{BC}=\sup_\xi\Omega(w_{BC}^{\,\xi}).
 \tag{66.9}
\]
For a family of admissible retained data, the suprema
also range over those data. They are quantities for
the actual conditional law; pinning the overlap \(O\)
would define a different problem.

\begin{proposition}[The conditional wing angle controls the anticommutator]
\label{prop:f12v62-angle}
For any bounds \(c_{BC}\ge\rho_{BC}\), \(0\le c_{BC}\le1\),
\[
 h_Bh_C+h_Ch_B\ge-c_{BC}(h_B+h_C).
 \tag{66.10}
\]
One permissible bound is
\(c_{BC}=\tanh(\Omega_{BC}/4)\).
If no effective potential term meets both of two
disjoint blocks, then \(c_{BC}=0\) is permissible.
Nested blocks also have \(c_{BC}=0\).
\end{proposition}
\begin{proof}
In a fixed exterior fibre, \(Q_B\) projects onto
functions of \(V\), and \(Q_C\) onto functions of \(U\).
Their common range consists of constants, by strict
positivity of the compact density. On the centered
space write \(p,q\) for the two restricted projections.
The norm of \(p\) on \(\operatorname{ran}q\) is exactly
\(\rho_{BC}^\xi\), since the overlap has been integrated
in (66.8). If \(a\in\operatorname{ran}p\) and
\(b\in\operatorname{ran}q\), then
\[
 \|a+b\|_2^2\le
 (1+\rho_{BC}^\xi)(\|a\|_2^2+\|b\|_2^2).
 \tag{66.11}
\]
Indeed the mixed inner product is bounded by
\(\rho_{BC}^\xi\|a\|_2\|b\|_2\), and
\(2\|a\|_2\|b\|_2\le\|a\|_2^2+\|b\|_2^2\).
Apply this to the map \((a,b)\mapsto a+b\);
its product with its adjoint is \(p+q\). Hence
\(\|p+q\|\le1+\rho_{BC}^\xi\). It follows that
\[
 h_B+h_C\ge(1-\rho_{BC}^\xi)(I-Q_D)
 \tag{66.12}
\]
on the fibre. Since \(\ker(h_B+h_C)=\operatorname{ran}Q_D\),
the spectral theorem gives
\((h_B+h_C)^2\ge(1-\rho_{BC}^\xi)(h_B+h_C)\).
Expand the square using \(h_B^2=h_B\), \(h_C^2=h_C\),
and integrate over \(\xi\) to prove (66.10).
Theorem~\ref{thm:f12v62-rectangle} gives the rectangle
choice, uniformly in exterior data.

For disjoint noninteracting blocks, the conditional
density factors, so the wing angle is zero. For nested
blocks, one exclusive part is empty and the two
conditional projections commute by the tower property.
\end{proof}

Thus overlap does not itself incur a negative
anticommutator cost. Product Haar updates commute even
when they overlap. Conversely, overlap does not itself
imply commutation for an interacting measure.

\subsection{Block assembly with every rate and comparison cost retained}

Let \((B_a)_{a\in\mathcal A}\) cover every interface
coordinate. Suppose their owner diameters are bounded
by a fixed block side \(b\), and let \(w_a>0\) be their
rates. Do not identify \(h_{B_a}\), which is one
conditional projection, with the sum of the
single-group operators inside \(B_a\). Put
\[
 G=\sum_a w_a h_{B_a},\qquad
 M_w=\max_{i\in I}\sum_{a:i\in B_a}w_a .
 \tag{66.13}
\]
Choose a symmetric matrix \(c_{ab}\) of the bounds in
Proposition~\ref{prop:f12v62-angle}, with \(c_{aa}=0\).
Pairs of disjoint noninteracting blocks may be assigned
zero, so each row has only a bounded number of
potentially nonzero entries for a bounded-overlap
geometric cover. Define
\[
                 \kappa=\min_a
                    \left(w_a-\sum_{b\ne a}w_b c_{ab}\right).
 \tag{66.14}
\]

\begin{theorem}[Rate-aware conditional block criterion]
\label{thm:f12v62-block}
If \(\kappa>0\), then
\[
 G^2\ge\kappa G,\qquad
 \operatorname{gap}(G)\ge\kappa .
 \tag{66.15}
\]
Suppose in addition that the actual all-exterior
single-group dynamics inside each \(B_a\) has Poincare
constant at most \(C_b(\beta)\). Then the original
rate-one interface heat baths satisfy
\[
                    h_N\ge\frac{\kappa}{M_w C_b(\beta)} .
 \tag{66.16}
\]
For fixed blocks, the existing unrestricted estimate
permits \(C_b(\beta)=e^{A(\beta+1)b^3}\), with V61's
effective-potential constant \(A\).
\end{theorem}
\begin{proof}
Expand and use (66.10):
\[
 \begin{aligned}
 G^2
 &=\sum_a w_a^2h_{B_a}
        +\sum_{a<b}w_aw_b\{h_{B_a},h_{B_b}\}\\
 &\ge\sum_a w_a
           \left(w_a-\sum_{b\ne a}w_bc_{ab}\right)h_{B_a}
 \ge\kappa G .
 \end{aligned}
 \tag{66.17}
\]
The common kernel consists of functions independent
of each block. Because the blocks cover all coordinates
and the density is equivalent to product Haar,
V61's kernel argument shows that these are precisely
constants. The spectral theorem proves the gap claim.

For each actual conditional block law, its
single-group Poincare estimate gives, after
disintegration,
\[
 \langle f,h_{B_a}f\rangle
 \le C_b(\beta)\sum_{i\in B_a}\langle f,h_i f\rangle .
 \tag{66.18}
\]
Conditional consistency identifies every \(h_i\)
here with the original unrestricted one-group
conditional operator. Multiply by \(w_a\), sum,
and use (66.13):
\[
 \langle f,Gf\rangle\le
          C_b(\beta)M_w\langle f,Hf\rangle .
 \tag{66.19}
\]
Combine this with
\(\kappa\operatorname{Var}_{\bar\nu}(f)\le\langle f,Gf\rangle\)
to prove (66.16). The value of \(C_b(\beta)\) follows
from V61's already checked effective compact
specification and V38's inherited recursion.
\end{proof}

Multiplying all rates by \(t>0\) multiplies
\(\kappa\) and \(M_w\) by the same \(t\). It cannot
improve (66.16). A single block projection's unit
gap on its nonzero spectral subspace is not a free
unit lower bound for its constituent single-group
generator. Equation (66.18) is the required conversion.
The comparison cost may be very large at large
\(\beta\), but is finite and volume-independent for
fixed \(b,\beta\). The remaining issue is a positive
uniform block margin \(\kappa\), not that conversion.

\subsection{A concrete buffered-cube count}

Take an owner torus of side \(M\), with \(M\) a
multiple of \(s\), \(s\ge R_*\), and \(M\ge10s\).
Use the side-\(3s\) cubes starting at \(s\mathbb Z^4\)
as update blocks, including every interface group
whose owner lies in each cube. These are buffered
versions of a side-\(s\) tiling. Use all block rates
equal to one. Each owner is contained in exactly
three intervals per coordinate, hence
\[
                             M_w=3^4=81
 \tag{66.20}
\]
whenever the interface is nonempty. Holes in the
interface do not increase this count.

If two such cubes can overlap or have an effective
potential term meeting both, their start displacement
\(sk\) satisfies
\[
 |k_a|s\le3s-1+R_*\le4s-1,\qquad |k_a|\le3
           \quad(a=1,\ldots,4).
 \tag{66.21}
\]
Use shortest periodic start displacements; the size
assumption avoids ambiguity in this count. There
are at most \(7^4-1=2400\) other block labels of
this kind for each block. Thus, if an all-exterior
bound \(\Omega_{ab}\le\Omega_*\) were proved for
every potentially interacting block pair, one
sufficient unoptimized margin would be
\[
 \kappa\ge1-2400\tanh(\Omega_*/4)>0,\qquad
 h_N\ge
 \frac{1-2400\tanh(\Omega_*/4)}
             {81e^{A(\beta+1)(3s)^3}} .
 \tag{66.22}
\]
The displayed positive-margin condition is equivalent
to \(\Omega_*<4\operatorname{artanh}(1/2400)\).
It is a target, not a bound already proved for YM.
Keeping the separate \(c_{ab}\) and rates in
(66.14) can be much better than replacing every
pair by the same worst bound.

\subsection{Why a large cube overlap does not supply that target}

There is a specific geometric failure in a tempting
application of the earlier buffered-annulus estimates.
In four dimensions let
\[
 B=\{0,\ldots,3s-1\}^4,\qquad C=B+(s,s,0,0).
 \tag{66.23}
\]
Their overlap has \(36s^4\) owner cells. Nonetheless,
with
\[
 \begin{aligned}
 u&=(s-1,3s-1,0,0)\in B\setminus C,\\
 z&=(s,3s-1,0,0)\in B\cap C,\\
 v&=(s,3s,0,0)\in C\setminus B ,
 \end{aligned}
 \tag{66.24}
\]
the nearest-neighbour path \(u,z,v\) has length two,
independent of \(s\). Thus the two exclusive parts
do not acquire separation proportional to the
overlap width or to the block side.

\begin{proposition}[A local positive corner model with a nonvanishing angle]
\label{prop:f12v62-corner}
There is a bounded-degree nearest-neighbour positive
spin law on the geometry (66.23), for every \(s\),
whose two exclusive-part maximal correlation is
the same fixed number \(r^2>0\), \(0<r<1\).
The resulting two-block heat-bath gap on their
three interacting spins is \(1-r^2\).
\end{proposition}
\begin{proof}
Put signs at the owners and make every interaction
zero except along \(u,z,v\). Relative to product
uniform signs, give these three spins the density
\[
                  (1+r x_u x_z)(1+r x_z x_v).
 \tag{66.25}
\]
It is strictly positive and its mean is one.
All other spins are independent spectators.
Summing over the overlap spin \(x_z\) gives the
exclusive-endpoint joint law
\[
                \pi(x_u,x_v)=\frac{1+r^2 x_u x_v}{4}.
 \tag{66.26}
\]
Every other exclusive spin is independent; conditional
averaging over these spectators has norm at most one.
The centered conditional channel in (66.26) is
multiplication by \(r^2\) on its single sign mode,
so the full exclusive-part angle is exactly \(r^2\).
It does not depend on \(s\).

On the three spins, \(Q_B\) retains only \(x_v\) and
\(Q_C\) only \(x_u\). Their centered ranges are the
one-dimensional sign spans, with inner product
\(r^2\). The two-block bound (66.12) gives gap at
least \(1-r^2\), and \(x_u+x_v\) is an eigenfunction
with that eigenvalue. This proves equality.
\end{proof}

The example is not a counterexample to a YM gap or
to mixing. Its nontrivial interacting component
has fixed size and a positive gap. It disproves
only the proposed geometric inference that the
pair angle must vanish as the cube overlap grows.
In fact (66.26) saturates the sharp rectangle
conversion, so that conversion cannot remove
this corner effect.

Even without diagonal shifts, some pairs in the
full buffered-cube cover merely touch or interact
across a thin boundary; the count (66.21) includes
them. They cannot be assigned the same large
separation as an aligned, deeply overlapping pair.

\subsection{What the existing rectangle estimates do and do not improve}

V54 proves for its legitimate guarded-annulus
geometry that the full compact endpoint joint law
has rectangle diameter at most
\[
                       \Lambda=\Omega-2\log(1-q).
 \tag{66.27}
\]
Theorem~\ref{thm:f12v62-rectangle} therefore improves
the endpoint maximal-correlation certificate there to
\[
                             \rho\le\tanh(\Lambda/4).
 \tag{66.28}
\]
This uses the same full law, the same repair geometry,
and the same restrictions on pins. No new guard
estimate is assumed. V56's already proved complete-record
Markov factorization then gives the stronger product
bound
\[
 \rho(X_0,X_n)\le\prod_{j=0}^{n-1}
                        \tanh(\Lambda_j/4).
 \tag{66.29}
\]
It is still restricted to V56's mesoscopic scale
window. A better numerical channel factor does not
extend that window or supply the missing all-exterior
block-pair estimates in (66.9).

For (66.22), every relevant block pair and all its
compact exterior values must be covered. The
corner geometry above is not a concentric
core/annulus/outer decomposition with a separating
annulus of width comparable to \(s\). In addition,
V53's repair patches must remain unfrozen by the
specified exterior. Consequently V53 cannot be
applied to all these block pairs merely because
their intersection contains many sites.

The generic compact comparison also fails to
provide a useful small pair bound. For example,
on a fixed block union contained in an owner
cube of side at most \(6s\), the effective
potential terms that depend on the union have
total oscillation bounded by
\(C(\beta+1)s^4\), by V61's bounded incidence and
term oscillation. Integration over \(O\) preserves
the corresponding ratio bound. Taking a rectangle
uses at most twice that logarithmic range, so it
yields only
\[
                    \Omega_{ab}\le C'(\beta+1)s^4.
 \tag{66.30}
\]
This is an available coarse upper certificate,
not a lower bound on the true rectangle diameter.
Its resulting certified \(c_{ab}\) approaches one,
not zero, at large \(\beta\) or large \(s\).
Inserting this same upper certificate for all
2400 pairs therefore gives no positive margin
in (66.22) in that regime.

\subsection{Exact checks and next direction}

The script
\begin{center}\small\ttfamily
scripts/verify\_ym\_f12\_v62\_block\_angles.py
\end{center}
checks four positive six-state joint laws obtained
from the matrix
\(\left(\begin{smallmatrix}1&2&3\\3&2&1\end{smallmatrix}\right)\)
by separate endpoint tilts. Their maximal rectangle
ratio is exactly nine, so the sharp angle bound is
\(1/2\). Both row total-variation bounds and both
weighted \(L^2\) squared-channel inequalities pass
exact rational positive-semidefinite tests.

For the eight-state corner model with \(r=3/4\),
the script verifies the exact wing angle \(9/16\),
two-block gap \(7/16\), and noncommutation. With
block rates \(1,5/4\), it verifies (66.17)--(66.19)
with
\[
 \kappa=19/64,\qquad M_w=9/4,\qquad C_b=49,\qquad
                  h_N\ge19/7056 .
 \tag{66.31}
\]
The factor 49 is the conservative product-comparison
constant \(((1+r)/(1-r))^2\). It also checks the
corner membership and length-two path at scales
\(s=1,2,4,9\), and the incidence and safe neighbour
counts in (66.20)--(66.21).
All checks pass. They check the exact algebra and
the stated local countermodel, not the missing YM
block-angle estimates.

\begin{decision}[Do not replace boundary analysis by overlap volume]
The block route has a rigorous all-exterior
conditional-angle criterion, including update-rate
and single-group comparison costs. No step proves
that its margin is positive for the unrestricted
YM law. The coarse large-overlap shortcut fails
at explicit corner paths.

The next useful upstream investigation is to keep
strong near-corner interactions inside finite
patches and test whether a multi-patch square
inequality can average their boundary losses,
rather than requiring every pair of blocks to
have a small angle. The exact pair-dependent
anticommutator bound (66.10) is available for
that analysis. Any improvement must address the
actual all-exterior patch gap left open in V61;
renaming a unit block projection as that patch
gap is not a solution.
The next companion audit remains V65, and the
next ten-version direction review remains V70.
\end{decision}

V61 had 972495 bytes and SHA--256
\begin{center}\scriptsize\ttfamily
5B82BEDDAB079F8188930FB05F27B884595356D54E1115A47987788553A78863.
\end{center}
Its body is retained in full apart from the title
version. Only active V61 is replaced by V62; all
archives, companion files, and previous scripts
are preserved.

\begin{firewall}[Position after V62]
The sharp compact rectangle-to-angle bound, actual
overlapping-block anticommutator estimate, and
rate-aware block-to-group gap criterion are proved.
The geometric corner obstruction is explicit,
and the earlier valid annular channel estimates
have a sharper constant. None of these supplies
the missing positive cofinal YM block or patch
margin. The continuum construction and physical
Yang--Mills mass gap remain unproved, and the
full research goal remains active.
\end{firewall}
\section{V63: weighted patch squares lower the sufficient threshold}
\label{sec:f12v63}

\subsection{What changes in the local target}

V61's indicator-weight patch cover gives a sufficient
unrestricted local gap threshold of order \(\ell^{-1}\).
V62 shows why replacing this by a requirement that all
overlapping blocks have small pair angles is not a
geometric shortcut: strong short corner paths remain.
Here we retain those interactions inside translated
patches and average \emph{weighted} patch squares.
The resulting sufficient unweighted local gap
threshold is of order \(\ell^{-3/2}\), with all
constants and local update rates kept explicit.

This is an improvement of the criterion, not a proof
that the actual compact YM patches satisfy it.
The already proved surface-order lower certificate
still does not cross the new threshold.
The new target is weaker, however: a local gap lower
bound \(c_\beta\ell^{-p}\) with \(p<3/2\) would suffice
at a sufficiently large fixed patch size, whereas
V61's corresponding power test requires \(p<1\).

Weighted finite-size methods and the exponent
\(3/2\) are not claimed as new principles.
For primary context see M.~Lemm and E.~Mozgunov,
\emph{Spectral gaps of frustration-free spin systems
with boundary},
\url{https://arxiv.org/abs/1801.08915}, and the
higher-dimensional finite-size work cited in V61.
Their stated spin-system hypotheses are not substituted
for this programme's conditional-expectation hypotheses.
The complete argument for the present four-dimensional
compact interface is given below. The archive search
found no previously completed weighted version of
V61 in the YM lineages.

\subsection{Translate a weight, not the underlying measure}

Use the actual unrestricted law \(\bar\nu\), the
one-group projections \(h_i=I-P_i\), and the rate-one
generator \(H=\sum_i h_i\) from V61. The commutation
graph \(E\) has degree at most \(\Delta\) and owner
range at most \(R_*\). The measure may be
inhomogeneous, including in its retained coarse data.
Only the deterministic patch weights are translated.

For a nonnegative finitely supported owner weight
\(a:\mathbb Z^4\to[0,\infty)\), periodize it on a
sufficiently large torus and write
\[
 \begin{aligned}
 m_1&=\sum_y a(y),&
 m_2&=\sum_y a(y)^2,\\
 q_a(\delta)&=\sum_y a(y)a(y-\delta),&
 d_a(\delta)&=m_2-q_a(\delta)
       =\tfrac12\sum_y|a(y)-a(y-\delta)|^2 .
 \end{aligned}
 \tag{67.1}
\]
All sums in the periodic algebra are over the torus.
For displacements of size at most \(R_*\), the
nonperiodic formula applies whenever the support and
torus size prevent a second wrap-around overlap.
Assume \(m_1,m_2>0\). For each owner translate \(x\),
put
\[
 A_x=\sum_i a(\operatorname{owner}(i)-x)h_i,\qquad
 B_x=\{i:a(\operatorname{owner}(i)-x)>0\}.
 \tag{67.2}
\]
The support \(B_x\) is an ordinary finite conditional
patch; the coefficients in \(A_x\) are its specified
one-group update rates. Holes in the interface and
several group variables at one owner are allowed.

For each edge choose a symmetric number \(c_{ij}\)
such that
\[
 \{h_i,h_j\}\ge-c_{ij}(h_i+h_j),\qquad 0\le c_{ij}\le1.
 \tag{67.3}
\]
The universal value \(c_{ij}=1\) is always available
by V61. V62 permits smaller values from actual
conditional pair angles, if they have been proved.
No small-angle bound will be needed for the explicit
criterion below. Define the unnormalized loss
\[
 D_a=\max_i\sum_{j:\{i,j\}\in E}
        c_{ij}\,d_a(\operatorname{owner}(i)-\operatorname{owner}(j)).
 \tag{67.4}
\]
Pairs with the same owner contribute zero.

\begin{theorem}[Weighted conditional patch-square criterion]
\label{thm:f12v63-weighted}
Suppose every \(A_x\), in its actual conditional
patch law with arbitrary compact exterior data,
has gap at least \(\gamma_a>0\). The same bound
must hold for all the admitted retained data and
patch positions. Then
\[
 H^2\ge
       \frac{m_1\gamma_a-D_a}{m_2}\,H .
 \tag{67.5}
\]
If \(m_1\gamma_a>D_a\), this proves
\[
                    h_N\ge\frac{m_1\gamma_a-D_a}{m_2}.
 \tag{67.6}
\]
\end{theorem}
\begin{proof}
The first and second moments of the translated
weight are independent of the owner:
\[
 \sum_xA_x=m_1H,\qquad
 \sum_xA_x^2=m_2H+
            \sum_{i<j}q_a(\operatorname{owner}(i)-\operatorname{owner}(j))
                            \{h_i,h_j\}.
 \tag{67.7}
\]
The diagonal coefficient in the second identity
is \(m_2\), not \(m_1\); \(A_x\) itself is generally
not a projection. Subtracting from \(m_2H^2\) gives
\[
 m_2H^2-\sum_xA_x^2
     =\sum_{i<j}d_a(\operatorname{owner}(i)-\operatorname{owner}(j))
                       \{h_i,h_j\}\ge-D_aH .
 \tag{67.8}
\]
To justify the inequality, (67.1) makes every
coefficient nonnegative. Nonedges commute and
have positive anticommutators. Apply (67.3) to
edges and collect their incident row sums.

Strict positivity of the compact density and
positivity of all rates on \(B_x\) give
\(\ker A_x=\operatorname{ran}Q_{B_x}\), using
V61's kernel argument. The assumed fibrewise gap
and spectral theorem imply
\(A_x^2\ge\gamma_a A_x\), after disintegration
over the exterior. Sum this inequality in (67.8),
use \(\sum_xA_x=m_1H\), and divide by \(m_2\).
The kernel of \(H\) consists of constants, so a
positive coefficient yields (67.6).
\end{proof}

\begin{corollary}[Conversion from a rate-one patch gap]
\label{cor:f12v63-unweighted}
Let
\(\alpha_a=\min\{a(y):a(y)>0\}>0\).
If the ordinary rate-one conditional patch
generator \(\sum_{i\in B_x}h_i\) has gap at least
\(\gamma_\ell\), uniformly in all the same exterior
data, then one may take
\(\gamma_a=\alpha_a\gamma_\ell\).
The exact sufficient unweighted threshold is therefore
\[
 \gamma_\ell>\tau_a,\qquad
 \tau_a=\frac{D_a}{\alpha_a m_1},\qquad
 h_N\ge\frac{\alpha_a m_1}{m_2}(\gamma_\ell-\tau_a).
 \tag{67.9}
\]
\end{corollary}
\begin{proof}
On the conditional patch,
\[
 A_x\ge\alpha_a\sum_{i\in B_x}h_i
       \ge\alpha_a\gamma_\ell(I-Q_{B_x}).
 \tag{67.10}
\]
This proves the gap statement. Applying its spectral
theorem gives the required local square inequality.
We do \emph{not} square an arbitrary operator
inequality by assuming that squaring preserves
positive-operator order.
\end{proof}

For indicator weights \(a=1_C\), \(m_1=m_2=|C|\)
and \(\alpha_a=1\); (67.9) recovers V61 with its
possible pair-angle improvement. Multiplying \(a\)
by a constant multiplies both \(D_a\) and
\(\alpha_a m_1\) by its square. Such a rescaling
cannot improve \(\tau_a\). This prevents a hidden
update-rate normalization from creating the gain.

\subsection{A floor plus a four-dimensional pyramid}

Choose an integer side \(\ell\) divisible by four,
\[
 \ell\ge4,\qquad \ell>2R_*,\qquad \ell\ge R_*^2,
 \qquad M\ge2\ell+2R_* ,
 \tag{67.11}
\]
on an owner torus of side \(M\) compatible with
the fixed packing. On \(C_\ell=\{1,\ldots,\ell\}^4\)
define
\[
 t_\ell(y)=\frac1\ell
       \min_{1\le b\le4}\{y_b,\ell+1-y_b\},\qquad
 \varepsilon=\ell^{-1/2},\qquad
 a_\ell(y)=
 \begin{cases}
   \varepsilon+t_\ell(y),&y\in C_\ell,\\
   0,&y\notin C_\ell .
 \end{cases}
 \tag{67.12}
\]
The minimum is taken over all eight distances.
This is not a tensor product of four one-dimensional
tapers. Its positive floor avoids multiplying four
small boundary rates at a corner. For even \(\ell\),
\[
 \alpha_a=\varepsilon+\ell^{-1},\qquad
 0<a_\ell\le1\ \hbox{on }C_\ell,\qquad
 m_2\le m_1,\qquad m_1\ge\ell^4/64 .
 \tag{67.13}
\]
Indeed \(t_\ell\le1/2\) and \(\varepsilon\le1/2\).
In the central cube with each coordinate between
\(\ell/4+1\) and \(3\ell/4\), there are
\((\ell/2)^4\) sites and \(t_\ell\ge1/4\).
These observations prove every bound in (67.13).

\begin{proposition}[Quadratic translation defect]
\label{prop:f12v63-defect}
For \(\|\delta\|_\infty\le R_*\),
\[
 d_{a_\ell}(\delta)\le
       B_{R_*}\ell^2,\qquad
 B_R=\tfrac12R^2+16R .
 \tag{67.14}
\]
In particular \(D_{a_\ell}\le\Delta B_{R_*}\ell^2\)
using only the universal anticommutator bound.
\end{proposition}
\begin{proof}
Where both \(y\) and \(y-\delta\) lie in \(C_\ell\),
the constant floors cancel and the minimum of
one-Lipschitz distance functions is one-Lipschitz
in the \(\ell^\infty\) distance. Thus
\[
 |a_\ell(y)-a_\ell(y-\delta)|\le R_*/\ell .
 \tag{67.15}
\]
There are at most \(\ell^4\) such sites, giving
a squared-difference sum at most \(R_*^2\ell^2\).

The symmetric difference of the two cubes has at
most \(2\ell^3\|\delta\|_1\le8R_*\ell^3\) sites.
For a site in that symmetric difference, the
point that lies inside its cube is within
\(R_*\) lattice layers of some face, so its
nonzero weight is at most
\(\varepsilon+R_*/\ell\le2\varepsilon\), by (67.11).
This part of the squared-difference sum is therefore
at most \(32R_*\ell^2\). Add the two estimates and
divide by two using (67.1). This proves (67.14).
The torus size in (67.11) allows this zero-extension
calculation for every relevant displacement.
\end{proof}

\begin{theorem}[An explicit \(\ell^{-3/2}\) sufficient patch threshold]
\label{thm:f12v63-threshold}
In the actual unrestricted compact interface law,
assume the rate-one patch gaps for owner cubes of
side \(\ell\) satisfy the all-exterior uniform bound
\(\gamma_\ell\), with (67.11). Define
\[
 C_* =64\Delta B_{R_*}
           =32\Delta(R_*^2+32R_*).
 \tag{67.16}
\]
If
\[
                         \gamma_\ell>C_*\ell^{-3/2},
 \tag{67.17}
\]
then, uniformly in the admitted large volumes,
\[
             h_N\ge\ell^{-1/2}
                          \bigl(\gamma_\ell-C_*\ell^{-3/2}\bigr)>0 .
 \tag{67.18}
\]
No small pair-angle estimate is required.
\end{theorem}
\begin{proof}
Corollary~\ref{cor:f12v63-unweighted} permits
\(\gamma_a\ge\varepsilon\gamma_\ell\), since
\(\alpha_a\ge\varepsilon\). Equations (67.13)--(67.14)
give
\[
 \frac{D_a}{\varepsilon m_1}
     \le
 \frac{\Delta B_{R_*}\ell^2}
                 {\ell^{-1/2}\ell^4/64}
     =C_*\ell^{-3/2}.
 \tag{67.19}
\]
Insert into (67.6):
\[
 h_N\ge
     \frac{\varepsilon m_1}{m_2}
           \bigl(\gamma_\ell-C_*\ell^{-3/2}\bigr).
 \tag{67.20}
\]
Under the strict positivity condition the parentheses
are positive; \(m_1/m_2\ge1\) then gives (67.18).
\end{proof}

Strong nearby and near-corner interactions are
retained in every relevant weighted patch square.
Only their small translation defects are charged
in (67.8). This differs from demanding a small
angle for each large-block pair, which V62 showed
cannot follow from overlap volume alone.
The constants in (67.16) are conservative: the
exact moment and displacement sums in (67.9)
are available whenever sharper geometric counts
are useful.

\subsection{Why the floor has size \(\ell^{-1/2}\)}

There are two competing costs. A discontinuity
of height \(\varepsilon\) at the support boundary
contributes squared differences of order
\(\varepsilon^2\ell^3\). The smoothly varying
interior contributes order \(\ell^2\).
The minimum-rate conversion divides by a factor
of order \(\varepsilon\ell^4\). This suggests the
two threshold terms
\[
                     \frac{\varepsilon}{\ell}
                         +\frac1{\varepsilon\ell^2}.
 \tag{67.21}
\]
The following exact axial calculation justifies
the scaling and prevents treating that heuristic
as a free improvement of the patch gap itself.

Keep (67.12) but for the moment let
\(\ell^{-1}\le\varepsilon\le1\), and let
\(\delta=e_1\). Then
\[
 \begin{aligned}
 \alpha_a&=\varepsilon+\ell^{-1},\\
 m_1&=\varepsilon\ell^4+
               \frac{16}{\ell}\sum_{k=1}^{\ell/2}k^4,\\
 d_a(e_1)&=\ell^3(\varepsilon+\ell^{-1})^2
                                      +\frac{(\ell-2)^2}{8}.
 \end{aligned}
 \tag{67.22}
\]
These identities concern the weights alone, not
the unknown YM conditional gap.

To prove the middle identity, use the layer-count
formula for the integer depth:
\[
 \sum_{y\in C_\ell}\ell t_\ell(y)
       =\sum_{d=1}^{\ell/2}(\ell-2d+2)^4
       =16\sum_{k=1}^{\ell/2}k^4.
 \tag{67.23}
\]
For the defect, the two boundary faces contribute
\(\ell^3\alpha_a^2\) after the factor \(1/2\) in
(67.1). In the interior, a one-step axial move
changes the depth by one exactly on the two
sets of layers where the other three coordinate
depths are sufficiently large. Their contribution is
\[
 \frac1{\ell^2}\sum_{k=1}^{\ell/2-1}(2k)^3
                         =\frac{(\ell-2)^2}{8}.
 \tag{67.24}
\]
One may also verify this by fixing the other
three coordinates and summing the squared
differences along their one-dimensional tent
with its flat central part.

For this one axial edge, let
\(\tau_{\rm ax}=d_a(e_1)/(\alpha_a m_1)\).
Since \(m_1\le(3/2)\ell^4\),
\(\alpha_a\le2\varepsilon\), and
\((\ell-2)^2/8\ge\ell^2/32\) for \(\ell\ge4\),
\[
 \tau_{\rm ax}\ge
        \frac{\varepsilon}{3\ell}
          +\frac1{96\varepsilon\ell^2}
       \ge\frac1{6\sqrt2}\ell^{-3/2}.
 \tag{67.25}
\]
The last step is the arithmetic-geometric mean
inequality. Thus, within this floor-plus-pyramid
family and this minimum-rate conversion, an axial
edge with the universal coefficient \(c_{ij}=1\)
already prevents improving the \emph{power}
beyond \(3/2\) merely by varying the floor in the
stated range. This is not an impossibility theorem
for other finite-size methods or sharper local
weighted-gap estimates.

More precisely, for any fixed \(c>0\), taking
\(\varepsilon=c\ell^{-1/2}\) in (67.22) gives
\[
 \lim_{\ell\to\infty}\ell^{3/2}\tau_{\rm ax}
                   =10\left(c+\frac1{8c}\right).
 \tag{67.26}
\]
Indeed \(m_1/\ell^4\to1/10\),
\(\alpha_a\sqrt\ell\to c\), and
\(d_a(e_1)/\ell^2\to c^2+1/8\).
These limits follow directly from (67.22) or
the corresponding Riemann sums.
The choice \(c=1\) used for the explicit theorem
has limit \(45/4\). The best constant in this
axial asymptotic expression is \(5\sqrt2\),
attained at \(c=1/\sqrt8\). We have not substituted
this optimized axial constant for the full
four-dimensional row bound (67.16).

\subsection{The available YM certificate still does not meet the new test}

V61 verified the effective compact term oscillations
and applied V38's inherited recursive comparison.
For the present all-exterior patches its available
certificate remains
\[
             \gamma_{\rm cert}(\ell,\beta)
                  =e^{-A(\beta+1)\ell^3}.
 \tag{67.27}
\]
It is a lower bound on the true local gap, not
an assertion of its actual size.
If \(A(\beta+1)\ge1\), then for every \(\ell\ge1\),
\[
                \gamma_{\rm cert}(\ell,\beta)
                    <\ell^{-3/2}.
 \tag{67.28}
\]
For example,
\(\ell^3>(3/2)\log\ell\) follows by differentiation
on \([1,\infty)\) and the value at one.
For an interacting graph with \(\Delta\ge1\),
the displayed constant \(C_*\) is larger than one,
so (67.27) cannot verify (67.17).

There is also an exact-weight version of this
deficit, not only the conservative constant
comparison. Use \(c_{ij}=1\) and assume the graph
has some edge joining distinct owners. For any
nonzero displacement \(\delta\) of that edge,
the symmetric difference of the support cubes
contains an entire face from each cube, of
\(\ell^3\) points each. On these faces the nonzero
weight equals \(\alpha_a\). Consequently
\[
 d_a(\delta)\ge\alpha_a^2\ell^3,\qquad
 \tau_a=\frac{D_a}{\alpha_a m_1}
        \ge\frac{\alpha_a\ell^3}{m_1}
        \ge\ell^{-3/2}
              \quad(\varepsilon=\ell^{-1/2}).
 \tag{67.29}
\]
Here \(m_1\le\ell^4\) follows from (67.13).
Thus the certificate (67.27), combined only with
the minimum-rate comparison, does not verify
even the exact universal-coefficient test (67.9).
Proved smaller \(c_{ij}\), a better actual patch
gap, or a direct weighted-gap estimate could
change that conclusion; none has been silently
assumed in (67.29).

\begin{corollary}[A weaker polynomial local target would now suffice]
\label{cor:f12v63-power}
At fixed \(\beta\) and fixed packing, suppose the
actual all-exterior patch gaps satisfy
\(\gamma_\ell\ge c_\beta\ell^{-p}\) for all
sufficiently large eligible \(\ell\), with
\(c_\beta>0\) and \(p<3/2\).
Then the unrestricted interface has a positive
volume-uniform rate-one heat-bath gap.
\end{corollary}
\begin{proof}
Choose one eligible fixed side so large that
\[
                 \ell^{\,3/2-p}>C_*/c_\beta .
 \tag{67.30}
\]
Then (67.17) holds and (67.18) gives a strictly
positive lower bound independent of the larger
ambient torus. No limit in \(\ell\) is needed
after making this choice.
\end{proof}

This corollary is conditional. It replaces neither
the all-exterior requirement nor the compact
patch law by a boxed reference law. If its
hypothesis were proved, V60 and V40 would then
convert the resulting \(h_N\) to the corresponding
auxiliary diffusion and full-fibre heat-bath
bounds. They would still not identify a physical
transfer Hamiltonian or prove continuum existence.

\subsection{Exact tests and next upstream step}

The script
\begin{center}\small\ttfamily
scripts/verify\_ym\_f12\_v63\_weighted\_patches.py
\end{center}
checks (67.7)--(67.10) in exact rational arithmetic
on the same positive five-site interacting sign
ring used in V61, now with cyclic weights
\((3/4,1,1,3/4,0)\). This finite ring tests the
periodic algebra, not the large-torus nonwrap
assumption in the geometric theorem. It finds
\[
 m_1=7/2,\quad m_2=25/8,\quad \alpha_a=3/4,\quad
 D_a=5/4,\quad
 \gamma_a=(3/4)(99/101)^5,
 \tag{67.31}
\]
and verifies the resulting exact positive
global floor
\[
                       \frac{94606905469}{262752512525}.
 \tag{67.32}
\]
The code explicitly rejects the incorrect
replacement of \(m_2\) by \(m_1\) in the
diagonal of the local-square sum.

For the four-dimensional weights it directly
enumerates the weights and translation defects
at sides four and sixteen, checking four and
six displacement cases respectively, including
diagonal and range-two shifts where eligible.
It verifies (67.1), the upper defect bound,
the face lower bound, and the axial formula.
Layer sums also check sides \(4,16,36,64,144,256\).
For example the one-axial-edge thresholds from
the exact sums at sides 16 and 256 are
\[
 \tau_{\rm ax}(16)=\frac{1698}{31445},
 \qquad
 \tau_{\rm ax}(256)=\frac{1750368}{1000683019},
 \tag{67.33}
\]
compared with indicator-weight thresholds
\(1/16\) and \(1/256\), respectively.
They are per-edge geometric thresholds, not
the full graph row sum and not measured YM gaps.
The script also checks a conservative consequence
of (67.25) for every floor \(k/\ell\),
\(1\le k\le\ell\), at sides 16 and 64.
All checks pass. The proofs, not these finite
tests, provide the general inequalities.

\begin{decision}[Return to the actual patch energy]
The universal finite-size threshold has been
lowered without discarding strong corner
interactions or rescaling the global dynamics.
Further adjustment of the floor within this
family cannot improve its asymptotic power
when using only the minimum-rate comparison.

The next upstream task is therefore analytic:
test whether the actual unrestricted weighted
patch generator \(A_x\) can be bounded below
at the scale required by
\(m_1\gamma_a>D_a\), which here asks for a
weighted local gap larger than a constant
times \(\ell^{-2}\). The equivalent sufficient
unweighted target in (67.17) is also available.
Buffered local rarity may help only if it
controls the test function's energy under
the same arbitrary compact exterior data;
an unconditional probability of a good patch
is not that estimate.
This is the next statement to investigate,
not an assumption that has already closed it.
The next companion audit is V65 and the
next ten-version direction review is V70.
\end{decision}

V62 had 993564 bytes and SHA--256
\begin{center}\scriptsize\ttfamily
0D545B8ED562CC52BE9AB6B0CC7CB895B3272BFDCF237244A2917DDA563E9F3A.
\end{center}
Its complete body is retained apart from the
title version. Only active V62 is replaced by
V63. All archives, companions, and earlier
verification scripts are preserved.

\begin{firewall}[Position after V63]
The weighted projection-square identity, explicit
four-dimensional translation-defect bound, and
\(\ell^{-3/2}\) sufficient compact patch-gap
criterion are proved. The existing YM surface
certificate does not meet that improved test.
No actual cofinal auxiliary gap, continuum
construction, or physical Yang--Mills mass gap
has been inferred. The full research goal
remains active.
\end{firewall}
\section{V64: boundary conditioning can defeat a sufficient patch criterion}
\label{sec:f12v64}

\subsection{Audit the all-exterior hypothesis before optimizing it further}

V63 lowers the sufficient local gap threshold to order
\(\ell^{-3/2}\), but does not prove the required all-exterior
patch estimate. We now test whether such an estimate is a
necessary feature of a globally gapped finite-range system.
It is not.

We construct a strictly positive finite-state model on a
four-dimensional owner lattice whose ordinary heat baths
have a proved volume-uniform gap. Certain frozen exteriors
of an owner cube nevertheless create a low-temperature
two-dimensional Ising component on its boundary. An explicitly
identified published equilibrium large-deviation theorem
then gives exponentially small conditional patch gaps.
Completing the patch by its incident auxiliary variables
removes this defect in the example.

This is a countermodel to a general inference, not a
counterexample to YM or to V61--V63's sufficient theorems.
It has discrete variables and does not satisfy all the
programme's gauge, curvature, and vertical-tether hypotheses.
The distinction between worst-exterior criteria and averaged
tensorization was already emphasized in archive f6 V82.
The explicit globally gapped embedding, its completed patches,
and its implications for the weighted criterion are the
additional work here.

\subsection{A local auxiliary representation with an exactly product reference}

Let \(\mathcal G=(V,E)\) be a finite simple graph of degree
at most \(d\). Put \(\sigma_v\in\{-1,1\}\) at each vertex and
\(\eta_e\in\{-1,0,1\}\) at each edge. For fixed \(J>0\), set
\[
 \mu_J(\sigma,\eta)
 =2^{-|V|}z_J^{-|E|}
       \exp\!\left(J\sum_{e=\{u,v\}}\eta_e\sigma_u\sigma_v\right),
 \qquad z_J=1+2\cosh J .
 \tag{68.1}
\]
Every state has positive probability. Each edge potential
has oscillation \(2J\). The coordinate change
\[
                         \tau_e=\eta_e\sigma_u\sigma_v
 \tag{68.2}
\]
is a bijection, with inverse \(\eta_e=\tau_e\sigma_u\sigma_v\).
In coordinates \((\sigma,\tau)\), (68.1) becomes exactly
\[
 \nu_J=\bigotimes_v\operatorname{Unif}\{-1,1\}
           \otimes\bigotimes_e p_J,\qquad
             p_J(t)=e^{Jt}/z_J,\quad t\in\{-1,0,1\}.
 \tag{68.3}
\]
This proves the normalization in (68.1), without dropping
an interaction-dependent normalizer. The full spin marginal
is product at every \(J\).

Write \(e_\alpha(f)\) for the ordinary conditional variance
when physical coordinate \(\alpha\) is refreshed, and \(e_B(f)\)
for the conditional block variance. All rates are one.

\begin{theorem}[A volume-uniform gap in the physical coordinates]
\label{thm:f12v64-global}
The ordinary heat baths of (68.1), updating individual
\(\sigma_v\) and \(\eta_e\), satisfy
\[
 \operatorname{Var}_{\mu_J}(f)
 \le(2e^{2Jd}+1)
       \left(\sum_v e_{\sigma_v}(f)+\sum_e e_{\eta_e}(f)\right).
 \tag{68.4}
\]
Their gap is therefore at least \((2e^{2Jd}+1)^{-1}\),
uniformly in the graph size.
\end{theorem}
\begin{proof}
Product variance tensorization in (68.3) gives
\[
 \operatorname{Var}_{\mu_J}(f)
          \le\sum_v\widetilde e_v(f)+\sum_e\widetilde e_e(f),
 \tag{68.5}
\]
where the tildes denote conjugates of reference-coordinate
heat-bath energies. A \(\tau_e\) refresh at fixed spins changes
only \(\eta_e\), according to its actual conditional law.
Thus \(\widetilde e_e=e_{\eta_e}\).

A reference \(\sigma_v\) refresh holds all \(\tau\) fixed.
When it flips \(\sigma_v\), it reverses every incident
\(\eta_e\), leaving zero values unchanged. Its support is
\[
                   S_v=\{\sigma_v\}\cup\{\eta_e:e\ni v\}.
 \tag{68.6}
\]
The sigma-field of all reference coordinates except
\(\sigma_v\) contains every physical coordinate outside
\(S_v\). Conditional variance decreases when conditioning
is enlarged, so \(\widetilde e_v\le e_{S_v}\).
The reference refresh is not being identified with a
full physical star refresh; it preserves the incident
\(\tau\) variables.

For every physical exterior of \(S_v\), its conditional
density relative to product uniform measure has potential
oscillation at most \(2J\deg(v)\le2Jd\). V38's proved
product-variance comparison, also valid on these finite
state spaces, gives
\[
 e_{S_v}(f)\le e^{2Jd}
          \left(e_{\sigma_v}(f)+\sum_{e\ni v}e_{\eta_e}(f)\right).
 \tag{68.7}
\]
The factor is the largest-to-smallest density ratio, not
its square. Summing, each spin occurs once and each
auxiliary in two stars. Equation (68.5) therefore yields
\[
 \operatorname{Var}_{\mu_J}(f)
   \le e^{2Jd}\sum_v e_{\sigma_v}(f)
           +(2e^{2Jd}+1)\sum_e e_{\eta_e}(f),
 \tag{68.8}
\]
which proves the result.
\end{proof}

The change of variables is local in both directions.
Physical observables whose reference-coordinate dependence
sets are disjoint are exactly independent. In a finite-range
owner embedding, sufficiently separated local observables
have this property. The slow conditional patch below is
not being hidden in a slow full equilibrium measure.

\subsection{A fixed-range four-dimensional owner embedding}

Take a periodic four-dimensional lattice of side \(M\),
much larger than the patch sizes. Use only graph edges
\[
                       (x,x+e_j),\qquad j=1,2,
 \tag{68.9}
\]
in each plane of fixed \((x_3,x_4)\), so \(d=4\).
Assign owners by the fixed rule
\[
 \operatorname{owner}(\sigma_x)=x,\qquad
 \operatorname{owner}(\eta_{x,j})=x+e_4.
 \tag{68.10}
\]
Each interaction involves owners \(x,x+e_j,x+e_4\),
of \(\ell^\infty\) diameter one. There are three variables
per owner and bounded local incidences. The owner rule
does not depend on the patch subsequently chosen.

Let \(P_\ell\) update precisely the physical coordinates
whose owners lie in \(C_\ell=\{1,\ldots,\ell\}^4\).
For \(\ell\ge2\), consider the spin plane
\[
 D_\ell=\{(x_1,x_2,1,\ell):1\le x_1,x_2\le\ell\}.
 \tag{68.11}
\]
Every auxiliary on an edge incident on this plane has
fourth owner coordinate \(\ell+1\), and is outside
\(P_\ell\). This includes edges whose two spin endpoints
both belong to the updated plane.

Choose an exterior by setting \(\eta_e=1\) on internal
edges of \(D_\ell\), \(\eta_e=0\) on edges with exactly
one endpoint in \(D_\ell\), and every remaining exterior
auxiliary to zero. Exterior spins are arbitrary.
These prescriptions are consistent and have positive
probability in every finite ambient volume.

\begin{proposition}[An exact free-boundary Ising component]
\label{prop:f12v64-frozen}
Under this exterior, the conditional patch density
factors into a density on the remaining variables and
the independent plane-spin density
\[
 \pi_{\ell,J}^{\,\mathrm f}(\sigma)\ \propto\
 \exp\!\left(J\sum_{\substack{\{u,v\}\subset D_\ell\\|u-v|_1=1}}
                                           \sigma_u\sigma_v\right).
 \tag{68.12}
\]
This is the zero-field, free-boundary, two-dimensional
nearest-neighbour Ising law. The conditional rate-one
patch gap is at most its heat-bath gap
\(g_{\ell,J}^{\,\mathrm f}\).
\end{proposition}
\begin{proof}
No graph edges run in directions three or four.
Every factor touching \(D_\ell\) is either internal,
with frozen auxiliary one, or crosses its boundary,
with auxiliary zero and hence factor one. No updated
auxiliary touches the plane. The density therefore
factors as stated.

Functions of the plane spins form a reducing subspace
of the patch generator. Its restriction is exactly the
rate-one Ising heat-bath generator; all other updates
fix those functions. The variational characterization
of the gap proves the upper bound.
\end{proof}

\subsection{The declared external input and the slow conditional gap}

We use one published Ising equilibrium result, rather
than claim a new proof of it: for a fixed sufficiently
large \(J\), there are \(C_J,c_J>0\) such that, for odd
\(\ell\),
\[
 \pi_{\ell,J}^{\,\mathrm f}
      \left(\left|\sum_{x\in D_\ell}\sigma_x\right|=1\right)
                           \le C_J e^{-c_J\ell}.
 \tag{68.13}
\]
See F.~Cesi, G.~Guadagni, F.~Martinelli and R.~H.~Schonmann,
\emph{On the two-dimensional stochastic Ising model in
the phase coexistence region near the critical point},
J.~Stat.~Phys.~85 (1996), 55--102,
\url{https://link.springer.com/article/10.1007/BF02175556}.
The magnetization limit used here is also recorded as
equation (1.3) in F.~Martinelli and F.~L.~Toninelli,
\emph{On the mixing time of the 2D stochastic Ising model
with plus boundary conditions at low temperature},
\url{https://arxiv.org/abs/0905.3040}.
Its Sections 1.1--1.2 specify the free-boundary measure
and rate-one heat-bath convention. Only (68.13) is
imported; no YM analogue is asserted.

The dynamical consequence is elementary. For odd
\(\ell\), take \(F=1_{\{\sum_{D_\ell}\sigma_x>0\}}\).
Spin-flip symmetry and the absence of zero magnetization
give \(\operatorname{Var}(F)=1/4\). A refresh at \(x\)
can change \(F\) only when the sum of the other spins
is zero, in which case its conditional variance is
at most \(1/4\). That event implies full magnetization
\(1\) or \(-1\). Thus
\[
 \begin{aligned}
 \sum_{x\in D_\ell}
     \pi_{\ell,J}^{\,\mathrm f}\operatorname{Var}_x(F)
  &\le\frac{\ell^2}{4}
      \pi_{\ell,J}^{\,\mathrm f}(|\textstyle\sum\sigma|=1),\\
 g_{\ell,J}^{\,\mathrm f}
  &\le C_J\ell^2e^{-c_J\ell}.
 \end{aligned}
 \tag{68.14}
\]
This proof keeps the number of rate-one updates and
uses (68.13) as its explicit external theorem.

\begin{corollary}[Full and worst-exterior patch gaps can separate]
\label{cor:f12v64-separation}
For this one fixed \(J\), the full owner model has gap
at least \((2e^{8J}+1)^{-1}\) in every ambient torus.
Its worst-exterior owner-patch gaps, along odd sides,
are at most \(C_J\ell^2e^{-c_J\ell}\). They cannot
admit an all-exterior lower bound \(c\ell^{-p}\),
with fixed \(c>0\) and finite \(p\), at all sufficiently
large patch sizes.
\end{corollary}
\begin{proof}
Combine Theorem~\ref{thm:f12v64-global} at degree four,
Proposition~\ref{prop:f12v64-frozen}, and (68.14).
An exponential upper bound eventually lies below
every fixed positive power-law lower bound.
\end{proof}

For even \(\ell\), use an odd square of side \(\ell-1\)
inside the top plane and freeze its crossing auxiliaries
to zero. This also covers V63's sides divisible by four.
Every spin in this smaller top-plane square has the
same V63 pyramid rate
\(\alpha_a=\ell^{-1/2}+\ell^{-1}\). The weighted
conditional patch gap is therefore at most
\[
                   \alpha_a C_J(\ell-1)^2 e^{-c_J(\ell-1)}.
 \tag{68.15}
\]
Weighting does not remove this frozen-boundary mode.
The example does not rule out some useful fixed-patch
criterion. It rules out inferring V63's all-large-patch
polynomial hypothesis merely from a uniform global gap.

\subsection{A variable-type completion fixes this countermodel}

For any spin vertex set \(S\), instead update
\[
 \widehat P(S)=\{\sigma_v:v\in S\}
                  \cup\{\eta_e:e\cap S\ne\varnothing\}.
 \tag{68.16}
\]
All exterior spins and all nonincident auxiliaries
remain fixed. This includes all incident auxiliary
dependencies without adding every spin in an enlarged
owner cube.

\begin{proposition}[Completed conditional patches are uniformly gapped]
\label{prop:f12v64-completed}
For every \(S\) and every exterior of \(\widehat P(S)\),
its ordinary one-coordinate conditional heat baths
have gap at least \((2e^{2Jd}+1)^{-1}\).
\end{proposition}
\begin{proof}
Every factor involving a spin in \(S\) has its
auxiliary included. Change coordinates to those
spins and \(\tau_e=\eta_e\sigma_u\sigma_v\) on the
incident edges, with exterior spins fixed. This
remains a bijection, and the conditional law is
product: uniform spins in \(S\) and independent
\(p_J\) incident variables. Nonincident factors
depend only on the exterior.

All reference spin refreshes have their entire
physical star inside the completed patch. Repeat
(68.5)--(68.8). Each incident auxiliary belongs
to at most two updated stars, including when
only one of its endpoints belongs to \(S\).
The same uniform constant follows.
\end{proof}

The completion extends owners by only a fixed collar
in (68.10). Simply enlarging an owner cube would
not be equivalent: it introduces new boundary spins
whose own incident auxiliaries are again frozen.
Variable types and dependencies matter. There is
no proved analogous product coordinate change for
the actual YM interface, so the conclusion cannot
be transferred to YM without additional work.

\subsection{Retain bad-exterior energy in the weighted criterion}

Return to V63's actual YM operators \(A_x,H\), weight
moments \(m_1,m_2\), and translation loss \(D_a\).
Let \(G_x\) be an event measurable in the exterior
of \(B_x\). Suppose the actual weighted conditional
patch gap is at least \(\gamma>0\) on \(G_x\);
no positive gap is assumed on \(G_x^c\). Define
the nonnegative form
\[
 \mathcal B(f)=\sum_x\langle f,1_{G_x^c}A_x f\rangle
   =\sum_{x,i}a(\operatorname{owner}(i)-x)
                         \|1_{G_x^c}h_i f\|_2^2.
 \tag{68.17}
\]
Only \(i\in B_x\) contribute. Multiplication by the
exterior event commutes with \(A_x\) and these
\(h_i\), so the identity and positivity are exact.

\begin{theorem}[Weighted squares with an exterior-energy defect]
\label{thm:f12v64-defect}
Under these good-exterior gap assumptions,
\[
 m_2\|Hf\|_2^2
  \ge(\gamma m_1-D_a)\langle f,Hf\rangle-\gamma\mathcal B(f).
 \tag{68.18}
\]
If an independent argument proves, for every \(f\),
\[
                  \mathcal B(f)\le\eta\langle f,Hf\rangle,
           \qquad \gamma(m_1-\eta)>D_a,
 \tag{68.19}
\]
then
\[
                    h_N\ge
                 \frac{\gamma(m_1-\eta)-D_a}{m_2}>0.
 \tag{68.20}
\]
\end{theorem}
\begin{proof}
On good exterior fibres the spectral theorem gives
\(A_x^2\ge\gamma A_x\); on bad fibres \(A_x^2\ge0\).
Exterior measurability therefore yields
\[
                A_x^2\ge\gamma A_x-\gamma1_{G_x^c}A_x.
 \tag{68.21}
\]
Sum this and use V63's proved identity and bound
\(m_2H^2\ge\sum_xA_x^2-D_aH\), together with
\(\sum_xA_x=m_1H\). This proves (68.18).
Insert (68.19) and apply the spectral theorem,
using the constant kernel of \(H\), to obtain
(68.20).
\end{proof}

This is the precise error a random-exterior argument
must control. A small value of \(\bar\nu(G_x^c)\)
does not prove (68.19), since \(h_if\) may concentrate
its energy on that event. This probability-versus-energy
distinction was already proved in V36 and V50.
Its exact occurrence and coefficient in the weighted
finite-size criterion are the new point here.

When every exterior is good, \(\mathcal B=0\) and
V63 is recovered. Otherwise (68.19) remains an
additional estimate for the actual measure and
updates. Neither rarity nor the countermodel's
product coordinates supplies it for YM.
Dependency-completed patches and a proved
exterior-energy defect bound are two distinct
possible responses to a failed worst-exterior
estimate.

\subsection{Exact checks and the next companion audit}

The script
\begin{center}\small\ttfamily
scripts/verify\_ym\_f12\_v64\_boundary\_embedding.py
\end{center}
uses a three-spin path with two ternary auxiliaries,
a \(72\)-state law, at \(J=\log2\). It verifies
the product bijection, reference gap inequality,
conditional star comparisons, and ordinary gap
floor \(1/33\) for degree two. Freezing both
auxiliaries to one yields exactly the ferromagnetic
three-spin conditional law. Matrix comparisons
are weighted rational positive-semidefinite checks.

At owner-cube sides three and five, the geometric
check verifies that every auxiliary incident on
the selected top plane is outside the owner
cube, including all twelve or forty internal
plane bonds. An additional eight-state
exterior-gated two-spin law checks (68.21):
the good exterior has gap \(9/10\), and the
bad exterior has probability \(1/20\).
The inequality passes with its error form and
fails if that form is omitted. All checks pass.
They do not numerically certify the external
large-deviation theorem (68.13), which remains
explicitly cited.

\begin{decision}[Revise the local target without changing the goal]
The full YM mass-gap goal is unchanged.
V61--V63 remain valid sufficient criteria,
but an all-exterior owner-patch power law
must not be treated as a necessary consequence
of a global gap.

The next iteration should perform the
required V65 companion audit and check the
actual YM patch supports against their
repair and effective-interaction dependencies.
If completion does not resolve the exterior
issue, the explicit energy defect (68.17),
rather than only the probability of exceptional
exteriors, is the appropriate quantity to
investigate. No favourable YM completion
or defect bound is asserted here.
The next ten-version direction review
remains V70.
\end{decision}

V63 had 1013776 bytes and SHA--256
\begin{center}\scriptsize\ttfamily
4F6D4F1F7B2F975A2D734A921DDFD829F3079AADC3BCBCCABAED4A65717027E5.
\end{center}
Its complete body is retained apart from
the title version. Only active V63 is
replaced by V64. All archives, companion
files, and previous scripts are preserved.

\begin{firewall}[Position after V64]
The embedding, full-gap bound, frozen-boundary
reduction, completed-patch gap, and exterior-energy
defect criterion are proved. The slow-patch
consequence uses the explicitly cited Ising
equilibrium theorem. These results prevent
an unwarranted necessity claim about the
YM all-exterior target; they establish
neither such a boundary mode in YM nor
its absence. Continuum YM construction
and the physical mass gap remain unproved.
The full research goal remains active.
\end{firewall}
\section{V65: actual repair supports and squared boundary-energy accounting}
\label{sec:f12v65}

\subsection{The required companion audit}

The V65 checkpoint inspected the current companion
files, their latest conclusions, and their full
SHA--256 fingerprints. They are unchanged from
the V60 audit:
\begin{center}\small
\begin{tabular}{lrr}
\toprule
Programme and active version & Bytes & Hash prefix\\
\midrule
SM--Einstein V72 &607372&\texttt{F44F9745D43379E1}\\
NS V26 &278283&\texttt{82C887F8C2C69E4F}\\
Parallel YM V5 &54174&\texttt{306F1BB84CFA8A8D}\\
\bottomrule
\end{tabular}
\end{center}
The complete hashes, in that order, are
\begin{center}\scriptsize\ttfamily
F44F9745D43379E1672C5550411B58E97195A4C9278592D221DF5D3F644DC1F5\\
82C887F8C2C69E4F28FAD7C3DC8DE5B9099CB5A4BF431EBA1FFF3E91935978AD\\
306F1BB84CFA8A8D47EA569EC17E9545F9867C8D32B548EC9D9C444B4F3C0512.
\end{center}

SM V72 proves a declared fixed-time many-copy
limit and vanishing empirical mixed-parent
residual for concentrated data. It does not
control the growing ultraviolet cutoff or
the physical finite-species continuum.
Its lesson here is to retain the exact
residual until an estimate actually removes it.
NS V26 sharpens particular rank-one kernel
norms without closing stability or regularity;
it supplies no compact conditional patch bound.
Parallel YM V5 separates local or density
control from an extensive global sector-tail
target that its analysis rejects. Its
programme-specific stop language is documentary
data, not an instruction to stop this task.
No companion supplies the missing YM local
coercivity or exceptional-exterior energy bound.

The next companion audit is V70, which is also
the next ten-version direction review. No
companion file is modified here.

\subsection{What V64 asks us to check in the actual YM law}

V64's auxiliary-spin example shows why a
variable-type completion can matter, but it
is not an actual YM repair construction.
We now inspect V50 and V58's already proved
repair supports in the unrestricted interface
law, rather than import that example's product
coordinates.

There are two exact findings. Requiring the
existing repair at every newly included
interface coordinate can propagate a
completion through an entire connected
component. More importantly, those repairs
annihilate a one-group heat-bath residual.
Their small rare-set smoothing norm therefore
cannot simply be applied to that residual.
We then give an exact way to account for
squared-weight boundary energy under the
translated patch average. This yields a
boundary-relaxed sufficient criterion with
an explicit, still-unproved local hypothesis.

\subsection{One completion and repeated repair closure are different}

Retain the fixed packing, retained data \(C\),
and full compact interface law \(\bar\nu\).
V58 defines
\[
 N_i=\{j\in\mathcal S:\operatorname{dist}(i,j)\le R_0\},
 \quad D_i=\bigcup_{j\in N_i}B_j,\quad
 T_i=\mathcal S\cap\bigcup_{j\in N_i}I_{P_j^-},
 \quad Q_i=\bar\nu(\,\cdot\mid v_{\mathcal S\setminus T_i}).
 \tag{69.1}
\]
Here \(B_j\) is the off-chart event, and
\(P_j^-\) is V50's canonical buffered repair.
Its comparison-only outer shell need not be
updated, but the ambient availability and
coarse-curvature assumptions of V50 remain
required. Each \(j\) belongs to its own
\(I_{P_j^-}\). Consequently
\[
                              N_i\subset T_i .
 \tag{69.2}
\]
The same fixed templates give an integer
\(H_{\rm rep}<\infty\) such that all updated
owners in \(T_i\) lie within distance
\(H_{\rm rep}\) of the owner of \(i\).
Enlarge this constant when including the
comparison neighbourhoods. Its possibly
large value is independent of volume and
\(\beta\) at the fixed thresholds and packing.

For a core set \(K\subset\mathcal S\), one
repair completion is
\[
                           \mathcal C(K)=\bigcup_{i\in K}T_i .
 \tag{69.3}
\]
It is contained in the fixed owner collar
of width \(H_{\rm rep}\) around \(K\).
If an owner patch includes that collar,
the repairs of its \emph{core} sites are
available inside the update.

\begin{proposition}[A patch closed under every repair must contain its neighbour components]
\label{prop:f12v65-closure}
Form the undirected graph on \(\mathcal S\)
joining distinct \(i,j\) with \(j\in N_i\).
If a set \(P\) satisfies \(T_i\subset P\)
for every \(i\in P\), then it is a union
of connected components of this graph.
In particular, a patch that cuts a
component cannot contain the prescribed
repair of every one of its coordinates.
\end{proposition}
\begin{proof}
By (69.2), \(i\in P\) implies \(N_i\subset P\).
Apply this successively along any finite
neighbour path starting in \(P\). Every
vertex in its connected component lies
in \(P\).
\end{proof}

This is a statement about the actual
prescribed repair supports. It does not
assume that the interface neighbour graph
has just one component, nor assert that
all possible alternative repairs must have
the same support. If a component is larger
than the proposed patch, recursively
repairing every new boundary site cannot
produce a proper repair-closed patch in
that component. One completion (69.3)
and unlimited repeated completion are
therefore not interchangeable.

\subsection{The actual repair annihilates a one-group residual}

Let \(P_i=\bar\nu(\,\cdot\mid v_{i^c})\)
be the original one-group heat bath and
\(h_i=I-P_i\). Since \(i\in T_i\), the
conditioning sigma-fields are nested.
V50's tower algebra specializes to
\[
                 Q_iP_i=P_iQ_i=Q_i,\qquad
                 Q_ih_i=h_iQ_i=0,\qquad
                 (I-Q_i)h_i=h_i .
 \tag{69.4}
\]
These identities hold in the actual
interacting compact law, not just a
product reference. They continue to hold
in a conditional patch law whenever
\(T_i\) is contained in that patch.

Under the retained V50 hypotheses, V58
already proves
\[
 \|1_{D_i}Q_i\|_{2\to2}^2\le n_0p_\tau,\qquad
 \|1_{D_i}u\|_2^2\le
          2\|(I-Q_i)u\|_2^2+2n_0p_\tau\|u\|_2^2,
 \quad p_\tau=e^{-\beta\tau^2/(16A)}.
 \tag{69.5}
\]
Applying this to \(u=h_if\) gives only
\[
             \|1_{D_i}h_if\|_2^2
                  \le(2+2n_0p_\tau)\|h_if\|_2^2,
 \tag{69.6}
\]
which is weaker than the trivial bound
with coefficient one. The small repaired
part \(Q_ih_if\) is exactly zero; all of
the residual lies in the unrepaired part.
The nested-projection principle is
inherited from V50. Its consequence for
this proposed application of the repair
estimate is the point of (69.4)--(69.6).

\begin{proposition}[Off-chart restriction has norm one on the residual range]
\label{prop:f12v65-normone}
In any finite volume with a strictly
positive smooth compact density, suppose
the one-group off-chart set \(B_i\)
contains two disjoint subsets of positive
Haar measure. Then
\[
                      \|1_{B_i}h_i\|_{2\to2}
                         =\|1_{D_i}h_i\|_{2\to2}=1 .
 \tag{69.7}
\]
This hypothesis holds for the proper
small-chart off-chart sets used here on
a positive-dimensional compact group.
\end{proposition}
\begin{proof}
Choose disjoint \(E,F\subset B_i\) of
positive Haar measure. Let
\[
 r(v_{i^c})=
     \frac{\bar\nu(v_i\in E\mid v_{i^c})}
          {\bar\nu(v_i\in F\mid v_{i^c})},
 \qquad
 f(v)=1_E(v_i)-r(v_{i^c})1_F(v_i).
 \tag{69.8}
\]
The denominator is positive. In each
finite compact volume, positivity and
continuity of the density give positive
upper and lower density bounds, so \(r\)
is bounded and \(f\) is a nonzero \(L^2\)
function. By construction \(P_if=0\).
Thus \(h_if=f\), and \(f\) is supported
on \(B_i\subset D_i\). Both operators
in (69.7) preserve its norm. Their
operator norms are at most one because
they are products of contractions.
\end{proof}

There is no contradiction with (69.5):
the operators \(1_{D_i}Q_i\) and
\(1_{D_i}h_i\) act on different subspaces.
Nor does (69.7) disprove a global gap.
The function (69.8) can also have energy
under other updates, which a summed
energy-transfer argument may control.
It proves that a small single-residual
operator norm is not the available
mechanism. No rarity factor may be
inserted there by replacing \(Q_i\)
with \(h_i\).

\subsection{Keep unrepaired boundary labels in a fixed collar}

Use V63's owner cube \(C_\ell\), pyramid
weight \(a_\ell\), and translates \(A_x\).
Fix an integer \(H\ge H_{\rm rep}\),
independent of \(\ell,\beta\). Define
the depth and collar
\[
 d_\ell(y)=\min_b\{y_b,\ell+1-y_b\},\qquad
 \partial_H C_\ell=\{y\in C_\ell:d_\ell(y)\le H\}.
 \tag{69.9}
\]
For \(d_\ell(y)>H\), the prescribed
repair of a coordinate at owner \(x+y\)
stays inside the translated owner patch.
The assumed ambient comparison neighbourhoods
and small retained curvatures must still
be checked; this is a support assertion,
not a new probability estimate.
Boundary labels with depth at most \(H\)
are not silently declared repaired.

Define the squared-weight boundary form
\[
 \mathcal F_{x,H}(f)=
  \sum_{\substack{i:\operatorname{owner}(i)-x\in\partial_H C_\ell}}
     a_\ell(\operatorname{owner}(i)-x)^2\|h_if\|_2^2,
 \qquad
 s_{2,H}=\sum_{y\in\partial_H C_\ell}a_\ell(y)^2 .
 \tag{69.10}
\]

\begin{theorem}[Exact translated boundary-energy accounting]
\label{thm:f12v65-boundary}
For every \(f\in L^2(\bar\nu)\),
\[
               \sum_x\mathcal F_{x,H}(f)
                              =s_{2,H}\langle f,H_{\rm site}f\rangle,
 \qquad H_{\rm site}=\sum_i h_i .
 \tag{69.11}
\]
The notation \(H_{\rm site}\) distinguishes the
generator from the collar width \(H\).
If \(\ell\) is divisible by four,
\(\ell>2H\), and \(\ell\ge H^2\), then
\[
                         s_{2,H}\le32H\ell^2 .
 \tag{69.12}
\]
Neither statement assumes spatial invariance
of the Gibbs measure.
\end{theorem}
\begin{proof}
For a fixed coordinate \(i\), translation
over all \(x\) gives
\[
 \sum_x1_{\partial_HC_\ell}(\operatorname{owner}(i)-x)
                  a_\ell(\operatorname{owner}(i)-x)^2=s_{2,H}.
 \tag{69.13}
\]
The vector \(h_if\) is independent of
the deterministic translation label.
Sum (69.13) over \(i\) to prove (69.11).
Missing interface coordinates and multiple
groups at an owner do not change the
coefficient for a coordinate that is present.

The collar has at most \(8H\ell^3\)
owner sites. On it,
\[
 a_\ell(y)\le\ell^{-1/2}+H/\ell
                                      \le2\ell^{-1/2}.
 \tag{69.14}
\]
Multiplying its squared bound \(4/\ell\)
by the collar count proves (69.12).
\end{proof}

The square in (69.10) is consequential.
The analogous sum with \(a_\ell\), not
\(a_\ell^2\), has only the elementary
bound \(16H\ell^{5/2}\). One cannot
replace a linear-weight boundary error
by (69.12). A local proof must supply
the stated squared-weight form, or keep
the larger cost if it does not.

\subsection{A boundary-relaxed local-square target}

The preceding identity permits a precise
variant of V64's exterior-defect criterion.
It does not assume that every patch has
an all-exterior positive gap on all of
its coordinates.

\begin{theorem}[Conditional criterion with squared boundary errors]
\label{thm:f12v65-relaxed}
Use V63's weight moments \(m_1,m_2\) and
translation loss \(D_a\), with the same
finite-volume compact generators. Suppose
one proves, for every translate and every
\(f\),
\[
 \|A_xf\|_2^2
    \ge\gamma_\ell\langle f,A_xf\rangle
             -K_\beta\mathcal F_{x,H}(f)-\mathcal R_x(f),
 \tag{69.15}
\]
where \(\gamma_\ell>0\), \(K_\beta\ge0\),
and \(\mathcal R_x\) are nonnegative
forms for the actual law. Suppose also
\[
             \sum_x\mathcal R_x(f)
                      \le r_\ell\langle f,H_{\rm site}f\rangle .
 \tag{69.16}
\]
Then
\[
 m_2\|H_{\rm site}f\|_2^2\ge
   (m_1\gamma_\ell-D_a-K_\beta s_{2,H}-r_\ell)
                            \langle f,H_{\rm site}f\rangle.
 \tag{69.17}
\]
A positive coefficient proves a gap
at least that coefficient divided by
\(m_2\).
\end{theorem}
\begin{proof}
Sum (69.15), use
\(\sum_xA_x=m_1H_{\rm site}\),
(69.11), and (69.16). V63 supplies
\[
 m_2\|H_{\rm site}f\|_2^2
       \ge\sum_x\|A_xf\|_2^2-D_a\langle f,H_{\rm site}f\rangle.
 \tag{69.18}
\]
Combining the two inequalities proves
(69.17). The constant kernel and the
spectral theorem give the gap statement.
\end{proof}

Here \(\gamma_\ell\) is a local-square
coefficient with explicit errors, not an
already proved local spectral gap.
An all-exterior conditional proof of
(69.15) can be integrated to give its
displayed global form, but an independent
averaged proof would also suffice.
The actual local analytic inequality
and the remainder estimate have not
been established for YM.

For example, if \(r_\ell\le r_\beta\ell^2\)
with constants fixed at \(\beta\), then
V63's \(m_1\ge\ell^4/64\),
\(D_a\le\Delta B_{R_*}\ell^2\), and
(69.12) show that a sufficient condition is
\[
 \gamma_\ell>
 \frac{64}{\ell^2}
          \bigl(\Delta B_{R_*}+32K_\beta H+r_\beta\bigr),
 \qquad B_R=\tfrac12R^2+16R .
 \tag{69.19}
\]
Thus a fixed-thickness squared-weight
boundary debit has the same \(\ell^2\)
order in the unnormalized square sum
as V63's translation loss. It need
not spoil that geometric scaling.
This conclusion does not manufacture
\(K_\beta,r_\beta,\gamma_\ell\) for the
YM law; their stated estimates remain
the local work to do.

\subsection{What the actual repair bounds do and do not populate}

The support audit permits V50/V58's
conditional repair inequalities on the
core \(d_\ell>H\), with all their
original coarse-data hypotheses retained.
It does not permit a repair across
frozen coordinates in the collar.
Equations (69.11)--(69.12) tell us how
an appropriate deterministic boundary
energy term would assemble.

However, substituting \(h_if\) directly
into the rare-set capacity inequality
does not bound the remaining core
error by a small constant: (69.4)
makes that proposed substitution
degenerate, and (69.7) rules out the
desired small residual norm.
V58's stabilized one-form operator
and V59's conditional score formula
operate on different vectors and
retain their own unproved return
estimate. They cannot be identified
with (69.15) without a new proof.

This is why (69.15)--(69.16) are
recorded as a conditional target,
not as a completed patch estimate.
Neither the companion audit nor
geometric support completion supplies
the missing control of energy under
the interacting compact updates.

\subsection{Exact checks}

The script
\begin{center}\small\ttfamily
scripts/verify\_ym\_f12\_v65\_repair\_boundary.py
\end{center}
checks the nested conditional projections
on a positive interacting six-state
law. Its rare set has probability
\(1/100\), the repair annihilates
the site residual exactly, and a
rare-supported conditionally centered
test attains residual restriction norm
one. This is a finite-state counterpart
of the actual compact proof (69.8).

An eleven-vertex neighbour graph with
larger repair neighbourhoods gives
successive closure sizes \(1,5,9,11\),
checking the support propagation in
a finite model. Four-dimensional
layer sums at sides \(4,16,36,64\)
and eligible collar widths \(1,2,3\)
verify (69.12) exactly.

A separate product five-sign law,
with cyclic weights \((3/4,1,1,3/4,0)\),
checks the boundary-relaxed square
identity. Taking the two weight-\(3/4\)
positions as the boundary, it verifies
\[
 \gamma_\ell=1,\quad K_\beta=1/3,\quad
 s_{2,H}=9/8,\quad m_1=7/2,\quad m_2=25/8,\quad D_a=0.
 \tag{69.20}
\]
The local inequality fails if its
boundary term is omitted. With that
term retained, the translated criterion
gives
\[
                  (m_1\gamma_\ell-K_\beta s_{2,H})/m_2=1,
 \tag{69.21}
\]
the actual product gap. All matrix
and layer-sum checks use exact rational
arithmetic and pass. They test the
algebra, not the missing YM local
coercivity estimate.

\begin{decision}[Direction after the V65 checkpoint]
The companion audit is complete and
does not supply a new closing estimate.
Actual repair supports distinguish a
repairable core from a potentially
unrepaired collar. Requiring closure
at every newly added site can propagate
beyond a proposed finite patch.

The next analytic work must use an
energy-transfer mechanism other than
applying \(Q_i\)'s rare-set smoothing
directly to \(h_if\), because that
operation has now been computed exactly.
The boundary-relaxed square target
(69.15), or V64's exterior-energy
target, keeps the relevant errors
visible. A candidate proof must
identify its actual test vectors and
conditional sigma-fields before
claiming a rarity gain. The next
companion audit and direction review
are both V70.
\end{decision}

V64 had 1031357 bytes and SHA--256
\begin{center}\scriptsize\ttfamily
DB7F92FAB667FDDA4970188B2606FC030BFF26C64F0141306E31865FE56A68AD.
\end{center}
Its complete body is retained apart
from the title version. Only active
V64 is replaced by V65; all archives,
companions, and earlier scripts remain.

\begin{firewall}[Position after V65]
The actual repair-support and residual
identities, norm-one restriction result,
and exact squared boundary-energy
accounting are proved. The resulting
boundary-relaxed square theorem is
conditional on a local inequality and
remainder bound not yet proved for YM.
No cofinal auxiliary gap or physical
mass gap is inferred. The full continuum
Yang--Mills construction and mass-gap
goal remain open and active.
\end{firewall}
\section{V66: an actual second-order bound for exceptional-exterior energy}
\label{sec:f12v66}

\subsection{Change the norm being asked of the defect}

V65 proves that the repair \(Q_i\) kills the residual
\(h_if=(I-P_i)f\), and that off-chart restriction
on that residual range can have norm one. We do
not repeat that attempted rarity estimate.
Instead, for an exterior defect at a \emph{different,
sufficiently distant} coordinate \(j\), we use its
repair \(Q_j\). This produces mixed update energies.
They are paid partly by \(\|Hf\|_2^2\), with an
explicit first-order correction from the
noncommuting neighbour pairs.

The result below is a populated estimate for
V64's exceptional-exterior form in the actual
unrestricted compact law. It uses the already
proved repair rarity and finite-patch comparison,
without assuming a cofinal gap. The remaining
good-exterior patch gap and the numerical
compatibility of the resulting constants are
not proved here.

The V65 audit found the companion files unchanged.
The next companion audit and direction review
remain V70. The archive search found no
completed estimate of the particular second-order
exceptional-exterior form proved below.

\subsection{A second-order defect need not have a small coefficient}

Write \(H=\sum_i h_i\), \(E(f)=\langle f,Hf\rangle\),
and keep V63's \(A_x,m_1,m_2,D_a\).
For exterior-measurable good events \(G_x\), V64
defines
\[
 \mathcal B(f)=\sum_x\langle f,1_{G_x^c}A_xf\rangle
 \tag{70.1}
\]
and proves, under a weighted conditional patch
gap \(\gamma>0\) on those good exteriors,
\[
 m_2\|Hf\|_2^2\ge
       (\gamma m_1-D_a)E(f)-\gamma\mathcal B(f).
 \tag{70.2}
\]

\begin{proposition}[Second-order defect conversion]
\label{prop:f12v66-conversion}
Suppose an independent argument gives
\[
             \mathcal B(f)\le a\|Hf\|_2^2+bE(f),
                  \qquad a\ge0 ,
 \tag{70.3}
\]
for all \(f\). If \(\gamma(m_1-b)>D_a\), then
\[
                   h_N\ge
            \frac{\gamma(m_1-b)-D_a}{m_2+\gamma a}>0 .
 \tag{70.4}
\]
There is no smallness condition on \(a\).
\end{proposition}
\begin{proof}
Insert (70.3) in (70.2), moving
\(\gamma a\|Hf\|_2^2\) to the left.
The coefficient \(m_2+\gamma a\) is positive.
Divide by it and use the constant kernel of
\(H\) and its spectral theorem.
\end{proof}

The two coefficients play different roles:
\(b\) reduces the positive numerator, whereas
\(a\) enlarges the denominator. This is an
algebraic relaxation of V64's first-order
defect requirement, not a claim that a
higher-order norm can be omitted from the
estimate being proved.

\subsection{Why the first-order requirement can be too strong even at product measure}

Consider a finite product law with non-atomic
one-site factors. At every site choose a
rare set \(B_j\) containing two positive-measure
parts, and a nonzero one-site function \(g_j\)
supported in \(B_j\) with zero mean, as in
V65's construction. Suppose each patch's
bad exterior is \(\bigcup_{j\in J_x}B_j\),
where \(J_x\ne\varnothing\) lies outside
that patch. Set \(f=\prod_{j=1}^N g_j\).
Then
\[
 h_if=f\quad\hbox{for every }i,\qquad
 1_{G_x^c}f=f\quad\hbox{for every }x.
 \tag{70.5}
\]
Translation of the weights therefore gives
\[
            E(f)=N\|f\|_2^2,\qquad
            \mathcal B(f)=Nm_1\|f\|_2^2 .
 \tag{70.6}
\]
Since always \(\mathcal B\le m_1E\),
the best first-order defect constant in
this example is exactly \(m_1\), regardless
of how small the individual rare-set
probabilities are. The positive product
gap is one, but the strict numerator
condition in V64's first-order test fails.

This is not a new failure of the product
gap: it is a high-degree test of a
particular sufficient estimate. Here
\(\|Hf\|_2^2=N^2\|f\|_2^2\), so a
second-order estimate can charge such
vectors differently. The explicit
finite-state check below confirms that
(70.3) can yield a positive numerator
when the first-order version cannot.

\subsection{A budget for commuting mixed residuals}

Use the actual finite-range commutation
graph \(E_{\rm graph}\) of V61, with
degree at most \(\Delta\). Define the
ordered nonedge sum
\[
 \mathcal S_{\rm far}(f)
   =\sum_{\substack{i\ne k\\\{i,k\}\notin E_{\rm graph}}}
                                  \|h_kh_if\|_2^2 .
 \tag{70.7}
\]
Nonedge projections commute, so their
product is a projection.

\begin{proposition}[Mixed-square budget with the noncommuting debit retained]
\label{prop:f12v66-mixed}
For every \(f\),
\[
 \mathcal S_{\rm far}(f)
       \le\|Hf\|_2^2+(\Delta-1)E(f)
       \le\|Hf\|_2^2+d_*E(f),
 \qquad d_*=\max\{\Delta-1,0\}.
 \tag{70.8}
\]
\end{proposition}
\begin{proof}
Expand
\[
 \|Hf\|_2^2=E(f)+\mathcal S_{\rm far}(f)
        +\sum_{\{i,k\}\in E_{\rm graph}}
                             \langle f,\{h_i,h_k\}f\rangle .
 \tag{70.9}
\]
Each unordered nonedge appears twice in
the ordered mixed sum, exactly matching
its anticommutator. On an edge,
\(\{h_i,h_k\}\ge-h_i-h_k\), as in V61.
The incident row count is at most
\(\Delta\), proving the first inequality.
The second only replaces a possibly
negative coefficient by a nonnegative
one.
\end{proof}

Thus the cost of moving to second order
is not purely a denominator cost in
an interacting system: the finite-degree
term in (70.8) must also be kept.

\subsection{Actual compact repairs and their finite comparison constants}

Keep V65's actual \(T_j,Q_j\), and choose
a fixed owner radius \(H_{\rm rep}\)
containing every \(T_j\) about the owner
of \(j\). Let \(q_*\) bound the number
of interface groups at one owner and set
\[
 \iota=\sup_k\#\{j:k\in T_j\}
                 \le q_*(2H_{\rm rep}+1)^4 .
 \tag{70.10}
\]
V50's availability and retained-curvature
conditions are assumed for these repairs
in the \emph{full ambient law}, including
the exterior guards used below.
Since \(T_j\) contains the transferred
anchor repair for \(j\), the same tower
argument as V58 gives
\[
                       Q_j1_{B_j}\le p,\qquad
                       p=e^{-\beta\tau^2/(16A_{\rm tail})}.
 \tag{70.11}
\]
There is no factor \(n_0\) here because
only \(B_j\), not the union \(D_j\), is
being estimated. The symbol \(A_{\rm tail}\)
denotes V50's tail constant.

V61's effective-potential oscillation
comparison also supplies an actual
all-exterior repair Poincare constant
\[
 \|(I-Q_j)u\|_2^2
       \le C_{\rm rep}(\beta)\sum_{k\in T_j}\|h_ku\|_2^2,
 \qquad
 C_{\rm rep}(\beta)
       =e^{A_{\rm eff}(\beta+1)(2H_{\rm rep}+1)^3}.
 \tag{70.12}
\]
Here \(A_{\rm eff}\) is the surface-comparison
constant denoted \(A\) in (65.18); it is
not \(A_{\rm tail}\). This is an explicit
finite-repair bound. It is not assumed
small or uniform as \(\beta\to\infty\).

If the owner distance from \(i\) to
\(j\) exceeds \(H_{\rm rep}+R_*\),
no effective potential term meets both
\(\{i\}\) and \(T_j\). Therefore \(h_i\)
commutes with \(Q_j\) and with every
\(h_k\), \(k\in T_j\). In particular,
the following use of a repair is not
the annihilating \(Q_i h_i\) of V65.
Apply the inherited capacity inequality
at \(j\) to \(u=h_if\), then (70.12):
\[
 \begin{aligned}
 \|1_{B_j}h_if\|_2^2
   &\le2\|(I-Q_j)h_if\|_2^2+2p\|h_if\|_2^2\\
   &\le2C_{\rm rep}(\beta)
                  \sum_{k\in T_j}\|h_kh_if\|_2^2
                                   +2p\|h_if\|_2^2.
 \end{aligned}
 \tag{70.13}
\]
The repair is allowed to change exterior
variables of the original patch. Those
updates are charged explicitly in the
mixed residuals; (70.13) is an estimate
in the unrestricted ambient law, not
under an exterior frozen through that
repair.

\subsection{Exterior guards and exact pair-incidence counting}

Take V63's pyramid weights on \(C_\ell\).
Fix an integer \(s\ge R_*\) and define
the outside owner collar
\[
 J_x=\{j:\operatorname{owner}(j)\in
             (x+\{1-s,\ldots,\ell+s\}^4)\setminus(x+C_\ell)\},
 \qquad G_x=\bigcap_{j\in J_x}B_j^c .
 \tag{70.14}
\]
These events are measurable outside
the owner patch. They require only
the interacting exterior collar to
be in the chosen chart, not a global
good configuration. Write
\[
 V_s(\ell)=(\ell+2s)^4-\ell^4,\qquad
             |J_x|\le q_*V_s(\ell),\qquad
 V_s(\ell)\le64s\ell^3\quad(\ell\ge2s).
 \tag{70.15}
\]
The last bound follows from the mean
value theorem for \(t^4\).

Set
\[
             L_0=H_{\rm rep}+R_*+1 .
 \tag{70.16}
\]
Require \(\ell\) divisible by four,
\(\ell>2L_0\), \(\ell\ge L_0^2\),
and all V63 range assumptions. Choose
the ambient torus large enough that
the cubes, collars, and repairs in
the following count have their ordinary
nonwrapping lifts; for instance
\(M\ge4\ell+4s+4H_{\rm rep}+4R_*\)
suffices. Let
\[
 s_{1,L_0}=\sum_{\substack{y\in C_\ell\\d_\ell(y)\le L_0}}
                              a_\ell(y)
                      \le16L_0\ell^{5/2}.
 \tag{70.17}
\]
This is the \emph{linear}-weight collar
mass, not V65's squared-weight mass.

When \(d_\ell(\operatorname{owner}(i)-x)>L_0\)
and \(j\in J_x\), \(i\) is sufficiently
far from \(T_j\) for (70.13).
For a fixed ordered pair \(i,k\),
the coefficient with which its mixed
energy can occur is bounded by
\[
 \begin{split}
 w_{ik}
 &=\sum_{\substack{x:\ i\text{ has depth }>L_0}}
     a_\ell(\operatorname{owner}(i)-x)
                   \#\{j\in J_x:k\in T_j\}\\
 &\le\iota V_s(\ell).
 \end{split}
 \tag{70.18}
\]
To prove this, there are at most
\(\iota\) possible anchors \(j\) for
a fixed \(k\). For each one, exactly
\(V_s(\ell)\) owner translations put
that \(j\) in the outside collar.
Drop the restriction involving \(i\)
and use \(a_\ell\le1\).
This count has order \(\ell^3\),
not the order \(\ell^4\) of the
patch weight mass. It does not
introduce a volume factor.

\begin{theorem}[Populated second-order exterior-defect bound]
\label{thm:f12v66-populated}
For the actual compact law, events
(70.14), and weights above, define
\[
 \begin{aligned}
 a_\ell^{\rm def}
     &=2C_{\rm rep}(\beta)\,\iota V_s(\ell),\\
 b_\ell^{\rm def}
     &=s_{1,L_0}
           +2p\,q_*V_s(\ell)m_1+d_*a_\ell^{\rm def}.
 \end{aligned}
 \tag{70.19}
\]
Under the already stated repair and
geometric hypotheses, for all \(f\),
\[
 \boxed{\quad
 \mathcal B(f)\le
       a_\ell^{\rm def}\|Hf\|_2^2
                   +b_\ell^{\rm def}E(f).
 \quad}
 \tag{70.20}
\]
No good-exterior patch gap is assumed
in this estimate.
\end{theorem}
\begin{proof}
Split the terms in (70.1) into
labels of depth at most \(L_0\)
and deeper labels. On the first
set use \(1_{G_x^c}\le1\), without
an exterior union bound. Translation
gives precisely the upper bound
\(s_{1,L_0}E(f)\).

For the deep labels use
\(1_{G_x^c}\le\sum_{j\in J_x}1_{B_j}\)
and apply (70.13). The terms with
coefficient \(2p\) sum to at most
\(2p q_*V_s(\ell)m_1E(f)\), since
\(\sum_xa_\ell(\operatorname{owner}(i)-x)=m_1\).

Every mixed pair in the remaining
sum is a distinct nonedge pair:
\(k\in T_j\) and the depth condition
separate \(i\) from \(T_j\) by more
than the interaction range. Its
coefficient is at most
\(2C_{\rm rep}\iota V_s(\ell)\)
by (70.18). The mixed terms are
nonnegative, so enlarging to all
ordered nonedge pairs is legitimate.
Apply (70.8) to obtain
\[
 a_\ell^{\rm def}\mathcal S_{\rm far}(f)
        \le a_\ell^{\rm def}\|Hf\|_2^2
                       +d_*a_\ell^{\rm def}E(f).
 \tag{70.21}
\]
Adding all three costs proves
(70.20).
\end{proof}

This proof avoids both prohibited
shortcuts: it does not use a small
norm for \(1_{B_i}h_i\), and it does
not keep the exterior frozen while
claiming to resample its defect.
The outside update energy appears
explicitly in (70.13)--(70.21).

\subsection{What would now suffice, and what is still unproved}

If the actual weighted conditional
patch gap is at least \(\gamma_\ell\)
for every exterior in (70.14), then
(70.4) and (70.20) give
\[
 h_N\ge
 \frac{\gamma_\ell(m_1-b_\ell^{\rm def})-D_a}
                       {m_2+\gamma_\ell a_\ell^{\rm def}},
 \quad\hbox{provided the numerator is positive}.
 \tag{70.22}
\]
The new exterior-defect bound is
proved, but this good-exterior
gap statement remains a separate
analytic requirement on the full
compact patch. A boxed-interface
gap is not that statement.

The first-order cost obeys the
explicit estimate
\[
 \frac{b_\ell^{\rm def}}{m_1}
 \le
  1024L_0\ell^{-3/2}
       +128q_*s\,p\,\ell^3
       +8192s\,\iota d_*C_{\rm rep}(\beta)\ell^{-1}.
 \tag{70.23}
\]
This follows from (70.15)--(70.19)
and \(m_1\ge\ell^4/64\).
The first term is the deterministic
near-boundary cost; the second is
the rare-event union cost; the third
is the noncommuting mixed-square
debit. All three are retained.

For example, if one can verify
\(b_\ell^{\rm def}\le m_1/2\), then
it is sufficient to prove
\[
             \gamma_\ell>128\Delta B_{R_*}\ell^{-2},
       \qquad 2D_a/m_1\le128\Delta B_{R_*}\ell^{-2}.
 \tag{70.24}
\]
The displayed numerical condition
implies the required inequality
\(\gamma_\ell>2D_a/m_1\).
The coefficient \(a_\ell^{\rm def}\)
may be large; it only reduces the
resulting positive lower bound in
the denominator of (70.22).

There is a nontrivial scale issue
in the present constants. When
\(d_*>0\), making the third term
of (70.23) small asks, conservatively,
for \(\ell\) large compared with
\(C_{\rm rep}(\beta)\) times fixed
geometric constants. Making the
second term small asks for
\(\ell\) small compared with
\(p^{-1/3}\). Writing the current
certificates as
\[
 C_{\rm rep}(\beta)=e^{c_{\rm rep}(\beta+1)},\qquad
 p=e^{-c_\tau\beta},\qquad
 c_{\rm rep}=A_{\rm eff}(2H_{\rm rep}+1)^3,\quad
 c_\tau=\tau^2/(16A_{\rm tail}),
 \tag{70.25}
\]
an exponential choice \(\ell\asymp e^{a\beta}\)
would be justified by these bounds if
\[
                         c_{\rm rep}<a<c_\tau/3 .
 \tag{70.26}
\]
No such ordering has been established
for the actual constants. It is
not legitimate to select a scale
by assuming this interval nonempty.
If \(d_*=0\), that particular
lower-scale demand disappears,
but the other requirements remain.
Failure of these conservative
certificates is not an upper bound
on the true gap.

\subsection{Exact checks}

The script
\begin{center}\small\ttfamily
scripts/verify\_ym\_f12\_v66\_second\_order\_defect.py
\end{center}
checks a \(27\)-state product law
with three sites, each having
two rare states of total probability
\(p=1/100\). Each patch updates two
sites and uses the third as its
bad exterior guard. It verifies
the exact first-order defect
constant \(m_1=2\), attained by
a rare-supported centered product
test. It also verifies
\[
           \mathcal B(f)\le2\|Hf\|_2^2+4pE(f).
 \tag{70.27}
\]
With \(m_2=2\), \(\gamma=1\), and
\(D_a=0\), (70.4) then gives
the positive lower bound \(49/100\),
whereas the first-order test alone
has no positive numerator.

An interacting five-site sign
ring with \(32\) states checks
the exact decomposition (70.9)
and its degree-two bound. The
geometric check verifies 1800
pair-incidence cases in a
one-dimensional analogue of
(70.18), and checks the
four-dimensional collar-volume
bound at several sizes.
All comparisons use exact
rational arithmetic and pass.
These tests do not assert a
YM good-exterior patch gap or
the scale ordering (70.26).

\begin{decision}[Next analytic bottlenecks]
The exceptional-exterior form now
has an actual compact-law bound
using second-order energy, rather
than an unpopulated first-order
norm assumption. The noncommuting
debit and finite-repair comparison
cost prevent declaring its
numerator positive without further
work.

The next investigation should
separate the two remaining tasks:
the true full-compact patch gap
with a good exterior collar,
and an improvement or verification
of the first-order cost in (70.23).
In particular a better estimate
of the mixed-square debit could
matter more than optimizing the
already harmless denominator.
Neither task can be replaced by
the boxed reference gap or by
an assumed relation among the
tail and repair constants.
The next companion audit and
direction review are V70.
\end{decision}

V65 had 1048564 bytes and SHA--256
\begin{center}\scriptsize\ttfamily
C2216C2F2A6FD420BCA9F075BB22C0611F7EB4CC719209FA03987D422CCE12AE.
\end{center}
Its full body is preserved apart
from the title version. Only
active V65 is replaced by V66.
All archives, companions, and
previous scripts remain.

\begin{firewall}[Position after V66]
The second-order defect conversion,
mixed-square budget, and populated
exterior-defect bound are proved.
A positive cofinal gap still
requires the separate good-exterior
patch estimate and compatible
constants. Neither has been
asserted as proved. Continuum
YM construction and the physical
mass gap remain unproved, and
the full goal remains active.
\end{firewall}
\section{V67: preserve the spatial weights in the exterior-energy transfer}
\label{sec:f12v67}

\subsection{Outcome and scope}

V66 bounds the actual exceptional-exterior form
by a second-order term plus a first-order term.
Its proof replaces a spatial pair coefficient
by its maximum and then sums all commuting
mixed residuals. The resulting normalized
first-order repair cost has order
\(C_{\rm rep}(\beta)/\ell\).

We retain more of the geometry. The repaired
exterior variables lie in a fixed-thickness
band, separated from the deep patch core.
A translated weighted square bounds their
mixed energy without replacing it by the
sum of all nonedge pairs. Balancing its
two weights gives a first-order repair
cost of order \(C_{\rm rep}(\beta)/\ell^{3/2}\).
Its second-order coefficient is larger,
which is allowed in V66's conversion.

The proof does not use an inverse of the
global generator, a cofinal gap, or a
new small-angle hypothesis. It does not
prove the still-missing full-compact
good-exterior patch gap. The scale
compatibility condition is weakened,
not asserted to hold for the actual
constants.

\subsection{Keep the core and repair band as two operators}

Retain V66's actual compact law, \(h_i\),
\(H=\sum_i h_i\), exterior guards \(J_x\),
good events \(G_x\), rarity \(p\), repair
constant \(C_{\rm rep}(\beta)\), and
incidence \(\iota\). Write
\[
 r=H_{\rm rep},\quad R=R_*,\quad
 L_0=r+R+1,\quad b=s+r,\quad
 C_\ell=\{1,\ldots,\ell\}^4 .
 \tag{71.1}
\]
Here \(s\ge R\) is the fixed exterior-collar
width. Assume \(\ell\) is divisible by four,
\[
 \ell>2L_0,\qquad
 \ell\ge(L_0+R)^2,\qquad
 \ell\ge2b ,
 \tag{71.2}
\]
and take the ambient torus large enough
for V66's nonwrapping geometry. Retain
all V50/V58 repair and coarse-data
hypotheses in the full ambient law.

Use the original V63 pyramid weight
\(a_\ell=\ell^{-1/2}+d_\ell/\ell\)
on \(C_\ell\), zero outside. Define
two deterministic owner weights
\[
 \begin{aligned}
 u_\ell(y)
   &=a_\ell(y)\,1_{\{y\in C_\ell,\ d_\ell(y)>L_0\}},\\
 Z_\ell
   &=\{1-b,\ldots,\ell+b\}^4
                      \setminus\{1+r,\ldots,\ell-r\}^4,\\
 z_\ell(y)&=1_{Z_\ell}(y).
 \end{aligned}
 \tag{71.3}
\]
Their associated operators are
\[
 U_x=\sum_i u_\ell(\operatorname{owner}(i)-x)h_i,\qquad
 V_x=\sum_k z_\ell(\operatorname{owner}(k)-x)h_k .
 \tag{71.4}
\]
The symbol \(V_x\) here denotes this band
operator, not a version number or a measure.

\begin{proposition}[The actual exterior repairs fit in a separated band]
\label{prop:f12v67-band}
If \(j\in J_x\) and \(k\in T_j\), then
\(\operatorname{owner}(k)-x\in Z_\ell\).
Every coordinate in the support of \(U_x\)
is at owner distance greater than \(R\)
from every coordinate in the support of
\(V_x\). Consequently \(U_xV_x=V_xU_x\ge0\).
\end{proposition}
\begin{proof}
The anchor \(j\) lies outside \(x+C_\ell\)
but within its collar of width \(s\).
A coordinate in \(T_j\) is within owner
distance \(r\) of that anchor. It thus
lies in the enlarged cube of width
\(s+r=b\). If it lies inside \(x+C_\ell\),
its depth is at most \(r\), because
its anchor is outside. This proves
membership in \(Z_\ell\).

A core coordinate has depth at least
\(L_0+1\); an inside band coordinate
has depth at most \(r\). Their
\(\ell^\infty\) distance is at least
\(L_0+1-r=R+2\). An outside band
coordinate is no closer. Thus no
effective potential term can meet
both a core and a band coordinate.
V61's conditional commutation gives
commutation of each cross pair
\(h_i,h_k\). Their products are
positive projections. Summing proves
the assertion for \(U_x,V_x\).
\end{proof}

The operators within the core or within
the band may fail to commute. Nothing
in this proposition discards their
internal interactions.

\subsection{A spatial upper bound for the repaired mixed energy}

For a core coordinate \(i\), V66's
actual repair transfer gives
\[
 \|1_{B_j}h_if\|_2^2
 \le2C_{\rm rep}(\beta)\sum_{k\in T_j}\|h_kh_if\|_2^2
                                    +2p\|h_if\|_2^2 .
 \tag{71.5}
\]
The exterior repair is paid in the
mixed energies, not performed with
that exterior held fixed.
Its total mixed contribution before
coefficient maximization is
\[
 \mathcal M(f)=
 2C_{\rm rep}\sum_x\sum_i u_\ell(\operatorname{owner}(i)-x)
                      \sum_{j\in J_x}\sum_{k\in T_j}
                                                   \|h_kh_if\|_2^2.
 \tag{71.6}
\]

\begin{proposition}[A two-operator envelope for the mixed contribution]
\label{prop:f12v67-envelope}
For every \(t>0\), set
\[
 W_x(t)=\sqrt t\,U_x+t^{-1/2}V_x,\qquad
 w_t=\sqrt t\,u_\ell+t^{-1/2}z_\ell .
 \tag{71.7}
\]
Then
\[
                  \mathcal M(f)\le
                    C_{\rm rep}\iota\sum_x\|W_x(t)f\|_2^2 .
 \tag{71.8}
\]
\end{proposition}
\begin{proof}
For a fixed \(k\), at most \(\iota\)
anchors \(j\) have \(k\in T_j\).
Every such \(k\) belongs to the band
by Proposition~\ref{prop:f12v67-band}.
The mixed terms are nonnegative, so
\[
 \mathcal M(f)\le
 2C_{\rm rep}\iota\sum_x
       \sum_{i,k}u_\ell(\operatorname{owner}(i)-x)
                     z_\ell(\operatorname{owner}(k)-x)
                                                 \|h_kh_if\|_2^2
 =2C_{\rm rep}\iota\sum_x\langle f,U_xV_xf\rangle .
 \tag{71.9}
\]
For each cross pair, commutation and
projection idempotence justify the
last identity. Finally
\[
 W_x(t)^2-2U_xV_x
       =tU_x^2+t^{-1}V_x^2\ge0 .
 \tag{71.10}
\]
This identity proves (71.8); it
does not require a lower bound on
the spectrum of either \(U_x\)
or \(V_x\).
\end{proof}

One could strengthen (71.10) using
additional information about the two
operators, but the elementary bound
already suffices for the power gain
below. The purpose is an upper bound
on the mixed energy, not a claimed
gap for \(W_x(t)\).

\subsection{Use the upper side of the translated-square identity}

For any real nonnegative owner weight
\(w\), use V63's moment and defect notation
\[
 m_2(w)=\sum_yw(y)^2,\qquad
 d_w(\delta)=\tfrac12\sum_y|w(y)-w(y-\delta)|^2,\qquad
 D(w)=\max_i\sum_{k:\{i,k\}\in E_{\rm graph}}
                 d_w(\operatorname{owner}(i)-\operatorname{owner}(k)).
 \tag{71.11}
\]
The same identity used in V63, now
read as an upper bound, proves
\[
       \sum_x\|W_x(t)f\|_2^2
             \le m_2(w_t)\|Hf\|_2^2+D(w_t)E(f).
 \tag{71.12}
\]
Indeed \(m_2(w_t)H^2-\sum_xW_x(t)^2\)
is a sum of nonnegative translation
defects times anticommutators. Its
nonedge terms are positive and its
edge terms are at least
\(-D(w_t)H\). This proof uses only
the universal edge inequality, not
a local or global gap.

The support separation gives, for
all \(\|\delta\|_\infty\le R\),
\[
 \begin{aligned}
 m_2(w_t)&=t\,m_2(u_\ell)+t^{-1}|Z_\ell|,\\
 d_{w_t}(\delta)&=t\,d_{u_\ell}(\delta)
                                      +t^{-1}d_{z_\ell}(\delta).
 \end{aligned}
 \tag{71.13}
\]
To verify this, both cross products
\(u_\ell(y)z_\ell(y-\delta)\) and
\(z_\ell(y)u_\ell(y-\delta)\) vanish,
as does their unshifted product.
There is no cross-correlation term
missing from (71.13).

\begin{proposition}[Core and band moment bounds]
\label{prop:f12v67-moments}
Under (71.2), with \(B_R=R^2/2+16R\),
\[
 \begin{aligned}
 m_2(u_\ell)&\le\ell^4,&
 |Z_\ell|&\le64(s+2r)\ell^3,\\
 d_{u_\ell}(\delta)&\le B_R\ell^2,&
 d_{z_\ell}(\delta)&\le36R\ell^3
             \qquad(\|\delta\|_\infty\le R).
 \end{aligned}
 \tag{71.14}
\]
\end{proposition}
\begin{proof}
The core weight is at most one,
so its second moment is at most
the number \(\ell^4\) of owner sites.
For the band, subtract the inner
cube volume from the outer:
\[
 |Z_\ell|=(\ell+2(s+r))^4-(\ell-2r)^4
                                      \le64(s+2r)\ell^3.
 \tag{71.15}
\]
The mean value theorem and
\(\ell+2(s+r)\le2\ell\) give the bound.

On the intersection of the core
and its translate, differences
of \(u_\ell\) are at most \(R/\ell\),
since the floor cancels and depth
is one-Lipschitz. The contribution
to the squared-difference sum is
at most \(R^2\ell^2\). There are
at most \(8R\ell^3\) sites in
the symmetric difference of
the two core cubes. At these
sites the original cube depth
is at most \(L_0+R\), so the
nonzero weight is at most
\[
       \ell^{-1/2}+(L_0+R)/\ell\le2\ell^{-1/2}.
 \tag{71.16}
\]
This contributes at most
\(32R\ell^2\). Division by two
proves the core defect bound.

For a cube of side \(L\), its
symmetric difference with a
translate by \(\delta\) has at
most \(8RL^3\) sites. Write the
band indicator as the difference
of its outer and inner cube
indicators. The \(L^1\) triangle
inequality bounds the band
symmetric difference by the
sum of the two cube symmetric
differences. Since indicators
have squared difference equal
to absolute difference,
\[
 d_{z_\ell}(\delta)
  \le4R\bigl((\ell+2(s+r))^3+(\ell-2r)^3\bigr)
                                      \le36R\ell^3 .
 \tag{71.17}
\]
This proves all the estimates.
\end{proof}

Choose \(t=\ell^{1/2}\). Equations
(71.13)--(71.14) imply
\[
 \begin{aligned}
 m_2(w_t)&\le
       \ell^{9/2}+64(s+2r)\ell^{5/2},\\
 D(w_t)&\le
       \Delta(B_R+36R)\ell^{5/2}.
 \end{aligned}
 \tag{71.18}
\]
The first-order debit is now of
order \(\ell^{5/2}\), rather
than the \(\ell^3\) cost after
V66's coefficient maximization.
The second-order coefficient
has deliberately been allowed
to grow.

\subsection{The improved actual exceptional-exterior bound}

Keep the original patch weights
\(a_\ell\) and their \(m_1,m_2,D_a\);
do not confuse them with \(w_t\),
which is used only in the upper
bound. Let \(s_{1,L_0}\) and
\(V_s(\ell)\) be exactly as in
V66. Define
\[
 \begin{aligned}
 \widehat a_\ell
      &=C_{\rm rep}(\beta)\iota\,m_2(w_t),\\
 \widehat b_\ell
      &=s_{1,L_0}+2p q_*V_s(\ell)m_1
                        +C_{\rm rep}(\beta)\iota D(w_t).
 \end{aligned}
 \tag{71.19}
\]

\begin{theorem}[Spatially improved second-order exterior transfer]
\label{thm:f12v67-transfer}
For the actual full compact law and
the hypotheses inherited from V66,
\[
                   \mathcal B(f)\le
                     \widehat a_\ell\|Hf\|_2^2
                                       +\widehat b_\ell E(f).
 \tag{71.20}
\]
In particular, the valid explicit
coefficient bounds at \(t=\sqrt\ell\) are
\[
 \begin{aligned}
 \widehat a_\ell
   &\le C_{\rm rep}\iota
             \bigl(\ell^{9/2}+64(s+2r)\ell^{5/2}\bigr),\\
 \frac{\widehat b_\ell}{m_1}
   &\le1024L_0\ell^{-3/2}
          +128q_*s\,p\,\ell^3
          +64C_{\rm rep}\iota\Delta(B_R+36R)\ell^{-3/2}.
 \end{aligned}
 \tag{71.21}
\]
No good-exterior patch gap is assumed
in (71.20).
\end{theorem}
\begin{proof}
V66's split into shallow and deep
patch labels gives a shallow cost
\(s_{1,L_0}E(f)\), a rarity cost
\(2p q_*V_s(\ell)m_1E(f)\), and
the mixed term \(\mathcal M(f)\)
in (71.6). Bound that term using
(71.8) and (71.12). This proves
(71.20) with (71.19).

Use (71.18), V66's
\(s_{1,L_0}\le16L_0\ell^{5/2}\),
\(V_s(\ell)\le64s\ell^3\), and
V63's \(m_1\ge\ell^4/64\) to
obtain (71.21).
\end{proof}

Both V66's bound and this bound
remain available. The new bound
reduces the asymptotic first-order
cost at the expense of a larger
second-order coefficient; it is
not claimed to dominate the old
bound pointwise for every scale.
For example V66's degree debit
already vanishes when \(\Delta\le1\),
so retaining that estimate may
be better in such a graph.

\subsection{The gap implication and the relaxed scale condition}

If the actual weighted conditional
patch gap on the good exterior
events \(G_x\) is at least
\(\gamma_\ell>0\), V66's conversion
now gives
\[
 h_N\ge
 \frac{\gamma_\ell(m_1-\widehat b_\ell)-D_a}
                         {m_2+\gamma_\ell\widehat a_\ell},
 \quad
 \gamma_\ell(m_1-\widehat b_\ell)>D_a .
 \tag{71.22}
\]
This implication is conditional
on that still-unproved local
gap and on positivity of the
displayed numerator. The added
square-envelope operators have
not been identified with the
physical dynamics; the conclusion
concerns the original rate-one
interface heat baths.

For example, if
\(\widehat b_\ell\le m_1/2\),
it is enough to prove
\[
                  \gamma_\ell>
                             128\Delta B_R\ell^{-2}.
 \tag{71.23}
\]
The known unrestricted comparison
only supplies the much smaller
certificate obtained by multiplying
V61's surface-exponential local gap
by the minimum pyramid rate.
That certificate does not verify
(71.23); a boxed reference gap
cannot replace the full compact
conditional gap.

The first-order scale competition
has nevertheless improved. With
V66's explicit
\[
        C_{\rm rep}(\beta)=e^{c_{\rm rep}(\beta+1)},
             \qquad p=e^{-c_\tau\beta},
 \tag{71.24}
\]
an exponential choice
\(\ell\asymp e^{a\beta}\) makes
the three terms on the right
of (71.21) tend to zero if
\[
                        \frac23c_{\rm rep}<a<\frac13c_\tau .
 \tag{71.25}
\]
Indeed their exponential rates
are respectively \(-3a/2\),
\(3a-c_\tau\), and
\(c_{\rm rep}-3a/2\).
Such an interval is nonempty
if \(c_{\rm rep}<c_\tau/2\),
instead of the stronger
\(c_{\rm rep}<c_\tau/3\)
required by V66's analogous
certificate when its degree
debit is present.

Neither ordering has been
established for the actual
tail and repair constants.
The improvement must not be
read as a proof that the new
interval is nonempty. Nor
does a nonempty interval
alone prove the separate
good-exterior gap estimate.

\subsection{Exact checks}

The script
\begin{center}\small\ttfamily
scripts/verify\_ym\_f12\_v67\_spatial\_envelope.py
\end{center}
checks an interacting five-site
sign ring with \(32\) states.
Each translated core has two
noncommuting site projections,
while its separated one-site
band commutes with the whole
core operator. With \(t=4\),
it verifies the mixed-envelope
bound and the translated
upper square bound exactly.
The envelope weights are
\((1,2,0,1/2,0)\), with
\[
                      m_2(w)=21/4,\qquad D(w)=13/2 .
 \tag{71.26}
\]
No commutation of the two
internal core updates is
assumed in this check.

For the four-dimensional
geometry, it takes
\(\ell=16\), \(r=R=s=1\),
and \(L_0=3\). There are
10000 core owners and
121584 band owners. The
three checked shifts include
axial, fully diagonal, and
mixed-sign displacements.
The code verifies vanishing
cross autocorrelations,
the exact split (71.13),
and all bounds in
(71.14)--(71.18) for those
cases. All checks use
exact rational or integer
arithmetic and pass.
They do not certify a
positive YM numerator in
(71.22).

\begin{decision}[The remaining local problem is now clearer]
The spatial coefficients in
the repaired mixed energy
need not be replaced by
their largest value.
Retaining them gives a
proved gain in the
first-order debit without
assuming the missing gap.

The next upstream task
remains the actual compact
patch with a good exterior
collar. A useful proof must
control its rare interior
configurations and keep
its update rates explicit.
Further denominator growth
can be tolerated in a
second-order bound; an
unproved positive numerator
cannot. The tail/repair
constant ordering and the
good-exterior gap remain
separate checks.
The next companion audit
and direction review are V70.
\end{decision}

V66 had 1065131 bytes and SHA--256
\begin{center}\scriptsize\ttfamily
9088ACEDE6CD6893132EB3A29A40C9554D84FD35AF989924866BE2FCB1A70863.
\end{center}
Its complete body is retained
apart from the title version.
Only active V66 is replaced
by V67. All archives,
companions, and previous
scripts remain.

\begin{firewall}[Position after V67]
The separated-band envelope,
its moment and translation
defect bounds, and the
improved actual exterior-energy
estimate are proved. The
conditional gap implication
still requires unproved local
coercivity and compatible
constants. No cofinal
auxiliary gap, continuum
construction, or physical
Yang--Mills mass gap is
claimed. The full goal
remains active.
\end{firewall}
\section{V68: tapered core and collar weights give a quadratic repair scale}
\label{sec:f12v68}

\subsection{Outcome and relation to the unresolved local problem}

V67 retained the spatial coefficients of the exterior-repair
form but used an indicator of a fixed-thickness band. Its
translation defect has surface order. We now replace that
indicator by a weight that equals one on every required
repair coordinate and tapers over a distance proportional
to the patch side. The retained core is tapered too.
The two supports remain separated by more than the
interaction range.

This proves an actual exceptional-exterior estimate whose
normalized repair debit is \(O(C_{\rm rep}(\beta)\ell^{-2})\),
in place of V67's \(O(C_{\rm rep}(\beta)\ell^{-3/2})\).
A fixed part of the original patch weight is discarded,
and is charged explicitly. A positive order-\(\ell^4\)
mass is retained. The second-order coefficient is
\(O(C_{\rm rep}(\beta)\ell^4)\).

This calculation does not prove the full-compact
good-exterior patch gap. In particular it does not
restrict an ambient repair through a frozen boundary,
nor substitute the boxed gap for the compact gap.
It sharpens one of the independently necessary
estimates before attempting that local problem.
No general novelty claim is made for weighted
finite-size methods: the context remains the
finite-size literature already cited in V63,
including Lemm--Mozgunov,
\href{https://arxiv.org/abs/1801.08915}
{Spectral gaps of frustration-free spin systems with boundary}.
The calculation below is for the present exceptional-exterior
form, not an application of their theorem to an
unverified Yang--Mills Hamiltonian.

\subsection{A macroscopic taper without enlarging the repairs}

Retain all the actual-law, coarse-data, bounded-range,
tail and conditional-repair hypotheses of V66--V67.
Thus \(h_i=I-P_i\), \(H=\sum_i h_i\),
\(E(f)=\langle f,Hf\rangle\), the interaction graph
has degree at most \(\Delta\) and owner range \(R=R_*\),
and the original patch operator is
\[
 A_x=\sum_i a_\ell(\operatorname{owner}(i)-x)h_i,
 \qquad a_\ell=\ell^{-1/2}+d_\ell/\ell
                 \quad\hbox{on }C_\ell=\{1,\ldots,\ell\}^4 .
 \tag{72.1}
\]
It is zero off that cube. Write \(m_1,m_2,D_a\)
for its original moments and defect. The exterior
guard \(J_x\) is outside the cube in a collar of
width \(s\ge R\). Put
\[
 r=H_{\rm rep},\qquad b=s+r,\qquad
 \ell\in16\mathbb N,\qquad
 \ell\ge16(b+R+1).
 \tag{72.2}
\]
Take the ambient torus side
\(M\ge4\ell+4(b+R+1)\), so the supports and
all range-\(R\) translates below are nonwrapping.
The sets \(T_j\) are the same bounded repairs
as before, with radius \(r\), incidence \(\iota\),
conditional Poincare constant \(C_{\rm rep}(\beta)\)
and bad-event probability bound \(p\).
Their sizes have not been enlarged with \(\ell\).

Extend the depth to the whole owner lattice by
\[
 d_\ell(y)=\min_{\alpha=1,\ldots,4}
                    \{y_\alpha,\ell+1-y_\alpha\}.
 \tag{72.3}
\]
This signed depth is one-Lipschitz for the
\(\ell^\infty\) metric; it is nonpositive outside
the original cube. Define
\[
 \begin{aligned}
 u(y)&=\ell^{-1}\bigl(d_\ell(y)-\ell/4\bigr)_+,\\
 z(y)&=\left(1-\frac{16}{\ell}
                \operatorname{dist}(d_\ell(y),[-b,r])\right)_+,\\
 w(y)&=u(y)+z(y),\qquad \mu_\ell=\sum_y u(y),\\
 U_x&=\sum_i u(\operatorname{owner}(i)-x)h_i,\qquad
 V_x=\sum_k z(\operatorname{owner}(k)-x)h_k .
 \end{aligned}
 \tag{72.4}
\]
The distance in the formula for \(z\) is distance
in the real line. The letter \(\mu_\ell\) here
denotes a deterministic weight mass, not a
probability measure.

\begin{proposition}[Support, retention and separation]
\label{prop:f12v68-support}
One has \(0\le u\le a_\ell\) and
\(\mu_\ell\ge\ell^4/2048\). Whenever \(j\in J_x\)
and \(k\in T_j\), the band weight at the
owner of \(k\), relative to \(x\), is one.
Every cross pair in the supports of \(U_x,V_x\)
has owner distance greater than \(R\).
Thus \(U_xV_x=V_xU_x\ge0\).
\end{proposition}
\begin{proof}
The support of \(u\) is the inner cube
\(\{\ell/4+1,\ldots,3\ell/4\}^4\).
There \(u\le d_\ell/\ell\le a_\ell\).
On the smaller cube
\(\{3\ell/8+1,\ldots,5\ell/8\}^4\),
of cardinality \((\ell/4)^4\), one has
\(u\ge1/8\). This proves the mass bound.

For a repaired coordinate anchored outside
\(C_\ell\), V67's support argument gives
\(-b\le d_\ell\le r\): an outside coordinate
is in the enlarged cube of width \(b\),
and an inside one has depth at most \(r\).
Hence \(z=1\) on every such repair coordinate.

If \(u(y)>0\) and \(z(y')>0\), then
\[
 d_\ell(y)>\ell/4,\qquad
 d_\ell(y')<r+\ell/16 .
 \tag{72.5}
\]
Their depth difference exceeds
\(3\ell/16-r>R\), by (72.2).
The Lipschitz property gives the same lower
bound on their owner distance. The actual
cross-coordinate conditional expectations
therefore commute by V61. Products of
the cross residual projections are positive;
summing with nonnegative coefficients proves
the final assertion.
\end{proof}

The taper is only a deterministic coefficient
in an upper bound on the ambient energy.
It is not a new heat-bath block, and it does
not change \(T_j\), the law, or the dynamics.

\subsection{Volume-order moments and quadratic translation defects}

For a finitely supported nonnegative weight \(v\), put
\[
 m_2(v)=\sum_y v(y)^2,\qquad
 d_v(\delta)=\frac12\sum_y|v(y)-v(y-\delta)|^2,\qquad
 D(v)=\max_i\sum_{k:\{i,k\}\in E_{\rm graph}}
          d_v(\operatorname{owner}(i)-\operatorname{owner}(k)).
 \tag{72.6}
\]

\begin{proposition}[Explicit taper bounds]
\label{prop:f12v68-moments}
Under (72.2), let \(K_R=2049R^2\).
For \(\|\delta\|_\infty\le R\),
\[
 \begin{aligned}
 m_2(w)&=m_2(u)+m_2(z)\le17\ell^4,\\
 d_w(\delta)&=d_u(\delta)+d_z(\delta),\\
 d_u(\delta)&\le\tfrac12R^2\ell^2,\qquad
 d_z(\delta)\le2048R^2\ell^2,\qquad
 D(w)\le\Delta K_R\ell^2 .
 \end{aligned}
 \tag{72.7}
\]
\end{proposition}
\begin{proof}
The core weight is at most \(1/4\), hence
\(m_2(u)\le\ell^4/16\).
The support of \(z\), together with any of
its range-\(R\) translates, lies in a cube
of side at most \(2\ell\). Indeed positivity
of \(z\) implies
\(d_\ell>-b-\ell/16\); its support is
contained in the corresponding enlarged cube.
Allowing two integer endpoints and a shift
of size \(R\) still gives a side bounded by
\(\ell+2b+\ell/8+2+2R\le2\ell\).
Since \(0\le z\le1\), \(m_2(z)\le16\ell^4\).
The supports are disjoint, proving the
first line of (72.7).

The function \(u\) is globally
\(\ell^{-1}\)-Lipschitz. The union of its
support with a range-\(R\) translate has
at most \(\ell^4\) sites, since the core
side is \(\ell/2\) and \(\ell/2+2R\le\ell\).
Its squared-difference sum is therefore
at most \(R^2\ell^2\).
The function \(z\) is globally
\(16/\ell\)-Lipschitz. Its corresponding
union has at most \(16\ell^4\) sites,
so its squared-difference sum is at most
\(4096R^2\ell^2\). Division by two proves
the asserted bounds on \(d_u,d_z\).

The support separation also implies that
both shifted cross products
\(u(y)z(y-\delta)\) and
\(z(y)u(y-\delta)\) vanish. Expanding
the square or the autocorrelation proves
the exact split for \(d_w\). Summing over
at most \(\Delta\) interacting neighbors
proves the stated \(D(w)\) bound.
\end{proof}

These constants are deliberately conservative.
The gain is the exponent: both profiles now
vary on scale \(\ell\), so their defects
have order \(\ell^2\), without a discontinuity
at the inner edge of the retained core.

\subsection{An actual exterior-form bound with explicit discarded mass}

As in V66, let
\[
 \mathcal B(f)=\sum_x
           \langle f,1_{G_x^c}A_x f\rangle,\qquad
 V_s(\ell)=(\ell+2s)^4-\ell^4 .
 \tag{72.8}
\]
The events \(G_x\) concern the original
exterior guards \(J_x\), not the enlarged
support of \(z\). In particular the
rarity count is still a fixed-width
surface count:
\[
 |J_x|\le q_*V_s(\ell),\qquad V_s(\ell)\le64s\ell^3 .
 \tag{72.9}
\]
Define
\[
 \begin{aligned}
 a_{\rm tap}&=C_{\rm rep}\iota\,m_2(w),\\
 b_{\rm tap}&=m_1-\mu_\ell+
          2p q_*V_s(\ell)\mu_\ell+C_{\rm rep}\iota D(w),\\
 \varepsilon_\ell&=2p q_*V_s(\ell)
                     +\frac{C_{\rm rep}\iota D(w)}{\mu_\ell}.
 \end{aligned}
 \tag{72.10}
\]

\begin{theorem}[Tapered actual exceptional-exterior transfer]
\label{thm:f12v68-transfer}
For the full compact ambient law under the inherited
V66 hypotheses, without a local or global gap assumption,
\[
 \mathcal B(f)\le
           a_{\rm tap}\|Hf\|_2^2+b_{\rm tap}E(f).
 \tag{72.11}
\]
Moreover
\[
 \begin{aligned}
 a_{\rm tap}&\le17C_{\rm rep}\iota\ell^4,\\
 m_1-b_{\rm tap}&=\mu_\ell(1-\varepsilon_\ell),\\
 \varepsilon_\ell&\le128q_*s\,p\ell^3
                 +2048\iota\Delta K_R C_{\rm rep}\ell^{-2}.
 \end{aligned}
 \tag{72.12}
\]
\end{theorem}
\begin{proof}
Because \(G_x\) is exterior-measurable, its
indicator commutes with every residual in
the patch. Decompose \(A_x=U_x+(A_x-U_x)\).
The latter operator is a nonnegative
weighted sum of residual projections.
Dropping the exterior indicator in this
part, then summing its translated
coefficients, gives exactly the upper
bound \((m_1-\mu_\ell)E(f)\).

For the retained core, use
\(1_{G_x^c}\le\sum_{j\in J_x}1_{B_j}\).
The core is separated from all these
repairs. Thus the actual V66 conditional
repair estimate applies:
\[
 \|1_{B_j}h_i f\|_2^2
 \le2C_{\rm rep}\sum_{k\in T_j}\|h_kh_i f\|_2^2
                                      +2p\|h_i f\|_2^2 .
 \tag{72.13}
\]
After summation the last term is bounded
by \(2p q_*V_s(\ell)\mu_\ell E(f)\).
It uses the retained mass \(\mu_\ell\),
not the original \(m_1\).

Every relevant \(k\) has band weight
\(z=1\). At most \(\iota\) anchors use
a given \(k\). Positivity and commutation
of cross residuals therefore bound
the mixed part by
\[
 2C_{\rm rep}\iota\sum_x\langle f,U_xV_x f\rangle
       \le C_{\rm rep}\iota\sum_x\|(U_x+V_x)f\|_2^2 .
 \tag{72.14}
\]
The second inequality follows from
\((U_x+V_x)^2-2U_xV_x=U_x^2+V_x^2\ge0\).
No commutation inside either support
is assumed.

V67's translated-square upper bound,
which uses only the universal
anticommutator inequality on edges,
now gives
\[
 \sum_x\|(U_x+V_x)f\|_2^2
          \le m_2(w)\|Hf\|_2^2+D(w)E(f).
 \tag{72.15}
\]
Combining the three costs proves
(72.11). The definitions give the
exact expression for \(m_1-b_{\rm tap}\).
Finally use (72.7), (72.9), and
\(\mu_\ell\ge\ell^4/2048\).
\end{proof}

The discarded mass is not being called
small. It can be most of the original
\(m_1\). What matters in the conversion
is that \(\mu_\ell\), not
\(m_1-\mu_\ell\), has a positive
volume-order lower bound, and that
the two additional debits are small
relative to \(\mu_\ell\).

\subsection{Conditional gap conversion and the improved constant ordering}

\begin{corollary}[A sufficient local target, still to be proved]
\label{cor:f12v68-gap}
In addition to the preceding hypotheses, suppose the
actual weighted conditional patch gap on every
good exterior \(G_x\) is at least \(\gamma_\ell>0\).
If
\[
 \gamma_\ell\mu_\ell(1-\varepsilon_\ell)>D_a ,
 \tag{72.16}
\]
then the original rate-one interface heat-bath gap satisfies
\[
 h_N\ge
 \frac{\gamma_\ell\mu_\ell(1-\varepsilon_\ell)-D_a}
                   {m_2+\gamma_\ell a_{\rm tap}}>0 .
 \tag{72.17}
\]
In particular, if \(\varepsilon_\ell\le1/2\), a
sufficient local target and resulting explicit bound are
\[
 \begin{aligned}
 \gamma_\ell&>4096\Delta B_R\ell^{-2},
                 \qquad B_R=R^2/2+16R,\\
 h_N&\ge
 \frac{\gamma_\ell/4096-\Delta B_R\ell^{-2}}
                   {1+17\gamma_\ell C_{\rm rep}\iota}.
 \end{aligned}
 \tag{72.18}
\]
\end{corollary}
\begin{proof}
Insert (72.11) in V66's second-order
conversion. This gives (72.17) exactly.
For (72.18), use \(\mu_\ell\ge\ell^4/2048\),
\(D_a\le\Delta B_R\ell^2\), \(m_2\le\ell^4\),
and \(a_{\rm tap}\le17C_{\rm rep}\iota\ell^4\).
The numerator used in the lower bound
is positive by the first line, so
replacing the denominator by an upper
bound is legitimate.
\end{proof}

As before these are auxiliary rate-one
heat baths, not a physical time normalization.
At a fixed \(\beta\), any verified choice
of a finite \(\ell\) in (72.16) would give
a bound independent of sufficiently large
ambient torus size. Its proof is not
supplied merely by the conditional statement.

With the unchanged actual certificates
\[
 C_{\rm rep}(\beta)=e^{c_{\rm rep}(\beta+1)},\qquad
 p=e^{-c_\tau\beta},
 \tag{72.19}
\]
and fixed geometric constants, a choice
\(\ell\asymp e^{a\beta}\), rounded to a
multiple of sixteen, makes the right
side of (72.12) tend to zero provided
\[
                    \frac{c_{\rm rep}}2<a<\frac{c_\tau}3 .
 \tag{72.20}
\]
The rates are \(c_{\rm rep}-2a\) and
\(3a-c_\tau\). This interval is nonempty
if
\[
                         c_{\rm rep}<\frac23c_\tau .
 \tag{72.21}
\]
This is weaker than V67's
\(c_{\rm rep}<c_\tau/2\); neither
ordering has been verified for the
actual tail and surface-comparison
constants. They have not been treated
as freely adjustable parameters.

The new criterion still requires a
patch gap at diffusive weighted scale
\(\ell^{-2}\). If only an unweighted
compact patch gap \(g_\ell\) is available,
the minimum-rate comparison supplies
\(\gamma_\ell\ge(\ell^{-1/2}+\ell^{-1})g_\ell\);
a sufficient unweighted target remains
\[
                   g_\ell>4096\Delta B_R\ell^{-3/2}.
 \tag{72.22}
\]
Thus the new repair estimate does not
by itself improve the exponent of the
unweighted finite-size gap target.
V61's known surface-exponential
certificate does not verify that target.

\subsection{A limited geometric check against indefinite exponent chasing}

For these specific profiles the quadratic
defect order is also bounded below if
the interaction graph has an axial edge.
This observation is not an optimality
theorem over all repair methods.

\begin{proposition}[Lower scale for this nonnegative envelope]
\label{prop:f12v68-lower-scale}
Suppose some graph edge has owner displacement
\(e_1=(1,0,0,0)\). For \(w=u+z\) in (72.4),
\[
                 \frac{D(w)}{\mu_\ell}
                           \ge\frac{1}{12288\ell^2}.
 \tag{72.23}
\]
\end{proposition}
\begin{proof}
For any real sequence \(v\) supported in an
interval of \(L\) sites, extend it by zero.
Writing \(v_j=\sum_{k=1}^j(v_k-v_{k-1})\)
and using Cauchy--Schwarz gives
\[
 \sum_{j=1}^L v_j^2
 \le\frac{L(L+1)}2
       \sum_{k\in\mathbb Z}(v_k-v_{k-1})^2 .
 \tag{72.24}
\]
Apply this on each line parallel to \(e_1\)
for \(w\), whose support is contained
in a cube of side at most \(L=2\ell\).
Then \(d_w(e_1)\ge m_2(w)/(L(L+1))\).
Also \(m_2(w)\ge m_2(u)\ge\mu_\ell^2/\ell^4\)
by Cauchy--Schwarz and the size of the
core support. Since \(L(L+1)\le6\ell^2\)
and \(\mu_\ell\ge\ell^4/2048\), the result
follows from \(D(w)\ge d_w(e_1)\).
\end{proof}

The condition on an axial edge is explicit;
it has not been asserted for every possible
owner graph. The proposition shows that
within this particular envelope one
cannot obtain an \(o(\ell^{-2})\)
normalized defect merely by sharpening
the constants in (72.7).

\subsection{Exact verification and the next upstream obligation}

The script
\begin{center}\small\ttfamily
scripts/verify\_ym\_f12\_v68\_tapered\_transfer.py
\end{center}
checks the four-dimensional profiles at
\((\ell,r,R,s)=(64,1,1,1)\) and
\((128,2,2,2)\), both satisfying (72.2).
It expands the weights in signed nested
cube indicators and evaluates shifted
intersections exactly. This avoids
finite sampling of lattice sites.
The retained masses are respectively
\(60962\) and \(905762\).
For three displacements at each scale,
including fully diagonal and mixed-sign
shifts, it checks vanishing cross
correlations, the exact moment and
defect splits, and all upper bounds.
The cube expansions are independently
checked at each integer depth against
the defining piecewise-linear formulas.

An additional \(81\)-state product model
has four ternary coordinates with
one-site probabilities \((49/50,1/100,1/100)\).
It uses patch weights \((1/2,1,1/2,0)\),
retained weights \((0,1/4,0,0)\), and
one exterior repair of weight one.
The script checks the full defect
inequality, including the discarded
mass, by exact positive-semidefinite
elimination in the probability-weighted
inner product. Its values are
\[
 b_{\rm tap}=44/25,\qquad
 m_1-b_{\rm tap}=6/25,\qquad
 h_{\rm converted}=96/1625,
 \tag{72.25}
\]
while the true product gap is one.
This validates a nonzero discarded
mass and a positive conversion in a
controlled model; it does not populate
the actual Yang--Mills local gap.

\begin{decision}[Do not infer a frozen-patch theorem from the ambient bound]
The taper removes the hard-band surface
loss in V67. The remaining route must
still prove compact patch coercivity
on a good exterior, or obtain a different
noncircular local input.
In (72.13) the actual repair updates
all of \(T_j\) in the ambient law.
When a patch exterior is frozen, this
does not authorize replacing \(T_j\)
by \(T_j\cap C_\ell\) with the same
tail and Poincare estimates.
V65's support-completion analysis remains
in force. Iterating a completion is
not a proved fixed-collar localization.

Accordingly, the next direction should
address the finite-patch conditional law
itself, including boundary-touching
defects, rather than reinterpret the
ambient second-order estimate as that
law's gap. Independently, numerical
compatibility in (72.21) must be derived
from the actual constants if this
exponential-scale route is used.
The next companion audit and direction
review remain V70.
\end{decision}

V67 had 1080900 bytes and SHA--256
\begin{center}\scriptsize\ttfamily
D909642515414E63F2A6D1B9609A5BB8CE4601030BDFF00BB2C588F047029930.
\end{center}
Its complete body is retained apart
from the title version. Only active
V67 is replaced by V68. All archives,
companions and prior scripts remain.

\begin{firewall}[Position after V68]
The tapered-support geometry, the
quadratic defect bounds, and the
actual exceptional-exterior estimate
are proved under the inherited explicit
repair hypotheses. The local-gap
conversion remains conditional on
unproved compact patch coercivity
and a positive numerator. There is
still no cofinal auxiliary gap,
continuum construction, or physical
Yang--Mills mass-gap theorem.
The full goal remains active.
\end{firewall}
\section{V69: the present repair certificate cannot meet the tail scale}
\label{sec:f12v69}

\subsection{A correction to the direction, not to the preceding inequalities}

V66--V68 improved the actual exceptional-exterior
estimate while leaving the compatibility of its
constants unverified. The present audit proves a
stronger conclusion: the particular oscillation-based
repair certificate used there cannot satisfy any
of those proposed large-\(\beta\) exponent orderings.
The same local oscillation bound that supplies this
certificate also limits how fast a fixed off-chart
event of positive Haar measure can become rare.

The finite-energy lower bound used here was already
proved in V33 and used in V40. It is not a new
probabilistic mechanism. The new step is to compare
its coefficient with the explicitly chosen effective
surface-comparison coefficient in V61 and V66.
This yields a certificate obstruction and changes
the next action. The V66--V68 exterior inequalities
remain valid; their numerical premise must no
longer be described as merely unchecked for the
certificates currently inserted.

The conclusion does not lower-bound the optimal
repair Poincare constant, disprove an auxiliary
gap, or disprove the Yang--Mills mass gap. It
identifies an overly costly comparison in this
route. A field-preserving product reference is
specified below as a possible replacement, with
its required estimate left explicit.

\subsection{Recover the finite-energy lower bound with its actual coefficient}

Use the full compact interface potential
\(W_\beta=U_{\mathcal S}-\sum_b\log Z_b\)
relative to product normalized Haar measure.
Retain V61's effective constants:
\[
 \#\{\hbox{terms incident on }j\}\le D_*,
 \quad \operatorname{osc}\Phi_a\le\widehat v_\beta,
 \quad \widehat v_\beta\le c_*(\beta+1).
 \tag{73.1}
\]
Let \(B_j\) be the complement of the fixed proper
small chart used in V58--V68. At fixed retained
data, its one-group Haar measure is
\[
                         \alpha_j=\operatorname{Haar}(B_j)>0 .
 \tag{73.2}
\]
The chart is fixed as \(\beta\) varies; translation
by the section does not change its Haar measure.
Only one such coordinate and an admissible
family at arbitrarily large \(\beta\) are needed
for the argument. If no such admissible family
exists, the earlier large-\(\beta\) route is
unavailable for that independent reason.

\begin{proposition}[Finite-energy constraint on the tail exponent]
\label{prop:f12v69-tail-coefficient}
For every compact exterior of the single group \(j\),
\[
 P_j1_{B_j}\ge
       \alpha_j e^{-D_*\widehat v_\beta}
       \ge\alpha_j e^{-D_*c_*(\beta+1)} .
 \tag{73.3}
\]
The same lower bound holds for \(Q_j1_{B_j}\)
whenever the repair \(T_j\) contains \(j\).
Consequently, if the inherited repair bound
\(Q_j1_{B_j}\le e^{-c_\tau\beta}\) holds for
arbitrarily large \(\beta\), with \(c_\tau>0\)
fixed, then
\[
                              c_\tau\le D_*c_* .
 \tag{73.4}
\]
\end{proposition}
\begin{proof}
With all other groups fixed, the one-group
conditional potential is the sum of terms
incident on \(j\). Its oscillation is at most
\(D_*\widehat v_\beta\); terms not incident
on \(j\) cancel from its conditional density.
If this potential has minimum \(m\) and
maximum \(M\), its normalized Haar density
is at least \(e^{-M}/e^{-m}=e^{-(M-m)}\).
Integration over \(B_j\) proves (73.3).

Since \(j\in T_j\), conditional expectation
onto the exterior of \(T_j\) gives
\(Q_jP_j=Q_j\). Apply this positive operator
to (73.3). Combining the resulting lower
bound with the assumed tail upper bound yields
\[
 c_\tau\beta
       \le D_*c_*(\beta+1)-\log\alpha_j .
 \tag{73.5}
\]
Divide by \(\beta\) along the admissible
unbounded sequence. The chart Haar mass
is fixed and positive, so its contribution
vanishes in the limit.
\end{proof}

This use of the tower property requires no
independence and no change to the compact
law. Small probability under that law
does not mean small Haar measure.
The fixed positive \(\alpha_j\) is exactly
what prevents the two exponential
coefficients from being chosen independently.

\subsection{Compare with the surface coefficient actually selected}

Set \(L_{\rm rep}=2H_{\rm rep}+1\).
The certificate used in V66--V68 is
\[
 C_{\rm rep}(\beta)=e^{c_{\rm rep}(\beta+1)},\qquad
 c_{\rm rep}=C_{\rm surf}^*c_*L_{\rm rep}^3,
 \tag{73.6}
\]
where V61 gives, without suppressing constants,
\[
 C_{\rm surf}^*
    =q_*D_*+\frac{64}{7}\,8q_*D_*(R_*+1).
 \tag{73.7}
\]
For a nonempty interacting specification use
\(q_*\ge1,D_*\ge1,L_{\rm rep}\ge1\).
Define
\[
 \kappa_{\rm surf}
    =q_*\left(1+\frac{512}{7}(R_*+1)\right)L_{\rm rep}^3
    \ge1 .
 \tag{73.8}
\]

\begin{theorem}[Incompatibility of the presently populated exponents]
\label{thm:f12v69-incompatible}
Under the preceding actual-law and fixed-chart
hypotheses,
\[
                c_{\rm rep}
                   =\kappa_{\rm surf}D_*c_*
                   \ge\kappa_{\rm surf}c_\tau
                   \ge c_\tau .
 \tag{73.9}
\]
In particular, the V68 requirement
\(c_{\rm rep}<2c_\tau/3\), and the stronger
V67 and V66 requirements, cannot hold
with this selected surface certificate.
\end{theorem}
\begin{proof}
Substitute (73.7) into (73.6), obtaining
the equality in (73.9). Apply (73.4)
and \(\kappa_{\rm surf}\ge1\).
Since \(c_\tau>0\), the resulting lower
bound contradicts \(c_{\rm rep}<2c_\tau/3\).
The requirements \(c_{\rm rep}<c_\tau/2\)
and \(c_{\rm rep}<c_\tau/3\) are stronger.
\end{proof}

The equality in (73.9) refers to the
chosen certificate, not to an optimal
constant. Increasing \(c_*\), enlarging
the repair radius, or weakening the tail
exponent does not cure this particular
conflict. The argument also explains
why assigning numerical values to the
two constants independently would have
been invalid.

\subsection{Optimizing over every patch size does not cure the certificate}

The upper expression in V68's (72.12) has
the form
\[
 F_\beta(\ell)
       =A e^{-c_\tau\beta}\ell^3
                      +B e^{c_{\rm rep}(\beta+1)}\ell^{-2},
 \quad
 A=128q_*s,\quad B=2048\iota\Delta K_R .
 \tag{73.10}
\]
Assume here \(A,B>0\), as in the
nontrivial spatial case. The vanishing
coefficient cases are not included in
the following minimization argument.

\begin{proposition}[Exact scale optimization]
\label{prop:f12v69-optimize}
For any \(a,b>0\),
\[
 \inf_{\ell>0}(a\ell^3+b\ell^{-2})
   =\frac5{2^{2/5}3^{3/5}}\,a^{2/5}b^{3/5},
 \qquad
 \ell_{\rm opt}=\left(\frac{2b}{3a}\right)^{1/5}.
 \tag{73.11}
\]
Consequently
\[
 \inf_{\ell>0}F_\beta(\ell)
 =\frac5{2^{2/5}3^{3/5}}\,
       A^{2/5}B^{3/5}e^{3c_{\rm rep}/5}
       e^{(3c_{\rm rep}-2c_\tau)\beta/5}
             \longrightarrow+\infty .
 \tag{73.12}
\]
\end{proposition}
\begin{proof}
The function diverges at zero and infinity.
Its derivative vanishes exactly when
\(3a\ell^5=2b\), is negative before that
point and positive after it. Substitution
gives (73.11), including its constant.
Apply this with the two coefficients
in (73.10). By (73.9),
\(3c_{\rm rep}-2c_\tau\ge c_\tau>0\).
\end{proof}

Restricting \(\ell\) to admissible integer
patch sizes cannot reduce this infimum.
The obstruction is therefore not confined
to exponential ansatzes
\(\ell\asymp e^{a\beta}\).
However, \(F_\beta\) is an upper bound
on a proof coefficient. Its divergence
alone would not prove that the exact
coefficient diverges, much less that
an actual spectral gap vanishes.
The next observation addresses only
the exact coefficient chosen in V68.

\subsection{The chosen exact numerator also fails on a cross-owner graph}

\begin{proposition}[Translation-orbit version of the taper lower bound]
\label{prop:f12v69-orbits}
For the V68 envelope \(w=u+z\), if the
interaction graph has any edge with a
nonzero owner displacement, then
\[
                         \frac{D(w)}{\mu_\ell}
                                   \ge\frac{1}{12288\ell^2}.
 \tag{73.13}
\]
The displacement need not be axial
or primitive.
\end{proposition}
\begin{proof}
Lift the finite support to \(\mathbb Z^4\);
the nonwrapping hypotheses in V68 identify
its autocorrelations for every graph
displacement with those of this lift.
Fix a nonzero integer displacement
\(\delta\) furnished by an edge.
Partition \(\mathbb Z^4\) into its
orbits \(y+\mathbb Z\delta\).
At least one coordinate of \(\delta\)
has absolute value at least one.
Since \(w\) is supported in a cube
of side at most \(2\ell\), on each
orbit its support fits into an interval
of at most \(2\ell\) consecutive orbit
indices. Missing sites in that interval
can be left with value zero.

Apply the one-dimensional telescoping
inequality (72.24) on every orbit and
sum. This gives
\[
 d_w(\delta)\ge
       \frac{m_2(w)}{(2\ell)(2\ell+1)} .
 \tag{73.14}
\]
The remaining steps are V68's
\(m_2(w)\ge\mu_\ell^2/\ell^4\),
\(\mu_\ell\ge\ell^4/2048\), and
\((2\ell)(2\ell+1)\le6\ell^2\).
Since this displacement occurs on
an edge, \(D(w)\ge d_w(\delta)\).
\end{proof}

In this case \(s\ge R_*>0\). The exact
coefficient defined in (72.10) satisfies
\[
 \begin{aligned}
 \varepsilon_\ell
   &=2p q_*V_s(\ell)
                     +C_{\rm rep}\iota\,D(w)/\mu_\ell\\
   &\ge16q_*s\,p\ell^3
                     +\frac{\iota}{12288}C_{\rm rep}\ell^{-2},
 \end{aligned}
 \tag{73.15}
\]
because
\((\ell+2s)^4-\ell^4\ge8s\ell^3\).
Proposition~\ref{prop:f12v69-optimize},
now applied to this lower bound, proves
\[
               \inf_{\ell\ {\rm admissible}}
                              \varepsilon_\ell\longrightarrow+\infty .
 \tag{73.16}
\]
Thus for all sufficiently large \(\beta\)
and every admissible size, the particular
V68 numerator
\(\gamma_\ell\mu_\ell(1-\varepsilon_\ell)-D_a\)
is nonpositive for every \(\gamma_\ell>0\).
Even a separately proved excellent
local patch gap would not rescue that
populated certificate in this regime.

If there is no cross-owner edge this
conclusion is not asserted: the defect
can vanish, as already noted in V61.
Nor does (73.16) assert a lower bound
on the smallest possible
exceptional-exterior coefficient.
It concerns exactly the coefficients
constructed with (73.6).

\subsection{Why concentration need not make the true repair expensive}

A short interacting example separates
the failure of the certificate from
failure of a heat-bath gap. On two
bits \(x,y\in\{0,1\}\), set
\[
 \nu_\beta(x,y)
       =Z_\beta^{-1}e^{-\beta(x+y)+Jxy},
                      \qquad J>0\ \hbox{fixed}.
 \tag{73.17}
\]
Compare with the product law
\(\pi_\beta(x,y)\propto e^{-\beta(x+y)}\).
The likelihood ratio \(\nu_\beta/\pi_\beta\)
has maximum divided by minimum equal
to \(e^J\), independently of \(\beta\).
If \(m\le d\nu_\beta/d\pi_\beta\le M\),
the variational formulas
\[
 \operatorname{Var}_\rho f=\inf_c\rho[(f-c)^2],\qquad
 e_i^\rho(f)=\inf_{g\ {\rm measurable\ off}\ i}
                                       \rho[(f-g)^2]
 \tag{73.18}
\]
give
\[
 \operatorname{Var}_{\nu_\beta}f
 \le M\operatorname{Var}_{\pi_\beta}f
 \le M\sum_i e_i^{\pi_\beta}(f)
 \le(M/m)\sum_i e_i^{\nu_\beta}(f).
 \tag{73.19}
\]
Here the middle inequality is product
variance tensorization; the last one
uses \(e_i^{\nu_\beta}\ge m e_i^{\pi_\beta}\).
The actual rate-one two-coordinate
heat-bath gap is therefore at least
\(e^{-J}\) uniformly in \(\beta\).
On the other hand,
\[
 \nu_\beta(x=1\mid y)
       =\frac{e^{-\beta+Jy}}{1+e^{-\beta+Jy}}
       \le e^{-\beta+J}\le e^{-\beta/2}
                         \quad(\beta\ge2J).
 \tag{73.20}
\]
There is genuine interaction and
exponentially rare bad coordinates,
yet a uniformly bounded repair
Poincare constant. Charging the
one-coordinate field of size \(\beta\)
as a perturbation of uniform measure
would discard this fact.
This finite model is not a
Yang--Mills specification and does
not provide its missing estimate.

\subsection{An upstream replacement target that retains the one-coordinate fields}

The following is a direct specialization
of the product comparison already
underlying V38, not a new tensorization
principle. It specifies what a less
costly fixed-repair argument could
attempt in the actual compact law.

For the actual repair conditional
potential \(W_{T_j,\beta}^{\,\xi}\)
relative to product Haar measure,
define
\[
 \Gamma_\beta=
 \sup_{j,\xi}\ \inf_{\{V_i:i\in T_j\}}
       \operatorname{osc}\left(
          W_{T_j,\beta}^{\,\xi}(v)
                          -\sum_{i\in T_j}V_i(v_i)\right).
 \tag{73.21}
\]
The supremum includes all required
retained data and compact exteriors.
In the infimum the \(V_i\) may be
arbitrary continuous one-group
potentials depending on these
parameters and on \(\beta\).
On compact groups their normalized
one-group reference laws exist.

\begin{proposition}[Actual repair comparison after retaining local fields]
\label{prop:f12v69-field-reference}
If \(\Gamma_\beta<\infty\), the actual
repair inequality is valid with
\[
                              C_{\rm rep}^{\rm new}(\beta)=e^{\Gamma_\beta}.
 \tag{73.22}
\]
\end{proposition}
\begin{proof}
Fix the repair and its exterior.
For any chosen \(\{V_i\}\), take the
product reference with one-group
densities proportional to \(e^{-V_i}\).
Its likelihood ratio with the actual
conditional law has maximum/minimum
ratio equal to the exponential of
the residual oscillation in (73.21).
The exact argument (73.18)--(73.19)
therefore gives the conditional
Poincare inequality for the actual
single-group updates.

For each exterior take an arbitrarily
close approximation to the infimum,
then let its error decrease to zero.
No measurable choice of optimizers
is required, because the resulting
inequality is pointwise and its
constant is bounded by the same
supremum \(\Gamma_\beta\).
Disintegration yields the ambient
repair inequality in the form used
in (70.12).
\end{proof}

For example an actual estimate
\[
                  \Gamma_\beta\le K+\kappa\beta,\qquad
                         \kappa<2c_\tau/3
 \tag{73.23}
\]
would avoid the exponent conflict
proved above: V68 could use
\(C_{\rm rep}^{\rm new}\le e^K e^{\kappa\beta}\)
and \(\kappa/2<a<c_\tau/3\).
Unlike \(c_{\rm rep}\) in (73.6),
\(\kappa\) is a residual interaction
coefficient after the one-coordinate
fields have been retained. The
finite-energy argument does not
imply \(\kappa\ge c_\tau\);
(73.17) has residual oscillation
at most \(J\), so \(\kappa=0\)
is possible in that model.

Equation (73.23) has not been proved
for the actual compact interface
repair. A boxed Hessian estimate
does not control the global
oscillation in (73.21). Even if
(73.23) were proved, the independent
good-exterior compact patch gap
would still be needed in (72.16).
This target is a possible upstream
replacement, not an asserted closure.

\subsection{Exact checks and version accounting}

The script
\begin{center}\small\ttfamily
scripts/verify\_ym\_f12\_v69\_certificate\_obstruction.py
\end{center}
checks the scale inequality
\[
          108(a\ell^3+b\ell^{-2})^5\ge3125a^2b^3
 \tag{73.24}
\]
in 36 rational cases and four exact
equality cases. It also checks the
diverging minimum for
\(p=2^{-k}\), \(C_{\rm rep}=2^{k+1}\),
the minimal conflicting exponent
ordering, at five integer values
of \(k\). These checks are examples;
the all-scale proof is (73.11).

Using V68's exact nested-cube
expansions, it verifies (73.13)--(73.14)
for four nonaxial or nonprimitive
displacements of its \(\ell=64\)
envelope. Finally it takes
\(J=\log2\), \(\beta=k\log2\) in
(73.17), for ten even integers
\(k=2,\ldots,20\). Exact weighted
positive-semidefinite checks give
\(H^2\ge H/2\), the conditional
tail bounds hold, and the actual
product-comparison density ratio
is exactly two in every case.
All checks pass. None tests a
Yang--Mills continuum Hamiltonian.

\begin{decision}[Change the repair input before further scale optimization]
The previous iterations' missing
constant comparison is now a proved
obstruction for the populated
surface certificate. Further tuning
of the patch size or the same
taper constants cannot repair it
at large \(\beta\).

The next upstream step is to obtain
a better estimate for the fixed
actual compact repair, retaining
its confining one-coordinate terms
or using an equally justified
non-oscillation argument. The
compact good-exterior patch problem
remains separate. At V70, perform
the scheduled companion audit and
direction review in light of this
obstruction; do not restore the
disproved exponent ordering as an
open numerical possibility.
\end{decision}

V68 had 1099108 bytes and SHA--256
\begin{center}\scriptsize\ttfamily
75E93E7E91604B32B3BB831D575DB7F05071832688EB1F7E56510A219B8A4818.
\end{center}
Its complete body is retained apart
from the title version. Only active
V68 is replaced by V69. Archives,
companions and earlier scripts remain.

\begin{firewall}[Position after V69]
The current repair certificate's
exponent incompatibility and its
all-size large-\(\beta\) obstruction
are proved under the explicit
actual-law hypotheses. No lower
bound on the optimal repair cost
is inferred. A replacement
field-preserving comparison is
specified but not verified for
Yang--Mills. The cofinal auxiliary
gap, continuum construction, and
physical mass gap remain open
within this programme. The full
goal remains active.
\end{firewall}
\section{V70: audit and a classical variational test for the replacement repair}
\label{sec:f12v70}

\subsection{Companion audit and ten-version direction review}

The scheduled audit read the current companion endings and
checked their full SHA--256 digests. They are unchanged
from V65:
\begin{center}\small
\begin{tabular}{lll}
programme & active version & bytes\\ \hline
SM--Einstein & V72 & 607372\\
NS stability & V26 & 278283\\
parallel YM & V5 & 54174
\end{tabular}
\end{center}
The files are respectively
\texttt{Renewal\_SM\_Einstein\_Continuum\_Research\_Notes\_V72.tex},
\texttt{Renewal\_NS\_Stability\_Research\_Notes\_V26.tex},
and \texttt{Renewal\_YM\_Parallel\_Research\_Notes\_V5.tex}.
Their digests, in that order, are
\begin{center}\scriptsize\ttfamily
F44F9745D43379E1672C5550411B58E97195A4C9278592D221DF5D3F644DC1F5\\
82C887F8C2C69E4F28FAD7C3DC8DE5B9099CB5A4BF431EBA1FFF3E91935978AD\\
306F1BB84CFA8A8D47EA569EC17E9545F9867C8D32B548EC9D9C444B4F3C0512.
\end{center}

SM V72 proves a fixed-time many-copy limit with fixed
mode sets and, for point-concentrated data, vanishing
empirical mixed-parent residual. It does not supply
a growing-cutoff continuum theorem or a compact YM
repair estimate. NS V26 computes pointwise rank-one
kernel norms, not a stability/global-regularity
theorem or a YM spectral estimate. Parallel YM V5
rules out its extensive global-charge tail target
and retains local or density-level topology targets.
Its internal stopping language is documentary
context, not an instruction to stop this task.
None of these updates changes V69's coefficient
obstruction or supplies the missing patch gap.

The main direction over V61--V70 is now:
finite-size identities and V66--V68 repair transfers
are proved; the old surface repair certificate
cannot populate their large-\(\beta\) numerator,
by V69. The next estimate must change that repair
input. V69 proposed retaining arbitrary one-coordinate
confining potentials in the comparison reference.
This note identifies the exact leading coefficient
of that replacement comparison and a certifiable
finite variational test for it.

V29 already proved a uniform compact
zero-temperature approximation, and V28 already
proved the exact bounded-block factorization.
They are reused, not rediscovered below.
The new conclusion concerns the oscillation
remaining after all additive one-coordinate
potentials have been removed.

\subsection{Localize the inherited compact energy minimum to one repair}

Fix a repair template \(T=T_j\) and its compact
exterior parameters \(\xi\), including the retained
data. All comparisons below use a fixed parameter
family \(\Xi\), admissible for every sufficiently
large \(\beta\). Uniform statements may also be
taken over finitely many fixed-packing repair
templates. If instead the admissible parameter
set changes with \(\beta\), the pointwise uniform
estimate remains true, but its supremum must be
taken over that changing set; a fixed limiting
supremum is not then automatic.

Discard potential terms independent of \(v_T\).
V28's exact assignment of plaquettes and pivot
factors, with the fixed smooth compact completion,
then gives a representative conditional potential
\[
 W_{T,\beta}^{\,\xi}(v_T)
  =\beta S_{{\rm dir},T}^{\,\xi}(v_T)
       +\Psi_{{\rm dir},T}^{\,\xi}(v_T)
       -\sum_{b:J_b\cap T\ne\varnothing}
                      \log Z_{b,\beta}(v_{J_b}),
 \tag{74.1}
\]
where
\[
 Z_{b,\beta}(v_{J_b})
    =\int_{G^{I_b}}
            e^{-\beta S_b(x_b,v_{J_b})}a_b(x_b,v_{J_b})\,dx_b .
 \tag{74.2}
\]
Here the amplitudes contain the full inverse-pivot
factors, are independent of \(\beta\), and have
fixed positive lower and upper bounds on each
compact template. The direct \(\Psi\) terms are
bounded. The action functions are the actual
smooth compact functions from the completion,
not their Taylor polynomials on a chart.

Set
\[
 m_b(v_{J_b})=\min_{x_b\in G^{I_b}}S_b(x_b,v_{J_b}),\qquad
 \mathcal E_T^{\,\xi}(v_T)
   =S_{{\rm dir},T}^{\,\xi}(v_T)
                 +\sum_{b:J_b\cap T\ne\varnothing}m_b(v_{J_b}).
 \tag{74.3}
\]
Every minimum exists. It need not be unique,
nondegenerate, or attained at the section.
Only a bounded number of blocks touch the
fixed repair, independently of ambient volume.

\begin{proposition}[Uniform local version of the inherited energy limit]
\label{prop:f12v70-local-laplace}
There are fixed \(K_T,B_T<\infty\), uniform
over the stated compact parameters, such that
for \(\beta\ge1\)
\[
 W_{T,\beta}^{\,\xi}
       =\beta\mathcal E_T^{\,\xi}+R_{T,\beta}^{\,\xi},
 \qquad
 \operatorname{osc}_{v_T}R_{T,\beta}^{\,\xi}
                              \le K_T+B_T\log\beta .
 \tag{74.4}
\]
The energy \(\mathcal E_T^{\,\xi}\) is continuous
and Lipschitz in its free and exterior arguments
with fixed constants. Differentiability of
the minimizing branches is not asserted.
\end{proposition}
\begin{proof}
Apply V29's compact minimum argument separately
to each block in (74.2). If
\(n_b=\dim(G)|I_b|\), the amplitude bounds,
normalized Haar volume, and the action minimum
give
\[
 \beta m_b-C_b\le-\log Z_{b,\beta}
                  \le\beta m_b+\frac{n_b}{2}\log\beta+C_b .
 \tag{74.5}
\]
For completeness, the upper inequality uses a
product normal neighborhood of radius
proportional to \(\beta^{-1/2}\) at any minimum.
At a minimum on the compact group the gradient
vanishes. The uniform action Hessian upper
bound controls its increase by a quadratic,
and the Haar density and amplitude have
uniform positive lower bounds there.
This is exactly the proof already used in V29;
it requires no invertible Hessian.

Sum (74.5) only over blocks touching \(T\),
and add the bounded direct amplitude potential.
The number and dimensions of these blocks
are fixed. Their possible contribution to
the oscillation is at most a constant plus
\(\frac12\sum_b n_b\log\beta\), proving (74.4).
No full-volume entropy error is inserted.

The smooth actions on the fixed compact
templates are uniformly Lipschitz in
their parameters. For any such action,
\[
 |\min_x S(x,v)-\min_x S(x,v')|
                      \le\sup_x|S(x,v)-S(x,v')| .
 \tag{74.6}
\]
This proves the Lipschitz assertion after
summing the bounded number of terms.
\end{proof}

\subsection{The additive-oscillation quotient}

For a continuous function \(F\) on the compact
product \(\prod_{i\in T}G\), define
\[
 \mathfrak q_T(F)=
    \inf_{V_i\in C(G)}
         \operatorname{osc}\left(F-\sum_{i\in T}V_i(v_i)\right).
 \tag{74.7}
\]
Constants can be absorbed into one of the
\(V_i\). Thus normalizing a conditional
density does not change this quantity.
It is nonnegative and insensitive to adding
any one-coordinate potentials.

\begin{proposition}[Stability and homogeneity of the quotient]
\label{prop:f12v70-quotient}
For continuous \(F,G\) and \(t\ge0\),
\[
 |\mathfrak q_T(F)-\mathfrak q_T(G)|
                \le\operatorname{osc}(F-G),\qquad
 \mathfrak q_T(tF)=t\mathfrak q_T(F).
 \tag{74.8}
\]
\end{proposition}
\begin{proof}
For any additive \(A=\sum_i V_i\),
\(\operatorname{osc}(F-A)\le
\operatorname{osc}(G-A)+\operatorname{osc}(F-G)\).
Take an arbitrarily close approximation
to the infimum on the right, and then
exchange \(F,G\). For \(t>0\), scale
the candidate \(V_i\) by \(t\) in both
directions of the infimum. The assertion
at \(t=0\) follows from \(\mathfrak q_T(0)=0\).
\end{proof}

Use V69's actual comparison coefficient
and its classical counterpart:
\[
 \Gamma_\beta=\sup_{T,\xi}\mathfrak q_T(W_{T,\beta}^{\,\xi}),\qquad
 \Gamma_0=\sup_{T,\xi}\mathfrak q_T(\mathcal E_T^{\,\xi}).
 \tag{74.9}
\]

\begin{theorem}[Exact leading exponent of the field-preserving comparison]
\label{thm:f12v70-exponent}
For the fixed parameter family described above,
\[
       |\Gamma_\beta-\beta\Gamma_0|
                              \le K+B\log\beta ,
       \qquad
       \lim_{\beta\to\infty}\frac{\Gamma_\beta}{\beta}=\Gamma_0 .
 \tag{74.10}
\]
The constants are independent of ambient
volume once the fixed repair templates fit.
\end{theorem}
\begin{proof}
Apply (74.8) to (74.4), obtaining the
uniform pointwise bound
\[
 \left|\mathfrak q_T(W_{T,\beta}^{\,\xi})
              -\beta\mathfrak q_T(\mathcal E_T^{\,\xi})\right|
                         \le K_T+B_T\log\beta .
 \tag{74.11}
\]
Take the supremum over the same family
on both sides, using the largest fixed
constants. Divide by \(\beta\).
This does not exchange a nonuniform
limit with an infimum or supremum:
the error was uniform before either
operation.
\end{proof}

Accordingly, a strict actual estimate
\[
                            \Gamma_0<2c_\tau/3
 \tag{74.12}
\]
would make V69's field-preserving
replacement compatible with V68's
rarity scale. Choose
\(\Gamma_0<\kappa<2c_\tau/3\);
the logarithmic error in (74.10)
is absorbed into \(K_\kappa+
(\kappa-\Gamma_0)\beta\).
Conversely, if \(\Gamma_0>2c_\tau/3\),
this particular optimal-product
oscillation comparison fails that
strict exponent test. Equality is
a borderline case not resolved by
the leading exponent alone.
None of the three cases has yet
been established for the actual
worst-exterior repair family.

\subsection{Finite certificates: additive potentials and balanced marginals}

Let \(X=\prod_i X_i\) be a finite product grid
and \(F:X\to\mathbb R\). Define \(\mathfrak q_X\)
by (74.7) with arbitrary functions on the
finite one-coordinate sets. It is the value
of the linear program
\[
 \begin{gathered}
 \text{minimize }t\quad\text{over }t,c,\{a_i(s):s\in X_i\},\\
       c\le F(x)-\sum_i a_i(x_i)\le c+t
                          \quad\text{for every }x\in X .
 \end{gathered}
 \tag{74.13}
\]
Adding \(t\ge0\) is redundant but harmless.

\begin{proposition}[Balanced-marginal dual certificate]
\label{prop:f12v70-dual}
For the finite grid,
\[
 \mathfrak q_X(F)
   =\max_{\substack{p,q\ {\rm probabilities\ on}\ X\\
                   p_i=q_i\ {\rm for\ each}\ i}}
                 \left(\sum_xp(x)F(x)-\sum_xq(x)F(x)\right).
 \tag{74.14}
\]
\end{proposition}
\begin{proof}
Use the equivalent upper and lower residual
constraints in (74.13), with nonnegative
dual multipliers \(p(x)\) and \(q(x)\).
The coefficient of the free variable \(t\)
forces \(\sum p=1\). The coefficient
of \(c\) forces \(\sum q=\sum p=1\).
For each free variable \(a_i(s)\),
vanishing of its coefficient forces
the \(i\)-th marginals of \(p,q\) to agree.
The remaining objective is the one
displayed in (74.14).
The primal is feasible and bounded below,
so finite linear-programming strong
duality applies. This is the standard
LP theorem, not a new duality principle;
see Boyd--Vandenberghe,
\href{https://web.stanford.edu/~boyd/cvxbook/bv_cvxbook.pdf}
{Convex Optimization}, Chapter 5.

The weak bound is also immediate without
duality: equal one-coordinate marginals
annihilate every \(\sum_i a_i(x_i)\),
and the difference of expectations of
any residual is at most its oscillation.
Thus any explicit primal and balanced
dual pair with matching values is an
exact certificate of optimality.
\end{proof}

For two varying coordinates and all
others fixed, the four-point rectangle
gives the particular lower bound
\[
 \mathfrak q_T(F)\ge\frac12
       |F(a,b)+F(a',b')-F(a,b')-F(a',b)| .
 \tag{74.15}
\]
The factor \(1/2\) makes the positive
and negative point masses each have
total one. Balanced witnesses with
more configurations can be stronger
than a single rectangle.

\subsection{A grid estimate needs a uniform error, not just sampled values}

\begin{proposition}[Continuous-to-grid enclosure]
\label{prop:f12v70-grid}
Suppose \(F\) on \(G^n\) is \(L\)-Lipschitz
for the sum of the group distances.
Choose a finite \(\delta\)-net \(X_i\)
in each group. Then
\[
 \mathfrak q_X(F|_X)\le\mathfrak q_T(F)
               \le\mathfrak q_X(F|_X)+4nL\delta .
 \tag{74.16}
\]
\end{proposition}
\begin{proof}
Restriction of any continuous additive
approximation proves the lower bound.
For the upper bound choose continuous
partitions of unity on each group
subordinate to the balls of radius
\(2\delta\) about its net points.
For any grid additive functions \(a_i\),
extend each by its one-coordinate
partition of unity. Interpolate \(F|_X\)
with the product partition of unity.

At any configuration this interpolant
differs from \(F\) by at most
\(2nL\delta\), because every grid point
having nonzero interpolation weight
is within that sum-distance.
The interpolated residual is a convex
combination of grid residuals, and
therefore has oscillation no larger
than the grid residual oscillation.
Adding the error increases oscillation
by at most \(4nL\delta\).
Take an arbitrarily close grid
minimizer to prove (74.16).
\end{proof}

There is a corresponding enclosure for
the worst-exterior coefficient \(\Gamma_0\).
For clarity consider one repair template;
take maxima over templates afterwards.
Suppose the energies are uniformly
\(L_\xi\)-Lipschitz in the exterior
parameter metric, and an admissible
\(\eta\)-net \(\Xi_\eta\) covers \(\Xi\).
At every \(\xi\in\Xi_\eta\), suppose
certified grid energy values \(\widetilde E^\xi\)
satisfy
\[
 \max_{x\in X}|\widetilde E^\xi(x)-\mathcal E_T^\xi(x)|
                         \le\zeta,\qquad
 Q=\max_{\xi\in\Xi_\eta}\mathfrak q_X(\widetilde E^\xi).
 \tag{74.17}
\]
Then (74.8) and (74.16) imply the rigorous enclosure
\[
 \max\{0,Q-2\zeta\}\le
       \sup_{\xi\in\Xi}\mathfrak q_T(\mathcal E_T^\xi)
       \le Q+2\zeta+4nL\delta+2L_\xi\eta .
 \tag{74.18}
\]
Indeed grid energy errors change the
quotient by at most \(2\zeta\), and
moving to an exterior net point
changes it by at most \(2L_\xi\eta\).

In the actual problem, a certified
evaluation of (74.3) requires lower
as well as upper bounds on the compact
interior minima. An optimizer's sampled
configuration provides only an upper
bound on \(m_b\). It is not the
two-sided \(\zeta\) required here.
Likewise the exterior net must have
a proved covering radius for the
admissible set; arbitrary exterior
samples cannot certify its supremum.
No computational tractability claim
is made for this potentially very
large finite product problem.

\subsection{Exact examples: unique minima do not remove the interaction cost}

Consider a hidden bit \(z\) and endpoint
bits \(x,y\), with classical action
\[
                      S(z,x,y)=(z-x)^2+(z-y)^2 .
 \tag{74.19}
\]
The minimized energy is zero at equal
endpoints and one at unequal endpoints.
Its additive quotient is exactly one:
zero additive fields give residual
range one, while the probability
uniform on \(\{(0,1),(1,0)\}\), minus
the probability uniform on
\(\{(0,0),(1,1)\}\), is a balanced
dual witness of value one.
Adding \(2(x+y)\) makes the energy
have its unique zero at \((0,0)\),
without changing this quotient.
Thus V29's unique section minimum
cannot, by itself, prove (74.12).

The exact normalized hidden-bit
integral is
\[
 Z_\beta(0,0)=Z_\beta(1,1)
       =\frac{1+e^{-2\beta}}2,\qquad
 Z_\beta(0,1)=Z_\beta(1,0)=e^{-\beta}.
 \tag{74.20}
\]
For \(W_\beta=-\log Z_\beta\), the same
primal and dual witnesses give
\[
 \mathfrak q(W_\beta)
       =\beta-\log2+\log(1+e^{-2\beta}),\qquad
              \beta^{-1}\mathfrak q(W_\beta)\longrightarrow1 .
 \tag{74.21}
\]
This model contains a degenerate
interior minimum at unequal endpoints,
which does not obstruct the limit.
It is an illustrative finite model,
not the actual gauge-group integral.

For three bits the energy
\(xy+xz+yz+2(x+y+z)\) also has
additive quotient one. Subtracting
\(3(x+y+z)\) gives residual values
zero at total occupancy zero or
three, and minus one at occupancy
one or two. The dual probability
uniform on the two extreme states
and the one uniform on the other
six states have identical
one-coordinate marginals and
energy-expectation difference one.
Every single two-coordinate
rectangle of this energy has
absolute value one, furnishing
only the weaker lower bound \(1/2\).

The exact script
\begin{center}\small\ttfamily
scripts/verify\_ym\_f12\_v70\_classical\_repair\_quotient.py
\end{center}
checks these matching primal--dual
witnesses. At \(\beta=k\log2\) for
\(k=1,2,4,8,16\), it checks (74.20)
and the exponential of (74.21)
using rational arithmetic, including
invariance under strong additive
endpoint fields. It also certifies
the quotient \(1/2\) of \(F(x,y)=xy\)
on four grids of \([0,1]^2\), and
checks its exact bilinear interpolation.
These checks pass. They neither
evaluate nor bound the actual
\(\Gamma_0\) in (74.9).

\begin{decision}[V70 direction review]
The repair problem has been moved
upstream from a \(\beta\)-dependent
integral estimate to the actual
classical minimized action modulo
one-coordinate fields. The leading
coefficient is now identified by a
proved uniform limit, and both
lower and upper certification
mechanisms are explicit.

The next useful calculation is a
certified bound on this actual
classical repair quantity, starting
with a genuine compact template and
its full exterior dependence.
A lower witness above \(2c_\tau/3\)
would rule out the product-reference
route; an upper bound below it would
validate its repair exponent but
still leave the good-exterior
compact patch gap. Merely assuming
small residual coupling, or invoking
the unique section minimum, would
repeat an unsupported step.
The next companion audit is V75,
and the next ten-version direction
review is V80.
\end{decision}

V69 had 1116295 bytes and SHA--256
\begin{center}\scriptsize\ttfamily
FF20385BCB262F90F6A5E5B5C22E5B9F67032DA6C6C08B305223E7E6E08C7DF8.
\end{center}
Its complete body is retained apart
from the title version. Only active
V69 is replaced by V70. All archives,
companions and previous scripts remain.

\begin{firewall}[Position after V70]
The uniform local classical
reduction, the leading additive
oscillation exponent, and the
finite certification inequalities
are proved under the stated fixed
compact-template hypotheses.
The actual value of \(\Gamma_0\)
and its required comparison with
the tail exponent remain unknown.
No cofinal auxiliary gap, continuum
construction, or physical YM mass
gap has been obtained. The full
objective remains active.
\end{firewall}
\section{V71: an actual compact fibre rectangle obstructs the product reference}
\label{sec:f12v71}

\subsection{Outcome and geometric scope}

V70 identified the classical additive-oscillation
quantity \(\Gamma_0\) governing V69's optimized
one-coordinate product reference. We now obtain
an explicit lower witness from the actual Wilson
action and the actual root-pivot fibre, rather
than from an illustrative bit model.

The witness uses two free boundary links and
their shared pivot over one coarse edge.
It stays in a small commuting family on the
exact coarse fibre. Its mixed action rectangle
is positive and can be calculated without
a Taylor remainder.

To transfer it unchanged through V28's compact
block integration, we impose one explicit
clearance condition: the two free coordinates
are interface coordinates outside every
integrated-block input set \(J_b\).
This is available in sufficiently wide
corridors of the permitted periodic packings.
It is not asserted for every packing satisfying
only V28's three separation conditions.
No new packing has been silently substituted
for an already fixed one.

Under that clearance condition and a sufficiently
small fixed tail threshold, the witness proves
\(\Gamma_0>2c_\tau/3\) for the presently retained
tail certificate. Thus even the optimized
one-coordinate product oscillation comparison,
not merely V69's old surface bound, fails that
repair scale test in this geometry.
It does not show that the actual repair
Poincare constant is large.

\subsection{Three physical links on the exact coarse fibre}

Use zero-based direction labels \(0,1,2,3\);
direction \(2\) is the third coordinate.
In the parity cell at the origin set
\[
 \begin{aligned}
 i&=((0,0,1,1),2),\\
 j&=((0,1,1,0),2),\\
 p&=((0,0,1,0),2).
 \end{aligned}
 \tag{75.1}
\]
Here \(i,j\) are nonpivot free boundary links
and \(p\) is the reference boundary pivot
over the same coarse edge \((0,2)\).
The inherited free-coordinate ordering
assigns indices \(7,16\) to \(i,j\),
respectively. All three links have the
same coarse-edge-tail owner.

Let \(G\) be the compact matrix group in the
programme, of positive Lie dimension, and choose
\[
 X\in\mathfrak g,\qquad \|X\|_{\rm HS}=1.
 \tag{75.2}
\]
For small real \(u,v\), define the fine field by
\[
 W_i=e^{uX},\qquad W_j=e^{vX},\qquad
 W_p=e^{-(u+v)X},\qquad W_l=I\quad(l\ne i,j,p).
 \tag{75.3}
\]
The tree links are all identity. The coarse
field is fixed to \(C=I\).

\begin{proposition}[Exact coarse identity and agreement with the actual chart]
\label{prop:f12v71-fibre}
There is \(t_0>0\), depending only on the fixed
local completion, such that for
\(|u|,|v|\le t_0\), (75.3) lies on the exact
coarse fibre \(K^0(W)=I\) and is the field
specified by the programme's local inverse
pivot chart with these free coordinates.
\end{proposition}
\begin{proof}
Only paths over the coarse edge \((0,2)\)
contain a changed link. Each of its eight
boundary links occurs in exactly two of
the sixteen chronological two-link paths.
All first internal links have value \(I\).
The three changed path angles are therefore
\(u,v,-u-v\), each repeated twice, with
the other ten angles zero.

The reference path has angle \(-u-v\).
All path products commute. On a sufficiently
small common logarithm neighborhood,
the logarithm of the relative path
product is exactly its angle difference
times \(X\). The root-reference block
therefore has exponent equal to the
average of the sixteen path angles:
\[
       \frac1{16}\bigl(2u+2v-2(u+v)\bigr)X=0 .
 \tag{75.4}
\]
Every other coarse output is identity.
This is the nonlinear block formula
evaluated on commuting matrices, not
its linearization.

The fixed completion agrees with this
good local branch near the section.
Choose \(t_0\) so all three links,
relative path products, and output
data remain in that neighborhood.
The inherited uniqueness of the local
inverse pivot map then identifies
(75.3) with the actual completed
field for the declared free data.
\end{proof}

This proposition uses only a small compact
neighborhood of the section. It does not
specify the completion outside that
neighborhood or suppress its other branches.

\subsection{The sixteen affected plaquettes}

Retain the actual action normalization
of (17.1),
\[
 s_p(W)=\tfrac12\|I-W_p\|_{\rm HS}^2,\qquad
 S_0(W)=\sum_p s_p(W).
 \tag{75.5}
\]
Write the eigenvalues of \(X\) as
\(\mathrm i\lambda_a\), so
\(\sum_a\lambda_a^2=1\) and
\(|\lambda_a|\le1\). Set
\[
 s_X(\theta)=\tfrac12\|I-e^{\theta X}\|_{\rm HS}^2
                    =\sum_a(1-\cos(\lambda_a\theta)).
 \tag{75.6}
\]
It is an even function.

\begin{proposition}[Exact two-coordinate Wilson action]
\label{prop:f12v71-action}
For (75.3), in a torus large enough that
the listed plaquettes are distinct,
\[
 \begin{aligned}
 S_0(u,v)=&
      5s_X(u)+5s_X(v)+4s_X(u+v)\\
      &+s_X(2u+v)+s_X(u+2v).
 \end{aligned}
 \tag{75.7}
\]
\end{proposition}
\begin{proof}
A link in four dimensions belongs to
six plaquettes. The pivot shares one
plaquette with \(i\), since their bases
differ by one step in direction \(3\),
and one with \(j\), since those bases
differ by one step in direction \(1\).
The two free links have no common
plaquette. No plaquette contains all
three changed links.

Each free link therefore has five
plaquettes not containing another
changed link, giving \(5s_X(u)\)
and \(5s_X(v)\). Four pivot plaquettes
contain neither free link and give
\(4s_X(u+v)\). The shared plaquettes
contain oppositely oriented parallel
links, so their angles, up to overall
sign, are \(2u+v\) and \(u+2v\).
All other link values are identity,
so there are no further contributions.
There are \(5+5+4+1+1=16\) distinct
affected plaquettes. Evenness of
\(s_X\) removes the harmless
orientation signs.
\end{proof}

The corresponding coefficient count is
\begin{center}\small
\begin{tabular}{c|ccccc}
angle up to sign & \(u\) & \(v\) & \(u+v\) & \(2u+v\) & \(u+2v\)\\ \hline
plaquette count & 5 & 5 & 4 & 1 & 1
\end{tabular}
\end{center}
This table is checked both by the
inherited Laurent fibre stencil and
by a separate direct enumeration of
physical plaquette words in the script.

\subsection{A positive balanced rectangle without an expansion error}

For \(0<t\le t_0\), define
\[
 \mathcal R_X(t)=S_0(t,t)+S_0(-t,-t)
                       -S_0(t,-t)-S_0(-t,t).
 \tag{75.8}
\]

\begin{proposition}[Actual mixed action lower witness]
\label{prop:f12v71-rectangle}
The rectangle is exactly
\[
 \mathcal R_X(t)
       =16\sum_a\sin^2(\lambda_a t)
                              (1+\cos(\lambda_a t)).
 \tag{75.9}
\]
For \(0<t\le\min\{t_0,1\}\),
\[
                         \tfrac12\mathcal R_X(t)\ge8t^2 .
 \tag{75.10}
\]
\end{proposition}
\begin{proof}
For one eigenvalue and one plaquette
with angle \(a u+b v\), its contribution
to (75.8) is
\[
        4\sin(a\lambda t)\sin(b\lambda t).
 \tag{75.11}
\]
The pure \(u\) and pure \(v\) terms
vanish. The four \((1,1)\) terms
give \(16\sin^2(\lambda t)\).
The \((2,1)\) and \((1,2)\) terms
give \(8\sin(\lambda t)\sin(2\lambda t)\).
Using \(\sin(2z)=2\sin z\cos z\)
proves (75.9) after summing eigenvalues.

For \(|z|\le1\), the elementary Taylor
inequalities give
\(|\sin z|\ge(5/6)|z|\) and
\(1+\cos z\ge3/2\). Consequently
\[
 \sin^2z(1+\cos z)\ge(25/24)z^2\ge z^2 .
 \tag{75.12}
\]
Apply this to \(z=\lambda_at\) and use
\(\sum_a\lambda_a^2=1\).
\end{proof}

For \(G=SU(2)\) with the unnormalized
generator \(\operatorname{diag}(\mathrm i,-\mathrm i)\),
the exact rectangle instead reads
\(32\sin^2t(1+\cos t)\).
The norm of that generator is \(\sqrt2\);
it must not be substituted into the
unit-Hilbert--Schmidt comparison without
rescaling. Formula (75.9) keeps this
normalization explicit.

\subsection{Transfer to the actual minimized repair action}

Let \(\mathcal S\) and \(J_b\) be exactly
V28's interface and the full input sets
of its integrated blocks. Require the
following clearance for the translate
of the motif under consideration:
\[
                   i,j\in\mathcal S,\qquad
                   \{i,j\}\cap\bigcup_b J_b=\varnothing .
 \tag{75.13}
\]
Choose the actual repair \(T_i\) of
V58--V69. Both coordinates have the
same owner, so \(j\in N_i\subset T_i\)
by V65; thus the repair contains
both varied free coordinates. All
required comparison neighborhoods
must fit in the ambient torus,
as in V50. Set the retained field
and all other interface values to
their flat section values.

\begin{theorem}[A lower bound on the actual classical repair quotient]
\label{thm:f12v71-actual-quotient}
Under (75.13) and the stated availability
conditions, the V70 worst-parameter
coefficient satisfies
\[
                       \Gamma_0\ge\tfrac12\mathcal R_X(t)\ge8t^2
               \qquad(0<t\le\min\{t_0,1\}).
 \tag{75.14}
\]
The actual finite-temperature coefficient
also obeys
\[
                          \Gamma_\beta\ge
                              \beta\,\tfrac12\mathcal R_X(t)-C_{\rm loc},
 \tag{75.15}
\]
with a fixed finite local constant.
\end{theorem}
\begin{proof}
By (75.13) no compact block integral
\(Z_{b,\beta}\) depends on \(u\) or \(v\).
Nor does any of its minimized actions
\(m_b\). The exact action assignment
in (74.1)--(74.3) therefore shows that
the rectangle of the actual
\(\mathcal E_{T_i}^{\,\xi}\), on these
four configurations, equals (75.8).
This step uses the full input sets,
not just the visible locations of
the integrated variables.

Put probability \(1/2\) on each of
\((u,v)=(t,t),(-t,-t)\), and put
probability \(1/2\) on each of
\((t,-t),(-t,t)\) for the second
measure. All other repair coordinates
are fixed to their section values
under both measures. The two
probabilities have identical
one-group marginals, so they
annihilate every additive
one-coordinate potential.
Their energy-expectation difference
is \(\mathcal R_X(t)/2\).
V70's balanced-witness inequality
and the supremum over repairs and
parameters prove (75.14).

At finite \(\beta\), all block
partition functions are still
constant on the witness. The
only additional varying terms
are the bounded local inverse-pivot
amplitude potentials. The difference
of their expectations is bounded
in absolute value by their local
oscillation \(C_{\rm loc}\).
Their number is bounded by the
fixed incidence of \(i,j\).
The same balanced witness proves
(75.15), independently of volume.
\end{proof}

The clearance condition is substantive.
If an integrated block has either
coordinate in \(J_b\), minimization
can change the rectangle, and this
proof gives no bound on that change.
V28 permits periodic parent cubes
with arbitrarily enlarged but fixed
corridors. Since every template
and every input set has bounded
owner extent, choosing a corridor
wide enough leaves a cell whose
two motif coordinates satisfy
(75.13). This establishes an
allowed family to which the theorem
applies. The three original V28
separation requirements alone do
not assert that every chosen
corridor has such an unused cell.

\subsection{Comparison with the present tail certificate}

When fixing the V50 tail threshold,
one may choose it sufficiently small
that
\[
          0<\tau\le\min\{2,8t_0\},\qquad
                    t=\tau/8 .
 \tag{75.16}
\]
All four witness configurations then
lie on the valid chart branch.
In particular their free-link
distances from the section are at
most \(t\), because
\(\|e^{uX}-I\|_{\rm HS}\le|u|\).
The witness does not require an
off-chart configuration.
The radius and constants of the
actual repairs must be fixed at
this chosen threshold, as before.

The tail used in V50 and V66 is
\[
 c_\tau=\frac{\tau^2}{16A_{\rm tail}}.
 \tag{75.17}
\]
For the inherited deterministic
choice from (17.3), or an enlargement,
\[
 A_{\rm tail}\ge
        2\bigl(153+64(9+6L_{\log})^2\bigr)
                           \ge10674 .
 \tag{75.18}
\]
This is the currently populated
tail constant, not an assertion
about the best possible tail rate.

\begin{corollary}[Product-reference scale obstruction in the cleared corridor]
\label{cor:f12v71-product-obstruction}
Under the preceding clearance and
small-threshold hypotheses,
\[
             \Gamma_0\ge\frac{\tau^2}{8}
                        >\frac{\tau^2}{24A_{\rm tail}}
                         =\frac23c_\tau .
 \tag{75.19}
\]
The lower bound is at least \(3A_{\rm tail}\),
and hence at least \(32022\), times
the threshold \(2c_\tau/3\).
\end{corollary}
\begin{proof}
Insert \(t=\tau/8\) into (75.14).
This gives \(8t^2=\tau^2/8\).
Divide by the positive threshold
\(\tau^2/(24A_{\rm tail})\), and
use (75.18).
\end{proof}

By V70 this rules out its strict
replacement comparison test
\(\Gamma_0<2c_\tau/3\) in the
stated geometry and parameter range.
Equation (75.15) supplies the
same leading obstruction directly
at finite temperature. The
confining one-coordinate potentials
can be chosen arbitrarily in the
reference: the balanced rectangle
annihilates every such choice.

The result is more limited than
a spectral-gap obstruction.
Oscillation comparison with a
one-coordinate product law charges
genuine local correlations even
when those correlations coexist
with strong convexity or a large
heat-bath gap. V39's boxed gap
and V69's interacting example
already illustrate that distinction.
Neither the small-chart interaction
found here nor its large ratio
to the current conservative tail
rate proves that the actual
compact repair is slow.
An improved tail theorem could
also change the comparison threshold;
no optimal tail exponent is proved
or ruled out here.

\subsection{Exact actual-geometry checks}

The script
\begin{center}\small\ttfamily
scripts/verify\_ym\_f12\_v71\_compact\_rectangle.py
\end{center}
reuses the exact root-tree paths,
free-coordinate ordering, and
Laurent curl--fibre maps already
checked in V5 and the inherited
action-jet scripts. It first
verifies the coarse linear kernel
identity and locates the pair
\((7,16)\). It then independently
enumerates physical plaquettes
with the three changed links in
(75.3), without periodic wrapping,
and verifies equality with the
Laurent stencil coefficient list.
The enumeration gives precisely
the sixteen plaquettes in (75.7).

The six mixed plaquette coefficient
pairs, including orientation signs,
are
\[
                    (1,1),\ (1,2),\ (-2,-1),\
                       (1,1),\ (-1,-1),\ (-1,-1).
 \tag{75.20}
\]
They have \(\sum ab=8\) and
\(\sum|ab|(a^2+b^2)=28\), consistent
with the quadratic and fourth-order
stocks of the inherited fibre
stencil. These stocks are checks,
not a Taylor approximation used
in (75.9).
The maximum simultaneous changed
link coefficient is two, at the
pivot. The script verifies the
zero coarse angle sums, the
exact trigonometric lower-bound
coefficient using rational
inequalities, and the ratio
in Corollary~\ref{cor:f12v71-product-obstruction}.
All checks pass.

The script does not claim to have
enumerated the \(J_b\) of a particular
numerically fixed packing, nor
computed the completion-dependent
radius \(t_0\). These enter as the
explicit geometric and chart
hypotheses of the theorem. Their
existence for the declared
wide-corridor family follows from
finite template range and the
local branch neighborhood, not
from the finite stencil test alone.

\begin{decision}[Retain correlations in the next repair argument]
The new witness is an actual
compact-fibre calculation and
identifies a concrete limitation
of the optimized one-coordinate
product reference. For a packing
with the cleared corridor, do
not continue seeking an upper
bound contradicted by (75.19).

The next repair argument must
either retain these correlations
in its reference or work directly
with the actual conditional
Dirichlet form, instead of
charging their full action
oscillation. For packings without
(75.13), an actual block-minimum
calculation is still needed;
the present rectangle must not
be transferred through an
uncontrolled interior minimization.
The independent good-exterior
compact patch gap and physical
continuum problem remain unpaid.
The next companion audit is V75
and direction review is V80.
\end{decision}

V70 had 1134326 bytes and SHA--256
\begin{center}\scriptsize\ttfamily
AA3883F72D94064DA71087614C4473EDEEAFC5BF724C5057ABD988D2805F16BD.
\end{center}
Its complete body is retained
apart from the title version.
Only active V70 is replaced by
V71. Archives, companions and
earlier scripts are preserved.

\begin{firewall}[Position after V71]
The exact coarse-fibre family,
Wilson rectangle, and conditional
product-reference obstruction
are proved with their normalization,
clearance and threshold hypotheses.
No actual repair gap has been
upper- or lower-bounded by this
rectangle. No cofinal auxiliary
gap, continuum construction, or
physical Yang--Mills mass gap
is claimed. The full goal remains
active.
\end{firewall}
\section{V72: the correlation obstruction also crosses owner blocks}
\label{sec:f12v72}

\subsection{Why grouping the V71 pair is not yet a repair theorem}

V71 ruled out a one-coordinate product oscillation
comparison in a cleared corridor. A natural
response is to retain all correlations within
each coarse owner, or within a bounded group
of owners, in a product reference. This note
tests that response on the actual fibre.

There is a second exact commuting configuration
with one varied free coordinate at each of
two neighboring owners. Its unique mixed
plaquette has a nonzero rectangle. The
rectangle annihilates every reference
potential whose factor contains at most
one of the two varied coordinates, even
if that factor contains many other
correlated coordinates.

Thus grouping by owners removes the
specific within-owner rectangle of V71,
but not this cross-owner obstruction.
For bounded owner clusters, a sufficiently
long cleared corridor path necessarily
has a cut between two factors. We prove
the resulting limitation of this class
of comparison certificates. It is not
a lower bound on the true repair cost,
and it is not a statement about arbitrary
nonproduct correlated references.

\subsection{An exact two-owner coarse-fibre configuration}

Continue to use zero-based direction
labels. Let \(e_0=(1,0,0,0)\) be the
first coarse-cell displacement. At
owners \(0,e_0\), use the four physical
boundary links
\[
 \begin{aligned}
 i&=((1,0,1,1),2),&
 p_0&=((0,0,1,0),2),\\
 j&=((2,0,1,1),2),&
 p_1&=((2,0,1,0),2).
 \end{aligned}
 \tag{76.1}
\]
The free-coordinate index of \(i\)
is \(30\) at owner \(0\); that of
\(j\) is \(7\) at owner \(e_0\).
The two \(p\)'s are the reference
pivots for the corresponding coarse
edges in direction \(2\).

Choose the same type of unit-Hilbert--Schmidt
Lie generator \(X\) as in V71 and set
\[
 W_i=e^{uX},\quad W_{p_0}=e^{-uX},\qquad
 W_j=e^{vX},\quad W_{p_1}=e^{-vX},
 \qquad W_l=I\ \hbox{otherwise}.
 \tag{76.2}
\]
There is a fixed \(t_\times>0\) for
which this lies on the actual good
inverse-pivot branch for
\(|u|,|v|\le t_\times\).
Indeed at each of the two coarse
edges its free boundary link and
pivot occur twice in the sixteen
paths. The respective angle sums
are \(2u-2u=0\) and \(2v-2v=0\).
On the common commuting logarithm
neighborhood this proves \(K^0(W)=I\)
exactly, just as in V71.
All other coarse outputs are
unchanged. Local inverse uniqueness
identifies (76.2) with the actual
completed field.

\begin{proposition}[The actual cross-owner action]
\label{prop:f12v72-action}
With \(s_X(\theta)\) as in (75.6),
the Wilson action of (76.2) is exactly
\[
 S_0(u,v)=11s_X(u)+9s_X(v)+s_X(2v)+s_X(u-v).
 \tag{76.3}
\]
\end{proposition}
\begin{proof}
All four changed links are parallel.
The only pairs sharing a plaquette
are \(i,j\), whose bases differ by
one step in direction \(0\), and
\(j,p_1\), whose bases differ by
one step in direction \(3\).
The other base differences have
at least two unit steps, so those
pairs cannot be opposite edges
of a single plaquette. No
plaquette contains three changed
links.

The link \(i\) has five pure-\(u\)
plaquettes and \(p_0\) has six,
giving \(11s_X(u)\). The link
\(j\) has four pure-\(v\)
plaquettes and \(p_1\) has five,
giving \(9s_X(v)\). Their shared
plaquette has angle \(2v\), up
to sign. The sole cross-owner
shared plaquette has angle
\(u-v\), up to sign. This
accounts for \(11+9+1+1=22\)
distinct plaquettes. All other
plaquette values are identity.
\end{proof}

\begin{corollary}[Nonzero cross-owner rectangle]
\label{cor:f12v72-rectangle}
For \(0<t\le\min\{t_\times,1\}\), write
\[
 R_\times(t)=S_0(t,t)+S_0(-t,-t)
                         -S_0(t,-t)-S_0(-t,t).
 \tag{76.4}
\]
Then
\[
 R_\times(t)=-4\sum_a\sin^2(\lambda_at),\qquad
                   -\tfrac12R_\times(t)\ge t^2,
 \tag{76.5}
\]
where \(\mathrm i\lambda_a\) are
the eigenvalues of \(X\).
\end{corollary}
\begin{proof}
Only \(s_X(u-v)\) survives in the
rectangle. Its contribution for
each eigenvalue is
\(4\sin(\lambda_at)\sin(-\lambda_at)\).
For \(|\lambda_at|\le1\),
\(|\sin(\lambda_at)|\ge(5/6)|\lambda_at|\).
Since \(\sum_a\lambda_a^2=1\),
\(-R_\times/2\ge(25/18)t^2\ge t^2\).
\end{proof}

\subsection{A product of correlated factors still has cut rectangles}

Fix an actual repair \(T\) and a partition
\(\mathcal P=\{D_1,\ldots,D_m\}\) of its
free interface coordinates. A factor
\(D_\alpha\) may contain arbitrarily
correlated coordinates. Define
\[
 \mathfrak q_{\mathcal P}(F)=
       \inf_{V_\alpha\in C(G^{D_\alpha})}
             \operatorname{osc}\left(F-\sum_\alpha V_\alpha(v_{D_\alpha})\right).
 \tag{76.6}
\]
This is not V70's single-coordinate
quotient unless all factors are
singletons. Its elementary stability
and homogeneity properties are
identical to (74.8).

Assume a translate of (76.1) satisfies
\[
 i,j\in T\subset\mathcal S,\qquad
 \{i,j\}\cap\bigcup_b J_b=\varnothing,\qquad
 i\in D_\alpha,\ j\in D_{\alpha'},\ \alpha\ne\alpha'.
 \tag{76.7}
\]
Take flat retained data and section
values for every other interface
coordinate. The full input-set
clearance in (76.7) is the same
substantive condition as in V71.
It cannot be replaced by mere
disjointness from the integrated
coordinate sets.

\begin{theorem}[Actual lower witness for a cut between reference factors]
\label{thm:f12v72-factor-cut}
Under (76.7), the actual classical
repair energy and conditional potential
satisfy
\[
 \begin{aligned}
 \mathfrak q_{\mathcal P}(\mathcal E_T^\xi)
                       &\ge-\tfrac12R_\times(t)\ge t^2,\\
 \mathfrak q_{\mathcal P}(W_{T,\beta}^{\,\xi})
                       &\ge\beta t^2-C_\times ,
 \end{aligned}
 \tag{76.8}
\]
for a fixed finite local constant
\(C_\times\).
\end{theorem}
\begin{proof}
Use the probability uniform on
\((u,v)=(t,-t),(-t,t)\) as the
positive witness, and the probability
uniform on \((t,t),(-t,-t)\) as
the negative witness. All other
coordinates are fixed under both.
Their marginals on every complete
factor \(D_\alpha\) agree, because
no factor contains both varied
coordinates. They therefore
annihilate every function
\(\sum_\alpha V_\alpha(v_{D_\alpha})\),
not just one-coordinate potentials.

By clearance, none of the integrated
block minima or partition functions
depends on \(u,v\). The classical
energy-expectation difference is
exactly \(-R_\times(t)/2\).
At finite temperature the remaining
amplitude-potential difference is
bounded by the oscillation of the
fixed number of incident terms.
This proves (76.8), using the
balanced-witness inequality.
\end{proof}

The theorem applies to an actual
anchor repair containing the pair.
For instance the inherited
\(N_i\subset T_i\) contains the
neighboring owner when its range
bound is at least one, as it is
in the chosen specification.
All comparison neighborhoods must
still fit in the ambient torus.

\subsection{Why bounded owner clusters cannot absorb a long cleared path}

Consider the special, natural class
of factor partitions in which each
\(D_\alpha\) is a union of complete
owner-coordinate sets \(T\cap\operatorname{owner}^{-1}(r)\).
Assume each factor contains at most
\(K\) distinct nonempty owners.
The owners need not form a connected
set; the following argument does
not assume connected factors.

Suppose \(T\) contains a chain of
\(K+1\) owners
\[
                   r,\ r+e_0,\ldots,r+Ke_0
 \tag{76.9}
\]
such that at every neighboring pair
the corresponding translated free
coordinates in (76.1) belong to \(T\)
and satisfy the clearance in (76.7).
Then at least one neighboring pair
belongs to different factors.
Otherwise equality of the factors
would propagate along the whole
chain, placing all \(K+1\) owners
in one factor, a contradiction.
Theorem~\ref{thm:f12v72-factor-cut}
therefore applies to every such
bounded-owner factor partition.

This is a condition on a path
inside the actual repair, not an
assertion that every repair or
packing contains a path of every
length. Sufficiently wide empty
transverse strips in periodic
cube packings can provide cleared
corridor lines; the available
segment must still lie inside
the fixed repair. The parameter
\(K\) cannot exceed the verified
segment length in this argument.
In particular the one-owner
product reference already fails
whenever one cleared neighboring
pair is present.

\subsection{The factor's own single-coordinate dynamics cannot be omitted}

A correlated-factor reference has
another cost besides its density
oscillation. Let
\[
                \rho=\prod_{\alpha=1}^m\rho_\alpha,\qquad
 d\rho_\alpha\propto e^{-V_\alpha}\prod_{i\in D_\alpha}dv_i ,
 \tag{76.10}
\]
and let \(C_\alpha\) be a Poincare
constant for the rate-one
single-coordinate heat baths
inside the factor \(\rho_\alpha\).
Put \(C_{\max}=\max_\alpha C_\alpha\).
Product variance tensorization
and the same bounded-density
comparison as in V69 give
\[
 \operatorname{Var}_{\nu_T^\xi}f
       \le e^{\Omega}C_{\max}
                   \sum_{i\in T}e_i^{\nu_T^\xi}(f),\qquad
 \Omega=\operatorname{osc}
           \left(W_{T,\beta}^{\,\xi}-\sum_\alpha V_\alpha\right).
 \tag{76.11}
\]
Indeed, if \(m\le d\nu_T^\xi/d\rho\le M\),
then \(\operatorname{Var}_{\nu_T^\xi}\le
M\operatorname{Var}_{\rho}\),
\(\operatorname{Var}_{\rho}\le C_{\max}\sum_i e_i^\rho\),
and \(e_i^{\nu_T^\xi}\ge m e_i^\rho\);
also \(M/m=e^\Omega\).

If each factor has at most \(n_*\)
coordinates, then every valid
\(C_\alpha\) satisfies
\[
                         C_\alpha\ge1/|D_\alpha|\ge1/n_* .
 \tag{76.12}
\]
To see this, each one-coordinate
residual is a projection with norm
one, so their sum has norm at most
\(|D_\alpha|\). Test the Poincare
inequality on any nonconstant
centered function in that factor.
Thus a factor gap cannot provide
an exponentially growing rate
that cancels a positive
oscillation exponent at fixed
factor size.

Combining (76.8), (76.11), and
(76.12) shows that the numerical
coefficient furnished by this
comparison, for every choice
of its factor potentials, is at least
\[
                    e^\Omega C_{\max}
                        \ge n_*^{-1}e^{\beta t^2-C_\times}.
 \tag{76.13}
\]
For owner clusters of capacity \(K\)
one may use \(n_*=q_*K\).
This is a lower bound on the
comparison certificate, not on
the optimal Poincare constant
of \(\nu_T^\xi\). Its derivation
must not be read as reversing
the inequality in (76.11).

A joint refresh of an entire
factor has gap one by definition.
That different generator does
not justify setting
\(C_\alpha=1\) for its internal
single-coordinate updates.
Taking \(T\) itself as one
factor can remove the density
oscillation entirely, but then
its required \(C_\alpha\) is
the original unresolved repair
constant. No gap has been proved
by that relabeling.

\subsection{Comparison with the retained tail rate}

Fix a threshold with
\[
 0<\tau\le\min\{2,8t_\times\},\qquad
 t=\tau/8,\qquad
 c_\tau=\tau^2/(16A_{\rm tail}).
 \tag{76.14}
\]
As in V71, for the inherited
deterministic choice
\(A_{\rm tail}\ge10674\).
The exponent floor in (76.13)
is now
\[
             t^2=\frac{\tau^2}{64}
                  >\frac{\tau^2}{24A_{\rm tail}}
                         =\frac23c_\tau .
 \tag{76.15}
\]
Its ratio to the required threshold
is \(3A_{\rm tail}/8\ge16011/4\).
At fixed \(\tau,K\), the prefactors
\(n_*^{-1}\) and \(e^{-C_\times}\)
do not change this exponent.
Hence such a product of bounded
owner-cluster references cannot
supply the repair exponent demanded
by V68, whenever the verified path
condition forces a cut.

This does not rule out a better
tail estimate, a packing without
the stated cleared path, a
nonproduct reference, or a direct
conditional Dirichlet-form argument.
It rules out the particular
proposal to solve the existing
comparison problem merely by
absorbing correlations inside
bounded owner factors.

\subsection{Exact checks and the next direction}

The script
\begin{center}\small\ttfamily
scripts/verify\_ym\_f12\_v72\_cross\_owner\_rectangle.py
\end{center}
checks the actual free indices
\(30,7\), with the second one
translated by \(e_0\). It compares
an independent physical-link
plaquette enumeration with the
inherited Laurent curl--fibre
stencil, including the displacement
of the second coordinate. Both
give exactly 22 affected plaquettes.
The unique mixed plaquette has
fine base \((1,0,1,1)\), directions
\((0,2)\), and oriented
coefficients \((-1,1)\).
Thus \(\sum ab=-1\), consistent
with the mixed quadratic term
of (76.3).

The script independently checks
the zero chronological coarse
angle sums at both owners and
the action multiplicities
\(11,9,1,1\). It verifies the
rational sine lower-bound
coefficient and the tail ratio.
It also checks the elementary
path-partition argument in all
278 admissible labeled cases
of capacities \(K=1,2,3\).
All checks pass.
The general path argument is
proved above, not inferred
from those small cases.
As in V71, the checker does
not silently certify clearance
for an unspecified packing or
an explicit completion radius.

\begin{decision}[Stop replacing local correlations by product oscillation]
The actual fibre now provides
both within-owner and cross-owner
balanced witnesses. On a verified
cleared path, bounded owner
grouping cannot eliminate the
interaction cost that defeats
the current rare-event scale.
This closes that specific
proposed response to V71.

The next useful repair work
must estimate its conditional
dynamics, or introduce a
correlated reference with an
independently proved
single-coordinate gap and a
controlled comparison that
does not pay the cut action
oscillation above. The actual
compact patch gap, auxiliary
to physical-time identification,
and continuum construction
remain separate requirements.
The next companion audit
remains V75; the next
direction review is V80.
\end{decision}

V71 had 1151066 bytes and SHA--256
\begin{center}\scriptsize\ttfamily
912E9A115F14EE06821FE348A13E87EC71B692387F2AF0CDA2A37B9ACC896D76.
\end{center}
Its complete body is retained
apart from the title version.
Only active V71 is replaced by
V72. Archives, companions and
previous scripts are preserved.

\begin{firewall}[Position after V72]
The exact cross-owner action
and the cut-factor comparison
obstruction are proved with
explicit clearance, factor
separation and threshold
hypotheses. They do not
establish a slow actual repair
or disprove a mass gap.
No full compact repair bound
meeting the required scale,
cofinal auxiliary gap, continuum
construction, or physical
Yang--Mills gap has been proved.
The full objective remains active.
\end{firewall}
\section{V73: a direct compact-patch criterion with an explicit escape cost}
\label{sec:f12v73}

\subsection{Change the estimate, not the target dynamics}

V69--V72 identified why several oscillation comparisons
cannot meet the chosen rare-event scale. They did
not show that the actual conditional dynamics is
slow. We now give a different, direct criterion:
combine the already available Poincare inequality
on a small chart with an escape estimate for the
same compact conditional diffusion outside a
smaller chart.

The local-to-global Lyapunov method is classical;
see Cattiaux--Guillin,
\href{https://arxiv.org/abs/1001.1822}
{Functional Inequalities via Lyapunov conditions},
and Cattiaux--Guillin--Zitt,
\href{https://arxiv.org/abs/1012.5274}
{Poincare inequalities and hitting times}.
No new general principle is claimed. The work
here records the constants needed by this
programme, constructs a cutoff whose gradient
cost is independent of the number of chart
coordinates, and checks precisely which actual
conditional laws have the local input.

The new escape estimate is not supplied by the
rarity bound or by the energy floor. It remains
an explicit missing hypothesis. The populated
local input below concerns good-exterior patches;
it does not silently become an all-exterior
repair estimate.

\subsection{A local Poincare inequality and a killed spectral floor}

Let \(M=G^P\) be a finite compact product, with
the product bi-invariant metric and a positive
smooth probability density
\(\nu=Z^{-1}e^{-W}dv\). Set
\[
 L=-\Delta+\nabla W\cdot\nabla,\qquad
 \mathcal E(f)=\int_M|\nabla f|^2\,d\nu .
 \tag{77.1}
\]
This is the positive diffusion operator in
the normalization used in V60.

Choose nested chart regions
\(\overline{\mathcal G_0}\subset\mathcal G_1\)
of positive measure, and a Lipschitz function
\[
 0\le\chi\le1,\qquad
 \chi=1\ \hbox{on }\mathcal G_0,\qquad
 \chi=0\ \hbox{off }\mathcal G_1,\qquad
                      |\nabla\chi|^2\le M_\chi .
 \tag{77.2}
\]
Assume the actual normalized restriction
\(\nu_1=\nu(\,\cdot\mid\mathcal G_1)\)
has a Neumann Poincare inequality
\[
 \operatorname{Var}_{\nu_1}(f)
           \le\gamma^{-1}\int_{\mathcal G_1}|\nabla f|^2\,d\nu_1,
                  \qquad\gamma>0 .
 \tag{77.3}
\]
Define the killed floor
\[
 \lambda_0=
 \inf_{\substack{0\ne v\in H^1(M,\nu)\\
                       v=0\ {\rm a.e.\ on}\ \mathcal G_0}}
          \frac{\mathcal E(v)}{\|v\|_{L^2(\nu)}^2}.
 \tag{77.4}
\]
For the regular chart regions used below,
this is the bottom of the diffusion
killed on reaching \(\mathcal G_0\).
Functions vanish on the smaller chart,
not outside the larger one.

\begin{theorem}[Local-chart plus escape assembly]
\label{thm:f12v73-assembly}
If a positive lower bound \(\lambda\le\lambda_0\)
has been established, then
\[
 \operatorname{Var}_{\nu}(f)
       \le C_{\rm esc}\mathcal E(f),\qquad
 C_{\rm esc}=
       \frac1\gamma+\frac2\lambda
                             +\frac{2M_\chi}{\gamma\lambda}.
 \tag{77.5}
\]
Equivalently, the full compact diffusion gap obeys
\[
                   g\ge
       \frac{\gamma\lambda}{\lambda+2\gamma+2M_\chi}.
 \tag{77.6}
\]
\end{theorem}
\begin{proof}
Subtract the larger-chart mean:
\(h=f-\nu_1f\), and set \(v=(1-\chi)h\).
This \(H^1\) function vanishes on
\(\mathcal G_0\), and equals \(h\)
outside \(\mathcal G_1\).
Multiplication by the Lipschitz
cutoff is legitimate in \(H^1\).
The product rule and
\(|a+b|^2\le2|a|^2+2|b|^2\) give
\[
 \mathcal E(v)
     \le2\mathcal E(f)
               +2M_\chi\int_{\mathcal G_1}h^2\,d\nu .
 \tag{77.7}
\]
The gradient of \(\chi\) is zero
off the larger chart, almost
everywhere. By the killed floor,
\(\|v\|_2^2\le\lambda^{-1}\mathcal E(v)\).
Also (77.3), multiplied by
\(\nu(\mathcal G_1)\), gives
\[
                    \int_{\mathcal G_1}h^2\,d\nu
                                   \le\gamma^{-1}\mathcal E(f).
 \tag{77.8}
\]
Finally,
\[
 \operatorname{Var}_\nu(f)
 \le\|h\|_2^2
 \le\int_{\mathcal G_1}h^2\,d\nu+\|v\|_2^2.
 \tag{77.9}
\]
Insert the preceding bounds. Approximation
extends the argument from smooth
functions to the full form domain.
\end{proof}

There is no loss proportional to
\(\nu(\mathcal G_1)^{-1}\) in this
assembly. This does not make a tiny
good region harmless: its difficulty
can appear in \(\lambda_0\).
Indeed if the full gap \(g\) were
already known, then any function
in (77.4) satisfies
\[
 \operatorname{Var}_\nu(v)
       \ge\nu(\mathcal G_0)\|v\|_2^2,\qquad
                 \lambda_0\ge g\,\nu(\mathcal G_0).
 \tag{77.10}
\]
The first inequality is Cauchy--Schwarz
on the complement of \(\mathcal G_0\).
Using (77.10) to populate (77.6)
would be circular. The purpose of
(77.4) is to isolate an independent
escape estimate, not to assume the
original gap under a different name.

\subsection{A cutoff without a volume factor}

Use product normal coordinates
\(x\in(-a,a)^N\), \(N=\dim(G)|P|\),
on the larger chart, with \(a>0\)
fixed and sufficiently small.
Let \(\mathcal G_0\) be the inner
cube \((-a/2,a/2)^N\).
Define, in that chart,
\[
 \chi(x)=\min\left\{1,\left(2-\frac{2}{a}
                                     \max_k|x_k|\right)_+\right\},
 \tag{77.11}
\]
and extend it by zero outside
the larger chart. It is Lipschitz
and has the required traces.
Since the maximum-coordinate norm
is one-Lipschitz for the Euclidean
norm, its weak gradient has norm
at most one almost everywhere.
If \(g^{-1}\le c_{\rm inv}I\)
in these product coordinates,
\[
                          M_\chi\le4c_{\rm inv}/a^2 .
 \tag{77.12}
\]
The metric constant is the maximum
of a one-group chart constant,
not a sum over groups.
The bound is independent of \(N\).
Nondifferentiability where maxima
tie causes no problem for the
Sobolev product rule.

For comparison, multiplying many
one-coordinate cutoffs and summing
their gradient bounds can introduce
an unnecessary factor \(N\).
Equation (77.11) avoids that
specific loss; it supplies no
escape estimate by itself.

\subsection{The actual local input and its exterior restriction}

Take an actual full-compact
interface conditional law on
a patch \(P\), with all its free
group variables untruncated.
Assume the retained data obey
the hypotheses of V28, and the
frozen interface variables in
every effective term touching
\(P\) lie in the chosen good chart.
Variables farther away do not
affect this conditional potential.
Thus the good-exterior collar
must cover the complete effective
support range, as in V61.

On \(\mathcal G_1\), V28's measured
coordinate potential, with the
Haar chart factors included,
has Hessian at least
\((m/2)\beta I\). This also holds
after freezing the good exterior:
it is a principal Hessian block,
and any irrelevant farther
coordinates can be set to the
section. The convex-box Poincare
inequality already used in V39
therefore gives (77.3) with
\[
                 \gamma=c_{\rm box}\beta,\qquad
                 c_{\rm box}=m/(2c_{\rm met})>0 ,
 \tag{77.13}
\]
where \(g\le c_{\rm met}I\)
on the product chart. This
accounts for converting the
coordinate gradient to the
Riemannian gradient in (77.1).
Both chart metric constants
are independent of patch size.
The exact compact block integrals
remain in the potential.

For this actual conditional
diffusion, V60's Bochner and
Poisson argument applies to the
finite set of unfrozen groups.
Its mixed-Hessian row majorant
\(K_\beta=O((\beta+1)^2)\)
is independent of the number
of free groups. If \(h_P\)
is the corresponding rate-one
single-group heat-bath gap,
then
\[
 h_P\ge\frac{g_P}{g_P+K_\beta}
       \ge\frac1{1+K_\beta C_{\rm esc}}
       =\frac{\gamma\lambda}
          {\gamma\lambda+K_\beta(\lambda+2\gamma+2M_\chi)} .
 \tag{77.14}
\]
The first inequality is inherited
from V60; it is not an
identification of the two
generators or their clocks.

If a weighted patch uses rates
bounded below by \(\alpha>0\),
its gap is at least
\[
                           \frac{\alpha}{1+K_\beta C_{\rm esc}}.
 \tag{77.15}
\]
For V63's pyramid,
\(\alpha=\ell^{-1/2}+\ell^{-1}\).
Equations (77.13)--(77.15)
make the needed escape accuracy
explicit when testing any of
the earlier local gap thresholds.
They do not restore the
incompatible repair certificates
ruled out in V69--V72.

There are two important limitations.
First, \(\lambda\) in (77.14)
has not been bounded at the
required scale for the actual
compact patch. Second, V66's
repair comparison needs all
compact exteriors, whereas the
local input just populated uses
a good exterior. A bad frozen
neighbor cannot be replaced by
a chart value when it occurs
in a term touching \(P\).
The theorem can be used for
other exteriors only after
proving an appropriate local
Poincare region for them.

\subsection{Two noncircular ways an escape bound could be certified}

\begin{proposition}[A killed-semigroup survival certificate]
\label{prop:f12v73-survival}
Let \(\tau_0\) be the first entrance time
into \(\mathcal G_0\) for the diffusion
with generator \(-L\). If for some
\(t_*>0\) and \(0<a_*<1\),
\[
 \operatorname*{ess\,sup}_{x\notin\mathcal G_0}
                 \mathbb P_x(\tau_0>t_*)\le1-a_* ,
 \tag{77.16}
\]
then
\[
                         \lambda_0\ge
                       -t_*^{-1}\log(1-a_*) .
 \tag{77.17}
\]
\end{proposition}
\begin{proof}
The killed semigroup \(S_t\) is
positive and self-adjoint in
\(L^2(\nu)\) on the complement.
Its row sums are the survival
probabilities. At \(t_*\) these
are at most \(r=1-a_*\).
Symmetry gives the same bound
on column integrals. Cauchy--Schwarz
for its positive kernel gives
\(|S_{t_*}f|^2\le rS_{t_*}|f|^2\);
integration bounds the squared
\(L^2\) norm by \(r^2\|f\|_2^2\).
By the spectral theorem
\(\|S_{t_*}\|_{2\to2}=e^{-t_*\lambda_0}\).
This proves the claim.
\end{proof}

This is a first-entrance event
under the specified diffusion,
not a claim that a large
joint heat-bath refresh can
be performed for free.
The probability and time in
(77.16) may depend on the
patch, exterior and \(\beta\);
the uniformities required
by (77.14) must be proved.

\begin{proposition}[A differential escape certificate]
\label{prop:f12v73-lyapunov}
Suppose a positive function
\(U\in C^2(M)\) satisfies
\[
                          LU/U\ge\lambda>0
                       \quad\hbox{off }\mathcal G_0 .
 \tag{77.18}
\]
Then \(\lambda_0\ge\lambda\).
\end{proposition}
\begin{proof}
For a test \(v\) vanishing on
\(\mathcal G_0\), integration by
parts gives the ground-state
identity
\[
 \mathcal E(v)-\int\frac{LU}{U}v^2\,d\nu
                       =\int U^2|\nabla(v/U)|^2\,d\nu\ge0 .
 \tag{77.19}
\]
The compact ambient manifold
has no boundary, and \(U\) has
a positive minimum there.
Approximation proves the identity
on the killed form domain. Apply
(77.18) and take the infimum.
\end{proof}

One possible trial function is
\(U=e^{\eta W}\), \(0<\eta<1\).
In the actual Haar metric its
drift expression is exactly
\[
 \frac{L(e^{\eta W})}{e^{\eta W}}
       =\eta(1-\eta)|\nabla W|^2-\eta\Delta W .
 \tag{77.20}
\]
All score and second-derivative
terms in \(W=U_{\mathcal S}-\sum_b\log Z_b\)
must be retained. Upper derivative
bounds from V58 do not provide
the lower bound in (77.18).
At an off-chart strict local
minimum with \(\Delta W>0\),
this trial expression is negative.
V29's global energy floor does
not exclude such a critical point.

\subsection{Why an energy floor and a rare bad region do not prove escape}

On the circle take the smooth potential
\[
 V(\theta)=(1-\cos2\theta)+(1-\cos\theta)
          =(1-\cos\theta)\bigl(3+2\cos\theta\bigr).
 \tag{77.21}
\]
It has the coercive floor
\(V\ge1-\cos\theta\) and its
unique zero is at \(\theta=0\).
Nevertheless \(\theta=\pi\) is
a local minimum of energy \(2\),
with \(V''(\pi)=3>0\).
The two intervening saddles
satisfy \(\cos\theta=-1/4\),
have energy \(25/8\), and
separate that secondary well
from the global minimum.
The barrier above the secondary
well is \(9/8\).

For \(\nu_\beta\propto e^{-\beta V}d\theta\),
choose \(\mathcal G_0\) to be a
small fixed neighborhood of zero.
For every small \(\varepsilon>0\),
take a fixed smooth test \(v\)
equal to one near \(\pi\), zero
near zero, and changing only
in small intervals around the
two saddles where
\(V\ge25/8-\varepsilon\).
Its unnormalized energy is
at most \(C_\varepsilon
e^{-\beta(25/8-\varepsilon)}\).
A neighborhood of width
\(\beta^{-1/2}\) at \(\pi\)
has \(V\le2+C/\beta\), and
there \(v=1\) for large \(\beta\).
Its unnormalized squared norm
is therefore at least
\(c\beta^{-1/2}e^{-2\beta}\).
Normalization cancels, proving
\[
                  \lambda_0(\beta)\le
           C_\varepsilon\sqrt\beta\,
                         e^{-\beta(9/8-\varepsilon)} .
 \tag{77.22}
\]
The secondary well has exponentially
small stationary mass, but that
does not prevent slow diffusion
escape. At its minimum the
trial (77.20), with \(W=\beta V\),
equals \(-3\eta\beta\).

This example is not a YM
conditional potential. It
shows why the stated energy
and rarity information alone
does not imply the new hypothesis.
It also has a further caution:
a one-group heat bath refreshes
the whole circle and has gap
one for every \(\beta\).
Thus a poor diffusion escape
bound does not prove slow
single-group heat-bath dynamics.
The diffusion route is sufficient,
not necessary, and its clock
and conversion (77.14) must
remain explicit.

\subsection{Exact checks and remaining work}

The script
\begin{center}\small\ttfamily
scripts/verify\_ym\_f12\_v73\_killed\_escape.py
\end{center}
checks the assembly in a reversible
five-vertex path model. Its
inner good set is vertex zero,
its larger good set is the
first three vertices, and the
cutoff is \((1,1/2,0,0,0)\).
The cutoff energy is supported
in the larger set. Exact
positive-semidefinite tests give
a killed floor \(1/10\), local
good gap \(1\), and cutoff
energy bound \(1/4\); (77.5)
then gives the checked full
gap bound \(1/26\).
This discrete check uses the
verified localization of its
cutoff energy, not an unproved
diffusion product rule for
arbitrary jump kernels.

The script checks the
dimension-independent Euclidean
cutoff inequality in 2500 rational
pairs across dimensions
\(1,2,7,45\). It also verifies
the exact critical energies,
second derivatives, energy-floor
identity and barrier in (77.21).
All checks pass. The general
proofs are above; none of these
checks supplies the actual YM
escape estimate (77.16) or
(77.18).

\begin{decision}[The next quantity is dynamical escape, not probability alone]
The local chart Poincare input
and cutoff are populated for
the actual good-exterior compact
patch. The missing term is now
an explicit killed-generator
lower bound, with a survival
or differential route for
certifying it. The survival
route must control all starting
configurations outside the
inner chart at the stated
conditional exterior. An
independent form estimate can
instead prove (77.4) directly.

Do not replace this by a
small stationary bad-set
probability, the unique
classical section minimum,
or the known full-refresh
probability of a larger repair.
For all-exterior repairs the
local good-region input also
remains to be established.
The next companion audit is
V75 and direction review V80.
\end{decision}

V72 had 1165567 bytes and SHA--256
\begin{center}\scriptsize\ttfamily
B3D80809D8A5F41EADEBC4CE531ED52B31569F6D1F86E5412A85FC391124ABB6.
\end{center}
Its complete body is retained
apart from the title version.
Only active V72 is replaced by
V73. Archives, companions and
prior scripts remain.

\begin{firewall}[Position after V73]
The direct chart-plus-escape
criterion, dimension-independent
cutoff, and conditional heat-bath
conversion are proved. The
actual killed escape bound is
not proved, and the application
does not cover arbitrary bad
exteriors without new input.
No cofinal auxiliary gap,
continuum construction, or
physical YM mass gap has
been obtained. The full
objective remains active.
\end{firewall}
\section{V74: escape by actual heat-bath energy improvement}
\label{sec:f12v74}

\subsection{Use the jump dynamics and expose its own obstruction}

V73 isolated a diffusion escape estimate, but its
circle example also showed why a whole-group
heat bath can escape where diffusion is slow.
We now work directly with the actual rate-one
single-group heat baths. The relevant
zero-temperature obstruction is not every
local differential minimum: it is a
configuration at which no single-group
replacement can lower the classical
conditional energy.

The new criterion uses the actual compact
conditional potential and V70's uniform
classical approximation. A quantitative
single-coordinate energy improvement outside
a target region implies uniform access to
that region in the random-scan chain, and
hence a lower bound for the killed
rate-one generator. The global improvement
hypothesis remains unproved for YM.
We do prove a fractional improvement for
the actual two-link action of V71 on
its specified small commuting box.
That calculation is not extended to
the full configuration space by assumption.

\subsection{Conditional heat baths select near-minimal one-group energies}

Fix a finite repair or patch \(T\), with
\(n=|T|\) group coordinates. Work in a
fixed admissible compact parameter family,
as in V70. Suppress the exterior parameter
\(\xi\) when convenient and write
\[
 W_\beta(v)=\beta E(v)+R_\beta(v),\qquad
 \operatorname{osc}R_\beta\le K+B\log\beta ,
 \quad \beta\ge1.
 \tag{78.1}
\]
Here \(E=\mathcal E_T^\xi\) is the actual
minimized classical action, and (78.1)
is the proved local compact approximation
(74.4), not a replacement measure.
The function \(E\) is uniformly Lipschitz
on this fixed compact product.

For coordinate \(i\), set
\[
 m_i(v_{-i})=\min_{z\in G}E(z,v_{-i}),\qquad
 d_i(v)=E(v)-m_i(v_{-i})\ge0 .
 \tag{78.2}
\]
The actual heat-bath projection \(P_i\)
has conditional density proportional
to \(e^{-W_\beta(z,v_{-i})}\) against
normalized Haar measure in \(z\).

\begin{proposition}[Uniform one-coordinate energy selection]
\label{prop:f12v74-selection}
For every fixed \(\eta>0\) there is
a finite \(C_\eta\), uniform in the
stated parameters, such that
\[
 P_i\!\left(E(z,v_{-i})>m_i(v_{-i})+\eta\right)
               \le C_\eta\beta^B e^{-\beta\eta/2}.
 \tag{78.3}
\]
\end{proposition}
\begin{proof}
Let \(L_E\) bound the one-group
Lipschitz constants of \(E\).
If \(L_E=0\) the event is empty.
Otherwise choose a fixed ball
radius \(r>0\) small enough that
\(L_Er\le\eta/2\). About a
minimizer of \(E(\,\cdot,v_{-i})\)
this ball has Haar measure at
least \(b_\eta>0\), uniformly
over its center on the compact
group, and energy at most
\(m_i+\eta/2\).

The conditional denominator is
therefore at least
\(b_\eta e^{-\beta(m_i+\eta/2)}
e^{-\sup_zR_\beta}\).
The numerator of the event in
(78.3) is at most
\(e^{-\beta(m_i+\eta)}
e^{-\inf_zR_\beta}\).
Their ratio is bounded by
\(b_\eta^{-1}e^K\beta^B
e^{-\beta\eta/2}\).
The restriction of the remainder
to one conditional coordinate
has no larger oscillation than
the full bound in (78.1).
\end{proof}

No minimizer needs to be unique
or nondegenerate. The estimate
uses the exact conditional
normalization. It charges neither
the full action oscillation nor
the coupling to a product reference.

\subsection{A direct killed-generator criterion}

Let \(\mathcal A\) be an open target
region and suppose
\[
 \operatorname{osc}E\le M_E,\qquad
 \max_{i\in T}d_i(v)\ge\delta>0
                       \quad(v\notin\mathcal A).
 \tag{78.4}
\]
Uniformity over the required exterior
parameters is part of this hypothesis.
Define
\[
 k_*=\left\lfloor\frac{2M_E}{\delta}\right\rfloor+1,\qquad
 \varepsilon_\beta=C_{\delta/2}\beta^B e^{-\beta\delta/4}.
 \tag{78.5}
\]
For sufficiently large \(\beta\),
\(\varepsilon_\beta\le1/2\).

The random-scan kernel and the target
rate-one generator are
\[
 P=\frac1n\sum_{i\in T}P_i,\qquad
 H=\sum_{i\in T}(I-P_i)=n(I-P).
 \tag{78.6}
\]
Let \(\mathcal B=\mathcal A^c\),
let \(P_{\mathcal B}=1_{\mathcal B}P1_{\mathcal B}\)
act on \(L^2(\mathcal B,\nu_\beta)\),
and set
\[
 \kappa_{\mathcal B}
  =\inf_{\substack{0\ne f\in L^2(\nu_\beta)\\f=0\ {\rm on}\ \mathcal A}}
                      \frac{\langle f,Hf\rangle}{\|f\|_2^2}.
 \tag{78.7}
\]
This is a killed heat-bath floor,
not V73's killed diffusion floor.

\begin{theorem}[Actual-clock escape from a coordinate-improvement bound]
\label{thm:f12v74-escape}
Under (78.1)--(78.5),
\[
 \begin{aligned}
 a_\beta&=\left(\frac{1-\varepsilon_\beta}{n}\right)^{k_*},\\
 \kappa_{\mathcal B}
       &\ge n\left(1-(1-a_\beta)^{1/k_*}\right)
         \ge\frac{na_\beta}{k_*}.
 \end{aligned}
 \tag{78.8}
\]
In particular, for all sufficiently
large \(\beta\),
\[
                 \kappa_{\mathcal B}\ge
                       \frac{n}{k_*(2n)^{k_*}}>0 .
 \tag{78.9}
\]
The constants may depend strongly
on the fixed patch and \(\delta\);
no uniformity in growing patch
size is claimed without those
dependencies being controlled.
\end{theorem}
\begin{proof}
At a state outside \(\mathcal A\),
choose the first coordinate attaining
\(\max_i d_i\). This is a measurable
choice: \(m_i\) and \(d_i\) are
continuous by compactness and
continuity of \(E\).
The actual random scan chooses
that coordinate with probability
\(1/n\). Conditional on this
choice, (78.3) with \(\eta=\delta/2\)
gives, with probability at least
\(1-\varepsilon_\beta\),
\[
                         E(v_{\rm new})\le E(v)-\delta/2 .
 \tag{78.10}
\]
This is an event in the original
random-scan chain; its scan rule
has not been changed to a greedy
or state-dependent generator.

Consider the event that each
step before the first entrance
into \(\mathcal A\) makes this
successful coordinate choice
and update. Its probability
through \(k_*\) steps is at
least \(a_\beta\); after entrance
there is no further condition.
Remaining outside for all those
successful steps would lower
the energy by \(k_*\delta/2>M_E\),
which is impossible. Therefore
\[
           \sup_{v\in\mathcal B}
                 P_{\mathcal B}^{\,k_*}1(v)\le1-a_\beta .
 \tag{78.11}
\]

The kernel \(P\) is a positive
self-adjoint contraction, since
it is an average of orthogonal
conditional-expectation projections.
Its compression \(P_{\mathcal B}\)
is positive and self-adjoint.
The reversible-kernel Schur bound,
as in V73, implies
\(\|P_{\mathcal B}^{k_*}\|_{2\to2}\le1-a_\beta\).
Consequently
\(\|P_{\mathcal B}\|\le(1-a_\beta)^{1/k_*}\).
On functions supported in
\(\mathcal B\), the generator
is \(n(I_{\mathcal B}-P_{\mathcal B})\),
giving the first floor in (78.8).
Concavity of \(x^{1/k_*}\) gives
the second. Use
\(\varepsilon_\beta\le1/2\)
to obtain (78.9).
\end{proof}

At fixed patch size and fixed
\(\delta\), (78.9) has no
positive exponential cost in
\(\beta\). It could therefore
avoid the exponent obstruction
of the previous oscillation
comparisons. Its hard premise
is (78.4), not an already
populated estimate.

On a compact fixed parameter
family, this premise is
equivalent to excluding all
coordinatewise minimizing
configurations outside the
target, with a uniform margin.
More precisely, if the compact
set of exterior parameters
times \(\mathcal A^c\) contains
no point with \(d_i=0\) for
every \(i\), continuity gives
some \(\delta>0\). A unique
global energy minimum does
not establish this exclusion.
For bad frozen exteriors even
a global conditional minimum
may lie outside the section
chart; a different target
region would then be required.

\subsection{A killed floor still needs within-target conductance}

For clarity, an actual heat-bath
version of the assembly can be
stated without any diffusion
comparison. Suppose, independently,
\[
 \langle f,Hf\rangle\ge
           \gamma_{\rm in}\nu_\beta(\mathcal A)
                 \operatorname{Var}_{\nu_\beta(\cdot\mid\mathcal A)}(f),
                 \qquad\gamma_{\rm in}>0.
 \tag{78.12}
\]
Then
\[
 \operatorname{gap}(H)\ge
       \frac{\kappa_{\mathcal B}\gamma_{\rm in}}
            {\kappa_{\mathcal B}+2n+2\gamma_{\rm in}} .
 \tag{78.13}
\]
To prove this, subtract the target-set
mean and split \(h=h_{\mathcal A}+h_{\mathcal B}\)
by indicators. Since \(0\le H\le nI\),
\[
 \begin{aligned}
 \langle h_{\mathcal B},Hh_{\mathcal B}\rangle
       &\le2\langle h,Hh\rangle
                          +2n\|h_{\mathcal A}\|_2^2,\\
 \|h_{\mathcal A}\|_2^2
       &\le\gamma_{\rm in}^{-1}\langle f,Hf\rangle .
 \end{aligned}
 \tag{78.14}
\]
Apply the killed floor to the first
line and use the orthogonal support
split to bound \(\operatorname{Var}f\).
This yields (78.13).

The within-target input is not
automatically V39's boxed gap.
For a product target
\(\mathcal A=\prod_i\mathcal A_i\),
V49's exact conductance normalization
(53.10) gives (78.12) with
\(\gamma_{\rm in}=q_{\min}\gamma_{\rm box}\)
if
\[
 \inf_{\substack{i,\ v_{-i}\in\prod_{j\ne i}\mathcal A_j}}
                 P_i(\mathcal A_i\mid v_{-i})\ge q_{\min}>0
 \tag{78.15}
\]
and the normalized restricted law
has heat-bath gap \(\gamma_{\rm box}\).
The acceptance factor cannot be
discarded, and its required
\(\beta\)-dependence is not
proved here. Equations
(78.4) and (78.12) are separate
inputs for a full repair gap.

\subsection{A positive actual calculation on the V71 family}

The large oscillation rectangle
in V71 does not create a
coordinate trap near the section.
This can be checked in its
actual action, rather than
assumed from an analogy.
Retain the cleared-pair geometry,
unit-Hilbert--Schmidt generator
and exact action
\[
 S(u,v)=5s_X(u)+5s_X(v)+4s_X(u+v)
                         +s_X(2u+v)+s_X(u+2v).
 \tag{78.16}
\]
Choose \(0<a\le\min\{t_0,1/6\}\)
and \((u,v)\in[-a,a]^2\).
All plaquette angle magnitudes
are at most \(3a\le1/2\).
Since \(\sum\lambda_r^2=1\),
\[
              7/8\le s_X''(\theta)\le1
                            \quad(|\theta|\le1/2).
 \tag{78.17}
\]
Summing the actual coefficient
outer products gives
\[
 H_0=\begin{pmatrix}14&8\\8&14\end{pmatrix},\qquad
             \frac78H_0\preceq S''\preceq H_0 .
 \tag{78.18}
\]
In particular
\[
 S''\succeq\frac{21}{4}I,\qquad
       \frac{49}{4}\le S_{uu},S_{vv}\le14,\qquad
                         |S_{uv}|\le8 .
 \tag{78.19}
\]

\begin{proposition}[Actual small-box coordinate improvement]
\label{prop:f12v74-pair}
Each one-variable minimum of \(S\)
in this box is interior. Its
minimizer satisfies
\[
          |u_*(v)|\le\frac{32}{49}|v|,\qquad
          |v_*(u)|\le\frac{32}{49}|u|.
 \tag{78.20}
\]
Moreover
\[
 \max\left\{S(u,v)-\min_{z\in[-a,a]}S(z,v),\
            S(u,v)-\min_{z\in[-a,a]}S(u,z)\right\}
                                 \ge\frac3{16}S(u,v).
 \tag{78.21}
\]
\end{proposition}
\begin{proof}
One has \(S_u(0,0)=0\) and
\(|S_u(0,v)|\le8|v|\).
The lower bound \(S_{uu}\ge49/4\)
shows opposite derivative signs
at the two endpoints of the
\(u\)-interval, and bounds the
unique root by \((32/49)|v|\).
The other coordinate is symmetric.

The gradient step
\(u'=u-S_u(u,v)/14\) lies between
\(u\) and \(u_*(v)\), by the
one-variable mean value theorem
and \(0<S_{uu}\le14\). It
therefore remains inside the
box and decreases \(S\) by
at least \(S_u(u,v)^2/28\).
The analogous statement holds
in the other coordinate.
Strong convexity with
\(m=21/4\), together with
\(S(0,0)=0=\nabla S(0,0)\),
implies
\(\|\nabla S\|^2\ge2mS\).
This follows by applying the
strong-convexity inequality
at \((u,v)\) with comparison
point zero and completing the
square in the gradient.
The larger squared partial
derivative is at least
\(\|\nabla S\|^2/2\ge mS\).
The better of the two steps
thus decreases the energy
by at least \(mS/28=3S/16\).
The exact coordinate minimum
does at least as well.
\end{proof}

For a full-group heat bath at
one of these actual coordinates,
minimization over \(G\) can do
no worse than these one-parameter
candidates. Under the V71
clearance condition, the
classical energy difference
of the actual repair equals
the displayed action difference.
Thus (78.21) genuinely supplies
an improvement for states in
this actual small family.

It does not verify (78.4) on
the full compact state space.
A generic heat-bath update
does not remain on this
one-parameter commuting
family, so the calculation
cannot be iterated as if
that family were invariant.
It proves that the V71
oscillation witness is not
itself a coordinate trap,
not that no other trap exists.

\subsection{A rare coordinate trap is a genuine jump-dynamics warning}

On two bits take the energy
\[
 E(0,0)=0,\quad E(1,0)=E(0,1)=3,\quad E(1,1)=2 .
 \tag{78.22}
\]
It has a unique global minimum
and the floor \(E(x,y)\ge x+y\).
But \((1,1)\) is minimizing
in each coordinate separately,
so \(d_1=d_2=0\) there.
For the actual Gibbs heat baths
at inverse temperature \(\beta\),
its probability is at most
\(e^{-2\beta}\), while its
total rate of leaving in the
rate-one two-coordinate dynamics
is
\[
                         \frac{2e^{-\beta}}{1+e^{-\beta}} .
 \tag{78.23}
\]
Testing the full Poincare
inequality with the indicator
of \((1,1)\) gives
\[
 \operatorname{gap}(H)
 \le
 \frac{2e^{-\beta}}
      {(1+e^{-\beta})(1-\nu_\beta(1,1))}
                         \le4e^{-\beta}\quad(\beta\ge0).
 \tag{78.24}
\]
Thus small stationary mass
does not imply fast heat-bath
escape either. Unlike the
one-circle diffusion example
in V73, this example is
slow for the actual coordinate
jump dynamics. It is not
a YM conditional law, and
does not establish the
existence of a YM coordinate
trap.

\subsection{Checks and the next acquisition}

The exact script
\begin{center}\small\ttfamily
scripts/verify\_ym\_f12\_v74\_heat\_bath\_descent.py
\end{center}
checks an eight-state product
law with \(E=\sum_{i=1}^3x_i\),
target \((0,0,0)\), and
\(\delta=1\). The criterion
uses \(k_*=7\), access
\(1/279936\), and rate-one
killed floor \(1/653184\).
At four rational Gibbs weights,
the script checks the killed
seven-step row bound and the
generator lower bound by
exact weighted positive-semidefinite
elimination.

It also verifies the rare
coordinate trap and Rayleigh
quotients in (78.22)--(78.24).
For the actual V71 plaquette
coefficients it checks \(H_0\),
the curvature and diagonal
bounds, the minimizer ratio
\(32/49\), and the descent
fraction \(3/16\).
All checks pass. The general
escape theorem is proved
above; the finite tests
do not supply the missing
global YM condition.

\begin{decision}[Classify actual coordinatewise minimizing configurations]
The next upstream target is
now concrete: determine
whether the actual classical
conditional energy has
coordinatewise minimizing
configurations outside an
appropriate target region,
uniformly over the required
exteriors. The compact
approximation and the
heat-bath normalization are
already controlled.

The positive V71-family
calculation covers one
small actual configuration
family, not the complement
of the chart. Any extension
must keep the full group
minimization and the
integrated block minima.
The within-target conductance
input remains separate.
The next companion audit
is V75 and direction review
V80.
\end{decision}

V73 had 1181448 bytes and SHA--256
\begin{center}\scriptsize\ttfamily
40F91CEE3C8960E5F0FE2BB05ECE69D70DE536C67D6F354C9D765D5D32A6B5F0.
\end{center}
Its complete body is retained
apart from the title version.
Only active V73 is replaced
by V74. Archives, companions
and previous scripts remain.

\begin{firewall}[Position after V74]
The actual-clock escape
criterion and the explicit
small-family YM energy
improvement are proved.
The global coordinatewise
minimum exclusion and the
required within-target
conductance are not proved.
No full repair gap, cofinal
auxiliary gap, continuum
construction, or physical
YM mass gap has been
established. The full
objective remains active.
\end{firewall}
\section{V75: actual low-energy coordinate descent after the companion audit}
\label{sec:f12v75}

\subsection{Audit, inherited inputs, and the limited new conclusion}

The scheduled five-version review checked the active companion
endings and complete file digests. They remain SM--Einstein V72
(607372 bytes), NS stability V26 (278283 bytes), and parallel
YM V5 (54174 bytes). Their SHA--256 digests are, respectively,
\begin{center}\scriptsize\ttfamily
F44F9745D43379E1672C5550411B58E97195A4C9278592D221DF5D3F644DC1F5\\
82C887F8C2C69E4F28FAD7C3DC8DE5B9099CB5A4BF431EBA1FFF3E91935978AD\\
306F1BB84CFA8A8D47EA569EC17E9545F9867C8D32B548EC9D9C444B4F3C0512.
\end{center}
These match V70. SM's fixed-mode many-copy limit, NS's
pointwise kernel norms, and parallel YM's local rather than
extensive charge-tail targets supply no new compact repair
gap estimate. Their internal directions are documentary
material, not instructions controlling this continuation.
The next companion review and ten-version direction review
are both V80.

V74 left the classification of coordinatewise minimizing
configurations open. Its explicit calculation concerned
only the two-link commuting family. Here we extend the
descent conclusion to every group coordinate in an actual
finite interface patch, on a specified low-energy sublevel.
The retained field is \(C=I\), and all exterior interface
variables are fixed at their section values \(I\).
This restriction is essential.

The argument combines, without reproving their lattice
construction, V28's exact compact marginal Hessian,
V29's minimized-action floor, and V70's localized compact
Laplace estimate. All integrated blocks remain compact.
No cleared corridor as in V71--V72 is required.
The new step is a block-coordinate descent inequality
and its application to a sublevel that the energy floor
places strictly inside the chart.
This is not a claim of novelty for general convex
coordinate descent.

The full objective remains the continuum construction
and physical Hamiltonian gap, as distinguished in the
\href{https://www.claymath.org/millennium/yang-mills-the-maths-gap/}
{Clay Mathematics Institute problem statement}
(consulted 13 September 2026). Neither a finite-patch
descent estimate nor the auxiliary killed generator
below supplies that objective.

\subsection{The actual conditional energy and the flat section}

Fix a nonempty finite patch \(T\subset\mathcal S\),
\(n=|T|\), and the packing of V28. The ambient periodic
volume is finite and sufficiently large for the fixed
templates under discussion. At \(C=I\) and section
exterior write
\[
 E(v_T)=\mathcal E_I(v_T,I_{\mathcal S\setminus T}),
 \qquad E(I_T)=0,\qquad
 D_T(v)=\sum_{i\in T}\|v_i-I\|_{\rm HS}^2 .
 \tag{79.1}
\]
Here \(\mathcal E_I\) is the full compact minimum
in (33.9). Equivalently \(E\) is the localized
energy in (74.3), with its constant chosen to vanish
at the section. To see the equality, minimize the
independent integrated blocks in V28's factorization.
Blocks whose input sets miss \(T\), and direct terms
independent of \(T\), have their section data and
zero minimum action. They contribute no classical
constant at \(C=I\). Blocks touching \(T\) retain
their full compact minima, including any off-chart
minimizers.

Thus (33.11) and (74.4) give
\[
 E(v)\ge\frac{D_T(v)}{2A}\ge0,\qquad
 W_{T,\beta}(v)=\beta E(v)+R_{T,\beta}(v),\qquad
 \operatorname{osc}R_{T,\beta}\le K_T+B_T\log\beta .
 \tag{79.2}
\]
The law is the actual conditional
\(\nu_\beta(dv)\propto e^{-W_{T,\beta}(v)}\,dv\),
where \(dv\) is product normalized Haar measure.
Constants in its potential are irrelevant.
In particular the section is its classical energy's
unique zero on the entire compact patch.
Uniqueness of this global minimum alone does not
exclude other coordinatewise minima; V74 already
provided the counterexample explaining that distinction.

\subsection{Passing the chart curvature to the minimized energy}

Choose an open convex product box
\(\mathcal C^+\subset\mathfrak g^T\), centered at zero,
strictly inside the fixed interface coordinate box
of V28, and a smaller such box
\(\overline{\mathcal C}\subset\mathcal C^+\).
Use a fixed orthonormal Lie basis and
\(v_i=\exp(x_i)\). All exterior coordinates remain zero.
The radii can be chosen per group coordinate,
independently of \(n\) and ambient volume.
Set \(F(x)=E(\exp x)\).

\begin{proposition}[Classical chart regularity with retained compact blocks]
\label{prop:f12v75-classical-convexity}
On \(\mathcal C\), \(F\) is \(C^{1,1}\),
\(F(0)=0\), \(\nabla F(0)=0\), and one can take
\(\mu=m/2>0\), \(L=K\ge\mu\), from (32.16), so that
\[
 \begin{split}
 F(y)&\ge F(x)+\langle\nabla F(x),y-x\rangle
                                      +\frac{\mu}{2}\|y-x\|^2,\\
 F(y)&\le F(x)+\langle\nabla F(x),y-x\rangle
                                      +\frac{L}{2}\|y-x\|^2,\\
 \|\nabla F(y)-\nabla F(x)\|&\le L\|y-x\|.
 \end{split}
 \tag{79.3}
\]
The curvature constants are independent of patch
cardinality and ambient volume under the inherited
fixed-packing hypotheses. No differentiability of \(E\)
outside this chart is asserted.
\end{proposition}
\begin{proof}
The Lebesgue-coordinate potential of the conditional
is the restriction of \(\Phi^{\rm c}\), up to a
constant. Its Hessian is a principal restriction
of the full interface Hessian. Hence its quotient
by \(\beta\) has Hessian between \(\mu I\) and
\(LI\) on \(\mathcal C^+\).

More explicitly, with \(j_G\) the local Haar density,
put
\[
 F_\beta(x)=\beta^{-1}\left[
 W_{T,\beta}(\exp x)-W_{T,\beta}(I)
 -\sum_{i\in T}\log\frac{j_G(x_i)}{j_G(0)}\right].
 \tag{79.4}
\]
Its Hessian has those bounds. By (79.2) and the
bounded smooth logarithm of \(j_G\) on the fixed
chart, \(F_\beta\to F\) uniformly on every compact
subset of \(\mathcal C^+\), for this fixed finite
patch. The finite-patch error is
\(O((1+\log\beta)/\beta)\), not an asserted
dimension-independent uniform convergence rate.

Here is a justification for passing derivatives,
rather than merely passing convexity. On any
compact box inside \(\mathcal C^+\), the gradients
\(\nabla F_\beta\) are \(L\)-Lipschitz.
They are bounded there: take small coordinate
finite differences in a slightly larger box and
use uniform bounds on \(F_\beta\) together with
the Hessian upper bound. In fixed finite dimension,
subsequences of these gradients therefore converge
uniformly on smaller compact boxes. Passing to the
limit in the line-integral identity for
\(F_\beta(y)-F_\beta(x)\) identifies the limit
as \(\nabla F\). This identification is unique,
and yields an \(L\)-Lipschitz continuous gradient.
The lower and upper Taylor inequalities pass
to the limit and prove (79.3).
Finally \(F\ge0=F(0)\), so its gradient at zero
vanishes. No selection or differentiation of
a minimizing integrated-block branch was used.
\end{proof}

This supplies local regularity not asserted globally
in V29. It does not say that the global compact
minimum has a unique smooth minimizing branch.

\subsection{An inward block update with a quantitative improvement}

Treat \(x_i\in\mathfrak g\) as one block:
the count \(n\) is the number of group updates,
not \(n\dim G\). Set \(s=\mu/L\in(0,1]\).
Let \(x^{[i]}\) agree with \(x\) except
\(x_i^{[i]}=(1-s)x_i\). These points and their
joining segments remain in the product chart.

\begin{theorem}[Actual full-group improvement throughout the chart]
\label{thm:f12v75-chart-descent}
For \(v=\exp x\) with \(x\in\mathcal C\),
define the actual compact improvement defects
as in (78.2):
\[
 d_i(v)=E(v)-\min_{z\in G}E(z,v_{-i}).
 \tag{79.5}
\]
Then
\[
 \begin{split}
 \sum_{i\in T}\bigl(F(x)-F(x^{[i]})\bigr)
                         &\ge\frac{\mu}{L}F(x),\\
 \max_{i\in T}d_i(v)&\ge\frac{\mu}{nL}E(v).
 \end{split}
 \tag{79.6}
\]
\end{theorem}
\begin{proof}
The upper inequality in (79.3), applied to the
single-block displacement \(-s x_i\), gives
\[
 F(x)-F(x^{[i]})
 \ge s\langle\nabla_iF(x),x_i\rangle
                              -\frac{Ls^2}{2}\|x_i\|^2 .
 \tag{79.7}
\]
The lower inequality in (79.3), compared with
the origin, gives
\(\langle\nabla F(x),x\rangle\ge
F(x)+(\mu/2)\|x\|^2\).
Summing (79.7), and using \(Ls=\mu\), yields
\[
 \sum_i(F(x)-F(x^{[i]}))
 \ge sF(x)+\frac{s\mu-Ls^2}{2}\|x\|^2=sF(x).
 \tag{79.8}
\]
At least one of the \(n\) summands is at least
\(sF(x)/n\). Its candidate
\(\exp((1-s)x_i)\) is among the full compact
replacements in (79.5). Minimization over the
entire group can only improve that candidate,
proving the second assertion.
\end{proof}

Individual inward updates need not all lower
the energy. The proof selects at least one
productive update and uses no absolute diagonal
dominance. Noncommuting Lie coordinates and
all interactions inherited through compact
block integration are included in \(F\).
This replaces V74's two-dimensional-family
restriction, but not its exterior restriction.

\subsection{A compact sublevel containing no nonsection coordinate trap}

Choose \(\tau_{\rm ch}>0\) sufficiently small that
the closed Hilbert--Schmidt ball
\(\{g:\|g-I\|_{\rm HS}\le\tau_{\rm ch}\}\)
is contained in the one-group chart used for
\(\mathcal C\). This follows from the local
normal-coordinate neighborhood of the identity;
shrink it if necessary so the inclusion is
strict. Define
\[
 e_*=\frac{\tau_{\rm ch}^2}{4A}>0,\qquad
 \mathcal K=\{v:E(v)\le e_*\},\qquad
 c_n=\frac{\mu}{nL}.
 \tag{79.9}
\]
The threshold is fixed and independent of patch
cardinality; it is a total-energy threshold,
not an energy density.

\begin{corollary}[All possible nonsection traps lie above a positive energy]
\label{cor:f12v75-low-energy-traps}
The compact sublevel \(\mathcal K\) is contained
strictly in the product chart. On it
\[
                  \max_i d_i(v)\ge c_n E(v).
 \tag{79.10}
\]
The section is the only configuration in
\(\mathcal K\) minimizing \(E\) separately in
every full group coordinate. Consequently any
nonsection coordinatewise minimum, if one
exists on this actual flat-exterior compact
patch, has \(E>e_*\).
\end{corollary}
\begin{proof}
For \(v\in\mathcal K\), (79.2) gives
\[
 \sum_i\|v_i-I\|_{\rm HS}^2
 \le2Ae_*=\tau_{\rm ch}^2/2.
 \tag{79.11}
\]
Each coordinate is therefore strictly inside
the stated chart. Apply (79.6). At a
coordinatewise minimum every \(d_i=0\), so
(79.10) implies \(E=0\); (79.2) then forces
every coordinate to be \(I\). Conversely the
zero-energy section minimizes every coordinate,
because the energy is globally nonnegative.
\end{proof}

This is an actual finite-patch classification
on a proved sublevel, rather than an assumed
global basin. It makes no statement excluding
stationary points above \(e_*\).

\subsection{Access to a smaller sublevel in the original heat-bath clock}

Fix \(0<\epsilon<e_*\), and set
\[
 \mathcal A_\epsilon=\{E<\epsilon\},\quad
 \mathcal B_\epsilon=\{\epsilon\le E\le e_*\},\quad
 \delta=c_n\epsilon,\quad
 k=\left\lfloor\frac{2e_*}{\delta}\right\rfloor+1 .
 \tag{79.12}
\]
Let \(P=n^{-1}\sum_iP_i\) be the original
random-scan kernel with its full compact
conditional replacements, and \(H=n(I-P)\).
No rejection or reflecting boundary is added.
Let \(\tau_{\mathcal A_\epsilon}\) and
\(\tau_{\mathcal K^c}\) be the first entrance
times, including time zero.
With the constant from (78.3), put
\[
 \eta_\beta=C_{\delta/2}\beta^{B_T}e^{-\beta\delta/4},
 \qquad a_\beta=\left(\frac{1-\eta_\beta}{n}\right)^k .
 \tag{79.13}
\]
Below \(\beta\) is large enough that
\(\eta_\beta\le1/2\).

\begin{theorem}[Low-energy access and the original killed energy-band floor]
\label{thm:f12v75-band-access}
For the actual patch at section exterior,
\[
 \inf_{v\in\mathcal K}
 \mathbb P_v\{\tau_{\mathcal A_\epsilon}\le k,\
          \tau_{\mathcal A_\epsilon}<\tau_{\mathcal K^c}\}
             \ge a_\beta\ge(2n)^{-k}.
 \tag{79.14}
\]
For the same rate-one generator \(H\),
\[
 \inf_{\substack{0\ne f\in L^2(\nu_\beta)\\
             f=0\ {\rm off}\ \mathcal B_\epsilon}}
 \frac{\langle f,Hf\rangle}{\|f\|_2^2}
 \ge n\bigl(1-(1-a_\beta)^{1/k}\bigr)
 \ge\frac{na_\beta}{k}
 \ge\frac{n}{k(2n)^k}.
 \tag{79.15}
\]
The left side kills both on entering the inner
target and on leaving the upper sublevel.
It is not the killed floor on the entire
complement of the inner target.
\end{theorem}
\begin{proof}
At a state in \(\mathcal B_\epsilon\),
(79.10) gives a coordinate with improvement
at least \(\delta\). Choose the first maximizing
coordinate; continuity of the minimum values
makes the choice measurable. The actual random
scan chooses it with probability \(1/n\).
Conditional on that choice, V74's selection
estimate with tolerance \(\delta/2\) gives,
with probability at least \(1-\eta_\beta\),
\[
                   E(v_{\rm new})\le E(v)-\delta/2.
 \tag{79.16}
\]
Such a successful full-group replacement
stays in \(\mathcal K\), by its energy,
even though it was not restricted to a
coordinate line or chart. Until the inner
target is reached, \(k\) successive successful
steps have probability at least \(a_\beta\).
If all \(k\) remained outside that target,
the energy would decrease by
\(k\delta/2>e_*\), contradicting \(E\ge0\).
After entrance impose no further condition.
For initial states already in the target
the desired event has probability one.
This proves (79.14) with no change to the
original scan rule.

Write \(P_{\mathcal B}=1_{\mathcal B_\epsilon}
P1_{\mathcal B_\epsilon}\).
The probability of surviving in the band for
\(k\) steps is at most \(1-a_\beta\).
Hence \(P_{\mathcal B}^k1\le1-a_\beta\).
Reversibility gives the same column bound,
and the Schur bound on the restricted measure
gives \(\|P_{\mathcal B}^k\|_{2\to2}\le1-a_\beta\).
Each \(P_i\) is an orthogonal projection,
so \(P\) and its band compression are
positive self-adjoint contractions.
The spectral theorem yields
\(\|P_{\mathcal B}\|\le(1-a_\beta)^{1/k}\).
For functions supported on the band the
quadratic form of \(H\) is exactly that
of \(n(I-P_{\mathcal B})\).
Finally concavity gives
\((1-a)^{1/k}\le1-a/k\) for \(0\le a\le1\).
These observations prove (79.15).
\end{proof}

The constants depend on the fixed patch,
\(\epsilon\), the chart, and the templates.
For example \(k\) grows with \(n\).
There is no volume-uniform conclusion when
the patch itself grows without further
control. Crucially, (79.15) cannot be
substituted for V74's required global
complement floor: the domain
\(\{E>e_*\}\) was killed, not analyzed.

\subsection{The unresolved region is rare, but its escape is still unknown}

\begin{proposition}[Fixed-patch high-energy mass bound]
\label{prop:f12v75-high-energy-mass}
Let \(N=n\dim G\). There is a fixed \(C_T<\infty\)
such that, for all sufficiently large \(\beta\),
\[
       \nu_\beta\{E>e_*\}
       \le C_T\beta^{B_T+N/2}e^{-\beta e_*}.
 \tag{79.17}
\]
\end{proposition}
\begin{proof}
The upper inequality in (79.3) at zero gives
\(F(x)\le(L/2)\|x\|^2\) locally. Integrate
the denominator in a Euclidean ball of
radius \(\beta^{-1/2}\) about zero, which
lies in the chart for sufficiently large
\(\beta\). Its product Haar volume is
at least \(b_T\beta^{-N/2}\), and
\(\beta E\le L/2\) there. Therefore
\[
 \int e^{-W_{T,\beta}}\,dv
 \ge b_T\beta^{-N/2}e^{-L/2}e^{-\sup R_{T,\beta}} .
 \tag{79.18}
\]
Normalized total Haar volume is one, so
the numerator on \(\{E>e_*\}\) is at most
\(e^{-\beta e_*}e^{-\inf R_{T,\beta}}\).
Divide and apply (79.2), absorbing \(e^{K_T+L/2}/b_T\)
into \(C_T\).
\end{proof}

The power \(N/2\) and the constant are retained.
This is a fixed-patch statement. It does not
show that every configuration of a growing
patch is in \(\mathcal K\) with useful
probability at a specified coupling.
More importantly, even exponentially small
mass of the remaining region cannot replace
an escape estimate there: the rare
coordinate trap of V74 is precisely a
counterexample to that inference.

\subsection{Exact checks, next acquisition, and version accounting}

The script
\begin{center}\small\ttfamily
scripts/verify\_ym\_f12\_v75\_low\_energy\_descent.py
\end{center}
checks (79.6) in 125 rational states for the
three-dimensional quadratic with Hessian
\(I+vv^\top\), \(v=(1,2,3)\), whose bounds
are \(\mu=1,L=15\). It also checks 729 states
in six dimensions, \(v=(1,2,3,4,5,6)\),
with three two-dimensional update blocks
and \(\mu=1,L=92\). The guaranteed
best-block fractions are \(1/45\) and
\(1/276\), respectively. These Hessians
are not absolutely diagonally dominant.

For an eight-state product Gibbs law with
\(E=\sum_{i=1}^3x_i\), the script takes
\(\mathcal K=\{E\le2\}\), inner target
\(\{E=0\}\), and \(\delta=1\).
Using the unchanged full conditional
kernels, it checks entrance to the core
before exit from \(\mathcal K\) in five
steps with lower bound \(1/7776\), and
the six-state band killed floor \(1/12960\).
The latter is verified by exact weighted
positive-semidefinite elimination at
\(\beta/\log2=1,2,4,8\).
All these auxiliary checks pass.
They check constants and finite examples,
not the inherited Yang--Mills construction
or a continuum spectral theorem.

\begin{decision}[Move beyond the established low-energy region]
The actual flat-exterior low-energy
coordinatewise-minimum exclusion is now
proved. Repeating chart convexity or the
two-link calculation cannot settle what
remains. The next acquisition is control
of full-group improvement or exit barriers
on \(\{E>e_*\}\), together with the
dependence on nonsection exteriors.
One must also retain V74's independent
within-target conductance requirement.

The bounds here isolate an unproved
region; they do not certify that a
descent path reaches this region from
every starting state. If the inherited
global energy floor and chart Hessian
alone give no high-energy information,
the next upstream analysis must use
additional structure of the actual
compact action and its minimizing blocks,
or identify precisely what that
structure fails to control.
\end{decision}

V74 had 1197176 bytes and SHA--256
\begin{center}\scriptsize\ttfamily
00416158C935D3C5F0B0F5AA9CACE2E54C5F940B4B69D38C98EE49277C047211.
\end{center}
Its entire body is preserved except
for the title version. Only active
V74 is replaced by V75; archives,
companions, and all earlier scripts
remain.

\begin{firewall}[Position after V75]
Under the inherited exact compact
interface hypotheses, low-energy
coordinate descent at the flat
section exterior, the stated
original-clock band killed floor,
and the fixed-patch high-energy
mass estimate are proved.
High-energy coordinate traps have
not been excluded, arbitrary
exteriors have not been handled,
and the required full repair gap
has not been established.
No cofinal auxiliary gap,
continuum construction, or physical
YM mass gap follows. The full
objective remains active.
\end{firewall}
\section{V76: an actual compact coordinatewise minimum that is a collective saddle}
\label{sec:f12v76}

\subsection{Outcome and the direction correction}

V75 established low-energy coordinate improvement at the flat
section exterior. An unrestricted extension based on excluding
every nonsection coordinatewise minimum is too strong for the
actual smooth completed fibre. We now construct an explicit
\(SU(2)\) configuration at \(C=I\) with the following properties:
every individual free-group replacement strictly increases its
classical action unless unchanged, but a simultaneous rotation
of a bounded set of free groups strictly decreases that action.
The point is a saddle, not a metastable local well.

This is an exact Wilson-action calculation for V18's smooth
completion, not the two-bit illustrative counterexample of V74.
The first theorem concerns the unintegrated completed fibre.
For the compactly integrated interface, a planar conditional
version transfers under an explicit clearance condition below.
No assertion is made that every prescribed V50 repair template
has this saddle after its integrated variables have been
minimized. That distinction is necessary.

The legacy f6 V27 centre-sheet erasure and probability estimates,
and f6 V41's identity-defect contour discussion, were checked for
overlap. They are not repeated here: neither calculation proves
the single-coordinate minimizing property, the pivot-completed
saddle, or the zero coordinate-path barrier established below.
The scheduled companion audit was completed at V75;
the next companion and direction reviews remain V80.

Throughout this section take \(G=SU(2)\) in its defining
representation. The result is not asserted for unrestricted
updates in \(SU(3)\) merely by embedding \(SU(2)\).
The physical gap required by the
\href{https://www.claymath.org/wp-content/uploads/2022/06/yangmills.pdf}
{Jaffe--Witten problem description}
remains a different, unproved continuum assertion.

\subsection{The sheet and its exact completed coarse field}

Use the fine four-torus \((\mathbb Z/(2M)\mathbb Z)^4\),
\(M\ge2\), and zero-based direction labels. Write
\(x_\perp=(x_0,x_1,x_3)\) and choose direction \(2\).
Define the nonpivot boundary-link set
\[
 T_0=\{(x,2):x_2=1,\ 
                   x_\perp\notin(2\mathbb Z/(2M)\mathbb Z)^3\}.
 \tag{80.1}
\]
There are \(7M^3\) such links. The \(M^3\) omitted
links at even \(x_\perp\) are exactly the reference
pivots for the coarse edges in this layer.
Put
\[
 W_l^*=\begin{cases}-I,&l\in T_0,\\ I,&l\notin T_0.\end{cases}
 \tag{80.2}
\]
All tree links are identity, and the links in \(T_0\)
are independent free variables in the actual chart.

Retain V18's fixed smooth cutoff, writing
\(\rho=\rho_1\). In addition to the smallness required
by Theorem~\ref{thm:f12v18-pivot}, assume
\[
                         0<\rho\le1/4 .
 \tag{80.3}
\]
This is compatible with that existence choice.
The results apply to any already fixed completion
satisfying (80.3); no cutoff is changed inside an
already specified measure. The cutoff's shape is
otherwise exactly as in (22.1)--(22.2).

\begin{proposition}[Exact sheet fibre and action]
\label{prop:f12v76-sheet-fibre}
The configuration (80.2) has \(\widetilde K(W^*)=I\)
and is the globally solved field \(W(Y^*,I)\).
Exactly \(6M^3\) plaquettes have holonomy \(-I\);
all others have holonomy \(I\). Consequently
\[
                         S_0(W^*)=24M^3 .
 \tag{80.4}
\]
\end{proposition}
\begin{proof}
For a direction-2 coarse edge in the layer,
each of its eight boundary links appears in two
of its sixteen chronological paths, whose other
internal link is identity. The two root paths
are \(I\); the other fourteen products are
\(-I\). Every path product for other coarse
edges is \(I\). Since \(g_\rho(I)=g_\rho(-I)=0\),
(22.14) evaluated at pivot \(B=I\) gives coarse
output \(I\). Global uniqueness of the pivot
solution identifies the actual completed field.
This uses the truncated logarithm, not the
original untruncated averaging formula.

A plaquette can be nontrivial only if it
compares two transverse neighboring direction-2
links in the layer. It is negative precisely
when one link is an omitted pivot and the other
belongs to \(T_0\). Each pivot has six such
neighbors, and no pair of pivots is adjacent.
Each negative plaquette therefore has exactly
one pivot endpoint and is counted once.
There are \(6M^3\) of them, including the
wraparound neighbors. The assumption \(2M\ge4\)
keeps opposite neighbors distinct.
Finally the actual normalization is
\(s_p(U)=\tfrac12\|I-U\|_{\rm HS}^2
=2-\Re\operatorname{Tr}U\), so \(s_p(-I)=4\).
\end{proof}

For a sheet link, let \(r_i\) count its pivot
neighbors in the transverse layer. If precisely
one component of \(x_\perp\) is odd then
\(r_i=2\); if two or three components are odd
then \(r_i=0\). Thus
\[
 h_i:=6-2r_i=
 \begin{cases}
 2,&3M^3\ \hbox{sheet coordinates},\\
 6,&4M^3\ \hbox{sheet coordinates}.
 \end{cases}
 \tag{80.5}
\]
Every physical link outside this direction-2
layer has at most one incident negative
plaquette. Indeed, for a transverse link a
negative plaquette can lie only between levels
\(x_2=1,2\), and for any other direction-2 link
there are none. Its signed plaquette count
is therefore at least four.

\subsection{Full compact minimization in a single free coordinate}

Let
\(\mathcal S(Y)=S_0(W(Y,I))\) be the actual
unintegrated completed-fibre action.
A replacement \(U\) below replaces the physical
nonpivot free link itself; its solved pivots
are recomputed by (22.6), not held fixed.
This concerns the classical action, not the
finite-\(\beta\) potential including the
inverse-pivot Jacobian.

Two uniform elementary bounds will be useful.
Every solved pivot at coarse value \(I\) obeys
\[
 B=\exp[-X(B)],\qquad
 \|\log B\|_{\rm HS}\le7\rho/8,\qquad
 \|B-I\|_{\rm HS}\le7\rho/8 .
 \tag{80.6}
\]
Here the logarithm is in the small injective
ball supplied by (22.8). Changing one pivot
while other links are arbitrary changes the
Wilson action by absolute value at most
\[
 6\sqrt2\,\|B-I\|_{\rm HS}
 \le(21\sqrt2/4)\rho<8\rho .
 \tag{80.7}
\]
Each pivot has six incident plaquettes;
for each, multiplication by other unitary
links preserves the Hilbert--Schmidt norm,
and \(|\operatorname{Tr}A|\le\sqrt2\|A\|_{\rm HS}\).
The same estimate can be summed while
changing two pivots successively.

\begin{theorem}[Strict coordinatewise global minimum on the completed fibre]
\label{thm:f12v76-coordinate-minimum}
For every one of the \(45M^4\) independent
free group coordinates \(i\), and every
\(U\in SU(2)\) different from \(Y_i^*\),
\[
                  \mathcal S(U,Y_{-i}^*)>\mathcal S(Y^*).
 \tag{80.8}
\]
Thus all actual full-group classical
improvement defects at \(Y^*\) vanish.
\end{theorem}
\begin{proof}
\emph{Sheet coordinates.}
First let \(i\in T_0\), whose current value
is \(-I\). Only its own coarse-edge pivot
can change. With all pivots held at \(I\),
the exact single-link action increase is
\[
                         h_i(2+\Re\operatorname{Tr}U).
 \tag{80.9}
\]
If \(g_\rho(U)=0\), every other path product
is \(-I\) or \(U\), so \(B=I\) solves the
pivot equation; uniqueness makes it the
actual solution. Equation (80.9) is then
strictly positive unless \(U=-I\).

If \(g_\rho(U)\ne0\), then
\(\|U-I\|_{\rm HS}<\rho\) and
\(\Re\operatorname{Tr}U\ge2-\rho^2/2\).
Using \(h_i\ge2\) and (80.7), the actual
increase is at least
\[
                 8-\rho^2-8\rho\ge95/16>0 .
 \tag{80.10}
\]
This proves the full-group statement
for every sheet coordinate, including
replacements near the section.

\emph{Identity coordinates away from the cutoff supports.}
Now let \(Y_i^*=I\). A physical internal link
appears only in the chronological paths for
its own coarse edge and possibly the previous
coarse edge in the same direction; a nonpivot
boundary link appears only for its own edge.
There are at most two affected pivots.
At pivot value \(I\), each nonroot path
product after replacement is \(I,-I,U\), or
\(-U\). If both \(g_\rho(U)\) and
\(g_\rho(-U)\) vanish, all pivots remain
\(I\) by uniqueness. The frozen-pivot
increase is exactly
\[
                   h_i^{\rm phys}(2-\Re\operatorname{Tr}U),
                   \qquad h_i^{\rm phys}\ge4,
 \tag{80.11}
\]
which is positive for \(U\ne I\).

\emph{Identity coordinates with replacement near \(-I\).}
If \(g_\rho(-U)\ne0\), then
\(\Re\operatorname{Tr}U\le-2+\rho^2/2\).
The two support balls are disjoint under
(80.3). Charging both possibly changed
pivots by (80.7), the increase is at least
\[
                        16-2\rho^2-16\rho\ge95/8>0 .
 \tag{80.12}
\]

\emph{Identity coordinates with replacement near \(I\).}
It remains to treat \(g_\rho(U)\ne0\).
Write \(U=\exp(uX)\), where
\(\|X\|_{\rm HS}=1\) and \(|u|<\rho\).
Every affected layer direction-2 pivot
remains \(I\): at \(B=I\) its nonroot
products are \(-I\) or \(-U\), both outside
the cutoff support. All other affected
pivots lie outside that layer.
By conjugation equivariance and uniqueness,
they commute with \(U\). In \(SU(2)\) the
centralizer of this noncentral \(U\) is its
one-dimensional torus, so those pivots
are \(\exp(b_jX)\), with \(|b_j|\le7\rho/8\).
The case \(U=I\) is already the unchanged
coordinate and needs no strict inequality.

Consider the action as a function of the
at most three independent physical angles
\(z=(u,b_1,b_2)\), keeping all other links
at the central background. These physical
links were all \(I\). At zero, its Hessian
has diagonal entries at least four.
For an internal free link the possible
pivots have distinct longitudinal base
coordinates \(x_\mu-1,x_\mu+1\), different
also from the free link's \(x_\mu\).
No plaquette contains two of these parallel
links, so there are no mixed entries.
For a boundary free link there is only
one possible pivot, and the pair shares
at most one plaquette. The mixed entry
therefore has absolute value at most one.
In both cases the physical-angle Hessian
at zero is bounded below by \(3I\).

On the angle box \(|z_j|\le\rho\), each
affected plaquette is
\(\sigma\exp((a\cdot z)X)\), with
\(\sigma=\pm1\) and at most two nonzero
entries of \(a\), each equal to \(1\)
or \(-1\). A unit-Hilbert--Schmidt
\(SU(2)\) generator has eigenvalues
\(\pm i/\sqrt2\), so its Hessian
contribution is
\(\sigma\cos((a\cdot z)/\sqrt2)aa^\top\).
There are at most eighteen affected
plaquettes,
\(\|a\|^2\le2\), and
\[
 |\cos((a\cdot z)/\sqrt2)-1|\le\rho^2 .
 \tag{80.13}
\]
The Hessian changes in operator norm
by at most \(36\rho^2\). It is therefore
at least \(3I/4\) throughout that box.
Its gradient at zero vanishes because
all plaquette holonomies there are
central. Integration on the straight
segment from zero to the actual vector
\((u,b_1,b_2)\) gives increase at least
\(3\|z\|^2/8>0\) when \(u\ne0\).
This proves the last case and (80.8).
\end{proof}

The near-\(I\) physical-angle argument
does not assume that the pivot is a
linear function of \(u\), nor that
the cutoff derivative is nonnegative.
The independent angle box is used only
to bound the action at the actual
nonlinear pivot solution.
All first derivatives of \(\mathcal S\)
at \(Y^*\) vanish, since all underlying
plaquette holonomies are central.

\subsection{A negative collective direction with bounded support}

Identify the layer with the transverse
three-torus. Call its even sites the
pivot sites and the other sites free.
For small \(t_i\) set the free sheet
links to \(-\exp(t_iX)\), leaving every
other free coordinate unchanged.
For \(|t_i|\le1/4\), all their path
products stay outside the cutoff
support near \(I\), so all pivots
remain exactly \(I\).

In one Lie direction, the Hessian
quadratic form on these sheet variables
is
\[
 \mathcal Q(t)=
       \sum_{\{i,j\}\ {\rm free\ neighbors}}(t_i-t_j)^2
                      -\sum_i r_it_i^2 .
 \tag{80.14}
\]
This follows by expanding the actual
plaquette action: free--free plaquettes
have sign \(+1\), while a free--pivot
plaquette has sign \(-1\).
Its diagonal is \(6-2r_i>0\).
For the constant vector on the entire
sheet,
\[
 \mathcal Q(\mathbf1)=-6M^3,\qquad
 \frac{\mathcal Q(\mathbf1)}{\|\mathbf1\|^2}=-6/7 .
 \tag{80.15}
\]
In fact the exact collective action is
\[
          S_0(t)=12M^3(1+\cos(t/\sqrt2)).
 \tag{80.16}
\]
Thus \(Y^*\) is not a local minimum,
despite (80.8).

\begin{proposition}[A fixed-size collective descent]
\label{prop:f12v76-local-collective}
Let \(1\le q<M\). Rotate together the
\(7q^3\) sheet coordinates belonging
to a transverse cube of \(q^3\)
coarse cells, and leave the rest
of the central sheet unchanged.
For \(|t|\le1/4\), the exact change is
\[
 S_0(t)-S_0(0)
       =6q^2(3-q)\,2(1-\cos(t/\sqrt2)).
 \tag{80.17}
\]
For \(q=4\), this is strictly negative
at every nonzero allowed \(t\).
It uses 448 independent free coordinates,
has Hessian quadratic value \(-96\),
and Hessian Rayleigh quotient \(-3/14\),
independently of the ambient size
when \(M>4\).
\end{proposition}
\begin{proof}
On each of the six coarse faces of the
selected cube, three of the four fine
transverse parity pairs are free--free;
the fourth has a pivot endpoint.
There are therefore \(18q^2\)
free--free edges crossing its boundary.
Those plaquettes change their action
by \(+2(1-\cos(t/\sqrt2))\).
Within the selected free set, each
of the \(3q^3\) parity-one sites has
two pivot neighbors, giving \(6q^3\)
free--pivot incidences. Those change
by the negative of the same quantity.
No other plaquette changes its action:
two simultaneously rotated free
neighbors have identical link values.
Subtracting these two counts proves
(80.17). The quadratic and Rayleigh
values follow by differentiation and
division by \(7q^3\).
\end{proof}

\subsection{The single-coordinate path barrier is zero, not positive}

The coordinatewise minimizing property
must not be turned into a metastability
claim. Fix \(q=4\), \(M>4\), and a
nonzero \(t_0\) of magnitude at most
\(1/4\). Let \(Y^\dagger\) be the
endpoint of (80.17). For finite paths
\(\gamma=(Y^0,\ldots,Y^k)\) with
\(Y^0=Y^*\), \(Y^k=Y^\dagger\), and
successive states differing in one
independent group coordinate, define
\[
 \mathfrak b(Y^*,Y^\dagger)=
       \inf_\gamma\max_{0\le j\le k}
                   [\mathcal S(Y^j)-\mathcal S(Y^*)].
 \tag{80.18}
\]

\begin{proposition}[Zero infimal coordinate communication barrier]
\label{prop:f12v76-zero-barrier}
For these actual completed-fibre endpoints,
\(\mathfrak b(Y^*,Y^\dagger)=0\).
Nevertheless every nonconstant finite
coordinate path out of \(Y^*\) has
a first action-increasing step.
\end{proposition}
\begin{proof}
The collective curve with parameter
between zero and \(t_0\) has action
at most its initial action, by (80.17).
Approximate it by a coordinate
staircase: divide the parameter
interval into \(N\) equal pieces,
and at each stage move each of the
448 coordinates to the next common
parameter value, one at a time.
Every intermediate state is within
product coordinate distance
\(\sqrt{448}|t_0|/N\) of a state
on the collective curve. The action
is uniformly continuous on this
fixed compact angle box. Hence the
maximum excess action of the
staircase tends to zero as
\(N\to\infty\). Its endpoints
are exact, and every step is an
allowed independent group replacement.
Since every path includes its
initial state, the infimum in
(80.18) is also at least zero.
Finally (80.8) forces the first
changed coordinate to increase
the action. Thus the infimum zero
need not be attained by a finite
non-increasing coordinate path.
\end{proof}

This result is not a stochastic
escape-time estimate. Exact paths
have Haar probability zero, and
using neighborhoods of these
staircases requires controlling
their heat-bath probabilities and
their dependence on \(\beta\).
It does show why a positive
one-coordinate improvement at
every state is an unnecessarily
strong condition for fast escape.
There is no positive classical
barrier for the particular descent
just exhibited.

\subsection{What transfers to the integrated interface}

For this paragraph only, assume that
the already chosen packing has
\[
                    T_0\subset\mathcal S,\qquad
                    T_0\cap J_b=\varnothing
                    \quad\hbox{for every integrated block }b .
 \tag{80.19}
\]
This is whole-sheet clearance,
stronger than the two-coordinate
clearance in V71. No assertion
that an arbitrary fixed packing
satisfies it is made.

Condition all interface coordinates
outside \(T_0\) to the section.
By the exact V28 factorization,
every integrated-block free energy
is independent of \(v_{T_0}\).
Consequently the actual minimized
classical patch energy equals the
sheet action, up to its constant
chosen to vanish at the section.
The sheet configuration is a
nonsection strict coordinatewise
minimum of this actual compact
conditional energy by the sheet
part of (80.8), but has the
collective descent (80.16)--(80.18).
The near-\(-I\) pivots are locally
the identity map in their input
pivot, so their Jacobian factors
are constant there as well.
No boxed replacement of a compact
block integral was used.

This conditional statement disproves
a blanket no-coordinatewise-minimum
assertion over all such compact
patches. It does not by itself
disprove the specific improvement
condition for every prescribed
repair \(T_j\): those patches
contain additional groups, may
have different exteriors, and may
couple to integrated blocks.
Nor does (80.8) for the unintegrated
action imply the same assertion
after those blocks are minimized;
minimizing a block allows a
collective move, exactly the
distinction at issue here.

\subsection{Exact checks and the next upstream estimate}

The independent script
\begin{center}\small\ttfamily
scripts/verify\_ym\_f12\_v76\_center\_sheet\_saddle.py
\end{center}
enumerates physical fine plaquettes
and all sixteen chronological
paths at \(M=2,3\). It verifies
the sheet's 56 and 189 free
coordinates, its 48 and 162
negative plaquettes, and actions
192 and 648. The signed physical
curl gives diagonal counts
\(24,32\) and \(81,108\) at
values \(2,6\), and collective
Rayleigh quotient \(-6/7\).

The same script independently
counts selected-cube boundaries
for \(q=1,2,3,4\) on a larger
torus. The quadratic values
are \(12,24,0,-96\); at \(q=4\)
the support is 448 and the
Rayleigh quotient is \(-3/14\).
It checks the exact cutoff
margins \(95/16,95/8,3/4\).

As an additional finite diagnostic,
it differentiates the actual
completed pivot equations at the
central background for all 720
free coordinates of the \(M=2\)
fibre. In a unit Lie direction
the Hessian diagonal histogram is
\[
 \begin{array}{c|rrrrrrrrr}
 \hbox{value}&2&6&7&15/2&8&9&10&12&14\\ \hline
 \hbox{count}&24&32&16&64&104&88&96&224&72
 \end{array}
\]
Every identity-coordinate physical
support has at most three links
and diagonal-dominance margin at
least three, as used in the
analytic proof. These rational
checks all pass. They supplement
the all-volume geometric arguments;
they are not a continuum mass-gap
certificate.

\begin{decision}[Allow collective saddle escape; keep the actual clock]
The universal single-coordinate
minimum exclusion is false for
the completed-fibre action.
The obstruction here is a
coordinatewise minimum with
negative collective curvature,
not a positive-barrier well.
Thus repeating the low-energy
convexity argument cannot supply
the missing global estimate.

The next meaningful input is a
quantitative escape argument for
the actual random-scan dynamics
near such saddles, or a
block-minimization argument that
removes them and is valid for
the already chosen repair
templates. Either route must
charge the update probabilities,
the dependence between interface
coordinates and minimized blocks, and the
within-target conductance.
Replacing the target clock by
a collective update without a
proved comparison would not
close the programme.
\end{decision}

V75 had 1215884 bytes and SHA--256
\begin{center}\scriptsize\ttfamily
04F123CE7072C578FDF6C65812C7A052E63EDD55519E187988408FDA72842DD3.
\end{center}
Its entire body is retained apart
from the title version. Only
active V75 is replaced by V76.
Archives, companion notes, and
previous scripts remain.

\begin{firewall}[Position after V76]
The explicit \(SU(2)\)
completed-fibre coordinatewise
minimum, its negative collective
direction, and its zero infimal
coordinate-path barrier are
proved under the stated cutoff
bound. The integrated-interface
version has the stated clearance
restriction. None proves slow
mixing or fast escape for the
actual heat bath. The required
full repair gap, cofinal
auxiliary gap, continuum
construction, and physical
Yang--Mills mass gap remain
unproved. The full objective
remains active.
\end{firewall}
\section{V77: a local saddle floor for the original heat-bath generator}
\label{sec:f12v77}

\subsection{Outcome, and what replaces a probability assigned to a path}

V76 was a progress turn: it constructed an actual compact
\(SU(2)\) coordinatewise minimum with negative collective
curvature, and thereby ruled out a blanket single-coordinate
descent argument. Its zero classical path barrier did not
bound probabilities of following those paths.

This note obtains a local spectral escape estimate by a
different argument. Negative curvature in one fixed
collective direction gives a killed diffusion floor in
a saddle neighborhood. A localization estimate then
transfers this information to the original full-group
heat-bath generator, without assuming a global diffusion
gap or replacing the clock by a collective update.

For the unintegrated completed fibre the resulting
local killed heat-bath floor is bounded below by a
positive constant for all sufficiently large \(\beta\).
For the compactly integrated interface, under V76's
explicit sheet-clearance condition, it is bounded
below by \(c/(\beta+1)\). All constants concern a
fixed local saddle pattern and are independent of
sufficiently large ambient volume.

These are bounds for functions supported in a
specified saddle neighborhood. They are not the
global complement estimate or the full repair
Poincare inequality still required by the programme.
The last companion audit was V75; the next
companion audit and direction review remain V80.

\subsection{The bounded-operator form of the inherited clock comparison}

On a finite product of the fixed compact connected
group, let \(\nu\propto e^{-W}\) relative to product
Haar, with \(W\) smooth. Retain (64.1):
\[
 L_i=-\Delta_i+\langle\nabla_iW,\nabla_i\,\cdot\,\rangle,
 \quad L=\sum_iL_i,\quad
 P_if=\nu(f\mid v_{i^c}),\quad H=\sum_i(I-P_i).
 \tag{81.1}
\]
Thus \(H\) still has rate one per original group.
Let \(K>0\) bound the mixed intrinsic Hessian row
as in (64.4). Increasing a nonnegative bound
slightly if necessary ensures \(K>0\), without
affecting the scaling conclusions below.
V60 already proved, without commuting the
coordinate generators,
\[
 \sum_i\|L_i u\|_2^2\le\|Lu\|_2^2+K E_L(u),
 \qquad P_iL_i u=0.
 \tag{81.2}
\]

\begin{proposition}[Global bounded comparison, not just a comparison of gaps]
\label{prop:f12v77-resolvent}
Under (81.2), on the entire weighted \(L^2\) space,
\[
                         H\succeq L(L+K)^{-1}.
 \tag{81.3}
\]
The right side is zero on constants.
No volume-uniform lower bound on \(L\) is assumed.
\end{proposition}
\begin{proof}
For smooth \(f\), put \(u=(L+K)^{-1}f\), which
is smooth by finite-volume elliptic regularity.
Set \(b=\langle f,L(L+K)^{-1}f\rangle\ge0\).
Conditional integration by parts and
Cauchy--Schwarz give
\[
 \begin{split}
 b^2
 &=\left|\sum_i\langle(I-P_i)f,L_i u\rangle\right|^2\\
 &\le E_H(f)\sum_i\|L_i u\|_2^2\\
 &\le E_H(f)\langle u,L(L+K)u\rangle
  =E_H(f)b .
 \end{split}
 \tag{81.4}
\]
Cancel \(b\) if positive; if zero there is
nothing to prove. Both operators in (81.3)
are bounded at fixed finite volume, so
density extends the inequality to all \(L^2\).
In particular the proof uses \((L+K)^{-1}\),
not a uniform bound on the centered inverse
\(L^{-1}\).
\end{proof}

This is the additional operator content of
V60's square estimate. A local killed
diffusion eigenvalue cannot simply be
substituted for the global \(g_N\) in
(64.9): a Dirichlet resolvent on a chart
has boundary conditions absent from
(81.2). The next two steps avoid that
unjustified substitution.

\subsection{A directional negative-curvature estimate with the correct density}

Let \(F\) be a fixed finite set of group
coordinates and use normal coordinates
\(x_F\) in a product chart, with all other
group variables denoted by \(z\).
Relative to \(dx_F\,dz\), write the density
as proportional to
\[
 e^{-\Phi(x_F,z)},\qquad
 \Phi=W-\sum_{i\in F}\log j_G(x_i),
 \tag{81.5}
\]
where translations of the chart centers
are understood. Only the Haar factors
of the charted groups occur in this
formula; the other reference factors
remain Haar.

Let \(a\) be a fixed Euclidean unit vector
supported in \(x_F\), and suppose on an
open product chart cylinder \(\mathcal U^+\)
that, uniformly in the uncharted exterior,
\[
                  \partial_a^2\Phi\le-\eta\beta,
                  \qquad \eta>0.
 \tag{81.6}
\]
Let \(c_g>0\) be a lower inverse-metric
bound on this chart. Then pointwise
\(\sum_i|\nabla_i h|^2\ge
c_g|\partial_a h|^2\).

\begin{proposition}[Local directional diffusion floor]
\label{prop:f12v77-directional}
For every smooth \(h\) supported compactly
in the chart coordinates of \(\mathcal U^+\),
\[
                 E_L(h)\ge\lambda_\beta\|h\|_2^2,
                 \qquad \lambda_\beta=c_g\eta\beta/2 .
 \tag{81.7}
\]
No condition is imposed on the other
Hessian directions or on \(\partial_a\Phi\).
\end{proposition}
\begin{proof}
For each fixed exterior, put
\(g=h e^{-\Phi/2}\), with the common
normalizing constant suppressed.
An integration by parts in the constant
direction \(a\) gives
\[
 \int|\partial_a h|^2e^{-\Phi}
 =
 \int\left[
 |\partial_a g|^2+
 \left(\frac{|\partial_a\Phi|^2}{4}
                    -\frac{\partial_a^2\Phi}{2}\right)|g|^2
 \right].
 \tag{81.8}
\]
The compact support removes boundary
terms. The last potential is at least
\(\eta\beta/2\) by (81.6). Integrate
the exterior and use the inverse-metric
lower bound. This proves (81.7).
\end{proof}

Equation (81.8) is the standard
ground-state conjugation underlying
the scalar Witten Laplacian; compare
L.~Michel,
\href{https://arxiv.org/pdf/1702.01837}
{\emph{About small eigenvalues of Witten Laplacian},
Section 1.1, p.~3}.
It is derived here with the normalization
and Haar correction actually used.
No Morse assumption, saddle eigenvalue
asymptotic, or Eyring--Kramers formula
from that work is being imported.

\subsection{Localization excludes low diffusion modes from the saddle region}

Choose an inner cylinder \(\mathcal U\)
whose closure is contained in
\(\mathcal U^+\), and a smooth cutoff
\(\chi\) depending only on \(F\), with
\[
 0\le\chi\le1,\qquad
 \chi=1\ \hbox{on }\mathcal U,\qquad
 \operatorname{supp}\chi\Subset\mathcal U^+,\qquad
 M_\chi=\sup\sum_i|\nabla_i\chi|^2<\infty .
 \tag{81.9}
\]
Compact support here refers to the
charted coordinates, not to a restriction
of the uncharted exterior.
The set \(F\) and the cutoff are fixed,
so \(M_\chi\) has no ambient-volume factor.

Apply (81.7) to \(\chi u\). The product
rule and \(|p+q|^2\le2|p|^2+2|q|^2\)
yield the global form inequality
\[
 \|1_{\mathcal U}u\|_2^2
 \le\frac{2}{\lambda_\beta}E_L(u)
              +\frac{2M_\chi}{\lambda_\beta}\|u\|_2^2 .
 \tag{81.10}
\]
This holds for every global smooth \(u\)
and then for the full diffusion form
domain by closure. It is stronger
than merely knowing a Dirichlet floor
for functions already supported in
the inner cylinder.

\begin{proposition}[Low-spectral-mass bound without commuting a cutoff]
\label{prop:f12v77-low-spectrum}
Let \(\Pi_t=1_{[0,t]}(L)\) for \(t>0\), and
put
\[
                         r_t=\frac{2(t+M_\chi)}{\lambda_\beta}.
 \tag{81.11}
\]
Then
\(\|1_{\mathcal U}\Pi_t\|^2\le r_t\).
Consequently every \(f=1_{\mathcal U}f\)
satisfies
\(\|\Pi_t f\|_2^2\le r_t\|f\|_2^2\).
\end{proposition}
\begin{proof}
Vectors in the spectral range of
\(\Pi_t\) have diffusion energy at most
\(t\) times their squared norm.
Insert them in (81.10), obtaining
the first operator-norm bound.
The adjoint of \(1_{\mathcal U}\Pi_t\)
is \(\Pi_t1_{\mathcal U}\) and has the
same norm. Applying it to
\(f=1_{\mathcal U}f\) proves the second
assertion. No commutation of the
spatial and spectral projections
is used.
\end{proof}

\subsection{The actual heat bath killed on exiting the saddle cylinder}

For this same original generator \(H\)
define
\[
 \kappa_{\mathcal U}^{H}
 =\inf_{\substack{0\ne f\in L^2(\nu)\\
                         f=0\ {\rm off}\ \mathcal U}}
                      \frac{E_H(f)}{\|f\|_2^2}.
 \tag{81.12}
\]
Functions need not be centered.
This is the bottom of the killed
operator \(H_{\mathcal U}=
1_{\mathcal U}H1_{\mathcal U}\) on
\(L^2(\mathcal U,\nu)\).
It is not the gap of a reflecting
or rejected-update chain in
\(\mathcal U\).

\begin{theorem}[Localized negative curvature gives an original-clock floor]
\label{thm:f12v77-killed-heat}
Under (81.2), (81.6), and (81.9),
for any \(t>0\) with \(r_t<1\),
\[
                 \kappa_{\mathcal U}^{H}
                 \ge(1-r_t)\frac{t}{t+K}.
 \tag{81.13}
\]
In particular, if
\(\lambda_\beta\ge8M_\chi\), then
\[
                  \kappa_{\mathcal U}^{H}
                  \ge\frac{\lambda_\beta}
                            {2(\lambda_\beta+8K)}.
 \tag{81.14}
\]
The proof assumes no positive
volume-uniform global diffusion
or heat-bath gap.
\end{theorem}
\begin{proof}
For \(f=1_{\mathcal U}f\), (81.3) and
the spectral theorem give
\[
 E_H(f)\ge
 \langle f,L(L+K)^{-1}f\rangle
 \ge\frac{t}{t+K}\|(I-\Pi_t)f\|_2^2.
 \tag{81.15}
\]
Proposition~\ref{prop:f12v77-low-spectrum}
bounds the omitted low-spectral norm
by \(r_t\|f\|_2^2\), proving (81.13).
For (81.14), choose \(t=\lambda_\beta/8\).
Then
\(r_t=1/4+2M_\chi/\lambda_\beta\le1/2\).
Substitution gives exactly (81.14).
\end{proof}

All diffusion operators in this proof
are the global ones on the identical
compact law. The chart cutoff enters
only through (81.10). This avoids
applying a boundary-free square
identity to a Dirichlet resolvent.
The resulting lower bound is for
the original full-group heat bath,
not for a newly chosen stochastic
flow along \(a\).

\subsection{Populate the hypotheses with the actual V76 saddle}

Return first to the unintegrated
\(SU(2)\) completed fibre at \(C=I\).
Its Haar-relative potential is
\[
          W_\beta(Y)=\beta\mathcal S(Y)-\log j(Y,I).
 \tag{81.16}
\]
Use V76's fixed \(q=4\) collective
direction, supported on 448 sheet
coordinates, and normalize its
coordinate vector to Euclidean norm
one. At the central background,
(80.17) gives
\[
                     \partial_a^2\mathcal S(Y^*)=-3/14.
 \tag{81.17}
\]
Choose a finite set \(F\) containing
these coordinates and the full
stencils of every action or
Jacobian term depending on them.
All chart centers in \(F\) are the
corresponding \(I\) or \(-I\)
values of the central pattern.
Take the ambient torus sufficiently
large that this fixed stencil is
represented without additional
wraparound identifications.

By V18's fixed-range smooth
completion, all derivatives needed
here are bounded on fixed compact
templates. The expression in
(81.17) depends only on \(F\).
Continuity therefore gives an
open product cylinder
\(\mathcal U^+\), of radius independent
of ambient volume, on which
\[
                     \partial_a^2\mathcal S\le-3/28,
 \tag{81.18}
\]
uniformly over all variables outside
\(F\). This is a neighborhood of a
local central pattern, not an event
requiring the entire fine volume
to stay near one configuration.

The directional second derivatives
of \(-\log j\) and of the coordinate
Haar correction in (81.5) are
bounded by a fixed constant.
After increasing a fixed threshold
\(\beta_0\), (81.18) yields
\[
                \partial_a^2\Phi_\beta\le-(3/56)\beta.
 \tag{81.19}
\]
Thus one may use \(\eta=3/56\)
and \(\lambda_\beta=c_0\beta\),
where \(c_0=3c_g/112>0\).
Choose a smaller cylinder and
cutoff as in (81.9), independently
of \(\beta\).

For (81.16), the global mixed
Hessian row has the stronger bound
\[
                         K_\beta\le C_{\rm f}(\beta+1).
 \tag{81.20}
\]
Indeed the action and logarithmic
Jacobian are sums of smooth
fixed-stencil terms. A given free
group belongs to only a bounded
number of terms, each with bounded
second derivatives before its
factor of \(\beta\). Summing a
mixed row therefore gives (81.20)
without a count of all cells.
No integrated covariance term
occurs in this unintegrated
potential.

\begin{corollary}[Actual completed-fibre local saddle floor]
\label{cor:f12v77-fibre}
For sufficiently large \(\beta\)
and ambient volume as just specified,
the original rate-one free-group
heat bath for (81.16) satisfies
\[
 \kappa_{\mathcal U}^{H_{\rm f}}
 \ge\frac{c_0\beta}
          {2(c_0\beta+8C_{\rm f}(\beta+1))}
 \ge\frac{c_0}{2(c_0+16C_{\rm f})}>0 .
 \tag{81.21}
\]
The constants are independent
of ambient volume and of
\(\beta\) above the fixed threshold.
\end{corollary}
\begin{proof}
Use (81.19)--(81.20) in (81.14),
increasing \(\beta_0\) so
\(c_0\beta\ge8M_\chi\) and
\(\beta\ge1\). Then
\(\beta+1\le2\beta\), proving
the displayed constant floor.
\end{proof}

Next consider the exact integrated
interface under V76's whole-sheet
clearance hypothesis (80.19).
Every integrated-block free energy
is independent of the collective
sheet direction. Its second
directional derivative is therefore
zero, not merely small.
The same central negative direction
and finite-stencil continuity
argument give (81.19), with a
possibly different fixed cylinder
and harmless constants.
All block integrals remain compact.

The global interface mixed-Hessian
row is the inherited bound
\(K_\beta\le C_{\rm I}(\beta+1)^2\),
which must be retained for the
global operator comparison.
Thus its actual rate-one
single-interface-group heat bath
has
\[
 \begin{split}
 \kappa_{\mathcal U}^{H_{\rm I}}
 &\ge\frac{c_0\beta}
       {2(c_0\beta+8C_{\rm I}(\beta+1)^2)}\\
 &\ge\frac{c_0}{2(c_0+32C_{\rm I}\beta)}
 \ge\frac{c}{\beta+1}
 \qquad(\beta\ge\beta_0),
 \end{split}
 \tag{81.22}
\]
where one may take
\(c=c_0/[2(c_0+32C_{\rm I})]\).
This is an interface conclusion
only under the stated clearance;
it is not asserted after arbitrary
block minimization or for an
arbitrary previously fixed packing.

\subsection{The precise stochastic meaning, including the initial-law limitation}

Let \(X_t\) be the continuous-time
heat bath with generator \(-H\)
and rate one per original group.
Let \(\tau_{\mathcal U^c}\) be its
first exit time from \(\mathcal U\).
Killing on this exit has semigroup
\(e^{-tH_{\mathcal U}}\); no proposals
are reflected or rejected at the
boundary. If \(\kappa\) is either
applicable lower bound above,
\[
                 \|e^{-tH_{\mathcal U}}\|_{2\to2}
                                      \le e^{-\kappa t}.
 \tag{81.23}
\]
For initial distribution
\(\nu(\,\cdot\mid\mathcal U)\), reversibility
and Cauchy--Schwarz give
\[
 \begin{split}
 \mathbb P_{\nu(\cdot\mid\mathcal U)}
        (\tau_{\mathcal U^c}>t)
 &=\frac{\langle1_{\mathcal U},
               e^{-tH_{\mathcal U}}1_{\mathcal U}\rangle}
          {\nu(\mathcal U)}
 \le e^{-\kappa t},\\
 \mathbb E_{\nu(\cdot\mid\mathcal U)}
          \tau_{\mathcal U^c}&\le\kappa^{-1}.
 \end{split}
 \tag{81.24}
\]
For any other initial probability
density \(r\) supported in the
cylinder and belonging to
\(L^2(\nu)\), the survival bound is
\(\|r\|_2\sqrt{\nu(\mathcal U)}
e^{-\kappa t}\).

There is no uniform pointwise
claim for every starting state,
in particular no such claim for
the exact central saddle as a
Dirac initial measure. A
conditional-equilibrium exit
bound and a killed \(L^2\)
floor do not provide that
pointwise bound without additional
regularization estimates.
Nor does exit from this cylinder
mean entry into the low-energy
target of V75: the exit could
lead to another high-energy region.

\subsection{Checks, remaining coverage, and version accounting}

The exact script
\begin{center}\small\ttfamily
scripts/verify\_ym\_f12\_v77\_saddle\_localization.py
\end{center}
checks the square and bounded
resolvent comparisons for an
interacting four-state conditional
law whose two heat-bath projections
do not commute.
For an eight-state product example,
the spatial and low-spectral
projections also do not commute.
It verifies (81.10), the low-mode
bound \(r_t=1/2\), and the
resulting killed floor \(1/8\);
the exact killed floor there is
\(3/2\).

Two-state weighted examples check
the distinction between linear
and quadratic \(K_\beta\): at
\(\beta=8,16,32\), the derived
floors are respectively \(1/4\)
and \(1/[2(1+\beta)]\), with
all operator inequalities checked
by exact weighted positive
semidefinite elimination.
Finally rational polynomial
differentiation checks the
weighted divergence identity
behind (81.8), with
\(\Phi=-a x^2/2\),
\(h=(1-x^2)^2\), and
\(a=1,3,7\). All checks pass.
These are auxiliary finite
checks; the actual saddle
application is the analytic
locality and curvature argument
above, using V76's physical
stencil calculation.

\begin{decision}[The strict saddle is locally controlled; the global coverage is not]
The identified coordinatewise
saddle no longer requires an
assumed pointwise monotone
descent path. Its actual
neighborhood has the local
killed heat-bath estimate
(81.21), and (81.22) where
the integrated clearance holds.

The remaining high-energy task
is to identify regions not
covered by productive group
updates or negative collective
curvature, including any genuine
local minima and the effect of
the actual integrated blocks
and required exteriors.
Even a collection of local
killed bounds does not
automatically yield a global
repair gap: transitions and
localization errors between
regions must be controlled.
That global assembly, and
V74's within-target conductance,
remain separate required inputs.
\end{decision}

V76 had 1236585 bytes and SHA--256
\begin{center}\scriptsize\ttfamily
17A05CEA0E1B1DE8CDF6E9D280EBDE453792513D6DFF89B6CD2CA969458AB577.
\end{center}
Its entire body is retained
apart from the title version.
Only active V76 is replaced
by V77. Archives, companions,
and previous scripts remain.

\begin{firewall}[Position after V77]
The original-clock local saddle
floor is proved for the
completed \(SU(2)\) fibre,
and with the stated restriction
for the compact interface.
The associated exit bound
has the initial-law scope
specified in (81.24).
Neither result is a full
repair gap or a classification
of all high-energy states.
No cofinal auxiliary gap,
continuum construction, or
physical Yang--Mills mass gap
has been established. The
full objective remains active.
\end{firewall}
\section{V78: bounded negative curvature for every one-direction central boundary pattern}
\label{sec:f12v78}

\subsection{The new coverage result and its scope}

V77 controlled a neighborhood of V76's fully occupied
central sheet. It did not establish that other central
patterns have a similar escape direction. We now prove
that every nonempty pattern of \(-I\) on nonpivot
boundary links in one fixed direction has a strictly
negative collective direction supported on at most
448 groups, with the same uniform Hessian Rayleigh
upper bound \(-3/14\).

The proof first treats one boundary layer and then
allows arbitrary patterns in several parallel layers.
It uses the actual \(SU(2)\) smooth completed fibre
at \(C=I\), with every internal link and every
boundary link in other directions fixed to \(I\)
at the central configuration. The negative direction
is nevertheless a direction in the full completed
fibre, not in an artificially altered action.
The compact completion and the cutoff assumptions
are those of V76.

This class contains far more configurations than
the full sheet, but does not include general
noncentral fields, negative internal links, or
simultaneous patterns in different directions.
No completeness claim for the full high-energy
region is made. The next scheduled companion
audit and direction review remain V80.

\subsection{Two cube decompositions of a boundary layer}

As in V76, identify the transverse layer with
\[
 \Lambda=(\mathbb Z/(2M)\mathbb Z)^3,\qquad
 P=(2\mathbb Z/(2M)\mathbb Z)^3,\qquad
 F_0=\Lambda\setminus P .
 \tag{82.1}
\]
The sites in \(P\) are the reference pivots,
and those in \(F_0\) are the \(7M^3\) free
boundary coordinates of this layer.
Assume \(M\ge2\).
All graph edges below are undirected
nearest-neighbor edges, counted once.

Split the edge set into \(\mathcal E_0\)
and \(\mathcal E_1\). The former contains
edges within the ordinary parity cells
\(2r+\{0,1\}^3\). The latter contains
edges crossing ordinary coarse-cell
boundaries. Equivalently,
\(\mathcal E_1\) is the union of the
edges within the shifted cells
\(2r+(1,1,1)+\{0,1\}^3\), with periodic
coordinates. Each of these two graphs
is a disjoint union of ordinary
eight-vertex cubes, each containing
exactly one pivot site. Every vertex
has degree three in each graph.
Together they give all six nearest
neighbors of the layer.

\begin{proposition}[Elementary pinned-cube cut bound]
\label{prop:f12v78-cube-cut}
If \(S\) is a subset of the vertices of
an ordinary three-cube excluding a
specified vertex, then its edge
boundary in that cube satisfies
\[
                         |\partial S|\ge(3/7)|S|.
 \tag{82.2}
\]
The constant is sharp for the seven
vertices other than the specified one.
\end{proposition}
\begin{proof}
The empty set is immediate. Every
nonempty proper subset of a three-cube
has at least three boundary edges.
To verify this directly, split the
cube into two square faces connected
by a matching. A nonempty proper
subset of a square has at least two
boundary edges within its face.
If both face subsets are proper
and nonempty, there are at least
four such edges. If just one is,
there are at least two within that
face and at least one in the
matching. If neither is, a
nonempty proper cube subset must
consist of one full face and one
empty face, giving four matching
edges. Thus the boundary is at
least three. Since \(1\le|S|\le7\),
this proves (82.2).
The complement of one vertex has
three boundary edges and seven
vertices, giving equality.
\end{proof}

Only this elementary edge-cut fact
is needed. No general hypercube
isoperimetric classification is
assumed.

For any \(A\subset F_0\), let
\(b_j(A)\) be its edge boundary
in \(\mathcal E_j\), and let
\(n_1(A)\) be the number of
\(\mathcal E_1\) edges with both
endpoints in \(A\).
Applying (82.2) in each cube of
the two decompositions gives
\[
 \begin{split}
 b_0(A),\,b_1(A)&\ge(3/7)|A|,\\
 b(A):=b_0(A)+b_1(A)&\ge(6/7)|A|,\\
 2n_1(A)+b_1(A)&=3|A|,\qquad
 n_1(A)\le(9/7)|A|.
 \end{split}
 \tag{82.3}
\]
The degree identity in the last
line counts the three
\(\mathcal E_1\) edges incident
to each vertex of \(A\).
All bounds are uniform in \(M\).
For \(A=F_0\), equality holds
in each of the displayed
inequalities.

\subsection{Arbitrary central patterns stay on the actual completed fibre}

Fix a nonempty \(A\subset F_0\).
In direction \(2\) and the layer
\(x_2=1\), give the free boundary
link at transverse site \(i\)
value \(-I\) if \(i\in A\) and
\(I\) otherwise. Give every
other physical link value \(I\).
Denote the resulting field by
\(W^A\), and its independent
coordinates by \(Y^A\).

All chronological two-link path
products are central. At root
pivot \(I\) they have value
\(I\) or \(-I\), so their
truncated logarithms vanish.
The global pivot uniqueness
of V18 therefore gives
\[
                   \widetilde K(W^A)=I,\qquad
                   W(Y^A,I)=W^A.
 \tag{82.4}
\]
Exactly the plaquettes across
the graph boundary of \(A\)
are negative. The actual
Wilson normalization consequently
gives
\[
                       \mathcal S(Y^A)=4b(A)>0 .
 \tag{82.5}
\]

Let \(B\subset A\), fix
\(X\in\mathfrak{su}(2)\) with
\(\|X\|_{\rm HS}=1\), and replace
the links in \(B\) by
\(-\exp(tX)\), leaving the other
free links unchanged.
For \(|t|\le1/4\), these changed
path products remain outside
the cutoff support near \(I\).
Every unchanged path product
is still \(I\) or \(-I\).
Thus every pivot remains exactly
\(I\), again by evaluating its
equation at \(I\) and using
uniqueness. In particular this
curve uses the actual compact
completion, not a linear pivot
substitution.

Write \(e(B,C)\) for the number
of layer edges joining disjoint
sets \(B,C\). Define
\[
             D_A(B)=e(B,A\setminus B)-e(B,\Lambda\setminus A).
 \tag{82.6}
\]
A plaquette between a rotated
negative link and an unrotated
negative link gains
\(s_X(t)=2(1-\cos(t/\sqrt2))\).
One between a rotated negative
link and a positive link loses
the same amount. A pair rotated
together contributes no change.
Therefore the exact change and
the corresponding Hessian
quadratic value are
\[
 \mathcal S(Y^{A,B}(t))-\mathcal S(Y^A)
                   =D_A(B)\,s_X(t),\qquad
 \left.\frac{d^2}{dt^2}\mathcal S(Y^{A,B}(t))\right|_{t=0}
                   =D_A(B).
 \tag{82.7}
\]
The raw direction has squared
coordinate norm \(|B|\).

All first derivatives of the full
completed action at \(Y^A\)
vanish: a first variation of a
Wilson plaquette at \(I\) or
\(-I\) has zero real trace
derivative, and the completed
chart is smooth. Thus a negative
value in (82.7) proves that
this full-fibre critical point
is not a local minimum.

\subsection{A translated-window identity selects a bounded negative direction}

For \(1\le q<M\), let \(Q_r(q)\)
be the union of a transverse
\(q\times q\times q\) cube of
ordinary coarse cells, with
coarse anchor
\(r\in(\mathbb Z/M\mathbb Z)^3\).
Put \(B_r=A\cap Q_r(q)\).
Then \(|B_r|\le7q^3\).

\begin{proposition}[Exact average of the physical rotation curvatures]
\label{prop:f12v78-window-average}
For every \(A\subset F_0\),
\[
 \begin{split}
 \sum_r |B_r|&=q^3|A|,\\
 \sum_r D_A(B_r)&=
             2q^2n_1(A)-q^3b(A)\\
 &\le\frac67 q^2(3-q)|A|.
 \end{split}
 \tag{82.8}
\]
\end{proposition}
\begin{proof}
Each site belongs to exactly
\(q^3\) translated windows,
giving the first identity.
An edge from \(A\) to its
complement contributes negatively
whenever its \(A\)-endpoint
lies in the window, hence
\(-q^3\) to the sum.

An \(A\)--\(A\) edge inside
one ordinary owner cell cannot
be cut by a union of owner cells.
An \(A\)--\(A\) edge crossing
an owner boundary has owners
separated by one coarse step.
There are exactly \(2q^2\)
translated windows containing
one owner but not the other.
Such an edge therefore contributes
\(2q^2\). This proves the second
identity, including periodic
windows because \(q<M\).
Insert \(n_1(A)\le9|A|/7\)
and \(b(A)\ge6|A|/7\) from
(82.3) for the last inequality.
\end{proof}

\begin{theorem}[Uniform bounded saddle direction for every nonempty pattern]
\label{thm:f12v78-bounded-descent}
If \(M>4\) and
\(\varnothing\ne A\subset F_0\),
there is a translated coarse
cube \(Q_r(4)\) such that
\[
 1\le|B_r|\le448,\qquad
                  \frac{D_A(B_r)}{|B_r|}\le-\frac3{14}.
 \tag{82.9}
\]
Consequently \(Y^A\) has a
unit full-fibre coordinate
direction \(a\), supported on
at most 448 groups, satisfying
\[
                         \partial_a^2\mathcal S(Y^A)\le-3/14 .
 \tag{82.10}
\]
It also has the exact decreasing
curve (82.7) on that support.
In particular no nonempty
one-layer central pattern is
a genuine local minimum of
the completed-fibre action.
\end{theorem}
\begin{proof}
Take \(q=4\) in (82.8).
The total support mass is
\(64|A|>0\), while the ratio
of the total curvature to
this mass is at most \(-3/14\).
It is the \(|B_r|\)-weighted
average of
\(D_A(B_r)/|B_r|\) over
nonempty \(B_r\). At least
one ratio is no larger.
The support bound is
\(7\cdot4^3=448\).
Normalize the corresponding
direction by \(|B_r|^{-1/2}\).
Equation (82.7) then gives
(82.10) and the strict
decrease for every sufficiently
small nonzero \(t\).
\end{proof}

The scale \(q=4\) and the
constant in (82.9) are sharp
for this particular selection
of directions: a uniform
rotation of the negative
links inside a coarse cube.
Indeed the full pattern
\(A=F_0\), already computed
in (80.17), has
\[
        \frac{D_{F_0}(F_0\cap Q_r(q))}
             {|F_0\cap Q_r(q)|}
                   =\frac67\left(\frac3q-1\right)
 \tag{82.11}
\]
for every anchor. It is
positive for \(q=1,2\), zero
for \(q=3\), and \(-3/14\)
for \(q=4\).
This does not assert that
448 is the smallest support
among all possible directions.

\subsection{Parallel layers and uniform local heat-bath bounds}

The same classification applies
if arbitrary nonpivot boundary
links in several layers of
one fixed direction carry
\(-I\), while all internal
links and all other directions
remain \(I\). A plaquette
containing this direction
compares two parallel links
at the same longitudinal
coordinate. Its action
therefore belongs to exactly
one such layer.
The completed coarse-edge
equations likewise use the
boundary links of their own
layer when internal links
are identity.
Choose any nonempty layer
and apply Theorem~\ref{thm:f12v78-bounded-descent}
there. Other layers do not
change (82.7).

Thus every nontrivial member
of this one-direction central
boundary class has the same
bounded strictly negative
direction. On an \(M^4\)
coarse torus the class allows
arbitrary choices on \(7M^4\)
boundary groups in the chosen
direction. This statement
does not extend by superposition
to different link directions:
their plaquettes can interact.

\begin{corollary}[Uniform local saddle estimates over the central-pattern class]
\label{cor:f12v78-uniform-local}
For the unintegrated completed
\(SU(2)\) fibre at \(C=I\),
there are fixed \(c_*>0\)
and \(\beta_0<\infty\)
such that every nontrivial
central pattern just described
lies in a local chart cylinder
\(\mathcal U_A\) with
\[
                    \kappa_{\mathcal U_A}^{H_{\rm f}}\ge c_*
                    \qquad(\beta\ge\beta_0).
 \tag{82.12}
\]
The cylinder restricts only
a fixed-size stencil around
a selected four-cell-side
cube. Its radii, stencil-size
bound, and the constants
can be chosen uniformly
over patterns and sufficiently
large ambient tori.
\end{corollary}
\begin{proof}
At the selected center the
unit-direction curvature is
at most \(-3/14\), exactly
the margin used in V77.
Enlarge the selected cube
by the fixed completion,
action, and Jacobian stencils.
Only finitely many local
central patterns are possible
on that fixed set, and
the selected direction has
entries \(X/\sqrt{k}\) on
a subset of size
\(1\le k\le448\).
Up to translation there
are finitely many such
pattern/direction pairs.

V18 bounds the needed local
derivatives uniformly on
these templates. Taking
a common smaller radius
over the finite collection
preserves the negative
curvature margins in
(81.18)--(81.19).
The inverse-metric bound,
the cutoff gradient cost,
and the amplitude derivative
bounds can also be chosen
uniformly. No condition on
variables outside the
enlarged stencil is needed,
by locality.
Finally the global mixed
Hessian row of the same
unintegrated law is still
\(C_{\rm f}(\beta+1)\).
Apply Corollary~\ref{cor:f12v77-fibre}.
\end{proof}

For the integrated interface,
the corresponding statement
holds for the \emph{single-layer}
family under the same
whole-layer clearance as
(80.19): all its nonpivot
boundary variables are
interface variables outside
every \(J_b\), and all
remaining interface variables
are fixed to the section
at the central configuration.
Every compact block free
energy is independent of
the selected direction.
The same finite-template
argument, with the inherited
global row bound
\(C_{\rm I}(\beta+1)^2\),
therefore gives
\[
                    \kappa_{\mathcal U_A}^{H_{\rm I}}
                                   \ge\frac{c_{\rm I}}{\beta+1}.
 \tag{82.13}
\]
No multilayer interface
conclusion is claimed if
other negative layers enter
the integrated blocks.
The clearance and exact
compact law have not been
dropped when applying V77.

The stochastic interpretation
of (82.12)--(82.13) is
precisely that of (81.23)--(81.24):
an original-clock killed
\(L^2\) floor and a survival
bound from equilibrium
conditioned on the particular
cylinder. It is not a
uniform pointwise exit-time
bound for arbitrary initial
states.

\subsection{Checks and the remaining high-energy coverage problem}

The script
\begin{center}\small\ttfamily
scripts/verify\_ym\_f12\_v78\_boundary\_patterns.py
\end{center}
checks all 127 nonempty
subsets of the seven
nonpivot cube vertices,
including the sharp cut
ratio \(3/7\).
It checks the two
decompositions and the
sharp full-layer ratio
\(6/7\) on all 127
nonempty periodic parity
patterns of a four-site-side
transverse torus.

For 25 structured or
deterministically seeded
patterns on a ten-site-side
layer, it checks all 125
translated side-four
coarse windows per pattern.
Both exact identities in
(82.8) hold, and every
pattern has a nonempty
window with normalized
curvature at most \(-3/14\).
The full pattern saturates
this worst-best bound.

Independent four-dimensional
physical plaquette and
chronological-path enumeration
checks three further patterns,
including one with a second
nonempty parallel layer.
Their actions are \(280,624,376\)
and the selected rotation
curvatures are \(-40,-48,-30\).
The changed path products
are always in the negative
central neighborhood, as
required by the exact pivot
argument. All checks pass.
The all-pattern, all-volume
statements are the proofs
above, not an inference from
the sampled patterns.

\begin{decision}[Central boundary saddles are classified; do not count local bounds as a global gap]
This step removes genuine
local minima from an entire
explicit central-pattern
class and supplies uniform
bounded-support negative
directions. It does not
classify noncentral phases,
internal-link excitations,
or patterns mixing link
directions.

Nor do the local killed
floors in (82.12) add
automatically to a floor
on their union. A trajectory
can exit one cylinder into
another, and summing
localizations can introduce
errors or volume factors.
The next step must either
control those transitions
and localization errors,
or identify genuinely
uncovered critical
configurations using the
actual action. Repeating
the central-sheet example
alone is no longer the
missing input.
\end{decision}

V77 had 1254658 bytes and SHA--256
\begin{center}\scriptsize\ttfamily
68B421944E1F1513DFB5812D37F59FC03FEB1A87F45D49068780577CF0B4EC8E.
\end{center}
Its entire body is retained
apart from the title version.
Only active V77 is replaced
by V78; archives, companions,
and previous scripts remain.

\begin{firewall}[Position after V78]
The bounded negative-direction
classification is proved for
the stated one-direction
central boundary class.
Its uniform local heat-bath
consequences retain the
specified law and clearance
restrictions. The global
repair estimate and general
high-energy classification
remain unproved, as do the
cofinal auxiliary gap,
continuum construction,
and physical Yang--Mills
mass gap. The full
objective remains active.
\end{firewall}
\section{V79: why local saddle floors cannot become a cofinal union-killing floor}
\label{sec:f12v79}

\subsection{Outcome and its relation to the earlier probability obstruction}

V78 proved uniform bounded negative directions for a
central-pattern class, with local killed heat-bath
floors supplied by V77. The next assembly step needs
a scope check. We now prove that a volume-uniform
first killed floor on the union of many such local
saddle cylinders is impossible at fixed coupling,
even for the actual completed-fibre law.

The obstruction is not slow exit from each local
cylinder. It is the rarity of the event that
\emph{none} of the selected cylinders occurs
anywhere. Killing only on that avoidance event
poses a different global hitting problem.

V40 already showed that the global all-chart
restriction is not a small total-variation
perturbation. That result is reused as context,
not reproved as a new claim. The new conclusions
here concern the first killed spectral value,
conditional-equilibrium exit times, and the
specific strict-saddle cylinders with proved
local rates in V77--V78.
The actual unrestricted Poincare target remains
unchanged. The next companion and direction
reviews are V80.

\subsection{A general variational bound retaining the constant mode}

Let \(\nu\) be a probability law and
\[
             H=\sum_{i=1}^N(I-P_i)
 \tag{83.1}
\]
its rate-one heat bath, with each \(P_i\)
an orthogonal conditional-expectation
projection. For a measurable region \(D\)
with \(0<\nu(D)<1\), write \(A=D^c\),
\(a=\nu(A)\), and
\[
 \kappa_D=\inf_{\substack{0\ne f\in L^2(\nu)\\f=0\ {\rm off}\ D}}
                   \frac{E_H(f)}{\|f\|_2^2},\qquad
 R_D=\frac{E_H(1_D)}{\nu(D)}.
 \tag{83.2}
\]
These infima are not restricted to
centered functions. They describe
killing on first exit from \(D\).
For the standard Dirichlet variational
formalism, compare
\href{https://www.stat.berkeley.edu/~aldous/RWG/Book_Ralph/Ch3.S6.html}
{Aldous--Fill, Section 3.6};
the bounded-operator argument needed
here is given explicitly below.

\begin{proposition}[The almost-constant indicator is an admissible killed test]
\label{prop:f12v79-indicator}
One has
\[
                  \kappa_D\le R_D\le\frac{Na}{1-a}.
 \tag{83.3}
\]
For the original continuous-time
heat bath, started from
\(\nu(\,\cdot\mid D)\), its exit time
\(\tau_A\) satisfies
\[
       \mathbb P_{\nu(\cdot\mid D)}(\tau_A>t)
                                  \ge e^{-R_Dt},\qquad
       \mathbb E_{\nu(\cdot\mid D)}\tau_A
                                  \ge R_D^{-1}.
 \tag{83.4}
\]
The usual infinite-mean interpretation
applies if \(R_D=0\).
\end{proposition}
\begin{proof}
Use \(1_D\) in the killed variational
infimum. Since \(H1=0\),
\(E_H(1_D)=E_H(1_A)\). Each
\(I-P_i\) is an orthogonal
projection and hence bounded
above by \(I\). Consequently
\[
 E_H(1_A)=\sum_i\|(I-P_i)1_A\|_2^2
                           \le N\|1_A\|_2^2=Na.
 \tag{83.5}
\]
Division by \(\nu(D)=1-a\)
proves (83.3).

On \(L^2(D,\nu)\), the killed
operator is \(H_D=1_DH1_D\).
It is bounded, nonnegative,
and self-adjoint. The spectral
measure of \(H_D\) for
\(1_D/\sqrt{\nu(D)}\) is a
probability measure with first
moment \(R_D\). The survival
probability equals its integral
of \(e^{-t\lambda}\). Convexity
of this function and Jensen's
inequality give the first
inequality in (83.4).
Integration over \(t\ge0\)
gives the second. No pointwise
initial-state bound is asserted.
\end{proof}

The small Rayleigh quotient in
(83.3) need not represent a
slow centered mode. Indeed
\[
 \begin{split}
 \left|\left\langle\frac{1_D}{\sqrt{\nu(D)}},1\right\rangle\right|^2
                            &=1-a,\\
 \|1_D-\nu(D)\|_2^2&=a(1-a),\\
 \frac{E_H(1_D-\nu(D))}
      {\|1_D-\nu(D)\|_2^2}&=\frac{R_D}{a}.
 \end{split}
 \tag{83.6}
\]
If \(a\) is tiny, the uncentered
killed test is almost entirely
the constant mode. Centering
costs the factor \(a\), which
cannot be omitted when comparing
to the genuine global gap.

\subsection{Many locally possible cylinders: no independence is needed}

Let \(U_1,\ldots,U_m\) depend,
respectively, on pairwise
disjoint finite coordinate
sets \(F_1,\ldots,F_m\).
Suppose a common \(p>0\) obeys
\[
            \nu(U_j\mid v_{F_j^c})\ge p
                       \quad\hbox{almost surely for every }j.
 \tag{83.7}
\]
Put
\[
                 D_m=\bigcup_{j=1}^m U_j,\qquad
                 A_m=\bigcap_{j=1}^m U_j^c .
 \tag{83.8}
\]

\begin{proposition}[Avoidance and the union-killing upper bound]
\label{prop:f12v79-union}
Under (83.7),
\[
 \nu(A_m)\le(1-p)^m,\qquad
 \kappa_{D_m}\le R_{D_m}
                 \le\frac{N(1-p)^m}{1-(1-p)^m}.
 \tag{83.9}
\]
In particular, for fixed \(p>0\)
and \(m\ge cN\), the right side
tends to zero as \(N\to\infty\).
\end{proposition}
\begin{proof}
The event
\(\bigcap_{j<k}U_j^c\) is measurable
outside \(F_k\), because the
coordinate sets are disjoint.
Conditioning on that exterior
and using (83.7) gives
\[
 \nu\left(\bigcap_{j\le k}U_j^c\right)
 \le(1-p)\,
       \nu\left(\bigcap_{j<k}U_j^c\right).
 \tag{83.10}
\]
Iteration proves the probability
bound. Apply (83.3) and the
monotonicity of \(a/(1-a)\)
to obtain the spectral bound.
Finally \((1-p)^m\le e^{-pm}\)
gives the limit. The events
may be strongly dependent;
independence was never used.
\end{proof}

\subsection{Application to the actual completed-fibre saddle cylinders}

Use the actual \(SU(2)\)
completed-fibre law at \(C=I\),
\[
             d\mu_{\beta,M}(Y)\propto
             e^{-\beta\mathcal S(Y)}j(Y,I)\,dY,
             \qquad N=45M^4 .
 \tag{83.11}
\]
Take one of V77's inner
strict-saddle cylinders \(U\),
with its fixed finite coordinate
stencil \(F\). Its chart centers
form the local central pattern
required for the bounded
negative direction. Outside
that stencil no restrictions
are imposed.
For all \(\beta\ge\beta_0\),
V77 supplies a fixed positive
local killed floor, uniformly
over sufficiently large
ambient tori.

The full potential
\(\beta\mathcal S-\log j\)
is a sum of smooth fixed-stencil
terms. When only the coordinates
of \(F\) vary, the sum of the
oscillations of incident terms
is at most
\(D_F(\beta+1)\), with a fixed
finite \(D_F\). Let
\(\alpha_F>0\) be the product
Haar measure of the inner
chart box defining \(U\).
Comparison of the conditional
numerator and denominator
therefore gives
\[
 \mu_{\beta,M}(U\mid Y_{F^c})
       \ge p_\beta:=
               \alpha_F e^{-D_F(\beta+1)}>0 .
 \tag{83.12}
\]
The chart is proper, so
\(p_\beta<1\). This is the
fixed-stencil form of the
inherited finite-energy bound;
it is uniform in the entire
compact exterior, not just
in a near-central exterior.

Choose a fixed coarse-cell
box side \(R\) large enough
to contain \(F\), with its
necessary translated margins.
One can place
\[
                       m=\lfloor M/R\rfloor^4
 \tag{83.13}
\]
disjoint translated copies
of this coordinate stencil
in the coarse torus.
The corresponding cylinder
events \(U_j\) obey (83.12)
with the same \(p_\beta\).
For \(M\ge2R\),
\[
                  m\ge\frac{M^4}{16R^4}
                    =\frac{N}{720R^4}.
 \tag{83.14}
\]
The local curvature and
cutoff estimates are unchanged
by these coarse translations.

\begin{theorem}[Actual local saddle floors coexist with a vanishing union floor]
\label{thm:f12v79-actual-union}
For the law (83.11) and
the disjoint translated
cylinders above, at any
fixed \(\beta\ge\beta_0\),
\[
 \begin{split}
 \inf_j\kappa_{U_j}^{H_{\rm f}}&\ge c_*>0,\\
 \kappa_{D_m}^{H_{\rm f}}
 &\le\frac{45M^4(1-p_\beta)^m}
                      {1-(1-p_\beta)^m}
                          \longrightarrow0
                          \quad(M\to\infty).
 \end{split}
 \tag{83.15}
\]
Moreover
\[
 \mathbb E_{\mu_{\beta,M}(\cdot\mid D_m)}\tau_{A_m}
       \ge\frac{1-(1-p_\beta)^m}
                   {45M^4(1-p_\beta)^m}
                          \longrightarrow\infty .
 \tag{83.16}
\]
This does not assert that
the unrestricted completed-fibre
gap tends to zero.
\end{theorem}
\begin{proof}
The local lower bound is
(81.21). The conditional
probability bound is (83.12),
and the supports in (83.13)
are disjoint. Apply
Proposition~\ref{prop:f12v79-union}
and then (83.4).
At fixed \(\beta\), the number
\(p_\beta\) is strictly positive;
(83.14) gives both limits.
\end{proof}

Each local cylinder can thus
be exited quickly in the
precise sense of V77, while
waiting until none of these
local patterns occurs anywhere
is slow in the different
sense of (83.16).
The avoidance event \(A_m\)
is not identified with the
vacuum or the all-good chart.
It merely excludes the
selected packed cylinders.

Any larger union containing
\(D_m\) has a no-larger
first killed floor, by
monotonicity of the variational
domain. Thus adding further
local saddle cylinders
cannot repair this cofinal
union-killing target.
No such conclusion is drawn
for the integrated interface
without constructing its
required clearance and
conditional-cylinder family.

The limits above keep
\(\beta\) fixed. For a joint
sequence \(\beta=\beta_M\),
the same upper bound tends
to zero under the sufficient
condition
\[
                       mp_{\beta_M}-\log(45M^4)\longrightarrow+\infty .
 \tag{83.17}
\]
No physical cutoff/coupling
scaling verifying that
condition is assumed.

\subsection{A product law separates the killed obstruction from a true gap}

Let \(N\) independent bits have
\(\nu(x_i=1)=p\in(0,1)\), and
use the same rate-one coordinate
heat bath. Its full gap is
exactly one by product
tensorization and a one-bit
eigenfunction.
For \(U_i=\{x_i=1\}\), the
killed operator on \(U_i\)
is
\((1-p)I+\sum_{j\ne i}(I-P_j)\).
Hence every local killed
floor is exactly \(1-p\).

For their union \(D=\{\sum_i x_i>0\}\),
the exact indicator flux is
\[
 \nu(D)=1-(1-p)^N,\qquad
 E_H(1_D)=Np(1-p)^N,\qquad
 R_D=\frac{Np(1-p)^N}{1-(1-p)^N}.
 \tag{83.18}
\]
Indeed coordinate \(i\) can
change \(1_D\) only when all
other coordinates are zero.
That event has probability
\((1-p)^{N-1}\), and the
remaining one-bit conditional
variance is \(p(1-p)\).
Sum over \(i\).
The union killed floor
therefore tends to zero,
although the unrestricted
gap stays exactly one.
Centering the same indicator
gives instead
\[
 \frac{E_H(1_D-\nu(D))}
      {\operatorname{Var}(1_D)}
                         =\frac{Np}{1-(1-p)^N},
 \tag{83.19}
\]
which does not vanish.
The bit example is illustrative;
the actual-law assertion is
Theorem~\ref{thm:f12v79-actual-union}.

\subsection{A direct audit of a proposed global spectral localization}

V77 used a global inequality
of the form
\[
                \|1_Du\|_2^2\le a_L E_L(u)+b_L\|u\|_2^2
 \tag{83.20}
\]
to exclude low diffusion modes
from a local cylinder.
If (83.20) holds for every
global \(u\), testing \(u=1\)
necessarily gives
\[
                              b_L\ge\nu(D).
 \tag{83.21}
\]
For \(D=D_m\), this forces
\(b_L\to1\) at fixed coupling.
The low-mode estimate used in
(81.13) then has
\(r_t=a_Lt+b_L\ge\nu(D_m)\).
Its available factor
\(1-r_t\) is at most
\(\nu(A_m)\), if positive
at all. Thus no sharper
calculation of the same
uncentered union localization
can produce a cofinal positive
margin: the constant mode
already rules it out.

This does not invalidate
(81.10) on a fixed local
cylinder. It prevents silently
extending that local inequality
to a union whose stationary
mass tends to one with the
same small coefficient \(b_L\).

\subsection{Exact checks and the corrected assembly target}

The script
\begin{center}\small\ttfamily
scripts/verify\_ym\_f12\_v79\_union\_killing.py
\end{center}
checks the full product gap
and the local killed floor
by exact weighted matrix
inequalities for four bits
with \(p=1/4\).
Its union indicator quotient
is \(81/175\), whereas the
centered quotient is \(256/175\).

The script also forms the
exact killed birth--death
generator for the number
of occupied bits, with rates
\(k(1-p)\) down and
\((N-k)p\) up. Rational
inversion verifies the
conditional-equilibrium
mean exit bound in seven
cases up to \(N=16\).
For \(N=16,p=1/2\), the
indicator upper bound is
\(8/65535\), and the mean
exit lower bound is
\(65535/8\), despite the
full gap remaining one.

A dependent six-bit Gibbs
example checks three
disjoint two-bit cylinders.
Their common conditional
probability lower bound is
\(1/15\), with the more
conservative oscillation
bound \(1/64\).
The exact avoidance
probability \(566/873\)
is at most \((14/15)^3\).
This verifies the conditioning
argument without an independence
assumption. All checks pass.

\begin{decision}[Do not use a global avoidance time as the cofinal gap criterion]
The cofinal first killed
floor on a union of local
saddle events is now a
proved unsuitable target
for this construction,
not merely an estimate
with uncomputed constants.
The relevant global
quantity remains the
centered unrestricted
Poincare inequality.

The original fixed repair
patches are not ruled out
by this result: their
cardinality is held fixed
as the ambient volume grows.
Their local assembly still
requires actual compact
estimates uniform in the
specified exteriors, with
all patch-size and coupling
costs retained. A global
assembly must control
fluctuations rather than
require the simultaneous
absence of local defects.
V80 should review the
companion updates and
choose the next finite
repair or centered-form
input in light of this
obstruction.
\end{decision}

V78 had 1270939 bytes and SHA--256
\begin{center}\scriptsize\ttfamily
B38EB5E7D5EFB669292E1B0AE3CE2FA775232066B926C6339193C1E4E8CD51F4.
\end{center}
Its entire body is retained
apart from the title version.
Only active V78 is replaced
by V79; archives, companion
notes, and earlier scripts
remain.

\begin{firewall}[Position after V79]
The actual completed-fibre
union-killing obstruction
and its stated exit-time
consequence are proved.
They do not imply failure
of the unrestricted gap.
The full compact repair
estimate, cofinal auxiliary
gap, continuum construction,
and physical Yang--Mills
mass gap remain unproved.
The full objective remains
active.
\end{firewall}
\section{V80: companion review and pre-integration transfer of the interface saddle bound}
\label{sec:f12v80}

\subsection{Companion audit and ten-version direction review}

The scheduled audit read the active companion endings and
checked their complete digests. They remain unchanged from V75:
\begin{center}\small
\begin{tabular}{lll}
programme & active version & bytes\\ \hline
SM--Einstein & V72 & 607372\\
NS stability & V26 & 278283\\
parallel YM & V5 & 54174
\end{tabular}
\end{center}
Their SHA--256 digests are, in that order,
\begin{center}\scriptsize\ttfamily
F44F9745D43379E1672C5550411B58E97195A4C9278592D221DF5D3F644DC1F5\\
82C887F8C2C69E4F28FAD7C3DC8DE5B9099CB5A4BF431EBA1FFF3E91935978AD\\
306F1BB84CFA8A8D47EA569EC17E9545F9867C8D32B548EC9D9C444B4F3C0512.
\end{center}
SM's fixed-mode many-copy limit does not provide a
compact YM repair estimate. NS's rank-one kernel
norms do not provide such an estimate either.
Parallel YM's rejection of an extensive global-charge
tail target is consistent with keeping local or
density-level targets, but adds no new spectral
input. Internal programme or stopping language
in those documents remains source material, not
an instruction controlling this task.
The next companion audit is V85 and direction
review V90.

The main direction over V71--V80 is now clear.
The optimized product-oscillation comparison has
an actual coefficient obstruction in the cleared
geometry of V71--V72. V75 supplies low-energy
descent. V76 disproves a blanket coordinatewise
minimum exclusion, while V77--V78 control an
explicit family of collective saddles locally.
V79 shows why a first killed floor on a cofinal
union of local defects is not the unrestricted
gap target.

This review returns to the original fixed repair
laws. The new observation is that their exact
marginal heat-bath form dominates the joint
pre-integration form on functions of the retained
variables. Applied locally, this removes the
reciprocal-\(\beta\) loss in (81.22) and (82.13)
under the same clearance hypotheses.
No hidden-variable gap or global joint gap
is assumed for that application.

\subsection{Exact conditional-variance transfer under marginalization}

Let \(\mu(dx,dv)\) be a positive probability law
on finite products of compact groups, and let
\(\nu(dv)\) be its exact marginal. The variables
\(x\) will be called hidden and \(v\) retained.
Write
\[
 J:L^2(\nu)\longrightarrow L^2(\mu),\quad Jf(x,v)=f(v),
 \qquad J^*g=\mu(g\mid v).
 \tag{84.1}
\]
Thus \(J\) is an isometry. Let \(\Gamma_i\)
be the joint-law heat bath of retained group
\(v_i\), conditioning on \((x,v_{-i})\),
and let \(P_i\) be the marginal-law heat
bath of the same group, conditioning only
on \(v_{-i}\). Include the hidden-coordinate
joint heat baths in \(H_\mu\), and set
\[
 H_\mu=\sum_{\text{all joint groups }k}(I-\Gamma_k),
 \qquad H_\nu=\sum_{\text{retained groups }i}(I-P_i).
 \tag{84.2}
\]
Both clocks have rate one per declared group.

\begin{proposition}[Marginal form domination and its exact difference]
\label{prop:f12v80-marginal-form}
For every \(f\in L^2(\nu)\),
\[
 \begin{split}
 E_{H_\nu}(f)-E_{H_\mu}(Jf)
  =\sum_i \nu\!\left[
     \operatorname{Var}_{\mu(dx\mid v_{-i})}
       \bigl(\mu(f(v)\mid x,v_{-i})\bigr)\right]
  \ge0 .
 \end{split}
 \tag{84.3}
\]
In particular
\[
                         H_\nu\succeq J^*H_\mu J.
 \tag{84.4}
\]
\end{proposition}
\begin{proof}
Given \(v_{-i}\), apply the conditional
variance decomposition to the joint
pair \((x,v_i)\):
\[
 \begin{split}
 \operatorname{Var}_{\nu(dv_i\mid v_{-i})}f
 ={}&\mu\!\left(
       \operatorname{Var}(f\mid x,v_{-i})
                             \mid v_{-i}\right)\\
   &+\operatorname{Var}_{\mu(dx\mid v_{-i})}
               \bigl(\mu(f\mid x,v_{-i})\bigr).
 \end{split}
 \tag{84.5}
\]
The conditioning law of \(x\) here is
\(\mu(dx\mid v_{-i})\), with \(v_i\)
already integrated; it is not
\(\mu(dx\mid v)\).
Integrate \(v_{-i}\) and sum over
retained groups. A hidden-coordinate
update leaves \(Jf\) unchanged,
so all its joint form contributions
vanish. This gives (84.3) exactly.
The quadratic-form interpretation
is (84.4).
\end{proof}

This is an application of the
elementary conditional-variance
identity, not a new general
variance theorem.
Equivalently
\(J^*\Gamma_iJ\succeq P_i\).
The compressed joint kernel
first samples hidden variables
from their posterior given \(v\)
and then updates \(v_i\) with
those hidden variables fixed.
It is generally not the ideal
marginal heat bath \(P_i\).
Neither equality of these
kernels nor a pathwise
identification of the two
processes is asserted.

\begin{corollary}[Global and visible-cylinder spectral consequences]
\label{cor:f12v80-gap-transfer}
Let \(g_\mu,g_\nu\) denote the
centered gaps of (84.2). Then
\[
                            g_\nu\ge g_\mu.
 \tag{84.6}
\]
More locally, for a measurable
retained-variable region \(D\),
put \(\widetilde D=G^x\times D\).
For the original killed operators,
\[
                         \kappa_D^{H_\nu}
                                  \ge\kappa_{\widetilde D}^{H_\mu}.
 \tag{84.7}
\]
There is no hidden dimension,
block multiplicity, or probability
factor in either comparison.
\end{corollary}
\begin{proof}
For centered \(f\), the lift \(Jf\)
has the same mean and squared norm.
Apply the joint Poincare inequality
to \(Jf\), then use (84.3).
For (84.7), a function supported
in \(D\) lifts to a function
supported in \(\widetilde D\),
again with the same norm.
The joint killed lower bound
and (84.3) give the marginal
killed lower bound after taking
the infimum.
\end{proof}

The product region
\(\widetilde D=G^x\times D\)
is essential. A joint estimate
requiring hidden variables to
remain inside a small chart
does not apply to \(Jf\) and
cannot be transferred this way.

For V40's full joint fibre and
interface marginal, (84.6)
sharpens the upper comparison
in (44.10) to
\[
 \frac{\lambda(H_{\mathcal S})}
       {M_{\rm inc}C_\beta}
 \le\lambda(H_{\rm site})
 \le\lambda(H_{\mathcal S}).
 \tag{84.8}
\]
The inherited lower comparison
and its exponential \(C_\beta\)
are unchanged. Only the
unnecessary factor \(n_*\)
in that upper gap bound is
removed; no unrestricted gap
is thereby proved.

\subsection{The exact smooth joint lift of a fixed repair law}

Fix an actual repair \(T\) and
its exterior parameters \(\xi\)
as in V70. Define
\[
 \mathcal B(T)=\{b:J_b\cap T\ne\varnothing\},
 \qquad X_T=\bigcup_{b\in\mathcal B(T)}I_b .
 \tag{84.9}
\]
The fixed packing makes these
a finite collection of compact
blocks for a fixed repair.
Introduce the joint law on
\(G^{X_T}\times G^T\)
with density proportional to
\[
 \begin{split}
 &\exp\left[-\beta\left(
 S_{{\rm dir},T}^{\,\xi}(v_T)
       +\sum_{b\in\mathcal B(T)}S_b(x_b,v_{J_b})\right)\right]\\
 &\hspace{12mm}\cdot
 e^{-\Psi_{{\rm dir},T}^{\,\xi}(v_T)}
                  \prod_{b\in\mathcal B(T)}a_b(x_b,v_{J_b}).
 \end{split}
 \tag{84.10}
\]
Its reference measure is
product normalized Haar.
All amplitudes are the positive
smooth factors already present
in (74.1)--(74.2).

Integrating the \(x_b\)'s in
(84.10) gives exactly
\(e^{-W_{T,\beta}^{\,\xi}(v_T)}\),
up to its normalization.
Blocks not in \(\mathcal B(T)\)
have no retained dependence
on \(v_T\) and contribute
only a conditional constant.
Thus (84.10) is an exact
joint lift of the prescribed
repair law, not a replacement
measure.

The joint potential is a smooth
classical action times \(\beta\),
plus a smooth amplitude potential.
Its mixed group-Hessian row
has a bound \(C(\beta+1)\)
from fixed-stencil derivative
bounds. This statement requires
no global differentiability of
the minimized energy in (74.3).
In particular the possible
nonsmooth switching of its
minimizing branches is not
silently ignored.

For every fixed exterior,
(84.6)--(84.7) apply to this
exact lift. A joint repair
Poincare bound would therefore
be a sufficient marginal
repair bound with the same
constant. Uniformity over
all required exterior
parameters must still be
proved; (84.10) alone supplies
no such spectral bound.

\subsection{A constant interface saddle floor from pre-integration curvature}

Return to the compact interface
of V28 and the single-layer
central-pattern family of
V78 under the whole-layer
clearance (80.19).
Let \(a\) be its selected
unit direction on at most
448 retained sheet groups.
The exact joint potential has
the block decomposition
\[
            V_\beta(x,v)=V_{{\mathcal S},\beta}(v)
                           +\sum_b V_{b,\beta}(x_b,v_{J_b}),
 \tag{84.11}
\]
where \(e^{-V_{b,\beta}}=
e^{-\beta S_b}a_b\).
Its marginal potential is
\(V_{{\mathcal S},\beta}
-\sum_b\log\int e^{-V_{b,\beta}}\).

Clearance means that no block
input \(J_b\) meets the support
of \(a\). Hence every block
potential and block integral
is independent of this
direction. In local retained
coordinates, including their
coordinate Haar factors,
\[
                 \partial_a^2\Phi_{\rm joint}(x,v)
                     =\partial_a^2\Phi_{\rm marginal}(v)
                     \quad\hbox{for every compact }x.
 \tag{84.12}
\]
This is an exact directional
identity. It does not assert
equality of the full Hessians.

The selected negative curvature
depends on a fixed finite
retained-coordinate stencil.
Choose a visible inner cylinder
\(D\) and larger visible
cylinder \(D^+\) as in V77,
uniformly over the permitted
central local patterns.
Equation (84.12) transfers
the bound
\[
                    \partial_a^2\Phi_{\rm joint}
                                       \le-(3/56)\beta
                    \quad\hbox{on }G^x\times D^+
 \tag{84.13}
\]
without restricting any hidden
group. The cutoff depends
only on the retained stencil.
The inverse-metric bound
and cutoff cost are therefore
fixed constants.

Apply V77's directional
localization proof to the
unintegrated joint law.
Its mixed-Hessian row obeys
\(K_{\rm joint}\le
C_{\rm f}(\beta+1)\), by
(81.20), rather than the
quadratic bound for the
full marginal Hessian.
With \(c_0=3c_g/112>0\),
the resulting joint killed
floor on \(G^x\times D\)
is at least
\[
             \frac{c_0\beta}
                    {2(c_0\beta+8C_{\rm f}(\beta+1))}.
 \tag{84.14}
\]
All hidden groups range
over their entire compact
product in this assertion.
No joint global gap is
used in the localization
or resolvent argument.

\begin{theorem}[Improved local floor for the actual interface heat bath]
\label{thm:f12v80-interface-saddle}
Under the same fixed completion,
central-pattern, and whole-layer
clearance hypotheses used in
(81.22) and (82.13), there are
fixed \(c_{\rm s}>0\) and
\(\beta_0<\infty\) such that
the actual rate-one interface
heat bath satisfies
\[
 \begin{split}
 \kappa_D^{H_{\mathcal S}}
 &\ge\frac{c_0\beta}
              {2(c_0\beta+8C_{\rm f}(\beta+1))}\\
 &\ge c_{\rm s}:=
       \frac{c_0}{2(c_0+16C_{\rm f})}>0,
       \qquad\beta\ge\beta_0 .
 \end{split}
 \tag{84.15}
\]
The constants and cylinder sizes
are uniform over the stated
local central patterns and
sufficiently large ambient
volume.
\end{theorem}
\begin{proof}
The uniform negative direction
and local stencil are provided
by V78. Clearance gives
(84.12), so the joint
localization hypothesis
holds uniformly in hidden
variables. V77 gives (84.14).
Apply the visible-cylinder
comparison (84.7) to the
same exact joint/marginal
pair. Finally enlarge
\(\beta_0\) to include
\(\beta\ge1\) and the
cutoff threshold from V77;
then \(\beta+1\le2\beta\).
\end{proof}

This improves the previously
proved \(c/(\beta+1)\) lower
bounds for these local
interface cylinders to a
positive constant. The
quadratic marginal-Hessian
bound itself has not been
improved. Its loss is
avoided by performing the
spectral comparison before
integration and then using
the exact marginal form
domination.

The conclusion concerns
the original interface
heat bath, not the
projected sample path of
the joint Gibbs process.
Its survival interpretation
is again (81.23)--(81.24),
with equilibrium conditioned
on \(D\), or the stated
\(L^2\) initial-density
factor. It does not
give a uniform pointwise
escape bound.

\subsection{Why a hidden global gap is not a necessary substitute}

The sufficient global
comparison (84.6) must
not be read as an
equivalence without further
information.
Take two hidden bits with
weights
\[
 \mu_x(00)=\mu_x(11)\ \propto1,\qquad
 \mu_x(01)=\mu_x(10)\ \propto\epsilon,
 \quad 0<\epsilon<1,
 \tag{84.16}
\]
and an independent uniform
retained bit \(v\).
For the two hidden heat
baths the gap is
\(2\epsilon/(1+\epsilon)\):
their sum has this eigenvalue
on \(x_1+x_2-1\), and its
other nonconstant eigenvalues
are \(2/(1+\epsilon)\) and
\(2\). Tensoring with the
retained heat bath therefore
gives
\[
        g_\mu=\frac{2\epsilon}{1+\epsilon},
        \qquad g_\nu=1.
 \tag{84.17}
\]
The marginal can remain
fast while an irrelevant
hidden global mode is slow.

For the retained event
\(D=\{v=1\}\), however,
the joint killed operator
on \(G^x\times D\) is
\(H_x+\tfrac12 I\).
Its floor is exactly
\(1/2\), equal to the
marginal killed floor,
independently of \(\epsilon\).
This illustrates why the
local transfer in (84.15)
does not need a hidden
global Poincare estimate.
It is a finite example,
not an assertion about
hidden YM metastability.

\subsection{Exact checks and the next repair input}

The script
\begin{center}\small\ttfamily
scripts/verify\_ym\_f12\_v80\_marginal\_transfer.py
\end{center}
uses a coupled sixteen-state
joint law with two hidden
and two retained bits.
It checks the isometric
lift, its conditional
adjoint, marginal form
domination, and the exact
variance differences
\(1125/8624\) and
\(2125/6272\) for two
retained updates and a
specified test function.
The visible event
\(\{v_0=1\}\) has a
verified joint and
marginal killed lower
bound \(1/11\), checked
by exact weighted positive
semidefinite elimination.

The independent hidden-mode
example is checked at
\(\epsilon=1/2,1/4,1/16\),
including the exact joint
gaps \(2/3,2/5,2/17\)
and the fixed visible
killed floor \(1/2\).
Rational scaling checks
also distinguish the
linear pre-integration
bound from the quadratic
post-integration bound.
All checks pass. They
verify auxiliary finite
identities, not a full
compact repair gap or
continuum theorem.

\begin{decision}[Use the exact fixed joint lift, with the visible support retained]
The improved local interface
estimate is a genuine
rate-one bound under the
existing clearance condition.
It does not remove V79's
cofinal union obstruction
and does not complete
the full repair inequality.

For non-cleared actual
repairs, (84.10) supplies
the precise smooth compact
joint law to analyze.
The next input must control
its relevant visible
fluctuations or a
visible-cylinder estimate
uniformly over the
required exteriors and
all integrated variables.
If a joint estimate
restricts hidden variables
to a favorable chart,
it cannot be inserted
in (84.7).
Nor should a global
hidden gap be declared
necessary merely because
it would be sufficient.
The remaining problem is
this coupled repair
geometry and its
centered form estimate,
not another optimization
of the already obstructed
global avoidance target.
\end{decision}

V79 had 1284943 bytes and SHA--256
\begin{center}\scriptsize\ttfamily
1F7771A23A5F4607B9A23D9F3DC115BD7F500E7E3D6E27BDFB66994C6F5E394E.
\end{center}
Its entire body is retained
apart from the title version.
Only active V79 is replaced
by V80; archives, companions,
and previous scripts remain.

\begin{firewall}[Position after V80]
The exact marginal form
transfer and the improved
constant local interface
saddle floor are proved
with the stated hypotheses.
The full repair Poincare
estimate, cofinal auxiliary
gap, continuum construction,
and physical Yang--Mills
mass gap remain unproved.
The full objective remains
active.
\end{firewall}
\section{V81: conditional saddle coverage beyond clearance}
\label{sec:f12v81}

V80 is retained in full. Its constant local interface
floor used clearance to make a negative joint direction
valid for every hidden configuration. This note replaces
that strong requirement by an explicitly quantified
conditional-mass condition. Hidden variables remain
unrestricted in the law and in the lifted test function.
A favorable hidden chart can be used in an intermediate
localization only after paying for its conditional coverage.

The resulting theorem permits mixed hidden--visible
directions. For a fixed actual repair, it reduces a
sufficient large-\(\beta\) escape estimate to the geometry
of the entire hidden conditional ground set. This is a
new sufficient criterion, not a verification of that
geometry for all required YM exteriors.

The non-duplication review included V51's Gaussian leakage
estimate, V70's compact Laplace principle,
V73's local/global diffusion gluing, and V77--V80's
saddle and marginal comparisons were inspected.
The ground-energy bounds in f11 V78 and the distinct
fixed-representation Wilson saddle expansion in f6 V75
do not prove the conditional-coverage assertion below.
V80 performed the companion and direction reviews;
the next are V85 and V90 respectively.

\subsection{The lift is not supported in a favorable hidden chart}

Use V80's exact joint/marginal pair
\(\mu(dx,dv)=\nu(dv)\pi(dx\mid v)\) and isometry
\(Jf(x,v)=f(v)\). Denote the joint rate-one group
heat bath by \(H_\mu\), its weighted diffusion by
\(L\), and the actual marginal heat bath by \(H_\nu\).
The comparisons already proved are
\[
 H_\nu\succeq J^*H_\mu J
       \succeq J^*L(L+K)^{-1}J,\qquad K>0 .
 \tag{85.1}
\]
The second inequality requires the joint mixed-Hessian
row bound of V77. It assumes no positive global gap of
\(L\) or \(H_\mu\).

Let \(D\) be a visible set and \(\mathcal G\) a
measurable joint set, not necessarily a visible cylinder.
Suppose
\[
 \begin{split}
 \pi(\mathcal G_v\mid v)&\ge\theta>0
       \quad\hbox{for \(\nu\)-almost every }v\in D,\\
 \mathcal G_v&=\{x:(x,v)\in\mathcal G\}.
 \end{split}
 \tag{85.2}
\]
Then every \(f=1_Df\) satisfies
\[
          \|1_{\mathcal G}Jf\|_\mu^2
                           \ge\theta\|f\|_\nu^2 .
 \tag{85.3}
\]
This follows directly by disintegration of \(|f(v)|^2\).
Unconditional small mass of \(\mathcal G^c\) does not
imply (85.2); a test function may concentrate on the
visible configurations having poor conditional coverage.

\begin{proposition}[Conditional-coverage spectral transfer]
\label{prop:f12v81-coverage}
Suppose (85.1)--(85.2) hold and, for
\(\Pi_t=1_{[0,t]}(L)\) with \(t>0\),
\[
                  \|1_{\mathcal G}\Pi_t\|^2\le r<\theta.
 \tag{85.4}
\]
The original marginal heat bath killed on leaving \(D\)
has floor
\[
 \kappa_D^{H_\nu}\ge
       \frac{t}{t+K}(\sqrt{\theta}-\sqrt r)^2 .
 \tag{85.5}
\]
\end{proposition}
\begin{proof}
Put \(F=Jf\) for \(f=1_Df\), and decompose \(F\)
using the global spectral projection \(\Pi_t\).
Equations (85.3)--(85.4) and the triangle inequality give
\[
 \begin{split}
 \sqrt{\theta}\|F\|_\mu
 &\le\|1_{\mathcal G}F\|_\mu\\
 &\le \sqrt r\,\|\Pi_t F\|_\mu
                   +\|(I-\Pi_t)F\|_\mu\\
 &\le \sqrt r\,\|F\|_\mu+\|(I-\Pi_t)F\|_\mu .
 \end{split}
 \tag{85.6}
\]
On the complementary spectral subspace,
\(L(L+K)^{-1}\succeq t/(t+K)\).
Consequently (85.1) and (85.6) prove (85.5)
for every supported \(L^2\) function. Taking the
Rayleigh infimum proves the assertion. Neither the
cutoff nor \(J\) has been commuted with \(\Pi_t\).
\end{proof}

This is not the invalid substitution
\(\| \Pi_tJf\|^2\le r\|f\|^2\): \(Jf\) need
not be supported in \(\mathcal G\).
The explicit conditional-coverage loss repairs
precisely that support failure.

\subsection{A finite collection of joint saddle charts}

Let \(\mathcal U_j\), \(1\le j\le m\), be joint
inner chart regions, with cutoffs \(0\le\chi_j\le1\),
\(\chi_j=1\) on \(\mathcal U_j\), compactly supported
in larger chart regions \(\mathcal U_j^+\).
Only charted variables are restricted. Suppose
\[
 \begin{aligned}
 E_L(h)&\ge\lambda\|h\|_\mu^2
       &&\text{if \(h\) is supported in }\mathcal U_j^+,\\
 \sup\sum_i|\nabla_i\chi_j|^2&\le M
       &&(1\le j\le m).
 \end{aligned}
 \tag{85.7}
\]
Set \(\mathcal G=\bigcup_j\mathcal U_j\).
The same product-rule calculation as V77 gives,
now with its finite multiplicity charged,
\[
 \|1_{\mathcal G}u\|_\mu^2
 \le\sum_j\|\chi_j u\|_\mu^2
 \le\frac{2m}{\lambda}E_L(u)
             +\frac{2mM}{\lambda}\|u\|_\mu^2 .
 \tag{85.8}
\]
Indeed, apply the first inequality of (85.7) to
\(\chi_j u\), use
\(|\nabla(\chi_j u)|^2\le
2\chi_j^2|\nabla u|^2+2|\nabla\chi_j|^2|u|^2\),
and sum. Closure extends the estimate from smooth
functions to the entire diffusion form domain.
In particular (85.4) holds with
\[
                         r=\frac{2m(t+M)}{\lambda}.
 \tag{85.9}
\]

\begin{corollary}[Quantified partial-hidden-coverage floor]
\label{cor:f12v81-partial}
If (85.2), (85.7), and
\(\lambda\ge16mM/\theta\) hold, then
\[
 \kappa_D^{H_\nu}\ge
       \frac{\theta^2\lambda}
                    {4(\theta\lambda+16mK)}.
 \tag{85.10}
\]
In particular, if \(\lambda=c\beta\),
\(K\le C(\beta+1)\), and \(m,M,c,C,\theta\)
are fixed, then for sufficiently large \(\beta\),
\[
 \kappa_D^{H_\nu}\ge
       \frac{\theta^2c}{4(\theta c+32mC)}>0 .
 \tag{85.11}
\]
\end{corollary}
\begin{proof}
Choose \(t=\theta\lambda/(16m)\).
Then (85.9) is at most \(\theta/4\).
Thus the squared difference in (85.5) is at least
\(\theta/4\), giving (85.10). For (85.11), also
take \(\beta\ge1\), so \(\beta+1\le2\beta\).
\end{proof}

No lower bound near one is required for the coverage:
any fixed \(\theta>0\) suffices. If coverage decreases
with \(\beta\), that dependence must instead be
retained in both the threshold and the floor.
The finite number \(m\) is not silently made
independent of the repair size or a growing volume.

\subsection{Coverage of conditional minimizers gives conditional probability}

Here is a sufficient geometric way to verify (85.2)
for the exact finite repair lift (84.10).
Write its law as
\[
 d\mu_\beta^\xi(x,v)
   =Z_{\beta,\xi}^{-1}
        a^\xi(x,v)e^{-\beta S^\xi(x,v)}\,dx\,dv ,
 \qquad X=G^{X_T},\quad V=G^T ,
 \tag{85.12}
\]
where \(dx,dv\) are normalized product Haar,
\(S^\xi=S_{{\rm dir},T}^{\,\xi}+\sum_b S_b\),
and \(a^\xi=e^{-\Psi_{{\rm dir},T}^{\,\xi}}\prod_ba_b\).
Thus this is the same smooth positive density, not a
boxed or branch-selected replacement.

For uniform statements let \(\xi\) range over a
declared compact parameter set. Assume the smooth
spatial derivatives used below are continuous in
\(\xi\) and uniformly bounded, and
\(0<a_-\le a^\xi\le a_+<\infty\).
These are fixed-repair properties of the smooth
completion, with constants allowed to depend on
the repair. Let \(D\subset V\) be closed, and define
\[
 \begin{split}
 E^\xi(v)&=\min_{x\in X}S^\xi(x,v),\\
 \mathcal M_D
 &=\{(\xi,x,v):v\in D,\ S^\xi(x,v)=E^\xi(v)\}.
 \end{split}
 \tag{85.13}
\]
Uniform continuity on the compact product proves
continuity of \(E^\xi(v)\): its change is at most
\(\sup_x|S^\xi(x,v)-S^{\xi'}(x,v')|\).
Hence \(\mathcal M_D\) is compact.
No differentiability or choice of minimizing branch
is asserted.

\begin{proposition}[Uniform mass of an open ground-set cover]
\label{prop:f12v81-ground-mass}
Let \(\mathcal G\) be an open neighborhood of
\(\mathcal M_D\) in parameter--joint space.
Unless its complement above \(D\) is empty, there
are \(\delta>0\) and \(C_0<\infty\) such that,
with \(d_X=\dim X\),
\[
 \sup_{\xi,\ v\in D}
 \pi_\beta^\xi(\mathcal G_{\xi,v}^{\,c}\mid v)
       \le C_0\beta^{d_X/2}e^{-\beta\delta},
       \qquad\beta\ge1 .
 \tag{85.14}
\]
If the complement is empty, its probability is zero.
In particular the coverage is at least \(1/2\)
for all sufficiently large \(\beta\).
\end{proposition}
\begin{proof}
The compact set above \(D\) outside \(\mathcal G\)
contains no conditional minimizer. Therefore the
continuous function \(S^\xi(x,v)-E^\xi(v)\)
has a strictly positive minimum \(\delta\) there.
The numerator of the conditional probability is
at most \(a_+e^{-\beta(E^\xi(v)+\delta)}\).

At any conditional minimizer \(x_*\), the hidden
gradient of \(S^\xi\) is zero. A uniform hidden
Hessian upper bound and Taylor's theorem in a
fixed-radius normal ball give
\[
 S^\xi(x,v)\le E^\xi(v)+B\,\operatorname{dist}_X(x,x_*)^2
       \quad\text{for }\operatorname{dist}_X(x,x_*)\le r_0 .
 \tag{85.15}
\]
Here \(\operatorname{dist}_X\) is product geodesic distance.
Compact Haar geometry gives ball volume at least \(b r^{d_X}\) for
\(r\le r_0\). Integrating over radius
\(r_0/\sqrt\beta\) bounds the conditional
normalizer from below by
\[
 a_- b r_0^{d_X}\beta^{-d_X/2}
           e^{-Br_0^2}e^{-\beta E^\xi(v)} .
 \tag{85.16}
\]
Division proves (85.14). In hidden dimension zero,
the same argument is interpreted with the one-point
hidden space and volume one.
\end{proof}

This is a finite-dimensional probability estimate,
not a Poincare estimate for the hidden conditional
law. It allows multiple or degenerate hidden
minimizers. The proof does not need their Hessians
to be positive definite.

\subsection{A noncritical-or-saddle criterion on the ground set}

All gradients and Hessians in the next theorem
refer to the classical joint action \(S^\xi\)
in its independent compact group variables.
At a critical point the sign of its intrinsic
Hessian agrees with that of the Hessian in
centered product normal coordinates.

\begin{theorem}[Fixed-repair ground-set escape criterion]
\label{thm:f12v81-ground-escape}
In (85.12)--(85.13), assume that each point of
\(\mathcal M_D\) satisfies one of the following:
\[
 \nabla_{x,v}S^\xi\ne0,
 \qquad\text{or}\qquad
 \nabla_{x,v}S^\xi=0
       \ \text{and}\ 
       \lambda_{\min}(\operatorname{Hess}_{x,v}S^\xi)<0 .
 \tag{85.17}
\]
Then there are \(\beta_0<\infty\) and \(\kappa_*>0\),
uniform in the declared compact parameter set, such that
the actual rate-one marginal repair heat bath satisfies
\[
                    \kappa_D^{H_\nu}\ge\kappa_*
                    \qquad(\beta\ge\beta_0).
 \tag{85.18}
\]
Constants can depend on the fixed repair and parameter
set. No joint global spectral gap is an assumption.
\end{theorem}
\begin{proof}
At a point with nonzero gradient choose a unit
constant chart direction \(q\) and a neighborhood
where \(|\partial_qS^\xi|\ge s>0\). In the correct
charted density write
\[
 \Phi_\beta^\xi=\beta S^\xi+\Psi^\xi,\qquad
 \Psi^\xi=-\log a^\xi-\sum_{\text{charted }i}\log j_G .
 \tag{85.19}
\]
On a smaller fixed chart the derivatives of \(\Psi^\xi\)
are bounded. If \(|\partial_q\Psi^\xi|\le B_1\),
\(|\partial_q^2 S^\xi|\le B_2\), and
\(|\partial_q^2\Psi^\xi|\le B_3\), then, for
\(\beta s\ge B_1\),
\[
 \frac{|\partial_q\Phi_\beta^\xi|^2}{4}
       -\frac{\partial_q^2\Phi_\beta^\xi}{2}
 \ge \frac{(\beta s-B_1)^2}{4}
                         -\frac{\beta B_2+B_3}{2}.
 \tag{85.20}
\]
For sufficiently large \(\beta\), this is at least
\(s^2\beta^2/16\). The ground-state identity (81.8),
with the inverse-metric lower bound, gives a local
diffusion floor of order \(\beta^2\).

At a critical point with a negative Hessian direction,
choose a unit constant chart direction \(q\) and a
neighborhood on which
\(\partial_q^2 S^\xi\le-\eta<0\).
For large \(\beta\),
\(\partial_q^2\Phi_\beta^\xi\le-\eta\beta/2\).
The same identity gives a local diffusion floor
at least \(c_g\eta\beta/4\). The direction may
involve both hidden and visible variables.

Choose nested chart neighborhoods and cutoffs at
these points. Compactness of \(\mathcal M_D\)
provides a finite collection of inner neighborhoods
covering it. Parameter neighborhoods may be chosen
alongside the spatial ones; at a given parameter use
the cutoffs whose parameter neighborhoods contain it.
There are at most a fixed \(m\) cutoffs, with uniform
spatial derivative bound \(M\). Taking minima of the
finitely many positive constants and increasing the
threshold, (85.7) holds with \(\lambda=c\beta\).

The union of the inner parameter--joint neighborhoods
is an open ground-set cover. Proposition
\ref{prop:f12v81-ground-mass} therefore gives
\(\theta=1/2\) uniformly for large \(\beta\).
The smooth fixed-repair joint potential obeys
\(K\le C(\beta+1)\), as in V80.
Corollary \ref{cor:f12v81-partial} proves (85.18).
All operators are those of the full compact law.
The chart cutoffs are used only in (85.8);
the lifted test functions are never chart-conditioned.
\end{proof}

At a conditional hidden minimizer the hidden gradient
vanishes. Thus the first alternative of (85.17)
is detected by a nonzero visible force. A negative
direction in the second alternative cannot be purely
hidden: the hidden Hessian at a hidden global minimum
is nonnegative. Mixed directions are nevertheless
allowed, and can be decisive.

\subsection{Relaxed hidden response supplies a mixed negative direction}

At a conditional minimizer use local coordinates
\((x,v)\) and suppose the hidden Hessian \(A\) is
positive definite. Write the classical joint Hessian
as
\[
       \begin{pmatrix}A&B\\B^T&C\end{pmatrix},
       \qquad A\succeq a_0 I,\quad a_0>0 .
 \tag{85.21}
\]
There is a local stationary hidden branch \(x(v)\)
by the implicit function theorem.
Differentiate its stationarity equation to obtain
\[
 Dx(v)=-A^{-1}B,\qquad
 D^2[S(x(v),v)]=C-B^TA^{-1}B
 \tag{85.22}
\]
at the base point. In the second formula the terms
containing \(D^2x(v)\) vanish because \(\partial_xS=0\).
This is a local branch identity, not a statement
that this branch is the global minimizer nearby.

If a unit visible vector \(a\) has
\(a^T(C-B^TA^{-1}B)a\le-\sigma<0\), set
\(z=(-A^{-1}Ba,a)\). Direct multiplication gives
\[
 \frac{z^T\operatorname{Hess}S\,z}{\|z\|^2}
 \le-\frac{\sigma}{1+\|A^{-1}B\|^2}.
 \tag{85.23}
\]
Indeed the numerator equals the Schur quadratic form,
and \(1\le\|z\|^2\le1+\|A^{-1}B\|^2\).
At a joint critical point this verifies the negative
alternative of (85.17). It need not be visible in
the unrelaxed visible Hessian \(C\).

For example, the rational matrices
\[
 A=\begin{pmatrix}2&1\\1&3\end{pmatrix},\qquad
 B=\begin{pmatrix}1\\2\end{pmatrix},\qquad C=\tfrac12
 \tag{85.24}
\]
give \(A^{-1}B=(1/5,3/5)^T\),
\(C-B^TA^{-1}B=-9/10\), and normalized joint
curvature \(-9/14\). Both the pure hidden block
and the pure visible curvature are positive.
The instability is in the coupled response.
This is an exact local quadratic diagnostic, not
an identification of the YM block Hessian.

\subsection{Why one selected hidden minimum is not enough}

The all-minimizer cover in (85.14) is sufficient,
not necessary: (85.2) permits any directly proved
positive conditional coverage. But a neighborhood
of just one hidden minimizer does not automatically
have such coverage.

A compact elementary diagnostic uses a hidden
disjoint union of two copies of \([-1,1]\), with
equal reference weights, and visible \(v\) near zero:
\[
 S_1(x,v)=x^2-v^2,\qquad S_2(x,v)=x^4+v^2 .
 \tag{85.25}
\]
At \(v=0\) both hidden minima have action zero.
The first is a joint saddle, whereas the second
is a degenerate joint local minimum. The first
component's conditional probability is
\[
 \frac{\int_{-1}^1 e^{-\beta x^2}\,dx}
      {\int_{-1}^1 e^{-\beta x^2}\,dx+
       \int_{-1}^1 e^{-\beta x^4}\,dx}
             \le \frac{3e}{2}\beta^{-1/4},
             \qquad\beta\ge1 .
 \tag{85.26}
\]
To verify the bound, rescale the quadratic integral:
it is at most \(3\beta^{-1/2}\), since
\(\int_{\mathbb R}e^{-t^2}dt<3\), also obtained
by splitting at \(|t|=1\) and using
\(e^{-t^2}\le |t|e^{-t^2}\) on the tails.
The quartic integral is at least
\(2e^{-1}\beta^{-1/4}\) by restricting to
\(|x|\le\beta^{-1/4}\).
Division proves (85.26).

Thus conditional mass of the selected saddle tends
to zero despite its having minimum hidden energy.
This diagnostic has disconnected interval hidden
space and is not an actual compact-group YM model.
Its purpose is solely to rule out the unqualified
inference from existence of one minimizer to
uniform positive conditional mass.

\subsection{The exact remaining geometric obligation}

For a fixed target \(A\subset V\), with closed
\(D=A^c\), Theorem \ref{thm:f12v81-ground-escape}
would control the entire complement, not just
an intermediate energy band. Together with an
independently established original-clock within-target
conductance
\[
 E_{H_\nu}(f)\ge
       \gamma_{\rm in}\nu(A)
                      \operatorname{Var}_{\nu(\cdot\mid A)}(f),
 \tag{85.27}
\]
the gluing inequality already proved in V74 gives,
with \(n=|T|\),
\[
       h_T\ge
       \frac{\kappa_*\gamma_{\rm in}}
                    {\kappa_*+2n+2\gamma_{\rm in}}.
 \tag{85.28}
\]
This last implication is recorded, not proved again.
It does not supply (85.17) or (85.27).

The ground-set hypothesis is a substantive classical
condition, with a simple necessary test. If the joint
action has a global minimum whose visible component
lies in \(D\), that point belongs to \(\mathcal M_D\),
has zero joint gradient and nonnegative joint Hessian.
It violates (85.17). Consequently a fixed section-centered
target cannot be imposed on arbitrary exterior data
without first showing that their joint global
minimizers lie inside it. Section interpolation alone
is not that proof; f11 V78 already distinguished
a section from a constrained minimizer.

For the actual non-cleared repairs, the next calculation
is therefore to locate the conditional hidden ground
set and analyze its stationary joint points, allowing
the response directions in (85.23). V78's central
sheet calculation controls specified unintegrated
central configurations. It does not assert that all
hidden conditional minima over the same visible
configuration are those central configurations.
That missing identification is not replaced by a
Gaussian branch approximation.

Nor does this theorem bypass V79's cofinal union
obstruction. Its finite cover and compactness constants
are fixed-repair constants. Applying it to an ambient
union growing with volume would require uniform
coverage, curvature, cutoff multiplicity, and energy
separation; none follows from the present argument.

\subsection{Exact verification and continuation record}

The script
\begin{center}\small\ttfamily
scripts/verify\_ym\_f12\_v81\_conditional\_saddle\_cover.py
\end{center}
checks the operator step on a coupled sixteen-state
joint law with one hidden and three visible bits.
For a visible set \(D\) with two states and a
strictly smaller joint cover, it finds
\[
       \theta=\tfrac23,\qquad r=\tfrac1{12},
       \qquad \kappa_{\rm bound}=\tfrac1{30},
       \qquad \kappa_D^{H_\nu}=\tfrac32 .
 \tag{85.29}
\]
Here the auxiliary operator is \(L=H_\mu\), with
\(K=1\) and \(t=1/4\); it is a finite operator
test of the transfer, not a discretization of the
compact diffusion. Exact weighted positive
semidefinite checks verify the resolvent compression
and the supported lower bounds. The cutoff and
low projection do not commute. For \(f=1_D\),
the squared low-projection ratio is \(1/8>r\),
explicitly detecting the invalid support shortcut
excluded after (85.6).

The script also checks (85.24), including the
relaxed minimizing response and curvature \(-9/14\).
With \(\theta=1/2,c=1/4,C=2,m=3,M=2\), it checks
the sufficient threshold \(\beta\ge768\) and
uniform floor \(1/12296\) in (85.11), at
\(\beta=768,1024,4096\). All exact checks pass.
They do not verify the continuum claim or the
actual YM ground-set hypothesis.

V80 had 1300714 bytes and SHA--256
\begin{center}\scriptsize\ttfamily
70E20D5B2FF8C5EBBE02301B3F75226E2F9ED2BC9167625D5FFF1EDB7323DDFF.
\end{center}
Its full body is retained apart from the title version.
Only active V80 is replaced by V81. Archives,
companions, and all earlier scripts remain.

\begin{firewall}[Position after V81]
The conditional-coverage transfer, its finite-cover
constant, and the fixed-repair conditional-ground-set
criterion are proved. They permit favorable hidden
charts and mixed directions without changing the
underlying compact law or requiring a hidden global gap.
The noncritical-or-saddle condition is not yet
verified for all required actual YM repairs and
exteriors. The unrestricted repair estimate,
cofinal auxiliary gap, continuum construction,
and physical Yang--Mills mass gap remain unproved.
The full objective remains active.
\end{firewall}
\section{V82: a 56-coordinate hidden descent and exclusion of the central block lift}
\label{sec:f12v82}

V81 gave a sufficient conditional-ground-set criterion,
but did not locate the actual hidden minimizers.
Here we test the central sheet lift against a genuine
compact integrated block. A nonuniform rotation in a
two-cell transverse cube already decreases the action.
When these variables are hidden, the central lift is
not a hidden conditional minimizer. A neighborhood
of it has exponentially small conditional mass.
Consequently it cannot simply be used as the
positive-coverage saddle chart in V81.

This calculation is on the actual \(SU(2)\) smooth
completion, with flat coarse field \(C=I\).
It is not a model with independently chosen pivots.
The new parts are the two-cell weighted direction,
its finite action decrease, its placement in a
canonical hidden block, the exact Jacobian cancellation,
and a necessary boundary-to-volume condition for
central hidden minimizers.

For non-duplication, V76's four-cell uniform rotation,
V78's arbitrary central-pattern cut formula, and V31's
all-exterior core-tail argument were inspected.
The last already supplies a different concentration
mechanism on sufficiently large guarded cores.
We do not re-prove that tail or count concentration
alone as a spectral gap. The present calculation
identifies an explicit wrong candidate for the
conditional ground set and a small hidden variation
which rules it out. The next companion audit remains
V85 and the next direction review V90.

\subsection{The exact sheet Hessian with the exterior frozen}

Use the full one-direction sheet of V76:
the nonpivot direction-\(2\) links at \(x_2=1\)
are \(-I\), while every other physical link is \(I\).
The transverse lattice is
\(\Lambda=(\mathbb Z/(2M)\mathbb Z)^3\).
Its pivot sites are
\(P=\{x:\text{all three coordinates are even}\}\);
write \(F=\Lambda\setminus P\).
For a non-wrapping transverse coarse cube of side
\(q<M\), let
\[
 H_q=F\cap\{0,\ldots,2q-1\}^3,\qquad |H_q|=7q^3 .
 \tag{86.1}
\]
All sheet variables outside \(H_q\) remain frozen.
Set \(u_x=0\) outside \(H_q\) and rotate its links to
\(-\exp(tu_xX)\), where
\(\|X\|_{\rm HS}=1\).
For \(|tu_x|\le1/4\), the changed chronological
paths stay near \(-I\), outside the cutoff support.
V18's exact pivot equation therefore leaves every
pivot at \(I\), as in V76.

Let \(r_x\) be the number of pivot neighbors of \(x\).
It is two when \(x\) has exactly one odd coordinate
and zero otherwise. The scalar sheet Hessian is
\[
 \begin{split}
 Q_q(u)&=\sum_{\{x,y\}\subset F,\ x\sim y}(u_x-u_y)^2
                                  -\sum_{x\in H_q}r_xu_x^2\\
       &=u^TK_qu,\qquad
 (K_q)_{xx}=6-2r_x,\quad
 (K_q)_{xy}=-1\quad(x,y\in H_q,\ x\sim y).
 \end{split}
 \tag{86.2}
\]
Off-diagonal entries not displayed are zero.
The first line is the already derived physical form
(80.14), with the frozen sheet variables set to zero
in the tangent vector. Thus it is the correct principal
form, including the positive edges to frozen negative
links and the negative edges to pivot holes.

For one coarse cell, the graph on the seven free
vertices of \(\{0,1\}^3\) is connected. Its internal
degrees are two at parity weight one and three
at parity weights two and three. Consequently
\[
 Q_1(u)=
   \sum_{\substack{x\sim y\\x,y\in H_1}}(u_x-u_y)^2
       +3\sum_{\substack{x\in H_1\\
                    \text{at least two odd coordinates}}}u_x^2>0
       \quad(u\ne0).
 \tag{86.3}
\]
Indeed equality would first make all values equal
by connectedness and then set that common value
to zero. This statement concerns sheet directions,
not variations of all the other block coordinates.

\subsection{An exact weighted direction in a two-cell cube}

Now \(q=2\). For \(j=0,1,2,3\) put
\(w_j=(1,2,2,1)_j\), and for \(x\in H_2\) define
\[
 p_x=w_{x_1}w_{x_2}w_{x_3},\qquad
 k(x)=\#\{a:x_a\text{ is odd}\},\qquad
 f_k(x)=1_{\{k(x)=k\}}p_x,\quad k=1,2,3 .
 \tag{86.4}
\]
The three disjointly supported basis vectors have
squared norms \(375,375,125\). To check this,
the even and odd sums of \(w_j^2\) both equal five;
multiply three such sums and count the choices
of the odd coordinates.

Their exact compressed quadratic matrix is
\[
       (f_i^TK_2f_j)_{i,j=1}^3
       =
       \begin{pmatrix}
         750&-1200&0\\
         -1200&2250&-600\\
         0&-600&750
       \end{pmatrix}.
 \tag{86.5}
\]
Here the diagonal values follow from \(6-2r_x\),
which is two on parity weight one and six on
the other two classes. A neighboring pair changes
parity weight by one. The one-dimensional edge
weight sum is
\(w_0w_1+w_1w_2+w_2w_3=8\).
For weights one and two there are three choices
of edge direction and two choices for the odd
coordinate among the remaining axes. Their squared
weight sums give \(3\cdot2\cdot8\cdot25=1200\).
For weights two and three the corresponding count
is \(3\cdot8\cdot25=600\).
The off-diagonal signs are negative, proving (86.5).

Take the explicit direction
\[
                  u=\frac{6f_1+4f_2+3f_3}{48}.
 \tag{86.6}
\]
It has \(1/16\le u_x\le1\) on all 56 selected sites.
Matrix multiplication and the disjoint norm sums give
\[
 \begin{split}
 (6,4,3)(f_i^TK_2f_j)(6,4,3)^T&=-2250,\\
 Q_2(u)&=-\frac{125}{128},\qquad
 \|u\|^2=\frac{6875}{768},\\
 \frac{Q_2(u)}{\|u\|^2}&=-\frac6{55}.
 \end{split}
 \tag{86.7}
\]

\begin{proposition}[A smaller actual sheet instability]
\label{prop:f12v82-weighted}
For \(M>2\), the full central sheet has a negative
completed-fibre direction supported on 56 free groups,
with normalized Hessian value \(-6/55\).
Among the integer coarse cubes in (86.1), side two
is the smallest permitting a negative direction
in this frozen-exterior sheet sector.
\end{proposition}
\begin{proof}
All first derivatives at the central field vanish
by the central-plaquette trace identity used in V78.
The exact pivot argument preceding (86.2) makes
its Hessian the actual completed-fibre Hessian.
Equations (86.6)--(86.7) prove negativity.
Equation (86.3) excludes the only smaller positive
integer side. There is no claim of minimum support
among all shapes or among all full-block directions.
\end{proof}

The earlier side-four statement is not contradicted:
V78 proved sharpness only for the specified
constant rotations on unions of coarse owner cells.
In fact \(Q_2(\mathbf1)=24>0\). The weights in
(86.6) are essential to this particular improvement.
The arbitrary-pattern theorem of V78 still has its
448-coordinate bound; the 56-coordinate result here
is for the full local negative-sheet pattern.

\subsection{A finite decrease, not only a Hessian sign}

Put \(s(z)=2(1-\cos(z/\sqrt2))\). With the actual
pivots fixed as above, the exact action change is
\[
 \Delta S(t)=
   \sum_{\{x,y\}\subset F,\ x\sim y}s(t(u_x-u_y))
                           -\sum_{x\in H_2}r_xs(tu_x).
 \tag{86.8}
\]
For the arguments used here Taylor's inequalities give
\[
              \frac{z^2}{2}-\frac{z^4}{48}
                         \le s(z)\le\frac{z^2}{2}.
 \tag{86.9}
\]
Using the upper bound on the positive terms and
the lower bound on the subtracted terms gives
\[
 \Delta S(t)\le
       \frac{t^2}{2}Q_2(u)
                    +\frac{t^4}{48}\sum_xr_xu_x^4 .
 \tag{86.10}
\]
There are \(3\cdot2^3=24\) parity-one sites,
each with two pivot neighbors, so
\(\sum r_xu_x^4\le48\).
At \(t_0=1/4\), equations (86.7) and (86.10) prove
\[
            \Delta S(t_0)\le-\delta_{56},
                  \qquad \delta_{56}:=\frac{109}{4096}>0.
 \tag{86.11}
\]
All angle magnitudes are at most \(1/4\), as required
by the pivot argument.

For an independently checkable stronger numerical
margin, the even and odd fourth-power sums of the
one-dimensional weights are both seventeen. Hence
\[
 \sum_xr_xu_x^4
       =2(6/48)^4\,3\cdot17^3=\frac{14739}{2048},
 \qquad
 -\Delta S(t_0)\ge\frac{251087}{8388608}>\delta_{56}.
 \tag{86.12}
\]
We use the simpler smaller margin in what follows.
No asymptotic-in-\(t\) remainder is left unbounded.

\subsection{The direction lies in an actual compact hidden block}

Let \(P_b\) be a canonical parent region as in V21,
embedded in V28's separated packing, and suppose
its free set \(I_b\) contains the 56 boundary
variables of (86.1) with \(q=2\).
This is a concrete geometric condition, not clearance.
It holds, for example, when the parent contains
the coarse cube \(\{0,1\}^4\) in the displayed
coordinates. Every selected sheet variable is a
nonpivot boundary variable on a direction-\(2\)
coarse edge based at \((r_1,r_2,0,r_3)\),
with each transverse \(r_a\in\{0,1\}\);
that edge has an endpoint in the parent.
The definition (25.1) therefore puts it in \(I_b\).
For the parent exactly equal to this cube,
\[
           |P_b|=16,\quad |E(P_b)|=96,\quad |I_b|=944 .
 \tag{86.13}
\]
These are checks of inclusion in the original
canonical update, not the definition of a new
56-group block kernel.

Condition the actual interface inputs of this block
to their values in the central sheet. Denote these
values by \(v^*\), and the corresponding central
hidden coordinates by \(x_b^*\).
All other hidden block variables are held at their
central values only to specify this candidate point,
not in the conditional measure. Let
\[
                E_b(v^*)=\min_{x_b\in G^{I_b}}S_b(x_b,v^*).
 \tag{86.14}
\]
Terms assigned to \(S_b\) differ from the full action
by constants under this hidden variation. Applying
the rotation entirely inside \(I_b\) gives
\[
                  S_b(x_b^*,v^*)-E_b(v^*)
                                    \ge\delta_{56}.
 \tag{86.15}
\]
In particular the proposed central hidden point
is not even a hidden local minimum. This conclusion
does not require knowing which other compact
configuration attains \(E_b(v^*)\).

An arbitrary unspecified packing is not asserted
to contain this geometry. The result applies to
the displayed canonical cube or any already chosen
parent containing these sheet coordinates.
Separated translates of such fixed cubes are
compatible with V28's packing conditions.

\subsection{Exact amplitude cancellation and exponential conditional exclusion}

The finite decrease extends to an actual neighborhood,
with a useful cancellation. Define a global
product-Haar-preserving map on the hidden variables by
\[
       (\mathcal R x)_i=x_i\exp(t_0u_iX)
           \quad\text{on the 56 selected boundary groups},
       \qquad
       (\mathcal R x)_i=x_i\quad\text{otherwise}.
 \tag{86.16}
\]
The multipliers are fixed, so this is an invertible
product of right translations. At the central point
it is exactly the rotation used in (86.11).

There are only eight coarse pivots whose factors
can depend on these boundary variables, one for
each of their owner edges. At both the reference
point and its image under \(\mathcal R\), all
fourteen nonroot chronological products for each
of these edges are near \(-I\), not near \(I\).
Consequently there are fixed neighborhoods of the
relevant independent coordinates on which, for
pivot input \(B\) near \(I\),
\[
 g_\rho(B^{-1}P_j)=0\quad(1\le j\le14),
        \qquad F(B)=B,\qquad j_e=1 .
 \tag{86.17}
\]
The same statements hold at both \(x\) and
\(\mathcal R x\). Global pivot uniqueness identifies
this local solution with the actual solved pivot.
The derivative of the pivot map is the identity,
so its inverse Haar Jacobian is exactly one.

A nonpivot boundary variable enters only its own
edge's pivot factor, as proved in V18. All other
pivot factors are unchanged by \(\mathcal R\).
Thus the full positive amplitude in the exact
block conditional satisfies
\[
                         a_b(\mathcal R x,v)=a_b(x,v)
 \tag{86.18}
\]
on a sufficiently small neighborhood. In particular
no unknown determinant prefactor or chart-Haar
ratio is discarded. The conditional reference
measure remains product Haar, and (86.16) preserves it.

Take a finite input closure of the action and pivot
factors changed by \(\mathcal R\). Smoothness and
finite range imply that it involves finitely many
hidden variables and retained inputs. By continuity
of the action difference in these variables, choose
a fixed visible neighborhood \(\mathcal V\) of
the required components of \(v^*\), and a hidden
chart cylinder \(\mathcal U\) about the corresponding
components of \(x_b^*\), such that
\[
 \begin{split}
 S_b(\mathcal R x,v)&\le S_b(x,v)-\delta_{56}/2,\\
 a_b(\mathcal R x,v)&=a_b(x,v),
        \qquad x\in\mathcal U,\quad v\in\mathcal V .
 \end{split}
 \tag{86.19}
\]
All hidden variables outside that finite closure
remain unrestricted. The coarse field remains \(I\).
The radius and closure are fixed independently
of \(\beta\) and ambient volume. For a larger
parent, the unused hidden variables do not enter
the changed local terms.

\begin{theorem}[The central hidden neighborhood is exponentially underweighted]
\label{thm:f12v82-hidden-exclusion}
For the actual full compact block conditional,
at every \(v\in\mathcal V\) and every \(\beta\ge0\),
\[
       \pi_{b,\beta}(\mathcal U\mid v)
              \le \exp\left(-\frac{109}{8192}\beta\right).
 \tag{86.20}
\]
Every \(x\in\mathcal U\) also satisfies
\[
              S_b(x,v)-\min_z S_b(z,v)
                                      \ge\frac{109}{8192}.
 \tag{86.21}
\]
\end{theorem}
\begin{proof}
The second assertion follows by comparing \(x\)
with the feasible hidden competitor \(\mathcal R x\)
in (86.19).
For the first, let \(Z_b(v)\) be the unrestricted
compact normalizer. Change variables by the
Haar-preserving map on the image of \(\mathcal U\):
\[
 \begin{split}
 Z_b(v)
 &\ge \int_{\mathcal R\mathcal U}
                    a_b(z,v)e^{-\beta S_b(z,v)}\,dz\\
 &=\int_{\mathcal U}
                    a_b(\mathcal R x,v)e^{-\beta S_b(\mathcal R x,v)}\,dx\\
 &\ge e^{\beta\delta_{56}/2}
             \int_{\mathcal U}a_b(x,v)e^{-\beta S_b(x,v)}\,dx .
 \end{split}
 \tag{86.22}
\]
Divide by \(Z_b(v)\) and use (86.11).
There is no replacement of the normalizer by
one favorable point, and no hidden-box conditioning.
\end{proof}

The conditional probability upper bound remains valid
after disintegrating the full exact joint law by its
visible interface: V28 identifies its hidden block
conditional with the same \(\pi_{b,\beta}\).
Therefore any proposed V81 cover contained in
this central hidden cylinder has conditional
coverage tending to zero uniformly over
\(\mathcal V\). It cannot be assigned a
\(\beta\)-independent positive \(\theta\).
This is a verified obstruction to that particular
choice of cover, not an obstruction to V81's theorem
or to a different cover of the true ground set.

\subsection{A boundary-to-volume constraint on central hidden minimizers}

There is also a necessary condition for other central
candidates in the same one-direction family.
Let \(H\subset F\) be the hidden sheet variables
of a fixed parent, \(A\subset F\) the negative
sites of a central candidate, and write
\[
                    A_{\rm h}=A\cap H,\qquad
                    A_{\rm e}=A\setminus H .
 \tag{86.23}
\]
All non-sheet free groups are at \(I\) for this
central-family statement. If the candidate is a
local minimum under the full hidden block
variation, then it is a local minimum under
the rotations of any \(B\subset A_{\rm h}\).
V78's exact physical formula therefore forces
\[
           D_A(B)=e(B,A\setminus B)
                      -e(B,\Lambda\setminus A)\ge0 .
 \tag{86.24}
\]
For \(B=A_{\rm h}\), writing \(b(A_{\rm h})\)
for its full transverse edge boundary, this becomes
\[
  0\le D_A(A_{\rm h})
      =-b(A_{\rm h})+2e(A_{\rm h},A_{\rm e}) .
 \tag{86.25}
\]
The two pinned-cube decompositions already proved
in (82.3) give \(b(A_{\rm h})\ge(6/7)|A_{\rm h}|\).
Thus every such hidden local minimum necessarily obeys
\[
                   |A_{\rm h}|
                     \le\frac73 e(A_{\rm h},A_{\rm e}) .
 \tag{86.26}
\]
No independence or low-temperature approximation
is used in this classical constraint.

For \(H=H_q\), there are exactly \(18q^2\) edges
from \(H_q\) to free sites outside it, by the same
face count used in V76. Hence
\[
                   |A_{\rm h}|\le42q^2,\qquad
                   \frac{|A_{\rm h}|}{|H_q|}\le\frac6q .
 \tag{86.27}
\]
The bound can be trivial for small \(q\); it is a
necessary surface-order constraint for large cubes,
not a sufficient test of stability. The weighted
side-two calculation already finds an instability
which rotating the entire \(H_2\) uniformly misses.

Every nonempty connected component of the hidden
negative set must also contact an exterior negative
site. If a component \(B\) did not, then
\(e(B,A\setminus B)=0\), so
\(D_A(B)=-b(B)\le-(6/7)|B|<0\), contradicting
(86.24). In particular with no exterior negative
sheet sites, no nonempty hidden negative central
pattern is a hidden local minimum in this family.

These statements do not classify noncentral
minimizers, patterns with other link directions,
or mixed noncentral interfaces. They cannot be
used to replace all integrated variables by their
central or section values.

\subsection{Verification, next calculation, and retained scope}

The new script is
\begin{center}\small\ttfamily
scripts/verify\_ym\_f12\_v82\_hidden\_sheet\_exclusion.py .
\end{center}
It verifies (86.5)--(86.12) in rational arithmetic.
On a fine torus of side six it independently
enumerates physical plaquettes and all chronological
paths. The sheet action is 648; the selected
56-coordinate physical Hessian is \(-125/128\);
48 changed negative plaquettes give the fourth-power
sum \(14739/2048\). Exactly 112 chronological paths
contain a moved link, each containing just one,
with negative central sign and angle magnitude
at most the declared \(1/4\).
It checks that all 56 selected variables belong
to the canonical 944-group parent in (86.13).
This small torus is a stencil check, not itself
a claim that this parent is a separated periodic
packing on that torus.

The one-cell principal matrix is also checked
exactly positive definite. Finally, for five
cube sizes and eleven patterns per size, the
script checks the signed-boundary identity,
the pinned-hole bound, and the necessary
surface-order implication (86.26)--(86.27).
All checks pass. Exploratory floating-point
eigenvectors were used to find a promising
shape; no floating-point eigenvalue appears
in the proof or final certificate.

\begin{decision}[Do not center the hidden conditional at an excluded saddle]
The first actual ground-set test has a definite
outcome: when this negative sheet crosses the
displayed hidden parent, its central hidden lift
lies above the hidden ground energy by an explicit
amount. Its local chart is exponentially
underweighted in the actual conditional law.
It is therefore the wrong reference chart for
positive conditional coverage in V81.

The next question is the boundary-driven hidden
minimizer after this descent. The classical
constraint (86.26) restricts any remaining central
candidate, but a genuine minimizer may be noncentral.
Its stationary equations and its mixed
hidden--visible Hessian, not another expansion
at \(x_b^*\), are the appropriate inputs for
V81's criterion. The present result supplies
neither that minimizer classification nor the
original-clock within-target conductance.
\end{decision}

V81 had 1321122 bytes and SHA--256
\begin{center}\scriptsize\ttfamily
BD0F7DBDA274926053512FD60EF5073263B9A67F3EE6EBE724B06476B18B368E.
\end{center}
Its entire body is preserved apart from the title
version. Only active V81 is replaced by V82.
Archives, companions, and earlier scripts remain.

\begin{firewall}[Position after V82]
Proved for the specified actual \(SU(2)\) completion:
a 56-coordinate negative sheet direction, a finite
action decrease, its hidden-block realization,
an exponential conditional exclusion with exact
amplitude cancellation, and necessary central
hidden-minimizer boundary constraints.
Not proved: classification of the noncentral
conditional ground set, the unrestricted repair
Poincare estimate, the cofinal auxiliary gap,
continuum construction, or the physical
Yang--Mills mass gap. The full objective remains active.
\end{firewall}
\section{V83: central-cut certification, noncentral conditional minima, and singular hidden response}
\label{sec:f12v83}

V82 excluded the full central sheet lift from the
ground set of a displayed hidden parent. It did not
exclude all other central candidates. We now solve
that central-candidate question exactly on its
56-group sheet slice. An integer routing identity
covers all \(2^{56}\) central replacements.
Every minimizer of this actual compact slice is
noncentral, and conjugation symmetry forces its
hidden Hessian to be singular.

The scope distinction matters. This slice is a
genuine conditional of the completed fibre, obtained
by fixing all other free groups to the displayed
central field. It is not the 944-group canonical
parent conditional. Its minimizers are not declared
to be minimizers of that larger parent. The result
does show, within the actual compact action, why
nondegenerate hidden Laplace branches cannot be
assumed for every local conditional.

We then replace V81's optional positive-definite
hidden Schur calculation by a kernel-compatible
response criterion. It applies to an arbitrary
hidden conditional minimizer, including singular
ones. V81's coverage theorem itself is unchanged:
it never required the hidden Hessian to be invertible.
The next companion audit remains V85; the next
direction review remains V90.

\subsection{A central-candidate problem on the exact compact slice}

Retain \(H_2\), \(K_2\), and the central exterior
from V82. Let \(S_{\rm sl}(x)\) be the actual
completed-fibre Wilson action when the 56 groups
in \(H_2\) vary over \(SU(2)^{56}\), all other
free groups are fixed at \(Y^*\), and \(C=I\).
The corresponding conditional density is the
positive smooth restriction of
\(j(Y,I)e^{-\beta S_0(Y,I)}\,dY_{H_2}\).
There is no frozen-pivot replacement in this
definition; the pivots are always solved by V18.

At any central configuration of these 56 groups,
every chronological path product is \(I\) or \(-I\).
The exact pivot equation at \(B=I\) and its global
uniqueness give all pivots equal to \(I\).
A central replacement is consequently specified
by the subset \(B\subset H_2\) of negative links
erased to \(+I\). Directly counting its central
plaquettes gives
\[
       S_{\rm sl}(Y^{*,B})-S_{\rm sl}(Y^*)
                     =4Q_2(1_B)=4\,1_B^TK_2\,1_B .
 \tag{87.1}
\]
Indeed each free--free boundary edge of \(B\)
creates one negative plaquette, while each edge
from \(B\) to a pivot hole removes one.
This is an endpoint calculation. It does not
extend the small-angle pivot identity along
a rotation all the way from \(-I\) to \(I\).

Let \(\mathcal E_{\rm in}\) be the undirected
nearest-neighbor edges with both ends in \(H_2\),
and define
\[
       m_x=(K_2)_{xx}-\deg_{\mathcal E_{\rm in}}(x).
 \tag{87.2}
\]
For binary values \(b_x=1_B(x)\),
\[
       Q_2(b)=\sum_{\{x,y\}\in\mathcal E_{\rm in}}|b_x-b_y|
                                    +\sum_xm_xb_x .
 \tag{87.3}
\]
We certify the sign of this expression for every
binary vector, rather than enumerate a sample
and infer its sign for untested configurations.

\subsection{A 24-unit routing identity for all central replacements}

For every vertex \(x\) having exactly one odd
coordinate, and every coordinate \(a\) with
\(x_a=2\), insert the oriented edge
\[
                 x\longrightarrow x+e_a .
 \tag{87.4}
\]
Its target has \(a\)-coordinate three and exactly
two odd coordinates. Denote these oriented
edges by \(\mathcal R\). There are 24 distinct
edges, all in \(\mathcal E_{\rm in}\).
For a chosen odd axis and odd value \(1\) or \(3\),
the two other coordinates are \(0\) or \(2\);
over their four possibilities the number of
twos totals four. Hence the count is
\(3\cdot2\cdot4=24\).

The following coefficients can be read directly
from the cube boundary and the pivot holes:
\[
 \begin{array}{c|c|c}
 k(x)&m_x&
 c_x:=m_x+\operatorname{out}_{\mathcal R}(x)
                         -\operatorname{in}_{\mathcal R}(x)\\
 \hline
 1&-\#\{a:x_a=2\}&0\\
 2&\#\{a:x_a=3\}+1_{\{\text{the even coordinate is }0\}}
          &1_{\{\text{the even coordinate is }0\}}\\
 3&\#\{a:x_a=3\}&\#\{a:x_a=3\}.
 \end{array}
 \tag{87.5}
\]
For \(k=1\), the two pivot neighbors account for
the negative terms; only an even coordinate
equal to zero contributes an exterior free edge.
For \(k=2,3\), there are no pivot neighbors,
so \(m_x\) counts exterior free edges.
Every odd coordinate equal to three at a
\(k=2\) vertex receives exactly one routing edge.
This verifies all entries of (87.5). In particular
\(c_x\ge0\), with \(\sum_xc_x=24\).

Let \(|\mathcal R|\) now denote its underlying
undirected edge set. Substituting the divergence
of this routing into (87.3) gives the exact identity
\[
 \begin{split}
 Q_2(b)
 &=\sum_{\{x,y\}\in\mathcal E_{\rm in}\setminus|\mathcal R|}
                                      |b_x-b_y|\\
 &\quad+\sum_{(x,y)\in\mathcal R}
                    \bigl(|b_x-b_y|-(b_x-b_y)\bigr)
                                      +\sum_xc_xb_x
 \ \ge0 .
 \end{split}
 \tag{87.6}
\]
Every summand on the right is nonnegative for
\(b_x\in\{0,1\}\). This is an integer flow/cut
certificate, with no numerical optimization
bound used as a theorem.

The graph obtained by removing the routing
edges is still connected. Here is an explicit
path to \((1,1,1)\) from any vertex.
First replace each even coordinate, zero or two,
by one. These edges never have the form (87.4),
and they do not pass through a pivot hole.
All coordinates are then odd. Replace each
three by one through \(3\to2\to1\).
The other two coordinates stay odd, so the
first edge joins parity weights three and two,
not weights two and one. The second edge ends
at one and is not a routing edge either.
The entire path is therefore in the graph
with routing edges removed.

If equality holds in (87.6), this connectedness
makes every \(b_x\) equal. The constant one
vector is excluded by \(\sum c_x=24>0\).
Thus equality holds only for \(b=0\).

\begin{theorem}[Exact central minimum and noncentral compact slice minimizers]
\label{thm:f12v83-slice-minimum}
Among all \(2^{56}\) central configurations of the
actual sheet slice, \(Y^*\) is the unique action
minimizer. Every other central configuration has
action at least \(S_{\rm sl}(Y^*)+4\).
Nevertheless every global compact minimizer of
\(S_{\rm sl}\) is noncentral.
\end{theorem}
\begin{proof}
Equations (87.1) and (87.6), with their equality
classification, prove the central assertion.
The strict integer quadratic value is at least one,
which gives the displayed action separation.

The continuous function \(S_{\rm sl}\) attains its
minimum on the compact product \(SU(2)^{56}\).
V82's feasible weighted rotation gives
\[
       \min S_{\rm sl}\le
              S_{\rm sl}(Y^*)-\frac{109}{4096}.
 \tag{87.7}
\]
No central configuration has action this low.
Every compact slice minimizer is therefore noncentral.
\end{proof}

The proof certifies a finite class with
\(2^{56}=72057594037927936\) members through
the universal identity (87.6), not through
exhaustive floating-point testing.
It is also stronger than merely finding a
noncentral local descent: it excludes every
central point from the slice ground set.

\subsection{Conjugation orbits force singular hidden Hessians}

The smooth completion is equivariant under simultaneous
conjugation, and the Wilson action is invariant.
All the frozen exterior data in this slice are
central, so conjugation preserves the slice itself:
\[
       S_{\rm sl}(hxh^{-1})=S_{\rm sl}(x),
                         \qquad h\in SU(2).
 \tag{87.8}
\]
In particular the conjugation orbit of any minimizer
consists entirely of minimizers.

\begin{corollary}[Unavoidable classical zero modes in this actual conditional]
\label{cor:f12v83-orbit-kernel}
At every compact minimizer of the 56-group slice,
the hidden classical Hessian has kernel dimension
at least two.
\end{corollary}
\begin{proof}
Theorem \ref{thm:f12v83-slice-minimum} makes at least
one group value noncentral. The centralizer of a
noncentral \(SU(2)\) element is a one-dimensional
circle subgroup. The common stabilizer of the
tuple is a subgroup of that centralizer, so
its conjugation orbit has dimension at least
\(3-1=2\).

Every point on this orbit is a minimum and hence
a critical point of the smooth slice action.
Differentiate its vanishing gradient along an
orbit curve. The resulting Hessian annihilates
the tangent to that curve. The independent orbit
tangents therefore lie in the hidden Hessian kernel.
\end{proof}

The orbit dimension is two when all noncentral
entries share a common conjugation axis and three
when their common centralizer is discrete.
For an elementary verification, write a unit
quaternion as \((a_i,v_i)\), with imaginary
part \(v_i\in\mathbb R^3\). An infinitesimal
conjugation with imaginary generator \(z\)
has velocity \(2z\times v_i\). Its orbit Gram
matrix in the Euclidean quaternion metric is
\[
          4\sum_i\bigl(|v_i|^2I-v_iv_i^T\bigr).
 \tag{87.9}
\]
Its kernel consists of vectors parallel to every
nonzero \(v_i\), proving the stated dimension
alternatives. Changing to the Hilbert--Schmidt
metric only multiplies this Gram matrix by a
positive fixed constant.

These are zero modes of a classical Hessian at
conditional minimizers. They do not imply zero
heat-bath gap or massless physical particles.
The compact symmetry orbit remains integrated
in the original law. No quotient measure,
gauge-volume division, or deletion of group
updates has been performed.

The conclusion cannot yet be transferred to every
minimizer of the 944-group parent: varying the
other hidden groups may change which central
configurations are competitive. What does transfer
without qualification is the symmetry statement:
for any actual block with central exterior, every
noncentral hidden minimizer has the corresponding
orbit directions in its hidden Hessian kernel.

\subsection{The response criterion when the hidden Hessian is singular}

Return to the actual fixed-repair joint action
\(S(x,v)\) of V81. At a hidden conditional
minimizer \(x_*\), the hidden Hessian is
nonnegative. At a joint critical point, write
the full classical Hessian in orthonormal
product normal coordinates as
\[
       \mathsf H=\begin{pmatrix}A&B\\B^T&C\end{pmatrix},
                  \qquad A\succeq0 .
 \tag{87.10}
\]
Let \(\Pi_0\) be the orthogonal projection onto
\(\ker A\), and let \(A^\dagger\) be the
self-adjoint inverse on \((\ker A)^\perp\),
extended by zero on \(\ker A\).
This defines the pseudoinverse without assuming
a uniform lower bound on the nonzero eigenvalues.

The following is standard generalized Schur
algebra; bibliographic background is
A.~Albert,
\href{https://epubs.siam.org/doi/10.1137/0117041}
{\emph{Conditions for Positive and Nonnegative
Definiteness in Terms of Pseudoinverses}}
(1969). The complete argument needed here follows,
so no inaccessible theorem or unverified regularity
assertion is imported.

\begin{proposition}[Kernel-compatible negative-direction dichotomy]
\label{prop:f12v83-singular-schur}
For (87.10), one has
\[
 \mathsf H\succeq0
 \quad\Longleftrightarrow\quad
 \Pi_0B=0\ \text{and}\ C-B^TA^\dagger B\succeq0 .
 \tag{87.11}
\]
If \(\Pi_0B\ne0\), there is a negative joint
direction without inverting \(A\).
If \(\Pi_0B=0\), negative visible relaxed
curvature in \(C-B^TA^\dagger B\) gives a
negative joint direction.
\end{proposition}
\begin{proof}
If \(k\in\ker A\) and \(a\) is visible, then
\[
        (tk,a)^T\mathsf H(tk,a)
                        =2t\,k^TBa+a^TCa .
 \tag{87.12}
\]
Nonnegativity for all real \(t\) forces
\(k^TBa=0\) for every \(k,a\), equivalently
\(\Pi_0B=0\). If that compatibility fails,
choose \(k,a\) with nonzero cross term and
the opposite sign of \(t\); the value is negative
for sufficiently large magnitude.

Under compatibility, \(AA^\dagger B=B\).
Expansion of the square then gives the exact identity
\[
 \begin{split}
 (z,a)^T\mathsf H(z,a)
 &=(z+A^\dagger Ba)^T A(z+A^\dagger Ba)\\
 &\quad+a^T(C-B^TA^\dagger B)a .
 \end{split}
 \tag{87.13}
\]
This proves sufficiency in (87.11).
Taking \(z=-A^\dagger Ba\) proves necessity
of the relaxed nonnegativity and supplies
the claimed negative direction when it fails.
\end{proof}

Two quantitative versions are useful for a finite
chart certificate. In the compatible case, if
\(\|a\|=1\) and
\(a^T(C-B^TA^\dagger B)a\le-\sigma<0\), then
\[
       z=(-A^\dagger Ba,a),\qquad
       \frac{z^T\mathsf H z}{\|z\|^2}
          \le-\frac{\sigma}{1+\|A^\dagger B\|^2}.
 \tag{87.14}
\]
This follows exactly as in (85.23), with the
kernel retained instead of discarded.

In the incompatible case, choose unit \(k\in\ker A\)
and unit \(a\), with \(b=k^TBa>0\) after a sign
choice. Put \(C_0=\|C\|\) and
\(\epsilon=b/(C_0+b)\). Then
\[
 \begin{split}
 z&=(k,-\epsilon a),\\
 z^T\mathsf H z
       &\le-2\epsilon b+\epsilon^2C_0
                    \le-\frac{b^2}{C_0+b},\\
 \frac{z^T\mathsf H z}{\|z\|^2}
       &\le-\frac{b^2}{2(C_0+b)}<0 .
 \end{split}
 \tag{87.15}
\]
Here \(0<\epsilon\le1\), so \(\|z\|^2\le2\).
The sign and normalization are explicit; a
diverging formal inverse is not used.

Compatibility is essential. For
\[
          A=(0),\quad B=(1),\quad C=(2),
          \qquad \mathsf H=\begin{pmatrix}0&1\\1&2\end{pmatrix},
 \tag{87.16}
\]
the naive pseudoinverse Schur expression is \(2>0\).
Nevertheless the vector \((1,-1/3)\) has normalized
curvature \(-2/5\). Thus replacing \(A^{-1}\)
by \(A^\dagger\) while omitting the kernel test
would produce an incorrect stability conclusion.

\subsection{How symmetry and minimized energy enter the criterion}

At a joint critical point with central frozen
exterior and central visible configuration,
the conjugation tangent has the form \((k,0)\):
conjugation fixes every visible central value.
The full action is invariant and the orbit of
a joint critical point remains joint critical.
Differentiating its gradient gives
\[
                     \mathsf H(k,0)=0,\qquad
                     Ak=0,\quad B^Tk=0 .
 \tag{87.17}
\]
Therefore exact symmetry modes automatically
satisfy the cross-term compatibility. Additional
hidden kernel directions, if present, need a
separate test. At a hidden minimizer with
nonzero visible force, V81's noncritical
alternative already applies; (87.17) is not
asserted there without the joint criticality.

Neither (87.13) nor (87.14) presupposes a smooth
branch of hidden global minimizers. To make this
precise, let \(x_*\) be a hidden global minimizer
at \(v_*\), let the joint point be critical,
and suppose either construction gives a joint
direction \((z,a)\) of strictly negative curvature.
Use the corresponding normal-coordinate curve
\((x(t),v(t))\). For the actual minimized energy
\(E(v)=\min_x S(x,v)\),
\[
 \begin{split}
 E(v(t))
 &\le S(x(t),v(t))\\
 &=E(v_*)+\tfrac12t^2
                  (z,a)^T\mathsf H(z,a)+o(t^2).
 \end{split}
 \tag{87.18}
\]
The hidden component alone cannot have negative
curvature because \(A\succeq0\), so \(a\ne0\).
Equation (87.18) gives an actual visible descent
upper bound even if minimizing branches switch.
It is not an equality of Hessians of \(E\),
nor a proof that \(E\) is twice differentiable.

Conversely, if \(v_*\) is a local minimum of \(E\),
then every hidden minimizer \(x_*\) above it
makes \((x_*,v_*)\) a joint local minimum:
\(S(x,v)\ge E(v)\ge E(v_*)\) for nearby \(v\).
Hence its Hessian satisfies (87.11).
The converse inference from a nonnegative Hessian
to a local minimum is not made.

For V81, (87.14) or (87.15) can supply the needed
negative joint chart direction at a verified
conditional ground point. Uniformity over a
ground set still requires positive margins.
In particular \(A^\dagger B\) can become large
when the nonzero spectrum of \(A\) approaches
zero or its rank changes. No uniform normal
coercivity, constant orbit type, Morse--Bott
structure, or Gaussian integration formula
has been proved by introducing the pseudoinverse.

\subsection{Exact checks and the remaining full-parent question}

The script
\begin{center}\small\ttfamily
scripts/verify\_ym\_f12\_v83\_central\_cut\_and\_singular\_response.py
\end{center}
verifies the 24 unit routing edges, their capacity
and residual coefficients, and an explicit nonrouted
path from each of the 56 vertices to \((1,1,1)\).
It checks the polynomial cut identity on the empty,
full, and 100 reproducible additional patterns.
The proof for all \(2^{56}\) patterns is the
nonnegative identity and connectivity argument,
not those 102 sample substitutions.

The compatible singular-matrix check adjoins an
exact hidden zero mode to the rational coupled
example of V81. The pseudoinverse identities,
kernel compatibility, completed square, relaxed
curvature \(-9/10\), and normalized joint value
\(-9/14\) are checked exactly.
The incompatible example (87.16) checks the
negative value \(-2/5\) despite its positive
naive Schur expression. Rational quaternion
examples verify orbit Gram ranks two and three.
All checks pass.

\begin{decision}[Analyze the full parent with the kernel retained]
The actual sheet slice now has a complete
central-candidate exclusion and forced symmetry
degeneracy at its compact minimizers. This is
evidence against assuming an isolated positive
hidden Hessian in the upstream calculation.
It is not the missing classification of the
full canonical parent: its other 888 hidden
groups may move and can change the minimizing
configuration.

The next full-parent calculation must either
identify that ground set with controlled
normal directions, or produce direct conditional
coverage for a different set of charts. At its
critical ground points the test is (87.11),
including its kernel compatibility, followed
by the actual joint negative-direction criterion
when applicable. None of this changes the
conditional law or turns a classical symmetry
zero mode into a claim about a physical spectrum.
\end{decision}

V82 had 1341183 bytes and SHA--256
\begin{center}\scriptsize\ttfamily
A3446689C4B372A54922249D875DA0F9EA738EDD755ED6949F865DC6597E5892.
\end{center}
Its complete body is retained apart from the title
version. Only active V82 is replaced by V83.
Archives, companions, and earlier scripts remain.

\begin{firewall}[Position after V83]
The central-cut certificate, noncentral compact
slice ground set, its symmetry zero modes, and
the kernel-compatible joint response criterion
are proved. The full canonical-parent ground
set, uniform repair estimate, cofinal auxiliary
gap, continuum construction, and physical
Yang--Mills mass gap remain unproved.
The full objective remains active.
\end{firewall}
\section{V84: convex reduction of the cutoff-safe sheet ground set}
\label{sec:f12v84}

V83 proved that the 56-group compact sheet conditional
has only noncentral minimizers, without identifying
them. We now classify that slice ground set as a
single conjugation sphere with a positive normal
Hessian. The proof first analyzes an auxiliary
identity-pivot spin action, then uses exact rational
supporting planes and a uniform pivot comparison
to prove that the true compact minimizers lie where
the auxiliary action is exactly the actual action.
The identity-pivot expression is not assumed valid
through the cutoff region.

There is an explicit additional hypothesis:
the already fixed V18 cutoff radius, denoted here
by \(\rho\), must satisfy its previous smallness
conditions and
\[
                              0<\rho\le1/64 .
 \tag{88.1}
\]
This section does not assert that every previously
admissible cutoff satisfies (88.1), or silently
change a specified numerical cutoff. The theorem
applies to the V18 family with this additional bound.
The conclusion is still about the actual 56-group
conditional slice, not the full 944-group parent.

The new finite certificate is stored as integer
data in the accompanying verifier. A numerical
search supplied candidate supporting planes, but
only exact rational inequalities are used below.
The archived one-axis determinant annulus in f9 V52
was checked and concerns a different regression
problem; it is not a previous sheet-ground-state
classification. The next companion audit is V85.

\subsection{The identity-pivot comparison action and its scalar reduction}

Keep the slice of V83, with all other independent
groups at the central sheet field and \(C=I\).
Let \(S_{\rm act}(U)\) be its actual action,
with every pivot solved by V18. Define
\(S_{\rm fp}(U)\) only for comparison by setting
those pivots to \(I\), at the same 56 free groups.
Write
\[
           U_i=a_iI+i\,v_i\cdot\sigma,\qquad
           a_i^2+|v_i|^2=1,\qquad
           s_i=\sqrt{1-a_i^2},
 \tag{88.2}
\]
using the usual Pauli matrices. Let
\(\mathcal E_{\rm in}\) be the 108 internal
free--free edges of \(H_2\). Let \(r_i\) count
pivot-hole neighbors and \(e_i\) count frozen
negative free neighbors outside \(H_2\).
Thus
\[
             h_i=r_i-e_i,\quad
             \sum_i r_i=48,\quad \sum_i e_i=72,\quad
             \sum_i h_i=-24 .
 \tag{88.3}
\]
The quaternion trace identity gives, relative
to \(S_*:=S_{\rm act}(Y^*)=S_{\rm fp}(Y^*)\),
\[
 \begin{split}
 S_{\rm fp}(U)-S_*
 &=2\bigl[\mathcal F(a,v)+132\bigr],\\
 \mathcal F(a,v)
 &=-\sum_{\{i,j\}\in\mathcal E_{\rm in}}
                       (a_i a_j+v_i\cdot v_j)-\sum_i h_i a_i .
 \end{split}
 \tag{88.4}
\]
The baseline value is
\(-108+\sum_i h_i=-132\).

For fixed \(a\), the inequalities
\(v_i\cdot v_j\le s_i s_j\) can be attained
simultaneously: take \(v_i=s_i n\) for one
common \(n\in S^2\). Hence
\[
 \begin{split}
 F(a)&=-\sum_{\{i,j\}\in\mathcal E_{\rm in}}
           \left(a_i a_j+\sqrt{1-a_i^2}\sqrt{1-a_j^2}\right)
                                      -\sum_i h_i a_i,\\
 \min_U\mathcal F(a,v)&=\min_{a\in[-1,1]^{56}}F(a).
 \end{split}
 \tag{88.5}
\]
In the second line the minimum on the left is over
all spin variables. Pointwise, \(\mathcal F(a,v)\ge F(a)\).
This is an exact minimization of the comparison
action; it does not yet identify \(S_{\rm act}\).

\subsection{Convexity and uniqueness of the scalar profile}

For \(a,b\in(-1,1)\), put
\(s=\sqrt{1-a^2}\), \(t=\sqrt{1-b^2}\), and
\[
                       f(a,b)=-ab-st .
 \tag{88.6}
\]
Direct differentiation gives
\[
 D^2f=
 \begin{pmatrix}
     t/s^3&-1-ab/(st)\\
     -1-ab/(st)&s/t^3
 \end{pmatrix},\qquad
 \det D^2f=\frac{1-(ab+st)^2}{s^2t^2}\ge0 .
 \tag{88.7}
\]
The last inequality is Cauchy--Schwarz for the
unit vectors \((a,s)\) and \((b,t)\).
The diagonal entries are positive. Equality
in the determinant holds only when \(a=b\),
since \(s,t>0\) exclude antipodal vectors.
Thus \(f\) is convex, strictly positive in its
Hessian when \(a\ne b\), and at \(a=b\) has
Hessian kernel spanned by \((1,1)\).
Continuity extends convexity to the closed square.
Summing edges and linear terms proves that \(F\)
is convex on \([-1,1]^{56}\).

Every minimizer of \(F\) is in the open cube.
To see this, first it cannot have all coordinates
equal to endpoints: V83's central-cut certificate
makes their minimum \(-132\), whereas V82's
feasible small rotation has smaller comparison
action. Thus at least one \(s_i>0\).
If an endpoint coordinate has a neighbor with
positive \(s_j\), perturb it so that its new
transverse length is \(\epsilon>0\) and its
longitudinal value retains the endpoint sign.
The neighboring transverse term decreases \(F\)
by \(\epsilon\sum_{j\sim i}s_j\), while changes
in all longitudinal terms are \(O(\epsilon^2)\).
For small positive \(\epsilon\) this is a strict
decrease. Connectivity rules out an interface
between endpoint and interior coordinates;
hence every coordinate is interior.

A constant interior vector cannot be a minimizer:
the edge gradients vanish there and
\(\nabla F=-h\ne0\), by (88.3).
At any nonconstant interior vector the Hessian
of \(F\) is positive definite. Indeed a connected
graph then has an edge with unequal endpoint
values. Its positive definite edge Hessian
annihilates no nonzero restriction of a null
vector. The null vector must vanish at both
ends. On any neighboring equal-value edge its
two components must agree, while on an
unequal-value edge they must both vanish.
Connectedness propagates zero to every vertex.

It follows that \(F\) has a unique minimizer,
denoted \(a^*\). Otherwise the line segment
between two minimizers would consist of
interior minimizers. None is constant and each
has positive definite Hessian, contradicting
constancy along the segment.
We have proved
\[
            a^*\in(-1,1)^{56},\qquad D^2F(a^*)>0 .
 \tag{88.8}
\]
No optimizer stopping criterion is used in this
existence, uniqueness, or Hessian argument.

\subsection{Exact rational separation from the near-identity region}

To compare with the true pivots, we need a
quantified energy exclusion, not a plot of
the scalar minimizer. We prove
\[
    \max_i a_i\ge24/25
                 \quad\Longrightarrow\quad F(a)\ge-131 .
 \tag{88.9}
\]
Here is the complete form of the finite certificate.
For a chosen index \(i\), choose an interior
rational Pythagorean tuple \(c,s\), so that
\(c_j^2+s_j^2=1\), \(s_j>0\), and \(c_i=24/25\).
Let \(g=\nabla F(c)\).
Convexity gives the supporting-plane inequality
\(F(a)\ge F(c)+g\cdot(a-c)\).
It follows also directly by applying the defining
convexity inequality on the segment from \(c\)
to \(a\), dividing by its length parameter,
and taking the limit at zero. See also
Boyd--Vandenberghe,
\href{https://web.stanford.edu/~boyd/cvxbook/bv_cvxbook.pdf}
{\emph{Convex Optimization}, Chapter 3},
for this standard first-order condition.

If \(g_i\ge0\), then for every \(a_i\ge24/25\),
\[
           F(a)\ge F(c)
                   -\sum_{j\ne i}|g_j|(1+|c_j|).
 \tag{88.10}
\]
This holds up to the boundary by continuity.
No constrained minimization algorithm is needed
to verify the bound.

The verifier stores \(d=100000\) and sixteen
lists of 56 positive integers \(k_j\).
For each list it defines exactly
\[
              c_j=\frac{k_j^2-d^2}{k_j^2+d^2},
              s_j=\frac{2k_jd}{k_j^2+d^2}.
 \tag{88.11}
\]
The distinguished integer is \(k_i=7d\),
giving \(c_i=24/25\), \(s_i=7/25\).
The value and every gradient component are rational:
\[
                 g_j=-h_j+
                   \sum_{\ell\sim j,\ \ell\in H_2}
                        \left(-c_\ell+\frac{c_j s_\ell}{s_j}\right).
 \tag{88.12}
\]
The stored data verify \(g_i>0\) and the
following lower bounds from (88.10).
Each displayed integer is a lower bound
multiplied by 1000.
\begin{center}\small
\begin{tabular}{rr|rr}
\toprule
Representative \(i\)&\(1000L_i\)&Representative \(i\)&\(1000L_i\)\\
\midrule
\((0,0,1)\)&\(-130623\)&\((1,1,2)\)&\(-128648\)\\
\((0,0,3)\)&\(-129778\)&\((1,1,3)\)&\(-125240\)\\
\((0,1,1)\)&\(-128666\)&\((1,2,2)\)&\(-130597\)\\
\((0,1,2)\)&\(-130642\)&\((1,2,3)\)&\(-126692\)\\
\((0,1,3)\)&\(-126325\)&\((1,3,3)\)&\(-123556\)\\
\((0,2,3)\)&\(-129882\)&\((2,2,3)\)&\(-129882\)\\
\((0,3,3)\)&\(-124644\)&\((2,3,3)\)&\(-125092\)\\
\((1,1,1)\)&\(-127877\)&\((3,3,3)\)&\(-122467\)\\
\bottomrule
\end{tabular}
\end{center}
Permuting the three transverse coordinates preserves
the graph and the fields. The sixteen sorted triples
listed are exactly the orbit representatives of
the 56 sites. Thus these inequalities cover
every possible distinguished index. Each
\(L_i\ge-131\), proving (88.9).

The full integer lists, their site ordering, and
the exact rational evaluation are in
\begin{center}\small\ttfamily
scripts/verify\_ym\_f12\_v84\_cutoff\_safe\_sheet\_ground.py .
\end{center}
The verifier imports no numerical optimizer.
This is a finite exact supporting-plane certificate,
not a claim that floating-point approximations
to the minimum are rigorous lower bounds.

\subsection{Comparison with the actual nonlinear pivots}

Only eight coarse pivots can depend on these
56 boundary variables. Their fourteen nonroot
path products consist of the seven corresponding
boundary groups, each appearing twice, since
all internal groups are fixed at \(I\).
For arbitrary values of the 56 groups, the
actual pivot equation at \(C=I\) gives
\[
           B_e=\exp(-X_e),\qquad
           \|X_e\|_{\rm HS}\le7\rho/8,\qquad
           \|B_e-I\|_{\rm HS}\le7\rho/8 .
 \tag{88.13}
\]
This uses the global cutoff value bound from V18,
not a small-chart assumption about the free groups.

Each physical pivot link belongs to six
plaquettes. Replacing one pivot by \(I\)
changes the sum of their Wilson actions by
at most \(6\sqrt2\,\|B_e-I\|_{\rm HS}\):
use unitary invariance of the Hilbert--Schmidt
norm and \(|\operatorname{Tr}A|\le\sqrt2\|A\|_{\rm HS}\).
Telescoping the eight replacements yields the
global slice comparison
\[
          |S_{\rm act}(U)-S_{\rm fp}(U)|
              \le42\sqrt2\,\rho<60\rho\le15/16 .
 \tag{88.14}
\]
Overlapping plaquettes cause no problem in this
telescoping bound. Other physical links stay fixed.

If every \(a_i<24/25\), then
\[
               \|U_i-I\|_{\rm HS}^2=4(1-a_i)>4/25.
 \tag{88.15}
\]
Their logarithmic distance from \(I\), whenever
the principal logarithm is used, is at least
their chord distance. Since \(2/5>\rho\),
all cutoff logarithms at pivot \(I\) vanish.
The exact pivot equation and uniqueness imply
\[
       B_e=I,\qquad S_{\rm act}=S_{\rm fp}
              \quad\text{on }\{\max_i a_i<24/25\}.
 \tag{88.16}
\]
The same is true in a neighborhood of each point
of this open set. There the eight varying inverse
pivot Jacobians equal one as well.

On its complement, (88.4), (88.9), and (88.14) give
\[
       S_{\rm act}(U)\ge S_*+2-15/16=S_*+17/16 .
 \tag{88.17}
\]
On the other hand the feasible V82 competitor has
actual action at most \(S_*-109/4096\).
Thus no actual global slice minimizer is in
the complement. Likewise no comparison-action
minimizer is there, by (88.9).
Both minima therefore occur in the open region
where the two actions are exactly equal.
Their complete global minimizing sets coincide.

\subsection{The actual ground set is one sphere with positive normal Hessian}

\begin{theorem}[Classified compact sheet ground manifold]
\label{thm:f12v84-ground-manifold}
Under (88.1), every global minimizer of the
actual 56-group slice has the form
\[
           U_i(n)=a_i^*I+i\sqrt{1-(a_i^*)^2}\,n\cdot\sigma,
                          \qquad n\in S^2 ,
 \tag{88.18}
\]
where \(a^*\) is the unique scalar minimizer in
(88.8) and \(a_i^*<24/25\) for every \(i\).
The ground set is an embedded \(S^2\).
The classical Hessian has exactly its two
tangent zero directions and is positive
definite on the normal space.
\end{theorem}
\begin{proof}
The preceding comparison identifies the actual
minimizing set with that of \(S_{\rm fp}\).
The scalar profile is unique and interior
by (88.8). Equality in
\(v_i\cdot v_j\le s_i s_j\) is necessary at
every internal edge, and all \(s_i>0\).
Connectedness therefore gives one common
unit vector \(n\). Conversely every such
choice attains the scalar minimum.
The formula is a smooth injective immersion
of the compact sphere, hence an embedding.
The cutoff bound follows from (88.9).

Use local coordinates \(a_i\) and \(n_i\in S^2\)
near one of these points. Longitudinal
second variations are governed by
\(2D^2F(a^*)>0\). If \(\eta_i\) is a tangent
variation of \(n_i\) at their common value,
the quadratic energy increase in the
orientation variables is
\[
          \sum_{\{i,j\}\in\mathcal E_{\rm in}}
                          s_i^*s_j^*|\eta_i-\eta_j|^2 .
 \tag{88.19}
\]
Indeed \(n_i\cdot n_j
=1-\tfrac12|\eta_i-\eta_j|^2+O(|\eta|^3)\).
All edge weights in (88.19) are positive.
Its kernel consists exactly of the two
constant tangent fields. Mixed longitudinal
and orientation second derivatives vanish
because the first orientation variation
is zero for every longitudinal profile
at aligned directions.

Thus there are precisely the two common
orientation zero modes and positive curvature
on a complementary normal space.
The actual action agrees with this action
on an open neighborhood by (88.16), so its
Hessian has the same property. Invariance
under common conjugation makes the normal
positivity uniform along the compact orbit.
\end{proof}

The theorem proves a Morse--Bott ground
manifold for this conditional: a smooth
minimum manifold with Hessian kernel
exactly its tangent space. It does not
classify all higher-energy stationary
points, which could still matter for
global stochastic mixing.

\subsection{A rigorous consequence for the compact normalization power}

Cancel all factors independent of the 56
groups in the slice conditional. Subtract
the constant \(S_*\) from the action and
write the resulting conditional normalizer as
\[
       Z_{\rm sl}(\beta)=
         \int a_{\rm sl}(U)
               e^{-\beta(S_{\rm act}(U)-S_*)}\,dU ,
       \qquad
       E_*=\min_U(S_{\rm act}(U)-S_*).
 \tag{88.20}
\]
Here \(a_{\rm sl}\) is the product of the
eight possibly varying inverse pivot
Jacobians. It is positive and bounded on
the whole compact slice and equals one
in the ground neighborhood.

\begin{corollary}[Normal rather than full-dimensional Laplace power]
\label{cor:f12v84-laplace-power}
For this fixed slice and cutoff there are
\(0<c<C<\infty\) and \(\beta_0<\infty\)
such that
\[
       c\,\beta^{-83}e^{-\beta E_*}
           \le Z_{\rm sl}(\beta)
           \le C\,\beta^{-83}e^{-\beta E_*},
                    \qquad \beta\ge\beta_0 .
 \tag{88.21}
\]
\end{corollary}
\begin{proof}
The slice dimension is \(3\cdot56=168\).
The ground manifold has dimension two,
so its normal dimension is 166.
Choose finitely many smooth tubular
coordinate charts \((y,z)\) about the
compact ground manifold, with
\(z\in\mathbb R^{166}\).
The positive normal Hessian and Taylor's
theorem give uniform constants on
sufficiently small charts:
\[
       c_1|z|^2\le
          S_{\rm act}(y,z)-S_*-E_*
                         \le C_1|z|^2 .
 \tag{88.22}
\]
The Haar-volume coordinate densities and
\(a_{\rm sl}\) have positive uniform
upper and lower bounds there.
Integrating a smaller orbit chart and
a normal ball of radius \(\beta^{-1/2}\)
proves the lower bound, with power
\(\beta^{-166/2}\).
For the upper bound sum the finitely
many chart estimates and integrate
the normal Gaussian over all
\(\mathbb R^{166}\).

The compact complement of a sufficiently
small ground neighborhood has an
action gap \(\delta>0\), since the
ground set is exactly (88.18).
Its contribution is at most a constant
times \(e^{-\beta(E_*+\delta)}\),
which is bounded by a constant times
\(\beta^{-83}e^{-\beta E_*}\).
This proves both bounds.
\end{proof}

Thus the isolated-nondegenerate power
\(\beta^{-84}\) is not the correct
two-sided order for this actual conditional.
No numerical normal determinant or
unproved Gaussian branch selection
is needed for (88.21). This does not
alter V81's conservative general
denominator lower bound: that bound
was not asserted to be sharp.

\subsection{Verification and the next full-parent test}

The new verifier freezes all sixteen
candidate tuples as positive integer lists.
It evaluates their rational scalar
values, gradients, and supporting-plane
debits exactly, certifying (88.9) on
all 56 possible near-identity sites.
It also checks 25 rational instances
of (88.7), the inequality
\(2\cdot42^2<60^2\), the cutoff
debit \(15/16\), the forbidden-region
margin \(17/16\), and normal dimension
166. All exact checks pass.
The global convexity, uniqueness,
orientation alignment, and normal
Hessian assertions are proved above;
they are not inferred from those
25 matrix substitutions.

\begin{decision}[Use the classified slice as a test point, not as a full-parent solution]
The compact sheet minimization is now
resolved under the stated cutoff bound:
there is one noncentral orbit, its
normal Hessian is positive, and its
normalization power is known.
The next question is whether its
criticality survives when the other
888 hidden groups of the canonical
parent are released. That requires
their actual first derivatives and
the coupled normal response; the
slice result alone does not answer it.

The positive normal Hessian is not
an unrestricted heat-bath gap.
Neither higher-energy stationary
points nor the full-parent conditional
ground set have been classified.
No scale-uniform repair inequality
or continuum conclusion is obtained
by treating the compact orbit as
an already integrated scalar variable.
\end{decision}

V83 had 1359792 bytes and SHA--256
\begin{center}\scriptsize\ttfamily
B72A009B3859BB21B8C67E3AD2A29B153AEC8C2208F8FB34309D43AF5E3268CD.
\end{center}
Its full body is retained apart from
the title version. Only active V83
is replaced by V84. Archives,
companions, and previous scripts remain.

\begin{firewall}[Position after V84]
Under the additional cutoff hypothesis,
the actual 56-group slice ground set,
its normal Hessian, and its two-sided
Laplace power are proved, using a
fully rational forbidden-region certificate.
The full-parent ground set,
unrestricted repair gap, cofinal
auxiliary gap, continuum construction,
and physical Yang--Mills mass gap
remain unproved. The full objective
remains active.
\end{firewall}
\section{V85: companion audit and the released-parent force}
\label{sec:f12v85}

V84 solved the compact sheet slice under
its additional cutoff bound \(\rho\le1/64\).
It explicitly left open stationarity when
the other hidden parent variables are released.
The answer is negative: a particular
non-tree internal transverse link has a
nonzero force at every point of the
classified slice ground sphere.
The induced pivot variations do not cancel it.
We derive a uniform finite decrease and
an exponential conditional exclusion of
a neighborhood of that entire embedded
sphere in the actual parent law.

This is an upstream correction to the
candidate ground set, not a full-parent
minimizer construction. The slice result
and its Laplace power remain valid for
the slice conditional.

\subsection{The scheduled companion review}

The active companions were read and hashed:
\begin{center}\small
\begin{tabular}{lrr}
\toprule
Programme&Version&Bytes\\
\midrule
SM--Einstein&72&607372\\
NS stability&26&278283\\
Parallel YM&5&54174\\
\bottomrule
\end{tabular}
\end{center}
\begin{center}\scriptsize\ttfamily
SM: F44F9745D43379E1672C5550411B58E97195A4C9278592D221DF5D3F644DC1F5\\
NS: 82C887F8C2C69E4F28FAD7C3DC8DE5B9099CB5A4BF431EBA1FFF3E91935978AD\\
YM: 306F1BB84CFA8A8D47EA569EC17E9545F9867C8D32B548EC9D9C444B4F3C0512
\end{center}
These are unchanged from V80. SM V72
controls its fixed-mode, fixed-time many-copy
limit, not the physical finite-species
continuum. NS V26 supplies its pointwise
rank-one formulas, not full stability or
global regularity. Parallel YM V5 diagnoses
the extensive global-charge tail target;
it does not provide a new local spectral
estimate. Document-internal continuation
and stop language is source material,
not a change to the present task.

No new companion input changes the next
calculation. We therefore continue the
promised full-parent stationarity test.
The next companion audit and ten-version
direction review are both V90.

\subsection{The parent coordinate that was frozen in the slice}

Use fine coordinates in directions \(0,1,2,3\),
as in the exact tree convention of V18.
Take the canonical parent
\(P_b=\{0,1\}^4\) in coarse coordinates,
with \(|I_b|=944\), embedded in a compatible
separated packing. The 56 sheet groups
are those of V84. Let \(X^*(n)\), \(n\in S^2\),
be its ground sphere embedded into the full
parent by leaving all other 888 hidden
groups at their reference value \(I\).
The exterior is the central sheet exterior,
denoted \(v^*\), and \(C=I\).

Select the fine link and neighboring sheet link
\[
 \begin{split}
       \ell&=((0,0,1,0),0),\\
       j&=((1,0,1,0),2).
 \end{split}
 \tag{89.1}
\]
The transverse sheet index of \(j\) is
\(j_\perp=(1,0,0)\).
The link \(\ell\) is internal in its owner cell:
its direction-\(0\) base coordinate is even.
It is not a tree link. Indeed the fixed
tree rule is
\[
 \text{\((x,\mu)\) is a tree link}
 \quad\Longleftrightarrow\quad
 x_\mu\bmod2=0,\quad
 x_k\bmod2=0\text{ for every }k>\mu,
 \tag{89.2}
\]
and \(x_2\bmod2=1\) at \(\ell\).
Thus \(\ell\in I_b\), among the hidden
groups which were frozen in the slice.

At the embedded sphere,
\[
          U_j=a_j^*I+i s_j^* n\cdot\sigma,
               \qquad s_j^*=\sqrt{1-(a_j^*)^2}>0,
               \qquad U_\ell=I .
 \tag{89.3}
\]
The strict positivity follows from the
interiority proved in V84, not from an
approximate numerical profile.

\subsection{The physical force and the induced pivots}

Among the six plaquettes incident on
\(\ell\), exactly one can have nonidentity
holonomy at this configuration. It is
the \((0,2)\) plaquette based at
\((0,0,1,0)\). Its oriented word is
\[
 U_0(0,0,1,0)\,
 U_2(1,0,1,0)\,
 U_0(0,0,2,0)^{-1}\,
 U_2(0,0,1,0)^{-1}.
 \tag{89.4}
\]
The last link is the direction-\(2\)
root pivot and equals \(I\); the
other transverse link also equals \(I\).
The word is therefore \(U_j\).
The plaquette on the other side in
direction \(2\) uses only direction-\(2\)
links based at \(x_2=0\), and has
identity holonomy. Plaquettes in the
other directions are identity as well.

For \(Z\in\mathfrak{su}(2)\), varying the
physical link \(\ell\) to \(\exp(tZ)\)
therefore has derivative
\[
             D_\ell S[Z]=-\operatorname{Re}\operatorname{Tr}(ZU_j).
 \tag{89.5}
\]
For the unit Hilbert--Schmidt generator
\(Z=i n\cdot\sigma/\sqrt2\), this equals
\(\sqrt2\,s_j^*>0\).

The independent-coordinate variation must
also solve the pivots. The only coarse
path families containing this internal
direction-\(0\) link are those of
the direction-\(0\) edges based at the
owner cell and the preceding cell
in direction \(0\). Each contains
\(\ell\) in one of its fourteen
nonroot chronological products.
At the reference configuration all
products in those two path families
are \(I\).
The derivative of their pivot map is
\(1-14/16=1/8\); the direct variation
contributes \(Z/16\). Thus the induced
pivot derivatives are
\[
                           \dot B_+=\dot B_-=-Z/2 .
 \tag{89.6}
\]
In particular the incoming coarse edge
is not omitted from the canonical parent.

Both physical pivot links have fine
\(x_2=0\). Every plaquette incident on
either is identity: their adjacent
direction-\(2\) links are based at
\(x_2=0\) or \(x_2=-1\), not at the
noncentral sheet \(x_2=1\).
Their physical Wilson gradients vanish.
Consequently their chain-rule contributions
to the first derivative are zero,
irrespective of the nonzero responses
in (89.6).

\begin{proposition}[The slice ground sphere is not parent-stationary]
\label{prop:f12v85-parent-force}
Every point \(X^*(n)\) of the embedded
slice sphere has a nonzero classical
gradient in the actual parent conditional.
More precisely, its independent
link derivative is exactly (89.5).
None of these points is a hidden
local minimum of the full parent.
\end{proposition}
\begin{proof}
The physical derivative is (89.5);
the only additional pivot derivatives
give zero physical force, as just
shown. Equation (89.3) then supplies
a nonzero derivative. The varied
coordinate belongs to \(I_b\), so
this is an allowed parent variation,
not a change in the exterior or
the retained coarse field.
\end{proof}

This is a statement about the classical
action. It is not an assertion that
an arbitrary path of physical link
changes preserves the coarse field;
the induced pivots were included
explicitly in this first variation.

\subsection{A rational uniform lower bound on a descent slope}

A very conservative explicit bound
avoids treating a numerical value
of \(s_j^*\) as a certificate.
Let \(m_s=\max_i s_i^*\) and
\(\delta_{56}=109/4096\).
The scalar minimum from V84 satisfies
\[
                       F(a^*)+132\le-\delta_{56}/2 .
 \tag{89.7}
\]
For any profile in the cutoff-safe
region \(a_i<24/25\),
\(1+a_i=s_i^2/(1-a_i)\le25s_i^2\).
Moreover
\[
 \begin{split}
 F(a)+132
 &=\sum_{\{i,k\}\in\mathcal E_{\rm in}}
                 (1-a_i a_k-s_i s_k)
                              -\sum_i h_i(a_i+1)\\
 &\ge-\sum_{\{i,k\}\in\mathcal E_{\rm in}}s_i s_k
                              -\sum_{h_i>0}h_i(a_i+1).
 \end{split}
 \tag{89.8}
\]
There are 108 edges and
\(\sum_{h_i>0}h_i=24\), as follows
from the explicit fields in V83--V84.
Thus (89.7)--(89.8) imply
\[
                         m_s^2\ge\delta_{56}/1416 .
 \tag{89.9}
\]

Interiority and scalar stationarity give
\[
        \lambda_i(a_i^*,s_i^*)
          =\left(h_i+\sum_{k\sim i}a_k^*,
                              \sum_{k\sim i}s_k^*\right),
        \qquad
        \lambda_i=\frac{\sum_{k\sim i}s_k^*}{s_i^*}>0 .
 \tag{89.10}
\]
The sums here include only internal
free--free neighbors. Since each
neighbor vector has length one,
\[
          \lambda_i\le |h_i|+\deg_{\rm in}(i)
            \le r_i+e_i+\deg_{\rm in}(i)=6 .
 \tag{89.11}
\]
Therefore \(s_i^*\ge s_k^*/6\) along
every internal edge.
Every vertex of \(H_2\) can reach
\((1,1,1)\) by the path used in V83
in at most six steps, and then reach
\(j_\perp=(1,0,0)\) through
\((1,1,0)\) in two more.
All vertices on those paths remain
in \(H_2\). Hence
\[
       (s_j^*)^2\ge6^{-16}m_s^2,\qquad
       2(s_j^*)^2\ge
       \gamma:=\frac{\delta_{56}}{708\cdot6^{16}}
                =\frac{109}{8181128456105361408}>0 .
 \tag{89.12}
\]
This small bound is used only for a
uniform positivity certificate, not
as an optimized physical rate.

\subsection{A smooth parent transport preserving product Haar}

Define the globally smooth equivariant
Lie-algebra function
\[
                         A(U)=\tfrac12(U-U^\dagger).
 \tag{89.13}
\]
For \(SU(2)\), it is traceless and
skew-Hermitian; if \(U=aI+i v\cdot\sigma\),
then \(A(U)=i v\cdot\sigma\) and
\(\|A(U)\|_{\rm HS}^2=2|v|^2\le2\).

On the full compact hidden parent define
\[
       (\mathcal R_tX)_\ell
               =X_\ell\exp[-tA(X_j)],\qquad
       (\mathcal R_tX)_i=X_i\quad(i\ne\ell).
 \tag{89.14}
\]
The controlling group \(j\) is distinct
from \(\ell\) and is unchanged by the
map. Thus the inverse is
\(\mathcal R_{-t}\). Conditional on
all groups except \(\ell\), (89.14)
is a right translation of that
group. Fubini and Haar invariance
prove that \(\mathcal R_t\) preserves
the entire product Haar measure.
The dependence on \(X_j\) introduces
no extra Jacobian.

At the embedded sphere, equation (89.5)
applied to \(-A(U_j)\) gives
\[
       \left.\frac{d}{dt}
          S_b(\mathcal R_tX^*(n),v^*)\right|_{t=0}
                       =-2(s_j^*)^2\le-\gamma .
 \tag{89.15}
\]
This choice covers the whole conjugation
orbit without choosing a preferred axis
or dividing by an orbit volume.

Let \(M_{\ell\ell}\) be a uniform bound
on the absolute intrinsic second derivative
of the actual classical regional action
in a unit direction of the independent
group \(\ell\), over its compact local
input space. V18's fixed-stencil
smoothness makes this finite and
independent of ambient size.
Set
\[
       C_2=\max\{1,2M_{\ell\ell}\},\qquad
       t_0=\gamma/C_2,\qquad
       \delta_{\rm p}=\gamma^2/(2C_2)>0 .
 \tag{89.16}
\]
The path in (89.14) is a geodesic
in its one varying group; \(X_j\)
is constant along it. Its squared
speed is at most two.
Thus the absolute second derivative
along this path is at most \(C_2\).
Since \(0<\gamma<1\), one has \(t_0\le1\).
Taylor's theorem and (89.15) prove
\[
       S_b(\mathcal R_{t_0}X^*(n),v^*)
           \le S_b(X^*(n),v^*)-\delta_{\rm p}
                 \quad\hbox{for every }n\in S^2 .
 \tag{89.17}
\]

\begin{corollary}[Strict lowering in the actual full parent]
\label{cor:f12v85-parent-energy}
The full parent ground energy at
\(v^*\) is smaller than the action
on the embedded slice ground sphere
by at least \(\delta_{\rm p}\).
\end{corollary}
\begin{proof}
The configurations in (89.17) are
feasible points of the same full
compact parent with the same exterior.
Take the infimum over all parent
configurations. The action on the
embedded sphere is constant by
simultaneous conjugation invariance.
\end{proof}

The bound does not require that
\(\mathcal R_{t_0}X^*(n)\) itself
be a parent stationary point.
It is a certified competitor.

\subsection{The embedded sphere has exponentially small parent weight}

The action difference in (89.17)
is smooth in the full parent variables
and their finite exterior inputs.
By compactness of the embedded sphere,
there are a fixed open neighborhood
\(\mathcal U\) of that sphere in the
parent and a fixed visible-input
neighborhood \(\mathcal V\) of \(v^*\)
such that
\[
        S_b(\mathcal R_{t_0}X,v)
             \le S_b(X,v)-\delta_{\rm p}/2,
                \qquad X\in\mathcal U,\quad v\in\mathcal V .
 \tag{89.18}
\]
The coarse field remains \(I\).
The neighborhoods are chosen at fixed
completion and parent geometry, not
shrunk as \(\beta\) grows.

Unlike V82's boundary-only rotation,
this internal-link move does not
have an exact amplitude cancellation.
It can change the two direction-\(0\)
inverse pivot Jacobians identified
above. No other pivot factor depends
on this internal direction-\(0\)
variable, by the explicit two-link
path structure. If the uniform
V18 bounds are \(j_-\le j_e\le j_+\),
then the actual block amplitude obeys
\[
       a_b(\mathcal R_{t_0}X,v)
          \ge e^{-B_0}a_b(X,v),\qquad
                B_0=2\log(j_+/j_-).
 \tag{89.19}
\]
This estimate is global on the
compact variables and has no factor
proportional to the number of
groups in the parent or ambient volume.
The controlling coordinate \(X_j\)
is unchanged, so its own pivot
factor does not acquire a new
argument under the map.

\begin{theorem}[Conditional exclusion after releasing the parent link]
\label{thm:f12v85-parent-exclusion}
For the actual unrestricted parent
conditional and \(v\in\mathcal V\),
\[
 \begin{split}
       \pi_{b,\beta}(\mathcal U\mid v)
          &\le \min\{1,\exp(B_0-\beta\delta_{\rm p}/2)\},\\
       S_b(X,v)-\min_ZS_b(Z,v)
          &\ge\delta_{\rm p}/2\qquad(X\in\mathcal U).
 \end{split}
 \tag{89.20}
\]
\end{theorem}
\begin{proof}
The second statement follows from
the feasible competitor in (89.18).
For the first, integrate over the
image of \(\mathcal U\) inside the
full compact normalizer and change
variables by (89.14):
\[
 \begin{split}
 Z_b(v)
 &\ge \int_{\mathcal R_{t_0}\mathcal U}
                      a_b(Z,v)e^{-\beta S_b(Z,v)}\,dZ\\
 &\ge e^{-B_0+\beta\delta_{\rm p}/2}
             \int_{\mathcal U}a_b(X,v)e^{-\beta S_b(X,v)}\,dX .
 \end{split}
 \tag{89.21}
\]
Use product-Haar preservation, (89.18),
and (89.19), then divide by \(Z_b(v)\).
The upper bound one is automatic.
\end{proof}

This transport is a comparison map,
not an original heat-bath trajectory.
Equation (89.20) is a conditional
probability and energy statement,
not a pointwise stochastic hitting
bound. In particular the result
does not prove a repair Poincare
inequality merely by exhibiting
one classical descent.

\subsection{What the release test changes}

The Morse--Bott sphere of V84 is
the correct ground manifold for
the 56-group slice conditional.
It is not even stationary for the
944-group parent, and its embedded
neighborhood is exponentially
underweighted there. Therefore
one cannot obtain the parent
normalizer by taking the
\(\beta^{-83}\) slice factor
and appending Gaussian powers
for the remaining 888 groups
at their old section values.
The leading parent energy is
already strictly lower by
(89.17), before such a prefactor
calculation could begin.

The released force is in an
actual internal transverse
connection, not just a rotation
of the common sheet axis.
The full-parent ground calculation
must relax those connections
and solve the coupled stationarity
equations before using the
kernel-compatible Hessian
response from V83.
At the current candidate, the
first derivative fails, so
another Hessian-only calculation
there would not address the
missing ground-set condition.

\subsection{Exact verification and continuation record}

The script
\begin{center}\small\ttfamily
scripts/verify\_ym\_f12\_v85\_released\_parent\_force.py
\end{center}
checks that \(\ell\) is a non-tree
internal parent coordinate.
It enumerates its two chronological
path families and verifies the
two inverse-pivot responses
\(-1/2\). On a fine torus of side
six it differentiates the affected
physical plaquette words in exact
quaternion arithmetic.

Three noncommuting rational sheet
configurations give released-link
derivatives
\[
                 -448/225,\qquad
                 -2800/1521,\qquad
                 -1008/625 .
 \tag{89.22}
\]
For each, the induced pivot-force
contribution is exactly zero.
Exactly one of the six plaquettes
at the released link is nonidentity,
and all six at each affected pivot
are identity. The test configurations
are not asserted to be ground
profiles: they verify the local
symbolic identity which was then
applied to V84's proved minimizers.

The script also verifies the
108-edge count, positive field
sum 24, paths of length at most
eight to \(j_\perp\), and the
rational constant \(\gamma\)
in (89.12). Rational substitutions
check the Taylor decrease in
(89.16)--(89.17); they do not
assign a numerical value to the
completion-dependent \(C_2\).
All checks pass.

V84 had 1378192 bytes and SHA--256
\begin{center}\scriptsize\ttfamily
FB09FE204111B30F90F45D12B87D5D761700072BFD45E019133695179E047D3B.
\end{center}
Its entire body is retained apart
from the title version. Only active
V84 is replaced by V85. Archives,
companions, and earlier scripts remain.

\begin{firewall}[Position after V85]
The actual parent first variation,
its uniform finite decrease, and
the conditional exclusion of the
embedded slice sphere are proved
under V84's cutoff and geometry
hypotheses. The companions provide
no new full-gap input.
The coupled full-parent minimizer,
unrestricted repair estimate,
cofinal auxiliary gap, continuum
construction, and physical
Yang--Mills mass gap remain unproved.
The full objective remains active.
\end{firewall}
\section{V86: Exact commuting equations for the whole parent,
the full pivot Jacobian, and the transverse stability test}
\label{sec:f12v86}

\subsection{Why the calculation must now involve the whole parent}

V84 classified the ground sphere of a compact 56-group slice.
V85 then proved that every point of its embedding into the
944-group parent has a nonzero force in a released internal
connection. The slice answer therefore cannot be used as the
parent's stationary centre. This section goes upstream to the
completed parent equations themselves. It does not freeze the
other 888 parent groups.

The setting is still the canonical parent
\(P=\{0,1\}^4\), coarse field \(C=I\), central frozen exterior
from V85, and the same globally completed map
\(\widetilde K\). Fix the cutoff from V84:
\(0<\rho\le 1/64\), also small enough for
Theorem~\ref{thm:f12v18-pivot}. There are
\[
       n=17|P|+7|E(P)|=944,\qquad |E(P)|=96.
 \tag{90.1}
\]
Here \emph{full parent} means all \(n\) hidden group coordinates at
these fixed visible data. It does not assert stationarity in
the visible coordinates of the joint lift.

We restrict each of the \(n\) independent groups to a common
maximal torus, but retain every independent angle. We prove
that a critical point of this restriction is an actual critical
point in all \(3n\) hidden directions. We also identify the
additional transverse test which a restricted minimum must
pass. This is distinct from the flat-toron stationary-centre
result in archive f8, Section 64: the present exterior is the
frustrated central sheet, and the completion and cutoff remain
nonlinear.

\subsection{A globally single-valued scalar pivot}

Choose quaternion units \(E_1,E_2,T\) with
\(E_1E_2=T\), \(T^2=-I\), and
\(\|E_1\|_{\rm HS}=\|E_2\|_{\rm HS}=\|T\|_{\rm HS}=\sqrt2\).
The torus is \(U(u)=\exp(uT)\), with angle period \(2\pi\).
In the cutoff of V18 define the smooth periodic function
\[
 h(u)=\chi\!\left(\frac{\sqrt2|\bar u|}{\rho}\right)\bar u,
 \qquad \bar u\in(-\pi,\pi],
 \qquad g_\rho(U(u))=h(u)T.
 \tag{90.2}
\]
The expression is zero in an open neighbourhood of the branch
cut. Thus the indicated zero extension is smooth, including
at the endpoints of the angle convention. At zero it is
exactly \(u\) in a neighbourhood. The monotonicity of \(\chi\)
gives
\[
       |h(u)|\le \rho/\sqrt2,\qquad h'(u)\le1.
 \tag{90.3}
\]
The second inequality is an upper bound, not a claim that the
derivative is nonnegative in the cutoff annulus.

Let the independent angles be
\(\theta\in(\mathbb R/2\pi\mathbb Z)^n\).
The fourteen nonroot path products for a coarse edge \(e\)
have scalar angles, in any local lift,
\[
 p_{ej}(\theta)=c_{ej}+\sum_i R_{ej,i}\theta_i,
                  \qquad 1\le j\le14.
 \tag{90.4}
\]
These are the actual two-link chronological words on the fixed
tree slice. Their integer incidence coefficients and frozen
constants include both internal and boundary coordinates.
The two omitted root words are the pivot itself. No other
completed pivot occurs in a nonroot word.

\begin{proposition}[Exact scalar completion, including the cutoff]
\label{prop:f12v86-scalar}
For each \(e\), the equation
\[
 b_e+\frac1{16}\sum_{j=1}^{14}h(p_{ej}-b_e)=0
 \tag{90.5}
\]
has exactly one real solution. It satisfies
\[
 |b_e|\le r:=\frac{7\rho}{8\sqrt2},\qquad
 D_e:=1-\frac1{16}\sum_j h'(p_{ej}-b_e)\ge\frac18.
 \tag{90.6}
\]
It is smooth and periodic in all the independent angles.
The actual SU(2) pivot at \(C_e=I\) is
\(B_e=\exp(b_eT)\).
\end{proposition}
\begin{proof}
For fixed \(p_{ej}\), the left side of (90.5), as a function
of \(b\), has derivative at least \(1-14/16=1/8\).
Its perturbation of \(b\) has absolute value at most \(r\).
It is nonpositive at \(-r\), nonnegative at \(r\), and
strictly increasing on the whole real line. Existence,
uniqueness, and (90.6) follow. The implicit function theorem
gives smoothness. Changing a lift of any \(\theta_i\) changes
each \(p_{ej}\) by an integer multiple of \(2\pi\), so leaves
the unique real solution unchanged.

Substitution into the actual map
\[
 B\longmapsto B\exp\!\left[
             \frac1{16}\sum_j g_\rho(B^{-1}P_{ej})\right]
\]
shows that \(\exp(b_eT)\) maps to \(I\). The globally unique
inverse from V18 identifies it with the full group pivot,
not a newly chosen branch. Alternatively, conjugation by
the torus fixes all the path products and the target; global
uniqueness forces the inverse to commute with that torus.
The small real representative is precisely (90.5).
\end{proof}

Writing \(q_{ej}=p_{ej}-b_e\) and differentiating this exact
equation gives
\[
 \partial_i b_e
   =-\frac{\sum_j h'(q_{ej})R_{ej,i}}{16D_e},
 \qquad
 r_{ej,i}:=R_{ej,i}-\partial_i b_e,
 \tag{90.7}
\]
and
\[
 \partial_i\partial_k b_e
   =-\frac{\sum_j h''(q_{ej})r_{ej,i}r_{ej,k}}{16D_e}.
 \tag{90.8}
\]
For completeness, differentiating (90.5) twice produces
\[
 D_e\partial_i\partial_k b_e+
 \frac1{16}\sum_j h''(q_{ej})
 (\partial_i p_{ej}-\partial_i b_e)
 (\partial_k p_{ej}-\partial_k b_e)=0;
\]
the path angles are affine. This proves (90.8), including
the terms arising from differentiation of \(D_e\).
The second derivatives are bounded at fixed \(\rho\), but
this statement is not a bound uniform as \(\rho\downarrow0\).

\subsection{The inverse Jacobian has two transverse factors}

The scalar factor \(D_e^{-1}\) is not the inverse Haar
Jacobian of the three-dimensional pivot. We compute that
Jacobian as well. Define the smooth, periodically extended
quantity
\[
 a_e=\frac1{16}\sum_j h(q_{ej})\cot q_{ej},
 \quad h(q)\cot q\big|_{q=0}=1,
 \quad h(q)\cot q=0\text{ near }q=\pi,
 \tag{90.9}
\]
where the value at every multiple of \(2\pi\) is interpreted
by continuity. On the support, \(0\le \bar q\cot\bar q\le1\);
hence \(0\le a_e\le7/8\). Put
\[
       k_e=\cos b_e-a_e\,\frac{\sin b_e}{b_e},
       \qquad \frac{\sin b_e}{b_e}\bigg|_{b_e=0}=1.
 \tag{90.10}
\]

\begin{theorem}[Full SU(2) pivot density on commuting data]
\label{thm:f12v86-jacobian}
At these commuting data, the forward pivot Jacobian relative
to Haar measures is exactly \(D_e k_e^2\). Therefore V18's
inverse-pivot density factor is
\[
                 j_e=\frac1{D_e k_e^2}.
 \tag{90.11}
\]
If \(C_\chi=\sup_u|h'(u)|\), then \(C_\chi\) can be chosen
independently of this small \(\rho\), and
\[
 \frac19\le k_e\le1,\qquad
 \frac1{1+7C_\chi/8}\le j_e\le648.
 \tag{90.12}
\]
When all fourteen \(P_{ej}=I\), one has \(j_e=512\).
When they are all outside the support of \(g_\rho\) in an
open neighbourhood of \(B=I\), one has \(j_e=1\).
\end{theorem}
\begin{proof}
Identify the transverse plane spanned by \(E_1,E_2\) with
\(\mathbb C\), with \(E_1\) corresponding to 1. Multiplication
by \(i\) is the rotation sending \(E_1\) to \(E_2\).
For the left perturbation \(e^{tE_1}U(q)\), its quaternion
vector has transverse derivative
\(\cos q\,E_1-\sin q\,E_2\).
The truncated logarithm multiplies this transverse vector
by \(h(q)/\sin q\). The scalar part and the axial angle have
zero first derivative in this perturbation. Consequently
the transverse complex multiplier of \(D_Lg_\rho(U(q))\) is
\[
                     h(q)(\cot q-i).
 \tag{90.13}
\]
It extends to 1 at \(q=0\) and to zero off the cutoff support.
The axial multiplier is \(h'(q)\).

In the notation of V18, set
\(X=16^{-1}\sum_j g_\rho(B^{-1}P_{ej})=-b_eT\) and
\(A=16^{-1}\sum_jD_Lg_\rho(U(q_{ej}))\).
For an input \(B e^{tv}\), the output right-trivialized
differential of \(B\exp X\) is
\(\operatorname{Ad}_B[I-\mathcal J(X)A]v\), where
\(\mathcal J(X)=\int_0^1\operatorname{Ad}_{e^{sX}}\,ds\).
The adjoint factor is an isometry of determinant one.
The axial multiplier of \(I-\mathcal J(X)A\) is \(D_e\).
On the transverse complex plane the two multipliers are
\[
    A=a_e+i b_e,\qquad
    \mathcal J(-b_eT)=e^{-ib_e}\frac{\sin b_e}{b_e}.
 \tag{90.14}
\]
Here (90.5) was used to replace the mean of \(h(q_{ej})\)
by \(-b_e\). Their combined multiplier is exactly
\[
  1-e^{-ib_e}\frac{\sin b_e}{b_e}(a_e+i b_e)
       =e^{-ib_e}\left(\cos b_e
                          -a_e\frac{\sin b_e}{b_e}\right).
 \tag{90.15}
\]
Its real two-dimensional determinant is \(k_e^2\).
This proves (90.11), including the transverse factors.

Since \(|b_e|\le r<\pi/2\), the sinc in (90.10) is positive
and at most one. Therefore
\[
  1\ge k_e\ge \cos b_e-\frac78
       \ge \frac18-\frac{b_e^2}{2}
       \ge \frac18-\frac{49\rho^2}{256}
       \ge\frac19.
 \tag{90.16}
\]
Also \(1/8\le D_e\le1+7C_\chi/8\). The derivative in (90.3)
is \(\chi(s)+s\chi'(s)\) in the support, which supplies the
claimed \(\rho\)-independent \(C_\chi\).
Equations (90.11)--(90.12) follow. At fourteen identity
products, \(b_e=0\), \(D_e=k_e=1/8\). In the stated inactive
case \(b_e=0\), \(D_e=k_e=1\). These give the two endpoints.
\end{proof}

In particular, the commuting restriction of the actual parent
amplitude is \(\prod_{e\in E(P)}(D_e k_e^2)^{-1}\), up to
factors independent of the parent variables. If finite-\(\beta\)
stationarity is sought, the restricted full potential is
\[
 \beta S(\theta)+\sum_{e\in E(P)}
                    \log(D_e(\theta)k_e(\theta)^2),
 \tag{90.17}
\]
up to an additive constant. Replacing that correction by
\(\sum_e\log D_e\) would omit the transverse density.
The maximal torus has Haar measure zero in SU(2); integrating
(90.17) against torus measure does not compute the parent
partition function.

\subsection{All parent scalar forces and second derivatives}

For every relevant physical plaquette \(p\), let its completed
angle be
\[
 \omega_p(\theta)=c_p+\sum_i L_{pi}\theta_i
                              +\sum_e K_{pe}b_e(\theta).
 \tag{90.18}
\]
The matrices \(L,K\) are signed physical word incidences,
not effective fitted coefficients. Fixed plaquettes may be
omitted by adding a constant to the action. Let
\(J_{pi}=L_{pi}+\sum_eK_{pe}\partial_i b_e\).
The exact restricted Wilson action and its derivatives are
\[
 \begin{split}
 S(\theta)&=S_{\rm const}+2\sum_p(1-\cos\omega_p),\\
 \partial_i S&=2\sum_p\sin\omega_p\,J_{pi},\\
 \partial_i\partial_k S
   &=2\sum_p\cos\omega_p\,J_{pi}J_{pk}
       +2\sum_e\left(\sum_pK_{pe}\sin\omega_p\right)
                                  \partial_i\partial_k b_e .
 \end{split}
 \tag{90.19}
\]
These follow by the ordinary chain rule from
\(\Re\Tr U(\omega)=2\cos\omega\) and (90.18).
Equations (90.5), (90.7)--(90.8), and (90.19) specify a
closed finite system involving all 944 scalar parent forces.
The second term in the Hessian is the physical pivot force
paired with the curvature of its eliminated coordinate.
It cannot be discarded merely because the parent scalar
gradient is zero. Parent stationarity is a constrained
statement, not independent stationarity of each physical pivot.

At central configurations \(b_e=0\), \(h''(q_{ej})=0\), and
every \(\sin\omega_p=0\). Thus the scalar first derivative
vanishes and the Hessian reduces to
\[
            S''=2J^{\mathsf T}
                         \operatorname{diag}(\cos\omega_p)J.
 \tag{90.20}
\]
Even in this case the nontrivial first pivot responses in
\(J\) must still be retained. At the V85 central sheet, for
example, 88 parent pivots have fourteen identity nonroot
products and eight have fourteen negative identity products.
The corresponding inverse Jacobians are respectively 512
and 1. A released internal link has the two responses
\(-1/2\) found in V85.

\subsection{Criticality transfers, but minimality does not}

\begin{theorem}[Elementary symmetric criticality for this parent]
\label{thm:f12v86-critical}
At the fixed central exterior and \(C=I\), a critical point
of \(S(\theta)\) is a critical point of the actual completed
classical action in all \(3n\) hidden group directions.
The same assertion holds for the actual finite-\(\beta\)
potential if its full restricted expression (90.17) is used.
\end{theorem}
\begin{proof}
Simultaneous conjugation of every parent group fixes the
central exterior and the coarse data. Equivariance and
uniqueness of the completion conjugate every completed link
by the same group element. The Wilson traces, Haar measures,
and pivot Jacobians are unchanged. Thus the relevant full
function is invariant under simultaneous conjugation.

At a commuting tuple \(Y_i=\exp(\theta_iT)\), conjugation by
\(\exp(tT)\) fixes the tuple. In right exponential coordinates
\(Y_i\exp(z_iT+x_iE_1+y_iE_2)\), it fixes every \(z_i\) and
rotates every pair \((x_i,y_i)\) by the same angle \(2t\).
An invariant linear functional on this tangent space has
zero coefficient in each transverse plane: invariance under
the rotation by \(\pi\) already forces this. Its only possible
coefficients are the scalar derivatives in the \(z_i\)
directions, namely the derivatives with respect to \(\theta_i\).
If those vanish, the full first derivative vanishes.
\end{proof}

\begin{corollary}[An actual parent critical value below the slice sphere]
\label{cor:f12v86-below}
Let \(E_{\rm emb}=S_b(X_*(\mathbf n),v_*)\) denote V85's
axis-independent embedded slice ground action, and let
\(\delta_{\rm p}>0\) be V85's proved full-parent decrease.
There exists a full hidden-parent critical point \(Y^\dagger\)
at these fixed visible data with
\[
                   S_b(Y^\dagger,v_*)
                         \le E_{\rm emb}-\delta_{\rm p}.
 \tag{90.21}
\]
It can be chosen to minimize the full-parent action over all
944 independent angles in one common maximal torus.
\end{corollary}
\begin{proof}
The smoothly completed action on this \(n\)-torus has a global
minimum by compactness. The embedded slice sphere has a
representative with axis \(T\). V85's map changes only
\(Y_\ell\) to \(Y_\ell\exp[-tA(Y_j)]\).
If the inputs are in this common torus, so is the output,
because \(A(Y_j)=(Y_j-Y_j^\dagger)/2\) is a multiple of \(T\).
The decreased configuration used in V85 is therefore an
admissible competitor in the full \(n\)-torus. Its action is
at most \(E_{\rm emb}-\delta_{\rm p}\).
The torus minimum is stationary there and
Theorem~\ref{thm:f12v86-critical} proves full hidden criticality.
\end{proof}

The corollary does not identify this point numerically,
prove it noncentral, or prove it a minimum in the transverse
directions. None of those claims follows from compactness.
In particular, a smaller central configuration elsewhere in
the full parent has not been excluded by V83's central
classification on the 56-group slice.

A direct warning about minimality is useful. On SU(2)\(^2\)
consider the conjugation-invariant frustrated plaquette function
\[
       F(U,V)=2+\Re\Tr(UVU^{-1}V^{-1}).
 \tag{90.22}
\]
It is identically 4 on a common torus, so every commuting
tuple is a global minimum of its torus restriction.
For \(U=\exp(\alpha T)\), \(V=\exp(tE_1)\), quaternion
multiplication gives
\[
 F(U,V)=4-4\sin^2\alpha\,\sin^2t,\qquad
 \left.\frac{d^2}{dt^2}F\right|_{t=0}=-8\sin^2\alpha .
 \tag{90.23}
\]
For \(\alpha\notin\pi\mathbb Z\) this is a strict transverse
descent. Indeed conjugation rotates the axis of \(V\) by
\(2\alpha\), so the scalar part of the commutator is
\(\cos^2t+\sin^2t\cos(2\alpha)\), which proves (90.23).
This example explains the logical gap; it is not a claim
that the same two-variable function equals the parent action.

\subsection{One real and one complex matrix constitute the stability test}

\begin{proposition}[Full hidden Hessian and its forced orbit kernel]
\label{prop:f12v86-hessian}
At a commuting full hidden critical point, the Hessian in
the right exponential coordinates used above has the form
\[
 H=
 \begin{pmatrix}
 H_\parallel&0&0\\
 0&A&-B\\
 0&B&A
 \end{pmatrix},
 \qquad A^{\mathsf T}=A,\quad B^{\mathsf T}=-B.
 \tag{90.24}
\]
Here \(H_\parallel\) is exactly the \(n\)-by-\(n\) Hessian
in (90.19). The transverse part is the real form of the
Hermitian matrix \(\mathcal H_\perp=A+iB\).
The full Hessian is nonnegative if and only if both
\(H_\parallel\ge0\) and \(\mathcal H_\perp\ge0\).

If at least one \(Y_i\) is noncentral, put
\[
                  w_i=e^{-2i\theta_i}-1.
 \tag{90.25}
\]
Then \(w\ne0\) and \(\mathcal H_\perp w=0\).
If all \(Y_i\) are central, instead \(B=0\) and
\(A=H_\parallel\).
\end{proposition}
\begin{proof}
Invariance under the stabilizing circle implies that the
Hessian commutes with every simultaneous transverse rotation.
The mixed longitudinal--transverse block must be zero,
as is seen by using the rotation by \(\pi\). A real linear
operator on the transverse coordinates commuting with the
complex structure has the form
\(\left(\begin{smallmatrix}A&-B\\B&A\end{smallmatrix}\right)\).
Symmetry of the real Hessian is precisely
\(A^{\mathsf T}=A\), \(B^{\mathsf T}=-B\).
For \(z=x+iy\) its quadratic form is
\(z^*(A+iB)z=x^{\mathsf T}Ax+y^{\mathsf T}Ay
-2x^{\mathsf T}By\). This proves the positivity equivalence.
At a critical point these coordinate second derivatives
represent the intrinsic Hessian; the common Hilbert--Schmidt
factor 2 in the coordinate metric changes neither signs nor
kernels.

Every point on the simultaneous conjugation orbit is critical.
Differentiating this statement shows that the Hessian kills
the orbit tangent. The tangent induced by \(E_1\) at coordinate
\(i\), expressed in the right frame, is
\(\operatorname{Ad}_{Y_i^{-1}}E_1-E_1\), whose complex value
is (90.25). The tangent induced by \(E_2\) is \(iw\).
This proves the stated kernel condition and its nonzero
character when some group is noncentral.

At an all-central tuple the stabilizer is the full SU(2).
Each \(3\)-by-\(3\) block of the Hessian commutes with every
adjoint rotation in \(\mathbb R^3\). Commutation with the
half-turns about the three coordinate axes first makes the
block diagonal; rotations interchanging axes make the
three diagonal entries equal. Thus every such block is a
scalar multiple of \(I_3\). Reordering coordinates proves
\(B=0\) and \(A=H_\parallel\).
\end{proof}

\begin{corollary}[A checkable sufficient local-minimum certificate]
\label{cor:f12v86-local}
At a noncentral critical point, suppose
\[
 H_\parallel>0,\qquad
 \mathcal H_\perp\ge0,\qquad
 \ker_{\mathbb C}\mathcal H_\perp=\mathbb Cw .
 \tag{90.26}
\]
Then its conjugation orbit is an embedded two-sphere of
full hidden local minima, with positive Hessian on all
\(3n-2=2830\) normal directions. On a sufficiently small
tube about this orbit there are constants \(c,C>0\) such that
\[
 c\,\operatorname{dist}(Y,\mathcal O)^2
  \le S_b(Y,v_*)-S_b(Y^\dagger,v_*)
  \le C\,\operatorname{dist}(Y,\mathcal O)^2.
 \tag{90.27}
\]
Conversely, a negative direction of \(\mathcal H_\perp\)
at a torus minimum is a genuine full-parent classical descent.
\end{corollary}
\begin{proof}
The noncentral commuting tuple has stabilizer the circle,
and its orbit is SU(2)/U(1). More explicitly it consists
of \(\cos\theta_i I+\sin\theta_i\,\mathbf n\cdot(E_1,E_2,T)\);
at least one sine is nonzero, so \(\mathbf n\in S^2\) gives
an injective smooth immersion and hence an embedding.
By (90.24)--(90.26) the only Hessian kernel is its
two-dimensional tangent space, and the normal Hessian is
positive definite. Equivariance and compactness of the
orbit give a uniform positive normal lower bound.
Taylor expansion in a small normal tube, whose first
derivative vanishes on the orbit, gives (90.27).
For the converse take the right exponential curve with the
negative transverse tangent. Its first derivative is zero
and its second derivative is strictly negative, so Taylor's
formula gives a decrease for sufficiently small nonzero time.
\end{proof}

Neither matrix inequality in (90.26) is asserted for the
existence point in (90.21). At an all-central point,
\(H_\parallel>0\) would instead give an isolated strict
full hidden minimum, while semidefiniteness alone would not
settle the higher-order terms. This keeps the central and
noncentral cases separate rather than assuming the two-sphere
structure found on the old slice.

If (90.26) is eventually certified, the integral over that
small parent tube with the actual smooth positive amplitude
has the two-sided order
\(\beta^{-1415}\exp[-\beta S_b(Y^\dagger,v_*)]\).
To see this, use (90.27), bounded positive volume and
amplitude densities in normal coordinates, and scale the
2830 normal coordinates by \(\beta^{-1/2}\); integrating
the compact orbit contributes constants. This is only a
local tube contribution. It becomes the leading full
parent normalizer only after global minimizers and the
off-tube energy gap are independently controlled.
It is not a reinterpretation of V84's slice exponent 83.

\subsection{Verification, nonduplication, and next certificate}

The new exact script is
\begin{center}\small\ttfamily
scripts/verify\_ym\_f12\_v86\_commuting\_parent.py.
\end{center}
It enumerates the 272 internal and 672 boundary coordinates
and all 96 pivot families on a fine torus of side six.
It verifies the two duplicate root paths, fourteen nonroot
paths without pivots, and the absence of parent hidden
coordinates from exterior pivot families. Fraction arithmetic
gives the central Jacobian histogram
\[
       \#\{j_e=512\}=88,\qquad \#\{j_e=1\}=8,
 \tag{90.28}
\]
and recovers the two \(-1/2\) responses of the V85 internal
link. The exact scalar curvature of the V82 weighted
56-link direction is \(-125/64\), twice its old
unit-Hilbert--Schmidt curvature because the present angle
generator has norm \(\sqrt2\). The central released-link
scalar curvature, including both induced pivots, is 14.
Eighteen rational algebra substitutions verify the complex
factorization (90.15); these substitutions test the identity,
not numerical values of trigonometric functions at rational
angles. A rational check also tests the Hermitian real form.
All these structural tests pass. They are not numerical
or interval certificates for a new minimizer.

The new work is the exact all-coordinate nonlinear reduction,
the full SU(2) density formula on it, the hidden critical-value
consequence (90.21), and the explicit transverse stability
criterion. The geometric and central-curvature checks are
regressions, not claimed as newly discovered results.
The next useful candidate calculation must solve (90.19)
with the pivot equations retained, then test the transverse
matrix in (90.24), respecting its forced kernel. A small
scalar residual or a positive scalar Hessian by itself is
not a certified parent ground. If the transverse test fails,
its negative direction supplies the next noncommuting release.
If it passes, global ground-set coverage is still a separate
requirement before invoking V81's conditional coverage theorem.

V85 had 1395032 bytes and SHA--256
\begin{center}\scriptsize\ttfamily
35835A9A713807659274D2F6ACC3312E7054E7E87DDF875CC115C596F3819911.
\end{center}
Its body is retained apart from the title version. Only the
active V85 file is replaced by V86. Archives, companion notes,
and previous scripts are preserved. The companion audit was
performed at V85; the next companion and direction review
is due at V90.

\begin{firewall}[Position after V86]
There is now an exact commuting full-parent system, with all
nonlinear pivots and their full Jacobians, and a full hidden
critical point below the excluded embedded slice ground
energy. The transverse stability and global-minimum
certificates for that point are not proved. No new
unrestricted repair inequality, cofinal auxiliary gap,
four-dimensional continuum construction, or physical
Yang--Mills mass gap is claimed. The full objective remains
active.
\end{firewall}
\section{V87: A rational whole-parent competitor more than
58 action units below the old slice ground}
\label{sec:f12v87}

\subsection{From equations to a certified coupled configuration}

V86 provided all 944 commuting parent equations, but did
not compute a candidate or test its value. We now use those
equations to find a coupled configuration and then discard
the floating-point optimizer as a source of proof. A frozen
rational configuration, exact plaquette arithmetic, and
analytic cutoff estimates prove a whole-parent action below
570. An existing V84 supporting plane puts every embedded
slice ground above 627.854 in the same normalization.
Thus the new coupled configuration is more than 58 action
units below the old ground sphere, without assuming that
either the numerical configuration or its rational version
is stationary.

Keep V86's geometry and fixed central exterior, with
\(0<\rho\le1/64\) satisfying V18's global completion theorem.
For the explicit certificate use the fine torus
\((\mathbb Z/6\mathbb Z)^4\), the coarse parent
\(P=\{0,1\}^4\), and \(C=I\).
Use the full fine-lattice Wilson action
\[
       \mathcal S(Y)=
         \sum_{p}\left(2-\Re\Tr U_p(Y)\right).
 \tag{91.1}
\]
There are \(6\cdot6^4=7776\) unoriented elementary
plaquettes, with one fixed orientation each. At these
visible data, \(\mathcal S\) differs from the conditional
parent action \(S_b\) only by a constant independent of the
944 parent variables. All energy differences and conditional
probability statements below are therefore unchanged.
At the full central one-direction sheet \(Y^*\),
\[
                  \mathcal S(Y^*)=648.
 \tag{91.2}
\]
This is an explicit fixed-parent certificate, not an
estimate uniform in cofinal volume.

\subsection{Exploration and its strict separation from certification}

The new exploratory program
\begin{center}\small\ttfamily
scripts/explore\_ym\_f12\_v87\_commuting\_parent.py
\end{center}
forms V86's actual signed incidence matrices, solves the
96 scalar pivot equations by safeguarded Newton steps,
and evaluates the full action and its 944 derivatives.
Its gradient includes the physical pivot forces through
the implicit responses (90.7). It uses the smooth cutoff
from (22.1), not an inactive-pivot approximation.

Starting from the V82 weighted sheet perturbation, a
floating-point L-BFGS search reported
\[
 \begin{split}
  \mathcal S_{\rm initial}&\simeq647.9698970394,\\
  \mathcal S_{\rm candidate}&\simeq569.6844118643,\\
  \|\nabla_\theta\mathcal S\|_\infty&\simeq1.69\cdot10^{-6},
       \qquad
  \|\nabla_\theta\mathcal S\|_2\simeq1.28\cdot10^{-5}.
 \end{split}
 \tag{91.3}
\]
It stopped after 85 iterations because of relative objective
reduction, not a certified stationarity test. Two deterministic
directional finite differences, including a configuration
with nonlinear cutoff activity, had absolute gradient
discrepancies below \(2.7\cdot10^{-8}\). The final numerical
pivot residual was below \(2.1\cdot10^{-17}\).
These are debugging and exploration records only. They
prove neither a critical point nor a positive Hessian.
No conclusion below relies on a numerical residual being zero.

\subsection{The frozen rational hidden field}

Index all parent free links by the lexicographically sorted
list in the V86 geometry verifier: first the four-coordinate
fine base point, then its direction. Let \(k_i\) and
\(\sigma_i\in\{-1,1\}\), \(0\le i<944\), be the fixed integer
and sign arrays in
\begin{center}\small\ttfamily
scripts/ym\_f12\_v87\_rational\_parent\_data.py.
\end{center}
This frozen data file is part of the finite arithmetic
certificate. Its SHA--256 is
\begin{center}\scriptsize\ttfamily
C2E4F462C1054350AB5DC3BD80AF314D48163B925759C28C71ED7D15B50788D7.
\end{center}
There are 944 integer parameters, exactly 25 negative
central signs, and common denominator \(d=1000\).
Define the independent parent groups by
\[
 \widehat Y_i=\sigma_i\left(
       \frac{d^2-k_i^2}{d^2+k_i^2}I+
       \frac{2dk_i}{d^2+k_i^2}T\right).
 \tag{91.4}
\]
The identity
\((d^2-k^2)^2+(2dk)^2=(d^2+k^2)^2\)
proves that each group lies exactly in SU(2).
The frozen exterior and tree links remain unchanged.
Completed pivots are not replaced in the definition of the
actual configuration: they are still the unique inverses
of V18's map at these independent groups.

For a comparison calculation only, let
\(\mathcal S_{\rm fp}(\widehat Y)\) be the physical action
obtained by setting the 96 parent pivots to \(I\).
This comparison field need not satisfy the coarse constraint.
The next argument bounds its difference from the actual
completed action, rather than treating it as the actual field.

\subsection{A small-pivot certificate using exact inverse pairing}

Identify a commuting group with its pair
\((a,s)\in\mathbb Q^2\), representing \(aI+sT\).
Multiplication and inversion are
\[
 (a,s)(c,t)=(ac-st,at+sc),\qquad
                  (a,s)^{-1}=(a,-s).
 \tag{91.5}
\]
For every coarse edge \(e\), compute its fourteen nonroot
products \(P_{ej}\) from (91.4) by (91.5), with the parent
pivots omitted as in V86. These products do not depend on
any pivot.

At \(B=I\), discard identity products and every product
\((a,s)\) with
\[
                    4(1-a)\ge\frac1{4096}.
 \tag{91.6}
\]
The latter inequality means
\(\|P-I\|_{\rm HS}\ge1/64\ge\rho\).
Because the chord distance is at most the length of the
principal exponential arc, these products are outside
the support of \(g_\rho\). Their truncated logarithms
vanish at \(B=I\). Among the remaining near-identity
products, cancel inverse pairs with multiplicity.
Oddness \(g_\rho(P^{-1})=-g_\rho(P)\) makes this cancellation
exact for every allowed \(\rho\). Denote the residual
multiset by \(\mathcal R_e\), and put
\[
        u_e=\sum_{(a,s)\in\mathcal R_e}|s|,
        \qquad
        \eta_e=\min\left\{\frac5{512},u_e\right\}.
 \tag{91.7}
\]
All quantities in this test are rational.

\begin{proposition}[Exact identity pivots and rational angle bounds]
\label{prop:f12v87-pivot-bound}
If \(\mathcal R_e\) is empty, the actual completed pivot
is exactly \(I\). In all cases, writing the actual pivot as
\(B_e=\exp(b_eT)\) with V86's small real representative gives
\[
                         |b_e|\le\eta_e .
 \tag{91.8}
\]
These statements hold for every fixed cutoff
\(0<\rho\le1/64\) satisfying the completion hypotheses.
\end{proposition}
\begin{proof}
Let
\(f_e(b)=b+16^{-1}\sum_jh(p_{ej}-b)\), so
\(f_e(b_e)=0\) and \(f'_e\ge1/8\) by V86.
The cancellations just described give
\[
 |b_e|\le8|f_e(0)|
      \le\frac12\sum_{(a,s)\in\mathcal R_e}|h(p)|,
               \qquad \exp(pT)=aI+sT.
 \tag{91.9}
\]
Every residual product has \(a>0\) and \(|s|<1/2\), as
follows from the failure of (91.6). Thus its principal angle
satisfies
\(|p|=\arcsin|s|\le2|s|\). One proof of the last inequality
is to integrate
\((1-t^2)^{-1/2}\le2\) on \(0\le t\le1/2\).
Since \(|h(p)|\le|p|\), (91.9) gives \(|b_e|\le u_e\).
The general bound from (90.6) also gives
\[
       |b_e|\le\frac{7\rho}{8\sqrt2}
          \le\frac7{512\sqrt2}<\frac5{512},
 \tag{91.10}
\]
where \(49<50\). Combining the bounds proves (91.8).
If there are no residual products, \(f_e(0)=0\);
uniqueness of the scalar and group inverses proves \(B_e=I\).
It is not necessary that every individual truncated
logarithm vanish before inverse-pair cancellation.
\end{proof}

For this particular frozen field, exact arithmetic proves
that 93 residual multisets are empty. The other three
have one product each. Outward-rounded rational upper
bounds for their pivot angles are as follows; a coarse
edge is given by its owner \(r\) and direction \(\mu\).
\[
\begin{array}{c|c|c}
 r&\mu&\text{upper bound for }|b_e|\\ \hline
 (0,0,0,1)&1&1/250\\
 (1,0,0,0)&1&1/250\\
 (1,0,0,1)&1&1/500
\end{array}
\qquad
                 \sum_e|b_e|\le\frac1{100}.
 \tag{91.11}
\]
The three remaining pivots are only potentially nonidentity;
we do not assert that these upper bounds are attained.

\begin{proposition}[Physical action debit for the actual pivots]
\label{prop:f12v87-action-debit}
For the rational field (91.4),
\[
 \left|\mathcal S(\widehat Y)
                -\mathcal S_{\rm fp}(\widehat Y)\right|
                       \le\frac3{25}.
 \tag{91.12}
\]
\end{proposition}
\begin{proof}
Replacing one physical pivot \(B_e\) by \(I\) changes each
incident plaquette trace by at most
\(\sqrt2\|B_e-I\|_{\rm HS}\).
This follows from the Hilbert--Schmidt Cauchy--Schwarz
inequality after multiplying by the other unitary factors;
inverse occurrences have the same bound.
For \(B_e=\exp(b_eT)\),
\(\|B_e-I\|_{\rm HS}\le\sqrt2|b_e|\), so the trace debit
is at most \(2|b_e|\) per plaquette. There are six plaquettes
incident to a physical link in dimension four.
Replace the pivots one at a time and telescope. The other
factors remain unitary at every replacement, even when
several changing pivots occur in one plaquette. The total
debit is at most
\(12\sum_e|b_e|\le12/100=3/25\), proving (91.12).
\end{proof}

\subsection{The exact finite arithmetic certificate}

The independent verification program
\begin{center}\small\ttfamily
scripts/verify\_ym\_f12\_v87\_rational\_parent.py
\end{center}
imports the frozen integers and earlier geometry, not the
optimizer. It checks all 944 unit-circle identities,
forms all 7776 physical plaquette words in rational
arithmetic, and performs the inverse-pair test above.
For a rational plaquette energy \(q_p\) let \(N=10^6\).
Instead of accumulating a large common denominator it uses
the exact enclosing sums
\[
 \frac1N\sum_p\lfloor Nq_p\rfloor
    \le\mathcal S_{\rm fp}(\widehat Y)
    \le\frac1N\sum_p\lceil Nq_p\rceil .
 \tag{91.13}
\]
Integer division of the numerator by the denominator of
each fraction computes both bounds without floating point.
The resulting enclosure is
\[
 \frac{142421221}{250000}
      \le\mathcal S_{\rm fp}(\widehat Y)
      \le\frac{142421761}{250000}.
 \tag{91.14}
\]
The same outward-rounding rule applied to (91.7) yields
(91.11). Equations (91.12)--(91.14) prove
\[
 \boxed{\displaystyle
  \mathcal S(\widehat Y)
     \le\frac{142451761}{250000}
       =569.807044<570.}
 \tag{91.15}
\]
The displayed decimal is an exact terminating decimal of
the rational upper bound, not a rounded numerical minimum.
This certificate is also uniform over the stated cutoff
range: it uses only the support, oddness, and monotonicity
properties already proved, not evaluations of the smooth
cutoff at transcendental angles.

\subsection{Comparison with the old slice ground without re-solving it}

Recall V84's convex scalar slice functional \(F\), its
unique minimizer \(a_*\), and
\[
        E_{\rm emb}
          =648+2\bigl(F(a_*)+132\bigr).
 \tag{91.16}
\]
This is the full-action normalization of the embedded
slice sphere. V84 already proved that its minimizing
pivots equal \(I\), so no new pivot debit is needed in
(91.16).

Reuse the frozen V84 supporting point \(c\) indexed by
the site \((0,0,3)\), with
\(c_i=(k_i^2-d^2)/(k_i^2+d^2)\),
\(s_i=2k_id/(k_i^2+d^2)\), and \(d=100000\).
Here the 56 \(k_i\)'s are the earlier V84 array, not the
new 944-element array. With \(g=\nabla F(c)\), convexity on
the closed cube gives the unrestricted supporting-plane bound
\[
 F(a)\ge F(c)+g\cdot(a-c)
       \ge F(c)-\sum_i|g_i|(1+|c_i|).
 \tag{91.17}
\]
Unlike the face-constrained V84 use of this point, the
debit now includes every index, including \((0,0,3)\).
The same rational value-and-gradient formulas from V84,
evaluated exactly in the new verifier, give
\[
 F(c)-\sum_i|g_i|(1+|c_i|)
             \ge-\frac{142073}{1000}.
 \tag{91.18}
\]
This is a reuse of a previously frozen supporting plane,
not a new minimization or a numerical estimate of \(a_*\).
Substituting into (91.16) yields
\[
                    E_{\rm emb}\ge\frac{313927}{500}
                                     =627.854.
 \tag{91.19}
\]

\begin{theorem}[Certified coupled energy separation and critical value]
\label{thm:f12v87-separation}
At the specified fixed parent and visible data,
\[
 E_{\rm emb}-\mathcal S(\widehat Y)
       \ge\frac{14511739}{250000}
          =58.046956>58.
 \tag{91.20}
\]
The full parent ground energy is at most
\(142451761/250000\). There also exists a commuting
full hidden-parent critical point of action at most that
number.
\end{theorem}
\begin{proof}
Subtract (91.15) from (91.19) to obtain (91.20).
The actual completed rational field is admissible, so it
bounds the infimum of the full compact parent action above.
It is also admissible in the torus of all 944 independent
angles. A global minimum on that compact torus exists and
has no greater action. Theorem~\ref{thm:f12v86-critical}
then gives criticality in all 2832 hidden group directions.
This last point need not equal the rational witness and
need not be a minimum in the transverse directions.
\end{proof}

\subsection{A stronger conditional exclusion, without a dynamics claim}

The action separation also sharpens the exclusion of the
old embedded orbit. Fix one allowed \(\rho\).
Let \(\mathcal O_{\rm emb}\) be that compact two-sphere
in the full parent hidden space \(\mathcal X=SU(2)^{944}\).
The actual conditional law is
\[
  \pi_\beta(dY\mid v)
     =Z_\beta(v)^{-1}a(Y,v)e^{-\beta S_b(Y,v)}\,dY,
 \tag{91.21}
\]
with its full positive inverse-pivot amplitude.

\begin{corollary}[Uniform local exclusion of the embedded slice sphere]
\label{cor:f12v87-exclusion}
There are a neighbourhood \(U\) of \(\mathcal O_{\rm emb}\),
a neighbourhood \(V\) of the fixed visible data \(v_*\),
and a finite constant \(C\), all independent of \(\beta\),
such that for \(v\in V\) and \(\beta\ge1\),
\[
             \pi_\beta(U\mid v)
                  \le\min\{1,C\beta^{1416}e^{-58\beta}\}.
 \tag{91.22}
\]
These neighbourhoods and constants are not asserted uniform
as \(\rho\downarrow0\), nor uniform in growing parent size.
\end{corollary}
\begin{proof}
At \(v_*\), (91.20) gives a strict margin above 58 between
the entire old orbit and the one fixed rational hidden
configuration \(\widehat Y\). Smooth global completion,
continuity, and compactness of the orbit give neighbourhoods
with compact closures, shrunk if necessary, on which
\[
        S_b(Y,v)-S_b(\widehat Y,v)\ge58
                      \quad(Y\in U,\ v\in V).
 \tag{91.23}
\]
The fixed compact parent amplitude is bounded above and below
by positive constants on these visible data.
Let \(E(v)=\min_Y S_b(Y,v)\). At any full hidden global
minimum the first derivative is zero; the uniform fixed-parent
second derivative bound gives a quadratic upper bound for
the action in a small geodesic ball about that minimum.
Product Haar volume of a radius \(r\) ball is bounded below
by a fixed positive multiple of \(r^{2832}\).
Integrate over a radius proportional to \(\beta^{-1/2}\)
to obtain the same elementary lower normalization bound
used in V81:
\[
             Z_\beta(v)\ge c\beta^{-1416}e^{-\beta E(v)}.
 \tag{91.24}
\]
No nondegeneracy or uniqueness of the minimizing point is
needed for this lower bound.
The numerator over \(U\) is at most
\(C_0e^{-\beta(S_b(\widehat Y,v)+58)}\) by (91.23).
Since \(E(v)\le S_b(\widehat Y,v)\), division by (91.24)
gives (91.22). The trivial probability upper bound gives
the minimum with 1.
\end{proof}

Equation (91.22) is conditional concentration in a fixed
finite parent. It gives no hitting-time estimate from the
old orbit, no original-clock heat-bath conductance, and no
cofinal spectral gap. In particular, V79's obstruction to
globally avoiding unions of local defects is not bypassed
by a larger local energy separation.

\subsection{Validation and remaining upstream task}

The exact V87 verifier passes. It checks the 944-group
field, all plaquettes, all pivot residual multisets,
the action enclosures, and the reused V84 lower bound.
V86's exact structural verifier also remains passing.
The numerical evaluator and the rational verifier serve
different purposes and have separate imports.
The certificate can be checked without SciPy or a rerun
of the optimizer.

V86 had 1418441 bytes and SHA--256
\begin{center}\scriptsize\ttfamily
079E20EF343045E12CD1E549C3C37347022EBDE6901E5E1906FE08C4F7FA92D7.
\end{center}
Its body is preserved apart from the title version. Only
active V86 is replaced by V87; frozen archives, companions,
and earlier scripts remain. The next companion audit and
direction review is due at V90.

\begin{firewall}[Position after V87]
The former slice ground is now separated from an explicit
actual whole-parent competitor by more than 58 action units.
The resulting critical-value existence and conditional
exclusion statements are proved in their stated finite
geometry. The floating-point candidate has not been certified
stationary, and its transverse Hessian has not been tested
or proved positive. The next upstream task is the full
transverse calculation at a properly certified coupled
critical point, or a certified transverse descent if it
fails. A global parent ground classification, unrestricted
repair estimate, cofinal auxiliary gap, continuum construction,
and physical Yang--Mills mass gap remain unproved.
The full objective remains active.
\end{firewall}
\section{V88: Four analytic inverse-root pivots, the full
transverse quadratic form, and a nonstationary symmetry defect}
\label{sec:f12v88}

\subsection{The transverse test must retain the pivot acceleration}

V87 supplied a rational whole-parent energy certificate,
not a critical-point certificate. We now calculate the
full transverse second variation near a small rational
adjustment of that field. All 944 independent parent
coordinates are included. Four induced pivots have a
particularly simple exact local description: each is an
inverse fifteenth root of one physical two-link product.
The resulting second variation contains both their first
responses and their axial accelerations.

The calculation also exposes a useful diagnostic trap.
A small negative eigenvalue occurs at the adjusted rational
point. An exact symmetry identity shows that this point is
not stationary, so the eigenvalue does not certify instability
of a nearby critical point. We certify a finite noncommuting
decrease from this particular point, while keeping the
stationary-centre and local-minimum questions separate.
The setting, full-action normalization \(\mathcal S\), and
cutoff assumptions remain those of V87.

\subsection{Three exact rational adjustments remove the cutoff ambiguity}

Keep the frozen V87 signs and denominator \(d=1000\).
Change only the following entries of its zero-based,
lexicographically indexed integer array:
\[
 \begin{array}{c|r|r|c}
 i&k_i\text{ in V87}&k_i\text{ now}&
              \text{physical link }(x,\mu)\\ \hline
 265&2&0&((1,1,0,3),1)\\
 703&-11&-9&((3,1,0,1),1)\\
 709&-1&0&((3,1,0,3),1).
 \end{array}
 \tag{92.1}
\]
All three signs are positive. Denote the adjusted independent
field by \(Y^\circ\). Formula (91.4) still defines each group
exactly; no new large parameter array is needed.

The exact geometry routine
\begin{center}\small\ttfamily
scripts/ym\_f12\_v88\_transverse\_geometry.py
\end{center}
checks all \(96\cdot14=1344\) nonroot path products.
Exactly four now equal \(I\), one on each coarse edge
\[
                   e=((a,0,0,b),1),
                       \qquad a,b\in\{0,1\}.
 \tag{92.2}
\]
The corresponding two physical links are
\[
 A_e=((2a+1,1,0,2b+1),1),\qquad
 B_e^{\rm path}=((2a+1,2,0,2b+1),1).
 \tag{92.3}
\]
Here \(B_e^{\rm path}\) names the second path link, not the
completed pivot; below we denote that pivot by \(\mathcal B_e\).
All other 1340 products have chord-squared distance from
\(I\) strictly greater than \(1/4096\). In fact the verifier
proves the uniform rational margin
\[
 \|P_{ej}(Y^\circ)-I\|_{\rm HS}^2-\frac1{4096}
       \ge\kappa:=\frac{14527}{100000000}>0
 \tag{92.4}
\]
for every one of those 1340 products.
Thus all 96 actual completed pivots at \(Y^\circ\)
are exactly \(I\), for every allowed \(\rho\).

\begin{proposition}[An explicit analytic branch of the actual completion]
\label{prop:f12v88-root}
Put
\[
                  \delta_\rho
                     =\min\{\rho/32,\kappa/12\}.
 \tag{92.5}
\]
On the groupwise neighbourhood
\(\max_i\|Y_i-Y_i^\circ\|_{\rm HS}<\delta_\rho\),
the 92 pivots not listed in (92.2) remain exactly \(I\).
For an edge in (92.2), let \(P_e(Y)\) be the product of
the two links in (92.3). The actual pivot is
\[
                    \mathcal B_e(Y)
                         =P_e(Y)^{-1/15}
                         :=\exp[-\tfrac1{15}\log P_e(Y)].
 \tag{92.6}
\]
The branch is principal near \(I\). These formulas hold
for arbitrary nearby SU(2) parent variations, not only
commuting ones. In particular the completed action is
real analytic on this neighbourhood.
\end{proposition}
\begin{proof}
Each two-link product changes in Hilbert--Schmidt norm
by at most \(2\delta_\rho\), by telescoping unitary factors.
For a product initially equal to \(I\), the principal
logarithm satisfies
\(\|\log P_e\|_{\rm HS}\le2\|P_e-I\|_{\rm HS}
 \le4\delta_\rho\).
The first inequality follows from the scalar exponential
chord formula on this small neighbourhood.
For the proposed pivot (92.6),
\[
 \|\log(\mathcal B_e^{-1}P_e)\|_{\rm HS}
        =\frac{16}{15}\|\log P_e\|_{\rm HS}
        \le\frac{64}{15}\delta_\rho<\rho/2.
 \tag{92.7}
\]
Thus this one truncated logarithm is the ordinary logarithm.

If \(d_{ej}\) is the initial chord distance of any other
product, then (92.4) gives \(d_{ej}^2-\rho^2\ge\kappa\).
The SU(2) chord diameter is \(2\sqrt2\), and
\(2\sqrt2+\rho<3\). Consequently
\[
                  d_{ej}-\rho\ge\kappa/3.
 \tag{92.8}
\]
The proposed pivot moves by at most \(4\delta_\rho/15\).
The new relative product therefore remains at chord
distance at least
\(d_{ej}-2\delta_\rho-4\delta_\rho/15>\rho\);
indeed \(34\delta_\rho/15<\kappa/3\) by (92.5).
All its truncated logarithms vanish.

The actual pivot equation on an edge in (92.2) therefore
reduces, for this proposed value, to
\[
 \mathcal B_e\exp\!\left[
        \frac1{16}\log(\mathcal B_e^{-1}P_e)\right]
       =P_e^{-1/15}\exp[\tfrac1{15}\log P_e]=I.
 \tag{92.9}
\]
Global uniqueness from V18 identifies this solution with
the actual completed pivot. On the other 92 edges the
proposed value \(I\) has all relative products outside
the cutoff, hence also solves the actual equation.
Principal logarithm and exponential are analytic here.
This proves all the assertions.
\end{proof}

At \(\rho=1/64\), (92.5) gives the explicit radius
\(14527/1200000000\). This is a conservative certified
neighbourhood, not the maximal extent of the analytic
branch. Its size is not a claim about a Newton basin.

\subsection{Exact first and second pivot responses}

At \(Y^\circ\), use the quaternion units from V86 and
right-geodesic variations
\[
        Y_i(t)=Y_i^\circ\exp(t v_i),\qquad
        v_i=x_iE_1+y_iE_2,\quad z_i=x_i+iy_i .
 \tag{92.10}
\]
Every transverse first derivative of the completed action
vanishes by the stabilizing-circle argument in V86, even
if its longitudinal derivative is nonzero.

For an active path write its base links as \(A\) and
\(B=A^{-1}\), with tangents \(v_a,v_b\); frozen link
tangents are interpreted as zero. Set
\[
        \alpha_e=e^{2i\theta_a},\qquad
        u_e=\alpha_e z_a+z_b,\qquad
        q_e=2\Im(\overline{\alpha_e z_a}\,z_b).
 \tag{92.11}
\]
All entries of \(\alpha_e\) are rational at \(Y^\circ\),
since \(A=aI+sT\) is rational and
\(\alpha_e=(a+is)^2\).

\begin{proposition}[The induced transverse velocity and axial acceleration]
\label{prop:f12v88-jets}
The completed pivot has transverse right velocity
\(-u_e/15\). In quaternion second-derivative notation,
\[
 \begin{split}
  \mathcal B_e(0)&=I,\\
  \mathcal B'_e(0)&=-\tfrac1{15}(\Re u_e\,E_1+\Im u_e\,E_2),\\
  \mathcal B''_e(0)&=
       -\frac{|u_e|^2}{225}I-\frac{q_e}{15}T .
 \end{split}
 \tag{92.12}
\]
The last axial term is an actual completion acceleration.
It is not part of the geodesic with the same initial
pivot velocity.
\end{proposition}
\begin{proof}
For the two-link product,
\[
 P_e(t)=A e^{t v_a}A^{-1}e^{t v_b}
       =e^{t\operatorname{Ad}_A v_a}e^{t v_b}.
 \tag{92.13}
\]
Its first derivative is the transverse vector \(u_e\).
Its second quaternion derivative has scalar part
\(-|u_e|^2\) and axial part
\(2(\operatorname{Ad}_A v_a\times v_b)_3T=q_eT\).
At \(P_e(0)=I\), the logarithm has the same first vector
derivative and the same second vector derivative.
Differentiate (92.6), then exponentiate the resulting
logarithmic jet. The exponential contributes
\(-|u_e|^2I/225\) to the scalar second derivative and
leaves the logarithmic axial acceleration \(-q_eT/15\).
This is (92.12).
\end{proof}

\subsection{An exact assembly rule for the whole transverse matrix}

Temporarily treat the 944 free groups and the four active
physical pivots as 948 independent physical variables.
At their base values define the Hermitian transverse
matrix \(H_{\rm phys}\) by the following elementary
plaquette rule. If an oriented plaquette word is
\(Q_1Q_2Q_3Q_4\), write \(c_p\) for the scalar part of
its product. Each variable occurrence adds \(2c_p\)
to the corresponding diagonal. For two distinct variable
occurrences, replace their factors by the respective
first derivatives in the \(E_1,E_1\) directions, and let
\(d_{11}\) be the scalar part of the resulting word.
Let \(d_{12}\) be the analogous scalar part for
\(E_1,E_2\). Their Hermitian off-diagonal contribution is
\[
                     -2d_{11}+2i d_{12},
 \tag{92.14}
\]
with the conjugate contribution in the transposed entry.
For a forward factor \(Q=Y\), the derivative is \(Y E_j\);
for a reverse factor \(Q=Y^{-1}\), it is \(-E_jY^{-1}\).
These rules include every physical plaquette with a
variable occurrence.

Let \(J_\perp:\mathbb C^{944}\to\mathbb C^{948}\) be
identity on free coordinates and have pivot row
\[
                (J_\perp z)_e
                      =-(\alpha_e z_a+z_b)/15.
 \tag{92.15}
\]
Let \(f_e\) be the physical axial pivot force, with the
other physical links temporarily independent:
\[
             f_e=2\sum_{p\ni\mathcal B_e}
                         \epsilon_{pe}\sin\omega_p .
 \tag{92.16}
\]
This is rational at \(Y^\circ\). Here
\(\epsilon_{pe}\) is the oriented plaquette incidence and
\(\omega_p\) its commuting angle at the base point.

\begin{theorem}[Full completed transverse quadratic form]
\label{thm:f12v88-transverse}
The actual completed second variation at \(Y^\circ\) is
\[
 \begin{split}
 Q_\perp(z)
     &=z^*J_\perp^*H_{\rm phys}J_\perp z
        -\frac2{15}\sum_{e\ {\rm active}}
             f_e\Im(\overline{\alpha_e z_a}\,z_b)\\
     &=z^*H_\perp z.
 \end{split}
 \tag{92.17}
\]
In particular every entry of the \(944\)-by-\(944\)
Hermitian matrix \(H_\perp\) is rational in its real
and imaginary parts. The acceleration term adds
\(i f_e\overline{\alpha_e}/15\) to entry \(ab\) and
its conjugate to entry \(ba\), with zero tangents
omitted when an endpoint is frozen.
\end{theorem}
\begin{proof}
For a physical right geodesic, each second derivative
of a group factor is \(-|z_i|^2Q_i\). The second product
rule gives the diagonal contribution \(2c_p|z_i|^2\).
A pair of first derivatives contributes \(-4d_{11}x_ix_j\)
or \(-4d_{12}x_iy_j\) to the real quadratic form.
The circle-invariant Hermitian representation from V86
therefore gives exactly (92.14).

The free coordinates follow right geodesics. Inserting
the first pivot responses (92.15) into this physical
geodesic Hessian gives its pullback
\(J_\perp^*H_{\rm phys}J_\perp\).
Equation (92.12) shows that the remaining acceleration
is axial, with coefficient \(-q_e/15\).
Pairing it with the physical force (92.16) gives the
second term in (92.17). Finally,
\(2\Re(i f_e\overline{\alpha_e}\,\overline z_a z_b/15)
 =-2 f_e\Im(\overline{\alpha_e z_a}z_b)/15\),
which proves the entry formula. Every operation uses
rational quaternion entries, integer incidences, and
division by 15, so the rationality assertion follows.
\end{proof}

There is an equally simple longitudinal formula on this
branch. If \(D\) is the physical signed plaquette incidence
matrix, let \(J_0\) be identity on free angles and have
entries \(-1/15\) at both endpoints of each active path.
With \(M=DJ_0\), local angle coordinates \(\xi\) satisfy
\[
 \begin{split}
  \mathcal S(\xi)&=\mathcal S_{\rm const}
                  +2\sum_p[1-\cos(\omega_p^\circ+(M\xi)_p)],\\
  \nabla\mathcal S(\xi)&=2M^{\mathsf T}
                        \sin(\omega^\circ+M\xi),\\
  \nabla^2\mathcal S(\xi)&=
         2M^{\mathsf T}\operatorname{diag}
                      (\cos(\omega^\circ+M\xi))M .
 \end{split}
 \tag{92.18}
\]
The scalar pivot angles are affine here, so their scalar
second derivatives vanish. This does not remove the
noncommuting acceleration in (92.17).
\subsection{The exact symmetry identity away from a critical point}

Write \(g_i=\partial_{\theta_i}\mathcal S(Y^\circ)\),
and define the transverse vector
\[
                  w_i=e^{-2i\theta_i}-1.
 \tag{92.19}
\]
At a critical point V86 gives \(H_\perp w=0\).
At a noncritical point even the quadratic form on \(w\)
need not vanish. The following identity specifies its defect.

\begin{proposition}[Orbit-geodesic curvature equals a longitudinal-force pairing]
\label{prop:f12v88-ward}
At any commuting parent point with this central exterior,
the completed right-geodesic transverse quadratic form obeys
\[
                 Q_\perp(w)
                    =2\sum_i g_i\sin(2\theta_i).
 \tag{92.20}
\]
In particular, a strictly negative value on \(w\) proves
that the point is not stationary; it does not prove
instability of a different, nearby critical point.
\end{proposition}
\begin{proof}
Consider the actual simultaneous conjugation curve
\(Y_i(s)=e^{sE_1}Y_i e^{-sE_1}\).
Its right logarithmic coordinate has expansion
\[
 \log(Y_i^{-1}Y_i(s))
   =s(\operatorname{Ad}_{Y_i^{-1}}E_1-E_1)
                         -s^2\sin(2\theta_i)T+O(s^3).
 \tag{92.21}
\]
Indeed the product is
\(e^{s\operatorname{Ad}_{Y_i^{-1}}E_1}e^{-sE_1}\);
the second-order logarithm term is half the commutator.
Using
\(\operatorname{Ad}_{Y_i^{-1}}E_1
 =\cos(2\theta_i)E_1-\sin(2\theta_i)E_2\)
and \([E_2,E_1]=-2T\) gives (92.21).
Its initial transverse tangent is \(w\), but its right
logarithmic acceleration is
\(-2\sin(2\theta_i)T\), not zero.

The completed action is constant on this conjugation
curve because the exterior and coarse data are central.
Twice differentiating gives
\(0=Q_\perp(w)-2\sum_i g_i\sin(2\theta_i)\).
This proves (92.20).
\end{proof}

\subsection{Exact curvature and energy certificates at the adjusted point}

The independent rational verifier
\begin{center}\small\ttfamily
scripts/verify\_ym\_f12\_v88\_full\_transverse.py
\end{center}
uses derivative triples
\((Q(0),Q'(0),Q''(0))\).
Quaternion multiplication obeys
\[
 \begin{split}
 (UV)'&=U'V+UV',\\
 (UV)''&=U''V+2U'V'+UV''.
 \end{split}
 \tag{92.22}
\]
Quaternion inversion along a unitary path is conjugation
of each derivative. For free right geodesics it inserts
\((Y,Yv,-|v|^2Y)\); for the four pivots it inserts
(92.12). It then differentiates all 7776 actual physical
plaquette words.

For the vector (92.19), both endpoints of every active
path have cancelling first responses. More strongly,
\(\operatorname{Ad}_A v_a=-v_b\), so (92.13) is identically
\(I\) along these particular free geodesics. All four
pivot derivative triples have zero first and second
derivatives. This special cancellation is not assumed
for general transverse vectors.

Using outward integer rounding at denominator \(10^9\)
for each rational plaquette contribution, the verifier gives
\[
 \begin{split}
 \frac{569686226279}{10^9}
      &\le\mathcal S(Y^\circ)
               \le\frac{14242155711}{25000000}<569.7,\\
 -\frac{8290081}{10^8}
      &\le Q_\perp(w)
               \le-\frac{82898649}{10^9}< -0.082 .
 \end{split}
 \tag{92.23}
\]
The first variation in this transverse direction is
exactly zero. The negative second variation proves a
local transverse descent from \(Y^\circ\), while
(92.20) proves \(g\ne0\).

The checker also verifies (92.20) separately on each
plaquette, using its completed scalar derivative.
This avoids comparing two rounded global quantities.
For independent probes which assign \(E_1\) and \(E_2\)
to the two endpoints of one active path, the acceleration
in (92.12) is nonzero. For example, on the edge
\(((0,0,0,0),1)\), the total action contribution omitted
by deleting that acceleration is enclosed by
\[
  -\frac{1791049}{10^9}
     \le \text{acceleration contribution}
     \le-\frac{1791043}{10^9}.
 \tag{92.24}
\]
The other three endpoint probes likewise have strictly
negative certified contributions. Thus the zero pivot
jet for the special vector \(w\) cannot justify dropping
the acceleration from the full matrix (92.17).

\subsection{A finite rational noncommuting decrease}

The strict infinitesimal descent can be realized by
an explicit rational endpoint. Let
\(v_i=\Re w_i E_1+\Im w_i E_2\), and define the Cayley path
\[
 R_i(t)=
   \frac{(1-t^2|w_i|^2/4)I+t v_i}
        {1+t^2|w_i|^2/4},
             \qquad Y_i(t)=Y_i^\circ R_i(t).
 \tag{92.25}
\]
Its quaternion norm is exactly one. Its first two
derivatives at zero agree with those of
\(\exp(t v_i)\), so (92.23) applies locally.
At the rational parameter
\[
                            t=\frac1{1024},
 \tag{92.26}
\]
all independent groups remain rational.

\begin{theorem}[An explicit noncommuting, fully completed endpoint]
\label{thm:f12v88-finite}
At (92.26), all 96 actual completed pivots equal \(I\),
the hidden configuration is noncommuting, and
\[
 -\frac{39071339}{10^{15}}
   \le \mathcal S(Y(t))-\mathcal S(Y^\circ)
   \le-\frac{39069178}{10^{15}}
       <-\frac{39}{10^9}.
 \tag{92.27}
\]
This is a finite energy decrease at the stated endpoint,
not a claim about a heat-bath trajectory or a barrier
along the entire intervening path.
\end{theorem}
\begin{proof}
For the four active products,
\(\operatorname{Ad}_A v_a=-v_b\).
Cayley transforms respect conjugation and inversion.
Thus
\[
          A R_a(t)A^{-1}R_b(t)
                       =R_{-v_b}(t)R_{v_b}(t)=I .
 \tag{92.28}
\]
This holds for every real \(t\).
For the other 1340 products the verifier forms their
exact rational quaternion values at (92.26) and checks
\(\|P_{ej}(Y(t))-I\|_{\rm HS}^2>1/4096\).
At this endpoint all truncated logarithms at proposed
pivot \(I\) therefore vanish, or are the logarithm of \(I\).
V18's global uniqueness gives the actual pivots \(I\).
This endpoint check does not assert the cutoff inequalities
at every intermediate value of \(t\).

All physical plaquettes at both endpoints are now rational
words. Subtract their exact energies first, then round
each rational difference outward with denominator \(10^{15}\).
Summing the resulting integer bounds gives (92.27).

To check noncommutation without a numerical tolerance,
write \(Y_i^\circ=a_iI+s_iT\). Since
\(w_i=-2s_i^2-2ia_is_i\), multiplication in (92.25)
gives the imaginary quaternion vector
\[
 \operatorname{Im}Y_i(t)
   =\frac{s_i}{1+t^2s_i^2}
                    (0,-2t,1-t^2s_i^2).
 \tag{92.29}
\]
Two such nonzero vectors have nonzero cross product
whenever \(s_i^2\ne s_j^2\) and \(t\ne0\).
The unchanged indices \(i=2,3\) have parameters
\(-358,-64\), so their nonzero squared sine values
are distinct by rational arithmetic. Their quaternion
products do not commute. This proves the noncommuting
assertion as well as the independently checked pair
test in the verifier.
\end{proof}

\subsection{What the numerical matrix tests do and do not say}

Before adjusting the field, a deliberately pivot-constant
transverse restriction froze 73 parent groups and varied
871. It excluded the three ambiguous pivot stencils and
all changing plaquettes incident to them. Its three
smallest numerical eigenvalues were approximately
\(0.09819,0.21048,0.24292\). This restricted computation
did not yield a descent and did not test all parent
directions.

The subsequent complete calculation at \(Y^\circ\),
using (92.17), reported the three smallest transverse
eigenvalues
\[
             -0.00018991,\qquad0.17347,\qquad0.19039.
 \tag{92.30}
\]
The two smallest longitudinal eigenvalues were approximately
\(0.25392,0.26363\), while the scalar gradient had
maximum component about \(0.05754\) and Euclidean norm
about \(0.23470\). Pinning transverse coordinate 156
gave a smallest numerical principal-minor eigenvalue
about \(0.01535\).
These are diagnostics from the two exploratory scripts
\begin{center}\small\ttfamily
scripts/explore\_ym\_f12\_v88\_transverse\_release.py\\
scripts/explore\_ym\_f12\_v88\_full\_transverse.py.
\end{center}
They are not verified eigenvalue enclosures and are not
used as proofs of positivity or a spectral gap.

The exact conclusion is more limited and more useful for
the next step: (92.23) and (92.20) prove that this adjusted
rational point has not been centred. The finite decrease
(92.27) does not establish that the torus minimum from
V86 is transversely unstable. Nor do the positive numerical
principal-minor values prove that it is stable.
One must first certify a zero of the longitudinal system
(92.18) within a verified completion branch and then test
the full transverse form at that same point. Near a
noncentral critical point, the exact forced orbit kernel
must be separated from the normal directions as in V86.

\subsection{Validation, preservation, and the next upstream certificate}

The exact V88 verifier passes without importing NumPy,
SciPy, or an optimizer. It checks the three integer
adjustments, all 1344 path products, the analytic-branch
margin, all 7776 Ward identities, all 7776 acceleration
identities for the orbit vector, four independent
acceleration probes, and the finite Cayley endpoint.
The negative-curvature and finite-decrease statements
are rational certificates at a specified nonstationary
point, not surrogate certificates for a nearby minimum.

V87 had 1435585 bytes and SHA--256
\begin{center}\scriptsize\ttfamily
2E39EB17622BF61698C247B121E2B9FD462E2FB6195E3B5251DAA41FC339D967.
\end{center}
Its body is retained apart from the title version. Only
active V87 is replaced by V88. The V87 frozen data,
archives, companion notes, and all earlier scripts remain.
The next companion and direction review is due at V90.

\begin{firewall}[Position after V88]
The full noncommuting pivot responses and completed
transverse quadratic form are now explicit on a certified
analytic neighbourhood. A rational low-energy point and
a finite noncommuting decrease from it are certified.
The exact symmetry defect shows why that point cannot
yet be used as a stationary centre.
A validated stationary-centre construction and its
transverse normal positivity remain open, as do global
parent ground control, unrestricted repair, a cofinal
auxiliary gap, continuum construction, and the physical
Yang--Mills mass gap. The full objective remains active.
\end{firewall}
\section{V89: A validated noncentral stationary centre
for all 944 hidden parent coordinates}
\label{sec:f12v89}

\subsection{What is certified in this iteration}

V88 showed why the rational test point could not be used
as a stationary centre. We now certify a zero of the
longitudinal system in a specified ball, with a positive
longitudinal Hessian throughout that ball. The stabilizing
circle then gives criticality in all 2832 hidden group
directions. This is a located, validated critical point,
not just the compactness existence statement of V86.

The proof has three separate parts: an explicit SU(2)
global pivot bound justifying the selected cutoff; rational
cutoff checks on the whole proposed ball; and an integer
matrix certificate combined with an outward-gradient
existence argument. A floating-point Newton calculation
only proposes the centre of the ball. It is not used as
a critical-point certificate.

The parent, central exterior, coarse field \(C=I\), and
full physical action \(\mathcal S\) are those of V87--V88.
In this section the cutoff is fixed at
\[
                              \rho=1/64.
 \tag{93.1}
\]
Unlike V87's energy comparison, the present stationary
point certificate is not asserted for every smaller
cutoff. All stationarity claims concern the classical
action at these fixed visible data, not the finite-\(\beta\)
potential including the logarithm of the pivot amplitude.

\subsection{An explicit SU(2) global pivot threshold}

V18 gave a qualitative sufficiently-small-cutoff theorem.
To use the definite value (93.1), we first strengthen that
theorem for SU(2). No derivative bound on the cutoff
annulus is silently assigned a numerical value.

\begin{theorem}[Uniform SU(2) pivot diffeomorphisms up to \(\rho=1/4\)]
\label{thm:f12v89-global-pivot}
For the monotone smooth cutoff of V18, using the principal
SU(2) logarithm, let \(0<\rho\le1/4\).
For arbitrary fourteen compact group parameters \(P_j\)
and \(0\le t\le1\), the map
\[
 F_t(B)=B\exp(tX(B)),\qquad
 X(B)=\frac1{16}\sum_{j=1}^{14}g_\rho(B^{-1}P_j)
 \tag{93.2}
\]
is a smooth diffeomorphism of SU(2). Every singular value
of its differential, in invariant Hilbert--Schmidt metrics,
is at least \(1/9\). Thus (93.1) is an admissible global
completion cutoff for the present SU(2) parent.
\end{theorem}
\begin{proof}
Every SU(2) input is conjugate to \(\exp(qT)\).
V86's direct quaternion derivative gives the axial
multiplier \(h'(q)\) and transverse multiplier
\(h(q)(\cot q-i)\) for \(D_Lg_\rho\).
The axial multiplier is at most 1 by monotonicity of
\(\chi\). On the support,
\[
 h(q)\cot q=\chi(\sqrt2|q|/\rho)\,q\cot q\le1.
 \tag{93.3}
\]
Here \(|q|\le\rho/\sqrt2<\pi/2\), so \(q\cot q\) is
nonnegative and at most 1; the last inequality follows
also by differentiating \(\sin q-q\cos q\) for \(q\ge0\).
Off the support the derivative is zero. Conjugation
is orthogonal in the invariant metric. Consequently,
for arbitrary, not necessarily commuting inputs,
\[
 \operatorname{Sym}D_Lg_\rho(U)\le I,\qquad
 A_B:=\frac1{16}\sum_jD_Lg_\rho(B^{-1}P_j),
          \quad\operatorname{Sym}A_B\le\frac78I.
 \tag{93.4}
\]

Up to an orthogonal output adjoint factor, the differential
of \(F_t(B e^{sv})\) is
\[
                    M_t=I-t\mathcal J(tX)A_B,
       \qquad \mathcal J(Z)=\int_0^1
                                    \operatorname{Ad}_{e^{sZ}}\,ds .
 \tag{93.5}
\]
Write \(r=\|tX\|_{\rm HS}/\sqrt2\). Since
\(\|X\|_{\rm HS}\le7\rho/8\),
\[
                        r^2\le49/2048.
 \tag{93.6}
\]
On the axis of \(X\), \(\mathcal J(tX)\) is identity.
On its transverse plane it is complex multiplication by
\(e^{ir}\sin r/r\), with the limiting interpretation at zero.
Thus its smallest singular value is at least
\(1-r^2/6\), while
\[
 \operatorname{Sym}\mathcal J(tX)^{-1}
       \ge r\cot r\,I\ge\cos r\,I\ge(1-r^2/2)I.
 \tag{93.7}
\]
All inequalities hold on the small range (93.6), where
\(\cos r>0\) and \(r/\sin r\ge1\).
Multiplying (93.5) on the left by the invertible
\(\mathcal J(tX)^{-1}\) gives
\[
 \operatorname{Sym}
       [\mathcal J(tX)^{-1}-tA_B]
                \ge\left(\frac18-\frac{49}{4096}\right)I.
 \tag{93.8}
\]
For a real operator \(L\), the lower bound
\(\operatorname{Sym}L\ge cI\) implies
\(\|Lv\|\ge c\|v\|\) by Cauchy--Schwarz.
Combining this fact with (93.7) gives
\[
 s_{\min}(M_t)\ge
 \left(1-\frac{49}{12288}\right)
 \left(\frac18-\frac{49}{4096}\right)>\frac19 .
 \tag{93.9}
\]
The last comparison is rational integer arithmetic.
It avoids multiplying an exponential-derivative error
by an unspecified absolute cutoff derivative bound.

Each \(F_t\) is therefore a local diffeomorphism of a
compact connected group, hence a covering map onto it.
The homotopy \(F_t\) begins at identity and has degree one.
The nonsingular differential has positive orientation
throughout the homotopy, so this covering has one sheet.
Thus \(F_t\) is a smooth global diffeomorphism.
The bound is uniform in all fourteen compact parameters.
\end{proof}

This theorem is specific to SU(2). It does not replace
the separate higher-rank group requirements of the full
Yang--Mills programme.

\subsection{The exact scalar system and its proposed centre}

Use the V88 adjusted field \(Y^\circ\) as reference and
write
\[
                 Y_i(\xi)=Y_i^\circ e^{\xi_iT},
                           \qquad \xi\in\mathbb R^{944}.
 \tag{93.10}
\]
For the four active paths \(e\), let \(r_e\xi\) be the
sum of the angles at their two free endpoints.
The proposed pivots are
\(\mathcal B_e(\xi)=\exp[-(r_e\xi)T/15]\);
the other 92 proposed pivots are \(I\).
The validity of these proposals on the entire ball will
be checked below, rather than inferred from V88's small
neighbourhood of the different point \(\xi=0\).

The exact incidence routine
\begin{center}\small\ttfamily
scripts/ym\_f12\_v89\_scalar\_system.py
\end{center}
retains 2376 plaquettes whose angles depend on \(\xi\).
It writes their effective incidence rows as
\(M_{pi}=m_{pi}/15\), with integer \(m_{pi}\), and
their base scalar and axial quaternion parts as rational
numbers \(c_p,s_p\), with \(c_p^2+s_p^2=1\).
Thus the angle-shifted cosine and sine are
\[
 \begin{split}
 C_p(\xi)&=c_p\cos(M_p\xi)-s_p\sin(M_p\xi),\\
 S_p(\xi)&=s_p\cos(M_p\xi)+c_p\sin(M_p\xi).
 \end{split}
 \tag{93.11}
\]
On any region where the pivot proposals are valid,
\[
       g(\xi)=2M^{\mathsf T}S(\xi),\qquad
       A(\xi)=2M^{\mathsf T}
                         \operatorname{diag}(C(\xi))M
 \tag{93.12}
\]
are the actual scalar gradient and Hessian.

Exact incidence enumeration gives
\[
 \max_p\sum_i|m_{pi}|\le60,\quad
 \max_p\sum_i m_{pi}^2\le900,\quad
 \max_i\sum_p|m_{pi}|\le96.
 \tag{93.13}
\]
Consequently
\(\|M\|_\infty\le4\), \(\|M\|_1\le32/5\),
\(\max_p\|M_p\|_2\le2\), and
\(\|M\|_2^2\le128/5\). It follows directly from
\(|\cos u-\cos v|\le|u-v|\) that
\[
             \|A(\xi)-A(\eta)\|_2
                    \le L\|\xi-\eta\|_2,\qquad
                           L=\frac{512}{5}.
 \tag{93.14}
\]
Indeed (93.12) bounds the left side by
\(2\|M\|_2^2\max_p|M_p(\xi-\eta)|\).

The exploratory Newton script proposed a centre
\[
             \widehat\xi_i=n_i/10^{16},
 \tag{93.15}
\]
whose 944 fixed integer numerators are stored in
\begin{center}\small\ttfamily
scripts/ym\_f12\_v89\_center\_data.py.
\end{center}
This frozen certificate data file has SHA--256
\begin{center}\scriptsize\ttfamily
9D2E9E4FA8E0C236FB8BCF88099865DD2117628BD3EA1800F322C698F60CE9AC.
\end{center}
Newton's numerical maximum gradient decreased from
approximately \(0.05753\) to \(6.7\cdot10^{-16}\) in
three steps. Those numbers are only an exploration record.
The exact datum (93.15), not the numerical stopping rule,
is submitted to the verifier.

\subsection{Rational cutoff checks on the entire proposed ball}

Fix
\[
  R=10^{-8},\qquad
  \mathbb B=\{\xi:\|\xi-\widehat\xi\|_2\le R\}.
 \tag{93.16}
\]
For a path with base product \(aI+sT\), its proposed
relative product at a pivot has the form
\(\exp((p_0+\ell\xi)T)\), where \(\ell\) is a rational
row of signed angle incidences and
\(\cos p_0=a,\sin p_0=s\).
Every coordinate perturbation on \(\mathbb B\) is at
most \(R\), so \(|\ell(\xi-\widehat\xi)|
\le\|\ell\|_1R\).

All trigonometric evaluations in the verifier use rational
Taylor enclosures. For \(|u|\le1/10\), define
\[
 C_8(u)=\sum_{k=0}^8\frac{(-1)^ku^{2k}}{(2k)!},\qquad
 S_8(u)=\sum_{k=0}^8\frac{(-1)^ku^{2k+1}}{(2k+1)!}.
 \tag{93.17}
\]
Taylor's remainder gives
\[
 |\cos u-C_8(u)|\le\frac{|u|^{18}}{18!},\qquad
 |\sin u-S_8(u)|\le\frac{|u|^{19}}{19!}.
 \tag{93.18}
\]
After multiplying these intervals by the rational base
parts \(a,s\), the endpoints are rounded outward by
integer division to denominator \(2^{80}\).
All arguments used by the certificate satisfy the
stated bound \(1/10\).

For the four active paths, exact interval arithmetic
on (93.15)--(93.16) gives
\[
   \max_{\xi\in\mathbb B,e}|r_e\xi|
            \le\frac{5065759777}{10^{12}}
                         <\frac{21\rho}{64}.
 \tag{93.19}
\]
Since \(\sqrt2<10/7\),
\[
  \|\log(\mathcal B_e^{-1}P_e)\|_{\rm HS}
      =\frac{16\sqrt2}{15}|r_e\xi|
      <\frac{32}{21}\frac{21\rho}{64}=\rho/2 .
 \tag{93.20}
\]
These are ordinary, fully active logarithms.
For every other relative product, the cosine enclosure
at \(\widehat\xi\), enlarged by \(\|\ell\|_1R\), gives
\[
 \|\mathcal B_e(\xi)^{-1}P_{ej}(\xi)-I\|_{\rm HS}^2
                   -\rho^2
       \ge\frac{80791219}{500000000000}>0
                       \quad(\xi\in\mathbb B).
 \tag{93.21}
\]
For an inactive edge the proposed pivot is \(I\).
For an active edge the row \(\ell\) includes the
\(r_e/15\) contribution of the inverse pivot; it is
not just the original path row.

\begin{proposition}[The verified scalar system is the actual parent system]
\label{prop:f12v89-branch}
Every proposed pivot on \(\mathbb B\) satisfies the
actual completion equation, hence is its unique global
inverse. Therefore (93.12) is the actual scalar gradient
and Hessian throughout \(\mathbb B\).
\end{proposition}
\begin{proof}
Equations (93.19)--(93.21) put the single active relative
logarithm on each active edge in the inner cutoff region
and put all other relative products outside the support.
The identity (92.9) therefore verifies the completion
equation on that edge. On inactive edges all truncated
logarithms vanish at proposed pivot \(I\).
Theorem~\ref{thm:f12v89-global-pivot} supplies global
uniqueness at the selected value \(\rho=1/64\).
Differentiating the exact affine angle formulas gives
(93.12).
\end{proof}
\subsection{A positive Hessian certificate using only integer arithmetic}

The main verifier is
\begin{center}\small\ttfamily
scripts/verify\_ym\_f12\_v89\_stationary\_center.py.
\end{center}
It encloses every entry of \(A(\widehat\xi)\) using
(93.17)--(93.18). After rational accumulation and outward
rounding, it obtains symmetric integer matrices \(A^-,A^+\)
such that, entry by entry,
\[
      \frac{A^-_{ij}}{s^2}
          \le A(\widehat\xi)_{ij}
          \le\frac{A^+_{ij}}{s^2},
                         \qquad s=2^{26}.
 \tag{93.22}
\]
The entrywise lower enclosure is not itself asserted
positive semidefinite.

Two rational triangular matrices \(R=R_{\mathbb Z}/s\)
and \(K=K_{\mathbb Z}/s\) are constructed by deterministic
integer recurrences. An integer Cholesky-type recurrence
applied to \(A^-\) proposes a lower triangular matrix:
diagonal values are integer square roots of the remaining
diagonal numerator, and off-diagonal values are rounded
integer quotients. Transposing gives \(R_{\mathbb Z}\).
A rounded triangular inverse recurrence proposes
\(K_{\mathbb Z}\). Their accuracy is not assumed.
Only the residual inequalities in the following proposition
are used to accept the resulting witnesses.

\begin{proposition}[Residual criterion for a positive Hessian]
\label{prop:f12v89-integer}
Suppose \(A\) is real symmetric and has the enclosure
(93.22). For rational matrices \(R,K\), set
\[
 \eta_\infty=\|I-KR\|_\infty,\quad
 \eta_1=\|I-KR\|_1,\quad
 b_\infty=\frac{\|K\|_\infty}{1-\eta_\infty},\quad
 b_1=\frac{\|K\|_1}{1-\eta_1}.
 \tag{93.23}
\]
Assume \(\eta_\infty,\eta_1<1\). If
\(\|A-R^{\mathsf T}R\|_\infty\le\epsilon\), then
\[
                   A\ge
             \left(\frac1{b_\infty b_1}-\epsilon\right)I .
 \tag{93.24}
\]
All quantities needed for (93.23)--(93.24) can be bounded
by exact integer products and integer row or column sums.
\end{proposition}
\begin{proof}
The Neumann series for \(KR\) is convergent in each of
the two induced norms. Since \(KR\) is a square invertible
matrix, both \(R\) and \(K\) are invertible, and
\[
          R^{-1}=(KR)^{-1}K,\qquad
          \|R^{-1}\|_\infty\le b_\infty,\quad
          \|R^{-1}\|_1\le b_1 .
 \tag{93.25}
\]
The elementary row-column inequality
\(\|T\|_2^2\le\|T\|_1\|T\|_\infty\) follows by
applying weighted Cauchy--Schwarz in each row and then
summing the column weights. Hence
\(R^{\mathsf T}R\ge(b_\infty b_1)^{-1}I\).
The error matrix is symmetric, so its operator norm is
at most its maximum absolute row sum, proving (93.24).

For explicit computation let
\(D=R_{\mathbb Z}^{\mathsf T}R_{\mathbb Z}\).
The enclosure (93.22) implies
\[
 \epsilon\le\frac1{s^2}\max_i
      \sum_j\max\{|A^-_{ij}-D_{ij}|,\ |A^+_{ij}-D_{ij}|\}.
 \tag{93.26}
\]
The inverse residuals are computed from the integer matrix
\(s^2I-K_{\mathbb Z}R_{\mathbb Z}\), divided by \(s^2\).
This supplies purely integer versions of every required
bound. No numerical eigenvalue is needed.
\end{proof}

For the fixed centre (93.15), the deterministic witnesses
pass the following rational bounds:
\[
 \begin{array}{c|c}
 \text{quantity}&\text{certified bound}\\ \hline
 \eta_\infty&\le4094323/250000000000\\
 \eta_1&\le3704247/250000000000\\
 b_\infty&\le4529213/1000000\\
 b_1&\le10186203/1000000\\
 (b_\infty b_1)^{-1}&\ge2167529/100000000\\
 \epsilon&\le10279847/10^{12}.
 \end{array}
 \tag{93.27}
\]
The displayed entries are outward-rounded reports of
the exact rational quantities used in the test.
Their unrounded values give
\[
       A(\widehat\xi)\ge
                  \frac{2166501}{100000000}I
                            \ge\frac1{50}I .
 \tag{93.28}
\]

The checker uses NumPy only for integer arrays and integer
products. Before each accelerated product it bounds the
absolute sum of every dot product, so no intermediate
partial sum can overflow a signed 64-bit integer.
For the factor recurrence the bound follows from tracked
integer squared row norms and Cauchy--Schwarz. For the
inverse recurrence and final residual products it uses
an integer row-sum bound times a maximum entry bound.
All these verified dot-product bounds are at most
\[
                       2^{58}
                          =288230376151711744<2^{63}-1.
 \tag{93.29}
\]
The additions and subtractions have separate guards.
Other arithmetic, including Taylor intervals, integer
square roots, and rational norm comparisons, uses
arbitrary-precision integers and fractions.
Thus floating-point Cholesky, floating-point inversion,
and floating-point eigenvalues are absent from the
accepted certificate.

\subsection{Residual, convexity, and an interior critical point}

The same trigonometric enclosures give
\[
       \|g(\widehat\xi)\|_2
          \le\frac{6327}{500000000000000000}
                       <\frac{13}{10^{15}}.
 \tag{93.30}
\]
To obtain this bound, the verifier encloses each gradient
component with denominator \(15\cdot2^{80}\).
It takes the larger absolute endpoint of each interval,
sums their integer squares, and rounds the integer square
root upward. No cancellation of interval errors is assumed.

Equations (93.14), (93.16), and (93.28) imply
\[
               A(\xi)\ge\frac1{100}I
                            \quad(\xi\in\mathbb B).
 \tag{93.31}
\]
For example,
\(1/50-(512/5)10^{-8}>1/100\).
The retained exact bound in the verifier is stronger,
\(10831993/500000000\), but the simpler bound is enough.

\begin{theorem}[Validated and unique scalar stationary centre]
\label{thm:f12v89-centre}
For the actual completed parent at \(\rho=1/64\), there
is exactly one scalar critical point \(\xi_*\) in
\(\mathbb B\). It is the strict scalar minimum on
\(\mathbb B\), and
\[
       \|\xi_*-\widehat\xi\|_2
              <\frac{13}{10^{13}}.
 \tag{93.32}
\]
Its full physical action satisfies
\[
           \mathcal S(\xi_*)
                \le\frac{569684412973}{10^9}<569.685.
 \tag{93.33}
\]
\end{theorem}
\begin{proof}
The actual action is smooth on a neighbourhood of the
closed ball by the strict cutoff inequalities above.
It attains a minimum on the ball. For a boundary point
\(\xi=\widehat\xi+h\), with \(\|h\|_2=R\), the fundamental
theorem of calculus and (93.31) give
\[
 \begin{split}
  h\cdot g(\widehat\xi+h)
    &=h\cdot g(\widehat\xi)
       +\int_0^1h^{\mathsf T}A(\widehat\xi+th)h\,dt\\
    &\ge-R\,\frac{13}{10^{15}}+\frac{R^2}{100}>0.
 \end{split}
 \tag{93.34}
\]
Moving a small distance radially inward therefore
decreases the action. No boundary point is a minimum,
so an interior minimizing point exists and has zero
gradient. The positive Hessian throughout the convex
ball gives strict convexity and uniqueness of the
critical point.

Put \(h_*=\xi_*-\widehat\xi\). Since \(g(\xi_*)=0\),
the same integrated Hessian bound gives
\[
       \frac{\|h_*\|_2^2}{100}
             \le-h_*\cdot g(\widehat\xi)
             \le\|h_*\|_2\|g(\widehat\xi)\|_2.
 \tag{93.35}
\]
If \(h_*=0\) the claim is immediate; otherwise divide
by its norm and use (93.30) to obtain (93.32).

The minimum has no greater action than the centre.
For the centre's energy the verifier starts with V88's
upper bound for \(\mathcal S(Y^\circ)\), then encloses
each difference \(2[c_p-C_p(\widehat\xi)]\) with outward
rational arithmetic over the 2376 variable plaquettes.
The resulting upper bound is (93.33).
This computation retains the four proposed pivots,
which have already been identified with the actual ones.
\end{proof}

\begin{corollary}[A located noncentral full hidden critical orbit]
\label{cor:f12v89-orbit}
The completed configuration \(Y(\xi_*)\) is noncentral
and critical in all 2832 hidden SU(2) directions.
Its simultaneous conjugation orbit is an embedded
two-sphere of critical points with the same action.
The longitudinal angular Hessian is at least \(I/100\).
The full Hessian has the two forced orbit tangent
directions in its kernel; additional transverse kernel
or negative directions are not excluded here.
\end{corollary}
\begin{proof}
All independent groups at \(\xi_*\) lie in the common
torus. Theorem~\ref{thm:f12v86-critical} transfers scalar
criticality to all hidden directions, because the coarse
field is \(I\) and the frozen exterior is central.

At index 2 of the adjusted reference field, rational
arithmetic gives \(|s_2^\circ|>3/5\). The frozen centre
and the ball satisfy
\(|\widehat\xi_2|+R<1/100\).
Since the sine function is 1-Lipschitz, the corresponding
group at \(\xi_*\) has sine magnitude greater than
\(3/5-1/100=59/100\). It is noncentral.
The stabilizer of a noncentral commuting tuple is its
circle, so its conjugation orbit is SU(2)/U(1), an
embedded two-sphere as proved in V86.
Invariance transports criticality along that orbit.
Differentiating the full criticality equation along
the orbit gives the two Hessian kernel directions.
Equation (93.31) gives the claimed longitudinal bound.
\end{proof}

\subsection{The next transverse test is at nonidentity pivots}

This iteration closes the local stationarity requirement
at a specifically located coupled configuration. It does
not prove a full-parent local minimum or a global minimum.
The previously tested V88 rational point remains
nonstationary; its negative orbit-geodesic value must
not be attributed to the point just certified.

There is a further concrete distinction. The exact
whole-ball checks prove that for at least one active
path,
\[
                    |r_e\xi|>1/1000
                                  \quad(\xi\in\mathbb B).
 \tag{93.36}
\]
Thus at least one of the four active pivots is genuinely
nonidentity at \(\xi_*\). V88's particular second-jet
formula (92.12) was evaluated at \(P_e=I\).
It cannot be used unchanged at this new centre.
The next normal-stability calculation must differentiate
the actual inverse-root map (92.6) at its nonidentity
argument, or equivalently its local algebraic relation
\[
                         \mathcal B_e^{15}P_e=I.
 \tag{93.37}
\]
Both the velocity and acceleration of that map must be
retained. Now that full hidden criticality has been
certified, the conjugation kernel may legitimately be
separated before testing transverse normal positivity.

\subsection{Validation and continuation record}

The exact verifier passes, including all whole-ball
cutoff inequalities, the complete 944-dimensional
gradient enclosure, all Hessian entry enclosures,
the deterministic integer factor and inverse tests,
overflow guards, the outward-gradient inequalities,
the energy upper bound, and noncentrality.
The Newton script is separate and is not imported by
the verifier. No numerical eigenvalue is used to infer
existence, uniqueness, or longitudinal positivity.

V88 had 1457392 bytes and SHA--256
\begin{center}\scriptsize\ttfamily
DE38D17A173653A09D6587EA7409117323C8FD333A7EBC9ACA526FAD053C97A6.
\end{center}
Its body is retained apart from the title version.
Only active V88 is replaced by V89. Earlier scripts,
frozen data, archives, and companion notes are preserved.
The next version, V90, includes the scheduled companion
audit and direction review.

\begin{firewall}[Position after V89]
At the specified finite SU(2) parent and cutoff, a
noncentral full hidden critical orbit has been located
and certified, with a unique scalar centre in the
specified ball and a strictly positive longitudinal
Hessian. The new explicit global SU(2) pivot threshold
justifies the chosen cutoff. Transverse normal positivity,
global ground-set coverage, unrestricted repair, the
cofinal auxiliary gap, continuum construction, and the
physical Yang--Mills mass gap remain unproved.
The full objective remains active.
\end{firewall}
\section{V90: nonidentity inverse-root curvature and a certified full hidden local minimum}
\label{sec:f12v90}

\subsection{Scope and the scheduled cross-programme audit}

This iteration finishes the local normal-stability question left
open by V89. The configuration is the same finite SU(2) parent,
with its specified central exterior, coarse field \(I\), cutoff
\(\rho=1/64\), and 944 independent hidden groups. No exterior
variable is released in the assertion proved below.
The new result is a full hidden Morse--Bott local minimum modulo
simultaneous conjugation. Global minimality is not asserted.
In particular, the normal Hessian gap below is a finite-dimensional
classical curvature bound, not the physical Yang--Mills mass gap.

The scheduled V90 companion audit found the following active
versions unchanged from the V85 audit:
\begin{itemize}
\item SM--Einstein V72 (607372 bytes): its quantum-to-covariance
result for fixed modes and times in the many-copy limit does not provide a uniform
ultraviolet limit, kinetic-time limit, or finite-species physical
continuum construction.
\item NS stability V26 (278283 bytes): the exact pointwise rank-one
norm formulas do not supply a full stability theorem or global
regularity.
\item Parallel YM V5 (54174 bytes): the extensive total-charge
exponential-tail target fails under tensorization. Its retained
fixed-window, density, and quasilocal alternatives do not provide
the missing cofinal Poincare estimate.
\end{itemize}
Their respective SHA--256 values are
\begin{center}\scriptsize\ttfamily
F44F9745D43379E1672C5550411B58E97195A4C9278592D221DF5D3F644DC1F5\\
82C887F8C2C69E4F28FAD7C3DC8DE5B9099CB5A4BF431EBA1FFF3E91935978AD\\
306F1BB84CFA8A8D47EA569EC17E9545F9867C8D32B548EC9D9C444B4F3C0512.
\end{center}
These are mathematical source documents, not new instructions
to the present programme. No new global-gap input was found.
The direction review appears at the end of this section.

\subsection{The exact critical-point box and what has already been proved}

Let \(\widehat\xi\) be the frozen rational centre of V89, and
let \(\xi_*\) be the unique scalar critical point proved in
Theorem~\ref{thm:f12v89-centre}. Set
\[
 r=\frac{13}{10^{13}},\qquad
 \mathcal Q=\{\xi:|\xi_i-\widehat\xi_i|\le r,\ 0\le i<944\}.
 \tag{94.1}
\]
Then \(\xi_*\in\mathcal Q\). Moreover,
\(\mathcal Q\subset\mathbb B\), the V89 Euclidean ball of
radius \(10^{-8}\), because
\(\sqrt{944}\,r<31r<10^{-8}\).
Consequently every cutoff gate, pivot identification, and smoothness
assertion already certified on \(\mathbb B\) holds on \(\mathcal Q\).
There are exactly four nontrivial local pivot maps, each
\[
             \mathcal B=P^\nu,\qquad \nu=-1/15,
 \tag{94.2}
\]
and the other 92 pivots are identically \(I\) in a neighbourhood
of the relevant configurations. Here the power is the local
analytic power near \(P=I\).
At least one of the four arguments is nonidentity at \(\xi_*\).

Corollary~\ref{cor:f12v89-orbit} already proves full hidden
criticality, a noncentral conjugation orbit \(\mathcal O\simeq S^2\),
and longitudinal angular Hessian at least \(I/100\).
We do not repeat the existence proof or replace \(\xi_*\) by
the nonstationary rational field used in V88.
The remaining local task is a transverse calculation at the
actual nonidentity roots.

\subsection{First and second jets of a nonidentity inverse root}

Use \(E_1E_2=T\), with \(E_1,E_2,T\) the unit quaternion
generators, each of Hilbert--Schmidt norm \(\sqrt2\).
Represent \(xE_1+yE_2\) by \(z=x+iy\).
For one active chronological path write
\[
 A=e^{\theta_aT},\quad B=e^{\theta_bT},\quad
 P=AB=e^{pT},\quad b=\nu p,\quad \mathcal B=e^{bT}.
 \tag{94.3}
\]
The small signed angle \(p\) is the local path angle; it is
not an arbitrary choice modulo \(2\pi\).
V88's snapped reference has \(P=I\), so on the present branch
\(p=\xi_a+\xi_b\).
Vary the independent endpoints by right geodesics
\(A(t)=A\exp(t v_a)\), \(B(t)=B\exp(t v_b)\), with transverse
complex representatives \(z_a,z_b\). Put
\[
 \alpha=e^{-2i\theta_b},\qquad
 u=\alpha z_a+z_b,\qquad
 h=\operatorname{Im}(\overline{\alpha z_a}\,z_b).
 \tag{94.4}
\]
In particular, away from \(P=I\) one must not replace
\(\alpha\) by \(e^{2i\theta_a}\).

\begin{proposition}[Nonidentity power response, including acceleration]
\label{lem:f12v90-root}
The right velocity of \(\mathcal B(t)=P(t)^\nu\) is transverse
with complex representative \(\ell(p)u\), where
\[
 \ell(p)=e^{i(1-\nu)p}\frac{\sin(\nu p)}{\sin p}
       =\nu e^{i(1-\nu)p}\frac{\operatorname{sinc}(\nu p)}
                                   {\operatorname{sinc}p}.
 \tag{94.5}
\]
Its second jet is
\[
 \mathcal B''(0)=\mathcal B
     \left(-|\ell(p)u|^2I+\zeta T\right),\qquad
 \zeta=2\nu h+\tau(p)|u|^2,
 \tag{94.6}
\]
where
\[
 \tau(p)=\nu\cot p-
       \frac{\sin(\nu p)\cos(\nu p)}{\sin^2p}
 =\nu\,\frac{\operatorname{sinc}(2p)-
                    \operatorname{sinc}(2\nu p)}
                  {p\,\operatorname{sinc}^2p}.
 \tag{94.7}
\]
All expressions extend analytically across \(p=0\), with
\(\ell(0)=\nu\) and \(\tau(0)=0\).
\end{proposition}
\begin{proof}
The product rule, after right trivialization, gives
\[
 P'(0)=P u,\qquad
 P''(0)=P(-|u|^2I+2hT).
 \tag{94.8}
\]
For example, the mixed product of the two transverse generators
has axial part \(hT\); it occurs twice in the second derivative.
The differential of the SU(2) power map at \(e^{pT}\)
multiplies the transverse quaternion component by
\(\sin(\nu p)/\sin p\). The input transverse quaternion component
is represented by \(e^{ip}u\), and right trivializing the output
multiplies it by \(e^{-i\nu p}\). This proves (94.5).

For \(p\ne0\) sufficiently small, use the signed polar angle
\(q(t)\) of \(P(t)\), with \(q(0)=p\).
Its first derivative vanishes because \(P'(0)\) is transverse.
Taking scalar parts in (94.8) yields
\[
                  q''(0)=\cot p\,|u|^2+2h .
 \tag{94.9}
\]
The scalar part of \((P(t)^\nu)''\) is therefore
\(-\nu\sin(\nu p)q''(0)\).
Equivariance under the circle fixing \(T\) shows that its
quadratic right-logarithmic acceleration is axial: rotation by
\(\pi\) reverses both transverse inputs, leaves a quadratic
expression unchanged, and reverses any transverse output.
The unit-quaternion constraint fixes the scalar right-trivialized
second derivative to \(-|\ell u|^2\).
Comparing scalar parts of (94.6) now gives (94.7) and (94.6).
This argument applies for both signs of the small nonzero \(p\).
The analytic power map and the sinc representations extend the
identity to zero. At zero it reduces to V88's identity-root
second jet, rather than assuming that jet remains unchanged.
\end{proof}

\subsection{The completed Hermitian Hessian}

Let \(H_{\rm phys}\) be the transverse Hermitian Hessian of the
physical plaquette action, first treating the 944 independent
links and four moving pivots as 948 independent physical links.
It is evaluated at their actual commuting values, including
the four nonidentity values \(e^{\nu pT}\).
V88's oriented-word differentiation rule applies at every
commuting tuple, not only at identity pivots.
Let \(J\) have identity rows for the 944 free links and, for
the pivot associated with \((a,b)\), entries
\[
                       J_{e,a}=\ell(p)\alpha,\qquad
                       J_{e,b}=\ell(p).
 \tag{94.10}
\]
Let \(f_e\) be the physical axial force at that pivot:
\[
             f_e=2\sum_{\text{plaquettes }q\ni e}
                            \epsilon_{q,e}\sin\omega_q .
 \tag{94.11}
\]
It need not vanish at a constrained critical point.

\begin{proposition}[Completed transverse Hessian]
\label{lem:f12v90-completed}
For every commuting configuration in \(\mathcal Q\), the
transverse quadratic form is \(z^*H_\perp z\), where
\[
 \begin{gathered}
 H_\perp=J^*H_{\rm phys}J+\sum_e C_e,\qquad
 (C_e)_{aa}=(C_e)_{bb}=f_e\tau(p),\\
 (C_e)_{ab}=f_e(\tau(p)-i\nu)\overline\alpha,\qquad
 (C_e)_{ba}=\overline{(C_e)_{ab}}.
 \end{gathered}
 \tag{94.12}
\]
All other entries of \(C_e\) vanish.
\end{proposition}
\begin{proof}
The second derivative of a composition is its physical Hessian
on the induced first velocities plus its physical differential
on the right-logarithmic accelerations.
The independent free links follow geodesics, so only pivot
accelerations contribute. Proposition~\ref{lem:f12v90-root} makes their
contribution
\[
       f_e\{2\nu\operatorname{Im}(\overline{\alpha z_a}z_b)
                      +\tau(p)|\alpha z_a+z_b|^2\}.
 \tag{94.13}
\]
Expanding this real quadratic expression gives exactly \(C_e\).
In particular the physical-force term is retained even though
the completed configuration will later be specialized to a
critical point. Omitting it would differentiate a different map.
\end{proof}

For clarity, the physical word rule used in the implementation
can be reproduced without a numerical derivative.
Represent a diagonal link by \(d=e^{i\theta}\).
Its forward oriented factor has derivative \(dE_1\); its reverse
factor has derivative \(-dE_1\), since
\(-E_1\overline d=-dE_1\).
For two differentiated positions \(k<l\), let
\(U,V,W\) be the products of the ordinary oriented diagonal
factors before, between, and after them, and let
\(\gamma_k,\gamma_l\) be the coefficients just specified.
The two-\(E_1\) word coefficient is
\[
                    C=-U\gamma_k\overline V\,
                                      \overline{\gamma_l}W.
 \tag{94.14}
\]
The Hermitian off-diagonal contribution for the action
\(2(1-\operatorname{Re}(\text{word}))\) is
\(-2\overline C\); the diagonal contribution at one position
is twice the real part of the word.
Replacing an input \(E_1\) by \(xE_1+yE_2\), or using
\(E_1d=\overline dE_1\), verifies the complex polarization.
Thus (94.12) can be assembled using only complex additions,
multiplications, conjugations, and the two functions in
(94.5)--(94.7).

\subsection{Outward intervals and a pinned integer positivity certificate}

The certificate uses the following new files:
\begin{center}\small\ttfamily
ym\_f12\_v90\_intervals.py\\
verify\_ym\_f12\_v90\_transverse\_minimum.py\\
verify\_ym\_f12\_v90\_root\_response.py.
\end{center}
They are in the repository's scripts directory. They reuse V88's
exact geometry, V89's frozen rational centre, and V89's checked
integer factor/inverse routines. No floating-point eigenvalue,
factorization, or inverse enters the accepted test.
The centre's SHA--256 remains
\begin{center}\scriptsize\ttfamily
9D2E9E4FA8E0C236FB8BCF88099865DD2117628BD3EA1800F322C698F60CE9AC.
\end{center}

Here are the enclosure rules, so the computational assertion
has a specified finite-arithmetic meaning. Real intervals have
integer endpoints divided by \(S=2^{70}\).
Addition is exact at that scale. Multiplication takes the
minimum and maximum of the four integer endpoint products,
then rounds the lower quotient down and the upper quotient up
after division by \(S\). For a positive interval \([l,u]/S\),
reciprocal endpoints are
\(\lfloor S^2/u\rfloor\) and \(\lceil S^2/l\rceil\).
Complex intervals use these rules componentwise.
Induction on the expression tree proves that every computed
interval contains every value of its expression on \(\mathcal Q\);
independence of the interval inputs is not required.

For each free group, the rational sine/cosine enclosure at
\(\widehat\xi_i\) from V89 is widened by \(r\) on both sides.
For each path angle the interval is
\(\widehat\xi_a+\widehat\xi_b+[-2r,2r]\); its pivot angle
has uncertainty \(2r/15\).
The functions \(\operatorname{sinc}x\) and \(\cos x\) use terms
through \(x^{16}\), with respective remainders
\[
          \frac{|x|^{18}}{19!},\qquad
          \frac{|x|^{18}}{18!},\qquad |x|\le1/10 .
 \tag{94.15}
\]
Sine is computed as \(x\operatorname{sinc}x\), with a further
conservative error of \(|x|^{19}/19!\).
The code checks all argument bounds and positivity of denominators.
For (94.7) no interval division by a possibly zero \(p\) is made.
Instead the code encloses the odd polynomial
\[
 \frac{\operatorname{sinc}(2p)-\operatorname{sinc}(2\nu p)}p
 =\sum_{k=1}^8
       \frac{(-1)^k2^{2k}(1-\nu^{2k})p^{2k-1}}{(2k+1)!}
       +\mathcal R(p),\quad
 |\mathcal R(p)|\le
       \frac{2^{18}(1+|\nu|^{18})|p|^{17}}{19!}.
 \tag{94.16}
\]
This follows by subtracting the two sinc Taylor formulas;
it also holds at zero by continuity.

Assembly of (94.12) gives 11080 potentially nonzero complex
entries of a \(944\)-dimensional Hermitian matrix.
Conjugate entries are replaced by their enclosing Hermitian
hulls. The diagonal imaginary part is exactly zero by Hermitian
symmetry, with interval containment of zero checked separately.
Pin the zero-based coordinate
\[
                             j=156 .
 \tag{94.17}
\]
Its sine magnitude is greater than \(1/2\) throughout the box.
Delete its complex row and column. For the resulting matrix
\(K=A+iB\), use the real symmetric representation
\[
                      K_{\mathbb R}=
                         \begin{pmatrix}A&-B\\B&A\end{pmatrix}
                             \quad\text{of size }1886.
 \tag{94.18}
\]
Each entry is rounded outward to denominator \(2^{52}\).

\begin{proposition}[Uniform pinned positivity, integer certificate]
\label{prop:f12v90-pinned}
For every \(\xi\in\mathcal Q\), the principal matrix obtained
by deleting complex coordinate 156 satisfies
\[
                 K(\xi)\ge
                   \frac{49641}{100000000}I
                         >\frac1{2500}I .
 \tag{94.19}
\]
\end{proposition}
\begin{proof}
Apply the deterministic integer construction of V89 to the
entry enclosure of (94.18), with dimension changed to 1886.
It constructs rational matrices \(R,L\) at denominator \(2^{26}\),
where \(R\) is triangular and \(L\) is an approximate inverse.
Write
\[
 \eta_\infty=\|I-LR\|_\infty,\quad
 \eta_1=\|I-LR\|_1,\quad
 b_\infty=\frac{\|L\|_\infty}{1-\eta_\infty},\quad
 b_1=\frac{\|L\|_1}{1-\eta_1}.
 \tag{94.20}
\]
The test verifies \(\eta_\infty,\eta_1<1\).
The Neumann series makes \(LR\), hence \(R\), invertible and
gives \(\|R^{-1}\|_\infty\le b_\infty\),
\(\|R^{-1}\|_1\le b_1\).
Therefore \(R^{\mathsf T}R\ge(b_\infty b_1)^{-1}I\).

If the entry endpoints are \(a^-_{kl},a^+_{kl}\), set
\[
 \epsilon=\max_k\sum_l
       \max\{|a^-_{kl}-(R^{\mathsf T}R)_{kl}|,
              |a^+_{kl}-(R^{\mathsf T}R)_{kl}|\}.
 \tag{94.21}
\]
Every enclosed real symmetric matrix differs from
\(R^{\mathsf T}R\) in operator norm by at most \(\epsilon\).
Indeed the displayed row bound controls the infinity norm,
and symmetry makes the one norm equal to it.
Consequently its smallest eigenvalue is at least
\((b_\infty b_1)^{-1}-\epsilon\).

The checked outward reports for this particular enclosure are
\[
\begin{array}{c|c}
 \eta_\infty\text{ upper}&33021737/10^{12}\\
 \eta_1\text{ upper}&403869/12500000000\\
 b_\infty\text{ upper}&7767817/500000\\
 b_1\text{ upper}&1549499/12500\\
 (b_\infty b_1)^{-1}\text{ lower}&25963/50000000\\
 \epsilon\text{ upper}&2285077/100000000000 .
\end{array}
 \tag{94.22}
\]
The unrounded rational quantities give, after rounding the
difference downward to denominator \(10^8\), the bound
\(49641/100000000\). This proves (94.19).
There is no sampling over \(\mathcal Q\): one interval enclosure
contains every matrix throughout that box.

For completeness the largest checked absolute dot-product
sum in the two final accelerated integer matrix products is
\[
                 387422008685988024<2^{59}<2^{63}-1 .
 \tag{94.23}
\]
The other final bound is \(242262362050304904\).
Factor-generation products have bound \(2^{58}\), and
inverse-generation products have bound \(56083336367269536\).
The integer routines check the additions and subtractions too.
Thus signed 64-bit accelerated products do not overflow;
all rounding bounds, norm comparisons, and Taylor coefficients
use arbitrary-precision integers or fractions.
The stronger \(2^{58}\) bound reported for all products in V89
is not asserted for this larger matrix.
\end{proof}

The independent root-response test is a useful implementation
check, not a substitute for Proposition~\ref{lem:f12v90-root}.
It chooses rational unit quaternions
\[
 z=\frac{1-t^2+2tT}{1+t^2},\quad
 t\in\{0,1/10000,-1/10000,3/20000\},\quad
 A=z^7,\ B=z^8,\ P=z^{15},\ \mathcal B=z^{-1}.
 \tag{94.24}
\]
For three rational transverse endpoint probes at each value,
it verifies exactly, through second order, the polynomial
quaternion-jet identity
\(\mathcal B(s)^{15}P(s)=I\), using (94.5)--(94.7).
It also checks that the interval coefficient formulas contain
the independently computed rational values at these four angles.
All twelve jet checks and four coefficient checks pass.
Finite probes alone would not establish the general identities;
their role is to catch sign, phase, and acceleration errors in
the implementation of the proved formulas.

\subsection{Removing the pin by the exact symmetry kernel}

\begin{proposition}[One-coordinate pinning with an exact kernel]
\label{lem:f12v90-pin}
Let \(H\) be Hermitian, \(Hw=0\), and \(w_j\ne0\).
If the principal matrix deleting \(j\) is at least \(\gamma I\),
then \(H\ge0\), \(\ker H=\mathbb Cw\), and
\[
                         H|_{w^\perp}\ge\gamma I .
 \tag{94.25}
\]
\end{proposition}
\begin{proof}
For arbitrary \(z\), put \(c=z_j/w_j\) and \(v=z-cw\).
Then \(v_j=0\) and \(z^*Hz=v^*Hv\ge\gamma\|v\|^2\).
This proves positivity and identifies its kernel.
If \(z\perp w\), then
\(\|v\|^2=\|z\|^2+|c|^2\|w\|^2\ge\|z\|^2\),
which proves the final assertion.
\end{proof}

\begin{theorem}[Full hidden Morse--Bott local minimum]
\label{thm:f12v90-minimum}
At the V89 critical point \(Y_*=Y(\xi_*)\), the completed
action has precisely the two conjugation tangent directions
in its Hessian kernel. Its other 2830 directions are positive.
On the normal space to \(\mathcal O\), using the product
Hilbert--Schmidt metric on the 944 independent hidden groups,
\[
             \operatorname{Hess}\mathcal S(Y_*)\ge
                                      \frac1{5000}I .
 \tag{94.26}
\]
The same bound holds along the conjugation orbit.
This orbit is a Morse--Bott local minimum of the full hidden
action, not merely of its commuting restriction.
\end{theorem}
\begin{proof}
At this full critical point, differentiating the invariant
criticality equation along the conjugation orbit gives the
exact transverse kernel vector
\[
                           w_i=e^{-2i\theta_i}-1,\qquad
                           H_\perp(\xi_*)w=0 .
 \tag{94.27}
\]
Its real and imaginary multiples are the two orbit directions.
The checked sine bound at index 156 ensures \(w_{156}\ne0\).
Proposition~\ref{prop:f12v90-pinned} and
Proposition~\ref{lem:f12v90-pin} give
\(H_\perp|_{w^\perp}\ge I/2500\) and no additional transverse
kernel.

The longitudinal block is at least \(I/100\) by V89.
The longitudinal--transverse block vanishes because the
stabilizing circle fixes longitudinal directions and rotates
transverse directions. These facts give normal positivity
at least \(1/2500\) in angular Euclidean coordinates.
Each quaternion generator has squared Hilbert--Schmidt norm
2, so converting to the stated product metric divides this
bound by 2, proving (94.26).
Conjugation is an isometry preserving the action and transports
the whole statement along the compact orbit.

In a small normal tubular neighbourhood of \(\mathcal O\),
the first derivative vanishes on the zero section, and the
normal Hessian has the uniform positive bound just established.
Smoothness and compactness give constants \(c,C>0\) such that
\[
        c\,d(Y,\mathcal O)^2
          \le\mathcal S(Y)-E_*
          \le C\,d(Y,\mathcal O)^2,\qquad
        E_*=\mathcal S(Y_*).
 \tag{94.28}
\]
This follows directly by Taylor's formula along normal
geodesics, after decreasing the tube radius if necessary.
The normal derivative is nonzero off the zero section in a
sufficiently small tube, by the same positive Hessian estimate.
Thus the critical set there is exactly the orbit, and it is a
Morse--Bott local minimum with the claimed normal dimension.
\end{proof}

The exact kernel in (94.27) uses full criticality.
The verifier also checks interval containment of every component
of \(H_\perp w\) over \(\mathcal Q\); this is only a consistency
diagnostic. It does not prove a zero identity at arbitrary
noncritical points in the box. In particular, pinned positivity
throughout \(\mathcal Q\) must not be misread as full positivity
at every point of \(\mathcal Q\); V88's noncritical Ward defect
explains why the criticality hypothesis matters.

\begin{corollary}[Local tube weight with the actual pivot amplitude]
\label{cor:f12v90-tube}
For a sufficiently small fixed tube \(\mathcal T\) around this
orbit, there are fixed-parent constants \(0<c_0\le C_0<\infty\)
such that, for all \(\beta\ge1\),
\[
 c_0\beta^{-1415}e^{-\beta E_*}
 \le\int_{\mathcal T}e^{-\beta\mathcal S(Y)}
                     a(Y)\,dY
 \le C_0\beta^{-1415}e^{-\beta E_*},
 \tag{94.29}
\]
where \(a\) is the actual product of inverse pivot Haar
Jacobians and \(dY\) is product Haar measure.
\end{corollary}
\begin{proof}
The amplitude and tubular volume density are smooth and
strictly positive; on a smaller tube their upper and lower
bounds are finite and positive. Apply (94.28) in the
2830-dimensional normal fibres and integrate over the compact
two-dimensional orbit. The upper Gaussian integral is bounded
by its integral over all normal space. For the lower bound,
integrate in a normal ball of radius a fixed constant times
\(\beta^{-1/2}\), contained in the tube for \(\beta\ge1\).
Both estimates have the power \(\beta^{-2830/2}\).
\end{proof}

This is a local integral. It is not an asymptotic formula for
the full parent normalizer: other minima may lie outside the
tube, may have lower action, and may be degenerate.
V89's \(E_*<569.685\) remains an upper bound for this orbit's
energy, not a lower bound for every hidden configuration.

\subsection{Direction review and continuation record}

The last ten versions separate three previously confusable
objects. The central sheet and the frozen 56-link slice do
not exhaust the full parent: V82 and V85--V87 supplied actual
escape or competitor arguments. V88's convenient rational
commuting tuple was not stationary. V89 located a genuine
full hidden critical orbit, and the present calculation proves
its full normal positivity, including nonidentity pivot
accelerations. The local problem for that orbit is now closed.

Repeating numerical refinements of its positive Hessian would
not address the remaining obstruction. The next upstream
question is global ground-set coverage and its joint visible
geometry. The V81 transfer criterion concerns all conditional
global minima on its visible domain, not one located hidden
local minimum at one central exterior.
Before using this branch for that criterion one must either
prove the required coverage or explicitly control all omitted
competitors. A joint visible force or Schur calculation on
this branch, by itself, would still not establish coverage.
Nor does a local Laplace power yield an original-clock heat
gap, unrestricted repair, or a cofinal estimate.

The accepted interval and integer tests and the independent
quaternion-jet tests pass. They provide a reproducible
computer-assisted finite-parent certificate using the specified
frozen data and formulas. They are not a formal proof-assistant
verification, and successful TeX compilation is not mathematical
certification.

V89 had 1479488 bytes and SHA--256
\begin{center}\scriptsize\ttfamily
98232EF29EB232E9FE18842EF49841E3283AB83D3A50717D3A449AB2B7D96D3E.
\end{center}
Its body is preserved apart from the title version.
Only active V89 is replaced by V90; frozen data, scripts,
archives, and companion files are retained.
The next scheduled companion audit is V95 and the next
direction review is V100. At V100 this lineage is to be
archived as family \(f12\), followed by a new compact V1
with a results-and-obligations summary.

\begin{firewall}[Position after V90]
The specified finite SU(2) parent has a certified noncentral
full hidden Morse--Bott local minimum with exactly two
symmetry zero modes and 2830 positive normal directions.
Global minimality, complete conditional ground-set coverage,
unrestricted repair, a cofinal auxiliary gap, continuum
construction, and the physical Yang--Mills mass gap remain
unproved. The full objective remains active.
\end{firewall}
\section{V91: a global central-sector certificate and uniform exclusion of all central tubes}
\label{sec:f12v91}

\subsection{Moving from a local branch to a whole competing sector}

V90 settles the local normal stability of one noncentral hidden
critical orbit. It does not establish global minimality.
The next question is therefore not another local eigenvalue
calculation, but which competing global minima can be ruled out.
This iteration treats \emph{every} central assignment of the
944 hidden SU(2) groups at the same fixed exterior and coarse
field used in V86--V90. There are \(2^{944}\) such assignments.
The result below is not restricted to the 56-link slice
classified in V83, nor to the one-direction pattern tested
in V76--V82.

Two exploratory checks motivated the calculation but are not
used as proofs. The all-identity hidden field has action 744,
and the original central field has action 648. Eight
deterministic commuting optimization starts did not produce
a value below the V89 branch; several converged to higher
values or reached their iteration limit. No completeness
or global-minimum conclusion follows from those searches.
The new proof instead gives a finite integer dual certificate
valid simultaneously on the entire central sector.

\subsection{The exact central reduction}

Write each independent central group as
\[
                         Y_i=(-I)^{x_i},\qquad
                         x_i\in\{0,1\},\quad 0\le i<944.
 \tag{95.1}
\]
All tree and exterior links are kept as in the earlier fixed
parent geometry. The 96 pivot values are not independently
chosen in (95.1).

\begin{proposition}[All central completions have identity pivots]
\label{prop:f12v91-central-pivots}
For every assignment (95.1), the actual 96 pivot values are \(I\).
\end{proposition}
\begin{proof}
Each of the fourteen nonroot path products for any pivot is
\(I\) or \(-I\). The truncated logarithm vanishes at both:
at \(I\) its logarithm is zero, and at \(-I\) its cutoff is zero.
Thus \(B=I\) solves the pivot equation
\[
 B\exp\left(\frac1{16}\sum_{j=1}^{14}
                         g(B^{-1}P_j)\right)=I.
 \tag{95.2}
\]
The global SU(2) diffeomorphism result of V89 applies at
\(\rho=1/64\), so this solution is the actual unique pivot.
This holds separately for all 96 pivots and every central
assignment, without choosing a local analytic branch.
\end{proof}

For each physical plaquette \(q\), let \(I_q\) be the set
of independent hidden indices in its oriented word and let
\(c_q\in\{0,1\}\) be the parity of its remaining fixed central
factors. By the preceding proposition the pivots contribute
zero parity. Orientation signs disappear modulo 2, and
the plaquette action is
\[
        4y_q,\qquad
        y_q=c_q+\sum_{i\in I_q}x_i\pmod2,\quad y_q\in\{0,1\}.
 \tag{95.3}
\]
The fine side is 6, so no plaquette has a repeated physical
link. The exact geometry enumeration checks that there are
5400 fixed plaquettes and 2376 plaquettes involving at least
one independent hidden group. The total fixed action is 408.
Combining equal pairs \((I_q,c_q)\) leaves 2292 distinct
variable parity terms. Write their multiplicities as \(m_k\),
so that \(\sum_km_k=2376\), and introduce one bit \(y_k\)
per term. Then, for every central assignment,
\[
              \mathcal S_{\rm cen}
                      =408+\sum_{k=1}^{2292}4m_ky_k.
 \tag{95.4}
\]

\subsection{Parity inequalities and a rational dual witness}

The following elementary inequality avoids enumerating
the \(2^{944}\) assignments.
For a local binary vector \(t=(x_i:i\in I_k,\ y_k)\),
the required parity is \(c_k\). For every binary vector \(b\)
of the wrong parity, impose
\[
            \sum_j(2b_j-1)t_j\le\sum_jb_j-1.
 \tag{95.5}
\]
Every valid binary \(t\) satisfies (95.5): its left side
equals \(\sum_jb_j-d_H(t,b)\), and the Hamming distance
is at least one because their parities differ.
The inequalities remain a valid relaxation when all variables
are allowed to lie in \([0,1]\). Only validity, not an
assertion of global integrality of the relaxation, is needed.

Collecting (95.5) over the ordered distinct terms gives
\[
       At\le b,\qquad 0\le t\le1,\qquad
       t\in\mathbb R^{3236},\quad
       A\in\mathbb Z^{15288\times3236}.
 \tag{95.6}
\]
The first 944 objective coefficients are zero; the remaining
ones are \(4m_k\). Call this objective vector \(c\).
Ordering is part of the certificate: independent links are
in V86's lexicographic order; pairs \((I_k,c_k)\) are sorted
lexicographically; for each pair, wrong-parity bit vectors
are taken in lexicographic order with \(0<1\).
The corresponding row has entries \(2b_j-1\) on the
indices of that term and its own auxiliary variable.

\begin{proposition}[Global central-sector action bound]
\label{prop:f12v91-central-bound}
Every one of the \(2^{944}\) central hidden configurations
at this fixed exterior and coarse field has actual action
\[
                            \mathcal S_{\rm cen}\ge616.
 \tag{95.7}
\]
This is a lower bound, not a claim that a central configuration
attaining 616 has been found.
\end{proposition}
\begin{proof}
The frozen file
\begin{center}\small\ttfamily
ym\_f12\_v91\_central\_dual\_data.py
\end{center}
specifies an integer vector \(u\) on the 15288 rows, with
637 nonzero entries, all negative. Its denominator is 23.
The independent integer verifier reconstructs (95.6) and checks
\[
      u\le0,\qquad
      r=23c-A^{\mathsf T}u\ge0,\qquad
      b^{\mathsf T}u=4784,\qquad
      0\le r_j\le184\quad\text{for every }j.
 \tag{95.8}
\]
All these are finite integer equalities or inequalities.
For a feasible \(t\), multiplication of \(At\le b\)
by \(u\le0\) reverses the inequality, whence
\[
 \begin{split}
  23c^{\mathsf T}t
      &=u^{\mathsf T}At+r^{\mathsf T}t\\
      &\ge u^{\mathsf T}b=4784.
 \end{split}
 \tag{95.9}
\]
Here \(r^{\mathsf T}t\ge0\) uses only \(t\ge0\);
the upper box bound is not needed for this dual witness.
Therefore every feasible relaxation point has
\[
                        408+c^{\mathsf T}t
                              \ge408+\frac{4784}{23}=616.
 \tag{95.10}
\]
Every central configuration gives such a feasible point
by (95.3)--(95.5), proving (95.7).

The frozen data file has SHA--256
\begin{center}\scriptsize\ttfamily
293ADBBEFA09A936B1C8ED50B26CAB23FDB340244763C1D11ED068D5D7CA51CA.
\end{center}
Its entries were proposed by a floating-point linear solver
and reconstructed at denominator 23. The proof uses neither
the solver's optimality flag nor its floating-point value.
The verifier imports the frozen integers and exact geometry,
not the solver, and checks (95.8) from scratch.
It also checks the 680 local valid-versus-forbidden parity
truth-table cases for local vector lengths 2 through 5.
The general Hamming-distance argument above proves validity
independently of those implementation checks.
\end{proof}

This is a computer-assisted finite certificate, not a
formal proof-assistant verification. Its arithmetic trust
boundary is simpler than the V90 matrix test: all accepted
arithmetic here uses Python arbitrary-precision integers
and fractions, with no accelerated machine-integer matrix
product or floating-point operation.
The enumeration and verifier are
\begin{center}\small\ttfamily
ym\_f12\_v91\_central\_parity.py\\
verify\_ym\_f12\_v91\_central\_exclusion.py.
\end{center}
The exploratory optimization and parity-relaxation scripts
are separate and are not imported by the verifier.

\subsection{A uniform tube around all central assignments}

Let \(d\) denote the bi-invariant SU(2) geodesic distance for
the Hilbert--Schmidt metric. Thus
\(d(I,-I)=\pi\sqrt2>4\).
Define the closed set
\[
 \mathcal N_\epsilon=
    \left\{Y:\max_{0\le i<944}d(Y_i,\{I,-I\})\le\epsilon\right\},
       \qquad \epsilon=\frac1{4096}.
 \tag{95.11}
\]
It is a neighbourhood, with boundary included, of every
central assignment at once. Its elements need not commute.
All comparisons in the next result use the actual pivots.

\begin{proposition}[Noncommutative pivot displacement near the centre]
\label{prop:f12v91-pivot-tube}
If \(Y\in\mathcal N_\epsilon\), choose a central assignment
\(Z\) with \(d(Y_i,Z_i)\le\epsilon\) for every independent
group. Every completed pivot obeys
\[
                            d(B_e(Y),I)\le14\epsilon.
 \tag{95.12}
\]
\end{proposition}
\begin{proof}
The product distance inequality gives
\[
                          d(P_j(Y),P_j(Z))\le2\epsilon
 \tag{95.13}
\]
for every nonroot path: each has two links and at most
two independent varying links. Its central reference
product is \(P_j(Z)=I\) or \(-I\).

The pivot equation (95.2) writes
\[
 B_e=\exp(-X_e),\qquad
 X_e=\frac1{16}\sum_{j=1}^{14}g(B_e^{-1}P_j(Y)).
 \tag{95.14}
\]
The cutoff bound \(\|g\|\le\rho\) gives
\(d(B_e,I)=\|X_e\|\le7\rho/8\); the logarithm is in
the injectivity ball at this value of \(\rho\).
If \(P_j(Z)=-I\), the triangle inequality gives
\[
 d(B_e,P_j(Y))
      \ge\pi\sqrt2-\frac{7\rho}{8}-2\epsilon>\rho.
 \tag{95.15}
\]
Its truncated logarithm is therefore zero. The strict
last inequality already follows from
\(4-7\rho/8-2\epsilon>\rho\), checked rationally at
\(\rho=1/64\).

There remain \(n\le14\) paths whose central product is \(I\).
Write \(b_e=d(B_e,I)\). Since \(0\le\chi\le1\), the
truncated logarithm satisfies
\[
       \|g(B_e^{-1}P_j(Y))\|
            \le d(B_e,P_j(Y))\le b_e+2\epsilon.
 \tag{95.16}
\]
This estimate also holds when its cutoff is zero.
Taking norms in (95.14) yields
\[
              b_e\le\frac n{16}(b_e+2\epsilon),
       \qquad
              b_e\le\frac{2n}{16-n}\epsilon\le14\epsilon.
 \tag{95.17}
\]
No commutativity, positive-path inner-cutoff assumption,
or derivative bound on the cutoff annulus was used.
\end{proof}

\begin{theorem}[Exclusion of the entire central tube from the ground set]
\label{thm:f12v91-tube}
For the actual completed action at the specified fixed parent,
\[
      \inf_{Y\in\mathcal N_\epsilon}\mathcal S(Y)
                         \ge\frac{156409}{256}
                         =610.97265625.
 \tag{95.18}
\]
In particular, with \(E_*\) the energy of the V89--V90 orbit,
\[
                  \inf_{\mathcal N_\epsilon}\mathcal S-E_*>41.
 \tag{95.19}
\]
No conditional global action minimizer belongs to
\(\mathcal N_\epsilon\).
\end{theorem}
\begin{proof}
For unitary factors, telescoping a plaquette word bounds
its Hilbert--Schmidt difference by the sum of the factor
differences. Taking an inverse preserves that norm.
Cauchy--Schwarz for the trace gives
\[
 |\mathcal S_q(Y)-\mathcal S_q(Z)|
     \le\sqrt2\sum_{\ell\in q}d(U_\ell(Y),U_\ell(Z)).
 \tag{95.20}
\]
There are exactly six plaquettes incident to each physical
link. The 944 independent links move by at most \(\epsilon\),
the 96 pivots by at most \(14\epsilon\), and all other links
are unchanged. Hence
\[
 \begin{split}
 |\mathcal S(Y)-\mathcal S(Z)|
      &\le6\sqrt2(944+96\cdot14)\epsilon\\
      &\le20592\epsilon,
 \end{split}
 \tag{95.21}
\]
where the last bound uses \(\sqrt2<3/2\).
Combining this with (95.7) and \(\epsilon=1/4096\) gives
\[
                    \mathcal S(Y)
                      \ge616-\frac{20592}{4096}
                      =\frac{156409}{256}.
 \tag{95.22}
\]
V89 proves
\(E_*\le569684412973/10^9<569.685\).
The rational difference between (95.22) and that upper bound
is greater than 41. A global minimum exists by compactness
and has energy no larger than \(E_*\), so it cannot lie in
the central tube.
\end{proof}

\subsection{Two consequences for all conditional ground states}

\begin{corollary}[All ground states have continuous conjugation degeneracy]
\label{cor:f12v91-ground-orbits}
Every conditional global action minimizer at this fixed
exterior is noncentral, with at least one independent link
at geodesic distance greater than \(1/4096\) from \(\{I,-I\}\).
Its simultaneous conjugation orbit has dimension 2 or 3
and is contained in the global ground set. Its hidden
Hessian has at least two zero directions.
\end{corollary}
\begin{proof}
The distance assertion is precisely exclusion from (95.11).
If all independent groups commute, at least one is noncentral
and their common stabilizer is a circle; the orbit is
SU(2)/U(1), of dimension 2. If the tuple contains two
noncommuting groups, the two corresponding imaginary
quaternion axes are not collinear. An adjoint rotation fixing
both is the identity, so the common stabilizer in SU(2)
is its centre and the orbit is SO(3), of dimension 3.
The fixed exterior and coarse field are central, so the
action is invariant under this conjugation.
Every orbit point is therefore a global minimum and is
critical. Differentiating criticality along the orbit puts
all its tangent directions in the Hessian kernel.
\end{proof}

This conclusion applies to all global minima, unlike the
normal positivity certificate for the one located orbit.
It does not prove that every global minimum is commuting,
that the remaining Hessian eigenvalues are positive, or that
kernel rank is constant across the ground set.
In particular, a ground-set analysis based on isolated
nondegenerate hidden minima is incompatible with this
central exterior. Symmetry must be retained or quotiented.

\begin{corollary}[Conditional probability of the entire central tube]
\label{cor:f12v91-probability}
For the same fixed conditional measure, there is a finite
constant \(C\), independent of \(\beta\), such that
\[
                 \pi_\beta(\mathcal N_\epsilon)
                      \le\min\{1,C\beta^{1415}e^{-41\beta}\},
                            \qquad \beta\ge1.
 \tag{95.23}
\]
\end{corollary}
\begin{proof}
Let \(a(Y)\) be the actual inverse-pivot Haar amplitude.
It is continuous and positive on the compact hidden space,
so \(A=\int a(Y)\,dY<\infty\).
The numerator obeys
\[
 \int_{\mathcal N_\epsilon}e^{-\beta\mathcal S(Y)}a(Y)\,dY
       \le A\exp(-\beta\inf_{\mathcal N_\epsilon}\mathcal S).
\]
The full normalizer is at least the local tube integral
of Corollary~\ref{cor:f12v90-tube}, namely
\(c_0\beta^{-1415}e^{-\beta E_*}\).
Divide and use (95.19). The normalization lower bound is
valid whether or not that local orbit is a global minimum;
lower-energy competitors only increase the denominator.
The bound by 1 is automatic.
\end{proof}

No factor \(2^{944}\) is omitted: the numerator above is
bounded using the total amplitude integral over the whole
hidden space, so it already includes the union of all
central neighbourhoods. The constants are fixed-parent
constants. Equation (95.23) is not an exit-rate estimate,
a conductance bound, or a spectral-gap lower bound.

\subsection{What this resolves and what remains upstream}

This iteration excludes a complete sector and a quantitative
noncommutative neighbourhood of it from the conditional global
ground set. It is a global-in-hidden-configurations result
at one specified exterior, not another restricted-slice
stability calculation. It also proves that every global
ground state at that exterior carries continuous conjugation
degeneracy.

The unresolved competitors now lie outside the central tube.
They include noncentral commuting orbits other than the
V89 orbit and potentially noncommuting configurations.
The integer parity inequalities do not apply directly to
arbitrary SU(2) plaquette products; replacing a quaternion
by a binary sign in (95.3) would be invalid.
A further relaxation or a ground-set argument must retain
that noncentral geometry.
Nor does the present single-exterior result verify V81's
coverage condition throughout a visible domain, its joint
force or negative-Hessian alternative, or the within-region
conductance needed for repair gluing.

The new exact verifier passes all geometry counts, dual
signs, 3236 residual inequalities, and quantitative tube
bounds. The V90 source had 1503861 bytes and SHA--256
\begin{center}\scriptsize\ttfamily
3A216824BAABAB59D046B7FCB9547E06CB03C891D20D1B8FB88F2FD14DC855FA.
\end{center}
Its body is preserved apart from the title version, and
only active V90 is replaced by V91. All prior scripts,
frozen data, and archives are retained. The V90 companion
audit is not repeated; the next is V95. The V100 direction
review and family-\(f12\) rollover remain scheduled.

\begin{firewall}[Position after V91]
All \(2^{944}\) central hidden assignments at the specified
finite SU(2) parent have action at least 616, and their
common radius-\(1/4096\) product-geodesic tube lies more
than 41 action units above an available noncentral competitor.
All conditional global ground states there are noncentral.
Complete noncentral ground-set coverage, unrestricted repair,
a cofinal auxiliary gap, continuum construction, and the
physical Yang--Mills mass gap remain unproved.
The full objective remains active.
\end{firewall}
\section{V92: the hidden minimum is joint-regular and has an exact coarse-link descent}
\label{sec:f12v92}

\subsection{The distinction needed for ground-set coverage}

V91 rules out all central hidden assignments and a uniform
neighbourhood of them as conditional ground states at the
specified exterior. V90 proves a full hidden local minimum
on the noncentral orbit \(\mathcal O\), but does not make it
a global minimum. The joint good-chart criterion of V81
asks a different question: are the relevant conditional
ground states regular or unstable when coarse variables
are also allowed to move?

This iteration answers that question for the entire located
orbit, not for all still-unknown conditional ground states.
One coarse link has a certified angular derivative greater
than \(49/5\) at the commuting representative. Its norm is
preserved along the conjugation orbit. Thus this hidden
minimum is \emph{joint-regular}, rather than a joint critical
point requiring a singular Schur test.
An explicit rational change of that coarse link decreases
the action by more than \(9/8\).

The coarse link is actually changed in this result.
The decreased configuration is not a competitor in the
same fixed-coarse-field hidden minimization problem.
There is consequently no contradiction with V90's local
minimum or with its possible, still-unproved global minimality.

\subsection{A coarse edge with a whole-ball inactive pivot gate}

Retain \(Y_*=Y(\xi_*)\), \(E_*=\mathcal S(Y_*)\), and the
critical-point box \(\mathcal Q\) from V89--V90. Select
\[
 e_0=((0,0,0,0),2),\qquad
 \ell_0=((0,0,1,0),2).
 \tag{96.1}
\]
Here \(e_0\) is a coarse edge and \(\ell_0\) its physical
pivot. It is not one of the four active inverse-root pivots.
For every commuting configuration in \(\mathcal Q\), the
integer interval evaluation gives, for all its fourteen
nonroot path products,
\[
        \|I-P_j\|_{\rm HS}^2
                    \ge\frac{6209511}{5000000}>1 .
 \tag{96.2}
\]
For a unitary \(P_j\), the squared chord used here is
\(4(1-(P_j)_0)\), with
\((P_j)_0\) its scalar quaternion component.
The same bound holds after simultaneously conjugating all
hidden groups and their completed pivots.

Let \(v(C)\) denote the coarse configuration obtained by
changing only \(e_0\) from \(I\) to \(C\), all other coarse
links and the fixed exterior remaining unchanged.

\begin{proposition}[Exact coarse pivot on a finite ball]
\label{prop:f12v92-coarse-pivot}
For every \(Y(\xi)\), \(\xi\in\mathcal Q\), and every
simultaneous conjugate of it, the actual pivot at \(e_0\)
for coarse input \(v(C)\) is
\[
                    B_{e_0}=C\qquad
                        \text{if }d(C,I)\le1/4.
 \tag{96.3}
\]
All other physical links are unchanged from their completed
values at coarse field \(I\).
\end{proposition}
\begin{proof}
Chordal distance is no greater than geodesic distance,
so (96.2) implies \(d(I,P_j)>1\).
For the proposed pivot \(B=C\),
\[
             d(C,P_j)\ge d(I,P_j)-d(I,C)>3/4>\rho.
 \tag{96.4}
\]
Every truncated logarithm in that pivot equation is zero.
Consequently its forward map sends \(B=C\) to \(C\).
V89's global SU(2) pivot diffeomorphism proves that this
is the unique actual solution, not just a proposal.
The nonroot path families contain no pivots, and changing
one coarse input does not change another family's input.
Thus all other pivots and independent physical links
remain as claimed. The same reasoning is invariant under
simultaneous conjugation.
\end{proof}

\subsection{The certified coarse force}

At a commuting field in \(\mathcal Q\), let \(\omega_q\)
be the physical plaquette phase at coarse field \(I\).
For the six plaquettes incident to \(\ell_0\), let
\(\sigma_q=\pm1\) be the occurrence orientation of that
pivot. Define
\[
          f(\xi)=2\sum_{q\ni\ell_0}\sigma_q\sin\omega_q,
       \qquad
          a(\xi)=2\sum_{q\ni\ell_0}\cos\omega_q.
 \tag{96.5}
\]
The simple bound \(|a(\xi)|\le12\) will suffice.
The exact interval assembly over the whole box gives
\[
\begin{aligned}
 \frac{196060493}{20000000}
          &\le f(\xi)\le\frac{490151233}{50000000},
             & f(\xi)&>\frac{49}{5},\\
 \frac{43848897}{10000000}
          &\le a(\xi)\le\frac{438488971}{100000000}.
\end{aligned}
 \tag{96.6}
\]
These are outward rational endpoints, not a numerical
derivative at an approximate stationary point.

The files implementing the accepted calculation are
\begin{center}\small\ttfamily
ym\_f12\_v92\_visible\_force\_intervals.py\\
verify\_ym\_f12\_v92\_joint\_regular\_orbit.py.
\end{center}
They reuse the fixed-point interval operations, frozen centre,
and whole-box trigonometric enclosures of V90.
Each free-link sine and cosine is enclosed over its interval
\(\widehat\xi_i+[-13/10^{13},13/10^{13}]\).
The four genuinely moving pivots are evaluated at their
actual angles \(-(\xi_a+\xi_b)/15\); the other pivots are \(I\).
The code multiplies the oriented plaquette factors, sums
their signed imaginary parts for \(f\), and their real
parts for \(a\). It also multiplies all fourteen path pairs
to verify (96.2). Every accepted operation is outward
integer/Fraction interval arithmetic. The optional decimal
scan is not used by the verifier.
Because \(\xi_*\in\mathcal Q\), all these bounds apply
to the exact critical point.

\begin{theorem}[Joint regularity of the entire hidden critical orbit]
\label{thm:f12v92-joint-regular}
At every point of \(\mathcal O\), with coarse field \(I\),
the joint action gradient has Hilbert--Schmidt product norm
strictly greater than 6. In particular, no point of
\(\mathcal O\) is joint critical. There is a joint
neighbourhood of this compact orbit on which that gradient
norm is greater than 3.
\end{theorem}
\begin{proof}
At the commuting representative, vary the selected coarse
input as \(C(t)=\exp(tT)\).
Proposition~\ref{prop:f12v92-coarse-pivot} identifies the
physical pivot with \(C(t)\) and leaves all other links fixed.
Direct differentiation of the six plaquettes gives
\[
        \left.\frac{d}{dt}\mathcal S(Y_*;v(e^{tT}))
                  \right|_{t=0}=f(\xi_*)>\frac{49}{5}.
 \tag{96.7}
\]
The circle fixing the commuting tuple forces the transverse
coarse derivatives at \(e_0\) to vanish.
In the Hilbert--Schmidt metric \(\|T\|=\sqrt2\), so the
norm of this one coarse-link gradient component is
\[
                     \frac{|f(\xi_*)|}{\sqrt2}
                         >\frac{98}{15}>6 .
 \tag{96.8}
\]
For \(Y_g=gY_*g^{-1}\), invariance gives the same component
rotated by \(\operatorname{Ad}_g\); its norm is unchanged.
The full joint gradient norm is no smaller.
The action and pivot maps are smooth, and \(\mathcal O\)
is compact. Continuity therefore supplies a joint open
neighbourhood on which the norm is greater than 3.
\end{proof}

The theorem concerns the physical action \(\mathcal S\).
For the local weighted potential
\(W_\beta=\beta\mathcal S-\log a_{\rm Haar}\), derivatives
of the positive smooth Haar amplitude are bounded on
a smaller compact neighbourhood. Thus the leading
joint force is not cancelled at large \(\beta\).
The regular-chart estimates already established in V81
may be applied on suitable smaller joint charts there;
their separate coverage hypothesis is not supplied by
this observation.

\subsection{An exact radial orbit envelope}

Write a coarse Lie-algebra vector as
\[
                 X=x_1E_1+x_2E_2+x_3T,\qquad
                 t=|x|=\|X\|_{\rm HS}/\sqrt2 .
 \tag{96.9}
\]
Define the \emph{orbit-restricted} energy envelope
\[
     E_{\mathcal O}(X)
         =\min_{g\in{\rm SU}(2)}
                    \mathcal S(gY_*g^{-1};v(e^X)).
 \tag{96.10}
\]
This is not defined by minimizing over all hidden fields.
In particular \(E_{\mathcal O}(0)=E_*\) regardless of
whether that orbit is a global hidden ground state.

\begin{theorem}[Exact orbit envelope and uniform radial descent]
\label{thm:f12v92-envelope}
Put \(f_*=f(\xi_*)\), \(a_*=a(\xi_*)\).
For \(0\le t\le1/8\),
\[
            E_{\mathcal O}(X)
                  =E_*+a_*(1-\cos t)-f_*\sin t
                  \le E_*-9t.
 \tag{96.11}
\]
At \(t=0\) the last inequality is an equality.
\end{theorem}
\begin{proof}
The bound \(\sqrt2\,t\le3/16<1/4\) permits the use of
(96.3) throughout this ball. Each incident plaquette
contains the moving pivot exactly once. All remaining
factors at the conjugated representative lie in the torus
with axis \(n=\operatorname{Ad}_gT\), which has unit
quaternion length.
Cyclic invariance of the trace combines the other factors
before taking the real part. Quaternion multiplication
therefore gives, for \(t>0\),
\[
 \mathcal S(Y_g;v(e^X))-E_*
        =a_*(1-\cos t)
                  +f_*\,\langle n,\widehat x\rangle\sin t,
       \qquad \widehat x=X/t .
 \tag{96.12}
\]
The pairing here is the Euclidean pairing of quaternion
generator coefficients, so both \(n\) and \(\widehat x\)
have length one.
Every direction \(n\in S^2\) is attained by conjugation.
Since \(f_*>0\) and \(\sin t\ge0\), minimizing over \(n\)
gives the equality in (96.11). The case \(t=0\) follows
by continuity.

For \(0\le t\le1/8\), the elementary Taylor bounds
\(1-\cos t\le t^2/2\) and \(\sin t\ge t-t^3/6\) yield
\[
 \begin{split}
  a_*(1-\cos t)-f_*\sin t
      &\le 6t^2-\frac{49}{5}(t-t^3/6)\\
      &=t\left(-\frac{49}{5}+6t+\frac{49}{30}t^2\right).
 \end{split}
 \tag{96.13}
\]
The increasing correction satisfies
\[
                6(1/8)+\frac{49}{30}(1/8)^2<\frac45.
 \tag{96.14}
\]
This gives (96.11), with strict inequality for \(t>0\).
\end{proof}

The envelope has a radial cusp at zero:
along every unit angular direction its one-sided
outward derivative equals \(-f_*\).
If, in a future step, \(E_*=E(I)\) is proved for the
full conditional ground energy \(E\), then
\[
                       E(v(e^X))\le E(I)-9|x|
 \tag{96.15}
\]
follows immediately. In that case \(E\) cannot be
differentiable at \(I\) in these three coarse directions:
the two opposite one-sided bounds are incompatible with
a single linear differential. Without that global-minimum
hypothesis, the unconditional statement is only
\(E(v(e^X))\le E_{\mathcal O}(X)\); no cusp assertion
about the full \(E\) has been established.

\begin{corollary}[A rational coarse-link descent endpoint]
\label{cor:f12v92-rational-descent}
Keeping the exact independent hidden groups \(Y_*\) fixed,
change only the coarse input at \(e_0\) to
\[
                        C_0=\frac{255}{257}I
                                       -\frac{32}{257}T.
 \tag{96.16}
\]
This is a unit quaternion, its actual physical pivot is
\(C_0\), and
\[
              \mathcal S(Y_*;v(C_0))
                    \le E_*-\frac{1448}{1285}
                    < E_*-\frac98.
 \tag{96.17}
\]
\end{corollary}
\begin{proof}
The identity \(255^2+32^2=257^2\) proves unitarity.
Its angular magnitude satisfies
\[
                         2\arctan(1/16)<1/8.
\]
Thus (96.3) applies. At the original axis \(T\), formula
(96.12) becomes
\[
       \mathcal S(Y_*;v(C_0))-E_*
                       =\frac2{257}a_*-\frac{32}{257}f_*.
 \tag{96.18}
\]
Use \(a_*\le12\) and \(f_*\ge49/5\).
The resulting decrease is
\((32\cdot49/5-24)/257=1448/1285>9/8\).
The hidden coordinates have not been numerically substituted
for \(Y_*\): (96.6) encloses its exact coefficients.
\end{proof}

\subsection{Noncommuting exploration and the adjoint it uses}

The new exploratory file
\begin{center}\small\ttfamily
explore\_ym\_f12\_v92\_full\_su2\_parent.py
\end{center}
allows all 944 hidden groups to move in SU(2), not only
in a common torus. It numerically solves all 96 actual
cutoff pivot equations at every evaluation.
This removes the commuting restriction from the exploratory
search, but does not turn it into a global optimizer.
The exact coarse-force certificate above does not import it.

For reproducibility, the noncommuting adjoint is as follows.
Use angular three-vector coordinates for the Lie algebra.
For \(Q=\cos r+\sin r\,n\), set
\[
 h(r)=\chi(\sqrt2r/\rho)r,\qquad
 L(Q)=h'(r)nn^{\mathsf T}
       +h(r)\cot r\,(I-nn^{\mathsf T})-h(r)[n]_\times.
 \tag{96.19}
\]
This is the differential of \(g\) under left multiplication
\(\exp(tv)Q\). Its right-multiplication differential is
\(L(Q)^{\mathsf T}\).
To see the signs directly, the transverse quaternion
variation under left multiplication has complex factor
\(e^{-ir}\); differentiation of the radial logarithm
multiplies it by \(h(r)/\sin r\).
The axial derivative is \(h'(r)\).
The formula extends as \(I\) at the identity and as zero
outside the cutoff, including a neighbourhood of \(-I\).

If \(B=\exp b\), the right-trivialized exponential Jacobian is
\[
 J_R(b)=nn^{\mathsf T}
       +\operatorname{sinc}r\,
            \{\cos r\,(I-nn^{\mathsf T})-\sin r\,[n]_\times\},
                   \quad r=|b|.
 \tag{96.20}
\]
At a solved pivot put
\[
 Q_j=B^{-1}P_j,\qquad
 A=\frac1{16}\sum_jL(Q_j),\qquad K=I-AJ_R(b).
 \tag{96.21}
\]
The small logarithmic pivot equation is
\[
                       b+\frac1{16}\sum_jg(Q_j)=0.
 \tag{96.22}
\]
Its \(b\)-derivative is \(K\): a change \(db\) gives
right pivot velocity \(J_R(b)db\), and changes \(Q_j\)
by left multiplication with its negative.
At an actual root \(K\) is invertible, equivalently to
the nonsingular pivot differential proved in V89.

For a physical right-angular pivot force \(F_B\), define
\[
        \lambda=-K^{-\mathsf T}J_R(b)^{\mathsf T}F_B,\qquad
        F_{P_j}=\frac1{16}L(Q_j)\lambda.
 \tag{96.23}
\]
This is the induced force on the right-angular path velocity.
Indeed a path variation \(v_j\) changes the left side of
(96.22) by \(L(Q_j)^{\mathsf T}v_j/16\);
implicit differentiation and pairing with \(F_B\) gives
exactly (96.23). For a two-link path \(P=UV\),
its right velocity is
\(\operatorname{Ad}_{V^{-1}}v_U+v_V\).
Consequently its endpoint forces are
\(\operatorname{Ad}_VF_P\) and \(F_P\).
Adding these induced forces to the direct physical
plaquette forces gives the complete independent-link
gradient. No pivot is frozen in that calculation.

The script uses normalized four-component quaternion
parameters. If a raw parameter is \(q_{\rm raw}\ne0\),
with \(q=q_{\rm raw}/\|q_{\rm raw}\|\) and right-angular
force \(F\), the corresponding Euclidean gradient is
\(q(0,F)/\|q_{\rm raw}\|\).
This follows because quaternion left multiplication is
orthogonal and \(q(0,F)\) is tangent to the unit sphere.

Directional finite-difference checks were made at the
V89 centre, a noncommuting perturbation of the central
field populating the cutoff transition, and a noncommuting
perturbation of the located orbit. The observed errors
were below \(5\cdot10^{-8}\); numerical pivot residuals
were below \(2\cdot10^{-14}\).
These are implementation diagnostics, not interval
proofs of arbitrary numerical root values or derivatives.
The algebra above specifies the exact differentiated map.

Three full-SU(2) starts were tested with a 160-iteration cap.
Two ended close to the known commuting orbit, at energies
approximately \(569.6844118644\) and \(569.6844118646\);
their two smaller imaginary-vector Gram eigenvalues
were below \(10^{-9}\).
The other stopped at its iteration cap near
\(569.6988188962\), still noncommuting and with a larger
gradient residual. No lower competitor was found.
Neither the returns to a commuting orbit nor the finite
list of starts excludes a different noncommuting global
ground state.

\subsection{The remaining coverage obligation}

The two latest iterations now have different, complementary
roles. V91 excludes all central tubes from the conditional
global ground set. The present theorem puts the entire
located noncentral orbit inside joint regular charts.
Its hidden symmetry kernel and positive normal Hessian
are therefore not an obstruction to joint regularity.
There is no need to calculate a Schur complement at
this orbit to decide whether its joint gradient vanishes:
it does not.

What remains is to show that \emph{all} conditional global
minima in the required visible region fall in suitable
joint good charts. One possible sufficient statement would
be complete ground-set containment in the neighbourhood
from Theorem~\ref{thm:f12v92-joint-regular}; a different
coverage proof could include further regular or unstable
branches. Neither statement is proved here.
The search diagnostics cannot replace it, and the rational
coarse descent does not prove global minimality at the
original coarse field. Upstream work must now control
the omitted noncentral competitors or prove a joint
regularity/instability alternative without classifying them.

The exact force verifier passes the whole-box force and
gate enclosures, Hilbert--Schmidt normalization, radial
Taylor inequality, and rational endpoint calculation.
V91 had 1520591 bytes and SHA--256
\begin{center}\scriptsize\ttfamily
C82FC2F210E8603DB442845DE7264179492255A15C2A685B42F904A8AFF550E2.
\end{center}
Its body is preserved apart from the title version.
Only active V91 is replaced by V92; archives and earlier
proof data are retained. The companion audit remains
scheduled for V95, and the direction review and
family-\(f12\) rollover for V100.

\begin{firewall}[Position after V92]
The located full hidden Morse--Bott minimum is rigorously
joint-regular on its entire conjugation orbit and admits
an explicit coarse-link descent. This verifies a useful
good-chart condition for that orbit, not coverage of all
conditional ground states. Complete noncentral coverage,
unrestricted repair, the cofinal auxiliary gap, continuum
construction, and the physical Yang--Mills mass gap remain
unproved. The full objective remains active.
\end{firewall}
\section{V93: nonabelian angle constraints, a spurious low-energy witness, and the cube identity that excludes it}
\label{sec:f12v93}

\subsection{Why the next global relaxation must couple plaquettes}

V91 excludes the whole central hidden sector, while V92
places the located noncentral orbit in joint regular charts.
Neither result covers every possible noncentral global
minimum. This iteration tests a genuinely global necessary
condition: rotation-angle inequalities valid for arbitrary
noncommuting SU(2) fields.

The local inequalities have an exact extension beyond the
binary sector. However, by themselves they admit a rational
point of true cosine cost below 529.095, more than 40 below
the located orbit. This point is not a lower-action gauge
configuration. An explicit elementary-cube identity proves
that it cannot be realized by any link field.
The missing constraint is therefore identified algebraically,
not merely inferred from an unsuccessful numerical search.

\subsection{The angle form of the parity inequalities}

For a unit quaternion \(U\), define its normalized angle
\[
                 \vartheta(U)=\frac1\pi\arccos(U_0)
                                   \in[0,1],
 \tag{97.1}
\]
where \(U_0\) is its scalar component.
It satisfies
\[
 \vartheta(U^{-1})=\vartheta(U),\quad
 \vartheta(gUg^{-1})=\vartheta(U),\quad
 \vartheta(-U)=1-\vartheta(U).
 \tag{97.2}
\]
The geodesic distance is
\(d(I,U)=\sqrt2\pi\vartheta(U)\).
Its product triangle inequality follows from bi-invariance.

\begin{proposition}[Parity inequalities for noncommuting products]
\label{prop:f12v93-angle-parity}
Suppose \(Q_1\cdots Q_m=(-I)^c\), with \(c\in\{0,1\}\),
and put \(t_j=\vartheta(Q_j)\).
For every \(b\in\{0,1\}^m\) whose parity differs from \(c\),
\[
                    \sum_{j=1}^m(2b_j-1)t_j
                              \le\sum_{j=1}^m b_j-1.
 \tag{97.3}
\]
No commutativity assumption is required.
\end{proposition}
\begin{proof}
Central signs commute with all factors, so
\[
                \prod_{j=1}^m(-I)^{b_j}Q_j=-I.
 \tag{97.4}
\]
The product distance inequality and (97.2) imply
\[
        1\le\sum_{j=1}^m
                       \vartheta((-I)^{b_j}Q_j)
          =\sum_{j=1}^m|t_j-b_j|.
 \tag{97.5}
\]
Expand the absolute values using \(0\le t_j\le1\).
The result is exactly (97.3).
\end{proof}

Apply this proposition to the variable oriented link factors
of a plaquette, followed by the inverse of its holonomy.
The fixed central factors can be moved out of the product
and contribute a parity \(c_q\).
With \(x_\ell=\vartheta(U_\ell)\) and
\(y_q=\vartheta(U_q)\), this gives V91's local parity
inequalities on \((x_\ell:\ell\in q,\ y_q)\).
Inverses have the same angle, so link orientation signs
do not enter those inequalities.
This is only a necessary-angle relaxation; it does not
identify link axes or impose compatibility between all
plaquette products.

\subsection{The actual moving pivots must be retained}

For arbitrary noncentral fields the pivots are not generally
\(I\). The present relaxation therefore has angle variables
for all 944 independent hidden links and all 96 pivots.
V89's global pivot estimate gives
\[
       d(B_e,I)\le\frac{7\rho}{8},
       \qquad
       \vartheta(B_e)\le\frac{7\rho}{8\sqrt2\pi}.
 \tag{97.6}
\]
At \(\rho=1/64\), the elementary lower bounds
\(\sqrt2>7/5\) and \(\pi>157/50\) give the convenient
rational cap
\[
                     0\le x_e\le\frac{125}{40192}
                                \quad\text{for each pivot}.
 \tag{97.7}
\]
All independent-link and plaquette angles lie in \([0,1]\).
Only the cap, not the full pivot equation, is imposed in
the relaxation.

Keeping these 1040 physical link variables leaves 5376
fixed plaquettes with total action 408 and 2400 variable
plaquettes. The resulting angle vector has dimension 3440.
The local parity system has 17184 inequalities
\[
                         At\le b,
 \tag{97.8}
\]
with integer coefficients. Its exact enumeration is in
\begin{center}\small\ttfamily
ym\_f12\_v93\_angle\_geometry.py.
\end{center}
Link order is the previous 944-link order followed by
the 96 pivots in the parent's edge order. Variable
plaquettes are ordered by their sorted active-link indices
and fixed parity, as in the script.

The true action associated with a relaxed angle assignment
is
\[
              \mathcal E_{\rm ang}(t)
                   =408+\sum_{q=1}^{2400}
                                    2(1-\cos(\pi y_q)).
 \tag{97.9}
\]
Every actual completed hidden configuration maps to a
feasible point and has exactly this cost.
The linear objective \(4y_q\) used for binary central
variables in V91 is not valid here: it generally exceeds
\(2(1-\cos(\pi y_q))\) near zero.
Thus the central action certificate cannot be carried
over by changing only the domain of its variables.

\subsection{An exact feasible point below the located orbit}

\begin{proposition}[A low-cost rational local-angle witness]
\label{prop:f12v93-witness}
The frozen angle vector
\[
                         t^{\rm w}=\frac{k}{10^9},
                                     \quad k\in\mathbb Z^{3440},
 \tag{97.10}
\]
specified in
\begin{center}\small\ttfamily
ym\_f12\_v93\_angle\_witness\_data.py
\end{center}
satisfies all 17184 inequalities (97.8), the ordinary
box bounds, and all 96 pivot caps (97.7). Its true cost
satisfies
\[
 \frac{529094296741}{10^9}
       \le\mathcal E_{\rm ang}(t^{\rm w})
       \le\frac{52909429837}{10^8}<529.095.
 \tag{97.11}
\]
The V89--V90 orbit has energy
\[
          \frac{113936882369}{200000000}
                 \le E_*\le\frac{569684411883}{10^9},
 \tag{97.12}
\]
so \(E_*-\mathcal E_{\rm ang}(t^{\rm w})>40\).
\end{proposition}
\begin{proof}
The verifier imports the frozen integers and reconstructs
the geometry. It checks
\[
 0\le k_j\le10^9,\qquad
 40192k_j\le125\cdot10^9\quad(944\le j<1040),
       \qquad
                  Ak-10^9b\le0 .
 \tag{97.13}
\]
The largest entry in the last vector is exactly zero.
These feasibility checks use arbitrary-precision integers.

For the cost, \(\pi\) is enclosed rationally using
\[
             \pi=16\arctan(1/5)-4\arctan(1/239).
 \tag{97.14}
\]
For completeness, the tangent double-angle formula gives
\(\tan(4\arctan(1/5))=120/119\), and hence
\[
 \tan(4\arctan(1/5)-\arctan(1/239))
       =\frac{120/119-1/239}{1+(120/119)/239}=1.
 \tag{97.15}
\]
The angle is in \((0,\pi/2)\), so it equals \(\pi/4\).
Each arctangent is bounded by its first 30 alternating
series terms, widened by \(x^{61}/61\).
The resulting interval also verifies \(\pi>157/50\).

For each \(z=\pi k_j/10^9\), the fixed-point interval
calculation evaluates
\[
          \cos z=\sum_{r=0}^{20}\frac{(-1)^rz^{2r}}{(2r)!}
                         +R(z),\qquad
          |R(z)|\le\frac{|z|^{42}}{42!}.
 \tag{97.16}
\]
The bound is Taylor's theorem through degree 41;
the degree-41 coefficient vanishes.
All arguments have absolute value less than 4.
The polynomial is evaluated with V90's outward interval
operations at denominator \(2^{70}\), including the
rational enclosure of \(\pi\) and all roundoff.
Summing the enclosed true costs proves (97.11).
No linearized plaquette cost is used in this check.

To obtain (97.12), the verifier separately reconstructs
the complete physical field over V90's critical-point box.
It retains the four inverse-root pivot angles and sums
all 7776 physical plaquette action intervals.
Since \(\xi_*\) lies in that box, the enclosure applies
to its exact energy. The two rational intervals imply
the stated strict separation.
\end{proof}

The data file has SHA--256
\begin{center}\scriptsize\ttfamily
7B756728095D8AE7F74A5DDEC8AD57CE169E6721CB1F77D3983C85AEAC4D68D3.
\end{center}
The accepted checker is
\begin{center}\small\ttfamily
verify\_ym\_f12\_v93\_angle\_obstruction.py.
\end{center}
The optimizer used to propose the vector is not imported
by this checker. As with the earlier finite certificates,
this is reproducible exact arithmetic plus proved enclosure
rules, not a formal proof-assistant certification.

\subsection{The six-face nonabelian cube identity}

The preceding witness does not establish an actual action
value below \(E_*\). The following identity detects its
failure to come from any field.

Let \(a<b<c\) be three lattice directions and \(x\) a
cube base. Denote its six positively oriented face
holonomies by
\[
 \begin{gathered}
 F_{ab}^0=U_{ab}(x),\quad F_{ab}^c=U_{ab}(x+e_c),\\
 F_{ac}^0=U_{ac}(x),\quad F_{ac}^b=U_{ac}(x+e_b),\\
 F_{bc}^0=U_{bc}(x),\quad F_{bc}^a=U_{bc}(x+e_a),
 \end{gathered}
 \tag{97.17}
\]
and put \(A_0=U_a(x)\), \(B_0=U_b(x)\), \(C_0=U_c(x)\).
The plaquette convention is the same four-factor
chronological convention used throughout these notes.

\begin{proposition}[Cube identity and angle consequences]
\label{prop:f12v93-cube}
For every lattice group field,
\[
 F_{ab}^0(B_0F_{ac}^bB_0^{-1})F_{bc}^0
   =(A_0F_{bc}^aA_0^{-1})F_{ac}^0
                       (C_0F_{ab}^cC_0^{-1}).
 \tag{97.18}
\]
For SU(2), its six normalized face angles \(y_1,\ldots,y_6\)
satisfy
\[
             \sum_{j=1}^6(2b_j-1)y_j
                         \le\sum_{j=1}^6b_j-1
             \quad\text{whenever }\sum_jb_j\text{ is odd}.
 \tag{97.19}
\]
In particular every face angle is no greater than the
sum of the other five.
\end{proposition}
\begin{proof}
Expand both sides of (97.18) using the oriented plaquette
definition. Adjacent inverse edge factors cancel, and
both sides reduce to the same six-edge path word
\[
 \begin{split}
 &U_a(x)\,U_b(x+e_a)\,U_c(x+e_a+e_b)\\
 &\hspace{1em}\cdot U_a(x+e_b+e_c)^{-1}
                  U_b(x+e_c)^{-1}U_c(x)^{-1}.
 \end{split}
 \tag{97.20}
\]
This proves the identity without commuting any factors.
Equivalently, multiply the left side by the inverse
of the right side: it is a product of six conjugated
or inverted face holonomies equal to \(I\).
Conjugation and inversion preserve their angles.
Proposition~\ref{prop:f12v93-angle-parity}, with \(c=0\),
therefore proves (97.19).
Choosing \(b\) with a single 1 gives the final assertion.
The verifier additionally expands this 30-letter group
word and checks its exact free reduction to the empty
word; that is an algebraic identity, not a numerical
test on selected matrices.
\end{proof}

\begin{theorem}[The rational witness is not realizable by link variables]
\label{thm:f12v93-no-lift}
No SU(2) lattice field has all the link and plaquette
angles of \(t^{\rm w}\). This obstruction already follows
from one elementary cube, independently of the nonlinear
pivot equations.
\end{theorem}
\begin{proof}
Choose the cube
\[
                         x=(2,0,1,3),\qquad
                         (a,b,c)=(0,2,3).
 \tag{97.21}
\]
In the face order of (97.17), the witness and its fixed
exterior assign the following normalized angles:
\[
 \frac1{10^9}
 (8303120,\ 1000000000,\ 500000,\
                  47318399,\ 8303120,\ 39516978).
 \tag{97.22}
\]
The second face is fixed by the exterior and has angle 1.
The other five sum to \(103941617/10^9\).
Thus the necessary single-face triangle inequality is
violated by
\[
                              \frac{896058383}{10^9}>0.
 \tag{97.23}
\]
The exact checker obtains these angles directly from
the physical plaquette-to-variable map, rather than
assuming that their order matches a solver output.
Equation (97.19) rules out any realizing group field.
\end{proof}

Consequently the local link/plaquette angle conditions
alone cannot prove an action lower bound \(E_*\).
This is not merely a shortcoming of piecewise linear
approximations to the cosine: the infeasible-as-a-field
point has true cosine cost below 529.095.
It is not a counterexample to global minimality of the
located orbit, nor to a Yang--Mills mass gap.
It identifies a genuine missing compatibility condition
in this proposed relaxation.

\subsection{The strictly strengthened relaxation}

Let \(\mathcal R_0\) be the compact local-angle feasible
set defined by (97.7)--(97.8) and the ordinary box bounds.
Add (97.19) for every elementary cube, substituting its
fixed exterior face angles and keeping the variable ones.
The resulting set is denoted \(\mathcal R_1\).
The exact enumeration produces 32768 additional nontrivial
inequalities. Cubes with all faces fixed already satisfy
their parity condition. The frozen witness violates 635
of the enumerated inequalities, with maximum violation
(97.23).

\begin{proposition}[Strict refinement without loss of physical fields]
\label{prop:f12v93-refinement}
Every actual completed hidden SU(2) configuration maps
into \(\mathcal R_1\), and
\[
                         \mathcal R_1\subsetneq\mathcal R_0.
 \tag{97.24}
\]
If \(L_i=\min_{\mathcal R_i}\mathcal E_{\rm ang}\) and
\(E_{\min}\) is the actual conditional ground energy, then
\[
                       L_0\le L_1\le E_{\min}\le E_*.
 \tag{97.25}
\]
No strict inequality between \(L_0\) and \(L_1\), or
equality between \(L_1\) and \(E_{\min}\), is asserted.
\end{proposition}
\begin{proof}
The pivot cap and local angle inequalities hold for
every actual configuration, and the cube inequalities
follow from its exact holonomies by (97.18)--(97.19).
Thus the strengthened set still contains all physical
angle data. It is nonempty and compact.
Theorem~\ref{thm:f12v93-no-lift} supplies a point of
\(\mathcal R_0\) not in \(\mathcal R_1\), proving strict
containment. Minimizing the same continuous objective
over these nested sets gives (97.25).
Strict containment alone does not imply a strict change
in the minimum value.
\end{proof}

The exploratory script
\begin{center}\small\ttfamily
explore\_ym\_f12\_v93\_angle\_relaxation.py
\end{center}
implements both versions. Its piecewise-linear search
model introduces one cost variable per variable plaquette
and uses 17 tangent proposals to the cosine cost, at
normalized angles \(j/64\), \(0\le j\le16\).
Without cube constraints it reported a floating-point
value near 520.4771902. Restoring the true cosine costs
of that numerical point gave approximately 528.5941371;
the slightly modified rational feasible witness has
the rigorously checked cost (97.11).
The reported linear-solver value is not an accepted
global lower certificate.

With all cube inequalities the search model has 5840
variables and 90752 inequalities. Bounded runs with
45-second and 120-second solver limits both terminated
at their stated time limits without an accepted optimum;
the latter reported 48180 iterations.
No optimized bound from either run is used here.
These are completed numerical attempts, not live jobs
being implicitly left running.

The further exploratory script
\begin{center}\small\ttfamily
explore\_ym\_f12\_v93\_angle\_lifts.py
\end{center}
used the proposed independent-link angles to seed actual
completed SU(2) fields with three choices of axes.
It discarded the independently assigned relaxed plaquette
angles and recomputed every actual holonomy and pivot.
Initial physical energies were approximately 636.163,
939.172, and 1438.374; all three searches returned near
the known value 569.684411864.
These results do not prove uniqueness or global minimality.
The rigorous nonrealizability result is instead the
single cube contradiction (97.23).

\subsection{Upstream direction and continuation record}

The new obstruction says exactly where this attempted
global argument lost information: independently feasible
plaquette angle data need not satisfy an elementary
nonabelian holonomy identity. Refining the local cosine
cost approximation cannot repair that loss.
The next relaxation must retain at least the cube
constraints now proved, and may need further shared-axis,
larger-loop, or actual pivot-equation information.
Even cube-angle consistency has not been proved sufficient
for realization by a link field.

This step therefore does not replace the remaining
ground-set coverage theorem with a low-dimensional
surrogate. It gives an exact necessary condition,
a verifiable failure of the weaker surrogate, and
a strictly strengthened feasible set that still contains
every physical configuration. A useful next lower
certificate must concern that stronger set or an
equally justified alternative. Any numerical optimum
must be converted into an independently checked bound
before it is used to exclude actual ground states.

The accepted verifier passes all rational bounds,
all 17184 local inequalities, the true cosine-cost
enclosure, the exact orbit-energy enclosure, the cube
mapping and violation checks, and free-word cancellation.
V92 had 1538157 bytes and SHA--256
\begin{center}\scriptsize\ttfamily
9C41235E65C02313FAFB1A0C988BB3C8812109BBC350010D10278DC05D8C6CEF.
\end{center}
Its body is preserved apart from the title version.
Only active V92 is replaced by V93. Prior scripts,
frozen data, archives, and companion notes are retained.
The next companion audit is V95. The V100 direction
review and family-\(f12\) rollover remain scheduled.

\begin{firewall}[Position after V93]
The global local-angle relaxation has been rigorously
shown to admit a nonphysical low-cost point, and an
exact nonabelian cube identity excludes that point.
The strengthened cube-angle constraints are valid for
every physical hidden field, but their optimized value
and their sufficiency for ground-set coverage are unproved.
Complete noncentral coverage, unrestricted repair,
the cofinal auxiliary gap, continuum construction,
and the physical Yang--Mills mass gap remain unproved.
The full objective remains active.
\end{firewall}
\section{V94: an exact global action bound and the larger-box obstruction beyond elementary cubes}
\label{sec:f12v94}

\subsection{A global certificate, with its remaining limitation}

The cube-consistent relaxation of V93 now gives an
independently checked lower bound for \emph{every} completed
hidden SU(2) field at the specified central exterior,
coarse field \(I\), and cutoff \(\rho=1/64\):
\[
                         \mathcal S(Y)\ge
                              \frac{21983613511}{40000000}
                              =549.590337775.
 \tag{98.1}
\]
This is a finite-parent action bound, not a spectral or
physical mass gap. Together with the existing competitor,
it brackets the conditional global ground energy, but
does not identify its minimizers.

The relaxation remains incomplete. A rational point
satisfies every local and elementary-cube angle inequality
and has true cosine cost below 564.652, yet it violates
a larger-box holonomy identity. Thus even elementary-cube
consistency alone does not characterize physical angle data.
Both the global lower bound and this further obstruction
are established below without trusting numerical solver
termination as a proof.

\subsection{Globally valid cosine minorants}

Let
\[
                         g(y)=2(1-\cos(\pi y)),
                                    \qquad 0\le y\le1.
 \tag{98.2}
\]
The linear model uses tangent minorants at
\(q_j=j/64\), \(0\le j\le16\). These must be valid on
the whole interval, including its concave part.

\begin{proposition}[A family of global tangent minorants]
\label{prop:f12v94-minorants}
For \(0\le q\le1/4\), put
\[
           \ell_q(y)=g(q)+g'(q)(y-q).
 \tag{98.3}
\]
Then \(\ell_q(y)\le g(y)\) for every \(0\le y\le1\).
\end{proposition}
\begin{proof}
The function \(g\) is convex on \([0,1/2]\), so the
tangent bound holds there. On \([1/2,1]\), the difference
\(g-\ell_q\) is concave, and its minimum is attained
at an endpoint. At \(1/2\) the desired inequality already
holds. It remains to check \(\ell_q(1)\le4\).

Differentiation with respect to \(q\) gives
\[
              \frac{d}{dq}\ell_q(1)=g''(q)(1-q)\ge0
                                      \quad(0\le q\le1/4).
 \tag{98.4}
\]
Its maximum is therefore at \(q=1/4\), where
\[
         \ell_{1/4}(1)=2+\sqrt2(3\pi/4-1)
             <2+\frac{10}{7}\left(\frac{3}{4}\frac{22}{7}-1\right)
             <4.
 \tag{98.5}
\]
The coefficient \(3\pi/4-1\) is positive, so the stated
upper bounds may be substituted in that order.
This proves global validity.
\end{proof}

The exact verifier encloses
\[
             \alpha_j=2\pi\sin(\pi q_j),\qquad
             \beta_j=g(q_j)-q_j\alpha_j
 \tag{98.6}
\]
using V93's rational enclosure of \(\pi\), outward
trigonometric intervals, and Taylor remainder bounds.
It rounds each lower endpoint down to denominator
\(P=10^9\), obtaining integers \(a_j,b_j\) with
\[
                  a_j/P\le\alpha_j,\qquad
                  b_j/P\le\beta_j.
 \tag{98.7}
\]
Since \(y\ge0\), the rational affine function
\((a_jy+b_j)/P\) remains a global minorant.
The verifier uses the degree-17 sine polynomial with
remainder \(|z|^{19}/19!\), at \(|z|<1\), and the
previous degree-40 cosine enclosure.
Rounding both coefficients downward is deliberate;
nearest rounding would not by itself preserve the bound.

\subsection{The finite linear programme and an exact residual certificate}

Retain all 3440 link and plaquette angle variables of V93,
including the 96 pivot angles. Introduce a variable \(z_q\)
for each of the 2400 plaquette costs and impose
\[
       a_j y_q-Pz_q\le-b_j,\qquad
       0\le z_q\le4,\qquad 0\le j\le16.
 \tag{98.8}
\]
Together with the 17184 local inequalities, 32768 cube
inequalities, ordinary box bounds, and pivot cap
\(125/40192\), this is a linear relaxation with 5840
variables and 90752 non-box inequalities.
Write it as
\[
             \mathcal A t\le d,\qquad
             0\le t\le U,\qquad
             \text{objective }408+c^{\mathsf T}t.
 \tag{98.9}
\]
Here \(c\) is zero on the angles and one on each cost
variable. The upper bound \(U_i\) is 1 on ordinary
angles, \(125/40192\) on pivot angles, and 4 on costs.
Every actual physical configuration supplies a feasible
point with \(z_q=g(y_q)\), by V93 and
Proposition~\ref{prop:f12v94-minorants}.

\begin{proposition}[Dual bound with exact box compensation]
\label{prop:f12v94-dual}
For any \(\lambda\le0\), define
\[
                     r=c-\mathcal A^{\mathsf T}\lambda.
 \tag{98.10}
\]
Every feasible \(t\) satisfies
\[
 408+c^{\mathsf T}t
       \ge408+d^{\mathsf T}\lambda+
                           \sum_i\min(r_i,0)U_i.
 \tag{98.11}
\]
\end{proposition}
\begin{proof}
Multiplication of \(\mathcal A t\le d\) by
\(\lambda\le0\) gives
\(\lambda^{\mathsf T}\mathcal A t\ge\lambda^{\mathsf T}d\).
Also \(r_it_i\ge\min(r_i,0)U_i\) on the given box.
Add these inequalities and use
\(c=\mathcal A^{\mathsf T}\lambda+r\).
\end{proof}

\begin{theorem}[Certified global action lower bound]
\label{thm:f12v94-global-bound}
For every actual completed hidden configuration of the
specified finite parent, (98.1) holds. In particular,
\[
                 \frac{1099}{2}<E_{\min}
                         \le E_*<569.685,
 \tag{98.12}
\]
where \(E_{\min}\) is its conditional global action minimum.
\end{theorem}
\begin{proof}
The frozen dual file
\begin{center}\small\ttfamily
ym\_f12\_v94\_angle\_dual\_data.py
\end{center}
specifies \(\lambda=v/10^8\) on the 90752 rows, with
3621 nonzero entries. The verifier checks every
\(v_j\le0\) and rebuilds every rational coefficient of
(98.9), including the downward-rounded minorants.
Its exact residual has 1037 negative coordinates.
They are not silently set to zero.

The outward reports for (98.11) are
\[
\begin{array}{c|c}
 408+d^{\mathsf T}\lambda\text{ lower}
                          &110322264321/200000000\\
 \sum_i\min(r_i,0)U_i\text{ lower}
                          &-202098383/100000000\\
 \text{combined lower}    &21983613511/40000000.
\end{array}
 \tag{98.13}
\]
The correction of about \(-2.021\) is essential:
it includes the effects of box bounds, not merely small
floating-point roundoff. Using the first row alone
would not be the certificate proved here.

For an entirely integer implementation, multiply matrix
rows and right sides by \(P=10^9\).
The residual numerators are
\[
        R_i=P\,10^8c_i-\sum_j(P\mathcal A_{ji})v_j.
 \tag{98.14}
\]
All upper bounds share denominator \(40192\), with
numerators \(40192\), \(125\), and \(4\cdot40192\)
on the three types of coordinates.
Consequently the complete lower numerator has denominator
\(P\,10^8\,40192\) and is evaluated using only
arbitrary-precision integer sums and products.
The final rational value, rounded downward as in
(98.13), proves (98.1) by
Proposition~\ref{prop:f12v94-dual}.
Physical angle data are feasible, so the bound applies
to all commuting and noncommuting hidden configurations.
The known competitor supplies the upper bound in (98.12).
\end{proof}

The dual file has SHA--256
\begin{center}\scriptsize\ttfamily
8FD47F32398F4AC5D59AB64CBFFDA22328F8A64723427C3F6E3D55E09344B3DA.
\end{center}
Its integers are stored by lossless little-endian packing,
compression, and base-85 encoding. The decoded length
and all certificate inequalities are checked; this is
serialization, not a numerical inference.
The accepted checker is
\begin{center}\small\ttfamily
verify\_ym\_f12\_v94\_cube\_bound.py.
\end{center}
An interior-point solver supplied the proposal in 36
iterations after V93's simplex attempts reached their
limits. Its value near 549.5903424 is not used as proof.
The rational bound (98.13) is independently checked.

\subsection{A rational point satisfying every elementary cube constraint}

The second frozen file,
\begin{center}\small\ttfamily
ym\_f12\_v94\_cube\_point\_data.py,
\end{center}
contains a 3440-dimensional rational angle point
\(\widehat t=k/10^{10}\).
It uses the same angle order as V93 and has SHA--256
\begin{center}\scriptsize\ttfamily
7847A1660DC79B4BF3C52A5BBFFD3894AFEB860E55FB297BC6490A44A63FCB7E.
\end{center}

\begin{proposition}[A low-cost cube-consistent angle point]
\label{prop:f12v94-point}
The point \(\widehat t\) satisfies all ordinary angle
bounds, the 96 pivot caps, all 17184 local inequalities,
and all 32768 elementary-cube inequalities. Its true cost
satisfies
\[
       \frac{282325994759}{500000000}
          \le\mathcal E_{\rm ang}(\widehat t)
          \le\frac{564651989587}{10^9}<564.652.
 \tag{98.15}
\]
It lies more than 5 action units below the located
orbit's energy \(E_*\).
\end{proposition}
\begin{proof}
The checker reconstructs all 49952 angle inequalities
and checks their integer defects
\(Ak-10^{10}b\le0\); the largest is zero.
It separately checks the box bounds and
\(40192k_i\le125\cdot10^{10}\) on pivot coordinates.
The true cosine cost is enclosed by the same rational
\(\pi\) and outward Taylor method as in V93.
This gives (98.15), without replacing the cosine by
its linear minorants.
The exact orbit-energy enclosure (97.12) proves the
separation greater than 5.
\end{proof}

This is still not a physical competitor. The next result
detects information not contained in the elementary-cube
angle inequalities.

\subsection{Rectangular-boundary identities}

Replace each edge of an elementary cube identity by the
ordered product of links along the corresponding edge
of an axis-aligned rectangular box.
V93's nonabelian identity remains valid because it is
a group-word identity. Each rectangular macroface can
then be factored into conjugates of its elementary
plaquette holonomies.

The factorization can be seen directly in a two-dimensional
rectangle. Let \(F_{ij}\) be its fine plaquette at \((i,j)\),
\(A_{00}\) its first horizontal edge, and \(B_{ij}\)
its vertical edge at \((i,j)\).
A \(2\)-by-\(1\) rectangle has holonomy
\[
                       (A_{00}F_{10}A_{00}^{-1})F_{00},
 \tag{98.16}
\]
and a \(1\)-by-\(2\) rectangle has holonomy
\[
                       F_{00}(B_{00}F_{01}B_{00}^{-1}).
 \tag{98.17}
\]
These follow by direct cancellation of the shared edge.
For a \(2\)-by-\(2\) macroface the resulting expression is
\[
 A_{00}\{F_{10}(B_{10}F_{11}B_{10}^{-1})\}A_{00}^{-1}
       \{F_{00}(B_{00}F_{01}B_{00}^{-1})\}.
 \tag{98.18}
\]
It is a product of four conjugated elementary holonomies.
Expanding its 22 fine-edge letters gives the eight-edge
rectangular boundary word. The verifier checks that
their quotient reduces freely to the empty word.

\begin{proposition}[Angle inequalities for a rectangular box]
\label{prop:f12v94-box}
For a rectangular box with side lengths 1 or 2 in
three lattice directions, let \(y_1,\ldots,y_m\) be
the normalized angles of all its elementary boundary
plaquettes. Then
\[
              \sum_{j=1}^m(2b_j-1)y_j
                         \le\sum_{j=1}^m b_j-1
                   \qquad\left(\sum_jb_j\text{ odd}\right).
 \tag{98.19}
\]
In particular \(y_i\le\sum_{j\ne i}y_j\) for every
boundary face.
\end{proposition}
\begin{proof}
Apply the cube group-word identity to the box's twelve
macroedge path products. Substitute (98.16)--(98.18)
for each macroface of size larger than one.
Conjugation distributes over products, and inverses
reverse the order while inverting the factors.
The result is a product of conjugated or inverted
elementary boundary holonomies equal to \(I\), with
each boundary plaquette appearing once.
Their angles are unchanged by these operations.
Apply Proposition~\ref{prop:f12v93-angle-parity} with
central parity zero to obtain (98.19).
The periodic side is 6 and the box sides are at most 2,
so these boxes have no self-identifications. The same
word identities apply when an edge crosses a periodic
coordinate seam.
\end{proof}

\begin{theorem}[The cube-consistent point has no physical realization]
\label{thm:f12v94-no-lift}
The point \(\widehat t\) cannot be realized by an SU(2)
link field, despite satisfying every elementary-cube
angle inequality.
\end{theorem}
\begin{proof}
Take the box based at
\[
                   x=(2,2,5,5),\qquad
                   (a,b,c)=(0,2,3),\qquad
                   (L_a,L_b,L_c)=(2,2,2).
 \tag{98.20}
\]
It has 24 elementary boundary plaquettes.
The face based at \((2,2,0,5)\) in directions \((2,3)\)
has witness angle \(4811811675/10^{10}\).
The sum of the other 23 boundary angles is
\(4309914329/10^{10}\).
Consequently its single-face inequality (98.19) is
violated by
\[
                         \frac{501897346}{10^{10}}
                           =\frac{250948673}{5000000000}>0.
 \tag{98.21}
\]
The verifier reconstructs all boundary face indices
with periodic shifts and checks these exact integers.
Proposition~\ref{prop:f12v94-box} excludes every physical
realization, independently of the nonlinear pivot equations.
\end{proof}

Thus elementary-cube angle consistency is a strictly
weaker requirement than consistency of all physical
holonomies. The cost in (98.15) is not an upper bound
on \(E_{\min}\). It is a feasible value for an incomplete
relaxation and is proved nonphysical.

\subsection{Why adjoining two cubes adds no new angle constraint}

The first larger surfaces tested were boundaries of
two adjacent cubes. The absence of a violation there
is explained algebraically, and should not be treated
as evidence for realizability.

\begin{proposition}[Redundancy for tree-like unions of checks]
\label{prop:f12v94-tree}
Suppose two even-parity angle-inequality systems have
exactly one shared variable \(s\). Then every
even-parity-system inequality on their combined external
variables is a sum of one inequality from each system.
Iterating this fact, boundary parity inequalities of a
collection of cubes whose face-adjacency graph is a tree
are implied by the individual cube inequalities.
\end{proposition}
\begin{proof}
Partition the external variables into the first and
second groups, and choose forbidden external bits
\(b^{(1)},b^{(2)}\) with odd total parity.
Their parities \(p_1,p_2\) differ.
In the first local inequality choose the shared forbidden
bit \(b_s^{(1)}=1-p_1\), and in the second choose
\(b_s^{(2)}=1-p_2\). Each complete local bit vector
has odd parity, so both inequalities are available.
The shared bits are complementary, and their coefficients
\(2b_s^{(i)}-1\) cancel.
The two right sides sum to
\[
        |b^{(1)}|+|b^{(2)}|
                   +b_s^{(1)}+b_s^{(2)}-2
             =|b^{(1)}|+|b^{(2)}|-1.
 \tag{98.22}
\]
This is precisely the desired external inequality.
For a tree, successively remove leaves or attach each
new cube along its single shared face and apply the
same argument inductively.
\end{proof}

The proof applies before fixed face angles are substituted;
substitution preserves the implication. It shows why
adding only two-cube boundaries cannot strengthen the
elementary-cube feasible set. Unions with cycles in
their cube-adjacency graph can add constraints:
Theorem~\ref{thm:f12v94-no-lift} gives an explicit example.

For complete finite testing of the chosen rectangular
surfaces, exponential enumeration of all odd subsets
is unnecessary. Given \(y_j\in[0,1]\), set
\[
 d_0=\sum_j\min(y_j,1-y_j),\qquad
 b_j^0=\mathbf1_{\{y_j>1/2\}}.
 \tag{98.23}
\]
The minimum of \(\sum_j|y_j-b_j|\) over odd-parity
binary vectors \(b\) is
\[
 d_{\rm odd}=
 \begin{cases}
 d_0,&\sum_jb_j^0\text{ odd},\\
 d_0+\min_j|1-2y_j|,&\sum_jb_j^0\text{ even}.
 \end{cases}
 \tag{98.24}
\]
Indeed the coordinatewise nearest bits minimize the
unconstrained distance. If their parity is wrong, an
odd number of bits must be changed, and each change
has nonnegative cost \(|1-2y_j|\). A cheapest single
change is optimal. Thus (98.19) fails precisely when
\(d_{\rm odd}<1\), including the case of ties at \(1/2\).

Using integer numerators in (98.23)--(98.24), the checker
tests all 36288 periodic boxes with side lengths in
\(\{1,2\}^3\), excluding elementary cubes.
It finds six violations: four for lengths \((1,2,2)\)
and two for \((2,2,2)\), in the ordered coordinate triples.
The 15552 two-cube boxes have no violations, as predicted
by Proposition~\ref{prop:f12v94-tree}.
The accepted point and box checker is
\begin{center}\small\ttfamily
verify\_ym\_f12\_v94\_cube\_point.py.
\end{center}

\subsection{Direction after the new global bound}

The present action lower bound applies to all hidden
configurations of the specified parent; it is not
another local Hessian test. Nevertheless it leaves
an interval of more than 20 action units below the
known competitor and does not identify the global
ground set. The larger-box obstruction also shows
that an elementary-cube-feasible low-cost point need
not be a physical competitor.

The next compatibility refinement should use genuinely
new closed-boundary constraints, such as the violated
boxes with adjacency cycles, rather than adding the
redundant tree-union inequalities. The exact separation
rule (98.24) permits adding only violated inequalities
without sacrificing validity.
Even a fully optimized angle relaxation may still
omit shared-axis information and the actual pivot
equations. No sufficiency claim is made for the larger
boxes, and the nonlinear cosine cost itself must not
be replaced by an uncertified linear objective.

These results do not show that every conditional
ground state is joint-regular or has an unstable
joint direction.
V92's verified regular orbit remains one covered
component, not a complete cover. Unrestricted repair
and the cofinal analytic estimates still require
their own proofs.

Both new exact verifiers pass, including the minorant
remainders, all dual signs and compensated residuals,
all angle and cube inequalities for the rational point,
its true cost enclosure, the larger-box violation,
and the rectangle/cube free-word identities.
V93 had 1555613 bytes and SHA--256
\begin{center}\scriptsize\ttfamily
E83D53AA2390395344E4766BE585056CECC524FBD4211741357ADC295D7C57DA.
\end{center}
Its body is preserved apart from the title version.
Only active V93 is replaced by V94; all prior scripts,
frozen data, archives, and companion notes are retained.
The next iteration V95 includes the scheduled companion
audit. The V100 direction review and family-\(f12\)
rollover remain scheduled.

\begin{firewall}[Position after V94]
A rational dual certificate gives a global actual-parent
action lower bound of \(549.590337775\).
Elementary-cube consistency is proved insufficient by
an exact larger-box obstruction, and a class of redundant
refinements is identified. None of these finite-parent
statements proves complete ground-set coverage,
unrestricted repair, the cofinal auxiliary gap,
continuum construction, or the physical Yang--Mills
mass gap. The full objective remains active.
\end{firewall}
\section{V95: surface cuts and the convex-cost obstruction}
\label{sec:f12v95}

\subsection{The scheduled companion audit}

The V95 audit rechecked the active companion sources.
They are unchanged from the V90 audit:
\begin{itemize}
\item SM--Einstein V72, 607372 bytes, proves a fixed-time,
fixed-mode many-copy limit for its covariance flow.
It still separates that auxiliary limit from ultraviolet,
kinetic-time, and physical finite-species continuum control.
\item NS stability V26, 278283 bytes, records exact
pointwise rank-one norm formulas and leaves full
stability and global regularity open.
\item Parallel YM V5, 54174 bytes, retains the obstruction
to a volume-uniform tail for an extensive total charge.
Its fixed-window, density, and quasilocal alternatives
do not supply a cofinal gap estimate.
\end{itemize}
Their respective SHA--256 values are
\begin{center}\scriptsize\ttfamily
F44F9745D43379E1672C5550411B58E97195A4C9278592D221DF5D3F644DC1F5\\
82C887F8C2C69E4F28FAD7C3DC8DE5B9099CB5A4BF431EBA1FFF3E91935978AD\\
306F1BB84CFA8A8D47EA569EC17E9545F9867C8D32B548EC9D9C444B4F3C0512.
\end{center}
No new transferable global coercivity or continuum theorem
was found. The companion documents are mathematical sources,
not additional instructions to terminate or redirect this
task. Their retained limitations remain applicable.

The present iteration first incorporates the genuinely
new surface cuts identified in V94. More importantly,
it finds a separate obstruction in the cost model:
globally valid separable convex plaquette minorants have
a certified deficit even at the known physical orbit.
Consequently adding angle compatibility constraints alone
cannot make that cost model certify the orbit's exact energy.
The final part gives interval-restricted minorants that
avoid this particular obstruction, while making explicit
the complete-cover obligation of any global branching proof.

\subsection{Nine valid surface cuts and a refined global bound}

The exploratory separation loop considers the 20736
rectangular boxes with lengths in \(\{1,2\}^3\) having
at least two sides of length 2. It omits the tree-like
two-cube unions proved redundant in V94.
For each numerical angle vector it uses the nearest-odd
bit rule (98.23)--(98.24) to select violated inequalities.
Every selected inequality is valid independently of the
numerical accuracy of that selection, because its bit
vector has odd parity and its underlying rectangular
boundary is verified.

The numerical progression was
\[
\begin{array}{c|c|c|c}
 \text{round}&\text{added cuts}&\text{reported objective}
                         &\text{largest tested violation}\\ \hline
 0&0&549.5903424&0.0547476\\
 1&8&549.5988457&0.0017836\\
 2&9&549.5988457&2.3\cdot10^{-16}\text{ or less}.
\end{array}
 \tag{99.1}
\]
These are exploratory reports, not exact feasibility,
optimality, or exhaustive physical-compatibility certificates.
The accepted bound below does not assume that separation
has converged or that every possible surface has been tested.

\begin{theorem}[Refined all-hidden action lower bound]
\label{thm:f12v95-bound}
For every actual completed hidden SU(2) field at the
specified central exterior, coarse field \(I\), and
\(\rho=1/64\),
\[
                  \mathcal S(Y)\ge
                       \frac{549598841167}{10^9}
                       =549.598841167.
 \tag{99.2}
\]
This improves V94's certified lower bound by more than
\(1/125\), without restricting the hidden groups to
commuting configurations.
\end{theorem}
\begin{proof}
The file
\begin{center}\small\ttfamily
ym\_f12\_v95\_surface\_dual\_data.py
\end{center}
contains nine surface descriptors and an integer dual
vector with denominator \(10^8\).
For each descriptor, the verifier reconstructs the
periodic box from its base, coordinate triple, side
lengths, and forbidden bits. It checks their ranges,
the number of boundary faces, and odd bit parity.
It then substitutes the fixed face angles in the
valid boundary inequality of
Proposition~\ref{prop:f12v94-box}.
Duplicate resulting rows are excluded.
Thus all nine appended rows are valid for every
physical field, regardless of how they were selected.

The resulting model has 90761 non-box rows.
The checker verifies all dual signs, including 3647
nonzero entries, rebuilds V94's rational cost minorants,
and recomputes the residual and box compensation in
(98.10)--(98.14). There are 1026 negative residual
coordinates, all retained in the compensation.
The outward lower reports are
\[
\begin{array}{c|c}
 408+d^{\mathsf T}\lambda&551621567243/10^9\\
 \sum_i\min(r_i,0)U_i&-505681519/250000000\\
 \text{combined bound}&549598841167/10^9.
\end{array}
 \tag{99.3}
\]
The exact integer calculation, followed by downward
rounding, gives the last row. The dual inequality
(98.11) proves (99.2).
The difference from \(21983613511/40000000\) is
greater than \(1/125\).
\end{proof}

The frozen data file has SHA--256
\begin{center}\scriptsize\ttfamily
5738FE756CB822B6898A93F9ED3BF4AB0591F8502E5559524C5E9E972C86BDAB.
\end{center}
The reconstruction and accepted checker are
\begin{center}\small\ttfamily
ym\_f12\_v95\_surface\_cuts.py\\
verify\_ym\_f12\_v95\_surface\_bound.py.
\end{center}
Compression in the data file is lossless serialization.
All certificate arithmetic after decoding is integer
or rational arithmetic. The numerical separation code
is not imported by the accepted verifier.

\subsection{A physical obstruction to globally convex separable costs}

The small improvement in (99.2) must not lead to an
indefinite attempt to close the entire ground-energy
interval using only more angle constraints.
There is an intrinsic limitation in the globally
convex cost surrogate itself.

At the exact physical orbit representative \(Y_*\),
consider the three plaquettes
\[
 \begin{array}{c|c}
 \text{base}&\text{directions}\\ \hline
 (0,0,1,3)&(2,3)\\
 (0,3,1,0)&(1,2)\\
 (3,0,1,0)&(0,2).
 \end{array}
 \tag{99.4}
\]
They are distinct variable plaquettes of the actual
completed field, not variables chosen freely in a relaxation.

\begin{proposition}[Three certified physical angles]
\label{prop:f12v95-physical-angles}
Each plaquette in (99.4) has normalized angle
\[
                              \frac12<y_q<\frac35
 \tag{99.5}
\]
and physical cost \(g(y_q)>13/5\).
\end{proposition}
\begin{proof}
The verifier reconstructs the physical field over
V90's entire box containing \(\xi_*\), including the
four moving inverse-root pivots.
For each selected plaquette its scalar holonomy
component \(c_q\) satisfies
\[
                           -\frac{309}{1000}<c_q<-\frac3{10}.
 \tag{99.6}
\]
The outward reports, in the order (99.4), are
\[
\begin{array}{c|c}
 \text{lower endpoint}&\text{upper endpoint}\\ \hline
 -30802199/10^8&-308021989/10^9\\
 -30802111/10^8&-308021109/10^9\\
 -61599569/200000000&-76999461/250000000.
\end{array}
 \tag{99.7}
\]
The same rational \(\pi\) and cosine enclosures check
\(\cos(3\pi/5)<-309/1000\).
Since cosine is strictly decreasing on \([0,\pi]\),
(99.6) implies (99.5).
Also \(g(y_q)=2(1-c_q)>13/5\).
These checks apply to the exact critical point because
they hold throughout its certified containing box.
\end{proof}

\begin{theorem}[A ceiling for globally valid separable convex minorants]
\label{thm:f12v95-convex-ceiling}
For each variable plaquette let \(\psi_q:[0,1]\to\mathbb R\)
be convex and satisfy \(\psi_q(y)\le g(y)\) on the
whole interval. Different plaquettes may use different
minorants. At the physical angle profile of \(Y_*\),
\[
                408+\sum_q\psi_q(y_q(Y_*))
                              <E_*-\frac{31}{25}.
 \tag{99.8}
\]
Therefore minimizing this separable surrogate over any
angle feasible set containing that physical profile
has value strictly below \(E_*-31/25\).
No refinement consisting only of valid angle constraints
can make such an objective certify a lower bound \(E_*\).
\end{theorem}
\begin{proof}
For \(1/4\le y\le1\), convexity and the endpoint
minorant bounds give
\[
 \begin{split}
 \psi_q(y)
    &\le \frac{1-y}{3/4}\psi_q(1/4)
                      +\frac{y-1/4}{3/4}\psi_q(1)\\
    &\le \frac{1-y}{3/4}(2-\sqrt2)
                      +\frac{y-1/4}{3/4}\,4 .
 \end{split}
 \tag{99.9}
\]
The final affine expression is increasing in \(y\).
For each selected plaquette, (99.5) therefore yields
\[
           \psi_q(y_q)
                 <\frac{44-8\sqrt2}{15}
                 <\frac{164}{75},
 \tag{99.10}
\]
using \(\sqrt2>7/5\).
Its true cost exceeds \(13/5\), so the loss is greater
than \(13/5-164/75=31/75\).
Sum over the three distinct plaquettes and use
\(\psi_q\le g\) on every other plaquette to obtain (99.8).
Any feasible set containing the physical profile has
surrogate infimum no greater than its value there.
Valid angle constraints cannot exclude a physical field.
\end{proof}

This theorem includes the maximum of any family of
globally valid affine minorants, with or without the
additional lower bound zero used in the cost epigraph.
It is not a limitation on general coupled cost
certificates, on the true nonlinear objective, or on
minorants valid only inside explicitly restricted
branches. It does not prove that \(E_*\) is the actual
global minimum. It proves that this particular global
convex separable strategy cannot establish that equality,
even if all missing angle compatibility constraints
were supplied.

The accepted checker is
\begin{center}\small\ttfamily
verify\_ym\_f12\_v95\_convex\_cost\_barrier.py.
\end{center}
Its finite calculations establish (99.6)--(99.7) and
the rational inequalities used above. The statement
for all convex \(\psi_q\) follows from the proof,
not from testing a finite list of cost functions.

\subsection{Interval-restricted costs with a uniform error certificate}

The upstream change is to restrict the domain before
convexifying the cost. This removes the specific obstruction
of Theorem~\ref{thm:f12v95-convex-ceiling}, but creates a
coverage obligation: a local cost plane cannot be applied
outside its certified interval.

\begin{proposition}[Certified local affine minorant]
\label{prop:f12v95-interval-cost}
Let \(I=[l,u]\subset[0,1]\), \(h=u-l>0\), and suppose
rational endpoint enclosures are available:
\[
 g_l^-\le g(l)\le g_l^+,\qquad
 g_u^-\le g(u)\le g_u^+,\qquad
 \epsilon_I=\max(g_l^+-g_l^-,g_u^+-g_u^-).
 \tag{99.11}
\]
Set \(C=121/49\) and
\[
 \begin{split}
 L_I^-(y)&=\frac{u-y}{h}g_l^-+\frac{y-l}{h}g_u^-,\\
 \ell_I(y)&=L_I^-(y)-Ch^2.
 \end{split}
 \tag{99.12}
\]
Then, for every \(y\in I\),
\[
             0\le g(y)-\ell_I(y)
                         \le 2Ch^2+\epsilon_I.
 \tag{99.13}
\]
For a singleton \(I=\{l\}\), the constant minorant
\(g_l^-\) has error at most \(g_l^+-g_l^-\).
\end{proposition}
\begin{proof}
The true endpoint chord \(L_I\) satisfies the interpolation
remainder bound
\[
 |g(y)-L_I(y)|\le
 \frac{\|g''\|_{\infty,I}}2 (y-l)(u-y)
 \le \frac{2\pi^2}{8}h^2<Ch^2,
 \tag{99.14}
\]
because \(g''(y)=2\pi^2\cos(\pi y)\) and \(\pi<22/7\).
For completeness, at an interior point \(y\), subtract from
\(g-L_I\) a multiple of \((t-l)(t-u)\) chosen to vanish
at \(t=y\). The resulting function has three zeros.
Two applications of Rolle's theorem give a point where
its second derivative is zero, yielding the usual
remainder formula. At the endpoints the error is zero.
The maximum of \((y-l)(u-y)\) is \(h^2/4\).

Both chord weights are nonnegative and sum to one, so
\(0\le L_I(y)-L_I^-(y)\le\epsilon_I\).
Adding \(Ch^2\) to the decomposition
\(g-L_I^-=(g-L_I)+(L_I-L_I^-)\) proves both inequalities
in (99.13). The singleton statement is the endpoint
enclosure itself.
\end{proof}

The helper
\begin{center}\small\ttfamily
ym\_f12\_v95\_interval\_cost.py
\end{center}
implements (99.11)--(99.13) using rational arithmetic.
It rounds V93's Machin-formula enclosure of \(\pi\)
outward to denominator \(10^{30}\). For a rational endpoint
\(y\), let \([a,b]\) enclose \(\pi y\), and put
\(m=(a+b)/2\). Its degree-40 cosine polynomial \(P_{40}\)
is evaluated exactly as a rational number. Since
\(0\le m<22/7\), the alternating cosine tail after
degree 40 is bounded by \(m^{42}/42!\).
The terms in this tail decrease in magnitude, as their
successive ratio is at most
\((22/7)^2/(43\cdot44)<1\).
Since cosine is 1-Lipschitz, the total error bound is
\[
 \left|\cos(\pi y)-P_{40}(m)\right|
           \le \frac{m^{42}}{42!}+\frac{b-a}{2}.
 \tag{99.15}
\]
The helper turns this into an enclosure of \(g(y)\),
rounds outward to denominator \(10^{20}\), and uses
the exact values \(g(0)=0\), \(g(1)=4\) at the endpoints.
No floating-point value is used for an accepted endpoint.
Higher Taylor order and finer rational reporting can
be used if a smaller endpoint error is required.

\subsection{What a complete branch calculation would certify}

\begin{proposition}[Coverage and objective-error accounting]
\label{prop:f12v95-branch-cover}
Let \(\mathcal R\) be a nonempty compact angle relaxation
containing every actual physical angle profile at the
fixed exterior and coarse input. Write
\[
 F(y)=408+\sum_{q=1}^{2400}g(y_q).
\]
Choose a finite family of closed product intervals
\(B_\alpha=\prod_q I_{q,\alpha}\) covering the plaquette
projection of \(\mathcal R\). Retain any auxiliary link
angles when intersecting with these intervals.
For each nonempty branch
\(\mathcal R_\alpha=\mathcal R\cap\{y\in B_\alpha\}\),
define
\[
 \begin{split}
 L_\alpha(y)&=408+\sum_q\ell_{I_{q,\alpha}}(y_q),\\
 D_\alpha&=\sum_q
       \bigl(2C|I_{q,\alpha}|^2+\epsilon_{q,\alpha}\bigr),\\
 b_\alpha&=\min_{\mathcal R_\alpha}L_\alpha,\qquad
 b=\min_\alpha b_\alpha,\qquad D=\max_\alpha D_\alpha .
 \end{split}
 \tag{99.16}
\]
Then
\[
 b\le \min_{\mathcal R}F\le b+D,\qquad
                     \min_Y\mathcal S(Y)\ge b.
 \tag{99.17}
\]
More generally, if only certified branch lower bounds
\(B_\alpha\le b_\alpha\) are available, the physical
lower bound remains valid with \(\min_\alpha B_\alpha\).
The upper error bound in (99.17) does not follow from
uncalibrated lower bounds alone.
\end{proposition}
\begin{proof}
On each nonempty branch, (99.13) gives
\(L_\alpha\le F\le L_\alpha+D_\alpha\).
Coverage implies that the minimum of \(F\) over
\(\mathcal R\) is the minimum of its branch minima,
which proves its lower bound by \(b\).
Choose a branch attaining \(b\) and a minimizer of
its continuous affine function on its compact domain.
At that point \(F\le b+D_\alpha\le b+D\), proving
the upper bound. Every physical profile belongs to
\(\mathcal R\), and its action equals \(F\), which
proves the physical lower bound. Replacing exact
branch minima by smaller certified values preserves
only that lower-bound argument.
\end{proof}

For example, intervals of width at most \(h\) and
endpoint errors at most \(\epsilon\) give
\[
          D\le \frac{580800}{49}h^2+2400\epsilon.
 \tag{99.18}
\]
At \(h=1/128\), the curvature term is
\(580800/802816<3/4\).
This is a uniform estimate of the objective
approximation error, not the cost of performing the
branch calculation. The naive product cover already
has \(128^{2400}\) boxes. No such enumeration or
certified pruning tree has been carried out.
Empty branches may be discarded only with a valid
infeasibility certificate. A branch restricted to
the vicinity of the known orbit is not a global cover.
One can instead cover a rigorously defined relevant
sublevel set, but then the excluded region needs
its own proved lower bound.

With endpoint errors also tending to zero, shrinking
intervals makes the objective approach the minimum
of the \emph{nonlinear angle relaxation}.
It does not prove that a relaxed profile is
realizable by a physical SU(2) link field.
Even exact branch optimization and vanishing errors
in (99.18) leave a possible compatibility gap between
\(\min_{\mathcal R}F\) and the physical minimum.
Thus this result repairs the convex-cost obstruction
without concealing the separate nonabelian realization
problem.

\subsection{Validation, direction, and continuation record}

The implementation checker
\begin{center}\small\ttfamily
verify\_ym\_f12\_v95\_interval\_cost.py
\end{center}
tests 20 nonsingleton intervals, nine rational sample
points in each interval, and three singleton intervals.
It checks the returned rational plane and reported
error against cost enclosures evaluated at those
sample points using the same rational cost helper.
These 180 sample checks are implementation diagnostics;
the all-points guarantee is
Proposition~\ref{prop:f12v95-interval-cost}, not finite
sampling. Together with the surface-dual and convex-cost
checkers, they separate exact accepted arithmetic from
the exploratory numerical cut search. They are not
machine-checked formalizations of all preceding analysis.

The next attack should therefore address interval
coverage and physical realizability jointly, rather
than seek arbitrarily small gains from additional
surface cuts with the unchanged global convex objective.
Useful next targets include a branch-level physical
infeasibility certificate or a coupled lower bound
that controls both the nonlinear plaquette cost and
shared link variables. Neither has been established
here. The exact interval helper is a component for
such work, not evidence that a complete branch proof
will be computationally feasible.

V95 preserves the complete V94 mathematical body,
apart from updating the cover-page version. The
predecessor had 1574318 bytes and SHA--256
\begin{center}\scriptsize\ttfamily
1A5FB686C33034EDA2572E7D45BA1620FDEB4FC0CBB7B8B178537CEAC9E5B21A.
\end{center}
Only V95 is retained as the active YM note.
Frozen witnesses and validation scripts are retained
for reproducibility. The five-version companion audit
is recorded at the start of this section; its next
scheduled audit and the next ten-version direction
review are V100. At that boundary the complete note
is to be preserved as \(\mathrm{V100\_f12}\), followed
by a new compact V1 with a checked summary.

\begin{firewall}
The new lower bound is an all-hidden action estimate
for one specified finite SU(2) parent problem. It is
not a spectral-gap estimate. The physical orbit remains
a certified local minimum, not a certified global one.
The interval-branch proposition does not supply a
completed cover calculation, a scale-uniform repair
inequality, a four-dimensional continuum construction,
or a physical Yang--Mills mass gap. The full goal
remains open in this programme.
\end{firewall}

\section{V96: shared-link Gram inequalities and a certified excluded branch}
\label{sec:f12v96}

\subsection{Continuation from the convex-cost obstruction}

The preceding turn was progress: V95 supplied a stronger
finite-parent lower bound, a proved ceiling for globally
convex separable plaquette costs, and a quantitative
interval-cost construction. The next unresolved issue is
not merely the approximation of the cosine. Scalar angles
must be realized by the \emph{same} physical link
quaternions wherever those links occur.

This version finds an explicit rational profile which
passes all previously tested local, cube, and rectangular
surface inequalities but fails such a shared-link condition.
The obstruction is a negative value assigned to a squared
quaternion norm. We then combine this condition with V95's
two-sided local cost enclosures to certify the exclusion
of a whole six-angle branch. This is the branch-level
physical infeasibility target left open at the end of V95,
not a complete branch-and-bound solution.

All statements remain at the specified central exterior,
coarse field \(I\), period-six fine lattice, and
\(\rho=1/64\). Physical pivots are still the actual
completed pivots. Nothing below freezes them or assumes
that the hidden links commute.

\subsection{Quaternion Gram constraints retain shared-link information}

Identify an SU(2) matrix with its unit quaternion
\(q=s+v\in\mathbb R\oplus\mathbb R^3\). Throughout this
section, \(|q|\) is the Euclidean quaternion norm, not
the Hilbert--Schmidt matrix norm. For unit quaternions,
\[
 \langle q,r\rangle=\operatorname{Re}(q\overline r),
 \qquad |q-r|^2=2\bigl(1-\operatorname{Re}(q\overline r)\bigr).
 \tag{100.1}
\]
Thus \(g(y)=2(1-\cos(\pi y))\) is exactly a squared
chord length when \(y\) is the relative normalized angle.

\begin{proposition}[Shared quaternion Gram test]
\label{prop:f12v96-gram}
For any finite family of physical quaternion words
\(q_0,\ldots,q_{m-1}\), its real Gram matrix
\(G_{ij}=\operatorname{Re}(q_i\overline q_j)\) is positive
semidefinite and has rank at most four.
For any real weights with \(\sum_i a_i=0\),
\[
  -\sum_{i<j}a_i a_j |q_i-q_j|^2
                         =\left|\sum_i a_iq_i\right|^2\ge0.
 \tag{100.2}
\]
The words may contain noncommuting physical links.
\end{proposition}
\begin{proof}
The columns \(q_i\) form a real \(4\)-by-\(m\) matrix,
whose transpose times itself is \(G\). This proves
positivity and the rank bound. Expanding the left
side of (100.2), the coefficients multiplying
\(|q_i|^2\) use \(\sum_i a_i=0\), and the remaining
cross terms give the displayed squared norm.
No commutation or independence assumption is involved.
\end{proof}

The initial exploration used all 24 monotone paths
across an elementary cube in four dimensions.
Swapping two consecutive steps relates two path
transporters by a plaquette, so it prescribes an entry
of their Gram matrix. A second exploration added entries
obtained by reducing fixed central links and matching
cyclic traces of physical words. These numerical matrix
completion searches did not themselves produce an
accepted infeasibility certificate. Their residuals and
iteration limits must not be reported as proofs.

Searching instead among the imaginary parts of shared
physical links yielded the three-link certificate below.
Only its exact word identities and rational inequalities
are used in the accepted result. The rank-four
restriction is not needed for this obstruction:
positive semidefiniteness alone already fails.

\subsection{Three physical links and three exact plaquette words}

In the existing zero-based physical-link enumeration,
let \(A=U_{35}\), \(B=U_{36}\), and \(C=U_{203}\).
Their actual locations are
\[
 \begin{array}{c|c|c}
  \text{index}&\text{base}&\text{direction}\\ \hline
  35&(0,0,3,3)&2\\
  36&(0,0,3,3)&3\\
  203&(0,5,3,3)&1 .
 \end{array}
 \tag{100.3}
\]
Write their normalized angles as \(x_A,x_B,x_C\).
The three plaquette angles used here have the following
locations, indices in the 3440-angle model, and reduced
physical words:
\[
 \begin{array}{c|c|c|c}
  \text{base}&\text{directions}&\text{angle index}&\text{word}\\ \hline
  (0,0,3,3)&(2,3)&1186&AB^{-1}\\
  (0,5,3,3)&(1,2)&1187&CA\\
  (0,5,3,3)&(1,3)&1190&CB .
 \end{array}
 \tag{100.4}
\]
Denote these angles by \(z_1,z_2,z_3\).
These are identities at the specified fixed exterior,
not identities imposed on freely chosen relaxation
variables. They follow by substituting the fixed
central links in the chronological plaquette words;
the remaining central sign is \(+1\) in all three cases.
The exact checker reconstructs the words from the
lattice and verifies this reduction.

For \(q_0=1,q_1=A,q_2=B,q_3=C^{-1}\), the six squared
chord lengths are consequently
\[
 \begin{array}{c|cccccc}
 (i,j)&(0,1)&(0,2)&(0,3)&(1,2)&(1,3)&(2,3)\\ \hline
 |q_i-q_j|^2&
 g(x_A)&g(x_B)&g(x_C)&g(z_1)&g(z_2)&g(z_3).
 \end{array}
 \tag{100.5}
\]
For example,
\(\operatorname{Re}(A\overline{C^{-1}})
=\operatorname{Re}(AC)=\operatorname{Re}(CA)\).
Only cyclic invariance of the scalar trace is used.

\begin{theorem}[An exact coupled link--plaquette cost inequality]
\label{thm:f12v96-cost-cut}
Every physical field in the fixed parent problem satisfies
\[
 \begin{split}
 K={}&5g(x_A)-2g(x_B)-2g(x_C)\\
    &+10g(z_1)+10g(z_2)-4g(z_3)\ge0.
 \end{split}
 \tag{100.6}
\]
More precisely,
\[
                  K=|-1+5A-2B-2C^{-1}|^2.
 \tag{100.7}
\]
\end{theorem}
\begin{proof}
Use weights \((-1,5,-2,-2)\), whose sum is zero,
in (100.2), and insert the six distances in (100.5).
The coefficients \(-a_i a_j\), in that order, are
\(5,-2,-2,10,10,-4\). This proves both statements.
\end{proof}

This is a globally valid \emph{coupled} constraint on
true link and plaquette costs. It is not a globally
convex separable minorant of the action, so V95's
convex-cost ceiling does not apply to it. It is also
not an extra energy term: the individual link costs
in (100.6) are auxiliary observables, not additions
to the Yang--Mills action.

There is a slightly stronger nonlinear form retaining
the scalar component of the quaternion sum. Put
\(s_A=\cos(\pi x_A)\), and similarly for \(s_B,s_C\).
Its imaginary component gives
\[
 \begin{split}
 Q={}&|-5v_A+2v_B-2v_C|^2\\
   ={}&33-(5s_A-2s_B-2s_C)^2\\
      &-20\cos(\pi z_1)-20\cos(\pi z_2)+8\cos(\pi z_3)
      \ge0,\\
 K={}&Q+(5s_A-2s_B-2s_C-1)^2 .
 \end{split}
 \tag{100.8}
\]
To check signs, (100.4) gives
\[
 \begin{split}
 v_A\cdot v_B&=\cos(\pi z_1)-s_As_B,\\
 v_A\cdot v_C&=s_As_C-\cos(\pi z_2),\\
 v_B\cdot v_C&=s_Bs_C-\cos(\pi z_3).
 \end{split}
 \tag{100.9}
\]
Expanding the squared norm proves (100.8).
The simpler linear true-cost inequality (100.6)
already suffices for an exact exclusion.

\subsection{A surface-consistent rational profile which is not physical}

\begin{theorem}[Exact separation beyond the tested scalar surfaces]
\label{thm:f12v96-profile}
The frozen rational profile in
\begin{center}\small\ttfamily
ym\_f12\_v96\_angle\_point\_data.py
\end{center}
has 3440 coordinates with denominator \(10^{10}\).
It satisfies all local and elementary-cube angle rows,
the nine V95 surface cuts, and all odd-subset boundary
inequalities for the 36288 rectangular boxes tested
in V94. Nevertheless, it cannot be the angle profile
of any physical field in the fixed parent problem.
\end{theorem}
\begin{proof}
The verifier decodes all integers and checks the
coordinate boxes, including
\(40192n_i\le125\cdot10^{10}\) on all 96 pivot angles.
It reconstructs 49961 scalar rows: the 49952 local
and cube rows and the nine V95 surface rows.
Every integer row defect is nonpositive.

For each of the 36288 boxes, the base ranges over
the period-six lattice, the coordinate triple is
arbitrary, and each side length is 1 or 2, with
\((1,1,1)\) omitted because elementary cubes have
already been included. If \(n_j\) are the boundary
angle numerators and \(D=10^{10}\), the checker
computes
\[
 d_0=\sum_j\min(n_j,D-n_j),\qquad
 d_{\mathrm{odd}}=
 \begin{cases}
 d_0,&\#\{j:2n_j>D\}\text{ is odd},\\
 d_0+\min_j|D-2n_j|,&\text{otherwise}.
 \end{cases}
 \tag{100.10}
\]
V94's exact nearest-odd-vertex argument shows that
all odd-subset inequalities hold if and only if
\(d_{\mathrm{odd}}\ge D\).
The largest integer defect \(D-d_{\mathrm{odd}}\)
is zero. This proves all of the stated scalar tests;
it does not assert tests of arbitrarily large surfaces.

The six selected numerators are
\[
 \begin{array}{c|rrrrrr}
 \text{coordinate}&x_A&x_B&x_C&z_1&z_2&z_3\\ \hline
 10^{10}y&
 270163139&348194339&348194339&
 83031200&83031200&129992632 .
 \end{array}
 \tag{100.11}
\]
Using V95's exact rational cosine-cost enclosures
and the signed sum (100.6), the checker obtains
\[
 -\frac{3813629}{781250000}
       \le K\le
 -\frac{4881445119}{10^{12}}<-\frac{48}{10000}.
 \tag{100.12}
\]
This contradicts Theorem~\ref{thm:f12v96-cost-cut}
for any physical realization.
\end{proof}

For completeness, the same checker also verifies the
16 odd-parity triangle rows on the four triples of
pairwise distances among \(1,A,B,C^{-1}\):
\[
 (x_A,x_B,z_1),\quad(x_A,x_C,z_2),\quad
 (x_B,x_C,z_3),\quad(z_1,z_2,z_3).
 \tag{100.13}
\]
They all hold. Thus even these four-point triangle
tests do not imply the shared quaternion Gram condition.
The imaginary-part test gives the additional enclosure
\[
 -\frac{4916367727}{10^{12}}
       \le Q\le-\frac{2458183863}{500000000000}
       <-\frac{49}{10000}.
 \tag{100.14}
\]
These are exact interval calculations, not negative
eigenvalues rounded from a numerical matrix.

The profile's \emph{relaxed true-cosine objective} lies in
\[
 \left[
 \frac{70586706231271}{125000000000},
 \frac{564693649850169}{10^{12}}
 \right]\subset(564,565).
 \tag{100.15}
\]
Since the profile is not physical, (100.15) is not
the energy of a competitor and is not an upper bound
on the physical ground energy. It displays a false
low-energy profile admitted by the stated scalar tests.

The frozen data file has SHA--256
\begin{center}\scriptsize\ttfamily
DC97C3EDD1074531DD05FE1454F02CFEEC32F5C67EB1607BCF99C8F88F05498E.
\end{center}
Its point was proposed by shrinking a numerical V95
relaxation solution by \(10^{-3}\) toward a central
barycentre and rounding to denominator \(10^{10}\).
Neither the optimizer nor that proposal rule proves
feasibility: every accepted claim above is rechecked
on the decoded rational point.

\subsection{An affine branch row with the correct signs}

V95 supplied lower affine cost enclosures. A Gram
exclusion with mixed positive and negative coefficients
requires an \emph{upper} enclosure of the entire signed
expression, not a lower enclosure substituted indiscriminately.

\begin{proposition}[Interval conversion of a coupled Gram constraint]
\label{prop:f12v96-affine-cut}
Let \(K(y)=\sum_i k_i g(y_i)\ge0\) be a physical
constraint. On a specified product interval branch,
suppose
\(\ell_i(y_i)\le g(y_i)\le\ell_i(y_i)+\delta_i\),
with \(\ell_i\) affine and \(\delta_i\ge0\).
Define
\[
 U(y)=\sum_i k_i\ell_i(y_i)
                          +\sum_{k_i>0}k_i\delta_i.
 \tag{100.16}
\]
Then throughout this branch,
\[
     0\le U(y)-K(y)\le\sum_i|k_i|\delta_i.
 \tag{100.17}
\]
Every physical point in the branch satisfies \(U(y)\ge0\).
Consequently, \(\max_{\mathrm{branch}}U<0\) is a
certificate that the entire branch is physically empty.
\end{proposition}
\begin{proof}
If \(k_i>0\), replace \(g\) by its upper enclosure;
the excess is in \([0,k_i\delta_i]\).
If \(k_i<0\), multiplication reverses order, so
replace \(g\) by its lower enclosure; the excess
is in \([0,|k_i|\delta_i]\).
Summing proves (100.17). The physical necessity
and the exclusion follow from \(K\ge0\).
\end{proof}

For the six coefficients in (100.6),
\(\sum_i|k_i|=33\). If all branch widths are at most
\(h\) and endpoint cost errors at most \(\epsilon\),
V95 gives
\[
        0\le U-K\le
          33\left(\frac{242}{49}h^2+\epsilon\right).
 \tag{100.18}
\]
Writing \(U(y)=b+\sum_i a_i y_i\), its exact maximum
over a box \(\prod_i[l_i,u_i]\) is
\[
 b+\sum_{a_i\ge0}a_i u_i+\sum_{a_i<0}a_i l_i .
 \tag{100.19}
\]
Thus the branch check requires only rational arithmetic
after constructing the endpoint enclosures.
The row \(-\sum_i a_i y_i\le b\) is a necessary
affine constraint \emph{on this branch}. It is not
licensed outside the intervals used to build it.

\begin{theorem}[A certified six-angle branch is physically empty]
\label{thm:f12v96-empty-branch}
Let \(y^0\) be the six coordinates in (100.11).
There is no physical field in the fixed parent problem
whose six corresponding angles all satisfy
\[
               |y_i-y_i^0|\le\frac1{2500}.
 \tag{100.20}
\]
All other hidden coordinates may vary arbitrarily.
\end{theorem}
\begin{proof}
Apply Proposition~\ref{prop:f12v96-affine-cut} with
V95's rational affine cost bounds on the six intervals
in (100.20). Each lies in \([0,1]\).
The exact checker reports the outward upper bounds
\[
 \begin{array}{c|c}
 \max_{\mathrm{branch}} U&
              -183639099/200000000000\\
 \sum_i|k_i|\delta_i&
               104306939/10^{12}.
 \end{array}
 \tag{100.21}
\]
In particular the first is less than \(-1/2000\).
It is obtained directly from (100.19); no optimizer
is involved. Every physical point in this branch
would satisfy \(0\le K\le U<-1/2000\), a contradiction.
The identity (100.7) depends only on these six
observables, so arbitrary choices of the remaining
coordinates cannot repair the contradiction.
\end{proof}

There is also a coarser independent continuity check.
Since \(|g'|\le2\pi\), (100.6) has sup-norm Lipschitz
constant at most \(66\pi<208\) in these six angles.
Combining this with (100.12) excludes radius \(1/50000\).
The affine interval calculation proves the twenty-times
larger radius in (100.20) and supplies an explicit
branch row rather than just a continuity estimate.

\subsection{Validation and the next upstream task}

The accepted verifier is
\begin{center}\small\ttfamily
verify\_ym\_f12\_v96\_link\_gram\_obstruction.py.
\end{center}
Its reusable helpers are
\begin{center}\small\ttfamily
ym\_f12\_v96\_path\_geometry.py\\
ym\_f12\_v96\_gram\_branch\_cut.py.
\end{center}
The verifier imports neither a floating-point
eigensolver nor an optimization routine. It reconstructs
the physical word reductions, the scalar inequalities,
the exact interval costs, and the affine box maximum.
Twelve rational noncommuting quaternion fixtures also
check the signs and coefficients in (100.7).
The general validity of (100.6) and (100.16) is
proved above; executable checks evaluate the stated
finite witnesses. This is not a formal proof-assistant
certification of the full programme.

The numerical path and link searches are retained
separately as reproducible exploration. A positive
completion residual is not used as evidence of
infeasibility. The accepted squared-norm witness has
small integer coefficients and does not depend on
any numerical eigenvector.

This changes the next action concretely. Shared-link
Gram inequalities can now be converted into certified
local affine cuts and branch deletions. The next task
is to incorporate such cuts into a complete, adaptively
refined cover, retaining every unresolved branch and
every actual pivot constraint. The present six-angle
box is only one certified deletion. Nothing here
shows that the remaining branches have energies above
the known orbit, nor that all necessary Gram conditions
would be sufficient for a global link-field realization.

V96 preserves the entire V95 mathematical body, apart
from updating the title version. V95 had 1592713 bytes
and SHA--256
\begin{center}\scriptsize\ttfamily
B264AAEBAF2E70D8C980A233B0592003D40184A31CDF58E6EB2A2486DE8EC398.
\end{center}
Only V96 is retained as the active note. Frozen data
and verification scripts are preserved. The companion
audit was completed at V95; the next audit and
ten-version direction review remain V100, followed
by preservation as \(\mathrm{V100\_f12}\) and a new
compact V1.

\begin{firewall}
The new achievement is a proved shared-link
compatibility obstruction and a certified physically
empty branch. It does not raise the global action
lower bound beyond V95's \(549.598841167\), identify
all finite-parent global minimizers, establish a
scale-uniform spectral or repair inequality, construct
the four-dimensional continuum theory, or prove
the physical Yang--Mills mass gap. Those remain
unproved obligations of the active programme.
\end{firewall}

\section{V97: coupled costs, pivot support, and exact two-link realization}
\label{sec:f12v97}

\subsection{What is being advanced}

V96 was a progress turn: it proved a shared-link
obstruction beyond the tested scalar surface constraints,
and certified an empty six-angle branch. The next test
is whether shared-link constraints can enter a global
lower-bound certificate rather than only delete a
neighbourhood of one false profile.

This version answers that limited question positively:
the new certified all-hidden action lower bound is
\(549.600501311\). The gain over V95 is small,
slightly more than \(0.00166\), and is not represented
as evidence that a mass-gap proof is close.
The exact dual support is more informative than the
size of the gain: all 13 supported Gram rows are
two-link rows, and 11 involve a physical pivot.
We also give an exact realization criterion for the
one- and two-link trace subsystem. Its remaining
three-link, four-link, and pivot-completion obligations
are kept explicit.

The finite parent, fixed central exterior, coarse
field \(I\), and \(\rho=1/64\) are unchanged.
The actual physical pivot values are not replaced by
identity values.

\subsection{Two-sided true-cost lifts}

The angle variables remain the 1040 physical link
angles and 2400 variable plaquette angles. Introduce
a cost variable for every one of these angles.
The 2400 plaquette costs already occurred in V94;
the 1040 link costs are new auxiliary variables.
Only plaquette costs enter the action objective.
Every physical field has the distinguished embedding
\[
              t_i=g(y_i)=2(1-\cos(\pi y_i)).
 \tag{101.1}
\]
The relaxed constraints do not require all feasible
lifted points to satisfy this equality.

Let \(P=10^9\), and let \(a_j/P,b_j/P\) be V94's
17 rational lower tangent coefficients. For every
link or plaquette angle we impose
\[
 \begin{split}
       t&\ge (a_j y+b_j)/P,\\
       t&\le 4-(a_j(1-y)+b_j)/P .
 \end{split}
 \tag{101.2}
\]
The first inequality is V94's global minorant;
the second follows from \(g(y)=4-g(1-y)\).
Thus both are valid for (101.1).
The nonnegative cost box also remains in force.

Put
\[
 L=\frac{\lceil 44P/7\rceil}{P}>\,2\pi .
\]
Since \(|g'|\le2\pi\), the endpoint values \(g(0)=0\),
\(g(1)=4\) give the additional valid inequalities
\[
                  4-L(1-y)\le t\le Ly .
 \tag{101.3}
\]
No numerical differentiation is used to justify them.

For a plaquette whose reduced physical word is a
single link or its inverse with central sign \(s\),
the true costs obey
\[
                   t_{\mathrm{face}}-s\,t_{\mathrm{link}}
                                      =2(1-s).
 \tag{101.4}
\]
The code reconstructs all 176 such relations and
imposes each equality as two inequalities.

Write \(u=125/40192\). All actual pivot angles are
at most \(u\), by V93. Since
\(2(1-\cos t)\le t^2\) and \(\pi^2<10\),
their auxiliary costs have the additional valid bound
\[
                0\le t_{\mathrm{pivot}}\le 10u^2 .
 \tag{101.5}
\]
The elementary inequality used here follows, for
example, from \(1-\cos t=2\sin^2(t/2)\) and
\(|\sin r|\le|r|\). This bound retains small,
potentially nonzero physical pivot costs.

Together with V95's scalar rows, these constraints
give 174153 non-box rows on 6880 variables. The
upper boxes are
\[
 \begin{array}{c|r|c}
 \text{variable type}&\text{number}&\text{upper bound}\\ \hline
 \text{free-link angle}&944&1\\
 \text{pivot angle}&96&u\\
 \text{plaquette angle}&2400&1\\
 \text{plaquette cost}&2400&4\\
 \text{free-link cost}&944&4\\
 \text{pivot cost}&96&10u^2 .
 \end{array}
 \tag{101.6}
\]
The 174153 rows consist of 90761 old scalar and
plaquette-minorant rows, 17680 new link-minorant rows,
58480 reflected upper rows, 6880 endpoint rows, and
352 one-link equality rows.

\subsection{Gram rows as globally valid coupled-cost constraints}

Write a two-link plaquette trace as
\[
 \cos(\pi z_{ij})=
       s_{ij}\operatorname{Re}(U_i^\alpha U_j^\beta),
 \qquad s_{ij}\in\{-1,1\},\quad
 o_{ij}=\alpha\beta\in\{-1,1\}.
 \tag{101.7}
\]
The exact word reductions determine \(s_{ij}\) and
\(o_{ij}\). There are 976 distinct such pairs.
Choose signs \(\sigma_i,\sigma_j\) satisfying
\(\sigma_i\sigma_j=-o_{ij}\). For
\(q_i=U_i^{\sigma_i}\), \(q_j=U_j^{\sigma_j}\),
\[
 \operatorname{Re}(q_i\overline q_j)
                          =s_{ij}\cos(\pi z_{ij}).
 \tag{101.8}
\]
Indeed inversion preserves a quaternion's scalar part
and reverses its imaginary part, so the sign condition
matches the two scalar-product formulas.

For a triangle of such pairs, a common choice of
three signs exists if and only if
\[
                    o_{ij}o_{ik}o_{jk}=-1.
 \tag{101.9}
\]
Necessity follows by multiplying the three sign
equations. For sufficiency set
\(\sigma_i=1,\sigma_j=-o_{ij},\sigma_k=-o_{ik}\);
then (101.9) gives the third equation.
All 128 triangles in the actual pair graph pass
this exact sign test.

Each pair therefore defines a Gram block on
\((1,q_i,q_j)\), and each of these triangles defines
one on \((1,q_i,q_j,q_k)\).
There are 1104 blocks in all. In true cost variables,
their entries are
\[
 G_{00}=G_{aa}=1,\quad
 G_{0a}=1-t_{i_a}/2,\quad
 G_{ab}=s_{i_ai_b}(1-t_{i_ai_b}/2).
 \tag{101.10}
\]
Here the edge cost is the corresponding plaquette
cost. It is not a new independent plaquette observable.

For any integer vector \(w=(w_0,\ldots,w_r)\), physical
positivity of \(w^{\mathsf T}Gw\) gives a linear cost row.
Writing \(D_w=10^4\) for the numerical proposal's
weight denominator, its exact normalized form is
\[
 \begin{split}
 &\sum_a \frac{w_0w_a}{D_w^2}t_{i_a}
  +\sum_{a<b}\frac{s_{i_ai_b}w_aw_b}{D_w^2}t_{i_ai_b}\\
 &\qquad\le
 \frac{\sum_{a=0}^r w_a^2+
       2w_0\sum_{a=1}^r w_a+
       2\sum_{1\le a<b\le r}s_{i_ai_b}w_aw_b}{D_w^2}.
 \end{split}
 \tag{101.11}
\]
This follows by expanding (101.10); the slack is
exactly \(w^{\mathsf T}Gw/D_w^2\).
Thus any rounded integer weight vector gives a valid
row. It need not be an exact eigenvector.
Since \(D_w^2\) divides \(P\), every such row has
integer coefficients and right side after multiplication
by \(P\), like the other accepted rows.

\subsection{The verified global action bound}

A bounded eight-round numerical search added at most
150 Gram rows per round. Selected numerical reports are
\[
 \begin{array}{r|c}
 \text{Gram rows present}&\text{reported LP objective}\\ \hline
 0&549.598844842\\
 600&549.598844841\\
 750&549.599516619\\
 900&549.600172704\\
 1050&549.600503993 .
 \end{array}
 \tag{101.12}
\]
Small last-digit fluctuations are numerical, not
violations of monotonicity for exact nested relaxations.
The final search still found many negative Gram
directions; it did not solve all 1104 semidefinite
constraints or establish convergence of separation.
None of those numerical claims is needed below.

\begin{theorem}[Coupled all-hidden lower bound]
\label{thm:f12v97-bound}
Every actual physical field in the fixed finite parent
problem satisfies
\[
                      \mathcal S(Y)\ge
             \frac{549600501311}{10^9}
                      =549.600501311.
 \tag{101.13}
\]
The increase over V95's certified bound is greater
than \(1/1000\).
\end{theorem}
\begin{proof}
Every physical field embeds by (101.1) in the boxes
(101.6), satisfies the old scalar rows, (101.2)--(101.5),
and every accepted row (101.11). Let these 175203
rows be \(Ax\le d\), with objective \(408+c^{\mathsf T}x\).
The 2400 plaquette multiplicities are all one, as
checked directly from the geometry.

The frozen dual consists of integers \(v_j\le0\)
with denominator \(D=10^8\).
Let \(\widetilde A=PA\), \(\widetilde d=Pd\), and set
\[
 R_i=PDc_i-\sum_j\widetilde A_{ji}v_j,\qquad
 B=40192^2,\qquad U_i=\widetilde U_i/B.
 \tag{101.14}
\]
All quantities with tildes, and every \(R_i\), are
integers. Multiplying \(Ax\le d\) by \(v/D\le0\),
and bounding each residual coordinate on its box,
gives the exact lower bound
\[
 408+\frac{\sum_j\widetilde d_jv_j}{PD}
       +\frac{\sum_i\min(R_i,0)\widetilde U_i}{PDB}.
 \tag{101.15}
\]
This is the same residual-compensated dual argument
as V94, now on the coupled model. Negative residuals
are not discarded.

The verifier reconstructs all 175203 rows, checks
every dual sign, and evaluates (101.15) with arbitrary
precision integers. There are 3666 nonzero dual entries
and 825 negative residual coordinates.
Independent outward reports are
\[
 \begin{array}{c|c}
 \text{dual part including }408&
                   275810795057/500000000\\
 \text{box compensation}&-505272201/250000000\\
 \text{combined exact expression, rounded down}&
                   549600501311/10^9 .
 \end{array}
 \tag{101.16}
\]
The last row is computed before rounding the two
parts separately; it is not obtained by adding
their independently rounded reports.
It proves (101.13). Subtracting
\(549598841167/10^9\) gives an increase of
\(1660144/10^9>1/1000\).
\end{proof}

The frozen files and their SHA--256 hashes are
\begin{center}\small\ttfamily
ym\_f12\_v97\_dual\_data.py\\
\scriptsize
8A4003B0E7AA8F27B2B224ED056E551548F71120CA4C8994718CB54A7E656591\\[3pt]
\small
ym\_f12\_v97\_gram\_cut\_data.py\\
\scriptsize
877D8C83581D1D1E498C888930048393040DD537A35926CCC0D00EBA04225C14.
\end{center}
The accepted checker is
\begin{center}\small\ttfamily
verify\_ym\_f12\_v97\_coupled\_bound.py.
\end{center}
Besides the dual calculation, it tests all 1050 row
identities on a deterministic rational noncommuting
quaternion fixture. Those tests check implementation
signs; the universal validity is the expansion
(101.11), not the fixture.

\subsection{Where the improvement actually comes from}

Only 13 of the 1050 Gram rows have nonzero weights
in the exact dual. All 13 are pair blocks on three
quaternions, not triangle blocks on four quaternions.
Eleven involve an index from 944 through 1039 and
therefore involve an actual physical pivot.
The other two involve pairs \((484,495)\) and
\((489,502)\) of free links.
These support statements are checked on the frozen
dual. They prevent attributing the improvement
to triple-link geometry that did not support this
particular certificate.

There is an exact way to understand the pair constraint.
Let \(h,k\in[0,4]\) be the two link costs. For central
sign \(s\), let \(e'=e\) if \(s=1\) and \(e'=4-e\)
if \(s=-1\), where \(e\) is the plaquette cost.

\begin{proposition}[Exact pair-cost feasibility interval]
\label{prop:f12v97-pair}
There exist two unit quaternions with these link
costs and the prescribed relative trace if and only if
\[
 \left|e'-h-k+\frac{hk}{2}\right|
 \le
 2\sqrt{h(1-h/4)\,k(1-k/4)}.
 \tag{101.17}
\]
The statement includes \(h=0,4\) or \(k=0,4\).
\end{proposition}
\begin{proof}
Put \(a=1-h/2\), \(b=1-k/2\), \(c=1-e'/2\).
The Gram matrix on the identity and the two unit
quaternions has first row \((1,a,b)\).
Taking the Schur complement of its leading \(1\),
positive semidefiniteness is equivalent to
\[
 \begin{pmatrix}
 1-a^2&c-ab\\ c-ab&1-b^2
 \end{pmatrix}\succeq0,
 \quad\text{or}\quad
 |c-ab|\le\sqrt{(1-a^2)(1-b^2)}.
 \tag{101.18}
\]
The diagonal entries are nonnegative because
\(h,k\in[0,4]\). Substitute their expressions in
terms of \(h,k,e'\) to obtain (101.17).

Conversely, a positive semidefinite matrix in (101.18)
is a Gram matrix of two real vectors, which can be
embedded in \(\mathbb R^3\). Adding scalar components
\(a,b\) gives unit quaternions with relative scalar
product \(c\). The inversion signs from (101.8)
then recover the original two-link word.
This construction uses no division by \(1-a^2\)
or \(1-b^2\), so the degenerate endpoints are included.
\end{proof}

For a pivot cost \(k\le\kappa=10u^2\), (101.17) yields
the useful necessary estimate
\[
                 |e'-h|\le \kappa+2\sqrt{h\kappa}.
 \tag{101.19}
\]
Indeed the centre differs from \(h\) by
\(k(1-h/2)\), whose absolute value is at most \(k\);
the square-root term is at most \(2\sqrt{hk}\).
At \(k=0\) the exact condition is \(e'=h\).
Thus near-identity pivots require consistency of
link and plaquette \emph{costs}, not merely broad
overlap of independent cost envelopes.
This is a supported next target for sharper
interval restrictions or analytic pair constraints.

\subsection{Exact realization of the two-link subsystem}

The preceding pair statement has an exact local
three-link counterpart.

\begin{proposition}[Local triangle Gram closure]
\label{prop:f12v97-local-realization}
For one of the 128 orientable triangles, specify
the three link angles and three two-link plaquette
angles in \([0,1]\). Form the \(4\)-by-\(4\) matrix
(101.10) using their \emph{true} cosine costs.
Those six observables can be realized by three
unit quaternions satisfying those three reduced
two-link words if and only if this matrix is
positive semidefinite.
\end{proposition}
\begin{proof}
Necessity is the physical Gram construction.
For sufficiency, any positive semidefinite matrix
of size four with diagonal one is a Gram matrix
of four unit vectors in \(\mathbb R^4\). Apply
an orthogonal transformation carrying the first
vector to the identity quaternion. The other
three vectors are unit quaternions with the
specified scalar and relative scalar products.
Undo the inversion signs whose existence was
proved in (101.9). Formula (101.8) then gives
the prescribed plaquette traces. Since cosine
is injective on \([0,\pi]\), the normalized angles
are the specified ones.
\end{proof}

Local realizations need not agree on shared links
when independently chosen. The exact global condition
for the two-link subsystem keeps that compatibility.

\begin{theorem}[Componentwise rank-three completion criterion]
\label{thm:f12v97-components}
Specify every physical link angle \(x_i\in[0,1]\)
and every two-link plaquette angle \(z_{ij}\in[0,1]\).
Let \(s_i=\cos(\pi x_i)\). On the graph with an edge
for each two-link word (101.7), prescribe a partial
real symmetric matrix by
\[
 M_{ii}=1-s_i^2,\qquad
 M_{ij}=o_{ij}\bigl(s_i s_j-s_{ij}\cos(\pi z_{ij})\bigr).
 \tag{101.20}
\]
There is an algebraic assignment of unit quaternions
realizing all these link angles and two-link traces
if and only if, on each connected component, this
partial matrix has a positive semidefinite completion
of rank at most three.
\end{theorem}
\begin{proof}
If \(U_i=s_i+v_i\), unit norm gives
\(|v_i|^2=1-s_i^2\).
The scalar part of \(U_i^\alpha U_j^\beta\) is
\(s_i s_j-o_{ij}v_i\cdot v_j\).
Thus the imaginary vectors supply the completion
in (101.20), positive semidefinite of rank at most
three.

Conversely, factor each prescribed component
completion as a Gram matrix of real vectors
in \(\mathbb R^3\), padding with zero coordinates
if its rank is smaller. Assign \(U_i=s_i+v_i\).
The diagonal entries make these unit quaternions,
and the edge entries give precisely (101.7).
Each component can be placed in the same
\(\mathbb R^3\) independently: there are no
prescribed edges between distinct components.
For an isolated vertex one may choose
\(v_i=\sqrt{1-s_i^2}(1,0,0)\).
No nonzero-angle or invertibility hypothesis is
needed. This proves the equivalence.
\end{proof}

The word \emph{algebraic} in this theorem is essential.
It does not impose the nonlinear coarse-field
pivot-completion equations, nor the plaquettes
with three or four unfixed links. Even a successful
rank-three completion would not by itself give an
actual completed hidden field in the original problem.
Likewise, ordinary positive semidefinite completion
without the rank restriction is only a relaxation.

The graph and word counts are exact:
\[
 \begin{array}{c|rrrrrrrrr}
 \text{component size}&1&2&3&4&8&20&28&31&32\\ \hline
 \text{number}&184&80&4&32&10&2&2&4&8 .
 \end{array}
 \tag{101.21}
\]
There are 326 components, 856 nonisolated link
vertices, and 976 edges. The largest component
has 32 vertices. The cyclically reduced physical
plaquette word lengths have distribution
\[
 \begin{array}{c|rrrrr}
 \text{unfixed letters}&0&1&2&3&4\\ \hline
 \text{plaquettes}&5376&176&976&880&368 .
 \end{array}
 \tag{101.22}
\]
Thus there are explicitly 1248 three- or four-link
plaquettes outside the realization criterion,
in addition to the pivot-completion obligation.
The one-link trace relations are (101.4) and
must also hold for a proposed full list of observables.
These counts are independently reconstructed by
\begin{center}\small\ttfamily
verify\_ym\_f12\_v97\_graph\_structure.py.
\end{center}

\subsection{Continuation record and proof boundary}

The exact row generator is
\begin{center}\small\ttfamily
ym\_f12\_v97\_coupled\_geometry.py.
\end{center}
The floating-point proposal search is retained
separately. Neither numerical optimality nor a
claim of full Gram feasibility enters the accepted
dual certificate. In particular, the final reported
cosine objective of the relaxation profile is not
the energy of a verified physical competitor.

The next direction should use the exact dual support:
tighten pivot-sensitive pair-cost consistency, and
test compatible component completions with actual
pivot equations and the remaining longer words.
Adding more triangle cuts indiscriminately is not
justified by this certificate. Complete branch
coverage and conditional global-ground coverage
are still absent; the present gain is a finite-parent
ingredient, not a substitute for them.

V97 preserves the entire V96 mathematical body,
apart from the title version. V96 had 1609155 bytes
and SHA--256
\begin{center}\scriptsize\ttfamily
759B2BCF81D1F63E6EFAA3E6DAB808CE546774E0F9DCA1B04D80EC4A521F3190.
\end{center}
Only V97 is the active note; frozen witnesses and
checkers are retained. The companion audit was
performed at V95. The next audit and ten-version
direction review remain V100, when this lineage
is to be archived as \(\mathrm{V100\_f12}\) and
followed by a new compact V1.

\begin{firewall}
The new action lower bound is global over the
hidden variables of one specified finite SU(2)
parent. It is not a spectral gap. The realization
theorems concern explicit trace subsystems and
do not construct a continuum gauge theory.
No scale-uniform repair inequality, complete
conditional-ground classification, or physical
four-dimensional Yang--Mills mass gap has been
proved here. The full goal remains active.
\end{firewall}

\section{V98: pivot source balance and the missing completion equation}
\label{sec:f12v98}

\subsection{An upstream change justified by the preceding certificate}

V97 was a progress turn: it raised the finite-parent
action lower bound and proved precise trace-realization
criteria. Its exact dual showed that 11 of the 13
supporting Gram rows involved pivots. We now go
upstream to the equation defining those pivots.
The small-angle and Gram restrictions used so far
are not replacements for that equation.

This version proves a strongly convex variational
description of each pivot, with a residual error bound
and an exact source-balance criterion. It gives a
complete local source disjunction and verifies that
71 pivots must be identity on a specified free-angle
branch. Finally, a small-cap counterexample shows
why even perfect algebraic link realization cannot
close the fixed-coarse problem if pivot completion
is still omitted.

The potential below is an auxiliary function for
solving a single pivot equation. It is not the
Yang--Mills action, and its strong convexity is not
a physical spectral-gap statement.

\subsection{A uniformly convex pivot potential}

Use the invariant Hilbert--Schmidt metric on SU(2).
In unit-quaternion coordinates it is the round
three-sphere metric of radius \(\sqrt2\).
Let \(0<\rho\le1/4\), and use the same smooth
nonincreasing cutoff as V18:
\(\chi:[0,\infty)\to[0,1]\), equal to one on
\([0,1/2]\) and zero on \([1,\infty)\).
Define
\[
 \Phi_\rho(r)=\int_0^r t\chi(t/\rho)\,dt,\qquad
 g_\rho(U)=\chi(d(I,U)/\rho)\log U ,
 \tag{102.1}
\]
with the smooth zero extension for \(g_\rho\)
outside its logarithm support.
Since \(\Phi_\rho\) is constant for \(r\ge\rho\),
the function \(\Phi_\rho(d(B,P))\) is smooth
even at the cut locus of \(P\). At \(B=P\)
it agrees locally with \(d(B,P)^2/2\).

Fix a coarse input \(C\) and fourteen nonroot path
products \(P_1,\ldots,P_{14}\). They may be arbitrary
noncommuting SU(2) matrices. Put
\[
 R=\frac{7\rho}{8},\qquad
 \kappa_\rho=\frac18-\frac{49\rho^2}{256},
 \qquad
 V_C(B)=\frac12d(B,C)^2
            -\frac1{16}\sum_{j=1}^{14}\Phi_\rho(d(B,P_j)).
 \tag{102.2}
\]
We use this potential on the closed ball
\(\overline B_R(C)\), where its first term is smooth.

\begin{theorem}[Strongly convex pivot characterization]
\label{thm:f12v98-potential}
On \(\overline B_R(C)\),
\[
                         \operatorname{Hess}V_C
                                         \ge\kappa_\rho I>0.
 \tag{102.3}
\]
Its right-trivialized gradient is
\[
 \mathcal R_C(B)=\log(C^{-1}B)
                  +\frac1{16}\sum_{j=1}^{14}g_\rho(B^{-1}P_j).
 \tag{102.4}
\]
The unique critical point of \(V_C\) is the unique
solution of the actual pivot equation
\[
 B\exp\!\left(\frac1{16}
               \sum_{j=1}^{14}g_\rho(B^{-1}P_j)\right)=C .
 \tag{102.5}
\]
It lies in the interior of the indicated ball.
At the working cutoff \(\rho=1/64\),
\[
                  \kappa_\rho=\frac{131023}{1048576}>\frac19 .
 \tag{102.6}
\]
\end{theorem}
\begin{proof}
We first record the radial Hessian calculation rather
than importing a bound on the absolute cutoff derivative.
For distance \(r\) on a sphere of radius \(\sqrt2\),
the Hessian of a radial function \(f(r)\) has radial
eigenvalue \(f''(r)\) and transverse eigenvalue
\[
                    \frac{f'(r)}{\sqrt2}
                                  \cot(r/\sqrt2).
 \tag{102.7}
\]
For a direct derivation, take a unit-speed geodesic
transverse to the radial direction. In quaternion
coordinates its scalar product with the centre is
\(\cos(r/\sqrt2)\cos(t/\sqrt2)\).
Twice differentiating the cosine of its distance
gives \(r''(0)=\cot(r/\sqrt2)/\sqrt2\).
The radial direction is an ordinary distance
coordinate. Rotational symmetry gives the stated
eigenvalues and vanishing mixed terms.

For \(f(r)=r^2/2\) and \(r\le R\), the transverse
eigenvalue obeys
\[
 (r/\sqrt2)\cot(r/\sqrt2)
       \ge\cos(r/\sqrt2)\ge1-r^2/4,
 \tag{102.8}
\]
using \(\sin t\le t\), \(\cos t>0\), and
\(\cos t\ge1-t^2/2\) on this small range.
Its radial eigenvalue is one.

For \(f=\Phi_\rho\), write \(s=r/\rho\).
The radial eigenvalue on the support is
\(\chi(s)+s\chi'(s)\le1\).
The transverse eigenvalue is
\(\chi(s)(r/\sqrt2)\cot(r/\sqrt2)\le1\):
on \(0\le r\le\rho\) the cotangent factor is
nonnegative and at most one. Outside the support
both eigenvalues vanish. The limits at zero and
at the support boundary are covered by smoothness.
Consequently
\(\operatorname{Hess}\Phi_\rho(d(\,\cdot\,,P_j))\le I\),
without any absolute lower bound on that Hessian.
Subtracting these fourteen upper bounds from
(102.8) gives
\[
 1-R^2/4-14/16=\kappa_\rho,
\]
which proves (102.3).
For all \(\rho\le1/4\) the lower bound is at least
\(463/4096>1/9\), and direct rational substitution
gives (102.6).

The right gradient of a squared distance to \(C\)
is \(-\log(B^{-1}C)=\log(C^{-1}B)\).
The right gradient of \(\Phi_\rho(d(B,P_j))\)
is \(-g_\rho(B^{-1}P_j)\).
This proves (102.4), including its signs.

The closed ball is geodesically convex. Indeed its
angular radius is less than \(\pi/2\); the short
great-circle interpolation of two points is a
positive sine-weighted combination whose scalar
product with the centre stays at least the cap
threshold. Therefore (102.3) implies strict
convexity along every connecting geodesic.

The continuous potential has a minimum on the
compact ball. At a boundary point its outward
radial derivative is at least
\[
 R-\frac1{16}\sum_j|g_\rho(B^{-1}P_j)|>R-\frac{14\rho}{16}=0.
 \tag{102.9}
\]
The inequality is strict because each truncated
logarithm has norm strictly less than \(\rho\):
on its support its argument has distance less
than \(\rho\), and outside it the value is zero.
Thus a minimum cannot lie on the boundary.
It is an interior critical point, and strict
convexity makes it unique.

On this ball, (102.4) vanishes exactly when (102.5)
holds, by exponentiating the small principal
logarithm. Conversely, any solution of (102.5),
wherever initially sought in SU(2), satisfies
\[
 d(B,C)\le\frac1{16}\sum_j|g_\rho(B^{-1}P_j)|<R,
\]
so there are no additional solutions outside the
ball. This also gives a variational proof of
uniqueness at fixed coarse input, consistent with
V89's already established global pivot diffeomorphism.
\end{proof}

\subsection{Residual control and exact source balance}

\begin{proposition}[Residual-to-root and source-to-root bounds]
\label{prop:f12v98-residual}
Let \(B_*\) be the pivot in Theorem~\ref{thm:f12v98-potential}.
For any proposed pivot \(\widetilde B\in\overline B_R(C)\),
\[
 d(\widetilde B,B_*)\le
                      \frac{|\mathcal R_C(\widetilde B)|}{\kappa_\rho}.
 \tag{102.10}
\]
Define the source vector at the coarse input by
\[
                         S_C=\sum_{j=1}^{14}g_\rho(C^{-1}P_j).
\]
Then
\[
 d(B_*,C)\le\frac{|S_C|}{16\kappa_\rho},
 \qquad B_*=C\ \Longleftrightarrow\ S_C=0.
 \tag{102.11}
\]
\end{proposition}
\begin{proof}
Integrate the Hessian lower bound along the
unit-speed minimizing geodesic from \(B_*\) to
\(\widetilde B\). The derivative of \(V_C\) is zero
at its first endpoint and at least
\(\kappa_\rho d(B_*,\widetilde B)\) at its last.
Cauchy--Schwarz bounds that last directional
derivative by the gradient norm, giving (102.10).
At \(\widetilde B=C\), the residual is \(S_C/16\),
which gives the distance estimate.
The equivalence follows from (102.4) and uniqueness
of its zero.
\end{proof}

The source condition allows cancellation between
nonzero truncated logarithms. It is not equivalent
to saying that every source is inactive.
However, the following sufficient identity test is
often easier to certify:
\[
       d(C,P_j)\ge\rho\ \text{for every }j
                       \quad\Longrightarrow\quad B_*=C.
 \tag{102.12}
\]
Unlike a test uniform in all possible \(B\), this
needs no subtraction of the pivot cap radius:
it tests the residual at the proposed solution \(B=C\)
and then uses uniqueness.

For the full parent, nonroot path families contain
no pivots. Thus all 96 residuals can be evaluated
with the free links fixed, without a coupled solve
between different pivots. There is also a useful
action error estimate. If every proposed pivot lies
in its admissible ball, let \(\widetilde{\mathcal S}\)
be the raw fine-link action before correcting them.
Then the completed action satisfies
\[
 |\mathcal S-\widetilde{\mathcal S}|
       \le\frac{6\sqrt2}{\kappa_\rho}
                          \sum_e|\mathcal R_{C_e}(\widetilde B_e)|.
 \tag{102.13}
\]
To prove this, each physical pivot occurs in six
plaquettes. A single Wilson plaquette cost has
Hilbert--Schmidt gradient norm at most \(\sqrt2\)
in any one link: differentiating its scalar
quaternion component gives that bound directly.
Move the pivots one at a time along their shortest
geodesics and sum the resulting \(6\sqrt2\) distance
bounds, then apply (102.10). This is a per-field
completion error estimate, not a new uniform
lower bound on the global ground energy.
The residual in these statements is the explicit
gradient (102.4); an unquantified forward group
equation error must not be substituted for it.

\subsection{A complete local source disjunction in free angles}

Return to coarse input \(I\) and \(\rho=1/64\).
For the specified parent, exact path enumeration
finds 1344 nonroot two-step paths. Of these, 800
contain one free link and 544 contain two free links.
Every other factor in these paths is \(+I\).
Thus their source products are either \(U_i\) or
\(U_iU_j\), with no pivot variables among them.

Write \(x_i=\arccos((U_i)_0)/\pi\in[0,1]\).
The normalized geodesic angle of a product satisfies
\[
                      \theta(U_iU_j)\ge |x_i-x_j|.
 \tag{102.14}
\]
Bi-invariance gives
\[
                 d(I,U_iU_j)=d(U_i^{-1},U_j).
\]
Apply the reverse triangle inequality and use
invariance of the angle under inversion.
The exact cutoff angle obeys
\[
       \frac{\rho}{\sqrt2\pi}
           <\frac{\rho}{(7/5)(157/50)}<\tau,
                         \qquad \tau=\frac1{250}.
 \tag{102.15}
\]

\begin{proposition}[Source cover for a nonidentity pivot]
\label{prop:f12v98-source-cover}
For each physical pivot at coarse input \(I\),
every completed field satisfies at least one of
the following closed branch conditions:
the pivot is identity; a one-free-link nonroot
path has \(x_i\le\tau\); or a two-free-link
nonroot path has \(|x_i-x_j|\le\tau\).
There are at most fifteen alternatives, counting
the identity branch and the fourteen source paths.
\end{proposition}
\begin{proof}
If the pivot is not identity, (102.11) implies
\(S_I\ne0\). Some summand \(g_\rho(P_j)\) is then
nonzero, so \(\theta(P_j)<\rho/(\sqrt2\pi)<\tau\).
For a one-link path this is \(x_i<\tau\).
For a two-link path, (102.14) gives
\(|x_i-x_j|<\tau\). Passing to closed inequalities
preserves coverage. The cases need not be disjoint.
\end{proof}

This is a complete \emph{local} source cover, not
an executed global branch tree. The identity
alternative must be retained even when some
sources are active, since they may cancel.
Conversely, a source-strip alternative is only
necessary for possible activity, not sufficient
for a nonidentity pivot.

For free-angle intervals \(x_i\in[l_i,u_i]\), the
following lower bounds require only rational arithmetic:
\[
 \begin{array}{c|c}
 \text{source product}&\text{normalized-angle lower bound}\\ \hline
 U_i&l_i\\
 U_iU_j&\max\{0,l_i-u_j,l_j-u_i\}.
 \end{array}
 \tag{102.16}
\]
If all fourteen lower bounds are at least \(\tau\),
(102.12) forces that pivot to be exactly identity
throughout the branch. These are derived constraints
on the actual completion, not an assumption that
small pivots can be frozen.

\begin{proposition}[Seventy-one certified identity consequences]
\label{prop:f12v98-seventy-one}
Let \(x^0\) be the first 944 coordinates of V96's
frozen rational profile. On the free-angle box
\[
      x_i\in
      [\,\max(0,x_i^0-10^{-4}),\min(1,x_i^0+10^{-4})\,],
                    \qquad 0\le i<944,
 \tag{102.17}
\]
the test (102.16) forces 71 of the 96 actual pivots
to be identity. The frozen profile assigns a
strictly positive angle to each of those 71 pivots.
\end{proposition}
\begin{proof}
The checker reconstructs the 96 ordered nonroot
families and the free-link indices. It verifies
the counts \(800+544=1344\) and that all fixed
path factors are \(+I\), then evaluates (102.16)
on the rational endpoints in (102.17).
Exactly 71 families have every lower bound at
least \(1/250\).
For example, the coarse edge \(e_0=((0,0,0,0),2)\)
has pivot-angle index 946 and minimum path lower bound
\[
                         \frac{11015453}{2000000000}>\frac1{250}.
 \tag{102.18}
\]
The frozen angle there is
\(1941851/625000000>0\).
All 71 detected pivots have positive frozen angles,
as checked individually. Applying (102.12) to
each certified family proves the assertion.
\end{proof}

The free-angle box itself has not been shown empty.
An actual field in that box must use the forced
identity values instead of the frozen profile's
71 nonzero pivot values. The other 25 pivots
remain unresolved by this test; failure of the
test is not evidence that their sources are active.
V96 already excluded its full profile by a Gram
condition. The new result is the family of
completion-based constraints valid on the entire
free-angle box, independently of the plaquette
angle choices or the link axes.

\subsection{A ceiling for exact algebraic realization with independent pivots}

The need for completion is not just a defect of
the frozen relaxed profile. Define the
\emph{cap-only fine-field class} to contain all
unit-quaternion fine-link assignments with the
specified fixed exterior and tree links, arbitrary
944 free links, and each of the 96 independent
pivots in its radius-\(R\) ball about \(I\).
In this class use the exact plaquette words and
the true action, but do not impose (102.5).
All actual fixed-coarse completed fields belong
to this larger class.

\begin{theorem}[Small-cap completion obstruction]
\label{thm:f12v98-cap-ceiling}
There is a genuine fine-link assignment in the
cap-only class whose exact action is less than
\[
                              E_*-\frac{19}{1000},
 \tag{102.19}
\]
where \(E_*\) is the action of the V89--V90
certified local-minimum orbit. It satisfies all
algebraic link-word, exact-cost, and quaternion
Gram identities, including those involving
three or four unfixed links, but it is not an
actual completed field at coarse input \(I\).
Consequently no relaxation containing the whole
cap-only class can certify a lower bound \(E_*\).
\end{theorem}
\begin{proof}
Use the exact commuting representative \(Y_*\)
and the V92 coarse edge
\(e_0=((0,0,0,0),2)\), whose pivot is
\(\ell_0=((0,0,1,0),2)\).
At \(Y_*\) this pivot is identity. V92 proved,
throughout the containing critical-point box,
that all fourteen source products at \(e_0\)
have Hilbert--Schmidt chord squared greater
than one. It also proved the incident-plaquette
force estimate \(f>49/5\) and \(|a|\le12\).
We use these proved estimates rather than
repeating their interval derivation.

Replace only this physical pivot by the rational
unit quaternion
\[
 B_\epsilon=cI-sT,\qquad
 c=\frac{1048575}{1048577},\quad
 s=\frac{2048}{1048577}.
 \tag{102.20}
\]
Indeed \(c^2+s^2=1\); its angular displacement is
\(2\arctan(1/1024)<1/512\). Hence
\[
 d(I,B_\epsilon)<\frac3{1024}<\frac7{512}=R,\qquad
 \theta(B_\epsilon)<\frac{25}{40192}<\frac{125}{40192}.
 \tag{102.21}
\]
Its cost is \(4/1048577<10(125/40192)^2\).
Thus this change stays inside both the geometric
pivot cap and the angle/cost boxes used by V97.
The other physical links, including the four
already completed moving pivots, are unchanged.

For the six incident plaquettes, V92's exact
trigonometric variation formula gives
\[
 \begin{split}
 \mathcal S_{\mathrm{raw}}-E_*
      &=a(1-c)-fs\\
      &\le 12(1-c)-\frac{49}{5}s
       =-\frac{100232}{5242885}
       <-\frac{19}{1000}.
 \end{split}
 \tag{102.22}
\]
This proves the energy statement inside the
cap-only class. It is a genuine unit-quaternion
assignment, so every algebraic word and Gram
identity holds automatically; it is not merely
an angle vector with guessed relative traces.

The source products do not contain the changed
pivot and remain unchanged. Chord distance is
no greater than geodesic distance, so
\(d(I,P_j)>1\). For every \(B\) in the original
cap,
\[
                    d(B,P_j)>1-R>\rho.
 \tag{102.23}
\]
All truncated logarithms vanish there. The
forward pivot map is therefore exactly
\(F_{e_0}(B)=B\) throughout that cap.
At \(B_\epsilon\ne I\) its coarse output is
\(B_\epsilon\), not \(I\). Thus (102.5) fails.
Equivalently, (102.12) requires the true pivot
to be identity.

Finally any feasible relaxation containing
every cap-only fine field also contains this
assignment, with its true action. Its infimum
is therefore below the right side of (102.19).
No collection of identities valid for all
cap-only fields can remove this witness.
\end{proof}

V92 already showed energy descent when a coarse
input is changed. Its larger displayed endpoint
was not inside the small independent-pivot cap
used in the global relaxations. The new endpoint
(102.20) is inside that cap, which is essential
for the present ceiling statement. It remains
a changed-coarse configuration, not a competitor
in the original fixed-coarse hidden minimization.
There is no contradiction with V90's local
minimality, and no claim that \(E_*\) is the
actual global minimum.

\subsection{Validation and next direction}

The exact checkers are
\begin{center}\small\ttfamily
verify\_ym\_f12\_v98\_pivot\_source.py\\
verify\_ym\_f12\_v98\_pivot\_gate\_branches.py .
\end{center}
The interval branch helper is
\begin{center}\small\ttfamily
ym\_f12\_v98\_pivot\_gates.py .
\end{center}
They check the rational constants, exact path
counts, all 71 identity consequences, the
small-cap quaternion, and the whole-box force
and source-distance estimates reused from V92.
The general Hessian, residual, and coverage
statements are proved above; finite executable
checks do not replace those analytic arguments.

The next global refinement should therefore
include completion-aware source branches or
certified residual restrictions, retaining every
unresolved alternative in Proposition~\ref{prop:f12v98-source-cover}.
Ordinary Gram completion, even with exact costs
and exact longer words, does not impose the
fixed coarse field. The cap-only obstruction
proves why pursuing that relaxation alone cannot
establish equality with \(E_*\).
No complete source branch tree or improved
global action lower bound is claimed in V98.

V98 preserves the entire V97 mathematical body
apart from its title version. V97 had 1627355 bytes
and SHA--256
\begin{center}\scriptsize\ttfamily
B11206D3D336CF7572FC8EF4AD54F6F264EBBCBD8E3EF4121025D902BB489B12.
\end{center}
Only V98 is retained as the active note.
The next companion audit and ten-version
direction review remain V100. At that boundary,
preserve this lineage as \(\mathrm{V100\_f12}\)
and restart with a checked compact V1.

\begin{firewall}
V97's action bound remains in force.
No stronger global bound is claimed here.
Strong convexity of the auxiliary pivot potential
does not imply convexity of the full gauge action
or a Markov or physical spectral gap.
Complete conditional-ground coverage, a
scale-uniform repair estimate, four-dimensional
continuum construction, and the physical
Yang--Mills mass gap remain unproved.
The full objective stays active.
\end{firewall}

\section{V99: a complete source cover and an energy-conditioned identity pivot}
\label{sec:f12v99}

\subsection{From conditional gates to a global certificate}

V98 was a progress turn: it proved a source-balance
description of completion, certified conditional identity
pivots, and exhibited an obstruction to omitting the
completion equation. We now execute one complete source
disjunction rather than choose only a favourable
identity branch. All eight cases are retained and
verified. Their weakest certified bound gives
\[
                      \mathcal S\ge549.617453988
 \tag{103.1}
\]
over the actual fixed-coarse hidden problem.

The source cases also give an energy threshold below
which one selected completion factor is exactly the
identity map on a uniform local neighbourhood.
The existence of fields below that threshold has
not been established. Thus this is not yet coverage
of the conditional global-ground set.

All statements use the same period-six SU(2) parent,
central exterior, coarse field \(I\), and
\(\rho=1/64\). The 944 free groups remain unrestricted
in their axes; there is no commuting-field assumption.

\subsection{The selected source family}

Choose
\[
 e_\dagger=((1,1,0,1),3),\qquad
 \ell_\dagger=((2,2,0,3),3).
 \tag{103.2}
\]
Its physical pivot-angle index is 1015 in the
existing 3440-angle enumeration. The exact geometry
shows that its fourteen nonroot paths consist of
seven distinct free-link matrices, each appearing
twice. Their free-link indices are
\[
                  J=(557,589,604,767,782,823,838).
 \tag{103.3}
\]
All other factors in these paths are \(+I\).
Consequently the actual pivot equation at this edge is
\[
 B\exp\!\left(\frac18\sum_{j\in J}g_\rho(B^{-1}U_j)\right)=I.
 \tag{103.4}
\]
This identity is reconstructed from the two-step
paths; it does not replace two-link products by
commuting surrogates. The special simplification
occurs because these paths have only one unfixed
link apiece.

Let \(y_B\) be the normalized pivot angle and
\(x_j\) the normalized angle of \(U_j\). Put
\(\tau=1/250\), as in V98.

\begin{proposition}[A complete eight-case cover]
\label{prop:f12v99-cover}
Every actual completed field at the specified
coarse input belongs to at least one of
\[
 \mathcal C_0=\{y_B=0\},\qquad
                  \mathcal C_j=\{x_j\le\tau\},\quad j\in J.
 \tag{103.5}
\]
\end{proposition}
\begin{proof}
If \(B=I\), the field belongs to \(\mathcal C_0\).
Otherwise V98's exact source-balance criterion
implies \(\sum_{j\in J}g_\rho(U_j)\ne0\).
At least one summand is nonzero, hence
\[
             x_j<\frac{\rho}{\sqrt2\pi}<\frac1{250}.
 \tag{103.6}
\]
It therefore belongs to the corresponding closed
source case. The duplicate multiplicity two
does not affect whether the sum vanishes.
\end{proof}

The cases overlap, which is harmless for lower
bounds. They are necessary conditions, not
equations solving the entire completion problem.
For example, the identity case does not itself
impose the source-vector cancellation required
for \(B=I\). Each branch relaxation is allowed
to contain such extra, nonphysical points.
What matters for (103.1) is that no actual
completed field is omitted from their union.

\subsection{All branch lower bounds are certified}

Use exactly V97's 6880-variable coupled model:
174153 base rows and the 1050 frozen Gram rows.
On case zero add the single row
\[
                              P y_B\le0,\qquad P=10^9.
 \tag{103.7}
\]
The nonnegative angle box then forces \(y_B=0\).
The existing endpoint cost bound \(t_B\le Ly_B\)
also forces \(t_B=0\); this does not require
another unproved constraint.
On a source case add
\[
                                  P x_j\le P/250=4000000 .
 \tag{103.8}
\]
Thus every case has 175204 non-box rows and
the same globally valid upper boxes as V97.
No equality is imposed on other pivots.

The exact branch results are listed below.
Each nine-place decimal in the third column
denotes the displayed rational number with
denominator \(10^9\), not an approximate optimum.
\[
 \begin{array}{c|c|c|r}
 k&\text{added condition}&\text{certified lower bound}&
                         \text{negative residuals}\\ \hline
 0&y_{1015}=0&549.617453988&845\\
 1&x_{557}\le1/250&550.616912146&878\\
 2&x_{589}\le1/250&549.647236309&810\\
 3&x_{604}\le1/250&549.969262071&862\\
 4&x_{767}\le1/250&549.642651322&780\\
 5&x_{782}\le1/250&549.893883248&811\\
 6&x_{823}\le1/250&549.647205753&804\\
 7&x_{838}\le1/250&549.718151597&802 .
 \end{array}
 \tag{103.9}
\]

\begin{theorem}[Global bound from the complete source cover]
\label{thm:f12v99-bound}
For every actual field in the fixed parent problem,
\[
                     \mathcal S(Y)\ge
                       \frac{137404363497}{250000000}.
 \tag{103.10}
\]
This increases V97's certified bound by
\[
                              \frac{16952677}{10^9}.
 \tag{103.11}
\]
\end{theorem}
\begin{proof}
For completeness, the numerical dual of each case
is treated only as a proposal. Its frozen entries
are integers \(v^{(k)}_r\le0\) with denominator
\(D=10^8\). The checker reconstructs the exact
scaled rows \(\widetilde A^{(k)}=PA^{(k)}\),
\(\widetilde d^{(k)}=Pd^{(k)}\), including the
correct source row or identity row for that case.
The objective is \(408+c^{\mathsf T}z\), with
unit coefficients on the 2400 plaquette costs
and zero on the other variables.

Let \(M=40192^2\), and write the common variable
upper bounds as \(U_i=\widetilde U_i/M\).
For every case define the integer residual
\[
 R_i^{(k)}=PDc_i-
                   \sum_r\widetilde A^{(k)}_{ri}v^{(k)}_r .
 \tag{103.12}
\]
Multiplication of each valid row by its nonpositive
dual weight, followed by minimization of the
residual on the coordinate boxes, proves
\[
 \mathcal S\ge
 408+\frac{\sum_r\widetilde d^{(k)}_r v^{(k)}_r}{PD}
       +\frac{\sum_i\min(R_i^{(k)},0)\widetilde U_i}{PDM}
                   \quad\text{on case }k.
 \tag{103.13}
\]
All negative residuals counted in (103.9) are
included in this expression.

The checker verifies that all eight case indices
occur, that each vector has exactly 175204 entries,
and that every entry is nonpositive. It evaluates
(103.13) using arbitrary-precision integers and
rounds each final bound downward to denominator
\(10^9\). The results are exactly (103.9).
The counts of nonzero dual entries, in case order,
are
\[
                (3672,3701,3665,3647,3653,3663,3675,3653).
 \tag{103.14}
\]
These finite checks establish a lower bound on
each branch independently of numerical optimality.

Proposition~\ref{prop:f12v99-cover} assigns every
actual field to at least one branch. Therefore
its action is at least the minimum of all eight
certified branch bounds. That minimum is the
case-zero value, proving (103.10).
Subtracting \(549600501311/10^9\) proves (103.11).
\end{proof}

As an arithmetic audit of the minimizing case,
its independently rounded dual and compensation
parts are
\[
 408+d^{\mathsf T}\lambda
       \ge\frac{551621700063}{10^9},\qquad
 \sum_i\min(r_i,0)U_i\ge-\frac{80169843}{40000000}.
 \tag{103.15}
\]
Their sum is the case-zero bound in (103.9).
The exact verifier always computes the combined
expression before reporting a final downward
rounding; it does not rely on a rounded optimizer
value.

The frozen eight-vector data file is
\begin{center}\small\ttfamily
ym\_f12\_v99\_source\_duals.py
\end{center}
with SHA--256
\begin{center}\scriptsize\ttfamily
03C5656D61EDC90FBCE3085E99A798A71ABAD5C0E91D0835DC130651C2BB5AFD.
\end{center}
The source identities and exact rows are generated by
\begin{center}\small\ttfamily
ym\_f12\_v99\_source\_branches.py,
\end{center}
and the accepted checker is
\begin{center}\small\ttfamily
verify\_ym\_f12\_v99\_source\_cover.py.
\end{center}
Lossless packing of the integer vectors is only
serialization. The numerical proposal code is
separate and is not imported by the checker.

\subsection{A precise energy-conditioned identity statement}

Let
\[
                          T_\dagger=
                       \frac{274821325661}{500000000}
                                  =549.642651322.
 \tag{103.16}
\]
This is the smallest bound among the seven
source cases, not among all eight cases.

\begin{proposition}[Source exclusion below a certified threshold]
\label{prop:f12v99-low-energy}
For any actual field with \(\mathcal S(Y)<T_\dagger\),
\[
                         x_j>\frac1{250}\quad(j\in J),
                              \qquad B_{e_\dagger}=I.
 \tag{103.17}
\]
Equivalently, every actual field with a nonidentity
selected pivot has action at least \(T_\dagger\).
\end{proposition}
\begin{proof}
If any \(x_j\le1/250\), the field belongs to a
source case and (103.9) gives
\(\mathcal S\ge T_\dagger\), a contradiction.
Thus all seven source angles exceed the threshold.
By (103.6), every source logarithm at \(I\) vanishes,
so V98's source-balance criterion forces \(B=I\).
Alternatively, a nonidentity pivot must belong
to a source case by the complete cover.
\end{proof}

The gap between the identity-case bound and the
source threshold is
\[
              T_\dagger-\frac{137404363497}{250000000}
                         =\frac{12598667}{500000000}>0.
 \tag{103.18}
\]
This is a gap between certified lower bounds,
not a spectral gap or a proved excitation energy
above the actual ground state. In particular,
we have not shown that the sublevel
\(\{\mathcal S<T_\dagger\}\) is nonempty.
The known orbit has energy about \(569.6844\),
above this threshold. Neither a ground-state
identity claim nor a low-temperature probability
estimate follows merely from (103.18).

\subsection{A uniform local completion factor on the conditional sublevel}

The strict source exclusion provides a quantitative
chart statement, conditional on the same sublevel.
All distances in the following theorem are
Hilbert--Schmidt geodesic distances.

\begin{theorem}[Uniform identity completion map near a low-energy field]
\label{thm:f12v99-flat-chart}
Take any actual field at coarse input \(I\) with
\(\mathcal S<T_\dagger\). Perturb each of its seven
source links \(U_j\), \(j\in J\), by distance at
most \(1/2000\), and change the selected coarse
input from \(I\) to \(C\) with \(d(I,C)\le1/1000\).
Then the actual selected pivot is exactly
\[
                                  B_{e_\dagger}=C.
 \tag{103.19}
\]
The selected forward pivot map is the identity
on \(d(I,B)\le1/1000\). Its Haar Jacobian is one
there, its inverse is independent of these seven
free links locally, and its corresponding
inverse-Jacobian amplitude factor is one.
These assertions concern this one completion factor.
\end{theorem}
\begin{proof}
By (103.17), each unperturbed source has
\[
 d(I,U_j)>\sqrt2\pi/250
                   >\frac{7}{5}\frac{157}{50}\frac1{250}
                    =\frac{1099}{62500}.
 \tag{103.20}
\]
Let \(U'_j\) be a perturbed source and let
\(d(I,B)\le1/1000\). The triangle inequality gives
\[
 \begin{split}
 d(B,U'_j)
   &>\frac{1099}{62500}-\frac1{1000}-\frac1{2000}\\
   &=\rho+\frac{459}{10^6}>\rho .
 \end{split}
 \tag{103.21}
\]
The exact word identities in (103.4) still apply:
the fixed tree factors stay \(+I\), and each
perturbed source occurs twice.
All truncated logarithms in that forward map
are therefore zero on the entire indicated
pivot ball. The map is \(F(B)=B\) there.
For the proposed value \(B=C\), it sends \(C\)
to \(C\). V98's uniqueness theorem proves that
this is the actual completed pivot.

The margin \(459/10^6\) is strict and uniform,
so the identity holds on an open neighbourhood
of each listed configuration. Its differential
in the pivot variable is identity and its
derivatives in the seven source variables vanish.
The Haar Jacobian and reciprocal amplitude
factor are consequently one.
\end{proof}

This does not say that the full parent action
has a positive Hessian on the chart, or that
the other 95 pivot amplitude factors are one.
It does not cover all conditional global
minimizers without a further energy upper bound
below \(T_\dagger\). It gives a quantitatively
controlled completion factor on the stated
conditional region.

\subsection{How the complete cover can be continued}

The identity branch is the present weakest leaf.
A further source split can be applied inside
that branch while retaining the other seven
leaves and their verified bounds. Every new
leaf must inherit all conditions of its parent,
and the source alternatives must cover every
actual field in that parent. The global lower
bound is always the minimum over the retained
leaves. In particular, a solved source branch
cannot be discarded just because another
branch appears more promising numerically.

This is a concrete continuation strategy, not
an executed multi-pivot tree in V99. Only the
eight leaves for (103.2) have been solved and
certified here. The other completion equations,
the full source-vector equations within each
necessary branch, and exact link realization
remain additional restrictions not fully
encoded in these relaxations.

The numerical solves all returned candidate
duals, but their stopping messages are not
part of the proof. Verification checks the
complete case list, every exact row and
dual sign, all box residuals, and the final
minimum. The analytic link from this finite
calculation to actual fields is
Proposition~\ref{prop:f12v99-cover}.

V99 preserves the complete V98 mathematical
body apart from the title version. V98 had
1646885 bytes and SHA--256
\begin{center}\scriptsize\ttfamily
9B7682305C8B5C9A25627CC144B98BA0BC08363288CF92E59BDEF930B92D3BBF.
\end{center}
Only V99 is retained as the active note.
The next version is the scheduled V100
companion audit and ten-version direction
review. After that version, retain the
completed lineage as \(\mathrm{V100\_f12}\)
and start a new compact V1 with checked
context and a summary of proved results.

\begin{firewall}
The new bound is a global action lower bound
for one finite parent, obtained from a complete
necessary source cover. The conditional
identity-pivot chart is not a spectral-gap
theorem and may describe an empty sublevel.
No complete conditional-ground coverage,
scale-uniform repair estimate, continuum
Yang--Mills construction, or physical mass
gap is proved here. The full goal remains active.
\end{firewall}

\section{V100: companion audit, direction review, and twelfth archive}
\label{sec:f12v100}

\subsection{Companion audit}

This is the scheduled five-version companion audit
and ten-version direction review. The current files
were inspected again, including their latest conclusions.
They are unchanged from the V95 audit:
\[
 \begin{array}{c|c|r}
 \text{programme}&\text{active version}&\text{bytes}\\ \hline
 \text{SM--Einstein}&72&607372\\
 \text{Navier--Stokes}&26&278283\\
 \text{parallel YM}&5&54174 .
 \end{array}
 \tag{104.1}
\]
Their SHA--256 values, in that order, are
\begin{center}\scriptsize\ttfamily
F44F9745D43379E1672C5550411B58E97195A4C9278592D221DF5D3F644DC1F5\\
82C887F8C2C69E4F28FAD7C3DC8DE5B9099CB5A4BF431EBA1FFF3E91935978AD\\
306F1BB84CFA8A8D47EA569EC17E9545F9867C8D32B548EC9D9C444B4F3C0512.
\end{center}

SM V72 proves the stated many-copy covariance flow
limit at fixed mode cutoff and bounded time.
It explicitly leaves growing ultraviolet cutoff,
kinetic time, and the physical finite-species
continuum uncontrolled. NS V26 gives pointwise
rank-one norm formulas, not global regularity.
Parallel YM V5 rules out the proposed uniform
tail for an extensive total charge under
tensorization, while retaining compact labels
and fixed-window or density formulations.
None supplies an all-exterior repair inequality,
an escape estimate from the present orbit tube,
or a physical continuum mass gap.

The parallel note's internal stop instruction
describes that document's failed route. It is
source material, not an instruction to terminate
the user's current programme. Its mathematical
warning about extensive charges remains relevant.

\subsection{What the last ten versions actually achieved}

V91 excluded the central hidden sector by an
exact action bound and separated a fixed
noncommuting neighbourhood of it. V92 established
joint regularity of the already located orbit,
not of all conditional global minima.
V93--V97 constructed and strengthened necessary
angle and Gram relaxations, with exact dual
certificates and explicit nonrealizability
obstructions. V98 returned to the actual pivot
equation and proved source balance and residual
control. V99 verified every case of one complete
source disjunction.

The current global finite-parent action bound is
\[
                     \mathcal S(Y)\ge549.617453988.
 \tag{104.2}
\]
The known orbit energy lies near \(569.6844\);
the orbit is a certified local minimum, not a
proved global minimum. The V99 conditional
identity threshold \(549.642651322\) is below
the known orbit energy. Nonemptiness of that
sublevel and coverage of actual ground states
by it have not been established.

There are now three proved reasons not to expect
unrestricted repetition of the same lower-bound
calculation to close the programme. Globally
convex separable cost minorants have V95's
objective ceiling. Scalar surface constraints
do not imply shared-link Gram realization,
as V96 demonstrates. Even exact algebraic
fine-link realization with small independent
pivots does not impose fixed-coarse completion,
as V98 demonstrates.

\subsection{The spectral input to acquire next}

The finite action inequalities do not estimate
the Dirichlet form of the actual conditional
diffusion. V90 did establish an actual
Morse--Bott orbit with a strictly positive
normal Hessian, and V73 already supplied a
local-Poincare-plus-escape assembly theorem.
What has not yet been written for this particular
orbit is the corresponding Neumann Poincare
inequality on its whole normal tube, with the
symmetry directions and actual amplitude retained.

The compact restart will derive that local
spectral input from the certified normal
Hessian, using an equivariant normal bundle.
This avoids assuming that the orbit directions
are strictly convex, or that a singular orbit
quotient has a smooth positive coordinate density.
The remaining quantitative escape estimate
outside a smaller tube must stay explicit.
No amount of local tube control alone supplies
that estimate, and a fixed-parent conditional
diffusion gap would still not be the physical
Hamiltonian mass gap.

The next acquisition target after the local
inequality is therefore an actual escape or
ground-coverage estimate, with its dependence
on exterior data and scale visible. A further
source split is justified when it certifies
such coverage or excludes a specific low-energy
obstruction, not merely because it raises
an unrelated action bound slightly.

\subsection{Archive record}

The complete V99 mathematical body is retained
above, apart from updating the title to V100.
V99 had 1661127 bytes and SHA--256
\begin{center}\scriptsize\ttfamily
D8C6E4F511FC70851A93C25AE86CC05122345977421904B4A4769125847EB547.
\end{center}
This completed lineage is preserved as
\begin{center}\small\ttfamily
Renewal\_Yang\_Mills\_Mass\_Gap\_Research\_Notes\_V100\_f12.tex.
\end{center}
The next active document restarts at V1,
with a scope-checked summary and new work.
The archives, frozen data, and exact checkers
are retained. This is an organizational
rollover, not a completion of the mass-gap goal.

\begin{firewall}
The full four-dimensional quantum construction
and physical Yang--Mills mass gap remain unproved.
Neither unchanged companion files nor successful
finite certificates close those obligations.
\end{firewall}

\end{document}
